% Auto-generated from data/segments.json + layout.json (+ transforms) — do not edit by hand.
% volume: 39Abhi11
% mode: printing
% segments: 7144
% transform_rules: 8
% toc_marks: 310

% Volume layout defaults (layout.json ``layout``).
\csromanlayoutapply{1}{1.5}{21.6pt}{8pt}{8pt}{65pt}{2.5em}%

\pitaka{อภิธมฺม\-ปิฏก}
\csromanheader{2}{\textbf{ธมฺมานุโลม}}
\csromanmatikaanchor{mtk.0}
\csromantocmarknopage{chapter}{ปฏฺฐานปาฬิ}{mtk.0}
\bookmark[dest={mtk.0},level=0]{ปฏฺฐานปาฬิ}
\gambhira{ทุกปฏฺฐาน\-ปาฬิ}
\namakkaram{นโม ตสฺส ภควโต อรหโต สมฺมา\-สมฺพุทฺธสฺส.}
\csromanmatikaanchor{mtk.1}
\csromanmatikaanchor{mtk.2}
\csromantocmarknopage{section}{จตุตฺถภาค}{mtk.1}
\bookmark[dest={mtk.1},level=1]{จตุตฺถภาค}
\csromantocmarknopage{chapter}{12. กิเลสโคจฺฉก}{mtk.2}
\bookmark[dest={mtk.2},level=0]{12. กิเลสโคจฺฉก}
\chapterhead{12. \textbf{กิเลสโคจฺฉก}}
\csromanmatikaanchor{mtk.3}
\csromantocmarkref{section}{75. กิเลสทุก}{mtk.3}
\bookmark[dest={mtk.3},level=1]{75. กิเลสทุก}
\csromanmarkh{75. กิเลสทุก}\csromanmarkhii{1. ปฏิจฺจวาร}\csromanheadernomark{1}{75. \textbf{กิเลสทุก 1. ปฏิจฺจวาร}}
\proseitem{1}{\csromanfolio{1}กิเลสํ ธมฺมํ ปฏิจฺจ กิเลโส ธมฺโม อุปฺปชฺชติ เหตุปจฺจยา. โลภํ ปฏิจฺจ โมโห, ทิฏฺฐิ, ถินํ\csromansharedfootnote{e0e469e6132e5391}{ถีนํ (สี, สฺยา)}, อุทฺธจฺจํ, อหิริกํ, อโนตฺตปฺปํ. โลภํ ปฏิจฺจ โมโห, ทิฏฺฐิ, อุทฺธจฺจํ, อหิริกํ, อโนตฺตปฺปํ. โลภํ ปฏิจฺจ โมโห, มาโน, ถินํ, อุทฺธจฺจํ, อหิริกํ, อโนตฺตปฺปํ. โลภํ ปฏิจฺจ โมโห, มาโน, อุทฺธจฺจํ, อหิริกํ, อโนตฺตปฺปํ. โลภํ ปฏิจฺจ โมโห, ถินํ, อุทฺธจฺจํ, อหิริกํ, อโนตฺตปฺปํ. โลภํ ปฏิจฺจ โมโห, อุทฺธจฺจํ, อหิริกํ, อโนตฺตปฺปํ. โทสํ ปฏิจฺจ โมโห, ถินํ, อุทฺธจฺจํ, อหิริกํ, อโนตฺตปฺปํ. โทสํ ปฏิจฺจ โมโห, อุทฺธจฺจํ, อหิริกํ, อโนตฺตปฺปํ. วิจิกิจฺฉํ ปฏิจฺจ โมโห, อุทฺธจฺจํ, อหิริกํ, อโนตฺตปฺปํ. อุทฺธจฺจํ ปฏิจฺจ โมโห, อหิริกํ, อโนตฺตปฺปํ. (1)}

\prose{\csromanfolio{2}กิเลสํ ธมฺมํ ปฏิจฺจ โนกิเลโส ธมฺโม อุปฺปชฺชติ เหตุปจฺจยา. กิเลสํ ปฏิจฺจ สมฺปยุตฺตกา ขนฺธา จิตฺต\-สมุฏฺฐานญฺจ รูปํ. (2)}

\prose{กิเลสํ ธมฺมํ ปฏิจฺจ กิเลโส จ โนกิเลโส จ ธมฺมา อุปฺปชฺชนฺติ เหตุปจฺจยา. โลภํ ปฏิจฺจ โมโห, ทิฏฺฐิ, ถินํ, อุทฺธจฺจํ, อหิริกํ, อโนตฺตปฺปํ สมฺปยุตฺตกา จ ขนฺธา จิตฺต\-สมุฏฺฐานญฺจ รูปํ. (จกฺกํ.) (3)}

\proseitem{2}{โนกิเลสํ ธมฺมํ ปฏิจฺจ โนกิเลโส ธมฺโม อุปฺปชฺชติ เหตุปจฺจยา. โนกิเลสํ เอกํ ขนฺธํ ปฏิจฺจ ตโย ขนฺธา จิตฺต\-สมุฏฺฐานญฺจ รูปํ ฯเปฯ เทฺว ขนฺเธ ฯเปฯ. ปฏิสนฺธิ\-กฺขเณ ฯเปฯ. ขนฺเธ ปฏิจฺจ วตฺถุ, วตฺถุํ ปฏิจฺจ ขนฺธา. เอกํ มหาภูตํ ฯเปฯ. (1)}

\prose{โนกิเลสํ ธมฺมํ ปฏิจฺจ กิเลโส ธมฺโม อุปฺปชฺชติ เหตุปจฺจยา. โนกิเลเส ขนฺเธ ปฏิจฺจ กิเลสา. (2)}

\prose{โนกิเลสํ ธมฺมํ ปฏิจฺจ กิเลโส จ โนกิเลโส จ ธมฺมา อุปฺปชฺชนฺติ เหตุปจฺจยา. โนกิเลสํ เอกํ ขนฺธํ ปฏิจฺจ ตโย ขนฺธา กิเลสา จ จิตฺต\-สมุฏฺฐานญฺจ รูปํ ฯเปฯ เทฺว ขนฺเธ ฯเปฯ. (3)}

\proseitem{3}{กิเลสญฺจ โนกิเลสญฺจ ธมฺมํ ปฏิจฺจ กิเลโส ธมฺโม อุปฺปชฺชติ เหตุปจฺจยา. โลภญฺจ สมฺปยุตฺตเก จ ขนฺเธ ปฏิจฺจ โมโห, ทิฏฺฐิ, ถินํ, อุทฺธจฺจํ, อหิริกํ, อโนตฺตปฺปํ. (จกฺกํ.) (1)}

\prose{กิเลสญฺจ โนกิเลสญฺจ ธมฺมํ ปฏิจฺจ โนกิเลโส ธมฺโม อุปฺปชฺชติ เหตุปจฺจยา. โนกิเลสํ เอกํ ขนฺธญฺจ กิเลเส จ ปฏิจฺจ ตโย ขนฺธา จิตฺต\-สมุฏฺฐานญฺจ รูปํ ฯเปฯ เทฺว ขนฺเธ จ ฯเปฯ. กิเลเส จ มหาภูเต จ ปฏิจฺจ จิตฺต\-สมุฏฺฐานํ รูปํ. (2)}

\prose{กิเลสญฺจ โนกิเลสญฺจ ธมฺมํ ปฏิจฺจ กิเลโส จ โนกิเลโส จ ธมฺมา อุปฺปชฺชนฺติ เหตุปจฺจยา. โนกิเลสํ เอกํ ขนฺธญฺจ โลภญฺจ ปฏิจฺจ ตโย ขนฺธา โมโห, ทิฏฺฐิ, ถินํ, อุทฺธจฺจํ, อหิริกํ, อโนตฺตปฺปํ จิตฺต\-สมุฏฺฐานญฺจ รูปํ ฯเปฯ เทฺว ขนฺเธ จ ฯเปฯ. (จกฺกํ, สํขิตฺตํ.) (3)}

\csromanmarkh{1. ปจฺจยานุโลม}\csromanmarkhii{2. สงฺขฺยาวาร}\csromanheadernomark{2}{1. \textbf{ปจฺจยานุโลม 2. สงฺขฺยาวาร}}
\proseitem{4}{เหตุยา นว, อารมฺมเณ นว, อธิปติยา นว, อนนฺตเร นว, สมนนฺตเร นว, (สพฺพตฺถ นว,) วิปาเก เอกํ ฯเปฯ อวิคเต นว.}

\csromanmarkh{2. ปจฺจย\-ปจฺจนีย}\csromanmarkhii{1. วิภงฺควาร}\csromanheadernomark{2}{2. ปจฺจย\-ปจฺจนีย 1. วิภงฺควาร}
\proseitem{5}{\csromanfolio{3}กิเลสํ ธมฺมํ ปฏิจฺจ กิเลโส ธมฺโม อุปฺปชฺชติ นเหตุปจฺจยา. วิจิกิจฺฉํ ปฏิจฺจ โมโห. อุทฺธจฺจํ ปฏิจฺจ โมโห. (1)}

\prose{โนกิเลสํ ธมฺมํ ปฏิจฺจ โนกิเลโส ธมฺโม อุปฺปชฺชติ นเหตุปจฺจยา. อเหตุกํ โนกิเลสํ เอกํ ขนฺธํ ปฏิจฺจ ตโย ขนฺธา จิตฺต\-สมุฏฺฐานญฺจ รูปํ ฯเปฯ เทฺว ขนฺเธ ฯเปฯ. อเหตุก\-ปฏิสนฺธิ\-กฺขเณ ฯเปฯ. (ยาว อสญฺญสตฺตา.) (1)}

\prose{โนกิเลสํ ธมฺมํ ปฏิจฺจ กิเลโส ธมฺโม อุปฺปชฺชติ นเหตุปจฺจยา. วิจิกิจฺฉา\-สหคเต อุทฺธจฺจ\-สหคเต ขนฺเธ ปฏิจฺจ วิจิกิจฺฉา\-สหคโต อุทฺธจฺจ\-สหคโต โมโห. (2)}

\prose{กิเลสญฺจ โนกิเลสญฺจ ธมฺมํ ปฏิจฺจ กิเลโส ธมฺโม อุปฺปชฺชติ นเหตุปจฺจยา. วิจิกิจฺฉา\-สหคเต ขนฺเธ จ วิจิกิจฺฉญฺจ ปฏิจฺจ วิจิกิจฺฉา\-สหคโต โมโห. อุทฺธจฺจ\-สหคเต ขนฺเธ จ อุทฺธจฺจญฺจ ปฏิจฺจ อุทฺธจฺจ\-สหคโต โมโห. (1)}

\prose{นอารมฺมณาทิ}

\proseitem{6}{กิเลสํ ธมฺมํ ปฏิจฺจ โนกิเลโส ธมฺโม อุปฺปชฺชติ นอารมฺมณ\-ปจฺจยา. กิเลเส ปฏิจฺจ จิตฺต\-สมุฏฺฐานํ รูปํ. (1)}

\prose{โนกิเลสํ ธมฺมํ ปฏิจฺจ โนกิเลโส ธมฺโม อุปฺปชฺชติ นอารมฺมณ\-ปจฺจยา. โนกิเลเส ขนฺเธ ปฏิจฺจ จิตฺต\-สมุฏฺฐานํ รูปํ. ปฏิสนฺธิ\-กฺขเณ ฯเปฯ. ขนฺเธ ปฏิจฺจ วตฺถุ. (ยาว อสญฺญสตฺตา.) (1)}

\prose{กิเลสญฺจ โนกิเลสญฺจ ธมฺมํ ปฏิจฺจ โนกิเลโส ธมฺโม อุปฺปชฺชติ นอารมฺมณ\-ปจฺจยา. กิเลเส จ สมฺปยุตฺตเก จ ขนฺเธ ปฏิจฺจ จิตฺต\-สมุฏฺฐานํ รูปํ. กิเลเส จ มหาภูเต จ ปฏิจฺจ จิตฺต\-สมุฏฺฐานํ รูปํ. (1)}

\prose{นอธิปติ\-ปจฺจยา. นอนนฺตร\-ปจฺจยา. นสมนนฺตร\-ปจฺจยา. นอญฺญมญฺญ\-ปจฺจยา. นอุปนิสฺสยปจฺจยา.}

\csromanheader{2}{\textbf{นปุเรชาตาทิ}}
\proseitem{7}{\csromanfolio{4}กิเลสํ ธมฺมํ ปฏิจฺจ กิเลโส ธมฺโม อุปฺปชฺชติ นปุเรชาต\-ปจฺจยา. อรูเป โลภํ ปฏิจฺจ โมโห, ทิฏฺฐิ, ถินํ, อุทฺธจฺจํ, อหิริกํ, อโนตฺตปฺปํ. โลภํ ปฏิจฺจ โมโห, ทิฏฺฐิ, อุทฺธจฺจํ, อหิริกํ, อโนตฺตปฺปํ. โลภํ ปฏิจฺจ โมโห, มาโน, ถินํ, อุทฺธจฺจํ, อหิริกํ, อโนตฺตปฺปํ. โลภํ ปฏิจฺจ โมโห, มาโน, อุทฺธจฺจํ, อหิริกํ, อโนตฺตปฺปํ. โลภํ ปฏิจฺจ โมโห, ถินํ, อุทฺธจฺจํ, อหิริกํ, อโนตฺตปฺปํ. โลภํ ปฏิจฺจ โมโห, อุทฺธจฺจํ, อหิริกํ, อโนตฺตปฺปํ. วิจิกิจฺฉํ ปฏิจฺจ โมโห, อุทฺธจฺจํ, อหิริกํ, อโนตฺตปฺปํ. อุทฺธจฺจํ ปฏิจฺจ โมโห, อหิริกํ, อโนตฺตปฺปํ. (อรูเป โทสมูลกํ นตฺถิ.) (1)}

\prose{กิเลสํ ธมฺมํ ปฏิจฺจ โนกิเลโส ธมฺโม อุปฺปชฺชติ นปุเรชาต\-ปจฺจยา. อรูเป กิเลเส ปฏิจฺจ สมฺปยุตฺตกา ขนฺธา. กิเลเส ปฏิจฺจ จิตฺต\-สมุฏฺฐานํ รูปํ. (เอวํ นวปิ ปญฺหา กาตพฺพา.) นปจฺฉาชาต\-ปจฺจยา. นอาเสวน\-ปจฺจยา.}

\csromanheader{2}{\textbf{นกมฺม}}
\proseitem{8}{กิเลสํ ธมฺมํ ปฏิจฺจ โนกิเลโส ธมฺโม อุปฺปชฺชติ นกมฺมปจฺจยา. กิเลเส ปฏิจฺจ สมฺปยุตฺตกา เจตนา. (1)}

\prose{โนกิเลสํ ธมฺมํ ปฏิจฺจ โนกิเลโส ธมฺโม อุปฺปชฺชติ นกมฺมปจฺจยา. โนกิเลเส ขนฺเธ ปฏิจฺจ สมฺปยุตฺตกา เจตนา. พาหิรํ. อาหารสมุฏฺฐานํ. อุตุสมุฏฺฐานํ ฯเปฯ. (1)}

\prose{กิเลสญฺจ โนกิเลสญฺจ ธมฺมํ ปฏิจฺจ โนกิเลโส ธมฺโม อุปฺปชฺชติ นกมฺมปจฺจยา. กิเลเส จ สมฺปยุตฺตเก จ ขนฺเธ ปฏิจฺจ สมฺปยุตฺตกา เจตนา. (1) (เอวํ สพฺเพ ปจฺจยา กาตพฺพา.)}

\csromanmarkh{2. ปจฺจย\-ปจฺจนีย}\csromanmarkhii{2. สงฺขฺยาวาร}\csromanheadernomark{2}{2. ปจฺจย\-ปจฺจนีย 2. สงฺขฺยาวาร}
\proseitem{9}{นเหตุยา จตฺตาริ, นอารมฺมเณ ตีณิ, นอธิปติยา นว, นอนนฺตเร ตีณิ, นสมนนฺตเร ตีณิ, นอญฺญมญฺเญ ตีณิ, นอุปนิสฺสเย ตีณิ, นปุเรชาเต นว, นปจฺฉาชาเต นว, นอาเสวเน นว, นกมฺเม ตีณิ, \csromanfolio{5}นวิปาเก นว, นอาหาเร เอกํ, นอินฺทฺริเย เอกํ, นฌาเน เอกํ, นมคฺเค เอกํ, นสมฺปยุตฺเต ตีณิ, นวิปฺปยุตฺเต นว, โนนตฺถิยา ตีณิ, โนวิคเต ตีณิ.}

\vspace{\tipitakaparskip}
\csromancenter{(เอวํ อิตเร เทฺว คณนาปิ สหชาตวาโรปิ กาตพฺโพ.)}
\csromansectionrule
\csromanmarkh{75. กิเลสทุก}\csromanmarkhii{3. ปจฺจยวาร}\csromanheadernomark{2}{75. กิเลสทุก 3. ปจฺจยวาร}
\proseitem{9}{กิเลสํ ธมฺมํ ปจฺจยา กิเลโส ธมฺโม อุปฺปชฺชติ เหตุปจฺจยา. ตีณิ. (ปฏิจฺจสทิสา.)}

\prose{โนกิเลสํ ธมฺมํ ปจฺจยา โนกิเลโส ธมฺโม อุปฺปชฺชติ เหตุปจฺจยา. โนกิเลสํ เอกํ ขนฺธํ ปจฺจยา ตโย ขนฺธา ฯเปฯ. (ยาว อชฺฌตฺติกา มหาภูตา.) วตฺถุํ ปจฺจยา โนกิเลสา ขนฺธา. (1)}

\prose{โนกิเลสํ ธมฺมํ ปจฺจยา กิเลโส ธมฺโม อุปฺปชฺชติ เหตุปจฺจยา. โนกิเลเส ขนฺเธ ปจฺจยา กิเลสา. วตฺถุํ ปจฺจยา กิเลสา. (2)}

\prose{โนกิเลสํ ธมฺมํ ปจฺจยา กิเลโส จ โนกิเลโส จ ธมฺมา อุปฺปชฺชนฺติ เหตุปจฺจยา. โนกิเลสํ เอกํ ขนฺธํ ปจฺจยา ตโย ขนฺธา กิเลสา จ จิตฺต\-สมุฏฺฐานญฺจ รูปํ ฯเปฯ เทฺว ขนฺเธ ฯเปฯ. วตฺถุํ ปจฺจยา กิเลสา. มหาภูเต ปจฺจยา จิตฺต\-สมุฏฺฐานํ รูปํ. วตฺถุํ ปจฺจยา กิเลสา จ สมฺปยุตฺตกา จ ขนฺธา. (3)}

\proseitem{11}{กิเลสญฺจ โนกิเลสญฺจ ธมฺมํ ปจฺจยา กิเลโส ธมฺโม อุปฺปชฺชติ เหตุปจฺจยา. โลภญฺจ สมฺปยุตฺตเก จ ขนฺเธ ปจฺจยา โมโห. ทิฏฺฐิ, ถินํ, อุทฺธจฺจํ, อหิริกํ, อโนตฺตปฺปํ. (จกฺกํ.) โลภญฺจ วตฺถุญฺจ ปจฺจยา กิเลสา. (1)}

\prose{กิเลสญฺจ โนกิเลสญฺจ ธมฺมํ ปจฺจยา โนกิเลโส ธมฺโม อุปฺปชฺชติ เหตุปจฺจยา. โนกิเลสํ เอกํ ขนฺธญฺจ กิเลสญฺจ ปจฺจยา ตโย ขนฺธา จิตฺต\-สมุฏฺฐานญฺจ รูปํ ฯเปฯ เทฺว ขนฺเธ จ ฯเปฯ. กิเลเส จ มหาภูเต จ ปจฺจยา จิตฺต\-สมุฏฺฐานํ รูปํ. กิเลเส จ วตฺถุญฺจ ปจฺจยา โนกิเลสา ขนฺธา. (2)}

\prose{\csromanfolio{6}กิเลสญฺจ โนกิเลสญฺจ ธมฺมํ ปจฺจยา กิเลโส จ โนกิเลโส จ ธมฺมา อุปฺปชฺชนฺติ เหตุปจฺจยา. โนกิเลสํ เอกํ ขนฺธญฺจ โลภญฺจ ปจฺจยา ตโย ขนฺธา โมโห, ทิฏฺฐิ, ถินํ, อุทฺธจฺจํ, อหิริกํ, อโนตฺตปฺปํ, จิตฺต\-สมุฏฺฐานญฺจ รูปํ ฯเปฯ เทฺว ขนฺเธ จ ฯเปฯ. (จกฺกํ.) โลภญฺจ วตฺถุญฺจ ปจฺจยา โมโห, ทิฏฺฐิ, ถินํ, อุทฺธจฺจํ, อหิริกํ, อโนตฺตปฺปํ สมฺปยุตฺตกา จ ขนฺธา. (จกฺกํ.) (3) (อารมฺมณ\-ปจฺจเย โนกิเลสมูเล ปญฺจ วิญฺญาณา กาตพฺพา.)}

\csromanmarkh{1. ปจฺจยานุโลม}\csromanmarkhii{2. สงฺขฺยาวาร}\csromanheadernomark{2}{1. \textbf{ปจฺจยานุโลม 2. สงฺขฺยาวาร}}
\proseitem{12}{เหตุยา นว, อารมฺมเณ นว, อธิปติยา นว, (สพฺพตฺถ นว,) วิปาเก เอกํ ฯเปฯ อวิคเต นว.}

\csromanmarkh{2. ปจฺจย\-ปจฺจนีย}\csromanmarkhii{1. วิภงฺควาร}\csromanheadernomark{2}{2. ปจฺจย\-ปจฺจนีย 1. วิภงฺควาร}
\proseitem{13}{กิเลสํ ธมฺมํ ปจฺจยา กิเลโส ธมฺโม อุปฺปชฺชติ นเหตุปจฺจยา. วิจิกิจฺฉํ ปจฺจยา วิจิกิจฺฉา\-สหคโต โมโห. อุทฺธจฺจํ ปจฺจยา อุทฺธจฺจหคโต โมโห. (1)}

\prose{โนกิเลสํ ธมฺมํ ปจฺจยา โนกิเลโส ธมฺโม อุปฺปชฺชติ นเหตุปจฺจยา. (ยาว อสญฺญสตฺตา.) จกฺขายตนํ ปจฺจยา จกฺขุวิญฺญาณํ ฯเปฯ. กายายตนํ ฯเปฯ. วตฺถุํ ปจฺจยา อเหตุกา โนกิเลสา ขนฺธา. (1)}

\prose{โนกิเลสํ ธมฺมํ ปจฺจยา กิเลโส ธมฺโม อุปฺปชฺชติ นเหตุปจฺจยา. วิจิกิจฺฉาคเต อุทฺธจฺจ\-สหคเต ขนฺเธ จ วตฺถุญฺจ ปจฺจยา วิจิกิจฺฉา\-สหคโต อุทฺธจฺจ\-สหคโต โมโห. (2)}

\prose{กิเลสญฺจ โนกิเลสญฺจ ธมฺมํ ปจฺจยา กิเลโส ธมฺโม อุปฺปชฺชติ นเหตุปจฺจยา. วิจิกิจฺฉญฺจ สมฺปยุตฺตเก จ ขนฺเธ วตฺถุญฺจ ปจฺจยา วิจิกิจฺฉา\-สหคโต โมโห. อุทฺธจฺจญฺจ สมฺปยุตฺตเก จ ขนฺเธ วตฺถุญฺจ ปจฺจยา อุทฺธจฺจ\-สหคโต โมโห. (สํขิตฺตํ.) (1)}

\csromanmarkh{2. ปจฺจย\-ปจฺจนีย}\csromanmarkhii{2. สงฺขฺยาวาร}\csromanheadernomark{2}{2. ปจฺจย\-ปจฺจนีย 2. สงฺขฺยาวาร}
\proseitem{14}{\csromanfolio{7}นเหตุยา จตฺตาริ, นอารมฺมเณ ตีณิ, นอธิปติยา นว ฯเปฯ นกมฺเม ตีณิ, นวิปาเก นว, นอาหาเร เอกํ ฯเปฯ โนวิคเต ตีณิ.}

\vspace{\tipitakaparskip}
\csromancenter{(เอวํ อิตเร เทฺว คณนาปิ นิสฺสยวาโรปิ กาตพฺโพ.)}
\csromanmarkh{75. กิเลสทุก}\csromanmarkhii{5. สํสฏฺฐวาร}\csromanheadernomark{2}{75. กิเลสทุก 5. สํสฏฺฐวาร}
\proseitem{15}{กิเลสํ ธมฺมํ สํสฏฺโฐ กิเลโส ธมฺโม อุปฺปชฺชติ เหตุปจฺจยา. โลภํ สํสฏฺโฐ โมโห, ทิฏฺฐิ, ถินํ, อุทฺธจฺจํ, อหิริกํ, อโนตฺตปฺปํ. (จกฺกํ.) (เอวํ นว ปญฺหา กาตพฺพา.)}

\prose{เหตุยา นว, อารมฺมเณ นว, (สพฺพตฺถ นว,) วิปาเก เอกํ ฯเปฯ อวิคเต นว. อนุโลมํ.}

\prose{กิเลสํ ธมฺมํ สํสฏฺโฐ กิเลโส ธมฺโม อุปฺปชฺชติ นเหตุปจฺจยา. (เอวํ นเหตุปญฺหา จตฺตาริ กาตพฺพา.)}

\prose{นเหตุยา จตฺตาริ, นอธิปติยา นว, นปุเรชาเต นว, นปจฺฉาชาเต นว, นอาเสวเน นว, นกมฺเม ตีณิ, นวิปาเก นว, นฌาเน เอกํ, นมคฺเค เอกํ, นวิปฺปยุตฺเต นว.}

\vspace{\tipitakaparskip}
\csromancenter{ปจฺจนียํ.}
\csromancenter{(เอวํ อิตเร เทฺว คณนาปิ สมฺปยุตฺตวาโรปิ กาตพฺโพ.)}
\csromanmarkh{75. กิเลสทุก}\csromanmarkhii{7. ปญฺหาวาร}\csromanheadernomark{2}{75. กิเลสทุก 7. ปญฺหาวาร}
\proseitem{16}{กิเลโส ธมฺโม กิเลสสฺส ธมฺมสฺส เหตุปจฺจเยน ปจฺจโย. กิเลสา เหตู สมฺปยุตฺตกานํ กิเลสานํ เหตุปจฺจเยน ปจฺจโย. \csromanfolio{8}(มูลํ ปุจฺฉิตพฺพํ.) กิเลสา เหตู สมฺปยุตฺตกานํ ขนฺธานํ จิตฺตสมุฏฺฐานานญฺจ รูปานํ เหตุปจฺจเยน ปจฺจโย. (มูลํ ปุจฺฉิตพฺพํ.) กิเลสา เหตู สมฺปยุตฺตกานํ ขนฺธานํ กิเลสานํ จิตฺตสมุฏฺฐานานญฺจ รูปานํ เหตุปจฺจเยน ปจฺจโย. (3)}

\prose{โนกิเลโส ธมฺโม โนกิเลสสฺส ธมฺมสฺส เหตุปจฺจเยน ปจฺจโย. โนกิเลสา เหตู สมฺปยุตฺตกานํ ขนฺธานํ จิตฺตสมุฏฺฐานานญฺจ รูปานํ เหตุปจฺจเยน ปจฺจโย. ปฏิสนฺธิ\-กฺขเณ ฯเปฯ. (1)}

\proseitem{17}{กิเลโส ธมฺโม กิเลสสฺส ธมฺมสฺส อารมฺมณ\-ปจฺจเยน ปจฺจโย. กิเลเส อารพฺภ กิเลสา อุปฺปชฺชนฺติ. (มูลํ ปุจฺฉิตพฺพํ.) กิเลเส อารพฺภ โนกิเลสา ขนฺธา อุปฺปชฺชนฺติ. (มูลํ ปุจฺฉิตพฺพํ.) กิเลเส อารพฺภ กิเลสา จ สมฺปยุตฺตกา จ ขนฺธา อุปฺปชฺชนฺติ. (3)}

\proseitem{18}{โนกิเลโส ธมฺโม โนกิเลสสฺส ธมฺมสฺส อารมฺมณ\-ปจฺจเยน ปจฺจโย. ทานํ ฯเปฯ สีลํ ฯเปฯ อุโปสถกมฺมํ ฯเปฯ. ปุพฺเพ สุจิณฺณานิ ฯเปฯ. ฌานา วุฏฺฐหิตฺวา ฌานํ ปจฺจเวกฺขติ, อสฺสาเทติ อภินนฺทติ, ตํ อารพฺภ ราโค ฯเปฯ ทิฏฺฐิ ฯเปฯ วิจิกิจฺฉา ฯเปฯ อุทฺธจฺจํ ฯเปฯ โทมนสฺสํ อุปฺปชฺชติ. อริยา มคฺคา วุฏฺฐหิตฺวา ฯเปฯ ผลสฺส, อาวชฺชนาย อารมฺมณ\-ปจฺจเยน ปจฺจโย. จกฺขุํ ฯเปฯ. วตฺถุํ โนกิเลเส ขนฺเธ อนิจฺจโต ฯเปฯ โทมนสฺสํ อุปฺปชฺชติ. ทิพฺเพน จกฺขุนา รูปํ ปสฺสติ. ทิพฺพาย โสตธาตุยา สทฺทํ สุณาติ ฯเปฯ. อนาคตํสญาณสฺส, อาวชฺชนาย อารมฺมณ\-ปจฺจเยน ปจฺจโย. (1)}

\prose{โนกิเลโส ธมฺโม กิเลสสฺส ธมฺมสฺส อารมฺมณ\-ปจฺจเยน ปจฺจโย. ทานํ ฯเปฯ. ฌานา วุฏฺฐหิตฺวา ฌานํ อสฺสาเทติ อภินนฺทติ, ตํ อารพฺภ ราโค ฯเปฯ ทิฏฺฐิ ฯเปฯ วิจิกิจฺฉา ฯเปฯ อุทฺธจฺจํ อุปฺปชฺชติ. ฌาเน ปริหีเน วิปฺปฏิสาริสฺส โทมนสฺสํ อุปฺปชฺชติ. จกฺขุํ ฯเปฯ. วตฺถุํ โนกิเลเส ขนฺเธ อสฺสาเทติ อภินนฺทติ, ตํ อารพฺภ ราโค ฯเปฯ โทมนสฺสํ อุปฺปชฺชติ. (2)}

\prose{โนกิเลโส ธมฺโม กิเลสสฺส จ โนกิเลสสฺส จ ธมฺมสฺส อารมฺมณ\-ปจฺจเยน ปจฺจโย. ทานํ ฯเปฯ. ฌานา วุฏฺฐหิตฺวา ฯเปฯ. จกฺขุํ ฯเปฯ. วตฺถุํ โนกิเลเส ขนฺเธ อสฺสาเทติ อภินนฺทติ, ตํ อารพฺภ กิเลสา จ สมฺปยุตฺตกา จ ขนฺธา อุปฺปชฺชนฺติ. (3)}

\prose{\csromanfolio{9}กิเลโส จ โนกิเลโส จ ธมฺมา กิเลสสฺส ธมฺมสฺส อารมฺมณ\-ปจฺจเยน ปจฺจโย. ตีณิ.}

\csromanheader{2}{\textbf{อธิปติ}}
\proseitem{19}{กิเลโส ธมฺโม กิเลสสฺส ธมฺมสฺส อธิปติ\-ปจฺจเยน ปจฺจโย. อารมฺมณา\-ธิปติ\csromandash{}กิเลเส ครุํ กตฺวา กิเลสา อุปฺปชฺชนฺติ. ตีณิ. (อารมฺมณา\-ธิปติเยว.) (3)}

\prose{โนกิเลโส ธมฺโม โนกิเลสสฺส ธมฺมสฺส อธิปติ\-ปจฺจเยน ปจฺจโย. อารมฺมณา\-ธิปติ, สหชาตา\-ธิปติ.  อารมฺมณา\-ธิปติ\csromandash{}ทานํ ฯเปฯ สีลํ ฯเปฯ อุโปสถกมฺมํ กตฺวา ตํ ครุํ กตฺวา ปจฺจเวกฺขติ, อสฺสาเทติ อภินนฺทติ. ตํ ครุํ กตฺวา ราโค อุปฺปชฺชติ. ทิฏฺฐิ อุปฺปชฺชติ. ปุพฺเพ ฯเปฯ. ฌานา ฯเปฯ. อริยา มคฺคา วุฏฺฐหิตฺวา มคฺคํ ครุํ กตฺวา ปจฺจเวกฺขนฺติ ฯเปฯ ผลสฺส อธิปติ\-ปจฺจเยน ปจฺจโย. จกฺขุํ ฯเปฯ. วตฺถุํ โนกิเลเส ขนฺเธ ครุํ กตฺวา อสฺสาเทติ อภินนฺทติ, ตํ ครุํ กตฺวา ราโค อุปฺปชฺชติ. ทิฏฺฐิ อุปฺปชฺชติ. สหชาตา\-ธิปติ\csromandash{}โนกิเลสา\-ธิปติ สมฺปยุตฺตกานํ ขนฺธานํ จิตฺตสมุฏฺฐานานญฺจ รูปานํ อธิปติ\-ปจฺจเยน ปจฺจโย. (1)}

\prose{โนกิเลโส ธมฺโม กิเลสสฺส ธมฺมสฺส อธิปติ\-ปจฺจเยน ปจฺจโย. อารมฺมณา\-ธิปติ, สหชาตา\-ธิปติ.  อารมฺมนาธิปติ\csromandash{}ทานํ ฯเปฯ. ฌานํ ฯเปฯ. จกฺขุํ ฯเปฯ. วตฺถุํ โนกิเลเส ขนฺเธ ครุํ กตฺวา อสฺสาเทติ อภินนฺทติ, ตํ ครุํ กตฺวา ราโค อุปฺปชฺชติ. ทิฏฺฐิ อุปฺปชฺชติ. สหชาตา\-ธิปติ\csromandash{}โนกิเลสา\-ธิปติ สมฺปยุตฺตกานํ กิเลสานํ อธิปติ\-ปจฺจเยน ปจฺจโย. (2)}

\prose{โนกิเลโส ธมฺโม กิเลสสฺส จ โนกิเลสสฺส จ ธมฺมสฺส อธิปติ\-ปจฺจเยน ปจฺจโย. อารมฺมณา\-ธิปติ, สหชาตา\-ธิปติ.  อารมฺมณา\-ธิปติ\csromandash{}ทานํ ฯเปฯ. ฌานํ ฯเปฯ. จกฺขุํ ฯเปฯ. วตฺถุํ โนกิเลเส ขนฺเธ ครุํ กตฺวา อสฺสาเทติ อภินนฺทติ, ตํ ครุํ กตฺวา กิเลสา จ สมฺปยุตฺตกา จ ขนฺธา อุปฺปชฺชนฺติ. สหชาตา\-ธิปติ\csromandash{}โนกิเลสา\-ธิปติ สมฺปยุตฺตกานํ ขนฺธานํ กิเลสานญฺจ จิตฺตสมุฏฺฐานานญฺจ รูปานํ อธิปติ\-ปจฺจเยน ปจฺจโย. (3)}

\prose{กิเลโส จ โนกิเลโส จ ธมฺมา กิเลสสฺส ธมฺมสฺส อธิปติ\-ปจฺจเยน ปจฺจโย. ตีณิ. (อารมฺมณา\-ธิปติเยว.) (3)}

\proseitem{20}{\csromanfolio{10}กิเลโส ธมฺโม กิเลสสฺส ธมฺมสฺส อนนฺตร\-ปจฺจเยน ปจฺจโย. ปุริมา ปุริมา กิเลสา ปจฺฉิมานํ ปจฺฉิมานํ กิเลสานํ อนนฺตร\-ปจฺจเยน ปจฺจโย. (มูลํ กาตพฺพํ.) ปุริมา ปุริมา กิเลสา ปจฺฉิมานํ ปจฺฉิมานํ โนกิเลสานํ ขนฺธานํ อนนฺตร\-ปจฺจเยน ปจฺจโย. กิเลสา วุฏฺฐานสฺส อนนฺตร\-ปจฺจเยน ปจฺจโย. (มูลํ กาตพฺพํ.) ปุริมา ปุริมา กิเลสา ปจฺฉิมานํ ปจฺฉิมานํ กิเลสานํ สมฺปยุตฺตกานญฺจ ขนฺธานํ อนนฺตร\-ปจฺจเยน ปจฺจโย. (3)}

\proseitem{21}{โนกิเลโส ธมฺโม โนกิเลสสฺส ธมฺมสฺส อนนฺตร\-ปจฺจเยน ปจฺจโย. ปุริมา ปุริมา โนกิเลสา ขนฺธา ปจฺฉิมานํ ปจฺฉิมานํ โนกิเลสานํ ขนฺธานํ อนนฺตร\-ปจฺจเยน ปจฺจโย ฯเปฯ ผลสมาปตฺติยา อนนฺตร\-ปจฺจเยน ปจฺจโย. (มูลํ กาตพฺพํ.) ปุริมา ปุริมา โนกิเลสา ขนฺธา ปจฺฉิมานํ ปจฺฉิมานํ กิเลสานํ อนนฺตร\-ปจฺจเยน ปจฺจโย. อาวชฺชนา กิเลสานํ อนนฺตร\-ปจฺจเยน ปจฺจโย. (มูลํ กาตพฺพํ.) ปุริมา ปุริมา โนกิเลสา ขนฺธา ปจฺฉิมานํ ปจฺฉิมานํ กิเลสานํ สมฺปยุตฺตกานญฺจ ขนฺธานํ อนนฺตร\-ปจฺจเยน ปจฺจโย. อาวชฺชนา กิเลสานํ สมฺปยุตฺตกานญฺจ ขนฺธานํ อนนฺตร\-ปจฺจเยน ปจฺจโย. (3)}

\proseitem{22}{กิเลโส จ โนกิเลโส จ ธมฺมา กิเลสสฺส ธมฺมสฺส อนนฺตร\-ปจฺจเยน ปจฺจโย. ปุริมา ปุริมา กิเลสา จ สมฺปยุตฺตกา จ ขนฺธา ปจฺฉิมานํ ปจฺฉิมานํ กิเลสานํ อนนฺตร\-ปจฺจเยน ปจฺจโย. (มูลํ กาตพฺพํ.) ปุริมา ปุริมา กิเลสา จ สมฺปยุตฺตกา จ ขนฺธา ปจฺฉิมานํ ปจฺฉิมานํ โนกิเลสานํ ขนฺธานํ อนนฺตร\-ปจฺจเยน ปจฺจโย. กิเลสา จ สมฺปยุตฺตกา จ ขนฺธา วุฏฺฐานสฺส อนนฺตร\-ปจฺจเยน ปจฺจโย. (มูลํ กาตพฺพํ.) ปุริมา ปุริมา กิเลสา จ สมฺปยุตฺตกา จ ขนฺธา ปจฺฉิมานํ ปจฺฉิมานํ กิเลสานํ สมฺปยุตฺตกานญฺจ ขนฺธานํ อนนฺตร\-ปจฺจเยน ปจฺจโย. (3)}

\prose{สมนนฺตร\-ปจฺจเยน ปจฺจโย. สหชาต\-ปจฺจเยน ปจฺจโย. อญฺญมญฺญ\-ปจฺจเยน ปจฺจโย. นิสฺสย\-ปจฺจเยน ปจฺจโย.}

\csromanheader{2}{\textbf{อุปนิสฺสย}}
\proseitem{23}{\csromanfolio{11}กิเลโส ธมฺโม กิเลสสฺส ธมฺมสฺส อุปนิสฺสย\-ปจฺจเยน ปจฺจโย. อารมฺมณู\-ปนิสฺสโย, อนนฺตรู\-ปนิสฺสโย, ปกตู\-ปนิสฺสโย ฯเปฯ. ปกตู\-ปนิสฺสโย\csromandash{}กิเลสา กิเลสานํ อุปนิสฺสย\-ปจฺจเยน ปจฺจโย. ตีณิ.}

\proseitem{24}{โนกิเลโส ธมฺโม โนกิเลสสฺส ธมฺมสฺส อุปนิสฺสย\-ปจฺจเยน ปจฺจโย. อารมฺมณู\-ปนิสฺสโย, อนนฺตรู\-ปนิสฺสโย, ปกตู\-ปนิสฺสโย ฯเปฯ. ปกตู\-ปนิสฺสโย\csromandash{}สทฺธํ อุปนิสฺสาย ทานํ เทติ ฯเปฯ. มานํ ชปฺเปติ. ทิฏฺฐิํ คณฺหาติ. สีลํ ฯเปฯ. ปญฺญํ. ราคํ. โทสํ. โมหํ. มานํ. ทิฏฺฐิํ. ปตฺถนํ. กายิกํ สุขํ ฯเปฯ. เสนาสนํ อุปนิสฺสาย ทานํ เทติ ฯเปฯ สํฆํ ภินฺทติ. สทฺธา ฯเปฯ. เสนาสนํ สทฺธาย ฯเปฯ ผลสมาปตฺติยา อุปนิสฺสย\-ปจฺจเยน ปจฺจโย. (1)}

\prose{โนกิเลโส ธมฺโม กิเลสสฺส ธมฺมสฺส อุปนิสฺสย\-ปจฺจเยน ปจฺจโย. อารมฺมณู\-ปนิสฺสโย, อนนฺตรู\-ปนิสฺสโย, ปกตู\-ปนิสฺสโย ฯเปฯ. ปกตู\-ปนิสฺสโย\csromandash{}สทฺธํ อุปนิสฺสาย มานํ ชปฺเปติ. ทิฏฺฐิํ คณฺหาติ. สีลํ ฯเปฯ. เสนาสนํ อุปนิสฺสาย ปาณํ หนติ ฯเปฯ สํฆํ ภินฺทติ, สทฺธา ฯเปฯ. เสนาสนํ กิเลสานํ อุปนิสฺสย\-ปจฺจเยน ปจฺจโย. (2)}

\prose{โนกิเลโส ธมฺโม กิเลสสฺส จ โนกิเลสสฺส จ ธมฺมสฺส อุปนิสฺสย\-ปจฺจเยน ปจฺจโย. อารมฺมณู\-ปนิสฺสโย, อนนฺตรู\-ปนิสฺสโย, ปกตู\-ปนิสฺสโย ฯเปฯ. ปกตู\-ปนิสฺสโย\csromandash{}สทฺธํ อุปนิสฺสาย มานํ ชปฺเปติ. ทิฏฺฐิํ คณฺหาติ. สีลํ ฯเปฯ. เสนาสนํ อุปนิสฺสาย ปาณํ หนติ ฯเปฯ สํฆํ ภินฺทติ. สทฺธา ฯเปฯ. เสนาสนํ กิเลสานํ สมฺปยุตฺตกานญฺจ ขนฺธานํ อุปนิสฺสย\-ปจฺจเยน ปจฺจโย. (3)}

\prose{กิเลโส จ โนกิเลโส จ ธมฺมา กิเลสสฺส ธมฺมสฺส อุปนิสฺสย\-ปจฺจเยน ปจฺจโย. ตีณิ.}

\csromanheader{2}{\textbf{ปุเรชาต}}
\proseitem{25}{โนกิเลโส ธมฺโม โนกิเลสสฺส ธมฺมสฺส ปุเรชาต\-ปจฺจเยน ปจฺจโย. อารมฺมณ\-ปุเรชาตํ, วตฺถุ\-ปุเรชาตํ.  อารมฺมณ\-ปุเรชาตํ\csromandash{} จกฺขุํ ฯเปฯ. วตฺถุํ อนิจฺจโต ฯเปฯ โทมนสฺสํ อุปฺปชฺชติ. ทิพฺเพน \csromanfolio{12}จกฺขุนา รูปํ ปสฺสติ. ทิพฺพาย โสตธาตุยา สทฺทํ สุณาติ. รูปายตนํ จกฺขุ\-วิญฺญาณสฺส ฯเปฯ. โผฏฺฐพฺพา\-ยตนํ กายวิญฺญาณสฺส ฯเปฯ. วตฺถุ\-ปุเรชาตํ\csromandash{}จกฺขายตนํ จกฺขุ\-วิญฺญาณสฺส ฯเปฯ. กายายตนํ กายวิญฺญาณสฺส ฯเปฯ. วตฺถุ โนกิเลสานํ ขนฺธานํ ปุเรชาต\-ปจฺจเยน ปจฺจโย. (1)}

\prose{โนกิเลโส ธมฺโม กิเลสสฺส ธมฺมสฺส ปุเรชาต\-ปจฺจเยน ปจฺจโย. อารมฺมณ\-ปุเรชาตํ, วตฺถุ\-ปุเรชาตํ.  อารมฺมณ\-ปุเรชาตํ\csromandash{}จกฺขุํ ฯเปฯ. วตฺถุํ อสฺสาเทติ อภินนฺทติ, ตํ อารพฺภ ราโค ฯเปฯ โทมนสฺสํ อุปฺปชฺชติ. วตฺถุ\-ปุเรชาตํ\csromandash{}วตฺถุ กิเลสานํ ปุเรชาต\-ปจฺจเยน ปจฺจโย. (2)}

\prose{โนกิเลโส ธมฺโม กิเลสสฺส จ โนกิเลสสฺส จ ธมฺมสฺส ปุเรชาต\-ปจฺจเยน ปจฺจโย. อารมฺมณ\-ปุเรชาตํ, วตฺถุ\-ปุเรชาตํ.  อารมฺมณ\-ปุเรชาตํ\csromandash{}จกฺขุํ ฯเปฯ. วตฺถุํ อสฺสาเทติ อภินนฺทติ, ตํ อารพฺภ กิเลสา จ สมฺปยุตฺตกา จ ขนฺธา อุปฺปชฺชนฺติ. วตฺถุ\-ปุเรชาตํ\csromandash{} วตฺถุ กิเลสานํ สมฺปยุตฺตกานญฺจ ขนฺธานํ ปุเรชาต\-ปจฺจเยน ปจฺจโย. (3)}

\csromanheader{2}{ปจฺฉาชาตา\-เสวน}
\proseitem{26}{กิเลโส ธมฺโม โนกิเลสสฺส ธมฺมสฺส ปจฺฉาชาต\-ปจฺจเยน ปจฺจโย. (สํขิตฺตํ.) (1)}

\prose{โนกิเลโส ธมฺโม โนกิเลสสฺส ธมฺมสฺส ปจฺฉาชาต\-ปจฺจเยน ปจฺจโย. (สํขิตฺตํ.) (1)}

\prose{กิเลโส จ โนกิเลโส จ ธมฺมา โนกิเลสสฺส ธมฺมสฺส ปจฺฉาชาต\-ปจฺจเยน ปจฺจโย. (สํขิตฺตํ.) อาเสวน\-ปจฺจเยน ปจฺจโย. นว.}

\csromanheader{2}{\textbf{กมฺม}}
\proseitem{27}{โนกิเลโส ธมฺโม โนกิเลสสฺส ธมฺมสฺส กมฺมปจฺจเยน ปจฺจโย. สหชาตา, นานากฺขณิกา.  สหชาตา โนกิเลสา เจตนา สมฺปยุตฺตกานํ ขนฺธานํ จิตฺตสมุฏฺฐานานญฺจ รูปานํ กมฺมปจฺจเยน ปจฺจโย. ปฏิสนฺธิ\-กฺขเณ ฯเปฯ. นานากฺขณิกา โนกิเลสา เจตนา วิปากานํ ขนฺธานํ กฏตฺตา จ รูปานํ กมฺมปจฺจเยน ปจฺจโย. (1)}

\prose{\csromanfolio{13}โนกิเลโส ธมฺโม กิเลสสฺส ธมฺมสฺส กมฺมปจฺจเยน ปจฺจโย. โนกิเลสา เจตนา กิเลสานํ กมฺมปจฺจเยน ปจฺจโย. (2)}

\prose{โนกิเลโส ธมฺโม กิเลสสฺส จ โนกิเลสสฺส จ ธมฺมสฺส กมฺมปจฺจเยน ปจฺจโย. โนกิเลสา เจตนา สมฺปยุตฺตกานํ ขนฺธานํ กิเลสานํ จิตฺตสมุฏฺฐานานญฺจ รูปานํ กมฺมปจฺจเยน ปจฺจโย. (3)}

\csromanheader{2}{\textbf{วิปากาทิ}}
\proseitem{28}{โนกิเลโส ธมฺโม โนกิเลสสฺส ธมฺมสฺส วิปาก\-ปจฺจเยน ปจฺจโย. เอกํ. อาหารปจฺจเยน ปจฺจโย. ตีณิ. อินฺทฺริย\-ปจฺจเยน ปจฺจโย. ตีณิ. ฌานปจฺจเยน ปจฺจโย. ตีณิ. มคฺคปจฺจเยน ปจฺจโย. นว. สมฺปยุตฺต\-ปจฺจเยน ปจฺจโย. นว.}

\csromanheader{2}{\textbf{วิปยุตฺต}}
\proseitem{29}{กิเลโส ธมฺโม โนกิเลสสฺส ธมฺมสฺส วิปฺปยุตฺต\-ปจฺจเยน ปจฺจโย. สหชาตํ, ปจฺฉาชาตํ. (สํขิตฺตํ.) (1)}

\prose{โนกิเลโส ธมฺโม โนกิเลสสฺส ธมฺมสฺส วิปฺปยุตฺต\-ปจฺจเยน ปจฺจโย. สหชาตํ, ปุเรชาตํ, ปจฺฉาชาตํ. (สํขิตฺตํ.) (1)}

\prose{โนกิเลโส ธมฺโม กิเลสสฺส ธมฺมสฺส วิปฺปยุตฺต\-ปจฺจเยน ปจฺจโย. ปุเรชาตํ วตฺถุ กิเลสานํ วิปฺปยุตฺต\-ปจฺจเยน ปจฺจโย. (2)}

\prose{โนกิเลโส ธมฺโม กิเลสสฺส จ โนกิเลสสฺส จ ธมฺมสฺส วิปฺปยุตฺต\-ปจฺจเยน ปจฺจโย. ปุเรชาตํ วตฺถุ กิเลสานํ สมฺปยุตฺตกานญฺจ ขนฺธานํ วิปฺปยุตฺต\-ปจฺจเยน ปจฺจโย. (3)}

\prose{กิเลโส จ โนกิเลโส จ ธมฺมา โนกิเลสสฺส ธมฺมสฺส วิปฺปยุตฺตปจเยน ปจฺจโย. สหชาตํ, ปจฺฉาชตํ. (สํขิตฺตํ. วิตฺถาเรตพฺพํ.)}

\prose{อตฺถิอาทิ}

\proseitem{30}{กิเลโส ธมฺโม กิเลสสฺส ธมฺมสฺส อตฺถิปจฺจเยน ปจฺจโย. เอกํ. (ปฏิจฺจสทิสํ.) (1)}

\prose{กิเลโส ธมฺโม โนกิเลสสฺส ธมฺมสฺส อตฺถิปจฺจเยน ปจฺจโย. สหชาตํ, ปจฺฉาชาตํ. (สํขิตฺตํ.) (2)}

\prose{\csromanfolio{14}กิเลโส ธมฺโม กิเลสสฺส จ โนกิเลสสฺส จ ธมฺมสฺส อตฺถิปจฺจเยน ปจฺจโย. (ปฏิจฺจสทิสํ.) (3)}

\prose{โนกิเลโส ธมฺโม โนกิเลสสฺส ธมฺมสฺส อตฺถิปจฺจเยน ปจฺจโย. สหชาตํ, ปุเรชาตํ, ปจฺฉาชาตํ, อาหารํ, อินฺทฺริยํ. (สํขิตฺตํ.) (1)}

\prose{โนกิเลโส ธมฺโม กิเลสสฺส ธมฺมสฺส อตฺถิปจฺจเยน ปจฺจโย. สหชาตํ, ปุเรชาตํ. (สํขิตฺตํ. สหชาตํ สหชาต\-สทิสํ. ปุเรชาตํ ปุเรชาต\-สทิสํ).  โนกิเลโส ธมฺโม กิเลสสฺส จ โนกิเลสสฺส จ ธมฺมสฺส อตฺถิปจฺจเยน ปจฺจโย. สหชาตํ, ปุเรชาตํ. (สหชาตํ สหชาต\-สทิสํ. ปุเรชาตํ ปุเรชาต\-สทิสํ.) (3)}

\proseitem{31}{กิเลโส จ โนกิเลโส จ ธมฺมา กิเลสสฺส ธมฺมสฺส อตฺถิปจฺจเยน ปจฺจโย. สหชาตํ, ปุเรชาตํ.  \textbf{สหชาโต} โลโภ จ สมฺปยุตฺตกา จ ขนฺธา โมหสฺส, ทิฏฺฐิยา, ถินสฺส, อุทฺธจฺจสฺส, อหิริกสฺส, อโนตฺตปฺปสฺส อตฺถิปจฺจเยน ปจฺจโย. (จกฺกํ.) \textbf{สหชาโต} โลโภ จ วตฺถุ จ โมหสฺส, ทิฏฺฐิยา, ถินสฺส, อุทฺธจฺจสฺส, อหิริกสฺส, อโนตฺตปฺปสฺส อตฺถิปจฺจเยน ปจฺจโย. (จกฺกํ.) (1)}

\prose{กิเลโส จ โนกิเลโส จ ธมฺมา โนกิเลสสฺส ธมฺมสฺส อตฺถิปจฺจเยน ปจฺจโย. สหชาตํ, ปุเรชาตํ, ปจฺฉาชาตํ, อาหารํ, อินฺทฺริยํ.  สหชาโต โนกิเลโส เอโก ขนฺโธ จ กิเลโส จ ติณฺณนฺนํ ขนฺธานํ จิตฺตสมุฏฺฐานานญฺจ รูปานํ อตฺถิปจฺจเยน ปจฺจโย ฯเปฯ เทฺว ขนฺธา จ ฯเปฯ. สหชาตา กิเลสา จ มหาภูตา จ จิตฺต\-สมุฏฺฐานานํ รูปานํ อตฺถิปจฺจเยน ปจฺจโย. สหชาตา กิเลสา จ วตฺถุ จ โนกิเลสานํ ขนฺธานํ อตฺถิปจฺจเยน ปจฺจโย. ปจฺฉาชาตา กิเลสา จ สมฺปยุตฺตกา ขนฺธา จ ปุเรชาตสฺส อิมสฺส กายสฺส อตฺถิปจฺจเยน ปจฺจโย. ปจฺฉาชาตา กิเลสา จ สมฺปยุตฺตกา จ ขนฺธา กพฬีกาโร อาหาโร จ อิมสฺส กายสฺส อตฺถิปจฺจเยน ปจฺจโย. ปจฺฉาชาตา กิเลสา จ สมฺปยุตฺตกา จ ขนฺธา รูปชีวิตินฺทฺริยญฺจ กฏตฺตารูปานํ อตฺถิปจฺจเยน ปจฺจโย. (2)}

\prose{กิเลโส จ โนกิเลโส จ ธมฺมา กิเลสสฺส จ โนกิเลสสฺส จ ธมฺมสฺส อตฺถิปจฺจเยน ปจฺจโย. สหชาตํ, ปุเรชาตํ.   \csromanfolio{15}\textbf{สหชาโต} โนกิเลโส เอโก ขนฺโธ จ โลโภ จ ติณฺณนฺนํ ขนฺธานํ จิตฺตสมุฏฺฐานานญฺจ รูปานํ โมหสฺส จ, ทิฏฺฐิยา, ถินสฺส, อุทฺธจฺจสฺส, อหิริกสฺส, อโนตฺตปฺปสฺส อตฺถิปจฺจเยน ปจฺจโย ฯเปฯ เทฺว ขนฺธา จ ฯเปฯ. \textbf{สหชาโต} โลโภ จ วตฺถุ จ โมหสฺส, ทิฏฺฐิยา, ถินสฺส, อุทฺธจฺจสฺส, อหิริกสฺส, อโนตฺตปฺปสฺส สมฺปยุตฺตกานญฺจ ขนฺธานํ อตฺถิปจฺจเยน ปจฺจโย. (จกฺกํ.) (3)}

\prose{นตฺถิ\-ปจฺจเยน ปจฺจโย. วิคต\-ปจฺจเยน ปจฺจโย. อวิคต\-ปจฺจเยน ปจฺจโย.}

\csromanmarkh{1. ปจฺจยานุโลม}\csromanmarkhii{2. สงฺขฺยาวาร}\csromanheadernomark{2}{1. \textbf{ปจฺจยานุโลม 2. สงฺขฺยาวาร}}
\proseitem{32}{เหตุยา จตฺตาริ, อารมฺมเณ นว, อธิปติยา นว, อนนฺตเร นว, สมนนฺตเร นว, สหชาเต นว, อญฺญมญฺเญ นว, นิสฺสเย นว, อุปนิสฺสเย นว, ปุเรชาเต ตีณิ, ปจฺฉาชาเต ตีณิ, อาเสวเน นว, กมฺเม ตีณิ, วิปาเก เอกํ, อาหาเร ตีณิ, อินฺทฺริเย ตีณิ, ฌาเน ตีณิ, มคฺเค นว, สมฺปยุตฺเต นว, วิปฺปยุตฺเต ปญฺจ, อตฺถิยา นว, นตฺถิยา นว, วิคเต นว, อวิคเต นว.}

\proseitem{33}{กิเลโส ธมฺโม กิเลสสฺส ธมฺมสฺส อารมฺมณ\-ปจฺจเยน ปจฺจโย. สหชตปจฺจเยน ปจฺจโย. อุปนิสฺสย\-ปจฺจเยน ปจฺจโย. (1)}

\prose{กิเลโส ธมฺโม โนกิเลสสฺส ธมฺมสฺส อารมฺมณ\-ปจฺจเยน ปจฺจโย. สหชาต\-ปจฺจเยน ปจฺจโย. อุปนิสฺสย\-ปจฺจเยน ปจฺจโย. ปจฺฉาชาต\-ปจฺจเยน ปจฺจโย. (2)}

\prose{กิเลโส ธมฺโม กิเลสสฺส จ โนกิเลสสฺส จ ธมฺมสฺส อารมฺมณ\-ปจฺจเยน ปจฺจโย. สหชาต\-ปจฺจเยน ปจฺจโย. อุปนิสฺสย\-ปจฺจเยน ปจฺจโย. (3)}

\proseitem{34}{โนกิเลโส ธมฺโม โนกิเลสสฺส ธมฺมสฺส อารมฺมณ\-ปจฺจเยน ปจฺจโย. สหชาต\-ปจฺจเยน ปจฺจโย. อุปนิสฺสย\-ปจฺจเยน \csromanfolio{16}ปจฺจโย. ปุเรชาต\-ปจฺจเยน ปจฺจโย. ปจฺฉาชาต\-ปจฺจเยน ปจฺจโย. กมฺมปจฺจเยน ปจฺจโย. อาหารปจฺจเยน ปจฺจโย. อินฺทฺริย\-ปจฺจเยน ปจฺจโย. (1)}

\prose{โนกิเลโส ธมฺโม กิเลสสฺส ธมฺมสฺส อารมฺมณ\-ปจฺจเยน ปจฺจโย. สหชาต\-ปจฺจเยน ปจฺจโย. อุปนิสฺสย\-ปจฺจเยน ปจฺจโย. ปุเรชาต\-ปจฺจเยน ปจฺจโย. (2)}

\prose{โนกิเลโส ธมฺโม กิเลสสฺส จ โนกิเลสสฺส จ ธมฺมสฺส อารมฺมณ\-ปจฺจเยน ปจฺจโย. สหชาต\-ปจฺจเยน ปจฺจโย. อุปนิสฺสย\-ปจฺจเยน ปจฺจโย. ปุเรชาต\-ปจฺจเยน ปจฺจโย. (3)}

\proseitem{35}{กิเลโส จ โนกิเลโส จ ธมฺมา กิเลสสฺส ธมฺมสฺส อารมฺมณ\-ปจฺจเยน ปจฺจโย. สหชาต\-ปจฺจเยน ปจฺจโย. อุปนิสฺสย\-ปจฺจเยน ปจฺจโย. (1)}

\prose{กิเลโส จ โนกิเลโส จ ธมฺมา โนกิเลสสฺส ธมฺมสฺส อารมฺมณ\-ปจฺจเยน ปจฺจโย. สหชาต\-ปจฺจเยน ปจฺจโย. อุปนิสฺสย\-ปจฺจเยน ปจฺจโย. ปจฺฉาชาต\-ปจฺจเยน ปจฺจโย. (2)}

\prose{กิเลโส จ โนกิเลโส จ ธมฺมา กิเลสสฺส จ โนกิเลสสฺส จ ธมฺมสฺส อารมฺมณ\-ปจฺจเยน ปจฺจโย. สหชาต\-ปจฺจเยน ปจฺจโย. อุปนิสฺสย\-ปจฺจเยน ปจฺจโย. (3)}

\csromanmarkh{2. ปจฺจย\-ปจฺจนีย}\csromanmarkhii{2. สงฺขฺยาวาร}\csromanheadernomark{2}{2. ปจฺจย\-ปจฺจนีย 2. สงฺขฺยาวาร}
\proseitem{36}{นเหตุยา นว, นอารมฺมเณ นว, นอธิปติยา นว, (สพฺพตฺถ นว,) โนอวิคเต นว.}

\proseitem{37}{เหตุปจฺจยา นอารมฺมเณ จตฺตาริ, นอธิปติยา จตฺตาริ, นอนนฺตเร จตฺตาริ, นสมนนฺตเร จตฺตาริ, นอญฺญมญฺเญ เทฺว, นอุปนิสฺสเย จตฺตาริ, (สพฺพตฺถ จตฺตาริ,) นสมฺปยุตฺเต เทฺว, นวิปฺปยุตฺเต จตฺตาริ, โนนตฺถิยา จตฺตาริ, โนวิคเต จตฺตาริ.}

\csromanmatikaanchor{mtk.4}
\csromanmatikaanchor{mtk.5}
\csromantocmarkref{section}{76. สํกิเลสทุก}{mtk.4}
\bookmark[dest={mtk.4},level=1]{76. สํกิเลสทุก}
\csromantocmarkref{section}{77. สํกิลิฏฺฐทุก}{mtk.5}
\bookmark[dest={mtk.5},level=1]{77. สํกิลิฏฺฐทุก}
\csromanmarkh{77. สํกิลิฏฺฐทุก}\csromanmarkhii{4. ปจฺจย\-ปจฺจนียา\-นุโลม}\csromanheadernomark{1}{77. สํกิลิฏฺฐทุก 4. ปจฺจย\-ปจฺจนียา\-นุโลม}
\proseitemwithcloser{38}{\csromanfolio{17}นเหตุปจฺจยา อารมฺมเณ นว, อธิปติยา นว, (อนุโลมมาติกา กาตพฺพา.) ฯเปฯ อวิคเต นว.}{กิเลสทุกํ นิฏฺฐิตํ.}
\csromanmarkh{76. สํกิเลสิกทุก}\csromanmarkhii{1. ปฏิจฺจวาร}\csromanheadernomark{2}{76. \textbf{สํกิเลสิกทุก 1. ปฏิจฺจวาร}}
\vspace{\tipitakaparskip}
\csromancenter{สํกิเลสิกํ ธมฺมํ ปฏิจฺจ สํกิเลสิโก ธมฺโม อุปฺปชฺชติ เหตุปจฺจยา.}
\csromancenter{(ยถา โลกิยทุกํ, เอวํ นินฺนานากรณํ.)}
\nitthitam{สํกิเลสิกทุกํ นิฏฺฐิตํ.}
\csromanmarkh{77. สํกิลิฏฺฐทุก}\csromanmarkhii{1. ปฏิจฺจวาร}\csromanheadernomark{2}{77. สํกิลิฏฺฐทุก 1. ปฏิจฺจวาร}
\proseitem{40}{สํกิลิฏฺฐํ ธมฺมํ ปฏิจฺจ สํกิลิฏฺโฐ ธมฺโม อุปฺปชฺชติ เหตุปจฺจยา. สํกิลิฏฺฐํ เอกํ ขนฺธํ ปฏิจฺจ ตโย ขนฺธา ฯเปฯ เทฺว ขนฺเธ ฯเปฯ. (1)}

\prose{สํกิลิฏฺฐํ ธมฺมํ ปฏิจฺจ อสํกิลิฏฺโฐ ธมฺโม อุปฺปชฺชติ เหตุปจฺจยา. สํกิลิฏฺเฐ ขนฺเธ ปฏิจฺจ จิตฺต\-สมุฏฺฐานํ รูปํ. (2)}

\prose{สํกิลิฏฺฐํ ธมฺมํ ปฏิจฺจ สํกิลิฏฺโฐ จ อสํกิลิฏฺโฐ จ ธมฺมา อุปฺปชฺชนฺติ เหตุปจฺจยา. สํกิลิฏฺฐํ เอกํ ขนฺธํ ปฏิจฺจ ตโย ขนฺธา จิตฺต\-สมุฏฺฐานญฺจ รูปํ ฯเปฯ เทฺว ขนฺเธ ฯเปฯ. (3)}

\prose{อสํกิลิฏฺฐํ ธมฺมํ ปฏิจฺจ อสํกิลิฏฺโฐ ธมฺโม อุปฺปชฺชติ เหตุปจฺจยา. อสํกิลิฏฺฐํ เอกํ ขนฺธํ ปฏิจฺจ ตโย ขนฺธา จิตฺต\-สมุฏฺฐานญฺจ รูปํ ฯเปฯ เทฺว ขนฺเธ ฯเปฯ. ปฏิสนฺธิ\-กฺขเณ ฯเปฯ. ขนฺเธ ปฏิจฺจ วตฺถุ, ปตฺถุํ ปฏิจฺจ ขนฺธา. เอกํ มหาภูตํ ฯเปฯ. (1)}

\prose{\csromanfolio{18}สํกิลิฏฺฐญฺจ อสํกิลิฏฺฐญฺจ ธมฺมํ ปฏิจฺจ อสํกิลิฏฺโฐ ธมฺโม อุปฺปชฺชติ เหตุปจฺจยา. สํกิลิฏฺเฐ ขนฺเธ จ มหาภูเต จ ปฏิจฺจ จิตฺต\-สมุฏฺฐานํ รูปํ. (สํขิตฺตํ.) (1)}

\csromanmarkh{1. ปจฺจยานุโลม}\csromanmarkhii{2. สงฺขฺยาวาร}\csromanheadernomark{2}{1. \textbf{ปจฺจยานุโลม 2. สงฺขฺยาวาร}}
\proseitem{41}{เหตุยา ปญฺจ, อารมฺมเณ เทฺว, อธิปติยา ปญฺจ, อนนฺตเร เทฺว, สมนนฺตเร เทฺว, สหชาเต ปญฺจ, อญฺญมญฺเญ เทฺว, นิสฺสเย ปญฺจ, อุปนิสฺสเย เทฺว, ปุเรชาเต เทฺว, อาเสวเน เทฺว, กมฺเม ปญฺจ, วิปาเก เอกํ อาหาเร ปญฺจ, (สํขิตฺตํ.) อวิคเต ปญฺจ.}

\csromanmarkh{2. ปจฺจย\-ปจฺจนีย}\csromanmarkhii{1. วิภงฺควาร}\csromanheadernomark{2}{2. ปจฺจย\-ปจฺจนีย 1. วิภงฺควาร}
\proseitem{42}{สํกิลิฏฺฐํ ธมฺมํ ปฏิจฺจ สํกิลิฏฺโฐ ธมฺโม อุปฺปชฺชติ นเหตุปจฺจยา. วิจิกิจฺฉา\-สหคเต อุทฺธจฺจ\-สหคเต ขนฺเธ ปฏิจฺจ วิจิกิจฺฉา\-สหคโต อุทฺธจฺจ\-สหคโต โมโห. (1)}

\prose{อสํกิลิฏฺฐํ ธมฺมํ ปฏิจฺจ อสํกิลิฏฺโฐ ธมฺโม อุปฺปชฺชติ นเหตุปจฺจยา. อเหตุกํ อสํกิลิฏฺฐํ เอกํ ขนฺธํ ปฏิจฺจ ตโย ขนฺธา จิตฺต\-สมุฏฺฐานญฺจ รูปํ ฯเปฯ เทฺว ขนฺเธ ฯเปฯ. อเหตุก\-ปฏิสนฺธิ\-กฺขเณ ฯเปฯ. (ยาว อสญฺญสตฺตา.) (1)}

\csromanmarkh{2. ปจฺจย\-ปจฺจนีย}\csromanmarkhii{2. สงฺขฺยาวาร}\csromanheadernomark{2}{2. ปจฺจย\-ปจฺจนีย 2. สงฺขฺยาวาร}
\vspace{\tipitakaparskip}
\csromancenter{\textbf{นเหตุ}ยา เทฺว, นอารมฺมเณ ตีณิ, นอธิปติยา ปญฺจ, นอนนฺตเร ตีณิ, นสมนนฺตเร ตีณิ, นอญฺญมญฺเญ ตีณิ, นฺเอาปนิสฺสเย ตีณิ, นปุเรชาเต จตฺตาริ, นปจฺฉาชาเต ปญฺจ, นอาเสวเน ปญฺจ, นกมฺเม เทฺว, นวิปาเก ปญฺจ, นอาหาเร เอกํ, ไนนฺทฺริเย เอกํ, นฌาเน เอกํ, นมคฺเค เอกํ, นสมฺปยุตฺเต ตีณิ, นวิปฺปยุตฺเต เทฺว, โนนตฺถิยา ตีณิ, โนวิคเต ตีณิ.}
\csromancenter{(เอวํ อิตเร เทฺว คณนาปิ สหชาตวาโรปิ กาตพฺโพ.)}
\csromanmarkh{77. สํกิลิฏฺฐทุก}\csromanmarkhii{77. สํกิลิฏฺฐทุก 3. ปจฺจยวาร}\csromanheadernomark{2}{77. สํกิลิฏฺฐทุก \textbf{77. สํกิลิฏฺฐทุก 3. ปจฺจยวาร}}
\proseitem{44}{\csromanfolio{19}สํกิลิฏฺฐิ ธมฺมํ ปจฺจยา สํกิลิฏฺโฐ ธมฺโม อุปฺปชฺชติ เหตุปจฺจยา. ตีณิ. (ปฏิจฺจสทิสํ.)}

\prose{อสํกิลิฏฺฐํ ธมฺมํ ปจฺจยา อสํกิลิฏฺโฐ ธมฺโม อุปฺปชฺชติ เหตุปจฺจยา. อสํกิลิฏฺฐํ เอกํ ขนฺธํ ปจฺจยา ตโย ขนฺธา จิตฺต\-สมุฏฺฐานญฺจ รูปํ ฯเปฯ เทฺว ขนฺเธ ฯเปฯ. ปฏิสนฺธิ\-กฺขเณ ฯเปฯ. (ยาว อชฺฌตฺติกา มหาภูตา.) วตฺถุํ ปจฺจยา อสํกิลิฏฺฐา ขนฺธา. (1)}

\prose{อสํกิลิฏฺฐํ ธมฺมํ ปจฺจยา สํกิลิฏฺโฐ ธมฺโม อุปฺปชฺชติ เหตุปจฺจยา. วตฺถุํ ปจฺจยา สํกิลิฏฺฐา ขนฺธา. (2)}

\prose{อสํกิลิฏฺฐํ ธมฺมํ ปจฺจยา สํกิลิฏฺโฐ จ อสํกิลิฏฺโฐ จ ธมฺมา อุปฺปชฺชนฺติ เหตุปจฺจยา. วตฺถุํ ปจฺจยา สํกิลิฏฺฐา ขนฺธา. มหาภูเต ปจฺจยา จิตฺต\-สมุฏฺฐานํ รูปํ. (3)}

\proseitem{45}{สํกิลิฏฺฐญฺจ อสํกิลิฏฺฐญฺจ ธมฺมํ ปจฺจยา สํกิลิฏฺโฐ ธมฺโม อุปฺปชฺชติ เหตุปจฺจยา. สํกิลิฏฺฐํ เอกํ ขนฺธญฺจ วตฺถุญฺจ ปจฺจยา ตโย ขนฺธา ฯเปฯ เทฺว ขนฺเธ จ ฯเปฯ. (1)}

\prose{สํกิลิฏฺฐญฺจ อสํกิลิฏฺฐญฺจ ธมฺมํ ปจฺจยา อสํกิลิฏฺโฐ ธมฺโม อุปฺปชฺชติ เหตุปจฺจยา. สํกิลิฏฺเฐ ขนฺเธ จ มหาภูเต จ ปจฺจยา จิตฺต\-สมุฏฺฐานํ รูปํ. (2)}

\prose{สํกิลิฏฺฐญฺจ อสํกิลิฏฺฐญฺจ ธมฺมํ ปจฺจยา สํกิลิฏฺโฐ จ อสํกิลิฏฺโฐ จ ธมฺมา อุปฺปชฺชนฺติ เหตุปจฺจยา. สํกิลิฏฺฐํ เอกํ ขนฺธญฺจ วตฺถุญฺจ ปจฺจยา ตโย ขนฺธา ฯเปฯ เทฺว ขนฺเธ จ ฯเปฯ. สํกิลิฏฺเฐ ขนฺเธ จ มหาภูเต จ ปจฺจยา จิตฺต\-สมุฏฺฐานํ รูปํ. (สํขิตฺตํ.) (3)}

\csromanmarkh{1. ปจฺจยานุโลม}\csromanmarkhii{2. สงฺขฺยาวาร}\csromanheadernomark{2}{1. \textbf{ปจฺจยานุโลม 2. สงฺขฺยาวาร}}
\proseitem{46}{\textbf{เหตุ}ยา นว, อารมฺมเณ จตฺตาริ, อธิปติยา นว, อนนฺตเร จตฺตาริ, สมนนฺตเร จตฺตาริ, สหชาเต นว, อญฺญมญฺเญ จตฺตาริ, นิสฺสเย นว, อุปนิสฺสเย จตฺตาริ, ปุเรชาเต จตฺตาริ, อาเสวเน \csromanfolio{20}\textbf{จตฺตาริ, กมฺเม นว, วิปาเก เอกํ, อาหาเร นว, อินฺทฺริเย นว ฯเปฯ วิคเต} จตฺตาริ, อวิคเต นว.}

\csromanmarkh{2. ปจฺจย\-ปจฺจนีย}\csromanmarkhii{1. วิภงฺควาร}\csromanheadernomark{2}{2. ปจฺจย\-ปจฺจนีย 1. วิภงฺควาร}
\proseitem{47}{สํกิลิฏฺฐํ ธมฺมํ ปจฺจยา สํกิลิฏฺโฐ ธมฺโม อุปฺปชฺชติ นเหตุปจฺจยา. วิจิกิจฺฉา\-สหคเต อุทฺธจฺจ\-สหคเต ขนฺเธ ปจฺจยา วิจิกิจฺฉา\-สหคโต อุทฺธจฺจ\-สหคโต โมโห. (1)}

\prose{อสํกิลิฏฺฐํ ธมฺมํ ปจฺจยา อสํกิลิฏฺโฐ ธมฺโม อุปฺปชฺชติ นเหตุปจฺจยา. อเหตุกํ อสํกิลิฏฺฐํ ฯเปฯ. (ยาว อสญฺญสตฺตา.) จกฺขายตนํ ปจฺจยา จกฺขุวิญฺญาณํ ฯเปฯ. กายายตนํ ปจฺจยา กายวิญฺญาณํ. วตฺถุํ ปจฺจยา อเหตุกา อสํกิลิฏฺฐา ขนฺธา. (1)}

\prose{อสํกิลิฏฺฐํ ธมฺมํ ปจฺจยา สํกิลิฏฺโฐ ธมฺโม อุปฺปชฺชติ นเหตุปจฺจยา. วตฺถุํ ปจฺจยา วิจิกิจฺฉา\-สหคโต อุทฺธจฺจ\-สหคโต โมโห. (2)}

\prose{สํกิลิฏฺฐญฺจ อสํกิลิฏฺฐญฺจ ธมฺมํ ปจฺจยา สํกิลิฏฺโฐ ธมฺโม อุปฺปชฺชติ นเหตุปจฺจยา. วิจิกิจฺฉา\-สหคเต อุทฺธจฺจ\-สหคเต ขนฺเธ จ วตฺถุญฺจ ปจฺจยา วิจิกิจฺฉา\-สหคโต อุทฺธจฺจสคโต โมโห. (1)}

\csromanmarkh{2. ปจฺจย\-ปจฺจนีย}\csromanmarkhii{2. สงฺขฺยาวาร}\csromanheadernomark{2}{2. ปจฺจย\-ปจฺจนีย 2. สงฺขฺยาวาร}
\vspace{\tipitakaparskip}
\csromancenter{\textbf{นเหตุ}ยา จตฺตาริ, นอารมฺมเณ ตีณิ, นอธิปติยา นว, นอนนฺตเร ตีณิ, นสมนนฺตเร ตีณิ, นอญฺญมญฺเญ ตีณิ, นฺเอาปนิสฺสเย ตีณิ, นปุเรชาเต จตฺตาริ, นปจฺฉาชาเต นว, นอาเสวเน นว, นกมฺเม จตฺตาริ, นวิปาเก นว, นอาหาเร เอกํ, ไนนฺทฺริเย เอกํ, นฌาเน เอกํ, นมคฺเค เอกํ, นสมฺปยุตฺเต ตีณิ, นวิปฺปยุตฺเต เทฺว, โนนตฺถิยา ตีณิ, โนวิคเต ตีณิ.}
\csromancenter{(เอวํ อิตเร เทฺว คณนาปิ นิสฺสยวาโรปิ กาตพฺโพ.)}
\csromanmarkh{77. สํกิลิฏฺฐทุก}\csromanmarkhii{77. สํกิลิฏฺฐทุก 5. สํสฏฺฐวาร}\csromanheadernomark{2}{77. สํกิลิฏฺฐทุก \textbf{77. สํกิลิฏฺฐทุก 5. สํสฏฺฐวาร}}
\prose{\csromanfolio{21}1. ปจฺจยานุโลมาทิเหตุ}

\proseitem{49}{สํกิลิฏฺฐํ ธมฺมํ สํสฏฺโฐ สํกิลิฏฺโฐ ธมฺโม อุปฺปชฺชติ เหตุปจฺจยา. สํกิลิฏฺฐํ เอกํ ขนฺธํ สํสฏฺฐา ตโย ขนฺธา ฯเปฯ เทฺว ขนฺเธ ฯเปฯ. (1)}

\prose{อสํกิลิฏฺฐํ ธมฺมํ สํสฏฺโฐ อสํกิลิฏฺโฐ ธมฺโม อุปฺปชฺชติ เหตุปจฺจยา. อสํกิลิฏฺฐํ เอกํ ขนฺธํ สํสฏฺฐา ตโย ขนฺธา ฯเปฯ เทฺว ขนฺเธ ฯเปฯ. ปฏิสนฺธิ\-กฺขเณ ฯเปฯ. (1)}

\prose{เหตุยา เทฺว, อารมฺมเณ เทฺว, อธิปติยา เทฺว, (สพฺพตฺถ เทฺว.) วิปาเก เอกํ ฯเปฯ อวิคเต เทฺว. อนุโลมํ.}

\proseitem{50}{สํกิลิฏฺฐํ ธมฺมํ สํสฏฺโฐ สํกิลิฏฺโฐ ธมฺโม อุปฺปชฺชติ นเหตุปจฺจยา. วิจิกิจฺฉา\-สหคเต อุทฺธจฺจ\-สหคเต ขนฺเธ สํสฏฺโฐ วิจิกิจฺฉา\-สหคโต อุทฺธจฺจ\-สหคโต โมโห. (1)}

\prose{อสํกิลิฏฺฐํ ธมฺมํ สํสฏฺโฐ อสํกิลิฏฺโฐ ธมฺโม อุปฺปชฺชติ นเหตุปจฺจยา. อเหตุกํ อสํกิลิฏฺฐํ เอกํ ขนฺธํ สํสฏฺฐา ตโย ขนฺธา ฯเปฯ เทฺว ขนฺเธ ฯเปฯ. อเหตุก\-ปฏิสนฺธิ\-กฺขเณ ฯเปฯ. (1)}

\prose{นเหตุยา เทฺว, นอธิปติยา เทฺว, นปุเรชาเต เทฺว, นปจฺฉาชาเต เทฺว, นอาเสวเน เทฺว, นกมฺเม เทฺว, นวิปาเก เทฺว, นฌาเน เอกํ, นมคฺเค เอกํ, นวิปฺปยุตฺเต เทฺว.}

\vspace{\tipitakaparskip}
\csromancenter{ปจฺจนียํ.}
\csromancenter{(เอวํ อิตเร เทฺว คณนาปิ สมฺปยุตฺตวาโรปิ กาตพฺโพ.)}
\csromanmarkh{77. สํกิลิฏฺฐทุก}\csromanmarkhii{7. ปญฺหาวาร}\csromanheadernomark{2}{77. สํกิลิฏฺฐทุก 7. ปญฺหาวาร}
\proseitem{51}{สํกิลิฏฺโฐ ธมฺโม สํกิลิฏฺฐสฺส ธมฺมสฺส เหตุปจฺจเยน ปจฺจโย. สํกิลิฏฺฐา เหตู สมฺปยุตฺตกานํ ขนฺธานํ เหตุปจฺจเยน ปจฺจโย. (1)}

\prose{\csromanfolio{22}สํกิลิฏฺโฐ ธมฺโม อสํกิลิฏฺฐสฺส ธมฺมสฺส เหตุปจฺจเยน ปจฺจโย. สํกิลิฏฺฐา เหตู จิตฺต\-สมุฏฺฐานานํ รูปานํ เหตุปจฺจเยน ปจฺจโย. (2)}

\prose{สํกิลิฏฺโฐ ธมฺโม สํกิลิฏฺฐสฺส จ อสํกิลิฏฺฐสฺส จ ธมฺมสฺส เหตุปจฺจเยน ปจฺจโย. สํกิลิฏฺฐา เหตู สมฺปยุตฺตกานํ ขนฺธานํ จิตฺตสมุฏฺฐานานญฺจ รูปานํ เหตุปจฺจเยน ปจฺจโย. (3)}

\prose{อสํกิลิฏฺโฐ ธมฺโม อสํกิลิฏฺฐสฺส ธมฺมสฺส เหตุปจฺจเยน ปจฺจโย. อสํกิลิฏฺฐา เหตู สมฺปยุตฺตกานํ ขนฺธานํ จิตฺตสมุฏฺฐานานญฺจ รูปานํ เหตุปจฺจเยน ปจฺจโย. ปฏิสนฺธิ\-กฺขเณ ฯเปฯ. (1)}

\proseitem{52}{สํกิลิฏฺโฐ ธมฺโม สํกิลิฏฺฐสฺส ธมฺมสฺส อารมฺมณ\-ปจฺจเยน ปจฺจโย. ราคํ อสฺสาเทติ อภินนฺทติ, ตํ อารพฺภ ราโค อุปฺปชฺชติ. ทิฏฺฐิ อุปฺปชฺชติ. วิจิกิจฺฉา ฯเปฯ อุทฺธจฺจํ ฯเปฯ โทมนสฺสํ อุปฺปชฺชติ. ทิฏฺฐิํ อสฺสาเทติ ฯเปฯ. (กุสล\-ตฺติก\-สทิสํ.) วิจิกิจฺฉํ อารพฺภ ฯเปฯ อุทฺธจฺจํ อารพฺภ ฯเปฯ โทมนสฺสํ อุปฺปชฺชติ. ทิฏฺฐิ ฯเปฯ วิจิกิจฺฉา ฯเปฯ อุทฺธจฺจํ อุปฺปชฺชติ. (1)}

\prose{สํกิลิฏฺโฐ ธมฺโม อสํกิลิฏฺฐสฺส ธมฺมสฺส อารมฺมณ\-ปจฺจเยน ปจฺจโย. อริยา ปหีเน กิเลเส ปจฺจเวกฺขนฺติ, วิกฺขมฺภิเต กิเลเส ฯเปฯ. ปุพฺเพ สมุทาจิณฺเณ ฯเปฯ. สํกิลิฏฺเฐ ขนฺเธ อนิจฺจโต ฯเปฯ วิปสฺสติ. เจโตปริย\-ญาเณน สํกิลิฏฺฐจิตฺต\-สมงฺคิสฺส จิตฺตํ ชานาติ. สํกิลิฏฺฐา ขนฺธา เจโตปริย\-ญาณสฺส, ปุพฺเพ\-นิวาสา\-นุสฺสติ\-ญาณสฺส, ยถากมฺมูปคญาณสฺส, อนาคตํสญาณสฺส, อาวชฺชนาย อารมฺมณ\-ปจฺจเยน ปจฺจโย. (2)}

\proseitem{53}{อสํกิลิฏฺโฐ ธมฺโม อสํกิลิฏฺฐสฺส ธมฺมสฺส อารมฺมณ\-ปจฺจเยน ปจฺจโย. ทานํ ฯเปฯ สีลํ ฯเปฯ อุโปสถกมฺมํ ฯเปฯ. ปุพฺเพ ฯเปฯ. ฌานา วุฏฺฐหิตฺวา ฌานํ ปจฺจเวกฺขติ. อริยา มคฺคา วุฏฺฐหิตฺวา มคฺคํ ปจฺจเวกฺขนฺติ ฯเปฯ อาวชฺชนาย อารมฺมณ\-ปจฺจเยน ปจฺจโย. จกฺขุํ ฯเปฯ. วตฺถุํ อสํกิลิฏฺเฐ ขนฺเธ อนิจฺจโต ฯเปฯ วิปสฺสติ. ทิพฺเพน จกฺขุนา รูปํ ปสฺสติ ฯเปฯ. อนาคตํสญาณสฺส, อาวชฺชนาย อารมฺมณ\-ปจฺจเยน ปจฺจโย. (1)}

\prose{อสํกิลิฏฺโฐ ธมฺโม สํกิลิฏฺฐสฺส ธมฺมสฺส อารมฺมณ\-ปจฺจเยน ปจฺจโย. ทานํ ฯเปฯ สีลํ ฯเปฯ. ฌานา วุฏฺฐหิตฺวา ฯเปฯ. จกฺขุํ ฯเปฯ. วตฺถุํ อสํกิลิฏฺเฐ \csromanfolio{23}ขนฺเธ อสฺสาเทติ อภินนฺทติ, ตํ อารพฺภ ราโค ฯเปฯ โทมนสฺสํ อุปฺปชฺชติ. (2)}

\csromanheader{2}{\textbf{อธิปติ}}
\proseitem{54}{สํกิลิฏฺโฐ ธมฺโม สํกิลิฏฺฐสฺส ธมฺมสฺส อธิปติ\-ปจฺจเยน ปจฺจโย. อารมฺมณา\-ธิปติ, สหชาตา\-ธิปติ.  อารมฺมณา\-ธิปติ\csromandash{}ราคํ ครุํ กตฺวา อสฺสาเทติ อภินนฺทติ, ตํ ครุํ กตฺวา ราโค อุปฺปชฺชติ. ทิฏฺฐิ อุปฺปชฺชติ. ทิฏฺฐิํ ครุํ กตฺวา อสฺสาเทติ อภินนฺทติ, ตํ ครุํ กตฺวา ราโค อุปฺปชฺชติ. ทิฏฺฐิ อุปฺปชฺชติ. สหชาตา\-ธิปติ\csromandash{}สํกิลิฏฺฐา\-ธิปติ สมฺปยุตฺตกานํ ขนฺธานํ อธิปติ\-ปจฺจเยน ปจฺจโย. (1)}

\prose{สํกิลิฏฺโฐ ธมฺโม อสํกิลิฏฺฐสฺส ธมฺมสฺส อธิปติ\-ปจฺจเยน ปจฺจโย. สหชาตา\-ธิปติ\csromandash{}สํกิลิฏฺฐา\-ธิปติ จิตฺต\-สมุฏฺฐานานํ รูปานํ อธิปติ\-ปจฺจเยน ปจฺจโย. (2)}

\prose{สํกิลิฏฺโฐ ธมฺโม สํกิลิฏฺฐสฺส จ อสํกิลิฏฺฐสฺส จ ธมฺมสฺส อธิปติ\-ปจฺจเยน ปจฺจโย. สหชาตา\-ธิปติ\csromandash{}สํกิลิฏฺฐา\-ธิปติ สมฺปยุตฺตกานํ ขนฺธานํ จิตฺตสมุฏฺฐานานญฺจ รูปานํ อธิปติ\-ปจฺจเยน ปจฺจโย. (3)}

\proseitem{55}{อสํกิลิฏฺโฐ ธมฺโม อสํกิลิฏฺฐสฺส ธมฺมสฺส อธิปติ\-ปจฺจเยน ปจฺจโย. อารมฺมณา\-ธิปติ, สหชาตา\-ธิปติ.  อารมฺมณา\-ธิปติ\csromandash{}ทานํ ฯเปฯ สีลํ ฯเปฯ อุโปสถกมฺมํ ฯเปฯ. ปุพฺเพ ฯเปฯ. ฌานา วุฏฺฐหิตฺวา ฌานํ ครุํ กตฺวา ปจฺจเวกฺขติ. อริยา มคฺคา วุฏฺฐหิตฺวา มคฺคํ ครุํ กตฺวา ปจฺจเวกฺขนฺติ ฯเปฯ. นิพฺพานํ ผลสฺส อธิปติ\-ปจฺจเยน ปจฺจโย. สหชาตา\-ธิปติ\csromandash{}อสํกิลิฏฺฐา\-ธิปติ สมฺปยุตฺตกานํ ขนฺธานํ จิตฺตสมุฏฺฐานานญฺจ รูปานํ อธิปติ\-ปจฺจเยน ปจฺจโย. (1)}

\prose{อสํกิลิฏฺโฐ ธมฺโม สํกิลิฏฺฐสฺส ธมฺมสฺส อธิปติ\-ปจฺจเยน ปจฺจโย. อารมฺมณา\-ธิปติ\csromandash{}ทานํ ฯเปฯ สีลํ ฯเปฯ อุโปสถกมฺมํ ฯเปฯ. ปุพฺเพ ฯเปฯ. ฌานา วุฏฺฐหิตฺวา ฯเปฯ. จกฺขุํ ฯเปฯ. วตฺถุํ อสํกิลิฏฺเฐ ขนฺเธ ครุํ กตฺวา อสฺสาเทติ อภินนฺทติ, ตํ ครุํ กตฺวา ราโค อุปฺปชฺชติ. ทิฏฺฐิ อุปฺปชฺชติ. (2)}

\proseitem{56}{\csromanfolio{24}สํกิลิฏฺโฐ ธมฺโม สํกิลิฏฺฐสฺส ธมฺมสฺส อนนฺตร\-ปจฺจเยน ปจฺจโย. ปุริมา ปุริมา สํกิลิฏฺฐา ขนฺธา ปจฺฉิมานํ ปจฺฉิมานํ สํกิลิฏฺฐานํ ขนฺธานํ อนนฺตร\-ปจฺจเยน ปจฺจโย. (1)}

\prose{สํกิลิฏฺโฐ ธมฺโม อสํกิลิฏฺฐสฺส ธมฺมสฺส อนนฺตร\-ปจฺจเยน ปจฺจโย. สํกิลิฏฺฐา ขนฺธา วุฏฺฐานสฺส อนนฺตร\-ปจฺจเยน ปจฺจโย. (2)}

\proseitem{57}{อสํกิลิฏฺโฐ ธมฺโม อสํกิลิฏฺฐสฺส ธมฺมสฺส อนนฺตร\-ปจฺจเยน ปจฺจโย. ปุริมา ปุริมา อสํกิลิฏฺฐา ขนฺธา ปจฺฉิมานํ ปจฺฉิมานํ อสํกิลิฏฺฐานํ ขนฺธานํ ฯเปฯ ผลสมาปตฺติยา อนนฺตร\-ปจฺจเยน ปจฺจโย. (1)}

\prose{อสํกิลิฏฺโฐ ธมฺโม สํกิลิฏฺฐสฺส ธมฺมสฺส อนนฺตร\-ปจฺจเยน ปจฺจโย. อาวชฺชนา สํกิลิฏฺฐานํ ขนฺธานํ อนนฺตร\-ปจฺจเยน ปจฺจโย. (2)}

\prose{สมนนฺตร\-ปจฺจเยน ปจฺจโย. จตฺตาริ. สหชาต\-ปจฺจเยน ปจฺจโย. ปญฺจ. อญฺญมญฺญ\-ปจฺจเยน ปจฺจโย. เทฺว. นิสฺสย\-ปจฺจเยน ปจฺจโย. สตฺต.}

\csromanheader{2}{\textbf{อุปนิสฺสย}}
\proseitem{58}{สํกิลิฏฺโฐ ธมฺโม สํกิลิฏฺฐสฺส ธมฺมสฺส อุปนิสฺสย\-ปจฺจเยน ปจฺจโย. อารมฺมณู\-ปนิสฺสโย, อนนฺตรู\-ปนิสฺสโย, ปกตู\-ปนิสฺสโย ฯเปฯ. ปกตู\-ปนิสฺสโย\csromandash{}ราคํ อุปนิสฺสาย ปาณํ หนติ ฯเปฯ สํฆํ ภินฺทติ. โทสํ ฯเปฯ. ปตฺถนํ อุปนิสฺสาย ปาณํ หนติ ฯเปฯ สํฆํ ภินฺทติ. ราโค ฯเปฯ. ปตฺถนา ราคสฺส ฯเปฯ ปตฺถนาย อุปนิสฺสย\-ปจฺจเยน ปจฺจโย. (1)}

\prose{สํกิลิฏฺโฐ ธมฺโม อสํกิลิฏฺฐสฺส ธมฺมสฺส อุปนิสฺสย\-ปจฺจเยน ปจฺจโย. อนนฺตรู\-ปนิสฺสโย, ปกตู\-ปนิสฺสโย ฯเปฯ. ปกตู\-ปนิสฺสโย\csromandash{} ราคํ อุปนิสฺสาย ทานํ เทติ ฯเปฯ สมาปตฺติํ อุปฺปาเทติ. โทสํ ฯเปฯ. ปตฺถนํ อุปนิสฺสาย ทานํ เทติ ฯเปฯ สมาปตฺติํ อุปฺปาเทติ. ราโค ฯเปฯ. ปตฺถนา สทฺธาย ฯเปฯ ปญฺญาย. กายิกสฺส สุขสฺส. กายิกสฺส ทุกฺขสฺส. มคฺคสฺส, ผลสมาปตฺติยา อุปนิสฺสย\-ปจฺจเยน ปจฺจโย. (2)}

\proseitem{59}{อสํกิลิฏฺโฐ ธมฺโม อสํกิลิฏฺฐสฺส ธมฺมสฺส อุปนิสฺสย\-ปจฺจเยน ปจฺจโย. อารมฺมณู\-ปนิสฺสโย, อนนฺตรู\-ปนิสฺสโย, ปกตู\-ปนิสฺสโย ฯเปฯ. ปกตู\-ปนิสฺสโย\csromandash{}สทฺธํ อุปนิสฺสาย ทานํ เทติ ฯเปฯ \csromanfolio{25}สมาปตฺติํ อุปฺปาเทติ. สีลํ ฯเปฯ. ปญฺญํ. กายิกํ สุขํ. กายิกํ ทุกฺขํ. อุตุํ. โภชนํ. เสนาสนํ อุปนิสฺสาย ทานํ เทติ ฯเปฯ สมาปตฺติํ อุปฺปาเทติ. สทฺธา ฯเปฯ. เสนาสนํ สทฺธาย ฯเปฯ ผลสมาปตฺติยา อุปานิสฺสยปจฺจเยน ปจฺจโย. (1)}

\prose{อสํกิลิฏฺโฐ ธมฺโม สํกิลิฏฺฐสฺส ธมฺมสฺส อุปนิสฺสย\-ปจฺจเยน ปจฺจโย. อารมฺมณู\-ปนิสฺสโย, อนนฺตรู\-ปนิสฺสโย, ปกตู\-ปนิสฺสโย ฯเปฯ. ปกตู\-ปนิสฺสโย\csromandash{}สทฺธํ อุปนิสฺสาย มานํ ชปฺเปติ. ทิฏฺฐิํ คณฺหาติ. สีลํ ฯเปฯ. ปญฺญํ. กายิกํ สุขํ. กายิกํ ทุกฺขํ. อุตุํ. โภชนํ. เสนาสนํ อุปนิสฺสาย ปาณํ หนติ ฯเปฯ สํฆํ ภินฺทติ. สทฺธา ฯเปฯ. เสนาสนํ ราคสฺส ฯเปฯ ปตฺถนาย อุปนิสฺสย\-ปจฺจเยน ปจฺจโย. (2)}

\csromanheader{2}{\textbf{ปุเรชาต}}
\proseitem{60}{อสํกิลิฏฺโฐ ธมฺโม อสํกิลิฏฺฐสฺส ธมฺมสฺส ปุเรชาต\-ปจฺจเยน ปจฺจโย. อารมฺมณ\-ปุเรชาตํ, วตฺถุ\-ปุเรชาตํ. (สํขิตฺตํ.) อสํกิลิฏฺโฐ ธมฺโม สํกิลิฏฺฐสฺส ธมฺมสฺส ปุเรชาต\-ปจฺจเยน ปจฺจโย. อารมฺมณ\-ปุเรชาตํ, วตฺถุ\-ปุเรชาตํ. (สํขิตฺตํ.) (2)}

\csromanheader{2}{ปจฺฉาชาตา\-เสวน}
\proseitem{61}{สํกิลิฏฺโฐ ธมฺโม อสํกิลิฏฺฐสฺส ธมฺมสฺส ปจฺฉาชาต\-ปจฺจเยน ปจฺจโย. (สํขิตฺตํ.) อสํกิลิฏฺโฐ ธมฺโม อสํกิลิฏฺฐสฺส ธมฺมสฺส ปจฺฉาชาต\-ปจฺจเยน ปจฺจโย. (สํขิตฺตํ.) อาเสวน\-ปจฺจเยน ปจฺจโย เทฺว.}

\csromanheader{2}{\textbf{กมฺม}}
\proseitem{62}{สํกิลิฏฺโฐ ธมฺโม สํกิลิฏฺฐสฺส ธมฺมสฺส กมฺมปจฺจเยน ปจฺจโย. สํกิลิฏฺฐา เจตนา สมฺปยุตฺตกานํ ขนฺธานํ กมฺมปจฺจเยน ปจฺจโย. (1)}

\prose{สํกิลิฏฺโฐ ธมฺโม อสํกิลิฏฺฐสฺส ธมฺมสฺส กมฺมปจฺจเยน ปจฺจโย. สหชาตา, นานากฺขณิกา.  สหชาตา สํกิลิฏฺฐา เจตนา จิตฺต\-สมุฏฺฐานานํ รูปานํ กมฺมปจฺจเยน ปจฺจโย. นานากฺขณิกา สํกิลิฏฺฐา เจตนา วิปากานํ ขนฺธานํ กฏตฺตา จ รูปานํ กมฺมปจฺจเยน ปจฺจโย. (มูลํ กาตพฺพํ.) สํกิลิฏฺฐา เจตนา สมฺปยุตฺตกานํ ขนฺธานํ จิตฺตสมุฏฺฐานานญฺจ รูปานํ กมฺมปจฺจเยน ปจฺจโย. (3)}

\prose{\csromanfolio{26}อสํกิลิฏฺโฐ ธมฺโม อสํกิลิฏฺฐสฺส ธมฺมสฺส กมฺมปจฺจเยน ปจฺจโย. สหชาตา, นานากฺขณิกา.  สหชาตา อสํกิลิฏฺฐา เจตนา สมฺปยุตฺตกานํ ขนฺธานํ จิตฺตสมุฏฺฐานานญฺจ รูปานํ กมฺมปจฺจเยน ปจฺจโย. ปฏิสนฺธิ\-กฺขเณ ฯเปฯ. นานากฺขณิกา อสํกิลิฏฺฐา เจตนา วิปากานํ ขนฺธานํ กฏตฺตา จ รูปานํ กมฺมปจฺจเยน ปจฺจโย. (1)}

\csromanheader{2}{\textbf{วิปากาทิ}}
\proseitem{63}{อสํกิลิฏฺโฐ ธมฺโม อสํกิลิฏฺฐสฺส ธมฺมสฺส วิปาก\-ปจฺจเยน ปจฺจโย. เอกํ.}

\prose{สํกิลิฏฺโฐ ธมฺโม สํกิลิฏฺฐสฺส ธมฺมสฺส อาหารปจฺจเยน ปจฺจโย. อินฺทฺริย\-ปจฺจเยน ปจฺจโย. ฌานปจฺจเยน ปจฺจโย. มคฺคปจฺจเยน ปจฺจโย. สมฺปยุตฺต\-ปจฺจเยน ปจฺจโย.}

\csromanheader{2}{\textbf{วิปฺปยุตฺต}}
\proseitem{64}{สํกิลิฏฺโฐ ธมฺโม อสํกิลิฏฺฐสฺส ธมฺมสฺส วิปฺปยุตฺต\-ปจฺจเยน ปจฺจโย. สหชาตํ, ปจฺฉาชาตํ. (สํขิตฺตํ.) (1)}

\prose{อสํกิลิฏฺโฐ ธมฺโม อสํกิลิฏฺฐสฺส ธมฺมสฺส วิปฺปยุตฺต\-ปจฺจเยน ปจฺจโย. สหชาตํ, ปุเรชาตํ, ปจฺฉาชาตํ. (สํขิตฺตํ.) (1)}

\prose{อสํกิลิฏฺโฐ ธมฺโม สํกิลิฏฺฐสฺส ธมฺมสฺส วิปฺปยุตฺต\-ปจฺจเยน ปจฺจโย. ปุเรชาตํ วตฺถุ สํกิลิฏฺฐานํ ขนฺธานํ วิปฺปยุตฺต\-ปจฺจเยน ปจฺจโย. (2)}

\prose{อตฺถิอาทิ}

\proseitem{65}{สํกิลิฏฺโฐ ธมฺโม สํกิลิฏฺฐสฺส ธมฺมสฺส อตฺถิปจฺจเยน ปจฺจโย. เอกํ. (ปฏิจฺจ\-วาร\-สทิสํ.) สํกิลิฏฺโฐ ธมฺโม อสํกิลิฏฺฐสฺส ธมฺมสฺส อตฺถิปจฺจเยน ปจฺจโย. สหชาตํ, ปจฺฉาชาตํ. (สํขิตฺตํ.) สํกิลิฏฺโฐ ธมฺโม สํกิลิฏฺฐสฺส จ อสํกิลิฏฺฐสฺส จ ธมฺมสฺส อตฺถิปจฺจเยน ปจฺจโย. (ปฏิจฺจสทิสํ.) (3)}

\prose{อสํกิลิฏฺโฐ ธมฺโม อสํกิลิฏฺฐสฺส ธมฺมสฺส อตฺถิปจฺจเยน ปจฺจโย. สหชาตํ, ปุเรชาตํ, ปจฺฉาชาตํ, อาหารํ, อินฺทฺริยํ. (สํขิตฺตํ.) อสํกิลิฏฺโฐ ธมฺโม สํกิลิฏฺฐสฺส ธมฺมสฺส อตฺถิปจฺจเยน ปจฺจโย. ปุเรชาตํ. (สํขิตฺตํ.) (2)}

\proseitem{66}{\csromanfolio{27}สํกิลิฏฺโฐ จ อสํกิลิฏฺโฐ จ ธมฺมา สํกิลิฏฺฐสฺส ธมฺมสฺส อตฺถิปจฺจเยน ปจฺจโย. สหชาตํ, ปุเรชาตํ.  สหชาโต สํกิลิฏฺโฐ เอโก ขนฺโธ จ วตฺถุ จ ติณฺณนฺนํ ขนฺธานํ อตฺถิปจฺจเยน ปจฺจโย ฯเปฯ เทฺว ขนฺธา จ ฯเปฯ. (สํขิตฺตํ.) (1)}

\prose{สํกิลิฏฺโฐ จ อสํกิลิฏฺโฐ จ ธมฺมา อสํกิลิฏฺฐสฺส ธมฺมสฺส อตฺถิปจฺจเยน ปจฺจโย. สหชาตํ, ปจฺฉาชาตํ, อาหารํ, อินฺทฺริยํ.  สหชาตา สํกิลิฏฺฐา ขนฺธา จ มหาภูตา จ จิตฺต\-สมุฏฺฐานานํ รูปานํ อตฺถิปจฺจเยน ปจฺจโย. ปจฺฉาชาตา สํกิลิฏฺฐา ขนฺธา จ กพฬีกาโร อาหาโร จ อิมสฺส กายสฺส อตฺถิปจฺจเยน ปจฺจโย. ปจฺฉาชาตา สํกิลิฏฺฐา ขนฺธา จ รูปชีวิตินฺทฺริยญฺจ กฏตฺตารูปานํ อตฺถิปจฺจเยน ปจฺจโย. (2)}

\prose{นตฺถิ\-ปจฺจเยน ปจฺจโย. วิคต\-ปจฺจเยน ปจฺจโย. อวิคต\-ปจฺจเยน ปจฺจโย.}

\csromanmarkh{1. ปจฺจยานุโลม}\csromanmarkhii{2. สงฺขฺยาวาร}\csromanheadernomark{2}{1. \textbf{ปจฺจยานุโลม 2. สงฺขฺยาวาร}}
\proseitem{67}{เหตุยา จตฺตาริ, อารมฺมเณ จตฺตาริ, อธิปติยา ปญฺจ, อนนฺตเร จตฺตาริ, สมนนฺตเร จตฺตาริ, สหชาเต ปญฺจ, อญฺญมญฺเญ เทฺว, นิสฺสเย สตฺต, อุปนิสฺสเย จตฺตาริ, ปุเรชาเต เทฺว, ปจฺฉาชาเต เทฺว, อาเสวเน เทฺว, กมฺเม จตฺตาริ, วิปาเก เอกํ, อาหาเร จตฺตาริ, อินฺทฺริเย จตฺตาริ, ฌาเน จตฺตาริ, มคฺเค จตฺตาริ, สมฺปยุตฺเต เทฺว, วิปฺปยุตฺเต ตีณิ, อตฺถิยา สตฺต, นตฺถิยา จตฺตาริ, วิคเต จตฺตาริ, อวิคเต สตฺต.}

\proseitem{68}{สํกิลิฏฺโฐ ธมฺโม สํกิลิฏฺฐสฺส ธมฺมสฺส อารมฺมณ\-ปจฺจเยน ปจฺจโย. สหชาต\-ปจฺจเยน ปจฺจโย. อุปนิสฺสย\-ปจฺจเยน ปจฺจโย. (1)}

\prose{สํกิลิฏฺโฐ ธมฺโม อสํกิลิฏฺฐสฺส ธมฺมสฺส อารมฺมณ\-ปจฺจเยน ปจฺจโย. สหชาต\-ปจฺจเยน ปจฺจโย. อุปนิสฺสย\-ปจฺจเยน ปจฺจโย. ปจฺฉาชาต\-ปจฺจเยน ปจฺจโย. กมฺมปจฺจเยน ปจฺจโย. (2)}

\prose{\csromanfolio{28}สํกิลิฏฺโฐ ธมฺโม สํกิลิฏฺฐสฺส จ อสํกิลิฏฺฐสฺส จ ธมฺมสฺส สหชาต\-ปจฺจเยน ปจฺจโย. (3)}

\proseitem{69}{อสํกิลิฏฺโฐ ธมฺโม อสํกิลิฏฺฐสฺส ธมฺมสฺส อารมฺมณ\-ปจฺจเยน ปจฺจโย. สหชาต\-ปจฺจเยน ปจฺจโย. อุปนิสฺสย\-ปจฺจเยน ปจฺจโย. ปุเรชาต\-ปจฺจเยน ปจฺจโย. ปจฺฉาชาต\-ปจฺจเยน ปจฺจโย. กมฺมปจฺจเยน ปจฺจโย. อาหารปจฺจเยน ปจฺจโย. อินฺทฺริย\-ปจฺจเยน ปจฺจโย. (1)}

\prose{อสํกิลิฏฺโฐ ธมฺโม สํกิลิฏฺฐสฺส ธมฺมสฺส อารมฺมณ\-ปจฺจเยน ปจฺจโย. อุปนิสฺสย\-ปจฺจเยน ปจฺจโย. ปุเรชาต\-ปจฺจเยน ปจฺจโย. (2)}

\prose{สํกิลิฏฺโฐ จ อสํกิลิฏฺโฐ จ ธมฺมา สํกิลิฏฺฐสฺส ธมฺมสฺส สหชาตํ. ปุเรชาตํ. (1)}

\prose{สํกิลิฏฺโฐ จ อสํกิลิฏฺโฐ จ ธมฺมา อสํกิลิฏฺฐสฺส ธมฺมสฺส สหชาตํ. ปจฺฉาชาตํ. อาหารํ. อินฺทฺริยํ. (2)}

\csromanmarkh{2. ปจฺจย\-ปจฺจนีย}\csromanmarkhii{2. สงฺขฺยาวาร}\csromanheadernomark{2}{2. ปจฺจย\-ปจฺจนีย 2. สงฺขฺยาวาร}
\proseitem{70}{นเหตุยา สตฺต, นอารมฺมเณ สตฺต, นอธิปติยา สตฺต, นอนนฺตเร สตฺต, นสมนนฺตเร สตฺต, นสหชาเต ปญฺจ, นอญฺญมญฺเญ ปญฺจ, นนิสฺสเย ปญฺจ, นอุปนิสฺสเย สตฺต, นปุเรชาเต ฉ, นปจฺฉาชาเต สตฺต ฯเปฯ นมคฺเค สตฺต, นสมฺปยุตฺเต ปญฺจ, นวิปฺปยุตฺเต จตฺตาริ, โนอตฺถิยา จตฺตาริ, โนนตฺถิยา สตฺต, โนวิคเต สตฺต, โนอวิคเต จตฺตาริ.}

\proseitem{71}{เหตุปจฺจยา นอารมฺมเณ จตฺตาริ, นอธิปติยา จตฺตาริ, นอนนฺตเร จตฺตาริ, นสมนนฺตเร จตฺตาริ, นอญฺญมญฺเญ เทฺว, นอุปนิสฺสเย จตฺตาริ ฯเปฯ นสมฺปยุตฺเต เทฺว, นวิปฺปยุตฺเต เทฺว, โนนตฺถิยา จตฺตาริ, โนวิคเต จตฺตาริ.}

\csromanmatikaanchor{mtk.6}
\csromanmatikaanchor{mtk.7}
\csromantocmarkref{section}{78. กิเลสสมฺปยุตฺตทุก}{mtk.6}
\bookmark[dest={mtk.6},level=1]{78. กิเลสสมฺปยุตฺตทุก}
\csromantocmarkref{section}{79. กิเลสสํกิเลสิกทุก}{mtk.7}
\bookmark[dest={mtk.7},level=1]{79. กิเลสสํกิเลสิกทุก}
\csromanmarkh{79. กิเลสสํกิเลสิกทุก}\csromanmarkhii{4. ปจฺจย\-ปจฺจนียา\-นุโลม}\csromanheadernomark{1}{79. กิเลสสํกิเลสิกทุก 4. ปจฺจย\-ปจฺจนียา\-นุโลม}
\proseitemwithcloser{72}{\csromanfolio{29}นเหตุปจฺจยา อารมฺมเณ จตฺตาริ, อธิปติยา ปญฺจ, (อนุโลมมาติกา.) ฯเปฯ อวิคเต สตฺต.}{สํกิลิฏฺฐทุกํ นิฏฺฐิตํ.}
\csromanmarkh{78. กิเลส\-สมฺปยุตฺตทุก}\csromanmarkhii{1. ปฏิจฺจวาร}\csromanheadernomark{1}{78. กิเลส\-สมฺปยุตฺตทุก 1. ปฏิจฺจวาร}
\proseitem{73}{กิเลส\-สมฺปยุตฺตํ ธมฺมํ ปฏิจฺจ กิเลส\-สมฺปยุตฺโต ธมฺโม อุปฺปชฺชติ เหตุปจฺจยา. กิเลส\-สมฺปยุตฺตํ เอกํ ขนฺธํ ปฏิจฺจ ตโย ธนฺธา ฯเปฯ เทฺว ขนฺเธ ฯเปฯ. (1)}

\prose{กิเลส\-สมฺปยุตฺตํ ธมฺมํ ปฏิจฺจ กิเลส\-วิปฺปยุตฺโต ธมฺโม อุปฺปชฺชติ เหตุปจฺจยา. กิเลส\-สมฺปยุตฺเต ขนฺเธ ปฏิจฺจ จิตฺต\-สมุฏฺฐานํ รูปํ. (2)}

\prosewithcloser{(กิเลส\-สมฺปยุตฺตทุกํ สํกิลิฏฺฐ\-ทุก\-สทิสํ, นินฺนานากรณํ.)}{กิเลส\-สมฺปยุตฺตทุกํ นิฏฺฐิตํ.}
\csromanmarkh{79. กิเลสสํกิเลสิกทุก}\csromanmarkhii{1. ปฏิจฺจวาร}\csromanheadernomark{2}{79. กิเลสสํกิเลสิกทุก 1. ปฏิจฺจวาร}
\proseitem{74}{กิเลสญฺเจว สํกิเลสิกญฺจ ธมฺมํ ปฏิจฺจ กิเลโส เจว สํกิเลสิโก จ ธมฺโม อุปฺปชฺชติ เหตุปจฺจยา. โลภํ ปฏิจฺจ โมโห, ทิฏฺฐิ, ถินํ, อุทฺธจฺจํ, อหิริกํ, อโนตฺตปฺปํ. (จกฺกํ.) (1)}

\prose{กิเลสญฺเจว สํกิเลสิกญฺจ ธมฺมํ ปฏิจฺจ สํกิเลสิโก เจว โน จ กิเลโส ธมฺโม อุปฺปชฺชติ เหตุปจฺจยา. กิเลเส ปฏิจฺจ สมฺปยุตฺตกา ขนฺธา จิตฺต\-สมุฏฺฐานญฺจ รูปํ. (2)}

\prose{กิเลสญฺเจว สํกิเลสิกญฺจ ธมฺมํ ปฏิจฺจ กิเลโส เจว สํกิเลสิโก จ สํกิเลสิโก เจว โน จ กิเลโส จ ธมฺมา อุปฺปชฺชนฺติ เหตุปจฺจยา. โลภํ ปฏิจฺจ โมโห, ทิฏฺฐิ, ถินํ, อุทฺธจฺจํ, อหิริกํ, อโนตฺตปฺปํ สมฺปยุตฺตกา จ ขนฺธา จิตฺต\-สมุฏฺฐานญฺจ รูปํ. (3)}

\prose{\csromanfolio{30}(เอวํ ปฏิจฺจวาโรปิ สหชาตวาโรปิ ปจฺจยวาโรปิ นิสฺสยวาโรปิ สํสฏฺฐ\-วาโรปิ สมฺปยุตฺตวาโรปิ กิเลส\-ทุก\-สทิสา. นินฺนานากรณํ. อามสนํ นานํ.)}

\csromanmarkh{79. กิเลสสํกิเลสิกทุก}\csromanmarkhii{7. ปญฺหาวาร}\csromanheadernomark{2}{79. \textbf{กิเลสสํกิเลสิกทุก 7. ปญฺหาวาร}}
\proseitem{75}{กิเลโส เจว สํกิเลสิโก จ ธมฺโม กิเลสสฺส เจว สํกิเลสิกสฺส จ ธมฺมสฺส เหตุปจฺจเยน ปจฺจโย. กิเลสา เจว สํกิเลสิกา จ เหตู สมฺปยุตฺตกานํ กิเลสานํ เหตุปจฺจเยน ปจฺจโย. (เอวํ จตฺตาริ, กิเลส\-ทุก\-สทิสํ.) (4)}

\proseitem{76}{กิเลโส เจว สํกิเลสิโก จ ธมฺโม กิเลสสฺส เจว สํกิเลสิกสฺส จ ธมฺมสฺส อารมฺมณ\-ปจฺจเยน ปจฺจโย. กิเลเส อารพฺภ กิเลสา อุปฺปชฺชนฺติ. (มูลํ กาตพฺพํ.) กิเลเส อารพฺภ สํกิเลสิกา เจว โน จ กิเลสา ขนฺธา อุปฺปชฺชนฺติ. (มูลํ กาตพฺพํ.) กิเลเส อารพฺภ กิเลสา จ สมฺปยุตฺตกา จ ขนฺธา อุปฺปชฺชนฺติ. (3)}

\proseitem{77}{สํกิเลสิโก เจว โน จ กิเลโส ธมฺโม สํกิเลสิกสฺส เจว โน จ กิเลสสฺส ธมฺมสฺส อารมฺมณ\-ปจฺจเยน ปจฺจโย. ทานํ ฯเปฯ สีลํ ฯเปฯ อุโปสถกมฺมํ ฯเปฯ. ปุพฺเพ สุจิณฺณานิ ฯเปฯ. ฌานา วุฏฺฐหิตฺวา ฌานํ ปจฺจเวกฺขติ, อสฺสาเทติ อภินนฺทติ, ตํ อารพฺภ ราโค อุปฺปชฺชติ ฯเปฯ อุทฺธจฺจํ อุปฺปชฺชติ. ฌาเน ปริหีเน วิปฺปฏิสาริสฺส โทมนสฺสํ อุปฺปชฺชติ. อริยา โคตฺรภุํ ปจฺจเวกฺขนฺติ, โวทานํ ปจฺจเวกฺขนฺติ. จกฺขุํ ฯเปฯ. วตฺถุํ สํกิเลสิเก เจว โน จ กิเลเส ขนฺเธ อนิจฺจโต ฯเปฯ โทมนสฺสํ อุปฺปชฺชติ. ทิพฺเพน จกฺขุนา รูปํ ปสฺสติ ฯเปฯ อาวชฺชนาย อารมฺมณ\-ปจฺจเยน ปจฺจโย.}

\vspace{\tipitakaparskip}
\csromancenter{(อิตเร เทฺว กิเลส\-ทุก\-สทิสา, ฆฏนา\-รมฺมณาปิ กิเลส\-ทุก\-สทิสา.)}
\csromanheader{2}{79. กิเลสสํกิเลสิกทุก \textbf{อธิปติ}}
\proseitem{78}{\csromanfolio{31}กิเลโส เจว สํกิเลสิโก จ ธมฺโม กิเลสสฺส เจว สํกิเลสิกสฺส จ ธมฺมสฺส อธิปติ\-ปจฺจเยน ปจฺจโย. อารมฺมณา\-ธิปติ. ตีณิ.}

\prose{สํกิเลสิโก เจว โน จ กิเลโส ธมฺโม สํกิเลสิกสฺส เจว โน จ กิเลสสฺส ธมฺมสฺส อธิปติ\-ปจฺจเยน ปจฺจโย. อารมฺมณา\-ธิปติ, สหชาตา\-ธิปติ.  อารมฺมณา\-ธิปติ\csromandash{}ทานํ ฯเปฯ สีลํ ฯเปฯ อุโปสถกมฺมํ ฯเปฯ. ปุพฺเพ สุจิณฺณานิ ฯเปฯ. ฌานา วุฏฺฐหิตฺวา ฌานํ ครุํ กตฺวา ปจฺจเวกฺขติ. เสกฺขา โคตฺรภุํ ครุํ กตฺวา ปจฺจเวกฺขนฺติ, โวทานํ ครุํ กตฺวา ปจฺจเวกฺขนฺติ. จกฺขุํ ฯเปฯ. วตฺถุํ สํกิเลสิเก เจว โน จ กิเลเส ขนฺเธ ครุํ กตฺวา สํกิเลสิกา เจว โน จ กิเลสา ขนฺธา อุปฺปชฺชนฺติ. สหชาตา\-ธิปติ\csromandash{}สํกิเลสิกา เจว โน จ กิเลสาธิปติ สมฺปยุตฺตกานํ ขนฺธานํ จิตฺตสมุฏฺฐานานญฺจ รูปานํ อธิปติ\-ปจฺจเยน ปจฺจโย. (อิตเร เทฺวปิ กิเลส\-ทุก\-สทิสา. ฆฏนาธิปติปิ.)}

\proseitem{79}{กิเลโส เจว สํกิเลสิโก จ ธมฺโม กิเลสสฺส เจว สํกิเลสิกสฺส จ ธมฺมสฺส อนนฺตร\-ปจฺจเยน ปจฺจโย. ตีณิ. (กิเลส\-ทุก\-สทิสา.)}

\prose{สํกิเลสิโก เจว โน จ กิเลโส ธมฺโม สํกิเลสิกสฺส เจว โน จ กิเลสสฺส ธมฺมสฺส อนนฺตร\-ปจฺจเยน ปจฺจโย. ปุริมา ปุริมา สํกิเลสิกา เจว โน จ กิเลสา ขนฺธา ปจฺฉิมานํ ปจฺฉิมานํ สํกิเลสิกาน\-ญฺเจว โน จ กิเลสานํ ขนฺธานํ อนนฺตร\-ปจฺจเยน ปจฺจโย. อนุโลมํ โคตฺรภุสฺส. อนุโลมํ โวทานสฺส. อาวชฺชนา สํกิเลสิกาน\-ญฺเจว โน จ กิเลสานํ ขนฺธานํ อนนฺตร\-ปจฺจเยน ปจฺจโย.}

\prosewithcloser{(อิตเร เทฺว อนนฺตรา กิเลส\-ทุก\-สทิสา, นินฺนานากรณา. ฆฏนานนฺตรมฺปิ สพฺเพ ปจฺจยา กิเลส\-ทุก\-สทิสา, นินฺนานากรณา. อุปนิสฺสเย โลกุตฺตรํ นตฺถิ, อิทํ ทุกํ กิเลส\-ทุก\-สทิสํ, นินฺนานากรณํ.)}{กิเลส\-สํกิเลสิกทุกํ นิฏฺฐิตํ.}
\csromanmatikaanchor{mtk.8}
\csromantocmarkref{section}{80. กิเลสสํกิลิฏฺฐทุก}{mtk.8}
\bookmark[dest={mtk.8},level=1]{80. กิเลสสํกิลิฏฺฐทุก}
\csromanmarkh{80. กิเลส\-สํกิลิฏฺฐทุก}\csromanmarkhii{1. ปฏิจฺจวาร}\csromanheadernomark{1}{80. กิเลส\-สํกิลิฏฺฐทุก 1. ปฏิจฺจวาร}
\proseitem{80}{\csromanfolio{32}กิเลสญฺเจว สํกิลิฏฺฐญฺจ ธมฺมํ ปฏิจฺจ กิเลโส เจว สํกิลิฏฺโฐ จ ธมฺโม อุปฺปชฺชติ เหตุปจฺจยา. โลภํ ปฏิจฺจ โมโห, ทิฏฺฐิ, ถินํ, อุทฺธจฺจํ, อหิริกํ, อโนตฺตปฺปํ. (จกฺกํ.) (1)}

\prose{กิเลสญฺจว สํกิลิฏฺฐญฺจ ธมฺมํ ปฏิจฺจ สํกิลิฏฺโฐ เจว โน จ กิเลโส ธมฺโม อุปฺปชฺชติ เหตุปจฺจยา. กิเลเส ปฏิจฺจ สมฺปยุตฺตกา ขนฺธา. (2)}

\prose{กิเลสญฺเจว สํกิลิฏฺฐญฺจ ธมฺมํ ปฏิจฺจ กิเลโส เจว สํกิลิฏฺโฐ จ สํกิลิฏฺโฐ เจว โน จ กิเลโส จ ธมฺมา อุปฺปชฺชติ เหตุปจฺจยา. โลภํ ปฏิจฺจ โมโห, ทิฏฺฐิ, ถินํ, อุทฺธจฺจํ, อหิริกํ, อโนตฺตปฺปํ สมฺปยุตฺตกา จ ขนฺธา. (จกฺกํ.) (3)}

\proseitem{81}{สํกิลิฏฺฐ\-ญฺเจว โน จ กิเลสํ ธมฺมํ ปฏิจฺจ สํกิลิฏฺโฐ เจว โน จ กิเลโส ธมฺโม อุปฺปชฺชติ เหตุปจฺจยา. สํกิลิฏฺฐ\-ญฺเจว โน จ กิเลสํ เอกํ ขนฺธํ ปฏิจฺจ ตโย ขนฺธา ฯเปฯ เทฺว ขนฺเธ ฯเปฯ. (1)}

\prose{สํกิลิฏฺฐ\-ญฺเจว โน จ กิเลสํ ธมฺมํ ปฏิจฺจ กิเลโส เจว สํกิลิฏฺโฐ จ ธมฺโม อุปฺปชฺชติ เหตุปจฺจยา. สํกิลิฏฺเฐ เจว โน จ กิเลเส ขนฺเธ ปฏิจฺจ กิเลสา. (2)}

\prose{สํกิลิฏฺฐ\-ญฺเจว โน จ กิเลสํ ธมฺมํ ปฏิจฺจ กิเลโส เจว สํกิลิฏฺโฐ จ สํกิลิฏฺโฐ เจว โน จ กิเลโส จ ธมฺมา อุปฺปชฺชนฺติ เหตุปจฺจยา. สํกิลิฏฺฐ\-ญฺเจว โน จ กิเลสํ เอกํ ขนฺธํ ปฏิจฺจ ตโย ขนฺธา กิเลสา จ ฯเปฯ เทฺว ขนฺเธ ฯเปฯ. (3)}

\proseitem{82}{กิเลสญฺเจว สํกิลิฏฺฐญฺจ สํกิลิฏฺฐ\-ญฺเจว โน จ กิเลสญฺจ ธมฺมํ ปฏิจฺจ กิเลโส เจว สํกิลิฏฺโฐ จ ธมฺโม อุปฺปชฺชติ เหตุปจฺจยา. โลภญฺจ สมฺปยุตฺตเก จ ขนฺเธ ปฏิจฺจ โมโห, ทิฏฺฐิ, ถินํ, อุทฺธจฺจํ, อหิริกํ, อโนตฺตปฺปํ. (จกฺกํ.) (1)}

\csromanheader{2}{80. กิเลส\-สํกิลิฏฺฐทุก}
\prose{\csromanfolio{33}กิเลสญฺเจว สํกิลิฏฺฐญฺจ สํกิลิฏฺฐ\-ญฺเจว โน จ กิเลสญฺจ ธมฺมํ ปฏิจฺจ สํกิลิฏฺโฐ เจว โน จ กิเลโส ธมฺโม อุปฺปชฺชติ เหตุปจฺจยา. สํกิลิฏฺฐ\-ญฺเจว โน จ กิเลสํ เอกํ ขนฺธญฺจ กิเลเส จ ปฏิจฺจ ตโย ขนฺธา ฯเปฯ เทฺว ขนฺเธ จ ฯเปฯ. (2)}

\prose{กิเลสญฺเจว สํกิลิฏฺฐญฺจ สํกิลิฏฺฐ\-ญฺเจว โน จ กิเลสญฺจ ธมฺมํ ปฏิจฺจ กิเลโส เจว สํกิลิฏฺโฐ จ สํกิลิฏฺโฐ เจว โน จ กิเลโส จ ธมฺมา อุปฺปชฺชนฺติ เหตุปจฺจยา. สํกิลิฏฺฐ\-ญฺเจว โน จ กิเลสํ เอกํ ขนฺธญฺจ โลภญฺจ ปฏิจฺจ ตโย ขนฺธา โมโห, ทิฏฺฐิ, ถินํ, อุทฺธจฺจํ, อหิริกํ, อโนตฺตปฺปํ ฯเปฯ เทฺว ขนฺเธ จ ฯเปฯ. (จกฺกํ.) (3)}

\csromanmarkh{1. ปจฺจยานุโลม}\csromanmarkhii{2. สงฺขฺยาวาร}\csromanheadernomark{2}{1. \textbf{ปจฺจยานุโลม 2. สงฺขฺยาวาร}}
\proseitem{83}{เหตุยา นว, อารมฺมเณ นว, (สพฺพตฺถ นว,) กมฺเม นว, อาหาเร นว ฯเปฯ อวิคเต นว.}

\csromanmarkh{2. ปจฺจย\-ปจฺจนีย}\csromanmarkhii{1. วิภงฺควาร}\csromanheadernomark{2}{2. ปจฺจย\-ปจฺจนีย 1. วิภงฺควาร}
\proseitem{84}{กิเลสญฺเจว สํกิลิฏฺฐญฺจ ธมฺมํ ปฏิจฺจ กิเลโส เจว สํกิลิฏฺโฐ จ ธมฺโม อุปฺปชฺชติ นเหตุปจฺจยา. วิจิกิจฺฉํ ปฏิจฺจ วิจิกิจฺฉา\-สหคโต โมโห. อุทฺธจฺจํ ปฏิจฺจ อุทฺธจฺจ\-สหคโต โมโห. (1)}

\prose{สํกิลิฏฺฐ\-ญฺเจว โน จ กิเลสํ ธมฺมํ ปฏิจฺจ กิเลโส เจว สํกิลิฏฺโฐ จ ธมฺโม อุปฺปชฺชติ นเหตุปจฺจยา. วิจิกิจฺฉา\-สหคเต อุทฺธจฺจ\-สหคเต ขนฺเธ ปฏิจฺจ วิจิกิจฺฉา\-สหคโต อุทฺธจฺจ\-สหคโต โมโห. (1)}

\prose{กิลิสญฺเจว สํกิลิฏฺฐญฺจ สํกิลิฏฺฐ\-ญฺเจว โน จ กิเลสญฺจ ธมฺมํ ปฏิจฺจ กิเลโส เจว สํกิลิฏฺโฐ จ ธมฺโม อุปฺปชฺชติ นเหตุปจฺจยา. วิจิกิจฺฉา\-สหคเต อุทฺธจฺจ\-สหคเต ขนฺเธ จ วิจิกิจฺฉญฺจ อุทฺธจฺจญฺจ ปฏิจฺจ วิจิกิจฺฉา\-สหคโต อุทฺธจฺจ\-สหคโต โมโห. (1)}

\csromanmarkh{2. ปจฺจย\-ปจฺจนีย}\csromanmarkhii{2. สงฺขฺยาวาร}\csromanheadernomark{2}{2. ปจฺจย\-ปจฺจนีย 2. สงฺขฺยาวาร}
\proseitem{85}{\csromanfolio{34}นเหตุยา ตีณิ, นอธิปติยา นว, นปุเรชาเต นว, นปจฺฉาชาเต นว, นอาเสวเน นว, นกมฺเม ตีณิ, นวิปาเก นว, นวิปฺปยุตฺเต นว.}

\prose{(เอวํ อิตเร เทฺว คณนาปิ สหชาตวาโรปิ ปจฺจยวาโรปิ นิสฺสยวาโรปิ สํสฏฺฐ\-วาโรปิ สมฺปยุตฺตวาโรปิ ปฏิจฺจ\-วาร\-สทิสา.)}

\csromanmarkh{80. กิเลส\-สํกิลิฏฺฐทุก}\csromanmarkhii{7. ปญฺหาวาร}\csromanheadernomark{2}{80. กิเลส\-สํกิลิฏฺฐทุก 7. ปญฺหาวาร}
\proseitem{86}{กิเลโส เจว สํกิลิฏฺโฐ จ ธมฺโม กิเลสสฺส เจว สํกิลิฏฺฐสฺส จ ธมฺมสฺส เหตุปจฺจเยน ปจฺจโย. กิเลสา เจว สํกิลิฏฺฐา จ เหตู สมฺปยุตฺตกานํ กิเลสานํ เหตุปจฺจเยน ปจฺจโย. (1)}

\prose{กิเลโส เจว สํกิลิฏฺโฐ จ ธมฺโม สํกิลิฏฺฐสฺส เจว โน จ กิเลสสฺส ธมฺมสฺส เหตุปจฺจเยน ปจฺจโย. กิเลสา เจว สํกิลิฏฺฐา จ เหตู สมฺปยุตฺตกานํ ขนฺธานํ เหตุปจฺจเยน ปจฺจโย. (2)}

\prose{กิเลโส เจว สํกิลิฏฺโฐ จ ธมฺโม กิเลสสฺส เจว สํกิลิฏฺฐสฺส จ สํกิลิฏฺฐสฺส เจว โน จ กิเลสสฺส จ ธมฺมสฺส เหตุปจฺจเยน ปจฺจโย. กิเลสา เจว สํกิลิฏฺฐา จ เหตู สมฺปยุตฺตกานํ ขนฺธานํ กิเลสานญฺจ เหตุปจฺจเยน ปจฺจโย. (3)}

\proseitem{87}{กิเลโส เจว สํกิลิฏฺโฐ จ ธมฺโม กิเลสสฺส เจว สํกิลิฏฺฐสฺส จ ธมฺมสฺส อารมฺมณ\-ปจฺจเยน ปจฺจโย. กิเลเส อารพฺภ กิเลสา อุปฺปชฺชนฺติ. (มูลํ กาตพฺพํ.) กิเลเส อารพฺภ สํกิลิฏฺฐา เจว โน จ กิเลสา ขนฺธา อุปฺปชฺชนฺติ. (มูลํ กาตพฺพํ.) กิเลเส อารพฺภ กิเลสา จ สมฺปยุตฺตกา ขนฺธา จ อุปฺปชฺชนฺติ. (3)}

\csromanheader{2}{80. กิเลส\-สํกิลิฏฺฐทุก}
\proseitem{88}{\csromanfolio{35}สํกิลิฏฺโฐ เจว โน จ กิเลโส ธมฺโม สํกิลิฏฺฐสฺส เจว โน จ กิเลสสฺส ธมฺมสฺส อารมฺมณ\-ปจฺจเยน ปจฺจโย. สํกิลิฏฺเฐ เจว โน จ กิเลเส ขนฺเธ อารพฺภ สํกิลิฏฺฐา เจว โน จ กิเลสา ขนฺธา อุปฺปชฺชนฺติ. (มูลํ กาตพฺพํ.) สํกิลิฏฺเฐ เจว โน จ กิเลเส ขนฺเธ อารพฺภ กิเลสา อุปฺปชฺชนฺติ. (มูลํ กาตพฺพํ.) สํกิลิฏฺเฐ เจว โน จ กิเลเส ขนฺเธ อารพฺภ กิเลสา จ สมฺปยุตฺตกา ขนฺธา จ อุปฺปชฺชนฺติ. (3) (อิตเรปิ ตีณิ กาตพฺพา.)}

\csromanheader{2}{\textbf{อธิปติ}}
\proseitem{89}{กิเลโส เจว สํกิลิฏฺโฐ จ ธมฺโม กิเลสสฺส เจว สํกิลิฏฺฐสฺส จ ธมฺมสฺส อธิปติ\-ปจฺจเยน ปจฺจโย. อารมฺมณา\-ธิปติ. ตีณิ.}

\prose{สํกิลิฏฺโฐ เจว โน จ กิเลโส ธมฺโม สํกิลิฏฺฐสฺส เจว โน จ กิเลสสฺส ธมฺมสฺส อธิปติ\-ปจฺจเยน ปจฺจโย. อารมฺมณา\-ธิปติ, สหชาตา\-ธิปติ.  อารมฺมณา\-ธิปติ\csromandash{}สํกิลิฏฺเฐ เจว โน จ กิเลเส ขนฺเธ ครุํ กตฺวา ฯเปฯ. ตีณิ. (เทฺว อธิปติ. ตีณิปิ กาตพฺพา. อิตเร เทฺวปิ ตีณิ กาตพฺพา.)}

\proseitem{90}{กิเลโส เจว สํกิลิฏฺโฐ จ ธมฺโม กิเลสสฺส เจว สํกิลิฏฺฐสฺส จ ธมฺมสฺส อนนฺตร\-ปจฺจเยน ปจฺจโย. (นวปิ กาตพฺพา. อาวชฺชนาปิ วุฏฺฐานมฺปิ นตฺถิ.) สมนนฺตร\-ปจฺจเยน ปจฺจโย. สหชาต\-ปจฺจเยน ปจฺจโย. อญฺญมญฺญ\-ปจฺจเยน ปจฺจโย. นิสฺสย\-ปจฺจเยน ปจฺจโย. อุปนิสฺสย\-ปจฺจเยน ปจฺจโย. นว ปญฺหา. (ปุเรชาต\-ปจฺจโย ปจฺฉาชาตปจฺจโยปิ นตฺถิ.) อาเสวน\-ปจฺจเยน ปจฺจโย.}

\csromanheader{2}{\textbf{กมฺมาทิ}}
\proseitem{91}{สํกิลิฏฺโฐ เจว โน จ กิเลโส ธมฺโม สํกิลิฏฺฐสฺส เจว โน จ กิเลสสฺส ธมฺมสฺส กมฺมปจฺจเยน ปจฺจโย. สํกิลิฏฺฐา เจว โน จ กิเลสา เจตนา สมฺปยุตฺตกานํ ขนฺธานํ กมฺมปจฺจเยน ปจฺจโย. (1)}

\prose{สํกิลิฏฺโฐ เจว โน จ กิเลโส ธมฺโม กิเลสสฺส เจว สํกิลิฏฺฐสฺส จ ธมฺมสฺส กมฺมปจฺจเยน ปจฺจโย. สํกิลิฏฺฐา เจว, \csromanfolio{36}โน จ กิเลสา เจตนา สมฺปยุตฺตกานํ กิเลสานํ กมฺมปจฺจเยน ปจฺจโย. (2)}

\prose{สํกิลิฏฺโฐ เจว โน จ กิเลโส ธมฺโม กิเลสสฺส เจว สํกิลิฏฺฐสฺส จ สํกิลิฏฺฐสฺส เจว โน จ กิเลสสฺส ธมฺมสฺส กมฺมปจฺจเยน ปจฺจโย. สํกิลิฏฺฐา เจว โน จ กิเลสา เจตนา สมฺปยุตฺตกานํ ขนฺธานํ กิเลสานญฺจ กมฺมปจฺจเยน ปจฺจโย. (3)}

\prose{อาหารปจฺจเยน ปจฺจโย. ตีณิ. อินฺทฺริย\-ปจฺจเยน ปจฺจโย. ตีณิ. ฌานปจฺจเยน ปจฺจโย. ตีณิ. มคฺคปจฺจเยน ปจฺจโย. นว. สมฺปยุตฺต\-ปจฺจเยน ปจฺจโย. นว. อตฺถิปจฺจเยน ปจฺจโย. นว. นตฺถิ\-ปจฺจเยน ปจฺจโย. วิคต\-ปจฺจเยน ปจฺจโย. อวิคต\-ปจฺจเยน ปจฺจโย. นว.}

\csromanmarkh{1. ปจฺจยานุโลม}\csromanmarkhii{2. สงฺขฺยาวาร}\csromanheadernomark{2}{1. \textbf{ปจฺจยานุโลม 2. สงฺขฺยาวาร}}
\proseitem{91}{เหตุยา ตีณิ, อารมฺมเณ นว, อธิปติยา นว, อนนฺตเร นว, สมนนฺตเร นว, สหชาเต นว, อญฺญมญฺเญ นว, นิสฺสเย นว, อุปนิสฺสเย นว, อาเสวเน นว, กมฺเม ตีณิ, อาหาเร ตีณิ, อินฺทฺริเย ตีณิ, ฌาเน ตีณิ, มคฺเค นว, สมฺปยุตฺเต นว, อตฺถิยา นว, นตฺถิยา นว, วิคเต นว, อวิคเต นว.}

\prose{กิเลโส เจว สํกิลิฏฺโฐ จ ธมฺโม กิเลสสฺส เจว สํกิลิฏฺฐสฺส จ ธมฺมสฺส อารมฺมณ\-ปจฺจเยน ปจฺจโย. สหชาต\-ปจฺจเยน ปจฺจโย. อุปนิสฺสย\-ปจฺจเยน ปจฺจโย. (นวปิ, ตีณิเยว ปทา กาตพฺพา.)}

\csromanmarkh{2. ปจฺจย\-ปจฺจนีย}\csromanmarkhii{2. สงฺขฺยาวาร}\csromanheadernomark{2}{2. ปจฺจย\-ปจฺจนีย 2. สงฺขฺยาวาร}
\vspace{\tipitakaparskip}
\csromancenter{นเหตุยา นว, นอารมฺมเณ นว, (สพฺพตฺถ นว,) โนอวิคเต นว.}
\csromanmatikaanchor{mtk.9}
\csromanmatikaanchor{mtk.10}
\csromantocmarkref{section}{81. กิเลสกิเลสสมฺปยุตฺตทุก}{mtk.9}
\bookmark[dest={mtk.9},level=1]{81. กิเลสกิเลสสมฺปยุตฺตทุก}
\csromantocmarkref{section}{82. กิเลสวิปฺปยุตฺตสํกิเลสิกทุก}{mtk.10}
\bookmark[dest={mtk.10},level=1]{82. กิเลสวิปฺปยุตฺตสํกิเลสิกทุก}
\csromanmarkh{82. กิเลส\-วิปฺปยุตฺต\-สํกิเลสิกทุก}\csromanmarkhii{3. ปจฺจยา\-นุโลม\-ปจฺจนีย}\csromanheadernomark{1}{82. กิเลส\-วิปฺปยุตฺต\-สํกิเลสิกทุก 3. ปจฺจยา\-นุโลม\-ปจฺจนีย}
\proseitem{95}{\csromanfolio{37}เหตุปจฺจยา นอารมฺมเณ ตีณิ, นอธิปติยา ตีณิ, นอนนฺตเร ตีณิ, นสมนนฺตเร ตีณิ, นอุปนิสฺสเย ตีณิ ฯเปฯ นมคฺเค ตีณิ ฯเปฯ นวิปฺปยุตฺเต ตีณิ, โนนตฺถิยา ตีณิ, โนวิคเต ตีณิ.}

\csromanheader{2}{4. ปจฺจย\-ปจฺจนียา\-นุโลม}
\csromanheader{2}{96. นเหตุปจฺจยา อารมฺมเณ นว, (อนุโลมมาติกา กาตพฺพา) ฯเปฯ อวิคเต นว.}
\nitthitam{กิเลส\-สํกิลิฏฺฐทุกํ นิฏฺฐิตํ.}
\csromanmarkh{81. กิเลส\-กิเลส\-สมฺปยุตฺตทุก}\csromanmarkhii{1. ปฏิจฺจวาร}\csromanheadernomark{1}{81. กิเลส\-กิเลส\-สมฺปยุตฺตทุก 1. ปฏิจฺจวาร}
\proseitem{97}{กิเลสญฺเจว กิเลส\-สมฺปยุตฺตญฺจ ธมฺมํ ปฏิจฺจ กิเลโส เจว กิเลส\-สมฺปยุตฺโต จ ธมฺโม อุปฺปชฺชติ เหตุปจฺจยา. โลภํ ปฏิจฺจ โมโห, ทิฏฺฐิ, ถินํ, อุทฺธจฺจํ, อหิริกํ, อโนตฺตปฺปํ.}

\vspace{\tipitakaparskip}
\csromancenter{(กิเลส\-สํกิลิฏฺฐ\-ทุก\-สทิสํ, นินฺนานากรณํ, สพฺเพ วารา.)}
\nitthitam{กิเลส\-กิเลส\-สมฺปยุตฺตทุกํ นิฏฺฐิตํ.}
\csromanmarkh{82. กิเลส\-วิปฺปยุตฺต\-สํกิเลสิกทุก}\csromanmarkhii{1. ปฏิจฺจวาร}\csromanheadernomark{2}{82. กิเลส\-วิปฺปยุตฺต\-สํกิเลสิกทุก 1. ปฏิจฺจวาร}
\proseitemwithcloser{98}{กิเลส\-วิปฺปยุตฺตํ สํกิเลสิกํ ธมฺมํ ปฏิจฺจ กิเลส\-วิปฺปยุตฺโต สํกิเลสิโก ธมฺโม อุปฺปชฺชติ เหตุปจฺจยา. กิเลส\-วิปฺปยุตฺตํ สํกิเลสิกํ เอกํ ขนฺธํ ปฏิจฺจ ตโย ขนฺธา จิตฺต\-สมุฏฺฐานญฺจ รูปํ ฯเปฯ เทฺว ขนฺเธ ฯเปฯ.}{(ยถา โลกิยทุกํ, เอวํ นินฺนานากรณํ.) กิเลส\-วิปฺปยุตฺต\-สํกิเลสิกทุกํ นิฏฺฐิตํ.}
\nitthitam{กิเลส\-โคจฺฉกํ นิฏฺฐิตํ.}
\csromanmatikaanchor{mtk.11}
\csromantocmarknopage{chapter}{13. ปิฏฺฐิทุก}{mtk.11}
\bookmark[dest={mtk.11},level=0]{13. ปิฏฺฐิทุก}
\chapterheadpageread{13. \textbf{ปิฏฺฐิทุก}}
\csromanmatikaanchor{mtk.12}
\csromantocmarkref{section}{83. ทสฺสเนนปหาตพฺพทุก}{mtk.12}
\bookmark[dest={mtk.12},level=1]{83. ทสฺสเนนปหาตพฺพทุก}
\csromanmarkh{83. ทสฺสเนน\-ปหาตพฺพทุก}\csromanmarkhii{1. ปฏิจฺจวาร}\csromanheadernomark{1}{83. ทสฺสเนน\-ปหาตพฺพทุก 1. ปฏิจฺจวาร}
\proseitem{1}{\csromanfolio{38}ทสฺสเนน ปหาตพฺพํ ธมฺมํ ปฏิจฺจ ทสฺสเนน ปหาตพฺโพ ธมฺโม อุปฺปชฺชติ เหตุปจฺจยา. ทสฺสเนน ปหาตพฺพํ เอกํ ขนฺธํ ปฏิจฺจ ตโย ขนฺธา ฯเปฯ เทฺว ขนฺเธ ฯเปฯ. (1)}

\prose{ทสฺสเนน ปหาตพฺพํ ธมฺมํ ปฏิจฺจ นทสฺสเนน ปหาตพฺโพ ธมฺโม อุปฺปชฺชติ เหตุปจฺจยา. ทสฺสเนน ปหาตพฺเพ ขนฺเธ ปฏิจฺจ จิตฺต\-สมุฏฺฐานํ รูปํ. (2)}

\prose{ทสฺสเนน ปหาตพฺพํ ธมฺมํ ปฏิจฺจ ทสฺสเนน ปหาตพฺโพ จ นทสฺสเนน ปหาตพฺโพ จ ธมฺมา อุปฺปชฺชนฺติ เหตุปจฺจยา. ทสฺสเนน ปหาตพฺพํ เอกํ ขนฺธํ ปฏิจฺจ ตโย ขนฺธา จิตฺต\-สมุฏฺฐานญฺจ ฯเปฯ เทฺว ขนฺเธ ฯเปฯ. (3)}

\prose{นทสฺสเนน ปหาตพฺพํ ธมฺมํ ปฏิจฺจ นทสฺสเนน ปหาตพฺโพ ธมฺโม อุปฺปชฺชติ เหตุปจฺจยา. นทสฺสเนน ปหาตพฺพํ เอกํ ขนฺธํ ปฏิจฺจ ตโย ขนฺธา จิตฺต\-สมุฏฺฐานญฺจ รูปํ ฯเปฯ เทฺว ขนฺเธ ฯเปฯ. ปฏิสนฺธิ\-กฺขเณ ฯเปฯ. ขนฺเธ ปฏิจฺจ วตฺถุ, วตฺถุํ ปฏิจฺจ ขนฺธา. เอกํ มหาภูตํ ฯเปฯ. (1)}

\prose{ทสฺสเนน ปหาตพฺพญฺจ นทสฺสเนน ปหาตพฺพญฺจ ธมฺมํ ปฏิจฺจ นทสฺสเนน ปหาตพฺโพ ธมฺโม อุปฺปชฺชติ เหตุปจฺจยา. ทสฺสเนน ปหาตพฺเพ ขนฺเธ จ มหาภูเต จ ปฏิจฺจ จิตฺต\-สมุฏฺฐานํ รูปํ. (สํขิตฺตํ.) (1)}

\csromanmarkh{1. ปจฺจยานุโลม}\csromanmarkhii{2. สงฺขฺยาวาร}\csromanheadernomark{2}{1. \textbf{ปจฺจยานุโลม 2. สงฺขฺยาวาร}}
\proseitem{2}{\textbf{เหตุ}ยา ปญฺจ, อารมฺมเณ เทฺว, อธิปติยา ปญฺจ, อนนฺตเร เทฺว, สมนนฺตเร เทฺว, สหชาเต ปญฺจ, อญฺญมญฺเญ เทฺว, นิสฺสเย ปญฺจ, อุปนิสฺสเย เทฺว, ปุเรชาเต เทฺว, อาเสวเน เทฺว, กมฺเม ปญฺจ, วิปาเก เอกํ, อาหาเร ปญฺจ ฯเปฯ อวิคเต ปญฺจ.}

\csromanmarkh{83. ทสฺสเนน\-ปหาตพฺพทุก}\csromanmarkhii{2. ปจฺจย\-ปจฺจนีย 1. วิภงฺควาร}\csromanheadernomark{2}{83. ทสฺสเนน\-ปหาตพฺพทุก 2. ปจฺจย\-ปจฺจนีย 1. วิภงฺควาร}
\proseitem{3}{\csromanfolio{39}ทสฺสเนน ปหาตพฺพํ ธมฺมํ ปฏิจฺจ ทสฺสเนน ปหาตพฺโพ ธมฺโม อุปฺปชฺชติ นเหตุปจฺจยา. วิจิกิจฺฉา\-สหคเต ขนฺเธ ปฏิจฺจ วิจิกิจฺฉา\-สหคโต โมโห. (1)}

\prose{นทสฺสเนน ปหาตพฺพํ ธมฺมํ ปฏิจฺจ นทสฺสเนน ปหาตพฺโพ ธมฺโม อุปชฺชติ นเหตุปจฺจยา. อเหตุกํ นทสฺสเนน ปหาตพฺพํ เอกํ ขนฺธํ ฯเปฯ. อเหตุก\-ปฏิสนฺธิ\-กฺขเณ ฯเปฯ. (ยาว อสญฺญสตฺตา.) อุทฺธจฺจ\-สหคเต ขนฺเธ ปฏิจฺจ อุทฺธจฺจ\-สหคโต โมโห. (1)}

\csromanmarkh{2. ปจฺจย\-ปจฺจนีย}\csromanmarkhii{2. สงฺขฺยาวาร}\csromanheadernomark{2}{2. ปจฺจย\-ปจฺจนีย 2. สงฺขฺยาวาร}
\vspace{\tipitakaparskip}
\csromancenter{\textbf{นเหตุ}ยา เทฺว, นอารมฺมเณ ตีณิ, นอธิปติยา ปญฺจ, นอนนฺตเร ตีณิ, นสมนนฺตเร ตีณิ, นอญฺญมญฺเญ ตีณิ, นฺเอาปนิสฺสเย ตีณิ, นปุเรชาเต จตฺตาริ, นปจฺฉาชาเต ปญฺจ, นอาเสวเน ปญฺจ, นกมฺเม เทฺว, นวิปาเก ปญฺจ, นอาหาเร เอกํ, ไนนฺทฺริเย เอกํ, นฌาเน เอกํ, นมคฺเค เอกํ, นสมฺปยุตฺเต ตีณิ, นวิปฺปยุตฺเต เทฺว, โนนตฺถิยา ตีณิ, โนวิคเต ตีณิ.}
\csromancenter{(อิตเร เทฺว คณนาปิ สหชาตวาโรปิ กาตพฺโพ.)}
\csromanmarkh{83. ทสฺสเนน\-ปหาตพฺพทุก}\csromanmarkhii{3. ปจฺจยวาร}\csromanheadernomark{2}{83. ทสฺสเนน\-ปหาตพฺพทุก 3. ปจฺจยวาร}
\proseitem{5}{ทสฺสเนน ปหาตพฺพํ ธมฺมํ ปจฺจยา ทสฺสเนน ปหาตพฺโพ ธมฺโม อุปฺปชฺชติ เหตุปจฺจยา. ตีณิ. (ปฏิจฺจสทิสา.)}

\proseitem{6}{นทสฺสเนน ปหาตพฺพํ ธมฺมํ ปจฺจยา นทสฺสเนน ปหาตพฺโพ ธมฺโม อุปฺปชฺชติ เหตุปจฺจยา. นทสฺสเนน ปหาตพฺพํ เอกํ ขนฺธํ ปจฺจยา ตโย ขนฺธา \csromanfolio{40}จิตฺต\-สมุฏฺฐานญฺจ รูปํ ฯเปฯ เทฺว ขนฺเธ ฯเปฯ. ปฏิสนฺธิ\-กฺขเณ ฯเปฯ. (ยาว อชฺฌตฺติกา มหาภูตา.) วตฺถุํ ปจฺจยา นทสฺสเนน ปหาตพฺพา ขนฺธา. (1)}

\prose{นทสฺสเนน ปหาตพฺพํ ธมฺมํ ปจฺจยา ทสฺสเนน ปหาตพฺโพ ธมฺโม อุปฺปชฺชติ เหตุปจฺจยา. วตฺถุํ ปจฺจยา ทสฺสเนน ปหาตพฺพา ขนฺธา. (2)}

\prose{นทสฺสเนน ปหาตพฺพํ ธมฺมํ ปจฺจยา ทสฺสเนน ปหาตพฺโพ จ นทสฺสเนน ปหาตพฺโพ จ ธมฺมา อุปฺปชฺชนฺติ เหตุปจฺจยา. วตฺถุํ ปจฺจยา ทสฺสเนน ปหาตพฺพา ขนฺธา. มหาภูเต ปจฺจยา จิตฺต\-สมุฏฺฐานํ. (3)}

\proseitem{7}{ทสฺสเนน ปหาตพฺพญฺจ นทสฺสเนน ปหาตพฺพญฺจ ธมฺมํ ปจฺจยา ทสฺสเนน ปหาตพฺโพ ธมฺโม อุปฺปชฺชติ เหตุปจฺจยา. ทสฺสเนน ปหาตพฺพํ เอกํ ขนฺธญฺจ วตฺถุญฺจ ปจฺจยา ตโย ขนฺธา ฯเปฯ เทฺว ขนฺเธ จ ฯเปฯ. (1)}

\prose{ทสฺสเนน ปหาตพฺพญฺจ นทสฺสเนน ปหาตพฺพญฺจ ธมฺมํ ปจฺจยา นทสฺสเนน ปหาตพฺโพ ธมฺโม อุปฺปชฺชติ เหตุปจฺจยา. ทสฺสเนน ปหาตพฺเพ ขนฺเธ จ มหาภูเต จ ปจฺจยา จิตฺต\-สมุฏฺฐานํ รูปํ. (2)}

\prose{ทสฺสเนน ปหาตพฺพญฺจ นทสฺสเนน ปหาตพฺพญฺจ ธมฺมํ ปจฺจยา ทสฺสเนน ปหาตพฺโพ จ นทสฺสเนน ปหาตพฺโพ จ ธมฺมา อุปฺปชฺชนฺติ เหตุปจฺจยา. ทสฺสเนน ปหาตพฺพํ เอกํ ขนฺธญฺจ วตฺถุญฺจ ปจฺจยา ตโย ขนฺธา ฯเปฯ เทฺว ขนฺเธ จ ฯเปฯ. ทสฺสเนน ปหาตพฺเพ ขนฺเธ จ มหาภูเต จ ปจฺจยา จิตฺต\-สมุฏฺฐานํ รูปํ. (3)}

\csromanmarkh{1. ปจฺจยานุโลม}\csromanmarkhii{2. สงฺขฺยาวาร}\csromanheadernomark{2}{1. \textbf{ปจฺจยานุโลม 2. สงฺขฺยาวาร}}
\vspace{\tipitakaparskip}
\csromancenter{เหตุยา นว, อารมฺมเณ จตฺตาริ, อธิปติยา นว ฯเปฯ อวิคเต นว.}
\csromanmarkh{2. ปจฺจย\-ปจฺจนีย}\csromanmarkhii{1. วิภงฺควาร}\csromanheadernomark{2}{2. ปจฺจย\-ปจฺจนีย 1. วิภงฺควาร}
\proseitem{9}{ทสฺสเนน ปหาตพฺพํ ธมฺมํ ปจฺจยา ทสฺสเนน ปหาตพฺโพ ธมฺโม อุปฺปชฺชติ นเหตุปจฺจยา. วิจิกิจฺฉา\-สหคเต ขนฺเธ ปจฺจยา วิจิกิจฺฉา\-สหคโต โมโห. (1)}

\prose{\csromanfolio{41}นทสฺสเนน ปหาตพฺพํ ธมฺมํ ปจฺจยา นทสฺสเนน ปหาตพฺโพ ธมฺโม อุปฺปชฺชติ นเหตุปจฺจยา. อเหตุกํ นทสฺสเนน ปหาตพฺพํ เอกํ ขนฺธํ ปจฺจยา ตโย ขนฺธา จิตฺต\-สมุฏฺฐานญฺจ รูปํ ฯเปฯ เทฺว ขนฺเธ ฯเปฯ. (ยาว อสญฺญสตฺตา.) จกฺขายตนํ ปจฺจยา จกฺขุวิญฺญาณํ ฯเปฯ. กายายตนํ ปจฺจยา กายวิญฺญาณํ. วตฺถุํ ปจฺจยา อเหตุกา นทสฺสเนน ปหาตพฺพา ขนฺธา. อุทฺธจฺจ\-สหคเต ขนฺเธ ปจฺจยา อุทฺธสหคโต โมโห. วตฺถุํ ปจฺจยา อุทฺธจฺจ\-สหคโต โมโห. (1)}

\prose{นทสฺสเนน ปหาตพฺพํ ธมฺมํ ปจฺจยา ทสฺสเนน ปหาตพฺโพ ธมฺโม อุปฺปชฺชติ นเหตุปจฺจยา. วตฺถุํ ปจฺจยา วิจิกิจฺฉา\-สหคโต โมโห. (2)}

\prose{ทสฺสเนน ปหาตพฺพญฺจ นทสฺสเนน ปหาตพฺพญฺจ ธมฺมํ ปจฺจยา ทสฺสเนน ปหาตพฺโพ ธมฺโม อุปฺปชฺชติ นเหตุปจฺจยา. วิจิกิจฺฉา\-สหคเต ขนฺเธ จ วตฺถุญฺจ ปจฺจยา วิจิกิจฺฉา\-สหคโต โมโห. (1)}

\csromanmarkh{2. ปจฺจย\-ปจฺจนีย}\csromanmarkhii{2. สงฺขฺยาวาร}\csromanheadernomark{2}{2. ปจฺจย\-ปจฺจนีย 2. สงฺขฺยาวาร}
\proseitem{10}{นเหตุยา จตฺตาริ, นอารมฺมเณ ตีณิ, นอธิปติยา นว, นอนนฺตเร ตีณิ, นสมนนฺตเร ตีณิ, นอญฺญมญฺเญ ตีณิ, นอุปนิสฺสเย ตีณิ, นปุเรชาเต จตฺตาริ, นปจฺฉาชาเต นว, นอาเสวเน นว, นกมฺเม จตฺตาริ, นวิปาเก นว, นอาหาเร เอกํ, นอินฺทฺริเย เอกํ, นฌาเน เอกํ, นมคฺเค เอกํ, นสมฺปยุตฺเต ตีณิ, นวิปฺปยุตฺเต เทฺว, โนนตฺถิยา ตีณิ, โนวิคเต ตีณิ.}

\vspace{\tipitakaparskip}
\csromancenter{(อิตเร เทฺว คณนาปิ นิสฺสยวาโรปิ กาตพฺโพ.)}
\csromanmarkh{83. ทสฺสเนน\-ปหาตพฺพทุก}\csromanmarkhii{5. สํสฏฺฐวาร}\csromanheadernomark{2}{83. ทสฺสเนน\-ปหาตพฺพทุก 5. สํสฏฺฐวาร}
\proseitem{11}{ทสฺสเนน ปหาตพฺพํ ธมฺมํ สํสฏฺโฐ ทสฺสเนน ปหาตพฺโพ ธมฺโม อุปฺปชฺชติ เหตุปจฺจยา. (สํขิตฺตํ.) (1)}

\prose{\csromanfolio{42}นทสฺสเนน ปหาตพฺพํ ธมฺมํ สํสฏฺโฐ นทสฺสเนน ปหาตพฺโพ ธมฺโม อุปฺปชฺชติ เหตุปจฺจยา. (สํขิตฺตํ.) (1)}

\prose{เหตุยา เทฺว, อารมฺมเณ เทฺว, อธิปติยา เทฺว, อวิคเต เทฺว. อนุโลมํ.}

\proseitem{12}{ทสฺสเนน ปหาตพฺพํ ธมฺมํ สํสฏฺโฐ ทสฺสเนน ปหาตพฺโพ ธมฺโม อุปฺปชฺชติ นเหตุปจฺจยา. วิจิกิจฺฉา\-สหคเต ขนฺเธ สํสฏฺโฐ วิจิกิจฺฉา\-สหคโต โมโห. (1)}

\prose{นทสฺสเนน ปหาตพฺพํ ธมฺมํ สํสฏฺโฐ นทสฺสเนน ปหาตพฺโพ ธมฺโม อุปฺปชฺชติ นเหตุปจฺจยา. (สํขิตฺตํ.) (1)}

\prose{นเหตุยา เทฺว, นอธิปติยา เทฺว, นปุเรชาเต เทฺว, นปจฺฉาชาเต เทฺว, นอาเสวเน เทฺว, นกมฺเม เทฺว, นวิปาเก เทฺว, นฌาเน เอกํ, นมคฺเค เอกํ, นวิปฺปยุตฺเต เทฺว.}

\vspace{\tipitakaparskip}
\csromancenter{ปจฺจนียํ.}
\csromancenter{(เอวํ อิตเร เทฺว คณนาปิ สมฺปยุตฺตวาโรปิ กาตพฺโพ.)}
\csromanmarkh{83. ทสฺสเนน\-ปหาตพฺพทุก}\csromanmarkhii{7. ปญฺหาวาร}\csromanheadernomark{2}{83. ทสฺสเนน\-ปหาตพฺพทุก 7. ปญฺหาวาร}
\proseitem{13}{ทสฺสเนน ปหาตพฺโพ ธมฺโม ทสฺสเนน ปหาตพฺพสฺส ธมฺมสฺส เหตุปจฺจเยน ปจฺจโย. ตีณิ.}

\prose{นทสฺสเนน ปหาตพฺโพ ธมฺโม นทสฺสเนน ปหาตพฺพสฺส ธมฺมสฺส เหตุปจฺจเยน ปจฺจโย. นทสฺสเนน ปหาตพฺพา เหตู สมฺปยุตฺตกานํ ขนฺธานํ จิตฺต\-สมุฏฺฐานญฺจ รูปานํ เหตุปจฺจเยน ปจฺจโย. ปฏิสนฺธิ\-กฺขเณ ฯเปฯ. (1)}

\csromanheader{2}{83. ทสฺสเนน\-ปหาตพฺพทุก อารมฺมณ}
\proseitem{14}{\csromanfolio{43}ทสฺสเนน ปหาตพฺโพ ธมฺโม ทสฺสเนน ปหาตพฺพสฺส ธมฺมสฺส อารมฺมณ\-ปจฺจเยน ปจฺจโย. ทสฺสเนน ปหาตพฺพํ ราคํ อสฺสาเทติ อภินนฺทติ, ตํ อารพฺภ ทสฺสเนน ปหาตพฺโพ ราโค อุปฺปชฺชติ. ทิฏฺฐิ ฯเปฯ วิจิกิจฺฉา อุปฺปชฺชติ. ทสฺสเนน ปหาตพฺพํ โทมนสฺสํ อุปฺปชฺชติ. ทสฺสเนน ปหาตพฺพํ ทิฏฺฐิํ อสฺสาเทติ อภินนฺทติ, ตํ อารพฺภ ทสฺสเนน ปหาตพฺโพ ราโค อุปฺปชฺชติ. ทิฏฺฐิ ฯเปฯ วิจิกิจฺฉา ฯเปฯ โทมนสฺสํ อุปฺปชฺชติ. วิจิกิจฺฉํ อารพฺภ วิจิกิจฺฉา อุปฺปชฺชติ. ทิฏฺฐิ ฯเปฯ โทมนสฺสํ อุปฺปชฺชติ. ทสฺสเนน ปหาตพฺพํ โทมนสฺสํ อารพฺภ ทสฺสเนน ปหาตพฺพํ โทมนสฺสํ อุปฺปชฺชติ. ทิฏฺฐิ อุปฺปชฺชติ. วิจิกิจฺฉา อุปฺปชฺชติ. (1)}

\prose{ทสฺสเนน ปหาตพฺโพ ธมฺโม นทสฺสเนน ปหาตพฺพสฺส ธมฺมสฺส อารมฺมณ\-ปจฺจเยน ปจฺจโย. อริยา ทสฺสเนน ปหาตพฺเพ ปหีเน กิเลเส ปจฺจเวกฺขนฺติ. ปุพฺเพ สมุทาจิณฺเณ ฯเปฯ. ทสฺสเนน ปหาตพฺเพ ขนฺเธ อนิจฺจโต ฯเปฯ วิปสฺสติ. เจโตปริย\-ญาเณน ทสฺสเนน ปหาตพฺพ\-จิตฺต\-สมงฺคิสฺส จิตฺตํ ชานาติ. ทสฺสเนน ปหาตพฺพา ขนฺธา เจโตปริย\-ญาณสฺส, ปุพฺเพ\-นิวาสา\-นุสฺสติ\-ญาณสฺส, ยถากมฺมูปคญาณสฺส, อนาคตํสญาณสฺส, อาวชฺชนาย อารมฺมณ\-ปจฺจเยน ปจฺจโย. (2)}

\proseitem{15}{นทสฺสเนน ปหาตพฺโพ ธมฺโม นทสฺสเนน ปหาตพฺพสฺส ธมฺมสฺส อารมฺมณ\-ปจฺจเยน ปจฺจโย. ทานํ ฯเปฯ สีลํ ฯเปฯ อุโปสถกมฺมํ กตฺวา ตํ ปจฺจเวกฺขติ, อสฺสาเทติ อภินนฺทติ, ตํ อารพฺภ นทสฺสเนน ปหาตพฺโพ ราโค อุปฺปชฺชติ. อุทฺธจฺจํ ฯเปฯ. นทสฺสเนน ปหาตพฺพํ โทมนสฺสํ อุปฺปชฺชติ. ปุพฺเพ สุจิณฺณานิ ฯเปฯ. ฌานา ฯเปฯ. อริยา มคฺคา วุฏฺฐหิตฺวา มคฺคํ ปจฺจเวกฺขนฺติ ฯเปฯ ผลสฺส, อาวชฺชนาย อารมฺมณ\-ปจฺจเยน ปจฺจโย. อริยา นทสฺสเนน ปหาตพฺเพ ปหีเน กิเลเส ปจฺจเวกฺขนฺติ. วิกฺขมฺภิเต กิเลเส ฯเปฯ. ปุพฺเพ ฯเปฯ. จกฺขุํ ฯเปฯ. วตฺถุํ นทสฺสเนน ปหาตพฺเพ ขนฺเธ อนิจฺจโต ฯเปฯ วิปสฺสติ, อสฺสาเทติ อภินนฺทติ, ตํ อารพฺภ นทสฺสเนน ปหาตพฺโพ ราโค อุปฺปชฺชติ. อุทฺธจฺจํ ฯเปฯ. นทสฺสเนน ปหาตพฺพํ โทมนสฺสํ อุปฺปชฺชติ. ทิพฺเพน จกฺขุนา รูปํ ปสฺสติ ฯเปฯ อนาคตํสญาณสฺส, อาวชฺชนาย อารมฺมณ\-ปจฺจเยน ปจฺจโย. (1)}

\prose{\csromanfolio{44}นทสฺสเนน ปหาตพฺโพ ธมฺโม ทสฺสเนน ปหาตพฺพสฺส ธมฺมสฺส อารมฺมณ\-ปจฺจเยน ปจฺจโย. ทานํ ฯเปฯ สีลํ ฯเปฯ อุโปสถกมฺมํ ฯเปฯ. ปุพฺเพ ฯเปฯ. ฌานา ฯเปฯ. จกฺขุํ ฯเปฯ. วตฺถุํ นทสฺสเนน ปหาตพฺเพ ขนฺเธ อสฺสาเทติ อภินนฺทติ, ตํ อารพฺภ ทสฺสเนน ปหาตพฺโพ ราโค อุปฺปชฺชติ. ทิฏฺฐิ ฯเปฯ วิจิกิจฺฉา ฯเปฯ ทสฺสเนน ปหาตพฺพํ โทมนสฺสํ อุปฺปชฺชติ. (2)}

\csromanheader{2}{\textbf{อธิปติ}}
\proseitem{16}{ทสฺสเนน ปหาตพฺโพ ธมฺโม ทสฺสเนน ปหาตพฺพสฺส ธมฺมสฺส อธิปติ\-ปจฺจเยน ปจฺจโย. อารมฺมณา\-ธิปติ, สหชาตา\-ธิปติ.  อารมฺมณา\-ธิปติ\csromandash{}ทสฺสเนน ปหาตพฺพํ ราคํ ครุํ กตฺวา อสฺสาเทติ อภินนฺทติ, ตํ ครุํ กตฺวา ทสฺสเนน ปหาตพฺโพ ราโค อุปฺปชฺชติ. ทิฏฺฐิ อุปฺปชฺชติ. ทิฏฺฐิํ ครุํ กตฺวา ฯเปฯ. สหชาตา\-ธิปติ\csromandash{}ทสฺสเนน ปหาตพฺพา\-ธิปติ สมฺปยุตฺตกานํ ขนฺธานํ อธิปติ\-ปจฺจเยน ปจฺจโย. (1)}

\prose{ทสฺสเนน ปหาตพฺโพ ธมฺโม นทสฺสเนน ปหาตพฺพสฺส ธมฺมสฺส อธิปติ\-ปจฺจเยน ปจฺจโย. สหชาตา\-ธิปติ\csromandash{}ทสฺสเนน ปหาตพฺพา\-ธิปติ จิตฺต\-สมุฏฺฐานานํ รูปานํ อธิปติ\-ปจฺจเยน ปจฺจโย. (2)}

\prose{ทสฺสเนน ปหาตพฺโพ ธมฺโม ทสฺสเนน ปหาตพฺพสฺส จ นทสฺสเนน ปหาตพฺพสฺส จ ธมฺมสฺส อธิปติ\-ปจฺจเยน ปจฺจโย. สหชาตา\-ธิปติ\csromandash{} ทสฺสเนน ปหาตพฺพา\-ธิปติ สมฺปยุตฺตกานํ ขนฺธานํ จิตฺตสมุฏฺฐานานญฺจ รูปานํ อธิปติ\-ปจฺจเยน ปจฺจโย. (3)}

\proseitem{17}{นทสฺสเนน ปหาตพฺโพ ธมฺโม นทสฺสเนน ปหาตพฺพสฺส ธมฺมสฺส อธิปติ\-ปจฺจเยน ปจฺจโย. อารมฺมณา\-ธิปติ, สหชาตา\-ธิปติ.  อารมฺมณา\-ธิปติ\csromandash{}ทานํ ฯเปฯ สีลํ ฯเปฯ อุโปสถกมฺมํ ฯเปฯ. ปุพฺเพ ฯเปฯ. ฌานา วุฏฺฐหิตฺวา ฌานํ ครุํ กตฺวา ปจฺจเวกฺขติ, อสฺสาเทติ อภินนฺทติ, ตํ ครุํ กตฺวา นทสฺสเนน ปหาตพฺโพ ราโค อุปฺปชฺชติ อริยา มคฺคา วุฏฺฐหิตฺวา มคฺคํ ครุํ กตฺวา ฯเปฯ ผลสฺส อธิปติ\-ปจฺจเยน ปจฺจโย. จกฺขุํ ฯเปฯ. วตฺถุํ นทสฺสเนน ปหาตพฺเพ ขนฺเธ ครุํ กตฺวา อสฺสาเทติ อภินนฺทติ, ตํ ครุํ กตฺวา นทสฺสเนน ปหาตพฺโพ ราโค อุปฺปชฺชติ.  สหชาตา\-ธิปติ\csromandash{}นทสฺสเนน ปหาตพฺพา\-ธิปติ สมฺปยุตฺตกานํ ขนฺธานํ จิตฺตสมุฏฺฐานานญฺจ รูปานํ อธิปติ\-ปจฺจเยน ปจฺจโย. (1)}

\prose{\csromanfolio{45}นทสฺสเนน ปหาตพฺโพ ธมฺโม ทสฺสเนน ปหาตพฺพสฺส ธมฺมสฺส อธิปติ\-ปจฺจเยน ปจฺจโย. อารมฺมณา\-ธิปติ\csromandash{}ทานํ ทตฺวา ฯเปฯ ฌานํ ฯเปฯ. จกฺขุํ ฯเปฯ. วตฺถุํ นทสฺสเนน ปหาตพฺเพ ขนฺเธ ครุํ กตฺวา อสฺสาเทติ อภินนฺทติ, ตํ ครุํ กตฺวา ทสฺสเนน ปหาตพฺโพ ราโค อุปฺปชฺชติ. ทิฏฺฐิ อุปฺปชฺชติ. (2)}

\csromanheader{2}{\textbf{อนนฺตร}}
\proseitem{18}{ทสฺสเนน ปหาตพฺโพ ธมฺโม ทสฺสเนน ปหาตพฺพสฺส ธมฺมสฺส อนนฺตร\-ปจฺจเยน ปจฺจโย. ปุริมา ปุริมา ทสฺสเนน ปหาตพฺพา ขนฺธา ปจฺฉิมานํ ปจฺฉิมานํ ทสฺสเนน ปหาตพฺพานํ ขนฺธานํ อนนฺตร\-ปจฺจเยน ปจฺจโย. (1)}

\prose{ทสฺสเนน ปหาตพฺโพ ธมฺโม นทสฺสเนน ปหาตพฺพสฺส ธมฺมสฺส อนนฺตร\-ปจฺจเยน ปจฺจโย. ทสฺสเนน ปหาตพฺพา ขนฺธา วุฏฺฐานสฺส อนนฺตร\-ปจฺจเยน ปจฺจโย. (มูลํ กาตพฺพํ.) ปุริมา ปุริมา นทสฺสเนน ปหาตพฺพา ขนฺธา ฯเปฯ ผลสมาปตฺติยา อนนฺตร\-ปจฺจเยน ปจฺจโย. (มูลํ กาตพฺพํ.) อาวชฺชานา ทสฺสเนน ปหาตพฺพานํ ขนฺธานํ อนนฺตร\-ปจฺจเยน ปจฺจโย. (2)}

\prose{สมนนฺตร\-ปจฺจเยน ปจฺจโย. สหชาต\-ปจฺจเยน ปจฺจโย. ปญฺจ. อญฺญมญฺญ\-ปจฺจเยน ปจฺจโย. เทฺว. นิสฺสย\-ปจฺจเยน ปจฺจโย. สตฺต.}

\csromanheader{2}{\textbf{อุปนิสฺสย}}
\proseitem{19}{ทสฺสเนน ปหาตพฺโพ ธมฺโม ทสฺสเนน ปหาตพฺพสฺส ธมฺมสฺส อุปนิสฺสย\-ปจฺจเยน ปจฺจโย. อารมฺมณู\-ปนิสฺสโย, อนนฺตรู\-ปนิสฺสโย, ปกตู\-ปนิสฺสโย ฯเปฯ. ปกตู\-ปนิสฺสโย\csromandash{}ทสฺสเนน ปหาตพฺพํ ราคํ อุปนิสฺสาย ปาณํ หนติ ฯเปฯ สํฆํ ภินฺทติ. ทสฺสเนน ปหาตพฺพํ โทสํ. โมหํ. ทิฏฺฐิํ. ปตฺถนํ อุปนิสฺสาย ปาณํ หนติ ฯเปฯ สํฆํ ภินฺทติ. ทสฺสเนน ปหาตพฺโพ ราโค ฯเปฯ ปตฺถนา ทสฺสเนน ปหาตพฺพสฺส ราคสฺส ฯเปฯ ปตฺถนาย อุปนิสฺสย\-ปจฺจเยน ปจฺจโย. (1)}

\prose{ทสฺสเนน ปหาตพฺโพ ธมฺโม นทสฺสเนน ปหาตพฺพสฺส ธมฺมสฺส อุปนิสฺสย\-ปจฺจเยน ปจฺจโย. อนนฺตรู\-ปนิสฺสโย, ปกตู\-ปนิสฺสโย ฯเปฯ. ปกตู\-ปนิสฺสโย\csromandash{}ทสฺสเนน ปหาตพฺพํ ราคํ อุปนิสฺสาย ทานํ เทติ ฯเปฯ สมาปตฺติํ อุปฺปาเทติ. ทสฺสเนน ปหาตพฺพํ โทสํ. โมหํ. ทิฏฺฐิํ. ปตฺถนํ \csromanfolio{46}อุปนิสฺสาย ทานํ เทติ ฯเปฯ สมาปตฺติํ อุปฺปาเทติ. ทสฺสเนน ปหาตพฺโพ ราโค ฯเปฯ ปตฺถนา สทฺธาย ฯเปฯ ปญฺญาย. นทสฺสเนน ปหาตพฺพสฺส ราคสฺส. โทสสฺส. โมหสฺส. มานสฺส. ปตฺถนาย. กายิกสฺส สุขสฺส ฯเปฯ ผลสมาปตฺติยา อุปนิสฺสย\-ปจฺจเยน ปจฺจโย. (2)}

\proseitem{20}{นทสฺสเนน ปหาตพฺโพ ธมฺโม นทสฺสเนน ปหาตพฺพสฺส ธมฺมสฺส อุปนิสฺสย\-ปจฺจเยน ปจฺจโย. อารมฺมณู\-ปนิสฺสโย, อนนฺตรู\-ปนิสฺสโย, ปกตู\-ปนิสฺสโย ฯเปฯ. ปกตู\-ปนิสฺสโย\csromandash{}สทฺธํ อุปนิสฺสาย ทานํ เทติ ฯเปฯ สมาปตฺติํ อุปฺปาเทติ. มานํ ชปฺเปติ. สีลํ ฯเปฯ. ปญฺญํ. นทสฺสเนน ปหาตพฺพํ ราคํ. โทสํ. โมหํ. มานํ. ปตฺถนํ. กายิกํ สุขํ ฯเปฯ. เสนาสนํ อุปนิสฺสาย ทานํ เทติ ฯเปฯ สมาปตฺติํ อุปฺปาเทติ. มานํ ชปฺเปติ. สทฺธา ฯเปฯ. ปญฺญา. นทสฺสเนน ปหาตพฺโพ ราโค ฯเปฯ ปตฺตนา. กายิกํ สุขํ ฯเปฯ. เสนาสนํ สทฺธาย ฯเปฯ ปญฺญาย. นทสฺสเนน ปหาตพฺพสฺส ราคสฺส ฯเปฯ ปตฺถนาย. กายิกสฺส สุขสฺส. กายิกสฺส ทุกฺขสฺส. มคฺคสฺส, ผลสมาปตฺติยา อุปนิสฺสย\-ปจฺจเยน ปจฺจโย. (1)}

\prose{นทสฺสเนน ปหาตพฺโพ ธมฺโม ทสฺสเนน ปหาตพฺพสฺส ธมฺมสฺส อุปนิสฺสย\-ปจฺจเยน ปจฺจโย. อารมฺมณู\-ปนิสฺสโย, อนนฺตรู\-ปนิสฺสโย, ปกตู\-ปนิสฺสโย ฯเปฯ. ปกตู\-ปนิสฺสโย\csromandash{}สทฺธํ อุปนิสฺสาย ทิฏฺฐิํ คณฺหาติ. สีลํ ฯเปฯ. ปญฺญํ อุปนิสฺสาย ทิฏฺฐิํ คณฺหาติ. นทสฺสเนน ปหาตพฺพํ ราคํ. โทสํ. โมหํ. มานํ. ปตฺถนํ. กายิกํ สุขํ ฯเปฯ. เสนาสนํ อุปนิสฺสาย ปาณํ หนติ ฯเปฯ สํฆํ ภินฺทติ. สทฺธา ฯเปฯ. เสนาสนํ ทสฺสเนน ปหาตพฺพสฺส ราคสฺส. โทสสฺส. โมหสฺส. ทิฏฺฐิยา. ปตฺถนาย อุปนิสฺสย\-ปจฺจเยน ปจฺจโย. (2)}

\csromanheader{2}{\textbf{ปุเรชาต}}
\proseitem{21}{นทสฺสเนน ปหาตพฺโพ ธมฺโม นทสฺสเนน ปหาตพฺพสฺส ธมฺมสฺส ปุเรชาต\-ปจฺจเยน ปจฺจโย. อารมฺมณ\-ปุเรชาตํ, วตฺถุ\-ปุเรชาตํ.  อารมฺมณ\-ปุเรชาตํ\csromandash{}จกฺขุํ ฯเปฯ. วตฺถุํ อนิจฺจโต ฯเปฯ วิปสฺสติ, อสฺสาเทติ อภินนฺทติ, ตํ อารพฺภ นทสฺสเนน ปหาตพฺโพ ราโค ฯเปฯ อุทฺธจฺจํ ฯเปฯ. นทสฺสเนน ปหาตพฺพํ โทมนสฺสํ อุปฺปชฺชติ. ทิพฺเพน จกฺขุนา รูปํ ปสฺสติ. ทิพฺพาย โสตธาตุยา สทฺทํ สุณาติ. รูปายตนํ จกฺขุ\-วิญฺญาณสฺส ฯเปฯ. โผฏฺฐพฺพา\-ยตนํ กายวิญฺญาณสฺส ฯเปฯ. วตฺถุ\-ปุเรชาตํ\csromandash{}จกฺขายตนํ \csromanfolio{47}จกฺขุ\-วิญฺญาณสฺส ฯเปฯ. กายายตนํ กายวิญฺญาณสฺส ฯเปฯ. วตฺถุ นทสฺสเนน ปหาตพฺพานํ ขนฺธานํ ปุเรชาต\-ปจฺจเยน ปจฺจโย. (1)}

\prose{นทสฺสเนน ปหาตพฺโพ ธมฺโม ทสฺสเนน ปหาตพฺพสฺส ธมฺมสฺส ปุเรชาต\-ปจฺจเยน ปจฺจโย. อารมฺมณ\-ปุเรชาตํ, วตฺถุ\-ปุเรชาตํ.  อารมฺมณ\-ปุเรชาตํ\csromandash{}จกฺขุํ ฯเปฯ. วตฺถุํ อสฺสาเทติ อภินนฺทติ, ตํ อารพฺภ ทสฺสเนน ปหาตพฺโพ ราโค ฯเปฯ ทิฏฺฐิ ฯเปฯ วิจิกิจฺฉา ฯเปฯ. ทสฺสเนน ปหาตพฺพํ โทมนสฺสํ อุปฺปชฺชติ. วตฺถุ\-ปุเรชาตํ\csromandash{}วตฺถุ ทสฺสเนน ปหาตพฺพานํ ขนฺธานํ ปุเรชาต\-ปจฺจเยน ปจฺจโย. (2)}

\csromanheader{2}{ปจฺฉาชาตา\-เสวน}
\proseitem{22}{ทสฺสเนน ปหาตพฺโพ ธมฺโม นทสฺสเนน ปหาตพฺพสฺส ธมฺมสฺส ปจฺฉาชาต\-ปจฺจเยน ปจฺจโย. (สํขิตฺตํ.) (1)}

\prose{นทสฺสเนน ปหาตพฺโพ ธมฺโม นทสฺสเนน ปหาตพฺพสฺส ธมฺมสฺส ปจฺฉาชาต\-ปจฺจเยน ปจฺจโย. (สํขิตฺตํ.) อาเสวน\-ปจฺจเยน ปจฺจโย. เทฺว.}

\csromanheader{2}{\textbf{กมฺม}}
\proseitem{23}{ทสฺสเนน ปหาตพฺโพ ธมฺโม ทสฺสเนน ปหาตพฺพสฺส ธมฺมสฺส กมฺมปจฺจเยน ปจฺจโย. ทสฺสเนน ปหาตพฺพา เจตนา สมฺปยุตฺตกานํ ขนฺธานํ กมฺมปจฺจเยน ปจฺจโย. (1)}

\prose{ทสฺสเนน ปหาตพฺโพ ธมฺโม นทสฺสเนน ปหาตพฺพสฺส ธมฺมสฺส กมฺมปจฺจเยน ปจฺจโย. สหชาตา, นานากฺขณิกา.  สหชาตา ทสฺสเนน ปหาตพฺพา เจตนา จิตฺต\-สมุฏฺฐานานํ รูปานํ กมฺมปจฺจเยน ปจฺจโย. นานากฺขณิกา ทสฺสเนน ปหาตพฺพา เจตนา วิปากานํ ขนฺธานํ กฏตฺตา จ รูปานํ กมฺมปจฺจเยน ปจฺจโย. (2)}

\prose{ทสฺสเนน ปหาตพฺโพ ธมฺโม ทสฺสเนน ปหาตพฺพสฺส จ นทสฺสเนน ปหาตพฺพสฺส จ ธมฺมสฺส กมฺมปจฺจเยน ปจฺจโย. ทสฺสเนน ปหาตพฺพา เจตนา สมฺปยุตฺตกานํ ขนฺธานํ จิตฺตสมุฏฺฐานานญฺจ รูปานํ กมฺมปจฺจเยน ปจฺจโย. (3)}

\proseitem{24}{\csromanfolio{48}นทสฺสเนน ปหาตพฺโพ ธมฺโม นทสฺสเนน ปหาตพฺพสฺส ธมฺมสฺส กมฺมปจฺจเยน ปจฺจโย. สหชาตา, นานากฺขณิกา.  สหชาตา นทสฺสเนน ปหาตพฺพา เจตนา สมฺปยุตฺตกานํ ขนฺธานํ จิตฺตสมุฏฺฐานานญฺจ รูปานํ กมฺมปจฺจเยน ปจฺจโย. ปฏิสนฺธิ\-กฺขเณ ฯเปฯ. นานากฺขณิกา นทสฺสเนน ปหาตพฺพา เจตนา วิปากานํ ขนฺธานํ กฏตฺตา จ รูปานํ กมฺมปจฺจเยน ปจฺจโย.}

\csromanheader{2}{\textbf{วิปากาทิ}}
\proseitem{25}{นทสฺสเนน ปหาตพฺโพ ธมฺโม นทสฺสเนน ปหาตพฺพสฺส ธมฺมสฺส วิปาก\-ปจฺจเยน ปจฺจโย. เอกํ. อาหารปจฺจเยน ปจฺจโย. จตฺตาริ. อินฺทฺริย\-ปจฺจเยน ปจฺจโย. จตฺตาริ. ฌานปจฺจเยน ปจฺจโย. จตฺตาริ. มคฺคปจฺจเยน ปจฺจโย. จตฺตาริ. สมฺปยุตฺต\-ปจฺจเยน ปจฺจโย. เทฺว.}

\csromanheader{2}{\textbf{วิปฺปยุตฺต}}
\proseitem{26}{ทสฺสเนน ปหาตพฺโพ ธมฺโม นทสฺสเนน ปหาตพฺพสฺส ธมฺมสฺส วิปฺปยุตฺต\-ปจฺจเยน ปจฺจโย. สหชาตํ, ปจฺฉาชาตํ. (สํขิตฺตํ.) (1)}

\prose{นทสฺสเนน ปหาตพฺโพ ธมฺโม นทสฺสเนน ปหาตพฺพสฺส ธมฺมสฺส วิปฺปยุตฺต\-ปจฺจเยน ปจฺจโย. สหชาตํ, ปุเรชาตํ, ปจฺฉาชาตํ. (สํขิตฺตํ.)}

\prose{นทสฺสเนน ปหาตพฺโพ ธมฺโม ทสฺสเนน ปหาตพฺพสฺส ธมฺมสฺส วิปฺปยุตฺต\-ปจฺจเยน ปจฺจโย. ปุเรชาตํ วตฺถุ ทสฺสเนน ปหาตพฺพานํ ขนฺธานํ วิปฺปยุตฺต\-ปจฺจเยน ปจฺจโย. (2)}

\prose{อตฺถิอาทิ}

\proseitem{27}{ทสฺสเนน ปหาตพฺโพ ธมฺโม ทสฺสเนน ปหาตพฺพสฺส ธมฺมสฺส อตฺถิปจฺจเยน ปจฺจโย. (ปฏิจฺจสทิสํ.) ทสฺสเนน ปหาตพฺโพ ธมฺโม นทสฺสเนน ปหาตพฺพสฺส ธมฺมสฺส อตฺถิปจฺจเยน ปจฺจโย. สหชาตํ, ปจฺฉาชาตํ. (สํขิตฺตํ.) ทสฺสเนน ปหาตพฺโพ ธมฺโม ทสฺสเนน ปหาตพฺพสฺส จ นทสฺสเนน ปหาตพฺพสฺส จ ธมฺมสฺส อตฺถิปจฺจเยน ปจฺจโย. (ปฏิจฺจ\-วาร\-สทิสา.) (3)}

\prose{\csromanfolio{49}นทสฺสเนน ปหาตพฺโพ ธมฺโม นทสฺสเนน ปหาตพฺพสฺส ธมฺมสฺส อตฺถิปจฺจเยน ปจฺจโย. สหชาตํ, ปุเรชาตํ, ปจฺฉาชาตํ, อาหารํ, อินฺทฺริยํ. (สํขิตฺตํ.) นทสฺสเนน ปหาตพฺโพ ธมฺโม ทสฺสเนน ปหาตพฺพสฺส ธมฺมสฺส อตฺถิปจฺจเยน ปจฺจโย. ปุเรชาตํ\csromandash{}จกฺขุํ ฯเปฯ. (สํขิตฺตํ. ปุเรชาต\-สทิสํ.) (2)}

\proseitem{28}{ทสฺสเนน ปหาตพฺโพ จ นทสฺสเนน ปหาตพฺโพ จ ธมฺมา ทสฺสเนน ปหาตพฺพสฺส ธมฺมสฺส อตฺถิปจฺจเยน ปจฺจโย. \textbf{สหชาตํ,} ปุเรชาตํ.  สหชาโต ทสฺสเนน ปหาตพฺโพ เอโก ขนฺโธ จ วตฺถุ จ ติณฺณนฺนํ ขนฺธานํ อตฺถิปจฺจเยน ปจฺจโย ฯเปฯ เทฺว ขนฺธา จ ฯเปฯ. (1)}

\prose{ทสฺสเนน ปหาตพฺโพ จ นทสฺสเนน ปหาตพฺโพ จ ธมฺมา นทสฺสเนน ปหาตพฺพสฺส ธมฺมสฺส อตฺถิปจฺจเยน ปจฺจโย. สหชาตํ, ปจฺฉาชาตํ, อาหารํ, อินฺทฺริยํ.  สหชาตา ทสฺสเนน ปหาตพฺพา ขนฺธา จ มหาภูตา จ จิตฺต\-สมุฏฺฐานานํ รูปานํ อตฺถิปจฺจเยน ปจฺจโย. ปจฺฉาชาตา ทสฺสเนน ปหาตพฺพา ขนฺธา จ กพฬีกาโร อาหาโร จ อิมสฺส กายสฺส อตฺถิปจฺจเยน ปจฺจโย. ปจฺฉาชาตา ทสฺสเนน ปหาตพฺพา ขนฺธา จ รูปชีวิตินฺทฺริยญฺจ กฏตฺตารูปานํ อตฺถิปจฺจเยน ปจฺจโย. (2)}

\prose{นตฺถิ\-ปจฺจเยน ปจฺจโย. วิคต\-ปจฺจเยน ปจฺจโย. อวิคต\-ปจฺจเยน ปจฺจโย.}

\csromanmarkh{1. ปจฺจยานุโลม}\csromanmarkhii{2. สงฺขฺยาวาร}\csromanheadernomark{2}{1. \textbf{ปจฺจยานุโลม 2. สงฺขฺยาวาร}}
\proseitem{29}{เหตุยา จตฺตาริ, อารมฺมเณ จตฺตาริ, อธิปติยา ปญฺจ, อนนฺตเร จตฺตาริ, สมนนฺตเร จตฺตาริ, สหชาเต ปญฺจ, อญฺญมญฺเญ เทฺว, นิสฺสเย สตฺต, อุปนิสฺสเย จตฺตาริ, ปุเรชาเต เทฺว, ปจฺฉาชาเต เทฺว, อาเสวเน เทฺว, กมฺเม จตฺตาริ, วิปาเก เอกํ, อาหาเร จตฺตาริ, อินฺทฺริเย จตฺตาริ, ฌาเน จตฺตาริ, มคฺเค จตฺตาริ, สมฺปยุตฺเต เทฺว, วิปฺปยุตฺเต ตีณิ, อตฺถิยา สตฺต, นตฺถิยา จตฺตาริ, วิคเต จตฺตาริ, อวิคเต สตฺต.}

\proseitem{30}{\csromanfolio{50}ทสฺสเนน ปหาตพฺโพ ธมฺโม ทสฺสเนน ปหาตพฺพสฺส ธมฺมสฺส อารมฺมณ\-ปจฺจเยน ปจฺจโย. สหชาต\-ปจฺจเยน ปจฺจโย. อุปนิสฺสย\-ปจฺจเยน ปจฺจโย. (1)}

\prose{ทสฺสเนน ปหาตพฺโพ ธมฺโม นทสฺสเนน ปหาตพฺพสฺส ธมฺมสฺส อารมฺมณ\-ปจฺจเยน ปจฺจโย. สหชาต\-ปจฺจเยน ปจฺจโย. อุปนิสฺสย\-ปจฺจเยน ปจฺจโย. ปจฺฉาชาต\-ปจฺจเยน ปจฺจโย. กมฺมปจฺจเยน ปจฺจโย. (2)}

\prose{ทสฺสเนน ปหาตพฺโพ ธมฺโม ทสฺสเนน ปหาตพฺพสฺส จ นทสฺสเนน ปหาตพฺพสฺส จ ธมฺมสฺส สหชาต\-ปจฺจเยน ปจฺจโย. (3)}

\proseitem{31}{นทสฺสเนน ปหาตพฺโพ ธมฺโม นทสฺสเนน ปหาตพฺพสฺส ธมฺมสฺส อารมฺมณ\-ปจฺจเยน ปจฺจโย. สหชาต\-ปจฺจเยน ปจฺจโย. อุปนิสฺสย\-ปจฺจเยน ปจฺจโย. ปุเรชาต\-ปจฺจเยน ปจฺจโย. ปจฺฉาชาต\-ปจฺจเยน ปจฺจโย. กมฺมปจฺจเยน ปจฺจโย. อาหารปจฺจเยน ปจฺจโย. อินฺทฺริย\-ปจฺจเยน ปจฺจโย. (1)}

\prose{นทสฺสเนน ปหาตพฺโพ ธมฺโม ทสฺสเนน ปหาตพฺพสฺส ธมฺมสฺส อารมฺมณ\-ปจฺจเยน ปจฺจโย. อุปนิสฺสย\-ปจฺจเยน ปจฺจโย. ปุเรชาต\-ปจฺจเยน ปจฺจโย. (2)}

\prose{ทสฺสเนน ปหาตพฺโพ จ นทสฺสเนน ปหาตพฺโพ จ ธมฺมา ทสฺสเนน ปหาตพฺพสฺส ธมฺมสฺส สหชาตํ, ปุเรชาตํ. (1)}

\prose{ทสฺสเนน ปหาตพฺโพ จ นทสฺสเนน ปหาตพฺโพ จ ธมฺมา นทสฺสเนน ปหาตพฺพสฺส ธมฺมสฺส สหชาตํ, ปจฺฉาชาตํ, อาหารํ, อินฺทฺริยํ. (2)}

\csromanmarkh{2. ปจฺจย\-ปจฺจนีย}\csromanmarkhii{2. สงฺขฺยาวาร}\csromanheadernomark{2}{2. ปจฺจย\-ปจฺจนีย 2. สงฺขฺยาวาร}
\proseitem{32}{นเหตุยา สตฺต, นอารมฺมเณ สตฺต, นอธิปติยา สตฺต, นอนนฺตเร สตฺต, นสมนนฺตเร สตฺต, นสหชาเต ปญฺจ, นอญฺญมญฺเญ ปญฺจ, นนิสฺสเย ปญฺจ, นอุปนิสฺสเย สตฺต, นปุเรชาเต ฉ, นปจฺฉาชาเต สตฺต ฯเปฯ นมคฺเค สตฺต, นสมฺปยุตฺเต ปญฺจ, นวิปฺปยุตฺเต จตฺตาริ, โนอตฺถิยา จตฺตาริ, โนนตฺถิยา สตฺต, โนวิคเต สตฺต, โนอวิคเต จตฺตาริ.}

\csromanmatikaanchor{mtk.13}
\csromantocmarkref{section}{84. ภาวนายปหาตพฺพทุก}{mtk.13}
\bookmark[dest={mtk.13},level=1]{84. ภาวนายปหาตพฺพทุก}
\csromanmarkh{84. ภาวนาย\-ปหาตพฺพทุก}\csromanmarkhii{3. ปจฺจยา\-นุโลม\-ปจฺจนีย}\csromanheadernomark{1}{84. ภาวนาย\-ปหาตพฺพทุก 3. ปจฺจยา\-นุโลม\-ปจฺจนีย}
\proseitem{33}{\csromanfolio{51}เหตุปจฺจยา นอารมฺมเณ จตฺตาริ, นอธิปติยา นว, นอนนฺตเร จตฺตาริ, นสมนนฺตเร จตฺตาริ, นอญฺญมญฺเญ เทฺว, นอุปนิสฺสเย จตฺตาริ, (สพฺพตฺถ จตฺตาริ.) นสมฺปยุตฺเต เทฺว, นวิปฺปยุตฺเต เทฺว, โนนตฺถิยา จตฺตาริ, โนวิคเต จตฺตาริ.}

\csromanheader{2}{4. ปจฺจย\-ปจฺจนียา\-นุโลม}
\proseitemwithcloser{34}{นเหตุปจฺจยา อารมฺมเณ จตฺตาริ, อธิปติยา ปญฺจ, (อนุโลมมาติกา กาตพฺพา.) ฯเปฯ อวิคเต สตฺต.}{ทสฺสเนน ปหาตพฺพทุกํ นิฏฺฐิตํ.}
\csromanmarkh{84. ภาวนาย\-ปหาตพฺพทุก}\csromanmarkhii{1. ปฏิจฺจวาร}\csromanheadernomark{2}{84. ภาวนาย\-ปหาตพฺพทุก 1. ปฏิจฺจวาร}
\proseitem{35}{ภาวนาย ปหาตพฺพํ ธมฺมํ ปฏิจฺจ ภาวนาย ปหาตพฺโพ ธมฺโม อุปฺปชฺชติ เหตุปจฺจยา.}

\vspace{\tipitakaparskip}
\csromancenter{(ยถา ทสฺสนทุกํ, เอวํ วิตฺถาเรตพฺพํ, นินฺนานากรณํ.)}
\prose{เหตุยา ปญฺจ ฯเปฯ อวิคเต ปญฺจ. อนุโลมํ.}

\prose{ภาวนาย ปหาตพฺพํ ธมฺมํ ปฏิจฺจ ภาวนาย ปหาตพฺโพ ธมฺโม อุปฺปชฺชติ นเหตุปจฺจยา. อุทฺธจฺจ\-สหคเต ขนฺเธ ปฏิจฺจ อุทฺธจฺจ\-สหคโต โมโห.}

\prose{นภาวนาย ปหาตพฺพํ ธมฺมํ ปฏิจฺจ นภาวนาย ปหาตพฺโพ ธมฺโม อุปฺปชฺชติ นเหตุปจฺจยา. อเหตุกํ นภาวนาย ปหาตพฺพํ เอกํ ขนฺธํ ฯเปฯ. (ยาว อสญฺญสตฺตา.) วิจิกิจฺฉา\-สหคเต ขนฺเธ ปฏิจฺจ วิจิกิจฺฉา\-สหคโต โมโห. (สํขิตฺตํ.)}

\prose{นเหตุยา เทฺว ฯเปฯ โนวิคเต ตีณิ.}

\vspace{\tipitakaparskip}
\csromancenter{ปจฺจนียํ.}
\prose{\csromanfolio{52}(ปจฺจยวาร\-ปจฺจนีเย นเหตุปจฺจเย อุทฺธจฺจ\-สหคเต ตีณิ, โมโห อุทฺธริตพฺโพ, สพฺเพปิ วารา ทสฺสน\-ทุก\-สทิสา, อุทฺธจฺจปจฺจนียมฺปิ นานํ.)}

\csromanmarkh{84. ภาวนาย\-ปหาตพฺพทุก}\csromanmarkhii{7. ปญฺหาวาร}\csromanheadernomark{2}{84. ภาวนาย\-ปหาตพฺพทุก 7. ปญฺหาวาร}
\proseitem{36}{ภาวนาย ปหาตพฺโพ ธมฺโม ภาวนาย ปหาตพฺพสฺส ธมฺมสฺส เหตุปจฺจเยน ปจฺจโย. ภาวนาย ปหาตพฺพา เหตู สมฺปยุตฺตกานํ ขนฺธานํ เหตุปจฺจเยน ปจฺจโย. ตีณิ.}

\prose{นภาวนาย ปหาตพฺโพ ธมฺโม นภาวนาย ปหาตพฺพสฺส ธมฺมสฺส เหตุปจฺจเยน ปจฺจโย. (สํขิตฺตํ.) (1)}

\proseitem{37}{ภาวนาย ปหาตพฺโพ ธมฺโม ภาวนาย ปหาตพฺพสฺส ธมฺมสฺส อารมฺมณ\-ปจฺจเยน ปจฺจโย. ภาวนาย ปหาตพฺพํ ราคํ อสฺสาเทติ อภินนฺทติ, ตํ อารพฺภ ภาวนาย ปหาตพฺโพ ราโค อุปฺปชฺชติ. อุทฺธจฺจํ อุปฺปชฺชติ. ภาวนาย ปหาตพฺพํ โทมนสฺสํ อุปฺปชฺชติ. อุทฺธจฺจํ อารพฺภ อุทฺธจฺจํ อุปฺปชฺชติ. ภาวนาย ปหาตพฺพํ โทมนสฺสํ อุปฺปชฺชติ. ภาวนาย ปหาตพฺพํ โทมนสฺสํ อารพฺภ ภาวนาย ปหาตพฺพํ โทมนสฺสํ อุปฺปชฺชติ. อุทฺธจฺจํ อุปฺปชฺชติ. (1)}

\prose{ภาวนาย ปหาตพฺโพ ธมฺโม นภาวนาย ปหาตพฺพสฺส ธมฺมสฺส อารมฺมณ\-ปจฺจเยน ปจฺจโย. อริยา ภาวนาย ปหาตพฺเพ ปหีเน กิเลเส ฯเปฯ. วิกฺขมฺภิเต กิเลเส ฯเปฯ. ปุพฺเพ สมุทาจิณฺเณ ฯเปฯ. ภาวนาย ปหาตพฺเพ ขนฺเธ อนิจฺจโต ฯเปฯ วิปสฺสติ, อสฺสาเทติ อภินนฺทติ, ตํ อารพฺภ นภาวนาย ปหาตพฺโพ ราโค อุปฺปชฺชติ. ทิฏฺฐิ ฯเปฯ วิจิกิจฺฉา ฯเปฯ. นภาวนาย ปหาตพฺพํ โทมนสฺสํ อุปฺปชฺชติ. เจโตปริย\-ญาเณน ภาวนาย ปหาตพฺพ\-จิตฺต\-สมงฺคิสฺส จิตฺตํ ชานาติ. ภาวนาย ปหาตพฺพา ขนฺธา เจโตปริย\-ญาณสฺส, ปุพฺเพ นิวาสา\-นุสฺสติ\-ญาณสฺส, \csromanfolio{53}ยถากมฺมูปคญาณสฺส, อนาคตํสญาณสฺส, อาวชฺชนาย อารมฺมณ\-ปจฺจเยน ปจฺจโย. (2)}

\proseitem{38}{นภาวนาย ปหาตพฺโพ ธมฺโม นภาวนาย ปหาตพฺพสฺส ธมฺมสฺส อารมฺมณ\-ปจฺจเยน ปจฺจโย. ทานํ ฯเปฯ สีลํ ฯเปฯ อุโปสถกมฺมํ กตฺวา ตํ ปจฺจเวกฺขติ, อสฺสาเทติ อภินนฺทติ, ตํ อารพฺภ นภาวนาย ปหาตพฺโพ ราโค อุปฺปชฺชติ. ทิฏฺฐิ ฯเปฯ วิจิกิจฺฉา ฯเปฯ. นภาวนาย ปหาตพฺพํ โทมนสฺสํ อุปฺปชฺชติ. ปุพฺเพ ฯเปฯ. ฌานา ฯเปฯ. อริยา มคฺคา วุฏฺฐหิตฺวา ฯเปฯ ผลสฺส, อาวชฺชนาย อารมฺมณ\-ปจฺจเยน ปจฺจโย. อริยา นภาวนาย ปหาตพฺเพ ปหีเน กิเลเส ฯเปฯ. จกฺขุํ ฯเปฯ. วตฺถุํ นภาวนาย ปหาตพฺเพ ขนฺเธ อนิจฺจโต ฯเปฯ วิปสฺสติ, อสฺสาเทติ อภินนฺทติ, ตํ อารพฺภ นภาวนาย ปหาตพฺโพ ราโค อุปฺปชฺชติ. ทิฏฺฐิ ฯเปฯ วิจิกิจฺฉา ฯเปฯ. นภาวนาย ปหาตพฺพํ โทมนสฺสํ อุปฺปชฺชติ. ทิพฺเพน จกฺขุนา รูปํ ปสฺสติ ฯเปฯ ยถากมฺมูปคญาณสฺส, อนาคตํสญาณสฺส, อาวชฺชนาย อารมฺมณ\-ปจฺจเยน ปจฺจโย. (1)}

\prose{นภาวนาย ปหาตพฺโพ ธมฺโม ภาวนาย ปหาตพฺพสฺส ธมฺมสฺส อารมฺมณ\-ปจฺจเยน ปจฺจโย. ทานํ ฯเปฯ สีลํ ฯเปฯ ฌานํ ฯเปฯ. จกฺขุํ ฯเปฯ. วตฺถุํ นภาวนาย ปหาตพฺเพ ขนฺเธ อสฺสาเทติ อภินนฺทติ, ตํ อารพฺภ ภาวนาย ปหาตพฺโพ ราโค อุปฺปชฺชติ. อุทฺธจฺจํ อุปฺปชฺชติ. ภาวนาย ปหาตพฺพํ โทมนสฺสํ อุปฺปชฺชติ. (2)}

\csromanheader{2}{\textbf{อธิปติ}}
\proseitem{39}{ภาวนาย ปหาตพฺโพ ธมฺโม ภาวนาย ปหาตพฺพสฺส ธมฺมสฺส อธิปติ\-ปจฺจเยน ปจฺจโย. อารมฺมณา\-ธิปติ, สหชาตา\-ธิปติ.  อารมฺมณา\-ธิปติ\csromandash{}ภาวนาย ปหาตพฺพํ ราคํ ครุํ กตฺวา อสฺสาเทติ อภินนฺทติ, ตํ ครุํ กตฺวา ภาวนาย ปหาตพฺโพ ราโค อุปฺปชฺชติ. สหชาตา\-ธิปติ\csromandash{}ภาวนาย ปหาตพฺพา\-ธิปติ สมฺปยุตฺตกานํ ขนฺธานํ อธิปติ\-ปจฺจเยน ปจฺจโย. (1)}

\prose{ภาวนาย ปหาตพฺโพ ธมฺโม นภาวนาย ปหาตพฺพสฺส ธมฺมสฺส อธิปติ\-ปจฺจเยน ปจฺจโย. อารมฺมณา\-ธิปติ, สหชาตา\-ธิปติ.  อารมฺมณา\-ธิปติ\csromandash{}ภาวนาย ปหาตพฺพํ ราคํ ครุํ กตฺวา อสฺสาเทติ อภินนฺทติ, ตํ ครุํ กตฺวา นภาวนาย ปหาตพฺโพ ราโค อุปฺปชฺชติ. \csromanfolio{54}ทิฏฺฐิ อุปฺปชฺชติ. สหชาตา\-ธิปติ\csromandash{}ภาวนาย ปหาตพฺพา\-ธิปติ จิตฺต\-สมุฏฺฐานานํ รูปานํ อธิปติ\-ปจฺจเยน ปจฺจโย. (2)}

\prose{ภาวนาย ปหาตพฺโพ ธมฺโม ภาวนาย ปหาตพฺพสฺส จ นภาวนาย ปหาตพฺพสฺส จ ธมฺมสฺส อธิปติ\-ปจฺจเยน ปจฺจโย. สหชาตา\-ธิปติ\csromandash{}ภาวนาย ปหาตพฺพา\-ธิปติ สมฺปยุตฺตกานํ ขนฺธานํ จิตฺตสมุฏฺฐานานญฺจ รูปานํ อธิปติ\-ปจฺจเยน ปจฺจโย. (3)}

\proseitem{40}{นภาวนาย ปหาตพฺโพ ธมฺโม นภาวนาย ปหาตพฺพสฺส ธมฺมสฺส อธิปติ\-ปจฺจเยน ปจฺจโย. อารมฺมณา\-ธิปติ, สหชาตา\-ธิปติ.  อารมฺมณา\-ธิปติ\csromandash{}ทานํ ทตฺวา, สีลํ ฯเปฯ อุโปสถกมฺมํ กตฺวา ตํ ครุํ กตฺวา ปจฺจเวกฺขติ, อสฺสาเทติ อภินนฺทติ, ตํ ครุํ กตฺวา นภาวนาย ปหาตพฺโพ ราโค อุปฺปชฺชติ. สหชาตา\-ธิปติ\csromandash{}นภาวนาย ปหาตพฺพา\-ธิปติ สมฺปยุตฺตกานํ ขนฺธานํ จิตฺตสมุฏฺฐานานญฺจ รูปานํ อธิปติ\-ปจฺจเยน ปจฺจโย. (1)}

\prose{นภาวนาย ปหาตพฺโพ ธมฺโม ภาวนาย ปหาตพฺพสฺส ธมฺมสฺส อธิปติ\-ปจฺจเยน ปจฺจโย. อารมฺมณา\-ธิปติ\csromandash{}ทานํ ฯเปฯ ฌานํ ฯเปฯ. จกฺขุํ ฯเปฯ. วตฺถุํ นภาวนาย ปหาตพฺเพ ขนฺเธ ครุํ กตฺวา อสฺสาเทติ อภินนฺทติ, ตํ ครุํ กตฺวา ภาวนาย ปหาตพฺโพ ราโค อุปฺปชฺชติ. (2)}

\proseitem{41}{ภาวนาย ปหาตพฺโพ ธมฺโม ภาวนาย ปหาตพฺพสฺส ธมฺมสฺส อนนฺตร\-ปจฺจเยน ปจฺจโย. จตฺตาริ. (ทสฺสน\-ทุก\-สทิสา ภาวนา นินฺนานากรณา.) สมนนฺตร\-ปจฺจเยน ปจฺจโย. จตฺตาริ. สหชาต\-ปจฺจเยน ปจฺจโย. ปญฺจ. อญฺญมญฺญ\-ปจฺจเยน ปจฺจโย. เทฺว. นิสฺสย\-ปจฺจเยน ปจฺจโย. สตฺต.}

\csromanheader{2}{\textbf{อุปนิสฺสย}}
\proseitem{42}{ภาวนาย ปหาตพฺโพ ธมฺโม ภาวนาย ปหาตพฺพสฺส ธมฺมสฺส อุปนิสฺสย\-ปจฺจเยน ปจฺจโย. อารมฺมณู\-ปนิสฺสโย, อนนฺตรู\-ปนิสฺสโย, ปกตู\-ปนิสฺสโย ฯเปฯ. ปกตู\-ปนิสฺสโย\csromandash{}ภาวนาย ปหาตพฺโพ ราโค. โทโส. โมโห. มาโน. ปตฺถนา ภาวนาย ปหาตพฺพสฺส \csromanfolio{55}ราคสฺส. โทสสฺส. โมหสฺส. มานสฺส. ปตฺถนาย อุปนิสฺสย\-ปจฺจเยน ปจฺจโย. (1)}

\prose{ภาวนาย ปหาตพฺโพ ธมฺโม นภาวนาย ปหาตพฺพสฺส ธมฺมสฺส อุปนิสฺสย\-ปจฺจเยน ปจฺจโย. อารมฺมณู\-ปนิสฺสโย, อนนฺตรู\-ปนิสฺสโย, ปกตู\-ปนิสฺสโย ฯเปฯ. ปกตู\-ปนิสฺสโย\csromandash{}ภาวนาย ปหาตพฺพํ ราคํ อุปนิสฺสาย ทานํ เทติ ฯเปฯ สมาปตฺติํ อุปฺปาเทติ. ปาณํ หนติ ฯเปฯ สํฆํ ภินฺทติ. ภาวนาย ปหาตพฺพํ โทสํ. โมหํ. มานํ. ปตฺถนํ อุปนิสฺสาย ทานํ เทติ ฯเปฯ สมาปตฺติํ อุปฺปาเทติ. ปาณํ หนติ ฯเปฯ สํฆํ ภินฺทติ. ภาวนาย ปหาตพฺโพ ราโค ฯเปฯ ปตฺถนา สทฺธาย ฯเปฯ ปญฺญาย. นภาวนาย ปหาตพฺพสฺส ราคสฺส. โทสสฺส. โมหสฺส. ทิฏฺฐิยา. ปตฺถนาย. กายิกสฺส สุขสฺส. กายิกสฺส ทุกฺขสฺส. มคฺคสฺส, ผลสมาปตฺติยา อุปนิสฺสย\-ปจฺจเยน ปจฺจโย. (2)}

\proseitem{43}{นภาวนาย ปหาตพฺโพ ธมฺโม นภาวนาย ปหาตพฺพสฺส ธมฺมสฺส อุปนิสฺสย\-ปจฺจเยน ปจฺจโย. อารมฺมณู\-ปนิสฺสโย, อนนฺตรู\-ปนิสฺสโย, ปกตู\-ปนิสฺสโย ฯเปฯ. ปกตู\-ปนิสฺสโย\csromandash{}สทฺธํ อุปนิสฺสาย ทานํ เทติ ฯเปฯ สมาปตฺติํ อุปฺปาเทติ. ทิฏฺฐิํ คณฺหาติ. สีลํ ฯเปฯ. ปญฺญํ. นภาวนาย ปหาตพฺพํ ราคํ. โทสํ. โมหํ. ทิฏฺฐิํ. ปตฺถนํ. กายิกํ สุขํ. กายิกํ ทุกฺขํ ฯเปฯ. เสนาสนํ อุปนิสฺสาย ทานํ เทติ ฯเปฯ สมาปตฺติํ อุปฺปาเทติ. ปาณํ หนติ ฯเปฯ สํฆํ ภินฺทติ. สทฺธา ฯเปฯ. เสนาสนํ สทฺธาย ฯเปฯ ปญฺญาย. นภาวนาย ปหาตพฺพสฺส ราคสฺส. โทสสฺส. โมหสฺส. ทิฏฺฐิยา. ปตฺถนาย. กายิกสฺส สุขสฺส. กายิกสฺส ทุกฺขสฺส. มคฺคสฺส, ผลสมาปตฺติยา อุปนิสฺสย\-ปจฺจเยน ปจฺจโย. (1)}

\prose{นภาวนาย ปหาตพฺโพ ธมฺโม ภาวนาย ปหาตพฺพสฺส ธมฺมสฺส อุปนิสฺสย\-ปจฺจเยน ปจฺจโย. อารมฺมณู\-ปนิสฺสโย, อนนฺตรู\-ปนิสฺสโย, ปกตู\-ปนิสฺสโย ฯเปฯ. ปกตู\-ปนิสฺสโย\csromandash{}สทฺธํ อุปนิสฺสาย มานํ ชปฺเปติ ฯเปฯ. สีลํ ฯเปฯ. ปญฺญํ. ราคํ ฯเปฯ. กายิกํ สุขํ. กายิกํ ทุกฺขํ. เสนาสนํ อุปนิสฺสาย มานํ ชปฺเปติ. สทฺธา ฯเปฯ. เสนาสนํ ภาวนาย ปหาตพฺพสฺส ราคสฺส. โทสสฺส. โมหสฺส. มานสฺส. ปตฺถนาย อุปนิสฺสย\-ปจฺจเยน ปจฺจโย. (2)}

\csromanheader{2}{\textbf{ปุเรชาตาทิ}}
\proseitem{44}{\csromanfolio{56}นภาวนาย ปหาตพฺโพ ธมฺโม นภาวนาย ปหาตพฺพสฺส ธมฺมสฺส ปุเรชาต\-ปจฺจเยน ปจฺจโย. อารมฺมณ\-ปุเรชาตํ, วตฺถุ\-ปุเรชาตํ.  อารมฺมณ\-ปุเรชาตํ\csromandash{}จกฺขุํ ฯเปฯ. วตฺถุํ อนิจฺจโต ฯเปฯ วิปสฺสติ, อสฺสาเทติ อภินนฺทติ, ตํ อารพฺภ นภาวนาย ปหาตพฺโพ ราโค อุปฺปชฺชติ. ทิฏฺฐิ อุปฺปชฺชติ. วิจิกิจฺฉา อุปฺปชฺชติ. นภาวนาย ปหาตพฺพํ โทมนสฺสํ อุปฺปชฺชติ. ทิพฺเพน จกฺขุนา รูปํ ปสฺสติ ฯเปฯ. โผฏฺฐพฺพา\-ยตนํ กายวิญฺญาณสฺส ฯเปฯ. วตฺถุ\-ปุเรชาตํ\csromandash{}จกฺขายตนํ จกฺขุ\-วิญฺญาณสฺส ฯเปฯ กายายตนํ กายวิญฺญาณสฺส ฯเปฯ. วตฺถุ นภาวนาย ปหาตพฺพานํ ขนฺธานํ ปุเรชาต\-ปจฺจเยน ปจฺจโย. (1)}

\prose{นภาวนาย ปหาตพฺโพ ธมฺโม ภาวนาย ปหาตพฺพสฺส ธมฺมสฺส ปุเรชาต\-ปจฺจเยน ปจฺจโย. อารมฺมณ\-ปุเรชาตํ, วตฺถุ\-ปุเรชาตํ.  อารมฺมณ\-ปุเรชาตํ\csromandash{}จกฺขุํ ฯเปฯ. วตฺถุํ อสฺสาเทติ อภินนฺทติ, ตํ อารพฺภ ภาวนาย ปหาตพฺโพ ราโค อุปฺปชฺชติ. อุทฺธจฺจํ อุปฺปชฺชติ. ภาวนาย ปหาตพฺพํ โทมนสฺสํ อุปฺปชฺชติ. วตฺถุ\-ปุเรชาตํ\csromandash{}วตฺถุ ภาวนาย ปหาตพฺพานํ ขนฺธานํ ปุเรชาต\-ปจฺจเยน ปจฺจโย. (2)}

\proseitem{45}{ปจฺฉาชาต\-ปจฺจเยน ปจฺจโย. เทฺว. อาเสวน\-ปจฺจเยน ปจฺจโย. เทฺว. กมฺมปจฺจเยน ปจฺจโย. ภาวนาย ปหาตพฺพา เจตนา สมฺปยุตฺตกานํ ขนฺธานํ กมฺมปจฺจเยน ปจฺจโย. (มูลํ กาตพฺพํ.) ภาวนาย ปหาตพฺพา เจตนา จิตฺต\-สมุฏฺฐานานํ รูปานํ กมฺมปจฺจเยน ปจฺจโย. (มูลํ กาตพฺพํ.) ภาวนาย ปหาตพฺพา เจตนา สมฺปยุตฺตกานํ ขนฺธานํ จิตฺตสมุฏฺฐานานญฺจ รูปานํ กมฺมปจฺจเยน ปจฺจโย. (3)}

\prose{นภาวนาย ปหาตพฺโพ ธมฺโม นภาวนาย ปหาตพฺพสฺส ธมฺมสฺส กมฺมปจฺจเยน ปจฺจโย. สหชาตา, นานากฺขณิกา.  สหชาตา นภาวนาย ปหาตพฺพา เจตนา สมฺปยุตฺตกานํ ขนฺธานํ จิตฺตสมุฏฺฐานานญฺจ รูปานํ กมฺมปจฺจเยน ปจฺจโย. ปฏิสนฺธิ\-กฺขเณ ฯเปฯ. นานากฺขณิกา นภาวนาย ปหาตพฺพา เจตนา วิปากานํ ขนฺธานํ กฏตฺตา จ รูปานํ กมฺมปจฺจเยน ปจฺจโย. วิปาก\-ปจฺจเยน ปจฺจโย. เอกํ ฯเปฯ. อวิคต\-ปจฺจเยน ปจฺจโย. (สพฺพปจฺจยา ทสฺสน\-ทุก\-สทิสา. ภาวนา นินฺนานากรณา.)}

\csromanmatikaanchor{mtk.14}
\csromantocmarkref{section}{85. ทสฺสเนนปหาตพฺพเหตุกทุก}{mtk.14}
\bookmark[dest={mtk.14},level=1]{85. ทสฺสเนนปหาตพฺพเหตุกทุก}
\csromanmarkh{85. ทสฺสเนน\-ปหาตพฺพ\-เหตุกทุก}\csromanmarkhii{1. ปจฺจยานุโลม 2. สงฺขฺยาวาร}\csromanheadernomark{1}{85. ทสฺสเนน\-ปหาตพฺพ\-เหตุกทุก 1. ปจฺจยานุโลม 2. สงฺขฺยาวาร}
\proseitem{46}{\csromanfolio{57}เหตุยา จตฺตาริ, อารมฺมเณ จตฺตาริ, อธิปติยา ปญฺจ, อนนฺตเร จตฺตาริ, สมนนฺตเร จตฺตาริ, สหชาเต ปญฺจ, อญฺญมญฺเญ เทฺว, นิสฺสเย สตฺต, อุปนิสฺสเย จตฺตาริ, ปุเรชาเต เทฺว, ปจฺฉาชาเต เทฺว, อาเสวเน เทฺว, กมฺเม จตฺตาริ, วิปาเก เอกํ, อาหาเร จตฺตาริ, อินฺทฺริเย จตฺตาริ, ฌาเน จตฺตาริ, มคฺเค จตฺตาริ, สมฺปยุตฺเต เทฺว, วิปฺปยุตฺเต ตีณิ, อตฺถิยา สตฺต, นตฺถิยา จตฺตาริ, วิคเต จตฺตาริ, อวิคเต สตฺต.}

\prosewithcloser{(ปจฺจนีย\-วิภงฺโค ทสฺสน\-ทุก\-สทิโส วิภชิตพฺโพ. เอวํ ตีณิ คณนาปิ คเณตพฺพา.)}{ภาวนาย ปหาตพฺพทุกํ นิฏฺฐิตํ.}
\csromanmarkh{85. ทสฺสเนน\-ปหาตพฺพ\-เหตุกทุก}\csromanmarkhii{1. ปฏิจฺจวาร}\csromanheadernomark{2}{85. ทสฺสเนน\-ปหาตพฺพ\-เหตุกทุก 1. ปฏิจฺจวาร}
\proseitem{47}{ทสฺสเนน ปหาตพฺพ\-เหตุกํ ธมฺมํ ปฏิจฺจ ทสฺสเนน ปหาตพฺพ\-เหตุโก ธมฺโม อุปฺปชฺชติ เหตุปจฺจยา. ทสฺสเนน ปหาตพฺพ\-เหตุกํ เอกํ ขนฺธํ ปฏิจฺจ ตโย ขนฺธา ฯเปฯ เทฺว ขนฺเธ ฯเปฯ. (1)}

\prose{ทสฺสเนน ปหาตพฺพ\-เหตุกํ ธมฺมํ ปฏิจฺจ นทสฺสเนน ปหาตพฺพ\-เหตุโก ธมฺโม อุปฺปชฺชติ เหตุปจฺจยา. ทสฺสเนน ปหาตพฺพ\-เหตุเก ขนฺเธ ปฏิจฺจ จิตฺต\-สมุฏฺฐานํ รูปํ. (2)}

\prose{ทสฺสเนน ปหาตพฺพ\-เหตุกํ ธมฺมํ ปฏิจฺจ ทสฺสเนน ปหาตพฺพ\-เหตุโก จ นทสฺสเนน ปหาตพฺพ\-เหตุโก จ ธมฺมา อุปฺปชฺชนฺติ เหตุปจฺจยา. ทสฺสเนน ปหาตพฺพ\-เหตุกํ เอกํ ขนฺธํ ปฏิจฺจ ตโย ขนฺธา จิตฺต\-สมุฏฺฐานญฺจ รูปํ ฯเปฯ เทฺว ขนฺเธ ฯเปฯ. (3)}

\proseitem{48}{\csromanfolio{58}นทสฺสเนน ปหาตพฺพ\-เหตุกํ ธมฺมํ ปฏิจฺจ นทสฺสเนน ปหาตพฺพ\-เหตุโก ธมฺโม อุปฺปชฺชติ เหตุปจฺจยา. นทสฺสเนน ปหาตพฺพ\-เหตุกํ เอกํ ขนฺธํ ปฏิจฺจ ตโย ขนฺธา จิตฺต\-สมุฏฺฐานญฺจ รูปํ ฯเปฯ เทฺว ขนฺเธ ฯเปฯ. วิจิกิจฺฉา\-สหคตํ โมหํ ปฏิจฺจ จิตฺต\-สมุฏฺฐานํ รูปํ. ปฏิสนฺธิ\-กฺขเณ ฯเปฯ. (ยาว อชฺฌตฺติกา มหาภูตา.) (1)}

\prose{นทสฺสเนน ปหาตพฺพ\-เหตุกํ ธมฺมํ ปฏิจฺจ ทสฺสเนน ปหาตพฺพ\-เหตุโก ธมฺโม อุปฺปชฺชติ เหตุปจฺจยา. วิจิกิจฺฉา\-สหคตํ โมหํ ปฏิจฺจ สมฺปยุตฺตกา ขนฺธา. (2)}

\prose{นทสฺสเนน ปหาตพฺพ\-เหตุกํ ธมฺมํ ปฏิจฺจ ทสฺสเนน ปหาตพฺพ\-เหตุโก จ นทสฺสเนน ปหาตพฺพ\-เหตุโก จ ธมฺมา อุปฺปชฺชนฺติ เหตุปจฺจยา. วิจิกิจฺฉา\-สหคตํ โมหํ ปฏิจฺจ สมฺปยุตฺตกา ขนฺธา จิตฺต\-สมุฏฺฐานญฺจ รูปํ. (3)}

\proseitem{49}{ทสฺสเนน ปหาตพฺพ\-เหตุ\-กญฺจ นทสฺสเนน ปหาตพฺพ\-เหตุ\-กญฺจ ธมฺมํ ปฏิจฺจ ทสฺสเนน ปหาตพฺพ\-เหตุโก ธมฺโม อุปฺปชฺชติ เหตุปจฺจยา. วิจิกิจฺฉา\-สหคตํ เอกํ ขนฺธญฺจ โมหญฺจ ปฏิจฺจ ตโย ขนฺธา ฯเปฯ เทฺว ขนฺเธ จ ฯเปฯ. (1)}

\prose{ทสฺสเนน ปหาตพฺพ\-เหตุ\-กญฺจ นทสฺสเนน ปหาตพฺพ\-เหตุ\-กญฺจ ธมฺมํ ปฏิจฺจ นทสฺสเนน ปหาตพฺพ\-เหตุโก ธมฺโม อุปฺปชฺชติ เหตุปจฺจยา. ทสฺสเนน ปหาตพฺพ\-เหตุเก ขนฺเธ จ มหาภูเต จ ปฏิจฺจ จิตฺต\-สมุฏฺฐานํ รูปํ. วิจิกิจฺฉา\-สหคเต ขนฺเธ จ โมหญฺจ ปฏิจฺจ จิตฺต\-สมุฏฺฐานํ รูปํ. (2)}

\prose{ทสฺสเนน ปหาตพฺพ\-เหตุ\-กญฺจ นทสฺสเนน ปหาตพฺพ\-เหตุ\-กญฺจ ธมฺมํ ปฏิจฺจ ทสฺสเนน ปหาตพฺพ\-เหตุโก จ นทสฺสเนน ปหาตพฺพ\-เหตุโก จ ธมฺมา อุปฺปชฺชนฺติ เหตุปจฺจยา. วิจิกิจฺฉา\-สหคตํ เอกํ ขนฺธญฺจ โมหญฺจ ปฏิจฺจ ตโย ขนฺธา จิตฺต\-สมุฏฺฐานญฺจ รูปํ ฯเปฯ เทฺว ขนฺเธ จ ฯเปฯ. (3)}

\proseitem{50}{ทสฺสเนน ปหาตพฺพ\-เหตุกํ ธมฺมํ ปฏิจฺจ ทสฺสเนน ปหาตพฺพ\-เหตุโก ธมฺโม อุปฺปชฺชติ อารมฺมณ\-ปจฺจยา. ทสฺสเนน ปหาตพฺพ\-เหตุกํ เอกํ ขนฺธํ ปฏิจฺจ ตโย ขนฺธา ฯเปฯ เทฺว ขนฺเธ ฯเปฯ. (1)}

\prose{\csromanfolio{59}ทสฺสเนน ปหาตพฺพ\-เหตุกํ ธมฺมํ ปฏิจฺจ นทสฺสเนน ปหาตพฺพ\-เหตุโก ธมฺโม อุปฺปชฺชติ อารมฺมณ\-ปจฺจยา. วิจิกิจฺฉา\-สหคเต ขนฺเธ ปฏิจฺจ วิจิกิจฺฉา\-สหคโต โมโห. (2)}

\prose{ทสฺสเนน ปหาตพฺพ\-เหตุกํ ธมฺมํ ปฏิจฺจ ทสฺสเนน ปหาตพฺพ\-เหตุโก จ นทสฺสเนน ปหาตพฺพ\-เหตุโก จ ธมฺมา อุปฺปชฺชนฺติ อารมฺมณ\-ปจฺจยา. วิจิกิจฺฉา\-สหคตํ เอกํ ขนฺธํ ปฏิจฺจ ตโย ขนฺธา โมโห จ ฯเปฯ เทฺว ขนฺเธ ฯเปฯ. (3)}

\proseitem{51}{นทสฺสเนน ปหาตพฺพ\-เหตุกํ ธมฺมํ ปฏิจฺจ นทสฺสเนน ปหาตพฺพ\-เหตุโก ธมฺโม อุปฺปชฺชติ อารมฺมณ\-ปจฺจยา. นทสฺสเนน ปหาตพฺพ\-เหตุกํ เอกํ ขนฺธํ ปฏิจฺจ ตโย ขนฺธา ฯเปฯ เทฺว ขนฺเธ ฯเปฯ. ปฏิสนฺธิ\-กฺขเณ ฯเปฯ. (1)}

\prose{นทสฺสเนน ปหาตพฺพ\-เหตุกํ ธมฺมํ ปฏิจฺจ ทสฺสเนน ปหาตพฺพ\-เหตุโก ธมฺโม อุปฺปชฺชติ อารมฺมณ\-ปจฺจยา. วิจิกิจฺฉา\-สหคตํ โมหํ ปฏิจฺจ สมฺปยุตฺตกา ขนฺธา. (2)}

\prose{ทสฺสเนน ปหาตพฺพ\-เหตุ\-กญฺจ นทสฺสเนน ปหาตพฺพ\-เหตุ\-กญฺจ ธมฺมํ ปฏิจฺจ ทสฺสเนน ปหาตพฺพ\-เหตุโก ธมฺโม อุปฺปชฺชติ อารมฺมณ\-ปจฺจยา. วิจิกิจฺฉา\-สหคตํ เอกํ ขนฺธญฺจ โมหญฺจ ปฏิจฺจ ตโย ขนฺธา ฯเปฯ เทฺว ขนฺเธ จ ฯเปฯ. (1) (สํขิตฺตํ.)}

\csromanmarkh{1. ปจฺจยานุโลม}\csromanmarkhii{2. สงฺขฺยาวาร}\csromanheadernomark{2}{1. \textbf{ปจฺจยานุโลม 2. สงฺขฺยาวาร}}
\proseitem{52}{เหตุยา นว, อารมฺมเณ ฉ, อธิปติยา ปญฺจ, อนนฺตเร ฉ, สมนนฺตเร ฉ, สหชาเต นว, อญฺญมญฺเญ ฉ, นิสฺสเย นว, อุปนิสฺสเย ฉ, ปุเรชาเต ฉ, อาเสวเน ฉ, กมฺเม นว, วิปาเก เอกํ, อาหาเร นว, อินฺทฺริเย นว, ฌาเน นว, มคฺเค นว, สมฺปยุตฺเต ฉ, วิปฺปยุตฺเต นว, อตฺถิยา นว, นตฺถิยา ฉ, วิคเต ฉ, อวิคเต นว.}

\csromanmarkh{2. ปจฺจย\-ปจฺจนีย}\csromanmarkhii{1. วิภงฺควาร}\csromanheadernomark{2}{2. ปจฺจย\-ปจฺจนีย 1. วิภงฺควาร}
\proseitem{53}{\csromanfolio{60}ทสฺสเนน ปหาตพฺพ\-เหตุกํ ธมฺมํ ปฏิจฺจ นทสฺสเนน ปหาตพฺพ\-เหตุโก ธมฺโม อุปฺปชฺชติ นเหตุปจฺจยา. วิจิกิจฺฉา\-สหคเต ขนฺเธ ปฏิจฺจ วิจิกิจฺฉา\-สหคโต โมโห. (1)}

\prose{นทสฺสเนน ปหาตพฺพ\-เหตุกํ ธมฺมํ ปฏิจฺจ นทสฺสเนน ปหาตพฺพ\-เหตุโก ธมฺโม อุปฺปชฺชติ นเหตุปจฺจยา. อเหตุกํ นทสฺสเนน ปหาตพฺพ\-เหตุกํ เอกํ ขนฺธํ ปฏิจฺจ ตโย ขนฺธา จิตฺต\-สมุฏฺฐานญฺจ รูปํ ฯเปฯ. (ยาว อสญฺญสตฺตา.) อุทฺธจฺจ\-สหคเต ขนฺเธ ปฏิจฺจ อุทฺธจฺจ\-สหคโต โมโห. (1)}

\csromanmarkh{2. ปจฺจย\-ปจฺจนีย}\csromanmarkhii{2. สงฺขฺยาวาร}\csromanheadernomark{2}{2. ปจฺจย\-ปจฺจนีย 2. สงฺขฺยาวาร}
\vspace{\tipitakaparskip}
\csromancenter{\textbf{นเหตุ}ยา เทฺว, นอารมฺมเณ ตีณิ, นอธิปติยา นว, นอนนฺตเร ตีณิ, นสมนนฺตเร ตีณิ, นอญฺญมญฺเญ ตีณิ, นฺเอาปนิสฺสเย ตีณิ, นปุเรชาเต สตฺต, นปจฺฉาชาเต นว, นอาเสวเน นว, นกมฺเม จตฺตาริ, นวิปาเก นว, นอาหาเร เอกํ, ไนนฺทฺริเย เอกํ, นฌาเน เอกํ, นมคฺเค เอกํ, นสมฺปยุตฺเต ตีณิ, นวิปฺปยุตฺเต ฉ, โนนตฺถิยา ตีณิ, โนวิคเต ตีณิ.}
\proseitem{55}{เหตุปจฺจยา นอารมฺมเณ ตีณิ, นอธิปติยา นว ฯเปฯ นปุเรชาเต สตฺต ฯเปฯ นวิปฺปยุตฺเต จตฺตาริ, โนนตฺถิยา ตีณิ, โนวิคเต ตีณิ.}

\csromanheader{2}{4. ปจฺจย\-ปจฺจนียา\-นุโลม}
\proseitem{56}{\textbf{นเหตุ}ปจฺจยา อารมฺมเณ เทฺว, อนนฺตเร เทฺว ฯเปฯ วิปาเก เอกํ ฯเปฯ มคฺเค เทฺว ฯเปฯ อวิคเต เทฺว.}

\vspace{\tipitakaparskip}
\csromancenter{(สหชาตวาโร ปฏิจฺจ\-วาร\-สทิโส.)}
\csromancenter{85. ทสฺสเนน\-ปหาตพฺพ\-เหตุกทุก 3. ปจฺจยวาร}
\proseitem{57}{\csromanfolio{61}ทสฺสเนน ปหาตพฺพ\-เหตุกํ ธมฺมํ ปจฺจยา ทสฺสเนน ปหาตพฺพ\-เหตุโก ธมฺโม อุปฺปชฺชติ เหตุปจฺจยา. ตีณิ.}

\prose{นทสฺสเนน ปหาตพฺพ\-เหตุกํ ธมฺมํ ปจฺจยา นทสฺสเนน ปหาตพฺพ\-เหตุโก ธมฺโม อุปฺปชฺชติ เหตุปจฺจยา. นทสฺสเนน ปหาตพฺพ\-เหตุกํ เอกํ ขนฺธํ ปจฺจยา ตโย ขนฺธา จิตฺต\-สมุฏฺฐานญฺจ รูปํ ฯเปฯ เทฺว ขนฺเธ ฯเปฯ. ปฏิสนฺธิ\-กฺขเณ ฯเปฯ. (ยาว อชฺฌตฺติกา มหาภูตา.) วตฺถุํ ปจฺจยา นทสฺสเนน ปหาตพฺพ\-เหตุกา ขนฺธา. (1)}

\prose{นทสฺสเนน ปหาตพฺพ\-เหตุกํ ธมฺมํ ปจฺจยา ทสฺสเนน ปหาตพฺพ\-เหตุโก ธมฺโม อุปฺปชฺชติ เหตุปจฺจยา. วตฺถุํ ปจฺจยา ทสฺสเนน ปหาตพฺพ\-เหตุกา ขนฺธา. วิจิกิจฺฉา\-สหคตํ โมหํ ปจฺจยา สมฺปยุตฺตกา ขนฺธา. (2)}

\prose{นทสฺสเนน ปหาตพฺพ\-เหตุกํ ธมฺมํ ปจฺจยา ทสฺสเนน ปหาตพฺพ\-เหตุโก จ นทสฺสเนน ปหาตพฺพ\-เหตุโก จ ธมฺมา อุปฺปชฺชนฺติ เหตุปจฺจยา. วตฺถุํ ปจฺจยา ทสฺสเนน ปหาตพฺพ\-เหตุกา ขนฺธา. มหาภูเต ปจฺจยา จิตฺต\-สมุฏฺฐานํ รูปํ. วิจิกิจฺฉา\-สหคตํ โมหํ ปจฺจยา สมฺปยุตฺตกา ขนฺธา จิตฺต\-สมุฏฺฐานญฺจ รูปํ. (3)}

\proseitem{58}{ทสฺสเนน ปหาตพฺพ\-เหตุ\-กญฺจ นทสฺสเนน ปหาตพฺพ\-เหตุ\-กญฺจ ธมฺมํ ปจฺจยา ทสฺสเนน ปหาตพฺพ\-เหตุโก ธมฺโม อุปฺปชฺชติ เหตุปจฺจยา. ทสฺสเนน ปหาตพฺพ\-เหตุกํ เอกํ ขนฺธญฺจ วตฺถุญฺจ ปจฺจยา ตโย ขนฺธา ฯเปฯ เทฺว ขนฺเธ จ ฯเปฯ. วิจิกิจฺฉา\-สหคตํ เอกํ ขนฺธญฺจ วตฺถุญฺจ ปจฺจยา ตโย ขนฺธา ฯเปฯ เทฺว ขนฺเธ จ ฯเปฯ. (1)}

\prose{ทสฺสเนน ปหาตพฺพ\-เหตุ\-กญฺจ นทสฺสเนน ปหาตพฺพ\-เหตุ\-กญฺจ ธมฺมํ ปจฺจยา นทสฺสเนน ปหาตพฺพ\-เหตุโก ธมฺโม อุปฺปชฺชติ เหตุปจฺจยา. ทสฺสเนน ปหาตพฺพ\-เหตุเก ขนฺเธ จ มหาภูเต จ ปจฺจยา จิตฺต\-สมุฏฺฐานํ รูปํ. วิจิกิจฺฉา\-สหคเต ขนฺเธ จ โมหญฺจ ปจฺจยา จิตฺต\-สมุฏฺฐานํ รูปํ. (2)}

\prose{\csromanfolio{62}ทสฺสเนน ปหาตพฺพ\-เหตุ\-กญฺจ นทสฺสเนน ปหาตพฺพ\-เหตุ\-กญฺจ ธมฺมํ ปจฺจยา ทสฺสเนน ปหาตพฺพ\-เหตุโก จ นทสฺสเนน ปหาตพฺพ\-เหตุโก จ ธมฺมา อุปฺปชฺชนฺติ เหตุปจฺจยา. ทสฺสเนน ปหาตพฺพ\-เหตุกํ เอกํ ขนฺธญฺจ วตฺถุญฺจ ปจฺจยา ตโย ขนฺธา ฯเปฯ เทฺว ขนฺเธ จ ฯเปฯ. ทสฺสเนน ปหาตพฺพ\-เหตุเก ขนฺเธ จ มหาภูเต จ ปจฺจยา จิตฺต\-สมุฏฺฐานํ รูปํ. วิจิกิจฺฉา\-สหคตํ เอกํ ขนฺธญฺจ โมหญฺจ ปจฺจยา ตโย ขนฺธา จิตฺต\-สมุฏฺฐานญฺจ รูปํ ฯเปฯ เทฺว ขนฺเธ จ ฯเปฯ. (3)}

\proseitem{59}{ทสฺสเนน ปหาตพฺพ\-เหตุกํ ธมฺมํ ปจฺจยา ทสฺสเนน ปหาตพฺพ\-เหตุโก ธมฺโม อุปฺปชฺชติ อารมฺมณ\-ปจฺจยา. ตีณิ. (ปฏิจฺจสทิสา.)}

\prose{นทสฺสเนน ปหาตพฺพ\-เหตุกํ ธมฺมํ ปจฺจยา นทสฺสเนน ปหาตพฺพ\-เหตุโก ธมฺโม อุปฺปชฺชติ อารมฺมณ\-ปจฺจยา. นทสฺสเนน ปหาตพฺพ\-เหตุกํ เอกํ ขนฺธํ ปจฺจยา ตโย ขนฺธา ฯเปฯ เทฺว ขนฺเธ ฯเปฯ. ปฏิสนฺธิ\-กฺขเณ ฯเปฯ. จกฺขายตนํ ปจฺจยา จกฺขุวิญฺญาณํ ฯเปฯ. วตฺถุํ ปจฺจยา นทสฺสเนน ปหาตพฺพ\-เหตุกา ขนฺธา. วตฺถุํ ปจฺจยา วิจิกิจฺฉา\-สหคโต โมโห.  นทสฺสเนน ปหาตพฺพ\-เหตุกํ ธมฺมํ ปจฺจยา ทสฺสเนน ปหาตพฺพ\-เหตุโก ธมฺโม อุปฺปชฺชติ อารมฺมณ\-ปจฺจยา. วตฺถุํ ปจฺจยา ทสฺสเนน ปหาตพฺพ\-เหตุกา ขนฺธา. วิจิกิจฺฉา\-สหคตํ โมหํ ปจฺจยา สมฺปยุตฺตกา ขนฺธา.  นทสฺสเนน ปหาตพฺพ\-เหตุกํ ธมฺมํ ปจฺจยา ทสฺสเนน ปหาตพฺพ\-เหตุโก จ นทสฺสเนน ปหาตพฺพ\-เหตุโก จ ธมฺมา อุปฺปชฺชนฺติ อารมฺมณ\-ปจฺจยา. วตฺถุํ ปจฺจยา วิจิกิจฺฉา\-สหคตา ขนฺธา จ โมโห จ. (3)}

\proseitem{60}{ทสฺสเนน ปหาตพฺพ\-เหตุ\-กญฺจ นทสฺสเนน ปหาตพฺพ\-เหตุ\-กญฺจ ธมฺมํ ปจฺจยา ทสฺสเนน ปหาตพฺพ\-เหตุโก ธมฺโม อุปฺปชฺชติ อารมฺมณ\-ปจฺจยา. ทสฺสเนน ปหาตพฺพ\-เหตุกํ เอกํ ขนฺธญฺจ วตฺถุญฺจ ปจฺจยา ตโย ขนฺธา ฯเปฯ เทฺว ขนฺเธ จ ฯเปฯ. วิจิกิจฺฉา\-สหคตํ เอกํ ขนฺธญฺจ โมหญฺจ ปจฺจยา ตโย ขนฺธา ฯเปฯ เทฺว ขนฺเธ จ ฯเปฯ. (1)}

\prose{ทสฺสเนน ปหาตพฺพ\-เหตุ\-กญฺจ นทสฺสเนน ปหาตพฺพ\-เหตุ\-กญฺจ ธมฺมํ ปจฺจยา นทสฺสเนน ปหาตพฺพ\-เหตุโก ธมฺโม อุปฺปชฺชติ อารมฺมณ\-ปจฺจยา. \csromanfolio{63}วิจิกิจฺฉา\-สหคเต ขนฺเธ จ วตฺถุญฺจ ปจฺจยา วิจิกิจฺฉา\-สหคโต โมโห. (2)}

\prose{ทสฺสเนน ปหาตพฺพ\-เหตุ\-กญฺจ นทสฺสเนน ปหาตพฺพ\-เหตุ\-กญฺจ ธมฺมํ ปจฺจยา ทสฺสเนน ปหาตพฺพ\-เหตุโก จ นทสฺสเนน ปหาตพฺพ\-เหตุโก จ ธมฺมา อุปฺปชฺชนฺติ อารมฺมณ\-ปจฺจยา. วิจิกิจฺฉา\-สหคตํ เอกํ ขนฺธญฺจ วตฺถุญฺจ ปจฺจยา ตโย ขนฺธา โมโห จ ฯเปฯ เทฺว ขนฺเธ จ ฯเปฯ. (3)}

\csromanmarkh{1. ปจฺจยานุโลม}\csromanmarkhii{2. สงฺขฺยาวาร}\csromanheadernomark{2}{1. \textbf{ปจฺจยานุโลม 2. สงฺขฺยาวาร}}
\proseitem{61}{เหตุยา นว, อารมฺมเณ นว, อธิปติยา นว, (สพฺพตฺถ นว.) วิปาเก เอกํ ฯเปฯ อวิคเต นว.}

\csromanmarkh{2. ปจฺจย\-ปจฺจนีย}\csromanmarkhii{1. วิภงฺควาร}\csromanheadernomark{2}{2. ปจฺจย\-ปจฺจนีย 1. วิภงฺควาร}
\proseitem{62}{ทสฺสเนน ปหาตพฺพ\-เหตุกํ ธมฺมํ ปจฺจยา นทสฺสเนน ปหาตพฺพ\-เหตุโก ธมฺโม อุปฺปชฺชติ นเหตุปจฺจยา. วิจิกิจฺฉา\-สหคเต ขนฺเธ ปจฺจยา วิจิกิจฺฉา\-สหคโต โมโห. (1)}

\prose{นทสฺสเนน ปหาตพฺพ\-เหตุกํ ธมฺมํ ปจฺจยา นทสฺสเนน ปหาตพฺพ\-เหตุโก ธมฺโม อุปฺปชฺชติ นเหตุปจฺจยา. อเหตุกํ นทสฺสเนน ปหาตพฺพ\-เหตุกํ ฯเปฯ. (ยาว อสญฺญสตฺตา.) จกฺขายตนํ ปจฺจยา จกฺขุวิญฺญาณํ ฯเปฯ. กายายตนํ ปจฺจยา กายวิญฺญาณํ. วตฺถุํ ปจฺจยา อเหตุกา นทสฺสเนน ปหาตพฺพ\-เหตุกา ขนฺธา. อุทฺธจฺจ\-สหคเต ขนฺเธ จ วตฺถุญฺจ ปจฺจยา อุทฺธจฺจ\-สหคโต โมโห. (1)}

\prose{ทสฺสเนน ปหาตพฺพ\-เหตุ\-กญฺจ นทสฺสเนน ปหาตพฺพ\-เหตุ\-กญฺจ ธมฺมํ ปจฺจยา นทสฺสเนน ปหาตพฺพ\-เหตุโก ธมฺโม อุปฺปชฺชติ นเหตุปจฺจยา. วิจิกิจฺฉา\-สหคเต ขนฺเธ จ วตฺถุญฺจ ปจฺจยา วิจิกิจฺฉา\-สหคโต โมโห. (สํขิตฺตํ.) (1)}

\csromanmarkh{2. ปจฺจย\-ปจฺจนีย}\csromanmarkhii{2. สงฺขฺยาวาร}\csromanheadernomark{2}{2. ปจฺจย\-ปจฺจนีย 2. สงฺขฺยาวาร}
\proseitem{63}{\csromanfolio{64}นเหตุยา ตีณิ, นอารมฺมเณ ตีณิ, นอธิปติยา นว, นอนนฺตเร ตีณิ, นสมนนฺตเร ตีณิ, นอญฺญมญฺเญ ตีณิ, นอุปนิสฺสเย ตีณิ, นปุเรชาเต สตฺต, นปจฺฉาชาเต นว, นอาเสวเน นว, นกมฺเม จตฺตาริ, นวิปาเก นว, นอาหาเร เอกํ, นอินฺทฺริเย เอกํ, นฌาเน เอกํ, นมคฺเค เอกํ, นสมฺปยุตฺเต ตีณิ, นวิปฺปยุตฺเต ฉ, โนนตฺถิยา ตีณิ, โนวิคเต ตีณิ.}

\vspace{\tipitakaparskip}
\csromancenter{(เอวํ อิตเร เทฺว คณนาปิ นิสฺสยวาโรปิ กาตพฺโพ.)}
\csromanmarkh{85. ทสฺสเนน\-ปหาตพฺพ\-เหตุกทุก}\csromanmarkhii{5. สํสฏฺฐวาร}\csromanheadernomark{2}{85. ทสฺสเนน\-ปหาตพฺพ\-เหตุกทุก 5. สํสฏฺฐวาร}
\csromanheader{2}{\textbf{ปจฺจยจตุกฺก}}
\proseitem{64}{ทสฺสเนน ปหาตพฺพ\-เหตุกํ ธมฺมํ สํสฏฺโฐ ทสฺสเนน ปหาตพฺพ\-เหตุโก ธมฺโม อุปฺปชฺชติ เหตุปจฺจยา. ทสฺสเนน ปหาตพฺพ\-เหตุกํ เอกํ ขนฺธํ สํสฏฺฐา ตโย ขนฺธา ฯเปฯ เทฺว ขนฺเธ ฯเปฯ. (1)}

\prose{นทสฺสเนน ปหาตพฺพ\-เหตุกํ ธมฺมํ สํสฏฺโฐ นทสฺสเนน ปหาตพฺพ\-เหตุโก ธมฺโม อุปฺปชฺชติ เหตุปจฺจยา. นทสฺสเนน ปหาตพฺพ\-เหตุกํ เอกํ ขนฺธํ สํสฏฺฐา ตโย ขนฺธา ฯเปฯ เทฺว ขนฺเธ ฯเปฯ. ปฏิสนฺธิ\-กฺขเณ ฯเปฯ. (1)}

\prose{นทสฺสเนน ปหาตพฺพ\-เหตุกํ ธมฺมํ สํสฏฺโฐ ทสฺสเนน ปหาตพฺพ\-เหตุโก ธมฺโม อุปฺปชฺชติ เหตุปจฺจยา. วิจิกิจฺฉา\-สหคตํ โมหํ สํสฏฺฐา สมฺปยุตฺตกา ขนฺธา. (2)}

\prose{ทสฺสเนน ปหาตพฺพ\-เหตุ\-กญฺจ นทสฺสเนน ปหาตพฺพ\-เหตุ\-กญฺจ ธมฺมํ สํสฏฺโฐ ทสฺสเนน ปหาตพฺพ\-เหตุโก ธมฺโม อุปฺปชฺชติ เหตุปจฺจยา. วิจิกิจฺฉา\-สหคตํ เอกํ ขนฺธญฺจ โมหญฺจ สํสฏฺฐา ตโย ขนฺธา ฯเปฯ เทฺว ขนฺเธ จ ฯเปฯ. (1)}

\csromanheader{2}{85. ทสฺสเนน\-ปหาตพฺพ\-เหตุกทุก อารมฺมณ}
\proseitem{65}{\csromanfolio{65}ทสฺสเนน ปหาตพฺพ\-เหตุกํ ธมฺมํ สํสฏฺโฐ ทสฺสเนน ปหาตพฺพ\-เหตุโก ธมฺโม อุปฺปชฺชติ อารมฺมณ\-ปจฺจยา. ตีณิ. (ปฏิจฺจสทิสา.)}

\prose{นทสฺสเนน ปหาตพฺพ\-เหตุกํ ธมฺมํ สํสฏฺโฐ นทสฺสเนน ปหาตพฺพ\-เหตุโก ธมฺโม อุปฺปชฺชติ อารมฺมณ\-ปจฺจยา. นทสฺสเนน ปหาตพฺพ\-เหตุกํ เอกํ ขนฺธํ สํสฏฺฐา ตโย ขนฺธา ฯเปฯ เทฺว ขนฺเธ ฯเปฯ. ปฏิสนฺธิ\-กฺขเณ ฯเปฯ. (1)}

\prose{นทสฺสเนน ปหาตพฺพ\-เหตุกํ ธมฺมํ สํสฏฺโฐ ทสฺสเนน ปหาตพฺพ\-เหตุโก ธมฺโม อุปฺปชฺชติ อารมฺมณ\-ปจฺจยา. (ปฏิจฺจสทิสํ.) (2)}

\prose{ทสฺสเนน ปหาตพฺพ\-เหตุ\-กญฺจ นทสฺสเนน ปหาตพฺพ\-เหตุ\-กญฺจ ธมฺมํ สํสฏฺโฐ ทสฺสเนน ปหาตพฺพ\-เหตุโก ธมฺโม อุปฺปชฺชติ อารมฺมณ\-ปจฺจยา. (ปฏิจฺจสทิสํ. สํขิตฺตํ.) (1)}

\proseitem{66}{เหตุยา จตฺตาริ, อารมฺมเณ ฉ, อธิปติยา เทฺว, อนนฺตเร ฉ, (สพฺพตฺถ ฉ,) วิปาเก เอกํ ฯเปฯ อวิคเต ฉ. อนุโลมํ.}

\proseitem{67}{ทสฺสเนน ปหาตพฺพ\-เหตุกํ ธมฺมํ สํสฏฺโฐ นทสฺสเนน ปหาตพฺพ\-เหตุโก ธมฺโม อุปฺปชฺชติ นเหตุปจฺจยา. วิจิกิจฺฉา\-สหคเต ขนฺเธ สํสฏฺโฐ วิจิกิจฺฉา\-สหคโต โมโห. (1)}

\prose{นทสฺสเนน ปหาตพฺพ\-เหตุกํ ธมฺมํ สํสฏฺโฐ นทสฺสเนน ปหาตพฺพ\-เหตุโก ธมฺโม อุปฺปชฺชติ นเหตุปจฺจยา. อเหตุกํ นทสฺสเนน ปหาตพฺพ\-เหตุกํ เอกํ ขนฺธํ สํสฏฺฐา ตโย ขนฺธา ฯเปฯ เทฺว ขนฺเธ ฯเปฯ. อเหตุก\-ปฏิสนฺธิ\-กฺขเณ ฯเปฯ. อุทฺธจฺจ\-สหคเต ขนฺเธ สํสฏฺโฐ อุทฺธจฺจ\-สหคโต โมโห. (1)}

\vspace{\tipitakaparskip}
\csromancenter{\textbf{นเหตุ}ยา เทฺว, นอธิปติยา ฉ, นปุเรชาเต ฉ, นปจฺฉาชาเต ฉ, นอาเสวเน ฉ, นกมฺเม จตฺตาริ, นวิปาเก ฉ, นฌาเน เอกํ, นมคฺเค เอกํ, นวิปฺปยุตฺเต ฉ.}
\csromancenter{ปจฺจนียํ.}
\csromanheader{2}{\textbf{เหตุทุก}}
\vspace{\tipitakaparskip}
\csromancenter{\textbf{เหตุ}ปจฺจยา นอธิปติยา จตฺตาริ, นปุเรชาเต จตฺตาริ ฯเปฯ นวิปฺปยุตฺเต จตฺตาริ.}
\csromanheader{2}{\textbf{นเหตุทุก}}
\vspace{\tipitakaparskip}
\csromancenter{นเหตุปจฺจยา อารมฺมเณ เทฺว, อนนฺตเร เทฺว ฯเปฯ วิปาเก เอกํ ฯเปฯ อวิคเต เทฺว.}
\csromancenter{(สมฺปยุตฺตวาโร สํสฏฺฐวาร\-สทิโส.)}
\csromanmarkh{85. ทสฺสเนน\-ปหาตพฺพ\-เหตุกทุก}\csromanmarkhii{7. ปญฺหาวาร}\csromanheadernomark{2}{85. ทสฺสเนน\-ปหาตพฺพ\-เหตุกทุก 7. ปญฺหาวาร}
\proseitem{68}{\csromanfolio{66}ทสฺสเนน ปหาตพฺพ\-เหตุโก ธมฺโม ทสฺสเนน ปหาตพฺพ\-เหตุกสฺส ธมฺมสฺส เหตุปจฺจเยน ปจฺจโย. ทสฺสเนน ปหาตพฺพ\-เหตุกา เหตู สมฺปยุตฺตกานํ ขนฺธานํ เหตุปจฺจเยน ปจฺจโย. (มูลํ กาตพฺพํ.) ทสฺสเนน ปหาตพฺพ\-เหตุกา เหตู จิตฺต\-สมุฏฺฐานานํ รูปานํ เหตุปจฺจเยน ปจฺจโย. (มูลํ กาตพฺพํ.) ทสฺสเนน ปหาตพฺพ\-เหตุกา เหตู สมฺปยุตฺตกานํ ขนฺธานํ จิตฺตสมุฏฺฐานานญฺจ รูปานํ เหตุปจฺจเยน ปจฺจโย. (3)}

\proseitem{69}{นทสฺสเนน ปหาตพฺพ\-เหตุโก ธมฺโม นทสฺสเนน ปหาตพฺพ\-เหตุกสฺส ธมฺมสฺส เหตุปจฺจเยน ปจฺจโย. นทสฺสเนน ปหาตพฺพ\-เหตุกา เหตู สมฺปยุตฺตกานํ ขนฺธานํ จิตฺตสมุฏฺฐานานญฺจ รูปานํ เหตุปจฺจเยน ปจฺจโย. วิจิกิจฺฉา\-สหคโต โมโห จิตฺต\-สมุฏฺฐานานํ รูปานํ เหตุปจฺจเยน ปจฺจโย. ปฏิสนฺธิ\-กฺขเณ ฯเปฯ. (มูลํ กาตพฺพํ.) วิจิกิจฺฉา\-สหคโต โมโห สมฺปยุตฺตกานํ ขนฺธานํ เหตุปจฺจเยน ปจฺจโย. (มูลํ กาตพฺพํ.) วิจิกิจฺฉา\-สหคโต โมโห สมฺปยุตฺตกานํ ขนฺธานํ จิตฺตสมุฏฺฐานานญฺจ รูปานํ เหตุปจฺจเยน ปจฺจโย. (3)}

\csromanheader{2}{85. ทสฺสเนน\-ปหาตพฺพ\-เหตุกทุก อารมฺมณ}
\proseitem{70}{\csromanfolio{67}ทสฺสเนน ปหาตพฺพ\-เหตุโก ธมฺโม ทสฺสเนน ปหาตพฺพ\-เหตุกสฺส ธมฺมสฺส อารมฺมณ\-ปจฺจเยน ปจฺจโย. ทสฺสเนน ปหาตพฺพ\-เหตุเก ขนฺเธ อารพฺภ ทสฺสเนน ปหาตพฺพ\-เหตุกา ขนฺธา อุปฺปชฺชนฺติ. (มูลํ กาตพฺพํ.) ทสฺสเนน ปหาตพฺพ\-เหตุเก ขนฺเธ อารพฺภ นทสฺสเนน ปหาตพฺพ\-เหตุกา ขนฺธา จ โมโห จ อุปฺปชฺชนฺติ. (มูลํ กาตพฺพํ.) ทสฺสเนน ปหาตพฺพ\-เหตุเก ขนฺเธ อารพฺภ วิจิกิจฺฉา\-สหคตา ขนฺธา จ โมโห จ อุปฺปชฺชนฺติ. (3)}

\proseitem{71}{นทสฺสเนน ปหาตพฺพ\-เหตุโก ธมฺโม นทสฺสเนน ปหาตพฺพ\-เหตุกสฺส ธมฺมสฺส อารมฺมณ\-ปจฺจเยน ปจฺจโย. ทานํ ฯเปฯ สีลํ ฯเปฯ อุโปสถกมฺมํ กตฺวา ตํ ปจฺจเวกฺขติ, อสฺสาเทติ อภินนฺทติ, ตํ อารพฺภ นทสฺสเนน ปหาตพฺพ\-เหตุโก ราโค อุปฺปชฺชติ. อุทฺธจฺจํ อุปฺปชฺชติ. นทสฺสเนน ปหาตพฺพ\-เหตุกํ โทมนสฺสํ อุปฺปชฺชติ. ปุพฺเพ สุจิณฺณานิ ฯเปฯ. ฌานา วุฏฺฐหิตฺวา ฯเปฯ. อริยา มคฺคา วุฏฺฐหิตฺวา มคฺคํ ปจฺจเวกฺขนฺติ ฯเปฯ. อริยา นทสฺสเนน ปหาตพฺพ\-เหตุเก ปหีเน กิเลเส ปจฺจเวกฺขนฺติ. วิกฺขมฺภิเต กิเลเส ฯเปฯ. ปุพฺเพ สมุทาจิณฺเณ กิเลเส ชานนฺติ. จกฺขุํ ฯเปฯ. วตฺถุํ นทสฺสเนน ปหาตพฺพ\-เหตุเก ขนฺเธ จ โมหญฺจ อนิจฺจโต ฯเปฯ วิปสฺสติ, อสฺสาเทติ อภินนฺทติ, ตํ อารพฺภ นทสฺสเนน ปหาตพฺพ\-เหตุโก ราโค อุปฺปชฺชติ. อุทฺธจฺจํ อุปฺปชฺชติ. นทสฺสเนน ปหาตพฺพ\-เหตุกํ โทมนสฺสํ อุปฺปชฺชติ. ทิพฺเพน จกฺขุนา รูปํ ปสฺสติ ฯเปฯ อนาคตํสญาณสฺส, อาวชฺชนาย, โมหสฺส จ อารมฺมณ\-ปจฺจเยน ปจฺจโย. (1)}

\prose{นทสฺสเนน ปหาตพฺพ\-เหตุโก ธมฺโม ทสฺสเนน ปหาตพฺพ\-เหตุกสฺส ธมฺมสฺส อารมฺมณ\-ปจฺจเยน ปจฺจโย. ทานํ ฯเปฯ. ฌานํ ฯเปฯ. จกฺขุํ ฯเปฯ. วตฺถุํ นทสฺสเนน ปหาตพฺพ\-เหตุเก ขนฺเธ จ โมหญฺจ อสฺสาเทติ อภินนฺทติ, ตํ อารพฺภ ทสฺสเนน ปหาตพฺพ\-เหตุโก ราโค อุปฺปชฺชติ. ทิฏฺฐิ ฯเปฯ วิจิกิจฺฉา ฯเปฯ. ทสฺสเนน ปหาตพฺพ\-เหตุกํ โทมนสฺสํ อุปฺปชฺชติ. (2)}

\prose{นทสฺสเนน ปหาตพฺพ\-เหตุโก ธมฺโม ทสฺสเนน ปหาตพฺพ\-เหตุกสฺส จ นทสฺสเนน ปหาตพฺพ\-เหตุกสฺส จ ธมฺมสฺส อารมฺมณ\-ปจฺจเยน ปจฺจโย. จกฺขุํ ฯเปฯ. วตฺถุํ นทสฺสเนน ปหาตพฺพ\-เหตุเก ขนฺเธ จ โมหญฺจ อารพฺภ วิจิกิจฺฉา\-สหคตา ขนฺธา จ โมโห จ อุปฺปชฺชนฺติ. (3)}

\proseitem{72}{\csromanfolio{68}ทสฺสเนน ปหาตพฺพ\-เหตุโก จ นทสฺสเนน ปหาตพฺพ\-เหตุโก จ ธมฺมา ทสฺสเนน ปหาตพฺพ\-เหตุกสฺส ธมฺมสฺส อารมฺมณ\-ปจฺจเยน ปจฺจโย. วิจิกิจฺฉา\-สหคเต ขนฺเธ จ โมหญฺจ อารพฺภ ทสฺสเนน ปหาตพฺพ\-เหตุกา ขนฺธา อุปฺปชฺชนฺติ. (มูลํ กาตพฺพํ.) วิจิกิจฺฉา\-สหคเต ขนฺเธ จ โมหญฺจ อารพฺภ นทสฺสเนน ปหาตพฺพ\-เหตุกา ขนฺธา จ โมโห จ อุปฺปชฺชนฺติ. (มูลํ กาตพฺพํ.) วิจิกิจฺฉา\-สหคเต ขนฺเธ จ โมหญฺจ อารพฺภ วิจิกิจฺฉา\-สหคตา ขนฺธา จ โมโห จ อุปฺปชฺชนฺติ. (3)}

\csromanheader{2}{\textbf{อธิปติ}}
\proseitem{73}{ทสฺสเนน ปหาตพฺพ\-เหตุโก ธมฺโม ทสฺสเนน ปหาตพฺพ\-เหตุกสฺส ธมฺมสฺส อธิปติ\-ปจฺจเยน ปจฺจโย. อารมฺมณา\-ธิปติ, สหชาตา\-ธิปติ.  อารมฺมณา\-ธิปติ\csromandash{}ทสฺสเนน ปหาตพฺพ\-เหตุเก ขนฺเธ ครุํ กตฺวา ทสฺสเนน ปหาตพฺพ\-เหตุกา ขนฺธา อุปฺปชฺชนฺติ. สหชาตา\-ธิปติ\csromandash{}ทสฺสเนน ปหาตพฺพ\-เหตุกา\-ธิปติ สมฺปยุตฺตกานํ ขนฺธานํ อธิปติ\-ปจฺจเยน ปจฺจโย. (1)}

\prose{ทสฺสเนน ปหาตพฺพ\-เหตุโก ธมฺโม นทสฺสเนน ปหาตพฺพ\-เหตุกสฺส ธมฺมสฺส อธิปติ\-ปจฺจเยน ปจฺจโย. สหชาตา\-ธิปติ\csromandash{}ทสฺสเนน ปหาตพฺพ\-เหตุกา\-ธิปติ จิตฺต\-สมุฏฺฐานานํ รูปานํ อธิปติ\-ปจฺจเยน ปจฺจโย. (2)}

\prose{ทสฺสเนน ปหาตพฺพ\-เหตุโก ธมฺโม ทสฺสเนน ปหาตพฺพ\-เหตุกสฺส จ นทสฺสเนน ปหาตพฺพ\-เหตุกสฺส จ ธมฺมสฺส อธิปติ\-ปจฺจเยน ปจฺจโย. สหชาตา\-ธิปติ\csromandash{}ทสฺสเนน ปหาตพฺพ\-เหตุกา\-ธิปติ สมฺปยุตฺตกานํ ขนฺธานํ จิตฺตสมุฏฺฐานานญฺจ รูปานํ อธิปติ\-ปจฺจเยน ปจฺจโย. (3)}

\proseitem{74}{นทสฺสเนน ปหาตพฺพ\-เหตุโก ธมฺโม นทสฺสเนน ปหาตพฺพ\-เหตุกสฺส ธมฺมสฺส อธิปติ\-ปจฺจเยน ปจฺจโย. อารมฺมณา\-ธิปติ, สหชาตา\-ธิปติ.  อารมฺมณา\-ธิปติ\csromandash{}ทานํ ฯเปฯ สีลํ ฯเปฯ อุโปสถกมฺมํ กตฺวา ตํ ครุํ กตฺวา ปจฺจเวกฺขติ, อสฺสาเทติ อภินนฺทติ, ตํ ครุํ กตฺวา นทสฺสเนน ปหาตพฺพ\-เหตุโก ราโค อุปฺปชฺชติ. ปุพฺเพ สุจิณฺณานิ ฯเปฯ. ฌานา ฯเปฯ. อริยา มคฺคา วุฏฺฐหิตฺวา มคฺคํ ครุํ กตฺวา ฯเปฯ \csromanfolio{69}ผลสฺส อธิปติ\-ปจฺจเยน ปจฺจโย. จกฺขุํ ฯเปฯ. วตฺถุํ นทสฺสเนน ปหาตพฺพ\-เหตุเก ขนฺเธ ครุํ กตฺวา อสฺสาเทติ อภินนฺทติ, ตํ ครุํ กตฺวา นทสฺสเนน ปหาตพฺพ\-เหตุโก ราโค อุปฺปชฺชติ. สหชาตา\-ธิปติ\csromandash{} นทสฺสเนน ปหาตพฺพ\-เหตุกา\-ธิปติ สมฺปยุตฺตกานํ ขนฺธานํ จิตฺตสมุฏฺฐานานญฺจ รูปานํ อธิปติ\-ปจฺจเยน ปจฺจโย. (1)}

\prose{นทสฺสเนน ปหาตพฺพ\-เหตุโก ธมฺโม ทสฺสเนน ปหาตพฺพ\-เหตุกสฺส ธมฺมสฺส อธิปติ\-ปจฺจเยน ปจฺจโย. อารมฺมณา\-ธิปติ\csromandash{}ทานํ ฯเปฯ สีลํ ฯเปฯ อุโปสถกมฺมํ กตฺวา ฯเปฯ. ฌานา ฯเปฯ. จกฺขุํ ฯเปฯ. วตฺถุํ นทสฺสเนน ปหาตพฺพ\-เหตุเก ขนฺเธ ครุํ กตฺวา อสฺสาเทติ อภินนฺทติ, ตํ ครุํ กตฺวา ทสฺสเนน ปหาตพฺพ\-เหตุโก ราโค อุปฺปชฺชติ. ทิฏฺฐิ อุปฺปชฺชติ. (2)}

\csromanheader{2}{\textbf{อนนฺตร}}
\proseitem{74}{ทสฺสเนน ปหาตพฺพ\-เหตุโก ธมฺโม ทสฺสเนน ปหาตพฺพ\-เหตุกสฺส ธมฺมสฺส อนนฺตร\-ปจฺจเยน ปจฺจโย. ปุริมา ปุริมา ทสฺสเนน ปหาตพฺพ\-เหตุกา ขนฺธา ปจฺฉิมานํ ปจฺฉิมานํ ทสฺสเนน ปหาตพฺพ\-เหตุกานํ ขนฺธานํ อนนฺตร\-ปจฺจเยน ปจฺจโย. (มูลํ กาตพฺพํ.) ปุริมา ปุริมา วิจิกิจฺฉา\-สหคตา ขนฺธา ปจฺฉิมสฺส ปจฺฉิมสฺส วิจิกิจฺฉา\-สหคตสฺส โมหสฺส อนนฺตร\-ปจฺจเยน ปจฺจโย. ทสฺสเนน ปหาตพฺพ\-เหตุกา ขนฺธา วุฏฺฐานสฺส อนนฺตร\-ปจฺจเยน ปจฺจโย. (มูลํ กาตพฺพํ.) ปุริมา ปุริมา วิจิกิจฺฉา\-สหคตา ขนฺธา ปจฺฉิมานํ ปจฺฉิมานํ วิจิกิจฺฉา\-สหคตานํ ขนฺธานํ โมหสฺส จ อนนฺตร\-ปจฺจเยน ปจฺจโย. (3)}

\proseitem{76}{นทสฺสเนน ปหาตพฺพ\-เหตุโก ธมฺโม นทสฺสเนน ปหาตพฺพ\-เหตุกสฺส ธมฺมสฺส อนนฺตร\-ปจฺจเยน ปจฺจโย. ปุริโม ปุริโม วิจิกิจฺฉา\-สหคโต โมโห ปจฺฉิมสฺส ปจฺฉิมสฺส วิจิกิจฺฉา\-สหคตสฺส โมหสฺส อนนฺตร\-ปจฺจเยน ปจฺจโย. ปุริมา ปุริมา นทสฺสเนน ปหาตพฺพ\-เหตุกา ขนฺธา ฯเปฯ ผลสมาปตฺติยา อนนฺตร\-ปจฺจเยน ปจฺจโย. (มูลํ กาตพฺพํ.) ปุริโม ปุริโม วิจิกิจฺฉา\-สหคโต โมโห ปจฺฉิมานํ ปจฺฉิมานํ วิจิกิจฺฉา\-สหคตานํ ขนฺธานํ อนนฺตร\-ปจฺจเยน ปจฺจโย. อาวชฺชนา ทสฺสเนน ปหาตพฺพ\-เหตุกานํ ขนฺธานํ อนนฺตร\-ปจฺจเยน ปจฺจโย. (มูลํ กาตพฺพํ.) ปุริโม ปุริโม วิจิกิจฺฉา\-สหคโต โมโห ปจฺฉิมานํ \csromanfolio{70}ปจฺฉิมานํ วิจิกิจฺฉา\-สหคตานํ ขนฺธานํ โมหสฺส จ อนนฺตร\-ปจฺจเยน ปจฺจโย. อาวชฺชนา วิจิกิจฺฉา\-สหคตานํ ขนฺธานํ โมหสฺส จ อนนฺตร\-ปจฺจเยน ปจฺจโย. (3)}

\prose{ทสฺสเนน ปหาตพฺพ\-เหตุโก จ นทสฺสเนน ปหาตพฺพ\-เหตุโก จ ธมฺมา ทสฺสเนน ปหาตพฺพ\-เหตุกสฺส ธมฺมสฺส อนนฺตร\-ปจฺจเยน ปจฺจโย. ปุริมา ปุริมา วิจิกิจฺฉา\-สหคตา ขนฺธา จ โมโห จ ปจฺฉิมานํ ปจฺฉิมานํ วิจิกิจฺฉา\-สหคตานํ ขนฺธานํ อนนฺตร\-ปจฺจเยน ปจฺจโย. (มูลํ กาตพฺพํ.) ปุริมา ปุริมา วิจิกิจฺฉา\-สหคตา ขนฺธา จ โมโห จ ปจฺฉิมสฺส ปจฺฉิมสฺส วิจิกิจฺฉา\-สหคตสฺส โมหสฺส อนนฺตร\-ปจฺจเยน ปจฺจโย. วิจิกิจฺฉา\-สหคตา ขนฺธา จ โมโห จ วุฏฺฐานสฺส อนนฺตร\-ปจฺจเยน ปจฺจโย. (มูลํ กาตพฺพํ.) ปุริมา ปุริมา วิจิกิจฺฉา\-สหคตา ขนฺธา จ โมโห จ ปจฺฉิมานํ ปจฺฉิมานํ วิจิกิจฺฉา\-สหคตานํ ขนฺธานํ โมหสฺส จ อนนฺตร\-ปจฺจเยน ปจฺจโย. (3)}

\prose{สมนนฺตร\-ปจฺจเยน ปจฺจโย. นว. สหชาต\-ปจฺจเยน ปจฺจโย. นว. อญฺญมญฺญ\-ปจฺจเยน ปจฺจโย. ฉ. นิสฺสย\-ปจฺจเยน ปจฺจโย. นว.}

\csromanheader{2}{\textbf{อุปนิสฺสย}}
\proseitem{78}{ทสฺสเนน ปหาตพฺพ\-เหตุโก ธมฺโม ทสฺสเนน ปหาตพฺพ\-เหตุกสฺส ธมฺมสฺส อุปนิสฺสย\-ปจฺจเยน ปจฺจโย. อารมฺมณู\-ปนิสฺสโย, อนนฺตรู\-ปนิสฺสโย, ปกตู\-ปนิสฺสโย ฯเปฯ. ปกตู\-ปนิสฺสโย\csromandash{}ทสฺสเนน ปหาตพฺพ\-เหตุกา ขนฺธา ทสฺสเนน ปหาตพฺพ\-เหตุกานํ ขนฺธานํ อุปนิสฺสย\-ปจฺจเยน ปจฺจโย. (อวเสเสสุ ทฺวีสุ อนนฺตรู\-ปนิสฺสโย, ปกตู\-ปนิสฺสโย.) (มูลํ กาตพฺพํ.) ทสฺสเนน ปหาตพฺพ\-เหตุกา ขนฺธา นทสฺสเนน ปหาตพฺพ\-เหตุกานํ ขนฺธานํ โมหสฺส จ อุปนิสฺสย\-ปจฺจเยน ปจฺจโย. (มูลํ กาตพฺพํ.) ทสฺสเนน ปหาตพฺพ\-เหตุกา ขนฺธา วิจิกิจฺฉา\-สหคตานํ ขนฺธานํ โมหสฺส จ อุปนิสฺสย\-ปจฺจเยน ปจฺจโย. (3)}

\proseitem{79}{นทสฺสเนน ปหาตพฺพ\-เหตุโก ธมฺโม นทสฺสเนน ปหาตพฺพ\-เหตุกสฺส ธมฺมสฺส อุปนิสฺสย\-ปจฺจเยน ปจฺจโย. อารมฺมณู\-ปนิสฺสโย, อนนฺตรู\-ปนิสฺสโย, ปกตู\-ปนิสฺสโย ฯเปฯ. ปกตู\-ปนิสฺสโย\csromandash{}สทฺธํ อุปนิสฺสาย ทานํ เทติ ฯเปฯ สมาปตฺติํ อุปฺปาเทติ. มานํ ชปฺเปติ. \csromanfolio{71}สีลํ ฯเปฯ. ปญฺญํ. นทสฺสเนน ปหาตพฺพ\-เหตุกํ ราคํ. โทสํ. โมหํ. มานํ. ปตฺถนํ. กายิกํ สุขํ. กายิกํ ทุกฺขํ. อุตุํ. โภชนํ. เสนาสนํ อุปนิสฺสาย ทานํ เทติ ฯเปฯ สมาปตฺติํ อุปฺปาเทติ. สทฺธา ฯเปฯ. เสนาสนํ สทฺธาย ฯเปฯ ปญฺญาย. นทสฺสเนน ปหาตพฺพ\-เหตุกสฺส ราคสฺส ฯเปฯ ปตฺถนาย ผลสมาปตฺติยา อุปนิสฺสย\-ปจฺจเยน ปจฺจโย. (1)}

\prose{นทสฺสเนน ปหาตพฺพ\-เหตุโก ธมฺโม ทสฺสเนน ปหาตพฺพ\-เหตุกสฺส ธมฺมสฺส อุปนิสฺสย\-ปจฺจเยน ปจฺจโย. อารมฺมณู\-ปนิสฺสโย, อนนฺตรู\-ปนิสฺสโย, ปกตู\-ปนิสฺสโย ฯเปฯ. ปกตู\-ปนิสฺสโย\csromandash{}สทฺธํ อุปนิสฺสาย ทิฏฺฐิํ คณฺหาติ. สีลํ ฯเปฯ. ปญฺญํ. นทสฺสเนน ปหาตพฺพ\-เหตุกํ ราคํ. โทสํ. โมหํ. มานํ. ปตฺถนํ. กายิกํ สุขํ. กายิกํ ทุกฺขํ. อุตุํ. โภชนํ. เสนาสนํ อุปนิสฺสาย ปาณํ หนติ ฯเปฯ สํฆํ ภินฺทติ. สทฺธา ฯเปฯ. เสนาสนํ ทสฺสเนน ปหาตพฺพ\-เหตุกสฺส ราคสฺส. โทสสฺส. โมหสฺส. ทิฏฺฐิยา. ปตฺถนาย อุปนิสฺสย\-ปจฺจเยน ปจฺจโย. (2)}

\prose{นทสฺสเนน ปหาตพฺพ\-เหตุโก ธมฺโม ทสฺสเนน ปหาตพฺพ\-เหตุกสฺส จ นทสฺสเนน ปหาตพฺพ\-เหตุกสฺส จ ธมฺมสฺส อุปนิสฺสย\-ปจฺจเยน ปจฺจโย. อนนฺตรู\-ปนิสฺสโย, ปกตู\-ปนิสฺสโย ฯเปฯ. ปกตู\-ปนิสฺสโย\csromandash{} สทฺธา ฯเปฯ ปญฺญา. นทสฺสเนน ปหาตพฺพ\-เหตุโก ราโค. โทโส. โมโห. มาโน. ปตฺถนา. กายิกํ สุขํ ฯเปฯ. เสนาสนํ. วิจิกิจฺฉา\-สหคตานํ ขนฺธานํ โมหสฺส จ อุปนิสฺสย\-ปจฺจเยน ปจฺจโย. (3)}

\proseitem{80}{ทสฺสเนน ปหาตพฺพ\-เหตุโก จ นทสฺสเนน ปหาตพฺพ\-เหตุโก จ ธมฺมา ทสฺสเนน ปหาตพฺพ\-เหตุกสฺส ธมฺมสฺส อุปนิสฺสย\-ปจฺจเยน ปจฺจโย. อนนฺตรู\-ปนิสฺสโย, ปกตู\-ปนิสฺสโย ฯเปฯ. ปกตู\-ปนิสฺสโย\csromandash{} วิจิกิจฺฉา\-สหคตา ขนฺธา จ โมโห จ ทสฺสเนน ปหาตพฺพ\-เหตุกานํ ขนฺธานํ อุปนิสฺสย\-ปจฺจเยน ปจฺจโย. (มูลํ กาตพฺพํ.) วิจิกิจฺฉา\-สหคตา ขนฺธา จ โมโห จ นทสฺสเนน ปหาตพฺพ\-เหตุกานํ ขนฺธานํ โมหสฺส จ อุปนิสฺสย\-ปจฺจเยน ปจฺจโย. (มูลํ กาตพฺพํ.) วิจิกิจฺฉา\-สหคตา ขนฺธา จ โมโห จ วิจิกิจฺฉา\-สหคตานํ ขนฺธานํ โมหสฺส จ อุปนิสฺสย\-ปจฺจเยน ปจฺจโย. (3)}

\csromanheader{2}{\textbf{ปุเรชาต}}
\proseitem{81}{\csromanfolio{72}นทสฺสเนน ปหาตพฺพ\-เหตุโก ธมฺโม นทสฺสเนน ปหาตพฺพ\-เหตุกสฺส ธมฺมสฺส ปุเรชาต\-ปจฺจเยน ปจฺจโย. อารมฺมณ\-ปุเรชาตํ, วตฺถุ\-ปุเรชาตํ.  อารมฺมณ\-ปุเรชาตํ\csromandash{}จกฺขุํ ฯเปฯ. วตฺถุํ อนิจฺจโต ฯเปฯ วฺปสฺสติ, อสฺสาเทตุ อภินนฺทติ, ตํ อารพฺภ นทสฺสเนน ปหาตพฺพ\-เหตุโก ราโค อุปฺปชฺชติ. อุทฺธจฺจํ อุปฺปชฺชติ. นทสฺสเนน ปหาตพฺพ\-เหตุกํ โทมนสฺสํ อุปฺปชฺชติ. ทิพฺเพน ฯเปฯ. (สํขิตฺตํ. วตฺถุ\-ปุเรชาตํ. สํขิตฺตํ.) (1)}

\prose{นทสฺสเนน ปหาตพฺพ\-เหตุโก ธมฺโม ทสฺสเนน ปหาตพฺพ\-เหตุกสฺส ธมฺมสฺส ปุเรชาต\-ปจฺจเยน ปจฺจโย. อารมฺมณ\-ปุเรชาตํ, วตฺถุ\-ปุเรชาตํ.  อารมฺมณ\-ปุเรชาตํ\csromandash{}จกฺขุํ ฯเปฯ. วตฺถุํ อสฺสาเทติ อภินนฺทติ, ตํ อารพฺภ ทสฺสเนน ปหาตพฺพ\-เหตุโก ราโค อุปฺปชฺชติ. ทิฏฺฐิ ฯเปฯ วิจิกิจฺฉา ฯเปฯ. ทสฺสเนน ปหาตพฺพ\-เหตุกํ โทมนสฺสํ อุปฺปชฺชติ. (วตฺถุ\-ปุเรชาตํ. สํขิตฺตํ.) (2)}

\prose{นทสฺสเนน ปหาตพฺพ\-เหตุโก ธมฺโม ทสฺสเนน ปหาตพฺพ\-เหตุกสฺส จ นทสฺสเนน ปหาตพฺพ\-เหตุกสฺส จ ธมฺมสฺส ปุเรชาต\-ปจฺจเยน ปจฺจโย. อารมฺมณ\-ปุเรชาตํ, วตฺถุ\-ปุเรชาตํ.  อารมฺมณ\-ปุเรชาตํ\csromandash{} จกฺขุํ ฯเปฯ. วตฺถุํ อารพฺภ วิจิกิจฺฉา\-สหคตา ขนฺธา จ โมโห จ อุปฺปชฺชนฺติ. (วตฺถุ\-ปุเรชาตํ. สํขิตฺตํ.) (3)}

\csromanheader{2}{ปจฺฉาชาตา\-เสวน}
\proseitem{82}{ทสฺสเนน ปหาตพฺพ\-เหตุโก ธมฺโม นทสฺสเนน ปหาตพฺพ\-เหตุกสฺส ธมฺมสฺส ปจฺฉาชาต\-ปจฺจเยน ปจฺจโย. (สํกฺขิตฺตํ.).  นทสฺสเนน ปหาตพฺพ\-เหตุโก ธมฺโม นทสฺสเนน ปหาตพฺพ\-เหตุกสฺส ธมฺมสฺส ปจฺฉาชาต\-ปจฺจเยน ปจฺจโย. (สํขิตฺตํ.).  ทสฺสเนน ปหาตพฺพ\-เหตุโก จ นทสฺสเนน ปหาตพฺพ\-เหตุโก จ ธมฺมา นทสฺสเนน ปหาตพฺพ\-เหตุกสฺส ธมฺมสฺส ปจฺฉาชาต\-ปจฺจเยน ปจฺจโย. (สํขิตฺตํ.) อาเสวน\-ปจฺจเยน ปจฺจโย.}

\csromanheader{2}{\textbf{กมฺมาทิ}}
\proseitem{83}{ทสฺสเนน ปหาตพฺพ\-เหตุโก ธมฺโม ทสฺสเนน ปหาตพฺพ\-เหตุกสฺส ธมฺมสฺส กมฺมปจฺจเยน ปจฺจโย. ทสฺสเนน ปหาตพฺพ\-เหตุกา \csromanfolio{73}เจตนา สมฺปยุตฺตกานํ ขนฺธานํ กมฺมปจฺจเยน ปจฺจโย. (มูลํ กาตพฺพํ.) สหชาตา, นานากฺขณิกา.  สหชาตา ทสฺสเนน ปหาตพฺพ\-เหตุกา เจตนา จิตฺต\-สมุฏฺฐานานํ รูปานํ กมฺมปจฺจเยน ปจฺจโย. นานากฺขณิกา ทสฺสเนน ปหาตพฺพ\-เหตุกา เจตนา วิปากานํ ขนฺธานํ กฏตฺตา จ รูปานํ กมฺมปจฺจเยน ปจฺจโย.  ทสฺสเนน ปหาตพฺพ\-เหตุโก ธมฺโม ทสฺสเนน ปหาตพฺพ\-เหตุกสฺส จ นทสฺสเนน ปหาตพฺพ\-เหตุกสฺส จ ธมฺมสฺส กมฺมปจฺจเยน ปจฺจโย. ทสฺสเนน ปหาตพฺพ\-เหตุกา เจตนา สมฺปยุตฺตกานํ ขนฺธานํ จิตฺตสมุฏฺฐานานญฺจ รูปานํ กมฺมปจฺจเยน ปจฺจโย. (3)}

\prose{นทสฺสเนน ปหาตพฺพ\-เหตุโก ธมฺโม นทสฺสเนน ปหาตพฺพ\-เหตุกสฺส ธมฺมสฺส กมฺมปจฺจเยน ปจฺจโย. สหชาตา, นานากฺขณิกา.  สหชาตา นทสฺสเนน ปหาตพฺพ\-เหตุกา เจตนา สมฺปยุตฺตกานํ ขนฺธานํ จิตฺตสมุฏฺฐานานญฺจ รูปานํ กมฺมปจฺจเยน ปจฺจโย. ปฏิสนฺธิ\-กฺขเณ ฯเปฯ. นานากฺขณิกา นทสฺสเนน ปหาตพฺพ\-เหตุกา เจตนา วิปากานํ ขนฺธานํ กฏตฺตา จ รูปานํ กมฺมปจฺจเยน ปจฺจโย. (1)}

\prose{วิปาก\-ปจฺจเยน ปจฺจโย. อาหารปจฺจเยน ปจฺจโย. อินฺทฺริย\-ปจฺจเยน ปจฺจโย. ฌานปจฺจเยน ปจฺจโย. มคฺคปจฺจเยน ปจฺจโย. สมฺปยุตฺต\-ปจฺจเยน ปจฺจโย. ฉ. วิปฺปยุตฺต\-ปจฺจเยน ปจฺจโย. ปญฺจ.}

\prose{อตฺถิอาทิ}

\proseitem{83}{ทสฺสเนน ปหาตพฺพ\-เหตุโก ธมฺโม ทสฺสเนน ปหาตพฺพ\-เหตุกสฺส ธมฺมสฺส อตฺถิปจฺจเยน ปจฺจโย. ตีณิ.}

\prose{นทสฺสเนน ปหาตพฺพ\-เหตุโก ธมฺโม นทสฺสเนน ปหาตพฺพ\-เหตุกสฺส ธมฺมสฺส อตฺถิปจฺจเยน ปจฺจโย. สหชาตํ, ปุเรชาตํ, ปจฺฉาชาตํ, อาหารํ, อินฺทฺริยํ. (สํขิตฺตํ.).  นทสฺสเนน ปหาตพฺพ\-เหตุโก ธมฺโม ทสฺสเนน ปหาตพฺพ\-เหตุกสฺส ธมฺมสฺส อตฺถิปจฺจเยน ปจฺจโย. สหชาตํ, ปุเรชาตํ. (สํขิตฺตํ.).  นทสฺสเนน ปหาตพฺพ\-เหตุโก ธมฺโม ทสฺสเนน ปหาตพฺพ\-เหตุกสฺส จ นทสฺสเนน ปหาตพฺพ\-เหตุกสฺส จ ธมฺมสฺส อตฺถิปจฺจเยน ปจฺจโย. สหชาตํ, ปุเรชาตํ. (สํขิตฺตํ.) (3)}

\proseitem{85}{\csromanfolio{74}ทสฺสเนน ปหาตพฺพ\-เหตุโก จ นทสฺสเนน ปหาตพฺพ\-เหตุโก จ ธมฺมา ทสฺสเนน ปหาตพฺพ\-เหตุกสฺส ธมฺมสฺส อตฺถิปจฺจเยน ปจฺจโย. สหชาตํ, ปุเรชาตํ.  สหชาโต ทสฺสเนน ปหาตพฺพ\-เหตุโก เอโก ขนฺโธ จ วตฺถุ จ ติณฺณนฺนํ ขนฺธานํ อตฺถิปจฺจเยน ปจฺจโย ฯเปฯ เทฺว ขนฺธา จ ฯเปฯ. วิจิกิจฺฉา\-สหคโต เอโก ขนฺโธ จ โมโห จ ติณฺณนฺนํ ขนฺธานํ อตฺถิปจฺจเยน ปจฺจโย ฯเปฯ เทฺว ขนฺธา จ ฯเปฯ. (1)}

\prose{ทสฺสเนน ปหาตพฺพ\-เหตุโก จ นทสฺสเนน ปหาตพฺพ\-เหตุโก จ ธมฺมา นทสฺสเนน ปหาตพฺพ\-เหตุกสฺส ธมฺมสฺส อตฺถิปจฺจเยน ปจฺจโย. สหชาตํ, ปุเรชาตํ, ปจฺฉาชาตํ, อาหารํ, อินฺทฺริยํ. (สํขิตฺตํ.) (2)}

\prose{ทสฺสเนน ปหาตพฺพ\-เหตุโก จ นทสฺสเนน ปหาตพฺพ\-เหตุโก จ ธมฺมา ทสฺสเนน ปหาตพฺพ\-เหตุกสฺส จ นทสฺสเนน ปหาตพฺพ\-เหตุกสฺส จ ธมฺมสฺส อตฺถิปจฺจเยน ปจฺจโย. สหชาตํ, ปุเรชาตํ. (สํขิตฺตํ.) (3)}

\prose{นตฺถิ\-ปจฺจเยน ปจฺจโย. วิคต\-ปจฺจเยน ปจฺจโย. อวิคต\-ปจฺจเยน ปจฺจโย.}

\csromanmarkh{1. ปจฺจยานุโลม}\csromanmarkhii{2. สงฺขฺยาวาร}\csromanheadernomark{2}{1. \textbf{ปจฺจยานุโลม 2. สงฺขฺยาวาร}}
\proseitem{86}{เหตุยา ฉ, อารมฺมเณ นว, อธิปติยา ปญฺจ, อนนฺตเร นว, สมนนฺตเร นว, สหชาเต นว, อญฺญมญฺเญ ฉ, นิสฺสเย นว, อุปนิสฺสเย นว, ปุเรชาเต ตีณิ, ปจฺฉาชาเต ตีณิ, อาเสวเน นว, กมฺเม จตฺตาริ, วิปาเก เอกํ, อาหาเร จตฺตาริ, อินฺทฺริเย จตฺตาริ, ฌาเน จตฺตาริ, มคฺเค จตฺตาริ, สมฺปยุตฺเต ฉ, วิปฺปยุตฺเต ปญฺจ, อตฺถิยา นว, นตฺถิยา นว, วิคเต นว, อวิคเต นว.}

\proseitem{87}{ทสฺสเนน ปหาตพฺพ\-เหตุโก ธมฺโม ทสฺสเนน ปหาตพฺพ\-เหตุกสฺส ธมฺมสฺส อารมฺมณ\-ปจฺจเยน ปจฺจโย. สหชาต\-ปจฺจเยน ปจฺจโย. อุปนิสฺสย\-ปจฺจเยน ปจฺจโย.  ทสฺสเนน ปหาตพฺพ\-เหตุโก ธมฺโม นทสฺสเนน ปหาตพฺพ\-เหตุกสฺส ธมฺมสฺส อารมฺมณ\-ปจฺจเยน ปจฺจโย. สหชาต\-ปจฺจเยน ปจฺจโย. อุปนิสฺสย\-ปจฺจเยน ปจฺจโย. \csromanfolio{75}ปจฺฉาชาต\-ปจฺจเยน ปจฺจโย. กมฺมปจฺจเยน ปจฺจโย.  ทสฺสเนน ปหาตพฺพ\-เหตุโก ธมฺโม ทสฺสเนน ปหาตพฺพ\-เหตุกสฺส จ นทสฺสเนน ปหาตพฺพ\-เหตุกสฺส จ ธมฺมสฺส อารมฺมณ\-ปจฺจเยน ปจฺจโย. สหชาต\-ปจฺจเยน ปจฺจโย. อุปนิสฺสย\-ปจฺจเยน ปจฺจโย. (3)}

\proseitem{88}{นทสฺสเนน ปหาตพฺพ\-เหตุโก ธมฺโม นทสฺสเนน ปหาตพฺพ\-เหตุกสฺส ธมฺมสฺส อารมฺมณ\-ปจฺจเยน ปจฺจโย. สหชาต\-ปจฺจเยน ปจฺจโย. อุปนิสฺสย\-ปจฺจเยน ปจฺจโย. ปุเรชาต\-ปจฺจเยน ปจฺจโย. ปจฺฉาชาต\-ปจฺจเยน ปจฺจโย. กมฺมปจฺจเยน ปจฺจโย. อาหารปจฺจเยน ปจฺจโย. อินฺทฺริย\-ปจฺจเยน ปจฺจโย.  นทสฺสเนน ปหาตพฺพ\-เหตุโก ธมฺโม ทสฺสเนน ปหาตพฺพ\-เหตุกสฺส ธมฺมสฺส อารมฺมณ\-ปจฺจเยน ปจฺจโย. สหชาต\-ปจฺจเยน ปจฺจโย. อุปนิสฺสย\-ปจฺจเยน ปจฺจโย. ปุเรชาต\-ปจฺจเยน ปจฺจโย.  นทสฺสเนน ปหาตพฺพ\-เหตุโก ธมฺโม ทสฺสเนน ปหาตพฺพ\-เหตุกสฺส จ นทสฺสเนน ปหาตพฺพ\-เหตุกสฺส จ ธมฺมสฺส อารมฺมณ\-ปจฺจเยน ปจฺจโย. สหชาต\-ปจฺจเยน ปจฺจโย. อุปนิสฺสย\-ปจฺจเยน ปจฺจโย. ปุเรชาต\-ปจฺจเยน ปจฺจโย. (3) ทสฺสเนน ปหาตพฺพ\-เหตุโก จ นทสฺสเนน ปหาตพฺพ\-เหตุโก จ ธมฺมา ทสฺสเนน ปหาตพฺพ\-เหตุกสฺส ธมฺมสฺส อารมฺมณ\-ปจฺจเยน ปจฺจโย. สหชาต\-ปจฺจเยน ปจฺจโย. อุปนิสฺสย\-ปจฺจเยน ปจฺจโย.  ทสฺสเนน ปหาตพฺพ\-เหตุโก จ นทสฺสเนน ปหาตพฺพ\-เหตุโก จ ธมฺมา นทสฺสเนน ปหาตพฺพ\-เหตุกสฺส ธมฺมสฺส อารมฺมณ\-ปจฺจเยน ปจฺจโย. สหชาต\-ปจฺจเยน ปจฺจโย. อุปนิสฺสย\-ปจฺจเยน ปจฺจโย. ปจฺฉาชาต\-ปจฺจเยน ปจฺจโย.  ทสฺสเนน ปหาตพฺพ\-เหตุโก จ นทสฺสเนน ปหาตพฺพ\-เหตุโก จ ธมฺมา ทสฺสเนน ปหาตพฺพ\-เหตุกสฺส จ นทสฺสเนน ปหาตพฺพ\-เหตุกสฺส จ ธมฺมสฺส อรมฺมณปจฺจเยน ปจฺจโย. สหชาต\-ปจฺจเยน ปจฺจโย. อุปนิสฺสย\-ปจฺจเยน ปจฺจโย. (3)}

\csromanmarkh{2. ปจฺจย\-ปจฺจนีย}\csromanmarkhii{2. สงฺขฺยาวาร}\csromanheadernomark{2}{2. ปจฺจย\-ปจฺจนีย 2. สงฺขฺยาวาร}
\vspace{\tipitakaparskip}
\csromancenter{นเหตุยา นว, นอารมฺมเณ นว, (สพฺพตฺถ นว,) โนอวิคเต นว.}
\csromancenter{\textbf{เหตุ}ปจฺจยา นอารมฺมเณ ฉ, นอธิปติยา ฉ, นอนนฺตเร ฉ, นสมนนฺตเร ฉ, นอญฺญมญฺเญ เทฺว, นฺเอาปนิสฺสเย ฉ ฯเปฯ นฺสมฺปยุตฺเต เทฺว, นวิปฺปยุตฺเต ตีณิ, โนนตฺถิยา ฉ, โนวิคเต ฉ.}
\csromanheader{2}{4. ปจฺจย\-ปจฺจนียา\-นุโลม}
\proseitemwithcloser{92}{\csromanfolio{76}นเหตุปจฺจยา อารมฺมเณ นว, อธิปติยา ปญฺจ, (อนุโลมมาติกา.) ฯเปฯ อวิคเต นว.}{ทสฺสเนน ปหาตพฺพ\-เหตุกทุกํ นิฏฺฐิตํ.}
\csromanmatikaanchor{mtk.15}
\csromantocmarkref{section}{86. ภาวนายปหาตพฺพเหตุกทุก}{mtk.15}
\bookmark[dest={mtk.15},level=1]{86. ภาวนายปหาตพฺพเหตุกทุก}
\csromanmarkh{86. ภาวนาย\-ปหาตพฺพ\-เหตุกทุก}\csromanmarkhii{1. ปฏิจฺจวาร}\csromanheadernomark{1}{86. ภาวนาย\-ปหาตพฺพ\-เหตุกทุก 1. ปฏิจฺจวาร}
\proseitem{93}{ภาวนาย ปหาตพฺพ\-เหตุกํ ธมฺมํ ปฏิจฺจ ภาวนาย ปหาตพฺพ\-เหตุโก ธมฺโม อุปฺปชฺชติ เหตุปจฺจยา. ภาวนาย ปหาตพฺพ\-เหตุกํ เอกํ ขนฺธํ ปฏิจฺจ ตโย ขนฺธา ฯเปฯ เทฺว ขนฺเธ ฯเปฯ. (เอวํ ปฏิจฺจวาโรปิ สหชาตวาโรปิ ปจฺจยวาโรปิ นิสฺสยวาโรปิ สํสฏฺฐ\-วาโรปิ สมฺปยุตฺตวาโรปิ ทสฺสเนน\-ปหาตพฺพ\-เหตุกทุก- สทิสา. อุทฺธจฺจ\-สหคโต โมโห วิจิกิจฺฉา\-สหคต\-โมห\-ฏฺฐาเน ฐเปตพฺโพ.)}

\csromanmarkh{86. ภาวนายปหาตพฺพฺอเหตุกทุก}\csromanmarkhii{7. ปญฺหาวาร}\csromanheadernomark{2}{86. \textbf{ภาวนายปหาตพฺพฺ}อเหตุกทุก 7. ปญฺหาวาร}
\proseitem{94}{ภาวนาย ปหาตพฺพ\-เหตุโก ธมฺโม ภาวนาย ปหาตพฺพ\-เหตุกสฺส ธมฺมสฺส เหตุปจฺจเยน ปจฺจโย. ฉ ฯเปฯ. (ทสฺสเนน ปหาตพฺพ\-เหตุก\-ทุก\-สทิสา.)}

\csromanheader{2}{86. ภาวนาย\-ปหาตพฺพ\-เหตุกทุก อารมฺมณ}
\proseitem{95}{\csromanfolio{77}ภาวนาย ปหาตพฺพ\-เหตุโก ธมฺโม ภาวนาย ปหาตพฺพ\-เหตุกสฺส ธมฺมสฺส อารมฺมณ\-ปจฺจเยน ปจฺจโย. ตีณิ. (อารพฺภ. ทสฺสเนน ปหาตพฺพ\-เหตุก\-ทุก\-สทิสา.)}

\proseitem{96}{นภาวนาย ปหาตพฺพ\-เหตุโก ธมฺโม นภาวนาย ปหาตพฺพ\-เหตุกสฺส ธมฺมสฺส อารมฺมณ\-ปจฺจเยน ปจฺจโย. ทานํ ฯเปฯ สีลํ ฯเปฯ อุโปสถกมฺมํ กตฺวา ตํ ปจฺจเวกฺขติ, อสฺสาเทติ อภินนฺทติ, ตํ อารพฺภ นภาวนาย ปหาตพฺพ\-เหตุโก ราโค อุปฺปชฺชติ. ทิฏฺฐิ ฯเปฯ วิจิกิจฺฉา ฯเปฯ. นภาวนาย ปหาตพฺพ\-เหตุกํ โทมนสฺสํ อุปฺปชฺชติ. ปุพฺเพ สุจิณฺณานิ ฯเปฯ. ฌานา ฯเปฯ. อริยา มคฺคา ฯเปฯ ผลสฺส, อาวชฺชนาย อารมฺมณ\-ปจฺจเยน ปจฺจโย. อริยา นภาวนาย ปหาตพฺพ\-เหตุเก ปหีเน กิเลเส ฯเปฯ. ปุพฺเพ สมุทาจิณฺเณ ฯเปฯ. จกฺขุํ ฯเปฯ. วตฺถุํ นภาวนาย ปหาตพฺพ\-เหตุเก ขนฺเธ จ โมหญฺจ อนิจฺจโต ฯเปฯ วิปสฺสติ, อสฺสาเทติ อภินนฺทติ, ตํ อารพฺภ นภาวนาย ปหาตพฺพ\-เหตุโก ราโค อุปฺปชฺชติ. ทิฏฺฐิ ฯเปฯ. วิจิกิจฺฉา ฯเปฯ. นภาวนาย ปหาตพฺพ\-เหตุกํ โทมนสฺสํ อุปฺปชฺชติ. ทิพฺเพน จกฺขุนา รูปํ ปสฺสติ ฯเปฯ อนาคตํสญาณสฺส, อาวชฺชนาย, โมหสฺส จ อารมฺมณ\-ปจฺจเยน ปจฺจโย. (1)}

\prose{นภาวนาย ปหาตพฺพ\-เหตุโก ธมฺโม ภาวนาย ปหาตพฺพ\-เหตุกสฺส ธมฺมสฺส อารมฺมณ\-ปจฺจเยน ปจฺจโย. ทานํ ฯเปฯ. สีลํ ฯเปฯ. ฌานํ ฯเปฯ. จกฺขุํ ฯเปฯ. วตฺถุํ นภาวนาย ปหาตพฺพ\-เหตุเก ขนฺเธ จ โมหญฺจ อสฺสาเทติ อภินนฺทติ, ตํ อารพฺภ ภาวนาย ปหาตพฺพ\-เหตุโก ราโค อุปฺปชฺชติ. อุทฺธจฺจํ ฯเปฯ. ภาวนาย ปหาตพฺพ\-เหตุกํ โทมนสฺสํ อุปฺปชฺชติ. (2)}

\prose{นภาวนาย ปหาตพฺพ\-เหตุโก ธมฺโม ภาวนาย ปหาตพฺพ\-เหตุกสฺส จ นภาวนาย ปหาตพฺพ\-เหตุกสฺส จ ธมฺมสฺส อารมฺมณ\-ปจฺจเยน ปจฺจโย. จกฺขุํ ฯเปฯ. วตฺถุํ นภาวนาย ปหาตพฺพ\-เหตุเก ขนฺเธ จ โมหญฺจ อารพฺภ อุทฺธจฺจ\-สหคตา ขนฺธา จ โมโห จ อุปฺปชฺชนฺติ. (ฆฏนารมฺมณา ตีณิปิ กาตพฺพา.) (3)}

\csromanheader{2}{\textbf{อธิปตฺยาทิ}}
\proseitem{97}{\csromanfolio{78}ภาวนาย ปหาตพฺพ\-เหตุโก ธมฺโม ภาวนาย ปหาตพฺพ\-เหตุกสฺส ธมฺมสฺส อธิปติ\-ปจฺจเยน ปจฺจโย. อารมฺมณา\-ธิปติ, สหชาตา\-ธิปติ.  อารมฺมณา\-ธิปติ\csromandash{}ภาวนาย ปหาตพฺพ\-เหตุกํ ราคํ ครุํ กตฺวา อสฺสาเทติ อภินนฺทติ, ตํ ครุํ กตฺวา ภาวนาย ปหาตพฺพ\-เหตุโก ราโค อุปฺปชฺชติ. สหชาตา\-ธิปติ\csromandash{} ภาวนาย ปหาตพฺพ\-เหตุกา\-ธิปติ สมฺปยุตฺตกานํ ขนฺธานํ อธิปติ\-ปจฺจเยน ปจฺจโย. (1)}

\prose{ภาวนาย ปหาตพฺพ\-เหตุโก ธมฺโม นภาวนาย ปหาตพฺพ\-เหตุกสฺส ธมฺมสฺส อธิปติ\-ปจฺจเยน ปจฺจโย. อารมฺมณา\-ธิปติ, สหชาตา\-ธิปติ.  อารมฺมณา\-ธิปติ\csromandash{}ภาวนาย ปหาตพฺพ\-เหตุกํ ราคํ ครุํ กตฺวา อสฺสาเทติ อภินนฺทติ, ตํ ครุํ กตฺวา นภาวนาย ปหาตพฺพ\-เหตุโก ราโค อุปฺปชฺชติ. ทิฏฺฐิ อุปฺปชฺชติ. สหชาตา\-ธิปติ\csromandash{}ภาวนาย ปหาตพฺพ\-เหตุกา\-ธิปติ จิตฺต\-สมุฏฺฐานานํ รูปานํ อธิปติ\-ปจฺจเยน ปจฺจโย. (2)}

\prose{ภาวนาย ปหาตพฺพ\-เหตุโก ธมฺโม ภาวนาย ปหาตพฺพ\-เหตุกสฺส จ นภาวนาย ปหาตพฺพ\-เหตุกสฺส จ ธมฺมสฺส อธิปติ\-ปจฺจเยน ปจฺจโย. สหชาตา\-ธิปติ\csromandash{}ภาวนาย ปหาตพฺพ\-เหตุกา\-ธิปติ สมฺปยุตฺตกานํ ขนฺธานํ จิตฺตสมุฏฺฐานานญฺจ รูปานํ อธิปติ\-ปจฺจเยน ปจฺจโย. (3)}

\proseitem{98}{นภาวนาย ปหาตพฺพ\-เหตุโก ธมฺโม นภาวนาย ปหาตพฺพ\-เหตุกสฺส ธมฺมสฺส อธิปติ\-ปจฺจเยน ปจฺจโย. อารมฺมณา\-ธิปติ, สหชาตา\-ธิปติ.  อารมฺมณา\-ธิปติ\csromandash{}ทานํ ฯเปฯ สีลํ ฯเปฯ อุโปสถกมฺมํ กตฺวา ตํ ครุํ กตฺวา ปจฺจเวกฺขติ, อสฺสาเทติ อภินนฺทติ, ตํ ครุํ กตฺวา นภาวนาย ปหาตพฺพ\-เหตุโก ราโค อุปฺปชฺชติ. ทิฏฺฐิ อุปฺปชฺชติ. ปุพฺเพ ฯเปฯ. ฌานา ฯเปฯ. อริยา มคฺคา วุฏฺฐหิตฺวา ฯเปฯ ผลสฺส อธิปติ\-ปจฺจเยน ปจฺจโย. จกฺขุํ ฯเปฯ. วตฺถุํ นภาวนาย ปหาตพฺพ\-เหตุเก ขนฺเธ ครุํ กตฺวา อสฺสาเทติ อภินนฺทติ, ตํ ครุํ กตฺวา นภาวนาย ปหาตพฺพ\-เหตุโก ราโค อุปฺปชฺชติ. ทิฏฺฐิ อุปฺปชฺชติ. สหชาตา\-ธิปติ\csromandash{}นภาวนาย ปหาตพฺพ\-เหตุกา\-ธิปติ สมฺปยุตฺตกานํ ขนฺธานํ จิตฺตสมุฏฺฐานานญฺจ รูปานํ อธิปติ\-ปจฺจเยน ปจฺจโย. (1)}

\csromanheader{2}{86. ภาวนาย\-ปหาตพฺพ\-เหตุกทุก}
\prose{\csromanfolio{79}นภาวนาย ปหาตพฺพ\-เหตุโก ธมฺโม ภาวนาย ปหาตพฺพ\-เหตุกสฺส ธมฺมสฺส อธิปติ\-ปจฺจเยน ปจฺจโย. อารมฺมณา\-ธิปติ\csromandash{}ทานํ ฯเปฯ ฌานํ ฯเปฯ. จกฺขุํ ฯเปฯ. วตฺถุํ นภาวนาย ปหาตพฺพ\-เหตุเก ขนฺเธ ครุํ กตฺวา อสฺสาเทติ อภินนฺทติ, ตํ ครุํ กตฺวา ภาวนาย ปหาตพฺพ\-เหตุโก ราโค อุปฺปชฺชติ. (2)}

\prose{(อนนฺตร\-ปจฺจเย นภาวนาย ปหาตพฺพ\-เหตุก\-การณา วิจิกิจฺฉา\-สหคโต โมโห น กาตพฺโพ, อุทฺธจฺจ\-สหคโต โมโห กาตพฺโพ.) สมนนฺตร\-ปจฺจเยน ปจฺจโย. สหชาต\-ปจฺจเยน ปจฺจโย. นว. อญฺญมญฺญ\-ปจฺจเยน ปจฺจโย ฉ. นิสฺสย\-ปจฺจเยน ปจฺจโย. นว.}

\csromanheader{2}{\textbf{อุปนิสฺสย}}
\proseitem{99}{ภาวนาย ปหาตพฺพ\-เหตุโก ธมฺโม ภาวนาย ปหาตพฺพ\-เหตุกสฺส ธมฺมสฺส อุปนิสฺสย\-ปจฺจเยน ปจฺจโย. อารมฺมณู\-ปนิสฺสโย, อนนฺตรู\-ปนิสฺสโย, ปกตู\-ปนิสฺสโย ฯเปฯ. ปกตู\-ปนิสฺสโย\csromandash{}ภาวนาย ปหาตพฺพ\-เหตุกา ขนฺธา ภาวนาย ปหาตพฺพ\-เหตุกานํ ขนฺธานํ อุปนิสยปจฺจเยน ปจฺจโย. (มูลํ กาตพฺพํ.) ภาวนาย ปหาตพฺพ\-เหตุกา ขนฺธา นภาวนาย ปหาตพฺพ\-เหตุกานํ ขนฺธานํ โมหสฺส จ อุปนิสฺสย\-ปจฺจเยน ปจฺจโย. สกภณฺเฑ ฉนฺทราโค ปรภณฺเฑ ฉนฺทราคสฺส อุปนิสฺสย\-ปจฺจเยน ปจฺจโย. สกปริคฺคเห ฉนฺทราโค ปรปริคฺคเห ฉนฺทราคสฺส อุปนิสฺสย\-ปจฺจเยน ปจฺจโย. (มูลํ กาตพฺพํ.) ภาวนาย ปหาตพฺพ\-เหตุกา ขนฺธา อุทฺธจฺจ\-สหคตานํ ขนฺธานํ โมหสฺส จ อุปนิสฺสย\-ปจฺจเยน ปจฺจโย. (3)}

\proseitem{100}{นภาวนาย ปหาตพฺพ\-เหตุโก ธมฺโม นภาวนาย ปหาตพฺพ\-เหตุกสฺส ธมฺมสฺส อุปนิสฺสย\-ปจฺจเยน ปจฺจโย. อารมฺมณู\-ปนิสฺสโย, อนนฺตรู\-ปนิสฺสโย, ปกตู\-ปนิสฺสโย ฯเปฯ. ปกตู\-ปนิสฺสโย\csromandash{}สทฺธํ อุปนิสฺสาย ทานํ เทติ ฯเปฯ. สมาปตฺติํ อุปฺปาเทติ. ทิฏฺฐิํ คณฺหาติ. สีลํ ฯเปฯ. ปญฺญํ. นภาวนาย ปหาตพฺพ\-เหตุกํ ราคํ. โทสํ. โมหํ. ทิฏฺฐิํ. ปตฺถนํ. กายิกํ สุขํ. กายิกํ ทุกฺขํ ฯเปฯ. เสนาสนํ อุปนิสฺสาย ทานํ เทติ ฯเปฯ ปาณํ หนติ ฯเปฯ สํฆํ ภินฺทติ. สทฺธา ฯเปฯ. เสนาสนํ สทฺธาย ฯเปฯ ปญฺญาย. นภาวนาย ปหาตพฺพ\-เหตุกสฺส \csromanfolio{80}ราคสฺส. โทสสฺส. โมหสฺส. ทิฏฺฐิยา. ปตฺถนาย กายิกสฺส สุขสฺส ฯเปฯ ผลสมาปตฺติยา อุปนิสฺสย\-ปจฺจเยน ปจฺจโย. (1)}

\prose{นภาวนาย ปหาตพฺพ\-เหตุโก ธมฺโม ภาวนาย ปหาตพฺพ\-เหตุกสฺส ธมฺมสฺส อุปนิสฺสย\-ปจฺจเยน ปจฺจโย. อารมฺมณู\-ปนิสฺสโย, อนนฺตรู\-ปนิสฺสโย, ปกตู\-ปนิสฺสโย ฯเปฯ. ปกตู\-ปนิสฺสโย\csromandash{}สทฺธํ อุปนิสฺสาย มานํ ชปฺเปติ ฯเปฯ. สทฺธา ฯเปฯ. เสนาสนํ ภาวนาย ปหาตพฺพ\-เหตุกสฺส ราคสฺส. โทสสฺส. โมหสฺส. มานสฺส. ปตฺถนาย อุปนิสฺสย\-ปจฺจเยน ปจฺจโย. (2)}

\prose{นภาวนาย ปหาตพฺพ\-เหตุโก ธมฺโม ภาวนาย ปหาตพฺพ\-เหตุกสฺส จ นภาวนาย ปหาตพฺพ\-เหตุกสฺส จ ธมฺมสฺส อุปนิสฺสย\-ปจฺจเยน ปจฺจโย. อารมฺมณู\-ปนิสฺสโย, อนนฺตรู\-ปนิสฺสโย, ปกตู\-ปนิสฺสโย ฯเปฯ. ปกตู\-ปนิสฺสโย\csromandash{}สทฺธา ฯเปฯ ปญฺญา. กายิกํ สุขํ ฯเปฯ. เสนาสนํ อุทฺธจฺจ\-สหคตานํ ขนฺธานํ โมหสฺส จ อุปนิสฺสย\-ปจฺจเยน ปจฺจโย. (ฆฏนู\-ปนิสฺสยาปิ ตีณิปิ กาตพฺพา.) (3)}

\csromanheader{2}{\textbf{ปุเรชาตาทิ}}
\proseitem{101}{นภาวนาย ปหาตพฺพ\-เหตุโก ธมฺโม นภาวนาย ปหาตพฺพ\-เหตุกสฺส ธมฺมสฺส ปุเรชาต\-ปจฺจเยน ปจฺจโย. ตีณิ. ปจฺฉาชาต\-ปจฺจเยน ปจฺจโย. ตีณิ. อาเสวน\-ปจฺจเยน ปจฺจโย. นว. กมฺมปจฺจเยน ปจฺจโย. (นภาวนาย ปหาตพฺพ\-ภาชน\-การเณ นานากฺขณิกา ลพฺภติ.) ฯเปฯ. โนวิคต\-ปจฺจเยน ปจฺจโย. (สํขิตฺตํ. ยถา ทสฺสเนน ปหาตพฺพ\-เหตุกทุกํ, เอวํ ภาวนาย ปหาตพฺพ\-เหตุ\-กปจฺจยาปิ ปจฺจนียาปิ วิภาโคปิ คณนาปิ นินฺนานากรณา.)}

\prose{นทสฺสเนน ปหาตพฺโพ ธมฺโม นทสฺสเนน ปหาตพฺพสฺส ธมฺมสฺส ฯเปฯ. (ปรนฺเตน สกภณฺฑฉนฺทราโคปิ กาตพฺโพ.)}

\vspace{\tipitakaparskip}
\csromancenter{(ปรนฺเตน สกภณฺฑฉนฺทราโคปิ กาตพฺโพ.)}
\prose{ภาวนาย ปหาตพฺโพ ธมฺโม นภาวนาย ปหาตพฺพสฺส ธมฺมสฺส ฯเปฯ. (ปรนฺเตน “สกภณฺฑฉนฺทราโค”ติ กาตพฺพํ.)}

\vspace{\tipitakaparskip}
\csromancenter{(ปรนฺเตน “สกภณฺฑฉนฺทราโค”ติ กาตพฺพํ.)}
\nitthitam{ภาวนาย ปหาตพฺพ\-เหตุกทุกํ นิฏฺฐิตํ.}
\csromanmatikaanchor{mtk.16}
\csromantocmarkref{section}{87. สวิตกฺกทุก}{mtk.16}
\bookmark[dest={mtk.16},level=1]{87. สวิตกฺกทุก}
\csromanmarkh{87. สวิตกฺกทุก}\csromanmarkhii{87. สวิตกฺกทุก 1. ปฏิจฺจวาร}\csromanheadernomark{1}{87. สวิตกฺกทุก \textbf{87. สวิตกฺกทุก 1. ปฏิจฺจวาร}}
\proseitem{102}{\csromanfolio{81}สวิตกฺกํ ธมฺมํ ปฏิจฺจ สวิตกฺโก ธมฺโม อุปฺปชฺชติ เหตุปจฺจยา. สวิตกฺกํ เอกํ ขนฺธํ ปฏิจฺจ ตโย ขนฺธา ฯเปฯ เทฺว ขนฺเธ ฯเปฯ. ปฏิสนฺธิ\-กฺขเณ ฯเปฯ. สวิตกฺกํ ธมฺมํ ปฏิจฺจ อวิตกฺโก ธมฺโม อุปฺปชฺชติ เหตุปจฺจยา. สวิตกฺเก ขนฺเธ ปฏิจฺจ วิตกฺโก จิตฺต\-สมุฏฺฐานญฺจ รูปํ. ปฏิสนฺธิ\-กฺขเณ ฯเปฯ. สวิตกฺกํ ธมฺมํ ปฏิจฺจ สวิตกฺโก จ อวิตกฺโก จ ธมฺมา อุปฺปชฺชนฺติ เหตุปจฺจยา. สวิตกฺกํ เอกํ ขนฺธํ ปฏิจฺจ ตโย ขนฺธา วิตกฺโก จ จิตฺต\-สมุฏฺฐานญฺจ รูปํ ฯเปฯ เทฺว ขนฺเธ ฯเปฯ. ปฏิสนฺธิ\-กฺขเณ ฯเปฯ. (3)}

\proseitem{103}{อวิตกฺกํ ธมฺมํ ปฏิจฺจ อวิตกฺโก ธมฺโม อุปฺปชฺชติ เหตุปจฺจยา. อวิตกฺกํ เอกํ ขนฺธํ ปฏิจฺจ ตโย ขนฺธา จิตฺต\-สมุฏฺฐานญฺจ รูปํ ฯเปฯ เทฺว ขนฺเธ ฯเปฯ. วิตกฺกํ ปฏิจฺจ จิตฺต\-สมุฏฺฐานํ รูปํ. ปฏิสนฺธิ\-กฺขเณ อวิตกฺกํ เอกํ ขนฺธํ ปฏิจฺจ ตโย ขนฺธา กฏตฺตา จ รูปํ ฯเปฯ เทฺว ขนฺเธ ฯเปฯ. วิตกฺกํ ปฏิจฺจ กฏตฺตารูปํ. ขนฺเธ ปฏิจฺจ วตฺถุ, วตฺถุํ ปฏิจฺจ ขนฺธา. วิตกฺกํ ปฏิจฺจ วตฺถุ, วตฺถุํ ปฏิจฺจ วิตกฺโก. เอกํ มหาภูตํ ฯเปฯ. (1)}

\prose{อวิตกฺกํ ธมฺมํ ปฏิจฺจ สวิตกฺโก ธมฺโม อุปฺปชฺชติ เหตุปจฺจยา. วิตกฺกํ ปฏิจฺจ สมฺปยุตฺตกา ขนฺธา. ปฏิสนฺธิ\-กฺขเณ วิตกฺกํ ปฏิจฺจ สมฺปยุตฺตกา ขนฺธา. ปฏิสนฺธิ\-กฺขเณ วตฺถุํ ปฏิจฺจ สวิตกฺกา ขนฺธา. (2)}

\prose{อวิตกฺกํ ธมฺมํ ปฏิจฺจ สวิตกฺโก จ อวิตกฺโก จ ธมฺมา อุปฺปชฺชนฺติ เหตุปจฺจยา. วิตกฺกํ ปฏิจฺจ สมฺปยุตฺตกา ขนฺธา จิตฺต\-สมุฏฺฐานญฺจ รูปํ. วิตกฺกํ ปฏิจฺจ สวิตกฺกา ขนฺธา. มหาภูเต ปฏิจฺจ จิตฺต\-สมุฏฺฐานํ รูปํ. ปฏิสนฺธิ\-กฺขเณ วิตกฺกํ ปฏิจฺจ สมฺปยุตฺตกา ขนฺธา กฏตฺตา จ รูปํ. ปฏิสนฺธิ\-กฺขเณ วิตกฺกํ ปฏิจฺจ สวิตกฺกา ขนฺธา. มหาภูเต ปฏิจฺจ กฏตฺตารูปํ. ปฏิสนฺธิ\-กฺขเณ วตฺถุํ ปฏิจฺจ สวิตกฺกา ขนฺธา. มหาภูเต ปฏิจฺจ กฏตฺตารูปํ. ปฏิสนฺธิ\-กฺขเณ วตฺถุํ ปฏิจฺจ วิตกฺโก สมฺปยุตฺตกา จ ขนฺธา. (3)}

\proseitem{104}{สวิตกฺกญฺจ อวิตกฺกญฺจ ธมฺมํ ปฏิจฺจ สวิตกฺโก ธมฺโม อุปฺปชฺชติ เหตุปจฺจยา. สวิตกฺกํ เอกํ ขนฺธญฺจ วิตกฺกญฺจ ปฏิจฺจ ตโย ขนฺธา ฯเปฯ เทฺว \csromanfolio{82}ขนฺเธ ฯเปฯ. ปฏิสนฺธิ\-กฺขเณ สวิตกฺกํ เอกํ ขนฺธญฺจ วิตกฺกญฺจ ปฏิจฺจ ตโย ขนฺธา ฯเปฯ เทฺว ขนฺเธ ฯเปฯ. ปฏิสนฺธิ\-กฺขเณ สวิตกฺกํ เอกํ ขนฺธญฺจ วตฺถุญฺจ ปฏิจฺจ ตโย ขนฺธา ฯเปฯ เทฺว ขนฺเธ ฯเปฯ. (1)}

\prose{สวิตกฺกญฺจ อวิตกฺกญฺจ ธมฺมํ ปฏิจฺจ อวิตกฺโก ธมฺโม อุปฺปชฺชติ เหตุปจฺจยา. สวิตกฺเก ขนฺเธ จ วิตกฺกญฺจ มหาภูเต จ ปฏิจฺจ จิตฺต\-สมุฏฺฐานํ รูปํ. ปฏิสนฺธิ\-กฺขเณ สวิตกฺเก ขนฺเธ จ วิตกฺกญฺจ มหาภูเต จ ปฏิจฺจ กฏตฺตารูปํ. ปฏิสนฺธิ\-กฺขเณ สวิตกฺเก ขนฺเธ จ วตฺถุญฺจ ปฏิจฺจ วิตกฺโก. (2)}

\prose{สวิตกฺกญฺจ อวิตกฺกญฺจ ธมฺมํ ปฏิจฺจ สวิตกฺโก จ อวิตกฺโก จ ธมฺมา อุปฺปชฺชนฺติ เหตุปจฺจยา. สวิตกฺกํ เอกํ ขนฺธญฺจ วิตกฺกญฺจ ปฏิจฺจ ตโย ขนฺธา จิตฺต\-สมุฏฺฐานญฺจ รูปํ ฯเปฯ เทฺว ขนฺเธ ฯเปฯ. สวิตกฺกํ เอกํ ขนฺธญฺจ วิตกฺกญฺจ ปฏิจฺจ ตโย ขนฺธา ฯเปฯ เทฺว ขนฺเธ ฯเปฯ. สวิตกฺเก ขนฺเธ จ วิตกฺกญฺจ มหาภูเต จ ปฏิจฺจ จิตฺต\-สมุฏฺฐานํ รูปํ. ปฏิสนฺธิ\-กฺขเณ สวิตกฺกํ เอกํ ขนฺธญฺจ วิตกฺกญฺจ ปฏิจฺจ ตโย ขนฺธา กฏตฺตา จ รูปํ ฯเปฯ เทฺว ขนฺเธ ฯเปฯ. ปฏิสนฺธิ\-กฺขเณ สวิตกฺกํ เอกํ ขนฺธญฺจ วิตกฺกญฺจ ปฏิจฺจ ตโย ขนฺธา ฯเปฯ เทฺว ขนฺเธ ฯเปฯ. สวิตกฺเก ขนฺเธ จ วิตกฺกญฺจ มหาภูเต จ ปฏิจฺจ กฏตฺตารูปํ. ปฏิสนฺธิ\-กฺขเณ สวิตกฺกํ เอกํ ขนฺธญฺจ วตฺถุญฺจ ปฏิจฺจ ตโย ขนฺธา วิตกฺโก จ ฯเปฯ เทฺว ขนฺเธ ฯเปฯ. (3) (สํขิตฺตํ.)}

\csromanmarkh{1. ปจฺจยานุโลม}\csromanmarkhii{2. สงฺขฺยาวาร}\csromanheadernomark{2}{1. \textbf{ปจฺจยานุโลม 2. สงฺขฺยาวาร}}
\proseitem{105}{เหตุยา นว, อารมฺมเณ นว, อธิปติยา นว ฯเปฯ อุปนิสฺสเย นว, ปุเรชาเต ฉ, อาเสวเน ฉ, กมฺเม นว, วิปาเก นว, (สพฺพตฺถ นว.) อวิคเต นว.}
\csromansectionrule

\csromanmarkh{2. ปจฺจย\-ปจฺจนีย}\csromanmarkhii{1. วิภงฺควาร}\csromanheadernomark{2}{2. ปจฺจย\-ปจฺจนีย 1. วิภงฺควาร}
\proseitem{106}{สวิตกฺกํ ธมฺมํ ปฏิจฺจ สวิตกฺโก ธมฺโม อุปฺปชฺชติ นเหตุปจฺจยา. อเหตุกํ สวิตกฺกํ เอกํ ขนฺธํ ปฏิจฺจ ตโย ขนฺธา ฯเปฯ เทฺว ขนฺเธ ฯเปฯ. อเหตุก\-ปฏิสนฺธิ\-กฺขเณ ฯเปฯ. วิจิกิจฺฉา\-สหคเต อุทฺธจฺจ\-สหคเต ขนฺเธ \csromanfolio{83}ปฏิจฺจ วิจิกิจฺฉา\-สหคโต อุทฺธจฺจ\-สหคโต โมโห. (สวิตกฺก\-มูลกา อวเสสา เทฺว ปญฺหา กาตพฺพา. อเหตุกํ นินฺนานํ.) (3)}

\prose{อวิตกฺกํ ธมฺมํ ปฏิจฺจ อวิตกฺโก ธมฺโม อุปฺปชฺชติ นเหตุปจฺจยา. อเหตุกํ อวิตกฺกํ เอกํ ขนฺธํ ปฏิจฺจ ตโย ขนฺธา ฯเปฯ เทฺว ขนฺเธ ฯเปฯ. อเหตุกํ วิตกฺกํ ปฏิจฺจ จิตฺต\-สมุฏฺฐานํ รูปํ. อเหตุก\-ปฏิสนฺธิ\-กฺขเณ วิตกฺกํ ปฏิจฺจ กฏตฺตารูปํ. วิตกฺกํ ปฏิจฺจ วตฺถุ, วตฺถุํ ปฏิจฺจ วิตกฺโก. เอกํ มหาภูตํ ฯเปฯ. (1)}

\prose{อวิตกฺกํ ธมฺมํ ปฏิจฺจ สวิตกฺโก ธมฺโม อุปฺปชฺชติ นเหตุปจฺจยา. อเหตุกํ วิตกฺกํ ปฏิจฺจ สมฺปยุตฺตกา ขนฺธา. อเหตุก\-ปฏิสนฺธิ\-กฺขเณ วิตกฺกํ ปฏิจฺจ สมฺปยุตฺตกา ขนฺธา. ปฏิสนฺธิ\-กฺขเณ วตฺถุํ ปฏิจฺจ อเหตุกา สวิตกฺกา ขนฺธา. วิตกฺกํ ปฏิจฺจ วิจิกิจฺฉา\-สหคโต อุทฺธจฺจ\-สหคโต โมโห. (2)}

\prose{อวิตกฺกํ ธมฺมํ ปฏิจฺจ สวิตกฺโก จ อวิตกฺโก จ ธมฺมา อุปฺปชฺชนฺติ นเหตุปจฺจยา. (สํขิตฺตํ. เหตุปจฺจย\-สทิสํ. “อเหตุกนฺ”ติ นิยาเมตพฺพํ.) (3)}

\prose{สวิตกฺกญฺจ อวิตกฺกญฺจ ธมฺมํ ปฏิจฺจ สวิตกฺโก ธมฺโม อุปฺปชฺชติ นเหตุปจฺจยา. อเหตุกํ สวิตกฺกํ เอกํ ขนฺธญฺจ วิตกฺกญฺจ ปฏิจฺจ ตโย ขนฺธา ฯเปฯ เทฺว ขนฺเธ ฯเปฯ. อเหตุก\-ปฏิสนฺธิ\-กฺขเณ สวิตกฺกํ เอกํ ขนฺธญฺจ วิตกฺกญฺจ ปฏิจฺจ ตโย ขนฺธา ฯเปฯ เทฺว ขนฺเธ ฯเปฯ. อเหตุก\-ปฏิสนฺธิ\-กฺขเณ สวิตกฺกํ เอกํ ขนฺธญฺจ วตฺถุญฺจ วิตกฺกญฺจ ปฏิจฺจ ตโย ขนฺธา ฯเปฯ เทฺว ขนฺเธ ฯเปฯ. วิจิกิจฺฉา\-สหคเต อุทฺธจฺจ\-สหคเต ขนฺเธ จ วิตกฺกญฺจ ปฏิจฺจ วิจิกิจฺฉา\-สหคโต อุทฺธจฺจ\-สหคโต โมโห. (อวเสสา เทฺว ปญฺหา เหตุปจฺจย\-สทิสา นินฺนานา, “อเหตุกนฺ”ติ นิยาเมตพฺพํ.) (3)}

\csromanmarkh{2. ปจฺจย\-ปจฺจนีย}\csromanmarkhii{2. สงฺขฺยาวาร}\csromanheadernomark{2}{2. ปจฺจย\-ปจฺจนีย 2. สงฺขฺยาวาร}
\proseitem{107}{นเหตุยา นว, นอารมฺมเณ ตีณิ, นอธิปติยา นว, นอนนฺตเร ตีณิ ฯเปฯ นอุปนิสฺสเย ตีณิ, นปุเรชาเต นว, นปจฺฉาชาเต นว, นอาเสวเน นว, นกมฺเม จตฺตาริ, นวิปาเก นว, นอาหาเร เอกํ, นอินฺทฺริเย เอกํ, นฌาเน เอกํ, นมคฺเค นว, นสมฺปยุตฺเต ตีณิ, นวิปฺปยุตฺเต ฉ, โนนตฺถิยา ตีณิ, โนวิคเต ตีณิ.}

\vspace{\tipitakaparskip}
\csromancenter{\textbf{เหตุ}ปจฺจยา นอารมฺมเณ ตีณิ, นอธิปติยา นว, นอนนฺตเร ตีณิ ฯเปฯ นกมฺเม จตฺตาริ, นวิปาเก นว, นสมฺปยุตฺเต ตีณิ, นวิปฺปยุตฺเต ฉ, โนนตฺถิยา ตีณิ, โนวิคเต ตีณิ.}
\csromanheader{2}{4. ปจฺจย\-ปจฺจนียา\-นุโลม}
\proseitem{109}{\csromanfolio{84}นเหตุปจฺจยา อารมฺมเณ นว ฯเปฯ อนนฺตเร นว ฯเปฯ ปุเรชาเต ฉ, อาเสวเน ปญฺจ, กมฺเม นว ฯเปฯ มคฺเค ตีณิ, สมฺปยุตฺเต นว. (สพฺพตฺถ นว.)}

\vspace{\tipitakaparskip}
\csromancenter{(สหชาตวาโร ปฏิจฺจ\-วาร\-สทิโส.)}
\csromanmarkh{87. สวิตกฺกทุก}\csromanmarkhii{3. ปจฺจยวาร}\csromanheadernomark{2}{87. \textbf{สวิตกฺกทุก 3. ปจฺจยวาร}}
\proseitem{110}{สวิตกฺกํ ธมฺมํ ปจฺจยา สวิตกฺโก ธมฺโม อุปฺปชฺชติ เหตุปจฺจยา. ตีณิ. (ปฏิจฺจ\-วาร\-สทิสา.)}

\prose{อวิตกฺกํ ธมฺมํ ปจฺจยา อวิตกฺโก ธมฺโม อุปฺปชฺชติ เหตุปจฺจยา. อวิตกฺกํ เอกํ ขนฺธํ ปจฺจยา ตโย ขนฺธา จิตฺต\-สมุฏฺฐานญฺจ รูปํ ฯเปฯ เทฺว ขนฺเธ ฯเปฯ. วิตฺตกฺกํ ปจฺจยา จิตฺต\-สมุฏฺฐานํ รูปํ. ปฏิสนฺธิ\-กฺขเณ อวิตกฺกํ เอกํ ขนฺธํ ปจฺจยา ตโย ขนฺธา กฏตฺตา จ รูปํ ฯเปฯ เทฺว ขนฺเธ ฯเปฯ. ปฏิสนฺธิ\-กฺขเณ วิตกฺกํ ปจฺจยา กฏตฺตารูปํ. ขนฺเธ ปจฺจยา วตฺถุ, วตฺถุํ ปจฺจยา ขนฺธา. วิตกฺกํ ปจฺจยา วตฺถุ, วตฺถุํ, ปจฺจยา วิตกฺโก. เอกํ มหาภูตํ ปจฺจยา ตโย มหาภูตา ฯเปฯ. วตฺถุํ ปจฺจยา อวิตกฺกา ขนฺธา. วตฺถุํ ปจฺจยา วิตกฺโก. (1)}

\prose{อวิตกฺกํ ธมฺมํ ปจฺจยา สวิตกฺโก ธมฺโม อุปฺปชฺชติ เหตุปจฺจยา. วิตกฺกํ ปจฺจยา สมฺปยุตฺตกา ขนฺธา. วตฺถุํ ปจฺจยา สวิตกฺกา ขนฺธา. (ปฏิสนฺธิยาปิ เทฺว.) (2)}

\prose{อวิตกฺกํ ธมฺมํ ปจฺจยา สวิตกฺโก จ อวิตกฺโก จ ธมฺมา อุปฺปชฺชนฺติ เหตุปจฺจยา. วิตกฺกํ ปจฺจยา สมฺปยุตฺตกา ขนฺธา จิตฺต\-สมุฏฺฐานญฺจ รูปํ. วิตกฺกํ \csromanfolio{85}ปจฺจยา สมฺปยุตฺตกา ขนฺธา. มหาภูเต ปจฺจยา จิตฺต\-สมุฏฺฐานํ รูปํ. วตฺถุํ ปจฺจยา สวิตกฺกา ขนฺธา. มหาภูเต ปจฺจยา จิตฺต\-สมุฏฺฐานํ รูปํ. วตฺถุํ ปจฺจยา วิตกฺโก สมฺปยุตฺตกา จ ขนฺธา. ปฏิสนฺธิ\-กฺขเณ ฯเปฯ. (ปฏิสนฺธิยาปิ ปวตฺติ\-สทิสาเยว.) (3)}

\proseitem{111}{สวิตกฺกญฺจ อวิตกฺกญฺจ ธมฺมํ ปจฺจยา สวิตกฺโก ธมฺโม อุปฺปชฺชติ เหตุปจฺจยา. สวิตกฺกํ เอกํ ขนฺธญฺจ วิตกฺกญฺจ ปจฺจยา ตโย ขนฺธา ฯเปฯ เทฺว ขนฺเธ ฯเปฯ. สวิตกฺกํ เอกํ ขนฺธญฺจ วตฺถุญฺจ ปจฺจยา ตโย ขนฺธา ฯเปฯ เทฺว ขนฺเธ ฯเปฯ. (ปฏิสนฺธิ\-กฺขเณ. เทฺว กาตพฺพา.) (1)}

\prose{สวิตกฺกญฺจ อวิตกฺกญฺจ ธมฺมํ ปจฺจยา อวิตกฺโก ธมฺโม อุปฺปชฺชติ เหตุปจฺจยา. สวิตกฺเก ขนฺเธ จ วิตกฺกญฺจ ปจฺจยา จิตฺต\-สมุฏฺฐานํ รูปํ. สวิตกฺเก ขนฺเธ จ วิตกฺกญฺจ มหาภูเต จ ปจฺจยา จิตฺต\-สมุฏฺฐานํ รูปํ. สวิตกฺเก ขนฺเธ จ วตฺถุญฺจ ปจฺจยา วิตกฺโก. ปฏิสนฺธิ\-กฺขเณ ฯเปฯ. (ตีณิ. ปฏิสนฺธิยาปิ.) (2)}

\prose{สวิตกฺกญฺจ อวิตกฺกญฺจ ธมฺมํ ปจฺจยา สวิตกฺโก จ อวิตกฺโก จ ธมฺมา อุปฺปชฺชนฺติ เหตุปจฺจยา. สวิตกฺกํ เอกํ ขนฺธญฺจ วิตกฺกญฺจ ปจฺจยา ตโย ขนฺธา จิตฺต\-สมุฏฺฐานญฺจ รูปํ ฯเปฯ เทฺว ขนฺเธ ฯเปฯ. สวิตกฺกํ เอกํ ขนฺธญฺจ วิตกฺกญฺจ วตฺถุญฺจ ปจฺจยา ตโย ขนฺธา ฯเปฯ เทฺว ขนฺเธ ฯเปฯ. สวิตกฺเก ขนฺเธ จ วิตกฺกญฺจ มหาภูเต จ ปจฺจยา จิตฺต\-สมุฏฺฐานํ รูปํ. สวิตกฺกํ เอกํ ขนฺธญฺจ วตฺถุญฺจ ปจฺจยา ตโย ขนฺธา วิตกฺโก จ ฯเปฯ เทฺว ขนฺเธ ฯเปฯ. ปฏิสนฺธิ\-กฺขเณ ฯเปฯ. (3)}

\csromanmarkh{1. ปจฺจยานุโลม}\csromanmarkhii{2. สงฺขฺยาวาร}\csromanheadernomark{2}{1. \textbf{ปจฺจยานุโลม 2. สงฺขฺยาวาร}}
\vspace{\tipitakaparskip}
\csromancenter{เหตุยา นว, อารมฺมเณ นว, อธิปติยา นว, (สพฺพตฺถ นว.) อวิคเต นว.}
\csromanmarkh{2. ปจฺจย\-ปจฺจนีย}\csromanmarkhii{1. วิภงฺควาร}\csromanheadernomark{2}{2. ปจฺจย\-ปจฺจนีย 1. วิภงฺควาร}
\proseitem{113}{สวิตกฺกํ ธมฺมํ ปจฺจยา สวิตกฺโก ธมฺโม อุปฺปชฺชติ นเหตุปจฺจยา. (นว ปญฺหา กาตพฺพา. “อเหตุกา”ติ นิยาเมตพฺพา ตีณิเยว. โมโห อุทฺธริตพฺโพ, ยถา ปฏิจฺจวาเร เหตุปจฺจย\-สทิสาเยว ปญฺหา ปญฺจวิญฺญาณา อติเรกา โมโห วิตกฺกํ.)}

\csromanmarkh{2. ปจฺจย\-ปจฺจนีย}\csromanmarkhii{2. สงฺขฺยาวาร}\csromanheadernomark{2}{2. ปจฺจย\-ปจฺจนีย 2. สงฺขฺยาวาร}
\proseitem{114}{\csromanfolio{86}นเหตุยา นว, นอารมฺมเณ ตีณิ, นอธิปติยา นว, นอนนฺตเร ตีณิ ฯเปฯ นอุปนิสฺสเย ตีณิ, นปุเรชาเต นว, นปจฺฉาชาเต นว, นอาเสวเน นว, นกมฺเม จตฺตาริ, นวิปาเก นว, นอาหาเร เอกํ, นอินฺทฺริเย เอกํ, นฌาเน เอกํ, นมคฺเค นว, นสมฺปยุตฺเต ตีณิ, นวิปฺปยุตฺเต ฉ, โนนตฺถิยา ตีณิ, โนวิคเต ตีณิ.}

\vspace{\tipitakaparskip}
\csromancenter{(เอวํ อิตเร เทฺว คณนาปิ นิสฺสยวาโรปิ กาตพฺโพ.)}
\csromanmarkh{87. สวิตกฺกทุก}\csromanmarkhii{5. สํสฏฺฐวาร 1. ปจฺจยา\-นุโลมาทิ}\csromanheadernomark{2}{87. สวิตกฺกทุก 5. สํสฏฺฐวาร 1. ปจฺจยา\-นุโลมาทิ}
\proseitem{115}{สวิตกฺกํ ธมฺมํ สํสฏฺโฐ สวิตกฺโก ธมฺโม อุปฺปชฺชติ เหตุปจฺจยา. สวิตกฺกํ เอกํ ขนฺธํ สํสฏฺฐา ตโย ขนฺธา ฯเปฯ เทฺว ขนฺเธ ฯเปฯ. ปฏิสนฺธิ\-กฺขเณ ฯเปฯ. สวิตกฺกํ ธมฺมํ สํสฏฺโฐ อวิตกฺโก ธมฺโม อุปฺปชฺชติ เหตุปจฺจยา. สวิตกฺเก ขนฺเธ สํสฏฺโฐ วิตกฺโก. ปฏิสนฺธิ\-กฺขเณ ฯเปฯ. สวิตกฺกํ ธมฺมํ สํสฏฺโฐ สวิตกฺโก จ อวิตกฺโก จ ธมฺมา อุปฺปชฺชนฺติ เหตุปจฺจยา. สวิตกฺกํ เอกํ ขนฺธํ สํสฏฺฐา ตโย ขนฺธา วิตกฺโก จ ฯเปฯ เทฺว ขนฺเธ ฯเปฯ. ปฏิสนฺธิ\-กฺขเณ ฯเปฯ. (3)}

\prose{อวิตกฺกํ ธมฺมํ สํสฏฺโฐ อวิตกฺโก ธมฺโม อุปฺปชฺชติ เหตุปจฺจยา. อวิตกฺกํ เอกํ ขนฺธํ สํสฏฺฐา ตโย ขนฺธา ฯเปฯ เทฺว ขนฺเธ ฯเปฯ. ปฏิสนฺธิ\-กฺขเณ ฯเปฯ. อวิตกฺกํ ธมฺมํ สํสฏฺโฐ สวิตกฺโก ธมฺโม อุปฺปชฺชติ เหตุปจฺจยา. วิตกฺกํ สํสฏฺฐา สมฺปยุตฺตกา ขนฺธา. ปฏิสนฺธิ\-กฺขเณ ฯเปฯ. (2)}

\prose{สวิตกฺกญฺจ อวิตกฺกญฺจ ธมฺมํ สํสฏฺโฐ สวิตกฺโก ธมฺโม อุปฺปชฺชติ เหตุปจฺจยา. สวิตกฺกํ เอกํ ขนฺธญฺจ วิตกฺกญฺจ สํสฏฺฐา ตโย ขนฺธา ฯเปฯ เทฺว ขนฺเธ ฯเปฯ. ปฏิสนฺธิ\-กฺขเณ ฯเปฯ. (1) (สํขิตฺตํ.)}

\prose{\textbf{เหตุ}ยา ฉ, อารมฺมเณ ฉ, อธิปติยา ฉ, (สพฺพตฺถ ฉ.) อวิคเต ฉ.}

\vspace{\tipitakaparskip}
\csromancenter{อนุโลมํ.}
\prose{\csromanfolio{87}สวิตกฺกํ ธมฺมํ สํสฏฺโฐ สวิตกฺโก ธมฺโม อุปฺปชฺชติ นเหตุปจฺจยา. (เอวํ ฉ ปญฺหา กาตพฺพา, อนุโลมสทิสา, “อเหตุกา”ติ นิยาเมตพฺพา. ตีณิเยว, โมโห อุทฺธริตพฺโพ.)}

\prose{นเหตุยา ฉ, นอธิปติยา ฉ, นปุเรชาเต ฉ, นปจฺฉาชาเต ฉ, นอาเสวเน ฉ, นกมฺเม จตฺตาริ, นวิปาเก ฉ, นฌาเน เอกํ, นมคฺเค ฉ, นวิปฺปยุตฺเต ฉ.}

\vspace{\tipitakaparskip}
\csromancenter{ปจฺจนียํ.}
\csromancenter{(เอวํ อิตเร เทฺว คณนาปิ สมฺปยุตฺตวาโรปิ กาตพฺโพ.)}
\csromanmarkh{87. สวิตกฺกทุก}\csromanmarkhii{7. ปญฺหาวาร}\csromanheadernomark{2}{87. สวิตกฺกทุก 7. ปญฺหาวาร}
\proseitem{116}{สวิตกฺโก ธมฺโม สวิตกฺกสฺส ธมฺมสฺส เหตุปจฺจเยน ปจฺจโย. สวิตกฺกา เหตู สมฺปยุตฺตกานํ ขนฺธานํ เหตุปจฺจเยน ปจฺจโย. ปฏิสนฺธิ\-กฺขเณ ฯเปฯ. สวิตกฺโก ธมฺโม อวิตกฺกสฺส ธมฺมสฺส เหตุปจฺจเยน ปจฺจโย. สวิตกฺกา เหตู วิตกฺกสฺส จิตฺตสมุฏฺฐานานญฺจ รูปานํ เหตุปจฺจเยน ปจฺจโย. ปฏิสนฺธิ\-กฺขเณ ฯเปฯ. (มูลํ กาตพฺพํ.) สวิตกฺกา เหตู สมฺปยุตฺตกานํ ขนฺธานํ วิตกฺกสฺส จ จิตฺตสมุฏฺฐานานญฺจ รูปานํ เหตุปจฺจเยน ปจฺจโย. ปฏิสนฺธิ\-กฺขเณ ฯเปฯ. (3)}

\prose{อวิตกฺโก ธมฺโม อวิตกฺกสฺส ธมฺมสฺส เหตุปจฺจเยน ปจฺจโย. อวิตกฺกา เหตู สมฺปยุตฺตกานํ ขนฺธานํ จิตฺตสมุฏฺฐานานญฺจ รูปานํ เหตุปจฺจเยน ปจฺจโย. ปฏิสนฺธิ\-กฺขเณ ฯเปฯ. (1)}

\proseitem{117}{สวิตกฺโก ธมฺโม สวิตกฺกสฺส ธมฺมสฺส อารมฺมณ\-ปจฺจเยน ปจฺจโย. สวิตกฺเก ขนฺเธ อารพฺภ สวิตกฺกา ขนฺธา อุปฺปชฺชนฺติ. (มูลํ กาตพฺพํ.) สวิตกฺเก ขนฺเธ อารพฺภ อวิตกฺกา ขนฺธา จ วิตกฺโก จ \csromanfolio{88}อุปฺปชฺชนฺติ. (มูลํ กาตพฺพํ.) สวิตกฺเก ขนฺเธ อารพฺภ สวิตฺตกฺกา ขนฺธา จ วิตกฺโก จ อุปฺปชฺชนฺติ. (3)}

\proseitem{118}{อวิตกฺโก ธมฺโม อวิตกฺกสฺส ธมฺมสฺส อารมฺมณ\-ปจฺจเยน ปจฺจโย. อริยา อวิตกฺกา ฌานา วิฏฺฐหิตฺวา อวิตกฺกํ ฌานํ ปจฺจเวกฺขนฺติ, มคฺคา วุฏฺฐหิตฺวา มคฺคํ ปจฺจเวกฺขนฺติ, ผลา วุฏฺฐหิตฺวา ผลํ ปจฺจเวกฺขนฺติ, นิพฺพานํ ปจฺจเวกฺขนฺติ. นิพฺพานํ อวิตกฺกสฺส มคฺคสฺส, ผลสฺส, วิตกฺกสฺส จ อารมฺมณ\-ปจฺจเยน ปจฺจโย. จกฺขุํ ฯเปฯ. วตฺถุํ อวิตกฺเก ขนฺเธ จ วิตกฺกญฺจ อนิจฺจโต ฯเปฯ วิปสฺสติ, อสฺสาเทติ อภินนฺทติ, ตํ อารพฺภ วิตกฺโก อุปฺปชฺชติ. ทิพฺเพน จกฺขุนา รูปํ ปสฺสติ. ทิพฺพาย โสตธาตุยา สทฺทํ สุณาติ. เจโตปริย\-ญาเณน อวิตกฺกจิตฺตสมงฺคีสฺส จิตฺตํ ชานาติ. อากาสา\-นญฺจา\-ยตนํ ฯเปฯ. อากิญฺจญฺญา\-ยตนํ ฯเปฯ. รูปายตนํ จกฺขุ\-วิญฺญาณสฺส ฯเปฯ. โผฏฺฐพฺพา\-ยตนํ ฯเปฯ. อวิตกฺกา ขนฺธา อิทฺธิวิธ\-ญาณสฺส, เจโตปริย\-ญาณสฺส, ปุพฺเพ\-นิวาสา\-นุสฺสติ\-ญาณสฺส, ยถากมฺมูปคญาณสฺส, อนาคตํสญาณสฺส, วิตกฺกสฺส จ อารมฺมณ\-ปจฺจเยน ปจฺจโย. อวิตกฺเก ขนฺเธ จ วิตกฺกญฺจ อารพฺภ อวิตกฺกา ขนฺธา จ วิตกฺโก จ อุปฺปชฺชนฺติ. (1)}

\prose{อวิตกฺโก ธมฺโม สวิตกฺกสฺส ธมฺมสฺส อารมฺมณ\-ปจฺจเยน ปจฺจโย. อริยา อวิตกฺกา ฌานา วุฏฺฐหิตฺวา ฯเปฯ มคฺคา วุฏฺฐหิตฺวา ฯเปฯ ผลา วุฏฺฐหิตฺวา ผลํ ปจฺจเวกฺขนฺติ, นิพฺพานํ ปจฺจเวกฺขนฺติ. นิพฺพานํ โคตฺรภุสฺส, โวทานสฺส, สวิตกฺกสฺส มคฺคสฺส, ผลสฺส, อาวชฺชนาย อารมฺมณ\-ปจฺจเยน ปจฺจโย. จกฺขุํ ฯเปฯ. วตฺถุํ อวิตกฺเก ขนฺเธ จ วิตกฺกญฺจ อนิจฺจโต ฯเปฯ วิปสฺสติ, อสฺสาเทติ อภินนฺทติ, ตํ อารพฺภ ราโค อุปฺปชฺชติ ฯเปฯ โทมนสฺสํ อุปฺปชฺชติ. อวิตกฺเก ขนฺเธ จ วิตกฺกญฺจ อารพฺภ สวิตกฺกา ขนฺธา อุปฺปชฺชนฺติ. (2)}

\prose{อวิตกฺโก ธมฺโม สวิตกฺกสฺส จ อวิตกฺกสฺส จ ธมฺมสฺส อารมฺมณ\-ปจฺจเยน ปจฺจโย. อริยา อวิตกฺกา ฌานา วุฏฺฐหิตฺวา ฯเปฯ มคฺคา วุฏฺฐหิตฺวา ฯเปฯ ผลา วุฏฺฐหิตฺวา ผลํ ปจฺจเวกฺขนฺติ, นิพฺพานํ ปจฺจเวกฺขนฺติ. นิพฺพานํ โคตฺรภุสฺส, โวทานสฺส, สวิตกฺกสฺส มคฺคสฺส, ผลสฺส, อาวชฺชนาย, วิตกฺกสฺส จ อารมฺมณ\-ปจฺจเยน ปจฺจโย. จกฺขุํ ฯเปฯ. วตฺถุํ อวิตกฺเก ขนฺเธ จ วิตกฺกญฺจ อนิจฺจโต ฯเปฯ วิปสฺสติ, อสฺสาเทติ อภินนฺทติ, ตํ อารพฺภ \csromanfolio{89}สวิตกฺกา ขนฺธา จ วิตกฺโก จ อุปฺปชฺชนฺติ. อวิตกฺเก ขนฺเธ จ วิตกฺกญฺจ อารพฺภ สวิตกฺกา ขนฺธา จ วิตกฺโก จ อุปฺปชฺชนฺติ. (3)}

\proseitem{119}{สวิตกฺโก จ อวิตกฺโก จ ธมฺมา สวิตกฺกสฺส ธมฺมสฺส อารมฺมณ\-ปจฺจเยน ปจฺจโย. สวิตกฺเก ขนฺเธ จ วิตกฺกญฺจ อารพฺภ สวิตกฺกา ขนฺธา อุปฺปชฺชนฺติ. (มูลํ กาตพฺพํ.) สวิตกฺเก ขนฺเธ จ วิตกฺกญฺจ อารพฺภ อวิตกฺกา ขนฺธา จ วิตกฺโก จ อุปฺปชฺชนฺติ. (มูลํ กาตพฺพํ.) สวิตกฺเก ขนฺเธ จ วิตกฺกญฺจ อารพฺภ สวิตกฺกา ขนฺธา จ วิตกฺโก จ อุปฺปชฺชนฺติ. (3)}

\csromanheader{2}{\textbf{อธิปติ}}
\proseitem{120}{สวิตกฺโก ธมฺโม สวิตกฺกสฺส ธมฺมสฺส อธิปติ\-ปจฺจเยน ปจฺจโย. อารมฺมณา\-ธิปติ, สหชาตา\-ธิปติ.  อารมฺมณา\-ธิปติ\csromandash{}สวิตกฺเก ขนฺเธ ครุํ กตฺวา สวิตกฺกา ขนฺธา อุปฺปชฺชนฺติ. สหชาตา\-ธิปติ\csromandash{} สวิตกฺกา\-ธิปติ สมฺปยุตฺตกานํ ขนฺธานํ อธิปติ\-ปจฺจเยน ปจฺจโย. (1)}

\prose{สวิตกฺโก ธมฺโม อวิตกฺกสฺส ธมฺมสฺส อธิปติ\-ปจฺจเยน ปจฺจโย. อารมฺมณา\-ธิปติ, สหชาตา\-ธิปติ.  อารมฺมณา\-ธิปติ\csromandash{}สวิตกฺเก ขนฺเธ ครุํ กตฺวา วิตกฺโก อุปฺปชฺชติ. สหชาตา\-ธิปติ\csromandash{}สวิตกฺกา\-ธิปติ วิตกฺกสฺส จ จิตฺตสมุฏฺฐานานญฺจ รูปานํ อธิปติ\-ปจฺจเยน ปจฺจโย. (2)}

\prose{สวิตกฺโก ธมฺโม สวิตกฺกสฺส จ อวิตกฺกสฺส จ ธมฺมสฺส อธิปติ\-ปจฺจเยน ปจฺจโย. อารมฺมณา\-ธิปติ, สหชาตา\-ธิปติ.  อารมฺมณา\-ธิปติ\csromandash{}สวิตกฺเก ขนฺเธ ครุํ กตฺวา สวิตกฺกา ขนฺธา จ วิตกฺโก จ อุปฺปชฺชนฺติ. สหชาตา\-ธิปติ\csromandash{}สวิตกฺกา\-ธิปติ สมฺปยุตฺตกานํ ขนฺธานํ วิตกฺกสฺส จ จิตฺตสมุฏฺฐานานญฺจ รูปานํ อธิปติ\-ปจฺจเยน ปจฺจโย. (3)}

\proseitem{121}{อวิตกฺโก ธมฺโม อวิตกฺกสฺส ธมฺมสฺส อธิปติ\-ปจฺจเยน ปจฺจโย. อารมฺมณา\-ธิปติ, สหชาตา\-ธิปติ.  อารมฺมณา\-ธิปติ\csromandash{}อริยา อวิตกฺกา ฌานา วุฏฺฐหิตฺวา ฯเปฯ มคฺคา วุฏฺฐหิตฺวา ฯเปฯ ผลา วุฏฺฐหิตฺวา ผลํ ครุํ กตฺวา ปจฺจเวกฺขนฺติ, นิพฺพานํ ครุํ กตฺวา ปจฺจเวกฺขนฺติ. นิพฺพานํ อวิตกฺกสฺส มคฺคสฺส, ผลสฺส, วิตกฺกสฺส จ อธิปติ\-ปจฺจเยน ปจฺจโย. จกฺขุํ ฯเปฯ. วตฺถุํ อวิตกฺเก ขนฺเธ จ วิตกฺกญฺจ ครุํ กตฺวา อสฺสาเทติ อภินนฺทติ, ตํ ครุํ กตฺวา วิตกฺโก อุปฺปชฺชติ. สหชาตา\-ธิปติ\csromandash{} \csromanfolio{90}อวิตกฺกา\-ธิปติ สมฺปยุตฺตกานํ ขนฺธานํ จิตฺตสมุฏฺฐานานญฺจ รูปานํ อธิปติ\-ปจฺจเยน ปจฺจโย. (1)}

\prose{อวิตกฺโก ธมฺโม สวิตกฺกสฺส ธมฺมสฺส อธิปติ\-ปจฺจเยน ปจฺจโย. อารมฺมณา\-ธิปติ\csromandash{}อริยา อวิตกฺกา ฌานา วุฏฺฐหิตฺวา ฯเปฯ มคฺคา วุฏฺฐหิตฺวา ฯเปฯ ผลา วุฏฺฐหิตฺวา ผลํ ครุํ กตฺวา ปจฺจเวกฺขนฺติ, นิพฺพานํ ครุํ กตฺวา ปจฺจเวกฺขนฺติ. นิพฺพานํ โคตฺรภุสฺส, โวทานสฺส, สวิตกฺกสฺส มคฺคสฺส, ผลสฺส อธิปติ\-ปจฺจเยน ปจฺจโย. จกฺขุํ ฯเปฯ. วตฺถุํ อวิตกฺเก ขนฺเธ จ วิตกฺกญฺจ ครุํ กตฺวา อสฺสาเทติ อภินนฺทติ, ตํ ครุํ กตฺวา ราโค อุปฺปชฺชติ, ทิฏฺฐิ อุปฺปชฺชติ. อวิตกฺเก ขนฺเธ จ วิตกฺกญฺจ ครุํ กตฺวา สวิตกฺกา ขนฺธา อุปฺปชฺชนฺติ. (2)}

\prose{อวิตกฺโก ธมฺโม สวิตกฺกสฺส จ อวิตกฺกสฺส จ ธมฺมสฺส อธิปติ\-ปจฺจเยน ปจฺจโย. อารมฺมณา\-ธิปติ\csromandash{}อริยา อวิตกฺกา ฌานา วุฏฺฐหิตฺวา ฯเปฯ มคฺคา วุฏฺฐหิตฺวา ฯเปฯ ผลา วุฏฺฐหิตฺวา ผลํ ครุํ กตฺวา ปจฺจเวกฺขนฺติ, นิพฺพานํ ครุํ กตฺวา ปจฺจเวกฺขนฺติ. นิพฺพานํ โคตฺรภุสฺส, โวทานสฺส, สวิตกฺกสฺส มคฺคสฺส, ผลสฺส, วิตกฺกสฺส จ อธิปติ\-ปจฺจเยน ปจฺจโย. จกฺขุํ ฯเปฯ. วตฺถุํ อวิตกฺเก ขนฺเธ จ วิตกฺกญฺจ ครุํ กตฺวา อสฺสาเทติ อภินนฺทติ, ตํ ครุํ กตฺวา ราโค อุปฺปชฺชติ, ทิฏฺฐิ อุปฺปชฺชติ. อวิตกฺเก ขนฺเธ จ วิตกฺกญฺจ ครุํ กตฺวา สวิตกฺกา ขนฺธา จ วิตกฺโก จ อุปฺปชฺชนฺติ. (3)}

\prose{สวิตกฺโก จ อวิตกฺโก จ ธมฺมา สวิตกฺกสฺส ธมฺมสฺส อธิปติ\-ปจฺจเยน ปจฺจโย. อารมฺมณา\-ธิปติ\csromandash{}สวิตกฺเก ขนฺเธ จ วิตกฺกญฺจ ครุํ กตฺวา สวิตกฺกา ขนฺธา อุปฺปชฺชนฺติ. (มูลํ กาตพฺพํ.) สวิตกฺเก ขนฺเธ จ วิตกฺกญฺจ ครุํ กตฺวา วิตกฺโก อุปฺปชฺชติ. (มูลํ กาตพฺพํ.) สวิตกฺเก ขนฺเธ จ วิตกฺกญฺจ ครุํ กตฺวา สวิตกฺกา ขนฺธา จ วิตกฺโก จ อุปฺปชฺชนฺติ. (3)}

\proseitem{122}{สวิตกฺโก ธมฺโม สวิตกฺกสฺส ธมฺมสฺส อนนฺตร\-ปจฺจเยน ปจฺจโย. ปุริมา ปุริมา สวิตกฺกา ขนฺธา ปจฺฉิมานํ ปจฺฉิมานํ สวิตกฺกานํ ขนฺธานํ อนนฺตร\-ปจฺจเยน ปจฺจโย. (1)}

\prose{\csromanfolio{91}สวิตกฺโก ธมฺโม อวิตกฺกสฺส ธมฺมสฺส อนนฺตร\-ปจฺจเยน ปจฺจโย. ปุริมา ปุริมา สวิตกฺกา ขนฺธา ปจฺฉิมสฺส ปจฺฉิมสฺส วิตกฺกสฺส อนนฺตร\-ปจฺจเยน ปจฺจโย. สวิตกฺกํ จุติจิตฺตํ อวิตกฺกสฺส อุปปตฺติ\-จิตฺตสฺส อนนฺตร\-ปจฺจเยน ปจฺจโย. อาวชฺชนา ปญฺจนฺนํ วิญฺญาณานํ อนนฺตร\-ปจฺจเยน ปจฺจโย. สวิตกฺกา ขนฺธา อวิตกฺกสฺส วุฏฺฐานสฺส อนนฺตร\-ปจฺจเยน ปจฺจโย. ทุติยสฺส ฌานสฺส ปริกมฺมํ ทุติยสฺส ฌานสฺส อนนฺตร\-ปจฺจเยน ปจฺจโย. ตติยสฺส ฌานสฺส ปริกมฺมํ ฯเปฯ. เนว\-สญฺญา\-นา\-สญฺญา\-ยตนสฺส ปริกมฺมํ เนว\-สญฺญา\-นา\-สญฺญา\-ยตนสฺส ฯเปฯ. ทิพฺพสฺส จกฺขุสฺส ปริกมฺมํ ฯเปฯ. ทิพฺพาย โสตธาตุยา ปริกมฺมํ ฯเปฯ. อิทฺธิวิธ\-ญาณสฺส ปริกมฺมํ ฯเปฯ. เจโตปริย\-ญาณสฺส ปริกมฺมํ ฯเปฯ. ปุพฺเพ\-นิวาสา\-นุสฺสติ\-ญาณสฺส ปริกมฺมํ ฯเปฯ. ยถากมฺมูปคญาณสฺส ปริกมฺมํ ยถากมฺมูปคญาณสฺส ฯเปฯ. อนาคตํสญาณสฺส ปริกมฺมํ อนาคตํสญาณสฺส อนนฺตร\-ปจฺจเยน ปจฺจโย. โคตฺรภุ อวิตกฺกสฺส มคฺคสฺส. โวทานํ อวิตกฺกสฺส มคฺคสฺส. อนุโลมํ อวิตกฺกาย ผลสมาปตฺติยา อนนฺตร\-ปจฺจเยน ปจฺจโย. (2)}

\prose{สวิตกฺโก ธมฺโม สวิตกฺกสฺส จ อวิตกฺกสฺส จ ธมฺมสฺส อนนฺตร\-ปจฺจเยน ปจฺจโย. ปุริมา ปุริมา สวิตกฺกา ขนฺธา ปจฺฉิมานํ ปจฺฉิมานํ สวิตกฺกานํ ขนฺธานํ วิตกฺกสฺส จ อนนฺตร\-ปจฺจเยน ปจฺจโย. (3)}

\proseitem{123}{อวิตกฺโก ธมฺโม อวิตกฺกสฺส ธมฺมสฺส อนนฺตร\-ปจฺจเยน ปจฺจโย. ปุริโม ปุริโม วิตกฺโก ปจฺฉิมสฺส ปจฺฉิมสฺส วิตกฺกสฺส อนนฺตร\-ปจฺจเยน ปจฺจโย. ปุริมา ปุริมา อวิตกฺกา ขนฺธา ปจฺฉิมานํ ปจฺฉิมานํ อวิตกฺกานํ ขนฺธานํ อนนฺตร\-ปจฺจเยน ปจฺจโย. อวิตกฺโก มคฺโค อวิตกฺกสฺส ผลสฺส ฯเปฯ. อวิตกฺกํ ผลํ อวิตกฺกสฺส ผลสฺส ฯเปฯ. นิโรธา วุฏฺฐหนฺตสฺส เนว\-สญฺญา\-นา\-สญฺญา\-ยตนํ อวิตกฺกาย ผลสมาปตฺติยา อนนฺตร\-ปจฺจเยน ปจฺจโย. (1)}

\prose{อวิตกฺโก ธมฺโม สวิตกฺกสฺส ธมฺมสฺส อนนฺตร\-ปจฺจเยน ปจฺจโย. ปุริโม ปุริโม วิตกฺโก ปจฺฉิมานํ ปจฺฉิมานํ สวิตกฺกานํ ขนฺธานํ อนนฺตร\-ปจฺจเยน ปจฺจโย. อวิตกฺกํ จุติจิตฺตํ สวิตกฺกสฺส อุปปตฺติ\-จิตฺตสฺส. อวิตกฺกํ ภวงฺคํ อาวชฺชนาย. อวิตกฺกา ขนฺธา สวิตกฺกสฺส วุฏฺฐานสฺส อนนฺตร\-ปจฺจเยน ปจฺจโย. (2)}

\prose{\csromanfolio{92}อวิตกฺโก ธมฺโม สวิตกฺกสฺส จ อวิตกฺกสฺส จ ธมฺมสฺส อนนฺตร\-ปจฺจเยน ปจฺจโย. ปุริโม ปุริโม วิตกฺโก ปจฺฉิมานํ ปจฺฉิมานํ สวิตกฺกานํ ขนฺธานํ วิตกฺกสฺส จ อนนฺตร\-ปจฺจเยน ปจฺจโย. (3)}

\proseitem{124}{สวิตกฺโก จ อวิตกฺโก จ ธมฺมา สวิตกฺกสฺส ธมฺมสฺส อนนฺตร\-ปจฺจเยน ปจฺจโย. ปุริมา ปุริมา สวิตกฺกา ขนฺธา จ วิตกฺโก จ ปจฺฉิมานํ ปจฺฉิมานํ สวิตกฺกานํ ขนฺธานํ อนนฺตร\-ปจฺจเยน ปจฺจโย. (1)}

\prose{สวิตกฺโก จ อวิตกฺโก จ ธมฺมา อวิตกฺกสฺส ธมฺมสฺส อนนฺตร\-ปจฺจเยน ปจฺจโย. ปุริมา ปุริมา สวิตกฺกา ขนฺธา จ วิตกฺโก จ ปจฺฉิมสฺส ปจฺฉิมสฺส วิตกฺกสฺส อนนฺตร\-ปจฺจเยน ปจฺจโย. สวิตกฺกํ จุติจิตฺตญฺจ วิตกฺโก จ อวิตกฺกสฺส อุปปตฺติ\-จิตฺตสฺส ฯเปฯ. อาวชฺชนา จ วิตกฺโก จ ปญฺจนฺนํ วิญฺญาณานํ ฯเปฯ. สวิตกฺกา ขนฺธา จ วิตกฺโก จ อวิตกฺกสฺส วุฏฺฐานสฺส อนนฺตร\-ปจฺจเยน ปจฺจโย. ทุติยสฺส ฌานสฺส ปริกมฺมญฺจ วิตกฺโก จ ฯเปฯ. (เหฏฺฐา ลิขิตํ เลขํ อิมินา การเณน ทฏฺฐพฺพํ.) อนุโลมญฺจ วิตกฺโก จ อวิตกฺกาย ผลสมาปตฺติยา อนนฺตร\-ปจฺจเยน ปจฺจโย. (2)}

\prose{สวิตกฺโก จ อวิตกฺโก จ ธมฺมา สวิตกฺกสฺส จ อวิตกฺกสฺส จ ธมฺมสฺส อนนฺตร\-ปจฺจเยน ปจฺจโย. ปุริมา ปุริมา สวิตกฺกา ขนฺธา จ วิตกฺโก จ ปจฺฉิมานํ ปจฺฉิมานํ สวิตกฺกานํ ขนฺธานํ วิตกฺกสฺส จ อนนฺตร\-ปจฺจเยน ปจฺจโย. (3)}

\prose{สมนนฺตร\-ปจฺจเยน ปจฺจโย. สหชาต\-ปจฺจเยน ปจฺจโย. นว. อญฺญมญฺญ\-ปจฺจเยน ปจฺจโย. นว. นิสฺสย\-ปจฺจเยน ปจฺจโย. นว.}

\csromanheader{2}{\textbf{อุปนิสฺสย}}
\proseitem{125}{สวิตกฺโก ธมฺโม สวิตกฺกสฺส ธมฺมสฺส อุปนิสฺสย\-ปจฺจเยน ปจฺจโย. อารมฺมณู\-ปนิสฺสโย, อนนฺตรู\-ปนิสฺสโย, ปกตู\-ปนิสฺสโย ฯเปฯ. ปกตู\-ปนิสฺสโย\csromandash{}สวิตกฺกา ขนฺธา สวิตกฺกานํ ขนฺธานํ อุปนิสฺสย\-ปจฺจเยน ปจฺจโย. (มูลํ กาตพฺพํ.) สวิตกฺกา ขนฺธา อวิตกฺกานํ ขนฺธานํ วิตกฺกสฺส จ อุปนิสฺสย\-ปจฺจเยน ปจฺจโย. (มูลํ กาตพฺพํ.) สวิตกฺกา ขนฺธา สวิตกฺกานํ ขนฺธานํ วิตกฺกสฺส จ อุปนิสฺสย\-ปจฺจเยน ปจฺจโย. (3)}

\proseitem{126}{\csromanfolio{93}อวิตกฺโก ธมฺโม อวิตกฺกสฺส ธมฺมสฺส อุปนิสฺสย\-ปจฺจเยน ปจฺจโย. อารมฺมณู\-ปนิสฺสโย, อนนฺตรู\-ปนิสฺสโย, ปกตู\-ปนิสฺสโย ฯเปฯ. ปกตู\-ปนิสฺสโย\csromandash{}อวิตกฺกํ สทฺธํ อุปนิสฺสาย อวิตกฺกํ ฌานํ อุปฺปาเทติ, มคฺคํ อุปฺปาเทติ, อภิญฺญํ อุปฺปาเทติ, สมาปตฺติํ อุปฺปาเทติ. อวิตกฺกํ สีลํ ฯเปฯ ปญฺญํ. กายิกํ สุขํ. กายิกํ ทุกฺขํ. อุตุํ. โภชนํ. เสนาสนํ วิตกฺกํ อุปนิสฺสาย อวิตกฺกํ ฌานํ อุปฺปาเทติ. มคฺคํ ฯเปฯ. อภิญฺญํ ฯเปฯ สมาปตฺติํ อุปฺปาเทติ. อวิตกฺกา สทฺธา ฯเปฯ. เสนาสนํ วิตกฺโก จ อวิตกฺกาย สทฺธาย ฯเปฯ ปญฺญาย. กายิกสฺส สุขสฺส. กายิกสฺส ทุกฺขสฺส. อวิตกฺกสฺส มคฺคสฺส, ผลสมาปตฺติยา อุปนิสฺสย\-ปจฺจเยน ปจฺจโย. (1)}

\prose{อวิตกฺโก ธมฺโม สวิตกฺกสฺส ธมฺมสฺส อุปนิสฺสย\-ปจฺจเยน ปจฺจโย. (ตีณิปิ อุปนิสฺสยา สพฺพตฺถ กาตพฺพา.) อวิตกฺกํ สทฺธํ อุปนิสฺสาย ทานํ เทติ, สีลํ สมาทิยติ, อุโปสถกมฺมํ กโรติ. สวิตกฺกํ ฌานํ อุปฺปาเทติ, วิปสฺสนํ ฯเปฯ มคฺคํ ฯเปฯ สมาปตฺติํ อุปฺปาเทติ, มานํ ชปฺเปติ, ทิฏฺฐิํ คณฺหาติ. อวิตกฺกํ สีลํ ฯเปฯ. เสนาสนํ วิตกฺกํ อุปนิสฺสาย ทานํ เทติ ฯเปฯ สมาปตฺติํ อุปฺปาเทติ, ปาณํ หนติ ฯเปฯ สํฆํ ภินฺทติ. อวิตกฺกา สทฺธา ฯเปฯ. เสนาสนํ วิตกฺโก จ สวิตกฺกาย สทฺธาย ฯเปฯ ปญฺญาย. ราคสฺส ฯเปฯ. ปตฺถนาย สวิตกฺกสฺส มคฺคสฺส, ผลสมาปตฺติยา อุปนิสฺสย\-ปจฺจเยน ปจฺจโย. (2)}

\prose{อวิตกฺโก ธมฺโม สวิตกฺกสฺส จ อวิตกฺกสฺส จ ธมฺมสฺส อุปนิสฺสย\-ปจฺจเยน ปจฺจโย. อวิตกฺกํ สทฺธํ อุปนิสฺสาย ทานํ เทติ ฯเปฯ. (ทุติยวาเร ลิขิตปทา สพฺเพ กาตพฺพา.) สมาปตฺติํ อุปฺปาเทติ, มานํ ชปฺเปติ, ทิฏฺฐิํ คณฺหาติ. สีลํ ฯเปฯ. ปญฺญํ ฯเปฯ. เสนาสนํ วิตกฺกํ อุปนิสฺสาย ทานํ เทติ ฯเปฯ ปาณํ หนติ ฯเปฯ สํฆํ ภินฺทติ. อวิตกฺกา สทฺธา ฯเปฯ. เสนาสนํ วิตกฺโก จ สวิตกฺกาย สทฺธาย ฯเปฯ ปญฺญาย. ราคสฺส ฯเปฯ. ปตฺถนาย สวิตกฺกสฺส มคฺคสฺส, ผลสมาปตฺติยา วิตกฺกสฺส จ อุปนิสฺสย\-ปจฺจเยน ปจฺจโย. (3)}

\proseitem{127}{สวิตกฺโก จ อวิตกฺโก จ ธมฺมา สวิตกฺกสฺส ธมฺมสฺส อุปนิสฺสย\-ปจฺจเยน ปจฺจโย. สวิตกฺกา ขนฺธา จ วิตกฺโก จ สวิตกฺกานํ ขนฺธานํ อุปนิสฺสย\-ปจฺจเยน ปจฺจโย. (มูลํ กาตพฺพํ.) สวิตกฺกา ขนฺธา จ วิตกฺโก จ อวิตกฺกานํ ขนฺธานํ วิตกฺกสฺส จ อุปนิสฺสย\-ปจฺจเยน ปจฺจโย. \csromanfolio{94}(มูลํ กาตพฺพํ.) สวิตกฺกา ขนฺธา จ วิตกฺโก จ สวิตกฺกานํ ขนฺธานํ วิตกฺกสฺส จ อุปนิสฺสย\-ปจฺจเยน ปจฺจโย. (3)}

\csromanheader{2}{\textbf{ปุเรชาต}}
\proseitem{128}{อวิตกฺโก ธมฺโม อวิตกฺกสฺส ธมฺมสฺส ปุเรชาต\-ปจฺจเยน ปจฺจโย. อารมฺมณ\-ปุเรชาตํ, วตฺถุ\-ปุเรชาตํ.  อารมฺมณ\-ปุเรชาตํ\csromandash{} จกฺขุํ ฯเปฯ. วตฺถุํ อนิจฺจโต ฯเปฯ วิปสฺสติ, อสฺสาเทติ อภินนฺทติ, ตํ อารพฺภ วิตกฺโก อุปฺปชฺชติ. ทิพฺเพน จกฺขุนา รูปํ ปสฺสติ ฯเปฯ. โผฏฺฐพฺพา\-ยตนํ กายวิญฺญาณสฺส ฯเปฯ. วตฺถุ\-ปุเรชาตํ\csromandash{}จกฺขายตนํ จกฺขุ\-วิญฺญาณสฺส ฯเปฯ. กายายตนํ กายวิญฺญาณสฺส ฯเปฯ. วตฺถุ อวิตกฺกานํ ขนฺธานํ วิตกฺกสฺส จ ปุเรชาต\-ปจฺจเยน ปจฺจโย. (1)}

\prose{อวิตกฺโก ธมฺโม สวิตกฺกสฺส ธมฺมสฺส ปุเรชาต\-ปจฺจเยน ปจฺจโย. อารมฺมณ\-ปุเรชาตํ, วตฺถุ\-ปุเรชาตํ.  อารมฺมณ\-ปุเรชาตํ\csromandash{}จกฺขุํ ฯเปฯ. วตฺถุํ อนิจฺจโต ฯเปฯ วิปสฺสติ, อสฺสาเทติ อภินนฺทติ, ตํ อารพฺภ สวิตกฺกา ขนฺธา อุปฺปชฺชนฺติ. วตฺถุ\-ปุเรชาตํ\csromandash{}วตฺถุ สวิตกฺกานํ ขนฺธานํ ปุเรชาต\-ปจฺจเยน ปจฺจโย. (2)}

\prose{อวิตกฺโก ธมฺโม สวิตกฺกสฺส จ อวิตกฺกสฺส จ ธมฺมสฺส ปุเรชาต\-ปจฺจเยน ปจฺจโย. อารมฺมณ\-ปุเรชาตํ, วตฺถุ\-ปุเรชาตํ.  อารมฺมณ\-ปุเรชาตํ\csromandash{}จกฺขุํ ฯเปฯ. วตฺถุํ อนิจฺจโต ฯเปฯ วิปสฺสติ, อสฺสาเทติ อภินนฺทติ, ตํ อารพฺภ สวิตกฺกา ขนฺธา จ วิตกฺโก จ อุปฺปชฺชนฺติ. วตฺถุ\-ปุเรชาตํ\csromandash{}วตฺถุ สวิตกฺกานํ ขนฺธานํ วิตกฺกสฺส จ ปุเรชาต\-ปจฺจเยน ปจฺจโย. (3)}

\csromanheader{2}{ปจฺฉาชาตา\-เสวน}
\proseitem{129}{สวิตกฺโก ธมฺโม อวิตกฺกสฺส ธมฺมสฺส ปจฺฉาชาต\-ปจฺจเยน ปจฺจโย. (ตีณิ, ปจฺฉาชาตา.) อาเสวน\-ปจฺจเยน ปจฺจโย. นว.}

\csromanheader{2}{\textbf{กมฺมาทิ}}
\proseitem{130}{สวิตกฺโก ธมฺโม สวิตกฺกสฺส ธมฺมสฺส กมฺมปจฺจเยน ปจฺจโย. สหชาตา, นานากฺขณิกา.  สหชาตา สวิตกฺกา เจตนา สมฺปยุตฺตกานํ ขนฺธานํ กมฺมปจฺจเยน ปจฺจโย. ปฏิสนฺธิ\-กฺขเณ ฯเปฯ. นานากฺขณิกา สวิตกฺกา เจตนา วิปากานํ สวิตกฺกานํ ขนฺธานํ กมฺมปจฺจเยน ปจฺจโย.}

\vspace{\tipitakaparskip}
\csromancenter{(เอวํ จตฺตาริ, สหชาตาปิ นานากฺขณิกาปิ กาตพฺพา.)}
\prose{\csromanfolio{95}วิปาก\-ปจฺจเยน ปจฺจโย. นว. อาหารปจฺจเยน ปจฺจโย. จตฺตาริ. อินฺทฺริย\-ปจฺจเยน ปจฺจโย. จตฺตาริ. ฌานปจฺจเยน ปจฺจโย. นว. มคฺคปจฺจเยน ปจฺจโย. นว. สมฺปยุตฺต\-ปจฺจเยน ปจฺจโย. ฉ.}

\proseitem{131}{สวิตกฺโก ธมฺโม สวิตกฺกสฺส ธมฺมสฺส วิปฺปยุตฺต\-ปจฺจเยน ปจฺจโย. สหชาตํ, ปจฺฉาชาตํ. (สํขิตฺตํ.) (1)}

\prose{อวิตกฺโก ธมฺโม อวิตกฺกสฺส ธมฺมสฺส วิปฺปยุตฺต\-ปจฺจเยน ปจฺจโย. สหชาตํ, ปุเรชาตํ, ปจฺฉาชาตํ. (สํขิตฺตํ.). อวิตกฺโก ธมฺโม สวิตกฺกสฺส ธมฺมสฺส วิปฺปยุตฺต\-ปจฺจเยน ปจฺจโย. สหชาตํ, ปุเรชาตํ. (สํขิตฺตํ.). อวิตกฺโก ธมฺโม สวิตกฺกสฺส จ อวิตกฺกสฺส จ ธมฺมสฺส วิปฺปยุตฺต\-ปจฺจเยน ปจฺจโย. สหชาตํ, ปุเรชาตํ.  สหชาตํ ปฏิสนฺธิ\-กฺขเณ วตฺถุ วิตกฺกสฺส สมฺปยุตฺตกานญฺจ ขนฺธานํ วิปฺปยุตฺต\-ปจฺจเยน ปจฺจโย. ปุเรชาตํ วตฺถุ วิตกฺกสฺส สมฺปยุตฺตกานญฺจ ขนฺธานํ วิปฺปยุตฺต\-ปจฺจเยน ปจฺจโย. (3)}

\prose{สวิตกฺโก จ อวิตกฺโก จ ธมฺมา อวิตกฺกสฺส ธมฺมสฺส วิปฺปยุตฺต\-ปจฺจเยน ปจฺจโย. สหชาตํ, ปจฺฉาชาตํ. (สํขิตฺตํ.)}

\prose{อตฺถิอาทิ}

\proseitem{132}{สวิตกฺโก ธมฺโม สวิตกฺกสฺส ธมฺมสฺส อตฺถิปจฺจเยน ปจฺจโย. (เอกํ, ปฏิจฺจสทิสํ.). สวิตกฺโก ธมฺโม อวิตกฺกสฺส ธมฺมสฺส อตฺถิปจฺจเยน ปจฺจโย. สหชาตํ, ปจฺฉาชาตํ. (สํขิตฺตํ.). สวิตกฺโก ธมฺโม สวิตกฺกสฺส จ อวิตกฺกสฺส จ ธมฺมสฺส อตฺถิปจฺจเยน ปจฺจโย. (ปฏิจฺจสทิสํ.) (3)}

\prose{อวิตกฺโก ธมฺโม อวิตกฺกสฺส ธมฺมสฺส อตฺถิปจฺจเยน ปจฺจโย. สหชาตํ, ปุเรชาตํ, ปจฺฉาชาตํ, อาหารํ, อินฺทฺริยํ. (สํขิตฺตํ.). อวิตกฺโก ธมฺโม สวิตกฺกสฺส ธมฺมสฺส อตฺถิปจฺจเยน ปจฺจโย. สหชาตํ, ปุเรชาตํ. (สํขิตฺตํ.). อวิตกฺโก ธมฺโม สวิตกฺกสฺส จ อวิตกฺกสฺส จ ธมฺมสฺส อตฺถิปจฺจเยน ปจฺจโย. สหชาตํ, ปุเรชาตํ.  สหชาโต วิตกฺโก สมฺปยุตฺตกานํ ขนฺธานํ จิตฺตสมุฏฺฐานานญฺจ รูปานํ อตฺถิปจฺจเยน ปจฺจโย. ปฏิสนฺธิ\-กฺขเณ วิตกฺโก สมฺปยุตฺตกานํ ขนฺธานํ กฏตฺตา \csromanfolio{96}จ รูปานํ อตฺถิปจฺจเยน ปจฺจโย. ปฏิสนฺธิ\-กฺขเณ วตฺถุ วิตกฺกสฺส สมฺปยุตฺตกานญฺจ ขนฺธานํ อตฺถิปจฺจเยน ปจฺจโย. ปุเรชาตํ จกฺขุํ ฯเปฯ. วตฺถุํ อนิจฺจโต ฯเปฯ วิปสฺสติ, อสฺสาเทติ อภินนฺทติ, ตํ อารพฺภ วิตกฺโก จ สมฺปยุตฺตกา จ ขนฺธา อุปฺปชฺชนฺติ. วตฺถุ วิตกฺกสฺส สมฺปยุตฺตกานญฺจ ขนฺธานํ อตฺถิปจฺจเยน ปจฺจโย. (3)}

\proseitem{133}{สวิตกฺโก จ อวิตกฺโก จ ธมฺมา สวิตกฺกสฺส ธมฺมสฺส อตฺถิปจฺจเยน ปจฺจโย. สหชาตํ, ปุเรชาตํ.  สหชาโต สวิตกฺโก เอโก ขนฺโธ จ วิตกฺโก จ ติณฺณนฺนํ ขนฺธานํ อตฺถิปจฺจเยน ปจฺจโย ฯเปฯ เทฺว ขนฺธา จ ฯเปฯ. สหชาโต สวิตกฺโก เอโก ขนฺโธ จ วตฺถุ จ ติณฺณนฺนํ ขนฺธานํ อตฺถิปจฺจเยน ปจฺจโย ฯเปฯ เทฺว ขนฺธา จ ฯเปฯ. (ปฏิสนฺธิ\-กฺขเณ สหชาตาปิ เทฺวปิ กาตพฺพา.) (1)}

\prose{สวิตกฺโก จ อวิตกฺโก จ ธมฺมา อวิตกฺกสฺส ธมฺมสฺส อตฺถิปจฺจเยน ปจฺจโย. สหชาตํ, ปุเรชาตํ, ปจฺฉาชาตํ, อาหารํ, อินฺทฺริยํ.  สหชาตา สวิตกฺกา ขนฺธา จ วิตกฺโก จ จิตฺต\-สมุฏฺฐานานํ รูปานํ อตฺถิปจฺจเยน ปจฺจโย. สหชาตา สวิตกฺกา ขนฺธา จ วตฺถุ จ วิตกฺกสฺส อตฺถิปจฺจเยน ปจฺจโย. (ปฏิสนฺธิ\-กฺขเณ. ตีณิ.) ปจฺฉาชาตา สวิตกฺกา ขนฺธา จ วิตกฺโก จ ปุเรชาตสฺส อิมสฺส กายสฺส อตฺถิปจฺจเยน ปจฺจโย. ปจฺฉาชาตา สวิตกฺกา ขนฺธา จ วิตกฺโก จ กพฬีกาโร อาหาโร จ อิมสฺส กายสฺส อตฺถิปจฺจเยน ปจฺจโย. ปจฺฉาชาตา สวิตกฺกา ขนฺธา จ วิตกฺโก จ รูปชีวิตินฺทฺริยญฺจ กฏตฺตารูปานํ อตฺถิปจฺจเยน ปจฺจโย. (2)}

\prose{สวิตกฺโก จ อวิตกฺโก จ ธมฺมา สวิตกฺกสฺส จ อวิตกฺกสฺส จ ธมฺมสฺส อตฺถิปจฺจเยน ปจฺจโย. สหชาตํ, ปุเรชาตํ.  \textbf{สหชาโต} สวิตกฺโก เอโก ขนฺโธ จ วิตกฺโก จ ติณฺณนฺนํ ขนฺธานํ จิตฺตสมุฏฺฐานานญฺจ รูปานํ อตฺถิปจฺจเยน ปจฺจโย ฯเปฯ เทฺว ขนฺธา จ ฯเปฯ. \textbf{สหชาโต} สวิตกฺโก เอโก ขนฺโธ จ วตฺถุ จ ติณฺณนฺนํ ขนฺธานํ วิตกฺกสฺส จ อตฺถิปจฺจเยน ปจฺจโย ฯเปฯ เทฺว ขนฺธา จ ฯเปฯ. (ปฏิสนฺธิยาปิ เทฺว.) (3)}

\prose{นตฺถิ\-ปจฺจเยน ปจฺจโย. วิคต\-ปจฺจเยน ปจฺจโย. อวิคต\-ปจฺจเยน ปจฺจโย.}

\csromanmarkh{87. สวิตกฺกทุก}\csromanmarkhii{1. ปจฺจยานุโลม 2. สงฺขฺยาวาร}\csromanheadernomark{2}{87. สวิตกฺกทุก \textbf{1. ปจฺจยานุโลม 2. สงฺขฺยาวาร}}
\proseitem{134}{\csromanfolio{97}เหตุยา จตฺตาริ, อารมฺมเณ นว, อธิปติยา นว, อนนฺตเร นว, สมนนฺตเร นว, สหชาเต นว, อญฺญมญฺเญ นว, นิสฺสเย นว, อุปนิสฺสเย นว, ปุเรชาเต ตีณิ, ปจฺฉาชาเต ตีณิ, อาเสวเน นว, กมฺเม จตฺตาริ, วิปาเก นว, อาหาเร จตฺตาริ, อินฺทฺริเย จตฺตาริ, ฌาเน นว, มคฺเค นว, สมฺปยุตฺเต ฉ, วิปฺปยุตฺเต ปญฺจ, อตฺถิยา นว, นตฺถิยา นว, วิคเต นว, อวิคเต นว.}

\vspace{\tipitakaparskip}
\csromancenter{อนุโลมํ.}
\proseitem{135}{สวิตกฺโก ธมฺโม สวิตกฺกสฺส ธมฺมสฺส อารมฺมณ\-ปจฺจเยน ปจฺจโย. สหชาต\-ปจฺจเยน ปจฺจโย. อุปนิสฺสย\-ปจฺจเยน ปจฺจโย. กมฺมปจฺจเยน ปจฺจโย.  สวิตกฺโก ธมฺโม อวิตกฺกสฺส ธมฺมสฺส อารมฺมณ\-ปจฺจเยน ปจฺจโย. สหชาต\-ปจฺจเยน ปจฺจโย. อุปนิสฺสย\-ปจฺจเยน ปจฺจโย. ปจฺฉาชาต\-ปจฺจเยน ปจฺจโย. กมฺมปจฺจเยน ปจฺจโย.  สวิตกฺโก ธมฺโม สวิตกฺกสฺส จ อวิตกฺกสฺส จ ธมฺมสฺส อารมฺมณ\-ปจฺจเยน ปจฺจโย. สหชาต\-ปจฺจเยน ปจฺจโย. อุปนิสฺสย\-ปจฺจเยน ปจฺจโย. กมฺมปจฺจเยน ปจฺจโย. (3)}

\prose{อวิตกฺโก ธมฺโม อวิตกฺกสฺส ธมฺมสฺส อารมฺมณ\-ปจฺจเยน ปจฺจโย. สหชาต\-ปจฺจเยน ปจฺจโย. อุปนิสฺสย\-ปจฺจเยน ปจฺจโย. ปุเรชาต\-ปจฺจเยน ปจฺจโย. ปจฺฉาชาต\-ปจฺจเยน ปจฺจโย. กมฺมปจฺจเยน ปจฺจโย. อาหารปจฺจเยน ปจฺจโย. อินฺทฺริย\-ปจฺจเยน ปจฺจโย.  อวิตกฺโก ธมฺโม สวิตกฺกสฺส ธมฺมสฺส อารมฺมณ\-ปจฺจเยน ปจฺจโย. สหชาต\-ปจฺจเยน ปจฺจโย. อุปนิสฺสย\-ปจฺจเยน ปจฺจโย. ปุเรชาต\-ปจฺจเยน ปจฺจโย.  อวิตกฺโก ธมฺโม สวิตกฺกสฺส จ อวิตกฺกสฺส จ ธมฺมสฺส อารมฺมณ\-ปจฺจเยน ปจฺจโย. สหชาต\-ปจฺจเยน ปจฺจโย. อุปนิสฺสย\-ปจฺจเยน ปจฺจโย. ปุเรชาต\-ปจฺจเยน ปจฺจโย. (3)}

\proseitem{137}{\csromanfolio{98}สวิตกฺโก จ อวิตกฺโก จ ธมฺมา สวิตกฺกสฺส ธมฺมสฺส อารมฺมณ\-ปจฺจเยน ปจฺจโย. สหชาต\-ปจฺจเยน ปจฺจโย. อุปนิสฺสย\-ปจฺจเยน ปจฺจโย.  สวิตกฺโก จ อวิตกฺโก จ ธมฺมา อวิตกฺกสฺส ธมฺมสฺส อารมฺมณ\-ปจฺจเยน ปจฺจโย. สหชาต\-ปจฺจเยน ปจฺจโย. อุปนิสฺสย\-ปจฺจเยน ปจฺจโย. ปจฺฉาชาต\-ปจฺจเยน ปจฺจโย.  สวิตกฺโก จ อวิตกฺโก จ ธมฺมา สวิตกฺกสฺส จ อวิตกฺกสฺส จ ธมฺมสฺส อารมฺมณ\-ปจฺจเยน ปจฺจโย. สหชาต\-ปจฺจเยน ปจฺจโย. อุปนิสฺสย\-ปจฺจเยน ปจฺจโย. (3)}

\csromanmarkh{2. ปจฺจย\-ปจฺจนีย}\csromanmarkhii{2. สงฺขฺยาวาร}\csromanheadernomark{2}{2. ปจฺจย\-ปจฺจนีย 2. สงฺขฺยาวาร}
\vspace{\tipitakaparskip}
\csromancenter{นเหตุยา นว, นอารมฺมเณ นว, (สพฺพตฺถ นว,) โนอวิคเต นว.}
\proseitem{139}{เหตุปจฺจยา นอารมฺมเณ จตฺตาริ ฯเปฯ นสมนนฺตเร จตฺตาริ, นอญฺญมญฺเญ เทฺว, นอุปนิสฺสเย จตฺตาริ ฯเปฯ นสมฺปยุตฺเต เทฺว, นวิปฺปยุตฺเต จตฺตาริ, โนนตฺถิยา จตฺตาริ, โนวิคเต จตฺตาริ.}

\csromanheader{2}{4. ปจฺจย\-ปจฺจนียา\-นุโลม}
\proseitemwithcloser{140}{นเหตุปจฺจยา อารมฺมเณ นว, อธิปติยา นว, (อนุโลมมาติกา วิตฺถาเรตพฺพา) ฯเปฯ อวิคเต นว.}{สวิตกฺกทุกํ นิฏฺฐิตํ.}
\csromanmatikaanchor{mtk.17}
\csromantocmarkref{section}{88. สวิจารทุก}{mtk.17}
\bookmark[dest={mtk.17},level=1]{88. สวิจารทุก}
\csromanmarkh{88. สวิจารทุก}\csromanmarkhii{1. ปฏิจฺจวาร}\csromanheadernomark{1}{88. \textbf{สวิจารทุก 1. ปฏิจฺจวาร}}
\proseitemwithcloser{141}{สวิจารํ ธมฺมํ ปฏิจฺจ สวิจาโร ธมฺโม อุปฺปชฺชติ เหตุปจฺจยา. สวิจารํ เอกํ ขนฺธํ ปฏิจฺจ ตโย ขนฺธา ฯเปฯ เทฺว ขนฺเธ ฯเปฯ. ปฏิสนฺธิ\-กฺขเณ ฯเปฯ. (ยถา สวิตกฺกทุกํ, เอวํ กาตพฺพํ, นินฺนานากรณํ. อิธ มคฺเค จตฺตาริ กาตพฺพานิ. สวิจารทุเก อิมํ นานากรณํ.)}{สวิจารทุกํ นิฏฺฐิตํ.}
\csromanmatikaanchor{mtk.18}
\csromantocmarkref{section}{89. สปฺปีติกทุก}{mtk.18}
\bookmark[dest={mtk.18},level=1]{89. สปฺปีติกทุก}
\csromanmarkh{89. สปฺปีติกทุก}\csromanmarkhii{89. สปฺปีติกทุก 1. ปฏิจฺจวาร}\csromanheadernomark{1}{89. สปฺปีติกทุก \textbf{89. สปฺปีติกทุก 1. ปฏิจฺจวาร}}
\proseitem{142}{\csromanfolio{99}สปฺปีติกํ ธมฺมํ ปฏิจฺจ สปฺปีติโก ธมฺโม อุปฺปชฺชติ เหตุปจฺจยา. สปฺปีติกํ เอกํ ขนฺธํ ปฏิจฺจ ตโย ขนฺธา ฯเปฯ เทฺว ขนฺเธ ฯเปฯ. ปฏิสนฺธิ\-กฺขเณ ฯเปฯ. สปฺปีติกํ ธมฺมํ ปฏิจฺจ อปฺปีติโก ธมฺโม อุปฺปชฺชติ เหตุปจฺจยา. สปฺปีติเก ขนฺเธ ปฏิจฺจ ปีติ จ จิตฺต\-สมุฏฺฐานญฺจ รูปํ. ปฏิสนฺธิ\-กฺขเณ ฯเปฯ. สปฺปีติกํ ธมฺมํ ปฏิจฺจ สปฺปีติโก จ อปฺปีติโก จ ธมฺมา อุปฺปชฺชนฺติ เหตุปจฺจยา. สปฺปีติกํ เอกํ ขนฺธํ ปฏิจฺจ ตโย ขนฺธา ปีติ จ จิตฺต\-สมุฏฺฐานญฺจ รูปํ ฯเปฯ เทฺว ขนฺเธ ฯเปฯ. ปฏิสนฺธิ\-กฺขเณ ฯเปฯ. (3)}

\prose{อปฺปีติกํ ธมฺมํ ปฏิจฺจ อปฺปีติโก ธมฺโม อุปฺปชฺชติ เหตุปจฺจยา. อปฺปีติกํ เอกํ ขนฺธํ ปฏิจฺจ ตโย ขนฺธา จิตฺต\-สมุฏฺฐานญฺจ รูปํ ฯเปฯ เทฺว ขนฺเธ ฯเปฯ. ปีติํ ปฏิจฺจ จิตฺต\-สมุฏฺฐานํ รูปํ. ปฏิสนฺธิ\-กฺขเณ อปฺปีติกํ เอกํ ขนฺธํ ปฏิจฺจ ตโย ขนฺธา กฏตฺตา จ รูปํ. เทฺว ขนฺเธ ฯเปฯ. ปีติํ ปฏิจฺจ กฏตฺตารูปํ. ขนฺเธ ปฏิจฺจ วตฺถุ, วตฺถุํ ปฏิจฺจ ขนฺธา. ปีติํ ปฏิจฺจ วตฺถุ, วตฺถุํ ปฏิจฺจ ปีติ. เอกํ มหาภูตํ ฯเปฯ. (ยถา สวิตกฺกทุกํ สพฺพตฺถ, เอวํ สปฺปีติกทุกํ กาตพฺพํ. สพฺพตฺถ ปวตฺติ\-ปฏิสนฺธิ นวปิ ปญฺหา.)}

\csromanmarkh{1. ปจฺจยานุโลม}\csromanmarkhii{2. สงฺขฺยาวาร}\csromanheadernomark{2}{1. \textbf{ปจฺจยานุโลม 2. สงฺขฺยาวาร}}
\proseitem{143}{\textbf{เหตุ}ยา นว, อารมฺมเณ นว, อธิปติยา นว ฯเปฯ ปุเรชาเต ฉ ฯเปฯ กมฺเม นว, วิปาเก นว ฯเปฯ อวิคเต นว.}
\csromansectionrule

\csromanmarkh{2. ปจฺจย\-ปจฺจนีย}\csromanmarkhii{1. วิภงฺควาร}\csromanheadernomark{2}{2. ปจฺจย\-ปจฺจนีย 1. วิภงฺควาร}
\proseitem{144}{สปฺปีติกํ ธมฺมํ ปฏิจฺจ สปฺปีติโก ธมฺโม อุปฺปชฺชติ นเหตุปจฺจยา. อเหตุกํ สปฺปีติกํ เอกํ ขนฺธํ ปฏิจฺจ ตโย ขนฺธา ฯเปฯ เทฺว ขนฺเธ ฯเปฯ. สปฺปีติกํ ธมฺมํ ปฏิจฺจ อปฺปีติโก ธมฺโม อุปฺปชฺชติ นเหตุปจฺจยา. อเหตุเก สปฺปีติเก ขนฺเธ ปฏิจฺจ ปีติ จ จิตฺต\-สมุฏฺฐานญฺจ รูปํ.   \csromanfolio{100}สปฺปีติกํ ธมฺมํ ปฏิจฺจ สปฺปีติโก จ อปฺปีติโก จ ธมฺมา อุปฺปชฺชนฺติ นเหตุปจฺจยา. อเหตุกํ สปฺปีติกํ เอกํ ขนฺธํ ปฏิจฺจ ตโย ขนฺธา ปีติ จ จิตฺต\-สมุฏฺฐานญฺจ รูปํ ฯเปฯ เทฺว ขนฺเธ ฯเปฯ. (3)}

\prose{อปฺปีติกํ ธมฺมํ ปฏิจฺจ อปฺปีติโก ธมฺโม อุปฺปชฺชติ นเหตุปจฺจยา. อเหตุกํ อปฺปีติกํ เอกํ ขนฺธํ ปฏิจฺจ ตโย ขนฺธา จิตฺต\-สมุฏฺฐานญฺจ รูปํ ฯเปฯ เทฺว ขนฺเธ ฯเปฯ. ปีติํ ปฏิจฺจ จิตฺต\-สมุฏฺฐานํ รูปํ. อเหตุก\-ปฏิสนฺธิ\-กฺขเณ อปฺปีติกํ เอกํ ขนฺธํ ปฏิจฺจ ตโย ขนฺธา กฏตฺตา จ รูปํ ฯเปฯ เทฺว ขนฺเธ ปฏิจฺจ เทฺว ขนฺธา กฏตฺตา จ รูปํ. ขนฺเธ ปฏิจฺจ วตฺถุ, วตฺถุํ ปฏิจฺจ ขนฺธา. เอกํ มหาภูตํ ฯเปฯ. (ยาว อสญฺญสตฺตาปิ.) วิจิกิจฺฉา\-สหคเต อุทฺธจฺจ\-สหคเต ขนฺเธ ปฏิจฺจ วิจิกิจฺฉา\-สหคโต อุทฺธจฺจ\-สหคโต โมโห.  อปฺปีติกํ ธมฺมํ ปฏิจฺจ สปฺปีติโก ธมฺโม อุปฺปชฺชติ นเหตุปจฺจยา. อเหตุกํ ปีติํ ปฏิจฺจ สปฺปีติกา ขนฺธา. (มูลํ กาตพฺพํ.) ปีติํ ปฏิจฺจ สปฺปีติกา ขนฺธา จิตฺต\-สมุฏฺฐานญฺจ รูปํ. (3)}

\proseitem{145}{สปฺปีติกญฺจ อปฺปีติกญฺจ ธมฺมํ ปฏิจฺจ สปฺปีติโก ธมฺโม อุปฺปชฺชติ นเหตุปจฺจยา. อเหตุกํ สปฺปีติกํ เอกํ ขนฺธญฺจ ปีติญฺจ ปฏิจฺจ ตโย ขนฺธา ฯเปฯ เทฺว ขนฺเธ จ ฯเปฯ. สปฺปีติกญฺจ อปฺปีติกญฺจ ธมฺมํ ปฏิจฺจ อปฺปีติโก ธมฺโม อุปฺปชฺชติ นเหตุปจฺจยา. อเหตุเก สปฺปีติเก ขนฺเธ จ ปีติญฺจ ปฏิจฺจ จิตฺต\-สมุฏฺฐานํ รูปํ. อเหตุเก สปฺปีติเก ขนฺเธ จ มหาภูเต จ ปฏิจฺจ จิตฺต\-สมุฏฺฐานํ รูปํ.  สปฺปีติกญฺจ อปฺปีติกญฺจ ธมฺมํ ปฏิจฺจ สปฺปีติโก จ อปฺปีติโก จ ธมฺมา อุปฺปชฺชนฺติ นเหตุปจฺจยา. อเหตุกํ สปฺปีติกํ เอกํ ขนฺธญฺจ ปีติญฺจ ปฏิจฺจ ตโย ขนฺธา จิตฺต\-สมุฏฺฐานญฺจ รูปํ ฯเปฯ เทฺว ขนฺเธ จ ฯเปฯ. อเหตุกํ สปฺปีติกํ เอกํ ขนฺธญฺจ ปีติญฺจ ปฏิจฺจ ตโย ขนฺธา ฯเปฯ เทฺว ขนฺเธ จ ฯเปฯ. อเหตุเก สปฺปีติเก ขนฺเธ จ ปีติญฺจ มหาภูเต จ ปฏิจฺจ จิตฺต\-สมุฏฺฐานํ รูปํ. (3)}

\csromanmarkh{2. ปจฺจย\-ปจฺจนีย}\csromanmarkhii{2. สงฺขฺยาวาร}\csromanheadernomark{2}{2. ปจฺจย\-ปจฺจนีย 2. สงฺขฺยาวาร}
\proseitem{146}{นเหตุยา นว, นอารมฺมเณ ตีณิ, นอธิปติยา นว, นอนนฺตเร ตีณิ ฯเปฯ นอุปนิสฺสเย ตีณิ, นปุเรชาเต นว, นปจฺฉาชาเต นว, นอาเสวเน นว, นกมฺเม จตฺตาริ, นวิปาเก นว, นอาหาเร เอกํ,}

\prose{\csromanfolio{101}ไนนฺทฺริเย เอกํ, นฌาเน เอกํ, นมคฺเค นว, นสมฺปยุตฺเต ตีณิ, นวิปฺปยุตฺเต ฉ, โนนตฺถิยา ตีณิ, โนวิคเต ตีณิ.}

\vspace{\tipitakaparskip}
\csromancenter{(เอวํ อิตเร เทฺว คณนาปิ สหชาตวาโรปิ กาตพฺโพ.)}
\csromanmarkh{89. สปฺปีติกทุก}\csromanmarkhii{3. ปจฺจยวาร}\csromanheadernomark{2}{89. สปฺปีติกทุก 3. ปจฺจยวาร}
\proseitem{147}{สปฺปีติกํ ธมฺมํ ปจฺจยา สปฺปีติโก ธมฺโม อุปฺปชฺชติ เหตุปจฺจยา. (สํขิตฺตํ. ยถา สวิตกฺกทุเก อนุโลม\-ปจฺจยวารํ, เอวํ ปวตฺติ\-ปฏิสนฺธิ นว ปญฺหา ปริปุณฺณา ปีติ นินฺนานากรณา.)}

\prose{เหตุยา นว, อารมฺมเณ นว, อธิปติยา นว ฯเปฯ อวิคเต นว. อนุโลมํ.}

\prose{สปฺปีติกํ ธมฺมํ ปจฺจยา สปฺปีติโก ธมฺโม อุปฺปชฺชติ นเหตุปจฺจยา. ตีณิ. (ปฏิจฺจสทิสา.)}

\prose{อปฺปีติกํ ธมฺมํ ปจฺจยา อปฺปีติโก ธมฺโม อุปฺปชฺชติ นเหตุปจฺจยา. (ปวตฺติ\-ปฏิสนฺธิ กาตพฺพา. ปฏิจฺจ\-วาร\-สทิสา. ยาว อสญฺญสตฺตา.) จกฺขายตนํ ปจฺจยา จกฺขุวิญฺญาณํ ฯเปฯ. กายายตนํ ปจฺจยา กายวิญฺญาณํ. วตฺถุํ ปจฺจยา อเหตุกา อปฺปีติกา ขนฺธา ปีติ จ. วิจิกิจฺฉา\-สหคเต อุทฺธจฺจ\-สหคเต ขนฺเธ จ วตฺถุญฺจ ปจฺจยา วิจิกิจฺฉา\-สหคโต อุทฺธจฺจ\-สหคโต โมโห. (อนุโลมสทิสา. นว ปญฺหา ปวตฺติเยว. ปฏิสนฺธิ นตฺถิ. เอโกเยว โมโห.)}

\prose{นเหตุยา นว, นอารมฺมเณ ตีณิ, นอธิปติยา นว, นอนนฺตเร ตีณิ ฯเปฯ นอุปนิสฺสเย ตีณิ, นปุเรชาเต นว, นปจฺฉาชาเต นว, นอาเสวเน นว, นกมฺเม นว, นวิปาเก นว, นอาหาเร เอกํ, นอินฺทฺริเย เอกํ, นฌาเน เอกํ, นมคฺเค นว, นสมฺปยุตฺเต ตีณิ, นวิปฺปยุตฺเต ฉ, โนนตฺถิยา ตีณิ, โนวิคเต ตีณิ.}

\vspace{\tipitakaparskip}
\csromancenter{ปจฺจนียํ.}
\csromancenter{(เอวํ อิตเร เทฺว คณนาปิ สิสฺสยวาโรปิ กาตพฺโพ.)}
\csromanmarkh{89. สปฺปีติกทุก}\csromanmarkhii{5. สํสฏฺฐวาร}\csromanheadernomark{2}{89. \textbf{สปฺปีติกทุก 5. สํสฏฺฐวาร}}
\proseitem{148}{\csromanfolio{102}สปฺปีติกํ ธมฺมํ สํสฏฺโฐ สปฺปีติโก ธมฺโม อุปฺปชฺชติ เหตุปจฺจยา ฯเปฯ.}

\prose{เหตุยา ฉ, อารมฺมเณ ฉ, (สพฺพตฺถ ฉ,) อวิคเต ฉ. อนุโลมํ.}

\prose{นเหตุยา ฉ, นอธิปติยา ฉ, นปุเรชาเต ฉ, นปจฺฉาชาเต ฉ, นอาเสวเน ฉ, นกมฺเม จตฺตาริ, นวิปาเก ฉ, นฌาเน เอกํ, นมคฺเค ฉ, นวิปฺปยุตฺเต ฉ.}

\vspace{\tipitakaparskip}
\csromancenter{ปจฺจนียํ.}
\csromancenter{(เอวํ อิตเร เทฺว คณนาปิ สมฺปยุตฺตวาโรปิ กาตพฺโพ.)}
\csromanmarkh{89. สปฺปีติกทุก}\csromanmarkhii{7. ปญฺหาวาร}\csromanheadernomark{2}{89. \textbf{สปฺปีติกทุก 7. ปญฺหาวาร}}
\proseitem{149}{สปฺปีติโก ธมฺโม สปฺปีติกสฺส ธมฺมสฺส เหตุปจฺจเยน ปจฺจโย. สปฺปีติกา เหตู สมฺปยุตฺตกานํ ขนฺธานํ เหตุปจฺจเยน ปจฺจโย. ปฏิสนฺธิ\-กฺขเณ ฯเปฯ. สปฺปีติโก ธมฺโม อปฺปีติกสฺส ธมฺมสฺส เหตุปจฺจเยน ปจฺจโย. สปฺปีติกา เหตู ปีติยา จ จิตฺตสมุฏฺฐานานญฺจ รูปานํ เหตุปจฺจเยน ปจฺจโย. ปฏิสนฺธิ\-กฺขเณ ฯเปฯ. (มูลํ กาตพฺพํ.) สปฺปีติกา เหตู สมฺปยุตฺตกานํ ขนฺธานํ ปีติยา จ จิตฺตสมุฏฺฐานานญฺจ รูปานํ เหตุปจฺจเยน ปจฺจโย. ปฏิสนฺธิ\-กฺขเณ ฯเปฯ. (3)}

\prose{อปฺปีติโก ธมฺโม อปฺปีติกสฺส ธมฺมสฺส เหตุปจฺจเยน ปจฺจโย. อปฺปีติกา เหตู สมฺปยุตฺตกานํ ขนฺธานํ จิตฺตสมุฏฺฐานานญฺจ รูปานํ เหตุปจฺจเยน ปจฺจโย. ปฏิสนฺธิ\-กฺขเณ ฯเปฯ. (1)}

\csromanheader{2}{89. สปฺปีติกทุก \textbf{อารมฺมณ}}
\proseitem{150}{\csromanfolio{103}สปฺปีติโก ธมฺโม สปฺปีติกสฺส ธมฺมสฺส อารมฺมณ\-ปจฺจเยน ปจฺจโย. สปฺปีติเก ขนฺเธ อารพฺภ สปฺปีติกา ขนฺธา อุปฺปชฺชนฺติ. (มูลํ กาตพฺพํ.) สปฺปีติเก ขนฺเธ อารพฺภ อปฺปีติกา ขนฺธา จ ปีติ จ อุปฺปชฺชนฺติ. (มูลํ กาตพฺพํ.) สปฺปีติเก ขนฺเธ อารพฺภ สปฺปีติกา ขนฺธา จ ปีติ จ อุปฺปชฺชนฺติ. (3)}

\proseitem{151}{อปฺปีติโก ธมฺโม อปฺปีติกสฺส ธมฺมสฺส อรมฺมณปจฺจเยน ปจฺจโย. อปฺปีติเกน จิตฺเตน ทานํ ทตฺวา, สีลํ ฯเปฯ อุโปสถกมฺมํ กตฺวา อปฺปีติเกน จิตฺเตน ปจฺจเวกฺขติ, อสฺสาเทติ อภินนฺทติ, ตํ อารพฺภ อปฺปีติโก ราโค อุปฺปชฺชติ. ทิฏฺฐิ อุปฺปชฺชติ, วิจิกิจฺฉา อุปฺปชฺชติ. อุทฺธจฺจํ อุปฺปชฺชติ. โทมนสฺสํ อุปฺปชฺชติ. อปฺปีติกา ฌานา วุฏฺฐหิตฺวา ฯเปฯ มคฺคา วุฏฺฐหิตฺวา ฯเปฯ ผลา วุฏฺฐหิตฺวา อปฺปีติเกน จิตฺเตน ผลํ ปจฺจเวกฺขติ. อริยา อปฺปีติเกน จิตฺเตน นิพฺพานํ ปจฺจเวกฺขนฺติ. นิพฺพานํ อปฺปีติกสฺส โคตฺรภุสฺส, โวทานสฺส, มคฺคสฺส, ผลสฺส, อาวชฺชนาย ปีติยา จ อารมฺมณ\-ปจฺจเยน ปจฺจโย. อริยา อปฺปีติเกน จิตฺเตน อปฺปีติเก ปหีเน กิเลเส ฯเปฯ. วิกฺขมฺภิเต กิเลเส ฯเปฯ. ปุพฺเพ ฯเปฯ. จกฺขุํ ฯเปฯ. วตฺถุํ อปฺปีติเก ขนฺเธ จ ปีติญฺจ อปฺปีติเกน จิตฺเตน อนิจฺจโต ฯเปฯ วิปสฺสติ, อสฺสาเทติ อภินนฺทติ, ตํ อารพฺภ อปฺปีติโก ราโค อุปฺปชฺชติ ฯเปฯ โทมนสฺสํ อุปฺปชฺชติ. ทิพฺเพน จกฺขุนา รูปํ ปสฺสติ ฯเปฯ. โผฏฺฐพฺพา\-ยตนํ กายวิญฺญาณสฺส อารมฺมณ\-ปจฺจเยน ปจฺจโย. อปฺปีติกา ขนฺธา อิทฺธิวิธ\-ญาณสฺส, เจโตปริย\-ญาณสฺส, ปุพฺเพ\-นิวาสา\-นุสฺสติ\-ญาณสฺส, ยถากมฺมูปคญาณสฺส, อนาคตํสญาณสฺส, อาวชฺชานาย, ปีติยา จ อารมฺมณ\-ปจฺจเยน ปจฺจโย. อปฺปีติเก ขนฺเธ จ ปีติญฺจ อารพฺภ อปฺปีติกา ขนฺธา จ ปีติ จ อุปฺปชฺชนฺติ. (1)}

\prose{อปฺปีติโก ธมฺโม สปฺปีติกสฺส ธมฺมสฺส อารมฺมณ\-ปจฺจเยน ปจฺจโย. อปฺปีติเกน จิตฺเตน ทานํ ทตฺวา, สีลํ ฯเปฯ อุโปสถกมฺมํ ฯเปฯ. สปฺปีติเกน จิตฺเตน ปจฺจเวกฺขติ, อสฺสาเทติ อภินนฺทติ, ตํ อารพฺภ สปฺปีติโก ราโค อุปฺปชฺชติ. ทิฏฺฐิ อุปฺปชฺชติ. อปฺปีติกา ฌานา วุฏฺฐหิตฺวา ฯเปฯ มคฺคา วุฏฺฐหิตฺวา ฯเปฯ ผลา วุฏฺฐหิตฺวา สปฺปีติเกน จิตฺเตน ผลํ ปจฺจเวกฺขติ. อริยา สปฺปีติเกน จิตฺเตน นิพฺพานํ ปจฺจเวกฺขนฺติ. นิพฺพานํ สปฺปีติกสฺส โคตฺรภุสฺส, โวทานสฺส, มคฺคสฺส, ผลสฺส อารมฺมณ\-ปจฺจเยน ปจฺจโย. อริยา สปฺปีติเกน จิตฺเตน อปฺปีติเก ปหีเน กิเลเส ฯเปฯ. วิกฺขมฺภิเต \csromanfolio{104}กิเลเส ฯเปฯ. ปุพฺเพ ฯเปฯ. จกฺขุํ ฯเปฯ. วตฺถุํ อปฺปีติเก ขนฺเธ จ ปีติญฺจ สปฺปีติเกน จิตฺเตน อนิจฺจโต ฯเปฯ วิปสฺสติ, อสฺสาเทติ อภินนฺทติ, ตํ อารพฺภ สปฺปีติโก ราโค อุปฺปชฺชติ. ทิฏฺฐิ อุปฺปชฺชติ. อปฺปีติเก ขนฺเธ จ ปีติญฺจ อารพฺภ สปฺปีติกา ขนฺธา อุปฺปชฺชนฺติ. (2)}

\prose{อปฺปีติโก ธมฺโม สปฺปีติกสฺส จ อปฺปีติกสฺส จ ธมฺมสฺส อารมฺมณ\-ปจฺจเยน ปจฺจโย. อปฺปีติเกน จิตฺเตน ทานํ ทตฺวา, สีลํ ฯเปฯ อุโปสถกมฺมํ ฯเปฯ. สปฺปีติเกน จิตฺเตน ปจฺจเวกฺขติ, อสฺสาเทติ อภินนฺทติ, ตํ อารพฺภ สปฺปีติกา ขนฺธา จ ปีติ จ อุปฺปชฺชนฺติ. อปฺปีติกา ฌานา ฯเปฯ มคฺคา ฯเปฯ ผลา วุฏฺฐหิตฺวา สปฺปีติเกน จิตฺเตน ผลํ ปจฺจเวกฺขติ. อริยา สปฺปีติเกน จิตฺเตน นิพฺพานํ ปจฺจเวกฺขนฺติ. นิพฺพานํ สปฺปีติกสฺส โคตฺรภุสฺส, โวทานสฺส, มคฺคสฺส, ผลสฺส, ปีติยา จ อารมฺมณ\-ปจฺจเยน ปจฺจโย. อริยา สปฺปีติเกน จิตฺเตน อปฺปีติเก ปหีเน กิเลเส ปจฺจเวกฺขนฺติ. วิกฺขมฺภิเต กิเลเส ฯเปฯ. ปุพฺเพ ฯเปฯ. จกฺขุํ ฯเปฯ. วตฺถุํ อปฺปีติเก ขนฺเธ จ ปีติญฺจ สปฺปีติเกน จิตฺเตน อนิจฺจโต ฯเปฯ วิปสฺสติ, อสฺสาเทติ อภินนฺทติ, ตํ อารพฺภ สปฺปีติกา ขนฺธา จ ปีติ จ อุปฺปชฺชนฺติ. อปฺปีติเก ขนฺเธ จ ปิติญฺจ อารพฺภ สปฺปีติกา ขนฺธา จ ปีติ จ อุปฺปชฺชนฺติ. (3)}

\proseitem{152}{สปฺปีติโก จ อปฺปีติโก จ ธมฺมา สปฺปีติกสฺส ธมฺมสฺส อารมฺมณ\-ปจฺจเยน ปจฺจโย. สปฺปีติเก ขนฺเธ จ ปีติญฺจ อารพฺภ สปฺปีติกา ขนฺธา อุปฺปชฺชนฺติ. (มูลํ กาตพฺพํ.) สปฺปีติเก ขนฺเธ จ ปีติญฺจ อารพฺภ อปฺปีติกา ขนฺธา จ ปีติ จ อุปฺปชฺชนฺติ. (มูลํ กาตพฺพํ.) สปฺปีติเก ขนฺเธ จ ปีติญฺจ อารพฺภ สปฺปีติกา ขนฺธา จ ปีติ จ อุปฺปชฺชนฺติ. (3)}

\csromanheader{2}{\textbf{อธิปติ}}
\proseitem{153}{สปฺปีติโก ธมฺโม สปฺปีติกสฺส ธมฺมสฺส อธิปติ\-ปจฺจเยน ปจฺจโย. อารมฺมณา\-ธิปติ, สหชาตา\-ธิปติ.  อารมฺมณา\-ธิปติ\csromandash{}สปฺปีติเก ขนฺเธ ครุํ กตฺวา สปฺปีติกา ขนฺธา อุปฺปชฺชนฺติ. สหชาตา\-ธิปติ\csromandash{} สปฺปีติกา\-ธิปติ สมฺปยุตฺตกานํ ขนฺธานํ อธิปติ\-ปจฺจเยน ปจฺจโย. (1)}

\prose{สปฺปีติโก ธมฺโม อปฺปีติกสฺส ธมฺมสฺส อธิปติ\-ปจฺจเยน ปจฺจโย. อารมฺมณา\-ธิปติ, สหชาตา\-ธิปติ.  อารมฺมณา\-ธิปติ\csromandash{}สปฺปีติเก ขนฺเธ ครุํ กตฺวา อปฺปีติกา ขนฺธา จ ปีติ จ อุปฺปชฺชนฺติ. สหชาตา\-ธิปติ\csromandash{} \csromanfolio{105}สปฺปีติกา\-ธิปติ ปีติยา จ จิตฺตสมุฏฺฐานานญฺจ รูปานํ อธิปติ\-ปจฺจเยน ปจฺจโย. (2)}

\prose{สปฺปีติโก ธมฺโม สปฺปีติกสฺส จ อปฺปีติกสฺส จ ธมฺมสฺส อธิปติ\-ปจฺจเยน ปจฺจโย. อารมฺมณา\-ธิปติ, สหชาตา\-ธิปติ.  อารมฺมณา\-ธิปติ\csromandash{}สปฺปีติเก ขนฺเธ ครุํ กตฺวา สปฺปีติกา ขนฺธา จ ปีติ จ อุปฺปชฺชนฺติ. สหชาตา\-ธิปติ\csromandash{}สปฺปีติกา\-ธิปติ สมฺปยุตฺตกานํ ขนฺธานํ ปีติยา จ จิตฺตสมุฏฺฐานานญฺจ รูปานํ อธิปติ\-ปจฺจเยน ปจฺจโย. (3)}

\proseitem{154}{อปฺปีติโก ธมฺโม อปฺปีติกสฺส ธมฺมสฺส อธิปติ\-ปจฺจเยน ปจฺจโย. อารมฺมณา\-ธิปติ, สหชาตา\-ธิปติ.  อารมฺมณา\-ธิปติ\csromandash{} อปฺปีติเกน จิตฺเตน ทานํ ฯเปฯ สีลํ ฯเปฯ อุโปสถกมฺมํ ฯเปฯ. อปฺปีติเกน จิตฺเตน ตํ ครุํ กตฺวา ปจฺจเวกฺขติ, อสฺสาเทติ อภินนฺทติ, ตํ ครุํ กตฺวา อปฺปีติโก ราโค อุปฺปชฺชติ, ทิฏฺฐิ อุปฺปชฺชติ. อปฺปีติกา ฌานา ฯเปฯ มคฺคา ฯเปฯ ผลา วุฏฺฐหิตฺวา อปฺปีติเกน จิตฺเตน ผลํ ครุํ กตฺวา ปจฺจเวกฺขติ. อริยา อปฺปีติเกน จิตฺเตน นิพฺพานํ ครุํ กตฺวา ปจฺจเวกฺขนฺติ. นิพฺพานํ อปฺปีติกสฺส โคตฺรภุสฺส, โวทานสฺส, มคฺคสฺส, ผลสฺส อธิปติ\-ปจฺจเยน ปจฺจโย. จกฺขุํ ฯเปฯ. วตฺถุํ อปฺปีติเก ขนฺเธ จ ปีติญฺจ อปฺปีติเกน จิตฺเตน ครุํ กตฺวา อสฺสาเทติ อภินนฺทติ, ตํ ครุํ กตฺวา อปฺปีติโก ราโค อุปฺปชฺชติ. ทิฏฺฐิ อุปฺปชฺชติ. อปฺปีติเก ขนฺเธ จ ปีติญฺจ ครุํ กตฺวา อปฺปีติกา ขนฺธา จ ปีติ จ อุปฺปชฺชนฺติ. สหชาตา\-ธิปติ\csromandash{}อปฺปีติกาธิปติ สมฺปยุตฺตกานํ ขนฺธานํ จิตฺตสมุฏฺฐานานญฺจ รูปานํ อธิปติ\-ปจฺจเยน ปจฺจโย. (1)}

\prose{อปฺปีติโก ธมฺโม สปฺปีติกสฺส ธมฺมสฺส อธิปติ\-ปจฺจเยน ปจฺจโย. อารมฺมณา\-ธิปติ\csromandash{}อปฺปีติเกน จิตฺเตน ทานํ ฯเปฯ สีลํ ฯเปฯ อุโปสถกมฺมํ ฯเปฯ. (สํขิตฺตํ.) นิพฺพานํ สปฺปีติกสฺส โคตฺรภุสฺส, โวทานสฺส, มคฺคสฺส, ผลสฺส อธิปติ\-ปจฺจเยน ปจฺจโย. จกฺขุํ ฯเปฯ. วตฺถุํ อปฺปีติเก ขนฺเธ จ ปีติญฺจ สปฺปีติเกน จิตฺเตน ครุํ กตฺวา อสฺสาเทติ อภินนฺทติ, ตํ ครุํ กตฺวา สปฺปีติโก ราโค อุปฺปชฺชติ. ทิฏฺฐิ อุปฺปชฺชติ. อปฺปีติเก ขนฺเธ จ ปีติญฺจ ครุํ กตฺวา สปฺปีติกา ขนฺธา อุปฺปชฺชนฺติ. (2)}

\prose{อปฺปีติโก ธมฺโม สปฺปีติกสฺส จ อปฺปีติกสฺส จ ธมฺมสฺส อธิปติ\-ปจฺจเยน ปจฺจโย. อารมฺมณา\-ธิปติ\csromandash{}ทานํ ฯเปฯ (สํขิตฺตํ.) นิพฺพานํ สปฺปีติกสฺส โคตฺรภุสฺส, โวทานสฺส, มคฺคสฺส, ผลสฺส, ปีติยา จ \csromanfolio{106}อธิปติ\-ปจฺจเยน ปจฺจโย. จกฺขุํ ฯเปฯ. วตฺถุํ อปฺปีติเก ขนฺเธ จ ปีติญฺจ สปฺปีติเกน จิตฺเตน ครุํ กตฺวา อสฺสาเทติ อภินนฺทติ, ตํ ครุํ กตฺวา ราโค อุปฺปชฺชติ. ทิฏฺฐิ อุปฺปชฺชติ, อปฺปีติเก ขนฺเธ จ ปีติญฺจ ครุํ กตฺวา สปฺปีติกา ขนฺธา จ ปีติ จ อุปฺปชฺชนฺติ. (3)}

\proseitem{155}{สปฺปีติโก จ อปฺปีติโก จ ธมฺมา สปฺปีติกสฺส ธมฺมสฺส อธิปติ\-ปจฺจเยน ปจฺจโย. อารมฺมณา\-ธิปติ\csromandash{}สปฺปีติเก ขนฺเธ จ ปีติญฺจ ครุํ กตฺวา สปฺปีติกา ขนฺธา อุปฺปชฺชนฺติ. (มูลํ กาตพฺพํ.) สปฺปีติเก ขนฺเธ จ ปีติญฺจ ครุํ กตฺวา อปฺปีติกา ขนฺธา จ ปีติ จ อุปฺปชฺชนฺติ. (มูลํ กาตพฺพํ.) สปฺปีติเก ขนฺเธ จ ปีติญฺจ ครุํ กตฺวา สปฺปีติกา ขนฺธา จ ปีติ จ อุปฺปชฺชนฺติ. (3)}

\proseitem{156}{สปฺปีติโก ธมฺโม สปฺปีติกสฺส ธมฺมสฺส อนนฺตร\-ปจฺจเยน ปจฺจโย. ปุริมา ปุริมา สปฺปีติกา ขนฺธา ปจฺฉิมานํ ปจฺฉิมานํ สปฺปีติกานํ ขนฺธานํ อนนฺตร\-ปจฺจเยน ปจฺจโย. (มูลํ กาตพฺพํ.) ปุริมา ปุริมา สปฺปีติกา ขนฺธา ปจฺฉิมาย ปจฺฉิมาย ปีติยา อนนฺตร\-ปจฺจเยน ปจฺจโย. สปฺปีติกํ จุติจิตฺตํ อปฺปีติกสฺส อุปปตฺติ\-จิตฺตสฺส. สปฺปีติกํ ภวงฺคํ อาวชฺชนาย. สปฺปีติกา ขนฺธา อปฺปีติกสฺส วุฏฺฐานสฺส. ปีติสหคตา วิปาก\-มโนวิญฺญาณ\-ธาตุ กิริย\-มโนวิญฺญาณ\-ธาตุยา. สปฺปีติกํ ภวงฺคํ อปฺปีติกสฺส ภวงฺคสฺส. สปฺปีติกํ กุสลากุสลํ อปฺปีติกสฺส วุฏฺฐานสฺส. กิริยํ วุฏฺฐานสฺส. ผลํ วุฏฺฐานสฺส อนนฺตร\-ปจฺจเยน ปจฺจโย. (มูลํ กาตพฺพํ.) ปุริมา ปุริมา สปฺปีติกา ขนฺธา ปจฺฉิมานํ ปจฺฉิมานํ สปฺปีติกานํ ขนฺธานํ ปีติยา จ อนนฺตร\-ปจฺจเยน ปจฺจโย. (3)}

\proseitem{157}{อปฺปีติโก ธมฺโม อปฺปีติกสฺส ธมฺมสฺส อนนฺตร\-ปจฺจเยน ปจฺจโย. ปุริมา ปุริมา ปีติ ปจฺฉิมาย ปจฺฉิมาย ปีติยา อนนฺตร\-ปจฺจเยน ปจฺจโย. ปุริมา ปุริมา อปฺปีติกา ขนฺธา ปจฺฉิมานํ ปจฺฉิมานํ อปฺปีติกานํ ขนฺธานํ อนนฺตร\-ปจฺจเยน ปจฺจโย. อนุโลมํ โคตฺรภุสฺส ฯเปฯ. อปฺปีติกาย ผลสมาปตฺติยา ปีติยา จ อนนฺตร\-ปจฺจเยน ปจฺจโย. (มูลํ กาตพฺพํ.) ปุริมา ปุริมา ปีติ ปจฺฉิมานํ ปจฺฉิมานํ สปฺปีติกานํ ขนฺธานํ อนนฺตร\-ปจฺจเยน ปจฺจโย. อปฺปีติกํ จุติจิตฺตํ สปฺปีติกสฺส อุปปตฺติ\-จิตฺตสฺส. อาวชฺชนา สปฺปีติกานํ ขนฺธานํ. อปฺปีติกา ขนฺธา สปฺปีติกสฺส วุฏฺฐานสฺส. วิปาก\-มโนธาตุ สปฺปีติกาย วิปาก\-มโนวิญฺญาณ\-ธาตุยา. อปฺปีติกํ \csromanfolio{107}ภวงฺคํ สปฺปีติกสฺส ภวงฺคสฺส. อปฺปีติกํ กุสลากุสลํ สปฺปีติกสฺส วุฏฺฐานสฺส. กิริยํ วุฏฺฐานสฺส. ผลํ วุฏฺฐานสฺส อนนฺตร\-ปจฺจเยน ปจฺจโย. นิโรธา วุฏฺฐหนฺตสฺส เนว\-สญฺญา\-นา\-สญฺญา\-ยตนํ สปฺปีติกาย ผลสมาปตฺติยา อนนฺตร\-ปจฺจเยน ปจฺจโย. (2)}

\prose{อปฺปีติโก ธมฺโม สปฺปีติกสฺส จ อปฺปีติกสฺส จ ธมฺมสฺส อนนฺตร\-ปจฺจเยน ปจฺจโย. ปุริมา ปุริมา ปีติ ปจฺฉิมานํ ปจฺฉิมานํ สปฺปีติกานํ ขนฺธานํ ปีติยา จ อนนฺตร\-ปจฺจเยน ปจฺจโย. (3)}

\proseitem{158}{สปฺปีติโก จ อปฺปีติโก จ ธมฺมา สปฺปีติกสฺส ธมฺมสฺส อนนฺตร\-ปจฺจเยน ปจฺจโย. ปุริมา ปุริมา สปฺปีติกา ขนฺธา จ ปีติ จ ปจฺฉิมานํ ปจฺฉิมานํ สปฺปีติกานํ ขนฺธานํ อนนฺตร\-ปจฺจเยน ปจฺจโย. (1)}

\prose{สปฺปีติโก จ อปฺปีติโก จ ธมฺมา อปฺปีติกสฺส ธมฺมสฺส อนนฺตร\-ปจฺจเยน ปจฺจโย. ปุริมา ปุริมา สปฺปีติกา ขนฺธา จ ปีติ จ ปจฺฉิมาย ปจฺฉิมาย ปีติยา อนนฺตร\-ปจฺจเยน ปจฺจโย. สปฺปีติกํ จุติจิตฺตญฺจ ปีติ จ อปฺปีติกสฺส อุปปตฺติ\-จิตฺตสฺส. สปฺปีติกํ ภวงฺคญฺจ ปีติ จ อาวชฺชนาย. สปฺปีติกา ขนฺธา จ ปีติ จ อปฺปีติกสฺส วุฏฺฐานสฺส. สปฺปีติกา วิปาก\-มโนวิญฺญาณ\-ธาตุ จ ปีติ จ กิริย\-มโนวิญฺญาณ\-ธาตุยา. สปฺปีติกํ ภวงฺคญฺจ ปีติ จ อปฺปีติกสฺส ภวงฺคสฺส. สปฺปีติกํ กุสลากุสลญฺจ ปีติ จ อปฺปีติกสฺส วุฏฺฐานสฺส. กิริยญฺจ ปีติ จ วุฏฺฐานสฺส. ผลญฺจ ปีติ จ วุฏฺฐานสฺส อนนฺตร\-ปจฺจเยน ปจฺจโย. (2)}

\prose{สปฺปีติโก จ อปฺปีติโก จ ธมฺมา สปฺปีติกสฺส จ อปฺปีติกสฺส จ ธมฺมสฺส อนนฺตร\-ปจฺจเยน ปจฺจโย. ปุริมา ปุริมา สปฺปีติกา ขนฺธา จ ปีติ จ ปจฺฉิมานํ ปจฺฉิมานํ สปฺปีติกานํ ขนฺธานํ ปีติยา จ อนนฺตร\-ปจฺจเยน ปจฺจโย. (3)}

\prose{สมนนฺตร\-ปจฺจเยน ปจฺจโย. นว. สหชาต\-ปจฺจเยน ปจฺจโย. นว. อญฺญมญฺญ\-ปจฺจเยน ปจฺจโย. นว. นิสฺสย\-ปจฺจเยน ปจฺจโย. นว.}

\csromanheader{2}{\textbf{อุปนิสฺสย}}
\proseitem{159}{สปฺปีติโก ธมฺโม สปฺปีติกสฺส ธมฺมสฺส อุปนิสฺสย\-ปจฺจเยน ปจฺจโย. อารมฺมณู\-ปนิสฺสโย, อนนฺตรู\-ปนิสฺสโย, ปกตู\-ปนิสฺสโย ฯเปฯ. ปกตู\-ปนิสฺสโย\csromandash{}สปฺปีติกา ขนฺธา สปฺปีติกานํ ขนฺธานํ \csromanfolio{108}อุปนิสฺสย\-ปจฺจเยน ปจฺจโย. (มูลํ กาตพฺพํ.) สปฺปีติกา ขนฺธา อปฺปีติกานํ ขนฺธานํ ปีติยา จ อุปนิสฺสย\-ปจฺจเยน ปจฺจโย. (มูลํ กาตพฺพํ.) สปฺปีติกา ขนฺธา สปฺปีติกานํ ขนฺธานํ ปีติยา จ อุปนิสฺสย\-ปจฺจเยน ปจฺจโย. (3)}

\proseitem{160}{อปฺปีติโก ธมฺโม อปฺปีติกสฺส ธมฺมสฺส อุปนิสฺสย\-ปจฺจเยน ปจฺจโย. อารมฺมณู\-ปนิสฺสโย, อนนฺตรู\-ปนิสฺสโย, ปกตู\-ปนิสฺสโย ฯเปฯ. ปกตู\-ปนิสฺสโย\csromandash{}อปฺปีติกํ สทฺธํ อุปนิสฺสาย อปฺปีติเกน จิตฺเตน ทานํ เทติ, สิลํ สมาทิยติ, อุโปสถกมฺมํ กโรติ. อปฺปีติกํ ฌานํ ฯเปฯ วิปสฺสนํ. มคฺคํ. อภิญฺญํ. สมาปตฺติํ อุปฺปาเทติ. มานํ ชปฺเปติ. ทิฏฺฐิํ คณฺหาติ. อปฺปีติกํ สีลํ ฯเปฯ ปญฺญํ. ราคํ. โทสํ. โมหํ. มานํ. ทิฏฺฐิํ. ปตฺถนํ. กายิกํ สุขํ. กายิกํ ทุกฺขํ. อุตุํ. โภชนํ. เสนาสนํ ปีติํ อุปนิสฺสาย อปฺปีติเกน จิตฺเตน ทานํ เทติ ฯเปฯ สมาปตฺติํ อุปฺปาเทติ. ปาณํ หนติ ฯเปฯ สํฆํ ภินฺทติ. อปฺปีติกา สทฺธา ฯเปฯ. เสนาสนํ ปีติ จ อปฺปีติกาย สทฺธาย ฯเปฯ ปตฺถนาย. กายิกสฺส สุขสฺส. กายิกสฺส ทุกฺขสฺส. มคฺคสฺส ผลสมาปตฺติยา ปีติยา จ อุปนิสฺสย\-ปจฺจเยน ปจฺจโย. (1)}

\prose{อปฺปีติโก ธมฺโม สปฺปีติกสฺส ธมฺมสฺส อุปนิสฺสย\-ปจฺจเยน ปจฺจโย. (ตีณิ อุปนิสฺสยา.) อปฺปีติกํ สทฺธํ อุปนิสฺสาย สปฺปีติเกน จิตฺเตน ทานํ เทติ ฯเปฯ. อปฺปีติกา ฌานา ฯเปฯ. มานํ ชปฺเปติ. ทิฏฺฐิํ คณฺหาติ. อปฺปีติกํ สีลํ ฯเปฯ. เสนาสนํ ปีติํ อุปนิสฺสาย สปฺปีติเกน จิตฺเตน ทานํ เทติ ฯเปฯ สมาปตฺติํ อุปฺปาเทติ. สปฺปีติเกน จิตฺเตน อทินฺนํ อาทิยติ. มุสา ฯเปฯ ปิสุณํ ฯเปฯ สมฺผํ ฯเปฯ สนฺธิํ ฯเปฯ นิลฺโลปํ ฯเปฯ เอกาคาริกํ ฯเปฯ ปริปนฺเถ ฯเปฯ ปรทารํ ฯเปฯ คามฆาตํ ฯเปฯ นิคมฆาตํ กโรติ. อปฺปีติกา สทฺธา ฯเปฯ. เสนาสนํ ปีติ จ สปฺปีติกาย สทฺธาย ฯเปฯ ปญฺญาย. ราคสฺส. โมหสฺส. มานสฺส. ทิฏฺฐิยา. ปตฺถนาย. มคฺคสฺส, ผลสมาปตฺติยา อุปนิสฺสย\-ปจฺจเยน ปจฺจโย. (2)}

\prose{อปฺปีติโก ธมฺโม สปฺปีติกสฺส จ อปฺปีติกสฺส จ ธมฺมสฺส อุปนิสฺสย\-ปจฺจเยน ปจฺจโย. (ตีณิ อุปนิสฺสยา.) อปฺปีติกํ สทฺธํ อุปนิสฺสาย สปฺปีติเกน จิตฺเตน ทานํ เทติ ฯเปฯ สมาปตฺติํ อุปฺปาเทติ. มานํ ชปฺเปติ. ทิฏฺฐิํ คณฺหาติ. อปฺปีติกํ สีลํ ฯเปฯ. เสนาสนํ ปีติํ อุปนิสฺสาย \csromanfolio{109}สปฺปีติเกน จิตฺเตน ทานํ เทติ ฯเปฯ สมาปตฺติํ อุปฺปาเทติ. สปฺปีติเกน จิตฺเตน อทินฺนํ อาทิยติ ฯเปฯ. (ทุติยวาร\-สทิสํ.) นิคมฆาตํ กโรติ. อปฺปีติกา สทฺธา ฯเปฯ. เสนาสนํ ปีติ จ สปฺปีติกาย สทฺธาย ฯเปฯ ปญฺญาย. ราคสฺส. โมหสฺส. มานสฺส. ทิฏฺฐิยา. ปตฺถนาย. มคฺคสฺส, ผลสมาปตฺติยา ปีติยา จ อุปนิสฺสย\-ปจฺจเยน ปจฺจโย. (3)}

\prose{สปฺปีติโก จ อปฺปีติโก จ ธมฺมา สปฺปีติกสฺส ธมฺมสฺส อุปนิสฺสย\-ปจฺจเยน ปจฺจโย. (ตีณิปิ อุปนิสฺสยา.) สปฺปีติกา ขนฺธา จ ปีติ จ สปฺปีติกานํ ขนฺธานํ อุปนิสฺสย\-ปจฺจเยน ปจฺจโย. (มูลํ กาตพฺพํ.) สปฺปีติกา ขนฺธา จ ปีติ จ อปฺปีติกานํ ขนฺธานํ ปีติยา จ อุปนิสฺสย\-ปจฺจเยน ปจฺจโย. (มูลํ กาตพฺพํ.) สปฺปีติกา ขนฺธา จ ปีติ จ สปฺปีติกานํ ขนฺธานํ ปีติยา จ อุปนิสฺสย\-ปจฺจเยน ปจฺจโย. (3)}

\csromanheader{2}{\textbf{ปุเรชาต}}
\proseitem{161}{อปฺปีติโก ธมฺโม อปฺปีติกสฺส ธมฺมสฺส ปุเรชาต\-ปจฺจเยน ปจฺจโย. อารมฺมณ\-ปุเรชาตํ, วตฺถุ\-ปุเรชาตํ.  อารมฺมณ\-ปุเรชาตํ\csromandash{} จกฺขุํ ฯเปฯ. วตฺถุํ อปฺปีติเกน จิตฺเตน อนิจฺจโต ฯเปฯ วิปสฺสติ, อสฺสาเทติ อภินนฺทติ, ตํ อารพฺภ อปฺปีติโก ราโค ฯเปฯ โทมนสฺสํ อุปฺปชฺชติ. ปีติ อุปฺปชฺชติ. ทิพฺเพน จกฺขุนา รูปํ ปสฺสาติ ฯเปฯ. โผฏฺฐพฺพา\-ยตนํ กายวิญฺญาณสฺส ฯเปฯ. วตฺถุ\-ปุเรชาตํ\csromandash{}จกฺขายตนํ จกฺขุ\-วิญฺญาณสฺส ฯเปฯ กายายตนํ กายวิญฺญาณสฺส ฯเปฯ. วตฺถุ อปฺปีติกานํ ขนฺธานํ ปีติยา จ ปุเรชาต\-ปจฺจเยน ปจฺจโย. (1)}

\prose{อปฺปีติโก ธมฺโม สปฺปีติกสฺส ธมฺมสฺส ปุเรชาต\-ปจฺจเยน ปจฺจโย. อารมฺมณ\-ปุเรชาตํ, วตฺถุ\-ปุเรชาตํ.  อารมฺมณ\-ปุเรชาตํ\csromandash{}จกฺขุํ ฯเปฯ. วตฺถุํ สปฺปีติเกน จิตฺเตน อนิจฺจโต ฯเปฯ วิปสฺสติ, อสฺสาเทติ อภินนฺทติ, ตํ อารพฺภ สปฺปีติโก ราโค อุปฺปชฺชติ. ทิฏฺฐิ อุปฺปชฺชติ. วตฺถุ\-ปุเรชาตํ\csromandash{}วตฺถุ สปฺปีติกานํ ขนฺธานํ ปุเรชาต\-ปจฺจเยน ปจฺจโย. (2)}

\prose{อปฺปีติโก ธมฺโม สปฺปีติกสฺส จ อปฺปีติกสฺส จ ธมฺมสฺส ปุเรชาต\-ปจฺจเยน ปจฺจโย. อารมฺมณ\-ปุเรชาตํ, วตฺถุ\-ปุเรชาตํ.  อารมฺมณ\-ปุเรชาตํ\csromandash{}จกฺขุํ ฯเปฯ. วตฺถุํ สปฺปีติเกน จิตฺเตน อนิจฺจโต ฯเปฯ วิปสฺสติ, อสฺสาเทติ อภินนฺทติ, ตํ อารพฺภ ปีติ จ สมฺปยุตฺตกา ขนฺธา จ อุปฺปชฺชนฺติ. \csromanfolio{110}วตฺถุ\-ปุเรชาตํ\csromandash{}วตฺถุ สปฺปีติกานํ ขนฺธานํ ปีติยา จ ปุเรชาต\-ปจฺจเยน ปจฺจโย. (3)}

\csromanheader{2}{\textbf{ปจฺฉาชาตาทิ}}
\proseitem{162}{สปฺปีติโก ธมฺโม อปฺปีติกสฺส ธมฺมสฺส ปจฺฉาชาต\-ปจฺจเยน ปจฺจโย. ตีณิ. อาเสวน\-ปจฺจเยน ปจฺจโย. นว. กมฺมปจฺจเยน ปจฺจโย. ฉ. (สหชาตาปิ นานากฺขณิกาปิ กาตพฺพา. เทฺว นานากฺขณิกา.) วิปาก\-ปจฺจเยน ปจฺจโย. นว. อาหารปจฺจเยน ปจฺจโย. จตฺตาริ. อินฺทฺริย\-ปจฺจเยน ปจฺจโย. จตฺตาริ. ฌานปจฺจเยน ปจฺจโย. นว. มคฺคปจฺจเยน ปจฺจโย. จตฺตาริ. สมฺปยุตฺต\-ปจฺจเยน ปจฺจโย. ฉ. วิปฺปยุตฺต\-ปจฺจเยน ปจฺจโย. ปญฺจ. อตฺถิปจฺจเยน ปจฺจโย. นว. (สํขิตฺตํ. สวิตกฺก\-ทุก\-สทิสํ กาตพฺพํ.) นตฺถิ\-ปจฺจเยน ปจฺจโย. วิคต\-ปจฺจเยน ปจฺจโย. อวิคต\-ปจฺจเยน ปจฺจโย. นว.}

\csromanmarkh{1. ปจฺจยานุโลม}\csromanmarkhii{2. สงฺขฺยาวาร}\csromanheadernomark{2}{1. \textbf{ปจฺจยานุโลม 2. สงฺขฺยาวาร}}
\proseitem{163}{เหตุยา จตฺตาริ, อารมฺมเณ นว, อธิปติยา นว, อนนฺตเร นว, สมนนฺตเร นว, สหชาเต นว, อญฺญมญฺเญ นว, นิสฺสเย นว, อุปนิสฺสเย นว, ปุเรชาเต ตีณิ, ปจฺฉาชาเต ตีณิ, อาเสวเน นว, กมฺเม ฉ, วิปาเก นว, อาหาเร จตฺตาริ, อินฺทฺริเย จตฺตาริ, ฌาเน นว, มคฺเค จตฺตาริ, สมฺปยุตฺเต ฉ, วิปฺปยุตฺเต ปญฺจ, อตฺถิยา นว, นตฺถิยา นว, วิคเต นว, อวิคเต นว.}

\vspace{\tipitakaparskip}
\csromancenter{อนุโลมํ.}
\proseitem{164}{สปฺปีติโก ธมฺโม สปฺปีติกสฺส ธมฺมสฺส อารมฺมณ\-ปจฺจเยน ปจฺจโย. สหชาต\-ปจฺจเยน ปจฺจโย. อุปนิสฺสย\-ปจฺจเยน ปจฺจโย. กมฺมปจฺจเยน ปจฺจโย.  สปฺปีติโก ธมฺโม อปฺปีติกสฺส ธมฺมสฺส อารมฺมณ\-ปจฺจเยน ปจฺจโย. สหชาต\-ปจฺจเยน ปจฺจโย. อุปนิสฺสย\-ปจฺจเยน ปจฺจโย. ปจฺฉาชาต\-ปจฺจเยน ปจฺจโย. กมฺมปจฺจเยน ปจฺจโย.  }

\vspace{\tipitakaparskip}
\csromancenter{ปีติสหคตทุก สปฺปีติโก ธมฺโม สปฺปีติกสฺส จ อปฺปีติกสฺส จ ธมฺมสฺส อารมฺมณ\-ปจฺจเยน ปจฺจโย. สหชาต\-ปจฺจเยน ปจฺจโย. อุปนิสฺสย\-ปจฺจเยน ปจฺจโย. กมฺมปจฺจเยน ปจฺจโย. (3)}
\proseitem{165}{\csromanfolio{111}อปฺปีติโก ธมฺโม อปฺปีติกสฺส ธมฺมสฺส อารมฺมณ\-ปจฺจเยน ปจฺจโย. สหชาต\-ปจฺจเยน ปจฺจโย. อุปนิสฺสย\-ปจฺจเยน ปจฺจโย. ปุเรชาต\-ปจฺจเยน ปจฺจโย. ปจฺฉาชาต\-ปจฺจเยน ปจฺจโย. กมฺมปจฺจเยน ปจฺจโย. อาหารปจฺจเยน ปจฺจโย. อินฺทฺริย\-ปจฺจเยน ปจฺจโย.  อปฺปีติโก ธมฺโม สปฺปีติกสฺส ธมฺมสฺส อารมฺมณ\-ปจฺจเยน ปจฺจโย. สหชาต\-ปจฺจเยน ปจฺจโย. อุปนิสฺสย\-ปจฺจเยน ปจฺจโย. ปุเรชาต\-ปจฺจเยน ปจฺจโย.  อปฺปีติโก ธมฺโม สปฺปีติกสฺส จ อปฺปีติกสฺส จ ธมฺมสฺส อารมฺมณ\-ปจฺจเยน ปจฺจโย. สหชาต\-ปจฺจเยน ปจฺจโย. อุปนิสฺสย\-ปจฺจเยน ปจฺจโย. ปุเรชาต\-ปจฺจเยน ปจฺจโย. (3)}

\proseitem{166}{สปฺปีติโก จ อปฺปีติโก จ ธมฺมา สปฺปีติกสฺส ธมฺมสฺส อารมฺมณ\-ปจฺจเยน ปจฺจโย. สหชาต\-ปจฺจเยน ปจฺจโย. อุปนิสฺสย\-ปจฺจเยน ปจฺจโย.  สปฺปีติโก จ อปฺปีติโก จ ธมฺมา อปฺปีติกสฺส ธมฺมสฺส อารมฺมณ\-ปจฺจเยน ปจฺจโย. สหชาต\-ปจฺจเยน ปจฺจโย. อุปนิสฺสย\-ปจฺจเยน ปจฺจโย. ปจฺฉาชาต\-ปจฺจเยน ปจฺจโย.  สปฺปีติโก จ อปฺปีติโก จ ธมฺมา สปฺปีติกสฺส จ อปฺปีติกสฺส จ ธมฺมสฺส อารมฺมณ\-ปจฺจเยน ปจฺจโย. สหชาต\-ปจฺจเยน ปจฺจโย. อุปนิสฺสย\-ปจฺจเยน ปจฺจโย. (3)}

\prosewithcloser{(ปจฺจนีย\-วิภงฺค\-คณนาปิ สวิตกฺก\-ทุก\-สทิสา. ยทิปิ น สเมติ อิมํ อนุโลมํ ปจฺจเวกฺขิตฺวา คเณตพฺพํ. อิตเร เทฺว คณนา คเณตพฺพา.)}{สปฺปีติกทุกํ นิฏฺฐิตํ.}
\csromanmatikaanchor{mtk.19}
\csromantocmarkref{section}{90. ปีติสหคตทุก}{mtk.19}
\bookmark[dest={mtk.19},level=1]{90. ปีติสหคตทุก}
\csromanmarkh{90. ปีติสหคตทุก}\csromanmarkhii{1. ปฏิจฺจวาร}\csromanheadernomark{1}{90. ปีติสหคตทุก 1. ปฏิจฺจวาร}
\proseitem{167}{ปีติสหคตํ ธมฺมํ ปฏิจฺจ ปีติสหคโต ธมฺโม อุปฺปชฺชติ เหตุปจฺจยา. ปีติสหคตํ เอกํ ขนฺธํ ปฏิจฺจ ตโย ขนฺธา ฯเปฯ เทฺว ขนฺเธ ฯเปฯ.}

\prosewithcloser{\csromanfolio{112}(เอวํ ปีติสหคตทุกํ วิตฺถาเรตพฺพํ. สปฺปีติก\-ทุก\-สทิสํ, นินฺนานากรณํ, อามสนํ นินฺนานํ.)}{ปีติสหคตทุกํ นิฏฺฐิตํ.}
\csromanmatikaanchor{mtk.20}
\csromantocmarkref{section}{91. สุขสหคตทุก}{mtk.20}
\bookmark[dest={mtk.20},level=1]{91. สุขสหคตทุก}
\csromanmarkh{91. สุขสหคตทุก}\csromanmarkhii{1. ปฏิจฺจวาร}\csromanheadernomark{1}{91. \textbf{สุขสหคตทุก 1. ปฏิจฺจวาร}}
\proseitem{168}{สุขสหคตํ ธมฺมํ ปฏิจฺจ สุขสหคโต ธมฺโม อุปฺปชฺชติ เหตุปจฺจยา. สุขสหคตํ เอกํ ขนฺธํ ปฏิจฺจ เทฺว ขนฺธา, เทฺว ขนฺเธ ปฏิจฺจ เอโก ขนฺโธ. ปฏิสนฺธิ\-กฺขเณ ฯเปฯ. สุขสหคตํ ธมฺมํ ปฏิจฺจ นสุขสหคโต ธมฺโม อุปฺปชฺชติ เหตุปจฺจยา. สุขสหคเต ขนฺเธ ปฏิจฺจ สุขํ จิตฺต\-สมุฏฺฐานญฺจ รูปํ. (เอวํ สุขสหคตทุกํ วิตฺถาเรตพฺพํ. ยถา สปฺปีติก\-ทุกสฺส อนุโลม\-ปฏิจฺจ\-วาโร.)}

\prose{\textbf{เหตุ}ยา นว, อารมฺมเณ นว ฯเปฯ ปุเรชาเต ฉ, อาเสวเน ฉ, กมฺเม นว ฯเปฯ อวิคเต นว. อนุโลมํ.}

\proseitem{169}{สุขสหคตํ ธมฺมํ ปฏิจฺจ สุขสหคโต ธมฺโม อุปฺปชฺชติ นเหตุปจฺจยา. อเหตุกํ สุขสหคตํ เอกํ ขนฺธํ ปฏิจฺจ เทฺว ขนฺธา, เทฺว ขนฺเธ ฯเปฯ. (มูลํ กาตพฺพํ.) อเหตุเก สุขสหคเต ขนฺเธ ปฏิจฺจ สุขญฺจ จิตฺต\-สมุฏฺฐานญฺจ รูปํ.  สุขสหคตํ ธมฺมํ ปฏิจฺจ สุขสหคโต จ นสุขสหคโต จ ธมฺมา อุปฺปชฺชนฺติ นเหตุปจฺจยา. อเหตุกํ สุขสหคตํ เอกํ ขนฺธํ ปฏิจฺจ เทฺว ขนฺธา สุขญฺจ จิตฺต\-สมุฏฺฐานญฺจ รูปํ. เทฺว ขนฺเธ ฯเปฯ. (3)}

\prose{นสุขสหคตํ ธมฺมํ ปฏิจฺจ นสุขสหคโต ธมฺโม อุปฺปชฺชติ นเหตุปจฺจยา. อเหตุกํ นสุขสหคตํ เอกํ ขนฺธํ ปฏิจฺจ ตโย ขนฺธา จิตฺต\-สมุฏฺฐานญฺจ รูปํ ฯเปฯ เทฺว ขนฺเธ ฯเปฯ. อเหตุกํ สุขํ ปฏิจฺจ จิตฺต\-สมุฏฺฐานํ รูปํ. อเหตุก\-ปฏิสนฺธิ\-กฺขเณ นสุขสหคตํ เอกํ ขนฺธํ ปฏิจฺจ ตโย ขนฺธา กฏตฺตา จ}

\vspace{\tipitakaparskip}
\csromancenter{สุขสหคตทุก รูปํ ฯเปฯ เทฺว ขนฺเธ ฯเปฯ. ขนฺเธ ปฏิจฺจ วตฺถุ, วตฺถุํ ปฏิจฺจ ขนฺธา. เอกํ มหาภูตํ ฯเปฯ. (ยาว อสญฺญสตฺตา.) วิจิกิจฺฉา\-สหคเต อุทฺธจฺจ\-สหคเต ขนฺเธ ปฏิจฺจ วิจิกิจฺฉา\-สหคโต อุทฺธจฺจ\-สหคโต โมโห.}
\csromancenter{(ยถา สปฺปีติเก นเหตุปจฺจย\-สทิสํ, นินฺนานํ, สพฺพตฺถเมว นว ปญฺหา.)}
\prose{\csromanfolio{113}นเหตุยา นว, นอารมฺมเณ ตีณิ ฯเปฯ นอุปนิสฺสเย ตีณิ, นปุเรชาเต นว, นปจฺฉาชาเต นว, นอาเสวเน นว, นกมฺเม จตฺตาริ, นวิปาเก นว, นอาหาเร เอกํ, นอินฺทฺริเย เอกํ, นฌาเน ฉ, นมคฺเค นว, นสมฺปยุตฺเต ตีณิ, นวิปฺปยุตฺเต ฉ, โนนตฺถิยา ตีณิ, โนวิคเต ตีณิ. ปจฺจนียํ. (เอวํ อิตเร เทฺว คณนาปิ สหชาตวาโรปิ ปฏิจฺจ\-วาร\-สทิสา. ปจฺจยวาเร ปวตฺติปิ ปฏิสนฺธิปิ วิตฺถาเรตพฺพา, ยถา สปฺปีติกทุกปจฺจวารปจฺจนีเยปิ ปวตฺเต วตฺถุ จ วิตฺถาเรตพฺพํ, ยถา สปฺปีติกทุเก เอโกเยว โมโห, เอวํ อิตเร เทฺว คณนาปิ นิสฺสยวาโรปิ สํสฏฺฐ\-วาโรปิ สมฺปยุตฺตวาโรปิ ยถา สปฺปีติกทุกํ, เอวํ กาตพฺพํ.)}

\csromanmarkh{91. สุขสหคตทุก}\csromanmarkhii{7. ปญฺหาวาร}\csromanheadernomark{2}{91. สุขสหคตทุก 7. ปญฺหาวาร}
\proseitem{170}{สุขสหคโต ธมฺโม สุขสหคตสฺส ธมฺมสฺส เหตุปจฺจเยน ปจฺจโย. จตฺตาริ.}

\vspace{\tipitakaparskip}
\csromancenter{(อารมฺมเณปิ อธิปติยาปิ สปฺปีติก\-ทุก\-สทิสา. “สุขนฺ”ติ นานากรณํ.)}
\proseitem{171}{สุขสหคโต ธมฺโม สุขสหคตสฺส ธมฺมสฺส อนนฺตร\-ปจฺจเยน ปจฺจโย. ปุริมา ปุริมา สุขสหคตา ขนฺธา ปจฺฉิมานํ ปจฺฉิมานํ สุขสหคตานํ ขนฺธานํ อนนฺตร\-ปจฺจเยน ปจฺจโย. (มูลํ กาตพฺพํ.) ปุริมา ปุริมา \csromanfolio{114}สุขสหคตา ขนฺธา ปจฺฉิมสฺส ปจฺฉิมสฺส สุขสฺส อนนฺตร\-ปจฺจเยน ปจฺจโย. สุขสหคตํ จุติจิตฺตํ นสุข\-สหคตสฺส อุปปตฺติ\-จิตฺตสฺส. สุขสหคตํ ภวงฺคํ อาวชฺชนาย. สุขสหคตํ กายวิญฺญาณํ วิปาก\-มโนธาตุยา. สุขสหคตา วิปาก\-มโนวิญฺญาณ\-ธาตุ กิริย\-มโนวิญฺญาณ\-ธาตุยา. สุขสหคตํ ภวงฺคํ นสุข\-สหคตสฺส ภวงฺคสฺส. สุขสหคตํ กุสลากุสลํ นสุข\-สหคตสฺส วุฏฺฐานสฺส. กิริยํ วุฏฺฐานสฺส. ผลํ วุฏฺฐานสฺส อนนฺตร\-ปจฺจเยน ปจฺจโย. (มูลํ กาตพฺพํ.) ปุริมา ปุริมา สุขสหคตา ขนฺธา ปจฺฉิมานํ ปจฺฉิมานํ สุขสหคตานํ ขนฺธานํ สุขสฺส จ อนนฺตร\-ปจฺจเยน ปจฺจโย. (3)}

\prose{นสุขสหคโต ธมฺโม นสุข\-สหคตสฺส ธมฺมสฺส อนนฺตร\-ปจฺจเยน ปจฺจโย. ปุริมา ปุริมา นสุขสหคตา ฯเปฯ. (มูลํ. ตีณิปิ สปฺปีติก\-ทุก\-สทิสา.)}

\proseitem{172}{สุขสหคโต จ นสุขสหคโต จ ธมฺมา สุขสหคตสฺส ธมฺมสฺส อนนฺตร\-ปจฺจเยน ปจฺจโย. ปุริมา ปุริมา สุขสหคตา ขนฺธา จ สุขญฺจ ปจฺฉิมานํ ปจฺฉิมานํ สุขสหคตานํ ขนฺธานํ อนนฺตร\-ปจฺจเยน ปจฺจโย. (มูลํ กาตพฺพํ.) ปุริมา ปุริมา สุขสหคตา ขนฺธา จ สุขญฺจ ปจฺฉิมสฺส ปจฺฉิมสฺส สุขสฺส อนนฺตร\-ปจฺจเยน ปจฺจโย. สุขสหคตํ จุติจิตฺตญฺจ สุขญฺจ นสุข\-สหคตสฺส อุปปตฺติ\-จิตฺตสฺส. สุขสหคตํ ภวงฺคญฺจ สุขญฺจ อาวชฺชนาย. สุขสหคตํ กายวิญฺญาณญฺจ สุขญฺจ วิปาก\-มโนธาตุยา. สุขสหคตา วิปาก\-มโนวิญฺญาณ\-ธาตุ จ สุขญฺจ กิริย\-มโนวิญฺญาณ\-ธาตุยา. สุขสหคตํ ภวงฺคญฺจ สุขญฺจ นสุข\-สหคตสฺส ภวงฺคสฺส. สุขสหคตํ กุสลากุสลญฺจ สุขญฺจ นสุข\-สหคตสฺส วุฏฺฐานสฺส. กิริยํ วุฏฺฐานสฺส. ผลํ วุฏฺฐานสฺส อนนฺตร\-ปจฺจเยน ปจฺจโย. (มูลํ กาตพฺพํ.) ปุริมา ปุริมา สุขสหคตา ขนฺธา จ สุขญฺจ ปจฺฉิมานํ ปจฺฉิมานํ สุขสหคตานํ ขนฺธานํ สุขสฺส จ อนนฺตร\-ปจฺจเยน ปจฺจโย. (3)}

\prose{สมนนฺตร\-ปจฺจเยน ปจฺจโย. สหชาต\-ปจฺจเยน ปจฺจโย. อญฺญมญฺญ\-ปจฺจเยน ปจฺจโย. นิสฺสย\-ปจฺจเยน ปจฺจโย.}

\csromanheader{2}{\textbf{อุปนิสฺสย}}
\proseitem{173}{สุขสหคโต ธมฺโม สุขสหคตสฺส ธมฺมสฺส อุปนิสฺสย\-ปจฺจเยน ปจฺจโย. ตีณิ.}

\csromanheader{2}{91. สุขสหคตทุก}
\prose{\csromanfolio{115}นสุขสหคโต ธมฺโม นสุข\-สหคตสฺส ธมฺมสฺส อุปนิสฺสย\-ปจฺจเยน ปจฺจโย. อารมฺมณู\-ปนิสฺสโย, อนนฺตรู\-ปนิสฺสโย, ปกตู\-ปนิสฺสโย ฯเปฯ. ปกตู\-ปนิสฺสโย\csromandash{}นสุขสหคตํ สทฺธํ อุปนิสฺสาย นสุข\-สหคเตน จิตฺเตน ทานํ เทติ. สีลํ ฯเปฯ สมาปตฺติํ อุปฺปาเทติ, มานํ ชปฺเปติ, ทิฏฺฐิํ คณฺหาติ. นสุขสหคตํ สีลํ ฯเปฯ ปญฺญํ. ราคํ ฯเปฯ. ปตฺถนํ. กายิกํ สุขํ. กายิกํ ทุกฺขํ. อุตุํ. โภชนํ. เสนาสนํ สุขํ อุปนิสฺสาย นสุข\-สหคเตน จิตฺเตน ทานํ เทติ ฯเปฯ สมาปตฺติํ อุปฺปาเทติ. ปาณํ หนติ ฯเปฯ สํฆํ ภินฺทติ. นสุขสหคตา สทฺธา ฯเปฯ. เสนาสนํ สุขญฺจ นสุข\-สหคตาย สทฺธาย ฯเปฯ ปญฺญาย. ราคสฺส. โทสสฺส ฯเปฯ ปตฺถนาย. กายิกสฺส สุขสฺส, กายิกสฺส ทุกฺขสฺส, มคฺคสฺส, ผลสมาปตฺติยา สุขสฺส จ อุปนิสฺสย\-ปจฺจเยน ปจฺจโย. (1)}

\prose{นสุขสหคโต ธมฺโม สุขสหคตสฺส ธมฺมสฺส อุปนิสฺสย\-ปจฺจเยน ปจฺจโย. อารมฺมณู\-ปนิสฺสโย, อนนฺตรู\-ปนิสฺสโย, ปกตู\-ปนิสฺสโย ฯเปฯ. ปกตู\-ปนิสฺสโย\csromandash{}นสุขสหคตํ สทฺธํ อุปนิสฺสาย สุขสหคเตน จิตฺเตน ทานํ เทติ ฯเปฯ สมาปตฺติํ อุปฺปาเทติ. มานํ ชปฺเปติ, ทิฏฺฐิํ คณฺหาติ. นสุขสหคตํ สีลํ ฯเปฯ ปญฺญํ. ราคํ ฯเปฯ. ปตฺถนํ. กายิกํ สุขํ. กายิกํ ทุกฺขํ. อุตุํ. โภชนํ. เสนาสนํ สุขํ อุปนิสฺสาย สุขสหคเตน จิตฺเตน ทานํ เทติ ฯเปฯ. สมาปตฺติํ อุปฺปาเทติ. สุขสหคเตน จิตฺเตน อทินฺนํ อาทิยติ. มุสา ฯเปฯ ปิสุณํ ฯเปฯ สมฺผํ ฯเปฯ สนฺธิํ ฯเปฯ นิลฺโลปํ ฯเปฯ เอกาคาริกํ ฯเปฯ ปริปนฺเถ ฯเปฯ ปรทารํ ฯเปฯ คามฆาตํ ฯเปฯ นิคมฆาตํ กโรติ. นสุขสหคตา สทฺธา ฯเปฯ. เสนาสนํ สุขญฺจ สุขสหคตาย สทฺธาย ฯเปฯ ปญฺญาย. ราคสฺส. มานสฺส. ทิฏฺฐิยา. ปตฺถนาย. กายิกสฺส สุขสฺส. มคฺคสฺส, ผลสมาปตฺติยา อุปนิสฺสย\-ปจฺจเยน ปจฺจโย. (2)}

\prose{นสุขสหคโต ธมฺโม สุขสหคตสฺส จ นสุข\-สหคตสฺส จ ธมฺมสฺส อุปนิสฺสย\-ปจฺจเยน ปจฺจโย. อารมฺมณู\-ปนิสฺสโย, อนนฺตรู\-ปนิสฺสโย, ปกตู\-ปนิสฺสโย ฯเปฯ. ปกตู\-ปนิสฺสโย\csromandash{} นสุขสหคตํ สทฺธํ อุปนิสฺสาย สุขสหคเตน จิตฺเตน ทานํ เทติ ฯเปฯ. (ทุติยคมน\-สทิสํ.) มานํ ชปฺเปติ, ทิฏฺฐิํ คณฺหาติ. นสุขสหคตํ สีลํ ฯเปฯ. เสนาสนํ สุขํ อุปนิสฺสาย ทานํ เทติ ฯเปฯ สมาปตฺติํ อุปฺปาเทติ. สุขสหคเตน จิตฺเตน อทินฺนํ อาทิยติ ฯเปฯ. นสุขสหคตา สทฺธา ฯเปฯ. เสนาสนํ สุขญฺจ}

\prose{\csromanfolio{116}สุขสหคตาย สทฺธาย ฯเปฯ ปตฺถนาย. กายิกสฺส สุขสฺส, มคฺคสฺส, ผลสมาปตฺติยา สุขสฺส จ อุปนิสฺสย\-ปจฺจเยน ปจฺจโย. (3)}

\prose{สุขสหคโต จ นสุขสหคโต จ ธมฺมา สุขสหคตสฺส ธมฺมสฺส อุปนิสฺสย\-ปจฺจเยน ปจฺจโย. ตีณิ.}

\csromanheader{2}{\textbf{ปุเรชาตาทิ}}
\proseitem{174}{นสุขสหคโต ธมฺโม นสุข\-สหคตสฺส ธมฺมสฺส ปุเรชาต\-ปจฺจเยน ปจฺจโย. อารมฺมณ\-ปุเรชาตํ, วตฺถุ\-ปุเรชาตํ.  อารมฺมณ\-ปุเรชาตํ\csromandash{}จกฺขุํ ฯเปฯ. วตฺถุํ นสุข\-สหคเตน จิตฺเตน อนิจฺจโต ฯเปฯ วิปสฺสติ, อสฺสาเทติ อภินนฺทติ, ตํ อารพฺภ นสุขสหคโต ราโค อุปฺปชฺชติ ฯเปฯ โทมนสฺสํ อุปฺปชฺชติ, สุขํ อุปฺปชฺชติ. ทิพฺเพน จกฺขุนา รูปํ ปสฺสติ ฯเปฯ. โผฏฺฐพฺพา\-ยตนํ กายวิญฺญาณสฺส ฯเปฯ. วตฺถุ\-ปุเรชาตํ\csromandash{}จกฺขายตนํ จกฺขุ\-วิญฺญาณสฺส ฯเปฯ. กายายตนํ กายวิญฺญาณสฺส ฯเปฯ. วตฺถุ นสุข\-สหคตานํ ขนฺธานํ สุขสฺส จ ปุเรชาต\-ปจฺจเยน ปจฺจโย. (1)}

\prose{นสุขสหคโต ธมฺโม สุขสหคตสฺส ธมฺมสฺส ปุเรชาต\-ปจฺจเยน ปจฺจโย. อารมฺมณ\-ปุเรชาตํ, วตฺถุ\-ปุเรชาตํ.  อารมฺมณ\-ปุเรชาตํ\csromandash{}จกฺขุํ ฯเปฯ. วตฺถุํ สุขสหคเตน จิตฺเตน อนิจฺจโต ฯเปฯ วิปสฺสติ, อสฺสาเทติ อภินนฺทติ, ตํ อารพฺภ สุขสหคโต ราโค อุปฺปชฺชติ. ทิฏฺฐิ อุปฺปชฺชติ. วตฺถุ\-ปุเรชาตํ\csromandash{}วตฺถุ สุขสหคตานํ ขนฺธานํ ปุเรชาต\-ปจฺจเยน ปจฺจโย. (2)}

\prose{นสุขสหคโต ธมฺโม สุขสหคตสฺส จ นสุข\-สหคตสฺส จ ธมฺมสฺส ปุเรชาต\-ปจฺจเยน ปจฺจโย. อารมฺมณ\-ปุเรชาตํ, วตฺถุ\-ปุเรชาตํ.  อารมฺมณ\-ปุเรชาตํ\csromandash{}จกฺขุํ ฯเปฯ. วตฺถุํ สุขสหคเตน จิตฺเตน อนิจฺจโต ฯเปฯ วิปสฺสติ, อสฺสาเทติ อภินนฺทติ, ตํ อารพฺภ สุขญฺจ สมฺปยุตฺตกา ขนฺธา จ อุปฺปชฺชนฺติ. วตฺถุ\-ปุเรชาตํ\csromandash{} วตฺถุ สุขสหคตานํ ขนฺธานํ สุขสฺส จ ปุเรชาต\-ปจฺจเยน ปจฺจโย. (3)}

\prose{ปจฺฉาชาต\-ปจฺจเยน ปจฺจโย. ตีณิ. อาเสวน\-ปจฺจเยน ปจฺจโย. นว.}

\csromanheader{2}{92. อุเปกฺขาสหคตทุก \textbf{กมฺมาทิ}}
\proseitem{175}{\csromanfolio{117}สุขสหคโต ธมฺโม สุขสหคตสฺส ธมฺมสฺส กมฺมปจฺจเยน ปจฺจโย. สหชาตา, นานากฺขณิกา.  (ฉ กมฺมานิ จตฺตาริ สหชาตาปิ นานากฺขณิกาปิ กาตพฺพา. เทฺว นานากฺขณิกา.) วิปาก\-ปจฺจเยน ปจฺจโย. นว. อาหารปจฺจเยน ปจฺจโย. จตฺตาริ. อินฺทฺริย\-ปจฺจเยน ปจฺจโย. นว. ฌานปจฺจเยน ปจฺจโย. นว. มคฺคปจฺจเยน ปจฺจโย. จตฺตาริ. สมฺปยุตฺต\-ปจฺจเยน ปจฺจโย. ฉ. วิปฺปยุตฺต\-ปจฺจเยน ปจฺจโย. ปญฺจ. อตฺถิปจฺจเยน ปจฺจโย. นว. นตฺถิ\-ปจฺจเยน ปจฺจโย. นว. วิคต\-ปจฺจเยน ปจฺจโย. นว. อวิคต\-ปจฺจเยน ปจฺจโย. นว.}

\csromanmarkh{1. ปจฺจยานุโลม}\csromanmarkhii{2. สงฺขฺยาวาร}\csromanheadernomark{2}{1. \textbf{ปจฺจยานุโลม 2. สงฺขฺยาวาร}}
\proseitem{176}{เหตุยา จตฺตาริ, อารมฺมเณ นว, อธิปติยา นว, อนนฺตเร นว, สมนนฺตเร นว, สหชาเต นว, อญฺญมญฺเญ นว, นิสฺสเย นว, อุปนิสฺสเย นว, ปุเรชาเต ตีณิ, ปจฺฉาชาเต ตีณิ, อาเสวเน นว, กมฺเม ฉ, วิปาเก นว, อาหาเร จตฺตาริ, อินฺทฺริเย นว, ฌาเน นว, มคฺเค จตฺตาริ, สมฺปยุตฺเต ฉ, วิปฺปยุตฺเต ปญฺจ, อตฺถิยา นว, นตฺถิยา นว, วิคเต นว, อวิคเต นว.}

\prosewithcloser{(เอวํ ปจฺจนียวิภงฺโคปิ คณนาปิ สปฺปีติก\-ทุก\-สทิสํ กาตพฺพํ, ยทิปิ วิมติ อตฺถิ, อนุโลมํ ปสฺสิตฺวา คเณตพฺพํ.)}{สุขสหคตทุกํ นิฏฺฐิตํ.}
\csromanmatikaanchor{mtk.21}
\csromantocmarkref{section}{92. อุเปกฺขาสหคตทุก}{mtk.21}
\bookmark[dest={mtk.21},level=1]{92. อุเปกฺขาสหคตทุก}
\csromanmarkh{92. อุเปกฺขาสหคตทุก}\csromanmarkhii{1. ปฏิจฺจวาร}\csromanheadernomark{1}{92. อุเปกฺขาสหคตทุก 1. ปฏิจฺจวาร}
\proseitem{177}{อุเปกฺขา\-สหคตํ ธมฺมํ ปฏิจฺจ อุเปกฺขา\-สหคโต ธมฺโม อุปฺปชฺชติ เหตุปจฺจยา. อุเปกฺขา\-สหคตํ เอกํ ขนฺธํ ปฏิจฺจ เทฺว ขนฺธา. เทฺว \csromanfolio{118}ขนฺเธ ฯเปฯ. ปฏิสนฺธิ\-กฺขเณ ฯเปฯ. อุเปกฺขา\-สหคตํ ธมฺมํ ปฏิจฺจ นอุเปกฺขา\-สหคโต ธมฺโม อุปฺปชฺชติ เหตุปจฺจยา. อุเปกฺขา\-สหคเต ขนฺเธ ปฏิจฺจ อุเปกฺขา จิตฺต\-สมุฏฺฐานญฺจ รูปํ. ปฏิสนฺธิ\-กฺขเณ ฯเปฯ. อุเปกฺขา\-สหคตํ ธมฺมํ ปฏิจฺจ อุเปกฺขา\-สหคโต จ นอุเปกฺขา\-สหคโต จ ธมฺมา อุปฺปชฺชนฺติ เหตุปจฺจยา. อุเปกฺขา\-สหคตํ เอกํ ขนฺธํ ปฏิจฺจ เทฺว ขนฺธา อุเปกฺขา จ จิตฺต\-สมุฏฺฐานญฺจ รูปํ, เทฺว ขนฺเธ ฯเปฯ. ปฏิสนฺธิ\-กฺขเณ ฯเปฯ. (3)}

\prose{นอุเปกฺขาสหคตํ ธมฺมํ ปฏิจฺจ นอุเปกฺขาสหคโต ธมฺโม อุปฺปชฺชติ เหตุปจฺจยา. นอุเปกฺขาสหคตํ เอกํ ขนฺธํ ปฏิจฺจ ตโย ขนฺธา จิตฺต\-สมุฏฺฐานญฺจ รูปํ ฯเปฯ เทฺว ขนฺเธ ฯเปฯ. อุเปกฺขํ ปฏิจฺจ จิตฺต\-สมุฏฺฐานํ รูปํ. ปฏิสนฺธิ\-กฺขเณ นอุเปกฺขาสหคตํ เอกํ ขนฺธํ ปฏิจฺจ ตโย ขนฺธา กฏตฺตา จ รูปํ ฯเปฯ เทฺว ขนฺเธ ฯเปฯ. ปฏิสนฺธิ\-กฺขเณ อุเปกฺขํ ปฏิจฺจ กฏตฺตารูปํ. ขนฺเธ ปฏิจฺจ วตฺถุ, วตฺถุํ ปฏิจฺจ ขนฺธา. อุเปกฺขํ ปฏิจฺจ วตฺถุ, วตฺถุํ ปฏิจฺจ อุเปกฺขา, เอกํ มหาภูตํ ฯเปฯ. (สปฺปีติก\-ทุก\-สทิสํ. อนุโลเม นวปิ ปญฺหา.)}

\csromanmarkh{1. ปจฺจยานุโลม}\csromanmarkhii{2. สงฺขฺยาวาร}\csromanheadernomark{2}{1. \textbf{ปจฺจยานุโลม 2. สงฺขฺยาวาร}}
\proseitem{178}{เหตุยา นว, อารมฺมเณ นว, อธิปติยา นว ฯเปฯ ปุเรชาเต ฉ, อาเสวเน ฉ, กมฺเม นว, (สพฺพตฺถ นว.) อวิคเต นว.}

\csromanmarkh{2. ปจฺจย\-ปจฺจนีย}\csromanmarkhii{1. วิภงฺควาร}\csromanheadernomark{2}{2. ปจฺจย\-ปจฺจนีย 1. วิภงฺควาร}
\proseitem{179}{อุเปกฺขา\-สหคตํ ธมฺมํ ปฏิจฺจ อุเปกฺขา\-สหคโต ธมฺโม อุปฺปชฺชติ นเหตุปจฺจยา. อเหตุกํ อุเปกฺขา\-สหคตํ เอกํ ขนฺธํ ปฏิจฺจ เทฺว ขนฺธา. เทฺว ขนฺเธ ฯเปฯ. อเหตุก\-ปฏิสนฺธิ\-ขเณ ฯเปฯ. วิจิกิจฺฉา\-สหคเต อุทฺธจฺจ\-สหคเต ขนฺเธ ปฏิจฺจ วิจิกิจฺฉา\-สหคโต อุทฺธจฺจ\-สหคโต โมโห.  อุเปกฺขา\-สหคตํ ธมฺมํ ปฏิจฺจ นอุเปกฺขา\-สหคโต ธมฺโม อุปฺปชฺชติ นเหตุปจฺจยา. อเหตุเก อุเปกฺขา\-สหคเต ขนฺเธ ปฏิจฺจ อุเปกฺขา จิตฺต\-สมุฏฺฐานญฺจ รูปํ. อเหตุก\-ปฏิสนฺธิ\-กฺขเณ ฯเปฯ. อุเปกฺขา\-สหคตํ \csromanfolio{119}ธมฺมํ ปฏิจฺจ อุเปกฺขา\-สหคโต จ นฺเอาเปกฺขา\-สหคโต จ ธมฺมา อุปฺปชฺชนฺติ นเหตุปจฺจยา. อเหตุกํ อุเปกฺขา\-สหคตํ เอกํ ขนฺธํ ปฏิจฺจ เทฺว ขนฺธา อุเปกฺขา จ จิตฺต\-สมุฏฺฐานญฺจ รูปํ. เทฺว ขนฺเธ ฯเปฯ. อเหตุก\-ปฏิสนฺธิ\-กฺขเณ ฯเปฯ. (3)}

\proseitem{180}{นอุเปกฺขาสหคตํ ธมฺมํ ปฏิจฺจ นอุเปกฺขาสหคโต ธมฺโม อุปฺปชฺชติ นเหตุปจฺจยา. อเหตุกํ นอุเปกฺขาสหคตํ เอกํ ขนฺธํ ปฏิจฺจ ตโย ขนฺธา จิตฺต\-สมุฏฺฐานญฺจ รูปํ ฯเปฯ เทฺว ขนฺเธ ฯเปฯ. อเหตุกํ อุเปกฺขํ ปฏิจฺจ จิตฺต\-สมุฏฺฐานํ รูปํ. อเหตุก\-ปฏิสนฺธิ\-กฺขเณ ฯเปฯ. อุเปกฺขํ ปฏิจฺจ วตฺถุ, วตฺถุํ ปฏิจฺจ อุเปกฺขา. เอกํ มหาภูตํ ฯเปฯ. (ยาว อสญฺญสตฺตา.) (1)}

\prose{นอุเปกฺขาสหคตํ ธมฺมํ ปฏิจฺจ อุเปกฺขา\-สหคโต ธมฺโม อุปฺปชฺชติ นเหตุปจฺจยา. อเหตุกํ อุเปกฺขํ ปฏิจฺจ สมฺปยุตฺตกา ขนฺธา. อเหตุก\-ปฏิสนฺธิ\-กฺขเณ อุเปกฺขํ ปฏิจฺจ สมฺปยุตฺตกา ขนฺธา. อเหตุก\-ปฏิสนฺธิ\-กฺขเณ วตฺถุํ ปฏิจฺจ อุเปกฺขา\-สหคตา ขนฺธา. วิจิกิจฺฉา\-สหคตํ อุทฺธจฺจ\-สหคตํ อุเปกฺขํ ปฏิจฺจ วิจิกิจฺฉา\-สหคโต อุทฺธจฺจ\-สหคโต โมโห. (2)}

\prose{นอุเปกฺขาสหคตํ ธมฺมํ ปฏิจฺจ อุเปกฺขา\-สหคโต จ นอุเปกฺขา\-สหคโต จ ธมฺมา อุปฺปชฺชนฺติ นเหตุปจฺจยา. อเหตุกํ อุเปกฺขํ ปฏิจฺจ สมฺปยุตฺตกา ขนฺธา จิตฺต\-สมุฏฺฐานญฺจ รูปํ. อเหตุกํ อุเปกฺขํ ปฏิจฺจ สมฺปยุตฺตกา ขนฺธา. มหาภูเต ปฏิจฺจ จิตฺต\-สมุฏฺฐานํ รูปํ. อเหตุก\-ปฏิสนฺธิ\-กฺขเณ อุเปกฺขํ ปฏิจฺจ สมฺปยุตฺตกา ขนฺธา กฏตฺตา จ รูปํ. อเหตุก\-ปฏิสนฺธิ\-กฺขเณ อุเปกฺขํ ปฏิจฺจ สมฺปยุตฺตกา ขนฺธา. มหาภูเต ปฏิจฺจ กฏตฺตารูปํ. อเหตุก\-ปฏิสนฺธิ\-กฺขเณ วตฺถุํ ปฏิจฺจ อุเปกฺขา\-สหคตา ขนฺธา. มหาภูเต ปฏิจฺจ กฏตฺตารูปํ. อเหตุก\-ปฏิสนฺธิ\-กฺขเณ วตฺถุํ ปฏิจฺจ อุเปกฺขา จ สมฺปยุตฺตกา จ ขนฺธา. (3)}

\proseitem{181}{อุเปกฺขาสหคตญฺจ นอุเปกฺขาสหคตญฺจ ธมฺมํ ปฏิจฺจ อุเปกฺขา\-สหคโต ธมฺโม อุปฺปชฺชติ นเหตุปจฺจยา. อเหตุกํ อุเปกฺขา\-สหคตํ เอกํ ขนฺธญฺจ อุเปกฺขญฺจ ปฏิจฺจ เทฺว ขนฺธา. เทฺว ขนฺเธ ฯเปฯ. อเหตุก\-ปฏิสนฺธิ\-กฺขเณ อุเปกฺขา\-สหคตํ เอกํ ขนฺธญฺจ อุเปกฺขญฺจ ปฏิจฺจ เทฺว ขนฺธา. เทฺว ขนฺเธ ฯเปฯ. อเหตุก\-ปฏิสนฺธิ\-กฺขเณ อุเปกฺขา\-สหคตํ เอกํ ขนฺธญฺจ วตฺถุญฺจ ปฏิจฺจ เทฺว ขนฺธา. เทฺว ขนฺเธ ฯเปฯ. วิจิกิจฺฉา\-สหคเต อุทฺธจฺจ\-สหคเต ขนฺเธ จ อุเปกฺขญฺจ ปฏิจฺจ วิจิกิจฺฉา\-สหคโต อุทฺธจฺจ\-สหคโต โมโห. (1)}

\prose{\csromanfolio{120}อุเปกฺขาสหคตญฺจ นอุเปกฺขาสหคตญฺจ ธมฺมํ ปฏิจฺจ นอุเปกฺขาสหคโต ธมฺโม อุปฺปชฺชติ นเหตุปจฺจยา. อเหตุเก อุเปกฺขา\-สหคเต ขนฺเธ จ อุเปกฺขญฺจ ปฏิจฺจ จิตฺต\-สมุฏฺฐานํ รูปํ. อเหตุเก อุเปกฺขา\-สหคเต ขนฺเธ จ มหาภูเต จ ปฏิจฺจ จิตฺต\-สมุฏฺฐานํ รูปํ. อเหตุก\-ปฏิสนฺธิ\-กฺขเณ อุเปกฺขา\-สหคเต ขนฺเธ จ อุเปกฺขญฺจ ปฏิจฺจ กฏตฺตารูปํ. อเหตุก\-ปฏิสนฺธิ\-กฺขเณ อุเปกฺขา\-สหคเต ขนฺเธ จ มหาภูเต จ ปฏิจฺจ กฏตฺตารูปํ. อเหตุก\-ปฏิสนฺธิ\-กฺขเณ อุเปกฺขา\-สหคเต ขนฺเธ จ วตฺถุญฺจ ปฏิจฺจ อุเปกฺขา. (2)}

\prose{อุเปกฺขาสหคตญฺจ นอุเปกฺขาสหคตญฺจ ธมฺมํ ปฏิจฺจ อุเปกฺขา\-สหคโต จ นอุเปกฺขา\-สหคโต จ ธมฺมา อุปฺปชฺชนฺติ นเหตุปจฺจยา. อเหตุกํ อุเปกฺขา\-สหคตํ เอกํ ขนฺธญฺจ อุเปกฺขญฺจ ปฏิจฺจ เทฺว ขนฺธา จิตฺต\-สมุฏฺฐานญฺจ รูปํ. เทฺว ขนฺเธ ฯเปฯ. อเหตุกํ อุเปกฺขา\-สหคตํ เอกํ ขนฺธญฺจ อุเปกฺขญฺจ ปฏิจฺจ เทฺว ขนฺธา. เทฺว ขนฺเธ ฯเปฯ. อเหตุเก อุเปกฺขา\-สหคเต ขนฺเธ จ อุเปกฺขญฺจ มหาภูเต จ ปฏิจฺจ จิตฺต\-สมุฏฺฐานํ รูปํ. อเหตุก\-ปฏิสนฺธิ\-กฺขเณ อุเปกฺขา\-สหคตํ เอกํ ขนฺธญฺจ อุเปกฺขญฺจ ปฏิจฺจ เทฺว ขนฺธา กฏตฺตา จ รูปํ. เทฺว ขนฺเธ ฯเปฯ. อเหตุก\-ปฏิสนฺธิ\-กฺขเณ อุเปกฺขา\-สหคตํ เอกํ ขนฺธญฺจ อุเปกฺขญฺจ ปฏิจฺจ เทฺว ขนฺธา. เทฺว ขนฺเธ ฯเปฯ. อเหตุเก อุเปกฺขา\-สหคเต ขนฺเธ จ อุเปกฺขญฺจ มหาภูเต จ ปฏิจฺจ กฏตฺตารูปํ. อเหตุก\-ปฏิสนฺธิ\-กฺขเณ อุเปกฺขา\-สหคตํ เอกํ ขนฺธญฺจ วตฺถุญฺจ ปฏิจฺจ เทฺว ขนฺธา อุเปกฺขา จ. เทฺว ขนฺเธ ฯเปฯ. (สํขิตฺตํ.) (3)}

\csromanmarkh{2. ปจฺจย\-ปจฺจนีย}\csromanmarkhii{2. สงฺขฺยาวาร}\csromanheadernomark{2}{2. ปจฺจย\-ปจฺจนีย 2. สงฺขฺยาวาร}
\proseitem{182}{นเหตุยา นว, นอารมฺมเณ ตีณิ, นอธิปติยา นว, นอนนฺตเร ตีณิ, นอญฺญมญฺเญ ตีณิ, นอุปนิสฺสเย ตีณิ, นปุเรชาเต นว, นปจฺฉาชาเต นว, นอาเสวเน นว, นกมฺเม จตฺตาริ, นวิปาเก นว, นอาหาเร เอกํ, นอินฺทฺริเย เอกํ, นฌาเน นว, นมคฺเค นว, นสมฺปยุตฺเต ตีณิ, นวิปฺปยุตฺเต ฉ, โนนตฺถิยา ตีณิ, โนวิคเต ตีณิ.}

\vspace{\tipitakaparskip}
\csromancenter{(เอวํ อิตเร เทฺว คณนาปิ สหชาตวาโรปิ กาตพฺโพ.)}
\csromanmarkh{92. อุเปกฺขาสหคตทุก}\csromanmarkhii{92. อุเปกฺขาสหคตทุก 3. ปจฺจยวาร}\csromanheadernomark{2}{92. อุเปกฺขาสหคตทุก \textbf{92. อุเปกฺขาสหคตทุก 3. ปจฺจยวาร}}
\proseitem{183}{\csromanfolio{121}อุเปกฺขา\-สหคตํ ธมฺมํ ปจฺจยา อุเปกฺขา\-สหคโต ธมฺโม อุปฺปชฺชติ เหตุปจฺจยา. อุเปกฺขา\-สหคตํ เอกํ ขนฺธํ ปจฺจยา เทฺว ขนฺธา. เทฺว ขนฺเธ ฯเปฯ. ปฏิสนฺธิ\-กฺขเณ ฯเปฯ. (ยถา สวิตกฺก\-ทุก\-สทิสํ, ปจฺจยวาเร นานากรณํ, “อุเปกฺขนฺ”ติ นวปิ ปญฺหา กาตพฺพา. ปฏิสนฺธิปวตฺติปิ วตฺถุปิ.)}

\prose{\textbf{เหตุ}ยา นว, อารมฺมเณ นว ฯเปฯ ปุเรชาเต นว, อาเสวเน นว, (สพฺพตฺถ นว.) อวิคเต นว. อนุโลมํ.}

\proseitem{184}{อุเปกฺขา\-สหคตํ ธมฺมํ ปจฺจยา อุเปกฺขา\-สหคโต ธมฺโม อุปฺปชฺชติ นเหตุปจฺจยา. อเหตุกํ อุเปกฺขา\-สหคตํ เอกํ ขนฺธํ ปจฺจยา เทฺว ขนฺธา. เทฺว ขนฺเธ ฯเปฯ. อเหตุก\-ปฏิสนฺธิ\-กฺขเณ ฯเปฯ. วิจิกิจฺฉา\-สหคเต อุทฺธจฺจ\-สหคเต ขนฺเธ ปจฺจยา วิจิกิจฺฉา\-สหคโต อุทฺธจฺจ\-สหคโต โมโห. (เอวํ นวปิ ปญฺหา ปวตฺติ\-ปฏิสนฺธิโย ยถา สวิตกฺก\-ทุกสฺส เอวํ กาตพฺพา. ตีณิเยว โมโห. ปวตฺเต วตฺถุปิ กาตพฺพา.)}

\prose{นเหตุยา นว, นอารมฺมเณ ตีณิ, นอธิปติยา นว, นอนนฺตเร ตีณิ, นสมนนฺตเร ตีณิ, นอุปนิสฺสเย ตีณิ, นปุเรชาเต นว, นปจฺฉาชาเต นว, นอาเสวเน นว, นกมฺเม จตฺตาริ, นวิปาเก นว, นอาหาเร เอกํ, นอินฺทฺริเย เอกํ, นฌาเน เอกํ, นมคฺเค นว, นสมฺปยุตฺเต ตีณิ, นวิปฺปยุตฺเต ฉ, โนนตฺถิยา ตีณิ, โนวิคเต ตีณิ.}

\vspace{\tipitakaparskip}
\csromancenter{ปจฺจนียํ.}
\csromancenter{(เอวํ อิตเร เทฺว คณนาปิ นิสฺสยวาโรปิ กาตพฺโพ.)}
\csromanmarkh{92. อุเปกฺขาสหคตทุก}\csromanmarkhii{5. สํสฏฺฐวาร}\csromanheadernomark{2}{92. อุเปกฺขาสหคตทุก 5. สํสฏฺฐวาร}
\proseitem{185}{อุเปกฺขา\-สหคตํ ธมฺมํ สํสฏฺโฐ อุเปกฺขา\-สหคโต ธมฺโม อุปฺปชฺชติ เหตุปจฺจยา. อุเปกฺขา\-สหคตํ เอกํ ขนฺธํ สํสฏฺฐา เทฺว ขนฺธา. เทฺว \csromanfolio{122}ขนฺเธ ฯเปฯ. ปฏิสนฺธิ\-กฺขเณ ฯเปฯ.}

\vspace{\tipitakaparskip}
\csromancenter{(ยถา สวิตกฺกทุกํ สมฺปยุตฺตวาโร, เอวํ กาตพฺโพ.)}
\prose{เหตุยา ฉ, อารมฺมเณ ฉ, (สพฺพตฺถ ฉ.) อวิคเต ฉ. อนุโลมํ.}

\prose{อุเปกฺขา\-สหคตํ ธมฺมํ สํสฏฺโฐ อุเปกฺขา\-สหคโต ธมฺโม อุปฺปชฺชติ นเหตุปจฺจยา. อเหตุกํ อุเปกฺขา\-สหคตํ เอกํ ขนฺธํ สํสฏฺฐา เทฺว ขนฺธา. เทฺว ขนฺเธ ฯเปฯ. อเหตุก\-ปฏิสนฺธิ\-กฺขเณ ฯเปฯ. วิจิกิจฺฉา\-สหคเต อุทฺธจฺจ\-สหคเต ขนฺเธ สํสฏฺโฐ วิจิกิจฺฉา\-สหคโต อุทฺธจฺจ\-สหคโต โมโห.}

\vspace{\tipitakaparskip}
\csromancenter{(เอวํ ปญฺจ ปญฺหา ยถา สวิตกฺกทุเก, เอวํ กาตพฺพา.)}
\prose{นเหตุยา ฉ, นอธิปติยา ฉ, นปุเรชาเต ฉ, นปจฺฉาชาเต ฉ, นอาเสวเน ฉ, นกมฺเม จตฺตาริ, นวิปาเก ฉ, นวิปฺปยุตฺเต ฉ.}

\vspace{\tipitakaparskip}
\csromancenter{ปจฺจนียํ.}
\csromancenter{(เอวํ อิตเร เทฺว คณนาปิ สมฺปยุตฺตวาโรปิ กาตพฺโพ.)}
\csromanmarkh{92. อุเปกฺขาสหคตทุก}\csromanmarkhii{7. ปญฺหาวาร}\csromanheadernomark{2}{92. \textbf{อุเปกฺขาสหคตทุก 7. ปญฺหาวาร}}
\prose{เหตุอาทิ}

\proseitem{186}{อุเปกฺขา\-สหคโต ธมฺโม อุเปกฺขา\-สหคตสฺส ธมฺมสฺส เหตุปจฺจเยน ปจฺจโย. อุเปกฺขา\-สหคตา เหตู สมฺปยุตฺตกานํ ขนฺธานํ เหตุปจฺจเยน ปจฺจโย. ปฏิสนฺธิ\-กฺขเณ ฯเปฯ. (เอวํ จตฺตาริ ปญฺหา ยถา สวิตกฺกทุกสฺส.)}

\vspace{\tipitakaparskip}
\csromancenter{(เอวํ จตฺตาริ ปญฺหา ยถา สวิตกฺก\-ทุกสฺส.)}
\prose{อุเปกฺขา\-สหคโต ธมฺโม อุเปกฺขา\-สหคตสฺส ธมฺมสฺส อารมฺมณ\-ปจฺจเยน ปจฺจโย. อธิปติ\-ปจฺจเยน ปจฺจโย.}

\vspace{\tipitakaparskip}
\csromancenter{(ยถา สปฺปีติกทุกํ, เอวํ อารมฺมณมฺปิ อธิปติปิ วิตฺถาเรตพฺพา, “อุเปกฺขา”ติ นานํ.)}
\csromanheader{2}{92. อุเปกฺขาสหคตทุก \textbf{อนนฺตราทิ}}
\proseitem{187}{\csromanfolio{123}อุเปกฺขา\-สหคโต ธมฺโม อุเปกฺขา\-สหคตสฺส ธมฺมสฺส อนนฺตร\-ปจฺจเยน ปจฺจโย. ปุริมา ปุริมา อุเปกฺขา\-สหคตา ขนฺธา ปจฺฉิมานํ ปจฺฉิมานํ อุเปกฺขา\-สหคตานํ ขนฺธานํ อนนฺตร\-ปจฺจเยน ปจฺจโย. (1)}

\prose{อุเปกฺขา\-สหคโต ธมฺโม นอุเปกฺขาสหคตสฺส ธมฺมสฺส อนนฺตร\-ปจฺจเยน ปจฺจโย. ปุริมา ปุริมา อุเปกฺขา\-สหคตา ขนฺธา ปจฺฉิมาย ปจฺฉิมาย อุเปกฺขาย อนนฺตร\-ปจฺจเยน ปจฺจโย. อุเปกฺขา\-สหคตํ จุติจิตฺตํ นอุเปกฺขาสหคตสฺส อุปปตฺติ\-จิตฺตสฺส. อาวชฺชนา นอุเปกฺขา\-สหคตานํ ขนฺธานํ. วิปาก\-มโนธาตุ นอุเปกฺขา\-สหคตาย วิปาก\-มโนวิญฺญาณ\-ธาตุยา. อุเปกฺขา\-สหคตํ ภวงฺคํ นอุเปกฺขาสหคตสฺส ภวงฺคสฺส. อุเปกฺขา\-สหคตํ กุสลากุสลํ นอุเปกฺขาสหคตสฺส วุฏฺฐานสฺส. กิริยํ วุฏฺฐานสฺส. ผลํ วุฏฺฐานสฺส. นิโรธา วุฏฺฐหนฺตสฺส เนว\-สญฺญา\-นา\-สญฺญา\-ยตนํ นอุเปกฺขา\-สหคตาย ผลสมาปตฺติยา อนนฺตร\-ปจฺจเยน ปจฺจโย. (2)}

\prose{อุเปกฺขา\-สหคโต ธมฺโม อุเปกฺขา\-สหคตสฺส จ นอุเปกฺขา\-สหคตสฺส จ ธมฺมสฺส อนนฺตร\-ปจฺจเยน ปจฺจโย. ปุริมา ปุริมา อุเปกฺขา\-สหคตา ขนฺธา ปจฺฉิมานํ ปจฺฉิมานํ อุเปกฺขา\-สหคตานํ ขนฺธานํ อุเปกฺขาย จ อนนฺตร\-ปจฺจเยน. ปจฺจโย. (3)}

\proseitem{188}{นอุเปกฺขาสหคโต ธมฺโม นอุเปกฺขา\-สหคตสฺส ธมฺมสฺส อนนฺตร\-ปจฺจเยน ปจฺจโย. ปุริมา ปุริมา อุเปกฺขา ปจฺฉิมาย ปจฺฉิมาย อุเปกฺขาย อนนฺตร\-ปจฺจเยน ปจฺจโย. ปุริมา ปุริมา นอุเปกฺขาสหคตา ขนฺธา ปจฺฉิมานํ ปจฺฉิมานํ นอุเปกฺขา\-สหคตานํ ขนฺธานํ อนนฺตร\-ปจฺจเยน ปจฺจโย. อนุโลมํ โคตฺรภุสฺส ฯเปฯ. อนุโลมํ ผลสมาปตฺติยา อนนฺตร\-ปจฺจเยน ปจฺจโย. (มูลํ กาตพฺพํ.) ปุริมา ปุริมา อุเปกฺขา ปจฺฉิมานํ ปจฺฉิมานํ อุเปกฺขา\-สหคตานํ ขนฺธานํ อนนฺตร\-ปจฺจเยน ปจฺจโย. นอุเปกฺขาสหคตํ จุติจิตฺตํ อุเปกฺขา\-สหคตสฺส อุปปตฺติ\-จิตฺตสฺส. นอุปกฺขาสหคตํ ภวงฺคํ อาวชฺชนาย. กายวิญฺญาณ\-ธาตุ วิปาก\-มโนธาตุยา. นอุเปกฺขาสหคตา วิปาก\-มโนวิญฺญาณ\-ธาตุ กิริย\-มโนวิญฺญาณ\-ธาตุยา. \csromanfolio{124}นอุเปกฺขาสหคตํ ภวงฺคํ อุเปกฺขา\-สหคตสฺส ภวงฺคสฺส. นอุเปกฺขาสหคตํ กุสลากุสลํ อุเปกฺขา\-สหคตสฺส วุฏฺฐานสฺส. กิริยํ วุฏฺฐานสฺส. ผลํ วุฏฺฐานสฺส อนนฺตร\-ปจฺจเยน ปจฺจโย. (2)}

\prose{นอุเปกฺขาสหคโต ธมฺโม อุเปกฺขา\-สหคตสฺส จ นอุเปกฺขา\-สหคตสฺส จ ธมฺมสฺส อนนฺตร\-ปจฺจเยน ปจฺจโย. ปุริมา ปุริมา อุเปกฺขา ปจฺฉิมานํ ปจฺฉิมานํ อุเปกฺขา\-สหคตานํ ขนฺธานํ อุเปกฺขาย จ อนนฺตร\-ปจฺจเยน ปจฺจโย. (3)}

\proseitem{189}{อุเปกฺขา\-สหคโต จ นอุเปกฺขา\-สหคโต จ ธมฺมา อุเปกฺขา\-สหคตสฺส ธมฺมสฺส อนนฺตร\-ปจฺจเยน ปจฺจโย. ปุริมา ปุริมา อุเปกฺขา\-สหคตา ขนฺธา จ อุเปกฺขา จ ปจฺฉิมานํ ปจฺฉิมานํ อุเปกฺขา\-สหคตานํ ขนฺธานํ อนนฺตร\-ปจฺจเยน ปจฺจโย. (มูลํ กาตพฺพํ.) ปุริมา ปุริมา อุเปกฺขา\-สหคตา ขนฺธา จ อุเปกฺขา จ ปจฺฉิมาย ปจฺฉิมาย อุเปกฺขาย อนนฺตร\-ปจฺจเยน ปจฺจโย. อุเปกฺขา\-สหคตํ จุติจิตฺตญฺจ อุเปกฺขา จ นอุเปกฺขา\-สหคตสฺส อุปปตฺติ\-จิตฺตสฺส. อาวชฺชนา จ อุเปกฺขา จ นอุเปกฺขา\-สหคตานํ ขนฺธานํ. วิปาก\-มโนธาตุ จ อุเปกฺขา จ นอุเปกฺขา\-สหคตาย วิปาก\-มโนวิญฺญาณ\-ธาตุยา. อุเปกฺขา\-สหคตํ ภวงฺคญฺจ อุเปกฺขา จ นอุเปกฺขา\-สหคตสฺส ภวงฺคสฺส. อุเปกฺขา\-สหคตํ กุสลากุสลญฺจ อุเปกฺขา จ นอุเปกฺขา\-สหคตสฺส วุฏฺฐานสฺส. กิริยญฺจ อุเปกฺขา จ วุฏฺฐานสฺส. ผลญฺจ อุเปกฺขา จ วุฏฺฐานสฺส. นิโรธา วุฏฺฐหนฺตสฺส เนว\-สญฺญา\-นา\-สญฺญา\-ยตนญฺจ อุเปกฺขา จ นอุเปกฺขา\-สหคตาย ผลสมาปตฺติยา อนนฺตร\-ปจฺจเยน ปจฺจโย. (มูลํ กาตพฺพํ.) ปุริมา ปุริมา อุเปกฺขา\-สหคตา ขนฺธา จ อุเปกฺขา จ ปจฺฉิมานํ ปจฺฉิมานํ อุเปกฺขา\-สหคตานํ ขนฺธานํ อุเปกฺขาย จ อนนฺตร\-ปจฺจเยน ปจฺจโย. (3)}

\prose{สมนนฺตร\-ปจฺจเยน ปจฺจโย. สหชาต\-ปจฺจเยน ปจฺจโย. นว. อญฺญมญฺญ\-ปจฺจเยน ปจฺจโย. นว. นิสฺสย\-ปจฺจเยน ปจฺจโย. นว.}

\csromanheader{2}{\textbf{อุปนิสฺสย}}
\proseitem{190}{อุเปกฺขา\-สหคโต ธมฺโม อุเปกฺขา\-สหคตสฺส ธมฺมสฺส อุปนิสฺสย\-ปจฺจเยน ปจฺจโย. ตีณิ.}

\prose{\csromanfolio{125}นอุเปกฺขาสหคโต ธมฺโม นอุเปกฺขาสหคตสฺส ธมฺมสฺส อุปนิสฺสย\-ปจฺจเยน ปจฺจโย. อารมฺมณู\-ปนิสฺสโย, อนนฺตรู\-ปนิสฺสโย, ปกตู\-ปนิสฺสโย ฯเปฯ. ปกตู\-ปนิสฺสโย\csromandash{}นอุเปกฺขาสหคตํ สทฺธํ อุปนิสฺสาย นอุเปกฺขาสหคเตน จิตฺเตน ทานํ เทติ. สีลํ ฯเปฯ อุโปสถกมฺมํ ฯเปฯ. นอุเปกฺขาสหคตํ ฌานํ ฯเปฯ วิปสฺสนํ ฯเปฯ มคฺคํ ฯเปฯ สมาปตฺติํ อุปฺปาเทติ. มานํ ชปฺเปติ. ทิฏฺฐิํ คณฺหาติ. นอุเปกฺขาสหคตํ สีลํ ฯเปฯ ปญฺญํ. ราคํ. โทสํ. โมหํ. มานํ. ทิฏฺฐิํ. ปตฺถนํ. กายิกํ สุขํ. กายิกํ ทุกฺขํ. อุตุํ. โภชนํ. เสนาสนํ อุเปกฺขํ อุปนิสฺสาย นอุเปกฺขาสหคเตน จิตฺเตน ทานํ เทติ ฯเปฯ สมาปตฺติํ อุปฺปาเทติ. ปาณํ หนติ ฯเปฯ สํฆํ ภินฺทติ. นอุเปกฺขาสหคตา สทฺธา ฯเปฯ. เสนาสนํ อุเปกฺขา จ นอุเปกฺขาสหคตาย สทฺธาย ฯเปฯ ปญฺญาย. ราคสฺส ฯเปฯ ปตฺถนาย. กายิกสฺส สุขสฺส. กายิกสฺส ทุกฺขสฺส. มคฺคสฺส, ผลสมาปตฺติยา อุเปกฺขาย จ อุปนิสฺสย\-ปจฺจเยน ปจฺจโย. (1)}

\prose{นอุเปกฺขาสหคโต ธมฺโม อุเปกฺขา\-สหคตสฺส ธมฺมสฺส อุปนิสฺสย\-ปจฺจเยน ปจฺจโย. (ตีณิปิ อุปนิสฺสยา.) นอุเปกฺขาสหคตํ สทฺธํ อุปนิสฺสาย อุเปกฺขา\-สหคเตน จิตฺเตน ทานํ เทติ ฯเปฯ สมาปตฺติํ อุปฺปาเทติ. มานํ ชปฺเปติ. ทิฏฺฐิํ คณฺหาติ. นอุเปกฺขาสหคตํ สีลํ ฯเปฯ. เสนาสนํ อุเปกฺขํ อุปนิสฺสาย อุเปกฺขา\-สหคเตน จิตฺเตน ทานํ เทติ ฯเปฯ สมาปตฺติํ อุปฺปาเทติ. อุเปกฺขา\-สหคเตน จิตฺเตน อทินฺนํ อาทิยติ. มุสา ภณติ. ปิสุณํ ฯเปฯ สมฺผํ ฯเปฯ สนฺธิํ ฯเปฯ นิลฺโลปํ ฯเปฯ เอกาคาริกํ ฯเปฯ ปริปนฺเถ ฯเปฯ ปรทารํ ฯเปฯ คามฆาตํ ฯเปฯ นิคมฆาตํ กโรติ. นอุเปกฺขาสหคตา สทฺธา ฯเปฯ. เสนาสนํ อุเปกฺขา จ อุเปกฺขา\-สหคตาย สทฺธาย ฯเปฯ ปตฺถนาย. มคฺคสฺส, ผลสมาปตฺติยา อุปนิสฺสย\-ปจฺจเยน ปจฺจโย.}

\prose{นอุเปกฺขาสหคโต ธมฺโม อุเปกฺขา\-สหคตสฺส จ นอุเปกฺขา\-สหคตสฺส จ ธมฺมสฺส อุเปนิสฺสยปจฺจเยน ปจฺจโย. (ตีณิ อุปนิสฺสยา, ทุติยคมน\-สทิสา.) (3)}

\prose{อุเปกฺขา\-สหคโต จ นอุเปกฺขา\-สหคโต จ ธมฺมา อุเปกฺขา\-สหคตสฺส ธมฺมสฺส อุปนิสฺสย\-ปจฺจเยน ปจฺจโย. ตีณิ.}

\csromanheader{2}{\textbf{ปุเรชาต}}
\proseitem{191}{\csromanfolio{126}นอุเปกฺขาสหคโต ธมฺโม นอุเปกฺขาสหคตสฺส ธมฺมสฺส ปุเรชาต\-ปจฺจเยน ปจฺจโย. ตีณิ. (สปฺปีติก\-ทุก\-สทิสา.)}

\csromanheader{2}{\textbf{ปจฺฉาชาตาทิ}}
\proseitem{192}{อุเปกฺขา\-สหคโต ธมฺโม นอุเปกฺขาสหคตสฺส ธมฺมสฺส ปจฺฉาชาต\-ปจฺจเยน ปจฺจโย. ตีณิ. อาเสวน\-ปจฺจเยน ปจฺจโย. นว. กมฺมปจฺจเยน ปจฺจโย. ฉ. (จตฺตาริ สหชาตา นานากฺขณิกา กาตพฺพา. เทฺว นานากฺขณิกา จ.) วิปาก\-ปจฺจเยน ปจฺจโย. นว. อาหารปจฺจเยน ปจฺจโย. จตฺตาริ. อินฺทฺริย\-ปจฺจเยน ปจฺจโย. นว. ฌานปจฺจเยน ปจฺจโย. นว. มคฺคปจฺจเยน ปจฺจโย. จตฺตาริ. สมฺปยุตฺต\-ปจฺจเยน ปจฺจโย. ฉ. วิปฺปยุตฺต\-ปจฺจเยน ปจฺจโย. ปญฺจ. อตฺถิปจฺจเยน ปจฺจโย. นว. นตฺถิ\-ปจฺจเยน ปจฺจโย. วิคต\-ปจฺจเยน ปจฺจโย. อวิคต\-ปจฺจเยน ปจฺจโย.}

\vspace{\tipitakaparskip}
\csromancenter{(อิมานิ ปจฺจยานิ สปฺปีติก\-กรเณน วิภชิตพฺพานิ.)}
\csromanmarkh{1. ปจฺจยานุโลม}\csromanmarkhii{2. สงฺขฺยาวาร}\csromanheadernomark{2}{1. \textbf{ปจฺจยานุโลม 2. สงฺขฺยาวาร}}
\proseitem{193}{เหตุยา จตฺตาริ, อารมฺมเณ นว, อธิปติยา นว, อนนฺตเร นว, สมนนฺตเร นว, สหชาเต นว, อญฺญมญฺเญ นว, นิสฺสเย นว, อุปนิสฺสเย นว, ปุเรชาเต ตีณิ, ปจฺฉาชาเต ตีณิ, อาเสวเน นว, กมฺเม ฉ, วิปาเก นว, อาหาเร จตฺตาริ, อินฺทฺริเย นว, ฌาเน นว, มคฺเค จตฺตาริ, สมฺปยุตฺเต ฉ, วิปฺปยุตฺเต ปญฺจ, อตฺถิยา นว, นตฺถิยา นว, วิคเต นว, อวิคเต นว.}

\prosewithcloser{(เอวํ ปจฺจนียวิภงฺโคปิ อิตเร ตีณิ คณนาปิ สปฺปีติก\-ทุก\-สทิสา กาตพฺพา.)}{อุเปกฺขาสหคตทุกํ นิฏฺฐิตํ.}
\csromanmatikaanchor{mtk.22}
\csromantocmarkref{section}{93. กามาวจรทุก}{mtk.22}
\bookmark[dest={mtk.22},level=1]{93. กามาวจรทุก}
\csromanmarkh{93. กามาวจรทุก}\csromanmarkhii{93. กามาวจรทุก 1. ปฏิจฺจวาร}\csromanheadernomark{1}{93. กามาวจรทุก \textbf{93. กามาวจรทุก 1. ปฏิจฺจวาร}}
\proseitem{194}{\csromanfolio{127}กามาวจรํ ธมฺมํ ปฏิจฺจ กามาวจโร ธมฺโม อุปฺปชฺชติ เหตุปจฺจยา. กามาวจรํ เอกํ ขนฺธํ ปฏิจฺจ ตโย ขนฺธา จิตฺต\-สมุฏฺฐานญฺจ รูปํ ฯเปฯ. เทฺว ขนฺเธ ฯเปฯ. ปฏิสนฺธิ\-กฺขเณ ฯเปฯ. เอกํ มหาภูตํ ฯเปฯ. (สํขิตฺตํ.). กามาวจรํ ธมฺมํ ปฏิจฺจ นกามาวจโร ธมฺโม อุปฺปชฺชติ เหตุปจฺจยา. ปฏิสนฺธิ\-กฺขเณ วตฺถุํ ปฏิจฺจ นกามาวจรา ขนฺธา.  กามาวจรํ ธมฺมํ ปฏิจฺจ กามาวจโร จ นกามาวจโร จ ธมฺมา อุปฺปชฺชนฺติ เหตุปจฺจยา. ปฏิสนฺธิ\-กฺขเณ วตฺถุํ ปฏิจฺจ นกามาวจรา ขนฺธา. มหาภูเต ปฏิจฺจ กฏตฺตารูปํ. (3)}

\proseitem{195}{นกามาวจรํ ธมฺมํ ปฏิจฺจ นกามาวจโร ธมฺโม อุปฺปชฺชติ เหตุปจฺจยา. นกามาวจรํ เอกํ ขนฺธํ ปฏิจฺจ ตโย ขนฺธา ฯเปฯ เทฺว ขนฺเธ ฯเปฯ. ปฏิสนฺธิ\-กฺขเณ ฯเปฯ. นกามาวจรํ ธมฺมํ ปฏิจฺจ กามาวจโร ธมฺโม อุปฺปชฺชติ เหตุปจฺจยา. นกามาวจเร ขนฺเธ ปฏิจฺจ จิตฺต\-สมุฏฺฐานํ รูปํ. ปฏิสนฺธิ\-กฺขเณ นกามาวจเร ขนฺเธ ปฏิจฺจ กฏตฺตารูปํ.  นกามาวจรํ ธมฺมํ ปฏิจฺจ กามาวจโร จ นกามาวจโร จ ธมฺมา อุปฺปชฺชนฺติ เหตุปจฺจยา. นกามาวจรํ เอกํ ขนฺธํ ปฏิจฺจ ตโย ขนฺธา จิตฺต\-สมุฏฺฐานญฺจ รูปํ ฯเปฯ เทฺว ขนฺเธ ฯเปฯ. ปฏิสนฺธิ\-กฺขเณ ฯเปฯ. (3)}

\proseitem{196}{กามาวจรญฺจ นกามาวจรญฺจ ธมฺมํ ปฏิจฺจ กามาวจโร ธมฺโม อุปฺปชฺชติ เหตุปจฺจยา. นกามาวจเร ขนฺเธ จ มหาภูเต จ ปฏิจฺจ จิตฺต\-สมุฏฺฐานํ รูปํ. ปฏิสนฺธิ\-กฺขเณ นกามาวจเร ขนฺเธ จ มหาภูเต จ ปฏิจฺจ กฏตฺตารูปํ.  กามาวจรญฺจ นกามาวจรญฺจ ธมฺมํ ปฏิจฺจ นกามาวจโร ธมฺโม อุปฺปชฺชติ เหตุปจฺจยา. ปฏิสนฺธิ\-กฺขเณ นกามาวจรํ เอกํ ขนฺธญฺจ วตฺถุญฺจ ปฏิจฺจ ตโย ขนฺธา ฯเปฯ เทฺว ขนฺเธ จ ฯเปฯ. กามาวจรญฺจ นกามาวจรญฺจ ธมฺมํ ปฏิจฺจ กามาวจโร จ นกามาวจโร จ ธมฺมา อุปฺปชฺชนฺติ เหตุปจฺจยา. ปฏิสนฺธิ\-กฺขเณ นกามาวจรํ เอกํ ขนฺธญฺจ วตฺถุญฺจ ปฏิจฺจ ตโย ขนฺธา ฯเปฯ เทฺว ขนฺเธ จ ฯเปฯ. นกามาวจเร ขนฺเธ จ มหาภูเต จ ปฏิจฺจ กฏตฺตารูปํ. (3) (สํขิตฺตํ.)}

\csromanmarkh{1. ปจฺจยานุโลม}\csromanmarkhii{2. สงฺขฺยาวาร}\csromanheadernomark{2}{1. \textbf{ปจฺจยานุโลม 2. สงฺขฺยาวาร}}
\proseitem{197}{\csromanfolio{128}เหตุยา นว, อารมฺมเณ จตฺตาริ, อธิปติยา ปญฺจ, อนนฺตเร จตฺตาริ, สมนนฺตเร จตฺตาริ, สหชาเต นว, อญฺญมญฺเญ ฉ, นิสฺสเย นว, อุปนิสฺสเย จตฺตาริ, ปุเรชาเต เทฺว, อาเสวเน เทฺว, กมฺเม นว, วิปาเก นว, อาหาเร นว, อินฺทฺริเย นว, ฌาเน นว, มคฺเค นว, สมฺปยุตฺเต จตฺตาริ, วิปฺปยุตฺเต นว, อตฺถิยา นว, นตฺถิยา จตฺตาริ, วิคเต จตฺตาริ, อวิคเต นว.}

\csromanmarkh{2. ปจฺจย\-ปจฺจนีย}\csromanmarkhii{1. วิภงฺควาร}\csromanheadernomark{2}{2. ปจฺจย\-ปจฺจนีย 1. วิภงฺควาร}
\proseitem{198}{กามาวจรํ ธมฺมํ ปฏิจฺจ กามาวจโร ธมฺโม อุปฺปชฺชติ นเหตุปจฺจยา. อเหตุกํ กามาวจรํ เอกํ ขนฺธํ ปฏิจฺจ ตโย ขนฺธา จิตฺต\-สมุฏฺฐานญฺจ รูปํ ฯเปฯ เทฺว ขนฺเธ ฯเปฯ. อเหตุก\-ปฏิสนฺธิ\-กฺขเณ ฯเปฯ. (ยาว อสญฺญสตฺตา.) วิจิกิจฺฉา\-สหคเต อุทฺธจฺจ\-สหคเต ขนฺเธ ปฏิจฺจ วิจิกิจฺฉา\-สหคโต อุทฺธจฺจ\-สหคโต โมโห. (1). นอารมฺมณ\-ปจฺจยา ตีณิ.}

\prose{นอธิปตฺยาทิ}

\proseitem{199}{กามาวจรํ ธมฺมํ ปฏิจฺจ กามาวจโร ธมฺโม อุปฺปชฺชติ นอธิปติ\-ปจฺจยา. ตีณิ.}

\prose{นกามาวจรํ ธมฺมํ ปฏิจฺจ นกามาวจโร ธมฺโม อุปฺปชฺชติ นอธิปติ\-ปจฺจยา. นกามาวจเร ขนฺเธ ปฏิจฺจ นกามาวจรา\-ธิปติ. วิปากํ นกามาวจรํ เอกํ ขนฺธํ ปฏิจฺจ ตโย ขนฺธา ฯเปฯ เทฺว ขนฺเธ ฯเปฯ. ปฏิสนฺธิ\-กฺขเณ ฯเปฯ. นกามาวจรํ ธมฺมํ ปฏิจฺจ กามาวจโร ธมฺโม อุปฺปชฺชติ นอธิปติ\-ปจฺจยา. วิปาเก นกามาวจเร ขนฺเธ ปฏิจฺจ จิตฺต\-สมุฏฺฐานํ รูปํ. ปฏิสนฺธิ\-กฺขเณ ฯเปฯ. นกามาวจรํ ธมฺมํ ปฏิจฺจ กามาวจโร จ นกามาวจโร จ ธมฺมา อุปฺปชฺชนฺติ นอธิปติ\-ปจฺจยา. วิปากํ นกามาวจรํ เอกํ ขนฺธํ ปฏิจฺจ ตโย ขนฺธา จิตฺต\-สมุฏฺฐานญฺจ รูปํ ฯเปฯ เทฺว ขนฺเธ ฯเปฯ. ปฏิสนฺธิ\-กฺขเณ ฯเปฯ. (3)}

\prose{\csromanfolio{129}กามาวจรญฺจ นกามาวจรญฺจ ธมฺมํ ปฏิจฺจ กามาวจโร ธมฺโม อุปฺปชฺชติ นอธิปติ\-ปจฺจยา. วิปาเก นกามาวจเร ขนฺเธ จ มหาภูเต จ ปฏิจฺจ จิตฺต\-สมุฏฺฐานํ รูปํ. ปฏิสนฺธิ\-กฺขเณ ฯเปฯ. (อิตเร เทฺว ปากติกา.) นอนนฺตร\-ปจฺจยา ฯเปฯ. นปุเรชาต\-ปจฺจยา. นปจฺฉาชาต\-ปจฺจยา.}

\prose{นอาเสวน}

\proseitem{200}{กามาวจรํ ธมฺมํ ปฏิจฺจ กามาวจโร ธมฺโม อุปฺปชฺชติ นอาเสวน\-ปจฺจยา. ตีณิ.}

\prose{นกามาวจรํ ธมฺมํ ปฏิจฺจ นกามาวจโร ธมฺโม อุปฺปชฺชติ นอาเสวน\-ปจฺจยา. วิปากํ นกามาวจรํ เอกํ ขนฺธํ ปฏิจฺจ ตโย ขนฺธา ฯเปฯ เทฺว ขนฺเธ ฯเปฯ. ปฏิสนฺธิ\-กฺขเณ ฯเปฯ. นกามาวจรํ ธมฺมํ ปฏิจฺจ กามาวจโร ธมฺโม อุปฺปชฺชติ นอาเสวน\-ปจฺจยา. นกามาวจเร ขนฺเธ ปฏิจฺจ จิตฺต\-สมุฏฺฐานํ รูปํ. ปฏิสนฺธิ\-กฺขเณ ฯเปฯ. (มูลํ กาตพฺพํ.) วิปากํ นกามาวจรํ เอกํ ขนฺธํ ปฏิจฺจ ตโย ขนฺธา จิตฺต\-สมุฏฺฐานญฺจ รูปํ ฯเปฯ เทฺว ขนฺเธ ฯเปฯ. ปฏิสนฺธิ\-กฺขเณ ฯเปฯ. (3) (อวเสสา ตีณิ ปากติกา. สํขิตฺตํ.)}

\csromanmarkh{2. ปจฺจย\-ปจฺจนีย}\csromanmarkhii{2. สงฺขฺยาวาร}\csromanheadernomark{2}{2. ปจฺจย\-ปจฺจนีย 2. สงฺขฺยาวาร}
\proseitem{201}{นเหตุยา เอกํ, นอารมฺมเณ ตีณิ, นอธิปติยา นว, นอนนฺตเร ตีณิ, นสมนนฺตเร ตีณิ, นอุปนิสฺสเย ตีณิ, นปุเรชาเต นว, นปจฺฉาชาเต นว, นอาเสวเน นว, นกมฺเม เทฺว, นวิปาเก ปญฺจ, นอาหาเร เอกํ, นอินฺทฺริเย เอกํ, นฌาเน เอกํ, นมคฺเค เอกํ, นสมฺปยุตฺเต ตีณิ, นวิปฺปยุตฺเต เทฺว, โนนตฺถิยา ตีณิ, โนวิคเต ตีณิ. (อวเสสา คณนาปิ สหชาตวาโรปิ กาตพฺโพ.)}

\vspace{\tipitakaparskip}
\csromancenter{(อวเสสา คณนาปิ สหชาตวาโรปิ กาตพฺโพ.)}
\csromanmarkh{93. กามาวจรทุก}\csromanmarkhii{3. ปจฺจยวาร}\csromanheadernomark{2}{93. กามาวจรทุก 3. ปจฺจยวาร}
\proseitem{202}{กามาวจรํ ธมฺมํ ปจฺจยา กามาวจโร ธมฺโม อุปฺปชฺชติ เหตุปจฺจยา. กามาวจรํ เอกํ ขนฺธํ ปจฺจยา ตโย ขนฺธา จิตฺต\-สมุฏฺฐานญฺจ \csromanfolio{130}รูปํ ฯเปฯ เทฺว ขนฺเธ ฯเปฯ. ปฏิสนฺธิ\-กฺขเณ ฯเปฯ. (ยาว อชฺฌตฺติกา มหาภูตา.) วตฺถุํ ปจฺจยา กามาวจรา ขนฺธา.  กามาวจรํ ธมฺมํ ปจฺจยา นกามาวจโร ธมฺโม อุปฺปชฺชติ เหตุปจฺจยา. วตฺถุํ ปจฺจยา นกามาวจรา ขนฺธา. ปฏิสนฺธิ\-กฺขเณ ฯเปฯ. กามาวจรํ ธมฺมํ ปจฺจยา กามาวจโร จ นกามาวจโร จ ธมฺมา อุปฺปชฺชนฺติ เหตุปจฺจยา. วตฺถุํ ปจฺจยา นกามาวจรา ขนฺธา. มหาภูเต ปจฺจยา จิตฺต\-สมุฏฺฐานํ รูปํ. ปฏิสนฺธิ\-กฺขเณ ฯเปฯ. (3)}

\prose{นกามาวจรํ ธมฺมํ ปจฺจยา นกามาวจโร ธมฺโม อุปฺปชฺชติ เหตุปจฺจยา. ตีณิ. (ปฏิจฺจสทิสา.)}

\prose{กามาวจรญฺจ นกามาวจรญฺจ ธมฺมํ ปจฺจยา กามาวจโร ธมฺโม อุปฺปชฺชติ เหตุปจฺจยา. (ปฏิจฺจสทิสา.). กามาวจรญฺจ นกามาวจรญฺจ ธมฺมํ ปจฺจยา นกามาวจโร ธมฺโม อุปฺปชฺชติ เหตุปจฺจยา. นกามาวจรํ เอกํ ขนฺธญฺจ วตฺถุญฺจ ปจฺจยา ตโย ขนฺธา ฯเปฯ เทฺว ขนฺเธ จ ฯเปฯ. ปฏิสนฺธิ\-กฺขเณ ฯเปฯ. กามาวจรญฺจ นกามาวจรญฺจ ธมฺมํ ปจฺจยา กามาวจโร จ นกามาวจโร จ ธมฺมา อุปฺปชฺชนฺติ เหตุปจฺจยา. นกามาวจรํ เอกํ ขนฺธญฺจ วตฺถุญฺจ ปจฺจยา ตโย ขนฺธา ฯเปฯ เทฺว ขนฺเธ จ ฯเปฯ. นกามาวจเร ขนฺเธ จ มหาภูเต จ ปจฺจยา จิตฺต\-สมุฏฺฐานํ รูปํ. ปฏิสนฺธิ\-กฺขเณ ฯเปฯ. (3) (สํขิตฺตํ.)}

\prose{เหตุยา นว, อารมฺมเณ จตฺตาริ, อธิปติยา นว, อนนฺตเร จตฺตาริ, สมนนฺตเร จตฺตาริ, สหชาเต นว, อญฺญมญฺเญ ฉ, นิสฺสเย นว, อุปนิสฺสเย จตฺตาริ, ปุเรชาเต จตฺตาริ, อาเสวเน จตฺตาริ, กมฺเม นว ฯเปฯ อวิคเต นว.}

\vspace{\tipitakaparskip}
\csromancenter{\textbf{อนุโลมํ.}}
\proseitem{203}{กามาวจรํ ธมฺมํ ปจฺจยา กามาวจโร ธมฺโม อุปฺปชฺชติ นเหตุปจฺจยา. อเหตุกํ กามาวจรํ เอกํ ขนฺธํ ปจฺจยา ฯเปฯ. อเหตุก\-ปฏิสนฺธิ\-กฺขเณ ฯเปฯ. (ยาว อสญฺญสตฺตา.) จกฺขายตนํ ปจฺจยา จกฺขุวิญฺญาณํ ฯเปฯ. กายายตนํ ปจฺจยา กายวิญฺญาณํ. วตฺถุํ ปจฺจยา อเหตุกา กามาวจรา ขนฺธา. วิจิกิจฺฉา\-สหคเต อุทฺธจฺจ\-สหคเต ขนฺเธ จ วตฺถุญฺจ ปจฺจยา วิจิกิจฺฉา\-สหคโต อุทฺธจฺจ\-สหคโต โมโห. (1) นอารมฺมณ\-ปจฺจยา. ตีณิ.}

\csromanheader{2}{93. กามาวจรทุก นอธิปตฺยาทิ}
\proseitem{204}{\csromanfolio{131}กามาวจรํ ธมฺมํ ปจฺจยา กามาวจโร ธมฺโม อุปฺปชฺชติ นอธิปติ\-ปจฺจยา. เอกํ. (ยาว อสญฺญสตฺตา.). กามาวจรํ ธมฺมํ ปจฺจยา นกามาวจโร ธมฺโม อุปฺปชฺชติ นอธิปติ\-ปจฺจยา. วตฺถุํ ปจฺจยา นกามาวจรา\-ธิปติ. วตฺถุํ ปจฺจยา วิปากา นกามาวจรา ขนฺธา. ปฏิสนฺธิ\-กฺขเณ ฯเปฯ. กามาวจรํ ธมฺมํ ปจฺจยา กามาวจโร จ นกามาวจโร จ ธมฺมา อุปฺปชฺชนฺติ นอธิปติ\-ปจฺจยา. วตฺถุํ ปจฺจยา วิปากา นกามาวจรา ขนฺธา. มหาภูเต ปจฺจยา จิตฺต\-สมุฏฺฐานํ รูปํ. ปฏิสนฺธิ\-กฺขเณ ฯเปฯ. (3)}

\prose{นกามาวจรํ ธมฺมํ ปจฺจยา นกามาวจโร ธมฺโม อุปฺปชฺชติ นอธิปติ\-ปจฺจยา. ตีณิ. (ปฏิจฺจสทิสา.)}

\prose{กามาวจรญฺจ นกามาวจรญฺจ ธมฺมํ ปจฺจยา กามาวจโร ธมฺโม อุปฺปชฺชติ นอธิปติ\-ปจฺจยา. (ปฏิจฺจสทิสา. มูลํ กาตพฺพํ.) นกามาวจเร ขนฺเธ จ วตฺถุญฺจ ปจฺจยา นกามาวจรา\-ธิปติ. วิปากํ นกามาวจรํ เอกํ ขนฺธญฺจ วตฺถุญฺจ ปจฺจยา ตโย ขนฺธา ฯเปฯ เทฺว ขนฺเธ จ ฯเปฯ. ปฏิสนฺธิ\-กฺขเณ ฯเปฯ. (มูลํ กาตพฺพํ.) วิปากํ นกามาวจรํ เอกํ ขนฺธญฺจ วตฺถุญฺจ ปจฺจยา ตโย ขนฺธา ฯเปฯ เทฺว ขนฺเธ จ ฯเปฯ. วิปาเก นกามาวจเร ขนฺเธ จ มหาภูเต จ ปจฺจยา จิตฺต\-สมุฏฺฐานํ รูปํ. ปฏิสนฺธิ\-กฺขเณ ฯเปฯ. (3) นอุปนิสฺสยปจฺจยา. ตีณิ. นอาเสวน\-ปจฺจยา. (สุทฺธเก อรูปมิสฺสเก จ “วิปากนฺ”ติ นิยาเมตพฺพํ. รูปมิสฺสเก นตฺถิ. สํขิตฺตํ.)}

\prose{นเหตุยา เอกํ, นอารมฺมเณ ตีณิ, นอธิปติยา นว, นอนนฺตเร ตีณิ, นสมนนฺตเร ตีณิ, นอุปนิสฺสเย ตีณิ, นปุเรชาเต นว, นปจฺฉาชาเต นว, นอาเสวเน นว, นกมฺเม จตฺตาริ, นวิปาเก นว, นอาหาเร เอกํ, นอินฺทฺริเย เอกํ, นฌาเน เอกํ, นมคฺเค เอกํ, นสมฺปยุตฺเต ตีณิ, นวิปฺปยุตฺเต เทฺว, โนนตฺถิยา ตีณิ, โนวิคเต ตีณิ.}

\vspace{\tipitakaparskip}
\csromancenter{ปจฺจนียํ.}
\csromancenter{(เอวํ อิตเร เทฺว คณนาปิ นิสฺสยวาโรปิ กาตพฺโพ.)}
\csromanmarkh{93. กามาวจรทุก}\csromanmarkhii{5. สํสฏฺฐวาร}\csromanheadernomark{2}{93. \textbf{กามาวจรทุก 5. สํสฏฺฐวาร}}
\proseitem{205}{\csromanfolio{132}กามาวจรํ ธมฺมํ สํสฏฺโฐ กามาวจโร ธมฺโม อุปฺปชฺชติ เหตุปจฺจยา. กามาวจรํ เอกํ ขนฺธํ สํสฏฺฐา ตโย ขนฺธา ฯเปฯ เทฺว ขนฺเธ ฯเปฯ. ปฏิสนฺธิ\-กฺขเณ ฯเปฯ. (1)}

\prose{นกามาวจรํ ธมฺมํ สํสฏฺโฐ นกามาวจโร ธมฺโม อุปฺปชฺชติ เหตุปจฺจยา. นกามาวจรํ เอกํ ขนฺธํ สํสฏฺฐา ตโย ขนฺธา ฯเปฯ เทฺว ขนฺเธ ฯเปฯ. ปฏิสนฺธิ\-กฺขเณ ฯเปฯ. (1)}

\prose{\textbf{เหตุ}ยา เทฺว, อารมฺมเณ เทฺว, อธิปติยา เทฺว, (สพฺพตฺถ เทฺว,) อวิคเต เทฺว. อนุโลมํ.}

\prose{กามาวจรํ ธมฺมํ สํสฏฺโฐ กามาวจโร ธมฺโม อุปฺปชฺชติ นเหตุปจฺจยา. อเหตุกํ กามาวจรํ เอกํ ขนฺธํ สํสฏฺฐา ตโย ขนฺธา ฯเปฯ เทฺว ขนฺเธ ฯเปฯ. อเหตุก\-ปฏิสนฺธิ\-กฺขเณ ฯเปฯ. วิจิกิจฺฉา\-สหคเต อุทฺธจฺจ\-สหคเต ขนฺเธ สํสฏฺโฐ วิจิกิจฺฉา\-สหคโต อุทฺธจฺจ\-สหคโต โมโห.}

\prose{นเหตุยา เอกํ, นอธิปติยา เทฺว, นปุเรชาเต เทฺว, นปจฺฉาชาเต เทฺว, นอาเสวเน เทฺว, นกมฺเม เทฺว, นวิปาเก เทฺว, นฌาเน เอกํ, นมคฺเค เอกํ, นวิปฺปยุตฺเต เทฺว.}

\vspace{\tipitakaparskip}
\csromancenter{ปจฺจนียํ.}
\csromancenter{(เอวํ อิตเร เทฺว คณนาปิ สมฺปยุตฺตวาโรปิ กาตพฺโพ.)}
\csromanmarkh{93. กามาวจรทุก}\csromanmarkhii{7. ปญฺหาวาร}\csromanheadernomark{2}{93. \textbf{กามาวจรทุก 7. ปญฺหาวาร}}
\proseitem{206}{กามาวจโร ธมฺโม กามาวจรสฺส ธมฺมสฺส เหตุปจฺจเยน ปจฺจโย. กามาวจรา เหตู สมฺปยุตฺตกานํ ขนฺธานํ จิตฺตสมุฏฺฐานานญฺจ รูปานํ เหตุปจฺจเยน ปจฺจโย. ปฏิสนฺธิ\-กฺขเณ ฯเปฯ. (1)}

\prose{\csromanfolio{133}นกามาวจโร ธมฺโม นกามาวจรสฺส ธมฺมสฺส เหตุปจฺจเยน ปจฺจโย. นกามาวจรา เหตู สมฺปยุตฺตกานํ ขนฺธานํ เหตุปจฺจเยน ปจฺจโย. ปฏิสนฺธิ\-กฺขเณ ฯเปฯ. นกามาวจโร ธมฺโม กามาวจรสฺส ธมฺมสฺส เหตุปจฺจเยน ปจฺจโย. นกามาวจรา เหตู จิตฺต\-สมุฏฺฐานานํ รูปานํ เหตุปจฺจเยน ปจฺจโย. ปฏิสนฺธิ\-กฺขเณ ฯเปฯ. (มูลํ กาตพฺพํ.) นกามาวจรา เหตู สมฺปยุตฺตกานํ ขนฺธานํ จิตฺตสมุฏฺฐานานญฺจ รูปานํ เหตุปจฺจเยน ปจฺจโย. ปฏิสนฺธิ\-กฺขเณ ฯเปฯ. (3)}

\proseitem{207}{กามาวจโร ธมฺโม กามาวจรสฺส ธมฺมสฺส อารมฺมณ\-ปจฺจเยน ปจฺจโย. ทานํ ทตฺวา, สีลํ ฯเปฯ อุโปสถกมฺมํ ฯเปฯ. ปุพฺเพ สุจิณฺณานิ ปจฺจเวกฺขติ, อสฺสาเทติ อภินนฺทติ, ตํ อารพฺภ ราโค ฯเปฯ โทมนสฺสํ อุปฺปชฺชติ. อาริยา โคตฺรภุํ ปจฺจเวกฺขนฺติ, โวทานํ ปจฺจเวกฺขนฺติ. ปหีเน กิเลเส ฯเปฯ. วิกฺขมฺภิเต กิเลเส ฯเปฯ. ปุพฺเพ สุจิณฺณานิ ฯเปฯ. จกฺขุํ ฯเปฯ. วตฺถุํ กามาวจเร ขนฺเธ อนิจฺจโต ฯเปฯ โทมนสฺสํ อุปฺปชฺชติ. รูปายตนํ จกฺขุ\-วิญฺญาณสฺส ฯเปฯ. โผฏฺฐพฺพา\-ยตนํ กายวิญฺญาณสฺส อารมฺมณ\-ปจฺจเยน ปจฺจโย. (1)}

\prose{กามาวจโร ธมฺโม นกามาวจรสฺส ธมฺมสฺส อารมฺมณ\-ปจฺจเยน ปจฺจโย. ทิพฺเพน จกฺขุนา รูปํ ปสฺสติ. ทิพฺพาย โสตธาตุยา สทฺทํ สุณาติ. เจโตปริย\-ญาเณน กามาวจร\-จิตฺต\-สมงฺคิสฺส จิตฺตํ ชานาติ. กามาวจรา ขนฺธา อิทฺธิวิธ\-ญาณสฺส, เจโตปริย\-ญาณสฺส, ปุพฺเพ\-นิวาสา\-นุสฺสติ\-ญาณสฺส, ยถากมฺมูปคญาณสฺส, อนาคตํสญาณสฺส อารมฺมณ\-ปจฺจเยน ปจฺจโย. (2)}

\proseitem{208}{นกามาวจโร ธมฺโม นกามาวจรสฺส ธมฺมสฺส อารมฺมณ\-ปจฺจเยน ปจฺจโย. นิพฺพานํ มคฺคสฺส, ผลสฺส อารมฺมณ\-ปจฺจเยน ปจฺจโย. เจโตปริย\-ญาเณน นกามาวจร\-จิตฺต\-สมงฺคิสฺส จิตฺตํ ชานาติ. อากาสา\-นญฺจา\-ยตนํ วิญฺญาณญฺจา\-ยตนสฺส ฯเปฯ. อากิญฺจญฺญา\-ยตนํ เนว\-สญฺญา\-นา\-สญฺญา\-ยตนสฺส ฯเปฯ. นกามาวจรา ขนฺธา อิทฺธิวิธ\-ญาณสฺส, เจโตปริย\-ญาณสฺส, ปุพฺเพ\-นิวาสา\-นุสฺสติ\-ญาณสฺส, ยถากมฺมูปคญาณสฺส, อนาคตํสญาณสฺส อารมฺมณ\-ปจฺจเยน ปจฺจโย. (1)}

\prose{\csromanfolio{134}นกามาวจโร ธมฺโม กามาวจรสฺส ธมฺมสฺส อารมฺมณ\-ปจฺจเยน ปจฺจโย. ฌานา วุฏฺฐหิตฺวา ฌานํ ปจฺจเวกฺขติ, อสฺสาเทติ อภินนฺทติ, ตํ อารพฺภ ราโค อุปฺปชฺชติ. ทิฏฺฐิ อุปฺปชฺชติ ฯเปฯ. อริยา มคฺคา วุฏฺฐหิตฺวา มคฺคํ ปจฺจเวกฺขนฺติ. ผลํ ฯเปฯ. นิพฺพานํ ปจฺจเวกฺขนฺติ. นิพฺพานํ โคตฺรภุสฺส, โวทานสฺส, อาวชฺชนาย อารมฺมณ\-ปจฺจเยน ปจฺจโย. อากาสา\-นญฺจา\-ยตนํ ปจฺจเวกฺขติ. วิญฺญาณญฺจา\-ยตนํ ปจฺจเวกฺขติ. อากิญฺจญฺญา\-ยตนํ ปจฺจเวกฺขติ. เนว\-สญฺญา\-นา\-สญฺญา\-ยตนํ ปจฺจเวกฺขติ. ทิพฺพํ จกฺขุํ ปจฺจเวกฺขติ. ทิพฺพํ โสตธาตุํ ปจฺจเวกฺขติ. อิทฺธิวิธ\-ญาณํ ปจฺจเวกฺขติ. เจโตปริย\-ญาณํ ฯเปฯ. ปุพฺเพ นิวาสา\-นุสฺสติ\-ญาณํ ฯเปฯ. ยถากมฺมูปคญาณํ ฯเปฯ. อนาคตํสญาณํ ปจฺจเวกฺขติ. นกามาวจเร ขนฺเธ อนิจฺจโต ฯเปฯ โทมนสฺสํ อุปฺปชฺชติ. (2)}

\csromanheader{2}{\textbf{อธิปติ}}
\proseitem{209}{กามาวจโร ธมฺโม กามาวจรสฺส ธมฺมสฺส อธิปติ\-ปจฺจเยน ปจฺจโย. อารมฺมณา\-ธิปติ, สหชาตา\-ธิปติ.  อารมฺมณา\-ธิปติ\csromandash{}ทานํ ฯเปฯ สีลํ ฯเปฯ อุโปสถกมฺมํ ฯเปฯ. ปุพฺเพ สุจิณฺณานิ ครุํ กตฺวา ปจฺจเวกฺขติ, อสฺสาเทติ อภินนฺทติ, ตํ ครุํ กตฺวา ราโค อุปฺปชฺชติ. ทิฏฺฐิ อุปฺปชฺชติ. เสกฺขา โคตฺรภุํ ครุํ กตฺวา ปจฺจเวกฺขนฺติ. โวทานํ ครุํ กตฺวา ปจฺจเวกฺขนฺติ. จกฺขุํ ฯเปฯ. วตฺถุํ กามาวจเร ขนฺเธ ครุํ กตฺวา อสฺสาเทติ อภินนฺทติ, ตํ ครุํ กตฺวา ราโค อุปฺปชฺชติ. ทิฏฺฐิ อุปฺปชฺชติ. สหชาตา\-ธิปติ\csromandash{}กามาวจรา\-ธิปติ สมฺปยุตฺตกานํ ขนฺธานํ จิตฺตสมุฏฺฐานานญฺจ รูปานํ อธิปติ\-ปจฺจเยน ปจฺจโย. (1)}

\proseitem{210}{นกามาวจโร ธมฺโม นกามาวจรสฺส ธมฺมสฺส อธิปติ\-ปจฺจเยน ปจฺจโย. อารมฺมณา\-ธิปติ, สหชาตา\-ธิปติ.  อารมฺมณา\-ธิปติ\csromandash{}นิพฺพานํ มคฺคสฺส, ผลสฺส อธิปติ\-ปจฺจเยน ปจฺจโย. สหชาตา\-ธิปติ\csromandash{}นกามาวจรา\-ธิปติ สมฺปยุตฺตกานํ ขนฺธานํ อธิปติ\-ปจฺจเยน ปจฺจโย.  นกามาวจโร ธมฺโม กามาวจรสฺส ธมฺมสฺส อธิปติ\-ปจฺจเยน ปจฺจโย. อารมฺมณา\-ธิปติ, สหชาตา\-ธิปติ.  อารมฺมณา\-ธิปติ\csromandash{}ฌานา วุฏฺฐหิตฺวา ฌานํ ครุํ กตฺวา ปจฺจเวกฺขติ, อสฺสาเทติ อภินนฺทติ, ตํ ครุํ กตฺวา ราโค อุปฺปชฺชติ. ทิฏฺฐิ อุปฺปชฺชติ. อริยา มคฺคา วุฏฺฐหิตฺวา มคฺคํ ครุํ กตฺวา ปจฺจเวกฺขนฺติ. ผลํ ฯเปฯ. นิพฺพานํ ครุํ กตฺวา ปจฺจเวกฺขนฺติ. นิพฺพานํ โคตฺรภุสฺส, โวทานสฺส อธิปติ\-ปจฺจเยน ปจฺจโย. อากาสา\-นญฺจา\-ยตนํ ครุํ กตฺวา \csromanfolio{135}ปจฺจเวกฺขติ. วิญฺญาณญฺจา\-ยตนํ ฯเปฯ. อากิญฺจญฺญา\-ยตนํ ฯเปฯ. เนว\-สญฺญา\-นา\-สญฺญา\-ยตนํ ครุํ กตฺวา ปจฺจเวกฺขติ. ทิพฺพํ จกฺขุํ ฯเปฯ. ทิพฺพํ โสตธาตุํ ฯเปฯ. อิทฺธิวิธ\-ญาณํ ฯเปฯ. อนาคตํสญาณํ ครุํ กตฺวา ปจฺจเวกฺขติ. นกามาวจเร ขนฺเธ ครุํ กตฺวา อสฺสาเทติ อภินนฺทติ, ตํ ครุํ กตฺวา ราโค อุปฺปชฺชติ. ทิฏฺฐิ อุปฺปชฺชติ. สหชาตา\-ธิปติ\csromandash{}นกามาวจรา\-ธิปติ จิตฺต\-สมุฏฺฐานานํ รูปานํ อธิปติ\-ปจฺจเยน ปจฺจโย.  นกามาวจโร ธมฺโม กามาวจรสฺส จ นกามาวจรสฺส จ ธมฺมสฺส อธิปติ\-ปจฺจเยน ปจฺจโย. สหชาตา\-ธิปติ\csromandash{}นกามาวจรา\-ธิปติ สมฺปยุตฺตกานํ ขนฺธานํ จิตฺตสมุฏฺฐานานญฺจ รูปานํ อธิปติ\-ปจฺจเยน ปจฺจโย. (3)}

\proseitem{211}{กามาวจโร ธมฺโม กามาวจรสฺส ธมฺมสฺส อนนฺตร\-ปจฺจเยน ปจฺจโย. ปุริมา ปุริมา กามาวจรา ขนฺธา ปจฺฉิมานํ ปจฺฉิมานํ กามาวจรานํ ขนฺธานํ อนนฺตร\-ปจฺจเยน ปจฺจโย. อนุโลมํ โคตฺรภุสฺส. อนุโลมํ โวทานสฺส. อาวชฺชนา กามาวจรานํ ขนฺธานํ อนนฺตร\-ปจฺจเยน ปจฺจโย. (1)}

\prose{กามาวจโร ธมฺโม นกามาวจรสฺส ธมฺมสฺส อนนฺตร\-ปจฺจเยน ปจฺจโย. กามาวจรํ จุติจิตฺตํ นกามาวจรสฺส อุปปตฺติ\-จิตฺตสฺส. อาวชฺชนา นกามาวจรานํ ขนฺธานํ อนนฺตร\-ปจฺจเยน ปจฺจโย. กามาวจรา ขนฺธา นกามาวจรสฺส วุฏฺฐานสฺส อนนฺตร\-ปจฺจเยน ปจฺจโย. ปฐมสฺส ฌานสฺส ปริกมฺมํ ปฐมสฺส ฌานสฺส อนนฺตร\-ปจฺจเยน ปจฺจโย ฯเปฯ. จตุตฺถสฺส ฌานสฺส ฯเปฯ. เนว\-สญฺญา\-นา\-สญฺญา\-ยตนสฺส ปริกมฺมํ ฯเปฯ. ทิพฺพสฺส จกฺขุสฺส ฯเปฯ. ทิพฺพาย โสตธาตุยา ฯเปฯ. อิธิวิธญาณสฺส ฯเปฯ. เจโตปริย\-ญาณสฺส ฯเปฯ. ปุพฺเพ\-นิวาสา\-นุสฺสติ\-ญาณสฺส ฯเปฯ. ยถากมฺมูปคญาณสฺส ฯเปฯ. อนาคตํสญาณสฺส ปริกมฺมํ อนาคตํสญาณสฺส อนนฺตร\-ปจฺจเยน ปจฺจโย. โคตฺรภุ มคฺคสฺส. โวทานํ มคฺคสฺส. อนุโลมํ ผลสมาปตฺติยา อนนฺตร\-ปจฺจเยน ปจฺจโย. (2)}

\proseitem{212}{นกามาวจโร ธมฺโม นกามาวจรสฺส ธมฺมสฺส อนนฺตร\-ปจฺจเยน ปจฺจโย. ปุริมา ปุริมา นกามาวจรา ขนฺธา ปจฺฉิมานํ ปจฺฉิมานํ นกามาวจรานํ ขนฺธานํ อนนฺตร\-ปจฺจเยน ปจฺจโย. มคฺโค ผลสฺส. ผลํ ผลสฺส. นิโรธา วุฏฺฐหนฺตสฺส เนว\-สญฺญา\-นา\-สญฺญา\-ยตนํ ผลสมาปตฺติยา อนนฺตร\-ปจฺจเยน ปจฺจโย.  นกามาวจโร ธมฺโม กามาวจรสฺส ธมฺมสฺส อนนฺตร\-ปจฺจเยน \csromanfolio{136}ปจฺจโย. นกามาวจรํ จุติจิตฺตํ กามาวจรสฺส อุปปตฺติ\-จิตฺตสฺส. นกามาวจรํ ภวงฺคํ อาวชฺชนาย. นกามาวจรา ขนฺธา กามาวจรสฺส วุฏฺฐานสฺส อนนฺตร\-ปจฺจเยน ปจฺจโย. (2)}

\prose{สมนนฺตร\-ปจฺจเยน ปจฺจโย. สหชาต\-ปจฺจเยน ปจฺจโย. สตฺต. อญฺญมญฺญ\-ปจฺจเยน ปจฺจโย. ฉ. นิสฺสย\-ปจฺจเยน ปจฺจโย. สตฺต.}

\csromanheader{2}{\textbf{อุปนิสฺสย}}
\proseitem{213}{กามาวจโร ธมฺโม กามาวจรสฺส ธมฺมสฺส อุปนิสฺสย\-ปจฺจเยน ปจฺจโย. อารมฺมณู\-ปนิสฺสโย, อนนฺตรู\-ปนิสฺสโย, ปกตู\-ปนิสฺสโย ฯเปฯ. ปกตู\-ปนิสฺสโย\csromandash{}กามาวจรํ สทฺธํ อุปนิสฺสาย ทานํ เทติ ฯเปฯ สีลํ ฯเปฯ อุโปสถกมฺมํ ฯเปฯ. วิปสฺสนํ อุปฺปาเทติ. มานํ ชปฺเปติ. ทิฏฺฐิํ คณฺหาติ. กามาวจรํ สีลํ ฯเปฯ ปญฺญํ. ราคํ ฯเปฯ. ปตฺถนํ. กายิกํ สุขํ. กายิกํ ทุกฺขํ. อุตุํ. โภชนํ. เสนาสนํ อุปนิสฺสาย ทานํ เทติ ฯเปฯ. วิปสฺสนํ อุปฺปาเทติ. ปาณํ หนติ ฯเปฯ สํฆํ ภินฺทติ. กามาวจรา สทฺธา ฯเปฯ. เสนาสนํ กามาวจราย สทฺธาย ฯเปฯ ปตฺถนาย. กายิกสฺส สุขสฺส. กายิกสฺส ทุกฺขสฺส อุปนิสฺสย\-ปจฺจเยน ปจฺจโย. (1)}

\prose{กามาวจโร ธมฺโม นกามาวจรสฺส ธมฺมสฺส อุปนิสฺสย\-ปจฺจเยน ปจฺจโย. อนนฺตรู\-ปนิสฺสโย, ปกตู\-ปนิสฺสโย ฯเปฯ. ปกตู\-ปนิสฺสโย\csromandash{}กามาวจรํ สทฺธํ อุปนิสฺสาย นกามาวจรํ ฌานํ อุปฺปาเทติ. มคฺคํ ฯเปฯ. อภิญฺญํ ฯเปฯ. สมาปตฺติํ อุปฺปาเทติ. กามาวจรํ สีลํ ฯเปฯ. เสนาสนํ อุปนิสฺสาย นกามาวจรํ ฌานํ อุปฺปาเทติ. มคฺคํ ฯเปฯ. อภิญฺญํ ฯเปฯ. สมาปตฺติํ อุปฺปาเทติ. กามาวจรา สทฺธา ฯเปฯ. เสนาสนํ นกามาวจราย สทฺธาย ฯเปฯ ปญฺญาย. มคฺคสฺส, ผลสมาปตฺติยา อุปนิสฺสย\-ปจฺจเยน ปจฺจโย. ปฐมสฺส ฌานสฺส ปริกมฺมํ ปฐมสฺส ฌานสฺส ฯเปฯ. จตุตฺถสฺส ฌานสฺส ฯเปฯ. อากาสา\-นญฺจา\-ยตนสฺส ฯเปฯ. ปฐมสฺส มคฺคสฺส ฯเปฯ. จตุตฺถสฺส มคฺคสฺส ปริกมฺมํ จตุตฺถสฺส มคฺคสฺส อุปนิสฺสย\-ปจฺจเยน ปจฺจโย. (2)}

\proseitem{214}{นกามาวจโร ธมฺโม นกามาวจรสฺส ธมฺมสฺส อุปนิสฺสย\-ปจฺจเยน ปจฺจโย. อารมฺมณู\-ปนิสฺสโย, อนนฺตรู\-ปนิสฺสโย, ปกตู\-ปนิสฺสโย ฯเปฯ. ปกตู\-ปนิสฺสโย\csromandash{}นกามาวจรํ สทฺธํ อุปนิสฺสาย ฌานํ อุปฺปาเทติ. มคฺคํ ฯเปฯ. อภิญฺญํ ฯเปฯ. สมาปตฺติํ อุปฺปาเทติ. นกามาวจรํ \csromanfolio{137}สีลํ ฯเปฯ ปญฺญํ อุปนิสฺสาย นกามาวจรํ ฌานํ อุปฺปาเทติ ฯเปฯ. มคฺคํ ฯเปฯ. อภิญฺญํ ฯเปฯ. สมาปตฺติํ อุปฺปาเทติ. นกามาวจรา สทฺธา ฯเปฯ ปญฺญา นกามาวจราย สทฺธาย ฯเปฯ ปญฺญาย. มคฺคสฺส, ผลสมาปตฺติยา อุปนิสฺสย\-ปจฺจเยน ปจฺจโย. ปฐมํ ฌานํ ทุติยสฺส ฌานสฺส อุปนิสฺสย\-ปจฺจเยน ปจฺจโย ฯเปฯ. ตติยํ ฌานํ ฯเปฯ. จตุตฺถํ ฌานํ ฯเปฯ. อากาสา\-นญฺจา\-ยตนํ วิญฺญาณญฺจา\-ยตนสฺส ฯเปฯ. อากิญฺจญฺญา\-ยตนํ เนว\-สญฺญา\-นา\-สญฺญา\-ยตนสฺส อุปนิสฺสย\-ปจฺจเยน ปจฺจโย. ปฐโม มคฺโค ทุติยสฺส มคฺคสฺส ฯเปฯ. ทุติโย มคฺโค ตติยสฺส มคฺคสฺส ฯเปฯ. ตติโย มคฺโค จตุตฺถสฺส มคฺคสฺส อุปนิสฺสย\-ปจฺจเยน ปจฺจโย. มคฺโค ผลสมาปตฺติยา อุปนิสฺสย\-ปจฺจเยน ปจฺจโย. (1)}

\prose{นกามาวจรา ธมฺโม กามาวจรสฺส ธมฺมสฺส อุปนิสฺสย\-ปจฺจเยน ปจฺจโย. อารมฺมณู\-ปนิสฺสโย, อนนฺตรู\-ปนิสฺสโย, ปกตู\-ปนิสฺสโย ฯเปฯ. ปกตู\-ปนิสฺสโย\csromandash{}นกามาวจรํ สทฺธํ อุปนิสฺสาย ทานํ เทติ. สีลํ ฯเปฯ อุโปสถกมฺมํ ฯเปฯ. วิปสฺสนํ อุปฺปาเทติ. มานํ ชปฺเปติ. ทิฏฺฐิํ คณฺหาติ. นกามาวจรํ สีลํ ฯเปฯ ปญฺญํ อุปนิสฺสาย ทานํ เทติ. สีลํ ฯเปฯ. อุโปสถกมฺมํ ฯเปฯ. วิปสฺสนํ อุปฺปาเทติ. มานํ ชปฺเปติ. ทิฏฺฐิํ คณฺหาติ. นกามาวจรา สทฺธา ฯเปฯ ปญฺญา กามาวจราย สทฺธาย ฯเปฯ ปญฺญาย. ราคสฺส ฯเปฯ. ปตฺถนาย. กายิกสฺส สุขสฺส. กายิกสฺส ทุกฺขสฺส อุปนิสฺสย\-ปจฺจเยน ปจฺจโย. อริยา มคฺคํ อุปนิสฺสาย สงฺขาเร อนิจฺจโต ฯเปฯ วิปสฺสนฺติ. มคฺโค อริยานํ อตฺถ\-ปฏิสมฺภิทาย. ธมฺม\-ปฏิสมฺภิทาย. นิรุตฺติ\-ปฏิสมฺภิทาย. ปฏิภาน\-ปฏิสมฺภิทาย. ฐานาฐาน\-โกสลฺลสฺส อุปนิสฺสย\-ปจฺจเยน ปจฺจโย. ผลสมาปตฺติ กายิกสฺส สุขสฺส อุปนิสฺสย\-ปจฺจเยน ปจฺจโย. (2)}

\csromanheader{2}{\textbf{ปุเรชาตาทิ}}
\proseitem{215}{กามาวจโร ธมฺโม กามาวจรสฺส ธมฺมสฺส ปุเรชาต\-ปจฺจเยน ปจฺจโย. อารมฺมณ\-ปุเรชาตํ, วตฺถุ\-ปุเรชาตํ.  อารมฺมณ\-ปุเรชาตํ\csromandash{}จกฺขุํ ฯเปฯ. วตฺถุํ อนิจฺจโต ฯเปฯ โทมนสฺสํ อุปฺปชฺชติ. รูปายตนํ จกฺขุ\-วิญฺญาณสฺส ฯเปฯ. โผฏฺฐพฺพา\-ยตนํ กายวิญฺญาณสฺส ฯเปฯ. วตฺถุ\-ปุเรชาตํ\csromandash{}จกฺขายตนํ จกฺขุ\-วิญฺญาณสฺส ฯเปฯ. กายายตนํ กายวิญฺญาณสฺส ฯเปฯ. วตฺถุ กามาวจรานํ ขนฺธานํ ปุเรชาต\-ปจฺจเยน ปจฺจโย.  กามาวจโร \csromanfolio{138}ธมฺโม นกามาวจรสฺส ธมฺมสฺส ปุเรชาต\-ปจฺจเยน ปจฺจโย. อารมฺมณ\-ปุเรชาตํ, วตฺถุ\-ปุเรชาตํ.  อารมฺมณ\-ปุเรชาตํ\csromandash{}ทิพฺเพน จกฺขุนา รูปํ ปสฺสติ. ทิพฺพาย โสตธาตุยา สทฺทํ สุณาติ. วตฺถุ\-ปุเรชาตํ\csromandash{}วตฺถุ นกามาวจรานํ ขนฺธานํ ปุเรชาต\-ปจฺจเยน ปจฺจโย. (2)}

\prose{ปจฺฉาชาต\-ปจฺจเยน ปจฺจโย. เทฺว. อาเสวน\-ปจฺจเยน ปจฺจโย. ตีณิ.}

\csromanheader{2}{\textbf{กมฺมาทิ}}
\proseitem{216}{กามาวจโร ธมฺโม กามาวจรสฺส ธมฺมสฺส กมฺมปจฺจเยน ปจฺจโย. สหชาตา, นานากฺขณิกา. (สํขิตฺตํ.)}

\prose{นกามาวจโร ธมฺโม นกามาวจรสฺส ธมฺมสฺส กมฺมปจฺจเยน ปจฺจโย. สหชาตา, นานากฺขณิกา. (สํขิตฺตํ.). นกามาวจโร ธมฺโม กามาวจรสฺส ธมฺมสฺส กมฺมปจฺจเยน ปจฺจโย. สหชาตา, นานากฺขณิกา. (สํขิตฺตํ.). นกามาวจโร ธมฺโม กามาวจรสฺส จ นกามาวจรสฺส จ ธมฺมสฺส กมฺมปจฺจเยน ปจฺจโย. สหชาตา, นานากฺขณิกา. (สํขิตฺตํ.) (3)}

\prose{วิปาก\-ปจฺจเยน ปจฺจโย. จตฺตาริ. อาหารปจฺจเยน ปจฺจโย. จตฺตาริ. อินฺทฺริย\-ปจฺจเยน ปจฺจโย. จตฺตาริ. ฌานปจฺจเยน ปจฺจโย. จตฺตาริ. มคฺคปจฺจเยน ปจฺจโย. จตฺตาริ. สมฺปยุตฺต\-ปจฺจเยน ปจฺจโย. เทฺว.}

\csromanheader{2}{\textbf{วิปฺปยุตฺต}}
\proseitem{217}{กามาวจโร ธมฺโม กามาวจรสฺส ธมฺมสฺส วิปฺปยุตฺต\-ปจฺจเยน ปจฺจโย. สหชาตํ, ปุเรชาตํ, ปจฺฉาชาตํ. (สํขิตฺตํ.). กามาวจโร ธมฺโม นกามาวจรสฺส ธมฺมสฺส วิปฺปยุตฺต\-ปจฺจเยน ปจฺจโย. สหชาตํ, ปุเรชาตํ.  สหชาตํ\csromandash{} (สํขิตฺตํ.) ปฏิสนฺธิ\-กฺขเณ วตฺถุ นกามาวจรานํ ขนฺธานํ วิปฺปยุตฺต\-ปจฺจเยน ปจฺจโย. ปุเรชาตํ วตฺถุ นกามาวจรานํ ขนฺธานํ วิปฺปยุตฺต\-ปจฺจเยน ปจฺจโย. (2)}

\prose{นกามาวจโร ธมฺโม กามาวจรสฺส ธมฺมสฺส วิปฺปยุตฺต\-ปจฺจเยน ปจฺจโย. สหชาตํ, ปจฺฉาชาตํ.  สหชาตา นกามาวจรา ขนฺธา จิตฺต\-สมุฏฺฐานานํ รูปานํ วิปฺปยุตฺต\-ปจฺจเยน ปจฺจโย. ปฏิสนฺธิ\-กฺขเณ ฯเปฯ. ปจฺฉาชาตา นกามาวจรา ขนฺธา ปุเรชาตสฺส อิมสฺส กายสฺส วิปฺปยุตฺต\-ปจฺจเยน ปจฺจโย. (1)}

\csromanheader{2}{93. กามาวจรทุก อตฺถิอาทิ}
\proseitem{218}{\csromanfolio{139}กามาวจโร ธมฺโม กามาวจรสฺส ธมฺมสฺส อตฺถิปจฺจเยน ปจฺจโย. สหชาตํ, ปุเรชาตํ, ปจฺฉาชาตํ, อาหารํ, อินฺทฺริยํ. (สํขิตฺตํ.). กามาวจโร ธมฺโม นกามาวจรสฺส ธมฺมสฺส อตฺถิปจฺจเยน ปจฺจโย. สหชาตํ, ปุเรชาตํ.  สหชาตํ ปฏิสนฺธิ\-กฺขเณ วตฺถุ นกามาวจรานํ ขนฺธานํ อตฺถิปจฺจเยน ปจฺจโย. ปุเรชาตํ ทิพฺเพน จกฺขุนา รูปํ ปสฺสติ ฯเปฯ. (ปุเรชาต\-สทิสํ.) (2)}

\prose{นกามาวจโร ธมฺโม นกามาวจรสฺส ธมฺมสฺส อตฺถิปจฺจเยน ปจฺจโย. สหชาตํ. (สํขิตฺตํ.). นกามาวจโร ธมฺโม กามาวจรสฺส ธมฺมสฺส อตฺถิปจฺจเยน ปจฺจโย. สหชาตํ, ปจฺฉาชาตํ. (วิปฺปยุตฺต\-สทิสํ.). นกามาวจโร ธมฺโม กามาวจรสฺส จ นกามาวจรสฺส จ ธมฺมสฺส อตฺถิปจฺจเยน ปจฺจโย. (ปฏิจฺจสทิสํ.) (3)}

\proseitem{219}{กามาวจโร จ นกามาวจโร จ ธมฺมา กามาวจรสฺส ธมฺมสฺส อตฺถิปจฺจเยน ปจฺจโย. สหชาตํ, ปจฺฉาชาตํ, อาหารํ, อินฺทฺริยํ.  สหชาตา นกามาวจรา ขนฺธา จ มหาภูตา จ จิตฺต\-สมุฏฺฐานานํ รูปานํ อตฺถิปจฺจเยน ปจฺจโย. ปฏิสนฺธิ\-กฺขเณ ฯเปฯ. ปจฺฉาชาตา นกามาวจรา ขนฺธา จ กพฬีกาโร อาหาโร จ อิมสฺส กายสฺส อตฺถิปจฺจเยน ปจฺจโย. ปจฺฉาชาตา นกามาวจรา ขนฺธา จ รูปชีวิตินฺทฺริยญฺจ กฏตฺตารูปานํ อตฺถิปจฺจเยน ปจฺจโย.  กามาวจโร จ นกามาวจโร จ ธมฺมา นกามาวจรสฺส ธมฺมสฺส อตฺถิปจฺจเยน ปจฺจโย. สหชาตํ, ปุเรชาตํ.  สหชาโต นกามาวจโร เอโก ขนฺโธ จ วตฺถุ จ ติณฺณนฺนํ ขนฺธานํ อตฺถิปจฺจเยน ปจฺจโย ฯเปฯ เทฺว ขนฺธา จ ฯเปฯ. ปฏิสนฺธิ\-กฺขเณ ฯเปฯ. (2)}

\prose{นตฺถิ\-ปจฺจเยน ปจฺจโย. วิคต\-ปจฺจเยน ปจฺจโย. อวิคต\-ปจฺจเยน ปจฺจโย.}

\csromanmarkh{1. ปจฺจยานุโลม}\csromanmarkhii{2. สงฺขฺยาวาร}\csromanheadernomark{2}{1. \textbf{ปจฺจยานุโลม 2. สงฺขฺยาวาร}}
\proseitem{220}{เหตุยา จตฺตาริ, อารมฺมเณ จตฺตาริ, อธิปติยา จตฺตาริ, อนนฺตเร จตฺตาริ, สมนนฺตเร จตฺตาริ, สหชาเต สตฺต, อญฺญมญฺเญ ฉ, นิสฺสเย สตฺต, อุปนิสฺสเย จตฺตาริ, ปุเรชาเต เทฺว, ปจฺฉาชาเต เทฺว, \csromanfolio{140}อาเสวเน ตีณิ, กมฺเม จตฺตาริ, วิปาเก จตฺตาริ, อาหาเร จตฺตาริ, อินฺทฺริเย จตฺตาริ, ฌาเน จตฺตาริ, มคฺเค จตฺตาริ, สมฺปยุตฺเต เทฺว, วิปฺปยุตฺเต ตีณิ, อตฺถิยา สตฺต, นตฺถิยา จตฺตาริ, วิคเต จตฺตาริ, อวิคเต สตฺต.}

\proseitem{221}{กามาวจโร ธมฺโม กามาวจรสฺส ธมฺมสฺส อารมฺมณ\-ปจฺจเยน ปจฺจโย. สหชาต\-ปจฺจเยน ปจฺจโย. อุปนิสฺสย\-ปจฺจเยน ปจฺจโย. ปุเรชาต\-ปจฺจเยน ปจฺจโย. ปจฺฉาชาต\-ปจฺจเยน ปจฺจโย. กมฺมปจฺจเยน ปจฺจโย. อาหารปจฺจเยน ปจฺจโย. อินฺทฺริย\-ปจฺจเยน ปจฺจโย.  กามาวจโร ธมฺโม นกามาวจรสฺส ธมฺมสฺส อารมฺมณ\-ปจฺจเยน ปจฺจโย. สหชาต\-ปจฺจเยน ปจฺจโย. อุปนิสฺสย\-ปจฺจเยน ปจฺจโย. ปุเรชาต\-ปจฺจเยน ปจฺจโย. (2)}

\proseitem{222}{นกามาวจโร ธมฺโม นกามาวจรสฺส ธมฺมสฺส อารมฺมณ\-ปจฺจเยน ปจฺจโย. สหชาต\-ปจฺจเยน ปจฺจโย. อุปนิสฺสย\-ปจฺจเยน ปจฺจโย.  นกามาวจโร ธมฺโม กามาวจรสฺส ธมฺมสฺส อารมฺมณ\-ปจฺจเยน ปจฺจโย. สหชาต\-ปจฺจเยน ปจฺจโย. อุปนิสฺสย\-ปจฺจเยน ปจฺจโย. ปจฺฉาชาต\-ปจฺจเยน ปจฺจโย. กมฺมปจฺจเยน ปจฺจโย.  นกามาวจโร ธมฺโม กามาวจรสฺส จ นกามาวจรสฺส จ ธมฺมสฺส สหชาต\-ปจฺจเยน ปจฺจโย. กมฺมปจฺจเยน ปจฺจโย. (3)}

\prose{กามาวจโร จ นกามาวจโร จ ธมฺมา กามาวจรสฺส ธมฺมสฺส สหชาตํ, ปจฺฉาชาตํ, อาหารํ, อินฺทฺริยํ.  กามาวจโร จ นกามาวจโร จ ธมฺมา นกามาวจรสฺส ธมฺมสฺส สหชาตํ, ปุเรชาตํ. (2)}

\csromanmarkh{2. ปจฺจย\-ปจฺจนีย}\csromanmarkhii{2. สงฺขฺยาวาร}\csromanheadernomark{2}{2. ปจฺจย\-ปจฺจนีย 2. สงฺขฺยาวาร}
\proseitem{223}{นเหตุยา สตฺต, นอารมฺมเณ สตฺต, นอธิปติยา สตฺต, นอนนฺตเร สตฺต, นสมนนฺตเร สตฺต, น สหชาเต ฉ, นอญฺญมญฺเญ ฉ, นนิสฺสเย ฉ, นอุปนิสฺสเย สตฺต, นปุเรชาเต ฉ ฯเปฯ นสมฺปยุตฺเต ฉ,}

\prose{\csromanfolio{141}นวิปฺปยุตฺเต ปญฺจ, โนอตฺถิยา ปญฺจ, โนนตฺถิยา สตฺต, โนวิคเต สตฺต, โนอวิคเต ปญฺจ.}

\vspace{\tipitakaparskip}
\csromancenter{\textbf{เหตุ}ปจฺจยา นอารมฺมเณ จตฺตาริ, นอธิปติยา จตฺตาริ, นอนนฺตเร จตฺตาริ, นสมนนฺตเร จตฺตาริ, นอญฺญมญฺเญ เทฺว, นฺเอาปนิสฺสเย จตฺตาริ ฯเปฯ นสมฺปยุตฺเต เทฺว, นวิปฺปยุตฺเต เทฺว, โนนตฺถิยา จตฺตาริ, โนวิคเต จตฺตาริ.}
\csromanheader{2}{4. ปจฺจย\-ปจฺจนียา\-นุโลม}
\proseitemwithcloser{225}{นเหตุปจฺจยา อารมฺมเณ จตฺตาริ, อธิปติยา จตฺตาริ, (อนุโลมมาติกา กาตพฺพา.) ฯเปฯ อวิคเต สตฺต.}{กามาวจรทุกํ นิฏฺฐิตํ.}
\csromanmatikaanchor{mtk.23}
\csromantocmarkref{section}{94. รูปาวจรทุก}{mtk.23}
\bookmark[dest={mtk.23},level=1]{94. รูปาวจรทุก}
\csromanmarkh{94. รูปาวจรทุก}\csromanmarkhii{1. ปฏิจฺจวาร}\csromanheadernomark{1}{94. \textbf{รูปาวจรทุก 1. ปฏิจฺจวาร}}
\proseitem{226}{รูปาวจรํ ธมฺมํ ปฏิจฺจ รูปาวจโร ธมฺโม อุปฺปชฺชติ เหตุปจฺจยา. รูปาวจรํ เอกํ ขนฺธํ ปฏิจฺจ ตโย ขนฺธา ฯเปฯ เทฺว ขนฺเธ ฯเปฯ. ปฏิสนฺธิ\-กฺขเณ ฯเปฯ. รูปาวจรํ ธมฺมํ ปฏิจฺจ นรูปาวจโร ธมฺโม อุปฺปชฺชติ เหตุปจฺจยา. รูปาวจเร ขนฺเธ ปฏิจฺจ จิตฺต\-สมุฏฺฐานํ รูปํ. ปฏิสนฺธิ\-กฺขเณ ฯเปฯ. รูปาวจรํ ธมฺมํ ปฏิจฺจ รูปาวจโร จ นรูปาวจโร จ ธมฺมา อุปฺปชฺชนฺติ เหตุปจฺจยา. รูปาวจรํ เอกํ ขนฺธํ ปฏิจฺจ ตโย ขนฺธา จิตฺต\-สมุฏฺฐานญฺจ รูปํ ฯเปฯ เทฺว ขนฺเธ ฯเปฯ. ปฏิสนฺธิ\-กฺขเณ ฯเปฯ. (3)}

\proseitem{227}{นรูปาวจรํ ธมฺมํ ปฏิจฺจ นรูปาวจโร ธมฺโม อุปฺปชฺชติ เหตุปจฺจยา. นรูปาวจรํ เอกํ ขนฺธํ ปฏิจฺจ ตโย ขนฺธา จิตฺต\-สมุฏฺฐานญฺจ รูปํ ฯเปฯ เทฺว ขนฺเธ ฯเปฯ. \csromanfolio{142}ปฏิสนฺธิ\-กฺขเณ ฯเปฯ. เอกํ มหาภูตํ ฯเปฯ. นรูปาวจรํ ธมฺมํ ปฏิจฺจ รูปาวจโร ธมฺโม อุปฺปชฺชติ เหตุปจฺจยา. ปฏิสนฺธิ\-กฺขเณ วตฺถุํ ปฏิจฺจ รูปาวจรา ขนฺธา.  นรูปาวจรํ ธมฺมํ ปฏิจฺจ รูปาวจโร จ นรูปาวจโร จ ธมฺมา อุปฺปชฺชนฺติ เหตุปจฺจยา. ปฏิสนฺธิ\-กฺขเณ วตฺถุํ ปฏิจฺจ รูปาวจรา ขนฺธา. มหาภูเต ปฏิจฺจ กฏตฺตารูปํ. (3)}

\proseitem{228}{รูปาวจรญฺจ นรูปาวจรญฺจ ธมฺมํ ปฏิจฺจ รูปาวจโร ธมฺโม อุปฺปชฺชติ เหตุปจฺจยา. ปฏิสนฺธิ\-กฺขเณ รูปาวจรํ เอกํ ขนฺธญฺจ วตฺถุญฺจ ปฏิจฺจ ตโย ขนฺธา ฯเปฯ เทฺว ขนฺเธ จ ฯเปฯ. รูปาวจรญฺจ นรูปาวจรญฺจ ธมฺมํ ปฏิจฺจ นรูปาวจโร ธมฺโม อุปฺปชฺชติ เหตุปจฺจยา. รูปาวจเร ขนฺเธ จ มหาภูเต จ ปฏิจฺจ จิตฺต\-สมุฏฺฐานํ รูปํ. ปฏิสนฺธิ\-กฺขเณ รูปาวจเร ขนฺเธ จ มหาภูเต จ ปฏิจฺจ กฏตฺตารูปํ.  รูปาวจรญฺจ นรูปาวจรญฺจ ธมฺมํ ปฏิจฺจ รูปาวจโร จ นรูปาวจโร จ ธมฺมา อุปฺปชฺชนฺติ เหตุปจฺจยา. ปฏิสนฺธิ\-กฺขเณ รูปาวจรํ เอกํ ขนฺธญฺจ วตฺถุญฺจ ปฏิจฺจ ตโย ขนฺธา ฯเปฯ เทฺว ขนฺเธ จ ฯเปฯ. รูปาวจเร ขนฺเธ จ มหาภูเต จ ปฏิจฺจ กฏตฺตารูปํ. (สํขิตฺตํ.) (3)}

\csromanmarkh{1. ปจฺจยานุโลม}\csromanmarkhii{2. สงฺขฺยาวาร}\csromanheadernomark{2}{1. \textbf{ปจฺจยานุโลม 2. สงฺขฺยาวาร}}
\proseitem{229}{เหตุยา นว, อารมฺมเณ จตฺตาริ, อธิปติยา ปญฺจ, อนนฺตเร จตฺตาริ, สมนนฺตเร จตฺตาริ, สหชาเต นว, อญฺญมญฺเญ ฉ, นิสฺสเย นว, อุปนิสฺสเย จตฺตาริ, ปุเรชาเต เทฺว, อาเสวเน เทฺว, กมฺเม นว, วิปาเก นว, ฌาเน นว, มคฺเค นว, สมฺปยุตฺเต จตฺตาริ, วิปฺปยุตฺเต นว, อตฺถิยา นว, นตฺถิยา จตฺตาริ, วิคเต จตฺตาริ, อวิคเต นว.}

\csromanmarkh{2. ปจฺจย\-ปจฺจนีย}\csromanmarkhii{1. วิภงฺควาร}\csromanheadernomark{2}{2. ปจฺจย\-ปจฺจนีย 1. วิภงฺควาร}
\prose{นเหตุ นอารมฺมณ}

\proseitem{230}{นรูปาวจรํ ธมฺมํ ปฏิจฺจ นรูปาวจโร ธมฺโม อุปฺปชฺชติ นเหตุปจฺจยา. อเหตุกํ นรูปาวจรํ เอกํ ขนฺธํ ปฏิจฺจ ตโย ขนฺธา จิตฺต\-สมุฏฺฐานญฺจ รูปํ ฯเปฯ เทฺว ขนฺเธ ฯเปฯ. อเหตุก\-ปฏิสนฺธิ\-กฺขเณ ฯเปฯ. ขนฺเธ ปฏิจฺจ วตฺถุ, วตฺถุํ ปฏิจฺจ ขนฺธา. เอกํ มหาภูตํ ฯเปฯ. (ยาว อสญฺญสตฺตา.) วิจิกิจฺฉา\-สหคเต \csromanfolio{143}อุทฺธจฺจ\-สหคเต ขนฺเธ ปฏิจฺจ วิจิกิจฺฉา\-สหคโต อุทฺธจฺจ\-สหคโต โมโห. (1).  นอารมฺมณ\-ปจฺจยา ตีณิ.}

\prose{นอธิปตฺยาทิ}

\proseitem{231}{รูปาวจรํ ธมฺมํ ปฏิจฺจ รูปาวจโร ธมฺโม อุปฺปชฺชติ นอธิปติ\-ปจฺจยา. รูปาวจเร ขนฺเธ ปฏิจฺจ รูปาวจรา\-ธิปติ. วิปากํ รูปาวจรํ เอกํ ขนฺธํ ปฏิจฺจ ตโย ขนฺธา ฯเปฯ เทฺว ขนฺเธ ฯเปฯ. ปฏิสนฺธิ\-กฺขเณ ฯเปฯ. รูปาวจรํ ธมฺมํ ปฏิจฺจ นรูปาวจโร ธมฺโม อุปฺปชฺชติ นอธิปติ\-ปจฺจยา. วิปาเก รูปาวจเร ขนฺเธ ปฏิจฺจ จิตฺต\-สมุฏฺฐานํ รูปํ. ปฏิสนฺธิ\-กฺขเณ ฯเปฯ. รูปาวจรํ ธมฺมํ ปฏิจฺจ รูปาวจโร จ นรูปาวจโร จ ธมฺมา อุปฺปชฺชนฺติ นอธิปติ\-ปจฺจยา. วิปากํ รูปาวจรํ เอกํ ขนฺธํ ปฏิจฺจ ตโย ขนฺธา จิตฺต\-สมุฏฺฐานญฺจ รูปํ ฯเปฯ เทฺว ขนฺเธ ฯเปฯ. ปฏิสนฺธิ\-กฺขเณ ฯเปฯ. (3)}

\prose{นรูปาวจรํ ธมฺมํ ปฏิจฺจ นรูปาวจโร ธมฺโม อุปฺปชฺชติ นอธิปติ\-ปจฺจยา. นรูปาวจรํ เอกํ ขนฺธํ ปฏิจฺจ ตโย ขนฺธา จิตฺต\-สมุฏฺฐานญฺจ รูปํ ฯเปฯ เทฺว ขนฺเธ ฯเปฯ. ตีณิ. (นกามาวจรํ ปฏิจฺจ\-วาร\-สทิสํ นินฺนานํ, อิธ สพฺเพ มหาภูตา กาตพฺพา.) (3)}

\proseitem{232}{รูปาวจรญฺจ นรูปาวจรญฺจ ธมฺมํ ปฏิจฺจ รูปาวจโร ธมฺโม อุปฺปชฺชติ นอธิปติ\-ปจฺจยา. ปฏิสนฺธิ\-กฺขเณ รูปาวจรํ เอกํ ขนฺธญฺจ วตฺถุญฺจ ปฏิจฺจ ตโย ขนฺธา ฯเปฯ เทฺว ขนฺเธ จ ฯเปฯ. รูปาวจรญฺจ นรูปาวจรญฺจ ธมฺมํ ปฏิจฺจ นรูปาวจโร ธมฺโม อุปฺปชฺชติ นอธิปติ\-ปจฺจยา. วิปาเก รูปาวจเร ขนฺเธ จ มหาภูเต จ ปฏิจฺจ จิตฺต\-สมุฏฺฐานํ รูปํ. ปฏิสนฺธิ\-กฺขเณ ฯเปฯ. รูปาวจรญฺจ นรูปาวจรญฺจ ธมฺมํ ปฏิจฺจ รูปาวจโร จ นรูปาวจโร จ ธมฺมา อุปฺปชฺชนฺติ นอธิปติ\-ปจฺจยา. ปฏิสนฺธิ\-กฺขเณ รูปาวจรํ เอกํ ขนฺธญฺจ วตฺถุญฺจ ปฏิจฺจ ตโย ขนฺธา ฯเปฯ เทฺว ขนฺเธ จ ฯเปฯ. รูปาวจเร ขนฺเธ จ มหาภูเต จ ปฏิจฺจ กฏตฺตารูปํ. (3).  นอนนฺตร\-ปจฺจยา ฯเปฯ. นอุปนิสฺสยปจฺจยา.}

\csromanheader{2}{\textbf{นปุเรชาตาทิ}}
\proseitem{233}{รูปาวจรํ ธมฺมํ ปฏิจฺจ รูปาวจโร ธมฺโม อุปฺปชฺชติ นปุเรชาต\-ปจฺจยา. ปฏิสนฺธิ\-กฺขเณ รูปาวจรํ เอกํ ขนฺธํ ปฏิจฺจ ตโย ขนฺธา ฯเปฯ เทฺว ขนฺเธ ฯเปฯ. รูปาวจรํ ธมฺมํ ปฏิจฺจ นรูปาวจโร ธมฺโม อุปฺปชฺชติ นปุเรชาต\-ปจฺจยา. รูปาวจเร ขนฺเธ ปฏิจฺจ จิตฺต\-สมุฏฺฐานํ รูปํ. ปฏิสนฺธิ\-กฺขเณ ฯเปฯ. \csromanfolio{144}รูปาวจรํ ธมฺมํ ปฏิจฺจ รูปาวจโร จ นรูปาวจโร จ ธมฺมา อุปฺปชฺชนฺติ นปุเรชาต\-ปจฺจยา. ปฏิสนฺธิ\-กฺขเณ รูปาวจรํ เอกํ ขนฺธํ ปฏิจฺจ ตโย ขนฺธา กฏตฺตา จ รูปํ ฯเปฯ เทฺว ขนฺเธ ฯเปฯ. (3)}

\prose{นรูปาวจรํ ธมฺมํ ปฏิจฺจ นรูปาวจโร ธมฺโม อุปฺปชฺชติ นปุเรชาต\-ปจฺจยา. อรูเป นรูปาวจรํ เอกํ ขนฺธํ ปฏิจฺจ ตโย ขนฺธา ฯเปฯ เทฺว ขนฺเธ ฯเปฯ. นรูปาวจเร ขนฺเธ ปฏิจฺจ จิตฺต\-สมุฏฺฐานํ รูปํ. ปฏิสนฺธิ\-กฺขเณ ฯเปฯ. (ยาว อสญฺญสตฺตา. อิตเร ปญฺจปิ ปญฺหา. อนุโลมํ กาตพฺพํ.) นปจฺฉาชาต\-ปจฺจยา. นว.}

\prose{นอาเสวนาทิ}

\proseitem{234}{รูปาวจรํ ธมฺมํ ปฏิจฺจ รูปาวจโร ธมฺโม อุปฺปชฺชติ นอาเสวน\-ปจฺจยา. วิปากํ รูปาวจรํ เอกํ ขนฺธํ ปฏิจฺจ ตโย ขนฺธา ฯเปฯ เทฺว ขนฺเธ ฯเปฯ. ปฏิสนฺธิ\-กฺขเณ ฯเปฯ. รูปาวจรํ ธมฺมํ ปฏิจฺจ นรูปาวจโร ธมฺโม อุปฺปชฺชติ นอาเสวน\-ปจฺจยา. รูปาวจเร ขนฺเธ ปฏิจฺจ จิตฺต\-สมุฏฺฐานํ รูปํ. ปฏิสนฺธิ\-กฺขเณ ฯเปฯ. รูปาวจรํ ธมฺมํ ปฏิจฺจ รูปาวจโร จ นรูปาวจโร จ ธมฺมา อุปฺปชฺชนฺติ นอาเสวน\-ปจฺจยา. วิปากํ รูปาวจรํ เอกํ ขนฺธํ ปฏิจฺจ ตโย ขนฺธา จิตฺต\-สมุฏฺฐานญฺจ รูปํ ฯเปฯ เทฺว ขนฺเธ ฯเปฯ. ปฏิสนฺธิ\-กฺขเณ ฯเปฯ. (3)}

\prose{นรูปาวจรํ ธมฺมํ ปฏิจฺจ นรูปาวจโร ธมฺโม อุปฺปชฺชติ นอาเสวน\-ปจฺจยา. ตีณิ.}

\prose{รูปาวจรญฺจ นรูปาวจรญฺจ ธมฺมํ ปฏิจฺจ รูปาวจโร ธมฺโม อุปฺปชฺชติ นอาเสวน\-ปจฺจยา. (สํขิตฺตํ. มูลํ. อิตเรปิ ปญฺหา กาตพฺพา.) นกมฺมปจฺจยา. เทฺว ฯเปฯ. นสมฺปยุตฺต\-ปจฺจยา.}

\proseitem{235}{นรูปาวจรํ ธมฺมํ ปฏิจฺจ นรูปาวจโร ธมฺโม อุปฺปชฺชติ นวิปฺปยุตฺต\-ปจฺจยา. อรูเป นรูปาวจรํ เอกํ ขนฺธํ ปฏิจฺจ ตโย ขนฺธา ฯเปฯ เทฺว ขนฺเธ ฯเปฯ. พาหิรํ. อาหารสมุฏฺฐานํ. อุตุสมุฏฺฐานํ. อสญฺญสตฺตานํ ฯเปฯ. (สํขิตฺตํ.)}

\csromanmarkh{2. ปจฺจย\-ปจฺจนีย}\csromanmarkhii{2. สงฺขฺยาวาร}\csromanheadernomark{2}{2. ปจฺจย\-ปจฺจนีย 2. สงฺขฺยาวาร}
\proseitem{236}{นเหตุยา เอกํ, นอารมฺมเณ ตีณิ, นอธิปติยา นว, นอนนฺตเร ตีณิ, นสมนนฺตเร ตีณิ, นอุปนิสฺสเย ตีณิ, นปุเรชาเต นว, นปจฺฉาชาเต นว, นอาเสวเน นว, นกมฺเม เทฺว, นวิปาเก}

\prose{\csromanfolio{145}ปญฺจ, นอาหาเร เอกํ, ไนนฺทฺริเย เอกํ, นฌาเน เอกํ, นมคฺเค เอกํ, นสมฺปยุตฺเต ตีณิ, นวิปฺปยุตฺเต เอกํ, โนนตฺถิยา ตีณิ, โนวิคเต ตีณิ.}

\vspace{\tipitakaparskip}
\csromancenter{(เอวํ อิตเร เทฺว คณนาปิ สหชาตวาโรปิ กาตพฺโพ.)}
\csromanmarkh{94. รูปาวจรทุก}\csromanmarkhii{3. ปจฺจยวาร}\csromanheadernomark{2}{94. รูปาวจรทุก 3. ปจฺจยวาร}
\proseitem{237}{รูปาวจรํ ธมฺมํ ปจฺจยา รูปาวจโร ธมฺโม อุปฺปชฺชติ เหตุปจฺจยา. ตีณิ. (ปฏิจฺจสทิสํ.)}

\prose{นรูปาวจรํ ธมฺมํ ปจฺจยา นรูปาวจโร ธมฺโม อุปฺปชฺชติ เหตุปจฺจยา. นรูปาวจรํ เอกํ ขนฺธํ ปจฺจยา ตโย ขนฺธา จิตฺต\-สมุฏฺฐานญฺจ รูปํ ฯเปฯ เทฺว ขนฺเธ ฯเปฯ. (ยาว อชฺฌตฺติกํ มหาภูตํ.) วตฺถุํ ปจฺจยา นรูปาวจรา ขนฺธา.  นรูปาวจรํ ธมฺมํ ปจฺจยา รูปาวจโร ธมฺโม อุปฺปชฺชติ เหตุปจฺจยา. วตฺถุํ ปจฺจยา รูปาวจรา ขนฺธา. ปฏิสนฺธิ\-กฺขเณ ฯเปฯ. นรูปาวจรํ ธมฺมํ ปจฺจยา รูปาวจโร จ นรูปาวจโร จ ธมฺมา อุปฺปชฺชนฺติ เหตุปจฺจยา. วตฺถุํ ปจฺจยา รูปาวจรา ขนฺธา. มหาภูเต ปจฺจยา จิตฺต\-สมุฏฺฐานํ รูปํ. ปฏิสนฺธิ\-กฺขเณ ฯเปฯ. (3)}

\proseitem{238}{รูปาวจรญฺจ นรูปาวจรญฺจ ธมฺมํ ปจฺจยา รูปาวจโร ธมฺโม อุปฺปชฺชติ เหตุปจฺจยา. รูปาวจรํ เอกํ ขนฺธญฺจ วตฺถุญฺจ ปจฺจยา ตโย ขนฺธา ฯเปฯ เทฺว ขนฺเธ จ ฯเปฯ. ปฏิสนฺธิ\-กฺขเณ ฯเปฯ. รูปาวจรญฺจ นรูปาวจรญฺจ ธมฺมํ ปจฺจยา นรูปาวจโร ธมฺโม อุปฺปชฺชติ เหตุปจฺจยา. รูปาวจเร ขนฺเธ จ มหาภูเต จ ปจฺจยา จิตฺต\-สมุฏฺฐานํ รูปํ. ปฏิสนฺธิ\-กฺขเณ ฯเปฯ. รูปาวจรญฺจ นรูปาวจรญฺจ ธมฺมํ ปจฺจยา รูปาวจโร จ นรูปาวจโร จ ธมฺมา อุปฺปชฺชนฺติ เหตุปจฺจยา. รูปาวจรํ เอกํ ขนฺธญฺจ วตฺถุญฺจ ปจฺจยา ตโย ขนฺธา ฯเปฯ เทฺว ขนฺเธ จ ฯเปฯ. รูปาวจเร ขนฺเธ จ มหาภูเต จ ปจฺจยา จิตฺต\-สมุฏฺฐานํ รูปํ. ปฏิสนฺธิ\-กฺขเณ ฯเปฯ. (3) (สํขิตฺตํ.)}

\csromanmarkh{1. ปจฺจยานุโลม}\csromanmarkhii{2. สงฺขฺยาวาร}\csromanheadernomark{2}{1. \textbf{ปจฺจยานุโลม 2. สงฺขฺยาวาร}}
\proseitem{239}{\csromanfolio{146}เหตุยา นว, อารมฺมเณ จตฺตาริ, อธิปติยา นว, อนนฺตเร จตฺตาริ, สมนนฺตเร จตฺตาริ, สหชาเต นว, อญฺญมญฺเญ ฉ, นิสฺสเย นว, อุปนิสฺสเย จตฺตาริ, ปุเรชาเต จตฺตาริ, อาเสวเน จตฺตาริ, กมฺเม นว ฯเปฯ อวิคเต นว.}

\csromanmarkh{2. ปจฺจย\-ปจฺจนีย}\csromanmarkhii{1. วิภงฺควาร}\csromanheadernomark{2}{2. ปจฺจย\-ปจฺจนีย 1. วิภงฺควาร}
\proseitem{240}{นรูปาวจรํ ธมฺมํ ปจฺจยา นรูปาวจโร ธมฺโม อุปฺปชฺชติ นเหตุปจฺจยา. อเหตุกํ นรูปาวจรํ เอกํ ขนฺธํ ปจฺจยา ตโย ขนฺธา จิตฺต\-สมุฏฺฐานญฺจ รูปํ ฯเปฯ เทฺว ขนฺเธ ฯเปฯ. อเหตุก\-ปฏิสนฺธิ\-กฺขเณ ฯเปฯ. (ยาว อสญฺญสตฺตา.) จกฺขายตนํ ปจฺจยา จกฺขุวิญฺญาณํ ฯเปฯ. กายายตนํ ปจฺจยา กายวิญฺญาณํ วตฺถุํ ปจฺจยา อเหตุกา นรูปาวจรา ขนฺธา. วิจิกิจฺฉา\-สหคเต อุทฺธจฺจ\-สหคเต ขนฺเธ จ วตฺถุญฺจ ปจฺจยา วิจิกิจฺฉา\-สหคโต อุทฺธจฺจ\-สหคโต โมโห. (1) (สํขิตฺตํ.)}

\prose{นอธิปติ}

\proseitem{241}{รูปาวจรํ ธมฺมํ ปจฺจยา รูปาวจโร ธมฺโม อุปฺปชฺชติ นอธิปติ\-ปจฺจยา. ตีณิ. (ปฏิจฺจ\-วาร\-สทิสํ.)}

\prose{นรูปาวจรํ ธมฺมํ ปจฺจยา นรูปาวจโร ธมฺโม อุปฺปชฺชติ นอธิปติ\-ปจฺจยา. (ปฏิจฺจ\-วาร\-สทิสํ.). นรูปาวจรํ ธมฺมํ ปจฺจยา รูปาวจโร ธมฺโม อุปฺปชฺชติ นอธิปติ\-ปจฺจยา. วตฺถุํ ปจฺจยา รูปาวจรา\-ธิปติ. วตฺถุํ ปจฺจยา วิปากา รูปาวจรา ขนฺธา. ปฏิสนฺธิ\-กฺขเณ ฯเปฯ. นรูปาวจรํ ธมฺมํ ปจฺจยา รูปาวจโร จ นรูปาวจโร จ ธมฺมา อุปฺปชฺชนฺติ นอธิปติ\-ปจฺจยา. วตฺถุํ ปจฺจยา วิปากา รูปาวจรา ขนฺธา. มหาภูเต ปจฺจยา จิตฺต\-สมุฏฺฐานํ รูปํ. ปฏิสนฺธิ\-กฺขเณ ฯเปฯ. (3)}

\proseitem{242}{รูปาวจรญฺจ นรูปาวจรญฺจ ธมฺมํ ปจฺจยา รูปาวจโร ธมฺโม อุปฺปชฺชติ นอธิปติ\-ปจฺจยา. รูปาวจเร ขนฺเธ จ วตฺถุญฺจ ปจฺจยา รูปาวจรา\-ธิปติ. \csromanfolio{147}วิปากํ รูปาวจรํ เอกํ ขนฺธญฺจ วตฺถุญฺจ ปจฺจยา ตโย ขนฺธา ฯเปฯ เทฺว ขนฺเธ จ ฯเปฯ. ปฏิสนฺธิ\-กฺขเณ รูปาวจรํ เอกํ ขนฺธญฺจ วตฺถุญฺจ ปจฺจยา ตโย ขนฺธา ฯเปฯ เทฺว ขนฺเธ จ ฯเปฯ. รูปาวจรญฺจ นรูปาวจรญฺจ ธมฺมํ ปจฺจยา นรูปาวจโร ธมฺโม อุปฺปชฺชติ นอธิปติ\-ปจฺจยา. วิปาเก รูปาวจเร ขนฺเธ จ มหาภูเต จ ปจฺจยา จิตฺต\-สมุฏฺฐานํ รูปํ. ปฏิสนฺธิ\-กฺขเณ ฯเปฯ. รูปาวจรญฺจ นรูปาวจรญฺจ ธมฺมํ ปจฺจยา รูปาวจโร จ นรูปาวจโร จ ธมฺมา อุปฺปชฺชนฺติ นอธิปติ\-ปจฺจยา. วิปากํ รูปาวจรํ เอกํ ขนฺธญฺจ วตฺถุญฺจ ปจฺจยา ตโย ขนฺธา ฯเปฯ เทฺว ขเธ จ ฯเปฯ. วิปาเก รูปาวจเร ขนฺเธ จ มหาภูเต จ ปจฺจยา จิตฺต\-สมุฏฺฐานํ รูปํ. ปฏิสนฺธิ\-กฺขเณ ฯเปฯ. (3) (สํขิตฺตํ.)}

\csromanmarkh{2. ปจฺจย\-ปจฺจนีย}\csromanmarkhii{2. สงฺขฺยาวาร}\csromanheadernomark{2}{2. ปจฺจย\-ปจฺจนีย 2. สงฺขฺยาวาร}
\proseitem{243}{นเหตุยา เอกํ, นอารมฺมเณ ตีณิ, นอธิปติยา นว, นอนนฺตเร ตีณิ ฯเปฯ นอุปนิสฺสเย ตีณิ, นปุเรชาเต นว, นปจฺฉาชาเต นว, นอาเสวเน นว, (สุทฺธิเก อรูเป จ มิสฺสเก จ “วิปากนฺ”ติ นิยาเมตพฺพํ.) นกมฺเม จตฺตาริ, นวิปาเก นว, นอาหาเร เอกํ, นอินฺทฺริเย เอกํ, นฌาเน เอกํ, นมคฺเค เอกํ, นสมฺปยุตฺเต ตีณิ, นวิปฺปยุตฺเต เอกํ, โนนตฺถิยา ตีณิ, โนวิคเต ตีณิ.}

\vspace{\tipitakaparskip}
\csromancenter{(เอวํ อิตเร เทฺว คณนาปิ นิสฺสยวาโรปิ กาตพฺโพ.)}
\csromanmarkh{94. รูปาวจรทุก}\csromanmarkhii{5. สํสฏฺฐวาร}\csromanheadernomark{2}{94. รูปาวจรทุก 5. สํสฏฺฐวาร}
\proseitem{244}{รูปาวจรํ ธมฺมํ สํสฏฺโฐ รูปาวจโร ธมฺโม อุปฺปชฺชติ เหตุปจฺจยา. รูปาวจรํ เอกํ ขนฺธํ สํสฏฺฐา ตโย ขนฺธา ฯเปฯ เทฺว ขนฺเธ ฯเปฯ. ปฏิสนฺธิ\-กฺขเณ ฯเปฯ. (1)}

\prose{\csromanfolio{148}นรูปาวจรํ ธมฺมํ สํสฏฺโฐ นรูปาวจโร ธมฺโม อุปฺปชฺชติ เหตุปจฺจยา. นรูปาวจรํ เอกํ ขนฺธํ สํสฏฺฐา ตโย ขนฺธา ฯเปฯ เทฺว ขนฺเธ ฯเปฯ. ปฏิสนฺธิ\-กฺขเณ ฯเปฯ. (1)}

\prose{เหตุยา เทฺว, อารมฺมเณ เทฺว, (สพฺพตฺถ เทฺว.) อวิคเต เทฺว. อนุโลมํ.}

\prose{นรูปาวจรํ ธมฺมํ สํสฏฺโฐ นรูปาวจโร ธมฺโม อุปฺปชฺชติ นเหตุปจฺจยา. (สํขิตฺตํ.)}

\prose{นเหตุยา เอกํ, นอธิปติยา เทฺว, นปุเรชาเต เทฺว, นปจฺฉาชาเต เทฺว, นอาเสวเน เทฺว, นกมฺเม เทฺว, นวิปาเก เทฺว, นฌาเน เอกํ, นมคฺเค เอกํ, นวิปฺปยุตฺเต เอกํ.}

\vspace{\tipitakaparskip}
\csromancenter{ปจฺจนียํ.}
\csromancenter{(เอวํ อิตเร เทฺว คณนาปิ สมฺปยุตฺตวาโรปิ กาตพฺโพ.)}
\csromanmarkh{94. รูปาวจรทุก}\csromanmarkhii{7. ปญฺหาวาร}\csromanheadernomark{2}{94. \textbf{รูปาวจรทุก 7. ปญฺหาวาร}}
\proseitem{245}{รูปาวจโร ธมฺโม รูปาวจรสฺส ธมฺมสฺส เหตุปจฺจเยน ปจฺจโย. รูปาวจรา เหตู สมฺปยุตฺตกานํ ขนฺธานํ เหตุปจฺจเยน ปจฺจโย. ปฏิสนฺธิ\-กฺขเณ ฯเปฯ. (มูลํ กาตพฺพํ.) รูปาวจรา เหตู จิตฺต\-สมุฏฺฐานานํ รูปานํ เหตุปจฺจเยน ปจฺจโย. ปฏิสนฺธิ\-กฺขเณ ฯเปฯ. (มูลํ กาตพฺพํ.) รูปาวจรา เหตู สมฺปยุตฺตกานํ ขนฺธานํ จิตฺตสมุฏฺฐานานญฺจ รูปานํ เหตุปจฺจเยน ปจฺจโย. ปฏิสนฺธิ\-กฺขเณ ฯเปฯ. (3)}

\prose{นรูปาวจโร ธมฺโม นรูปาวจรสฺส ธมฺมสฺส เหตุปจฺจเยน ปจฺจโย. นรูปาวจรา เหตู สมฺปยุตฺตกานํ ขนฺธานํ จิตฺตสมุฏฺฐานานญฺจ รูปานํ เหตุปจฺจเยน ปจฺจโย. ปฏิสนฺธิ\-กฺขเณ ฯเปฯ. (1)}

\csromanheader{2}{94. รูปาวจรทุก \textbf{อารมฺมณ}}
\proseitem{246}{\csromanfolio{149}รูปาวจโร ธมฺโม รูปาวจรสฺส ธมฺมสฺส อารมฺมณ\-ปจฺจเยน ปจฺจโย. เจโตปริย\-ญาเณน รูปาวจร\-จิตฺต\-สมงฺคิสฺส จิตฺตํ ชานาติ. รูปาวจรา ขนฺธา อิทฺธิวิธ\-ญาณสฺส, เจโตปริย\-ญาณสฺส, ปุพฺเพ\-นิวาสา\-นุสฺสติ\-ญาณสฺส, ยถากมฺมูปคญาณสฺส, อนาคตํสญาณสฺส อารมฺมณ\-ปจฺจเยน ปจฺจโย.  รูปาวจโร ธมฺโม นรูปาวจรสฺส ธมฺมสฺส อารมฺมณ\-ปจฺจเยน ปจฺจโย. ปฐมํ ฌานํ ปจฺจเวกฺขติ ฯเปฯ. จตุตฺถํ ฌานํ ปจฺจเวกฺขติ. ทิพฺพํ จกฺขุํ ปจฺจเวกฺขติ. ทิพฺพํ โสตธาตุํ ฯเปฯ. อิทฺธิวิธ\-ญาณํ ฯเปฯ. เจโตปริย\-ญาณํ ฯเปฯ. ปุพฺเพ\-นิวาสา\-นุสฺสติ\-ญาณํ ฯเปฯ. ยถากมฺมูปคญาณํ ฯเปฯ. อนาคตํสญาณํ ปจฺจเวกฺขติ. รูปาวจเร ขนฺเธ อนิจฺจโต ฯเปฯ โทมนสฺสํ อุปฺปชฺชติ. (2)}

\proseitem{247}{นรูปาวจโร ธมฺโม นรูปาวจรสฺส ธมฺมสฺส อารมฺมณ\-ปจฺจเยน ปจฺจโย. ทานํ ทตฺวา, สีลํ ฯเปฯ อุโปสถกมฺมํ ฯเปฯ. ปุพฺเพ สุจิณฺณานิ ปจฺจเวกฺขติ, อสฺสาเทติ อภินนฺทติ, ตํ อารพฺภ ราโค ฯเปฯ โทมนสฺสํ อุปฺปชฺชติ. อริยา มคฺคา วุฏฺฐหิตฺวา มคฺคํ ปจฺจเวกฺขนฺติ, ผลํ ปจฺจเวกฺขนฺติ, นิพฺพานํ ปจฺจเวกฺขนฺติ. นิพฺพานํ โคตฺรภุสฺส, โวทานสฺส, มคฺคสฺส, ผลสฺส, อาวชฺชนาย อารมฺมณ\-ปจฺจเยน ปจฺจโย. อริยา ปหีเน กิเลเส ฯเปฯ. วิกฺขมฺภิเต กิเลเส ฯเปฯ. ปุพฺเพ ฯเปฯ. จกฺขุํ ฯเปฯ. วตฺถุํ นรูปาวจเร ขนฺเธ อนิจฺจโต ฯเปฯ โทมนสฺสํ อุปฺปชฺชติ. อากาสา\-นญฺจา\-ยตนํ วิญฺญาณญฺจา\-ยตนสฺส ฯเปฯ. อากิญฺจญฺญา\-ยตนํ เนว\-สญฺญา\-นา\-สญฺญา\-ยตนสฺส ฯเปฯ. รูปายตนํ จกฺขุ\-วิญฺญาณสฺส ฯเปฯ. โผฏฺฐพฺพา\-ยตนํ กายวิญฺญาณสฺส อารมฺมณ\-ปจฺจเยน ปจฺจโย. (1)}

\prose{นรูปาวจโร ธมฺโม รูปาวจรสฺส ธมฺมสฺส อารมฺมณ\-ปจฺจเยน ปจฺจโย. ทิพฺเพน จกฺขุนา รูปํ ปสฺสติ, ทิพฺพาย โสตธาตุยา สทฺทํ สุณาติ. เจโตปริย\-ญาเณน นรูปาวจร\-จิตฺต\-สมงฺคิสฺส จิตฺตํ ชานาติ. นรูปาวจรา ขนฺธา อิทฺธิวิธ\-ญาณสฺส, เจโตปริย\-ญาณสฺส, ปุพฺเพ\-นิวาสา\-นุสฺสติ\-ญาณสฺส, ยถากมฺมูปคญาณสฺส, อนาคตํสญาณสฺส อารมฺมณ\-ปจฺจเยน ปจฺจโย. (2)}

\csromanheader{2}{\textbf{อธิปติ}}
\proseitem{248}{รูปาวจโร ธมฺโม รูปาวจรสฺส ธมฺมสฺส อธิปติ\-ปจฺจเยน ปจฺจโย. สหชาตา\-ธิปติ\csromandash{}รูปาวจรา\-ธิปติ สมฺปยุตฺตกานํ ขนฺธานํ อธิปติ\-ปจฺจเยน \csromanfolio{150}ปจฺจโย.  รูปาวจโร ธมฺโม นรูปาวจรสฺส ธมฺมสฺส อธิปติ\-ปจฺจเยน ปจฺจโย. อารมฺมณา\-ธิปติ, สหชาตา\-ธิปติ.  อารมฺมณา\-ธิปติ\csromandash{}ปฐมํ ฌานํ ครุํ กตฺวา ปจฺจเวกฺขติ, อสฺสาเทติ อภินนฺทติ, ตํ ครุํ กตฺวา ราโค อุปฺปชฺชติ, ทิฏฺฐิ อุปฺปชฺชติ ฯเปฯ. จตุตฺถํ ฌานํ ฯเปฯ. ทิพฺพํ จกฺขุํ ฯเปฯ. ทิพฺพํ โสตธาตุํ ฯเปฯ. อิทฺธิวิธ\-ญาณํ ฯเปฯ. เจโตปริย\-ญาณํ ฯเปฯ. ปุพฺเพ\-นิวาสา\-นุสฺสติ\-ญาณํ ฯเปฯ. ยถากมฺมูปคญาณํ ฯเปฯ. อนาคตํสญาณํ ครุํ กตฺวา ปจฺจเวกฺขติ, อสฺสาเทติ อภินนฺทติ, ตํ ครุํ กตฺวา ราโค อุปฺปชฺชติ, ทิฏฺฐิ อุปฺปชฺชติ. รูปาวจเร ขนฺเธ ครุํ กตฺวา อสฺสาเทติ อภินนฺทติ, ตํ ครุํ กตฺวา ราโค อุปฺปชฺชติ. ทิฏฺฐิ อุปฺปชฺชติ. สหชาตา\-ธิปติ\csromandash{}รูปาวจรา\-ธิปติ จิตฺต\-สมุฏฺฐานานํ รูปานํ อธิปติ\-ปจฺจเยน ปจฺจโย.  รูปาวจโร ธมฺโม รูปาวจรสฺส จ นรูปาวจรสฺส จ ธมฺมสฺส อธิปติ\-ปจฺจเยน ปจฺจโย. สหชาตา\-ธิปติ\csromandash{}รูปาวจรา\-ธิปติ สมฺปยุตฺตกานํ ขนฺธานํ จิตฺตสมุฏฺฐานานญฺจ รูปานํ อธิปติ\-ปจฺจเยน ปจฺจโย. (3)}

\proseitem{249}{นรูปาวจโร ธมฺโม นรูปาวจรสฺส ธมฺมสฺส อธิปติ\-ปจฺจเยน ปจฺจโย. อารมฺมณา\-ธิปติ, สหชาตา\-ธิปติ.  อารมฺมณา\-ธิปติ\csromandash{}ทานํ ฯเปฯ สีลํ ฯเปฯ อุโปสถกมฺมํ ฯเปฯ. ปุพฺเพ สุจิณฺณานิ ครุํ กตฺวา ปจฺจเวกฺขติ, อสฺสาเทติ อภินนฺทติ, ตํ ครุํ กตฺวา ราโค อุปฺปชฺชติ. ทิฏฺฐิ อุปฺปชฺชติ. อริยา มคฺคา วุฏฺฐหิตฺวา มคฺคํ ครุํ กตฺวา ฯเปฯ ผลํ ฯเปฯ นิพฺพานํ ครุํ กตฺวา ปจฺจเวกฺขนฺติ. นิพฺพานํ โคตฺรภุสฺส, โวทานสฺส, มคฺคสฺส, ผลสฺส อธิปติ\-ปจฺจเยน ปจฺจโย. จกฺขุํ ฯเปฯ. วตฺถุํ นรูปาวจเร ขนฺเธ ครุํ กตฺวา อสฺสาเทติ อภินนฺทติ, ตํ ครุํ กตฺวา ราโค อุปฺปชฺชติ. ทิฏฺฐิ อุปฺปชฺชติ. อากาสา\-นญฺจา\-ยตนํ ครุํ กตฺวา ปจฺจเวกฺขติ ฯเปฯ. เนว\-สญฺญา\-นา\-สญฺญา\-ยตนํ ครุํ กตฺวา ปจฺจเวกฺขติ. สหชาตา\-ธิปติ\csromandash{} นรูปาวจรา\-ธิปติ สมฺปยุตฺตกานํ ขนฺธานํ จิตฺตสมุฏฺฐานานญฺจ รูปานํ อธิปติ\-ปจฺจเยน ปจฺจโย. (1)}

\proseitem{250}{รูปาวจโร ธมฺโม รูปาวจรสฺส ธมฺมสฺส อนนฺตร\-ปจฺจเยน ปจฺจโย. ปุริมา ปุริมา รูปาวจรา ขนฺธา ปจฺฉิมานํ ปจฺฉิมานํ รูปาวจรานํ ขนฺธานํ อนนฺตร\-ปจฺจเยน ปจฺจโย.  รูปาวจโร ธมฺโม นรูปาวจรสฺส ธมฺมสฺส อนนฺตร\-ปจฺจเยน ปจฺจโย. รูปาวจรํ จุติจิตฺตํ นรูปาวจรสฺส \csromanfolio{151}อุปปตฺติ\-จิตฺตสฺส อนนฺตร\-ปจฺจเยน ปจฺจโย. รูปาวจรํ ภวงฺคํ อาวชฺชนาย. รูปาวจรา ขนฺธา นรูปาวจรสฺส วุฏฺฐานสฺส อนนฺตร\-ปจฺจเยน ปจฺจโย. (2)}

\proseitem{251}{นรูปาวจโร ธมฺโม นรูปาวจรสฺส ธมฺมสฺส อนนฺตร\-ปจฺจเยน ปจฺจโย. ปุริมา ปุริมา นรูปาวจรา ขนฺธา ปจฺฉิมานํ ปจฺฉิมานํ นรูปาวจรานํ ขนฺธานํ อนนฺตร\-ปจฺจเยน ปจฺจโย. อนุโลมํ โคตฺรภุสฺส ฯเปฯ. เนว\-สญฺญา\-นา\-สญฺญา\-ยตนํ ผลสมาปตฺติยา อนนฺตร\-ปจฺจเยน ปจฺจโย.  นรูปาวจโร ธมฺโม รูปาวจรสฺส ธมฺมสฺส อนนฺตร\-ปจฺจเยน ปจฺจโย. นรูปาวจรํ จุติจิตฺตํ รูปาวจรสฺส อุปปตฺติ\-จิตฺตสฺส อนนฺตร\-ปจฺจเยน ปจฺจโย. นรูปาวจรา ขนฺธา รูปาวจรสฺส วุฏฺฐานสฺส อนนฺตร\-ปจฺจเยน ปจฺจโย. ปฐมสฺส ฌานสฺส ปริกมฺมํ ปฐมสฺส ฌานสฺส อนนฺตร\-ปจฺจเยน ปจฺจโย ฯเปฯ. จตุตฺถสฺส ฌานสฺส ฯเปฯ. ทิพฺพสฺส จกฺขุสฺส ฯเปฯ. ทิพฺพาย โสตธาตุยา ฯเปฯ. อิทฺธิวิธ\-ญาณสฺส ฯเปฯ. เจโตปริย\-ญาณสฺส ฯเปฯ. ปุพฺเพ\-นิวาสา\-นุสฺสติ\-ญาณสฺส ฯเปฯ. ยถากมฺมูปคญาณสฺส ฯเปฯ. อนาคตํสญาณสฺส ปริกมฺมํ อนาคตํสญาณสฺส อนนฺตร\-ปจฺจเยน ปจฺจโย. (2)}

\prose{สมนนฺตร\-ปจฺจเยน ปจฺจโย. สหชาต\-ปจฺจเยน ปจฺจโย. สตฺต. อญฺญมญฺญ\-ปจฺจเยน ปจฺจโย. ฉ. นิสฺสย\-ปจฺจเยน ปจฺจโย. สตฺต.}

\csromanheader{2}{\textbf{อุปนิสฺสย}}
\proseitem{252}{รูปาวจโร ธมฺโม รูปาวจรสฺส ธมฺมสฺส อุปนิสฺสย\-ปจฺจเยน ปจฺจโย. อนนฺตรู\-ปนิสฺสโย, ปกตู\-ปนิสฺสโย ฯเปฯ. ปกตู\-ปนิสฺสโย\csromandash{}รูปาวจรํ สทฺธํ อุปนิสฺสาย รูปาวจรํ ฌานํ อุปฺปาเทติ, อภิญฺญํ ฯเปฯ สมาปตฺติํ อุปฺปาเทติ. รูปาวจรํ สีลํ ฯเปฯ ปญฺญํ อุปนิสฺสาย รูปาวจรํ ฌานํ ฯเปฯ อภิญฺญํ ฯเปฯ สมาปตฺติํ อุปฺปาเทติ. รูปาวจรา สทฺธา ฯเปฯ ปญฺญา รูปาวจราย สทฺธาย ฯเปฯ ปญฺญาย อุปนิสฺสย\-ปจฺจเยน ปจฺจโย. ปฐมํ ฌานํ ทุติยสฺส ฌานสฺส อุปนิสฺสย\-ปจฺจเยน ปจฺจโย. ทุติยํ ฌานํ ตติยสฺส ฌานสฺส ฯเปฯ. ตติยํ ฌานํ จตุตฺถสฺส ฌานสฺส อุปนิสฺสย\-ปจฺจเยน ปจฺจโย. (1)}

\prose{รูปาวจโร ธมฺโม นรูปาวจรสฺส ธมฺมสฺส อุปนิสฺสย\-ปจฺจเยน ปจฺจโย. อารมฺมณู\-ปนิสฺสโย, อนนฺตรู\-ปนิสฺสโย, ปกตู\-ปนิสฺสโย ฯเปฯ. ปกตู\-ปนิสฺสโย\csromandash{}รูปาวจรํ สทฺธํ อุปนิสฺสาย ทานํ เทติ. สีลํ ฯเปฯ \csromanfolio{152}อุโปสถกมฺมํ กโรติ. นรูปาวจรํ ฌานํ ฯเปฯ วิปสฺสนํ ฯเปฯ มคฺคํ ฯเปฯ สมาปตฺติํ อุปฺปาเทติ. มานํ ชปฺเปติ, ทิฏฺฐิํ คณฺหาติ. รูปาวจรํ สีลํ ฯเปฯ ปญฺญํ อุปนิสฺสาย ทานํ เทติ ฯเปฯ มานํ ชปฺเปติ, ทิฏฺฐิํ คณฺหาติ. รูปาวจรา สทฺธา ฯเปฯ ปญฺญา นรูปาวจราย สทฺธาย ฯเปฯ ปญฺญาย. ราคสฺส ฯเปฯ ปตฺถนาย. กายิกสฺส สุขสฺส. กายิกสฺส ทุกฺขสฺส. มคฺคสฺส, ผลสมาปตฺติยา อุปนิสฺสย\-ปจฺจเยน ปจฺจโย. (2)}

\proseitem{253}{นรูปาวจโร ธมฺโม นรูปาวจรสฺส ธมฺมสฺส อุปนิสฺสย\-ปจฺจเยน ปจฺจโย. อารมฺมณู\-ปนิสฺสโย, อนนฺตรู\-ปนิสฺสโย, ปกตู\-ปนิสฺสโย ฯเปฯ. ปกตู\-ปนิสฺสโย\csromandash{}นรูปาวจรํ สทฺธํ อุปนิสฺสาย ทานํ ฯเปฯ สีลํ ฯเปฯ อุโปสถกมฺมํ กโรติ. นรูปาวจรํ ฌานํ ฯเปฯ วิปสฺสนํ ฯเปฯ มคฺคํ ฯเปฯ สมาปตฺติํ อุปฺปาเทติ. มานํ ชปฺเปติ, ทิฏฺฐิํ คณฺหาติ. นรูปาวจรํ สีลํ ฯเปฯ ปญฺญํ. ราคํ ฯเปฯ. โภชนํ. เสนาสนํ อุปนิสฺสาย ทานํ เทติ ฯเปฯ ปาณํ หนติ ฯเปฯ สํฆํ ภินฺทติ. นรูปาวจรา สทฺธา ฯเปฯ ปญฺญา. ราโค ฯเปฯ. เสนาสนํ นรูปาวจราย สทฺธาย ฯเปฯ ปญฺญาย. ราคสฺส ฯเปฯ ปตฺถนาย. กายิกสฺส สุขสฺส. กายิกสฺส ทุกฺขสฺส. มคฺคสฺส, ผลสมาปตฺติยา อุปนิสฺสย\-ปจฺจเยน ปจฺจโย. (1)}

\prose{นรูปาวจโร ธมฺโม รูปาวจรสฺส ธมฺมสฺส อุปนิสฺสย\-ปจฺจเยน ปจฺจโย. อนนฺตรู\-ปนิสฺสโย, ปกตู\-ปนิสฺสโย ฯเปฯ. ปกตู\-ปนิสฺสโย\csromandash{} นรูปาวจรํ สทฺธํ อุปนิสฺสาย รูปาวจรํ ฌานํ ฯเปฯ อภิญฺญํ ฯเปฯ สมาปตฺติํ อุปฺปาเทติ. นรูปาวจรํ สีลํ ฯเปฯ. เสนาสนํ อุปนิสฺสาย รูปาวจรํ ฌานํ ฯเปฯ อภิญฺญํ ฯเปฯ สมาปตฺติํ อุปฺปาเทติ. นรูปาวจรา สทฺธา ฯเปฯ. เสนาสนํ รูปาวจราย สทฺธาย ฯเปฯ ปญฺญาย อุปนิสฺสย\-ปจฺจเยน ปจฺจโย. ปฐมสฺส ฌานสฺส ปริกมฺมํ ปฐมสฺส ฌานสฺส อุปนิสฺสย\-ปจฺจเยน ปจฺจโย ฯเปฯ. จตุตฺถสฺส ฌานสฺส ฯเปฯ. ทิพฺพสฺส จกฺขุสฺส ปริกมฺมํ ฯเปฯ. ทิพฺพาย โสตธาตุยา ฯเปฯ. อิทฺธิวิธ\-ญาณสฺส ฯเปฯ. เจโตปริย\-ญาณสฺส ฯเปฯ. ปุพฺเพ\-นิวาสา\-นุสฺสติ\-ญาณสฺส ฯเปฯ. ยถากมฺมูปคญาณสฺส ฯเปฯ. อนาคตํสญาณสฺส ปริกมฺมํ อนาคตํสญาณสฺส อุปนิสฺสย\-ปจฺจเยน ปจฺจโย. (2)}

\csromanheader{2}{\textbf{ปุเรชาตาทิ}}
\proseitem{254}{นรูปาวจโร ธมฺโม นรูปาวจรสฺส ธมฺมสฺส ปุเรชาต\-ปจฺจเยน ปจฺจโย. อารมฺมณ\-ปุเรชาตํ, วตฺถุ\-ปุเรชาตํ.  อารมฺมณ\-ปุเรชาตํ\csromandash{} \csromanfolio{153}จกฺขุํ ฯเปฯ. วตฺถุํ อนิจฺจโต ฯเปฯ โทมนสฺสํ อุปฺปชฺชติ. รูปายตนํ จกฺขุ\-วิญฺญาณสฺส ฯเปฯ. โผฏฺฐพฺพา\-ยตนํ กายวิญฺญาณสฺส ปุเรชาต\-ปจฺจเยน ปจฺจโย. วตฺถุ\-ปุเรชาตํ\csromandash{}จกฺขายตนํ จกฺขุ\-วิญฺญาณสฺส ฯเปฯ. กายายตนํ กายวิญฺญาณสฺส ปุเรชาต\-ปจฺจเยน ปจฺจโย. วตฺถุ นรูปาวจรานํ ขนฺธานํ ปุเรชาต\-ปจฺจเยน ปจฺจโย.  นรูปาวจโร ธมฺโม รูปาวจรสฺส ธมฺมสฺส ปุเรชาต\-ปจฺจเยน ปจฺจโย. อารมฺมณ\-ปุเรชาตํ, วตฺถุ\-ปุเรชาตํ.  อารมฺมณ\-ปุเรชาตํ\csromandash{} ทิพฺเพน จกฺขุนา รูปํ ปสฺสติ. ทิพฺพาย โสตธาตุยา สทฺทํ สุณาติ. วตฺถุ\-ปุเรชาตํ\csromandash{}วตฺถุ รูปาวจรานํ ขนฺธานํ ปุเรชาต\-ปจฺจเยน ปจฺจโย. (2)}

\prose{ปจฺฉาชาต\-ปจฺจเยน ปจฺจโย. เทฺว. อาเสวน\-ปจฺจเยน ปจฺจโย. ตีณิ.}

\csromanheader{2}{\textbf{กมฺมาทิ}}
\proseitem{255}{รูปาวจโร ธมฺโม รูปาวจรสฺส ธมฺมสฺส กมฺมปจฺจเยน ปจฺจโย. สหชาตา, นานากฺขณิกา.  สหชาตา รูปาวจรา เจตนา สมฺปยุตฺตกานํ ขนฺธานํ กมฺมปจฺจเยน ปจฺจโย. ปฏิสนฺธิ\-กฺขเณ ฯเปฯ. นานากฺขณิกา รูปาวจรา เจตนา วิปากานํ รูปาวจรานํ ขนฺธานํ กมฺมปจฺจเยน ปจฺจโย.  รูปาวจโร ธมฺโม นรูปาวจรสฺส ธมฺมสฺส กมฺมปจฺจเยน ปจฺจโย. สหชาตา, นานากฺขณิกา.  สหชาตา รูปาวจรา เจตนา จิตฺต\-สมุฏฺฐานานํ รูปานํ กมฺมปจฺจเยน ปจฺจโย. ปฏิสนฺธิ\-กฺขเณ ฯเปฯ. นานากฺขณิกา รูปาวจรา เจตนา กฏตฺตารูปานํ กมฺมปจฺจเยน ปจฺจโย.  รูปาวจโร ธมฺโม รูปาวจรสฺส จ นรูปาวจรสฺส จ ธมฺมสฺส กมฺมปจฺจเยน ปจฺจโย. สหชาตา, นานากฺขณิกา.  สหชาตา รูปาวจรา เจตนา สมฺปยุตฺตกานํ ขนฺธานํ จิตฺตสมุฏฺฐานานญฺจ รูปานํ กมฺมปจฺจเยน ปจฺจโย. ปฏิสนฺธิ\-กฺขเณ ฯเปฯ. นานากฺขณิกา รูปาวจรา เจตนา วิปากานํ รูปาวจรานํ ขนฺธานํ กฏตฺตา จ รูปานํ กมฺมปจฺจเยน ปจฺจโย. (3)}

\prose{นรูปาวจโร ธมฺโม นรูปาวจรสฺส ธมฺมสฺส กมฺมปจฺจเยน ปจฺจโย. สหชาตา, นานากฺขณิกา.  สหชาตา นรูปาวจรา เจตนา สมฺปยุตฺตกานํ ขนฺธานํ จิตฺตสมุฏฺฐานานญฺจ รูปานํ กมฺมปจฺจเยน ปจฺจโย. ปฏิสนฺธิ\-กฺขเณ ฯเปฯ. นานากฺขณิกา นรูปาวจรา เจตนา วิปากานํ นรูปาวจรานํ ขนฺธานํ กฏตฺตา จ รูปานํ กมฺมปจฺจเยน ปจฺจโย. (1)}

\prose{\csromanfolio{154}วิปาก\-ปจฺจเยน ปจฺจโย. จตฺตาริ. อาหารปจฺจเยน ปจฺจโย. จตฺตาริ. อินฺทฺริย\-ปจฺจเยน ปจฺจโย. จตฺตาริ. ฌานปจฺจเยน ปจฺจโย. จตฺตาริ. มคฺคปจฺจเยน ปจฺจโย. จตฺตาริ. สมฺปยุตฺต\-ปจฺจเยน ปจฺจโย. เทฺว.}

\csromanheader{2}{\textbf{วิปฺปยุตฺต}}
\proseitem{256}{รูปาวจโร ธมฺโม นรูปาวจรสฺส ธมฺมสฺส วิปฺปยุตฺต\-ปจฺจเยน ปจฺจโย. สหชาตํ, ปจฺฉาชาตํ. (สํขิตฺตํ.) (1)}

\prose{นรูปาวจโร ธมฺโม นรูปาวจรสฺส ธมฺมสฺส วิปฺปยุตฺต\-ปจฺจเยน ปจฺจโย. สหชาตํ, ปุเรชาตํ, ปจฺฉาชาตํ. (สํขิตฺตํ.). นรูปาวจโร ธมฺโม รูปาวจรสฺส ธมฺมสฺส วิปฺปยุตฺต\-ปจฺจเยน ปจฺจโย. สหชาตํ, ปุเรชาตํ.  สหชาตํ ปฏิสนฺธิ\-กฺขเณ วตฺถุ รูปาวจรานํ ขนฺธานํ วิปฺปยุตฺต\-ปจฺจเยน ปจฺจโย. ปุเรชาตํ วตฺถุ รูปาวจรานํ ขนฺธานํ วิปฺปยุตฺต\-ปจฺจเยน ปจฺจโย. (2)}

\prose{อตฺถิอาทิ}

\proseitem{257}{รูปาวจโร ธมฺโม รูปาวจรสฺส ธมฺมสฺส อตฺถิปจฺจเยน ปจฺจโย. (ปฏิจฺจสทิสํ.). รูปาวจโร ธมฺโม นรูปาวจรสฺส ธมฺมสฺส อตฺถิปจฺจเยน ปจฺจโย. สหชาตํ, ปุเรชาตํ, ปจฺฉาชาตํ.  รูปาวจโร ธมฺโม รูปาวจรสฺส จ นรูปาวจรสฺส จ ธมฺมสฺส อตฺถิปจฺจเยน ปจฺจโย. (ปฏิจฺจสทิสํ.) (3)}

\prose{นรูปาวจโร ธมฺโม นรูปาวจรสฺส ธมฺมสฺส อตฺถิปจฺจเยน ปจฺจโย. สหชาตํ, ปุเรชาตํ, ปจฺฉาชาตํ, อาหารํ, อินฺทฺริยํ. (สํขิตฺตํ.). นรูปาวจโร ธมฺโม รูปาวจรสฺส ธมฺมสฺส อตฺถิปจฺจเยน ปจฺจโย. สหชาตํ, ปุเรชาตํ.  สหชาตํ ปฏิสนฺธิ\-กฺขเณ วตฺถุ รูปาวจรานํ ขนฺธานํ อตฺถิปจฺจเยน ปจฺจโย. ปุเรชาตํ ทิพฺเพน จกฺขุนา รูปํ ปสฺสติ. ทิพฺพาย โสตธาตุยา สทฺทํ สุณาติ. วตฺถุ รูปาวจรานํ ขนฺธานํ อตฺถิปจฺจเยน ปจฺจโย. (2)}

\proseitem{258}{รูปาวจโร จ นรูปาวจโร จ ธมฺมา รูปาวจรสฺส ธมฺมสฺส อตฺถิปจฺจเยน ปจฺจโย. สหชาตํ, ปุเรชาตํ.  สหชาโต รูปาวจโร เอโก ขนฺโธ จ วตฺถุ จ ติณฺณนฺนํ ขนฺธานํ อตฺถิปจฺจเยน ปจฺจโย ฯเปฯ เทฺว ขนฺธา จ ฯเปฯ. ปฏิสนฺธิ\-กฺขเณ ฯเปฯ. รูปาวจโร จ นรูปาวจโร จ ธมฺมา \csromanfolio{155}นรูปาวจรสฺส ธมฺมสฺส อตฺถิปจฺจเยน ปจฺจโย. สหชาตํ, ปจฺฉาชาตํ, อาหารํ, อินฺทฺริยํ.  สหชาตา รูปาวจรา ขนฺธา จ มหาภูตา จ จิตฺต\-สมุฏฺฐานานํ รูปานํ อตฺถิปจฺจเยน ปจฺจโย. ปฏิสนฺธิ\-กฺขเณ ฯเปฯ. ปจฺฉาชาตา รูปาวจรา ขนฺธา จ กพฬีกาโร อาหาโร จ ปุเรชาตสฺส อิมสฺส กายสฺส อตฺถิปจฺจเยน ปจฺจโย. ปจฺฉาชาตา รูปาวจรา ขนฺธา จ รูปาชีวิตินฺทฺริยญฺจ กฏตฺตารูปานํ อตฺถิปจฺจเยน ปจฺจโย. (2)}

\prose{นตฺถิ\-ปจฺจเยน ปจฺจโย. วิคต\-ปจฺจเยน ปจฺจโย. อวิคต\-ปจฺจเยน ปจฺจโย.}

\csromanmarkh{1. ปจฺจยานุโลม}\csromanmarkhii{2. สงฺขฺยาวาร}\csromanheadernomark{2}{1. \textbf{ปจฺจยานุโลม 2. สงฺขฺยาวาร}}
\proseitem{259}{เหตุยา จตฺตาริ, อารมฺมเณ จตฺตาริ, อธิปติยา จตฺตาริ, อนนฺตเร จตฺตาริ, สมนนฺตเร จตฺตาริ, สหชาเต สตฺต, อญฺญมญฺเญ ฉ, นิสฺสเย สตฺต, อุปนิสฺสเย จตฺตาริ, ปุเรชาเต เทฺว, ปจฺฉาชาเต เทฺว, อาเสวเน ตีณิ, กมฺเม จตฺตาริ, วิปาเก จตฺตาริ, อาหาเร จตฺตาริ, อินฺทฺริเย จตฺตาริ, ฌาเน จตฺตาริ, มคฺเค จตฺตาริ, สมฺปยุตฺเต เทฺว, วิปฺปยุตฺเต ตีณิ, อตฺถิยา สตฺต, นตฺถิยา จตฺตาริ, วิคเต จตฺตาริ, อวิคเต สตฺต.}

\proseitem{260}{รูปาวจโร ธมฺโม รูปาวจรสฺส ธมฺมสฺส อารมฺมณ\-ปจฺจเยน ปจฺจโย. สหชาต\-ปจฺจเยน ปจฺจโย. อุปนิสฺสย\-ปจฺจเยน ปจฺจโย. รูปาวจโร ธมฺโม นรูปาวจรสฺส ธมฺมสฺส อารมฺมณ\-ปจฺจเยน ปจฺจโย. สหชาต\-ปจฺจเยน ปจฺจโย. อุปนิสฺสย\-ปจฺจเยน ปจฺจโย. ปจฺฉาชาต\-ปจฺจเยน ปจฺจโย. กมฺมปจฺจเยน ปจฺจโย.  รูปาวจโร ธมฺโม รูปาวจรสฺส จ นรูปาวจรสฺส จ ธมฺมสฺส สหชาต\-ปจฺจเยน ปจฺจโย. กมฺมปจฺจเยน ปจฺจโย. (3)}

\proseitem{261}{นรูปาวจโร ธมฺโม นรูปาวจรสฺส ธมฺมสฺส อารมฺมณ\-ปจฺจเยน ปจฺจโย. สหชาต\-ปจฺจเยน ปจฺจโย. อุปนิสฺสย\-ปจฺจเยน ปจฺจโย. \csromanfolio{156}ปุเรชาต\-ปจฺจเยน ปจฺจโย. ปจฺฉาชาต\-ปจฺจเยน ปจฺจโย. กมฺมปจฺจเยน ปจฺจโย. อาหารปจฺจเยน ปจฺจโย. อินฺทฺริย\-ปจฺจเยน ปจฺจโย.  นรูปาวจโร ธมฺโม รูปาวจรสฺส ธมฺมสฺส อารมฺมณ\-ปจฺจเยน ปจฺจโย. สหชาต\-ปจฺจเยน ปจฺจโย. อุปนิสฺสย\-ปจฺจเยน ปจฺจโย. ปุเรชาต\-ปจฺจเยน ปจฺจโย. (2)}

\prose{รูปาวจโร จ นรูปาวจโร จ ธมฺมา รูปาวจรสฺส ธมฺมสฺส สหชาตํ, ปุเรชาตํ.  รูปาวจโร จ นรูปาวจโร จ ธมฺมา นรูปาวจรสฺส ธมฺมสฺส สหชาตํ, ปจฺฉาชาตํ, อาหารํ, อินฺทฺริยํ. (2)}

\csromanmarkh{2. ปจฺจย\-ปจฺจนีย}\csromanmarkhii{2. สงฺขฺยาวาร}\csromanheadernomark{2}{2. ปจฺจย\-ปจฺจนีย 2. สงฺขฺยาวาร}
\proseitem{262}{นเหตุยา สตฺต, นอารมฺมเณ สตฺต, นอธิปติยา สตฺต, นอนนฺตเร สตฺต, นสมนนฺตเร สตฺต, นสหชาเต ฉ, นอญฺญมญฺเญ ฉ, นนิสฺสเย ฉ, นอุปนิสฺสเย สตฺต, นปุเรชาเต สตฺต, นปจฺฉาชาเต สตฺต ฯเปฯ นสมฺปยุตฺเต ฉ, นวิปฺปยุตฺเต ปญฺจ, โนอตฺถิยา ปญฺจ, โนนตฺถิยา สตฺต, โนวิคเต สตฺต, โนอวิคเต ปญฺจ.}

\proseitem{263}{เหตุปจฺจยา นอารมฺมเณ จตฺตาริ, นอธิปติยา จตฺตาริ, นอนนฺตเร นสมนนฺตเร จตฺตาริ, นอญฺญมญฺเญ เทฺว, นอุปนิสฺสเย จตฺตาริ ฯเปฯ นสมฺปยุตฺเต เทฺว, นวิปฺปยุตฺเต เทฺว, โนนตฺถิยา จตฺตาริ, โนวิคเต จตฺตาริ.}

\csromanheader{2}{4. ปจฺจย\-ปจฺจนียา\-นุโลม}
\proseitemwithcloser{264}{นเหตุปจฺจยา อารมฺมเณ จตฺตาริ, อธิปติยา จตฺตาริ, (อนุโลมมาติกา กาตพฺพา.) ฯเปฯ อวิคเต สตฺต.}{รูปาวจรทุกํ นิฏฺฐิตํ.}
\csromanmatikaanchor{mtk.24}
\csromantocmarkref{section}{95. อรูปาวจรทุก}{mtk.24}
\bookmark[dest={mtk.24},level=1]{95. อรูปาวจรทุก}
\csromanmarkh{95. อรูปาวจรทุก}\csromanmarkhii{95. อรูปาวจรทุก 1. ปฏิจฺจวาร}\csromanheadernomark{1}{95. อรูปาวจรทุก \textbf{95. อรูปาวจรทุก 1. ปฏิจฺจวาร}}
\proseitem{265}{\csromanfolio{157}อรูปาวจรํ ธมฺมํ ปฏิจฺจ อรูปาวจโร ธมฺโม อุปฺปชฺชติ เหตุปจฺจยา. อรูปาวจรํ เอกํ ขนฺธํ ปฏิจฺจ ตโย ขนฺธา ฯเปฯ เทฺว ขนฺเธ ฯเปฯ. ปฏิสนฺธิ\-กฺขเณ ฯเปฯ. อรูปาวจรํ ธมฺมํ ปฏิจฺจ นอรูปาวจโร ธมฺโม อุปฺปชฺชติ เหตุปจฺจยา. อรูปาวจเร ขนฺเธ ปฏิจฺจ จิตฺต\-สมุฏฺฐานํ รูปํ.  อรูปาวจรํ ธมฺมํ ปฏิจฺจ อรูปาวจโร จ นอรูปาวจโร จ ธมฺมา อุปฺปชฺชนฺติ เหตุปจฺจยา. อรูปาวจรํ เอกํ ขนฺธํ ปฏิจฺจ ตโย ขนฺธา จิตฺต\-สมุฏฺฐานญฺจ รูปํ ฯเปฯ เทฺว ขนฺเธ ฯเปฯ. (3)}

\proseitem{266}{นอรูปาวจรํ ธมฺมํ ปฏิจฺจ นอรูปาวจโร ธมฺโม อุปฺปชฺชติ เหตุปจฺจยา. นอรูปาวจรํ เอกํ ขนฺธํ ปฏิจฺจ ตโย ขนฺธา จิตฺต\-สมุฏฺฐานญฺจ รูปํ ฯเปฯ เทฺว ขนฺเธ ฯเปฯ. ปฏิสนฺธิ\-กฺขเณ นอรูปาวจรํ เอกํ ขนฺธํ ปฏิจฺจ ตโย ขนฺธา กฏตฺตา จ รูปํ ฯเปฯ เทฺว ขนฺเธ ฯเปฯ. ขนฺเธ ปฏิจฺจ วตฺถุ, วตฺถุํ ปฏิจฺจ ขนฺธา. เอกํ มหาภูตํ ฯเปฯ. (1)}

\prose{อรูปาวจรญฺจ นอรูปาวจรญฺจ ธมฺมํ ปฏิจฺจ นอรูปาวจโร ธมฺโม อุปฺปชฺชติ เหตุปจฺจยา. อรูปาวจเร ขนฺเธ จ มหาภูเต จ ปฏิจฺจ จิตฺต\-สมุฏฺฐานํ รูปํ. (1) (สํขิตฺตํ.)}

\csromanmarkh{1. ปจฺจยานุโลม}\csromanmarkhii{2. สงฺขฺยาวาร}\csromanheadernomark{2}{1. \textbf{ปจฺจยานุโลม 2. สงฺขฺยาวาร}}
\proseitem{267}{\textbf{เหตุ}ยา ปญฺจ, อารมฺมเณ เทฺว, อธิปติยา ปญฺจ, อนนฺตเร เทฺว, สมนนฺตเร เทฺว, สหชาเต ปญฺจ, อญฺญมญฺเญ เทฺว, นิสฺสเย ปญฺจ, อุปนิสฺสเย เทฺว, ปุเรชาเต เทฺว, อาเสวเน เทฺว, กมฺเม ปญฺจ, วิปาเก เทฺว, อาหาเร ปญฺจ ฯเปฯ อวิคเต ปญฺจ.}

\csromanmarkh{2. ปจฺจย\-ปจฺจนีย}\csromanmarkhii{1. วิภงฺควาร}\csromanheadernomark{2}{2. ปจฺจย\-ปจฺจนีย 1. วิภงฺควาร}
\prose{นเหตุ นอารมฺมณ}

\proseitem{268}{นอรูปาวจรํ ธมฺมํ ปฏิจฺจ นอรูปาวจโร ธมฺโม อุปฺปชฺชติ นเหตุปจฺจยา. อเหตุกํ นอรูปาวจรํ เอกํ ขนฺธํ ปฏิจฺจ ตโย ขนฺธา จิตฺต\-สมุฏฺฐานญฺจ \csromanfolio{158}รูปํ ฯเปฯ เทฺว ขนฺเธ ฯเปฯ. อเหตุก\-ปฏิสนฺธิ\-กฺขเณ ฯเปฯ. (ยาว อสญฺญสตฺตา.) วิจิกิจฺฉา\-สหคเต อุทฺธจฺจ\-สหคเต ขนฺเธ ปฏิจฺจ วิจิกิจฺฉา\-สหคโต อุทฺธจฺจ\-สหคโต โมโห. (1). นอารมฺมณ\-ปจฺจยา. ตีณิ.}

\prose{นอธิปตฺยาทิ}

\proseitem{269}{อรูปาวจรํ ธมฺมํ ปฏิจฺจ อรูปาวจโร ธมฺโม อุปฺปชฺชติ นอธิปติ\-ปจฺจยา. อรูปาวจเร ขนฺเธ ปฏิจฺจ อรูปาวจรา\-ธิปติ. วิปากํ อรูปาวจรํ เอกํ ขนฺธํ ปฏิจฺจ ตโย ขนฺธา ฯเปฯ เทฺว ขนฺเธ ฯเปฯ. ปฏิสนฺธิ\-กฺขเณ ฯเปฯ. (1)}

\prose{นอรูปาวจรํ ธมฺมํ ปฏิจฺจ นอรูปาวจโร ธมฺโม อุปฺปชฺชติ นอธิปติ\-ปจฺจยา. นอรูปาวจรํ เอกํ ขนฺธํ ปฏิจฺจ ตโย ขนฺธา จิตฺต\-สมุฏฺฐานญฺจ รูปํ ฯเปฯ เทฺว ขนฺเธ ฯเปฯ. ปฏิสนฺธิ\-กฺขเณ ฯเปฯ. (ยาว อสญฺญสตฺตา.) (1). นอนนฺตร\-ปจฺจยา ฯเปฯ. นอุปนิสฺสยปจฺจยา.}

\csromanheader{2}{\textbf{นปุเรชาตาทิ}}
\proseitem{270}{อรูปาวจรํ ธมฺมํ ปฏิจฺจ อรูปาวจโร ธมฺโม อุปฺปชฺชติ นปุเรชาต\-ปจฺจยา. อรูเป อรูปาวจรํ เอกํ ขนฺธํ ปฏิจฺจ ตโย ขนฺธา ฯเปฯ เทฺว ขนฺเธ ฯเปฯ. (มูลํ กาตพฺพํ.) อรูปาวจเร ขนฺเธ ปฏิจฺจ จิตฺต\-สมุฏฺฐานํ รูปํ. (2)}

\prose{นอรูปาวจรํ ธมฺมํ ปฏิจฺจ นอรูปาวจโร ธมฺโม อุปฺปชฺชติ นปุเรชาต\-ปจฺจยา. อรูเป นอรูปาวจรํ เอกํ ขนฺธํ ปฏิจฺจ ตโย ขนฺธา ฯเปฯ เทฺว ขนฺเธ ฯเปฯ. ปฏิสนฺธิ\-กฺขเณ ฯเปฯ. (ยาว อสญฺญสตฺตา.) (1)}

\prose{อรูปาวจรญฺจ นอรูปาวจรญฺจ ธมฺมํ ปฏิจฺจ นอรูปาวจโร ธมฺโม อุปฺปชฺชติ นปุเรชาต\-ปจฺจยา. อรูปาวจเร ขนฺเธ จ มหาภูเต จ ปฏิจฺจ จิตฺต\-สมุฏฺฐานํ รูปํ. (1). นปจฺฉาชาต\-ปจฺจยา.}

\prose{นอาเสวน}

\proseitem{271}{อรูปาวจรํ ธมฺมํ ปฏิจฺจ อรูปาวจโร ธมฺโม อุปฺปชฺชติ นอาเสวน\-ปจฺจยา. วิปากํ อรูปาวจรํ เอกํ ขนฺธํ ปฏิจฺจ ตโย ขนฺธา ฯเปฯ เทฺว ขนฺเธ ฯเปฯ. ปฏิสนฺธิ\-กฺขเณ ฯเปฯ. อรูปาวจรํ ธมฺมํ ปฏิจฺจ นอรูปาวจโร ธมฺโม อุปฺปชฺชติ นอาเสวน\-ปจฺจยา. อรูปาวจเร ขนฺเธ ปฏิจฺจ จิตฺต\-สมุฏฺฐานํ รูปํ. (2)}

\prose{\csromanfolio{159}นอรูปาวจรํ ธมฺมํ ปฏิจฺจ นอรูปาวจโร ธมฺโม อุปฺปชฺชติ นอาเสวน\-ปจฺจยา. นอรูปาวจรํ เอกํ ขนฺธํ ปฏิจฺจ ตโย ขนฺธา จิตฺต\-สมุฏฺฐานญฺจ รูปํ ฯเปฯ เทฺว ขนฺเธ ฯเปฯ. (ยาว อสญฺญสตฺตา.) (1)}

\prose{อรูปาวจรญฺจ นอรูปาวจรญฺจ ธมฺมํ ปฏิจฺจ นอรูปาวจโร ธมฺโม อุปฺปชฺชติ นอาเสวน\-ปจฺจยา. อรูปาวจเร ขนฺเธ จ มหาภูเต จ ปฏิจฺจ จิตฺต\-สมุฏฺฐานํ รูปํ. (1) (สํขิตฺตํ.)}

\csromanmarkh{2. ปจฺจย\-ปจฺจนีย}\csromanmarkhii{2. สงฺขฺยาวาร}\csromanheadernomark{2}{2. ปจฺจย\-ปจฺจนีย 2. สงฺขฺยาวาร}
\proseitem{272}{นเหตุยา เอกํ, นอารมฺมเณ ตีณิ, นอธิปติยา เทฺว, นอนนฺตเร ตีณิ, นสมนนฺตเร ตีณิ, นอญฺญมญฺเญ ตีณิ, นอุปนิสฺสเย ตีณิ, นปุเรชาเต จตฺตาริ, นปจฺฉาชาเต ปญฺจ, นอาเสวเน จตฺตาริ, นกมฺเม เทฺว, นวิปาเก ปญฺจ, นอาหาเร เอกํ, นอินฺทฺริเย เอกํ, นฌาเน เอกํ, นมคฺเค เอกํ, นสมฺปยุตฺเต ตีณิ, นวิปฺปยุตฺเต เทฺว, โนนตฺถิยา ตีณิ, โนวิคเต ตีณิ.}

\vspace{\tipitakaparskip}
\csromancenter{(เอวํ อิตเร เทฺว คณนาปิ สหชาตวาโรปิ กาตพฺโพ.)}
\csromanmarkh{95. อรูปาวจรทุก}\csromanmarkhii{3. ปจฺจยวาร}\csromanheadernomark{2}{95. อรูปาวจรทุก 3. ปจฺจยวาร}
\proseitem{273}{อรูปาวจรํ ธมฺมํ ปจฺจยา อรูปาวจโร ธมฺโม อุปฺปชฺชติ เหตุปจฺจยา. ตีณิ. (ปฏิจฺจสทิสา.)}

\prose{นอรูปาวจรํ ธมฺมํ ปจฺจยา นอรูปาวจโร ธมฺโม อุปฺปชฺชติ เหตุปจฺจยา. นอรูปาวจรํ เอกํ ขนฺธํ ฯเปฯ. (ยาว อชฺฌตฺติกา มหาภูตา.) วตฺถุํ ปจฺจยา นอรูปาวจรา ขนฺธา.  นอรูปาวจรํ ธมฺมํ ปจฺจยา อรูปาวจโร ธมฺโม อุปฺปชฺชติ เหตุปจฺจยา. วตฺถุํ ปจฺจยา อรูปาวจรา ขนฺธา.  นอรูปาวจรํ ธมฺมํ ปจฺจยา อรูปาวจโร จ นอรูปาวจโร จ ธมฺมา อุปฺปชฺชนฺติ เหตุปจฺจยา. วตฺถุํ ปจฺจยา อรูปาวจรา ขนฺธา. มหาภูเต ปจฺจยา จิตฺต\-สมุฏฺฐานํ รูปํ. (3)}

\proseitem{274}{\csromanfolio{160}อรูปาวจรญฺจ นอรูปาวจรญฺจ ธมฺมํ ปจฺจยา อรูปาวจโร ธมฺโม อุปฺปชฺชติ เหตุปจฺจยา. อรูปาวจรํ เอกํ ขนฺธญฺจ วตฺถุญฺจ ปจฺจยา ตโย ขนฺธา ฯเปฯ เทฺว ขนฺเธ จ ฯเปฯ. อรูปาวจรญฺจ นอรูปาวจรญฺจ ธมฺมํ ปจฺจยา นอรูปาวจโร ธมฺโม อุปฺปชฺชติ เหตุปจฺจยา. อรูปาวจเร ขนฺเธ จ มหาภูเต จ ปจฺจยา จิตฺต\-สมุฏฺฐานํ รูปํ.  อรูปาวจรญฺจ นอรูปาวจรญฺจ ธมฺมํ ปจฺจยา อรูปาวจโร จ นอรูปาวจโร จ ธมฺมา อุปฺปชฺชนฺติ เหตุปจฺจยา. อรูปาวจรํ เอกํ ขนฺธญฺจ วตฺถุญฺจ ปจฺจยา ตโย ขนฺธา ฯเปฯ เทฺว ขนฺเธ จ ฯเปฯ. อรูปาวจเร ขนฺเธ จ มหาภูเต จ ปจฺจยา จิตฺต\-สมุฏฺฐานํ รูปํ. (3)}

\csromanmarkh{1. ปจฺจยานุโลม}\csromanmarkhii{2. สงฺขฺยาวาร}\csromanheadernomark{2}{1. \textbf{ปจฺจยานุโลม 2. สงฺขฺยาวาร}}
\proseitem{275}{เหตุยา นว, อารมฺมเณ จตฺตาริ, อธิปติยา นว, อนนฺตเร จตฺตาริ, สมนนฺตเร จตฺตาริ, สหชาเต นว, อญฺญมญฺเญ จตฺตาริ, นิสฺสเย นว, อุปนิสฺสเย จตฺตาริ, ปุเรชาเต จตฺตาริ, อาเสวเน จตฺตาริ, กมฺเม นว, วิปาเก เทฺว ฯเปฯ อวิคเต นว.}

\csromanmarkh{2. ปจฺจย\-ปจฺจนีย}\csromanmarkhii{1. วิภงฺควาร}\csromanheadernomark{2}{2. ปจฺจย\-ปจฺจนีย 1. วิภงฺควาร}
\csromanheader{2}{\textbf{นเหตฺวาทิ}}
\proseitem{276}{นอรูปาวจรํ ธมฺมํ ปจฺจยา นอรูปาวจโร ธมฺโม อุปฺปชฺชติ นเหตุปจฺจยา. อเหตุกํ นอรูปาวจรํ เอกํ ขนฺธํ ฯเปฯ. (ยาว อสญฺญสตฺตา.) จกฺขายตนํ ปจฺจยา จกฺขุวิญฺญาณํ ฯเปฯ. กายายตนํ ปจฺจยา กายวิญฺญาณํ. วตฺถุํ ปจฺจยา อเหตุกา นอรูปาวจรา ขนฺธา. วิจิกิจฺฉา\-สหคเต อุทฺธจฺจ\-สหคเต ขนฺเธ จ วตฺถุญฺจ ปจฺจยา วิจิกิจฺฉา\-สหคโต อุทฺธจฺจ\-สหคโต โมโห. (1). นอารมฺมณ\-ปจฺจยา. ตีณิ.}

\prose{อรูปาวจรํ ธมฺมํ ปจฺจยา อรูปาวจโร ธมฺโม อุปฺปชฺชติ นอธิปติ\-ปจฺจยา. อรูปาวจเร ขนฺเธ ปจฺจยา อรูปาวจรา\-ธิปติ. วิปากํ อรูปาวจรํ เอกํ ขนฺธํ ฯเปฯ. ปฏิสนฺธิ\-กฺขเณ ฯเปฯ. (1)}

\prose{นอรูปาวจรํ ธมฺมํ ปจฺจยา นอรูปาวจโร ธมฺโม อุปฺปชฺชติ นอธิปติ\-ปจฺจยา. นอรูปาวจรํ เอกํ ขนฺธํ ฯเปฯ. ปฏิสนฺธิ\-กฺขเณ ฯเปฯ. (ยาว อสญฺญสตฺตา.) จกฺขายตนํ ปจฺจยา จกฺขุวิญฺญาณํ ฯเปฯ. กายายตนํ ปจฺจยา กายวิญฺญาณํ. วตฺถุํ ปจฺจยา นอรูปาวจรา ขนฺธา. (1)}

\csromanmarkh{95. อรูปาวจรทุก}\csromanmarkhii{2. ปจฺจย\-ปจฺจนีย 2. สงฺขฺยาวาร}\csromanheadernomark{2}{95. อรูปาวจรทุก 2. ปจฺจย\-ปจฺจนีย 2. สงฺขฺยาวาร}
\proseitem{277}{\csromanfolio{161}นเหตุยา เอกํ, นอารมฺมเณ ตีณิ, นอธิปติยา จตฺตาริ, นอนนฺตเร ตีณิ, นสมนนฺตเร นอญฺญมญฺเญ นอุปนิสฺสเย ตีณิ, นปุเรชาเต จตฺตาริ, นปจฺฉาชาเต นว, นอาเสวเน จตฺตาริ, นกมฺเม จตฺตาริ, นวิปาเก นว, นอาหาเร เอกํ, นอินฺทฺริเย เอกํ, นฌาเน เอกํ, นมคฺเค เอกํ, นสมฺปยุตฺเต ตีณิ, นวิปฺปยุตฺเต เทฺว, โนนตฺถิยา ตีณิ, โนวิคเต ตีณิ.}

\vspace{\tipitakaparskip}
\csromancenter{(เอวํ อิตเร เทฺว คณนาปิ นิสฺสยวาโรปิ กาตพฺโพ.)}
\csromanmarkh{95. อรูปาวจรทุก}\csromanmarkhii{5. สํสฏฺฐวาร}\csromanheadernomark{2}{95. อรูปาวจรทุก 5. สํสฏฺฐวาร}
\proseitem{278}{อรูปาวจรํ ธมฺมํ สํสฏฺโฐ อรูปาวจโร ธมฺโม อุปฺปชฺชติ เหตุปจฺจยา. อรูปาวจรํ เอกํ ขนฺธํ สํสฏฺฐา ตโย ขนฺธา ฯเปฯ เทฺว ขนฺเธ ฯเปฯ. ปฏิสนฺธิ\-กฺขเณ ฯเปฯ. (1)}

\prose{นอรูปาวจรํ ธมฺมํ สํสฏฺโฐ นอรูปาวจโร ธมฺโม อุปฺปชฺชติ เหตุปจฺจยา. นอรูปาวจรํ เอกํ ขนฺธํ สํสฏฺฐา ตโย ขนฺธา ฯเปฯ เทฺว ขนฺเธ ฯเปฯ. ปฏิสนฺธิ\-กฺขเณ ฯเปฯ. (1)}

\prose{\textbf{เหตุ}ยา เทฺว ฯเปฯ อวิคเต เทฺว. (อนุโลมํ.)}

\prose{นเหตุยา เอกํ, นอธิปติยา เทฺว, นปุเรชาเต เทฺว, นปจฺฉาชาเต เทฺว, นอาเสวเน เทฺว, นกมฺเม เทฺว, นวิปาเก เทฺว, นฌาเน เอกํ, นมคฺเค เอกํ, นวิปฺปยุตฺเต เทฺว. (ปจฺจนียํ.)}

\vspace{\tipitakaparskip}
\csromancenter{(เอวํ อิตเร เทฺว คณนาปิ สมฺปยุตฺตวาโรปิ กาตพฺโพ.)}
\csromanmarkh{95. อรูปาวจรทุก}\csromanmarkhii{7. ปญฺหาวาร}\csromanheadernomark{2}{95. \textbf{อรูปาวจรทุก 7. ปญฺหาวาร}}
\proseitem{279}{\csromanfolio{162}อรูปาวจโร ธมฺโม อรูปาวจรสฺส ธมฺมสฺส เหตุปจฺจเยน ปจฺจโย. อรูปาวจรา เหตู สมฺปยุตฺตกานํ ขนฺธานํ เหตุปจฺจเยน ปจฺจโย. ปฏิสนฺธิ\-กฺขเณ ฯเปฯ. (มูลํ กาตพฺพํ.) อรูปาวจรา เหตู จิตฺต\-สมุฏฺฐานานํ รูปานํ เหตุปจฺจเยน ปจฺจโย. (มูลํ กาตพฺพํ.) อรูปาวจรา เหตู สมฺปยุตฺตกานํ ขนฺธานํ จิตฺตสมุฏฺฐานานญฺจ รูปานํ เหตุปจฺจเยน ปจฺจโย. (3)}

\prose{นอรูปาวจโร ธมฺโม นอรูปาวจรสฺส ธมฺมสฺส เหตุปจฺจเยน ปจฺจโย. นอรูปาวจรา เหตู สมฺปยุตฺตกานํ ขนฺธานํ จิตฺตสมุฏฺฐานานญฺจ รูปานํ เหตุปจฺจเยน ปจฺจโย. ปฏิสนฺธิ\-กฺขเณ ฯเปฯ. (1)}

\proseitem{280}{อรูปาวจโร ธมฺโม อรูปาวจรสฺส ธมฺมสฺส อารมฺมณ\-ปจฺจเยน ปจฺจโย. อากาสา\-นญฺจา\-ยตนํ วิญฺญาณญฺจา\-ยตนสฺส อารมฺมณ\-ปจฺจเยน ปจฺจโย. อากิญฺจญฺญา\-ยตนํ เนว\-สญฺญา\-นา\-สญฺญา\-ยตนสฺส อารมฺมณ\-ปจฺจเยน ปจฺจโย.  อรูปาวจโร ธมฺโม นอรูปาวจรสฺส ธมฺมสฺส อารมฺมณ\-ปจฺจเยน ปจฺจโย. อากาสา\-นญฺจา\-ยตนํ ปจฺจเวกฺขติ. วิญฺญาณญฺจา\-ยตนํ ฯเปฯ. อากิญฺจญฺญา\-ยตนํ ฯเปฯ. เนว\-สญฺญา\-นา\-สญฺญา\-ยตนํ ปจฺจเวกฺขติ. อรูปาวจเร ขนฺเธ อนิจฺจโต ฯเปฯ โทมนสฺสํ อุปฺปชฺชติ. เจโตปริย\-ญาเณน อรูปาวจร\-จิตฺต\-สมงฺคิสฺส จิตฺตํ ชานาติ. อรูปาวจรา ขนฺธา เจโตปริย\-ญาณสฺส, ปุพฺเพ\-นิวาสา\-นุสฺสติ\-ญาณสฺส, ยถากมฺมูปคญาณสฺส, อนาคตํสญาณสฺส, อาวชฺชนาย อารมฺมณ\-ปจฺจเยน ปจฺจโย. (2)}

\prose{นอรูปาวจโร ธมฺโม นอรูปาวจรสฺส ธมฺมสฺส อารมฺมณ\-ปจฺจเยน ปจฺจโย. ทานํ ฯเปฯ สีลํ ฯเปฯ อุโปสถกมฺมํ กตฺวา ตํ ปจฺจเวกฺขติ, อสฺสาเทติ อภินนฺทติ, ตํ อารพฺภ ราโค ฯเปฯ โทมนสฺสํ อุปฺปชฺชติ. ปุพฺเพ สุจิณฺณานิ ฯเปฯ. ฌานา ฯเปฯ. อริยา มคฺคา วุฏฺฐหิตฺวา มคฺคํ ปจฺจเวกฺขนฺติ, ผลํ \csromanfolio{163}ปจฺจเวกฺขนฺติ, นิพฺพานํ ปจฺจเวกฺขนฺติ. นิพฺพานํ โคตฺรภุสฺส, โวทานสฺส, มคฺคสฺส, ผลสฺส, อาวชฺชนาย อารมฺมณ\-ปจฺจเยน ปจฺจโย. อริยา ปหีเน กิเลเส ปจฺจเวกฺขนฺติ, วิกฺขมฺภิเต กิเลเส ฯเปฯ. ปุพฺเพ ฯเปฯ. จกฺขุํ ฯเปฯ. วตฺถุํ นอรูปาวจเร ขนฺเธ อนิจฺจโต ฯเปฯ โทมนสฺสํ อุปฺปชฺชติ. ทิพฺเพน จกฺขุนา รูปํ ปสฺสติ. ทิพฺพาย โสตธาตุยา สทฺทํ สุณาติ. เจโตปริย\-ญาเณน นอรูปาวจร\-จิตฺต\-สมงฺคิสฺส จิตฺตํ ชานาติ. รูปายตนํ จกฺขุ\-วิญฺญาณสฺส ฯเปฯ. โผฏฺฐพฺพา\-ยตนํ กายวิญฺญาณสฺส ฯเปฯ. นอรูปาวจรา ขนฺธา อิทฺธิวิธ\-ญาณสฺส, เจโตปริย\-ญาณสฺส, ปุพฺเพ\-นิวาสา\-นุสฺสติ\-ญาณสฺส, ยถากมฺมูปคญาณสฺส, อนาคตํสญาณสฺส, อาวชฺชนาย อารมฺมณ\-ปจฺจเยน ปจฺจโย. (1)}

\csromanheader{2}{\textbf{อธิปติ}}
\proseitem{281}{อรูปาวจโร ธมฺโม อรูปาวจรสฺส ธมฺมสฺส อธิปติ\-ปจฺจเยน ปจฺจโย. สหชาตา\-ธิปติ\csromandash{}อรูปาวจรา\-ธิปติ สมฺปยุตฺตกานํ ขนฺธานํ อธิปติ\-ปจฺจเยน ปจฺจโย.  อรูปาวจโร ธมฺโม นอรูปาวจรสฺส ธมฺมสฺส อธิปติ\-ปจฺจเยน ปจฺจโย. อารมฺมณา\-ธิปติ, สหชาตา\-ธิปติ.  อารมฺมณา\-ธิปติ\csromandash{} อากาสา\-นญฺจา\-ยตนํ ครุํ กตฺวา ปจฺจเวกฺขติ ฯเปฯ. เนว\-สญฺญา\-นา\-สญฺญา\-ยตนํ ครุํ กตฺวา ปจฺจเวกฺขติ. อรูปาวจเร ขนฺเธ ครุํ กตฺวา อสฺสาเทติ อภินนฺทติ, ตํ ครุํ กตฺวา ราโค อุปฺปชฺชติ. ทิฏฺฐิ อุปฺปชฺชติ. สหชาตา\-ธิปติ\csromandash{}อรูปาวจรา\-ธิปติ จิตฺต\-สมุฏฺฐานานํ รูปานํ อธิปติ\-ปจฺจเยน ปจฺจโย. (มูลํ กาตพฺพํ.) อรูปาวจรา\-ธิปติ สมฺปยุตฺตกานํ ขนฺธานํ จิตฺตสมุฏฺฐานานญฺจ รูปานํ อธิปติ\-ปจฺจเยน ปจฺจโย. (3)}

\proseitem{282}{นอรูปาวจโร ธมฺโม นอรูปาวจรสฺส ธมฺมสฺส อธิปติ\-ปจฺจเยน ปจฺจโย. อารมฺมณา\-ธิปติ, สหชาตา\-ธิปติ.  อารมฺมณา\-ธิปติ\csromandash{}ทานํ ฯเปฯ สีลํ ฯเปฯ อุโปสถกมฺมํ กตฺวา ตํ ครุํ กตฺวา ปจฺจเวกฺขติ, อสฺสาเทติ อภินนฺทติ, ตํ ครุํ กตฺวา ราโค อุปฺปชฺชติ. ทิฏฺฐิ อุปฺปชฺชติ. ปุพฺเพ สุจิณฺณานิ ฯเปฯ. ฌานา ฯเปฯ. อริยา มคฺคา วุฏฺฐหิตฺวา มคฺคํ ครุํ กตฺวา ปจฺจเวกฺขนฺติ, ผลํ ฯเปฯ. นิพฺพานํ ครุํ กตฺวา ปจฺจเวกฺขนฺติ. นิพฺพานํ โคตฺรภุสฺส, โวทานสฺส, มคฺคสฺส, ผลสฺส อธิปติ\-ปจฺจเยน ปจฺจโย. จกฺขุํ ฯเปฯ. วตฺถุํ นอรูปาวจเร ขนฺเธ ครุํ กตฺวา อสฺสาเทติ อภินนฺทติ, ตํ ครุํ กตฺวา ราโค อุปฺปชฺชติ. ทิฏฺฐิ อุปฺปชฺชติ. สหชาตา\-ธิปติ\csromandash{}นอรูปาวจรา\-ธิปติ สมฺปยุตฺตกานํ ขนฺธานํ จิตฺตสมุฏฺฐานานญฺจ รูปานํ อธิปติ\-ปจฺจเยน ปจฺจโย. (1)}

\proseitem{283}{\csromanfolio{164}อรูปาวจโร ธมฺโม อรูปาวจรสฺส ธมฺมสฺส อนนฺตร\-ปจฺจเยน ปจฺจโย. ปุริมา ปุริมา อรูปาวจรา ขนฺธา ปจฺฉิมานํ ปจฺฉิมานํ อรูปาวจรานํ ขนฺธานํ อนนฺตร\-ปจฺจเยน ปจฺจโย.  อรูปาวจโร ธมฺโม นอรูปาวจรสฺส ธมฺมสฺส อนนฺตร\-ปจฺจเยน ปจฺจโย. อรูปาวจรํ จุติจิตฺตํ นอรูปาวจรสฺส อุปปตฺติ\-จิตฺตสฺส. อรูปาวจรํ ภวงฺคํ อาวชฺชนาย. อรูปาวจรา ขนฺธา นอรูปาวจรสฺส วุฏฺฐานสฺส. นิโรธา วุฏฺฐหนฺตสฺส เนว\-สญฺญา\-นา\-สญฺญา\-ยตนํ ผลสมาปตฺติยา อนนฺตร\-ปจฺจเยน ปจฺจโย. (2)}

\proseitem{284}{นอรูปาวจโร ธมฺโม นอรูปาวจรสฺส ธมฺมสฺส อนนฺตร\-ปจฺจเยน ปจฺจโย. ปุริมา ปุริมา นอรูปาวจรา ขนฺธา ปจฺฉิมานํ ปจฺฉิมานํ นอรูปาวจรานํ ขนฺธานํ อนนฺตร\-ปจฺจเยน ปจฺจโย. อนุโลมํ โคตฺรภุสฺส ฯเปฯ. อนุโลมํ ผลสมาปตฺติยา อนนฺตร\-ปจฺจเยน ปจฺจโย.  นอรูปาวจโร ธมฺโม อรูปาวจรสฺส ธมฺมสฺส อนนฺตร\-ปจฺจเยน ปจฺจโย. นอรูปาวจรํ จุติจิตฺตํ อรูปาวจรสฺส อุปปตฺติ\-จิตฺตสฺส อนนฺตร\-ปจฺจเยน ปจฺจโย. นอรูปาวจรา ขนฺธา อรูปาวจรสฺส วุฏฺฐานสฺส อนนฺตร\-ปจฺจเยน ปจฺจโย. อากาสา\-นญฺจา\-ยตนสฺส ปริกมฺมํ อากาสา\-นญฺจา\-ยตนสฺส อนนฺตร\-ปจฺจเยน ปจฺจโย. วิญฺญาณญฺจา\-ยตนสฺส ฯเปฯ. อากิญฺจญฺญา\-ยตนสฺส ฯเปฯ. เนว\-สญฺญา\-นา\-สญฺญา\-ยตนสฺส ปริกมฺมํ เนว\-สญฺญา\-นา\-สญฺญา\-ยตนสฺส อนนฺตร\-ปจฺจเยน ปจฺจโย. (2)}

\prose{สมนนฺตร\-ปจฺจเยน ปจฺจโย. สหชาต\-ปจฺจเยน ปจฺจโย. ปญฺจ. อญฺญมญฺญ\-ปจฺจเยน ปจฺจโย. เทฺว. นิสฺสย\-ปจฺจเยน ปจฺจโย. สตฺต.}

\csromanheader{2}{\textbf{อุปนิสฺสย}}
\proseitem{285}{อรูปาวจโร ธมฺโม อรูปาวจรสฺส ธมฺมสฺส อุปนิสฺสย\-ปจฺจเยน ปจฺจโย. อนนฺตรู\-ปนิสฺสโย, ปกตู\-ปนิสฺสโย ฯเปฯ. ปกตู\-ปนิสฺสโย\csromandash{}อากาสา\-นญฺจา\-ยตนํ วิญฺญาณญฺจา\-ยตนสฺส อุปนิสฺสย\-ปจฺจเยน ปจฺจโย. วิญฺญาณญฺจา\-ยตนํ อากิญฺจญฺญา\-ยตนสฺส ฯเปฯ. อากิญฺจญฺญา\-ยตนํ เนว\-สญฺญา\-นา\-สญฺญา\-ยตนสฺส อุปนิสฺสย\-ปจฺจเยน ปจฺจโย.  อรูปาวจโร ธมฺโม นอรูปาวจรสฺส ธมฺมสฺส อุปนิสฺสย\-ปจฺจเยน ปจฺจโย. อารมฺมณู\-ปนิสฺสโย, อนนฺตรู\-ปนิสฺสโย, ปกตู\-ปนิสฺสโย ฯเปฯ. \csromanfolio{165}ปกตู\-ปนิสฺสโย\csromandash{}อรูปาวจรํ สทฺธํ อุปนิสฺสาย ทานํ เทติ. สีลํ ฯเปฯ อุโปสถกมฺมํ กโรติ. นอรูปาวจรํ ฌานํ อุปฺปาเทติ. วิปสฺสนํ ฯเปฯ มคฺคํ ฯเปฯ อภิญฺญํ ฯเปฯ สมาปตฺติํ อุปฺปาเทติ. มานํ ชปฺเปติ. ทิฏฺฐิํ คณฺหาติ. อรูปาวจรํ สีลํ ฯเปฯ. ปญฺญํ อุปนิสฺสาย ทานํ เทติ ฯเปฯ. สมาปตฺติํ อุปฺปาเทติ. มานํ ชปฺเปติ. ทิฏฺฐิํ คณฺหาติ. อรูปาวจรา สทฺธา ฯเปฯ ปญฺญา นอรูปาวจราย สทฺธาย ฯเปฯ ปญฺญาย. ราคสฺส ฯเปฯ ปตฺถนาย. กายิกสฺส สุขสฺส. กายิกสฺส ทุกฺขสฺส, มคฺคสฺส, ผลสมาปตฺติยา อุปนิสฺสย\-ปจฺจเยน ปจฺจโย. (2)}

\proseitem{286}{นอรูปาวจโร ธมฺโม นอรูปาวจรสฺส ธมฺมสฺส อุปนิสฺสย\-ปจฺจเยน ปจฺจโย. อารมฺมณู\-ปนิสฺสโย, อนนฺตรู\-ปนิสฺสโย, ปกตู\-ปนิสฺสโย ฯเปฯ. ปกตู\-ปนิสฺสโย\csromandash{}นอรูปาวจรํ สทฺธํ อุปนิสฺสาย ทานํ เทติ. สีลํ ฯเปฯ. อุโปสถกมฺมํ กโรติ. นอรูปาวจรํ ฌานํ อุปฺปาเทติ. วิปสฺสนํ ฯเปฯ. มคฺคํ ฯเปฯ. อภิญฺญํ ฯเปฯ. สมาปตฺติํ อุปฺปาเทติ. มานํ ชปฺเปติ. ทิฏฺฐิํ คณฺหาติ. นอรูปาวจรํ สีลํ ฯเปฯ ปญฺญํ ฯเปฯ ราคํ ฯเปฯ ปตฺถนํ ฯเปฯ กายิกํ สุขํ. กายิกํ ทุกฺขํ. อุตุํ. โภชนํ. เสนาสนํ อุปนิสฺสาย ทานํ เทติ ฯเปฯ. สมาปตฺติํ อุปฺปาเทติ. ปาณํ หนติ ฯเปฯ สํฆํ ภินฺทติ. นอรูปาวจรา สทฺธา ฯเปฯ. เสนาสนํ นอรูปาวจราย สทฺธาย ฯเปฯ ปตฺถนาย. กายิกสฺส สุขสฺส. กายิกสฺส ทุกฺขสฺส, มคฺคสฺส, ผลสมาปตฺติยา อุปนิสฺสย\-ปจฺจเยน ปจฺจโย. (1)}

\prose{นอรูปาวจโร ธมฺโม อรูปาวจรสฺส ธมฺมสฺส อุปนิสฺสย\-ปจฺจเยน ปจฺจโย. อนนฺตรู\-ปนิสฺสโย, ปกตู\-ปนิสฺสโย ฯเปฯ. ปกตู\-ปนิสฺสโย\csromandash{}อากาสา\-นญฺจา\-ยตนสฺส ปริกมฺมํ อากาสา\-นญฺจา\-ยตนสฺส อุปนิสฺสย\-ปจฺจเยน ปจฺจโย ฯเปฯ. เนว\-สญฺญา\-นา\-สญฺญา\-ยตนสฺส ปริกมฺมํ เนว\-สญฺญา\-นา\-สญฺญา\-ยตนสฺส อุปนิสฺสย\-ปจฺจเยน ปจฺจโย. (2)}

\csromanheader{2}{\textbf{ปุเรชาต}}
\proseitem{287}{นอรูปาวจโร ธมฺโม นอรูปาวจรสฺส ธมฺมสฺส ปุเรชาต\-ปจฺจเยน ปจฺจโย. อารมฺมณ\-ปุเรชาตํ, วตฺถุ\-ปุเรชาตํ.  อารมฺมณ\-ปุเรชาตํ\csromandash{}จกฺขุํ ฯเปฯ. วตฺถุํ อนิจฺจโต ฯเปฯ โทมนสฺสํ อุปฺปชฺชติ. ทิพฺเพน จกฺขุนา รูปํ ปสฺสติ. ทิพฺพาย โสตธาตุยา สทฺทํ สุณาติ. รูปายตนํ จกฺขุ\-วิญฺญาณสฺส ฯเปฯ. โผฏฺฐพฺพา\-ยตนํ กายวิญฺญาณสฺส ฯเปฯ. วตฺถุ\-ปุเรชาตํ\csromandash{} \csromanfolio{166}จกฺขายตนํ จกฺขุ\-วิญฺญาณสฺส ฯเปฯ. กายายตนํ กายวิญฺญาณสฺส ฯเปฯ. วตฺถุ นอรูปาวจรานํ ขนฺธานํ ปุเรชาต\-ปจฺจเยน ปจฺจโย.  นอรูปาวจโร ธมฺโม อรูปาวจรสฺส ธมฺมสฺส ปุเรชาต\-ปจฺจเยน ปจฺจโย. วตฺถุ\-ปุเรชาตํ\csromandash{}วตฺถุ อรูปาวจรานํ ขนฺธานํ ปุเรชาต\-ปจฺจเยน ปจฺจโย. (2)}

\prose{ปจฺฉาชาต\-ปจฺจเยน ปจฺจโย. เทฺว. อาเสวน\-ปจฺจเยน ปจฺจโย. ตีณิ.}

\csromanheader{2}{\textbf{กมฺม}}
\proseitem{288}{อรูปาวจโร ธมฺโม อรูปาวจรสฺส ธมฺมสฺส กมฺมปจฺจเยน ปจฺจโย. สหชาตา, นานากฺขณิกา ฯเปฯ. อรูปาวจโร ธมฺโม นอรูปาวจรสฺส ธมฺมสฺส กมฺมปจฺจเยน ปจฺจโย. อรูปาวจรา เจตนา จิตฺต\-สมุฏฺฐานานํ รูปานํ กมฺมปจฺจเยน ปจฺจโย. (มูลํ กาตพฺพํ.) อรูปาวจรา เจตนา สมฺปยุตฺตกานํ ขนฺธานํ จิตฺตสมุฏฺฐานานญฺจ รูปานํ กมฺมปจฺจเยน ปจฺจโย. (3)}

\prose{นอรูปาวจโร ธมฺโม นอรูปาวจรสฺส ธมฺมสฺส กมฺมปจฺจเยน ปจฺจโย. สหชาตา, นานากฺขณิกา.  สหชาตา นอรูปาวจรา เจตนา สมฺปยุตฺตกานํ ขนฺธานํ จิตฺตสมุฏฺฐานานญฺจ รูปานํ กมฺมปจฺจเยน ปจฺจโย. (สํขิตฺตํ.) (1)}

\csromanheader{2}{\textbf{วิปากาทิ}}
\proseitem{289}{อรูปาวจโร ธมฺโม อรูปาวจรสฺส ธมฺมสฺส วิปาก\-ปจฺจเยน ปจฺจโย. (สํขิตฺตํ.)}

\prose{นอรูปาวจโร ธมฺโม นอรูปาวจรสฺส ธมฺมสฺส วิปาก\-ปจฺจเยน ปจฺจโย. (สํขิตฺตํ.) อาหารปจฺจเยน ปจฺจโย. จตฺตาริ. อินฺทฺริย\-ปจฺจเยน ปจฺจโย. จตฺตาริ. ฌานปจฺจเยน ปจฺจโย. จตฺตาริ. มคฺคปจฺจเยน ปจฺจโย. จตฺตาริ. สมฺปยุตฺต\-ปจฺจเยน ปจฺจโย. เทฺว.}

\csromanheader{2}{\textbf{วิปฺปยุตฺต}}
\proseitem{290}{อรูปาวจโร ธมฺโม นอรูปาวจรสฺส ธมฺมสฺส วิปฺปยุตฺต\-ปจฺจเยน ปจฺจโย. สหชาตํ, ปจฺฉาชาตํ. (สํขิตฺตํ.) (1)}

\prose{นอรูปาวจโร ธมฺโม นอรูปาวจรสฺส ธมฺมสฺส วิปฺปยุตฺต\-ปจฺจเยน ปจฺจโย. สหชาตํ, ปุเรชาตํ, ปจฺฉาชาตํ. (สํขิตฺตํ.) นอรูปาวจโร \csromanfolio{167}ธมฺโม อรูปาวจรสฺส ธมฺมสฺส วิปฺปยุตฺต\-ปจฺจเยน ปจฺจโย. ปุเรชาตํ วตฺถุ อรูปาวจรานํ ขนฺธานํ วิปฺปยุตฺต\-ปจฺจเยน ปจฺจโย. (2)}

\prose{อตฺถิอาทิ}

\proseitem{291}{อรูปาวจโร ธมฺโม อรูปาวจรสฺส ธมฺมสฺส อตฺถิปจฺจเยน ปจฺจโย. สหชาตํ ฯเปฯ. อรูปาวจโร ธมฺโม นอรูปาวจรสฺส ธมฺมสฺส อตฺถิปจฺจเยน ปจฺจโย. สหชาตํ, ปจฺฉาชาตํ ฯเปฯ. อรูปาวจโร ธมฺโม อรูปาวจรสฺส จ นอรูปาวจรสฺส จ ธมฺมสฺส อตฺถิปจฺจเยน ปจฺจโย. สหชาตํ. (สํขิตฺตํ.) (3)}

\prose{นอรูปาวจโร ธมฺโม นอรูปาวจรสฺส ธมฺมสฺส อตฺถิปจฺจเยน ปจฺจโย. สหชาตํ, ปุเรชาตํ, ปจฺฉาชาตํ, อาหารํ, อินฺทฺริยํ. (สํขิตฺตํ).  นอรูปาวจโร ธมฺโม อรูปาวจรสฺส ธมฺมสฺส อตฺถิปจฺจเยน ปจฺจโย. \textbf{ปุเรชาตํ}\csromandash{}วตฺถุ อรูปาวจรานํ ขนฺธานํ อตฺถิปจฺจเยน ปจฺจโย. (2)}

\proseitem{292}{อรูปาวจโร จ นอรูปาวจโร จ ธมฺมา อรูปาวจรสฺส ธมฺมสฺส อตฺถิปจฺจเยน ปจฺจโย. สหชาตํ, ปุเรชาตํ.  สหชาโต\csromandash{} อรูปาวจโร เอโก ขนฺโธ จ วตฺถุ จ ติณฺณนฺนํ ขนฺธานํ อตฺถิปจฺจเยน ปจฺจโย ฯเปฯ เทฺว ขนฺธา จ ฯเปฯ.  อรูปาวจโร จ นอรูปาวจโร จ ธมฺมา นอรูปาวจรสฺส ธมฺมสฺส อตฺถิปจฺจเยน ปจฺจโย. สหชาตํ, ปจฺฉาชาตํ, อาหารํ, อินฺทฺริยํ.  สหชาตา\csromandash{} อรูปาวจรา ขนฺธา จ มหาภูตา จ จิตฺต\-สมุฏฺฐานานํ รูปานํ อตฺถิปจฺจเยน ปจฺจโย. ปจฺฉาชาตา\csromandash{}อรูปาวจรา ขนฺธา จ กพฬีกาโร อาหาโร จ อิมสฺส กายสฺส อตฺถิปจฺจเยน ปจฺจโย. ปจฺฉาชาตา\csromandash{}อรูปาวจรา ขนฺธา จ รูปชีวิตินฺทฺริยญฺจ กฏตฺตารูปานํ อตฺถิปจฺจเยน ปจฺจโย. (2)}

\prose{นตฺถิ\-ปจฺจเยน ปจฺจโย. วิคต\-ปจฺจเยน ปจฺจโย. อวิคต\-ปจฺจเยน ปจฺจโย.}

\csromanmarkh{1. ปจฺจยานุโลม}\csromanmarkhii{2. สงฺขฺยาวาร}\csromanheadernomark{2}{1. \textbf{ปจฺจยานุโลม 2. สงฺขฺยาวาร}}
\proseitem{293}{เหตุยา จตฺตาริ, อารมฺมเณ ตีณิ, อธิปติยา จตฺตาริ, อนนฺตเร จตฺตาริ, สมนนฺตเร จตฺตาริ, สหชาเต ปญฺจ, อญฺญมญฺเญ เทฺว, นิสฺสเย สตฺต, อุปนิสฺสเย จตฺตาริ, ปุเรชาเต เทฺว, ปจฺฉาชาเต เทฺว, อาเสวเน ตีณิ, กมฺเม จตฺตาริ, วิปาเก เทฺว, อาหาเร จตฺตาริ, \csromanfolio{168}อินฺทฺริเย จตฺตาริ, ฌาเน จตฺตาริ, มคฺเค จตฺตาริ, สมฺปยุตฺเต เทฺว, วิปฺปยุตฺเต ตีณิ, \textbf{อตฺถิยา สตฺต, นตฺถิยา จตฺตาริ, วิคเต จตฺตาริ, อวิคเต สตฺต.}}

\proseitem{294}{อรูปาวจโร ธมฺโม อรูปาวจรสฺส ธมฺมสฺส อารมฺมณ\-ปจฺจเยน ปจฺจโย. สหชาต\-ปจฺจเยน ปจฺจโย. อุปนิสฺสย\-ปจฺจเยน ปจฺจโย.  อรูปาวจโร ธมฺโม นอรูปาวจรสฺส ธมฺมสฺส อารมฺมณ\-ปจฺจเยน ปจฺจโย. สหชาต\-ปจฺจเยน ปจฺจโย. อุปนิสยปจฺจเยน ปจฺจโย. ปจฺฉาชาต\-ปจฺจเยน ปจฺจโย.  อรูปาวจโร ธมฺโม อรูปาวจรสฺส จ นอรูปาวจรสฺส จ ธมฺมสฺส สหชาต\-ปจฺจเยน ปจฺจโย. (3)}

\proseitem{295}{นอรูปาวจโร ธมฺโม นอรูปาวจรสฺส ธมฺมสฺส อารมฺมณ\-ปจฺจเยน ปจฺจโย. สหชาต\-ปจฺจเยน ปจฺจโย. อุปนิสฺสย\-ปจฺจเยน ปจฺจโย. ปุเรชาต\-ปจฺจเยน ปจฺจโย. ปจฺฉาชาต\-ปจฺจเยน ปจฺจโย. กมฺมปจฺจเยน ปจฺจโย. อาหารปจฺจเยน ปจฺจโย. อินฺทฺริย\-ปจฺจเยน ปจฺจโย.  นอรูปาวจโร ธมฺโม อรูปาวจรสฺส ธมฺมสฺส อุปนิสฺสย\-ปจฺจเยน ปจฺจโย. ปุเรชาต\-ปจฺจเยน ปจฺจโย. (2)}

\prose{อรูปาวจโร จ นอรูปาวจโร จ ธมฺมา อรูปาวจรสฺส ธมฺมสฺส สหชาตํ. ปุเรชาตํ.  อรูปาวจโร จ นอรูปาวจโร จ ธมฺมา นอรูปาวจรสฺส ธมฺมสฺส สหชาตํ, ปจฺฉาชาตํ, อาหารํ, อินฺทฺริยํ. (2)}

\csromanmarkh{2. ปจฺจย\-ปจฺจนีย}\csromanmarkhii{2. สงฺขฺยาวาร}\csromanheadernomark{2}{2. ปจฺจย\-ปจฺจนีย 2. สงฺขฺยาวาร}
\proseitem{296}{นเหตุยา สตฺต, นอารมฺมเณ สตฺต, นอธิปติยา สตฺต, นอนนฺตเร สตฺต, นสมนนฺตเร สตฺต, นสหชาเต ปญฺจ, นอญฺญมญฺเญ ปญฺจ, นนิสฺสเย ปญฺจ, นอุปนิสฺสเย สตฺต, นปุเรชาเต ฉ, นปจฺฉาชาเต สตฺต ฯเปฯ นสมฺปยุตฺเต ปญฺจ, นวิปฺปยุตฺเต จตฺตาริ, โนอตฺถิยา จตฺตาริ, โนนตฺถิยา สตฺต, โนวิคเต สตฺต, โนอวิคเต จตฺตาริ.}

\csromanmatikaanchor{mtk.25}
\csromanmatikaanchor{mtk.26}
\csromantocmarkref{section}{96. ปริยาปนฺนทุก}{mtk.25}
\bookmark[dest={mtk.25},level=1]{96. ปริยาปนฺนทุก}
\csromantocmarkref{section}{97. นิยฺยานิกทุก}{mtk.26}
\bookmark[dest={mtk.26},level=1]{97. นิยฺยานิกทุก}
\csromanmarkh{97. นิยฺยานิกทุก}\csromanmarkhii{3. ปจฺจยา\-นุโลม\-ปจฺจนีย}\csromanheadernomark{1}{97. นิยฺยานิกทุก 3. ปจฺจยา\-นุโลม\-ปจฺจนีย}
\proseitem{297}{\csromanfolio{169}เหตุปจฺจยา นอารมฺมเณ จตฺตาริ, นอธิปติยา จตฺตาริ, นอนนฺตเร นสมนนฺตเร จตฺตาริ, นอญฺญมญฺเญ เทฺว, นอุปนิสฺสเย จตฺตาริ ฯเปฯ นสมฺปยุตฺเต เทฺว, นวิปฺปยุตฺเต เทฺว, โนนตฺถิยา จตฺตาริ, โนวิคเต จตฺตาริ.}

\csromanheader{2}{4. ปจฺจย\-ปจฺจนียา\-นุโลม}
\proseitemwithcloser{298}{นเหตุปจฺจยา อารมฺมเณ ตีณิ, อธิปติยา จตฺตาริ, (อนุโลมมาติกา กาตพฺพา) ฯเปฯ อวิคเต สตฺต.}{อรูปาวจรทุกํ นิฏฺฐิตํ.}
\csromanmarkh{96. ปริยาปนฺนทุก}\csromanmarkhii{1. ปฏิจฺจวาร}\csromanheadernomark{1}{96. \textbf{ปริยาปนฺนทุก 1. ปฏิจฺจวาร}}
\proseitemwithcloser{299}{ปริยาปนฺนํ ธมฺมํ ปฏิจฺจ ปริยาปนฺโน ธมฺโม อุปฺปชฺชติ เหตุปจฺจยา. ปริยาปนฺนํ เอกํ ขนฺธํ ปฏิจฺจ ตโย ขนฺธา จิตฺต\-สมุฏฺฐานญฺจ รูปํ ฯเปฯ เทฺว ขนฺเธ ฯเปฯ. ปฏิสนฺธิ\-กฺขเณ ฯเปฯ. (ยถา จูฬนฺตรทุเก โลกิยทุกํ. เอวํ อิมมฺปิ ทุกํ กาตพฺพํ, นินฺนานากรณํ.)}{ปริยาปนฺนทุกํ นิฏฺฐิตํ.}
\csromanmarkh{97. นิยฺยานิกทุก}\csromanmarkhii{1. ปฏิจฺจวาร}\csromanheadernomark{2}{97. นิยฺยานิกทุก 1. ปฏิจฺจวาร}
\proseitem{300}{นิยฺยานิกํ ธมฺมํ ปฏิจฺจ นิยฺยานิโก ธมฺโม อุปฺปชฺชติ เหตุปจฺจยา. นิยฺยานิกํ เอกํ ขนฺธํ ปฏิจฺจ ตโย ขนฺธา ฯเปฯ เทฺว ขนฺเธ ฯเปฯ. นิยฺยานิกํ ธมฺมํ ปฏิจฺจ อนิยฺยานิโก ธมฺโม อุปฺปชฺชติ เหตุปจฺจยา. นิยฺยานิเก ขนฺเธ ปฏิจฺจ จิตฺต\-สมุฏฺฐานํ รูปํ.  นิยฺยานิกํ ธมฺมํ ปฏิจฺจ นิยฺยานิโก จ อนิยฺยานิโก จ ธมฺมา อุปฺปชฺชนฺติ เหตุปจฺจยา. นิยฺยานิกํ เอกํ ขนฺธํ ปฏิจฺจ ตโย ขนฺธา จิตฺต\-สมุฏฺฐานญฺจ รูปํ ฯเปฯ เทฺว ขนฺเธ ฯเปฯ. (3)}

\prose{\csromanfolio{170}อนิยฺยานิกํ ธมฺมํ ปฏิจฺจ อนิยฺยานิโก ธมฺโม อุปฺปชฺชติ เหตุปจฺจยา. อนิยฺยานิกํ เอกํ ขนฺธํ ปฏิจฺจ ตโย ขนฺธา จิตฺต\-สมุฏฺฐานญฺจ รูปํ ฯเปฯ เทฺว ขนฺเธ ฯเปฯ. ปฏิสนฺธิ\-กฺขเณ ฯเปฯ. ขนฺเธ ปฏิจฺจ วตฺถุ, วตฺถุํ ปฏิจฺจ ขนฺธา. เอกํ มหาภูตํ ฯเปฯ. (1)}

\prose{นิยฺยานิกญฺจ อนิยฺยานิกญฺจ ธมฺมํ ปฏิจฺจ อนิยฺยานิโก ธมฺโม อุปฺปชฺชติ เหตุปจฺจยา. นิยฺยานิเก ขนฺเธ จ มหาภูเต จ ปฏิจฺจ จิตฺต\-สมุฏฺฐานํ รูปํ. (1)}

\csromanmarkh{1. ปจฺจยานุโลม}\csromanmarkhii{2. สงฺขฺยาวาร}\csromanheadernomark{2}{1. \textbf{ปจฺจยานุโลม 2. สงฺขฺยาวาร}}
\proseitem{301}{เหตุยา ปญฺจ, อารมฺมเณ เทฺว, อธิปติยา ปญฺจ, อนนฺตเร สมนนฺตเร เทฺว, สหชาเต ปญฺจ, อญฺญมญฺเญ เทฺว, นิสฺสเย ปญฺจ, อุปนิสฺสเย เทฺว, ปุเรชาเต เทฺว, อาเสวเน เทฺว, กมฺเม ปญฺจ, วิปาเก เอกํ, อาหาเร ปญฺจ ฯเปฯ อวิคเต ปญฺจ.}

\csromanmarkh{2. ปจฺจย\-ปจฺจนีย}\csromanmarkhii{1. วิภงฺควาร}\csromanheadernomark{2}{2. ปจฺจย\-ปจฺจนีย 1. วิภงฺควาร}
\prose{นเหตุ นอารมฺมณ}

\proseitem{302}{อนิยฺยานิกํ ธมฺมํ ปฏิจฺจ อนิยฺยานิโก ธมฺโม อุปฺปชฺชติ นเหตุปจฺจยา. อเหตุกํ อนิยฺยานิกํ เอกํ ขนฺธํ ปฏิจฺจ ตโย ขนฺธา จิตฺต\-สมุฏฺฐานญฺจ รูปํ ฯเปฯ เทฺว ขนฺเธ ฯเปฯ. อเหตุก\-ปฏิสนฺธิ\-กฺขเณ ฯเปฯ. ขนฺเธ ปฏิจฺจ วตฺถุ, วตฺถุํ ปฏิจฺจ ขนฺธา. เอกํ มหาภูตํ ฯเปฯ. (ยาว อสญฺญสตฺตา.) วิจิกิจฺฉา\-สหคเต อุทฺธจฺจ\-สหคเต ขนฺเธ ปฏิจฺจ วิจิกิจฺฉา\-สหคโต อุทฺธจฺจ\-สหคโต โมโห. (1). นอารมฺมณ\-ปจฺจยา. ตีณิ.}

\prose{นอธิปตฺยาทิ}

\proseitem{303}{นิยฺยานิกํ ธมฺมํ ปฏิจฺจ นิยฺยานิโก ธมฺโม อุปฺปชฺชติ นอธิปติ\-ปจฺจยา. นิยฺยานิเก ขนฺเธ ปฏิจฺจ นิยฺยานิกา\-ธิปติ. (1)}

\prose{อนิยฺยานิกํ ธมฺมํ ปฏิจฺจ อนิยฺยานิโก ธมฺโม อุปฺปชฺชติ นอธิปติ\-ปจฺจยา. อนิยฺยานิกํ เอกํ ขนฺธํ ปฏิจฺจ ตโย ขนฺธา จิตฺต\-สมุฏฺฐานญฺจ รูปํ ฯเปฯ เทฺว ขนฺเธ ฯเปฯ. ปฏิสนฺธิ\-กฺขเณ ฯเปฯ. (ยาว อสญฺญสตฺตา กาตพฺพา.) (1)}

\prose{นอนนฺตร\-ปจฺจยา. นสมนนฺตร\-ปจฺจยา. นอญฺญมญฺญ\-ปจฺจยา ฯเปฯ.}

\csromanheader{2}{97. นิยฺยานิกทุก \textbf{นปุเรชาต}}
\proseitem{304}{\csromanfolio{171}นิยฺยานิกํ ธมฺมํ ปฏิจฺจ นิยฺยานิโก ธมฺโม อุปฺปชฺชติ นปุเรชาต\-ปจฺจยา. อรูเป นิยฺยานิกํ เอกํ ขนฺธํ ฯเปฯ. นิยฺยานิกํ ธมฺมํ ปฏิจฺจ อนิยฺยานิโก ธมฺโม อุปฺปชฺชติ นปุเรชาต\-ปจฺจยา. นิยฺยานิเก ขนฺเธ ปฏิจฺจ จิตฺต\-สมุฏฺฐานํ รูปํ. (2)}

\prose{อนิยฺยานิกํ ธมฺมํ ปฏิจฺจ อนิยฺยานิโก ธมฺโม อุปฺปชฺชติ นปุเรชาต\-ปจฺจยา. อรูเป อนิยฺยานิกํ เอกํ ขนฺธํ ปฏิจฺจ ตโย ขนฺธา ฯเปฯ เทฺว ขนฺเธ ฯเปฯ. อนิยฺยานิเก ขนฺเธ ปฏิจฺจ จิตฺต\-สมุฏฺฐานํ รูปํ. ปฏิสนฺธิ\-กฺขเณ ฯเปฯ. (ยาว อสญฺญสตฺตา.) (1)}

\prose{นิยฺยานิกญฺจ อนิยฺยานิกญฺจ ธมฺมํ ปฏิจฺจ อนิยฺยานิโก ธมฺโม อุปฺปชฺชติ นปุเรชาต\-ปจฺจยา. นิยฺยานิเก ขนฺเธ จ มหาภูเต จ ปฏิจฺจ จิตฺต\-สมุฏฺฐานํ รูปํ. (1)}

\csromanmarkh{2. ปจฺจย\-ปจฺจนีย}\csromanmarkhii{2. สงฺขฺยาวาร}\csromanheadernomark{2}{2. ปจฺจย\-ปจฺจนีย 2. สงฺขฺยาวาร}
\proseitem{305}{นเหตุยา เอกํ, นอารมฺมเณ ตีณิ, นอธิปติยา เทฺว, นอนนฺตเร ตีณิ, นอุปนิสฺสเย ตีณิ, นปุเรชาเต จตฺตาริ, นปจฺฉาชาเต ปญฺจ, นอาเสวเน เอกํ, นกมฺเม เทฺว, นวิปาเก ปญฺจ, นอาหาเร เอกํ, นอินฺทฺริเย เอกํ, นฌาเน เอกํ, นมคฺเค เอกํ, นสมฺปยุตฺเต ตีณิ, นวิปฺปยุตฺเต เทฺว, โนนตฺถิยา ตีณิ, โนวิคเต ตีณิ.}

\vspace{\tipitakaparskip}
\csromancenter{(เอวํ อิตเร เทฺว คณนาปิ สหชาตวาโรปิ กาตพฺโพ.)}
\csromanmarkh{97. นิยฺยานิกทุก}\csromanmarkhii{3. ปจฺจยวาร}\csromanheadernomark{2}{97. นิยฺยานิกทุก 3. ปจฺจยวาร}
\proseitem{306}{นิยฺยานิกํ ธมฺมํ ปจฺจยา นิยฺยานิโก ธมฺโม อุปฺปชฺชติ เหตุปจฺจยา. ตีณิ. (ปฏิจฺจสทิสา.)}

\prose{\csromanfolio{172}อนิยฺยานิกํ ธมฺมํ ปจฺจยา อนิยฺยานิโก ธมฺโม อุปฺปชฺชติ เหตุปจฺจยา. อนิยฺยานิกํ เอกํ ขนฺธํ ปจฺจยา ตโย ขนฺธา จิตฺต\-สมุฏฺฐานญฺจ รูปํ ฯเปฯ เทฺว ขนฺเธ ฯเปฯ. (ยาว อชฺฌตฺติกา มหาภูตา.) วตฺถุํ ปจฺจยา อนิยฺยานิกา ขนฺธา.  อนิยฺยานิกํ ธมฺมํ ปจฺจยา นิยฺยานิโก ธมฺโม อุปฺปชฺชติ เหตุปจฺจยา. วตฺถุํ ปจฺจยา นิยฺยานิกา ขนฺธา.  อนิยฺยานิกํ ธมฺมํ ปจฺจยา นิยฺยานิโก จ อนิยฺยานิโก จ ธมฺมา อุปฺปชฺชนฺติ เหตุปจฺจยา. วตฺถุํ ปจฺจยา นิยฺยานิกา ขนฺธา. มหาภูเต ปจฺจยา จิตฺต\-สมุฏฺฐานํ รูปํ. (3)}

\proseitem{307}{นิยฺยานิกญฺจ อนิยฺยานิกญฺจ ธมฺมํ ปจฺจยา นิยฺยานิโก ธมฺโม อุปฺปชฺชติ เหตุปจฺจยา. นิยฺยานิกํ เอกํ ขนฺธญฺจ วตฺถุญฺจ ปจฺจยา ตโย ขนฺธา ฯเปฯ เทฺว ขนฺเธ จ ฯเปฯ. นิยฺยานิกญฺจ อนิยฺยานิกญฺจ ธมฺมํ ปจฺจยา อนิยฺยานิโก ธมฺโม อุปฺปชฺชติ เหตุปจฺจยา. นิยฺยานิเก ขนฺเธ จ มหาภูเต จ ปจฺจยา จิตฺต\-สมุฏฺฐานํ รูปํ.  นิยฺยานิกญฺจ อนิยฺยานิกญฺจ ธมฺมํ ปจฺจยา นิยฺยานิโก จ อนิยฺยานิโก จ ธมฺมา อุปฺปชฺชนฺติ เหตุปจฺจยา. นิยฺยานิกํ เอกํ ขนฺธญฺจ วตฺถุญฺจ ปจฺจยา ตโย ขนฺธา ฯเปฯ เทฺว ขนฺเธ จ ฯเปฯ. นิยฺยานิเก ขนฺเธ จ มหาภูเต จ ปจฺจยา จิตฺต\-สมุฏฺฐานํ รูปํ. (3)}

\csromanmarkh{1. ปจฺจยานุโลม}\csromanmarkhii{2. สงฺขฺยาวาร}\csromanheadernomark{2}{1. \textbf{ปจฺจยานุโลม 2. สงฺขฺยาวาร}}
\proseitem{308}{เหตุยา นว, อารมฺมเณ จตฺตาริ, อธิปติยา นว, อนนฺตเร จตฺตาริ, สมนนฺตเร จตฺตาริ, สหชาเต นว, อญฺญมญฺเญ จตฺตาริ, นิสฺสเย นว, อุปนิสฺสเย จตฺตาริ, ปุเรชาเต จตฺตาริ, อาเสวเน จตฺตาริ, กมฺเม นว, วิปาเก เอกํ ฯเปฯ อวิคเต นว.}

\csromanmarkh{2. ปจฺจย\-ปจฺจนีย}\csromanmarkhii{1. วิภงฺควาร}\csromanheadernomark{2}{2. ปจฺจย\-ปจฺจนีย 1. วิภงฺควาร}
\proseitem{309}{อนิยฺยานิกํ ธมฺมํ ปจฺจยา อนิยฺยานิโก ธมฺโม อุปฺปชฺชติ นเหตุปจฺจยา. อเหตุกํ อนิยฺยานิกํ เอกํ ขนฺธํ ปจฺจยา ตโย ขนฺธา จิตฺต\-สมุฏฺฐานญฺจ รูปํ ฯเปฯ เทฺว ขนฺเธ ฯเปฯ. (ยาว อสญฺญสตฺตา.) จกฺขายตนํ ปจฺจยา จกฺขุวิญฺญาณํ ฯเปฯ. กายายตนํ ปจฺจยา กายวิญฺญาณํ. วตฺถุํ ปจฺจยา อเหตุกา อนิยฺยานิกา ขนฺธา. วิจิกิจฺฉา\-สหคเต อุทฺธจฺจ\-สหคเต ขนฺเธ จ วตฺถุญฺจ ปจฺจยา วิจิกิจฺฉา\-สหคโต อุทฺธจฺจ\-สหคโต โมโห. (1)}

\csromanheader{2}{97. นิยฺยานิกทุก นอารมฺมณาทิ}
\proseitem{310}{\csromanfolio{173}นิยฺยานิกํ ธมฺมํ ปจฺจยา อนิยฺยานิโก ธมฺโม อุปฺปชฺชติ นอารมฺมณ\-ปจฺจยา. ตีณิ.}

\prose{นิยฺยานิกํ ธมฺมํ ปจฺจยา นิยฺยานิโก ธมฺโม อุปฺปชฺชติ นอธิปติ\-ปจฺจยา. นิยฺยานิเก ขนฺเธ ปจฺจยา นิยฺยานิกา\-ธิปติ. (1)}

\prose{อนิยฺยานิกํ ธมฺมํ ปจฺจยา อนิยฺยานิโก ธมฺโม อุปฺปชฺชติ นอธิปติ\-ปจฺจยา. อนิยฺยานิกํ เอกํ ขนฺธํ ปจฺจยา ฯเปฯ. (ยาว อสญฺญสตฺตา.) จกฺขายตนํ ปจฺจยา จกฺขุวิญฺญาณํ ฯเปฯ. กายายตนํ ปจฺจยา กายวิญฺญาณํ. วตฺถุํ ปจฺจยา อนิยฺยานิกา ขนฺธา.  อนิยฺยานิกํ ธมฺมํ ปจฺจยา นิยฺยานิโก ธมฺโม อุปฺปชฺชติ นอธิปติ\-ปจฺจยา. วตฺถุํ ปจฺจยา นิยฺยานิกา\-ธิปติ. (2)}

\prose{นิยฺยานิกญฺจ อนิยฺยานิกญฺจ ธมฺมํ ปจฺจยา นิยฺยานิโก ธมฺโม อุปฺปชฺชติ นอธิปติ\-ปจฺจยา. นิยฺยานิเก ขนฺเธ จ วตฺถุญฺจ ปจฺจยา นิยฺยานิกา\-ธิปติ. (1)}

\csromanmarkh{2. ปจฺจย\-ปจฺจนีย}\csromanmarkhii{2. สงฺขฺยาวาร}\csromanheadernomark{2}{2. ปจฺจย\-ปจฺจนีย 2. สงฺขฺยาวาร}
\proseitem{311}{นเหตุยา เอกํ, นอารมฺมเณ ตีณิ, นอธิปติยา จตฺตาริ, นอนนฺตเร ตีณิ ฯเปฯ นอุปนิสฺสเย ตีณิ, นปุเรชาเต จตฺตาริ, นปจฺฉาชาเต นว, นอาเสวเน เอกํ, นกมฺเม จตฺตาริ, นวิปาเก นว, นอาหาเร เอกํ, นอินฺทฺริเย เอกํ, นฌาเน เอกํ, นมคฺเค เอกํ, นสมฺปยุตฺเต ตีณิ, นวิปฺปยุตฺเต เทฺว, โนนตฺถิยา ตีณิ, โนวิคเต ตีณิ.}

\vspace{\tipitakaparskip}
\csromancenter{(เอวํ อิตเร เทฺว คณนาปิ นิสฺสยวาโรปิ กาตพฺโพ.)}
\csromanmarkh{97. นิยฺยานิกทุก}\csromanmarkhii{5. สํสฏฺฐวาร}\csromanheadernomark{2}{97. นิยฺยานิกทุก 5. สํสฏฺฐวาร}
\proseitem{312}{นิยฺยานิกํ ธมฺมํ สํสฏฺโฐ นิยฺยานิโก ธมฺโม อุปฺปชฺชติ เหตุปจฺจยา. นิยฺยานิกํ เอกํ ขนฺธํ สํสฏฺฐา ตโย ขนฺธา ฯเปฯ เทฺว ขนฺเธ ฯเปฯ. (1)}

\prose{\csromanfolio{174}อนิยฺยานิกํ ธมฺมํ สํสฏฺโฐ อนิยฺยานิโก ธมฺโม อุปฺปชฺชติ เหตุปจฺจยา. อนิยฺยานิกํ เอกํ ขนฺธํ สํสฏฺฐา ตโย ขนฺธา ฯเปฯ เทฺว ขนฺเธ ฯเปฯ. ปฏิสนฺธิ\-กฺขเณ ฯเปฯ. (1)}

\prose{เหตุยา เทฺว, อารมฺมเณ เทฺว, (สพฺพตฺถ เทฺว,) วิปาเก เอกํ ฯเปฯ อวิคเต เทฺว. (อนุโลมํ.)}

\prose{นเหตุยา เอกํ, นอธิปติยา เทฺว, นปุเรชาเต เทฺว, นปจฺฉาชาเต เทฺว, นอาเสวเน เอกํ, นกมฺเม เทฺว, นวิปาเก เทฺว, นฌาเน เอกํ, นมคฺเค เอกํ, นวิปฺปยุตฺเต เทฺว. (ปจฺจนียํ.)}

\vspace{\tipitakaparskip}
\csromancenter{(เอวํ อิตเร เทฺว คณนาปิ สมฺปยุตฺตวาโรปิ กาตพฺโพ.)}
\csromanmarkh{97. นิยฺยานิกทุก}\csromanmarkhii{7. ปญฺหาวาร}\csromanheadernomark{2}{97. \textbf{นิยฺยานิกทุก 7. ปญฺหาวาร}}
\proseitem{313}{นิยฺยานิโก ธมฺโม นิยฺยานิกสฺส ธมฺมสฺส เหตุปจฺจเยน ปจฺจโย. นิยฺยานิกา เหตู สมฺปยุตฺตกานํ ขนฺธานํ เหตุปจฺจเยน ปจฺจโย. (มูลํ กาตพฺพํ.) นิยฺยานิกา เหตู จิตฺต\-สมุฏฺฐานานํ รูปานํ เหตุปจฺจเยน ปจฺจโย. (มูลํ กาตพฺพํ.) นิยฺยานิกา เหตู สมฺปยุตฺตกานํ ขนฺธานํ จิตฺตสมุฏฺฐานานญฺจ รูปานํ เหตุปจฺจเยน ปจฺจโย. (3)}

\prose{อนิยฺยานิโก ธมฺโม อนิยฺยานิกสฺส ธมฺมสฺส เหตุปจฺจเยน ปจฺจโย. อนิยฺยานิกา เหตู สมฺปยุตฺตกานํ ขนฺธานํ จิตฺตสมุฏฺฐานานญฺจ รูปานํ เหตุปจฺจเยน ปจฺจโย. ปฏิสนฺธิ\-กฺขเณ ฯเปฯ. (1)}

\proseitem{314}{นิยฺยานิโก ธมฺโม อนิยฺยานิกสฺส ธมฺมสฺส อารมฺมณ\-ปจฺจเยน ปจฺจโย. อริยา มคฺคา วุฏฺฐหิตฺวา มคฺคํ ปจฺจเวกฺขนฺติ. เจโตปริย\-ญาเณน นิยฺยานิก\-จิตฺต\-สมงฺคิสฺส จิตฺตํ ชานาติ. นิยฺยานิกา ขนฺธา เจโตปริย\-ญาณสฺส, ปุพฺเพ\-นิวาสา\-นุสฺสติ\-ญาณสฺส, อนาคตํสญาณสฺส, อาวชฺชนาย อารมฺมณ\-ปจฺจเยน ปจฺจโย. (1)}

\prose{\csromanfolio{175}อนิยฺยานิโก ธมฺโม อนิยฺยานิกสฺส ธมฺมสฺส อารมฺมณ\-ปจฺจเยน ปจฺจโย. ทานํ ฯเปฯ. สีลํ ฯเปฯ อุโปสถกมฺมํ กตฺวา ตํ ปจฺจเวกฺขติ, อสฺสาเทติ อภินนฺทติ, ตํ อารพฺภ ราโค อุปฺปชฺชติ ฯเปฯ โทมนสฺสํ อุปฺปชฺชติ. ปุพฺเพ สุจิณฺณานิ ฯเปฯ. ฌานา ฯเปฯ. อริยา ผลํ ปจฺจเวกฺขนฺติ, นิพฺพานํ ปจฺจเวกฺขนฺติ. นิพฺพานํ โคตฺรภุสฺส, โวทานสฺส, ผลสฺส, อาวชฺชนาย อารมฺมณ\-ปจฺจเยน ปจฺจโย. อริยา ปหีเน กิเลเส ฯเปฯ. วิกฺขมฺภิเต ฯเปฯ. ปุพฺเพ สมุทาจิณฺเณ ฯเปฯ. จกฺขุํ ฯเปฯ. วตฺถุํ อนิยฺยานิเก ขนฺเธ อนิจฺจโต ฯเปฯ โทมนสฺสํ อุปฺปชฺชติ. ทิพฺเพน จกฺขุนา รูปํ ปสฺสติ. ทิพฺพาย โสตธาตุยา สทฺทํ สุณาติ. เจโตปริย\-ญาเณน อนิยฺยานิก\-จิตฺต\-สมงฺคิสฺส จิตฺตํ ชานาติ. อากาสา\-นญฺจา\-ยตนํ วิญฺญาณญฺจา\-ยตนสฺส ฯเปฯ. อากิญฺจญฺญา\-ยตนํ เนว\-สญฺญา\-นา\-สญฺญา\-ยตนสฺส ฯเปฯ. รูปายตนํ จกฺขุ\-วิญฺญาณสฺส ฯเปฯ. โผฏฺฐพฺพา\-ยตนํ กายวิญฺญาณสฺส ฯเปฯ. อนิยฺยานิกา ขนฺธา อิทฺธิวิธ\-ญาณสฺส, เจโตปริย\-ญาณสฺส, ปุพฺเพ\-นิวาสา\-นุสฺสติ\-ญาณสฺส, ยถากมฺมูปคญาณสฺส, อนาคตํสญาณสฺส, อาวชฺชนาย อารมฺมณ\-ปจฺจเยน ปจฺจโย. (1)}

\prose{อนิยฺยานิโก ธมฺโม นิยฺยานิกสฺส ธมฺมสฺส อารมฺมณ\-ปจฺจเยน ปจฺจโย. นิพฺพานํ มคฺคสฺส อารมฺมณ\-ปจฺจเยน ปจฺจโย. (2)}

\csromanheader{2}{\textbf{อธิปติ}}
\proseitem{315}{นิยฺยานิโก ธมฺโม นิยฺยานิกสฺส ธมฺมสฺส อธิปติ\-ปจฺจเยน ปจฺจโย. สหชาตา\-ธิปติ\csromandash{}นิยฺยานิกา\-ธิปติ สมฺปยุตฺตกานํ ขนฺธานํ อธิปติ\-ปจฺจเยน ปจฺจโย.  นิยฺยานิโก ธมฺโม อนิยฺยานิกสฺส ธมฺมสฺส อธิปติ\-ปจฺจเยน ปจฺจโย. อารมฺมณา\-ธิปติ, สหชาตา\-ธิปติ.  อารมฺมณา\-ธิปติ\csromandash{}อริยา มคฺคา วุฏฺฐหิตฺวา มคฺคํ ครุํ กตฺวา ปจฺจเวกฺขนฺติ. สหชาตา\-ธิปติ\csromandash{}นิยฺยานิกา\-ธิปติ จิตฺต\-สมุฏฺฐานานํ รูปานํ อธิปติ\-ปจฺจเยน ปจฺจโย. (มูลํ กาตพฺพํ.) นิยฺยานิกา\-ธิปติ สมฺปยุตฺตกานํ ขนฺธานํ จิตฺตสมุฏฺฐานานญฺจ รูปานํ อธิปติ\-ปจฺจเยน ปจฺจโย. (3)}

\proseitem{316}{อนิยฺยานิโก ธมฺโม อนิยฺยานิกสฺส ธมฺมสฺส อธิปติ\-ปจฺจเยน ปจฺจโย. อารมฺมณา\-ธิปติ, สหชาตา\-ธิปติ.  อารมฺมณา\-ธิปติ\csromandash{}ทานํ ฯเปฯ สีลํ ฯเปฯ อุโปสถกมฺมํ กตฺวา ตํ ครุํ กตฺวา ปจฺจเวกฺขติ, อสฺสาเทติ อภินนฺทติ, ตํ ครุํ กตฺวา ราโค อุปฺปชฺชติ. ทิฏฺฐิ อุปฺปชฺชติ. ปุพฺเพ สุจิณฺณานิ ฯเปฯ. ฌานา วุฏฺฐหิตฺวา ฯเปฯ. อริยา ผลํ ครุํ กตฺวา \csromanfolio{176}ปจฺจเวกฺขนฺติ, นิพฺพานํ ครุํ กตฺวา ปจฺจเวกฺขนฺติ. นิพฺพานํ โคตฺรภุสฺส, โวทานสฺส, ผลสฺส อธิปติ\-ปจฺจเยน ปจฺจโย. จกฺขุํ ฯเปฯ. วตฺถุํ อนิยฺยานิเก ขนฺเธ ครุํ กตฺวา อสฺสาเทติ อภินนฺทติ, ตํ ครุํ กตฺวา ราโค อุปฺปชฺชติ. ทิฏฺฐิ อุปฺปชฺชติ. สหชาตา\-ธิปติ\csromandash{} อนิยฺยานิกา\-ธิปติ สมฺปยุตฺตกานํ ขนฺธานํ จิตฺตสมุฏฺฐานานญฺจ รูปานํ อธิปติ\-ปจฺจเยน ปจฺจโย.  อนิยฺยานิโก ธมฺโม นิยฺยานิกสฺส ธมฺมสฺส อธิปติ\-ปจฺจเยน ปจฺจโย. อารมฺมณา\-ธิปติ\csromandash{}นิพฺพานํ มคฺคสฺส อธิปติ\-ปจฺจเยน ปจฺจโย. (2)}

\proseitem{317}{นิยฺยานิโก ธมฺโม อนิยฺยานิกสฺส ธมฺมสฺส อนนฺตร\-ปจฺจเยน ปจฺจโย. มคฺโค ผลสฺส อนนฺตร\-ปจฺจเยน ปจฺจโย. (1)}

\prose{อนิยฺยานิโก ธมฺโม อนิยฺยานิกสฺส ธมฺมสฺส อนนฺตร\-ปจฺจเยน ปจฺจโย. ปุริมา ปุริมา อนิยฺยานิกา ขนฺธา ปจฺฉิมานํ ปจฺฉิมานํ อนิยฺยานิกานํ ขนฺธานํ อนนฺตร\-ปจฺจเยน ปจฺจโย. อนุโลมํ โคตฺรภุสฺส. อนุโลมํ โวทานสฺส. ผลํ ผลสฺส. อนุโลมํ ผลสมาปตฺติยา. นิโรธา วุฏฺฐหนฺตสฺส เนว\-สญฺญา\-นา\-สญฺญา\-ยตนํ ผลสมาปตฺติยา อนนฺตร\-ปจฺจเยน ปจฺจยา.  อนิยฺยานิโก ธมฺโม นิยฺยานิกสฺส ธมฺมสฺส อนนฺตร\-ปจฺจเยน ปจฺจโย. โคตฺรภุ มคฺคสฺส. โวทานํ มคฺคสฺส อนนฺตร\-ปจฺจเยน ปจฺจโย. (2)}

\prose{สมนนฺตร\-ปจฺจเยน ปจฺจโย. สหชาต\-ปจฺจเยน ปจฺจโย. ปญฺจ. อญฺญมญฺญ\-ปจฺจเยน ปจฺจโย. เทฺว. นิสฺสย\-ปจฺจเยน ปจฺจโย. สตฺต.}

\csromanheader{2}{\textbf{อุปนิสฺสย}}
\proseitem{318}{นิยฺยานิโก ธมฺโม นิยฺยานิกสฺส ธมฺมสฺส อุปนิสฺสย\-ปจฺจเยน ปจฺจโย. ปกตู\-ปนิสฺสโย\csromandash{}ปฐโม มคฺโค ทุติยสฺส มคฺคสฺส อุปนิสฺสย\-ปจฺจเยน ปจฺจโย ฯเปฯ. ตติโย มคฺโค จตุตฺถสฺส มคฺคสฺส อุปนิสฺสย\-ปจฺจเยน ปจฺจโย.  นิยฺยานิโก ธมฺโม อนิยฺยานิกสฺส ธมฺมสฺส อุปนิสฺสย\-ปจฺจเยน ปจฺจโย. อารมฺมณู\-ปนิสฺสโย, อนนฺตรู\-ปนิสฺสโย, ปกตู\-ปนิสฺสโย ฯเปฯ. ปกตู\-ปนิสฺสโย\csromandash{}อริยา มคฺคํ อุปนิสฺสาย อนุปฺปนฺนํ สมาปตฺติํ อุปฺปาเทนฺติ, อุปฺปนฺนํ สมาปชฺชนฺติ, สงฺขาเร อนิจฺจโต ฯเปฯ วิปสฺสนฺติ. มคฺโค อริยานํ อตฺถ\-ปฏิสมฺภิทาย ฯเปฯ. ปฏิภาน\-ปฏิสมฺภิทาย. \csromanfolio{177}ฐานาฐาน\-โกสลฺลสฺส อุปนิสฺสย\-ปจฺจเยน ปจฺจโย. มคฺโค ผลสมาปตฺติยา อุปนิสฺสย\-ปจฺจเยน ปจฺจโย. (2)}

\proseitem{319}{อนิยฺยานิโก ธมฺโม อนิยฺยานิกสฺส ธมฺมสฺส อุปนิสฺสย\-ปจฺจเยน ปจฺจโย. อารมฺมณู\-ปนิสฺสโย, อนนฺตรู\-ปนิสฺสโย, ปกตู\-ปนิสฺสโย ฯเปฯ. ปกตู\-ปนิสฺสโย\csromandash{}อนิยฺยานิกํ สทฺธํ อุปนิสฺสาย ทานํ เทติ. สีลํ ฯเปฯ. อุโปสถกมฺมํ กโรติ. ฌานํ อุปฺปาเทติ. วิปสฺสนํ. อภิญฺญํ. สมาปตฺติํ อุปฺปาเทติ. มานํ ชปฺเปติ. ทิฏฺฐิํ คณฺหาติ. อนิยฺยานิกํ สีลํ ฯเปฯ ปญฺญํ. ราคํ ฯเปฯ ปตฺถนํ. อุตุํ. โภชนํ. เสนาสนํ อุปนิสฺสาย ทานํ เทติ ฯเปฯ. สมาปตฺติํ อุปฺปาเทติ. ปาณํ หนติ ฯเปฯ สํฆํ ภินฺทติ. อนิยฺยานิกา สทฺธา ฯเปฯ. เสนาสนํ อนิยฺยานิกาย สทฺธาย ฯเปฯ ปตฺถนาย. กายิกสฺส สุขสฺส. กายิกสฺส ทุกฺขสฺส. ผลสมาปตฺติยา อุปนิสฺสย\-ปจฺจเยน ปจฺจโย.  อนิยฺยานิโก ธมฺโม นิยฺยานิกสฺส ธมฺมสฺส อุปนิสฺสย\-ปจฺจเยน ปจฺจโย. อารมฺมณู\-ปนิสฺสโย, อนนฺตรู\-ปนิสฺสโย, ปกตู\-ปนิสฺสโย ฯเปฯ. ปกตู\-ปนิสฺสโย\csromandash{}ปฐมสฺส มคฺคสฺส ปริกมฺมํ ปฐมสฺส มคฺคสฺส อุปนิสฺสย\-ปจฺจเยน ปจฺจโย ฯเปฯ. จตุตฺถสฺส มคฺคสฺส ปริกมฺมํ จตุตฺถสฺส มคฺคสฺส อุปนิสฺสย\-ปจฺจเยน ปจฺจโย. (2)}

\csromanheader{2}{\textbf{ปุเรชาตาทิ}}
\proseitem{320}{อนิยฺยานิโก ธมฺโม อนิยฺยานิกสฺส ธมฺมสฺส ปุเรชาต\-ปจฺจเยน ปจฺจโย. เทฺว. (อรูป\-ทุก\-สทิสา กาตพฺพา.) ปจฺฉาชาต\-ปจฺจเยน ปจฺจโย. เทฺว. อาเสวน\-ปจฺจเยน ปจฺจโย. เทฺว.}

\csromanheader{2}{\textbf{กมฺม}}
\proseitem{321}{นิยฺยานิโก ธมฺโม นิยฺยานิกสฺส ธมฺมสฺส กมฺมปจฺจเยน ปจฺจโย. นิยฺยานิกา เจตนา สมฺปยุตฺตกานํ ขนฺธานํ กมฺมปจฺจเยน ปจฺจโย.  นิยฺยานิโก ธมฺโม อนิยฺยานิกสฺส ธมฺมสฺส กมฺมปจฺจเยน ปจฺจโย. สหชาตา, นานากฺขณิกา.  สหชาตา นิยฺยานิกา เจตนา จิตฺต\-สมุฏฺฐานานํ รูปานํ กมฺมปจฺจเยน ปจฺจโย. นานากฺขณิกา นิยฺยานิกา เจตนา ผลสฺส กมฺมปจฺจเยน ปจฺจโย. (มูลํ กาตพฺพํ.) นิยฺยานิกา เจตนา สมฺปยุตฺตกานํ ขนฺธานํ จิตฺตสมุฏฺฐานานญฺจ รูปานํ กมฺมปจฺจเยน ปจฺจโย. (3)}

\proseitem{322}{\csromanfolio{178}อนิยฺยานิโก ธมฺโม อนิยฺยานิกสฺส ธมฺมสฺส กมฺมปจฺจเยน ปจฺจโย. สหชาตา, นานากฺขณิกา.  สหชาตา (สํขิตฺตํ.) นานากฺขณิกา อนิยฺยานิกา เจตนา วิปากานํ ขนฺธานํ กฏตฺตา จ รูปานํ กมฺมปจฺจเยน ปจฺจโย. วิปาก\-ปจฺจเยน ปจฺจโย. เอกํเยว. อาหารปจฺจเยน ปจฺจโย. จตฺตาริ. อินฺทฺริย\-ปจฺจเยน ปจฺจโย. จตฺตาริ. ฌานปจฺจเยน ปจฺจโย. จตฺตาริ. มคฺคปจฺจเยน ปจฺจโย. จตฺตาริ. สมฺปยุตฺต\-ปจฺจเยน ปจฺจโย. เทฺว. วิปฺปยุตฺต\-ปจฺจเยน ปจฺจโย. ตีณิ. (อรูป\-ทุก\-สทิสา กาตพฺพา.) อตฺถิปจฺจเยน ปจฺจโย. สตฺต. (อรูป\-ทุก\-สทิสา กาตพฺพา. อามสนา นานาปทาเยว.) นตฺถิ\-ปจฺจเยน ปจฺจโย. ตีณิ. วิคต\-ปจฺจเยน ปจฺจโย. ตีณิ. อวิคต\-ปจฺจเยน ปจฺจโย. สตฺต.}

\csromanmarkh{1. ปจฺจยานุโลม}\csromanmarkhii{2. สงฺขฺยาวาร}\csromanheadernomark{2}{1. \textbf{ปจฺจยานุโลม 2. สงฺขฺยาวาร}}
\proseitem{323}{เหตุยา จตฺตาริ, อารมฺมเณ ตีณิ, อธิปติยา ปญฺจ, อนนฺตเร ตีณิ, สมนนฺตเร ตีณิ, สหชาเต ปญฺจ, อญฺญมญฺเญ เทฺว, นิสฺสเย สตฺต, อุปนิสฺสเย จตฺตาริ, ปุเรชาเต เทฺว, ปจฺฉาชาเต เทฺว, อาเสวเน เทฺว, กมฺเม จตฺตาริ, วิปาเก เอกํ, อาหาเร จตฺตาริ, อินฺทฺริเย จตฺตาริ, ฌาเน จตฺตาริ, มคฺเค จตฺตาริ, สมฺปยุตฺเต เทฺว, วิปฺปยุตฺเต ตีณิ, อตฺถิยา สตฺต, นตฺถิยา ตีณิ, วิคเต ตีณิ, อวิคเต สตฺต.}

\vspace{\tipitakaparskip}
\csromancenter{อนุโลมํ.}
\proseitem{324}{นิยฺยานิโก ธมฺโม นิยฺยานิกสฺส ธมฺมสฺส สหชาต\-ปจฺจเยน ปจฺจโย. อุปนิสฺสย\-ปจฺจเยน ปจฺจโย.  นิยฺยานิโก ธมฺโม อนิยฺยานิกสฺส ธมฺมสฺส อารมฺมณ\-ปจฺจเยน ปจฺจโย. สหชาต\-ปจฺจเยน ปจฺจโย. อุปนิสฺสย\-ปจฺจเยน ปจฺจโย. ปจฺฉาชาต\-ปจฺจเยน ปจฺจโย.  นิยฺยานิโก ธมฺโม นิยฺยานิกสฺส จ อนิยฺยานิกสฺส จ ธมฺมสฺส สหชาต\-ปจฺจเยน ปจฺจโย. (3)}

\proseitem{325}{อนิยฺยานิโก ธมฺโม อนิยฺยานิกสฺส ธมฺมสฺส อารมฺมณ\-ปจฺจเยน ปจฺจโย. สหชาต\-ปจฺจเยน ปจฺจโย. อุปนิสฺสย\-ปจฺจเยน ปจฺจโย. \csromanfolio{179}ปุเรชาต\-ปจฺจเยน ปจฺจโย. ปจฺฉาชาต\-ปจฺจเยน ปจฺจโย. กมฺมปจฺจเยน ปจฺจโย. อาหารปจฺจเยน ปจฺจโย. อินฺทฺริย\-ปจฺจเยน ปจฺจโย.  อนิยฺยานิโก ธมฺโม นิยฺยานิกสฺส ธมฺมสฺส อุปนิสฺสย\-ปจฺจเยน ปจฺจโย. ปุเรชาต\-ปจฺจเยน ปจฺจโย. (2)}

\prose{นิยฺยานิโก จ อนิยฺยานิโก จ ธมฺมา นิยฺยานิกสฺส ธมฺมสฺส สหชาตํ, ปุเรชาตํ.  นิยฺยานิโก จ อนิยฺยานิโก จ ธมฺมา อนิยฺยานิกสฺส ธมฺมสฺส สหชาตํ, ปจฺฉาชาตํ, อาหารํ, อินฺทฺริยํ. (2)}

\csromanmarkh{2. ปจฺจย\-ปจฺจนีย}\csromanmarkhii{2. สงฺขฺยาวาร}\csromanheadernomark{2}{2. ปจฺจย\-ปจฺจนีย 2. สงฺขฺยาวาร}
\proseitem{326}{นเหตุยา สตฺต, นอารมฺมเณ สตฺต, นอธิปติยา สตฺต, นอนนฺตเร นสมนนฺตเร สตฺต, นสหชาเต ปญฺจ, นอญฺญมญฺเญ ปญฺจ, นนิสฺสเย ปญฺจ, นอุปนิสฺสเย สตฺต, นปุเรชาเต ฉ, นปจฺฉาชาเต สตฺต ฯเปฯ นสมฺปยุตฺเต ปญฺจ, นวิปฺปยุตฺเต จตฺตาริ, โนอตฺถิยา จตฺตาริ, โนนตฺถิยา สตฺต, โนวิคเต สตฺต, โนอวิคเต จตฺตาริ.}

\proseitem{327}{เหตุปจฺจยา นอารมฺมเณ จตฺตาริ, นอธิปติยา จตฺตาริ, นอนนฺตเร นสมนนฺตเร จตฺตาริ, นอญฺญมญฺเญ เทฺว, นอุปนิสฺสเย จตฺตาริ ฯเปฯ นสมฺปยุตฺเต เทฺว, นวิปฺปยุตฺเต เทฺว, โนนตฺถิยา จตฺตาริ, โนวิคเต จตฺตาริ.}

\csromanheader{2}{4. ปจฺจย\-ปจฺจนียา\-นุโลม}
\proseitemwithcloser{328}{นเหตุปจฺจยา อารมฺมเณ ตีณิ, อธิปติยา ปญฺจ, (อนุโลมมาติกา วิตฺถาเรตพฺพา.) ฯเปฯ อวิคเต สตฺต.}{นิยฺยานิกทุกํ นิฏฺฐิตํ.}
\csromanheader{2}{98. \textbf{นิยตทุก 1-6. ปฏิจฺจาทิวาร}}
\proseitem{329}{\csromanfolio{180}นิยตํ ธมฺมํ ปฏิจฺจ นิยโต ธมฺโม อุปฺปชฺชติ เหตุปจฺจยา. นิยตํ เอกํ ขนฺธํ ปฏิจฺจ ตโย ขนฺธา ฯเปฯ เทฺว ขนฺเธ ฯเปฯ. นิยตํ ธมฺมํ ปฏิจฺจ อนิยโต ธมฺโม อุปฺปชฺชติ เหตุปจฺจยา. นิยเต ขนฺเธ ปฏิจฺจ จิตฺต\-สมุฏฺฐานํ รูปํ. (สํขิตฺตํ. ปญฺจปิ ปญฺหา กาตพฺพา, ยถา นิยฺยานิกทุกํ. เอวํ ปฏิจฺจวาโรปิ สหชาตวาโรปิ ปจฺจยวาโรปิ นิสฺสยวาโรปิ สํสฏฺฐ\-วาโรปิ สมฺปยุตฺตวาโรปิ กาตพฺพา.  นินฺนานากรณํ อามสนํ นานํ.)}

\csromanmatikaanchor{mtk.27}
\csromantocmarkref{section}{98. นิยตทุก}{mtk.27}
\bookmark[dest={mtk.27},level=1]{98. นิยตทุก}
\csromanmarkh{98. นิยตทุก}\csromanmarkhii{7. ปญฺหาวาร}\csromanheadernomark{1}{98. \textbf{นิยตทุก 7. ปญฺหาวาร}}
\proseitem{330}{นิยโต ธมฺโม นิยตสฺส ธมฺมสฺส เหตุปจฺจเยน ปจฺจโย. จตฺตาริ. (นิยฺยานิก\-ทุก\-สทิสา นินฺนานากรณา.)}

\proseitem{331}{นิยโต ธมฺโม อนิยตสฺส ธมฺมสฺส อารมฺมณ\-ปจฺจเยน ปจฺจโย. อริยา มคฺคา วุฏฺฐหิตฺวา มคฺคํ ปจฺจเวกฺขนฺติ. นิยเต ปหีเน กิเลเส ปจฺจเวกฺขนฺติ. ปุพฺเพ สมุทาจิณฺเณ ฯเปฯ. นิยเต ขนฺเธ อนิจฺจโต ฯเปฯ วิปสฺสติ. เจโตปริย\-ญาเณน นิยต\-จิตฺต\-สมงฺคิสฺส จิตฺตํ ชานาติ. นิยตา ขนฺธา เจโตปริย\-ญาณสฺส, ปุพฺเพ\-นิวาสา\-นุสฺสติ\-ญาณสฺส, ยถากมฺมูปคญาณสฺส, อนาคตํสญาณสฺส, อาวชฺชนาย อารมฺมณ\-ปจฺจเยน ปจฺจโย. (1)}

\prose{อนิยโต ธมฺโม อนิยตสฺส ธมฺมสฺส อารมฺมณ\-ปจฺจเยน ปจฺจโย. ทานํ ทตฺวา, สีลํ ฯเปฯ อุโปสถกมฺมํ ฯเปฯ. ปุพฺเพ สุจิณฺณานิ ฯเปฯ. ฌานา วุฏฺฐหิตฺวา ฌานํ ปจฺจเวกฺขติ, อสฺสาเทติ อภินนฺทติ, ตํ อารพฺภ อนิยโต ราโค อุปฺปชฺชติ. ทิฏฺฐิ ฯเปฯ วิจิกิจฺฉา ฯเปฯ อุทฺธจฺจํ ฯเปฯ. อนิยตํ โทมนสฺสํ อุปฺปชฺชติ. อริยา ผลํ ปจฺจเวกฺขนฺติ, นิพฺพานํ ปจฺจเวกฺขนฺติ. นิพฺพานํ \csromanfolio{181}โคตฺรภุสฺส, โวทานสฺส, ผลสฺส, อาวชฺชนาย อารมฺมณ\-ปจฺจเยน ปจฺจโย. อริยา อนิยเต ปหีเน กิเลเส ปจฺจเวกฺขนฺติ. วิกฺขมฺภิเต กิเลเส ฯเปฯ. ปุพฺเพ สมุทาจิณฺเณ ฯเปฯ. จกฺขุํ ฯเปฯ. วตฺถุํ อนิยเต ขนฺเธ อนิจฺจโต ฯเปฯ วิปสฺสติ, อสฺสาเทติ อภินนฺทติ, ตํ อารพฺภ อนิยโต ราโค ฯเปฯ โทมนสฺสํ อุปฺปชฺชติ. ทิพฺเพน จกฺขุนา รูปํ ปสฺสติ ฯเปฯ. อาวชฺชนาย อารมฺมณ\-ปจฺจเยน ปจฺจโย.  อนิยโต ธมฺโม นิยตสฺส ธมฺมสฺส อารมฺมณ\-ปจฺจเยน ปจฺจโย. นิพฺพานํ มคฺคสฺส อารมฺมณ\-ปจฺจเยน ปจฺจโย. รูปชีวิตินฺทฺริยํ มาตุ\-ฆาติ\-กมฺมสฺส. ปิตุฆาติ\-กมฺมสฺส. อรหนฺต\-ฆาติ\-กมฺมสฺส. รุหิรุ\-ปฺปาท\-กมฺมสฺส อารมฺมณ\-ปจฺจเยน ปจฺจโย. ยํ วตฺถุํ อามสนฺตสฺส มิจฺฉตฺต\-นิยตา ขนฺธา อุปฺปชฺชนฺติ, ตํ วตฺถุ มิจฺฉตฺต\-นิยตานํ ขนฺธานํ อารมฺมณ\-ปจฺจเยน ปจฺจโย. (2)}

\csromanheader{2}{\textbf{อธิปติ}}
\proseitem{332}{นิยโต ธมฺโม นิยตสฺส ธมฺมสฺส อธิปติ\-ปจฺจเยน ปจฺจโย. สหชาตา\-ธิปติ\csromandash{}นิยตาธิปติ สมฺปยุตฺตกานํ ขนฺธานํ อธิปติ\-ปจฺจเยน ปจฺจโย.  นิยโต ธมฺโม อนิยตสฺส ธมฺมสฺส อธิปติ\-ปจฺจเยน ปจฺจโย. อารมฺมณา\-ธิปติ, สหชาตา\-ธิปติ.  อารมฺมณา\-ธิปติ\csromandash{}อริยา มคฺคา วุฏฺฐหิตฺวา มคฺคํ ครุํ กตฺวา ปจฺจเวกฺขนฺติ. สหชาตา\-ธิปติ\csromandash{}นิยตาธิปติ จิตฺต\-สมุฏฺฐานานํ รูปานํ อธิปติ\-ปจฺจเยน ปจฺจโย.  นิยโต ธมฺโม นิยตสฺส จ อนิยตสฺส จ ธมฺมสฺส อธิปติ\-ปจฺจเยน ปจฺจโย. สหชาตา\-ธิปติ\csromandash{} นิยตาธิปติ สมฺปยุตฺตกานํ ขนฺธานํ จิตฺตสมุฏฺฐานานญฺจ รูปานํ อธิปติ\-ปจฺจเยน ปจฺจโย. (3)}

\proseitem{333}{อนิยโต ธมฺโม อนิยตสฺส ธมฺมสฺส อธิปติ\-ปจฺจเยน ปจฺจโย. อารมฺมณา\-ธิปติ, สหชาตา\-ธิปติ.  อารมฺมณา\-ธิปติ\csromandash{}ทานํ ทตฺวา, สีลํ ฯเปฯ อุโปสถกมฺมํ กตฺวา ตํ ครุํ กตฺวา ปจฺจเวกฺขติ, อสฺสาเทติ อภินนฺทติ, ตํ ครุํ กตฺวา อนิยโต ราโค อุปฺปชฺชติ. ทิฏฺฐิ อุปฺปชฺชติ. ปุพฺเพ ฯเปฯ. ฌานา ฯเปฯ. อริยา ผลํ ครุํ กตฺวา ปจฺจเวกฺขนฺติ, นิพฺพานํ ครุํ กตฺวา ปจฺจเวกฺขนฺติ. นิพฺพานํ โคตฺรภุสฺส, โวทานสฺส, ผลสฺส อธิปติ\-ปจฺจเยน ปจฺจโย. จกฺขุํ ฯเปฯ. วตฺถุํ อนิยเต ขนฺเธ ครุํ กตฺวา อสฺสาเทติ อภินนฺทติ, ตํ ครุํ กตฺวา อนิยโต ราโค อุปฺปชฺชติ. ทิฏฺฐิ อุปฺปชฺชติ. สหชาตา\-ธิปติ\csromandash{} อนิยตาธิปติ สมฺปยุตฺตกานํ ขนฺธานํ \csromanfolio{182}จิตฺตสมุฏฺฐานานญฺจ รูปานํ อธิปติ\-ปจฺจเยน ปจฺจยา.  อนิยโต ธมฺโม นิยตสฺส ธมฺมสฺส อธิปติ\-ปจฺจเยน ปจฺจโย. อารมฺมณา\-ธิปติ\csromandash{}นิพฺพานํ มคฺคสฺส อธิปติ\-ปจฺจเยน ปจฺจโย. (2)}

\proseitem{334}{นิยโต ธมฺโม อนิยตสฺส ธมฺมสฺส อนนฺตร\-ปจฺจเยน ปจฺจโย. มคฺโค ผลสฺส อนนฺตร\-ปจฺจเยน ปจฺจโย. นิยตา ขนฺธา วุฏฺฐานสฺส อนนฺตร\-ปจฺจเยน ปจฺจโย. (1)}

\prose{อนิยโต ธมฺโม อนิยตสฺส ธมฺมสฺส อนนฺตร\-ปจฺจเยน ปจฺจโย. ปุริมา ปุริมา อนิยตา ขนฺธา ปจฺฉิมานํ ปจฺฉิมานํ อนิยตานํ ขนฺธานํ อนนฺตร\-ปจฺจเยน ปจฺจโย. อนุโลมํ โคตฺรภุสฺส. อนุโลมํ โวทานสฺส. ผลํ ผลสฺส. อนุโลมํ ผลสมาปตฺติยา. นิโรธา วุฏฺฐหนฺตสฺส เนว\-สญฺญา\-นา\-สญฺญา\-ยตนํ ผลสมาปตฺติยา อนนฺตร\-ปจฺจเยน ปจฺจโย.  อนิยโต ธมฺโม นิยตสฺส ธมฺมสฺส อนนฺตร\-ปจฺจเยน ปจฺจโย. อนิยตํ โทมนสฺสํ นิยตสฺส โทมนสฺสสฺส. อนิยตา มิจฺฉาทิฏฺฐิ นิยต\-มิจฺฉาทิฏฺฐิยา อนนฺตร\-ปจฺจเยน ปจฺจโย. โคตฺรภุ มคฺคสฺส. โวทานํ มคฺคสฺส อนนฺตร\-ปจฺจเยน ปจฺจโย. (2)}

\prose{สมนนฺตร\-ปจฺจเยน ปจฺจโย. สหชาต\-ปจฺจเยน ปจฺจโย. ปญฺจ. อญฺญมญฺญ\-ปจฺจเยน ปจฺจโย. เทฺว. นิสฺสย\-ปจฺจเยน ปจฺจโย. สตฺต.}

\csromanheader{2}{\textbf{อุปนิสฺสย}}
\proseitem{335}{นิยโต ธมฺโม นิยตสฺส ธมฺมสฺส อุปนิสฺสย\-ปจฺจเยน ปจฺจโย. ปฺอกตูปนิสฺสโย\csromandash{}มาตุ\-ฆาติ\-กมฺมํ มาตุ\-ฆาติ\-กมฺมสฺส. ปิตุฆาติ\-กมฺมสฺส. อรหนฺต\-ฆาติ\-กมฺมสฺส. รุหิรุ\-ปฺปาท\-กมฺมสฺส. สํฆเภทกมฺมสฺส. นิยต\-มิจฺฉาทิฏฺฐิยา อุปนิสฺสย\-ปจฺจเยน ปจฺจโย. (จกฺกํ.) ปฐโม มคฺโค ทุติยสฺส มคฺคสฺส ฯเปฯ. ตติโย มคฺโค จตุตฺถสฺส มคฺคสฺส อุปนิสฺสย\-ปจฺจเยน ปจฺจโย.  นิยโต ธมฺโม อนิยตสฺส ธมฺมสฺส อุปนิสฺสย\-ปจฺจเยน ปจฺจโย. อารมฺมณู\-ปนิสฺสโย, อนนฺตรู\-ปนิสฺสโย, ปกตู\-ปนิสฺสโย ฯเปฯ. ปกตู\-ปนิสฺสโย\csromandash{}มาตรํ ชีวิตา โวโรเปตฺวา ฯเปฯ สํฆํ ภินฺทิตฺวา, ตสฺส ปฏิฆาต\-ตฺถาย ทานํ เทติ. สีลํ สมาทิยติ. \csromanfolio{183}อุโปสถกมฺมํ กโรติ. อริยา มคฺคํ อุปนิสฺสาย อนุปฺปนฺนํ สมาปตฺติํ อุปฺปาเทนฺติ. อุปฺปนฺนํ สมาปชฺชนฺติ ฯเปฯ. ฐานาฐาน\-โกสลฺลสฺส อุปนิสฺสย\-ปจฺจเยน ปจฺจโย. มคฺโค ผลสมาปตฺติยา อุปนิสฺสย\-ปจฺจเยน ปจฺจโย. (2)}

\proseitem{336}{อนิยโต ธมฺโม อนิยตสฺส ธมฺมสฺส อุปนิสฺสย\-ปจฺจเยน ปจฺจโย. อารมฺมณู\-ปนิสฺสโย, อนนฺตรู\-ปนิสฺสโย, ปกตู\-ปนิสฺสโย ฯเปฯ. ปกตู\-ปนิสฺสโย\csromandash{}อนิยตํ สทฺธํ อุปนิสฺสาย ทานํ เทติ. สีลํ ฯเปฯ. อุโปสถกมฺมํ กโรติ. ฌานํ อุปฺปาเทติ. วิปสฺสนํ ฯเปฯ. อภิญฺญํ ฯเปฯ. สมาปตฺติํ อุปฺปาเทติ. มานํ ชปฺเปติ. ทิฏฺฐิํ คณฺหาติ. อนิยตํ สีลํ ฯเปฯ ปญฺญํ. ราคํ ฯเปฯ. ปตฺถนํ. กายิกํ สุขํ. กายิกํ ทุกฺขํ. อุตุํ. โภชนํ. เสนาสนํ อุปนิสฺสาย ทานํ เทติ ฯเปฯ. สมาปตฺติํ อุปฺปาเทติ. ปาณํ หนติ ฯเปฯ นิคมฆาตํ กโรติ. อนิยตา สทฺธา ฯเปฯ. เสนาสนํ อนิยตาย สทฺธาย ฯเปฯ ปตฺถนาย. กายิกสฺส สุขสฺส. กายิกสฺส ทุกฺขสฺส, ผลสมาปตฺติยา อุปนิสฺสย\-ปจฺจเยน ปจฺจโย. (1)}

\prose{อนิยโต ธมฺโม นิยตสฺส ธมฺมสฺส อุปนิสฺสย\-ปจฺจเยน ปจฺจโย. อารมฺมณู\-ปนิสฺสโย, อนนฺตรู\-ปนิสฺสโย, ปกตู\-ปนิสฺสโย ฯเปฯ. ปกตู\-ปนิสฺสโย\csromandash{}อนิยตํ ราคํ อุปนิสฺสาย มาตรํ ชีวิตา โวโรเปติ ฯเปฯ สํฆํ ภินฺทติ. อนิยตํ โทสํ ฯเปฯ. เสนาสนํ อุปนิสฺสาย มาตรํ ชีวิตา โวโรเปติ ฯเปฯ สํฆํ ภินฺทติ. อนิยโต ราโค. โทสํ ฯเปฯ. เสนาสนํ มาตุ\-ฆาติ\-กมฺมสฺส ฯเปฯ สํฆเภทกมฺมสฺส อุปนิสฺสย\-ปจฺจเยน ปจฺจโย. ปฐมสฺส มคฺคสฺส ปริกมฺมํ ปฐมสฺส มคฺคสฺส ฯเปฯ. จตุตฺถสฺส มคฺคสฺส ปริกมฺมํ จตุตฺถสฺส มคฺคสฺส อุปนิสฺสย\-ปจฺจเยน ปจฺจโย. (2)}

\csromanheader{2}{\textbf{ปุเรชาตาทิ}}
\proseitem{337}{อนิยโต ธมฺโม อนิยตสฺส ธมฺมสฺส ปุเรชาต\-ปจฺจเยน ปจฺจโย. อารมฺมณ\-ปุเรชาตํ, วตฺถุ\-ปุเรชาตํ. (สํขิตฺตํ.) อนิยโต ธมฺโม นิยตสฺส ธมฺมสฺส ปุเรชาต\-ปจฺจเยน ปจฺจโย. อารมฺมณ\-ปุเรชาตํ, วตฺถุ\-ปุเรชาตํ.  อารมฺมณ\-ปุเรชาตํ\csromandash{} รูปชีวิตินฺทฺริยํ มาตุ\-ฆาติ\-กมฺมสฺส. ปิตุฆาติ\-กมฺมสฺส. อรหนฺต\-ฆาติ\-กมฺมสฺส. รุหิรุ\-ปฺปาท\-กมฺมสฺส ปุเรชาต\-ปจฺจเยน ปจฺจโย. วตฺถุ\-ปุเรชาตํ\csromandash{}วตฺถุ นิยตานํ ขนฺธานํ ปุเรชาต\-ปจฺจเยน ปจฺจโย. (2)}

\prose{ปจฺฉาชาต\-ปจฺจเยน ปจฺจโย. เทฺว. อาเสวน\-ปจฺจเยน ปจฺจโย. เทฺว.}

\csromanheader{2}{\textbf{กมฺม}}
\proseitem{338}{\csromanfolio{184}นิยโต ธมฺโม นิยตสฺส ธมฺมสฺส กมฺมปจฺจเยน ปจฺจโย. นิยตา เจตนา สมฺปยุตฺตกานํ ขนฺธานํ กมฺมปจฺจเยน ปจฺจโย.  นิยโต ธมฺโม อนิยตสฺส ธมฺมสฺส กมฺมปจฺจเยน ปจฺจโย. สหชาตา, นานากฺขณิกา.  สหชาตา นิยตา เจตนา จิตฺต\-สมุฏฺฐานานํ รูปานํ กมฺมปจฺจเยน ปจฺจโย. นานากฺขณิกา นิยตา เจตนา วิปากานํ ขนฺธานํ กฏตฺตา จ รูปานํ กมฺมปจฺจเยน ปจฺจโย. (มูลํ กาตพฺพํ.) นิยตา เจตนา สมฺปยุตฺตกานํ ขนฺธานํ จิตฺตสมุฏฺฐานานญฺจ รูปานํ กมฺมปจฺจเยน ปจฺจโย. (3)}

\prose{อนิยโต ธมฺโม อนิยตสฺส ธมฺมสฺส กมฺมปจฺจเยน ปจฺจโย. สหชาตา, นานากฺขณิกา. (สํขิตฺตํ.)}

\prose{วิปาก\-ปจฺจเยน ปจฺจโย. เอกํ. อาหารปจฺจเยน ปจฺจโย. จตฺตาริ. อินฺทฺริย\-ปจฺจเยน ปจฺจโย. จตฺตาริ. ฌานปจฺจเยน ปจฺจโย. จตฺตาริ. มคฺคปจฺจเยน ปจฺจโย. จตฺตาริ. สมฺปยุตฺต\-ปจฺจเยน ปจฺจโย. เทฺว. วิปฺปยุตฺต\-ปจฺจเยน ปจฺจโย. ตีณิ. (อรูป\-ทุก\-สทิสํ.) อตฺถิปจฺจเยน ปจฺจโย. สตฺต. (อรูปาวจร\-ทุก\-สทิสํ.) นตฺถิ\-ปจฺจเยน ปจฺจโย. วิคต\-ปจฺจเยน ปจฺจโย. อวิคต\-ปจฺจเยน ปจฺจโย. สตฺต.}

\csromanmarkh{1. ปจฺจยานุโลม}\csromanmarkhii{2. สงฺขฺยาวาร}\csromanheadernomark{2}{1. \textbf{ปจฺจยานุโลม 2. สงฺขฺยาวาร}}
\proseitem{339}{เหตุยา จตฺตาริ, อารมฺมเณ ตีณิ, อธิปติยา ปญฺจ, อนนฺตเร ตีณิ, สมนนฺตเร ตีณิ, สหชาเต ปญฺจ, อญฺญมญฺเญ เทฺว, นิสฺสเย สตฺต, อุปนิสฺสเย จตฺตาริ, ปุเรชาเต เทฺว, ปจฺฉาชาเต เทฺว, อาเสวเน เทฺว, กมฺเม จตฺตาริ, วิปาเก เอกํ, อาหาเร จตฺตาริ, อินฺทฺริเย ฌาเน จตฺตาริ, มคฺเค จตฺตาริ, สมฺปยุตฺเต เทฺว, วิปฺปยุตฺเต ตีณิ, อตฺถิยา สตฺต, นตฺถิยา ตีณิ, วิคเต ตีณิ, อวิคเต สตฺต.}

\proseitem{340}{นิยโต ธมฺโม นิยตสฺส ธมฺมสฺส สหชาต\-ปจฺจเยน ปจฺจโย. อุปนิสฺสย\-ปจฺจเยน ปจฺจโย.  นิยโต ธมฺโม อนิยตสฺส ธมฺมสฺส \csromanfolio{185}อารมฺมณ\-ปจฺจเยน ปจฺจโย. สหชาต\-ปจฺจเยน ปจฺจโย. อุปนิสฺสย\-ปจฺจเยน ปจฺจโย. ปจฺฉาชาต\-ปจฺจเยน ปจฺจโย. กมฺมปจฺจเยน ปจฺจโย.  นิยโต ธมฺโม นิยตสฺส จ อนิยตสฺส จ ธมฺมสฺส สหชาต\-ปจฺจเยน ปจฺจโย. (3)}

\proseitem{341}{อนิยโต ธมฺโม อนิยตสฺส ธมฺมสฺส อารมฺมณ\-ปจฺจเยน ปจฺจโย. สหชาต\-ปจฺจเยน ปจฺจโย. อุปนิสฺสย\-ปจฺจเยน ปจฺจโย. ปุเรชาต\-ปจฺจเยน ปจฺจโย. ปจฺฉาชาต\-ปจฺจเยน ปจฺจโย. กมฺมปจฺจเยน ปจฺจโย. อาหารปจฺจเยน ปจฺจโย. อินฺทฺริย\-ปจฺจเยน ปจฺจโย.  อนิยโต ธมฺโม นิยตสฺส ธมฺมสฺส อารมฺมณ\-ปจฺจเยน ปจฺจโย. อุปนิสฺสย\-ปจฺจเยน ปจฺจโย. ปุเรชาต\-ปจฺจเยน ปจฺจโย. (2)}

\prose{นิยโต จ อนิยโต จ ธมฺมา นิยตสฺส ธมฺมสฺส สหชาตํ, ปุเรชาตํ.  นิยโต จ อนิยโต จ ธมฺมา อนิยตสฺส ธมฺมสฺส สหชาตํ, ปจฺฉาชาตํ, อาหารํ, อินฺทฺริยํ. (2)}

\csromanmarkh{2. ปจฺจย\-ปจฺจนีย}\csromanmarkhii{2. สงฺขฺยาวาร}\csromanheadernomark{2}{2. ปจฺจย\-ปจฺจนีย 2. สงฺขฺยาวาร}
\proseitem{342}{นเหตุยา สตฺต, นอารมฺมเณ สตฺต, นอธิปติยา สตฺต, นอนนฺตเร นสมนนฺตเร สตฺต, นสหชาเต ปญฺจ, นอญฺญมญฺเญ ปญฺจ, นนิสฺสเย ปญฺจ, นอุปนิสฺสเย สตฺต, นปุเรชาเต ฉ, นปจฺฉาชาเต สตฺต ฯเปฯ นสมฺปยุตฺเต ปญฺจ, นวิปฺปยุตฺเต จตฺตาริ, โนอตฺถิยา จตฺตาริ, โนนตฺถิยา สตฺต, โนวิคเต สตฺต, โนอวิคเต จตฺตาริ.}

\proseitem{343}{เหตุปจฺจยา นอารมฺมเณ จตฺตาริ, นอธิปติยา จตฺตาริ, นอนนฺตเร นสมนนฺตเร จตฺตาริ, นอญฺญมญฺเญ เทฺว, นอุปนิสฺสเย จตฺตาริ ฯเปฯ นสมฺปยุตฺเต เทฺว, นวิปฺปยุตฺเต เทฺว, โนนตฺถิยา จตฺตาริ, โนวิคเต จตฺตาริ.}

\csromanheader{2}{4. ปจฺจย\-ปจฺจนียา\-นุโลม}
\proseitemwithcloser{344}{นเหตุปจฺจยา อารมฺมเณ ตีณิ, อธิปติยา ปญฺจ, (อนุโลมมาติกา วิตฺถาเรตพฺพา.) ฯเปฯ อวิคเต สตฺต.}{นิยตทุกํ นิฏฺฐิตํ.}
\csromanmatikaanchor{mtk.28}
\csromantocmarkref{section}{99. สอุตฺตรทุก}{mtk.28}
\bookmark[dest={mtk.28},level=1]{99. สอุตฺตรทุก}
\prose{\csromanfolio{186}99. สอุตฺตรทุก 1. ปฏิจฺจวาร}

\proseitemwithcloser{345}{สอุตฺตรํ ธมฺมํ ปฏิจฺจ สอุตฺตโร ธมฺโม อุปฺปชฺชติ เหตุปจฺจยา. สอุตฺตรํ เอกํ ขนฺธํ ปฏิจฺจ ตโย ขนฺธา จิตฺต\-สมุฏฺฐานญฺจ รูปํ ฯเปฯ เทฺว ขนฺเธ ฯเปฯ. ปฏิสนฺธิ\-กฺขเณ ฯเปฯ. (ยาว อชฺฌตฺติกา มหาภูตา. ยถา จูฬนฺตรทุเก โลกิย\-ทุก\-สทิสํ, นินฺนานากรณํ.)}{สอุตฺตรทุกํ นิฏฺฐิตํ.}
\csromanmatikaanchor{mtk.29}
\csromantocmarkref{subsection}{100. สรณทุก}{mtk.29}
\bookmark[dest={mtk.29},level=2]{100. สรณทุก}
\csromanmarkh{100. สรณทุก}\csromanmarkhii{1. ปฏิจฺจวาร}\csromanheadernomark{2}{100. \textbf{สรณทุก 1. ปฏิจฺจวาร}}
\proseitem{346}{สรณํ ธมฺมํ ปฏิจฺจ สรโณ ธมฺโม อุปฺปชฺชติ เหตุปจฺจยา. สรณํ เอกํ ขนฺธํ ปฏิจฺจ ตโย ขนฺธา ฯเปฯ เทฺว ขนฺเธ ฯเปฯ.}

\prose{(ปญฺจ ปญฺหา อรูปาวจร\-ทุก\-สทิสา.  อนุโลม\-ปฏิจฺจ\-สทิสา.)}

\prose{\textbf{เหตุ}ยา ปญฺจ, อารมฺมเณ เทฺว, อธิปติยา ปญฺจ ฯเปฯ อวิคเต ปญฺจ. อนุโลมํ.}

\prose{สรณํ ธมฺมํ ปฏิจฺจ สรโณ ธมฺโม อุปฺปชฺชติ นเหตุปจฺจยา. วิจิกิจฺฉา\-สหคเต อุทฺธจฺจ\-สหคเต ขนฺเธ ปฏิจฺจ วิจิกิจฺฉา\-สหคโต อุทฺธจฺจ\-สหคโต โมโห. (1)}

\prose{อรณํ ธมฺมํ ปฏิจฺจ อรโณ ธมฺโม อุปฺปชฺชติ นเหตุปจฺจยา. อเหตุกํ อรณํ เอกํ ขนฺธํ ปฏิจฺจ ตโย ขนฺธา จิตฺต\-สมุฏฺฐานญฺจ รูปํ ฯเปฯ เทฺว ขนฺเธ ฯเปฯ. อเหตุก\-ปฏิสนฺธิ\-กฺขเณ ฯเปฯ. (ยาว อสญฺญสตฺตา.) (1)}

\prose{นเหตุยา เทฺว, นอารมฺมเณ ตีณิ, นอธิปติยา ปญฺจ, นอนนฺตเร ตีณิ, นอุปนิสฺสเย ตีณิ, นปุเรชาเต จตฺตาริ, นปจฺฉาชาเต ปญฺจ, นอาเสวเน ปญฺจ, นกมฺเม เทฺว, นวิปาเก ปญฺจ, นอาหาเร เอกํ,}

\prose{\csromanfolio{187}ไนนฺทฺริเย เอกํ, นฌาเน เอกํ, นมคฺเค เอกํ, นสมฺปยุตฺเต ตีณิ, นวิปฺปยุตฺเต เทฺว, โนนตฺถิยา ตีณิ, โนวิคเต ตีณิ.}

\vspace{\tipitakaparskip}
\csromancenter{ปจฺจนียํ.}
\csromancenter{(เอวํ อิตเร เทฺว คณนาปิ สหชาตวาโรปิ กาตพฺโพ.)}
\csromanmarkh{100. สรณทุก}\csromanmarkhii{3. ปจฺจยวาร}\csromanheadernomark{2}{100. สรณทุก 3. ปจฺจยวาร}
\csromanheader{2}{\textbf{ปจฺจยจตุกฺก}}
\proseitem{347}{สรณํ ธมฺมํ ปจฺจยา สรโณ ธมฺโม อุปฺปชฺชติ เหตุปจฺจยา. สรณํ เอกํ ขนฺธํ ปจฺจยา ตโย ขนฺธา ฯเปฯ เทฺว ขนฺเธ ฯเปฯ.}

\vspace{\tipitakaparskip}
\csromancenter{(ยถา อรูปาวจร\-ทุกสฺส ปจฺจยวาโรปิ เอวํ กาตพฺโพ.)}
\prose{เหตุยา นว, อารมฺมเณ จตฺตาริ, อธิปติยา นว ฯเปฯ อวิคเต นว. อนุโลมํ.}

\proseitem{348}{สรณํ ธมฺมํ ปจฺจยา สรโณ ธมฺโม อุปฺปชฺชติ นเหตุปจฺจยา. วิจิกิจฺฉา\-สหคเต อุทฺธจฺจ\-สหคเต ขนฺเธ ปจฺจยา วิจิกิจฺฉา\-สหคโต อุทฺธจฺจ\-สหคโต โมโห. (1)}

\prose{อรณํ ธมฺมํ ปจฺจยา อรโณ ธมฺโม อุปฺปชฺชติ นเหตุปจฺจยา. อเหตุกํ อรณํ เอกํ ขนฺธํ ปจฺจยา ตโย ขนฺธา จิตฺต\-สมุฏฺฐานญฺจ รูปํ ฯเปฯ เทฺว ขนฺเธ ฯเปฯ. อเหตุก\-ปฏิสนฺธิ\-กฺขเณ ฯเปฯ. (ยาว อสญฺญสตฺตา.) จกฺขายตนํ ปจฺจยา จกฺขุวิญฺญาณํ ฯเปฯ. กายายตนํ ปจฺจยา กายวิญฺญาณํ. วตฺถุํ ปจฺจยา อเหตุกา ขนฺธา.  อรณํ ธมฺมํ ปจฺจยา สรโณ ธมฺโม อุปฺปชฺชติ นเหตุปจฺจยา. วตฺถุํ ปจฺจยา วิจิกิจฺฉา\-สหคโต อุทฺธจฺจ\-สหคโต โมโห. (2)}

\prose{สรณญฺจ อรณญฺจ ธมฺมํ ปจฺจยา สรโณ ธมฺโม อุปฺปชฺชติ นเหตุปจฺจยา. วิจิกิจฺฉา\-สหคเต อุทฺธจฺจ\-สหคเต ขนฺเธ จ วตฺถุญฺจ ปจฺจยา วิจิกิจฺฉา\-สหคโต อุทฺธจฺจ\-สหคโต โมโห. (1)}

\prose{\csromanfolio{188}นเหตุยา จตฺตาริ, นอารมฺมเณ ตีณิ, นอธิปติยา นว, นอนนฺตเร ตีณิ, นสมนนฺตเร ตีณิ, นอญฺญมญฺเญ ตีณิ, นอุปนิสฺสเย ตีณิ, นปุเรชาเต จตฺตาริ, นปจฺฉาชาเต นว, นอาเสวเน นว, นกมฺเม จตฺตาริ, นวิปาเก นว, นอาหาเร เอกํ, นอินฺทฺริเย เอกํ, นฌาเน เอกํ, นมคฺเค เอกํ, นสมฺปยุตฺเต ตีณิ, นวิปฺปยุตฺเต เทฺว, โนนตฺถิยา ตีณิ, โนวิคเต ตีณิ. ปจฺจนียํ.}

\prose{(เอวํ อิตเร เทฺว คณนาปิ นิสฺสยวาโรปิ กาตพฺโพ, สํสฏฺฐ\-วาเรปิ เทฺว ปญฺหา กาตพฺพา. สพฺพตฺถ.)}

\prose{เหตุยา เทฺว, อารมฺมเณ เทฺว, (สพฺพตฺถ เทฺว.) วิปาเก เอกํ, อวิคเต เทฺว. (อนุโลมํ.)}

\prose{นเหตุยา เทฺว, นอธิปติยา เทฺว, นปุเรชาเต เทฺว, นปจฺฉาชาเต เทฺว, นอาเสวเน เทฺว, นกมฺเม เทฺว, นวิปาเก เทฺว, นฌาเน เอกํ, นมคฺเค เอกํ, นวิปฺปยุตฺเต เทฺว. (ปจฺจนียํ.)}

\vspace{\tipitakaparskip}
\csromancenter{(เอวํ อิตเร เทฺว คณนาปิ สมฺปยุตฺตวาโรปิ กาตพฺโพ.)}
\csromanmarkh{100. สรณทุก}\csromanmarkhii{7. ปญฺหาวาร}\csromanheadernomark{2}{100. \textbf{สรณทุก 7. ปญฺหาวาร}}
\proseitem{349}{สรโณ ธมฺโม สรณสฺส ธมฺมสฺส เหตุปจฺจเยน ปจฺจโย. (อรูป\-ทุก\-สทิสํ, จตฺตาริ.)}

\proseitem{350}{สรโณ ธมฺโม สรณสฺส ธมฺมสฺส อารมฺมณ\-ปจฺจเยน ปจฺจโย. ราคํ อสฺสาเทติ อภินนฺทติ, ตํ อรพฺภ ราโค อุปฺปชฺชติ. ทิฏฺฐิ ฯเปฯ วิจิกิจฺฉา ฯเปฯ อุทฺธจฺจํ ฯเปฯ โทมนสฺสํ อุปฺปชฺชติ. ทิฏฺฐิํ อสฺสาเทติ อภินนฺทติ, ตํ อารพฺภ ราโค อุปฺปชฺชติ ฯเปฯ โทมนสฺสํ อุปฺปชฺชติ. วิจิกิจฺฉํ \csromanfolio{189}อารพฺภ ฯเปฯ อุทฺธจฺจํ อารพฺภ ฯเปฯ. โทมนสฺสํ อารพฺภ โทมนสฺสํ อุปฺปชฺชติ. ทิฏฺฐิ ฯเปฯ วิจิกิจฺฉา ฯเปฯ อุทฺธจฺจํ อุปฺปชฺชติ.  สรโณ ธมฺโม อรณสฺส ธมฺมสฺส อรมฺมณปจฺจเยน ปจฺจโย. อริยา ปหีเน กิเลเส ฯเปฯ. วิกฺขมฺภิเต กิเลเส ฯเปฯ. ปุพฺเพ สมุทาจิณฺเณ ฯเปฯ. สรเณ ขนฺเธ อนิจฺจโต ฯเปฯ วิปสฺสติ. เจโตปริย\-ญาเณน สรณ\-จิตฺต\-สมงฺคิสฺส จิตฺตํ ชานาติ. สรณา ขนฺธา เจโตปริย\-ญาณสฺส, ปุพฺเพ\-นิวาสา\-นุสฺสติ\-ญาณสฺส, ยถากมฺมูปคญาณสฺส, อนาคตํสญาณสฺส อาวชฺชนาย อารมฺมณ\-ปจฺจเยน ปจฺจโย. (2)}

\proseitem{351}{อรโณ ธมฺโม อรณสฺส ธมฺมสฺส อารมฺมณ\-ปจฺจเยน ปจฺจโย. ทานํ ฯเปฯ. สีลํ ฯเปฯ. อุโปสถกมฺมํ กตฺวา ตํ ปจฺจเวกฺขติ. ปุพฺเพ ฯเปฯ. ฌานา ฯเปฯ. อริยา มคฺคา วุฏฺฐหิตฺวา มคฺคํ ปจฺจเวกฺขนฺติ. ผลํ ฯเปฯ. นิพฺพานํ ปจฺจเวกฺขนฺติ. นิพฺพานํ โคตฺรภุสฺส, โวทานสฺส, มคฺคสฺส, ผลสฺส, อาวชฺชนาย อารมฺมณ\-ปจฺจเยน ปจฺจโย. จกฺขุํ ฯเปฯ. วตฺถุํ อรเณ ขนฺเธ อนิจฺจโต ฯเปฯ วิปสฺสติ. ทิพฺเพน จกฺขุนา รูปํ ปสฺสติ ฯเปฯ. อนาคตํสญาณสฺส, อาวชฺชนาย อารมฺมณ\-ปจฺจเยน ปจฺจโย.  อรโณ ธมฺโม สรณสฺส ธมฺมสฺส อารมฺมณ\-ปจฺจเยน ปจฺจโย. ทานํ ฯเปฯ. สีลํ ฯเปฯ. อุโปสถกมฺมํ กตฺวา ตํ อสฺสาเทติ อภินนฺทติ, ตํ อารพฺภ ราโค อุปฺปชฺชติ ฯเปฯ. ปุพฺเพ ฯเปฯ. ฌานา ฯเปฯ. จกฺขุํ ฯเปฯ. วตฺถุํ. อรเณ ขนฺเธ อสฺสาเทติ อภินนฺทติ, ตํ อารพฺภ ราโค อุปฺปชฺชติ ฯเปฯ โทมนสฺสํ อุปฺปชฺชติ. (2)}

\csromanheader{2}{\textbf{อธิปติ}}
\proseitem{352}{สรโณ ธมฺโม สรณสฺส ธมฺมสฺส อธิปติ\-ปจฺจเยน ปจฺจโย. อารมฺมณา\-ธิปติ, สหชาตา\-ธิปติ.  อารมฺมณา\-ธิปติ\csromandash{}ราคํ ครุํ กตฺวา อสฺสาเทติ อภินนฺทติ, ตํ ครุํ กตฺวา ราโค อุปฺปชฺชติ. ทิฏฺฐิ อุปฺปชฺชติ. ทิฏฺฐิํ ครุํ กตฺวา อสฺสาเทติ อภินนฺทติ. สหชาตา\-ธิปติ\csromandash{}สรณาธิปติ สมฺปยุตฺตกานํ ขนฺธานํ อธิปติ\-ปจฺจเยน ปจฺจโย.  สรโณ ธมฺโม อรณสฺส ธมฺมสฺส อธิปติ\-ปจฺจเยน ปจฺจโย. สหชาตา\-ธิปติ\csromandash{}สรณาธิปติ จิตฺต\-สมุฏฺฐานานํ รูปานํ อธิปติ\-ปจฺจเยน ปจฺจโย.  สรโณ ธมฺโม สรณสฺส จ อรณสฺส จ ธมฺมสฺส อธิปติ\-ปจฺจเยน ปจฺจโย. สหชาตา\-ธิปติ\csromandash{}สรณาธิปติ สมฺปยุตฺตกานํ ขนฺธานํ จิตฺตสมุฏฺฐานานญฺจ รูปานํ อธิปติ\-ปจฺจเยน ปจฺจโย. (3)}

\proseitem{353}{\csromanfolio{190}อรโณ ธมฺโม อรณสฺส ธมฺมสฺส อธิปติ\-ปจฺจเยน ปจฺจโย. อารมฺมณ\-ธิปติ, สหชาตา\-ธิปติ.  อารมฺมณา\-ธิปติ\csromandash{}ทานํ ฯเปฯ สีลํ ฯเปฯ อุโปสถกมฺมํ กตฺวา ตํ ครุํ กตฺวา ปจฺจเวกฺขติ. ปุพฺเพ ฯเปฯ. ฌานา ฯเปฯ. อริยา มคฺคา วุฏฺฐหิตฺวา มคฺคํ ครุํ กตฺวา ปจฺจเวกฺขนฺติ. ผลํ ฯเปฯ. นิพฺพานํ ครุํ กตฺวา ปจฺจเวกฺขนฺติ. นิพฺพานํ โคตฺรภุสฺส, โวทานสฺส, มคฺคสฺส, ผลสฺส อธิปติ\-ปจฺจเยน ปจฺจโย. สหชาตา\-ธิปติ\csromandash{}อรณาธิปติ สมฺปยุตฺตกานํ ขนฺธานํ จิตฺตสมุฏฺฐานานญฺจ รูปานํ อธิปติ\-ปจฺจเยน ปจฺจโย.  อรโณ ธมฺโม สรณสฺส ธมฺมสฺส อธิปติ\-ปจฺจเยน ปจฺจโย. อารมฺมณา\-ธิปติ\csromandash{}ทานํ ฯเปฯ สีลํ ฯเปฯ อุโปสถกมฺมํ กตฺวา ตํ ครุํ กตฺวา อสฺสาเทติ อภินนฺทติ, ตํ ครุํ กตฺวา ราโค อุปฺปชฺชติ. ทิฏฺฐิ อุปฺปชฺชติ. ปุพฺเพ สุจิณฺณานิ ฯเปฯ. ฌานา ฯเปฯ. จกฺขุํ ฯเปฯ. วตฺถุํ อรเณ ขนฺเธ ครุํ กตฺวา อสฺสาเทติ อภินนฺทติ, ตํ ครุํ กตฺวา ราโค อุปฺปชฺชติ. ทิฏฺฐิ อุปฺปชฺชติ. (2)}

\proseitem{354}{สรโณ ธมฺโม สรณสฺส ธมฺมสฺส อนนฺตร\-ปจฺจเยน ปจฺจโย. ปุริมา ปุริมา สรณา ขนฺธา ปจฺฉิมานํ ปจฺฉิมานํ สรณานํ ขนฺธานํ อนนฺตร\-ปจฺจเยน ปจฺจโย.  สรโณ ธมฺโม อรณสฺส ธมฺมสฺส อนนฺตร\-ปจฺจเยน ปจฺจโย. สรณา ขนฺธา วุฏฺฐานสฺส อนนฺตร\-ปจฺจเยน ปจฺจโย. (2)}

\proseitem{355}{อรโณ ธมฺโม อรณสฺส ธมฺมสฺส อนนฺตร\-ปจฺจเยน ปจฺจโย. ปุริมา ปุริมา อรณา ขนฺธา ปจฺฉิมานํ ปจฺฉิมานํ อรณานํ ขนฺธานํ อนนฺตร\-ปจฺจเยน ปจฺจโย. อนุโลมํ โคตฺรภุสฺส ฯเปฯ ผลสมาปตฺติยา อนนฺตร\-ปจฺจเยน ปจฺจโย.  อรโณ ธมฺโม สรณสฺส ธมฺมสฺส อนนฺตร\-ปจฺจเยน ปจฺจโย. อาวชฺชนา สรณานํ ขนฺธานํ อนนฺตร\-ปจฺจเยน ปจฺจโย. (2)}

\prose{สมนนฺตร\-ปจฺจเยน ปจฺจโย. สหชาต\-ปจฺจเยน ปจฺจโย. ปญฺจ. อญฺญมญฺญ\-ปจฺจเยน ปจฺจโย. เทฺว. นิสฺสย\-ปจฺจเยน ปจฺจโย. สตฺต.}

\csromanheader{2}{\textbf{อุปนิสฺสย}}
\proseitem{356}{สรโณ ธมฺโม สรณสฺส ธมฺมสฺส อุปนิสฺสย\-ปจฺจเยน ปจฺจโย. อารมฺมณู\-ปนิสฺสโย, อนนฺตรู\-ปนิสฺสโย, ปกตู\-ปนิสฺสโย ฯเปฯ. \csromanfolio{191}ปกตู\-ปนิสฺสโย ราคํ อุปนิสฺสาย ปาณํ หนติ ฯเปฯ สํฆํ ภินฺทติ. โทสํ ฯเปฯ. ปตฺถนํ อุปนิสฺสาย ปาณํ หนติ ฯเปฯ สํฆํ ภินฺทติ. ราโค ฯเปฯ. ปตฺถนา ราคสฺส ฯเปฯ ปตฺถนาย อุปนิสฺสย\-ปจฺจเยน ปจฺจโย.  สรโณ ธมฺโม อรณสฺส ธมฺมสฺส อุปนิสฺสย\-ปจฺจเยน ปจฺจโย. อนนฺตรู\-ปนิสฺสโย, ปกตู\-ปนิสฺสโย ฯเปฯ. อกตูปนิสฺสโย\csromandash{}ราคํ อุปนิสฺสาย ทานํ เทติ. สีลํ ฯเปฯ อุโปสถกมฺมํ กโรติ. ฌานํ อุปฺปาเทติ. วิปสฺสนํ ฯเปฯ. มคฺคํ ฯเปฯ. อภิญฺญํ ฯเปฯ. สมาปตฺติํ อุปฺปาเทติ. โทสํ ฯเปฯ. ปตฺถนํ อุปนิสฺสาย ทานํ เทติ ฯเปฯ. สมาปตฺติํ อุปฺปาเทติ. ราโค ฯเปฯ. ปตฺถนา สทฺธาย ฯเปฯ ปญฺญาย. กายิกสฺส สุขสฺส. กายิกสฺส ทุกฺขสฺส. มคฺคสฺส. ผลสมาปตฺติยา อุปนิสฺสย\-ปจฺจเยน ปจฺจโย. (2)}

\proseitem{357}{อรโณ ธมฺโม อรณสฺส ธมฺมสฺส อุปนิสฺสย\-ปจฺจเยน ปจฺจโย. อารมฺมณู\-ปนิสฺสโย, อนนฺตรู\-ปนิสฺสโย, ปกตู\-ปนิสฺสโย ฯเปฯ. ปกตู\-ปนิสฺสโย\csromandash{}สทฺธํ อุปนิสฺสาย ทานํ เทติ ฯเปฯ. สมาปตฺติํ อุปฺปาเทติ. สีลํ ฯเปฯ. ปญฺญํ. กายิกํ สุขํ. กายิกํ ทุกฺขํ. อุตุํ. โภชนํ. เสนาสนํ อุปนิสฺสาย ทานํ เทติ ฯเปฯ. สมาปตฺติํ อุปฺปาเทติ. สทฺธา ฯเปฯ. เสนาสนํ สทฺธาย ฯเปฯ ปญฺญาย. กายิกสฺส สุขสฺส. กายิกสฺส ทุกฺขสฺส. มคฺคสฺส. ผลสมาปตฺติยา อุปนิสฺสย\-ปจฺจเยน ปจฺจโย. (1)}

\prose{อรโณ ธมฺโม สรณสฺส ธมฺมสฺส อุปนิสฺสย\-ปจฺจเยน ปจฺจโย. อารมฺมณู\-ปนิสฺสโย, อนนฺตรู\-ปนิสฺสโย, ปกตู\-ปนิสฺสโย ฯเปฯ. ปกตู\-ปนิสฺสโย\csromandash{}สทฺธํ อุปนิสฺสาย มานํ ชปฺเปติ. ทิฏฺฐิํ คณฺหาติ. สีลํ ฯเปฯ. เสนาสนํ อุปนิสฺสาย ปาณํ หนติ ฯเปฯ สํฆํ ภินฺทติ. สทฺธา ฯเปฯ. เสนาสนํ ราคสฺส ฯเปฯ ปตฺถนาย อุปนิสฺสย\-ปจฺจเยน ปจฺจโย. (2)}

\csromanheader{2}{\textbf{ปุเรชาตาทิ}}
\proseitem{358}{อรโณ ธมฺโม อรณสฺส ธมฺมสฺส ปุเรชาต\-ปจฺจเยน ปจฺจโย. อารมฺมณ\-ปุเรชาตํ, วตฺถุ\-ปุเรชาตํ.  อารมฺมณ\-ปุเรชาตํ\csromandash{}จกฺขุํ ฯเปฯ. วตฺถุํ อนิจฺจโต ฯเปฯ วิปสฺสติ ฯเปฯ. ทิพฺเพน จกฺขุนา รูปํ ปสฺสติ. ทิพฺพาย โสตธาตุยา สทฺทํ สุณาติ. รูปายตนํ จกฺขุ\-วิญฺญาณสฺส ฯเปฯ. โผฏฺฐพฺพา\-ยตนํ กายวิญฺญาณสฺส ฯเปฯ. วตฺถุ\-ปุเรชาตํ\csromandash{}จกฺขายตนํ \csromanfolio{192}จกฺขุ\-วิญฺญาณสฺส ฯเปฯ. กายายตนํ กายวิญฺญาณสฺส ฯเปฯ. วตฺถุ อรณานํ ขนฺธานํ ปุเรชาต\-ปจฺจเยน ปจฺจโย.  อรโณ ธมฺโม สรณสฺส ธมฺมสฺส ปุเรชาต\-ปจฺจเยน ปจฺจโย. อารมฺมณ\-ปุเรชาตํ, วตฺถุ\-ปุเรชาตํ.  อารมฺมณ\-ปุเรชาตํ\csromandash{}จกฺขุํ ฯเปฯ. วตฺถุํ อสฺสาเทติ อภินนฺทติ, ตํ อารพฺภ ราโค อุปฺปชฺชติ ฯเปฯ โทมนสฺสํ อุปฺปชฺชติ. วตฺถุ\-ปุเรชาตํ\csromandash{}วตฺถุ สรณานํ ขนฺธานํ ปุเรชาต\-ปจฺจเยน ปจฺจโย. (2)}

\prose{ปจฺฉาชาต\-ปจฺจเยน ปจฺจโย. เทฺว. อาเสวน\-ปจฺจเยน ปจฺจโย. เทฺว.}

\csromanheader{2}{\textbf{กมฺม}}
\proseitem{359}{สรโณ ธมฺโม สรณสฺส ธมฺมสฺส กมฺมปจฺจเยน ปจฺจโย. สรณา เจตนา สมฺปยุตฺตกานํ ขนฺธานํ กมฺมปจฺจเยน ปจฺจโย. (มูลํ กาตพฺพํ.) สหชาตา, นานากฺขณิกา.  สหชาตา สรณา เจตนา จิตฺต\-สมุฏฺฐานานํ รูปานํ กมฺมปจฺจเยน ปจฺจโย. นานากฺขณิกา สรณา เจตนา วิปากานํ ขนฺธานํ กฏตฺตา จ รูปานํ กมฺมปจฺจเยน ปจฺจโย. (มูลํ กาตพฺพํ.) สรณา เจตนา สมฺปยุตฺตกานํ ขนฺธานํ จิตฺตสมุฏฺฐานานญฺจ รูปานํ กมฺมปจฺจเยน ปจฺจโย. (3)}

\prose{อรโณ ธมฺโม อรณสฺส ธมฺมสฺส กมฺมปจฺจเยน ปจฺจโย. สหชาตา, นานากฺขณิกา.  สหชาตา\csromandash{}(สํขิตฺตํ.) (1)}

\prose{วิปาก\-ปจฺจเยน ปจฺจโย. เอกํ. อาหารปจฺจเยน ปจฺจโย. จตฺตาริ. อินฺทฺริย\-ปจฺจเยน ปจฺจโย. จตฺตาริ. ฌานปจฺจเยน ปจฺจโย. จตฺตาริ. มคฺคปจฺจเยน ปจฺจโย. จตฺตาริ. สมฺปยุตฺต\-ปจฺจเยน ปจฺจโย. เทฺว. วิปฺปยุตฺต\-ปจฺจเยน ปจฺจโย. ตีณิ. (อรูป\-ทุก\-สทิสา.)}

\csromanheader{2}{\textbf{อตฺถิ}}
\proseitem{360}{สรโณ ธมฺโม สรณสฺส ธมฺมสฺส อตฺถิปจฺจเยน ปจฺจโย. (สํขิตฺตํ.). สรโณ ธมฺโม อรณสฺส ธมฺมสฺส อตฺถิปจฺจเยน ปจฺจโย. สหชาตํ, ปจฺฉาชาตํ. (สํขิตฺตํ.). สรโณ ธมฺโม สรณสฺส จ อรณสฺส จ ธมฺมสฺส อตฺถิปจฺจเยน ปจฺจโย. (สํขิตฺตํ.) (3)}

\proseitem{361}{อรโณ ธมฺโม อรณสฺส ธมฺมสฺส อตฺถิปจฺจเยน ปจฺจโย. สหชาตํ, ปุเรชาตํ, ปจฺฉาชาตํ, อาหารํ, อินฺทฺริยํ. (สํขิตฺตํ.). อรโณ ธมฺโม สรณสฺส ธมฺมสฺส อตฺถิปจฺจเยน ปจฺจโย. ปุเรชาตํ \csromanfolio{193}จกฺขุํ ฯเปฯ. วตฺถุํ อสฺสาเทติ อภินนฺทติ, ตํ อารพฺภ ราโค อุปฺปชฺชติ ฯเปฯ โทมนสฺสํ อุปฺปชฺชติ. วตฺถุ สรณานํ ขนฺธานํ อตฺถิปจฺจเยน ปจฺจโย. (2)}

\prose{สรโณ จ อรโณ จ ธมฺมา สรณสฺส ธมฺมสฺส อตฺถิปจฺจเยน ปจฺจโย. สหชาตํ, ปุเรชาตํ. (สํขิตฺตํ.). สรโณ จ อรโณ จ ธมฺมา อรณสฺส ธมฺมสฺส อตฺถิปจฺจเยน ปจฺจโย. สหชาตํ, ปจฺฉาชาตํ, อาหารํ, อินฺทฺริยํ.  สหชาตา สรณา ขนฺธา จ มหาภูตา จ จิตฺต\-สมุฏฺฐานานํ รูปานํ อตฺถิปจฺจเยน ปจฺจโย. ปจฺฉาชาตา สรณา ขนฺธา จ กพฬีกาโร อาหาโร จ อิมสฺส กายสฺส อตฺถิปจฺจเยน ปจฺจโย. ปจฺฉาชาตา สรณา ขนฺธา จ รูปชีวิตินฺทฺริยญฺจ กฏตฺตารูปานํ อตฺถิปจฺจเยน ปจฺจโย. (2)}

\csromanmarkh{1. ปจฺจยานุโลม}\csromanmarkhii{2. สงฺขฺยาวาร}\csromanheadernomark{2}{1. \textbf{ปจฺจยานุโลม 2. สงฺขฺยาวาร}}
\proseitem{362}{เหตุยา จตฺตาริ, อารมฺมเณ จตฺตาริ, อธิปติยา ปญฺจ, อนนฺตเร จตฺตาริ, สมนนฺตเร จตฺตาริ, สหชาเต ปญฺจ, อญฺญมญฺเญ เทฺว, นิสฺสเย สตฺต, อุปนิสฺสเย จตฺตาริ, ปุเรชาเต เทฺว, ปจฺฉาชาเต เทฺว, อาเสวเน เทฺว, กมฺเม จตฺตาริ, วิปาเก เอกํ, อาหาเร จาตฺตาริ, อินฺทฺริเย จตฺตาริ, ฌาเน จตฺตาริ, มคฺเค จตฺตาริ, สมฺปยุตฺเต เทฺว, วิปฺปยุตฺเต ตีณิ, อตฺถิยา สตฺต, นตฺถิยา จตฺตาริ, วิคเต จตฺตาริ, อวิคเต สตฺต.}

\proseitem{363}{สรโณ ธมฺโม สรณสฺส ธมฺมสฺส อารมฺมณ\-ปจฺจเยน ปจฺจโย. สหชาต\-ปจฺจเยน ปจฺจโย. อุปนิสฺสย\-ปจฺจเยน ปจฺจโย.  สรโณ ธมฺโม อรณสฺส ธมฺมสฺส อารมฺมณ\-ปจฺจเยน ปจฺจโย. สหชาต\-ปจฺจเยน ปจฺจโย. อุปนิสฺสย\-ปจฺจเยน ปจฺจโย. ปจฺฉาชาต\-ปจฺจเยน ปจฺจโย. กมฺมปจฺจเยน ปจฺจโย.  สรโณ ธมฺโม สรณสฺส จ อรณสฺส จ ธมฺมสฺส สหชาต\-ปจฺจเยน ปจฺจโย. (3)}

\proseitem{364}{อรโณ ธมฺโม อรณสฺส ธมฺมสฺส อารมฺมณ\-ปจฺจเยน ปจฺจโย. สหชาต\-ปจฺจเยน ปจฺจโย. อุปนิสฺสย\-ปจฺจเยน ปจฺจโย. ปุเรชาต\-ปจฺจเยน ปจฺจโย. ปจฺฉาชาต\-ปจฺจเยน ปจฺจโย. กมฺมปจฺจเยน ปจฺจโย. \csromanfolio{194}อาหารปจฺจเยน ปจฺจโย. อินฺทฺริย\-ปจฺจเยน ปจฺจโย.  อรโณ ธมฺโม สรณสฺส ธมฺมสฺส อารมฺมณ\-ปจฺจเยน ปจฺจโย. อุปนิสฺสย\-ปจฺจเยน ปจฺจโย. ปุเรชาต\-ปจฺจเยน ปจฺจโย. (2)}

\prose{สรโณ จ อรโณ จ ธมฺมา สรณสฺส ธมฺมสฺส สหชาตํ, ปุเรชาตํ.  สรโณ จ อรโณ จ ธมฺมา อรณสฺส ธมฺมสฺส สหชาตํ, ปจฺฉาชาตํ, อาหารํ, อินฺทฺริยํ. (2)}

\csromanmarkh{2. ปจฺจย\-ปจฺจนีย}\csromanmarkhii{2. สงฺขฺยาวาร}\csromanheadernomark{2}{2. ปจฺจย\-ปจฺจนีย 2. สงฺขฺยาวาร}
\proseitem{365}{นเหตุยา สตฺต, นสมนนฺตเร สตฺต, นสหชาเต ปญฺจ, นอญฺญมญฺเญ ปญฺจ, นนิสฺสเย ปญฺจ, นอุปนิสฺสเย สตฺต, นปุเรชาเต ฉ, นปจฺฉาชาเต สตฺต ฯเปฯ นสมฺปยุตฺเต ปญฺจ, นวิปฺปยุตฺเต จตฺตาริ, โนอตฺถิยา จตฺตาริ, โนนตฺถิยา สตฺต, โนวิคเต สตฺต, โนอวิคเต จตฺตาริ.}

\proseitem{366}{เหตุปจฺจยา นอารมฺมเณ จตฺตาริ, นอธิปติยา จตฺตาริ ฯเปฯ นสมนนฺตเร จตฺตาริ, นอญฺญมญฺเญ เทฺว, นอุปนิสฺสเย จตฺตาริ ฯเปฯ นสมฺปยุตฺเต เทฺว, นวิปฺปยุตฺเต เทฺว, โนนตฺถิยา จตฺตาริ, โนวิคเต จตฺตาริ.}

\csromanheader{2}{4. ปจฺจย\-ปจฺจนียา\-นุโลม}
\proseitemwithcloser{367}{นเหตุปจฺจยา อารมฺมเณ จตฺตาริ, อธิปติยา ปญฺจ, อนนฺตเร จตฺตาริ, (อนุโลมมาติกา คเหตพฺพา.) ฯเปฯ อวิคเต สตฺต.}{สรณทุกํ นิฏฺฐิตํ.}
\nitthitam{ปิฏฺฐิทุกํ นิฏฺฐิตํ.}
\nitthitam{\textbf{ธมฺมานุโลเม ทุกปฏฺฐานํ นิฏฺฐิตํ.}}
\csromanpalirecto
\pitaka{อภิธมฺม\-ปิฏก}
\csromanheader{2}{\textbf{ธมฺมานุโลม}}
\csromanmatikaanchor{mtk.30}
\csromantocmarknopage{chapter}{ทุกติกปฏฺฐานปาฬิ}{mtk.30}
\bookmark[dest={mtk.30},level=0]{ทุกติกปฏฺฐานปาฬิ}
\gambhira{ทุกติก\-ปฏฺฐาน\-ปาฬิ}
\namakkaram{นโม ตสฺส ภควโต อรหโต สมฺมา\-สมฺพุทฺธสฺส.}
\csromanmatikaanchor{mtk.31}
\csromanmatikaanchor{mtk.32}
\csromanmatikaanchor{mtk.92}
\csromantocmarknopage{chapter}{1. เหตุทุก}{mtk.31}
\bookmark[dest={mtk.31},level=0]{1. เหตุทุก}
\csromantocmarkref{section}{1. กุสลตฺติก}{mtk.32}
\bookmark[dest={mtk.32},level=1]{1. กุสลตฺติก}
\csromanmarkh{1. เหตุทุก}\csromanmarkhii{1. กุสลตฺติก}\csromanheadernomark{1}{1. \textbf{เหตุทุก 1. กุสลตฺติก}}
\csromanmarkh{1. กุสลปท}\csromanmarkhii{1. ปฏิจฺจวาร}\csromanheadernomark{2}{1. \textbf{กุสลปท 1. ปฏิจฺจวาร}}
\prose{\csromanfolio{195}ปจฺจยจตุกฺกเหตุ}

\proseitem{1}{เหตุํ กุสลํ ธมฺมํ ปฏิจฺจ เหตุ กุสโล ธมฺโม อุปฺปชฺชติ เหตุปจฺจยา.  เหตุํ กุสลํ ธมฺมํ ปฏิจฺจ นเหตุ กุสโล ธมฺโม อุปฺปชฺชติ เหตุปจฺจยา.  เหตุํ กุสลํ ธมฺมํ ปฏิจฺจ เหตุ กุสโล จ นเหตุ กุสโล จ ธมฺมา อุปฺปชฺชนฺติ เหตุปจฺจยา. (3)}

\proseitem{2}{นเหตุํ กุสลํ ธมฺมํ ปฏิจฺจ นเหตุ กุสโล ธมฺโม อุปฺปชฺชติ เหตุปจฺจยา.  นเหตุํ กุสลํ ธมฺมํ ปฏิจฺจ เหตุ กุสโล ธมฺโม อุปฺปชฺชติ เหตุปจฺจยา.  นเหตุํ กุสลํ ธมฺมํ ปฏิจฺจ เหตุ กุสโล จ นเหตุ กุสโล จ ธมฺมา อุปฺปชฺชนฺติ เหตุปจฺจยา. (3)}

\proseitem{3}{เหตุํ กุสลญฺจ นเหตุํ กุสลญฺจ ธมฺมํ ปฏิจฺจ เหตุ กุสโล ธมฺโม อุปฺปชฺชติ เหตุปจฺจยา.  เหตุํ กุสลญฺจ นเหตุํ กุสลญฺจ ธมฺมํ ปฏิจฺจ นเหตุ กุสโล ธมฺโม อุปฺปชฺชติ เหตุปจฺจยา.  เหตุํ กุสลญฺจ นเหตุํ กุสลญฺจ ธมฺมํ ปฏิจฺจ เหตุ กุสโล จ นเหตุ กุสโล จ ธมฺมา อุปฺปชฺชนฺติ เหตุปจฺจยา. (3)}

\csromanheader{2}{ธมฺมานุโลม ทุกติก\-ปฏฺฐาน\-ปาฬิ อารมฺมณ}
\proseitem{4}{\csromanfolio{196}เหตุํ กุสลํ ธมฺมํ ปฏิจฺจ เหตุ กุสโล ธมฺโม อุปฺปชฺชติ อารมฺมณ\-ปจฺจยา. (สํขิตฺตํ.)}

\proseitem{5}{เหตุยา นว, อารมฺมเณ นว, อธิปติยา นว, อนนฺตเร นว, สมนนฺตเร นว, สหชาเต นว, อญฺญมญฺเญ นว, นิสฺสเย นว, อุปนิสฺสเย นว, ปุเรชาเต นว, อาเสวเน นว, กมฺเม นว, อาหาเร นว, อินฺทฺริเย นว, ฌาเน นว, มคฺเค นว, สมฺปยุตฺเต นว, วิปฺปยุตฺเต นว, อตฺถิยา นว, นตฺถิยา นว, วิคเต นว, อวิคเต นว. (อนุโลมํ.)}

\prose{นอธิปติ}

\proseitem{6}{เหตุํ กุสลํ ธมฺมํ ปฏิจฺจ เหตุ กุสโล ธมฺโม อุปฺปชฺชติ นอธิปติ\-ปจฺจยา.  เหตุํ กุสลํ ธมฺมํ ปฏิจฺจ นเหตุ กุสโล ธมฺโม อุปฺปชฺชติ นอธิปติ\-ปจฺจยา.  เหตุํ กุสลํ ธมฺมํ ปฏิจฺจ เหตุ กุสโล จ นเหตุ กุสโล จ ธมฺมา อุปฺปชฺชนฺติ นอธิปติ\-ปจฺจยา. (3)}

\proseitem{7}{นเหตุํ กุสลํ ธมฺมํ ปฏิจฺจ นเหตุ กุสโล ธมฺโม อุปฺปชฺชติ นอธิปติ\-ปจฺจยา. ตีณิ.}

\prose{เหตุํ กุสลญฺจ นเหตุํ กุสลญฺจ ธมฺมํ ปฏิจฺจ เหตุ กุสโล ธมฺโม อุปฺปชฺชติ นอธิปติ\-ปจฺจยา. ตีณิ.}

\csromanheader{2}{\textbf{นปุเรชาตาทิ}}
\proseitem{8}{เหตุํ กุสลํ ธมฺมํ ปฏิจฺจ เหตุ กุสโล ธมฺโม อุปฺปชฺชติ นปุเรชาต\-ปจฺจยา. นว. นปจฺฉาชาต\-ปจฺจยา. นว. นอาเสวน\-ปจฺจยา. นว.}

\csromanheader{2}{\textbf{นกมฺม}}
\proseitem{9}{เหตุํ กุสลํ ธมฺมํ ปฏิจฺจ นเหตุ กุสโล ธมฺโม อุปฺปชฺชติ นกมฺมปจฺจยา. (1)}

\prose{นเหตุํ กุสลํ ธมฺมํ ปฏิจฺจ นเหตุ กุสโล ธมฺโม อุปฺปชฺชติ นกมฺมปจฺจยา. (1)}

\csromanmarkh{1. เหตุทุก}\csromanmarkhii{1. กุสลตฺติก}\csromanheadernomark{2}{1. เหตุทุก 1. กุสลตฺติก}
\prose{\csromanfolio{197}เหตุํ กุสลญฺจ นเหตุํ กุสลญฺจ ธมฺมํ ปฏิจฺจ นเหตุ กุสโล ธมฺโม อุปฺปชฺชติ นกมฺมปจฺจยา. (1)}

\csromanheader{2}{\textbf{นวิปากาทิ}}
\proseitem{10}{เหตุํ กุสลํ ธมฺมํ ปฏิจฺจ เหตุ กุสโล ธมฺโม อุปฺปชฺชติ นวิปาก\-ปจฺจยา. นว ฯเปฯ. นวิปฺปยุตฺต\-ปจฺจยา. นว.}

\proseitem{11}{นอธิปติยา นว, นปุเรชาเต นว, นปจฺฉาชาเต นว, นอาเสวเน นว, นกมฺเม ตีณิ, นวิปาเก นว, นวิปฺปยุตฺเต นว. (ปจฺจนียํ.)}

\prose{เหตุปจฺจยา นอธิปติยา นว. (สํขิตฺตํ. อนุโลม ปจฺจนียํ.)}

\prose{นอธิปติ\-ปจฺจยา เหตุยา นว. อารมฺมเณ นว. (สํขิตฺตํ. ปจฺจนียา\-นุโลมํ.)}

\prose{(สหชาตวารมฺปิ ปจฺจยวารมฺปิ นิสฺสยวารมฺปิ สํสฏฺฐ\-วารมฺปิ สมฺปยุตฺตวารมฺปิ ปฏิจฺจ\-วาร\-สทิสา วิตฺถาเรตพฺพา.)}

\csromanmarkh{1. กุสลปท}\csromanmarkhii{7. ปญฺหาวาร}\csromanheadernomark{2}{1. \textbf{กุสลปท 7. ปญฺหาวาร}}
\prose{ปจฺจยจตุกฺกเหตุ}

\proseitem{12}{เหตุ กุสโล ธมฺโม เหตุสฺส กุสลสฺส ธมฺมสฺส เหตุปจฺจเยน ปจฺจโย.  เหตุ กุสโล ธมฺโม นเหตุสฺส กุสลสฺส ธมฺมสฺส เหตุปจฺจเยน ปจฺจโย.  เหตุ กุสโล ธมฺโม เหตุสฺส กุสลสฺส จ นเหตุสฺส กุสลสฺส จ ธมฺมสฺส เหตุปจฺจเยน ปจฺจโย. (3)}

\csromanheader{2}{\textbf{อารมฺมณาทิ}}
\proseitem{13}{เหตุ กุสโล ธมฺโม เหตุสฺส กุสลสฺส ธมฺมสฺส อารมฺมณ\-ปจฺจเยน ปจฺจโย. ตีณิ.}

\prose{กุสโล ธมฺโม นเหตุสฺส กุสลสฺส ธมฺมสฺส อารมฺมณ\-ปจฺจเยน ปจฺจโย. ตีณิ.}

\csromanpalirecto
\gambhira{ธมฺมานุโลม ทุกติก\-ปฏฺฐาน\-ปาฬิ}
\prose{\csromanfolio{198}เหตุ กุสโล จ นเหตุ กุสโล จ ธมฺมา เหตุสฺส กุสลสฺส ธมฺมสฺส อารมฺมณ\-ปจฺจเยน ปจฺจโย. ตีณิ.}

\proseitem{14}{เหตุ กุสโล ธมฺโม เหตุสฺส กุสลสฺส ธมฺมสฺส อธิปติ\-ปจฺจเยน ปจฺจโย. อารมฺมณา\-ธิปติ, สหชาตา\-ธิปติ. ตีณิ.}

\prose{กุสโล ธมฺโม นเหตุสฺส กุสลสฺส ธมฺมสฺส อธิปติ\-ปจฺจเยน ปจฺจโย. อารมฺมณา\-ธิปติ, สหชาตา\-ธิปติ. ตีณิ.}

\prose{เหตุ กุสโล จ นเหตุ กุสโล จ ธมฺมา เหตุสฺส กุสลสฺส ธมฺมสฺส อธิปติ\-ปจฺจเยน ปจฺจโย. อารมฺมณา\-ธิปติ. ตีณิ.}

\proseitem{15}{เหตุ กุสโล ธมฺโม เหตุสฺส กุสลสฺส ธมฺมสฺส อนนฺตร\-ปจฺจเยน ปจฺจโย. สมนนฺตร\-ปจฺจเยน ปจฺจโย. สหชาต\-ปจฺจเยน ปจฺจโย. อญฺญมญฺญ\-ปจฺจเยน ปจฺจโย. นิสฺสย\-ปจฺจเยน ปจฺจโย.}

\proseitem{14}{เหตุ กุสโล ธมฺโม เหตุสฺส กุสลสฺส ธมฺมสฺส อุปนิสฺสย\-ปจฺจเยน ปจฺจโย. อารมฺมณู\-ปนิสฺสโย, อนนฺตรู\-ปนิสฺสโย, ปกตู\-ปนิสฺสโย. ตีณิ.}

\prose{กุสโล ธมฺโม นเหตุสฺส กุสลสฺส ธมฺมสฺส อุปนิสฺสย\-ปจฺจเยน ปจฺจโย. อารมฺมณู\-ปนิสฺสโย, อนนฺตรู\-ปนิสฺสโย, ปกตู\-ปนิสฺสโย. ตีณิ.}

\prose{เหตุ กุสโล จ นเหตุ กุสโล จ ธมฺมา เหตุสฺส กุสลสฺส ธมฺมสฺส อุปนิสฺสย\-ปจฺจเยน ปจฺจโย. อารมฺมณู\-ปนิสฺสโย, อนนฺตรู\-ปนิสฺสโย, ปกตู\-ปนิสฺสโย. ตีณิ. อาเสวน\-ปจฺจเยน ปจฺจโย. นว.}

\proseitem{17}{กุสโล ธมฺโม นเหตุสฺส กุสลสฺส ธมฺมสฺส กมฺมปจฺจเยน ปจฺจโย.  นเหตุ กุสโล ธมฺโม เหตุสฺส กุสลสฺส ธมฺมสฺส กมฺมปจฺจเยน ปจฺจโย.  นเหตุ กุสโล ธมฺโม เหตุสฺส กุสลสฺส จ นเหตุสฺส กุสลสฺส จ ธมฺมสฺส กมฺมปจฺจเยน ปจฺจโย. (3)}

\csromanmarkh{1. เหตุทุก}\csromanmarkhii{1. กุสลตฺติก อาหาราทิ}\csromanheadernomark{2}{1. เหตุทุก 1. กุสลตฺติก \textbf{อาหาราทิ}}
\proseitem{18}{\csromanfolio{199}กุสโล ธมฺโม นเหตุสฺส กุสลสฺส ธมฺมสฺส อาหารปจฺจเยน ปจฺจโย. ตีณิ.}

\proseitem{19}{เหตุ กุสโล ธมฺโม เหตุสฺส กุสลสฺส ธมฺมสฺส อินฺทฺริย\-ปจฺจเยน ปจฺจโย. นว.}

\proseitem{20}{กุสโล ธมฺโม นเหตุสฺส กุสลสฺส ธมฺมสฺส ฌานปจฺจเยน ปจฺจโย. ตีณิ.}

\proseitem{21}{เหตุ กุสโล ธมฺโม นเหตุสฺส กุสลสฺส ธมฺมสฺส มคฺคปจฺจเยน ปจฺจโย. นว.}

\proseitem{22}{เหตุยา ตีณิ, อารมฺมเณ นว, อธิปติยา นว, อนนฺตเร นว, สมนนฺตเร นว, สหชาเต นว, อญฺญมญฺเญ นว, นิสฺสเย นว, อุปนิสฺสเย นว, อาเสวเน นว, กมฺเม ตีณิ, อาหาเร ตีณิ, อินฺทฺริเย นว, ฌาเน ตีณิ, มคฺเค นว, สมฺปยุตฺเต นว, อตฺถิยา นว, นตฺถิยา นว, วิคเต นว, อวิคเต นว.}

\proseitem{23}{เหตุ กุสโล ธมฺโม เหตุสฺส กุสลสฺส ธมฺมสฺส อารมฺมณ\-ปจฺจเยน ปจฺจโย. สหชาต\-ปจฺจเยน ปจฺจโย. อุปนิสฺสย\-ปจฺจเยน ปจฺจโย. ตีณิ. (3)}

\proseitem{24}{กุสโล ธมฺโม นเหตุสฺส กุสลสฺส ธมฺมสฺส อารมฺมณ\-ปจฺจเยน ปจฺจโย. สหชาต\-ปจฺจเยน ปจฺจโย. อุปนิสฺสย\-ปจฺจเยน ปจฺจโย.  นเหตุ กุสโล ธมฺโม เหตุสฺส กุสลสฺส ธมฺมสฺส อารมฺมณ\-ปจฺจเยน ปจฺจโย. สหชาต\-ปจฺจเยน ปจฺจโย. อุปนิสฺสย\-ปจฺจเยน ปจฺจโย.  นเหตุ กุสโล ธมฺโม เหตุสฺส กุสลสฺส จ นเหตุสฺส กุสลสฺส จ ธมฺมสฺส อารมฺมณ\-ปจฺจเยน ปจฺจโย. สหชาต\-ปจฺจเยน ปจฺจโย. อุปนิสฺสย\-ปจฺจเยน ปจฺจโย. (3)}

\proseitem{25}{\csromanfolio{200}เหตุ กุสโล จ นเหตุ กุสโล จ ธมฺมา เหตุสฺส กุสลสฺส ธมฺมสฺส อารมฺมณ\-ปจฺจเยน ปจฺจโย. สหชาต\-ปจฺจเยน ปจฺจโย. อุปนิสฺสย\-ปจฺจเยน ปจฺจโย. ตีณิ. (3)}

\proseitem{26}{นเหตุยา นว, นอารมฺมเณ นว, นอธิปติยา นว ฯเปฯ โนอวิคเต นว. (สํขิตฺตํ. ปจฺจนียํ.)}

\prose{เหตุปจฺจยา นอารมฺมเณ ตีณิ ฯเปฯ โนนตฺถิยา ตีณิ, โนวิคเต ตีณิ. (สํขิตฺตํ. อนุโลม\-ปจฺจนียํ.)}

\prose{นเหตุปจฺจยา อารมฺมเณ นว ฯเปฯ อวิคเต นว. (สํขิตฺตํ. ปจฺจนียา\-นุโลมํ.)}

\prose{ยถา กุสลตฺติเก ปญฺหาวารสฺส อนุโลมมฺปิ ปจฺจนียมฺปิ อนุโลมปจฺจนียมฺปิ ปจฺจนียานุโลมมฺปิ คณิตํ, เอวํ คเณตพฺพํ.)}

\csromanmarkh{2. อกุสลปท}\csromanmarkhii{1. ปฏิจฺจวาร}\csromanheadernomark{2}{2. \textbf{อกุสลปท 1. ปฏิจฺจวาร}}
\prose{ปจฺจยจตุกฺกเหตุ}

\proseitem{27}{เหตุํ อกุสลํ ธมฺมํ ปฏิจฺจ เหตุ อกุสโล ธมฺโม อุปฺปชฺชติ เหตุปจฺจยา.  เหตุํ อกุสลํ ธมฺมํ ปฏิจฺจ นเหตุ อกุสโล ธมฺโม อุปฺปชฺชติ เหตุปจฺจยา.  เหตุํ อกุสลํ ธมฺมํ ปฏิจฺจ เหตุ อกุสโล จ นเหตุ อกุสโล จ ธมฺมา อุปฺปชฺชนฺติ เหตุปจฺจยา. (3)}

\proseitem{28}{นเหตุํ อกุสลํ ธมฺมํ ปฏิจฺจ นเหตุ อกุสโล ธมฺโม อุปฺปชฺชติ เหตุปจฺจยา.  นเหตุํ อกุสลํ ธมฺมํ ปฏิจฺจ เหตุ อกุสโล ธมฺโม อุปฺปชฺชติ เหตุปจฺจยา.  นเหตุํ อกุสลํ ธมฺมํ ปฏิจฺจ เหตุ อกุสโล จ นเหตุ อกุสโล จ ธมฺมา อุปฺปชฺชนฺติ เหตุปจฺจยา. (3)}

\proseitem{29}{เหตุํ อกุสลญฺจ นเหตุํ อกุสลญฺจ ธมฺมํ ปฏิจฺจ เหตุ อกุสโล ธมฺโม อุปฺปชฺชติ เหตุปจฺจยา.  เหตุํ อกุสลญฺจ นเหตุํ อกุสลญฺจ ธมฺมํ ปฏิจฺจ นเหตุ อกุสโล ธมฺโม อุปฺปชฺชติ เหตุปจฺจยา.  เหตุํ อกุสลญฺจ นเหตุํ อกุสลญฺจ ธมฺมํ ปฏิจฺจ เหตุ อกุสโล จ นเหตุ อกุสโล จ ธมฺมา อุปฺปชฺชนฺติ เหตุปจฺจยา. (3)}

\csromanmarkh{1. เหตุทุก}\csromanmarkhii{1. กุสลตฺติก}\csromanheadernomark{2}{1. เหตุทุก 1. กุสลตฺติก}
\proseitem{30}{\csromanfolio{201}เหตุยา นว, อารมฺมเณ นว, อธิปติยา นว ฯเปฯ กมฺเม นว, อาหาเร นว ฯเปฯ อวิคเต นว. (สํขิตฺตํ. อนุโลมํ.)}

\proseitem{31}{\textbf{นเหตุ}ํ อกุสลํ ธมฺมํ ปฏิจฺจ เหตุ อกุสโล ธมฺโม อุปฺปชฺชติ นเหตุปจฺจยา. (1)}

\prose{นอธิปตฺยาทิ}

\proseitem{32}{เหตุํ อกุสลํ ธมฺมํ ปฏิจฺจ เหตุ อกุสโล ธมฺโม อุปฺปชฺชติ นอธิปติ\-ปจฺจยา. นว ฯเปฯ.}

\prose{เหตุํ อกุสลํ ธมฺมํ ปฏิจฺจ นเหตุ อกุสโล ธมฺโม อุปฺปชฺชติ นกมฺมปจฺจยา. (1)}

\prose{\textbf{นเหตุ}ํ อกุสลํ ธมฺมํ ปฏิจฺจ นเหตุ อกุสโล ธมฺโม อุปฺปชฺชติ นกมฺมปจฺจยา. (1)}

\prose{เหตุํ อกุสลญฺจ นเหตุํ อกุสลญฺจ ธมฺมํ ปฏิจฺจ นเหตุ อกุสโล ธมฺโม อุปฺปชฺชติ นกมฺมปจฺจยา. (1) (สํขิตฺตํ.)}

\vspace{\tipitakaparskip}
\csromancenter{\textbf{นเหตุ}ยา เอกํ, นอธิปติยา นว, นปุเรชาเต นว, นปจฺฉาชาเต นว, นอาเสวเน นว, นกมฺเม ตีณิ, นวิปาเก นว, นวิปฺปยุตฺเต นว. (ปจฺจนียํ.)}
\prose{เหตุปจฺจยา นอธิปติยา นว, (สํขิตฺตํ. อนุโลม\-ปจฺจนียํ.)}

\prose{นเหตุปจฺจยา อารมฺมเณ เอกํ. (สํขิตฺตํ. ปจฺจนียา\-นุโลมํ.)}

\prose{(สหชาตวารมฺปิ ปจฺจยวารมฺปิ นิสฺสยวารมฺปิ สํสฏฺฐ\-วารมฺปิ สมฺปยุตฺตวารมฺปิ ปฏิจฺจ\-วาร\-สทิสา วิตฺถาเรตพฺพา.)}

\csromanmarkh{2. อกุสลปท}\csromanmarkhii{7. ปญฺหาวาร}\csromanheadernomark{2}{2. \textbf{อกุสลปท 7. ปญฺหาวาร}}
\prose{ปจฺจยจตุกฺกเหตุ}

\proseitem{34}{เหตุ อกุสโล ธมฺโม เหตุสฺส อกุสลสฺส ธมฺมสฺส เหตุปจฺจเยน ปจฺจโย.  เหตุ อกุสโล ธมฺโม นเหตุสฺส \csromanfolio{202}อกุสลสฺส ธมฺมสฺส เหตุปจฺจเยน ปจฺจโย.  เหตุ อกุสโล ธมฺโม เหตุสฺส อกุสลสฺส จ นเหตุสฺส อกุสลสฺส จ ธมฺมสฺส เหตุปจฺจเยน ปจฺจโย. (3)}

\csromanheader{2}{\textbf{อารมฺมณาทิ}}
\proseitem{35}{เหตุ อกุสโล ธมฺโม เหตุสฺส อกุสลสฺส ธมฺมสฺส อารมฺมณ\-ปจฺจเยน ปจฺจโย. ตีณิ.}

\prose{อกุสโล ธมฺโม นเหตุสฺส อกุสลสฺส ธมฺมสฺส อารมฺมณ\-ปจฺจเยน ปจฺจโย. ตีณิ.}

\prose{เหตุ อกุสโล จ นเหตุ อกุสโล จ ธมฺมา เหตุสฺส อกุสลสฺส ธมฺมสฺส อารมฺมณ\-ปจฺจเยน ปจฺจโย. ตีณิ.}

\proseitem{36}{เหตุ อกุสโล ธมฺโม เหตุสฺส อกุสลสฺส ธมฺมสฺส อธิปติ\-ปจฺจเยน ปจฺจโย. อารมฺมณา\-ธิปติ. ตีณิ.}

\prose{อกุสโล ธมฺโม นเหตุสฺส อกุสลสฺส ธมฺมสฺส อธิปติ\-ปจฺจเยน ปจฺจโย. อารมฺมณา\-ธิปติ, สหชาตา\-ธิปติ. ตีณิ.}

\prose{เหตุ อกุสโล จ นเหตุ อกุสโล จ ธมฺมา เหตุสฺส อกุสลสฺส ธมฺมสฺส อธิปติ\-ปจฺจเยน ปจฺจโย. อารมฺมณา\-ธิปติ. ตีณิ ฯเปฯ.}

\proseitem{37}{อกุสโล ธมฺโม นเหตุสฺส อกุสลสฺส ธมฺมสฺส กมฺมปจฺจเยน ปจฺจโย.  นเหตุ อกุสโล ธมฺโม เหตุสฺส อกุสลสฺส ธมฺมสฺส กมฺมปจฺจเยน ปจฺจโย.  นเหตุ อกุสโล ธมฺโม เหตุสฺส อกุสลสฺส จ นเหตุสฺส อกุสลสฺส จ ธมฺมสฺส กมฺมปจฺจเยน ปจฺจโย. (3)}

\csromanheader{2}{\textbf{อาหาราทิ}}
\proseitem{38}{อกุสโล ธมฺโม นเหตุสฺส อกุสลสฺส ธมฺมสฺส อาหารปจฺจเยน ปจฺจโย. ตีณิ.}

\proseitem{39}{อกุสโล ธมฺโม นเหตุสฺส อกุสลสฺส ธมฺมสฺส อินฺทฺริย\-ปจฺจเยน ปจฺจโย. ตีณิ.}

\csromanmarkh{1. เหตุทุก}\csromanmarkhii{1. กุสลตฺติก}\csromanheadernomark{2}{1. เหตุทุก 1. กุสลตฺติก}
\proseitem{40}{\csromanfolio{203}อกุสโล ธมฺโม นเหตุสฺส อกุสลสฺส ธมฺมสฺส ฌานปจฺจเยน ปจฺจโย. ตีณิ. (สํขิตฺตํ.)}

\proseitem{41}{เหตุยา ตีณิ, อารมฺมเณ นว, อธิปติยา นว, อนนฺตเร นว, สมนนฺตเร นว, สหชาเต นว, อญฺญมญฺเญ นว, นิสฺสเย นว, อุปนิสฺสเย นว, อาเสวเน นว, กมฺเม ตีณิ, อาหาเร ตีณิ, อินฺทฺริเย ตีณิ, ฌาเน ตีณิ, มคฺเค ตีณิ, สมฺปยุตฺเต นว, อตฺถิยา นว, นตฺถิยา นว, วิคเต นว, อวิคเต นว.}

\proseitem{42}{เหตุ อกุสโล ธมฺโม เหตุสฺส อกุสลสฺส ธมฺมสฺส อารมฺมณ\-ปจฺจเยน ปจฺจโย. สหชาต\-ปจฺจเยน ปจฺจโย. อุปนิสฺสย\-ปจฺจเยน ปจฺจโย. (สํขิตฺตํ.)}

\proseitem{43}{นเหตุยา นว, นอารมฺมเณ นว, นอธิปติยา นว, นอนนฺตเร นว, นสมนนฺตเร นว, นสหชาเต นว, นอญฺญมญฺเญ นว ฯเปฯ โนอวิคเต นว. (ปจฺจนียํ.)}

\prose{เหตุปจฺจยา นอารมฺมเณ ตีณิ. (สํขิตฺตํ. อนุโลม\-ปจฺจนียํ.)}

\prose{นเหตุปจฺจยา อารมฺมเณ นว. (สํขิตฺตํ. ปจฺจนียา\-นุโลมํ.)}

\prose{(ยถา กุสลตฺติเก ปญฺหาวารสฺส อนุโลมมฺปิ ปจฺจนียมฺปิ อนุโลมปจฺจนียมฺปิ ปจฺจนียานุโลมมฺปิ คณิตํ, เอวํ คเณตพฺพํ.)}

\csromanmarkh{3. อพฺยากตปท}\csromanmarkhii{1. ปฏิจฺจวาร}\csromanheadernomark{2}{3. \textbf{อพฺยากตปท 1. ปฏิจฺจวาร}}
\prose{ปจฺจยจตุกฺกเหตฺวาทิ}

\proseitem{44}{เหตุํ อพฺยากตํ ธมฺมํ ปฏิจฺจ เหตุ อพฺยากโต ธมฺโม อุปฺปชฺชติ เหตุปจฺจยา.  เหตุํ อพฺยากตํ ธมฺมํ ปฏิจฺจ นเหตุ อพฺยากโต ธมฺโม อุปฺปชฺชติ เหตุปจฺจยา.  เหตุํ อพฺยากตํ ธมฺมํ ปฏิจฺจ เหตุ อพฺยากโต จ นเหตุ อพฺยากโต จ ธมฺมา อุปฺปชฺชนฺติ เหตุปจฺจยา. (3)}

\proseitem{45}{\csromanfolio{204}นเหตุํ อพฺยากตํ ธมฺมํ ปฏิจฺจ นเหตุ อพฺยากโต ธมฺโม อุปฺปชฺชติ เหตุปจฺจยา. ตีณิ.}

\prose{เหตุํ อพฺยากตญฺจ นเหตุํ อพฺยากตญฺจ ธมฺมํ ปฏิจฺจ เหตุ อพฺยากโต ธมฺโม อุปฺปชฺชติ เหตุปจฺจยา. ตีณิ.}

\proseitem{46}{เหตุํ อพฺยากตํ ธมฺมํ ปฏิจฺจ เหตุ อพฺยากโต ธมฺโม อุปฺปชฺชติ อารมฺมณ\-ปจฺจยา. (สํขิตฺตํ.)}

\proseitem{47}{เหตุยา นว, อารมฺมเณ นว, อธิปติยา นว, อนนฺตเร นว, สมนนฺตเร นว, สหชาเต นว, อญฺญมญฺเญ นว, นิสฺสเย นว, อุปนิสฺสเย นว, ปุเรชาเต นว, อาเสวเน นว, กมฺเม นว, วิปาเก นว, อาหาเร นว, อินฺทฺริเย นว, ฌาเน นว, มคฺเค นว, สมฺปยุตฺเต นว, วิปฺปยุตฺเต นว, อตฺถิยา นว, นตฺถิยา นว, วิคเต นว, อวิคเต นว. (อนุโลมํ.)}

\prose{นเหตุ นอารมฺมณ}

\proseitem{48}{นเหตุํ อพฺยากตํ ธมฺมํ ปฏิจฺจ นเหตุ อพฺยากโต ธมฺโม อุปฺปชฺชติ นเหตุปจฺจยา. (1)}

\proseitem{49}{เหตุํ อพฺยากตํ ธมฺมํ ปฏิจฺจ นเหตุ อพฺยากโต ธมฺโม อุปฺปชฺชติ นอารมฺมณ\-ปจฺจยา. (1)}

\prose{นเหตุํ อพฺยากตํ ธมฺมํ ปฏิจฺจ นเหตุ อพฺยากโต ธมฺโม อุปฺปชฺชติ นอารมฺมณ\-ปจฺจยา. (1)}

\prose{เหตุํ อพฺยากตญฺจ นเหตุํ อพฺยากตญฺจ ธมฺมํ ปฏิจฺจ นเหตุ อพฺยากโต ธมฺโม อุปฺปชฺชติ นอารมฺมณ\-ปจฺจยา. (1)}

\prose{นอธิปตฺยาทิ}

\proseitem{50}{เหตุํ อพฺยากตํ ธมฺมํ ปฏิจฺจ เหตุ อพฺยากโต ธมฺโม อุปฺปชฺชติ นอธิปติ\-ปจฺจยา. นว ฯเปฯ.}

\proseitem{51}{เหตุํ อพฺยากตํ ธมฺมํ ปฏิจฺจ นเหตุ อพฺยากโต ธมฺโม อุปฺปชฺชติ นกมฺมปจฺจยา. (1)}

\csromanmarkh{1. เหตุทุก}\csromanmarkhii{1. กุสลตฺติก}\csromanheadernomark{2}{1. เหตุทุก 1. กุสลตฺติก}
\prose{\csromanfolio{205}นเหตุํ อพฺยากตํ ธมฺมํ ปฏิจฺจ นเหตุ อพฺยากโต ธมฺโม อุปฺปชฺชติ นกมฺมปจฺจยา. (1)}

\prose{เหตุํ อพฺยากตญฺจ นเหตุํ อพฺยากตญฺจ ธมฺมํ ปฏิจฺจ นเหตุ อพฺยากโต ธมฺโม อุปฺปชฺชติ นกมฺมปจฺจยา. (1)}

\proseitem{52}{นเหตุํ อพฺยากตํ ธมฺมํ ปฏิจฺจ นเหตุ อพฺยากโต ธมฺโม อุปฺปชฺชติ นอาหารปจฺจยา. นอินฺทฺริยปจฺจยา. นฌานปจฺจยา. (สํขิตฺตํ.)}

\proseitem{53}{นหุตุยา เอกํ, นอารมฺมเณ ตีณิ, นอธิปติยา นว, นอนนฺตเร ตีณิ, นสมนนฺตเร ตีณิ, นอญฺญมญฺเญ ตีณิ. นอุปนิสฺสเย ตีณิ, นปุเรชาเต นว, นปจฺฉาชาเต นว, นอาเสวเน นว, นกมฺเม ตีณิ, นวิปาเก นว, นอาหาเร เอกํ, นอินฺทฺริเย เอกํ, นฌาเน เอกํ, นมคฺเค เอกํ, นสมฺปยุตฺเต ตีณิ, นวิปฺปยุตฺเต นว, โนนตฺถิยา ตีณิ, โนวิคเต ตีณิ. (ปจฺจนียํ.)}

\prose{เหตุปจฺจยา นอารมฺมเณ ตีณิ, นอธิปติยา นว. (สํขิตฺตํ. อนุโลม\-ปจฺจนียํ.)}

\prose{นเหตุปจฺจยา อารมฺมเณ เอกํ. (สํขิตฺตํ. ปจฺจนียา\-นุโลมํ.)}

\prose{(สหชาตวาโรปิ ปจฺจยวาโรปิ นิสฺสยวาโรปิ สํสฏฺฐ\-วาโรปิ สมฺปยุตฺตวาโรปิ ปฏิจฺจ\-วาร\-สทิสา วิตฺถาเรตพฺพา.)}

\csromanmarkh{3. อพฺยากตปท}\csromanmarkhii{7. ปญฺหาวาร}\csromanheadernomark{2}{3. \textbf{อพฺยากตปท 7. ปญฺหาวาร}}
\prose{ปจฺจยจตุกฺกเหตุ}

\proseitem{54}{เหตุ อพฺยากโต ธมฺโม เหตุสฺส อพฺยากตสฺส ธมฺมสฺส เหตุปจฺจเยน ปจฺจโย.  เหตุ อพฺยากโต ธมฺโม นเหตุสฺส อพฺยากตสฺส ธมฺมสฺส เหตุปจฺจเยน ปจฺจโย.  เหตุ อพฺยากโต ธมฺโม เหตุสฺส อพฺยากตสฺส จ นเหตุสฺส อพฺยากตสฺส จ ธมฺมสฺส เหตุปจฺจเยน ปจฺจโย. (3)}

\csromanheader{2}{ธมฺมานุโลม ทุกติก\-ปฏฺฐาน\-ปาฬิ อารมฺมณาทิ}
\proseitem{55}{\csromanfolio{206}เหตุ อพฺยากโต ธมฺโม เหตุสฺส อพฺยากตสฺส ธมฺมสฺส อารมฺมณ\-ปจฺจเยน ปจฺจโย. นว.}

\proseitem{56}{เหตุ อพฺยากโต ธมฺโม เหตุสฺส อพฺยากตสฺส ธมฺมสฺส อธิปติ\-ปจฺจเยน ปจฺจโย. อารมฺมณา\-ธิปติ, สหชาตา\-ธิปติ. ตีณิ.}

\prose{อพฺยากโต ธมฺโม นเหตุสฺส อพฺยากตสฺส ธมฺมสฺส อธิปติ\-ปจฺจเยน ปจฺจโย. อารมฺมณา\-ธิปติ, สหชาตา\-ธิปติ. ตีณิ.}

\prose{เหตุ อพฺยากโต จ นเหตุ อพฺยากโต จ ธมฺมา เหตุสฺส อพฺยากตสฺส ธมฺมสฺส อธิปติ\-ปจฺจเยน ปจฺจโย. อารมฺมณา\-ธิปติ. ตีณิ ฯเปฯ.}

\csromanheader{2}{\textbf{ปุเรชาตาทิ}}
\proseitem{57}{อพฺยากโต ธมฺโม นเหตุสฺส อพฺยากตสฺส ธมฺมสฺส ปุเรชาต\-ปจฺจเยน ปจฺจโย. ตีณิ.}

\proseitem{58}{เหตุ อพฺยากโต ธมฺโม นเหตุสฺส อพฺยากตสฺส ธมฺมสฺส ปจฺฉาชาต\-ปจฺจเยน ปจฺจโย. (1)}

\prose{อพฺยากโต ธมฺโม นเหตุสฺส อพฺยากตสฺส ธมฺมสฺส ปจฺฉาชาต\-ปจฺจเยน ปจฺจโย. (1)}

\prose{เหตุ อพฺยากโต จ นเหตุ อพฺยากโต จ ธมฺมา นเหตุสฺส อพฺยากตสฺส ธมฺมสฺส ปจฺฉาชาต\-ปจฺจเยน ปจฺจโย. (1)}

\proseitem{59}{อพฺยากโต ธมฺโม นเหตุสฺส อพฺยากตสฺส ธมฺมสฺส กมฺมปจฺจเยน ปจฺจโย. ตีณิ.}

\proseitem{60}{เหตุ อพฺยากโต ธมฺโม เหตุสฺส อพฺยากตสฺส ธมฺมสฺส วิปาก\-ปจฺจเยน ปจฺจโย. นว.}

\proseitem{61}{อพฺยากโต ธมฺโม นเหตุสฺส อพฺยากตสฺส ธมฺมสฺส อาหารปจฺจเยน ปจฺจโย. อินฺทฺริย\-ปจฺจเยน ปจฺจโย. ฌานปจฺจเยน ปจฺจโย. มคฺคปจฺจเยน ปจฺจโย. สมฺปยุตฺต\-ปจฺจเยน ปจฺจโย.}

\csromanmarkh{1. เหตุทุก}\csromanmarkhii{1. กุสลตฺติก วิปฺปยุตฺต}\csromanheadernomark{2}{1. เหตุทุก 1. กุสลตฺติก \textbf{วิปฺปยุตฺต}}
\proseitem{62}{\csromanfolio{207}เหตุ อพฺยากโต ธมฺโม นเหตุสฺส อพฺยากตสฺส ธมฺมสฺส วิปฺปยุตฺต\-ปจฺจเยน ปจฺจโย. (1)}

\prose{อพฺยากโต ธมฺโม นเหตุสฺส อพฺยากตสฺส ธมฺมสฺส วิปฺปยุตฺต\-ปจฺจเยน ปจฺจโย. ตีณิ.}

\prose{เหตุ อพฺยากโต จ นเหตุ อพฺยากโต จ ธมฺมา นเหตุสฺส อพฺยากตสฺส ธมฺมสฺส วิปฺปยุตฺต\-ปจฺจเยน ปจฺจโย. (1) (สํขิตฺตํ.)}

\proseitem{63}{เหตุยา ตีณิ, อารมฺมเณ นว, อธิปติยา นว, อนนฺตเร นว, สมนนฺตเร นว, สหชาเต นว, อญฺญมญฺเญ นว, นิสฺสเย นว, อุปนิสฺสเย นว, ปุเรชาเต ตีณิ, ปจฺฉาชาเต ตีณิ, อาเสวเน นว, กมฺเม ตีณิ, วิปาเก นว, อาหาเร ตีณิ, อินฺทฺริเย นว, ฌาเน ตีณิ, มคฺเค นว, สมฺปยุตฺเต นว, วิปฺปยุตฺเต ปญฺจ, อตฺถิยา นว, นตฺถิยา นว, วิคเต นว, อวิคเต นว.}

\proseitem{64}{เหตุ อพฺยากโต ธมฺโม เหตุสฺส อพฺยากตสฺส ธมฺมสฺส อารมฺมณ\-ปจฺจเยน ปจฺจโย. สหชาต\-ปจฺจเยน ปจฺจโย. อุปนิสฺสย\-ปจฺจเยน ปจฺจโย.  เหตุ อพฺยากโต ธมฺโม นเหตุสฺส อพฺยากตสฺส ธมฺมสฺส อารมฺมณ\-ปจฺจเยน ปจฺจโย. สหชาต\-ปจฺจเยน ปจฺจโย. อุปนิสฺสย\-ปจฺจเยน ปจฺจโย. ปจฺฉาชาต\-ปจฺจเยน ปจฺจโย.  เหตุ อพฺยากโต ธมฺโม เหตุสฺส อพฺยากตสฺส จ นเหตุสฺส อพฺยากตสฺส จ ธมฺมสฺส อารมฺมณ\-ปจฺจเยน ปจฺจโย. สหชาต\-ปจฺจเยน ปจฺจโย. อุปนิสฺสย\-ปจฺจเยน ปจฺจโย. (3)}

\proseitem{65}{อพฺยากโต ธมฺโม นเหตุสฺส อพฺยากตสฺส ธมฺมสฺส อารมฺมณ\-ปจฺจเยน ปจฺจโย. สหชาต\-ปจฺจเยน ปจฺจโย. อุปนิสฺสย\-ปจฺจเยน ปจฺจโย. ปุเรชาต\-ปจฺจเยน ปจฺจโย. ปจฺฉาชาต\-ปจฺจเยน ปจฺจโย. อาหารปจฺจเยน ปจฺจโย. อินฺทฺริย\-ปจฺจเยน ปจฺจโย.}

\proseitem{66}{นเหตุยา นว, นอารมฺมเณ นว, นอธิปติยา นว. (สํขิตฺตํ. ปจฺจนียํ.)}

\prose{\csromanfolio{208}เหตุปจฺจยา นอารมฺมเณ ตีณิ. (สํขิตฺตํ. อนุโลม\-ปจฺจนียํ.)}

\prose{นเหตุปจฺจยา อารมฺมเณ นว. (สํขิตฺตํ. ปจฺจนียา\-นุโลมํ.)}

\prosewithcloser{(ยถา กุสลตฺติเก ปญฺหาวารสฺส อนุโลมมฺปิ ปจฺจนียมฺปิ อนุโลมปจฺจนียมฺปิ ปจฺจนียานุโลมมฺปิ คณิตํ, เอวํ คเณตพฺพํ.)}{เหตุทุกสลตฺติกํ นิฏฺฐิตํ.}
\csromanmatikaanchor{mtk.33}
\csromantocmarkref{section}{2. เวทนาตฺติก}{mtk.33}
\bookmark[dest={mtk.33},level=1]{2. เวทนาตฺติก}
\csromanmarkh{1. เหตุทุก}\csromanmarkhii{2. เวทนาตฺติก}\csromanheadernomark{1}{1. \textbf{เหตุทุก 2. เวทนาตฺติก}}
\csromanmarkh{1. สุขาย\-เวทนาย\-สมฺปยุตฺตปท}\csromanmarkhii{1. ปฏิจฺจวาร}\csromanheadernomark{2}{1. สุขาย\-เวทนาย\-สมฺปยุตฺตปท 1. ปฏิจฺจวาร}
\prose{ปจฺจยจตุกฺกเหตุ}

\proseitem{67}{เหตุํ สุขาย เวทนาย สมฺปยุตฺตํ ธมฺมํ ปฏิจฺจ เหตุ สุขาย เวทนาย สมฺปยุตฺโต ธมฺโม อุปฺปชฺชติ เหตุปจฺจยา.  เหตุํ สุขาย เวทนาย สมฺปยุตฺตํ ธมฺมํ ปฏิจฺจ นเหตุ สุขาย เวทนาย สมฺปยุตฺโต ธมฺโม อุปฺปชฺชติ เหตุปจฺจยา.  เหตุํ สุขาย เวทนาย สมฺปยุตฺตํ ธมฺมํ ปฏิจฺจ เหตุ สุขาย เวทนาย สมฺปยุตฺโต จ นเหตุ สุขาย เวทนาย สมฺปยุตฺโต จ ธมฺมา อุปฺปชฺชติ เหตุปจฺจยา. (3)}

\proseitem{68}{นเหตุํ สุขาย เวทนาย สมฺปยุตฺตํ ธมฺมํ ปฏิจฺจ นเหตุ สุขาย เวทนาย สมฺปยุตฺโต ธมฺโม อุปฺปชฺชติ เหตุปจฺจยา.  นเหตุํ สุขาย เวทนาย สมฺปยุตฺตํ ธมฺมํ ปฏิจฺจ เหตุ สุขาย เวทนาย สมฺปยุตฺโต ธมฺโม อุปฺปชฺชติ เหตุปจฺจยา.  นเหตุํ สุขาย เวทนาย สมฺปยุตฺตํ ธมฺมํ ปฏิจฺจ เหตุ สุขาย เวทนาย สมฺปยุตฺโต จ นเหตุ สุขาย เวทนาย สมฺปยุตฺโต จ ธมฺโม อุปฺปชฺชนฺติ เหตุปจฺจยา. (3)}

\proseitem{69}{เหตุํ สุขาย เวทนาย สมฺปยุตฺตญฺจ นเหตุํ สุขาย เวทนาย สมฺปยุตฺตญฺจ ธมฺมํ ปฏิจฺจ เหตุ สุขาย เวทนาย สมฺปยุตฺโต ธมฺโม อุปฺปชฺชติ เหตุปจฺจยา.  เหตุํ สุขาย เวทนาย สมฺปยุตฺตญฺจ นเหตุํ สุขาย เวทนาย สมฺปยุตฺตญฺจ ธมฺมํ ปฏิจฺจ นเหตุ สุขาย เวทนาย สมฺปยุตฺโต ธมฺโม อุปฺปชฺชติ เหตุปจฺจยา.  เหตุํ สุขาย เวทนาย}

\vspace{\tipitakaparskip}
\csromancenter{เหตุทุก 2. เวทนาตฺติก สมฺปยุตฺตญฺจ นเหตุํ สุขาย เวทนาย สมฺปยุตฺตญฺจ ธมฺมํ ปฏิจฺจ เหตุ สุขาย เวทนาย สมฺปยุตฺโต จ นเหตุ สุขาย เวทนาย สมฺปยุตฺโต จ ธมฺมา อุปฺปชฺชนฺติ เหตุปจฺจยา. (3) (สํขิตฺตํ.)}
\proseitem{70}{\csromanfolio{209}เหตุยา นว, อารมฺมเณ นว, อธิปติยา นว, อนนฺตเร นว, สมนนฺตเร นว, สหชาเต นว, อญฺญมญฺเญ นว, นิสฺสเย นว, อุปนิสฺสเย นว, ปุเรชาเต นว, อาเสวเน นว, กมฺเม นว, วิปาเก นว, อาหาเร นว ฯเปฯ อวิคเต นว. (อนุโลมํ.)}

\prose{นเหตุ นอธิปติ}

\proseitem{71}{นเหตุํ สุขาย เวทนาย สมฺปยุตฺตํ ธมฺมํ ปฏิจฺจ นเหตุ สุขาย เวทนาย สมฺปยุตฺโต ธมฺโม อุปฺปชฺชติ นเหตุปจฺจยา. (1)}

\proseitem{72}{เหตุํ สุขาย เวทนาย สมฺปยุตฺตํ ธมฺมํ ปฏิจฺจ เหตุ สุขาย เวทนาย สมฺปยุตฺโต ธมฺโม อุปฺปชฺชติ นอธิปติ\-ปจฺจยา.  เหตุํ สุขาย เวทนาย สมฺปยุตฺตํ ธมฺมํ ปฏิจฺจ นเหตุ สุขาย เวทนาย สมฺปยุตฺโต ธมฺโม อุปฺปชฺชติ นอธิปติ\-ปจฺจยา. (สํขิตฺตํ.)}

\proseitem{73}{นเหตุยา เอกํ, นอธิปติยา นว, นปุเรชาเต นว, นปจฺฉาชาเต นว, นอาเสวเน นว, นกมฺเม ตีณิ, นวิปาเก นว, นฌาเน เอกํ, นมคฺเค เอกํ, นวิปฺปยุตฺเต นว. (ปจฺจนียํ.)}

\prose{เหตุปจฺจยา นอธิปติยา นว. (สํขิตฺตํ. อนุโลม\-ปจฺจนียํ.)}

\prose{นเหตุปจฺจยา อารมฺมเณ เอกํ. (สํขิตฺตํ. ปจฺจนียา\-นุโลมํ.)}

\prose{(สหชาตวาโรปิ ปจฺจยวาโรปิ นิสฺสยวาโรปิ สํสฏฺฐ\-วาโรปิ สมฺปยุตฺตวาโรปิ ปฏิจฺจ\-วาร\-สทิสา วิตฺถาเรตพฺพา.)}

\csromanmarkh{1. สุขาย\-เวทนาย\-สมฺปยุตฺตปท}\csromanmarkhii{7. ปญฺหาวาร}\csromanheadernomark{2}{1. สุขาย\-เวทนาย\-สมฺปยุตฺตปท 7. ปญฺหาวาร}
\prose{ปจฺจยจตุกฺกเหตุ}

\proseitem{74}{เหตุ สุขาย เวทนาย สมฺปยุตฺโต ธมฺโม เหตุสฺส สุขาย เวทนาย สมฺปยุตฺตสฺส ธมฺมสฺส เหตุปจฺจเยน ปจฺจโย. ตีณิ.}

\csromanheader{2}{ธมฺมานุโลม ทุกติก\-ปฏฺฐาน\-ปาฬิ อารมฺมณาทิ}
\proseitem{75}{\csromanfolio{210}เหตุ สุขาย เวทนาย สมฺปยุตฺโต ธมฺโม เหตุสฺส สุขาย เวทนาย สมฺปยุตฺตสฺส ธมฺมสฺส อารมฺมณ\-ปจฺจเยน ปจฺจโย. ตีณิ.}

\prose{สุขาย เวทนาย สมฺปยุตฺโต ธมฺโม นเหตุสฺส สุขาย เวทนาย สมฺปยุตฺตสฺส ธมฺมสฺส อารมฺมณ\-ปจฺจเยน ปจฺจโย. ตีณิ.}

\prose{เหตุ สุขาย เวทนาย สมฺปยุตฺโต จ นเหตุ สุขาย เวทนาย สมฺปยุตฺโต จ ธมฺมา เหตุสฺส สุขาย เวทนาย สมฺปยุตฺตสฺส ธมฺมสฺส อารมฺมณ\-ปจฺจเยน ปจฺจโย. ตีณิ.}

\proseitem{76}{เหตุ สุขาย เวทนาย สมฺปยุตฺโต ธมฺโม เหตุสฺส สุขาย เวทนาย สมฺปยุตฺตสฺส ธมฺมสฺส อธิปติ\-ปจฺจเยน ปจฺจโย. อารมฺมณา\-ธิปติ, สหชาตา\-ธิปติ. ตีณิ.}

\prose{สุขาย เวทนาย สมฺปยุตฺโต ธมฺโม นเหตุสฺส สุขาย เวทนาย สมฺปยุตฺตสฺส ธมฺมสฺส อธิปติ\-ปจฺจเยน ปจฺจโย. อารมฺมณา\-ธิปติ, สหชาตา\-ธิปติ. ตีณิ.}

\prose{เหตุ สุขาย เวทนาย สมฺปยุตฺโต จ นเหตุ สุขาย เวทนาย สมฺปยุตฺโต จ ธมฺมา เหตุสฺส สุขาย เวทนาย สมฺปยุตฺตสฺส ธมฺมสฺส อธิปติ\-ปจฺจเยน ปจฺจโย. อารมฺมณา\-ธิปติ. ตีณิ ฯเปฯ.}

\csromanheader{2}{\textbf{อุปนิสฺสยาทิ}}
\proseitem{76}{เหตุ สุขาย เวทนาย สมฺปยุตฺโต ธมฺโม เหตุสฺส สุขาย เวทนาย สมฺปยุตฺตสฺส ธมฺมสฺส อุปนิสฺสย\-ปจฺจเยน ปจฺจโย. อารมฺมณู\-ปนิสฺสโย, อนนฺตรู\-ปนิสฺสโย, ปกตู\-ปนิสฺสโย. นว. (สํขิตฺตํ.)}

\proseitem{78}{สุขาย เวทนาย สมฺปยุตฺโต ธมฺโม นเหตุสฺส สุขาย เวทนาย สมฺปยุตฺตสฺส ธมฺมสฺส กมฺมปจฺจเยน ปจฺจโย.  นเหตุ สุขาย เวทนาย สมฺปยุตฺโต ธมฺโม เหตุสฺส สุขาย เวทนาย สมฺปยุตฺตสฺส ธมฺมสฺส กมฺมปจฺจเยน ปจฺจโย.  นเหตุ สุขาย เวทนาย สมฺปยุตฺโต ธมฺโม เหตุสฺส สุขาย เวทนาย สมฺปยุตฺตสฺส จ}

\vspace{\tipitakaparskip}
\csromancenter{เหตุทุก 2. เวทนาตฺติก นเหตุสฺส สุขาย เวทนาย สมฺปยุตฺตสฺส จ ธมฺมสฺส กมฺมปจฺจเยน ปจฺจโย. (3). วิปาก\-ปจฺจเยน ปจฺจโย.}
\proseitem{79}{\csromanfolio{211}สุขาย เวทนาย สมฺปยุตฺโต ธมฺโม นเหตุสฺส สุขาย เวทนาย สมฺปยุตฺตสฺส ธมฺมสฺส อาหารปจฺจเยน ปจฺจโย. ตีณิ ฯเปฯ. เหตุ สุขาย เวทนาย สมฺปยุตฺโต ธมฺโม เหตุสฺส สุขาย เวทนาย สมฺปยุตฺตสฺส ธมฺมสฺส อวิคต\-ปจฺจเยน ปจฺจโย.}

\proseitem{80}{เหตุยา ตีณิ, อารมฺมเณ นว, อธิปติยา นว, อนนฺตเร นว, สมนนฺตเร นว, สหชาเต นว, อญฺญมญฺเญ นว, นิสฺสเย นว, อุปนิสฺสเย นว, อาเสวเน นว, กมฺเม ตีณิ, วิปาเก นว, อาหาเร ตีณิ, อินฺทฺริเย นว, ฌาเน ตีณิ, มคฺเค นว, สมฺปยุตฺเต นว, อตฺถิยา นว, นตฺถิยา นว, วิคเต นว, อวิคเต นว. เหตุ สุขาย เวทนาย สมฺปยุตฺโต ธมฺโม เหตุสฺส สุขาย เวทนาย สมฺปยุตฺตสฺส ธมฺมสฺส อารมฺมณ\-ปจฺจเยน ปจฺจโย. สหชาต\-ปจฺจเยน ปจฺจโย. อุปนิสฺสย\-ปจฺจเยน ปจฺจโย. (สํขิตฺตํ.)}

\proseitem{82}{นเหตุยา นว, นอารมฺมเณ นว ฯเปฯ โนอวิคเต นว. (ปจฺจนียํ.)}

\prose{เหตุปจฺจยา นอารมฺมเณ ตีณิ. (สํขิตฺตํ. อนุโลม\-ปจฺจนียํ.)}

\prose{นเหตุปจฺจยา อารมฺมเณ นว. (สํขิตฺตํ. ปจฺจนียา\-นุโลมํ.)}

\prose{(ยถา กุสลตฺติเก ปญฺหาวารสฺส อนุโลมมฺปิ ปจฺจนียมฺปิ อนุโลมปจฺจนียมฺปิ ปจฺจนียานุโลมมฺปิ คณิตํ, เอวํ คเณตพฺพํ.)}

\prose{2. ทุกฺขาย\-เวทนาย\-สมฺปยุตฺตปท 1. ปฏิจฺจวาร}

\prose{ปจฺจยจตุกฺกเหตุ}

\proseitem{79}{เหตุํ ทุกฺขาย เวทนาย สมฺปยุตฺตํ ธมฺมํ ปฏิจฺจ เหตุ ทุกฺขาย เวทนาย สมฺปยุตฺโต ธมฺโม อุปฺปชฺชติ เหตุปจฺจยา.  เหตุํ ทุกฺขาย \csromanfolio{212}เวทนาย สมฺปยุตฺตํ ธมฺมํ ปฏิจฺจ นเหตุ ทุกฺขาย เวทนาย สมฺปยุตฺโต ธมฺโม อุปฺปชฺชติ เหตุปจฺจยา.  เหตุํ ทุกฺขาย เวทนาย สมฺปยุตฺตํ ธมฺมํ ปฏิจฺจ เหตุ ทุกฺขาย เวทนาย สมฺปยุตฺโต จ นเหตุ ทุกฺขาย เวทนาย สมฺปยุตฺโต จ ธมฺมา อุปฺปชฺชนฺติ เหตุปจฺจยา. (3)}

\proseitem{83}{นเหตุํ ทุกฺขาย เวทนาย สมฺปยุตฺตํ ธมฺมํ ปฏิจฺจ นเหตุ ทุกฺขาย เวทนาย สมฺปยุตฺโต ธมฺโม อุปฺปชฺชติ เหตุปจฺจยา. ตีณิ.}

\prose{เหตุํ ทุกฺขาย เวทนาย สมฺปยุตฺตญฺจ นเหตุํ ทุกฺขาย เวทนาย สมฺปยุตฺตญฺจ ธมฺมํ ปฏิจฺจ เหตุ ทุกฺขาย เวทนาย สมฺปยุตฺโต ธมฺโม อุปฺปชฺชติ เหตุปจฺจยา. ตีณิ. (สํขิตฺตํ.)}

\prose{เหตุยา นว, อารมฺมเณ นว, อธิปติยา นว, อนนฺตเร นว, สมนนฺตเร นว, สหชาเต นว, อญฺญมญฺเญ นว, นิสฺสเย นว, อุปนิสฺสเย นว, ปุเรชาเต นว, อาเสวเน นว, กมฺเม นว, วิปาเก เอกํ, อาหาเร นว, อินฺทฺริเย นว, ฌาเน นว, มคฺเค นว, สมฺปยุตฺเต นว, วิปฺปยุตฺเต นว, อตฺถิยา นว, นตฺถิยา นว, วิคเต นว, อวิคเต นว. (อนุโลมํ.)}

\prose{นเหตุ นอธิปติ}

\proseitem{83}{นเหตุํ ทุกฺขาย เวทนาย สมฺปยุตฺตํ ธมฺมํ ปฏิจฺจ นเหตุ ทุกฺขาย เวทนาย สมฺปยุตฺโต ธมฺโม อุปฺปชฺชติ นเหตุปจฺจยา. (1)}

\proseitem{86}{เหตุํ ทุกฺขาย เวทนาย สมฺปยุตฺตํ ธมฺมํ ปฏิจฺจ เหตุ ทุกฺขาย เวทนาย สมฺปยุตฺโต ธมฺโม อุปฺปชฺชติ นอธิปติ\-ปจฺจยา. (สํขิตฺตํ.)}

\proseitem{83}{นเหตุยา เอกํ, นอธิปติยา นว, นปจฺฉาชาเต นว, นอาเสวเน นว, นกมฺเม ตีณิ, นวิปาเก นว, นฌาเน เอกํ, นมคฺเค เอกํ. (ปจฺจนียํ.)}

\prose{เหตุปจฺจยา นอธิปติยา นว. (สํขิตฺตํ. อนุโลม\-ปจฺจนียํ.)}

\prose{นเหตุปจฺจยา อารมฺมเณ เอกํ. (สํขิตฺตํ. ปจฺจนียา\-นุโลมํ)}

\prose{(สหชาตวาโรปิ ปจฺจยวาโรปิ นิสฺสยวาโรปิ สํสฏฺฐ\-วาโรปิ สมฺปยุตฺตวาโรปิ ปฏิจฺจ\-วาร\-สทิสา วิตฺถาเรตพฺพา.)}

\vspace{\tipitakaparskip}
\csromancenter{เหตุทุก 2. เวทนาตฺติก 2. ทุกฺขาย\-เวทนาย\-สมฺปยุตฺตปท 7. ปญฺหาวาร}
\prose{\csromanfolio{213}ปจฺจยจตุกฺกเหตุ}

\proseitem{88}{เหตุ ทุกฺขาย เวทนาย สมฺปยุตฺโต ธมฺโม เหตุสฺส ทุกฺขาย เวทนาย สมฺปยุตฺตสฺส ธมฺมสฺส เหตุปจฺจเยน ปจฺจโย.  เหตุ ทุกฺขาย เวทนาย สมฺปยุตฺโต ธมฺโม นเหตุสฺส ทุกฺขาย เวทนาย สมฺปยุตฺตสฺส ธมฺมสฺส เหตุปจฺจเยน ปจฺจโย.  เหตุ ทุกฺขาย เวทนาย สมฺปยุตฺโต ธมฺโม เหตุสฺส ทุกฺขาย เวทนาย สมฺปยุตฺตสฺส จ นเหตุสฺส ทุกฺขาย เวทนาย สมฺปยุตฺตสฺส จ ธมฺมสฺส เหตุปจฺจเยน ปจฺจโย. (3)}

\csromanheader{2}{\textbf{อารมฺมณาทิ}}
\proseitem{89}{เหตุ ทุกฺขาย เวทนาย สมฺปยุตฺโต ธมฺโม เหตุสฺส ทุกฺขาย เวทนาย สมฺปยุตฺตสฺส ธมฺมสฺส อารมฺมณ\-ปจฺจเยน ปจฺจโย. ตีณิ. ทุกฺขาย เวทนาย สมฺปยุตฺโต ธมฺโม นเหตุสฺส ทุกฺขาย เวทนาย สมฺปยุตฺตสฺส ธมฺมสฺส อารมฺมณ\-ปจฺจเยน ปจฺจโย. ตีณิ.}

\prose{เหตุ ทุกฺขาย เวทนาย สมฺปยุตฺโต จ นเหตุ ทุกฺขาย เวทนาย สมฺปยุตฺโต จ ธมฺมา เหตุสฺส ทุกฺขาย เวทนาย สมฺปยุตฺตสฺส ธมฺมสฺส อารมฺมณ\-ปจฺจเยน ปจฺจโย. ตีณิ. (สํขิตฺตํ.) ทุกฺขาย เวทนาย สมฺปยุตฺโต ธมฺโม นเหตุสฺส ทุกฺขาย เวทนาย สมฺปยุตฺตสฺส ธมฺมสฺส กมฺมปจฺจเยน ปจฺจโย. ตีณิ. (สํขิตฺตํ.)}

\proseitem{91}{เหตุยา ตีณิ, อารมฺมเณ นว, อธิปติยา ตีณิ, อนนฺตเร นว, สมนนฺตเร นว, สหชาเต นว, อญฺญมญฺเญ นว, นิสฺสเย นว, อุปนิสฺสเย นว, อาเสวเน นว, กมฺเม ตีณิ, วิปาเก เอกํ, อาหาเร ตีณิ, อินฺทฺริเย ตีณิ, ฌาเน ตีณิ, มคฺเค ตีณิ, สมฺปยุตฺเต นว, อตฺถิยา นว, นตฺถิยา นว, วิคเต นว, อวิคเต นว. (อนุโลมํ.)}

\csromanheader{2}{ธมฺมานุโลม ทุกติก\-ปฏฺฐาน\-ปาฬิ ปจฺจนียุ\-ทฺธาร}
\proseitem{92}{\csromanfolio{214}เหตุ ทุกฺขาย เวทนาย สมฺปยุตฺโต ธมฺโม เหตุสฺส ทุกฺขาย เวทนาย สมฺปยุตฺตสฺส ธมฺมสฺส อารมฺมณ\-ปจฺจเยน ปจฺจโย. สหชาต\-ปจฺจเยน ปจฺจโย. อุปนิสฺสย\-ปจฺจเยน ปจฺจโย. (สํขิตฺตํ.)}

\proseitem{93}{นเหตุยา นว, นอารมฺมเณ นว, นอธิปติยา นว. (สํขิตฺตํ. ปจฺจนียํ.)}

\prose{เหตุปจฺจยา นอารมฺมเณ ตีณิ. (สํขิตฺตํ. อนุโลม\-ปจฺจนียํ.)}

\prose{นเหตุปจฺจยา อารมฺมเณ นว. (สํขิตฺตํ. ปจฺจนียา\-นุโลมํ.)}

\prose{(ยถา กุสลตฺติเก ปญฺหาวารสฺส อนุโลมมฺปิ ปจฺจนียมฺปิ อนุโลมปจฺจนียมฺปิ ปจฺจนียานุโลมมฺปิ คณิตํ, เอวํ คเณตพฺพํ.)}

\prose{3. อทุกฺขมสุขาย\-เวทนาย\-สมฺปยุตฺตปท 1. ปฏิจฺจวาร}

\prose{ปจฺจยจตุกฺกเหตุ}

\proseitem{94}{เหตุํ อทุกฺขม\-สุขาย เวทนาย สมฺปยุตฺตํ ธมฺมํ ปฏิจฺจ เหตุ อทุกฺขม\-สุขาย เวทนาย สมฺปยุตฺโต ธมฺโม อุปฺปชฺชติ เหตุปจฺจยา. ตีณิ.}

\prose{นเหตุํ อทุกฺขม\-สุขาย เวทนาย สมฺปยุตฺตํ ธมฺมํ ปฏิจฺจ นเหตุ อทุกฺขม\-สุขาย เวทนาย สมฺปยุตฺโต ธมฺโม อุปฺปชฺชติ เหตุปจฺจยา. ตีณิ.}

\prose{เหตุํ อทุกฺขม\-สุขาย เวทนาย สมฺปยุตฺตญฺจ นเหตุํ อทุกฺขม\-สุขาย เวทนาย สมฺปยุตฺตญฺจ ธมฺมํ ปฏิจฺจ เหตุ อทุกฺขม\-สุขาย เวทนาย สมฺปยุตฺโต ธมฺโม อุปฺปชฺชติ เหตุปจฺจยา. ตีณิ. (สํขิตฺตํ.)}

\proseitem{95}{เหตุยา นว, อารมฺมเณ นว, อธิปติยา นว, อนนฺตเร นว, สมนนฺตเร นว, สหชาเต นว, อญฺญมญฺเญ นว, นิสฺสเย นว, อุปนิสฺสเย นว, ปุเรชาเต นว, อาเสวเน นว, กมฺเม นว, วิปาเก นว, อาหาเร}

\vspace{\tipitakaparskip}
\csromancenter{เหตุทุก 2. เวทนาตฺติก นว, อินฺทฺริเย นว, ฌาเน นว, มคฺเค นว, สมฺปยุตฺเต นว, วิปฺปยุตฺเต นว, อตฺถิยา นว, นตฺถิยา นว, วิคเต นว, อวิคเต นว. (อนุโลมํ.)}
\prose{\csromanfolio{215}นเหตุ นอธิปติ}

\proseitem{96}{นเหตุํ อทุกฺขม\-สุขาย เวทนาย สมฺปยุตฺตํ ธมฺมํ ปฏิจฺจ นเหตุ อทุกฺขม\-สุขาย เวทนาย สมฺปยุตฺโต ธมฺโม อุปฺปชฺชติ นเหตุปจฺจยา. เทฺว.}

\proseitem{97}{เหตุํ อทุกฺขม\-สุขาย เวทนาย สมฺปยุตฺตํ ธมฺมํ ปฏิจฺจ เหตุ อทุกฺขม\-สุขาย เวทนาย สมฺปยุตฺโต ธมฺโม อุปฺปชฺชติ นอธิปติ\-ปจฺจยา. (สํขิตฺตํ.)}

\proseitem{98}{นเหตุยา เทฺว, นอธิปติยา นว, นปุเรชาเต นว, นปจฺฉาชาเต นว, นอาเสวเน นว, นกมฺเม ตีณิ, นวิปาเก นว, นฌาเน เอกํ, นมคฺเค เอกํ, นวิปฺปยุตฺเต นว. (สํขิตฺตํ. ปจฺจนียํ.)}

\prose{เหตุปจฺจยา นอธิปติยา นว. (สํขิตฺตํ. อนุโลม\-ปจฺจนียํ.)}

\prose{นเหตุปจฺจยา อารมฺมเณ เทฺว. (สํขิตฺตํ. ปจฺจนียา\-นุโลมํ.)}

\prose{(สหชาตวาโรปิ ปจฺจยวาโรปิ นิสฺสยวาโรปิ สํสฏฺฐ\-วาโรปิ สมฺปยุตฺตวาโรปิ ปฏิจฺจ\-วาร\-สทิสา วิตฺถาเรตพฺพา.)}

\prose{3. อทุกฺขมสุขาย\-เวทนาย\-สมฺปยุตฺตปท 7. ปญฺหาวาร}

\prose{ปจฺจยจตุกฺกเหตฺวาทิ}

\proseitem{99}{เหตุ อทุกฺขม\-สุขาย เวทนาย สมฺปยุตฺโต ธมฺโม เหตุสฺส อทุกฺขม\-สุขาย เวทนาย สมฺปยุตฺตสฺส ธมฺมสฺส เหตุปจฺจเยน ปจฺจโย. ตีณิ.}

\proseitem{100}{เหตุ อทุกฺขม\-สุขาย เวทนาย สมฺปยุตฺโต ธมฺโม เหตุสฺส อทุกฺขม\-สุขาย เวทนาย สมฺปยุตฺตสฺส ธมฺมสฺส อารมฺมณ\-ปจฺจเยน ปจฺจโย. นว.}

\proseitem{101}{\csromanfolio{216}เหตุ อทุกฺขม\-สุขาย เวทนาย สมฺปยุตฺโต ธมฺโม เหตุสฺส อทุกฺขม\-สุขาย เวทนาย สมฺปยุตฺตสฺส ธมฺมสฺส อธิปติ\-ปจฺจเยน ปจฺจโย. อารมฺมณา\-ธิปติ, สหชาตา\-ธิปติ. ตีณิ.}

\prose{อทุกฺขม\-สุขาย เวทนาย สมฺปยุตฺโต ธมฺโม นเหตุสฺส อทุกฺขม\-สุขาย เวทนาย สมฺปยุตฺตสฺส ธมฺมสฺส อธิปติ\-ปจฺจเยน ปจฺจโย. อารมฺมณา\-ธิปติ, สหชาตา\-ธิปติ. ตีณิ.}

\prose{เหตุ อทุกฺขม\-สุขาย เวทนาย สมฺปยุตฺโต จ นเหตุ อทุกฺขม\-สุขาย เวทนาย สมฺปยุตฺโต จ ธมฺมา เหตุสฺส อทุกฺขม\-สุขาย เวทนาย สมฺปยุตฺตสฺส ธมฺมสฺส อธิปติ\-ปจฺจเยน ปจฺจโย. อารมฺมณา\-ธิปติ. ตีณิ. (อารมฺมณา\-ธิปติเยว.) ฯเปฯ.}

\csromanheader{2}{\textbf{อุปนิสฺสยาทิ}}
\proseitem{102}{เหตุ อทุกฺขม\-สุขาย เวทนาย สมฺปยุตฺโต ธมฺโม เหตุสฺส อทุกฺขม\-สุขาย เวทนาย สมฺปยุตฺตสฺส ธมฺมสฺส อุปนิสฺสย\-ปจฺจเยน ปจฺจโย. อารมฺมณู\-ปนิสฺสโย, อนนฺตรู\-ปนิสฺสโย, ปกตู\-ปนิสฺสโย. นว. อาเสวน\-ปจฺจเยน ปจฺจโย. นว.}

\proseitem{103}{อทุกฺขม\-สุขาย เวทนาย สมฺปยุตฺโต ธมฺโม นเหตุสฺส อทุกฺขม\-สุขาย เวทนาย สมฺปยุตฺตสฺส ธมฺมสฺส กมฺมปจฺจเยน ปจฺจโย. ตีณิ.}

\proseitem{104}{เหตุ อทุกฺขม\-สุขาย เวทนาย สมฺปยุตฺโต ธมฺโม เหตุสฺส อทุกฺขม\-สุขาย เวทนาย สมฺปยุตฺตสฺส ธมฺมสฺส วิปาก\-ปจฺจเยน ปจฺจโย. นว ฯเปฯ. อวิคต\-ปจฺจเยน ปจฺจโย. นว.}

\proseitem{105}{เหตุยา ตีณิ, อารมฺมเณ นว, อธิปติยา นว, อนนฺตเร นว, สมนนฺตเร นว, สหชาเต นว, อญฺญมญฺเญ นว, นิสฺสเย นว, อุปนิสฺสเย นว, อาเสวเน นว, กมฺเม ตีณิ, วิปาเก นว, อาหาเร ตีณิ, อินฺทฺริเย นว, ฌาเน ตีณิ, มคฺเค นว, สมฺปยุตฺเต นว, อตฺถิยา นว, นตฺถิยา นว, วิคเต นว, อวิคเต นว. (อนุโลมํ.)}

\csromanmarkh{1. เหตุทุก}\csromanmarkhii{3. วิปากตฺติก ปจฺจนียุ\-ทฺธาร}\csromanheadernomark{2}{1. เหตุทุก 3. วิปากตฺติก ปจฺจนียุ\-ทฺธาร}
\proseitem{106}{\csromanfolio{217}เหตุ อทุกฺขม\-สุขาย เวทนาย สมฺปยุตฺโต ธมฺโม เหตุสฺส อทุกฺขม\-สุขาย เวทนาย สมฺปยุตฺตสฺส ธมฺมสฺส อารมฺมณ\-ปจฺจเยน ปจฺจโย. สหชาต\-ปจฺจเยน ปจฺจโย. อุปนิสฺสย\-ปจฺจเยน ปจฺจโย. (สํขิตฺตํ.)}

\proseitem{107}{นเหตุยา นว, นอารมฺมเณ นว. (สํขิตฺตํ. ปจฺจนียํ.)}

\prose{เหตุปจฺจยา นอารมฺมเณ ตีณิ. (สํขิตฺตํ. อนุโลม\-ปจฺจนียํ.)}

\prose{นเหตุปจฺจยา อารมฺมเณ นว. (สํขิตฺตํ. ปจฺจนียา\-นุโลมํ.)}

\prosewithcloser{(ยถา กุสลตฺติเก ปญฺหาวารสฺส อนุโลมมฺปิ ปจฺจนียมฺปิ อนุโลมปจฺจนียมฺปิ ปจฺจนียานุโลมมฺปิ คณิตํ, เอวํ คเณตพฺพํ.)}{เหตุ\-ทุก\-เวทนา\-ตฺติกํ นิฏฺฐิตํ.}
\csromanmatikaanchor{mtk.34}
\csromantocmarkref{section}{3. วิปากตฺติก}{mtk.34}
\bookmark[dest={mtk.34},level=1]{3. วิปากตฺติก}
\csromanmarkh{1. เหตุทุก}\csromanmarkhii{3. วิปากตฺติก}\csromanheadernomark{1}{1. เหตุทุก 3. วิปากตฺติก}
\csromanmarkh{1. วิปากปท}\csromanmarkhii{1. ปฏิจฺจวาร}\csromanheadernomark{2}{1. \textbf{วิปากปท 1. ปฏิจฺจวาร}}
\prose{ปจฺจยจตุกฺกเหตุ}

\proseitem{108}{เหตุํ วิปากํ ธมฺมํ ปฏิจฺจ เหตุ วิปาโก ธมฺโม อุปฺปชฺชติ เหตุปจฺจยา.  เหตุํ วิปากํ ธมฺมํ ปฏิจฺจ นเหตุ วิปาโก ธมฺโม อุปฺปชฺชติ เหตุปจฺจยา.  เหตุํ วิปากํ ธมฺมํ ปฏิจฺจ เหตุ วิปาโก จ นเหตุ วิปาโก จ ธมฺมา อุปฺปชฺชนฺติ เหตุปจฺจยา. (3)}

\prose{นเหตุํ วิปากํ ธมฺมํ ปฏิจฺจ นเหตุ วิปาโก ธมฺโม อุปฺปชฺชติ เหตุปจฺจยา. ตีณิ.}

\prose{เหตุํ วิปากญฺจ นเหตุํ วิปากญฺจ ธมฺมํ ปฏิจฺจ เหตุ วิปาโก ธมฺโม อุปฺปชฺชติ เหตุปจฺจยา. ตีณิ. (สํขิตฺตํ.)}

\proseitem{109}{เหตุยา นว, อารมฺมเณ นว, อธิปติยา นว, อนนฺตเร นว, สมนนฺตเร นว, สหชาเต นว, อญฺญมญฺเญ นว, นิสฺสเย นว, อุปนิสฺสเย \csromanfolio{218}นว, ปุเรชาเต นว, กมฺเม นว, วิปาเก นว, อาหาเร นว, อินฺทฺริเย นว, ฌาเน นว, มคฺเค นว, สมฺปยุตฺเต นว, วิปฺปยุตฺเต นว, อตฺถิยา นว, นตฺถิยา นว, วิคเต นว, อวิคเต นว. (อนุโลมํ.)}

\prose{นเหตุ นอธิปติ}

\proseitem{110}{นเหตุํ วิปากํ ธมฺมํ ปฏิจฺจ นเหตุ วิปาโก ธมฺโม อุปฺปชฺชติ นเหตุปจฺจยา. (1)}

\proseitem{111}{เหตุํ วิปากํ ธมฺมํ ปฏิจฺจ เหตุ วิปาโก ธมฺโม อุปฺปชฺชติ นอธิปติ\-ปจฺจยา. (สํขิตฺตํ.)}

\proseitem{112}{นเหตุยา เอกํ, นอธิปติยา นว, นปุเรชาเต นว, นปจฺฉาชาเต นว, นอาเสวเน นว, นฌาเน เอกํ, นมคฺเค เอกํ, นวิปฺปยุตฺเต นว. (สํขิตฺตํ. ปจฺจนียํ.)}

\prose{เหตุปจฺจยา นอธิปติยา นว. (สํขิตฺตํ. อนุโลมฺมปจฺจนียํ.)}

\prose{นเหตุปจฺจยา อารมฺมเณ เอกํ. (สํขิตฺตํ. ปจฺจนียา\-นุโลมํ.)}

\prose{(สหชาตวาโรปิ ปจฺจยวาโรปิ นิสฺสยวาโรปิ สํสฏฺฐ\-วาโรปิ สมฺปยุตฺตวาโรปิ ปฏิจฺจ\-วาร\-สทิสา วิตฺถาเรตพฺพา.)}

\csromanmarkh{1. วิปากปท}\csromanmarkhii{7. ปญฺหาวาร}\csromanheadernomark{2}{1. \textbf{วิปากปท 7. ปญฺหาวาร}}
\prose{ปจฺจยจตุกฺกเหตฺวาทิ}

\proseitem{113}{เหตุ วิปาโก ธมฺโม เหตุสฺส วิปากสฺส ธมฺมสฺส เหตุปจฺจเยน ปจฺจโย. ตีณิ.}

\proseitem{114}{เหตุ วิปาโก ธมฺโม เหตุสฺส วิปากสฺส ธมฺมสฺส อารมฺมณ\-ปจฺจเยน ปจฺจโย. ตีณิ.}

\prose{วิปาโก ธมฺโม นเหตุสฺส วิปากสฺส ธมฺมสฺส อารมฺมณ\-ปจฺจเยน ปจฺจโย. ตีณิ.}

\csromanmarkh{1. เหตุทุก}\csromanmarkhii{3. วิปากตฺติก}\csromanheadernomark{2}{1. เหตุทุก 3. วิปากตฺติก}
\prose{\csromanfolio{219}เหตุ วิปาโก จ นเหตุ วิปาโก จ ธมฺมา เหตุสฺส วิปากสฺส ธมฺมสฺส อารมฺมณ\-ปจฺจเยน ปจฺจโย. ตีณิ. (ตทารมฺมณาเยว ลพฺภนฺติ.)}

\proseitem{115}{เหตุ วิปาโก ธมฺโม เหตุสฺส วิปากสฺส ธมฺมสฺส อธิปติ\-ปจฺจเยน ปจฺจโย. ตีณิ. (สหชาตา\-ธิปติเยว ลพฺภติ. อารมฺมณา\-ธิปติ นตฺถิ.) ฯเปฯ.}

\csromanheader{2}{\textbf{อุปนิสฺสยาทิ}}
\proseitem{116}{เหตุ วิปาโก ธมฺโม เหตุสฺส วิปากสฺส ธมฺมสฺส อุปนิสฺสย\-ปจฺจเยน ปจฺจโย. อนนฺตรู\-ปนิสฺสโย, ปกตู\-ปนิสฺสโย.  เหตุ วิปาโก ธมฺโม นเหตุสฺส วิปากสฺส ธมฺมสฺส อุปนิสฺสย\-ปจฺจเยน ปจฺจโย. อนนฺตรู\-ปนิสฺสโย, ปกตู\-ปนิสฺสโย.  เหตุ วิปาโก ธมฺโม เหตุสฺส วิปากสฺส จ นเหตุสฺส วิปากสฺส จ ธมฺมสฺส อุปนิสฺสย\-ปจฺจเยน ปจฺจโย. อนนฺตรู\-ปนิสฺสโย, ปกตู\-ปนิสฺสโย. (3) วิปาโก ธมฺโม นเหตุสฺส วิปากสฺส ธมฺมสฺส อุปนิสฺสย\-ปจฺจเยน ปจฺจโย. อนนฺตรู\-ปนิสฺสโย, ปกตู\-ปนิสฺสโย ฯเปฯ.}

\vspace{\tipitakaparskip}
\csromancenter{(อิตเร เทฺว อนนฺตรู\-ปนิสฺสโย ปกตูปนิสฺสโยเยว.)}
\prosecont{วิปาโก ธมฺโม นเหตุสฺส วิปากสฺส ธมฺมสฺส กมฺมปจฺจเยน ปจฺจโย. ตีณิ. (สหชาต\-กมฺมเมว. สํขิตฺตํ.)}

\proseitem{117}{เหตุ วิปาโก ธมฺโม เหตุสฺส วิปากสฺส ธมฺมสฺส วิปาก\-ปจฺจเยน ปจฺจโย. นว. วิปาโก ธมฺโม นเหตุสฺส วิปากสฺส ธมฺมสฺส อาหารปจฺจเยน ปจฺจโย. (สํขิตฺตํ.)}

\proseitem{118}{เหตุยา ตีณิ, อารมฺมเณ นว, อธิปติยา ฉ, อนนฺตเร นว, สมนนฺตเร นว, สหชาเต นว, อญฺญมญฺเญ นว, นิสฺสเย นว, อุปนิสฺสเย นว, กมฺเม ตีณิ, วิปาเก นว, อาหาเร ตีณิ, อินฺทฺริเย นว, ฌาเน ตีณิ, มคฺเค นว, สมฺปยุตฺเต นว, อตฺถิยา นว, นตฺถิยา นว, วิคเต นว, อวิคเต นว. (อนุโลมํ.)}

\csromanheader{2}{ธมฺมานุโลม ทุกติก\-ปฏฺฐาน\-ปาฬิ ปจฺจนียุ\-ทฺธาร}
\proseitem{119}{\csromanfolio{220}เหตุ วิปาโก ธมฺโม เหตุสฺส วิปากสฺส ธมฺมสฺส อารมฺมณ\-ปจฺจเยน ปจฺจโย. สหชาต\-ปจฺจเยน ปจฺจโย. อุปนิสฺสย\-ปจฺจเยน ปจฺจโย. (สํขิตฺตํ.)}

\proseitem{120}{นเหตุยา นว, นอารมฺมเณ นว. (สํขิตฺตํ. ปจฺจนียํ.)}

\prose{เหตุปจฺจยา นอารมฺมเณ ตีณิ. (สํขิตฺตํ. อนุโลม\-ปจฺจนียํ.)}

\prose{นเหตุปจฺจยา อารมฺมเณ นว. (สํขิตฺตํ. ปจฺจนียา\-นุโลมํ.)}

\prose{(ยถา กุสลตฺติเก ปญฺหาวารสฺส อนุโลมมฺปิ ปจฺจนียมฺปิ อนุโลมปจฺจนียมฺปิ ปจฺจนียานุโลมมฺปิ คณิตํ, เอวํ คเณตพฺพํ.)}

\csromanmarkh{2. วิปาก\-ธมฺมปท}\csromanmarkhii{1. ปฏิจฺจวาร}\csromanheadernomark{2}{2. วิปาก\-ธมฺมปท 1. ปฏิจฺจวาร}
\prose{ปจฺจยจตุกฺกเหตุ}

\proseitem{121}{เหตุํ วิปาก\-ธมฺม\-ธมฺมํ ปฏิจฺจ เหตุ วิปาก\-ธมฺม\-ธมฺโม อุปฺปชฺชติ เหตุปจฺจยา.  เหตุํ วิปาก\-ธมฺม\-ธมฺมํ ปฏิจฺจ นเหตุ วิปาก\-ธมฺม\-ธมฺโม อุปฺปชฺชติ เหตุปจฺจยา.  เหตุํ วิปาก\-ธมฺม\-ธมฺมํ ปฏิจฺจ เหตุ วิปาก\-ธมฺม\-ธมฺโม จ นเหตุ วิปาก\-ธมฺม\-ธมฺโม จ ธมฺมา อุปฺปชฺชนฺติ เหตุปจฺจยา. (3)}

\prose{นเหตุํ วิปาก\-ธมฺม\-ธมฺมํ ปฏิจฺจ นเหตุ วิปาก\-ธมฺม\-ธมฺโม อุปฺปชฺชติ เหตุปจฺจยา. ตีณิ.}

\prose{เหตุํ วิปาก\-ธมฺม\-ธมฺมญฺจ นเหตุํ วิปาก\-ธมฺม\-ธมฺมญฺจ ปฏิจฺจ เหตุ วิปาก\-ธมฺม\-ธมฺโม อุปฺปชฺชติ เหตุปจฺจยา. ตีณิ. (สํขิตฺตํ.)}

\proseitem{122}{เหตุยา นว, อารมฺมเณ นว, อธิปติยา นว ฯเปฯ กมฺเม นว, อาหาเร นว, อินฺทฺริเย นว, ฌาเน นว, มคฺเค นว, สมฺปยุตฺเต นว, วิปฺปยุตฺเต นว, อตฺถิยา นว, นตฺถิยา นว, วิคเต นว, อวิคเต นว.}

\prose{นเหตุ นอธิปติ}

\proseitem{123}{นเหตุํ วิปาก\-ธมฺม\-ธมฺมํ ปฏิจฺจ เหตุ วิปาก\-ธมฺม\-ธมฺโม อุปฺปชฺชติ นเหตุปจฺจยา. (1)}

\csromanmarkh{1. เหตุทุก}\csromanmarkhii{3. วิปากตฺติก}\csromanheadernomark{2}{1. เหตุทุก 3. วิปากตฺติก}
\proseitem{124}{\csromanfolio{221}เหตุํ วิปาก\-ธมฺม\-ธมฺมํ ปฏิจฺจ เหตุ วิปาก\-ธมฺม\-ธมฺโม อุปฺปชฺชติ นอธิปติ\-ปจฺจยา. นว.}

\csromanheader{2}{\textbf{นปุเรชาตาทิ}}
\proseitem{125}{เหตุํ วิปาก\-ธมฺม\-ธมฺมํ ปฏิจฺจ เหตุ วิปาก\-ธมฺม\-ธมฺโม อุปฺปชฺชติ นปุเรชาต\-ปจฺจยา. นว. นปจฺฉาชาต\-ปจฺจยา. นว. นอาเสวน\-ปจฺจยา. นว.}

\proseitem{126}{เหตุํ วิปาก\-ธมฺม\-ธมฺมํ ปฏิจฺจ นเหตุ วิปาก\-ธมฺม\-ธมฺโม อุปฺปชฺชติ นกมฺมปจฺจยา. (1)}

\prose{นเหตุํ วิปาก\-ธมฺม\-ธมฺมํ ปฏิจฺจ นเหตุ วิปาก\-ธมฺม\-ธมฺโม อุปฺปชฺชติ นกมฺมปจฺจยา. (1)}

\prose{เหตุํ วิปาก\-ธมฺม\-ธมฺมญฺจ นเหตุํ วิปาก\-ธมฺม\-ธมฺมญฺจ ปฏิจฺจ นเหตุ วิปาก\-ธมฺม\-ธมฺโม อุปฺปชฺชติ นกมฺมปจฺจยา. (1) (สํขิตฺตํ.)}

\proseitem{127}{นเหตุยา เอกํ, นอธิปติยา นว, นปุเรชาเต นว, นปจฺฉาชาเต นว, นอาเสวเน นว, นกมฺเม ตีณิ, นวิปาเก นว, นวิปฺปยุตฺเต นว. (สํขิตฺตํ. ปจฺจนียํ.)}

\prose{เหตุปจฺจยา นอธิปติยา นว. (สํขิตฺตํ. อนุโลม\-ปจฺจนียํ.)}

\prose{นเหตุปจฺจยา อารมฺมเณ เอกํ. (สํขิตฺตํ. ปจฺจนียา\-นุโลมํ.)}

\prose{(สหชาตวาโรปิ ปจฺจยวาโรปิ นิสฺสยวาโรปิ สํสฏฺฐ\-วาโรปิ สมฺปยุตฺตวาโรปิ ปฏิจฺจ\-วาร\-สทิสา วิตฺถาเรตพฺพา.)}

\csromanmarkh{2. วิปาก\-ธมฺมปท}\csromanmarkhii{7. ปญฺหาวาร}\csromanheadernomark{2}{2. วิปาก\-ธมฺมปท 7. ปญฺหาวาร}
\prose{ปจฺจยจตุกฺกเหตฺวาทิ}

\proseitem{128}{เหตุ วิปาก\-ธมฺม\-ธมฺโม เหตุสฺส วิปาก\-ธมฺม\-ธมฺมสฺส เหตุปจฺจเยน ปจฺจโย. ตีณิ.}

\proseitem{129}{เหตุ วิปาก\-ธมฺม\-ธมฺโม เหตุสฺส วิปาก\-ธมฺม\-ธมฺมสฺส อารมฺมณ\-ปจฺจเยน ปจฺจโย. นว.}

\proseitem{130}{\csromanfolio{222}เหตุ วิปาก\-ธมฺม\-ธมฺโม เหตุสฺส วิปาก\-ธมฺม\-ธมฺมสฺส อธิปติ\-ปจฺจเยน ปจฺจโย. อารมฺมณา\-ธิปติ, สหชาตา\-ธิปติ. นว.}

\proseitem{131}{เหตุ วิปาก\-ธมฺม\-ธมฺโม เหตุสฺส วิปาก\-ธมฺม\-ธมฺมสฺส อนนฺตร\-ปจฺจเยน ปจฺจโย ฯเปฯ อุปนิสฺสย\-ปจฺจเยน ปจฺจโย. อารมฺมณู\-ปนิสฺสโย, อนนฺตรู\-ปนิสฺสโย, ปกตู\-ปนิสฺสโย. นว. อาเสวน\-ปจฺจเยน ปจฺจโย. นว.}

\proseitem{132}{วิปาก\-ธมฺม\-ธมฺโม นเหตุสฺส วิปาก\-ธมฺม\-ธมฺมสฺส กมฺมปจฺจเยน ปจฺจโย. ตีณิ.}

\prose{วิปาก\-ธมฺม\-ธมฺโม นเหตุสฺส วิปาก\-ธมฺม\-ธมฺมสฺส อาหารปจฺจเยน ปจฺจโย. (สํขิตฺตํ.)}

\proseitem{133}{เหตุยา ตีณิ, อารมฺมเณ นว, อธิปติยา นว, อนนฺตเร นว, สมนนฺตเร นว, สหชาเต นว, อญฺญมญฺเญ นว, นิสฺสเย นว, อุปนิสฺสเย นว, อาเสวเน นว, กมฺเม ตีณิ, อาหาเร ตีณิ, อินฺทฺริเย นว, ฌาเน ตีณิ, มคฺเค นว, สมฺปยุตฺเต นว, อตฺถิยา นว, นตฺถิยา นว, วิคเต นว, อวิคเต นว. (สํขิตฺตํ. อนุโลมํ.)}

\proseitem{134}{เหตุ วิปาก\-ธมฺม\-ธมฺโม เหตุสฺส วิปาก\-ธมฺม\-ธมฺมสฺส อารมฺมณ\-ปจฺจเยน ปจฺจโย. สหชาต\-ปจฺจเยน ปจฺจโย. อุปนิสฺสย\-ปจฺจเยน ปจฺจโย. (สํขิตฺตํ.)}

\proseitem{135}{นเหตุยา นว, นอารมฺมเณ นว. (สํขิตฺตํ. ปจฺจนียํ.)}

\prose{เหตุปจฺจยา นอารมฺมเณ ตีณิ. (สํขิตฺตํ. อนุโลม\-ปจฺจนียํ.)}

\prose{นเหตุปจฺจยา อารมฺมเณ นว. (สํขิตฺตํ. ปจฺจนียา\-นุโลมํ.)}

\prose{(ยถา กุสลตฺติเก ปญฺหาวารสฺส อนุโลมมฺปิ ปจฺจนียมฺปิ อนุโลมปจฺจนียมฺปิ ปจฺจนียานุโลมมฺปิ คณิตํ, เอวํ คเณตพฺพํ.)}

\vspace{\tipitakaparskip}
\csromancenter{เหตุทุก 3. วิปากตฺติก 3. \textbf{ เนววิปากนวิปากธมฺมปท 1.} ปฏิจฺจวาร}
\prose{\csromanfolio{223}ปจฺจยจตุกฺกเหตุ}

\proseitem{136}{เหตุํ เนว\-วิปาก\-น\-วิปากธมฺม\-ธมฺมํ ปฏิจฺจ เหตุ เนว\-วิปาก\-น\-วิปาก\-ธมฺม\-ธมฺโม อุปฺปชฺชติ เหตุปจฺจยา. ตีณิ.}

\prose{นเหตุํ เนว\-วิปาก\-น\-วิปากธมฺม\-ธมฺมํ ปฏิจฺจ นเหตุ เนว\-วิปาก\-น\-วิปาก\-ธมฺม\-ธมฺโม อุปฺปชฺชติ เหตุปจฺจยา. ตีณิ.}

\prose{เหตุํ เนว\-วิปาก\-น\-วิปากธมฺม\-ธมฺมญฺจ นเหตุํ เนว\-วิปาก\-น\-วิปากธมฺม\-ธมฺมญฺจ ปฏิจฺจ เหตุ เนว\-วิปาก\-น\-วิปาก\-ธมฺม\-ธมฺโม อุปฺปชฺชติ เหตุปจฺจยา. ตีณิ. (สํขิตฺตํ.)}

\proseitem{137}{เหตุยา นว, อารมฺมเณ นว, อธิปติยา นว, อนนฺตเร นว, สมนนฺตเร นว, สหชาเต นว, อญฺญมญฺเญ นว, นิสฺสเย นว, อุปนิสฺสเย นว, ปุเรชาเต นว, อาเสวเน นว, กมฺเม นว, วิปาเก เอกํ, อาหาเร นว, อินฺทฺริเย นว, ฌาเน นว, มคฺเค นว, สมฺปยุตฺเต นว, วิปฺปยุตฺเต นว, อตฺถิยา นว, นตฺถิยา นว, วิคเต นว, อวิคเต นว.}

\csromanheader{2}{\textbf{นเหตฺวาทิ}}
\proseitem{138}{นเหตุํ เนว\-วิปาก\-น\-วิปากธมฺม\-ธมฺมํ ปฏิจฺจ นเหตุ เนว\-วิปาก\-น\-วิปาก\-ธมฺม\-ธมฺโม อุปฺปชฺชติ นเหตุปจฺจยา. (1)}

\proseitem{139}{เหตุํ เนว\-วิปาก\-น\-วิปากธมฺม\-ธมฺมํ ปฏิจฺจ นเหตุ เนว\-วิปาก\-น\-วิปาก\-ธมฺม\-ธมฺโม อุปฺปชฺชติ นอารมฺมณ\-ปจฺจยา. (1)}

\prose{นเหตุํ เนว\-วิปาก\-น\-วิปากธมฺม\-ธมฺมํ ปฏิจฺจ นเหตุ เนว\-วิปาก\-น\-วิปาก\-ธมฺม\-ธมฺโม อุปฺปชฺชติ นอารมฺมณ\-ปจฺจยา. (1)}

\prose{เหตุํ เนว\-วิปาก\-น\-วิปากธมฺม\-ธมฺมญฺจ นเหตุํ เนว\-วิปาก\-น\-วิปากธมฺม\-ธมฺมญฺจ ปฏิจฺจ นเหตุ เนว\-วิปาก\-น\-วิปาก\-ธมฺม\-ธมฺโม อุปฺปชฺชติ นอารมฺมณ\-ปจฺจยา. (1)}

\proseitem{140}{เหตุํ เนว\-วิปาก\-น\-วิปากธมฺม\-ธมฺมํ ปฏิจฺจ เหตุ เนว\-วิปาก\-น\-วิปาก\-ธมฺม\-ธมฺโม อุปฺปชฺชติ นอธิปติ\-ปจฺจยา. นว ฯเปฯ.}

\proseitem{141}{\csromanfolio{224}เหตุํ เนว\-วิปาก\-น\-วิปากธมฺม\-ธมฺมํ ปฏิจฺจ นเหตุ เนว\-วิปาก\-น\-วิปาก\-ธมฺม\-ธมฺโม อุปฺปชฺชติ นกมฺมปจฺจยา. ตีณิ ฯเปฯ.}

\proseitem{142}{นเหตุํ เนว\-วิปาก\-น\-วิปากธมฺม\-ธมฺมํ ปฏิจฺจ นเหตุ เนว\-วิปาก\-น\-วิปาก\-ธมฺม\-ธมฺโม อุปฺปชฺชติ นอาหารปจฺจยา. (1) (สํขิตฺตํ.)}

\proseitem{143}{นเหตุยา เอกํ, นอารมฺมเณ ตีณิ, นอธิปติยา นว, นอนนฺตเร ตีณิ, นสมนนฺตเร ตีณิ, นอญฺญมญฺเญ ตีณิ, นอุปนิสฺสเย ตีณิ, นปุเรชาเต นว, นปจฺฉาชาเต นว, นอาเสวเน นว, นกมฺเม ตีณิ, นวิปาเก นว, นอาหาเร เอกํ, นอินฺทฺริเย เอกํ, นฌาเน เอกํ, นมคฺเค เอกํ, นสมฺปยุตฺเต ตีณิ, นวิปฺปยุตฺเต นว, โนนตฺถิยา ตีณิ, โนวิคเต ตีณิ. (สํขิตฺตํ. ปจฺจนียํ.)}

\prose{เหตุปจฺจยา นอารมฺมเณ ตีณิ. (สํขิตฺตํ. อนุโลม\-ปจฺจนียํ.)}

\prose{นเหตุปจฺจยา อารมฺมเณ เอกํ. (สํขิตฺตํ. ปจฺจนียา\-นุโลมํ.)}

\prose{(สหชาตวาโรปิ ปจฺจยวาโรปิ นิสฺสยวาโรปิ สํสฏฺฐ\-วาโรปิ สมฺปยุตฺตวาโรปิ ปฏิจฺจ\-วาร\-สทิสา วิตฺถาเรตพฺพา.)}

\csromanmarkh{3. เนววิปากนวิปากธมฺมปท}\csromanmarkhii{7. ปญฺหาวาร}\csromanheadernomark{2}{3. \textbf{เนววิปากนวิปากธมฺมปท 7. ปญฺหาวาร}}
\prose{ปจฺจยจตุกฺกเหตฺวาทิ}

\proseitem{144}{เหตุ เนว\-วิปาก\-น\-วิปาก\-ธมฺม\-ธมฺโม เหตุสฺส เนว\-วิปาก\-น\-วิปากธมฺม\-ธมฺมสฺส เหตุปจฺจเยน ปจฺจโย. ตีณิ ฯเปฯ.}

\proseitem{145}{เหตุ เนว\-วิปาก\-น\-วิปาก\-ธมฺม\-ธมฺโม เหตุสฺส เนว\-วิปาก\-น\-วิปากธมฺม\-ธมฺมสฺส อธิปติ\-ปจฺจเยน ปจฺจโย. สหชาตา\-ธิปติ. ตีณิ.}

\prose{เนว\-วิปาก\-น\-วิปาก\-ธมฺม\-ธมฺโม นเหตุสฺส เนว\-วิปาก\-น\-วิปากธมฺม\-ธมฺมสฺส อธิปติ\-ปจฺจเยน ปจฺจโย. อารมฺมณา\-ธิปติ, สหชาตา\-ธิปติ. ตีณิ ฯเปฯ.}

\csromanmarkh{1. เหตุทุก}\csromanmarkhii{3. วิปากตฺติก ปุเรชาตาทิ}\csromanheadernomark{2}{1. เหตุทุก 3. วิปากตฺติก \textbf{ปุเรชาตาทิ}}
\proseitem{146}{\csromanfolio{225}เนว\-วิปาก\-น\-วิปาก\-ธมฺม\-ธมฺโม นเหตุสฺส เนว\-วิปาก\-น\-วิปากธมฺม\-ธมฺมสฺส ปุเรชาต\-ปจฺจเยน ปจฺจโย. อารมฺมณ\-ปุเรชาตํ, วตฺถุ\-ปุเรชาตํ. ตีณิ.}

\proseitem{147}{เหตุ เนว\-วิปาก\-น\-วิปาก\-ธมฺม\-ธมฺโม นเหตุสฺส เนว\-วิปาก\-น\-วิปากธมฺม\-ธมฺมสฺส ปจฺฉาชาต\-ปจฺจเยน ปจฺจโย. (1)}

\prose{เนว\-วิปาก\-น\-วิปาก\-ธมฺม\-ธมฺโม นเหตุสฺส เนว\-วิปาก\-น\-วิปากธมฺม\-ธมฺมสฺส ปจฺฉาชาต\-ปจฺจเยน ปจฺจโย. (1)}

\prose{เหตุ เนว\-วิปาก\-น\-วิปาก\-ธมฺม\-ธมฺโม จ นเหตุ เนว\-วิปาก\-น\-วิปาก\-ธมฺม\-ธมฺโม จ นเหตุสฺส เนว\-วิปาก\-น\-วิปากธมฺม\-ธมฺมสฺส ปจฺฉาชาต\-ปจฺจเยน ปจฺจโย. (1) อาเสวน\-ปจฺจเยน ปจฺจโย. นว.}

\proseitem{148}{เนว\-วิปาก\-น\-วิปาก\-ธมฺม\-ธมฺโม นเหตุสฺส เนว\-วิปาก\-น\-วิปากธมฺม\-ธมฺมสฺส กมฺมปจฺจเยน ปจฺจโย. ตีณิ. (สํขิตฺตํ.)}

\proseitem{149}{เหตุยา ตีณิ, อารมฺมเณ นว, อธิปติยา ฉ, อนนฺตเร นว, สมนนฺตเร นว, สหชาเต นว, อญฺญมญฺเญ นว, นิสฺสเย นว, อุปนิสฺสเย นว, ปุเรชาเต ตีณิ, ปจฺฉาชาเต ตีณิ, อาเสวเน นว, กมฺเม ตีณิ, อาหาเร ตีณิ, อินฺทฺริเย นว, ฌาเน ตีณิ, มคฺเค นว, สมฺปยุตฺเต นว, วิปฺปยุตฺเต ปญฺจ, อตฺถิยา นว, นตฺถิยา นว, วิคเต นว, อวิคเต นว.}

\proseitem{150}{เหตุ เนว\-วิปาก\-น\-วิปาก\-ธมฺม\-ธมฺโม เหตุสฺส เนว\-วิปาก\-น\-วิปากธมฺม\-ธมฺมสฺส อารมฺมณ\-ปจฺจเยน ปจฺจโย. สหชาต\-ปจฺจเยน ปจฺจโย. อุปนิสฺสย\-ปจฺจเยน ปจฺจโย. (สํขิตฺตํ.)}

\proseitem{151}{นเหตุยา นว, นอารมฺมเณ นว, นอธิปติยา นว. (สํขิตฺตํ.)}

\prose{เหตุปจฺจยา นอารมฺมเณ ตีณิ. (สํขิตฺตํ.)}

\prose{นเหตุปจฺจยา อารมฺมเณ นว. (สํขิตฺตํ.)}

\prosewithcloser{\csromanfolio{226}(ยถา กุสลตฺติเก ปญฺหาวารสฺส อนุโลมมฺปิ ปจฺจนียมฺปิ อนุโลมปจฺจนียมฺปิ ปจฺจนียานุโลมมฺปิ คณิตํ, เอวํ คเณตพฺพํ.)}{เหตุ\-ทุก\-วิปากตฺติกํ นิฏฺฐิตํ.}
\csromanmatikaanchor{mtk.35}
\csromantocmarkref{section}{4. อุปาทินฺนตฺติก}{mtk.35}
\bookmark[dest={mtk.35},level=1]{4. อุปาทินฺนตฺติก}
\csromanmarkh{1. เหตุทุก}\csromanmarkhii{4. อุปาทินฺนตฺติก}\csromanheadernomark{1}{1. \textbf{เหตุทุก 4. อุปาทินฺนตฺติก}}
\csromanmarkh{1. อุปาทินฺนุ\-ปาทานิยปท}\csromanmarkhii{1. ปฏิจฺจวาร}\csromanheadernomark{2}{1. อุปาทินฺนุ\-ปาทานิยปท 1. ปฏิจฺจวาร}
\prose{ปจฺจยจตุกฺกเหตุ}

\proseitem{152}{เหตุํ อุปาทินฺนุ\-ปาทานิยํ ธมฺมํ ปฏิจฺจ เหตุ อุปาทินฺนุ\-ปาทานิโย ธมฺโม อุปฺปชฺชติ เหตุปจฺจยา. ตีณิ.}

\prose{นเหตุํ อุปาทินฺนุ\-ปาทานิยํ ธมฺมํ ปฏิจฺจ นเหตุ อุปาทินฺนุ\-ปาทานิโย ธมฺโม อุปฺปชฺชติ เหตุปจฺจยา. ตีณิ.}

\prose{เหตุํ อุปาทินฺนุ\-ปาทานิยญฺจ นเหตุํ อุปาทินฺนุ\-ปาทานิยญฺจ ธมฺมํ ปฏิจฺจ เหตุ อุปาทินฺนุ\-ปาทานิโย ธมฺโม อุปฺปชฺชติ เหตุปจฺจยา. ตีณิ. (สํขิตฺตํ.)}

\proseitem{153}{เหตุยา นว, อารมฺมเณ นว, อนนฺตเร นว, สมนนฺตเร นว ฯเปฯ กมฺเม นว, วิปาเก นว ฯเปฯ อวิคเต นว. (สํขิตฺตํ.)}

\csromanheader{2}{\textbf{นเหตฺวาทิ}}
\proseitem{154}{นเหตุํ อุปาทินฺนุ\-ปาทานิยํ ธมฺมํ ปฏิจฺจ นเหตุ อุปาทินฺนุ\-ปาทานิโย ธมฺโม อุปฺปชฺชติ นเหตุปจฺจยา. (1)}

\proseitem{155}{เหตุํ อุปาทินฺนุ\-ปาทานิยํ ธมฺมํ ปฏิจฺจ นเหตุ อุปาทินฺนุ\-ปาทานิโย ธมฺโม อุปฺปชฺชติ นอารมฺมณ\-ปจฺจยา. (1)}

\prose{นเหตุํ อุปาทินฺนุ\-ปาทานิยํ ธมฺมํ ปฏิจฺจ นเหตุ อุปาทินฺนุ\-ปาทานิโย ธมฺโม อุปฺปชฺชติ นอารมฺมณ\-ปจฺจยา. (1)}

\prose{เหตุํ อุปาทินฺนุ\-ปาทานิยญฺจ นเหตุํ อุปาทินฺนุ\-ปาทานิยญฺจ ธมฺมํ ปฏิจฺจ นเหตุ อุปาทินฺนุ\-ปาทานิโย ธมฺโม อุปฺปชฺชติ นอารมฺมณ\-ปจฺจยา. (1)}

\csromanmarkh{1. เหตุทุก}\csromanmarkhii{4. อุปาทินฺนตฺติก}\csromanheadernomark{2}{1. เหตุทุก 4. อุปาทินฺนตฺติก}
\proseitem{156}{\csromanfolio{227}เหตุํ อุปาทินฺนุ\-ปาทานิยํ ธมฺมํ ปฏิจฺจ เหตุ อุปาทินฺนุ\-ปาทานิโย ธมฺโม อุปฺปชฺชติ นอธิปติ\-ปจฺจยา. นว ฯเปฯ.}

\proseitem{157}{นเหตุํ อุปาทินฺนุ\-ปาทานิยํ ธมฺมํ ปฏิจฺจ นเหตุ อุปาทินฺนุ\-ปาทานิโย ธมฺโม อุปฺปชฺชติ นวิปาก\-ปจฺจยา. นอาหารปจฺจยา. (สํขิตฺตํ.)}

\proseitem{158}{นเหตุยา เอกํ, นอารมฺมเณ ตีณิ, นอธิปติยา นว, นอนนฺตเร ตีณิ, นสมนนฺตเร ตีณิ, นอญฺญมญฺเญ ตีณิ, นอุปนิสฺสเย ตีณิ, นปุเรชาเต นว, นปจฺฉาชาเต นว, นอาเสวเน นว, นวิปาเก เอกํ, นอาหาเร เอกํ, นอินฺทฺริเย เอกํ, นฌาเน เอกํ, นมคฺเค เอกํ, นสมฺปยุตฺเต ตีณิ, นวิปฺปยุตฺเต นว, โนนตฺถิยา ตีณิ, โนวิคเต ตีณิ.}

\prose{เหตุปจฺจยา นอารมฺมเณ ตีณิ. (สํขิตฺตํ.)}

\prose{นเหตุปจฺจยา อารมฺมเณ เอกํ. (สํขิตฺตํ.)}

\prose{(สหชาตวาโรปิ ปจฺจยวาโรปิ นิสฺสยวาโรปิ สํสฏฺฐ\-วาโรปิ สมฺปยุตฺตวาโรปิ ปฏิจฺจ\-วาร\-สทิสา วิตฺถาเรตพฺพา.)}

\csromanmarkh{1. อุปาทินฺนุ\-ปาทานิยปท}\csromanmarkhii{7. ปญฺหาวาร}\csromanheadernomark{2}{1. อุปาทินฺนุ\-ปาทานิยปท 7. ปญฺหาวาร}
\prose{ปจฺจยจตุกฺกเหตุ อารมฺมณ}

\proseitem{159}{เหตุ อุปาทินฺนุ\-ปาทานิโย ธมฺโม เหตุสฺส อุปาทินฺนุ\-ปาทานิยสฺส ธมฺมสฺส เหตุปจฺจเยน ปจฺจโย. ตีณิ.}

\proseitem{160}{เหตุ อุปาทินฺนุ\-ปาทานิโย ธมฺโม เหตุสฺส อุปาทินฺนุ\-ปาทานิยสฺส ธมฺมสฺส อารมฺมณ\-ปจฺจเยน ปจฺจโย. ตีณิ.}

\prose{อุปาทินฺนุ\-ปาทานิโย ธมฺโม นเหตุสฺส อุปาทินฺนุ\-ปาทานิยสฺส ธมฺมสฺส อารมฺมณ\-ปจฺจเยน ปจฺจโย. ตีณิ.}

\prose{เหตุ อุปาทินฺนุ\-ปาทานิโย จ นเหตุ อุปาทินฺนุ\-ปาทานิโย จ ธมฺมา เหตุสฺส อุปาทินฺนุ\-ปาทานิยสฺส ธมฺมสฺส อารมฺมณ\-ปจฺจเยน ปจฺจโย. ตีณิ. (สํขิตฺตํ.)}

\proseitem{161}{\csromanfolio{228}เหตุยา ตีณิ, อารมฺมเณ นว, อนนฺตเร นว, สมนนฺตเร นว, สหชาเต นว, อญฺญมญฺเญ นว, นิสฺสเย นว, อุปนิสฺสเย นว, ปุเรชาเต ตีณิ, ปจฺฉาชาเต ตีณิ, กมฺเม ตีณิ, วิปาเก นว, อาหาเร ตีณิ, อินฺทฺริเย นว, ฌาเน ตีณิ, มคฺเค นว, สมฺปยุตฺเต นว, วิปฺปยุตฺเต ปญฺจ, อตฺถิยา นว, นตฺถิยา นว, วิคเต นว, อวิคเต นว.}

\proseitem{162}{เหตุ อุปาทินฺนุ\-ปาทานิโย ธมฺโม เหตุสฺส อุปาทินฺนุ\-ปาทานิยสฺส ธมฺมสฺส อารมฺมณ\-ปจฺจเยน ปจฺจโย. สหชาต\-ปจฺจเยน ปจฺจโย. อุปนิสฺสย\-ปจฺจเยน ปจฺจโย. ตีณิ.}

\prose{อุปาทินฺนุ\-ปาทานิโย ธมฺโม นเหตุสฺส อุปาทินฺนุ\-ปาทานิยสฺส ธมฺมสฺส อารมฺมณ\-ปจฺจเยน ปจฺจโย. สหชาต\-ปจฺจเยน ปจฺจโย. อุปนิสฺสย\-ปจฺจเยน ปจฺจโย. ปุเรชาต\-ปจฺจเยน ปจฺจโย. ปจฺฉาชาต\-ปจฺจเยน ปจฺจโย. อาหารปจฺจเยน ปจฺจโย. อินฺทฺริย\-ปจฺจเยน ปจฺจโย. (สํขิตฺตํ.)}

\proseitem{163}{นเหตุยา นว, นอารมฺมเณ นว, นอธิปติยา นว. (สํขิตฺตํ.)}

\prose{เหตุปจฺจยา นอารมฺมเณ ตีณิ. (สํขิตฺตํ.)}

\prose{นเหตุปจฺจยา อารมฺมเณ นว. (สํขิตฺตํ.)}

\prose{(ยถา กุสลตฺติเก ปญฺหาวารสฺส อนุโลมมฺปิ ปจฺจนียมฺปิ อนุโลมปจฺจนียมฺปิ ปจฺจนียานุโลมมฺปิ คณิตํ, เอวํ คเณตพฺพํ.)}

\csromanmarkh{2. อนุปาทินฺนุ\-ปาทานิยปท}\csromanmarkhii{1. ปฏิจฺจวาร}\csromanheadernomark{2}{2. อนุปาทินฺนุ\-ปาทานิยปท 1. ปฏิจฺจวาร}
\prose{ปจฺจยจตุกฺกเหตุ}

\proseitem{164}{เหตุํ อนุปาทินฺนุ\-ปาทานิยํ ธมฺมํ ปฏิจฺจ เหตุ อนุปาทินฺนุ\-ปาทานิโย ธมฺโม อุปฺปชฺชติ เหตุปจฺจยา. ตีณิ.}

\prose{นเหตุํ อนุปาทินฺนุ\-ปาทานิยํ ธมฺมํ ปฏิจฺจ นเหตุ อนุปาทินฺนุ\-ปาทานิโย ธมฺโม อุปฺปชฺชติ เหตุปจฺจยา. ตีณิ.}

\csromanmarkh{1. เหตุทุก}\csromanmarkhii{4. อุปาทินฺนตฺติก}\csromanheadernomark{2}{1. เหตุทุก 4. อุปาทินฺนตฺติก}
\prose{\csromanfolio{229}เหตุํ อนุปาทินฺนุ\-ปาทานิยญฺจ นเหตุํ อนุปาทินฺนุ\-ปาทานิยญฺจ ธมฺมํ ปฏิจฺจ เหตุ อนุปาทินฺนุ\-ปาทานิโย ธมฺโม อุปฺปชฺชติ เหตุปจฺจยา. ตีณิ. (สํขิตฺตํ.)}

\proseitem{165}{เหตุยา นว, อารมฺมเณ นว, อธิปติยา นว, อนนฺตเร นว, สมนนฺตเร นว, สหชาเต นว, อญฺญมญฺเญ นว, นิสฺสเย นว, อุปนิสฺสเย นว, ปุเรชาเต นว, อาเสวเน นว, กมฺเม นว, วิปาเก เอกํ, อาหาเร นว, อินฺทฺริเย นว, สมฺปยุตฺเต นว, วิปฺปยุตฺเต นว, อตฺถิยา นว, นตฺถิยา นว, วิคเต นว, อวิคเต นว.}

\prose{นเหตุ นอารมฺมณ}

\proseitem{166}{นเหตุํ อนุปาทินฺนุ\-ปาทานิยํ ธมฺมํ ปฏิจฺจ นเหตุ อนุปาทินฺนุ\-ปาทานิโย ธมฺโม อุปฺปชฺชติ นเหตุปจฺจยา. เทฺว. (2)}

\proseitem{167}{เหตุํ อนุปาทินฺนุ\-ปาทานิยํ ธมฺมํ ปฏิจฺจ นเหตุ อนุปาทินฺนุ\-ปาทานิโย ธมฺโม อุปฺปชฺชติ นอารมฺมณ\-ปจฺจยา. (1)}

\prose{นเหตุํ อนุปาทินฺนุ\-ปาทานิยํ ธมฺมํ ปฏิจฺจ นเหตุ อนุปาทินฺนุ\-ปาทานิโย ธมฺโม อุปฺปชฺชติ นอารมฺมณ\-ปจฺจยา. (1)}

\prose{เหตุํ อนุปาทินฺนุ\-ปาทานิยญฺจ นเหตุํ อนุปาทินฺนุ\-ปาทานิยญฺจ ธมฺมํ ปฏิจฺจ นเหตุ อนุปาทินฺนุ\-ปาทานิโย ธมฺโม อุปฺปชฺชติ นอารมฺมณ\-ปจฺจยา. (1) (สํขิตฺตํ.)}

\proseitem{168}{นเหตุยา เทฺว, นอารมฺมเณ ตีณิ, นอธิปติยา นว, นอนนฺตเร ตีณิ, นสมนนฺตเร ตีณิ, นอญฺญมญฺเญ ตีณิ, นอุปนิสฺสเย ตีณิ, นปุเรชาเต นว, นปจฺฉาชาเต นว, นอาเสวเน นว, นกมฺเม ตีณิ, นวิปาเก นว, นอาหาเร เอกํ, นอินฺทฺริเย เอกํ, นฌาเน เอกํ, นมคฺเค เอกํ, นสมฺปยุตฺเต ตีณิ, นวิปฺปยุตฺเต นว, โนนตฺถิยา ตีณิ, โนวิคเต ตีณิ. (สํขิตฺตํ.)}

\prose{เหตุปจฺจยา นอารมฺมเณ ตีณิ. (สํขิตฺตํ.)}

\prose{นเหตุปจฺจยา อารมฺมเณ เทฺว. (สํขิตฺตํ.)}

\prose{(สหชาตวาโรปิ ปจฺจยวาโรปิ นิสฺสยวาโรปิ สํสฏฺฐ\-วาโรปิ สมฺปยุตฺตวาโรปิ ปฏิจฺจ\-วาร\-สทิสา วิตฺถาเรตพฺพา.)}

\csromanmarkh{ธมฺมานุโลม ทุกติก\-ปฏฺฐาน\-ปาฬิ}\csromanmarkhii{2. อนุปาทินฺนุ\-ปาทานิยปท 7. ปญฺหาวาร}\csromanheadernomark{2}{ธมฺมานุโลม ทุกติก\-ปฏฺฐาน\-ปาฬิ 2. อนุปาทินฺนุ\-ปาทานิยปท 7. ปญฺหาวาร}
\prose{\csromanfolio{230}ปจฺจยจตุกฺกเหตฺวาทิ}

\proseitem{169}{เหตุ อนุปาทินฺนุ\-ปาทานิโย ธมฺโม เหตุสฺส อนุปาทินฺนุ\-ปาทานิยสฺส ธมฺมสฺส เหตุปจฺจเยน ปจฺจโย. ตีณิ.}

\proseitem{170}{เหตุ อนุปาทินฺนุ\-ปาทานิโย ธมฺโม เหตุสฺส อนุปาทินฺนุ\-ปาทานิยสฺส ธมฺมสฺส อารมฺมณ\-ปจฺจเยน ปจฺจโย. ตีณิ.}

\prose{อนุปาทินฺนุ\-ปาทานิโย ธมฺโม นเหตุสฺส อนุปาทินฺนุ\-ปาทานิยสฺส ธมฺมสฺส อารมฺมณ\-ปจฺจเยน ปจฺจโย. นว.}

\proseitem{171}{เหตุ อนุปาทินฺนุ\-ปาทานิโย ธมฺโม เหตุสฺส อนุปาทินฺนุ\-ปาทานิยสฺส ธมฺมสฺส อธิปติ\-ปจฺจเยน ปจฺจโย. อารมฺมณา\-ธิปติ, สหชาตา\-ธิปติ. ตีณิ.}

\prose{อนุปาทินฺนุ\-ปาทานิโย ธมฺโม นเหตุสฺส อนุปาทินฺนุ\-ปาทานิยสฺส ธมฺมสฺส อธิปติ\-ปจฺจเยน ปจฺจโย. อารมฺมณา\-ธิปติ, สหชาตา\-ธิปติ. นว ฯเปฯ.}

\proseitem{172}{เหตุ อนุปาทินฺนุ\-ปาทานิโย ธมฺโม เหตุสฺส อนุปาทินฺนุ\-ปาทานิยสฺส ธมฺมสฺส อุปนิสฺสย\-ปจฺจเยน ปจฺจโย. อารมฺมณู\-ปนิสฺสโย, อนนฺตรู\-ปนิสฺสโย, ปกตู\-ปนิสฺสโย. นว.}

\proseitem{173}{อนุปาทินฺนุ\-ปาทานิโย ธมฺโม นเหตุสฺส อนุปาทินฺนุ\-ปาทานิยสฺส ธมฺมสฺส ปุเรชาต\-ปจฺจเยน ปจฺจโย. อารมฺมณ\-ปุเรชาตํ. ตีณิ. (สํขิตฺตํ.)}

\proseitem{174}{เหตุยา ตีณิ, อารมฺมเณ นว, อธิปติยา นว, อนนฺตเร นว, สมนนฺตเร นว, สหชาเต นว, อญฺญมญฺเญ นว, นิสฺสเย นว, อุปนิสฺสเย นว, ปุเรชาเต ตีณิ, ปจฺฉาชาเต ตีณิ, อาเสวเน นว, กมฺเม ตีณิ, อาหาเร ตีณิ, อินฺทฺริเย นว, ฌาเน ตีณิ, มคฺเค นว, สมฺปยุตฺเต นว, วิปฺปยุตฺเต ตีณิ, อตฺถิยา นว, นตฺถิยา นว, วิคเต นว, อวิคเต นว.}

\csromanmarkh{1. เหตุทุก}\csromanmarkhii{4. อุปาทินฺนตฺติก ปจฺจนียุ\-ทฺธาร}\csromanheadernomark{2}{1. เหตุทุก 4. อุปาทินฺนตฺติก ปจฺจนียุ\-ทฺธาร}
\proseitem{175}{\csromanfolio{231}เหตุ อนุปาทินฺนุ\-ปาทานิโย ธมฺโม เหตุสฺส อนุปาทินฺนุ\-ปาทานิยสฺส ธมฺมสฺส อารมฺมณ\-ปจฺจเยน ปจฺจโย. สหชาต\-ปจฺจเยน ปจฺจโย. อุปนิสฺสย\-ปจฺจเยน ปจฺจโย. (สํขิตฺตํ.)}

\proseitem{176}{นเหตุยา นว, นอารมฺมเณ นว. (สํขิตฺตํ.)}

\prose{เหตุปจฺจยา นอารมฺมเณ ตีณิ. (สํขิตฺตํ.)}

\prose{นเหตุปจฺจยา อารมฺมเณ นว. (สํขิตฺตํ.)}

\prose{(ยถา กุสลตฺติเก ปญฺหาวารสฺส อนุโลมมฺปิ ปจฺจนียมฺปิ อนุโลมปจฺจนียมฺปิ ปจฺจนียานุโลมมฺปิ คณิตํ, เอวํ คเณตพฺพํ.)}

\prose{3. อนุปาทินฺน\-อนุปาทานิยปท 1. ปฏิจฺจวาร}

\prose{ปจฺจยจตุกฺกเหตุ}

\proseitem{177}{เหตุํ อนุปาทินฺน\-อนุปาทานิยํ ธมฺมํ ปฏิจฺจ เหตุ อนุปาทินฺน\-อนุปาทานิโย ธมฺโม อุปฺปชฺชติ เหตุปจฺจยา. ตีณิ.}

\prose{นเหตุํ อนุปาทินฺน\-อนุปาทานิยํ ธมฺมํ ปฏิจฺจ นเหตุ อนุปาทินฺน\-อนุปาทานิโย ธมฺโม อุปฺปชฺชติ เหตุปจฺจยา. ตีณิ.}

\prose{เหตุํ อนุปาทินฺน\-อนุปาทานิยญฺจ นเหตุํ อนุปาทินฺน\-อนุปาทานิยญฺจ ธมฺมํ ปฏิจฺจ เหตุ อนุปาทินฺน\-อนุปาทานิโย ธมฺโม อุปฺปชฺชติ เหตุปจฺจยา. ตีณิ. (สํขิตฺตํ.)}

\proseitem{178}{เหตุยา นว, อารมฺมเณ นว, อธิปติยา นว, อนนฺตเร นว, สมนนฺตเร นว, สหชาเต นว, อญฺญมญฺเญ นว, นิสฺสเย นว, อุปนิสฺสเย นว, ปุเรชาเต นว, อาเสวเน นว, กมฺเม นว, วิปาเก นว, อาหาเร นว, อินฺทฺริเย นว, ฌาเน นว, มคฺเค นว, สมฺปยุตฺเต นว, วิปฺปยุตฺเต นว, อตฺถิยา นว, นตฺถิยา นว, วิคเต นว, อวิคเต นว.}

\csromanheader{2}{ธมฺมานุโลม ทุกติก\-ปฏฺฐาน\-ปาฬิ นอธิปติ}
\proseitem{179}{\csromanfolio{232}เหตุํ อนุปาทินฺน\-อนุปาทานิยํ ธมฺมํ ปฏิจฺจ เหตุ อนุปาทินฺน\-อนุปาทานิโย ธมฺโม อุปฺปชฺชติ นอธิปติ\-ปจฺจยา.  เหตุํ อนุปาทินฺน\-อนุปาทานิยํ ธมฺมํ ปฏิจฺจ นเหตุ อนุปาทินฺน\-อนุปาทานิโย ธมฺโม อุปฺปชฺชติ นอธิปติ\-ปจฺจยา. (2)}

\prose{นเหตุํ อนุปาทินฺน\-อนุปาทานิยํ ธมฺมํ ปฏิจฺจ นเหตุ อนุปาทินฺน\-อนุปาทานิโย ธมฺโม อุปฺปชฺชติ นอธิปติ\-ปจฺจยา.  นเหตุํ อนุปาทินฺนอนุปาทานิย ธมฺมํ ปฏิจฺจ เหตุ อนุปาทินฺน\-อนุปาทานิโย ธมฺโม อุปฺปชฺชติ นอธิปติ\-ปจฺจยา. (2)}

\prose{เหตุํ อนุปาทินฺน\-อนุปาทานิยญฺจ นเหตุํ อนุปาทินฺน\-อนุปาทานิยญฺจ ธมฺมํ ปฏิจฺจ เหตุ อนุปาทินฺน\-อนุปาทานิโย ธมฺโม อุปฺปชฺชติ นอธิปติ\-ปจฺจยา.  เหตุํ อนุปาทินฺน\-อนุปาทานิยญฺจ นเหตุํ อนุปาทินฺน\-อนุปาทานิยญฺจ ธมฺมํ ปฏิจฺจ นเหตุ อนุปาทินฺน\-อนุปาทานิโย ธมฺโม อุปฺปชฺชติ นอธิปติ\-ปจฺจยา. (2) (สํขิตฺตํ.)}

\proseitem{180}{นอธิปติยา ฉ, นปุเรชาเต นว, นปจฺฉาชาเต นว, นอาเสวเน นว, นกมฺเม ตีณิ, นวิปาเก นว, นวิปฺปยุตฺเต นว. (สํขิตฺตํ.)}

\prose{เหตุปจฺจยา นอธิปติยา ฉ. (สํขิตฺตํ.)}

\prose{นอธิปติ\-ปจฺจยา เหตุยา ฉ. (สํขิตฺตํ.)}

\prose{(สหชาตวาโรปิ ปจฺจยวาโรปิ นิสฺสยวาโรปิ สํสฏฺฐ\-วาโรปิ สมฺปยุตฺตวาโรปิ ปฏิจฺจ\-วาร\-สทิสา วิตฺถาเรตพฺพา.)}

\prose{3. อนุปาทินฺน\-อนุปาทานิยปท 7. ปญฺหาวาร}

\prose{ปจฺจยจตุกฺกเหตุ อารมฺมณ}

\proseitem{181}{เหตุ อนุปาทินฺน\-อนุปาทานิโย ธมฺโม เหตุสฺส อนุปาทินฺน\-อนุปาทานิยสฺส ธมฺมสฺส เหตุปจฺจเยน ปจฺจโย. ตีณิ.}

\csromanmatikaanchor{mtk.36}
\csromantocmarkref{section}{5. สํกิลิฏฺฐตฺติก}{mtk.36}
\bookmark[dest={mtk.36},level=1]{5. สํกิลิฏฺฐตฺติก}
\csromanmarkh{1. เหตุทุก}\csromanmarkhii{5. สํกิลิฏฺฐ\-ตฺติก}\csromanheadernomark{1}{1. เหตุทุก 5. สํกิลิฏฺฐ\-ตฺติก}
\proseitem{182}{\csromanfolio{233}อนุปาทินฺน\-อนุปาทานิโย ธมฺโม นเหตุสฺส อนุปาทินฺน\-อนุปาทานิยสฺส ธมฺมสฺส อารมฺมณ\-ปจฺจเยน ปจฺจโย. ตีณิ. (สํขิตฺตํ.)}

\proseitem{183}{เหตุยา ตีณิ, อารมฺมเณ ตีณิ, อธิปติยา ฉ, อนนฺตเร นว, สมนนฺตเร นว, สหชาเต นว, อญฺญมญฺเญ นว, นิสฺสเย นว, อุปนิสฺสเย นว, กมฺเม ตีณิ, วิปาเก นว, อาหาเร ตีณิ, อินฺทฺริเย นว, ฌาเน ตีณิ, มคฺเค นว, สมฺปยุตฺเต นว, อตฺถิยา นว, นตฺถิยา นว, วิคเต นว, อวิคเต นว.}

\proseitem{184}{เหตุ อนุปาทินฺน\-อนุปาทานิโย ธมฺโม เหตุสฺส อนุปาทินฺน\-อนุปาทานิยสฺส ธมฺมสฺส สหชาต\-ปจฺจเยน ปจฺจโย. อุปนิสฺสย\-ปจฺจเยน ปจฺจโย. (สํขิตฺตํ.)}

\proseitem{185}{นเหตุยา นว, นอารมฺมเณ นว. (สํขิตฺตํ.)}

\prose{เหตุปจฺจยา นอารมฺมเณ ตีณิ. (สํขิตฺตํ.)}

\prose{นเหตุปจฺจยา อารมฺมเณ ตีณิ. (สํขิตฺตํ.)}

\prosewithcloser{(ยถา กุสลตฺติเก ปญฺหาวารสฺส อนุโลมมฺปิ ปจฺจนียมฺปิ อนุโลมปจฺจนียมฺปิ ปจฺจนียานุโลมมฺปิ คณิตํ, เอวํ คเณตพฺพํ.)}{เหตุทุกอุปาทินฺนตฺติกํ นิฏฺฐิตํ.}
\csromanmarkh{1. เหตุทุก}\csromanmarkhii{5. สํกิลิฏฺฐ\-ตฺติก}\csromanheadernomark{2}{1. เหตุทุก 5. สํกิลิฏฺฐ\-ตฺติก}
\csromanmarkh{1. สํกิลิฏฺฐสํกิเลสิกปท}\csromanmarkhii{1. ปฏิจฺจวาร}\csromanheadernomark{2}{1. \textbf{สํกิลิฏฺฐสํกิเลสิกปท 1. ปฏิจฺจวาร}}
\prose{ปจฺจยจตุกฺกเหตุ}

\proseitem{186}{เหตุํ สํกิลิฏฺฐ\-สํกิเลสิกํ ธมฺมํ ปฏิจฺจ เหตุ สํกิลิฏฺฐ\-สํกิเลสิโก ธมฺโม อุปฺปชฺชติ เหตุปจฺจยา. ตีณิ.}

\prose{นเหตุํ สํกิลิฏฺฐ\-สํกิเลสิกํ ธมฺมํ ปฏิจฺจ นเหตุ สํกิลิฏฺฐ\-สํกิเลสิโก ธมฺโม อุปฺปชฺชติ เหตุปจฺจยา. ตีณิ.}

\prose{\csromanfolio{234}เหตุํ สํกิลิฏฺฐ\-สํกิเลสิกญฺจ นเหตุํ สํกิลิฏฺฐ\-สํกิเลสิกญฺจ ธมฺมํ ปฏิจฺจ เหตุ สํกิลิฏฺฐ\-สํกิเลสิโก ธมฺโม อุปฺปชฺชติ เหตุปจฺจยา. ตีณิ. (สํขิตฺตํ.)}

\proseitem{187}{เหตุยา นว, อารมฺมเณ นว, อธิปติยา นว ฯเปฯ กมฺเม นว, อาหาเร นว, อวิคเต นว. (สํขิตฺตํ.)}

\prose{นเหตุ นอธิปติ}

\proseitem{188}{นเหตุํ สํกิลิฏฺฐ\-สํกิเลสิกํ ธมฺมํ ปฏิจฺจ เหตุ สํกิลิฏฺฐ\-สํกิเลสิโก ธมฺโม อุปฺปชฺชติ นเหตุปจฺจยา.}

\prose{เหตุํ สํกิลิฏฺฐ\-สํกิเลสิกํ ธมฺมํ ปฏิจฺจ เหตุ สํกิลิฏฺฐ\-สํกิเลสิโก ธมฺโม อุปฺปชฺชติ นอธิปติ\-ปจฺจยา. (สํขิตฺตํ.)}

\proseitem{189}{นเหตุยา เอกํ, นอธิปติยา นว, นปุเรชาเต นว, นปจฺฉาชาเต นว, นอาเสวเน นว, นกมฺเม ตีณิ, นวิปาเก นว, นวิปฺปยุตฺเต นว. (สํขิตฺตํ.)}

\prose{เหตุปจฺจยา นอธิปติยา นว. (สํขิตฺตํ.)}

\prose{นเหตุปจฺจยา อารมฺมเณ เอกํ. (สํขิตฺตํ.)}

\prose{(สหชาตวาโรปิ ปจฺจยวาโรปิ นิสฺสยวาโรปิ สํสฏฺฐ\-วาโรปิ สมฺปยุตฺตวาโรปิ ปฏิจฺจ\-วาร\-สทิสา วิตฺถาเรตพฺพา.)}

\csromanmarkh{1. สํกิลิฏฺฐสํกิเลสิกปท}\csromanmarkhii{7. ปญฺหาวาร}\csromanheadernomark{2}{1. \textbf{สํกิลิฏฺฐสํกิเลสิกปท 7. ปญฺหาวาร}}
\prose{ปจฺจยจตุกฺกเหตฺวาทิ}

\proseitem{190}{เหตุ สํกิลิฏฺฐ\-สํกิเลสิโก ธมฺโม เหตุสฺส สํกิลิฏฺฐ\-สํกิเลสิกสฺส ธมฺมสฺส เหตุปจฺจเยน ปจฺจโย. ตีณิ.}

\prose{เหตุ สํกิลิฏฺฐ\-สํกิเลสิโก ธมฺโม เหตุสฺส สํกิลิฏฺฐ\-สํกิเลสิกสฺส ธมฺมสฺส อารมฺมณ\-ปจฺจเยน ปจฺจโย. นว.}

\prose{เหตุ สํกิลิฏฺฐ\-สํกิเลสิโก ธมฺโม เหตุสฺส สํกิลิฏฺฐ\-สํกิเลสิกสฺส ธมฺมสฺส อธิปติ\-ปจฺจเยน ปจฺจโย. อารมฺมณา\-ธิปติ. นว. (สํขิตฺตํ.)}

\csromanmarkh{1. เหตุทุก}\csromanmarkhii{5. สํกิลิฏฺฐ\-ตฺติก}\csromanheadernomark{2}{1. เหตุทุก 5. สํกิลิฏฺฐ\-ตฺติก}
\proseitem{191}{\csromanfolio{235}เหตุยา ตีณิ, อารมฺมเณ นว, อธิปติยา นว, อนนฺตเร นว, สมนนฺตเร นว, สหชาเต นว, อญฺญมญฺเญ นว, นิสฺสเย นว, อุปนิสฺสเย นว, อาเสวเน นว, กมฺเม ตีณิ, อาหาเร ตีณิ, อินฺทฺริเย ตีณิ, ฌาเน ตีณิ, มคฺเค ตีณิ, สมฺปยุตฺเต นว, อตฺถิยา นว, นตฺถิยา นว, วิคเต นว, อวิคเต นว. (สํขิตฺตํ.)}

\proseitem{192}{เหตุ สํกิลิฏฺฐ\-สํกิเลสิโก ธมฺโม เหตุสฺส สํกิลิฏฺฐ\-สํกิเลสิกสฺส ธมฺมสฺส อารมฺมณ\-ปจฺจเยน ปจฺจโย. สหชาต\-ปจฺจเยน ปจฺจโย. อุปนิสฺสย\-ปจฺจเยน ปจฺจโย. (สํขิตฺตํ.)}

\proseitem{193}{นเหตุยา นว, นอารมฺมเณ นว. (สํขิตฺตํ.)}

\prose{เหตุปจฺจยา นอารมฺมเณ ตีณิ. (สํขิตฺตํ.)}

\prose{นเหตุปจฺจยา อารมฺมเณ นว. (สํขิตฺตํ.)}

\prose{(ยถา กุสลตฺติเก ปญฺหาวารสฺส อนุโลมมฺปิ ปจฺจนียมฺปิ อนุโลมปจฺจนียมฺปิ ปจฺจนียานุโลมมฺปิ คณิตํ, เอวํ คเณตพฺพํ.)}

\csromanmarkh{2. อสํกิลิฏฺฐสํกิเลสิกปท}\csromanmarkhii{1. ปฏิจฺจวาร}\csromanheadernomark{2}{2. \textbf{อสํกิลิฏฺฐสํกิเลสิกปท 1. ปฏิจฺจวาร}}
\prose{ปจฺจยจตุกฺกเหตุ}

\proseitem{194}{เหตุํ อสํกิลิฏฺฐ\-สํกิเลสิกํ ธมฺมํ ปฏิจฺจ เหตุ อสํกิลิฏฺฐ\-สํกิเลสิโก ธมฺโม อุปฺปชฺชติ เหตุปจฺจยา. ตีณิ.}

\prose{นเหตุํ อสํกิลิฏฺฐ\-สํกิเลสิกํ ธมฺมํ ปฏิจฺจ นเหตุ อสํกิลิฏฺฐ\-สํกิเลสิโก ธมฺโม อุปฺปชฺชติ เหตุปจฺจยา. ตีณิ.}

\prose{เหตุํ อสํกิลิฏฺฐ\-สํกิเลสิกญฺจ นเหตุํ อสํกิลิฏฺฐ\-สํกิเลสิกญฺจ ธมฺมํ ปฏิจฺจ เหตุ อสํกิลิฏฺฐ\-สํกิเลสิโก ธมฺโม อุปฺปชฺชติ เหตุปจฺจยา. ตีณิ. (สํขิตฺตํ.)}

\proseitem{195}{เหตุยา นว, อารมฺมเณ นว, อธิปติยา นว ฯเปฯ อวิคเต นว. (สํขิตฺตํ.)}

\csromanpalirecto
\gambhira{ธมฺมานุโลม ทุกติก\-ปฏฺฐาน\-ปาฬิ นเหตุ}
\proseitem{196}{\csromanfolio{236}นเหตุํ อสํกิลิฏฺฐ\-สํกิเลสิกํ ธมฺมํ ปฏิจฺจ นเหตุ อสํกิลิฏฺฐ\-สํกิเลสิโก ธมฺโม อุปฺปชฺชติ นเหตุปจฺจยา. (สํขิตฺตํ.)}

\vspace{\tipitakaparskip}
\csromancenter{\textbf{นเหตุ}ยา เอกํ, นอารมฺมเณ ตีณิ, นอธิปติยา นว, นอนนฺตเร ตีณิ, นสมนนฺตเร ตีณิ, นอญฺญมญฺเญ ตีณิ, นฺเอาปนิสฺสเย ตีณิ, นปุเรชาเต นว, นปจฺฉาชาเต นว, นอาเสวเน นว, นกมฺเม ตีณิ, นวิปาเก นว, นอาหาเร เอกํ, ไนนฺทฺริเย เอกํ, นฌาเน เอกํ, นมคฺเค เอกํ, นสมฺปยุตฺเต ตีณิ, นวิปฺปยุตฺเต นว, โนนตฺถิยา ตีณิ, โนวิคเต ตีณิ. (สํขิตฺตํ.)}
\prose{เหตุปจฺจยา นอารมฺมเณ ตีณิ. (สํขิตฺตํ.)}

\prose{\textbf{นเหตุ}ปจฺจยา อารมฺมเณ เอกํ. (สํขิตฺตํ.)}

\prose{(สหชาตวาโรปิ ปจฺจยวาโรปิ นิสฺสยวาโรปิ สํสฏฺฐ\-วาโรปิ สมฺปยุตฺตวาโรปิ ปฏิจฺจ\-วาร\-สทิสา วิตฺถาเรตพฺพา.)}

\csromanmarkh{2. อสํกิลิฏฺฐสํกิเลสิกปท}\csromanmarkhii{7. ปญฺหาวาร}\csromanheadernomark{2}{2. \textbf{อสํกิลิฏฺฐสํกิเลสิกปท 7. ปญฺหาวาร}}
\prose{ปจฺจยจตุกฺกเหตฺวาทิ}

\proseitem{198}{เหตุ อสํกิลิฏฺฐ\-สํกิเลสิโก ธมฺโม เหตุสฺส อสํกิลิฏฺฐ\-สํกิเลสิกสฺส ธมฺมสฺส เหตุปจฺจเยน ปจฺจโย. ตีณิ.}

\prose{เหตุ อสํกิลิฏฺฐ\-สํกิเลสิโก ธมฺโม เหตุสฺส อสํกิลิฏฺฐ\-สํกิเลสิกสฺส ธมฺมสฺส อารมฺมณ\-ปจฺจเยน ปจฺจโย. นว.}

\prose{เหตุ อสํกิลิฏฺฐ\-สํกิเลสิโก ธมฺโม เหตุสฺส อสํกิลิฏฺฐ\-สํกิเลสิกสฺส ธมฺมสฺส อธิปติ\-ปจฺจเยน ปจฺจโย. อารมฺมณา\-ธิปติ, สหชาตา\-ธิปติ. ตีณิ.}

\prose{อสํกิลิฏฺฐ\-สํกิเลสิโก ธมฺโม นเหตุสฺส อสํกิลิฏฺฐ\-สํกิเลสิกสฺส ธมฺมสฺส อธิปติ\-ปจฺจเยน ปจฺจโย. อารมฺมณา\-ธิปติ, สหชาตา\-ธิปติ. ตีณิ.}

\csromanmarkh{1. เหตุทุก}\csromanmarkhii{5. สํกิลิฏฺฐ\-ตฺติก}\csromanheadernomark{2}{1. เหตุทุก 5. สํกิลิฏฺฐ\-ตฺติก}
\prose{\csromanfolio{237}เหตุ อสํกิลิฏฺฐ\-สํกิเลสิโก จ นเหตุ อสํกิลิฏฺฐ\-สํกิเลสิโก จ ธมฺมา เหตุสฺส อสํกิลิฏฺฐ\-สํกิเลสิกสฺส ธมฺมสฺส อธิปติ\-ปจฺจเยน ปจฺจโย. อารมฺมณา\-ธิปติ. ตีณิ. (สํขิตฺตํ.)}

\proseitem{199}{เหตุยา ตีณิ, อารมฺมเณ นว, อธิปติยา นว, อนนฺตเร นว, สมนนฺตเร นว, สหชาเต นว, อญฺญมญฺเญ นว, นิสฺสเย นว, อุปนิสฺสเย นว, ปุเรชาเต ตีณิ, ปจฺฉาชาเต ตีณิ, อาเสวเน นว, กมฺเม ตีณิ, วิปาเก นว, อาหาเร ตีณิ, อินฺทฺริเย นว, ฌาเน ตีณิ, มคฺเค นว, สมฺปยุตฺเต นว, วิปฺปยุตฺเต ปญฺจ, อตฺถิยา นว, นตฺถิยา นว, วิคเต นว, อวิคเต นว. (สํขิตฺตํ.)}

\proseitem{200}{เหตุ อสํกิลิฏฺฐ\-สํกิเลสิโก ธมฺโม เหตุสฺส อสํกิลิฏฺฐ\-สํกิเลสิกสฺส ธมฺมสฺส อารมฺมณ\-ปจฺจเยน ปจฺจโย. สหชาต\-ปจฺจเยน ปจฺจโย. อุปนิสฺสย\-ปจฺจเยน ปจฺจโย. (สํขิตฺตํ.)}

\proseitem{201}{นเหตุยา นว, นอารมฺมเณ นว. (สํขิตฺตํ.)}

\prose{เหตุปจฺจยา นอารมฺมเณ ตีณิ. (สํขิตฺตํ.)}

\prose{นเหตุปจฺจยา อารมฺมเณ นว. (สํขิตฺตํ.)}

\prose{(ยถา กุสลตฺติเก ปญฺหาวารสฺส อนุโลมมฺปิ ปจฺจนียมฺปิ อนุโลมปจฺจนียมฺปิ ปจฺจนียานุโลมมฺปิ คณิตํ, เอวํ คเณตพฺพํ.)}

\prose{3. อสํกิลิฏฺฐอสํกิเลสิกปท 1. ปฏิจฺจวาร}

\prose{ปจฺจยจตุกฺกเหตุ}

\proseitem{202}{เหตุํ อสํกิลิฏฺฐ\-อสํกิเลสิกํ ธมฺมํ ปฏิจฺจ เหตุ อสํกิลิฏฺฐ\-อสํกิเลสิโก ธมฺโม อุปฺปชฺชติ เหตุปจฺจยา. (สํขิตฺตํ.)}

\proseitem{203}{เหตุยา นว, อารมฺมเณ นว, อธิปติยา นว, อนนฺตเร นว, สมนนฺตเร นว, สหชาเต นว, อญฺญมญฺเญ นว, นิสฺสเย นว, \csromanfolio{238}อุปนิสฺสเย นว, ปุเรชาเต นว, อาเสวเน นว, กมฺเม นว, วิปาเก นว, อาหาเร นว, อินฺทฺริเย นว, ฌาเน นว, มคฺเค นว, สมฺปยุตฺเต นว, วิปฺปยุตฺเต นว, อตฺถิยา นว, นตฺถิยา นว, วิคเต นว, อวิคเต นว.}

\prose{นอธิปติ}

\proseitem{204}{เหตุํ อสํกิลิฏฺฐ\-อสํกิเลสิกํ ธมฺมํ ปฏิจฺจ เหตุ อสํกิลิฏฺฐ\-อสํกิเลสิโก ธมฺโม อุปฺปชฺชติ นอธิปติ\-ปจฺจยา.  เหตุํ อสํกิลิฏฺฐ\-อสํกิเลสิกํ ธมฺมํ ปฏิจฺจ นเหตุ อสํกิลิฏฺฐ\-อสํกิเลสิโก ธมฺโม อุปฺปชฺชติ นอธิปติ\-ปจฺจยา. (2)}

\prose{นเหตุํ อสํกิลิฏฺฐ\-อสํกิเลสิกํ ธมฺมํ ปฏิจฺจ นเหตุ อสํกิลิฏฺฐ\-อสํกิเลสิโก ธมฺโม อุปฺปชฺชติ นอธิปติ\-ปจฺจยา.  นเหตุํ อสํกิลิฏฺฐ\-อสํกิเลสิกํ ธมฺมํ ปฏิจฺจ เหตุ อสํกิลิฏฺฐ\-อสํกิเลสิโก ธมฺโม อุปฺปชฺชติ นอธิปติ\-ปจฺจยา. (2) (สํขิตฺตํ.)}

\proseitem{205}{นอธิปติยา ฉ, นปุเรชาเต นว, นปจฺฉาชาเต นว, นอาเสวเน นว, นกมฺเม ตีณิ, นวิปาเก นว, นวิปฺปยุตฺเต นว. (สํขิตฺตํ.)}

\prose{เหตุปจฺจยา นอธิปติยา ฉ. (สํขิตฺตํ.)}

\prose{นอธิปติ\-ปจฺจยา เหตุยา ฉ. (สํขิตฺตํ.)}

\prose{(สหชาตวาโรปิ ปจฺจยวาโรปิ นิสฺสยวาโรปิ สํสฏฺฐ\-วาโรปิ สมฺปยุตฺตวาโรปิ ปฏิจฺจ\-วาร\-สทิสา วิตฺถาเรตพฺพา.)}

\prose{3. อสํกิลิฏฺฐอสํกิเลสิกปท 7. ปญฺหาวาร}

\prose{ปจฺจยจตุกฺกเหตุ อารมฺมณ}

\proseitem{206}{เหตุ อสํกิลิฏฺฐ\-อสํกิเลสิโก ธมฺโม เหตุสฺส อสํกิลิฏฺฐ\-อสํกิเลสิกสฺส ธมฺมสฺส เหตุปจฺจเยน ปจฺจโย. ตีณิ.}

\prose{อสํกิลิฏฺฐ\-อสํกิเลสิโก ธมฺโม นเหตุสฺส อสํกิลิฏฺฐ\-อสํกิเลสิกสฺส ธมฺมสฺส อารมฺมณ\-ปจฺจเยน ปจฺจโย. ตีณิ. (สํขิตฺตํ.)}

\csromanmatikaanchor{mtk.37}
\csromantocmarkref{section}{6. วิตกฺกตฺติก}{mtk.37}
\bookmark[dest={mtk.37},level=1]{6. วิตกฺกตฺติก}
\csromanmarkh{1. เหตุทุก}\csromanmarkhii{6. วิตกฺกตฺติก}\csromanheadernomark{1}{1. เหตุทุก 6. วิตกฺกตฺติก}
\proseitem{207}{\csromanfolio{239}เหตุยา ตีณิ, อารมฺมเณ ตีณิ, อธิปติยา ฉ, อนนฺตเร นว, สมนฺตเร นว, สหชาเต นว, อญฺญมญฺเญ นว, นิสฺสเย นว, อุปนิสฺสเย นว, กมฺเม ตีณิ, วิปาเก นว, อาหาเร ตีณิ, อินฺทฺริเย นว, ฌาเน ตีณิ, มคฺเค นว, สมฺปยุตฺเต นว, อตฺถิยา นว, นตฺถิยา นว, วิคเต นว, อวิคเต นว. (สํขิตฺตํ.)}

\proseitem{208}{เหตุ อสํกิลิฏฺฐ\-อสํกิเลสิโก ธมฺโม เหตุสฺส อสํกิลิฏฺฐ\-อสํกิเลสิกสฺส ธมฺมสฺส สหชาต\-ปจฺจเยน ปจฺจโย. อุปนิสฺสย\-ปจฺจเยน ปจฺจโย. (สํขิตฺตํ.)}

\proseitem{209}{นเหตุยา นว, นอารมฺมเณ นว. (สํขิตฺตํ.)}

\prose{เหตุปจฺจยา นอารมฺมเณ ตีณิ. (สํขิตฺตํ.)}

\prose{นเหตุปจฺจยา อารมฺมเณ ตีณิ. (สํขิตฺตํ.)}

\prosewithcloser{(ยถา กุสลตฺติเก ปญฺหาวารสฺส อนุโลมมฺปิ ปจฺจนียมฺปิ อนุโลมปจฺจนียมฺปิ ปจฺจนียานุโลมมฺปิ คณิตํ, เอวํ คเณตพฺพํ.)}{เหตุ\-ทุก\-สํกิลิฏฺฐตฺติกํ นิฏฺฐิตํ.}
\csromanmarkh{1. เหตุทุก}\csromanmarkhii{6. วิตกฺกตฺติก}\csromanheadernomark{2}{1. เหตุทุก 6. วิตกฺกตฺติก}
\csromanmarkh{1. สวิตกฺกสวิจารปท}\csromanmarkhii{1. ปฏิจฺจวาร}\csromanheadernomark{2}{1. \textbf{สวิตกฺกสวิจารปท 1. ปฏิจฺจวาร}}
\prose{ปจฺจยจตุกฺกเหตุ}

\proseitem{210}{เหตุํ สวิตกฺก\-สวิจารํ ธมฺมํ ปฏิจฺจ เหตุ สวิตกฺก\-สวิจาโร ธมฺโม อุปฺปชฺชติ เหตุปจฺจยา.  เหตุํ สวิตกฺก\-สวิจารํ ธมฺมํ ปฏิจฺจ นเหตุ สวิตกฺก\-สวิจาโร ธมฺโม อุปฺปชฺชติ เหตุปจฺจยา.  เหตุํ สวิตกฺก\-สวิจารํ ธมฺมํ ปฏิจฺจ เหตุ สวิตกฺก\-สวิจาโร จ นเหตุ สวิตกฺก\-สวิจาโร จ ธมฺมา อุปฺปชฺชนฺติ เหตุปจฺจยา. (3)}

\prose{นเหตุํ สวิตกฺก\-สวิจารํ ธมฺมํ ปฏิจฺจ นเหตุ สวิตกฺก\-สวิจาโร ธมฺโม อุปฺปชฺชติ เหตุปจฺจยา. ตีณิ.}

\prose{\csromanfolio{240}เหตุํ สวิตกฺกสวิจารญฺจ นเหตุํ สวิตกฺกสวิจารญฺจ ธมฺมํ ปฏิจฺจ เหตุ สวิตกฺก\-สวิจาโร ธมฺโม อุปฺปชฺชติ เหตุปจฺจยา. ตีณิ. (สํขิตฺตํ.)}

\proseitem{211}{เหตุยา นว, อารมฺมเณ นว, อธิปติยา นว, อนนฺตเร นว, สมนนฺตเร นว ฯเปฯ กมฺเม นว, วิปาเก นว ฯเปฯ อวิคเต นว. (สํขิตฺตํ.)}

\prose{นเหตุนอธิปติ}

\proseitem{212}{นเหตุํ สวิตกฺก\-สวิจารํ ธมฺมํ ปฏิจฺจ นเหตุ สวิตกฺก\-สวิจาโร ธมฺโม อุปฺปชฺชติ นเหตุปจฺจยา. เทฺว.}

\prose{เหตุํ สวิตกฺก\-สวิจารํ ธมฺมํ ปฏิจฺจ เหตุ สวิตกฺก\-สวิจาโร ธมฺโม อุปฺปชฺชติ นอธิปติ\-ปจฺจยา. (สํขิตฺตํ.)}

\proseitem{213}{นเหตุยา เทฺว, นอธิปติยา นว, นปุเรชาเต นว, นปจฺฉาชาเต นว, นอาเสวเน นว, นกมฺเม ตีณิ, นวิปาเก นว, นมคฺเค เอกํ, นวิปฺปยุตฺเต นว. (สํขิตฺตํ.)}

\prose{เหตุปจฺจยา นอธิปติยา นว. (สํขิตฺตํ.)}

\prose{นเหตุปจฺจยา อารมฺมเณ เทฺว. (สํขิตฺตํ.)}

\prose{(สหชาตวาโรปิ ปจฺจยวาโรปิ นิสฺสยวาโรปิ สํสฏฺฐ\-วาโรปิ สมฺปยุตฺตวาโรปิ ปฏิจฺจ\-วาร\-สทิสา วิตฺถาเรตพฺพา.)}

\csromanmarkh{1. สวิตกฺกสวิจารปท}\csromanmarkhii{7. ปญฺหาวาร}\csromanheadernomark{2}{1. \textbf{สวิตกฺกสวิจารปท 7. ปญฺหาวาร}}
\prose{ปจฺจยจตุกฺกเหตฺวาทิ}

\proseitem{214}{เหตุ สวิตกฺก\-สวิจาโร ธมฺโม เหตุสฺส สวิตกฺก\-สวิจารสฺส ธมฺมสฺส เหตุปจฺจเยน ปจฺจโย. ตีณิ.}

\prose{เหตุ สวิตกฺก\-สวิจาโร ธมฺโม เหตุสฺส สวิตกฺก\-สวิจารสฺส ธมฺมสฺส อารมฺมณ\-ปจฺจเยน ปจฺจโย. นว.}

\prose{เหตุ สวิตกฺก\-สวิจาโร ธมฺโม เหตุสฺส สวิตกฺก\-สวิจารสฺส ธมฺมสฺส อธิปติ\-ปจฺจเยน ปจฺจโย. อารมฺมณา\-ธิปติ, สหชาตา\-ธิปติ. ตีณิ.}

\csromanmarkh{1. เหตุทุก}\csromanmarkhii{6. วิตกฺกตฺติก}\csromanheadernomark{2}{1. เหตุทุก 6. วิตกฺกตฺติก}
\prose{\csromanfolio{241}สวิตกฺก\-สวิจาโร ธมฺโม นเหตุสฺส สวิตกฺก\-สวิจารสฺส ธมฺมสฺส อธิปติ\-ปจฺจเยน ปจฺจโย. อารมฺมณา\-ธิปติ, สหชาตา\-ธิปติ. ตีณิ.}

\prose{เหตุ สวิตกฺก\-สวิจาโร จ นเหตุ สวิตกฺก\-สวิจาโร จ ธมฺมา เหตุสฺส สวิตกฺก\-สวิจารสฺส ธมฺมสฺส อธิปติ\-ปจฺจเยน ปจฺจโย. อารมฺมณา\-ธิปติ. ตีณิ. (สํขิตฺตํ.)}

\proseitem{215}{เหตุยา ตีณิ, อารมฺมเณ นว, อธิปติยา นว, อนนฺตเร นว, สมนนฺตเร นว, สหชาเต นว, อญฺญมญฺเญ นว, นิสฺสเย นว, อุปนิสฺสเย นว, อาเสวเน นว, กมฺเม ตีณิ, วิปาเก นว, อาหาเร ตีณิ, อินฺทฺริเย นว, ฌาเน ตีณิ, มคฺเค นว, สมฺปยุตฺเต นว, อตฺถิยา นว, นตฺถิยา นว, วิคเต นว, อวิคเต นว. (สํขิตฺตํ.)}

\proseitem{216}{เหตุ สวิตกฺก\-สวิจาโร ธมฺโม เหตุสฺส สวิตกฺก\-สวิจารสฺส ธมฺมสฺส อารมฺมณ\-ปจฺจเยน ปจฺจโย. สหชาต\-ปจฺจเยน ปจฺจโย. อุปนิสฺสย\-ปจฺจเยน ปจฺจโย. (สํขิตฺตํ.)}

\proseitem{217}{นเหตุยา นว, นอารมฺมเณ นว. (สํขิตฺตํ.)}

\prose{เหตุปจฺจยา นอารมฺมเณ ตีณิ. (สํขิตฺตํ.)}

\prose{นเหตุปจฺจยา อารมฺมเณ นว. (สํขิตฺตํ.)}

\prose{(ยถา กุสลตฺติเก ปญฺหาวารสฺส อนุโลมมฺปิ ปจฺจนียมฺปิ อนุโลมปจฺจนียมฺปิ ปจฺจนียานุโลมมฺปิ คณิตํ, เอวํ คเณตพฺพํ.)}

\csromanmarkh{2. อวิตกฺกวิจารมตฺตปท}\csromanmarkhii{1. ปฏิจฺจวาร}\csromanheadernomark{2}{2. \textbf{อวิตกฺกวิจารมตฺตปท 1. ปฏิจฺจวาร}}
\prose{ปจฺจยจตุกฺกเหตุ}

\proseitem{218}{เหตุํ อวิตกฺก\-วิจารมตฺตํ ธมฺมํ ปฏิจฺจ เหตุ อวิตกฺก\-วิจารมตฺโต ธมฺโม อุปฺปชฺชติ เหตุปจฺจยา. ตีณิ.}

\prose{\csromanfolio{242}นเหตุํ อวิตกฺก\-วิจารมตฺตํ ธมฺมํ ปฏิจฺจ นเหตุ อวิตกฺก\-วิจารมตฺโต ธมฺโม อุปฺปชฺชติ เหตุปจฺจยา. ตีณิ.}

\prose{เหตุํ อวิตกฺกวิจารมตฺตญฺจ นเหตุํ อวิตกฺกวิจารมตฺตญฺจ ธมฺมํ ปฏิจฺจ เหตุ อวิตกฺก\-วิจารมตฺโต ธมฺโม อุปฺปชฺชติ เหตุปจฺจยา. ตีณิ. (สํขิตฺตํ.)}

\proseitem{219}{เหตุยา นว, อารมฺมเณ นว, อธิปติยา นว ฯเปฯ กมฺเม นว, วิปาเก นว ฯเปฯ อวิคเต นว. (สํขิตฺตํ.)}

\prose{นอธิปติ}

\proseitem{220}{เหตุํ อวิตกฺก\-วิจารมตฺตํ ธมฺมํ ปฏิจฺจ เหตุ อวิตกฺก\-วิจารมตฺโต ธมฺโม อุปฺปชฺชติ นอธิปติ\-ปจฺจยา. (สํขิตฺตํ.)}

\proseitem{221}{นอธิปติยา นว, นปุเรชาเต นว, นปจฺฉาชาเต นว, นอาเสวเน นว, นกมฺเม ตีณิ, นวิปาเก นว, นวิปฺปยุตฺเต นว. (สํขิตฺตํ.)}

\prose{เหตุปจฺจยา นอธิปติยา นว. (สํขิตฺตํ.)}

\prose{นอธิปติ\-ปจฺจยา เหตุยา นว. (สํขิตฺตํ.)}

\prose{(สหชาตวาโรปิ ปจฺจยวาโรปิ นิสฺสยวาโรปิ สํสฏฺฐ\-วาโรปิ สมฺปยุตฺตวาโรปิ ปฏิจฺจ\-วาร\-สทิสา วิตฺถาเรตพฺพา.)}

\csromanmarkh{2. อวิตกฺกวิจารมตฺตปท}\csromanmarkhii{7. ปญฺหาวาร}\csromanheadernomark{2}{2. \textbf{อวิตกฺกวิจารมตฺตปท 7. ปญฺหาวาร}}
\prose{ปจฺจยจตุกฺกเหตฺวาทิ}

\proseitem{222}{เหตุ อวิตกฺก\-วิจารมตฺโต ธมฺโม เหตุสฺส อวิตกฺก\-วิจารมตฺตสฺส ธมฺมสฺส เหตุปจฺจเยน ปจฺจโย. ตีณิ.}

\prose{เหตุ อวิตกฺก\-วิจารมตฺโต ธมฺโม นเหตุสฺส อวิตกฺก\-วิจารมตฺตสฺส ธมฺมสฺส อารมฺมณ\-ปจฺจเยน ปจฺจโย. ตีณิ.}

\prose{เหตุ อวิตกฺก\-วิจารมตฺโต ธมฺโม เหตุสฺส อวิตกฺก\-วิจารมตฺตสฺส ธมฺมสฺส อธิปติ\-ปจฺจเยน ปจฺจโย. สหชาตา\-ธิปติ. ตีณิ.}

\csromanmarkh{1. เหตุทุก}\csromanmarkhii{6. วิตกฺกตฺติก}\csromanheadernomark{2}{1. เหตุทุก 6. วิตกฺกตฺติก}
\prose{\csromanfolio{243}อวิตกฺก\-วิจารมตฺโต ธมฺโม นเหตุสฺส อวิตกฺก\-วิจารมตฺตสฺส ธมฺมสฺส อธิปติ\-ปจฺจเยน ปจฺจโย. อารมฺมณา\-ธิปติ, สหชาตา\-ธิปติ. ตีณิ.}

\prose{เหตุ อวิตกฺก\-วิจารมตฺโต จ นเหตุ อวิตกฺก\-วิจารมตฺโต จ ธมฺมา นเหตุสฺส อวิตกฺก\-วิจารมตฺตสฺส ธมฺมสฺส อธิปติ\-ปจฺจเยน ปจฺจโย. อารมฺมณา\-ธิปติ. (สํขิตฺตํ.)}

\proseitem{223}{เหตุยา ตีณิ, อารมฺมเณ ตีณิ, อธิปติยา สตฺต, อนนฺตเร นว, สมนนฺตเร นว, สหชาเต นว, อญฺญมญฺเญ นว, นิสฺสเย นว, อุปนิสฺสเย นว, อาเสวเน นว, กมฺเม ตีณิ, วิปาเก นว, อาหาเร ตีณิ, อินฺทฺริเย นว, ฌาเน ตีณิ, มคฺเค นว, สมฺปยุตฺเต นว, อตฺถิยา นว, นตฺถิยา นว, วิคเต นว, อวิคเต นว. (สํขิตฺตํ.)}

\proseitem{224}{เหตุ อวิตกฺก\-วิจารมตฺโต ธมฺโม เหตุสฺส อวิตกฺก\-วิจารมตฺตสฺส ธมฺมสฺส สหชาต\-ปจฺจเยน ปจฺจโย. อุปนิสฺสย\-ปจฺจเยน ปจฺจโย. (สํขิตฺตํ.)}

\proseitem{225}{นเหตุยา นว, นอารมฺมเณ นว. (สํขิตฺตํ.)}

\prose{เหตุปจฺจยา นอารมฺมเณ ตีณิ. (สํขิตฺตํ.)}

\prose{นเหตุปจฺจยา อารมฺมเณ ตีณิ. (สํขิตฺตํ.)}

\prose{(ยถา กุสลตฺติเก ปญฺหาวารสฺส อนุโลมมฺปิ ปจฺจนียมฺปิ อนุโลมปจฺจนียมฺปิ ปจฺจนียานุโลมมฺปิ คณิตํ, เอวํ คเณตพฺพํ.)}

\prose{3. อวิตกฺกอวิจารปท 1. ปฏิจฺจวาร}

\prose{ปจฺจยจตุกฺกเหตุ}

\proseitem{226}{เหตุํ อวิตกฺก\-อวิจารํ ธมฺมํ ปฏิจฺจ เหตุ อวิตกฺก\-อวิจาโร ธมฺโม อุปฺปชฺชติ เหตุปจฺจยา. (สํขิตฺตํ.)}

\proseitem{227}{\csromanfolio{244}เหตุยา นว, อารมฺมเณ นว, อธิปติยา นว ฯเปฯ กมฺเม นว, วิปาเก นว ฯเปฯ วิคเต นว, อวิคเต นว. (สํขิตฺตํ.)}

\csromanheader{2}{\textbf{นเหตฺวาทิ}}
\proseitem{228}{นเหตุํ อวิตกฺก\-อวิจารํ ธมฺมํ ปฏิจฺจ นเหตุ อวิตกฺก\-อวิจาโร ธมฺโม อุปฺปชฺชติ นเหตุปจฺจยา. (1)}

\prose{เหตุํ อวิตกฺก\-อวิจารํ ธมฺมํ ปฏิจฺจ นเหตุ อวิตกฺก\-อวิจาโร ธมฺโม อุปฺปชฺชติ นอารมฺมณ\-ปจฺจยา. (1)}

\prose{นเหตุํ อวิตกฺก\-อวิจารํ ธมฺมํ ปฏิจฺจ นเหตุ อวิตกฺก\-อวิจาโร ธมฺโม อุปฺปชฺชติ นอารมฺมณ\-ปจฺจยา. (1)}

\prose{เหตุํ อวิตกฺกอวิจารญฺจ นเหตุํ อวิตกฺกอวิจารญฺจ ธมฺมํ ปฏิจฺจ นเหตุ อวิตกฺก\-อวิจาโร ธมฺโม อุปฺปชฺชติ นอารมฺมณ\-ปจฺจยา. (1)}

\prose{เหตุํ อวิตกฺก\-อวิจารํ ธมฺมํ ปฏิจฺจ เหตุ อวิตกฺก\-อวิจาโร ธมฺโม อุปฺปชฺชติ นอธิปติ\-ปจฺจยา. นว. (สํขิตฺตํ.)}

\proseitem{229}{เอกํ, นอารมฺมเณ ตีณิ, นอธิปติยา นว, นอนนฺตเร ตีณิ, นสมนนฺตเร ตีณิ, นอญฺญมญฺเญ ตีณิ, นฺเอาปนิสฺสเย ตีณิ, นปุเรชาเต นว, นปจฺฉาชาเต นว, นอาเสวเน นว, นกมฺเม ตีณิ, นวิปาเก นว, นอาหาเร เอกํ, ไนนฺทฺริเย เอกํ, นฌาเน เอกํ, นมคฺเค เอกํ, นสมฺปยุตฺเต ตีณิ, นวิปฺปยุตฺเต นว, โนนตฺถิยา ตีณิ, โนวิคเต ตีณิ. (สํขิตฺตํ.)}

\prose{เหตุปจฺจยา นอารมฺมเณ ตีณิ. (สํขิตฺตํ.)}

\prose{นเหตุปจฺจยา อารมฺมเณ เอกํ. (สํขิตฺตํ.)}

\prose{(สหชาตวาโรปิ ปจฺจยวาโรปิ นิสฺสยวาโรปิ สํสฏฺฐ\-วาโรปิ สมฺปยุตฺตวาโรปิ ปฏิจฺจ\-วาร\-สทิสา วิตฺถาเรตพฺพา.)}

\vspace{\tipitakaparskip}
\csromancenter{เหตุทุก 6. วิตกฺกตฺติก 3. อวิตกฺกอวิจารปท 7. ปญฺหาวาร}
\prose{\csromanfolio{245}ปจฺจยจตุกฺกเหตุ อารมฺมณ}

\proseitem{230}{เหตุ อวิตกฺก\-อวิจาโร ธมฺโม เหตุสฺส อวิตกฺก\-อวิจารสฺส ธมฺมสฺส เหตุปจฺจเยน ปจฺจโย. ตีณิ.}

\prose{เหตุ อวิตกฺก\-อวิจาโร ธมฺโม เหตุสฺส อวิตกฺก\-อวิจารสฺส ธมฺมสฺส อารมฺมณ\-ปจฺจเยน ปจฺจโย. นว. (สํขิตฺตํ.)}

\proseitem{231}{เหตุยา ตีณิ, อารมฺมเณ นว, อธิปติยา ฉ, อนนฺตเร นว, สมนนฺตเร นว, สหชาเต นว, อญฺญมญฺเญ นว, นิสฺสเย นว, อุปนิสฺสเย นว, ปุเรชาเต ตีณิ, ปจฺฉาชาเต ตีณิ, อาเสวเน นว, กมฺเม ตีณิ, วิปาเก นว, อาหาเร ตีณิ, อินฺทฺริเย นว, ฌาเน ตีณิ, มคฺเค นว, สมฺปยุตฺเต นว, วิปฺปยุตฺเต ปญฺจ, อตฺถิยา นว, นตฺถิยา นว, วิคเต นว, อวิคเต นว. (สํขิตฺตํ.)}

\proseitem{232}{เหตุ อวิตกฺก\-อวิจาโร ธมฺโม เหตุสฺส อวิตกฺก\-อวิจารสฺส ธมฺมสฺส อารมฺมณ\-ปจฺจเยน ปจฺจโย. สหชาต\-ปจฺจเยน ปจฺจโย. อุปนิสฺสย\-ปจฺจเยน ปจฺจโย. (สํขิตฺตํ.)}

\proseitem{233}{นเหตุยา นว, นอารมฺมเณ นว, นอธิปติยา นว. (สํขิตฺตํ.)}

\prose{เหตุปจฺจยา นอารมฺมเณ ตีณิ. (สํขิตฺตํ.)}

\prose{นเหตุปจฺจยา อารมฺมเณ นว. (สํขิตฺตํ.)}

\prosewithcloser{(ยถา กุสลตฺติเก ปญฺหาวารสฺส อนิโลมมฺปิ ปจฺจนียมฺปิ อนุโลมปจฺจนียมฺปิ ปจฺจนียานุโลมมฺปิ คณิตํ, เอวํ คเณตพฺพํ.)}{เหตุ\-ทุ\-กวิต\-กฺกตฺติกํ นิฏฺฐิตํ.}
\csromanmatikaanchor{mtk.38}
\csromantocmarkref{section}{7. ปีติตฺติก}{mtk.38}
\bookmark[dest={mtk.38},level=1]{7. ปีติตฺติก}
\csromanmarkh{ธมฺมานุโลม ทุกติก\-ปฏฺฐาน\-ปาฬิ}\csromanmarkhii{1. เหตุทุก 7. ปีติตฺติก}\csromanheadernomark{1}{ธมฺมานุโลม ทุกติก\-ปฏฺฐาน\-ปาฬิ 1. เหตุทุก 7. ปีติตฺติก}
\csromanmarkh{1. ปีติสหคตปท}\csromanmarkhii{1. ปฏิจฺจวาร}\csromanheadernomark{2}{1. \textbf{ปีติสหคตปท 1. ปฏิจฺจวาร}}
\prose{\csromanfolio{246}ปจฺจยจตุกฺกเหตุ}

\proseitem{234}{เหตุํ ปีติสหคตํ ธมฺมํ ปฏิจฺจ เหตุ ปีติสหคโต ธมฺโม อุปฺปชฺชติ เหตุปจฺจยา. ตีณิ.}

\prose{นเหตุํ ปีติสหคตํ ธมฺมํ ปฏิจฺจ นเหตุ ปีติสหคโต ธมฺโม อุปฺปชฺชติ เหตุปจฺจยา. (สํขิตฺตํ.)}

\proseitem{235}{เหตุยา นว, อารมฺมเณ นว, อธิปติยา นว ฯเปฯ กมฺเม นว, วิปาเก นว, อาหาเร นว ฯเปฯ อวิคเต นว. (สํขิตฺตํ.)}

\prose{นเหตุ นอธิปติ}

\proseitem{236}{นเหตุํ ปีติสหคตํ ธมฺมํ ปฏิจฺจ นเหตุ ปีติสหคโต ธมฺโม อุปฺปชฺชติ นเหตุปจฺจยา.}

\prose{เหตุํ ปีติสหคตํ ธมฺมํ ปฏิจฺจ เหตุ ปีติสหคโต ธมฺโม อุปฺปชฺชติ นอธิปติ\-ปจฺจยา. (สํขิตฺตํ.)}

\proseitem{237}{นเหตุยา เอกํ, นอธิปติยา นว, นปุเรชาเต นว, นปจฺฉาชาเต นว, นอาเสวเน นว, นกมฺเม ตีณิ, นวิปาเก นว, นมคฺเค เอกํ, นวิปฺปยุตฺเต นว. (สํขิตฺตํ.)}

\prose{เหตุปจฺจยา นอธิปติยา นว. (สํขิตฺตํ.)}

\prose{นเหตุปจฺจยา อารมฺมเณ เอกํ. (สํขิตฺตํ.)}

\prose{(สหชาตวาโรปิ ปจฺจยวาโรปิ นิสฺสยวาโรปิ สํสฏฺฐ\-วาโรปิ สมฺปยุตฺตวาโรปิ ปฏิจฺจ\-วาร\-สทิสา วิตฺถาเรตพฺพา.)}

\csromanmarkh{1. เหตุทุก}\csromanmarkhii{7. ปีติตฺติก 1. ปีติสหคตปท 7. ปญฺหาวาร}\csromanheadernomark{2}{1. เหตุทุก 7. ปีติตฺติก \textbf{1. ปีติสหคตปท 7. ปญฺหาวาร}}
\prose{\csromanfolio{247}ปจฺจยจตุกฺกเหตฺวาทิ}

\proseitem{238}{เหตุ ปีติสหคโต ธมฺโม เหตุสฺส ปีติสหคตสฺส ธมฺมสฺส เหตุปจฺจเยน ปจฺจโย. ตีณิ.}

\prose{เหตุ ปีติสหคโต ธมฺโม เหตุสฺส ปีติสหคตสฺส ธมฺมสฺส อารมฺมณ\-ปจฺจเยน ปจฺจโย. ตีณิ.}

\prose{ปีติสหคโต ธมฺโม นเหตุสฺส ปีติสหคตสฺส ธมฺมสฺส อารมฺมณ\-ปจฺจเยน ปจฺจโย. ตีณิ.}

\prose{เหตุ ปีติสหคโต จ นเหตุ ปีติสหคโต จ ธมฺมา เหตุสฺส ปีติสหคตสฺส ธมฺมสฺส อารมฺมณ\-ปจฺจเยน ปจฺจโย. ตีณิ.}

\prose{เหตุ ปีติสหคโต ธมฺโม เหตุสฺส ปีติสหคตสฺส ธมฺมสฺส อธิปติ\-ปจฺจเยน ปจฺจโย. อารมฺมณา\-ธิปติ, สหชาตา\-ธิปติ. ตีณิ.}

\prose{ปีติสหคโต ธมฺโม นเหตุสฺส ปีติสหคตสฺส ธมฺมสฺส อธิปติ\-ปจฺจเยน ปจฺจโย. อารมฺมณา\-ธิปติ, สหชาตา\-ธิปติ. ตีณิ.}

\prose{เหตุ ปีติสหคโต จ นเหตุ ปีติสหคโต จ ธมฺมา เหตุสฺส ปีติสหคตสฺส ธมฺมสฺส อธิปติ\-ปจฺจเยน ปจฺจโย. อารมฺมณา\-ธิปติ. ตีณิ. (สํขิตฺตํ.)}

\proseitem{239}{เหตุยา ตีณิ, อารมฺมเณ นว, อธิปติยา นว, อนนฺตเร นว, สมนนฺตเร นว, สหชาเต นว, อญฺญมญฺเญ นว, นิสฺสเย นว, อุปนิสฺสเย นว, อาเสวเน นว, กมฺเม ตีณิ, วิปาเก นว, อาหาเร ตีณิ, อินฺทฺริเย นว, ฌาเน ตีณิ, มคฺเค นว, สมฺปยุตฺเต นว, อตฺถิยา นว, นตฺถิยา นว, วิคเต นว, อวิคเต นว. (สํขิตฺตํ.)}

\proseitem{240}{เหตุ ปีติสหคโต ธมฺโม เหตุสฺส ปีติสหคตสฺส ธมฺมสฺส อารมฺมณ\-ปจฺจเยน ปจฺจโย. สหชาต\-ปจฺจเยน ปจฺจโย. อุปนิสฺสย\-ปจฺจเยน ปจฺจโย. (สํขิตฺตํ.)}

\proseitem{241}{\csromanfolio{248}นเหตุยา นว, นอารมฺมเณ นว. (สํขิตฺตํ.)}

\prose{เหตุปจฺจยา นอารมฺมเณ ตีณิ. (สํขิตฺตํ.)}

\prose{นเหตุปจฺจยา อารมฺมเณ นว. (สํขิตฺตํ.)}

\prose{(ยถา กุสลตฺติเก ปญฺหาวารสฺส อนุโลมมฺปิ ปจฺจนียมฺปิ อนุโลมปจฺจนียมฺปิ ปจฺจนียานุโลมมฺปิ คณิตํ, เอวํ คเณตพฺพํ.)}

\csromanmarkh{2. สุขสหคตปท}\csromanmarkhii{1. ปฏิจฺจวาร}\csromanheadernomark{2}{2. \textbf{สุขสหคตปท 1. ปฏิจฺจวาร}}
\prose{ปจฺจยจตุกฺกเหตุ}

\proseitem{242}{เหตุํ สุขสหคตํ ธมฺมํ ปฏิจฺจ เหตุ สุขสหคโต ธมฺโม อุปฺปชฺชติ เหตุปจฺจยา. ตีณิ.}

\prose{นเหตุํ สุขสหคตํ ธมฺมํ ปฏิจฺจ นเหตุ สุขสหคโต ธมฺโม อุปฺปชฺชติ เหตุปจฺจยา. ตีณิ.}

\prose{เหตุํ สุขสหคตญฺจ นเหตุํ สุขสหคตญฺจ ธมฺมํ ปฏิจฺจ เหตุ สุขสหคโต ธมฺโม อุปฺปชฺชติ เหตุปจฺจยา. ตีณิ. (สํขิตฺตํ.)}

\proseitem{243}{เหตุยา นว, อารมฺมเณ นว, อธิปติยา นว ฯเปฯ อวิคเต นว. (สํขิตฺตํ.)}

\prose{นเหตุ นอธิปติ}

\proseitem{244}{นเหตุํ สุขสหคตํ ธมฺมํ ปฏิจฺจ นเหตุ สุขสหคโต ธมฺโม อุปฺปชฺชติ นเหตุปจฺจยา.}

\prose{เหตุํ สุขสหคตํ ธมฺมํ ปฏิจฺจ เหตุ สุขสหคโต ธมฺโม อุปฺปชฺชติ นอธิปติ\-ปจฺจยา. นว. (สํขิตฺตํ.)}

\proseitem{245}{นเหตุยา เอกํ, นอธิปติยา นว, นปุเรชาเต นว, นปจฺฉาชาเต นว, นอาเสวเน นว, นกมฺเม ตีณิ, นวิปาเก นว, นฌาเน เอกํ, นมคฺเค เอกํ, นวิปฺปยุตฺเต นว. (สํขิตฺตํ.)}

\csromanmarkh{1. เหตุทุก}\csromanmarkhii{7. ปีติตฺติก}\csromanheadernomark{2}{1. เหตุทุก 7. ปีติตฺติก}
\prose{\csromanfolio{249}เหตุปจฺจยา นอธิปติยา นว. (สํขิตฺตํ.)}

\prose{นเหตุปจฺจยา อารมฺมเณ เอกํ. (สํขิตฺตํ.)}

\prose{(สหชาตวาโรปิ ปจฺจยวาโรปิ นิสฺสยวาโรปิ สํสฏฺฐ\-วาโรปิ สมฺปยุตฺตวาโรปิ ปฏิจฺจ\-วาร\-สทิสา วิตฺถาเรตพฺพา.)}

\csromanmarkh{2. สุขสหคตปท}\csromanmarkhii{7. ปญฺหาวาร}\csromanheadernomark{2}{2. \textbf{สุขสหคตปท 7. ปญฺหาวาร}}
\prose{ปจฺจยจตุกฺกเหตฺวาทิ}

\proseitem{246}{เหตุ สุขสหคโต ธมฺโม เหตุสฺส สุขสหคตสฺส ธมฺมสฺส เหตุปจฺจเยน ปจฺจโย. ตีณิ.}

\prose{เหตุ สุขสหคโต ธมฺโม เหตุสฺส สุขสหคตสฺส ธมฺมสฺส อารมฺมณ\-ปจฺจเยน ปจฺจโย. นว.}

\prose{เหตุ สุขสหคโต ธมฺโม เหตุสฺส สุขสหคตสฺส ธมฺมสฺส อธิปติ\-ปจฺจเยน ปจฺจโย. นว. (สํขิตฺตํ.)}

\proseitem{247}{เหตุยา ตีณิ, อารมฺมเณ นว, อธิปติยา นว, อนนฺตเร นว, สมนนฺตเร นว, สหชาเต นว, อญฺญมญฺเญ นว, นิสฺสเย นว, อุปนิสฺสเย นว, อาเสวเน นว, กมฺเม ตีณิ, วิปาเก นว, อาหาเร ตีณิ, อินฺทฺริเย นว, ฌาเน ตีณิ, มคฺเค นว, สมฺปยุตฺเต นว, อตฺถิยา นว, นตฺถิยา นว, วิคเต นว, อวิคเต นว. (สํขิตฺตํ.)}

\proseitem{248}{เหตุ สุขสหคโต ธมฺโม เหตุสฺส สุขสหคตสฺส ธมฺมสฺส อารมฺมณ\-ปจฺจเยน ปจฺจโย. สหชาต\-ปจฺจเยน ปจฺจโย. อุปนิสฺสย\-ปจฺจเยน ปจฺจโย. (สํขิตฺตํ.)}

\proseitem{249}{นเหตุยา นว, นอารมฺมเณ นว. (สํขิตฺตํ.)}

\prose{เหตุปจฺจยา นอารมฺมเณ ตีณิ. (สํขิตฺตํ.)}

\prose{นเหตุปจฺจยา อารมฺมเณ นว. (สํขิตฺตํ.)}

\prose{\csromanfolio{250}(ยถา กุสลตฺติเก ปญฺหาวารสฺส อนุโลมมฺปิ ปจฺจนียมฺปิ อนุโลมปจฺจนียมฺปิ ปจฺจนียานุโลมมฺปิ คณิตํ, เอวํ คเณตพฺพํ.)}

\csromanmarkh{3. อุเปกฺขาสหคตปท}\csromanmarkhii{1. ปฏิจฺจวาร}\csromanheadernomark{2}{3. \textbf{อุเปกฺขาสหคตปท 1. ปฏิจฺจวาร}}
\prose{ปจฺจยจตุกฺกเหตุ}

\proseitem{250}{เหตุํ อุเปกฺขา\-สหคตํ ธมฺมํ ปฏิจฺจ เหตุ อุเปกฺขา\-สหคโต ธมฺโม อุปฺปชฺชติ เหตุปจฺจยา. ตีณิ.}

\prose{นเหตุํ อุเปกฺขา\-สหคตํ ธมฺมํ ปฏิจฺจ นเหตุ อุเปกฺขา\-สหคโต ธมฺโม อุปฺปชฺชติ เหตุปจฺจยา. ตีณิ.}

\prose{เหตุํ อุเปกฺขาสหคตญฺจ นเหตุํ อุเปกฺขาสหคตญฺจ ธมฺมํ ปฏิจฺจ เหตุ อุเปกฺขา\-สหคโต ธมฺโม อุปฺปชฺชติ เหตุปจฺจยา. ตีณิ. (สํขิตฺตํ.)}

\proseitem{251}{เหตุยา นว, อารมฺมเณ นว, อธิปติยา นว ฯเปฯ กมฺเม นว, วิปาเก นว ฯเปฯ อวิคเต นว. (สํขิตฺตํ.)}

\csromanheader{2}{\textbf{นเหตฺวาทิ}}
\proseitem{252}{นเหตุํ อุเปกฺขา\-สหคตํ ธมฺมํ ปฏิจฺจ นเหตุ อุเปกฺขา\-สหคโต ธมฺโม อุปฺปชฺชติ นเหตุปจฺจยา. เทฺว.}

\prose{เหตุํ อุเปกฺขา\-สหคตํ ธมฺมํ ปฏิจฺจ เหตุ อุเปกฺขา\-สหคโต ธมฺโม อุปฺปชฺชติ นอธิปติ\-ปจฺจยา. นว.}

\prose{เหตุํ อุเปกฺขา\-สหคตํ ธมฺมํ ปฏิจฺจ เหตุ อุเปกฺขา\-สหคโต ธมฺโม อุปฺปชฺชติ นปุเรชาต\-ปจฺจยา. นว.}

\prose{เหตุํ อุเปกฺขา\-สหคตํ ธมฺมํ ปฏิจฺจ เหตุ อุเปกฺขา\-สหคโต ธมฺโม อุปฺปชฺชติ นปจฺฉาชาต\-ปจฺจยา. นว. (สํขิตฺตํ.)}

\proseitem{253}{นเหตุยา เทฺว, นอธิปติยา นว, นปุเรชาเต นว, นปจฺฉาชาเต นว, นอาเสวเน นว, นกมฺเม ตีณิ, นวิปาเก นว, นฌาเน เอกํ, นมคฺเค เอกํ, นวิปฺปยุตฺเต นว. (สํขิตฺตํ.)}

\csromanmarkh{1. เหตุทุก}\csromanmarkhii{7. ปีติตฺติก}\csromanheadernomark{2}{1. เหตุทุก 7. ปีติตฺติก}
\prose{\csromanfolio{251}เหตุปจฺจยา นอธิปติยา นว. (สํขิตฺตํ.)}

\prose{นเหตุปจฺจยา อารมฺมเณ เทฺว. (สํขิตฺตํ.)}

\prose{(สหชาตวาโรปิ ปจฺจยวาโรปิ นิสฺสยวาโรปิ สํสฏฺฐ\-วาโรปิ สมฺปยุตฺตวาโรปิ ปฏิจฺจ\-วาร\-สทิสา วิตฺถาเรตพฺพา.)}

\csromanmarkh{3. อุเปกฺขาสหคตปท}\csromanmarkhii{7. ปญฺหาวาร}\csromanheadernomark{2}{3. \textbf{อุเปกฺขาสหคตปท 7. ปญฺหาวาร}}
\prose{ปจฺจยจตุกฺกเหตฺวาทิ}

\proseitem{254}{เหตุ อุเปกฺขา\-สหคโต ธมฺโม เหตุสฺส อุเปกฺขา\-สหคตสฺส ธมฺมสฺส เหตุปจฺจเยน ปจฺจโย. ตีณิ.}

\prose{เหตุ อุเปกฺขา\-สหคโต ธมฺโม เหตุสฺส อุเปกฺขา\-สหคตสฺส ธมฺมสฺส อารมฺมณ\-ปจฺจเยน ปจฺจโย. นว.}

\prose{เหตุ อุเปกฺขา\-สหคโต ธมฺโม เหตุสฺส อุเปกฺขา\-สหคตสฺส ธมฺมสฺส อธิปติ\-ปจฺจเยน ปจฺจโย. นว. (สํขิตฺตํ.)}

\proseitem{255}{เหตุยา ตีณิ, อารมฺมเณ นว, อธิปติยา นว, อนนฺตเร นว, สมนนฺตเร นว, สหชาเต นว, อญฺญมญฺเญ นว, นิสฺสเย นว, อุปนิสฺสเย นว, อาเสวเน นว, กมฺเม ตีณิ, วิปาเก นว, อาหาเร ตีณิ, อินฺทฺริเร นว, ฌาเน ตีณิ, มคฺเค นว, สมฺปยุตฺเต นว, อตฺถิยา นว, นตฺถิยา นว, วิคเต นว. อวิคเต นว. (สํขิตฺตํ.)}

\proseitem{256}{เหตุ อุเปกฺขา\-สหคโต ธมฺโม เหตุสฺส อุเปกฺขา\-สหคตสฺส ธมฺมสฺส อารมฺมณ\-ปจฺจเยน ปจฺจโย. สหชาต\-ปจฺจเยน ปจฺจโย. อุปนิสฺสย\-ปจฺจเยน ปจฺจโย. (สํขิตฺตํ.)}

\proseitem{257}{นเหตุยา นว, นอารมฺมเณ นว. (สํขิตฺตํ.)}

\prose{เหตุปจฺจยา นอารมฺมเณ ตีณิ. (สํขิตฺตํ.)}

\prose{นเหตุปจฺจยา อารมฺมเณ นว. (สํขิตฺตํ.)}

\prosewithcloser{\csromanfolio{252}(ยถา กุสลตฺติเก ปญฺหาวารสฺส อนุโลมมฺปิ ปจฺจนียมฺปิ อนุโลมปจฺจนียมฺปิ ปจฺจนียานุโลมมฺปิ คณิตํ, เอวํ คเณตพฺพํ.)}{เหตุ\-ทุ\-กปี\-ติตฺติกํ นิฏฺฐิตํ.}
\csromanmatikaanchor{mtk.39}
\csromantocmarkref{section}{8. ทสฺสเนนปหาตพฺพตฺติก}{mtk.39}
\bookmark[dest={mtk.39},level=1]{8. ทสฺสเนนปหาตพฺพตฺติก}
\csromanmarkh{1. เหตุทุก}\csromanmarkhii{8. ทสฺสเนน\-ปหาตพฺพ\-ตฺติก}\csromanheadernomark{1}{1. เหตุทุก 8. ทสฺสเนน\-ปหาตพฺพ\-ตฺติก}
\csromanmarkh{1. ทสฺสเนน\-ปหาตพฺพปท}\csromanmarkhii{1. ปฏิจฺจวาร}\csromanheadernomark{2}{1. ทสฺสเนน\-ปหาตพฺพปท 1. ปฏิจฺจวาร}
\prose{ปจฺจยจตุกฺกเหตุ}

\proseitem{258}{เหตุํ ทสฺสเนน ปหาตพฺพํ ธมฺมํ ปฏิจฺจ เหตุ ทสฺสเนน ปหาตพฺโพ ธมฺโม อุปฺปชฺชติ เหตุปจฺจยา. ตีณิ.}

\prose{นเหตุํ ทสฺสเนน ปหาตพฺพํ ธมฺมํ ปฏิจฺจ นเหตุ ทสฺสเนน ปหาตพฺโพ ธมฺโม อุปฺปชฺชติ เหตุปจฺจยา. ตีณิ.}

\prose{เหตุํ ทสฺสเนน ปหาตพฺพญฺจ นเหตุํ ทสฺสเนน ปหาตพฺพญฺจ ธมฺมํ ปฏิจฺจ เหตุ ทสฺสเนน ปหาตพฺโพ ธมฺโม อุปฺปชฺชติ เหตุปจฺจยา. ตีณิ. (สํขิตฺตํ.)}

\proseitem{259}{เหตุยา นว, อารมฺมเณ นว, อธิปติยา นว ฯเปฯ กมฺเม นว, อาหาเร นว ฯเปฯ อวิคเต นว. (สํขิตฺตํ.)}

\prose{นเหตุ นอธิปติ}

\proseitem{260}{นเหตุํ ทสฺสเนน ปหาตพฺพํ ธมฺมํ ปฏิจฺจ เหตุ ทสฺสเนน ปหาตพฺโพ ธมฺโม อุปฺปชฺชติ นเหตุปจฺจยา. (1)}

\prose{เหตุํ ทสฺสเนน ปหาตพฺพํ ธมฺมํ ปฏิจฺจ เหตุ ทสฺสเนน ปหาตพฺโพ ธมฺโม อุปฺปชฺชติ นอธิปติ\-ปจฺจยา. นว.}

\proseitem{261}{นเหตุยา เอกํ, นอธิปติยา นว, นปุเรชาเต นว, นปจฺฉาชาเต นว, นอาเสวเน นว, นกมฺเม ตีณิ, นวิปาเก นว, นวิปฺปยุตฺเต นว. (สํขิตฺตํ.)}

\prose{เหตุปจฺจยา นอธิปติยา นว. (สํขิตฺตํ.)}

\csromanmarkh{1. เหตุทุก}\csromanmarkhii{8. ทสฺสเนน\-ปหาตพฺพ\-ตฺติก}\csromanheadernomark{2}{1. เหตุทุก 8. ทสฺสเนน\-ปหาตพฺพ\-ตฺติก}
\prose{\csromanfolio{253}นเหตุปจฺจยา อารมฺมเณ เอกํ. (สํขิตฺตํ.)}

\prose{(สหชาตวาโรปิ ปจฺจยวาโรปิ นิสฺสยวาโรปิ สํสฏฺฐ\-วาโรปิ สมฺปยุตฺตวาโรปิ ปฏิจฺจ\-วาร\-สทิสา วิตฺถาเรตพฺพา.)}

\csromanmarkh{1. ทสฺสเนน\-ปหาตพฺพปท}\csromanmarkhii{7. ปญฺหาวาร}\csromanheadernomark{2}{1. ทสฺสเนน\-ปหาตพฺพปท 7. ปญฺหาวาร}
\prose{ปจฺจยจตุกฺกเหตุ อารมฺมณ}

\proseitem{262}{เหตุ ทสฺสเนน ปหาตพฺโพ ธมฺโม เหตุสฺส ทสฺสเนน ปหาตพฺพสฺส ธมฺมสฺส เหตุปจฺจเยน ปจฺจโย. ตีณิ.}

\prose{เหตุ ทสฺสเนน ปหาตพฺโพ ธมฺโม เหตุสฺส ทสฺสเนน ปหาตพฺพสฺส ธมฺมสฺส อารมฺมณ\-ปจฺจเยน ปจฺจโย. นว. (สํขิตฺตํ.)}

\proseitem{263}{เหตุยา ตีณิ, อารมฺมเณ นว, อธิปติยา นว, อนนฺตเร นว, สมนนฺตเร นว, สหชาเต นว, อญฺญมญฺเญ นว, นิสฺสเย นว, อุปนิสฺสเย นว, อาเสวเน นว, กมฺเม ตีณิ, อาหาเร ตีณิ, อินฺทฺริเย ตีณิ, ฌาเน ตีณิ, มคฺเค ตีณิ, สมฺปยุตฺเต นว ฯเปฯ อวิคเต นว. (สํขิตฺตํ.)}

\proseitem{264}{เหตุ ทสฺสเนน ปหาตพฺโพ ธมฺโม เหตุสฺส ทสฺสเนน ปหาตพฺพสฺส ธมฺมสฺส อารมฺมณ\-ปจฺจเยน ปจฺจโย. สหชาต\-ปจฺจเยน ปจฺจโย. อุปนิสฺสย\-ปจฺจเยน ปจฺจโย. (สํขิตฺตํ.)}

\proseitem{265}{นเหตุยา นว, นอารมฺมเณ นว. (สํขิตฺตํ.)}

\prose{เหตุปจฺจยา นอารมฺมเณ ตีณิ. (สํขิตฺตํ.)}

\prose{นเหตุปจฺจยา อารมฺมเณ นว. (สํขิตฺตํ.)}

\prose{(ยถา กุสลตฺติเก ปญฺหาวารสฺส อนุโลมมฺปิ ปจฺจนียมฺปิ อนุโลมปจฺจนียมฺปิ ปจฺจนียานุโลมมฺปิ คณิตํ, เอวํ คเณตพฺพํ.)}

\csromanmarkh{ธมฺมานุโลม ทุกติก\-ปฏฺฐาน\-ปาฬิ}\csromanmarkhii{2. ภาวนาย\-ปหาตพฺพปท 1. ปฏิจฺจวาร}\csromanheadernomark{2}{ธมฺมานุโลม ทุกติก\-ปฏฺฐาน\-ปาฬิ 2. ภาวนาย\-ปหาตพฺพปท 1. ปฏิจฺจวาร}
\prose{\csromanfolio{254}ปจฺจยจตุกฺกเหตุ}

\proseitem{266}{เหตุํ ภาวนาย ปหาตพฺพํ ธมฺมํ ปฏิจฺจ เหตุ ภาวนาย ปหาตพฺโพ ธมฺโม อุปฺปชฺชติ เหตุปจฺจยา. ตีณิ.}

\prose{นเหตุํ ภาวนาย ปหาตพฺพํ ธมฺมํ ปฏิจฺจ นเหตุ ภาวนาย ปหาตพฺโพ ธมฺโม อุปฺปชฺชติ เหตุปจฺจยา. ตีณิ.}

\prose{เหตุํ ภาวนาย ปหาตพฺพญฺจ นเหตุํ ภาวนาย ปหาตพฺพญฺจ ธมฺมํ ปฏิจฺจ เหตุ ภาวนาย ปหาตพฺโพ ธมฺโม อุปฺปชฺชติ เหตุปจฺจยา. ตีณิ. (สํขิตฺตํ.)}

\proseitem{267}{เหตุยา นว, อารมฺมเณ นว ฯเปฯ อวิคเต นว. (สํขิตฺตํ.)}

\vspace{\tipitakaparskip}
\csromancenter{\textbf{นเหตุ}}
\proseitem{268}{นเหตุํ ภาวนาย ปหาตพฺพํ ธมฺมํ ปฏิจฺจ เหตุ ภาวนาย ปหาตพฺโพ ธมฺโม อุปฺปชฺชติ นเหตุปจฺจยา. (1) (สํขิตฺตํ.)}

\vspace{\tipitakaparskip}
\csromancenter{\textbf{นเหตุ}ยา เอกํ, นอธิปติยา นว, นปุเรชาเต นว, นปจฺฉาชาเต นว, นอาเสวเน นว, นกมฺเม ตีณิ, นวิปาเก นว, นวิปฺปยุตฺเต นว. (สํขิตฺตํ.)}
\prose{เหตุปจฺจยา นอธิปติยา นว. (สํขิตฺตํ.)}

\prose{\textbf{นเหตุ}ปจฺจยา อารมฺมเณ เอกํ. (สํขิตฺตํ.)}

\prose{(สหชาตวาโรปิ ปจฺจยวาโรปิ นิสฺสยวาโรปิ สํสฏฺฐ\-วาโรปิ สมฺปยุตฺตวาโรปิ ปฏิจฺจ\-วาร\-สทิสา วิตฺถาเรตพฺพา.)}

\csromanmarkh{2. ภาวนาย\-ปหาตพฺพปท}\csromanmarkhii{7. ปญฺหาวาร}\csromanheadernomark{2}{2. ภาวนาย\-ปหาตพฺพปท 7. ปญฺหาวาร}
\prose{ปจฺจยจตุกฺกเหตุ อารมฺมณ}

\proseitem{270}{เหตุ ภาวนาย ปหาตพฺโพ ธมฺโม เหตุสฺส ภาวนาย ปหาตพฺพสฺส ธมฺมสฺส เหตุปจฺจเยน ปจฺจโย. ตีณิ.}

\csromanmarkh{1. เหตุทุก}\csromanmarkhii{8. ทสฺสเนน\-ปหาตพฺพ\-ตฺติก}\csromanheadernomark{2}{1. เหตุทุก 8. ทสฺสเนน\-ปหาตพฺพ\-ตฺติก}
\prose{\csromanfolio{255}เหตุ ภาวนาย ปหาตพฺโพ ธมฺโม เหตุสฺส ภาวนาย ปหาตพฺพสฺส ธมฺมสฺส อารมฺมณ\-ปจฺจเยน ปจฺจโย. นว. (สํขิตฺตํ.)}

\proseitem{271}{เหตุยา ตีณิ, อารมฺมเณ นว, อธิปติยา นว, อนนฺตเร นว, สมนนฺตเร นว, สหชาเต นว, อญฺญมญฺเญ นว, นิสฺสเย นว, อุปนิสฺสเย นว, อาเสวเน นว, กมฺเม ตีณิ, อาหาเร ตีณิ, อินฺทฺริเย ตีณิ ฯเปฯ อวิคเต นว. (สํขิตฺตํ.)}

\proseitem{272}{เหตุ ภาวนาย ปหาตพฺโพ ธมฺโม เหตุสฺส ภาวนาย ปหาตพฺพสฺส ธมฺมสฺส อารมฺมณ\-ปจฺจเยน ปจฺจโย. สหชาต\-ปจฺจเยน ปจฺจโย. อุปนิสฺสย\-ปจฺจเยน ปจฺจโย. (สํขิตฺตํ.)}

\proseitem{273}{นเหตุยา นว, นอารมฺมเณ นว. (สํขิตฺตํ.)}

\prose{เหตุปจฺจยา นอารมฺมเณ ตีณิ. (สํขิตฺตํ.)}

\prose{นเหตุปจฺจยา อารมฺมเณ นว. (สํขิตฺตํ.)}

\prose{(ยถา กุสลตฺติเก ปญฺหาวารสฺส อนุโลมมฺปิ ปจฺจนียมฺปิ อนุโลมปจฺจนียมฺปิ ปจฺจนียานุโลมมฺปิ คณิตํ, เอวํ คเณตพฺพํ.)}

\prose{3. \textbf{เนวทสฺสเนนภาวนายปหาตพฺพปท 1. ปฏิจฺจวาร}}

\prose{ปจฺจยจตุกฺกเหตุ}

\proseitem{274}{เหตุํ เนวทสฺสเนน นภาวนาย ปหาตพฺพํ ธมฺมํ ปฏิจฺจ เหตุ เนวทสฺสเนน นภาวนาย ปหาตพฺโพ ธมฺโม อุปฺปชฺชติ เหตุปจฺจยา. ตีณิ.}

\prose{นเหตุํ เนวทสฺสเนน นภาวนาย ปหาตพฺพํ ธมฺมํ ปฏิจฺจ นเหตุ เนวทสฺสเนน นภาวนาย ปหาตพฺโพ ธมฺโม อุปฺปชฺชติ เหตุปจฺจยา. (สํขิตฺตํ.)}

\proseitem{275}{\csromanfolio{256}เหตุยา นว, อารมฺมเณ นว, อธิปติยา นว ฯเปฯ กมฺเม นว, วิปาเก นว ฯเปฯ อวิคเต นว. (สํขิตฺตํ.)}

\csromanheader{2}{\textbf{นเหตฺวาทิ}}
\proseitem{276}{นเหตุํ เนวทสฺสเนน นภาวนาย ปหาตพฺพํ ธมฺมํ ปฏิจฺจ นเหตุ เนวทสฺสเนน นภาวนาย ปหาตพฺโพ ธมฺโม อุปฺปชฺชติ นเหตุปจฺจยา.}

\prose{เหตุํ เนวทสฺสเนน นภาวนาย ปหาตพฺพํ ธมฺมํ ปฏิจฺจ นเหตุ เนวทสฺสเนน นภาวนาย ปหาตพฺโพ ธมฺโม อุปฺปชฺชติ นอารมฺมณ\-ปจฺจยา. (1)}

\prose{นเหตุํ เนวทสฺสเนน นภาวนาย ปหาตพฺพํ ธมฺมํ ปฏิจฺจ นเหตุ เนวทสฺสเนน นภาวนาย ปหาตพฺโพ ธมฺโม อุปฺปชฺชติ นอารมฺมณ\-ปจฺจยา. (1)}

\prose{เหตุํ เนวทสฺสเนน นภาวนาย ปหาตพฺพญฺจ นเหตุํ เนวทสฺสเนน นภาวนาย ปหาตพฺพญฺจ ธมฺมํ ปฏิจฺจ นเหตุ เนวทสฺสเนน นภาวนาย ปหาตพฺโพ ธมฺโม อุปฺปชฺชติ นอารมฺมณ\-ปจฺจยา. (1)}

\prose{เหตุํ เนวทสฺสเนน นภาวนาย ปหาตพฺพํ ธมฺมํ ปฏิจฺจ เหตุ เนวทสฺสเนน นภาวนาย ปหาตพฺโพ ธมฺโม อุปฺปชฺชติ นอธิปติ\-ปจฺจยา. (สํขิตฺตํ.)}

\proseitem{277}{นเหตุยา เอกํ, นอารมฺมเณ ตีณิ, นอธิปติยา นว, นอนนฺตเร ตีณิ, นสมนนฺตเร ตีณิ, นอญฺญมญฺเญ ตีณิ, นอุปนิสฺสเย ตีณิ, นปุเรชาเต นว, นปจฺฉาชาเต นว, นอาเสวเน นว, นกมฺเม ตีณิ, นวิปาเก นว, นอาหาเร เอกํ, นอินฺทฺริเย เอกํ, นฌาเน เอกํ, นมคฺเค เอกํ, นสมฺปยุตฺเต ตีณิ, นวิปฺปยุตฺเต นว, โนนตฺถิยา ตีณิ, โนวิคเต ตีณิ. (สํขิตฺตํ.)}

\prose{เหตุปจฺจยา นอารมฺมเณ ตีณิ. (สํขิตฺตํ.)}

\prose{นเหตุปจฺจยา อารมฺมเณ เอกํ. (สํขิตฺตํ.)}

\prose{(สหชาตวาโรปิ ปจฺจยวาโรปิ นิสฺสยวาโรปิ สํสฏฺฐ\-วาโรปิ สมฺปยุตฺตวาโรปิ ปฏิจฺจ\-วาร\-สทิสา วิตฺถาเรตพฺพา.)}

\vspace{\tipitakaparskip}
\csromancenter{เหตุทุก 8. ทสฺสเนน\-ปหาตพฺพ\-ตฺติก 3. เนวทสฺสเนนนภาวนายปหาตพฺพปท 7. ปญฺหาวาร}
\prose{\csromanfolio{257}ปจฺจยจตุกฺกเหตุ อารมฺมณ}

\proseitem{278}{เหตุ เนวทสฺสเนน นภาวนาย ปหาตพฺโพ ธมฺโม เหตุสฺส เนวทสฺสเนน นภาวนาย ปหาตพฺพสฺส ธมฺมสฺส เหตุปจฺจเยน ปจฺจโย. ตีณิ.}

\prose{เหตุ เนวทสฺสเนน นภาวนาย ปหาตพฺโพ ธมฺโม เหตุสฺส เนวทสฺสเนน นภาวนาย ปหาตพฺพสฺส ธมฺมสฺส อารมฺมณ\-ปจฺจเยน ปจฺจโย. นว. (สํขิตฺตํ.)}

\proseitem{279}{เหตุยา ตีณิ, อารมฺมเณ นว, อธิปติยา นว, อนนฺตเร นว, สมนนฺตเร นว, สหชาเต นว, อญฺญมญฺเญ นว, นิสฺสเย นว, อุปนิสฺสเย นว, ปุเรชาเต ตีณิ, ปจฺฉาชาเต ตีณิ, อาเสวเน นว, กมฺเม ตีณิ, วิปาเก นว, อาหาเร ตีณิ, อินฺทฺริเย นว, ฌาเน ตีณิ, มคฺเค นว, สมฺปยุตฺเต นว, วิปฺปยุตฺเต ปญฺจ, อตฺถิยา นว, นตฺถิยา นว, วิคเต นว, อวิคเต นว.}

\proseitem{280}{เหตุ เนวทสฺสเนน นภาวนาย ปหาตพฺโพ ธมฺโม เหตุสฺส เนวทสฺสเนน นภาวนาย ปหาตพฺพสฺส ธมฺมสฺส อารมฺมณ\-ปจฺจเยน ปจฺจโย. สหชาต\-ปจฺจเยน ปจฺจโย. อุปนิสฺสย\-ปจฺจเยน ปจฺจโย. (สํขิตฺตํ.)}

\proseitem{281}{นเหตุยา นว, นอารมฺมเณ นว. (สํขิตฺตํ.)}

\prose{เหตุปจฺจยา นอารมฺมเณ ตีณิ. (สํขิตฺตํ.)}

\prose{นเหตุปจฺจยา อารมฺมเณ นว. (สํขิตฺตํ.)}

\prosewithcloser{(ยถา กุสลตฺติเก ปญฺหาวารสฺส อนุโลมมฺปิ ปจฺจนียมฺปิ อนุโลมปจฺจนียมฺปิ ปจฺจนียานุโลมมฺปิ คณิตํ, เอวํ คเณตพฺพํ.)}{เหตุทุก ทสฺสเนน ปหาตพฺพตฺติกํ นิฏฺฐิตํ.}
\csromanmatikaanchor{mtk.40}
\csromantocmarkref{section}{9. ทสฺสเนนปหาตพฺพเหตุกตฺติก}{mtk.40}
\bookmark[dest={mtk.40},level=1]{9. ทสฺสเนนปหาตพฺพเหตุกตฺติก}
\csromanmarkh{ธมฺมานุโลม ทุกติก\-ปฏฺฐาน\-ปาฬิ}\csromanmarkhii{1. เหตุทุก 9. ทสฺสเนน\-ปหาตพฺพ\-เหตุ\-กตฺติก}\csromanheadernomark{1}{ธมฺมานุโลม ทุกติก\-ปฏฺฐาน\-ปาฬิ 1. เหตุทุก 9. ทสฺสเนน\-ปหาตพฺพ\-เหตุ\-กตฺติก}
\csromanmarkh{1. ทสฺสเนน\-ปหาตพฺพ\-เหตุกปท}\csromanmarkhii{1. ปฏิจฺจวาร}\csromanheadernomark{2}{1. ทสฺสเนน\-ปหาตพฺพ\-เหตุกปท 1. ปฏิจฺจวาร}
\prose{\csromanfolio{258}ปจฺจยจตุกฺกเหตุ}

\proseitem{282}{เหตุํ ทสฺสเนน ปหาตพฺพ\-เหตุกํ ธมฺมํ ปฏิจฺจ เหตุ ทสฺสเนน ปหาตพฺพ\-เหตุโก ธมฺโม อุปฺปชฺชติ เหตุปจฺจยา. ตีณิ.}

\prose{นเหตุํ ทสฺสเนน ปหาตพฺพ\-เหตุกํ ธมฺมํ ปฏิจฺจ นเหตุ ทสฺสเนน ปหาตพฺพ\-เหตุโก ธมฺโม อุปฺปชฺชติ เหตุปจฺจยา. ตีณิ.}

\prose{เหตุํ ทสฺสเนน ปหาตพฺพ\-เหตุ\-กญฺจ นเหตุํ ทสฺสเนน ปหาตพฺพ\-เหตุ\-กญฺจ ธมฺมํ ปฏิจฺจ เหตุ ทสฺสเนน ปหาตพฺพ\-เหตุโก ธมฺโม อุปฺปชฺชติ เหตุปจฺจยา. ตีณิ. (สํขิตฺตํ.)}

\proseitem{283}{เหตุยา นว, อารมฺมเณ นว ฯเปฯ อวิคเต นว. (สํขิตฺตํ.)}

\prose{นอธิปติ}

\proseitem{284}{เหตุํ ทสฺสเนน ปหาตพฺพ\-เหตุกํ ธมฺมํ ปฏิจฺจ เหตุ ทสฺสเนน ปหาตพฺพ\-เหตุโก ธมฺโม อุปฺปชฺชติ นอธิปติ\-ปจฺจยา. (สํขิตฺตํ.)}

\proseitem{285}{นอธิปติยา นว, นปุเรชาเต นว, นปจฺฉาชาเต นว, นอาเสวเน นว, นกมฺเม ตีณิ, นวิปาเก นว, นวิปฺปยุตฺเต นว. (สํขิตฺตํ.)}

\prose{เหตุปจฺจยา นอธิปติยา นว. (สํขิตฺตํ.)}

\prose{นอธิปติ\-ปจฺจยา เหตุยา นว. (สํขิตฺตํ.)}

\prose{(สหชาตวาโรปิ ปจฺจยวาโรปิ นิสฺสยวาโรปิ สํสฏฺฐ\-วาโรปิ สมฺปยุตฺตวาโรปิ ปฏิจฺจ\-วาร\-สทิสา วิตฺถาเรตพฺพา.)}

\csromanmarkh{1. ทสฺสเนนปหาตพฺพฺอเหตุกปท}\csromanmarkhii{7. ปญฺหาวาร}\csromanheadernomark{2}{1. \textbf{ทสฺสเนนปหาตพฺพฺ}อเหตุกปท 7. ปญฺหาวาร}
\prose{ปจฺจยจตุกฺกเหตุ อารมฺมณ}

\proseitem{286}{เหตุ ทสฺสเนน ปหาตพฺพ\-เหตุโก ธมฺโม เหตุสฺส ทสฺสเนน ปหาตพฺพ\-เหตุกสฺส ธมฺมสฺส เหตุปจฺจเยน ปจฺจโย. ตีณิ.}

\csromanmarkh{1. เหตุทุก}\csromanmarkhii{9. ทสฺสเนน\-ปหาตพฺพ\-เหตุ\-กตฺติก}\csromanheadernomark{2}{1. เหตุทุก 9. ทสฺสเนน\-ปหาตพฺพ\-เหตุ\-กตฺติก}
\prose{\csromanfolio{259}เหตุ ทสฺสเนน ปหาตพฺพ\-เหตุโก ธมฺโม เหตุสฺส ทสฺสเนน ปหาตพฺพ\-เหตุกสฺส ธมฺมสฺส อารมฺมณ\-ปจฺจเยน ปจฺจโย. นว. (สํขิตฺตํ.)}

\proseitem{287}{เหตุยา ตีณิ, อารมฺมเณ นว, อธิปติยา นว ฯเปฯ อุปนิสฺสเย นว, อาเสวเน นว, กมฺเม ตีณิ, อาหาเร ตีณิ ฯเปฯ มคฺเค ตีณิ, สมฺปยุตฺเต นว ฯเปฯ อวิคเต นว. (สํขิตฺตํ.)}

\proseitem{288}{เหตุ ทสฺสเนน ปหาตพฺพ\-เหตุโก ธมฺโม เหตุสฺส ทสฺสเนน ปหาตพฺพ\-เหตุกสฺส ธมฺมสฺส อารมฺมณ\-ปจฺจเยน ปจฺจโย. สหชาต\-ปจฺจเยน ปจฺจโย. อุปนิสฺสย\-ปจฺจเยน ปจฺจโย. (สํขิตฺตํ.)}

\proseitem{289}{นเหตุยา นว, นอารมฺมเณ นว. (สํขิตฺตํ.)}

\prose{เหตุปจฺจยา นอารมฺมเณ ตีณิ. (สํขิตฺตํ.)}

\prose{นเหตุปจฺจยา อารมฺมเณ นว. (สํขิตฺตํ.)}

\prose{(ยถา กุสลตฺติเก ปญฺหาวารสฺส อนุโลมมฺปิ ปจฺจนียมฺปิ อนุโลมปจฺจนียมฺปิ ปจฺจนียานุโลมมฺปิ คณิตํ, เอวํ คเณตพฺพํ.)}

\csromanmarkh{2. ภาวนาย\-ปหาตพฺพ\-เหตุกปท}\csromanmarkhii{1. ปฏิจฺจวาร}\csromanheadernomark{2}{2. ภาวนาย\-ปหาตพฺพ\-เหตุกปท 1. ปฏิจฺจวาร}
\prose{ปจฺจยจตุกฺกเหตุ}

\proseitem{290}{เหตุํ ภาวนาย ปหาตพฺพ\-เหตุกํ ธมฺมํ ปฏิจฺจ เหตุ ภาวนาย ปหาตพฺพ\-เหตุโก ธมฺโม อุปฺปชฺชติ เหตุปจฺจยา. ตีณิ.}

\prose{นเหตุํ ภาวนาย ปหาตพฺพ\-เหตุกํ ธมฺมํ ปฏิจฺจ นเหตุ ภาวนาย ปหาตพฺพ\-เหตุโก ธมฺโม อุปฺปชฺชติ เหตุปจฺจยา. (สํขิตฺตํ.)}

\proseitem{291}{เหตุยา นว, อารมฺมเณ นว, อธิปติยา นว ฯเปฯ อวิคเต นว. (สํขิตฺตํ.)}

\csromanheader{2}{ธมฺมานุโลม ทุกติก\-ปฏฺฐาน\-ปาฬิ นอธิปติ}
\proseitem{292}{\csromanfolio{260}เหตุํ ภาวนาย ปหาตพฺพ\-เหตุกํ ธมฺมํ ปฏิจฺจ เหตุ ภาวนาย ปหาตพฺพ\-เหตุโก ธมฺโม อุปฺปชฺชติ นอธิปติ\-ปจฺจยา. นว. (สํขิตฺตํ.)}

\proseitem{293}{นอธิปติยา นว, นปุเรชาเต นว, นปจฺฉาชาเต นว, นอาเสวเน นว, นกมฺเม ตีณิ, นวิปาเก นว, นวิปฺปยุตฺเต นว. (สํขิตฺตํ.)}

\prose{เหตุปจฺจยา นอธิปติยา นว. (สํขิตฺตํ.)}

\prose{นอธิปติ\-ปจฺจยา เหตุยา นว. (สํขิตฺตํ.)}

\prose{(สหชาตวาโรปิ ปจฺจยวาโรปิ นิสฺสยวาโรปิ สํสฏฺฐ\-วาโรปิ สมฺปยุตฺตวาโรปิ ปฏิจฺจ\-วาร\-สทิสา วิตฺถาเรตพฺพา.)}

\csromanmarkh{2. ภาวนาย\-ปหาตพฺพ\-เหตุกปท}\csromanmarkhii{7. ปญฺหาวาร}\csromanheadernomark{2}{2. ภาวนาย\-ปหาตพฺพ\-เหตุกปท 7. ปญฺหาวาร}
\prose{ปจฺจยจตุกฺกเหตุ อารมฺมณ}

\proseitem{294}{เหตุ ภาวนาย ปหาตพฺพ\-เหตุโก ธมฺโม เหตุสฺส ภาวนาย ปหาตพฺพ\-เหตุกสฺส ธมฺมสฺส เหตุปจฺจเยน ปจฺจโย. ตีณิ.}

\prose{เหตุ ภาวนาย ปหาตพฺพ\-เหตุโก ธมฺโม เหตุสฺส ภาวนาย ปหาตพฺพ\-เหตุกสฺส ธมฺมสฺส อารมฺมณ\-ปจฺจเยน ปจฺจโย. นว. (สํขิตฺตํ.)}

\proseitem{295}{เหตุยา ตีณิ, อารมฺมเณ นว, อธิปติยา นว, อนนฺตเร นว, สมนนฺตเร นว, สหชาเต นว, อญฺญมญฺเญ นว, นิสฺสเย นว, อุปนิสฺสเย นว, อาเสวเน นว, กมฺเม ตีณิ, อาหาเร ตีณิ ฯเปฯ มคฺเค ตีณิ, สมฺปยุตฺเต นว ฯเปฯ อวิคเต นว. (สํขิตฺตํ.)}

\proseitem{296}{เหตุ ภาวนาย ปหาตพฺพ\-เหตุโก ธมฺโม เหตุสฺส ภาวนาย ปหาตพฺพ\-เหตุกสฺส ธมฺมสฺส อารมฺมณ\-ปจฺจเยน ปจฺจโย. สหชาต\-ปจฺจเยน ปจฺจโย. อุปนิสฺสย\-ปจฺจเยน ปจฺจโย. (สํขิตฺตํ.)}

\csromanmarkh{1. เหตุทุก}\csromanmarkhii{9. ทสฺสเนน\-ปหาตพฺพ\-เหตุ\-กตฺติก}\csromanheadernomark{2}{1. เหตุทุก 9. ทสฺสเนน\-ปหาตพฺพ\-เหตุ\-กตฺติก}
\vspace{\tipitakaparskip}
\csromancenter{\textbf{นเหตุ}ยา นว, นอารมฺมเณ นว. (สํขิตฺตํ.)}
\prose{\csromanfolio{261}เหตุปจฺจยา นอารมฺมเณ ตีณิ. (สํขิตฺตํ.)}

\prose{นเหตุปจฺจยา อารมฺมเณ นว. (สํขิตฺตํ.)}

\prose{(ยถา กุสลตฺติเก ปญฺหาวารสฺส อนุโลมมฺปิ ปจฺจนียมฺปิ อนุโลมปจฺจนียมฺปิ ปจฺจนียานุโลมมฺปิ คณิตํ, เอวํ คเณตพฺพํ.)}

\prose{3. \textbf{เนวทสฺสเนนนภาวนายปหาตพฺพเหตุกปท 1. ปฏิจฺจวาร}}

\prose{ปจฺจยจตุกฺกเหตุ}

\proseitem{298}{เหตุํ เนวทสฺสเนน นภาวนาย ปหาตพฺพ\-เหตุกํ ธมฺมํ ปฏิจฺจ เหตุ เนวทสฺสเนน นภาวนาย ปหาตพฺพ\-เหตุโก ธมฺโม อุปฺปชฺชติ เหตุปจฺจยา. (สํขิตฺตํ.)}

\proseitem{299}{เหตุยา นว, อารมฺมเณ นว ฯเปฯ อวิคเต นว. (สํขิตฺตํ.)}

\vspace{\tipitakaparskip}
\csromancenter{\textbf{นเหตุ}}
\proseitem{300}{เหตุํ เนวทสฺสเนน นภาวนาย ปหาตพฺพ\-เหตุกํ ธมฺมํ ปฏิจฺจ นเหตุ เนวทสฺสเนน นภาวนาย ปหาตพฺพ\-เหตุโก ธมฺโม อุปฺปชฺชติ นเหตุปจฺจยา. (สํขิตฺตํ.)}

\vspace{\tipitakaparskip}
\csromancenter{\textbf{นเหตุ}ยา เอกํ, นอารมฺมเณ ตีณิ, นอธิปติยา นว, นอนนฺตเร ตีณิ, นสมนนฺตเร ตีณิ, นอญฺญมญฺเญ ตีณิ, นฺเอาปนิสฺสเย ตีณิ, นปุเรชาเต นว, นปจฺฉาชาเต นว, นอาเสวเน นว, นกมฺเม ตีณิ, นวิปาเก นว, นอาหาเร เอกํ, ไนนฺทฺริเย เอกํ, นฌาเน เอกํ, นมคฺเค เอกํ, นสมฺปยุตฺเต ตีณิ, นวิปฺปยุตฺเต นว, โนนตฺถิยา ตีณิ, โนวิคเต ตีณิ. (สํขิตฺตํ.)}
\prose{เหตุปจฺจยา นอารมฺมเณ ตีณิ. (สํขิตฺตํ.)}

\prose{\textbf{นเหตุ}ปจฺจยา อารมฺมเณ เอกํ. (สํขิตฺตํ.)}

\prose{(สหชาตวาโรปิ ปจฺจยวาโรปิ นิสฺสยวาโรปิ สํสฏฺฐ\-วาโรปิ สมฺปยุตฺตวาโรปิ ปฏิจฺจ\-วาร\-สทิสา วิตฺถาเรตพฺพา.)}

\prose{\csromanfolio{262}3. เนวทสฺสเนนนภาวนายปหาตพฺพเหตุกปท 7. ปญฺหาวาร}

\prose{ปจฺจยจตุกฺกเหตุ}

\proseitem{302}{เหตุ เนวทสฺสเนน นภาวนาย ปหาตพฺพ\-เหตุโก ธมฺโม เหตุสฺส เนวทสฺสเนน นภาวนาย ปหาตพฺพ\-เหตุกสฺส ธมฺมสฺส เหตุปจฺจเยน ปจฺจโย. ตีณิ. (สํขิตฺตํ.)}

\proseitem{303}{เหตุยา ตีณิ, อารมฺมเณ นว, อธิปติยา นว, อนนฺตเร นว, สมนนฺตเร นว, สหชาเต นว, อญฺญมญฺเญ นว, นิสฺสเย นว, อุปนิสฺสเย นว, ปุเรชาเต ตีณิ, ปจฺฉาชาเต ตีณิ, อาเสวเน นว, กมฺเม ตีณิ, วิปาเก นว, อาหาเร ตีณิ, อินฺทฺริเย นว, ฌาเน ตีณิ, มคฺเค นว, สมฺปยุตฺเต นว, วิปฺปยุตฺเต ปญฺจ, อตฺถิยา นว, นตฺถิยา นว, วิคเต นว, อวิคเต นว. (สํขิตฺตํ.)}

\proseitem{304}{เหตุ เนวทสฺสเนน นภาวนาย ปหาตพฺพ\-เหตุโก ธมฺโม เหตุสฺส เนวทสฺสเนน นภาวนาย ปหาตพฺพ\-เหตุกสฺส ธมฺมสฺส อารมฺมณ\-ปจฺจเยน ปจฺจโย. สหชาต\-ปจฺจเยน ปจฺจโย. อุปนิสฺสย\-ปจฺจเยน ปจฺจโย. (สํขิตฺตํ.)}

\proseitem{305}{นเหตุยา นว, นอารมฺมเณ นว. (สํขิตฺตํ.)}

\prose{เหตุปจฺจยา นอารมฺมเณ ตีณิ. (สํขิตฺตํ.)}

\prose{นเหตุปจฺจยา อารมฺมเณ นว. (สํขิตฺตํ.)}

\prosewithcloser{(ยถา กุสลตฺติเก ปญฺหาวารสฺส อนุโลมมฺปิ ปจฺจนียมฺปิ อนุโลมปจฺจนียมฺปิ ปจฺจนียานุโลมมฺปิ คณิตํ, เอวํ คเณตพฺพํ.)}{เหตุทุก ทสฺสเนน ปหาตพฺพเหตุกตฺติกํ นิฏฺฐิตํ.}
\csromanmatikaanchor{mtk.41}
\csromantocmarkref{section}{10. อาจยคามิตฺติก}{mtk.41}
\bookmark[dest={mtk.41},level=1]{10. อาจยคามิตฺติก}
\csromanmarkh{1. เหตุทุก}\csromanmarkhii{10. อาจยคามิตฺติก 1. เหตุทุก 10. อาจยคามิตฺติก}\csromanheadernomark{1}{1. เหตุทุก 10. อาจยคามิตฺติก \textbf{1. เหตุทุก 10. อาจยคามิตฺติก}}
\csromanmarkh{1. อาจยคามิปท}\csromanmarkhii{1. ปฏิจฺจวาร}\csromanheadernomark{2}{1. \textbf{อาจยคามิปท 1. ปฏิจฺจวาร}}
\prose{\csromanfolio{263}ปจฺจยจตุกฺกเหตุ}

\proseitem{306}{เหตุํ อาจยคามิํ ธมฺมํ ปฏิจฺจ เหตุ อาจยคามี ธมฺโม อุปฺปชฺชติ เหตุปจฺจยา. ตีณิ.}

\prose{นเหตุํ อาจยคามิํ ธมฺมํ ปฏิจฺจ นเหตุ อาจยคามี ธมฺโม อุปฺปชฺชติ เหตุปจฺจยา. ตีณิ.}

\prose{เหตุํ อาจยคามิญฺจ นเหตุํ อาจยคามิญฺจ ธมฺมํ ปฏิจฺจ เหตุ อาจยคามี ธมฺโม อุปฺปชฺชติ เหตุปจฺจยา. ตีณิ. (สํขิตฺตํ.)}

\proseitem{307}{เหตุยา นว, อารมฺมเณ นว, อธิปติยา นว, อนนฺตเร นว, สมนนฺตเร นว ฯเปฯ อวิคเต นว. (สํขิตฺตํ.)}

\prose{นเหตุ นอธิปติ}

\proseitem{308}{นเหตุํ อาจยคามิํ ธมฺมํ ปฏิจฺจ เหตุ อาจยคามี ธมฺโม อุปฺปชฺชติ นเหตุปจฺจยา. (1)}

\prose{เหตุํ อาจยคามิํ ธมฺมํ ปฏิจฺจ เหตุ อาจยคามี ธมฺโม อุปฺปชฺชติ นอธิปติ\-ปจฺจยา. (สํขิตฺตํ.)}

\proseitem{309}{นเหตุยา เอกํ, นอธิปติยา นว, นปุเรชาเต นว, นปจฺฉาชาเต นว, นอาเสวเน นว, นกมฺเม ตีณิ, นวิปาเก นว, นวิปฺปยุตฺเต นว. (สํขิตฺตํ.)}

\prose{เหตุปจฺจยา นอธิปติยา นว. (สํขิตฺตํ.)}

\prose{นเหตุปจฺจยา อารมฺมเณ เอกํ. (สํขิตฺตํ.)}

\prose{(สหชาตวาโรปิ ปจฺจยวาโรปิ นิสฺสยวาโรปิ สํสฏฺฐ\-วาโรปิ สมฺปยุตฺตวาโรปิ ปฏิจฺจ\-วาร\-สทิสา วิตฺถาเรตพฺพา.)}

\csromanmarkh{ธมฺมานุโลม ทุกติก\-ปฏฺฐาน\-ปาฬิ}\csromanmarkhii{1. อาจยคามิปท 7. ปญฺหาวาร}\csromanheadernomark{2}{ธมฺมานุโลม ทุกติก\-ปฏฺฐาน\-ปาฬิ 1. อาจยคามิปท 7. ปญฺหาวาร}
\prose{\csromanfolio{264}ปจฺจยจตุกฺกเหตฺวาทิ}

\proseitem{310}{เหตุ อาจยคามี ธมฺโม เหตุสฺส อาจยคามิสฺส ธมฺมสฺส เหตุปจฺจเยน ปจฺจโย. ตีณิ.}

\prose{เหตุ อาจยคามี ธมฺโม เหตุสฺส อาจยคามิสฺส ธมฺมสฺส อารมฺมณ\-ปจฺจเยน ปจฺจโย. นว.}

\prose{เหตุ อาจยคามี ธมฺโม เหตุสฺส อาจยคามิสฺส ธมฺมสฺส อธิปติ\-ปจฺจเยน ปจฺจโย. อารมฺมณา\-ธิปติ, สหชาตา\-ธิปติ. ตีณิ.}

\prose{อาจยคามี ธมฺโม นเหตุสฺส อาจยคามิสฺส ธมฺมสฺส อธิปติ\-ปจฺจเยน ปจฺจโย. อารมฺมณา\-ธิปติ, สหชาตา\-ธิปติ. ตีณิ.}

\prose{เหตุ อาจยคามี จ นเหตุ อาจยคามี จ ธมฺมา เหตุสฺส อาจยคามิสฺส ธมฺมสฺส อธิปติ\-ปจฺจเยน ปจฺจโย. อารมฺมณา\-ธิปติ. ตีณิ. (สํขิตฺตํ.)}

\proseitem{311}{เหตุยา ตีณิ, อารมฺมเณ นว, อธิปติยา นว ฯเปฯ อุปนิสฺสเย นว, อาเสวเน นว, กมฺเม ตีณิ, อาหาเร ตีณิ ฯเปฯ อวิคเต นว. (สํขิตฺตํ.)}

\proseitem{312}{เหตุ อาจยคามี ธมฺโม เหตุสฺส อาจยคามิสฺส ธมฺมสฺส อารมฺมณ\-ปจฺจเยน ปจฺจโย. สหชาต\-ปจฺจเยน ปจฺจโย. อุปนิสฺสย\-ปจฺจเยน ปจฺจโย. (สํขิตฺตํ.)}

\proseitem{313}{นเหตุยา นว, นอารมฺมเณ นว. (สํขิตฺตํ.)}

\prose{เหตุปจฺจยา นอารมฺมเณ ตีณิ. (สํขิตฺตํ.)}

\prose{นเหตุปจฺจยา อารมฺมเณ นว. (สํขิตฺตํ.)}

\prose{(ยถา กุสลตฺติเก ปญฺหาวารสฺส อนุโลมมฺปิ ปจฺจนียมฺปิ อนุโลมปจฺจนียมฺปิ ปจฺจนียานุโลมมฺปิ คณิตํ, เอวํ คเณตพฺพํ.)}

\vspace{\tipitakaparskip}
\csromancenter{เหตุทุก 10. อาจยคามิตฺติก 2. อปจยคามิปท 1. ปฏิจฺจวาร}
\prose{\csromanfolio{265}ปจฺจยจตุกฺกเหตุ}

\proseitem{314}{เหตุํ อปจยคามิํ ธมฺมํ ปฏิจฺจ เหตุ อปจยคามี ธมฺโม อุปฺปชฺชติ เหตุปจฺจยา. ตีณิ.}

\prose{นเหตุํ อปจยคามิํ ธมฺมํ ปฏิจฺจ นเหตุ อปจยคามี ธมฺโม อุปฺปชฺชติ เหตุปจฺจยา. ตีณิ.}

\prose{เหตุํ อปจยคามิญฺจ นเหตุํ อปจยคามิญฺจ ธมฺมํ ปฏิจฺจ เหตุ อปจยคามี ธมฺโม อุปฺปชฺชติ เหตุปจฺจยา. ตีณิ. (สํขิตฺตํ.)}

\proseitem{315}{เหตุยา นว, อารมฺมเณ นว ฯเปฯ กมฺเม นว, อาหาเร นว ฯเปฯ อวิคเต นว. (สํขิตฺตํ.)}

\prose{นอธิปติ}

\proseitem{316}{เหตุํ อปจยคามิํ ธมฺมํ ปฏิจฺจ เหตุ อปจยคามี ธมฺโม อุปฺปชฺชติ นอธิปติ\-ปจฺจยา. (สํขิตฺตํ.)}

\proseitem{317}{นอธิปติยา ฉ, นปุเรชาเต นว, นปจฺฉาชาเต นว, นกมฺเม ตีณิ, นวิปาเก นว, นวิปฺปยุตฺเต นว. (สํขิตฺตํ.)}

\prose{เหตุปจฺจยา นอธิปติยา ฉ. (สํขิตฺตํ.)}

\prose{นอธิปติ\-ปจฺจยา เหตุยา ฉ. (สํขิตฺตํ.)}

\prose{(สหชาตวาโรปิ ปจฺจยวาโรปิ นิสฺสยวาโรปิ สํสฏฺฐ\-วาโรปิ สมฺปยุตฺตวาโรปิ ปฏิจฺจ\-วาร\-สทิสา วิตฺถาเรตพฺพา.)}

\csromanmarkh{2. อปจยคามิปท}\csromanmarkhii{7. ปญฺหาวาร}\csromanheadernomark{2}{2. \textbf{อปจยคามิปท 7. ปญฺหาวาร}}
\prose{ปจฺจยจตุกฺกเหตุ อธิปติ}

\proseitem{318}{เหตุ อปจยคามี ธมฺโม เหตุสฺส อปจยคามิสฺส ธมฺมสฺส เหตุปจฺจเยน ปจฺจโย. ตีณิ.}

\prose{\csromanfolio{266}เหตุ อปจยคามี ธมฺโม เหตุสฺส อปจยคามิสฺส ธมฺมสฺส อธิปติ\-ปจฺจเยน ปจฺจโย. สหชาตา\-ธิปติ. ตีณิ.}

\prose{อปจยคามี ธมฺโม นเหตุสฺส อปจยคามิสฺส ธมฺมสฺส อธิปติ\-ปจฺจเยน ปจฺจโย. สหชาตา\-ธิปติ. ตีณิ. (สํขิตฺตํ.)}

\proseitem{319}{เหตุยา ตีณิ, อธิปติยา ฉ, สหชาเต นว ฯเปฯ อุปนิสฺสเย นว, กมฺเม ตีณิ, อาหาเร ตีณิ, อินฺทฺริเย นว, ฌาเน ตีณิ, มคฺเค นว, สมฺปยุตฺเต นว, อตฺถิยา นว, อวิคเต นว. (สํขิตฺตํ.)}

\proseitem{320}{เหตุ อปจยคามี ธมฺโม เหตุสฺส อปจยคามิสฺส ธมฺมสฺส สหชาต\-ปจฺจเยน ปจฺจโย. อุปนิสฺสย\-ปจฺจเยน ปจฺจโย. (สํขิตฺตํ.)}

\proseitem{321}{นเหตุยา นว, นอารมฺมเณ นว. (สํขิตฺตํ.)}

\prose{เหตุปจฺจยา นอารมฺมเณ ตีณิ. (สํขิตฺตํ.)}

\prose{นเหตุปจฺจยา อธิปติยา ตีณิ. (สํขิตฺตํ.)}

\prose{(ยถา กุสลตฺติเก ปญฺหาวารสฺส อนุโลมมฺปิ ปจฺจนียมฺปิ อนุโลมปจฺจนียมฺปิ ปจฺจนียานุโลมมฺปิ คณิตํ, เอวํ คเณตพฺพํ.)}

\csromanmarkh{3. เนวาจยคามินาปจยคามิปท}\csromanmarkhii{1. ปฏิจฺจวาร}\csromanheadernomark{2}{3. \textbf{เนวาจยคามินาปจยคามิปท 1. ปฏิจฺจวาร}}
\prose{ปจฺจยจตุกฺกเหตุ}

\proseitem{322}{เหตุํ เนวา\-จยคามิ\-นา\-ปจยคามิํ ธมฺมํ ปฏิจฺจ เหตุ เนวา\-จยคามิ\-นา\-ปจยคามี ธมฺโม อุปฺปชฺชติ เหตุปจฺจยา. ตีณิ.}

\prose{นเหตุํ เนวา\-จยคามิ\-นา\-ปจยคามิํ ธมฺมํ ปฏิจฺจ นเหตุ เนวา\-จยคามิ\-นา\-ปจยคามี ธมฺโม อุปฺปชฺชติ เหตุปจฺจยา. ตีณิ.}

\prose{เหตุํ เนวา\-จยคามิ\-นา\-ปจยคามิญฺจ นเหตุํ เนวา\-จยคามิ\-นา\-ปจยคามิญฺจ ธมฺมํ ปฏิจฺจ เหตุ เนวา\-จยคามิ\-นา\-ปจยคามี ธมฺโม อุปฺปชฺชติ เหตุปจฺจยา. ตีณิ. (สํขิตฺตํ.)}

\csromanmarkh{1. เหตุทุก}\csromanmarkhii{10. อาจยคามิตฺติก}\csromanheadernomark{2}{1. เหตุทุก 10. อาจยคามิตฺติก}
\proseitem{323}{\csromanfolio{267}เหตุยา นว, อารมฺมเณ นว, อธิปติยา นว ฯเปฯ อุปนิสฺสเย นว, ปุเรชาเต นว, อาเสวเน นว, กมฺเม นว, วิปาเก นว, อาหาเร นว, อินฺทฺริเย นว ฯเปฯ อวิคเต นว. (สํขิตฺตํ.)}

\csromanheader{2}{\textbf{นเหตฺวาทิ}}
\proseitem{324}{นเหตุํ เนวา\-จยคามิ\-นา\-ปจยคามิํ ธมฺมํ ปฏิจฺจ นเหตุ เนวา\-จยคามิ\-นา\-ปจยคามี ธมฺโม อุปฺปชฺชติ นเหตุปจฺจยา. (1)}

\prose{เหตุํ เนวา\-จยคามิ\-นา\-ปจยคามิํ ธมฺมํ ปฏิจฺจ นเหตุ เนวา\-จยคามิ\-นา\-ปจยคามี ธมฺโม อุปฺปชฺชติ นอารมฺมณ\-ปจฺจยา. (1)}

\prose{นเหตุํ เนวา\-จยคามิ\-นา\-ปจยคามิํ ธมฺมํ ปฏิจฺจ นเหตุ เนวา\-จยคามิ\-นา\-ปจยคามี ธมฺโม อุปฺปชฺชติ นอารมฺมณ\-ปจฺจยา. (1)}

\prose{เหตุํ เนวา\-จยคามิ\-นา\-ปจยคามิญฺจ นเหตุํ เนวา\-จยคามิ\-นา\-ปจยคามิญฺจ ธมฺมํ ปฏิจฺจ นเหตุ เนวา\-จยคามิ\-นา\-ปจยคามี ธมฺโม อุปฺปชฺชติ นอารมฺมณ\-ปจฺจยา. (1)}

\prose{เหตุํ เนวา\-จยคามิ\-นา\-ปจยคามิํ ธมฺมํ ปฏิจฺจ เหตุ เนวา\-จยคามิ\-นา\-ปจยคามี ธมฺโม อุปฺปชฺชติ นอธิปติ\-ปจฺจยา ฯเปฯ.}

\prose{เหตุํ เนวา\-จยคามิ\-นา\-ปจยคามิํ ธมฺมํ ปฏิจฺจ เหตุ เนวา\-จยคามิ\-นา\-ปจยคามี ธมฺโม อุปฺปชฺชติ นปุเรชาต\-ปจฺจยา. ตีณิ.}

\prose{เหตุํ เนวา\-จยคามิ\-นา\-ปจยคามิํ ธมฺมํ ปฏิจฺจ เหตุ เนวา\-จยคามิ\-นา\-ปจยคามี ธมฺโม อุปฺปชฺชติ นปจฺฉาชาต\-ปจฺจยา ฯเปฯ.}

\prose{เหตุํ เนวา\-จยคามิ\-นา\-ปจยคามิํ ธมฺมํ ปฏิจฺจ นเหตุ เนวา\-จยคามิ\-นา\-ปจยคามี ธมฺโม อุปฺปชฺชติ นกมฺมปจฺจยา. (สํขิตฺตํ.)}

\proseitem{325}{นเหตุยา เอกํ, นอารมฺมเณ ตีณิ, นอธิปติยา นว, นอนนฺตเร ตีณิ, นสมนนฺตเร ตีณิ, นอญฺญมญฺเญ ตีณิ, นอุปนิสฺสเย ตีณิ, นปุเรชาเต นว, นปจฺฉาชาเต นว, นอาเสวเน นว, นกมฺเม ตีณิ, นวิปาเก นว, นอาหาเร เอกํ, นอินฺทฺริเย เอกํ, นฌาเน เอกํ, นมคฺเค เอกํ, นสมฺปยุตฺเต ตีณิ, นวิปฺปยุตฺเต นว, โนนตฺถิยา ตีณิ, โนวิคเต ตีณิ. (สํขิตฺตํ.)}

\prose{\csromanfolio{268}เหตุปจฺจยา นอารมฺมเณ ตีณิ. (สํขิตฺตํ.)}

\prose{นเหตุปจฺจยา อารมฺมเณ เอกํ. (สํขิตฺตํ.)}

\prose{(สหชาตวาโรปิ ปจฺจยวาโรปิ นิสฺสยวาโรปิ สํสฏฺฐ\-วาโรปิ สมฺปยุตฺตวาโรปิ ปฏิจฺจ\-วาร\-สทิสา วิตฺถาเรตพฺพา.)}

\csromanmarkh{3. เนวาจยคามินาปจยคามิปท}\csromanmarkhii{7. ปญฺหาวาร}\csromanheadernomark{2}{3. \textbf{เนวาจยคามินาปจยคามิปท 7. ปญฺหาวาร}}
\prose{ปจฺจยจตุกฺกเหตุ อารมฺมณ}

\proseitem{326}{เหตุ เนวา\-จยคามิ\-นา\-ปจยคามี ธมฺโม เหตุสฺส เนวา\-จยคามิ\-นา\-ปจยคามิสฺส ธมฺมสฺส เหตุปจฺจเยน ปจฺจโย. ตีณิ.}

\prose{เหตุ เนวา\-จยคามิ\-นา\-ปจยคามี ธมฺโม เหตุสฺส เนวา\-จยคามิ\-นา\-ปจยคามิสฺส ธมฺมสฺส อารมฺมณ\-ปจฺจเยน ปจฺจโย. (สํขิตฺตํ.)}

\proseitem{327}{เหตุยา ตีณิ, อารมฺมเณ นว, อธิปติยา นว, อนนฺตเร นว, สมนนฺตเร นว, สหชาเต นว, อญฺญมญฺเญ นว, นิสฺสเย นว, อุปนิสฺสเย นว, ปุเรชาเต ตีณิ, ปจฺฉาชาเต ตีณิ, อาเสวเน นว, กมฺเม ตีณิ, วิปาเก นว, อาหาเร ตีณิ, อินฺทฺริเย นว, ฌาเน ตีณิ, มคฺเค นว, สมฺปยุตฺเต นว, วิปฺปยุตฺเต ปญฺจ, อตฺถิยา นว, นตฺถิยา นว, วิคเต นว, อวิคเต นว. (สํขิตฺตํ.)}

\proseitem{328}{เหตุ เนวา\-จยคามิ\-นา\-ปจยคามี ธมฺโม เหตุสฺส เนวา\-จยคามิ\-นา\-ปจยคามิสฺส ธมฺมสฺส อารมฺมณ\-ปจฺจเยน ปจฺจโย. สหชาต\-ปจฺจเยน ปจฺจโย. อุปนิสฺสย\-ปจฺจเยน ปจฺจโย. ตีณิ.}

\prose{เนวา\-จยคามิ\-นา\-ปจยคามี ธมฺโม นเหตุสฺส เนวา\-จยคามิ\-นา\-ปจยคามิสฺส ธมฺมสฺส อารมฺมณ\-ปจฺจเยน ปจฺจโย. สหชาต\-ปจฺจเยน ปจฺจโย. อุปนิสฺสย\-ปจฺจเยน ปจฺจโย. ปุเรชาต\-ปจฺจเยน ปจฺจโย. ปจฺฉาชาต\-ปจฺจเยน ปจฺจโย. อาหารปจฺจเยน ปจฺจโย. อินฺทฺริย\-ปจฺจเยน ปจฺจโย. (สํขิตฺตํ.)}

\csromanmatikaanchor{mtk.42}
\csromantocmarkref{section}{11. เสกฺขตฺติก}{mtk.42}
\bookmark[dest={mtk.42},level=1]{11. เสกฺขตฺติก}
\csromanmarkh{1. เหตุทุก}\csromanmarkhii{11. เสกฺขตฺติก}\csromanheadernomark{1}{1. เหตุทุก 11. เสกฺขตฺติก}
\proseitem{329}{\csromanfolio{269}นเหตุยา นว, นอารมฺมเณ นว. (สํขิตฺตํ.)}

\prose{เหตุปจฺจยา นอารมฺมเณ ตีณิ. (สํขิตฺตํ.)}

\prose{นเหตุปจฺจยา อารมฺมเณ นว. (สํขิตฺตํ.)}

\prosewithcloser{(ยถา กุสลตฺติเก ปญฺหาวารสฺส อนุโลมมฺปิ ปจฺจนียมฺปิ อนุโลมปจฺจนียมฺปิ ปจฺจนียานุโลมมฺปิ คณิตํ, เอวํ คเณตพฺพํ.)}{เหตุ\-ทุก\-อาจยคามิ\-ตฺติกํ นิฏฺฐิตํ.}
\csromanmarkh{1. เหตุทุก}\csromanmarkhii{11. เสกฺขตฺติก}\csromanheadernomark{2}{1. เหตุทุก 11. เสกฺขตฺติก}
\csromanmarkh{1. เสกฺขปท}\csromanmarkhii{1. ปฏิจฺจวาร}\csromanheadernomark{2}{1. \textbf{เสกฺขปท 1. ปฏิจฺจวาร}}
\prose{ปจฺจยจตุกฺกเหตุ}

\proseitem{330}{เหตุํ เสกฺขํ ธมฺมํ ปฏิจฺจ เหตุ เสกฺโข ธมฺโม อุปฺปชฺชติ เหตุปจฺจยา. ตีณิ.}

\prose{นเหตุํ เสกฺขํ ธมฺมํ ปฏิจฺจ นเหตุ เสกฺโข ธมฺโม อุปฺปชฺชติ เหตุปจฺจยา. ตีณิ.}

\prose{เหตุํ เสกฺขญฺจ นเหตุํ เสกฺขญฺจ ธมฺมํ ปฏิจฺจ เหตุ เสกฺโข ธมฺโม อุปฺปชฺชติ เหตุปจฺจยา. ตีณิ. (สํขิตฺตํ.)}

\proseitem{331}{เหตุยา นว, อารมฺมเณ นว, อธิปติยา นว, อนนฺตเร นว, สมนนฺตเร นว, สหชาเต นว, อญฺญมญฺเญ นว, นิสฺสเย นว, อุปนิสฺสเย นว, ปุเรชาเต นว, อาเสวเน นว, กมฺเม นว, วิปาเก นว, อาหาเร นว ฯเปฯ อวิคเต นว. (สํขิตฺตํ.)}

\prose{นอธิปติ}

\proseitem{332}{เหตุํ เสกฺขํ ธมฺมํ ปฏิจฺจ เหตุ เสกฺโข ธมฺโม อุปฺปชฺชติ นอธิปติ\-ปจฺจยา. (สํขิตฺตํ.)}

\proseitem{333}{\csromanfolio{270}นอธิปติยา ฉ, นปุเรชาเต นว, นปจฺฉาชาเต นว, นอาเสวเน นว, นกมฺเม ตีณิ, นวิปาเก นว, นวิปฺปยุตฺเต นว.}

\prose{เหตุปจฺจยา นอธิปติยา ฉ. (สํขิตฺตํ.)}

\prose{นอธิปติ\-ปจฺจยา เหตุยา ฉ. (สํขิตฺตํ.)}

\prose{(สหชาตวาโรปิ ปจฺจยวาโรปิ นิสฺสยวาโรปิ สํสฏฺฐ\-วาโรปิ สมฺปยุตฺตวาโรปิ ปฏิจฺจ\-วาร\-สทิสา วิตฺถาเรตพฺพา.)}

\csromanmarkh{1. เสกฺขปท}\csromanmarkhii{7. ปญฺหาวาร}\csromanheadernomark{2}{1. \textbf{เสกฺขปท 7. ปญฺหาวาร}}
\prose{ปจฺจยจตุกฺกเหตฺวาทิ}

\proseitem{334}{เหตุ เสกฺโข ธมฺโม เหตุสฺส เสกฺขสฺส ธมฺมสฺส เหตุปจฺจเยน ปจฺจโย. ตีณิ.}

\prose{เหตุ เสกฺโข ธมฺโม เหตุสฺส เสกฺขสฺส ธมฺมสฺส อธิปติ\-ปจฺจเยน ปจฺจโย. สหชาตา\-ธิปติ. ตีณิ.}

\prose{เสกฺโข ธมฺโม นเหตุสฺส เสกฺขสฺส ธมฺมสฺส อธิปติ\-ปจฺจเยน ปจฺจโย. สหชาตา\-ธิปติ. ตีณิ.}

\prose{เหตุ เสกฺโข ธมฺโม เหตุสฺส เสกฺขสฺส ธมฺมสฺส อนนฺตร\-ปจฺจเยน ปจฺจโย. (สํขิตฺตํ.)}

\proseitem{335}{เหตุยา ตีณิ, อธิปติยา ฉ, อนนฺตเร นว, สมนนฺตเร นว, สหชาเต นว, อญฺญมญฺเญ นว, นิสฺสเย นว, อุปนิสฺสเย นว, กมฺเม ตีณิ, วิปาเก นว, อาหาเร ตีณิ, อินฺทฺริเย นว, ฌาเน ตีณิ, มคฺเค นว, สมฺปยุตฺเต นว, อตฺถิยา นว, นตฺถิยา นว, วิคเต นว, อวิคเต นว. (สํขิตฺตํ.)}

\proseitem{336}{เหตุ เสกฺโข ธมฺโม เหตุสฺส เสกฺขสฺส ธมฺมสฺส สหชาต\-ปจฺจเยน ปจฺจโย. อุปนิสฺสย\-ปจฺจเยน ปจฺจโย. (สํขิตฺตํ.)}

\csromanmarkh{1. เหตุทุก}\csromanmarkhii{11. เสกฺขตฺติก}\csromanheadernomark{2}{1. เหตุทุก 11. เสกฺขตฺติก}
\proseitem{337}{\csromanfolio{271}นเหตุยา นว, นอารมฺมเณ นว. (สํขิตฺตํ.)}

\prose{เหตุปจฺจยา นอารมฺมเณ ตีณิ. (สํขิตฺตํ.)}

\prose{นเหตุปจฺจยา อธิปติยา ตีณิ. (สํขิตฺตํ.)}

\prose{(ยถา กุสลตฺติเก ปญฺหาวารสฺส อนุโลมมฺปิ ปจฺจนียมฺปิ อนุโลมปจฺจนียมฺปิ ปจฺจนียานุโลมมฺปิ คณิตํ, เอวํ คเณตพฺพํ.)}

\csromanmarkh{2. อเสกฺขปท}\csromanmarkhii{1. ปฏิจฺจวาร}\csromanheadernomark{2}{2. \textbf{อเสกฺขปท 1. ปฏิจฺจวาร}}
\prose{ปจฺจยจตุกฺกเหตุ}

\proseitem{338}{เหตุํ อเสกฺขํ ธมฺมํ ปฏิจฺจ เหตุ อเสกฺโข ธมฺโม อุปฺปชฺชติ เหตุปจฺจยา. ตีณิ.}

\prose{นเหตุํ อเสกฺขํ ธมฺมํ ปฏิจฺจ นเหตุ อเสกฺโข ธมฺโม อุปฺปชฺชติ เหตุปจฺจยา. ตีณิ.}

\prose{เหตุํ อเสกฺขญฺจ นเหตุํ อเสกฺขญฺจ ธมฺมํ ปฏิจฺจ เหตุ อเสกฺโข ธมฺโม อุปฺปชฺชติ เหตุปจฺจยา. ตีณิ. (สํขิตฺตํ.)}

\proseitem{339}{เหตุยา นว, อารมฺมเณ นว, อธิปติยา นว, อนนฺตเร นว, สมนนฺตเร นว, สหชาเต นว, อญฺญมญฺเญ นว, นิสฺสเย นว, อุปนิสฺสเย นว, ปุเรชาเต นว, กมฺเม นว, วิปาเก นว ฯเปฯ อวิคเต นว. (สํขิตฺตํ.)}

\prose{นอธิปติ}

\proseitem{340}{เหตุํ อเสกฺขํ ธมฺมํ ปฏิจฺจ เหตุ อเสกฺโข ธมฺโม อุปฺปชฺชติ นอธิปติ\-ปจฺจยา. (สํขิตฺตํ.)}

\proseitem{341}{นอธิปติยา ฉ, นปุเรชาเต นว, นปจฺฉาชาเต นว, นอาเสวเน นว, นวิปฺปยุตฺเต นว. (สํขิตฺตํ.)}

\prose{เหตุปจฺจยา นอธิปติยา ฉ. (สํขิตฺตํ.)}

\prose{นอธิปติ\-ปจฺจยา เหตุยา ฉ. (สํขิตฺตํ.)}

\prose{\csromanfolio{272}(สหชาตวาโรปิ ปจฺจยวาโรปิ นิสฺสยวาโรปิ สํสฏฺฐ\-วาโรปิ สมฺปยุตฺตวาโรปิ ปฏิจฺจ\-วาร\-สทิสา วิตฺถาเรตพฺพา.)}

\csromanmarkh{2. อเสกฺขปท}\csromanmarkhii{7. ปญฺหาวาร}\csromanheadernomark{2}{2. \textbf{อเสกฺขปท 7. ปญฺหาวาร}}
\prose{ปจฺจยจตุกฺกเหตฺวาทิ}

\proseitem{342}{เหตุ อเสกฺโข ธมฺโม เหตุสฺส อเสกฺขสฺส ธมฺมสฺส เหตุปจฺจเยน ปจฺจโย. ตีณิ.}

\prose{เหตุ อเสกฺโข ธมฺโม เหตุสฺส อเสกฺขสฺส ธมฺมสฺส อธิปติ\-ปจฺจเยน ปจฺจโย. สหชาตา\-ธิปติ. ตีณิ.}

\prose{อเสกฺโข ธมฺโม นเหตุสฺส อเสกฺขสฺส ธมฺมสฺส อธิปติ\-ปจฺจเยน ปจฺจโย. สหชาตา\-ธิปติ. ตีณิ ฯเปฯ.}

\prose{เหตุ อเสกฺโข ธมฺโม เหตุสฺส อเสกฺขสฺส ธมฺมสฺส อุปนิสฺสย\-ปจฺจเยน ปจฺจโย. อนนฺตรู\-ปนิสฺสโย. (สํขิตฺตํ.)}

\proseitem{343}{เหตุยา ตีณิ, อธิปติยา ฉ, อนนฺตเร นว, สมนนฺตเร นว, สหชาเต นว, อญฺญมญฺเญ นว, นิสฺสเย นว, อุปนิสฺสเย นว, กมฺเม ตีณิ, วิปาเก นว, อาหาเร ตีณิ, อินฺทฺริเย นว, ฌาเน ตีณิ, มคฺเค นว, สมฺปยุตฺเต นว, อตฺถิยา นว, นตฺถิยา นว, วิคเต นว, อวิคเต นว.}

\proseitem{344}{เหตุ อเสกฺโข ธมฺโม เหตุสฺส อเสกฺขสฺส ธมฺมสฺส สหชาต\-ปจฺจเยน ปจฺจโย. อุปนิสฺสย\-ปจฺจเยน ปจฺจโย. (สํขิตฺตํ.)}

\proseitem{345}{นเหตุยา นว, นอารมฺมเณ นว. (สํขิตฺตํ.)}

\prose{เหตุปจฺจยา นอารมฺมเณ ตีณิ. (สํขิตฺตํ.)}

\prose{นเหตุปจฺจยา อธิปติยา ตีณิ. (สํขิตฺตํ.)}

\prose{(ยถา กุสลตฺติเก ปญฺหาวารสฺส อนุโลมมฺปิ ปจฺจนียมฺปิ อนุโลมปจฺจนียมฺปิ ปจฺจนียานุโลมมฺปิ คณิตํ, เอวํ คเณตพฺพํ.)}

\vspace{\tipitakaparskip}
\csromancenter{เหตุทุก 11. เสกฺขตฺติก 3. เนวเสกฺขนาเสกฺขปท 1. ปฏิจฺจวาร}
\prose{\csromanfolio{273}ปจฺจยจตุกฺกเหตุ}

\proseitem{346}{เหตุํ เนว\-เสกฺข\-นาเสกฺขํ ธมฺมํ ปฏิจฺจ เหตุ เนว\-เสกฺข\-นา\-เสกฺโข ธมฺโม อุปฺปชฺชติ เหตุปจฺจยา. ตีณิ.}

\prose{นเหตุํ เนว\-เสกฺข\-นาเสกฺขํ ธมฺมํ ปฏิจฺจ นเหตุ เนว\-เสกฺข\-นา\-เสกฺโข ธมฺโม อุปฺปชฺชติ เหตุปจฺจยา. ตีณิ.}

\prose{เหตุํ เนว\-เสกฺข\-นาเสกฺขญฺจ นเหตุํ เนว\-เสกฺข\-นาเสกฺขญฺจ ธมฺมํ ปฏิจฺจ เหตุ เนว\-เสกฺข\-นา\-เสกฺโข ธมฺโม อุปฺปชฺชติ เหตุปจฺจยา. ตีณิ. (สํขิตฺตํ.)}

\proseitem{347}{เหตุยา นว, อารมฺมเณ นว, อธิปติยา นว, อนนฺตเร นว, สมนนฺตเร นว, สหชาเต นว, อญฺญมญฺเญ นว, นิสฺสเย นว, อุปนิสฺสเย นว, ปุเรชาเต นว, อาเสวเน นว, กมฺเม นว, วิปาเก นว ฯเปฯ อวิคเต นว. (สํขิตฺตํ.)}

\prose{นเหตุ นอารมฺมณาทิ}

\proseitem{348}{นเหตุํ เนว\-เสกฺข\-นาเสกฺขํ ธมฺมํ ปฏิจฺจ นเหตุ เนว\-เสกฺข\-นา\-เสกฺโข ธมฺโม อุปฺปชฺชติ นเหตุปจฺจยา.  นเหตุํ เนว\-เสกฺข\-นาเสกฺขํ ธมฺมํ ปฏิจฺจ เหตุ เนว\-เสกฺข\-นา\-เสกฺโข ธมฺโม อุปฺปชฺชติ นเหตุปจฺจยา. (2)}

\prose{เหตุํ เนว\-เสกฺข\-นาเสกฺขํ ธมฺมํ ปฏิจฺจ นเหตุ เนว\-เสกฺข\-นา\-เสกฺโข ธมฺโม อุปฺปชฺชติ นอารมฺมณ\-ปจฺจยา. (1)}

\prose{นเหตุํ เนว\-เสกฺข\-นาเสกฺขํ ธมฺมํ ปฏิจฺจ นเหตุ เนว\-เสกฺข\-นา\-เสกฺโข ธมฺโม อุปฺปชฺชติ นอารมฺมณ\-ปจฺจยา. (1)}

\prose{เหตุํ เนว\-เสกฺข\-นาเสกฺขญฺจ นเหตุํ เนว\-เสกฺข\-นาเสกฺขญฺจ ธมฺมํ ปฏิจฺจ นเหตุ เนว\-เสกฺข\-นา\-เสกฺโข ธมฺโม อุปฺปชฺชติ นอารมฺมณ\-ปจฺจยา. (1)}

\proseitem{349}{เหตุํ เนว\-เสกฺข\-นาเสกฺขํ ธมฺมํ ปฏิจฺจ เหตุ เนว\-เสกฺข\-นา\-เสกฺโข ธมฺโม อุปฺปชฺชติ นอธิปติ\-ปจฺจยา ฯเปฯ.}

\prose{\csromanfolio{274}เหตุํ เนว\-เสกฺข\-นาเสกฺขํ ธมฺมํ ปฏิจฺจ เหตุ เนว\-เสกฺข\-นา\-เสกฺโข ธมฺโม อุปฺปชฺชติ นปุเรชาต\-ปจฺจยา. ตีณิ. (สํขิตฺตํ.)}

\proseitem{350}{นเหตุยา เทฺว, นอารมฺมเณ ตีณิ, นอธิปติยา นว, นอนนฺตเร ตีณิ, นสมนนฺตเร ตีณิ, นอญฺญมญฺเญ ตีณิ, นอุปนิสฺสเย ตีณิ, นปุเรชาเต นว, นปจฺฉาชาเต นว, นอาเสวเน นว, นกมฺเม ตีณิ, นวิปาเก นว, นอาหาเร เอกํ, นอินฺทฺริเย เอกํ, นฌาเน เอกํ, นมคฺเค เอกํ, นสมฺปยุตฺเต ตีณิ, นวิปฺปยุตฺเต นว, โนนตฺถิยา ตีณิ, โนวิคเต ตีณิ. (สํขิตฺตํ.)}

\prose{เหตุปจฺจยา นอารมฺมเณ ตีณิ. (สํขิตฺตํ.)}

\prose{นเหตุปจฺจยา อารมฺมเณ เทฺว. (สํขิตฺตํ.)}

\prose{(สหชาตวาโรปิ ปจฺจยวาโรปิ นิสฺสยวาโรปิ สํสฏฺฐ\-วาโรปิ สมฺปยุตฺตวาโรปิ ปฏิจฺจ\-วาร\-สทิสา วิตฺถาเรตพฺพา.)}

\csromanmarkh{3. เนวเสกฺขนาเสกฺขปท}\csromanmarkhii{7. ปญฺหาวาร}\csromanheadernomark{2}{3. \textbf{เนวเสกฺขนาเสกฺขปท 7. ปญฺหาวาร}}
\prose{ปจฺจยจตุกฺกเหตฺวาทิ}

\proseitem{351}{เหตุ เนว\-เสกฺข\-นา\-เสกฺโข ธมฺโม เหตุสฺส เนว\-เสกฺข\-นาเสกฺขสฺส ธมฺมสฺส เหตุปจฺจเยน ปจฺจโย. ตีณิ.}

\prose{เหตุ เนว\-เสกฺข\-นา\-เสกฺโข ธมฺโม เหตุสฺส เนว\-เสกฺข\-นาเสกฺขสฺส ธมฺมสฺส อารมฺมณ\-ปจฺจเยน ปจฺจโย. นว.}

\prose{เหตุ เนว\-เสกฺข\-นา\-เสกฺโข ธมฺโม เหตุสฺส เนว\-เสกฺข\-นาเสกฺขสฺส ธมฺมสฺส อธิปติ\-ปจฺจเยน ปจฺจโย. อารมฺมณา\-ธิปติ, สหชาตา\-ธิปติ. นว ฯเปฯ.}

\prose{เนว\-เสกฺข\-นา\-เสกฺโข ธมฺโม นเหตุสฺส เนว\-เสกฺข\-นาเสกฺขสฺส ธมฺมสฺส ปุเรชาต\-ปจฺจเยน ปจฺจโย. (สํขิตฺตํ.)}

\proseitem{352}{เหตุยา ตีณิ, อารมฺมเณ นว, อธิปติยา นว, อนนฺตเร นว, สมนนฺตเร นว, สหชาเต นว, อญฺญมญฺเญ นว, นิสฺสเย นว,}

\vspace{\tipitakaparskip}
\csromancenter{เหตุทุก 12. ปริตฺตตฺติก อุปนิสฺสเย นว, ปุเรชาเต ตีณิ, ปจฺฉาชาเต ตีณิ, อาเสวเน นว, กมฺเม ตีณิ, วิปาเก นว, อาหาเร ตีณิ, อินฺทฺริเย นว, ฌาเน ตีณิ, มคฺเค นว, สมฺปยุตฺเต นว, วิปฺปยุตฺเต ปญฺจ, อตฺถิยา นว, นตฺถิยา นว, วิคเต นว, อวิคเต นว. (สํขิตฺตํ.)}
\proseitem{353}{\csromanfolio{275}เหตุ เนว\-เสกฺข\-นา\-เสกฺโข ธมฺโม เหตุสฺส เนว\-เสกฺข\-นาเสกฺขสฺส ธมฺมสฺส อารมฺมณ\-ปจฺจเยน ปจฺจโย. สหชาต\-ปจฺจเยน ปจฺจโย. อุปนิสฺสย\-ปจฺจเยน ปจฺจโย. (สํขิตฺตํ.)}

\proseitem{354}{นเหตุยา นว, นอารมฺมเณ นว. (สํขิตฺตํ.)}

\prose{เหตุปจฺจยา นอารมฺมเณ ตีณิ. (สํขิตฺตํ.)}

\prose{นเหตุปจฺจยา อารมฺมเณ นว. (สํขิตฺตํ.)}

\prosewithcloser{(ยถา กุสลตฺติเก ปญฺหาวารสฺส อนุโลมมฺปิ ปจฺจนียมฺปิ อนุโลมปจฺจนียมฺปิ ปจฺจนียานุโลมมฺปิ คณิตํ, เอวํ คเณตพฺพํ.)}{เหตุ\-ทุก\-เสกฺข\-ตฺติกํ นิฏฺฐิตํ.}
\csromanmatikaanchor{mtk.43}
\csromantocmarkref{section}{12. ปริตฺตตฺติก}{mtk.43}
\bookmark[dest={mtk.43},level=1]{12. ปริตฺตตฺติก}
\csromanmarkh{1. เหตุทุก}\csromanmarkhii{12. ปริตฺตตฺติก}\csromanheadernomark{1}{1. เหตุทุก 12. ปริตฺตตฺติก}
\csromanmarkh{1. ปริตฺตปท}\csromanmarkhii{1. ปฏิจฺจวาร}\csromanheadernomark{2}{1. \textbf{ปริตฺตปท 1. ปฏิจฺจวาร}}
\prose{ปจฺจยจตุกฺกเหตุ}

\proseitem{355}{เหตุํ ปริตฺตํ ธมฺมํ ปฏิจฺจ เหตุ ปริตฺโต ธมฺโม อุปฺปชฺชติ เหตุปจฺจยา. ตีณิ.}

\prose{นเหตุํ ปริตฺตํ ธมฺมํ ปฏิจฺจ นเหตุ ปริตฺโต ธมฺโม อุปฺปชฺชติ เหตุปจฺจยา. ตีณิ.}

\prose{เหตุํ ปริตฺตญฺจ นเหตุํ ปริตฺตญฺจ ธมฺมํ ปฏิจฺจ เหตุ ปริตฺโต ธมฺโม อุปฺปชฺชติ เหตุปจฺจยา. ตีณิ. (สํขิตฺตํ.)}

\proseitem{356}{\csromanfolio{276}เหตุยา นว, อารมฺมเณ นว, อธิปติยา นว ฯเปฯ อวิคเต นว. (สํขิตฺตํ.)}

\prose{นเหตุ นอารมฺมณาทิ}

\proseitem{357}{นเหตุํ ปริตฺตํ ธมฺมํ ปฏิจฺจ นเหตุ ปริตฺโต ธมฺโม อุปฺปชฺชติ นเหตุปจฺจยา.  นเหตุํ ปริตฺตํ ธมฺมํ ปฏิจฺจ เหตุ ปริตฺโต ธมฺโม อุปฺปชฺชติ นเหตุปจฺจยา. (2)}

\prose{เหตุํ ปริตฺตํ ธมฺมํ ปฏิจฺจ นเหตุ ปริตฺโต ธมฺโม อุปฺปชฺชติ นอารมฺมณ\-ปจฺจยา. (1)}

\prose{นเหตุํ ปริตฺตํ ธมฺมํ ปฏิจฺจ นเหตุ ปริตฺโต ธมฺโม อุปฺปชฺชติ นอารมฺมณ\-ปจฺจยา. (1)}

\prose{เหตุํ ปริตฺตญฺจ นเหตุํ ปริตฺตญฺจ ธมฺมํ ปฏิจฺจ นเหตุ ปริตฺโต ธมฺโม อุปฺปชฺชติ นอารมฺมณ\-ปจฺจยา. (1)}

\prose{เหตุํ ปริตฺตํ ธมฺมํ ปฏิจฺจ เหตุ ปริตฺโต ธมฺโม อุปฺปชฺชติ นอธิปติ\-ปจฺจยา. (สํขิตฺตํ.)}

\proseitem{358}{นเหตุยา เทฺว, นอารมฺมเณ ตีณิ, นอธิปติยา นว, นอนนฺตเร ตีณิ, นสมนนฺตเร ตีณิ, นอญฺญมญฺเญ ตีณิ, นอุปนิสฺสเย ตีณิ, นปุเรชาเต นว, นปจฺฉาชาเต นว, นอาเสวเน นว, นกมฺเม ตีณิ, นวิปาเก นว, นอาหาเร เอกํ, นอินฺทฺริเย เอกํ, นฌาเน เอกํ, นมคฺเค เอกํ, นสมฺปยุตฺเต ตีณิ, นวิปฺปยุตฺเต นว, โนนตฺถิยา ตีณิ, โนวิคเต ตีณิ. (สํขิตฺตํ.)}

\prose{เหตุปจฺจยา นอารมฺมเณ ตีณิ. (สํขิตฺตํ.)}

\prose{นเหตุปจฺจยา อารมฺมเณ เทฺว. (สํขิตฺตํ.)}

\prose{(สหชาตวาโรปิ ปจฺจยวาโรปิ นิสฺสยวาโรปิ สํสฏฺฐ\-วาโรปิ สมฺปยุตฺตวาโรปิ ปฏิจฺจ\-วาร\-สทิสา วิตฺถาเรตพฺพา.)}

\csromanmarkh{1. เหตุทุก}\csromanmarkhii{12. ปริตฺตตฺติก 1. ปริตฺตปท 7. ปญฺหาวาร}\csromanheadernomark{2}{1. เหตุทุก 12. ปริตฺตตฺติก \textbf{1. ปริตฺตปท 7. ปญฺหาวาร}}
\prose{\csromanfolio{277}ปจฺจยจตุกฺกเหตุ อารมฺมณ}

\proseitem{359}{เหตุ ปริตฺโต ธมฺโม เหตุสฺส ปริตฺตสฺส ธมฺมสฺส เหตุปจฺจเยน ปจฺจโย. ตีณิ.}

\prose{เหตุ ปริตฺโต ธมฺโม เหตุสฺส ปริตฺตสฺส ธมฺมสฺส อารมฺมณ\-ปจฺจเยน ปจฺจโย. (สํขิตฺตํ.)}

\proseitem{360}{เหตุยา ตีณิ, อารมฺมเณ นว, อธิปติยา นว, อนนฺตเร นว, สมนนฺตเร นว ฯเปฯ อวิคเต นว. (สํขิตฺตํ.)}

\proseitem{361}{เหตุ ปริตฺโต ธมฺโม เหตุสฺส ปริตฺตสฺส ธมฺมสฺส อารมฺมณ\-ปจฺจเยน ปจฺจโย. สหชาต\-ปจฺจเยน ปจฺจโย. อุปนิสฺสย\-ปจฺจเยน ปจฺจโย. (สํขิตฺตํ.)}

\proseitem{362}{นเหตุยา นว, นอารมฺมเณ นว. (สํขิตฺตํ.)}

\prose{เหตุปจฺจยา นอารมฺมเณ ตีณิ. (สํขิตฺตํ.)}

\prose{นเหตุปจฺจยา อารมฺมเณ นว. (สํขิตฺตํ.)}

\prose{(ยถา กุสลตฺติเก ปญฺหาวารสฺส อนุโลมมฺปิ ปจฺจนียมฺปิ อนุโลมปจฺจนียมฺปิ ปจฺจนียานุโลมมฺปิ คณิตํ, เอวํ คเณตพฺพํ.)}

\csromanmarkh{2. มหคฺคตปท}\csromanmarkhii{1. ปฏิจฺจวาร}\csromanheadernomark{2}{2. \textbf{มหคฺคตปท 1. ปฏิจฺจวาร}}
\prose{ปจฺจยจตุกฺกเหตุ}

\proseitem{363}{เหตุํ มหคฺคตํ ธมฺมํ ปฏิจฺจ เหตุ มหคฺคโต ธมฺโม อุปฺปชฺชติ เหตุปจฺจยา. ตีณิ.}

\prose{\csromanfolio{278}นเหตุํ มหคฺคตํ ธมฺมํ ปฏิจฺจ นเหตุ มหคฺคโต ธมฺโม อุปฺปชฺชติ เหตุปจฺจยา. ตีณิ.}

\prose{เหตุํ มหคฺคตญฺจ นเหตุํ มหคฺคตญฺจ ธมฺมํ ปฏิจฺจ เหตุ มหคฺคโต ธมฺโม อุปฺปชฺชติ เหตุปจฺจยา. ตีณิ. (สํขิตฺตํ.)}

\proseitem{364}{เหตุยา นว, อารมฺมเณ นว, อธิปติยา นว, อนนฺตเร นว, สมนนฺตเร นว, สหชาเต นว, อญฺญมญฺเญ นว, นิสฺสเย นว, อุปนิสฺสเย นว, ปุเรชาเต นว, อาเสวเน นว, กมฺเม นว, วิปาเก นว, อาหาเร นว ฯเปฯ อวิคเต นว. (สํขิตฺตํ.)}

\prose{นอธิปติ}

\proseitem{365}{เหตุํ มหคฺคตํ ธมฺมํ ปฏิจฺจ เหตุ มหคฺคโต ธมฺโม อุปฺปชฺชติ นอธิปติ\-ปจฺจยา. (สํขิตฺตํ.)}

\proseitem{366}{นอธิปติยา นว, นปุเรชาเต นว, นปจฺฉาชาเต นว, นอาเสวเน นว, นกมฺเม ตีณิ, นวิปาเก นว, นวิปฺปยุตฺเต นว. (สํขิตฺตํ.)}

\prose{เหตุปจฺจยา นอธิปติยา นว. (สํขิตฺตํ.)}

\prose{นอธิปติ\-ปจฺจยา เหตุยา นว. (สํขิตฺตํ.)}

\prose{(สหชาตวาโรปิ ปจฺจยวาโรปิ นิสฺสยวาโรปิ สํสฏฺฐ\-วาโรปิ สมฺปยุตฺตวาโรปิ ปฏิจฺจ\-วาร\-สทิสา วิตฺถาเรตพฺพา.)}

\csromanmarkh{2. มหคฺคตปท}\csromanmarkhii{7. ปญฺหาวาร}\csromanheadernomark{2}{2. \textbf{มหคฺคตปท 7. ปญฺหาวาร}}
\prose{ปจฺจยจตุกฺกเหตฺวาทิ}

\proseitem{367}{เหตุ มหคฺคโต ธมฺโม เหตุสฺส มหคฺคตสฺส ธมฺมสฺส เหตุปจฺจเยน ปจฺจโย. ตีณิ.}

\prose{เหตุ มหคฺคโต ธมฺโม เหตุสฺส มหคฺคตสฺส ธมฺมสฺส อารมฺมณ\-ปจฺจเยน ปจฺจโย. ตีณิ.}

\prose{มหคฺคโต ธมฺโม นเหตุสฺส มหคฺคตสฺส ธมฺมสฺส อารมฺมณ\-ปจฺจเยน ปจฺจโย. ตีณิ.}

\csromanmarkh{1. เหตุทุก}\csromanmarkhii{12. ปริตฺตตฺติก}\csromanheadernomark{2}{1. เหตุทุก 12. ปริตฺตตฺติก}
\prose{\csromanfolio{279}เหตุ มหคฺคโต จ นเหตุ มหคฺคโต จ ธมฺมา เหตุสฺส มหคฺคตสฺส ธมฺมสฺส อารมฺมณ\-ปจฺจเยน ปจฺจโย. ตีณิ.}

\prose{เหตุ มหคฺคโต ธมฺโม เหตุสฺส มหคฺคตสฺส ธมฺมสฺส อธิปติ\-ปจฺจเยน ปจฺจโย. สหชาตา\-ธิปติ. ตีณิ.}

\prose{มหคฺคโต ธมฺโม นเหตุสฺส มหคฺคตสฺส ธมฺมสฺส อธิปติ\-ปจฺจเยน ปจฺจโย. สหชาตา\-ธิปติ. ตีณิ.}

\prose{เหตุ มหคฺคโต ธมฺโม เหตุสฺส มหคฺคตสฺส ธมฺมสฺส อนนฺตร\-ปจฺจเยน ปจฺจโย. (สํขิตฺตํ.)}

\proseitem{368}{เหตุยา ตีณิ, อารมฺมเณ นว, อธิปติยา ฉ, อนนฺตเร นว, สมนนฺตเร นว, สหชาเต นว, อญฺญมญฺเญ นว, นิสฺสเย นว, อุปนิสฺสเย นว, อาเสวเน นว, กมฺเม ตีณิ, วิปาเก นว, อาหาเร ตีณิ, อินฺทฺริเย นว, ฌาเน ตีณิ, มคฺเค นว, สมฺปยุตฺเต นว, อตฺถิยา นว ฯเปฯ อวิคเต นว. (สํขิตฺตํ.)}

\proseitem{369}{เหตุ มหคฺคโต ธมฺโม เหตุสฺส มหคฺคตสฺส ธมฺมสฺส อรมฺมณปจฺจเยน ปจฺจโย. สหชาต\-ปจฺจเยน ปจฺจโย. อุปนิสฺสย\-ปจฺจเยน ปจฺจโย. (สํขิตฺตํ.)}

\proseitem{370}{นเหตุยา นว, นอารมฺมเณ นว. (สํขิตฺตํ.)}

\prose{เหตุปจฺจยา นอารมฺมเณ ตีณิ. (สํขิตฺตํ.)}

\prose{นเหตุปจฺจยา อารมฺมเณ นว. (สํขิตฺตํ.)}

\prose{(ยถา กุสลตฺติเก ปญฺหาวารสฺส อนุโลมมฺปิ ปจฺจนียมฺปิ อนุโลมปจฺจนียมฺปิ ปจฺจนียานุโลมมฺปิ คณิตํ, เอวํ คเณตพฺพํ.)}

\csromanmarkh{3. อปฺปมาณปท}\csromanmarkhii{1. ปฏิจฺจวาร}\csromanheadernomark{2}{3. \textbf{อปฺปมาณปท 1. ปฏิจฺจวาร}}
\prose{ปจฺจยจตุกฺกเหตุ}

\proseitem{371}{เหตุํ อปฺปมาณํ ธมฺมํ ปฏิจฺจ เหตุ อปฺปมาโณ ธมฺโม อุปฺปชฺชติ เหตุปจฺจยา. ตีณิ.}

\prose{\csromanfolio{280}นเหตุํ อปฺปมาณํ ธมฺมํ ปฏิจฺจ นเหตุ อปฺปมาโณ ธมฺโม อุปฺปชฺชติ เหตุปจฺจยา. ตีณิ.}

\prose{เหตุํ อปฺปมาณญฺจ นเหตุํ อปฺปมาณญฺจ ธมฺมํ ปฏิจฺจ เหตุ อปฺปมาโณ ธมฺโม อุปฺปชฺชติ เหตุปจฺจยา. ตีณิ. (สํขิตฺตํ.)}

\proseitem{372}{เหตุยา นว, อารมฺมเณ นว, อธิปติยา นว ฯเปฯ อุปนิสฺสเย นว ฯเปฯ กมฺเม นว, วิปาเก นว ฯเปฯ อวิคเต นว. (สํขิตฺตํ.)}

\prose{นอธิปติ}

\proseitem{373}{เหตุํ อปฺปมาณํ ธมฺมํ ปฏิจฺจ เหตุ อปฺปมาโณ ธมฺโม อุปฺปชฺชติ นอธิปติยา. (สํขิตฺตํ.)}

\proseitem{374}{นอธิปติยา ฉ, นปุเรชาเต นว, นปจฺฉาชาเต นว, นอาเสวเน นว, นกมฺเม ตีณิ, นวิปาเก นว, นวิปฺปยุตฺเต นว. (สํขิตฺตํ.)}

\prose{เหตุปจฺจยา นอธิปติยา ฉ. (สํขิตฺตํ.)}

\prose{นอธิปติ\-ปจฺจยา เหตุยา ฉ. (สํขิตฺตํ.)}

\prose{(สหชาตวาโรปิ ปจฺจยวาโรปิ นิสฺสยวาโรปิ สํสฏฺฐ\-วาโรปิ สมฺปยุตฺตวาโรปิ ปฏิจฺจ\-วาร\-สทิสา วิตฺถาเรตพฺพา.)}

\csromanmarkh{3. อปฺปมาณปท}\csromanmarkhii{7. ปญฺหาวาร}\csromanheadernomark{2}{3. \textbf{อปฺปมาณปท 7. ปญฺหาวาร}}
\prose{ปจฺจยจตุกฺกเหตฺวาทิ}

\proseitem{375}{เหตุ อปฺปมาโณ ธมฺโม เหตุสฺส อปฺปมาณสฺส ธมฺมสฺส เหตุปจฺจเยน ปจฺจโย. ตีณิ.}

\prose{อปฺปมาโณ ธมฺโม นเหตุสฺส อปฺปมาณสฺส ธมฺมสฺส อารมฺมณ\-ปจฺจเยน ปจฺจโย. ตีณิ.}

\proseitem{376}{เหตุ อปฺปมาโณ ธมฺโม เหตุสฺส อปฺปมาณสฺส ธมฺมสฺส อธิปติ\-ปจฺจเยน ปจฺจโย. สหชาตา\-ธิปติ. ตีณิ.}

\csromanmarkh{1. เหตุทุก}\csromanmarkhii{12. ปริตฺตตฺติก}\csromanheadernomark{2}{1. เหตุทุก 12. ปริตฺตตฺติก}
\prose{\csromanfolio{281}อปฺปมาโณ ธมฺโม นเหตุสฺส อปฺปมาณสฺส ธมฺมสฺส อธิปติ\-ปจฺจเยน ปจฺจโย. อารมฺมณา\-ธิปติ, สหชาตา\-ธิปติ. ตีณิ.}

\proseitem{377}{เหตุ อปฺปมาโณ ธมฺโม เหตุสฺส อปฺปมาณสฺส ธมฺมสฺส อนนฺตร\-ปจฺจเยน ปจฺจโย. นว ฯเปฯ.}

\prose{เหตุ อปฺปมาโณ ธมฺโม เหตุสฺส อปฺปมาณสฺส ธมฺมสฺส อุปนิสฺสย\-ปจฺจเยน ปจฺจโย. อนนฺตรู\-ปนิสฺสโย, ปกตู\-ปนิสฺสโย. นว. (สํขิตฺตํ.)}

\proseitem{378}{เหตุยา ตีณิ, อารมฺมเณ ตีณิ, อธิปติยา ฉ, อนนฺตเร นว, สมนนฺตเร นว, สหชาเต นว, อญฺญมญฺเญ นว, นิสฺสเย นว, อุปนิสฺสเย นว, กมฺเม ตีณิ, วิปาเก นว, อาหาเร ตีณิ, อินฺทฺริเย นว, ฌาเน ตีณิ, มคฺเค นว, สมฺปยุตฺเต นว, อตฺถิยา นว, นตฺถิยา นว, วิคเต นว, อวิคเต นว. (สํขิตฺตํ.)}

\proseitem{379}{เหตุ อปฺปมาโณ ธมฺโม เหตุสฺส อปฺปมาณสฺส ธมฺมสฺส สหชาต\-ปจฺจเยน ปจฺจโย. อุปนิสฺสย\-ปจฺจเยน ปจฺจโย. (สํขิตฺตํ.)}

\proseitem{380}{นเหตุยา นว, นอารมฺมเณ นว. (สํขิตฺตํ.)}

\prose{เหตุปจฺจยา นอารมฺมเณ ตีณิ. (สํขิตฺตํ.)}

\prose{นเหตุปจฺจยา อารมฺมเณ ตีณิ. (สํขิตฺตํ.)}

\prosewithcloser{(ยถา กุสลตฺติเก ปญฺหาวารสฺส อนุโลมมฺปิ ปจฺจนียมฺปิ อนุโลมปจฺจนียมฺปิ ปจฺจนียานุโลมมฺปิ คณิตํ, เอวํ คเณตพฺพํ.)}{เหตุ\-ทุก\-ปริตฺตตฺติกํ นิฏฺฐิตํ.}
\csromanmatikaanchor{mtk.44}
\csromantocmarkref{section}{13. ปริตฺตารมฺมณตฺติก}{mtk.44}
\bookmark[dest={mtk.44},level=1]{13. ปริตฺตารมฺมณตฺติก}
\csromanmarkh{ธมฺมานุโลม ทุกติก\-ปฏฺฐาน\-ปาฬิ}\csromanmarkhii{1. เหตุทุก 13. ปริตฺตา\-รมฺมณ\-ตฺติก}\csromanheadernomark{1}{ธมฺมานุโลม ทุกติก\-ปฏฺฐาน\-ปาฬิ 1. เหตุทุก 13. ปริตฺตา\-รมฺมณ\-ตฺติก}
\csromanmarkh{1. ปริตฺตา\-รมฺมณปท}\csromanmarkhii{1. ปฏิจฺจวาร}\csromanheadernomark{2}{1. ปริตฺตา\-รมฺมณปท 1. ปฏิจฺจวาร}
\prose{\csromanfolio{282}ปจฺจยจตุกฺกเหตุ}

\proseitem{381}{เหตุํ ปริตฺตา\-รมฺมณํ ธมฺมํ ปฏิจฺจ เหตุ ปริตฺตารมฺมโณ ธมฺโม อุปฺปชฺชติ เหตุปจฺจยา. ตีณิ.}

\prose{นเหตุํ ปริตฺตา\-รมฺมณํ ธมฺมํ ปฏิจฺจ นเหตุ ปริตฺตารมฺมโณ ธมฺโม อุปฺปชฺชติ เหตุปจฺจยา. ตีณิ.}

\prose{เหตุํ ปริตฺตา\-รมฺมณญฺจ นเหตุํ ปริตฺตา\-รมฺมณญฺจ ธมฺมํ ปฏิจฺจ เหตุ ปริตฺตารมฺมโณ ธมฺโม อุปฺปชฺชติ เหตุปจฺจยา. ตีณิ. (สํขิตฺตํ.)}

\proseitem{382}{เหตุยา นว, อารมฺมเณ นว, อธิปติยา นว, อนนฺตเร นว, สมนนฺตเร นว, สหชาเต นว, อญฺญมญฺเญ นว, นิสฺสเย นว, อุปนิสฺสเย นว, ปุเรชาเต นว, อาเสวเน นว, กมฺเม นว, วิปาเก นว ฯเปฯ อวิคเต นว. (สํขิตฺตํ.)}

\prose{นเหตุ นอธิปติ}

\proseitem{383}{นเหตุํ ปริตฺตา\-รมฺมณํ ธมฺมํ ปฏิจฺจ นเหตุ ปริตฺตารมฺมโณ ธมฺโม อุปฺปชฺชติ นเหตุปจฺจยา.  นเหตุํ ปริตฺตา\-รมฺมณํ ธมฺมํ ปฏิจฺจ เหตุ ปริตฺตารมฺมโณ ธมฺโม อุปฺปชฺชติ นเหตุปจฺจยา. (2)}

\prose{เหตุํ ปริตฺตา\-รมฺมณํ ธมฺมํ ปฏิจฺจ เหตุ ปริตฺตารมฺมโณ ธมฺโม อุปฺปชฺชติ นอธิปติ\-ปจฺจยา. (สํขิตฺตํ.)}

\proseitem{384}{นเหตุยา เทฺว, นอธิปติยา นว, นปุเรชาเต นว, นปจฺฉาชาเต นว, นอาเสวเน นว, นกมฺเม ตีณิ, นวิปาเก นว, นฌาเน เอกํ, นมคฺเค เอกํ, นวิปฺปยุตฺเต นว. (สํขิตฺตํ.)}

\prose{เหตุปจฺจยา นอธิปติยา นว. (สํขิตฺตํ.)}

\prose{นเหตุปจฺจยา อารมฺมเณ เทฺว. (สํขิตฺตํ.)}

\prose{(สหชาตวาโรปิ ปจฺจยวาโรปิ นิสฺสยวาโรปิ สํสฏฺฐ\-วาโรปิ สมฺปยุตฺตวาโรปิ ปฏิจฺจ\-วาร\-สทิสา วิตฺถาเรตพฺพา.)}

\vspace{\tipitakaparskip}
\csromancenter{เหตุทุก 13. ปริตฺตา\-รมฺมณ\-ตฺติก 1. ปริตฺตา\-รมฺมณปท 7. ปญฺหาวาร}
\prose{\csromanfolio{283}ปจฺจยจตุกฺกเหตฺวาทิ}

\proseitem{385}{เหตุ ปริตฺตารมฺมโณ ธมฺโม เหตุสฺส ปริตฺตา\-รมฺมณสฺส ธมฺมสฺส เหตุปจฺจเยน ปจฺจโย. ตีณิ.}

\prose{เหตุ ปริตฺตารมฺมโณ ธมฺโม เหตุสฺส ปริตฺตา\-รมฺมณสฺส ธมฺมสฺส อารมฺมณ\-ปจฺจเยน ปจฺจโย. นว.}

\prose{เหตุ ปริตฺตารมฺมโณ ธมฺโม เหตุสฺส ปริตฺตา\-รมฺมณสฺส ธมฺมสฺส อธิปติ\-ปจฺจเยน ปจฺจโย. (สํขิตฺตํ.)}

\proseitem{368}{เหตุยา ตีณิ, อารมฺมเณ นว, อธิปติยา นว, อนนฺตเร นว, สมนนฺตเร นว, สหชาเต นว, อญฺญมญฺเญ นว, นิสฺสเย นว, อุปนิสฺสเย นว, อาเสวเน นว, กมฺเม ตีณิ, วิปาเก นว, อาหาเร ตีณิ, อินฺทฺริเย นว, ฌาเน ตีณิ, มคฺเค นว, สมฺปยุตฺเต นว, อตฺถิยา นว, นตฺถิยา นว, วิคเต นว, อวิคเต นว. (สํขิตฺตํ.)}

\proseitem{387}{เหตุ ปริตฺตารมฺมโณ ธมฺโม เหตุสฺส ปริตฺตา\-รมฺมณสฺส ธมฺมสฺส อารมฺมณ\-ปจฺจเยน ปจฺจโย. สหชาต\-ปจฺจเยน ปจฺจโย. อุปนิสฺสย\-ปจฺจเยน ปจฺจโย. (สํขิตฺตํ.)}

\proseitem{388}{นเหตุยา นว, นอารมฺมเณ นว. (สํขิตฺตํ.)}

\prose{เหตุปจฺจยา นอารมฺมเณ ตีณิ. (สํขิตฺตํ.)}

\prose{นเหตุปจฺจยา อารมฺมเณ นว. (สํขิตฺตํ.)}

\prose{(ยถา กุสลตฺติเก ปญฺหาวารสฺส อนุโลมมฺปิ ปจฺจนียมฺปิ อนุโลมปจฺจนียมฺปิ ปจฺจนียานุโลมมฺปิ คณิตํ, เอวํ คเณตพฺพํ.)}

\csromanmarkh{ธมฺมานุโลม ทุกติก\-ปฏฺฐาน\-ปาฬิ}\csromanmarkhii{2. มหคฺคตา\-รมฺมณปท 1. ปฏิจฺจวาร}\csromanheadernomark{2}{ธมฺมานุโลม ทุกติก\-ปฏฺฐาน\-ปาฬิ 2. มหคฺคตา\-รมฺมณปท 1. ปฏิจฺจวาร}
\prose{\csromanfolio{284}ปจฺจยจตุกฺกเหตุ}

\proseitem{389}{เหตุํ มหคฺคตา\-รมฺมณํ ธมฺมํ ปฏิจฺจ เหตุ มหคฺคตา\-รมฺมโณ ธมฺโม อุปฺปชฺชติ เหตุปจฺจยา. ตีณิ.}

\prose{นเหตุํ มหคฺคตา\-รมฺมณํ ธมฺมํ ปฏิจฺจ นเหตุ มหคฺคตา\-รมฺมโณ ธมฺโม อุปฺปชฺชติ เหตุปจฺจยา ตีณิ.}

\prose{เหตุํ มหคฺคตา\-รมฺมณญฺจ นเหตุํ มหคฺคตา\-รมฺมณญฺจ ธมฺมํ ปฏิจฺจ เหตุ มหคฺคตา\-รมฺมโณ ธมฺโม อุปฺปชฺชติ เหตุปจฺจยา. ตีณิ. (สํขิตฺตํ.)}

\proseitem{390}{เหตุยา นว, อารมฺมเณ นว, อธิปติยา นว ฯเปฯ กมฺเม นว, วิปาเก นว, อาหาเร นว ฯเปฯ อวิคเต นว. (สํขิตฺตํ.)}

\prose{นเหตุ นอธิปติ}

\proseitem{391}{นเหตุํ มหคฺคตา\-รมฺมณํ ธมฺมํ ปฏิจฺจ นเหตุ มหคฺคตา\-รมฺมโณ ธมฺโม อุปฺปชฺชติ นเหตุปจฺจยา.  นเหตุํ มหคฺคตา\-รมฺมณํ ธมฺมํ ปฏิจฺจ เหตุ มหคฺคตา\-รมฺมโณ ธมฺโม อุปฺปชฺชติ นเหตุปจฺจยา. (2)}

\prose{เหตุํ มหคฺคตา\-รมฺมณํ ธมฺมํ ปฏิจฺจ เหตุ มหคฺคตา\-รมฺมโณ ธมฺโม อุปฺปชฺชติ นอธิปติ\-ปจฺจยา. (สํขิตฺตํ.)}

\proseitem{392}{นเหตุยา เทฺว, นอธิปติยา นว, นปุเรชาเต นว, นปจฺฉาชาเต นว, นอาเสวเน นว, นกมฺเม ตีณิ, นวิปาเก นว, นมคฺเค เอกํ, นวิปฺปยุตฺเต นว. (สํขิตฺตํ.)}

\prose{เหตุปจฺจยา นอธิปติยา นว. (สํขิตฺตํ.)}

\prose{นเหตุปจฺจยา อารมฺมเณ เทฺว. (สํขิตฺตํ.)}

\prose{(สหชาตวาโรปิ ปจฺจยวาโรปิ นิสฺสยวาโรปิ สํสฏฺฐ\-วาโรปิ สมฺปยุตฺตวาโรปิ ปฏิจฺจ\-วาร\-สทิสา วิตฺถาเรตพฺพา.)}

\vspace{\tipitakaparskip}
\csromancenter{เหตุทุก 13. ปริตฺตา\-รมฺมณ\-ตฺติก 2. มหคฺคตา\-รมฺมณปท 7. ปญฺหาวาร}
\prose{\csromanfolio{285}ปจฺจยจตุกฺกเหตฺวาทิ}

\proseitem{393}{เหตุ มหคฺคตา\-รมฺมโณ ธมฺโม เหตุสฺส มหคฺคตา\-รมฺมณสฺส ธมฺมสฺส เหตุปจฺจเยน ปจฺจโย. ตีณิ.}

\prose{เหตุ มหคฺคตา\-รมฺมโณ ธมฺโม เหตุสฺส มหคฺคตา\-รมฺมณสฺส ธมฺมสฺส อารมฺมณ\-ปจฺจเยน ปจฺจโย. นว.}

\prose{เหตุ มหคฺคตา\-รมฺมโณ ธมฺโม เหตุสฺส มหคฺคตา\-รมฺมณสฺส ธมฺมสฺส อธิปติ\-ปจฺจเยน ปจฺจโย. อารมฺมณา\-ธิปติ, สหชาตา\-ธิปติ. ตีณิ.}

\prose{มหคฺคตา\-รมฺมโณ ธมฺโม นเหตุสฺส มหคฺคตา\-รมฺมณสฺส ธมฺมสฺส อธิปติ\-ปจฺจเยน ปจฺจโย. อารมฺมณา\-ธิปติ, สหชาตา\-ธิปติ. ตีณิ.}

\prose{เหตุ มหคฺคตา\-รมฺมโณ จ นเหตุ มหคฺคตา\-รมฺมโณ จ ธมฺมา เหตุสฺส มหคฺคตา\-รมฺมณสฺส ธมฺมสฺส อธิปติ\-ปจฺจเยน ปจฺจโย. อารมฺมณา\-ธิปติ. ตีณิ. (สํขิตฺตํ.)}

\proseitem{394}{เหตุยา ตีณิ, อารมฺมเณ นว, อธิปติยา นว ฯเปฯ นิสฺสเย นว, อุปนิสฺสเย นว, อาเสวเน นว, กมฺเม ตีณิ, วิปาเก นว, อาหาเร ตีณิ ฯเปฯ อวิคเต นว. (สํขิตฺตํ.)}

\proseitem{395}{เหตุ มหคฺคตา\-รมฺมโณ ธมฺโม เหตุสฺส มหคฺคตา\-รมฺมณสฺส ธมฺมสฺส อารมฺมณ\-ปจฺจเยน ปจฺจโย. สหชาต\-ปจฺจเยน ปจฺจโย. อุปนิสฺสย\-ปจฺจเยน ปจฺจโย. (สํขิตฺตํ.)}

\proseitem{396}{นเหตุยา นว, นอารมฺมเณ นว. (สํขิตฺตํ.)}

\prose{เหตุปจฺจยา นอารมฺมเณ ตีณิ. (สํขิตฺตํ.)}

\prose{นเหตุปจฺจยา อารมฺมเณ นว. (สํขิตฺตํ.)}

\prose{(ยถา กุสลตฺติเก ปญฺหาวารสฺส อนุโลมมฺปิ ปจฺจนียมฺปิ อนุโลมปจฺจนียมฺปิ ปจฺจนียานุโลมมฺปิ คณิตํ, เอวํ คเณตพฺพํ.)}

\csromanmarkh{ธมฺมานุโลม ทุกติก\-ปฏฺฐาน\-ปาฬิ}\csromanmarkhii{3. อปฺปามาณารมฺมณปท 1. ปฏิจฺจวาร}\csromanheadernomark{2}{ธมฺมานุโลม ทุกติก\-ปฏฺฐาน\-ปาฬิ 3. อปฺปามาณารมฺมณปท 1. ปฏิจฺจวาร}
\prose{\csromanfolio{286}ปจฺจยจตุกฺกเหตุ}

\proseitem{397}{เหตุํ อปฺปมาณา\-รมฺมณํ ธมฺมํ ปฏิจฺจ เหตุ อปฺปมาณารมฺมโณ ธมฺโม อุปฺปชฺชติ เหตุปจฺจยา. ตีณิ.}

\prose{นเหตุํ อปฺปมาณา\-รมฺมณํ ธมฺมํ ปฏิจฺจ นเหตุ อปฺปมาณารมฺมโณ ธมฺโม อุปฺปชฺชติ เหตุปจฺจยา. ตีณิ.}

\prose{เหตุํ อปฺปมาณารมฺมณญฺจ นเหตุํ อปฺปมาณารมฺมณญฺจ ธมฺมํ ปฏิจฺจ เหตุ อปฺปมาณารมฺมโณ ธมฺโม อุปฺปชฺชติ เหตุปจฺจยา. ตีณิ. (สํขิตฺตํ.)}

\proseitem{398}{เหตุยา นว, อารมฺมเณ นว, อธิปติยา นว, อนนฺตเร นว, สมนนฺตเร นว, สหชาเต นว, อญฺญมญฺเญ นว, นิสฺสเย นว, อุปนิสฺสเย นว, ปุเรชาเต นว, อาเสวเน นว, กมฺเม นว, วิปาเก นว ฯเปฯ อวิคเต นว. (สํขิตฺตํ.)}

\vspace{\tipitakaparskip}
\csromancenter{\textbf{นเหตุ}}
\proseitem{399}{นเหตุํ อปฺปมาณา\-รมฺมณํ ธมฺมํ ปฏิจฺจ นเหตุ อปฺปมาณารมฺมโณ ธมฺโม อุปฺปชฺชติ นเหตุปจฺจยา. (สํขิตฺตํ.)}

\vspace{\tipitakaparskip}
\csromancenter{\textbf{นเหตุ}ยา เอกํ, นอธิปติยา นว, นปุเรชาเต นว, นปจฺฉาชาเต นว, นอาเสวเน นว, นกมฺเม ตีณิ, นวิปาเก นว, นมคฺเค เอกํ, นวิปฺปยุตฺเต นว. (สํขิตฺตํ.)}
\prose{เหตุปจฺจยา นอธิปติยา นว. (สํขิตฺตํ.)}

\prose{\textbf{นเหตุ}ปจฺจยา อารมฺมเณ เอกํ. (สํขิตฺตํ.)}

\prose{(สหชาตวาโรปิ ปจฺจยวาโรปิ นิสฺสยวาโรปิ สํสฏฺฐ\-วาโรปิ สมฺปยุตฺตวาโรปิ ปฏิจฺจ\-วาร\-สทิสา วิตฺถาเรตพฺพา.)}

\vspace{\tipitakaparskip}
\csromancenter{เหตุทุก 13. ปริตฺตา\-รมฺมณ\-ตฺติก 3. อปฺปมาณารมฺมณปท 7. ปญฺหาวาร}
\prose{\csromanfolio{287}ปจฺจยจตุกฺกเหตฺวาทิ}

\proseitem{401}{เหตุ อปฺปมาณารมฺมโณ ธมฺโม เหตุสฺส อปฺปมาณา\-รมฺมณสฺส ธมฺมสฺส เหตุปจฺจเยน ปจฺจโย. ตีณิ.}

\prose{เหตุ อปฺปมาณารมฺมโณ ธมฺโม เหตุสฺส อปฺปมาณา\-รมฺมณสฺส ธมฺมสฺส อารมฺมณ\-ปจฺจเยน ปจฺจโย. นว.}

\prose{เหตุ อปฺปมาณารมฺมโณ ธมฺโม เหตุสฺส อปฺปมาณา\-รมฺมณสฺส ธมฺมสฺส อธิปติ\-ปจฺจเยน ปจฺจโย. อารมฺมณา\-ธิปติ, สหชาตา\-ธิปติ. ตีณิ.}

\prose{อปฺปมาณารมฺมโณ ธมฺโม นเหตุสฺส อปฺปมาณา\-รมฺมณสฺส ธมฺมสฺส อธิปติ\-ปจฺจเยน ปจฺจโย. อารมฺมณา\-ธิปติ, สหชาตา\-ธิปติ. ตีณิ.}

\prose{เหตุ อปฺปมาณารมฺมโณ จ นเหตุ อปฺปมาณารมฺมโณ จ ธมฺมา เหตุสฺส อปฺปมาณา\-รมฺมณสฺส ธมฺมสฺส อธิปติ\-ปจฺจเยน ปจฺจโย. อารมฺมณา\-ธิปติ. ตีณิ. (สํขิตฺตํ.)}

\proseitem{402}{เหตุยา ตีณิ, อารมฺมเณ นว, อธิปติยา นว, อนนฺตเร นว, สมนนฺตเร นว, สหชาเต นว, อญฺญมญฺเญ นว, นิสฺสเย นว, อุปนิสฺสเย นว, อาเสวเน นว, กมฺเม ตีณิ, วิปาเก นว, อาหาเร ตีณิ, อินฺทฺริเร นว, ฌาเน ตีณิ, มคฺเค นว, สมฺปยุตฺเต นว, อตฺถิยา นว ฯเปฯ อวิคเต นว. (สํขิตฺตํ.)}

\proseitem{403}{เหตุ อปฺปมาณารมฺมโณ ธมฺโม เหตุสฺส อปฺปมาณา\-รมฺมณสฺส ธมฺมสฺส อารมฺมณ\-ปจฺจเยน ปจฺจโย. สหชาต\-ปจฺจเยน ปจฺจโย. อุปนิสฺสย\-ปจฺจเยน ปจฺจโย. (สํขิตฺตํ.)}

\proseitem{404}{นเหตุยา นว, นอารมฺมเณ นว. (สํขิตฺตํ.)}

\prose{เหตุปจฺจยา นอารมฺมเณ ตีณิ. (สํขิตฺตํ.)}

\prose{\csromanfolio{288}นเหตุปจฺจยา อารมฺมเณ นว. (สํขิตฺตํ.)}

\prosewithcloser{(ยถา กุสลตฺติเก ปญฺหาวารสฺส อนุโลมมฺปิ ปจฺจนียมฺปิ อนุโลมปจฺจนียมฺปิ ปจฺจนียานุโลมมฺปิ คณิตํ, เอวํ คเณตพฺพํ.)}{เหตุทุกปริตฺตารมฺมณตฺติกํ นิฏฺฐิตํ.}
\csromanmatikaanchor{mtk.45}
\csromanmatikaanchor{mtk.118}
\csromantocmarkref{section}{14. หินตฺติก}{mtk.45}
\bookmark[dest={mtk.45},level=1]{14. หินตฺติก}
\csromanmarkh{1. เหตุทุก}\csromanmarkhii{14. หีนตฺติก}\csromanheadernomark{6}{1. เหตุทุก 14. หีนตฺติก}
\csromanmarkh{1. หีนปท}\csromanmarkhii{1. ปฏิจฺจวาร}\csromanheadernomark{2}{1. \textbf{หีนปท 1. ปฏิจฺจวาร}}
\prose{ปจฺจยจตุกฺกเหตุ}

\proseitem{405}{เหตุํ หีนํ ธมฺมํ ปฏิจฺจ เหตุ หีโน ธมฺโม อุปฺปชฺชติ เหตุปจฺจยา. ตีณิ.}

\prose{นเหตุํ หีนํ ธมฺมํ ปฏิจฺจ นเหตุ หีโน ธมฺโม อุปฺปชฺชติ เหตุปจฺจยา. ตีณิ.}

\prose{เหตุํ หีนญฺจ นเหตุํ หีนญฺจ ธมฺมํ ปฏิจฺจ เหตุ หีโน ธมฺโม อุปฺปชฺชติ เหตุปจฺจยา. ตีณิ. (สํขิตฺตํ.)}

\proseitem{406}{เหตุยา นว, อารมฺมเณ นว, อธิปติยา นว, อนนฺตเร นว, สมนนฺตเร นว, สหชาเต นว, อญฺญมญฺเญ นว, นิสฺสเย นว, อุปนิสฺสเย นว, ปุเรชาเต นว, อาเสวเน นว, กมฺเม นว, อาหาเร นว ฯเปฯ อวิคเต นว. (สํขิตฺตํ.)}

\prose{นเหตุ นอธิปติ}

\proseitem{407}{นเหตุํ หีนํ ธมฺมํ ปฏิจฺจ เหตุ หีโน ธมฺโม อุปฺปชฺชติ นเหตุปจฺจยา. (1)}

\prose{เหตุํ หีนํ ธมฺมํ ปฏิจฺจ เหตุ หีโน ธมฺโม อุปฺปชฺชติ นอธิปติ\-ปจฺจยา. ตีณิ.}

\prose{นเหตุํ หีนํ ธมฺมํ ปฏิจฺจ นเหตุ หีโน ธมฺโม อุปฺปชฺชติ นอธิปติ\-ปจฺจยา. ตีณิ.}

\csromanmarkh{1. เหตุทุก}\csromanmarkhii{14. หีนตฺติก}\csromanheadernomark{2}{1. เหตุทุก 14. หีนตฺติก}
\prose{\csromanfolio{289}เหตุํ หีนญฺจ นเหตุํ หีนญฺจ ธมฺมํ ปฏิจฺจ เหตุ หีโน ธมฺโม อุปฺปชฺชติ นอธิปติ\-ปจฺจยา. ตีณิ. (สํขิตฺตํ.)}

\proseitem{408}{นเหตุยา เอกํ, นอธิปติยา นว, นปุเรชาเต นว, นปจฺฉาชาเต นว, นอาเสวเน นว, นกมฺเม ตีณิ, นวิปาเก นว, นวิปฺปยุตฺเต นว. (สํขิตฺตํ.)}

\prose{เหตุปจฺจยา นอธิปติยา นว. (สํขิตฺตํ.)}

\prose{นเหตุปจฺจยา อารมฺมเณ เอกํ. (สํขิตฺตํ.)}

\prose{(สหชาตวาโรปิ ปจฺจยวาโรปิ นิสฺสยวาโรปิ สํสฏฺฐ\-วาโรปิ สมฺปยุตฺตวาโรปิ ปฏิจฺจ\-วาร\-สทิสา วิตฺถาเรตพฺพา.)}

\csromanmarkh{1. หีนปท}\csromanmarkhii{7. ปญฺหาวาร}\csromanheadernomark{2}{1. \textbf{หีนปท 7. ปญฺหาวาร}}
\prose{ปจฺจยจตุกฺกเหตฺวาทิ}

\proseitem{409}{เหตุ หีโน ธมฺโม เหตุสฺส หีนสฺส ธมฺมสฺส เหตุปจฺจเยน ปจฺจโย. ตีณิ.}

\prose{เหตุ หีโน ธมฺโม เหตุสฺส หีนสฺส ธมฺมสฺส อารมฺมณ\-ปจฺจเยน ปจฺจโย ฯเปฯ อธิปติ\-ปจฺจเยน ปจฺจโย. อารมฺมณา\-ธิปติ, สหชาตา\-ธิปติ. (สํขิตฺตํ.)}

\proseitem{410}{เหตุยา ตีณิ, อารมฺมเณ นว, อธิปติยา นว, อนนฺตเร นว, สมนนฺตเร นว, สหชาเต นว, อญฺญมญฺเญ นว, นิสฺสเย นว, อุปนิสฺสเย นว, อาเสวเน นว, กมฺเม ตีณิ, อาหาเร ตีณิ, อินฺทฺริเย ฌาเน มคฺเค ตีณิ, สมฺปยุตฺเต นว ฯเปฯ อวิคเต นว. (สํขิตฺตํ.)}

\proseitem{411}{เหตุ หีโน ธมฺโม เหตุสฺส หีนสฺส ธมฺมสฺส อารมฺมณ\-ปจฺจเยน ปจฺจโย. สหชาต\-ปจฺจเยน ปจฺจโย. อุปนิสฺสย\-ปจฺจเยน ปจฺจโย. (สํขิตฺตํ.)}

\vspace{\tipitakaparskip}
\csromancenter{\textbf{นเหตุ}ยา นว, นอารมฺมเณ นว. (สํขิตฺตํ.)}
\prose{\csromanfolio{290}เหตุปจฺจยา นอารมฺมเณ ตีณิ. (สํขิตฺตํ.)}

\prose{นเหตุปจฺจยา อารมฺมเณ นว. (สํขิตฺตํ.)}

\prose{(ยถา กุสลตฺติเก ปญฺหาวารสฺส อนุโลมมฺปิ ปจฺจนียมฺปิ อนุโลมปจฺจนียมฺปิ ปจฺจนียานุโลมมฺปิ คณิตํ, เอวํ คเณตพฺพํ.)}

\csromanmarkh{2. มชฺฌิมปท}\csromanmarkhii{1. ปฏิจฺจวาร}\csromanheadernomark{2}{2. \textbf{มชฺฌิมปท 1. ปฏิจฺจวาร}}
\prose{ปจฺจยจตุกฺกเหตุ}

\proseitem{413}{เหตุํ มชฺฌิมํ ธมฺมํ ปฏิจฺจ เหตุ มชฺฌิโม ธมฺโม อุปฺปชฺชติ เหตุปจฺจยา. (สํขิตฺตํ.)}

\proseitem{414}{เหตุยา นว, อารมฺมเณ นว ฯเปฯ อวิคเต นว. (สํขิตฺตํ.)}

\vspace{\tipitakaparskip}
\csromancenter{\textbf{นเหตุ}}
\proseitem{415}{\textbf{นเหตุ}ํ มชฺฌิมํ ธมฺมํ ปฏิจฺจ นเหตุ มชฺฌิโม ธมฺโม อุปฺปชฺชติ นเหตุปจฺจยา. (สํขิตฺตํ.)}

\vspace{\tipitakaparskip}
\csromancenter{\textbf{นเหตุ}ยา เอกํ, นอารมฺมเณ ตีณิ, นอธิปติยา นว ฯเปฯ โนวิคเต ตีณิ. (สํขิตฺตํ.)}
\prose{เหตุปจฺจยา นอารมฺมเณ ตีณิ. (สํขิตฺตํ.)}

\prose{\textbf{นเหตุ}ปจฺจยา อารมฺมเณ เอกํ. (สํขิตฺตํ.)}

\prose{(สหชาตวาโรปิ ปจฺจยวาโรปิ นิสฺสยวาโรปิ สํสฏฺฐ\-วาโรปิ สมฺปยุตฺตวาโรปิ ปฏิจฺจ\-วาร\-สทิสา วิตฺถาเรตพฺพา.)}

\csromanmarkh{1. เหตุทุก}\csromanmarkhii{14. หีนตฺติก 2. มชฺฌิมปท 7. ปญฺหาวาร}\csromanheadernomark{2}{1. เหตุทุก 14. หีนตฺติก \textbf{2. มชฺฌิมปท 7. ปญฺหาวาร}}
\prose{\csromanfolio{291}ปจฺจยจตุกฺกเหตฺวาทิ}

\proseitem{417}{เหตุ มชฺฌิโม ธมฺโม เหตุสฺส มชฺฌิมสฺส ธมฺมสฺส เหตุปจฺจเยน ปจฺจโย. ตีณิ.}

\prose{เหตุ มชฺฌิโม ธมฺโม เหตุสฺส มชฺฌิมสฺส ธมฺมสฺส อารมฺมณ\-ปจฺจเยน ปจฺจโย. อธิปติ\-ปจฺจเยน ปจฺจโย. อารมฺมณา\-ธิปติ, สหชาตา\-ธิปติ. (สํขิตฺตํ.)}

\proseitem{418}{เหตุยา ตีณิ, อารมฺมเณ นว, อธิปติยา นว, อนนฺตเร นว, สมนนฺตเร นว ฯเปฯ กมฺเม ตีณิ, วิปาเก นว, อาหาเร ตีณิ ฯเปฯ อวิคเต นว. (สํขิตฺตํ.)}

\proseitem{419}{เหตุ มชฺฌิโม ธมฺโม เหตุสฺส มชฺฌิมสฺส ธมฺมสฺส อารมฺมณ\-ปจฺจเยน ปจฺจโย. สหชาต\-ปจฺจเยน ปจฺจโย. อุปนิสฺสย\-ปจฺจเยน ปจฺจโย. (สํขิตฺตํ.)}

\proseitem{420}{นเหตุยา นว, นอารมฺมเณ นว. (สํขิตฺตํ.)}

\prose{เหตุปจฺจยา นอารมฺมเณ ตีณิ. (สํขิตฺตํ.)}

\prose{นเหตุปจฺจยา อารมฺมเณ นว. (สํขิตฺตํ.)}

\prose{(ยถา กุสลตฺติเก ปญฺหาวารสฺส อนุโลมมฺปิ ปจฺจนียมฺปิ อนุโลมปจฺจนียมฺปิ ปจฺจนียานุโลมมฺปิ คณิตํ, เอวํ คเณตพฺพํ.)}

\csromanmarkh{3. ปณีตปท}\csromanmarkhii{1. ปฏิจฺจวาร}\csromanheadernomark{2}{3. \textbf{ปณีตปท 1. ปฏิจฺจวาร}}
\prose{ปจฺจยจตุกฺกเหตุ}

\proseitem{421}{เหตุํ ปณีตํ ธมฺมํ ปฏิจฺจ เหตุ ปณีโต ธมฺโม อุปฺปชฺชติ เหตุปจฺจยา. ตีณิ.}

\prose{\csromanfolio{292}นเหตุํ ปณีตํ ธมฺมํ ปฏิจฺจ นเหตุ ปณีโต ธมฺโม อุปฺปชฺชติ เหตุปจฺจยา. ตีณิ.}

\prose{เหตุํ ปณีตญฺจ นเหตุํ ปณีตญฺจ ธมฺมํ ปฏิจฺจ เหตุ ปณีโต ธมฺโม อุปฺปชฺชติ เหตุปจฺจยา. ตีณิ. (สํขิตฺตํ.)}

\proseitem{422}{เหตุยา นว, อารมฺมเณ นว, อธิปติยา นว ฯเปฯ กมฺเม นว, วิปาเก นว, อาหาเร นว ฯเปฯ อวิคเต นว. (สํขิตฺตํ.)}

\prose{นอธิปติ}

\proseitem{423}{เหตุํ ปณีตํ ธมฺมํ ปฏิจฺจ เหตุ ปณีโต ธมฺโม อุปฺปชฺชติ นอธิปติ\-ปจฺจยา. (สํขิตฺตํ.)}

\proseitem{424}{นอธิปติยา ฉ, นปุเรชาเต นว, นปจฺฉาชาเต นว, นอาเสวเน นว, นกมฺเม ตีณิ, นวิปาเก นว, นวิปฺปยุตฺเต นว. (สํขิตฺตํ.)}

\prose{เหตุปจฺจยา นอธิปติยา ฉ. (สํขิตฺตํ.)}

\prose{นอธิปติ\-ปจฺจยา เหตุยา ฉ. (สํขิตฺตํ.)}

\prose{(สหชาตวาโรปิ ปจฺจยวาโรปิ นิสฺสยวาโรปิ สํสฏฺฐ\-วาโรปิ สมฺปยุตฺตวาโรปิ ปฏิจฺจ\-วาร\-สทิสา วิตฺถาเรตพฺพา.)}

\csromanmarkh{3. ปณีตปท}\csromanmarkhii{7. ปญฺหาวาร}\csromanheadernomark{2}{3. \textbf{ปณีตปท 7. ปญฺหาวาร}}
\prose{ปจฺจยจตุกฺกเหตฺวาทิ}

\proseitem{425}{เหตุ ปณีโต ธมฺโม เหตุสฺส ปณีตสฺส ธมฺมสฺส เหตุปจฺจเยน ปจฺจโย. ตีณิ.}

\prose{ปณีโต ธมฺโม นเหตุสฺส ปณีตสฺส ธมฺมสฺส อารมฺมณ\-ปจฺจเยน ปจฺจโย. ตีณิ.}

\prose{เหตุ ปณีโต ธมฺโม เหตุสฺส ปณีตสฺส ธมฺมสฺส อธิปติ\-ปจฺจเยน ปจฺจโย. อารมฺมณา\-ธิปติ, สหชาตา\-ธิปติ. ตีณิ. (สํขิตฺตํ.)}

\csromanmatikaanchor{mtk.46}
\csromantocmarkref{section}{15. มิจฺฉตฺตนิยตตฺติก}{mtk.46}
\bookmark[dest={mtk.46},level=1]{15. มิจฺฉตฺตนิยตตฺติก}
\csromanmarkh{1. เหตุทุก}\csromanmarkhii{15. มิจฺฉตฺตนิยต\-ตฺติก}\csromanheadernomark{1}{1. เหตุทุก 15. มิจฺฉตฺตนิยต\-ตฺติก}
\proseitem{426}{\csromanfolio{293}เหตุยา ตีณิ, อารมฺมเณ ตีณิ, อธิปติยา ฉ, อนนฺตเร นว, สมนนฺตเร นว, สหชาเต นว, อญฺญมญฺเญ นว, นิสฺสเย นว, อุปนิสฺสเย นว, กมฺเม ตีณิ, วิปาเก นว, อาหาเร ตีณิ, อินฺทฺริเย นว, ฌาเน ตีณิ, มคฺเค นว, สมฺปยุตฺเต นว, อตฺถิยา นว ฯเปฯ อวิคเต นว. (สํขิตฺตํ.)}

\proseitem{427}{เหตุ ปณีโต ธมฺโม เหตุสฺส ปณีตสฺส ธมฺมสฺส สหชาต\-ปจฺจเยน ปจฺจโย. อุปนิสฺสย\-ปจฺจเยน ปจฺจโย. (สํขิตฺตํ.)}

\proseitem{428}{นเหตุยา นว, นอารมฺมเณ นว. (สํขิตฺตํ.)}

\prose{เหตุปจฺจยา นอารมฺมเณ ตีณิ. (สํขิตฺตํ.)}

\prose{นเหตุปจฺจยา อารมฺมเณ ตีณิ. (สํขิตฺตํ.)}

\prosewithcloser{(ยถา กุสลตฺติเก ปญฺหาวารสฺส อนุโลมมฺปิ ปจฺจนียมฺปิ อนุโลมปจฺจนียมฺปิ ปจฺจนียานุโลมมฺปิ คณิตํ, เอวํ คเณตพฺพํ.)}{เหตุ\-ทุก\-หีน\-ตฺติกํ นิฏฺฐิตํ.}
\csromanmarkh{1. เหตุทุก}\csromanmarkhii{15. มิจฺฉตฺตนิยต\-ตฺติก}\csromanheadernomark{2}{1. เหตุทุก 15. มิจฺฉตฺตนิยต\-ตฺติก}
\csromanmarkh{1. มิจฺฉตฺตนิยตปท}\csromanmarkhii{1. ปฏิจฺจวาร}\csromanheadernomark{2}{1. \textbf{มิจฺฉตฺตนิยตปท 1. ปฏิจฺจวาร}}
\prose{ปจฺจยจตุกฺกเหตุ อารมฺมณ}

\proseitem{429}{เหตุํ มิจฺฉตฺต\-นิยตํ ธมฺมํ ปฏิจฺจ เหตุ มิจฺฉตฺต\-นิยโต ธมฺโม อุปฺปชฺชติ เหตุปจฺจยา. ตีณิ.}

\prose{นเหตุํ มิจฺฉตฺต\-นิยตํ ธมฺมํ ปฏิจฺจ นเหตุ มิจฺฉตฺต\-นิยโต ธมฺโม อุปฺปชฺชติ เหตุปจฺจยา. ตีณิ.}

\prose{เหตุํ มิจฺฉตฺตนิยตญฺจ นเหตุํ มิจฺฉตฺตนิยตญฺจ ธมฺมํ ปฏิจฺจ เหตุ มิจฺฉตฺต\-นิยโต ธมฺโม อุปฺปชฺชติ เหตุปจฺจยา. ตีณิ.}

\prose{\csromanfolio{294}เหตุํ มิจฺฉตฺต\-นิยตํ ธมฺมํ ปฏิจฺจ เหตุ มิจฺฉตฺต\-นิยโต ธมฺโม อุปฺปชฺชติ อารมฺมณ\-ปจฺจยา. (สํขิตฺตํ.)}

\proseitem{430}{เหตุยา นว, อารมฺมเณ นว, อธิปติยา นว, อนนฺตเร นว, สมนนฺตเร นว, สหชาเต นว, อญฺญมญฺเญ นว, นิสฺสเย นว, อุปนิสฺสเย นว, ปุเรชาเต นว, อาเสวเน นว, กมฺเม นว, อาหาเร นว ฯเปฯ อวิคเต นว. (สํขิตฺตํ.)}

\prose{นอธิปติ}

\proseitem{431}{เหตุํ มิจฺฉตฺต\-นิยตํ ธมฺมํ ปฏิจฺจ นเหตุ มิจฺฉตฺต\-นิยโต ธมฺโม อุปฺปชฺชติ นอธิปติ\-ปจฺจยา. (สํขิตฺตํ.)}

\proseitem{432}{นอธิปติยา ตีณิ, นปจฺฉาชาเต นว, นกมฺเม ตีณิ, นวิปาเก นว. (สํขิตฺตํ.)}

\prose{เหตุปจฺจยา นอธิปติยา ตีณิ. (สํขิตฺตํ.)}

\prose{นอธิปติ\-ปจฺจยา เหตุยา ตีณิ. (สํขิตฺตํ.)}

\prose{(สหชาตวาโรปิ ปจฺจยวาโรปิ นิสฺสยวาโรปิ สํสฏฺฐ\-วาโรปิ สมฺปยุตฺตวาโรปิ ปฏิจฺจ\-วาร\-สทิสา วิตฺถาเรตพฺพา.)}

\csromanmarkh{1. มิจฺฉตฺตนิยตปท}\csromanmarkhii{7. ปญฺหาวาร}\csromanheadernomark{2}{1. \textbf{มิจฺฉตฺตนิยตปท 7. ปญฺหาวาร}}
\prose{ปจฺจยจตุกฺกเหตุ อธิปติ}

\proseitem{433}{เหตุ มิจฺฉตฺต\-นิยโต ธมฺโม เหตุสฺส มิจฺฉตฺต\-นิยตสฺส ธมฺมสฺส เหตุปจฺจเยน ปจฺจโย. ตีณิ.}

\prose{มิจฺฉตฺต\-นิยโต ธมฺโม นเหตุสฺส มิจฺฉตฺต\-นิยตสฺส ธมฺมสฺส อธิปติ\-ปจฺจเยน ปจฺจโย. สหชาตา\-ธิปติ. ตีณิ. (สํขิตฺตํ.)}

\proseitem{434}{เหตุยา ตีณิ, อธิปติยา ตีณิ, สหชาเต นว, อญฺญมญฺเญ นว, นิสฺสเย นว, อุปนิสฺสเย นว, กมฺเม ตีณิ, อาหาเร ตีณิ, อินฺทฺริเย ตีณิ, ฌาเน ตีณิ, มคฺเค ตีณิ, สมฺปยุตฺเต นว, อตฺถิยา นว, อวิคเต นว. (สํขิตฺตํ.)}

\csromanmarkh{1. เหตุทุก}\csromanmarkhii{15. มิจฺฉตฺตนิยต\-ตฺติก ปจฺจนียุ\-ทฺธาร}\csromanheadernomark{2}{1. เหตุทุก 15. มิจฺฉตฺตนิยต\-ตฺติก ปจฺจนียุ\-ทฺธาร}
\proseitem{435}{\csromanfolio{295}เหตุ มิจฺฉตฺต\-นิยโต ธมฺโม เหตุสฺส มิจฺฉตฺต\-นิยตสฺส ธมฺมสฺส สหชาต\-ปจฺจเยน ปจฺจโย. อุปนิสฺสย\-ปจฺจเยน ปจฺจโย. (สํขิตฺตํ.)}

\proseitem{436}{นเหตุยา นว, นอารมฺมเณ นว. (สํขิตฺตํ.)}

\prose{เหตุปจฺจยา นอธิปติยา ตีณิ. (สํขิตฺตํ.)}

\prose{นเหตุยา อธิปติยา ตีณิ. (สํขิตฺตํ.)}

\prose{(ยถา กุสลตฺติเก ปญฺหาวารสฺส อนุโลมมฺปิ ปจฺจนียมฺปิ อนุโลมปจฺจนียมฺปิ ปจฺจนียานุโลมมฺปิ คณิตํ, เอวํ คเณตพฺพํ.)}

\csromanmarkh{2. สมฺมตฺตนิยตปท}\csromanmarkhii{1. ปฏิจฺจวาร}\csromanheadernomark{2}{2. \textbf{สมฺมตฺตนิยตปท 1. ปฏิจฺจวาร}}
\prose{ปจฺจยจตุกฺกเหตุ}

\proseitem{437}{เหตุํ สมฺมตฺต\-นิยตํ ธมฺมํ ปฏิจฺจ เหตุ สมฺมตฺตนิยโต ธมฺโม อุปฺปชฺชติ เหตุปจฺจยา. ตีณิ.}

\prose{นเหตุํ สมฺมตฺต\-นิยตํ ธมฺมํ ปฏิจฺจ นเหตุ สมฺมตฺตนิยโต ธมฺโม อุปฺปชฺชติ เหตุปจฺจยา. ตีณิ.}

\prose{เหตุํ สมฺมตฺตนิยตญฺจ นเหตุํ สมฺมตฺตนิยตญฺจ ธมฺมํ ปฏิจฺจ เหตุ สมฺมตฺตนิยโต ธมฺโม อุปฺปชฺชติ เหตุปจฺจยา. ตีณิ. (สํขิตฺตํ.)}

\proseitem{438}{เหตุยา นว, อารมฺมเณ นว, อธิปติยา นว ฯเปฯ ปุเรชาเต นว, อาเสวเน นว, กมฺเม นว, อาหาเร นว ฯเปฯ อวิคเต นว. (สํขิตฺตํ.)}

\prose{นอธิปติ}

\proseitem{439}{เหตุํ สมฺมตฺต\-นิยตํ ธมฺมํ ปฏิจฺจ เหตุ สมฺมตฺตนิยโต ธมฺโม อุปฺปชฺชติ นอธิปติ\-ปจฺจยา. (สํขิตฺตํ.)}

\proseitem{440}{นอธิปติยา ฉ, นปุเรชาเต นว, นปจฺฉาชาเต นว, นกมฺเม ตีณิ, นวิปาเก นว, นวิปฺปยุตฺเต นว. (สํขิตฺตํ.)}

\prose{\csromanfolio{296}เหตุปจฺจยา นอธิปติยา ฉ. (สํขิตฺตํ.)}

\prose{นอธิปติ\-ปจฺจยา เหตุยา ฉ. (สํขิตฺตํ.)}

\prose{(สหชาตวาโรปิ ปจฺจยวาโรปิ นิสฺสยวาโรปิ สํสฏฺฐ\-วาโรปิ สมฺปยุตฺตวาโรปิ ปฏิจฺจ\-วาร\-สทิสา วิตฺถาเรตพฺพา.)}

\csromanmarkh{2. สมฺมตฺตนิยตปท}\csromanmarkhii{7. ปญฺหาวาร}\csromanheadernomark{2}{2. \textbf{สมฺมตฺตนิยตปท 7. ปญฺหาวาร}}
\prose{ปจฺจยจตุกฺกเหตุ อธิปติ}

\proseitem{441}{เหตุ สมฺมตฺตนิยโต ธมฺโม เหตุสฺส สมฺมตฺต\-นิยตสฺส ธมฺมสฺส เหตุปจฺจเยน ปจฺจโย. ตีณิ.}

\prose{เหตุ สมฺมตฺตนิยโต ธมฺโม เหตุสฺส สมฺมตฺต\-นิยตสฺส ธมฺมสฺส อธิปติ\-ปจฺจเยน ปจฺจโย. สหชาตา\-ธิปติ. ตีณิ.}

\prose{สมฺมตฺตนิยโต ธมฺโม นเหตุสฺส สมฺมตฺต\-นิยตสฺส ธมฺมสฺส อธิปติ\-ปจฺจเยน ปจฺจโย. สหชาตา\-ธิปติ. ตีณิ. (สํขิตฺตํ.)}

\proseitem{442}{เหตุยา ตีณิ, อธิปติยา ฉ, สหชาเต นว, อญฺญมญฺเญ นว, นิสฺสเย นว, อุปนิสฺสเย นว, กมฺเม ตีณิ, อาหาเร ตีณิ, อินฺทฺริเย นว, ฌาเน ตีณิ, มคฺเค นว, สมฺปยุตฺเต นว, อตฺถิยา นว, อวิคเต นว. (สํขิตฺตํ.)}

\proseitem{443}{เหตุ สมฺมตฺตนิยโต ธมฺโม เหตุสฺส สมฺมตฺต\-นิยตสฺส ธมฺมสฺส สหชาต\-ปจฺจเยน ปจฺจโย. อุปนิสฺสย\-ปจฺจเยน ปจฺจโย. (สํขิตฺตํ.)}

\proseitem{444}{นเหตุยา นว, นอารมฺมเณ นว. (สํขิตฺตํ.)}

\prose{เหตุปจฺจยา นอารมฺมเณ ตีณิ. (สํขิตฺตํ.)}

\prose{นเหตุปจฺจยา อธิปติยา ตีณิ. (สํขิตฺตํ.)}

\prose{(ยถา กุสลตฺติเก ปญฺหาวารสฺส อนุโลมมฺปิ ปจฺจนียมฺปิ อนุโลมปจฺจนียมฺปิ ปจฺจนียานุโลมมฺปิ คณิตํ, เอวํ คเณตพฺพํ.)}

\vspace{\tipitakaparskip}
\csromancenter{เหตุทุก 15. มิจฺฉตฺตนิยต\-ตฺติก 3. อนิยตปท 1. ปฏิจฺจวาร}
\prose{\csromanfolio{297}ปจฺจยจตุกฺกเหตุ}

\proseitem{445}{เหตุํ อนิยตํ ธมฺมํ ปฏิจฺจ เหตุ อนิยโต ธมฺโม อุปฺปชฺชติ เหตุปจฺจยา. ตีณิ.}

\prose{นเหตุํ อนิยตํ ธมฺมํ ปฏิจฺจ นเหตุ อนิยโต ธมฺโม อุปฺปชฺชติ เหตุปจฺจยา. ตีณิ.}

\prose{เหตุํ อนิยตญฺจ นเหตุํ อนิยตญฺจ ธมฺมํ ปฏิจฺจ เหตุ อนิยโต ธมฺโม อุปฺปชฺชติ เหตุปจฺจยา. ตีณิ. (สํขิตฺตํ.)}

\proseitem{446}{เหตุยา นว, อารมฺมเณ นว, อธิปติยา นว ฯเปฯ กมฺเม นว, วิปาเก นว ฯเปฯ อวิคเต นว. (สํขิตฺตํ.)}

\prose{นเหตุ นอารมฺมณ}

\proseitem{447}{นเหตุํ อนิยตํ ธมฺมํ ปฏิจฺจ นเหตุ อนิยโต ธมฺโม อุปฺปชฺชติ นเหตุปจฺจยา.  นเหตุํ อนิยตํ ธมฺมํ ปฏิจฺจ เหตุ อนิยโต ธมฺโม อุปฺปชฺชติ นเหตุปจฺจยา. (2)}

\prose{เหตุํ อนิยตํ ธมฺมํ ปฏิจฺจ นเหตุ อนิยโต ธมฺโม อุปฺปชฺชติ นอารมฺมณ\-ปจฺจยา. (1)}

\prose{นเหตุํ อนิยตํ ธมฺมํ ปฏิจฺจ นเหตุ อนิยโต ธมฺโม อุปฺปชฺชติ นอารมฺมณ\-ปจฺจยา. (1)}

\prose{เหตุํ อนิยตญฺจ นเหตุํ อนิยตญฺจ ธมฺมํ ปฏิจฺจ นเหตุ อนิยโต ธมฺโม อุปฺปชฺชติ นอารมฺมณ\-ปจฺจยา. (1) (สํขิตฺตํ.)}

\proseitem{448}{นเหตุยา เทฺว, นอารมฺมเณ ตีณิ, นอธิปติยา นว, นอนนฺตเร ตีณิ, นสมนนฺตเร ตีณิ, นอญฺญมญฺเญ ตีณิ, นอุปนิสฺสเย ตีณิ, นปุเรชาเต นว, นปจฺฉาชาเต นว, นอาเสวเน นว, นกมฺเม ตีณิ, นวิปาเก นว, นอาหาเร เอกํ, นอินฺทฺริเย เอกํ, นฌาเน เอกํ, นมคฺเค เอกํ, นสมฺปยุตฺเต ตีณิ, นวิปฺปยุตฺเต นว, โนนตฺถิยา ตีณิ, โนวิคเต ตีณิ. (สํขิตฺตํ.)}

\prose{\csromanfolio{298}เหตุปจฺจยา นอารมฺมเณ ตีณิ. (สํขิตฺตํ.)}

\prose{นเหตุปจฺจยา อารมฺมเณ เทฺว. (สํขิตฺตํ.)}

\prose{(สหชาตวาโรปิ ปจฺจยวาโรปิ นิสฺสยวาโรปิ สํสฏฺฐ\-วาโรปิ สมฺปยุตฺตวาโรปิ ปฏิจฺจ\-วาร\-สทิสา วิตฺถาเรตพฺพา.)}

\csromanmarkh{3. อนิยตปท}\csromanmarkhii{7. ปญฺหาวาร}\csromanheadernomark{2}{3. \textbf{อนิยตปท 7. ปญฺหาวาร}}
\prose{ปจฺจยจตุกฺกเหตฺวาทิ}

\proseitem{449}{เหตุ อนิยโต ธมฺโม เหตุสฺส อนิยตสฺส ธมฺมสฺส เหตุปจฺจเยน ปจฺจโย. ตีณิ.}

\prose{เหตุ อนิยโต ธมฺโม เหตุสฺส อนิยตสฺส ธมฺมสฺส อารมฺมณ\-ปจฺจเยน ปจฺจโย. นว.}

\prose{เหตุ อนิยโต ธมฺโม เหตุสฺส อนิยตสฺส ธมฺมสฺส อธิปติ\-ปจฺจเยน ปจฺจโย. อารมฺมณา\-ธิปติ, สหชาตา\-ธิปติ. ตีณิ.}

\prose{อนิยโต ธมฺโม นเหตุสฺส อนิยตสฺส ธมฺมสฺส อธิปติ\-ปจฺจเยน ปจฺจโย. อารมฺมณา\-ธิปติ, สหชาตา\-ธิปติ. ตีณิ.}

\prose{เหตุ อนิยโต จ นเหตุ อนิยโต จ ธมฺมา เหตุสฺส อนิยตสฺส ธมฺมสฺส อธิปติ\-ปจฺจเยน ปจฺจโย. อารมฺมณา\-ธิปติ, ตีณิ. (สํขิตฺตํ.)}

\proseitem{450}{เหตุยา ตีณิ, อารมฺมเณ นว, อธิปติยา นว, อนนฺตเร นว, สมนนฺตเร นว, สหชาเต นว, อญฺญมญฺเญ นว, นิสฺสเย นว, อุปนิสฺสเย นว, ปุเรชาเต ตีณิ, ปจฺฉาชาเต ตีณิ, อาเสวเน นว, กมฺเม ตีณิ, วิปาเก นว, อาหาเร ตีณิ, อินฺทฺริเย นว, ฌาเน ตีณิ, มคฺเค นว, สมฺปยุตฺเต นว, วิปฺปยุตฺเต ปญฺจ, อตฺถิยา นว, นตฺถิยา นว, วิคเต นว, อวิคเต นว.}

\csromanmarkh{1. เหตุทุก}\csromanmarkhii{16. มคฺคา\-รมฺมณ\-ตฺติก ปจฺจนียุ\-ทฺธาร}\csromanheadernomark{2}{1. เหตุทุก 16. มคฺคา\-รมฺมณ\-ตฺติก ปจฺจนียุ\-ทฺธาร}
\proseitem{451}{\csromanfolio{299}เหตุ อนิยโต ธมฺโม เหตุสฺส อนิยตสฺส ธมฺมสฺส อารมฺมณ\-ปจฺจเยน ปจฺจโย. สหชาต\-ปจฺจเยน ปจฺจโย. อุปนิสฺสย\-ปจฺจเยน ปจฺจโย. (สํขิตฺตํ.)}

\proseitem{452}{นเหตุยา นว, นอารมฺมเณ นว. (สํขิตฺตํ.)}

\prose{เหตุปจฺจยา นอารมฺมเณ ตีณิ. (สํขิตฺตํ.)}

\prose{นเหตุปจฺจยา อารมฺมเณ นว. (สํขิตฺตํ.)}

\prosewithcloser{(ยถา กุสลตฺติเก ปญฺหาวารสฺส อนุโลมมฺปิ ปจฺจนียมฺปิ อนุโลมปจฺจนียมฺปิ ปจฺจนียานุโลมมฺปิ คณิตํ, เอวํ คเณตพฺพํ.)}{เหตุ\-ทุก\-มิจฺฉตฺตนิยต\-ตฺติกํ นิฏฺฐิตํ.}
\csromanmatikaanchor{mtk.47}
\csromantocmarkref{section}{16. มคฺคารมฺมณตฺติก}{mtk.47}
\bookmark[dest={mtk.47},level=1]{16. มคฺคารมฺมณตฺติก}
\csromanmarkh{1. เหตุทุก}\csromanmarkhii{16. มคฺคา\-รมฺมณ\-ตฺติก}\csromanheadernomark{1}{1. เหตุทุก 16. มคฺคา\-รมฺมณ\-ตฺติก}
\csromanmarkh{1. มคฺคา\-รมฺมณปท}\csromanmarkhii{1. ปฏิจฺจวาร}\csromanheadernomark{2}{1. มคฺคา\-รมฺมณปท 1. ปฏิจฺจวาร}
\prose{ปจฺจยจตุกฺกเหตุ}

\proseitem{453}{เหตุํ มคฺคารมฺมณํ ธมฺมํ ปฏิจฺจ เหตุ มคฺคารมฺมโณ ธมฺโม อุปฺปชฺชติ เหตุปจฺจยา. ตีณิ.}

\prose{นเหตุํ มคฺคารมฺมณํ ธมฺมํ ปฏิจฺจ นเหตุ มคฺคารมฺมโณ ธมฺโม อุปฺปชฺชติ เหตุปจฺจยา. ตีณิ.}

\prose{เหตุํ มคฺคา\-รมฺมณญฺจ นเหตุํ มคฺคา\-รมฺมณญฺจ ธมฺมํ ปฏิจฺจ เหตุ มคฺคารมฺมโณ ธมฺโม อุปฺปชฺชติ เหตุปจฺจยา. ตีณิ. (สํขิตฺตํ.)}

\proseitem{454}{เหตุยา นว, อารมฺมเณ นว, อธิปติยา นว ฯเปฯ กมฺเม นว ฯเปฯ อวิคเต นว. (สํขิตฺตํ.)}

\proseitem{455}{\csromanfolio{300}\textbf{นเหตุ}ํ มคฺคารมฺมณํ ธมฺมํ ปฏิจฺจ นเหตุ มคฺคารมฺมโณ ธมฺโม อุปฺปชฺชติ นเหตุปจฺจยา. (สํขิตฺตํ.)}

\vspace{\tipitakaparskip}
\csromancenter{\textbf{นเหตุ}ยา เอกํ, นอธิปติยา นว, นปุเรชาเต นว, นปจฺฉาชาเต นว, นอาเสวเน นว, นกมฺเม ตีณิ, นวิปาเก นว, นมคฺเค เอกํ, นวิปฺปยุตฺเต นว. (สํขิตฺตํ.)}
\prose{เหตุปจฺจยา นอธิปติยา นว. (สํขิตฺตํ.)}

\prose{\textbf{นเหตุ}ปจฺจยา อารมฺมเณ เอกํ. (สํขิตฺตํ.)}

\prose{(สหชาตวาโรปิ ปจฺจยวาโรปิ นิสฺสยวาโรปิ สํสฏฺฐ\-วาโรปิ สมฺปยุตฺตวาโรปิ ปฏิจฺจ\-วาร\-สทิสา วิตฺถาเรตพฺพา.)}

\csromanmarkh{1. มคฺคา\-รมฺมณปท}\csromanmarkhii{7. ปญฺหาวาร}\csromanheadernomark{2}{1. มคฺคา\-รมฺมณปท 7. ปญฺหาวาร}
\prose{ปจฺจยจตุกฺกเหตฺวาทิ}

\proseitem{457}{เหตุ มคฺคารมฺมโณ ธมฺโม เหตุสฺส มคฺคา\-รมฺมณสฺส ธมฺมสฺส เหตุปจฺจเยน ปจฺจโย. ตีณิ.}

\prose{เหตุ มคฺคารมฺมโณ ธมฺโม เหตุสฺส มคฺคา\-รมฺมณสฺส ธมฺมสฺส อธิปติ\-ปจฺจเยน ปจฺจโย. สหชาตา\-ธิปติ. ตีณิ.}

\prose{มคฺคารมฺมโณ ธมฺโม นเหตุสฺส มคฺคา\-รมฺมณสฺส ธมฺมสฺส อธิปติ\-ปจฺจเยน ปจฺจโย. สหชาตา\-ธิปติ. ตีณิ.}

\prose{เหตุ มคฺคารมฺมโณ ธมฺโม เหตุสฺส มคฺคา\-รมฺมณสฺส ธมฺมสฺส อนนฺตร\-ปจฺจเยน ปจฺจโย. (สํขิตฺตํ.)}

\proseitem{458}{เหตุยา ตีณิ, อธิปติยา ฉ, อนนฺตเร นว, สมนนฺตเร นว, สหชาเต นว, อญฺญมญฺเญ นว, นิสฺสเย นว, อุปนิสฺสเย นว, อาเสวเน นว, กมฺเม ตีณิ, อาหาเร ตีณิ, อินฺทฺริเย นว, ฌาเน ตีณิ, มคฺเค นว, สมฺปยุตฺเต นว, อตฺถิยา นว, นตฺถิยา นว, วิคเต นว, อวิคเต นว. (สํขิตฺตํ.)}

\csromanmarkh{1. เหตุทุก}\csromanmarkhii{16. มคฺคา\-รมฺมณ\-ตฺติก ปจฺจนียุ\-ทฺธาร}\csromanheadernomark{2}{1. เหตุทุก 16. มคฺคา\-รมฺมณ\-ตฺติก ปจฺจนียุ\-ทฺธาร}
\proseitem{459}{\csromanfolio{301}เหตุ มคฺคารมฺมโณ ธมฺโม เหตุสฺส มคฺคา\-รมฺมณสฺส ธมฺมสฺส สหชาต\-ปจฺจเยน ปจฺจโย. อุปนิสฺสย\-ปจฺจเยน ปจฺจโย. (สํขิตฺตํ.)}

\proseitem{460}{นเยตุยา นว, นอารมฺมเณ นว. (สํขิตฺตํ.)}

\prose{เหตุปจฺจยา นอารมฺมเณ ตีณิ. (สํขิตฺตํ.)}

\prose{นเหตุปจฺจยา อธิปติยา ตีณิ. (สํขิตฺตํ.)}

\prose{(ยถา กุสลตฺติเก ปญฺหาวารสฺส อนุโลมมฺปิ ปจฺจนียมฺปิ อนุโลมปจฺจนียมฺปิ ปจฺจนียานุโลมมฺปิ คณิตํ, เอวํ คเณตพฺพํ.)}

\csromanmarkh{2. มคฺค\-เหตุกปท}\csromanmarkhii{1. ปฏิจฺจวาร}\csromanheadernomark{2}{2. มคฺค\-เหตุกปท 1. ปฏิจฺจวาร}
\prose{ปจฺจยจตุกฺกเหตุ}

\proseitem{461}{เหตุํ มคฺคเหตุกํ ธมฺมํ ปฏิจฺจ เหตุ มคฺคเหตุโก ธมฺโม อุปฺปชฺชติ เหตุปจฺจยา. ตีณิ.}

\prose{นเหตุํ มคฺคเหตุกํ ธมฺมํ ปฏิจฺจ นเหตุ มคฺคเหตุโก ธมฺโม อุปฺปชฺชติ เหตุปจฺจยา. ตีณิ.}

\prose{เหตุํ มคฺคเหตุกญฺจ นเหตุํ มคฺคเหตุกญฺจ ธมฺมํ ปฏิจฺจ เหตุ มคฺคเหตุโก ธมฺโม อุปฺปชฺชติ เหตุปจฺจยา. ตีณิ. (สํขิตฺตํ.)}

\proseitem{462}{เหตุยา นว, อารมฺมเณ นว, อธิปติยา นว ฯเปฯ อาเสวเน นว, กมฺเม นว, อาหาเร นว, อินฺทฺริเย นว ฯเปฯ อวิคเต นว. (สํขิตฺตํ.)}

\prose{นอธิปติ}

\proseitem{463}{เหตุํ มคฺคเหตุกํ ธมฺมํ ปฏิจฺจ เหตุ มคฺคเหตุโก ธมฺโม อุปฺปชฺชติ นอธิปติ\-ปจฺจยา. (สํขิตฺตํ.)}

\proseitem{464}{นอธิปติยา ฉ, นปุเรชาเต นว, นปจฺฉาชาเต นว, นกมฺเม ตีณิ, นวิปาเก นว, นวิปฺปยุตฺเต นว. (สํขิตฺตํ.)}

\prose{\csromanfolio{302}เหตุปจฺจยา นอธิปติยา ฉ. (สํขิตฺตํ.)}

\prose{นอธิปติ\-ปจฺจยา เหตุยา ฉ. (สํขิตฺตํ.)}

\prose{(สหชาตวาโรปิ ปจฺจยวาโรปิ นิสฺสยวาโรปิ สํสฏฺฐ\-วาโรปิ สมฺปยุตฺตวาโรปิ ปฏิจฺจ\-วาร\-สทิสา วิตฺถาเรตพฺพา.)}

\csromanmarkh{2. มคฺค\-เหตุกปท}\csromanmarkhii{7. ปญฺหาวาร}\csromanheadernomark{2}{2. มคฺค\-เหตุกปท 7. ปญฺหาวาร}
\prose{ปจฺจยจตุกฺกเหตุ อธิปติ}

\proseitem{465}{เหตุ มคฺคเหตุโก ธมฺโม เหตุสฺส มคฺคเหตุกสฺส ธมฺมสฺส เหตุปจฺจเยน ปจฺจโย. ตีณิ.}

\prose{เหตุ มคฺคเหตุโก ธมฺโม เหตุสฺส มคฺคเหตุกสฺส ธมฺมสฺส อธิปติ\-ปจฺจเยน ปจฺจโย. สหชาตา\-ธิปติ. ตีณิ.}

\prose{มคฺคเหตุโก ธมฺโม นเหตุสฺส มคฺคเหตุกสฺส ธมฺมสฺส อธิปติ\-ปจฺจเยน ปจฺจโย. สหชาตา\-ธิปติ. ตีณิ. (สํขิตฺตํ.)}

\proseitem{466}{เหตุยา ตีณิ, อธิปติยา ฉ, สหชาเต นว, อญฺญมญฺเญ นว, นิสฺสเย นว, อุปนิสฺสเย นว, กมฺเม ตีณิ, อาหาเร ตีณิ, อินฺทฺริเย นว, ฌาเน ตีณิ, มคฺเค นว, สมฺปยุตฺเต นว, อตฺถิยา นว, อวิคเต นว. (สํขิตฺตํ.)}

\proseitem{467}{เหตุ มคฺคเหตุโก ธมฺโม เหตุสฺส มคฺคเหตุกสฺส ธมฺมสฺส สหชาต\-ปจฺจเยน ปจฺจโย. อุปนิสฺสย\-ปจฺจเยน ปจฺจโย. (สํขิตฺตํ.)}

\proseitem{468}{นเหตุยา นว, นอารมฺมเณ นว. (สํขิตฺตํ.)}

\prose{เหตุปจฺจยา นอารมฺมเณ ตีณิ. (สํขิตฺตํ.)}

\prose{นเหตุปจฺจยา อธิปติยา ตีณิ. (สํขิตฺตํ.)}

\prose{(ยถา กุสลตฺติเก ปญฺหาวารสฺส อนุโลมมฺปิ ปจฺจนียมฺปิ อนุโลมปจฺจนียมฺปิ ปจฺจนียานุโลมมฺปิ คณิตํ, เอวํ คเณตพฺพํ.)}

\vspace{\tipitakaparskip}
\csromancenter{เหตุทุก 16. มคฺคา\-รมฺมณ\-ตฺติก 3. มคฺคา\-ธิป\-ติปท 1. ปฏิจฺจวาร}
\prose{\csromanfolio{303}ปจฺจยจตุกฺกเหตุ}

\proseitem{469}{เหตุํ มคฺคาธิปติํ ธมฺมํ ปฏิจฺจ เหตุ มคฺคาธิปติ ธมฺโม อุปฺปชฺชติ เหตุปจฺจยา. ตีณิ.}

\prose{นเหตุํ มคฺคาธิปติํ ธมฺมํ ปฏิจฺจ นเหตุ มคฺคาธิปติ ธมฺโม อุปฺปชฺชติ เหตุปจฺจยา. ตีณิ.}

\prose{เหตุํ มคฺคา\-ธิปติญฺจ นเหตุํ มคฺคา\-ธิปติญฺจ ธมฺมํ ปฏิจฺจ เหตุ มคฺคาธิปติ ธมฺโม อุปฺปชฺชติ เหตุปจฺจยา. ตีณิ. (สํขิตฺตํ.)}

\proseitem{470}{เหตุยา นว, อารมฺมเณ นว ฯเปฯ กมฺเม นว, อาหาเร นว ฯเปฯ อวิคเต นว. (สํขิตฺตํ.)}

\prose{นอธิปติ}

\proseitem{471}{เหตุํ มคฺคาธิปติํ ธมฺมํ ปฏิจฺจ เหตุ มคฺคาธิปติ ธมฺโม อุปฺปชฺชติ นอธิปติ\-ปจฺจยา. (สํขิตฺตํ.)}

\proseitem{472}{นอธิปติยา นว, นปุเรชาเต นว, นปจฺฉาชาเต นว, นอาเสวเน นว, นกมฺเม ตีณิ, นวิปาเก นว, นวิปฺปยุตฺเต นว. (สํขิตฺตํ.)}

\prose{เหตุปจฺจยา นอธิปติยา นว. (สํขิตฺตํ.)}

\prose{นอธิปติ\-ปจฺจยา เหตุยา นว. (สํขิตฺตํ.)}

\prose{(สหชาตวาโรปิ ปจฺจยวาโรปิ นิสฺสยวาโรปิ สํสฏฺฐ\-วาโรปิ สมฺปยุตฺตวาโรปิ ปฏิจฺจ\-วาร\-สทิสา วิตฺถาเรตพฺพา.)}

\csromanmarkh{3. มคฺคา\-ธิป\-ติปท}\csromanmarkhii{7. ปญฺหาวาร}\csromanheadernomark{2}{3. มคฺคา\-ธิป\-ติปท 7. ปญฺหาวาร}
\prose{ปจฺจยจตุกฺกเหตฺวาทิ}

\proseitem{473}{เหตุ มคฺคาธิปติ ธมฺโม เหตุสฺส มคฺคา\-ธิปติสฺส ธมฺมสฺส เหตุปจฺจเยน ปจฺจโย. ตีณิ.}

\prose{\csromanfolio{304}เหตุ มคฺคาธิปติ ธมฺโม เหตุสฺส มคฺคา\-ธิปติสฺส ธมฺมสฺส อารมฺมณ\-ปจฺจเยน ปจฺจโย. นว.}

\prose{เหตุ มคฺคาธิปติ ธมฺโม เหตุสฺส มคฺคา\-ธิปติสฺส ธมฺมสฺส อธิปติ\-ปจฺจเยน ปจฺจโย. อารมฺมณา\-ธิปติ, สหชาตา\-ธิปติ. ตีณิ.}

\prose{มคฺคาธิปติ ธมฺโม นเหตุสฺส มคฺคา\-ธิปติสฺส ธมฺมสฺส อธิปติ\-ปจฺจเยน ปจฺจโย. อารมฺมณา\-ธิปติ, สหชาตา\-ธิปติ. ตีณิ.}

\prose{เหตุ มคฺคาธิปติ จ นเหตุ มคฺคาธิปติ จ ธมฺมา เหตุสฺส มคฺคา\-ธิปติสฺส ธมฺมสฺส อธิปติ\-ปจฺจเยน ปจฺจโย. อารมฺมณา\-ธิปติ. ตีณิ. (สํขิตฺตํ.)}

\proseitem{474}{เหตุยา ตีณิ, อารมฺมเณ นว, อธิปติยา นว, อนนฺตเร นว, สมนนฺตเร นว, สหชาเต นว, อญฺญมญฺเญ นว, นิสฺสเย นว, อุปนิสฺสเย นว, อาเสวเน นว, กมฺเม ตีณิ, อาหาเร ตีณิ, อินฺทฺริเย นว, ฌาเน ตีณิ, มคฺเค นว, สมฺปยุตฺเต นว, อตฺถิยา นว ฯเปฯ อวิคเต นว. (สํขิตฺตํ.)}

\proseitem{475}{เหตุ มคฺคาธิปติ ธมฺโม เหตุสฺส มคฺคา\-ธิปติสฺส ธมฺมสฺส อารมฺมณ\-ปจฺจเยน ปจฺจโย. สหชาต\-ปจฺจเยน ปจฺจโย. อุปนิสฺสย\-ปจฺจเยน ปจฺจโย. (สํขิตฺตํ.)}

\proseitem{476}{นเหตุยา นว, นอารมฺมเณ นว. (สํขิตฺตํ.)}

\prose{เหตุปจฺจยา นอารมฺมเณ ตีณิ. (สํขิตฺตํ.)}

\prose{นเหตุปจฺจยา อารมฺมเณ นว. (สํขิตฺตํ.)}

\prosewithcloser{(ยถา กุสลตฺติเก ปญฺหาวารสฺส อนุโลมมฺปิ ปจฺจนียมฺปิ อนุโลมปจฺจนียมฺปิ ปจฺจนียานุโลมมฺปิ คณิตํ, เอวํ คเณตพฺพํ.)}{เหตุทุกมคฺคารมฺมณตฺติกํ นิฏฺฐิตํ.}
\csromanmatikaanchor{mtk.48}
\csromantocmarkref{section}{17. อุปฺปนฺนตฺติก}{mtk.48}
\bookmark[dest={mtk.48},level=1]{17. อุปฺปนฺนตฺติก}
\csromanmarkh{1. เหตุทุก}\csromanmarkhii{17. อุปฺปนฺนตฺติก 1. เหตุทุก 17. อุปฺปนฺนตฺติก}\csromanheadernomark{1}{1. เหตุทุก 17. อุปฺปนฺนตฺติก \textbf{1. เหตุทุก 17. อุปฺปนฺนตฺติก}}
\csromanmarkh{1. อุปฺปนฺนปท}\csromanmarkhii{7. ปญฺหาวาร}\csromanheadernomark{2}{1. \textbf{อุปฺปนฺนปท 7. ปญฺหาวาร}}
\prose{\csromanfolio{305}ปจฺจยจตุกฺกเหตุ อุปนิสฺสย}

\proseitem{477}{เหตุ อุปฺปนฺโน ธมฺโม เหตุสฺส อุปฺปนฺนสฺส ธมฺมสฺส เหตุปจฺจเยน ปจฺจโย. ตีณิ.}

\prose{อุปฺปนฺโน ธมฺโม นเหตุสฺส อุปฺปนฺนสฺส ธมฺมสฺส อารมฺมณ\-ปจฺจเยน ปจฺจโย. ตีณิ.}

\prose{เหตุ อุปฺปนฺโน ธมฺโม เหตุสฺส อุปฺปนฺนสฺส ธมฺมสฺส อธิปติ\-ปจฺจเยน ปจฺจโย. สหชาตา\-ธิปติ. ตีณิ.}

\prose{อุปฺปนฺโน ธมฺโม นเหตุสฺส อุปฺปนฺนสฺส ธมฺมสฺส อธิปติ\-ปจฺจเยน ปจฺจโย. อารมฺมณา\-ธิปติ, สหชาตา\-ธิปติ. ตีณิ.}

\prose{เหตุ อุปฺปนฺโน ธมฺโม เหตุสฺส อุปฺปนฺนสฺส ธมฺมสฺส สหชาต\-ปจฺจเยน ปจฺจโย. (สํขิตฺตํ.)}

\prose{อุปฺปนฺโน ธมฺโม นเหตุสฺส อุปฺปนฺนสฺส ธมฺมสฺส อุปนิสฺสย\-ปจฺจเยน ปจฺจโย. อารมฺมณู\-ปนิสฺสโย, ปกตู\-ปนิสฺสโย ฯเปฯ. ปกตู\-ปนิสฺสโย\csromandash{}อุปฺปนฺนํ อุตุํ อุปนิสฺสาย ฌานํ อุปฺปาเทติ, วิปสฺสนํ. มคฺคํ. อภิญฺญํ. สมาปตฺติํ อุปฺปาเทติ. (สํขิตฺตํ.)}

\proseitem{478}{เหตุยา ตีณิ, อารมฺมเณ ตีณิ, อธิปติยา ฉ, สหชาเต นว, อญฺญมญฺเญ นว, นิสฺสเย นว, อุปนิสฺสเย ตีณิ, ปุเรชาเต ตีณิ, ปจฺฉาชาเต ตีณิ, กมฺเม ตีณิ, วิปาเก นว, อาหาเร ตีณิ, อินฺทฺริเย นว, ฌาเน ตีณิ, มคฺเค นว, สมฺปยุตฺเต นว, วิปฺปยุตฺเต ปญฺจ, อตฺถิยา นว, อวิคเต นว. (สํขิตฺตํ.)}

\proseitem{479}{เหตุ อุปฺปนฺโน ธมฺโม เหตุสฺส อุปฺปนฺนสฺส ธมฺมสฺส สหชาต\-ปจฺจเยน ปจฺจโย. (สํขิตฺตํ.)}

\proseitem{480}{\csromanfolio{306}นเหตุยา นว, นอารมฺมเณ นว. (สํขิตฺตํ.)}

\prose{เหตุปจฺจยา นอารมฺมเณ ตีณิ. (สํขิตฺตํ.)}

\prosewithcloser{นเหตุปจฺจยา อารมฺมเณ ตีณิ. (สํขิตฺตํ.) (ยถา กุสลตฺติเก ปญฺหาวารสฺส อนุโลมมฺปิ ปจฺจนียมฺปิ อนุโลมปจฺจนียมฺปิ ปจฺจนียานุโลมมฺปิ คณิตํ, เอวํ คเณตพฺพํ. อิมมฺหิ ทุกติเก ปฏิจฺจวารมฺปิ สหชาตวารมฺปิ ปจฺจยวารมฺปิ นิสฺสยวารมฺปิ สํสฏฺฐ\-วารมฺปิ สมฺปยุตฺตวารมฺปิ อนุปฺปนฺนมฺปิ อุปฺปาทีปิ นตฺถิ.)}{เหตุทุกอุปฺปนฺนตฺติกํ นิฏฺฐิตํ.}
\csromanmatikaanchor{mtk.49}
\csromantocmarkref{section}{18. อตีตตฺติก}{mtk.49}
\bookmark[dest={mtk.49},level=1]{18. อตีตตฺติก}
\csromanmarkh{1. เหตุทุก}\csromanmarkhii{18. อตีตตฺติก}\csromanheadernomark{1}{1. เหตุทุก 18. อตีตตฺติก}
\csromanmarkh{3. ปจฺจุปฺปนฺนปท}\csromanmarkhii{7. ปญฺหาวาร}\csromanheadernomark{2}{3. \textbf{ปจฺจุปฺปนฺนปท 7. ปญฺหาวาร}}
\prose{ปจฺจยจตุกฺกเหตฺวาทิ}

\proseitem{481}{เหตุ ปจฺจุปฺปนฺโน ธมฺโม เหตุสฺส ปจฺจุปฺปนฺนสฺส ธมฺมสฺส เหตุปจฺจเยน ปจฺจโย. ตีณิ.}

\prose{ปจฺจุปฺปนฺโน ธมฺโม นเหตุสฺส ปจฺจุปฺปนฺนสฺส ธมฺมสฺส อารมฺมณ\-ปจฺจเยน ปจฺจโย. ตีณิ.}

\prose{เหตุ ปจฺจุปฺปนฺโน ธมฺโม เหตุสฺส ปจฺจุปฺปนฺนสฺส ธมฺมสฺส อธิปติ\-ปจฺจเยน ปจฺจโย. สหชาตา\-ธิปติ. ตีณิ.}

\prose{ปจฺจุปฺปนฺโน ธมฺโม นเหตุสฺส ปจฺจุปฺปนฺนสฺส ธมฺมสฺส อธิปติ\-ปจฺจเยน ปจฺจโย. อารมฺมณา\-ธิปติ, สหชาตา\-ธิปติ. ตีณิ.}

\prose{เหตุ ปจฺจุปฺปนฺโน ธมฺโม เหตุสฺส ปจฺจุปฺปนฺนสฺส ธมฺมสฺส สหชาต\-ปจฺจเยน ปจฺจโย. (สํขิตฺตํ.)}

\proseitem{482}{เหตุยา ตีณิ, อารมฺมเณ ตีณิ, อธิปติยา ฉ, สหชาเต นว, อญฺญมญฺเญ นว, นิสฺสเย นว, อุปนิสฺสเย นว, ปุเรชาเต ตีณิ, ปจฺฉาชาเต ตีณิ, กมฺเม ตีณิ, วิปาเก นว, อาหาเร ตีณิ, อินฺทฺริเย}

\vspace{\tipitakaparskip}
\csromancenter{เหตุทุก 19. อตีตา\-รมฺมณ\-ตฺติก นว, ฌาเน ตีณิ, มคฺเค นว, สมฺปยุตฺเต นว, วิปฺปยุตฺเต ปญฺจ, อตฺถิยา นว, อวิคเต นว. (สํขิตฺตํ.)}
\proseitem{483}{\csromanfolio{307}เหตุ ปจฺจุปฺปนฺโน ธมฺโม เหตุสฺส ปจฺจุปฺปนฺนสฺส ธมฺมสฺส สหชาต\-ปจฺจเยน ปจฺจโย. (สํขิตฺตํ.)}

\proseitem{484}{นเหตุยา นว, นอารมฺมเณ นว. (สํขิตฺตํ.)}

\prose{เหตุปจฺจยา นอารมฺมเณ ตีณิ. (สํขิตฺตํ.)}

\prosewithcloser{นเหตุปจฺจยา อารมฺมเณ ตีณิ. (สํขิตฺตํ.) (ยถา กุสลตฺติเก ปญฺหาวารสฺส อนุโลมมฺปิ ปจฺจนียมฺปิ อนุโลมปจฺจนียมฺปิ ปจฺจนียานุโลมมฺปิ คณิตํ, เอวํ คเณตพฺพํ. อิมมฺหิ ทุกติเก ปฏิจฺจวารมฺปิ สหชาตวารมฺปิ ปจฺจยวารมฺปิ นิสฺสยวารมฺปิ สํสฏฺฐ\-วารมฺปิ สมฺปยุตฺตวารมฺปิ อตีตมฺปิ อนาคตมฺปิ นตฺถิ.)}{เหตุ\-ทุก\-อตีตตฺติกํ นิฏฺฐิตํ.}
\csromanmatikaanchor{mtk.50}
\csromantocmarkref{section}{19. อตีตารมฺมณตฺติก}{mtk.50}
\bookmark[dest={mtk.50},level=1]{19. อตีตารมฺมณตฺติก}
\csromanmarkh{1. เหตุทุก}\csromanmarkhii{19. อตีตา\-รมฺมณ\-ตฺติก}\csromanheadernomark{1}{1. เหตุทุก 19. อตีตา\-รมฺมณ\-ตฺติก}
\csromanmarkh{1. อตีตา\-รมฺมณปท}\csromanmarkhii{1. ปฏิจฺจวาร}\csromanheadernomark{2}{1. อตีตา\-รมฺมณปท 1. ปฏิจฺจวาร}
\prose{ปจฺจยจตุกฺกเหตุ}

\proseitem{485}{เหตุํ อตีตารมฺมณํ ธมฺมํ ปฏิจฺจ เหตุ อตีตารมฺมโณ ธมฺโม อุปฺปชฺชติ เหตุปจฺจยา. ตีณิ.}

\prose{นเหตุํ อตีตารมฺมณํ ธมฺมํ ปฏิจฺจ นเหตุ อตีตารมฺมโณ ธมฺโม อุปฺปชฺชติ เหตุปจฺจยา. ตีณิ.}

\prose{เหตุํ อตีตา\-รมฺมณญฺจ นเหตุํ อตีตา\-รมฺมณญฺจ ธมฺมํ ปฏิจฺจ เหตุ อตีตารมฺมโณ ธมฺโม อุปฺปชฺชติ เหตุปจฺจยา. ตีณิ. (สํขิตฺตํ.)}

\proseitem{486}{\csromanfolio{308}เหตุยา นว, อารมฺมเณ นว, อธิปติยา นว, อนนฺตเร นว, สมนนฺตเร นว ฯเปฯ กมฺเม นว ฯเปฯ อวิคเต นว. (สํขิตฺตํ.)}

\prose{นเหตุ นอธิปติ}

\proseitem{487}{นเหตุํ อตีตารมฺมณํ ธมฺมํ ปฏิจฺจ นเหตุ อตีตารมฺมโณ ธมฺโม อุปฺปชฺชติ นเหตุปจฺจยา.  นเหตุํ อตีตารมฺมณํ ธมฺมํ ปฏิจฺจ เหตุ อตีตารมฺมโณ ธมฺโม อุปฺปชฺชติ นเหตุปจฺจยา. (2)}

\prose{เหตุํ อตีตารมฺมณํ ธมฺมํ ปฏิจฺจ เหตุ อตีตารมฺมโณ ธมฺโม อุปฺปชฺชติ นอธิปติ\-ปจฺจยา. (สํขิตฺตํ.)}

\proseitem{488}{นเหตุยา เทฺว, นอธิปติยา นว, นปุเรชาเต นว, นปจฺฉาชาเต นว, นอาเสวเน นว, นกมฺเม ตีณิ, นวิปาเก นว, นมคฺเค เอกํ, นวิปฺปยุตฺเต นว. (สํขิตฺตํ.)}

\prose{เหตุปจฺจยา นอธิปติยา นว. (สํขิตฺตํ.)}

\prose{นเหตุปจฺจยา อารมฺมเณ เทฺว. (สํขิตฺตํ.)}

\prose{(สหชาตวาโรปิ ปจฺจยวาโรปิ นิสฺสยวาโรปิ สํสฏฺฐ\-วาโรปิ สมฺปยุตฺตวาโรปิ ปฏิจฺจ\-วาร\-สทิสา วิตฺถาเรตพฺพา.)}

\csromanmarkh{1. อตีตา\-รมฺมณปท}\csromanmarkhii{7. ปญฺหาวาร}\csromanheadernomark{2}{1. อตีตา\-รมฺมณปท 7. ปญฺหาวาร}
\prose{ปจฺจยจตุกฺกเหตฺวาทิ}

\proseitem{489}{เหตุ อตีตารมฺมโณ ธมฺโม เหตุสฺส อตีตา\-รมฺมณสฺส ธมฺมสฺส เหตุปจฺจเยน ปจฺจโย. ตีณิ.}

\prose{เหตุ อตีตารมฺมโณ ธมฺโม เหตุสฺส อตีตา\-รมฺมณสฺส ธมฺมสฺส อารมฺมณ\-ปจฺจเยน ปจฺจโย. นว.}

\prose{เหตุ อตีตารมฺมโณ ธมฺโม เหตุสฺส อตีตา\-รมฺมณสฺส ธมฺมสฺส อธิปติ\-ปจฺจเยน ปจฺจโย. อารมฺมณา\-ธิปติ, สหชาตา\-ธิปติ. ตีณิ.}

\csromanmarkh{1. เหตุทุก}\csromanmarkhii{19. อตีตา\-รมฺมณ\-ตฺติก}\csromanheadernomark{2}{1. เหตุทุก 19. อตีตา\-รมฺมณ\-ตฺติก}
\prose{\csromanfolio{309}อตีตารมฺมโณ ธมฺโม นเหตุสฺส อตีตา\-รมฺมณสฺส ธมฺมสฺส อธิปติ\-ปจฺจเยน ปจฺจโย. อารมฺมณา\-ธิปติ, สหชาตา\-ธิปติ. ตีณิ.}

\prose{เหตุ อตีตารมฺมโณ จ นเหตุ อตีตารมฺมโณ จ ธมฺมา เหตุสฺส อตีตา\-รมฺมณสฺส ธมฺมสฺส อธิปติ\-ปจฺจเยน ปจฺจโย. อารมฺมณา\-ธิปติ. ตีณิ. (สํขิตฺตํ.)}

\proseitem{490}{เหตุยา ตีณิ, อารมฺมเณ นว, อธิปติยา นว, อนนฺตเร นว, สมนนฺตเร นว, สหชาเต นว, อญฺญมญฺเญ นว, นิสฺสเย นว, อุปนิสฺสเย นว, อาเสวเน นว, กมฺเม ตีณิ, วิปาเก นว, อาหาเร ตีณิ, อินฺทฺริเย นว, ฌาเน ตีณิ, มคฺเค นว, สมฺปยุตฺเต นว, อตฺถิยา นว ฯเปฯ อวิคเต นว. (สํขิตฺตํ.)}

\proseitem{491}{เหตุ อตีตารมฺมโณ ธมฺโม เหตุสฺส อตีตา\-รมฺมณสฺส ธมฺมสฺส อารมฺมณ\-ปจฺจเยน ปจฺจโย. สหชาต\-ปจฺจเยน ปจฺจโย. อุปนิสฺสย\-ปจฺจเยน ปจฺจโย. (สํขิตฺตํ.)}

\proseitem{492}{นเหตุยา นว, นอารมฺมเณ นว. (สํขิตฺตํ.)}

\prose{เหตุปจฺจยา นอารมฺมเณ ตีณิ. (สํขิตฺตํ.)}

\prose{นเหตุปจฺจยา อารมฺมเณ นว. (สํขิตฺตํ.)}

\prose{(ยถา กุสลตฺติเก ปญฺหาวารสฺส อนุโลมมฺปิ ปจฺจนียมฺปิ อนุโลมปจฺจนียมฺปิ ปจฺจนียานุโลมมฺปิ คณิตํ, เอวํ คเณตพฺพํ.)}

\csromanmarkh{2. อนาคตา\-รมฺมณปท}\csromanmarkhii{1. ปฏิจฺจวาร}\csromanheadernomark{2}{2. อนาคตา\-รมฺมณปท 1. ปฏิจฺจวาร}
\prose{ปจฺจยจตุกฺกเหตุ}

\proseitem{493}{เหตุํ อนาคตา\-รมฺมณํ ธมฺมํ ปฏิจฺจ เหตุ อนาคตารมฺมโณ ธมฺโม อุปฺปชฺชติ เหตุปจฺจยา. ตีณิ.}

\prose{\csromanfolio{310}นเหตุํ อนาคตา\-รมฺมณํ ธมฺมํ ปฏิจฺจ นเหตุ อนาคตารมฺมโณ ธมฺโม อุปฺปชฺชติ เหตุปจฺจยา. ตีณิ.}

\prose{เหตุํ อนาคตา\-รมฺมณญฺจ นเหตุํ อนาคตา\-รมฺมณญฺจ ธมฺมํ ปฏิจฺจ เหตุ อนาคตารมฺมโณ ธมฺโม อุปฺปชฺชติ เหตุปจฺจยา. ตีณิ. (สํขิตฺตํ.)}

\proseitem{494}{เหตุยา นว, อารมฺมเณ นว ฯเปฯ กมฺเม นว ฯเปฯ อวิคเต นว. (สํขิตฺตํ.)}

\prose{นเหตุ นอธิปติ}

\proseitem{495}{นเหตุํ อนาคตา\-รมฺมณํ ธมฺมํ ปฏิจฺจ นเหตุ อนาคตารมฺมโณ ธมฺโม อุปฺปชฺชติ นเหตุปจฺจยา.  นเหตุํ อนาคตา\-รมฺมณํ ธมฺมํ ปฏิจฺจ เหตุ อนาคตารมฺมโณ ธมฺโม อุปฺปชฺชติ นเหตุปจฺจยา. (2)}

\prose{เหตุํ อนาคตา\-รมฺมณํ ธมฺมํ ปฏิจฺจ เหตุ อนาคตารมฺมโณ ธมฺโม อุปฺปชฺชติ นอธิปติ\-ปจฺจยา. (สํขิตฺตํ.)}

\proseitem{496}{นเหตุยา เทฺว, นอธิปติยา นว, นปุเรชาเต นว, นปจฺฉาชาเต นว, นอาเสวเน นว, นกมฺเม ตีณิ, นวิปาเก นว, นมคฺเค เอกํ, นวิปฺปยุตฺเต นว. (สํขิตฺตํ.)}

\prose{เหตุปจฺจยา นอธิปติยา นว. (สํขิตฺตํ.)}

\prose{นเหตุปจฺจยา อารมฺมเณ เทฺว. (สํขิตฺตํ.)}

\prose{(สหชาตวาโรปิ ปจฺจยวาโรปิ นิสฺสยวาโรปิ สํสฏฺฐ\-วาโรปิ สมฺปยุตฺตวาโรปิ ปฏิจฺจ\-วาร\-สทิสา วิตฺถาเรตพฺพา.)}

\csromanmarkh{2. อนาคตา\-รมฺมณปท}\csromanmarkhii{7. ปญฺหาวาร}\csromanheadernomark{2}{2. อนาคตา\-รมฺมณปท 7. ปญฺหาวาร}
\prose{ปจฺจยจตุกฺกเหตฺวาทิ}

\proseitem{497}{เหตุ อนาคตารมฺมโณ ธมฺโม เหตุสฺส อนาคตา\-รมฺมณสฺส ธมฺมสฺส เหตุปจฺจเยน ปจฺจโย. ตีณิ.}

\prose{เหตุ อนาคตารมฺมโณ ธมฺโม เหตุสฺส อนาคตา\-รมฺมณสฺส ธมฺมสฺส อารมฺมณ\-ปจฺจเยน ปจฺจโย. นว.}

\csromanmarkh{1. เหตุทุก}\csromanmarkhii{19. อตีตา\-รมฺมณ\-ตฺติก}\csromanheadernomark{2}{1. เหตุทุก 19. อตีตา\-รมฺมณ\-ตฺติก}
\prose{\csromanfolio{311}เหตุ อนาคตารมฺมโณ ธมฺโม เหตุสฺส อนาคตา\-รมฺมณสฺส ธมฺมสฺส อธิปติ\-ปจฺจเยน ปจฺจโย. อารมฺมณา\-ธิปติ, สหชาตา\-ธิปติ. ตีณิ.}

\prose{นเหตุํ อนาคตารมฺมโณ ธมฺโม นเหตุสฺส อนาคตา\-รมฺมณสฺส ธมฺมสฺส อธิปติ\-ปจฺจเยน ปจฺจโย. อารมฺมณา\-ธิปติ, สหชาตา\-ธิปติ. ตีณิ.}

\prose{เหตุ อนาคตารมฺมโณ จ นเหตุ อนาคตารมฺมโณ จ ธมฺมา เหตุสฺส อนาคตา\-รมฺมณสฺส ธมฺมสฺส อธิปติ\-ปจฺจเยน ปจฺจโย. อารมฺมณา\-ธิปติ. ตีณิ. (สํขิตฺตํ.)}

\proseitem{498}{เหตุยา ตีณิ, อารมฺมเณ นว, อธิปติยา นว, อนนฺตเร นว, สมนนฺตเร นว, สหชาเต นว, อญฺญมญฺเญ นว, นิสฺสเย นว, อุปนิสฺสเย นว, อาเสวเน นว, กมฺเม ตีณิ, วิปาเก นว, อาหาเร ตีณิ ฯเปฯ อวิคเต นว. (สํขิตฺตํ.)}

\proseitem{499}{เหตุ อนาคตารมฺมโณ ธมฺโม เหตุสฺส อนาคตา\-รมฺมณสฺส ธมฺมสฺส อารมฺมณ\-ปจฺจเยน ปจฺจโย. สหชาต\-ปจฺจเยน ปจฺจโย. อุปนิสฺสย\-ปจฺจเยน ปจฺจโย. (สํขิตฺตํ.)}

\proseitem{500}{นเหตุยา นว, นอารมฺมเณ นว. (สํขิตฺตํ.)}

\prose{เหตุปจฺจยา นอารมฺมเณ ตีณิ. (สํขิตฺตํ.)}

\prose{นเหตุปจฺจยา อารมฺมเณ นว. (สํขิตฺตํ.)}

\prose{(ยถา กุสลตฺติเก ปญฺหาวารสฺส อนุโลมมฺปิ ปจฺจนียมฺปิ อนุโลมปจฺจนียมฺปิ ปจฺจนียานุโลมมฺปิ คณิตํ, เอวํ คเณตพฺพํ.)}

\csromanmarkh{ธมฺมานุโลม ทุกติก\-ปฏฺฐาน\-ปาฬิ}\csromanmarkhii{3. ปจฺจุปฺปนฺนา\-รมฺมณปท 1. ปฏิจฺจวาร}\csromanheadernomark{2}{ธมฺมานุโลม ทุกติก\-ปฏฺฐาน\-ปาฬิ 3. ปจฺจุปฺปนฺนา\-รมฺมณปท 1. ปฏิจฺจวาร}
\prose{\csromanfolio{312}ปจฺจยจตุกฺกเหตุ}

\proseitem{501}{เหตุํ ปจฺจุปฺปนฺนา\-รมฺมณํ ธมฺมํ ปฏิจฺจ เหตุ ปจฺจุปฺปนฺนา\-รมฺมโณ ธมฺโม อุปฺปชฺชติ เหตุปจฺจยา. ตีณิ.}

\prose{นเหตุํ ปจฺจุปฺปนฺนา\-รมฺมณํ ธมฺมํ ปฏิจฺจ นเหตุ ปจฺจุปฺปนฺนา\-รมฺมโณ ธมฺโม อุปฺปชฺชติ เหตุปจฺจยา. ตีณิ.}

\prose{เหตุํ ปจฺจุปฺปนฺนา\-รมฺมณญฺจ นเหตุํ ปจฺจุปฺปนฺนา\-รมฺมณญฺจ ธมฺมํ ปฏิจฺจ เหตุ ปจฺจุปฺปนฺนา\-รมฺมโณ ธมฺโม อุปฺปชฺชติ เหตุปจฺจยา. ตีณิ. (สํขิตฺตํ.)}

\proseitem{502}{เหตุยา นว, อารมฺมเณ นว, อธิปติยา นว ฯเปฯ กมฺเม นว ฯเปฯ อวิคเต นว. (สํขิตฺตํ.)}

\prose{นเหตุ นอธิปติ}

\proseitem{503}{นเหตุํ ปจฺจุปฺปนฺนา\-รมฺมณํ ธมฺมํ ปฏิจฺจ นเหตุ ปจฺจุปฺปนฺนา\-รมฺมโณ ธมฺโม อุปฺปชฺชติ นเหตุปจฺจยา.  นเหตุํ ปจฺจุปฺปนฺนา\-รมฺมณํ ธมฺมํ ปฏิจฺจ เหตุ ปจฺจุปฺปนฺนา\-รมฺมโณ ธมฺโม อุปฺปชฺชติ นเหตุปจฺจยา. (2)}

\prose{เหตุํ ปจฺจุปฺปนฺนา\-รมฺมณํ ธมฺมํ ปฏิจฺจ เหตุ ปจฺจุปฺปนฺนา\-รมฺมโณ ธมฺโม อุปฺปชฺชติ นอธิปติ\-ปจฺจยา. (สํขิตฺตํ.)}

\proseitem{504}{นเหตุยา เทฺว, นอธิปติยา นว, นปุเรชาเต นว, นปจฺฉาชาเต นว, นอาเสวเน นว, นกมฺเม ตีณิ, นวิปาเก นว, นฌาเน เอกํ, นมคฺเค เอกํ, นวิปฺปยุตฺเต นว. (สํขิตฺตํ.)}

\prose{เหตุปจฺจยา นอธิปติยา นว. (สํขิตฺตํ.)}

\prose{นเหตุปจฺจยา อารมฺมเณ เทฺว. (สํขิตฺตํ.)}

\prose{(สหชาตวาโรปิ ปจฺจยวาโรปิ นิสฺสยวาโรปิ สํสฏฺฐ\-วาโรปิ สมฺปยุตฺตวาโรปิ ปฏิจฺจ\-วาร\-สทิสา วิตฺถาเรตพฺพา.)}

\vspace{\tipitakaparskip}
\csromancenter{เหตุทุก 19. อตีตา\-รมฺมณ\-ตฺติก 3. ปจฺจุปฺปนฺนา\-รมฺมณปท 7. ปญฺหาวาร}
\prose{\csromanfolio{313}ปจฺจยจตุกฺกเหตฺวาทิ}

\proseitem{505}{เหตุ ปจฺจุปฺปนฺนา\-รมฺมโณ ธมฺโม เหตุสฺส ปจฺจุปฺปนฺนา\-รมฺมณสฺส ธมฺมสฺส เหตุปจฺจเยน ปจฺจโย. ตีณิ.}

\prose{เหตุ ปจฺจุปฺปนฺนา\-รมฺมโณ ธมฺโม เหตุสฺส ปจฺจุปฺปนฺนา\-รมฺมณสฺส ธมฺมสฺส อารมฺมณ\-ปจฺจเยน ปจฺจโย. นว.}

\prose{เหตุ ปจฺจุปฺปนฺนา\-รมฺมโณ ธมฺโม เหตุสฺส ปจฺจุปฺปนฺนา\-รมฺมณสฺส ธมฺมสฺส อธิปติ\-ปจฺจเยน ปจฺจโย. อารมฺมณา\-ธิปติ, สหชาตา\-ธิปติ. ตีณิ.}

\prose{ปจฺจุปฺปนฺนา\-รมฺมโณ ธมฺโม นเหตุสฺส ปจฺจุปฺปนฺนา\-รมฺมณสฺส ธมฺมสฺส อธิปติ\-ปจฺจเยน ปจฺจโย. อารมฺมณา\-ธิปติ, สหชาตา\-ธิปติ. ตีณิ.}

\prose{เหตุ ปจฺจุปฺปนฺนา\-รมฺมโณ จ นเหตุ ปจฺจุปฺปนฺนา\-รมฺมโณ จ ธมฺมา เหตุสฺส ปจฺจุปฺปนฺนา\-รมฺมณสฺส ธมฺมสฺส อธิปติ\-ปจฺจเยน ปจฺจโย. อารมฺมณา\-ธิปติ. ตีณิ. (สํขิตฺตํ.)}

\proseitem{506}{เหตุยา ตีณิ, อารมฺมเณ นว, อธิปติยา นว, อนนฺตเร นว, สมนนฺตเร นว, สหชาเต นว, อญฺญมญฺเญ นว, นิสฺสเย นว, อุปนิสฺสเย นว, อาเสวเน นว, กมฺเม ตีณิ, วิปาเก นว, อาหาเร ตีณิ, อินฺทฺริเย นว ฯเปฯ อวิคเต นว. (สํขิตฺตํ.)}

\proseitem{507}{เหตุ ปจฺจุปฺปนฺนา\-รมฺมโณ ธมฺโม เหตุสฺส ปจฺจุปฺปนฺนา\-รมฺมณสฺส ธมฺมสฺส อารมฺมณ\-ปจฺจเยน ปจฺจโย. สหชาต\-ปจฺจเยน ปจฺจโย. อุปนิสฺสย\-ปจฺจเยน ปจฺจโย. (สํขิตฺตํ.)}

\proseitem{508}{นเหตุยา นว, นอารมฺมเณ นว. (สํขิตฺตํ.)}

\prose{เหตุปจฺจยา นอารมฺมเณ ตีณิ. (สํขิตฺตํ.)}

\prose{นเหตุปจฺจยา อารมฺมเณ นว. (สํขิตฺตํ.)}

\prosewithcloser{\csromanfolio{314}(ยถา กุสลตฺติเก ปญฺหาวารสฺส อนุโลมมฺปิ ปจฺจนียมฺปิ อนุโลมปจฺจนียมฺปิ ปจฺจนียานุโลมมฺปิ คณิตํ, เอวํ คเณตพฺพํ.)}{เหตุทุกอตีตารมฺมณตฺติกํ นิฏฺฐิตํ.}
\csromanmatikaanchor{mtk.51}
\csromantocmarkref{section}{20. อชฺฌตฺตตฺติก}{mtk.51}
\bookmark[dest={mtk.51},level=1]{20. อชฺฌตฺตตฺติก}
\csromanmarkh{1. เหตุทุก}\csromanmarkhii{20. อชฺฌตฺตตฺติก}\csromanheadernomark{1}{1. เหตุทุก 20. อชฺฌตฺตตฺติก}
\csromanmarkh{1. อชฺฌตฺตปท}\csromanmarkhii{1. ปฏิจฺจวาร}\csromanheadernomark{2}{1. \textbf{อชฺฌตฺตปท 1. ปฏิจฺจวาร}}
\prose{ปจฺจยจตุกฺกเหตุ}

\proseitem{509}{เหตุํ อชฺฌตฺตํ ธมฺมํ ปฏิจฺจ เหตุ อชฺฌตฺโต ธมฺโม อุปฺปชฺชติ เหตุปจฺจยา. ตีณิ.}

\prose{นเหตุํ อชฺฌตฺตํ ธมฺมํ ปฏิจฺจ นเหตุ อชฺฌตฺโต ธมฺโม อุปฺปชฺชติ เหตุปจฺจยา. ตีณิ.}

\prose{เหตุํ อชฺฌตฺตญฺจ นเหตุํ อชฺฌตฺตญฺจ ธมฺมํ ปฏิจฺจ เหตุ อชฺฌตฺโต ธมฺโม อุปฺปชฺชติ เหตุปจฺจยา. ตีณิ. (สํขิตฺตํ.)}

\proseitem{510}{เหตุยา นว, อารมฺมเณ นว, อธิปติยา นว ฯเปฯ กมฺเม นว, วิปาเก นว, อาหาเร นว ฯเปฯ อวิคเต นว. (สํขิตฺตํ.)}

\csromanheader{2}{\textbf{นเหตฺวาทิ}}
\proseitem{511}{นเหตุํ อชฺฌตฺตํ ธมฺมํ ปฏิจฺจ นเหตุ อชฺฌตฺโต ธมฺโม อุปฺปชฺชติ นเหตุปจฺจยา.  นเหตุํ อชฺฌตฺตํ ธมฺมํ ปฏิจฺจ เหตุ อชฺฌตฺโต ธมฺโม อุปฺปชฺชติ นเหตุปจฺจยา. (2)}

\prose{เหตุํ อชฺฌตฺตํ ธมฺมํ ปฏิจฺจ นเหตุ อชฺฌตฺโต ธมฺโม อุปฺปชฺชติ นอารมฺมณ\-ปจฺจยา. (1)}

\prose{นเหตุํ อชฺฌตฺตํ ธมฺมํ ปฏิจฺจ นเหตุ อชฺฌตฺโต ธมฺโม อุปฺปชฺชติ นอารมฺมณ\-ปจฺจยา. (1)}

\prose{เหตุํ อชฺฌตฺตญฺจ นเหตุํ อชฺฌตฺตญฺจ ธมฺมํ ปฏิจฺจ นเหตุ อชฺฌตฺโต ธมฺโม อุปฺปชฺชติ นอารมฺมณ\-ปจฺจยา. (1)}

\csromanmarkh{1. เหตุทุก}\csromanmarkhii{20. อชฺฌตฺตตฺติก}\csromanheadernomark{2}{1. เหตุทุก 20. อชฺฌตฺตตฺติก}
\prose{\csromanfolio{315}เหตุํ อชฺฌตฺตํ ธมฺมํ ปฏิจฺจ เหตุ อชฺฌตฺโต ธมฺโม อุปฺปชฺชติ นอธิปติ\-ปจฺจยา. (สํขิตฺตํ.)}

\proseitem{512}{นเหตุยา เทฺว, นอารมฺมเณ ตีณิ, นอธิปติยา นว, นอนนฺตเร ตีณิ, นสมนนฺตเร ตีณิ, นอญฺญมญฺเญ ตีณิ, นอุปนิสฺสเย ตีณิ, นปุเรชาเต นว, นปจฺฉาชาเต นว, นอาเสวเน นว, นกมฺเม ตีณิ, นวิปาเก นว, นอาหาเร เอกํ, นอินฺทฺริเย เอกํ, นฌาเน เอกํ, นมคฺเค เอกํ, นสมฺปยุตฺเต ตีณิ, นวิปฺปยุตฺเต นว, โนนตฺถิยา ตีณิ, โนวิคเต ตีณิ. (สํขิตฺตํ.)}

\prose{เหตุปจฺจยา นอารมฺมเณ ตีณิ. (สํขิตฺตํ.)}

\prose{นเหตุปจฺจยา อารมฺมเณ เทฺว. (สํขิตฺตํ.)}

\prose{(สหชาตวาโรปิ ปจฺจยวาโรปิ นิสฺสยวาโรปิ สํสฏฺฐ\-วาโรปิ สมฺปยุตฺตวาโรปิ ปฏิจฺจ\-วาร\-สทิสา วิตฺถาเรตพฺพา.)}

\csromanmarkh{1. อชฺฌตฺตปท}\csromanmarkhii{7. ปญฺหาวาร}\csromanheadernomark{2}{1. \textbf{อชฺฌตฺตปท 7. ปญฺหาวาร}}
\prose{ปจฺจยจตุกฺกเหตฺวาทิ}

\proseitem{513}{เหตุ อชฺฌตฺโต ธมฺโม เหตุสฺส อชฺฌตฺตสฺส ธมฺมสฺส เหตุปจฺจเยน ปจฺจโย. ตีณิ.}

\prose{เหตุ อชฺฌตฺโต ธมฺโม เหตุสฺส อชฺฌตฺตสฺส ธมฺมสฺส อารมฺมณ\-ปจฺจเยน ปจฺจโย. นว.}

\prose{เหตุ อชฺฌตฺโต ธมฺโม เหตุสฺส อชฺฌตฺตสฺส ธมฺมสฺส อธิปติ\-ปจฺจเยน ปจฺจโย. อารมฺมณา\-ธิปติ, สหชาตา\-ธิปติ. ตีณิ.}

\prose{อชฺฌตฺโต ธมฺโม นเหตุสฺส อชฺฌตฺตสฺส ธมฺมสฺส อธิปติ\-ปจฺจเยน ปจฺจโย. อารมฺมณา\-ธิปติ, สหชาตา\-ธิปติ. ตีณิ.}

\prose{เหตุ อชฺฌตฺโต จ นเหตุ อชฺฌตฺโต จ ธมฺมา เหตุสฺส อชฺฌตฺตสฺส ธมฺมสฺส อธิปติ\-ปจฺจเยน ปจฺจโย. อารมฺมณา\-ธิปติ. ตีณิ. (สํขิตฺตํ.)}

\proseitem{514}{\csromanfolio{316}เหตุยา ตีณิ, อารมฺมเณ นว, อธิปติยา นว, อนนฺตเร นว, สมนนฺตเร นว, สหชาเต นว, อญฺญมญฺเญ นว, นิสฺสเย นว, อุปนิสฺสเย นว, ปุเรชาเต ตีณิ, ปจฺฉาชาเต ตีณิ, อาเสวเน นว, กมฺเม ตีณิ, วิปาเก นว, อาหาเร ตีณิ, อินฺทฺริเย นว, ฌาเน ตีณิ, มคฺเค นว, สมฺปยุตฺเต นว, วิปฺปยุตฺเต ปญฺจ, อตฺถิยา นว, นตฺถิยา นว, วิคเต นว, อวิคเต นว. (สํขิตฺตํ.)}

\proseitem{515}{เหตุ อชฺฌตฺโต ธมฺโม เหตุสฺส อชฺฌตฺตสฺส ธมฺมสฺส อารมฺมณ\-ปจฺจเยน ปจฺจโย. สหชาต\-ปจฺจเยน ปจฺจโย. อุปนิสฺสย\-ปจฺจเยน ปจฺจโย. (สํขิตฺตํ.)}

\vspace{\tipitakaparskip}
\csromancenter{\textbf{นเหตุ} นว, นอารมฺมเณ นว. (สํขิตฺตํ.)}
\prose{เหตุปจฺจยา นอารมฺมเณ ตีณิ. (สํขิตฺตํ.)}

\prose{นเหตุปจฺจยา อารมฺมเณ นว. (สํขิตฺตํ.)}

\prose{(ยถา กุสลตฺติเก ปญฺหาวารสฺส อนุโลมมฺปิ ปจฺจนียมฺปิ อนุโลมปจฺจนียมฺปิ ปจฺจนียานุโลมมฺปิ คณิตํ, เอวํ คเณตพฺพํ.)}

\csromanmarkh{2. พหิทฺธาปท}\csromanmarkhii{1. ปฏิจฺจวาร}\csromanheadernomark{2}{2. \textbf{พหิทฺธาปท 1. ปฏิจฺจวาร}}
\prose{ปจฺจยจตุกฺกเหตุ}

\proseitem{517}{เหตุํ พหิทฺธา ธมฺมํ ปฏิจฺจ เหตุ พหิทฺธา ธมฺโม อุปฺปชฺชติ เหตุปจฺจยา. (สํขิตฺตํ.)}

\proseitem{518}{เหตุยา นว, อารมฺมเณ นว ฯเปฯ กมฺเม นว, วิปาเก นว ฯเปฯ อวิคเต นว. (สํขิตฺตํ.)}

\vspace{\tipitakaparskip}
\csromancenter{\textbf{นเหตุ}}
\proseitem{519}{\textbf{นเหตุ}ํ พหิทฺธา ธมฺมํ ปฏิจฺจ นเหตุ พหิทฺธา ธมฺโม อุปฺปชฺชติ นเหตุปจฺจยา.  นเหตุํ พหิทฺธา ธมฺมํ ปฏิจฺจ เหตุ พหิทฺธา ธมฺโม อุปฺปชฺชติ นเหตุปจฺจยา. (2) (สํขิตฺตํ.)}

\csromanmarkh{1. เหตุทุก}\csromanmarkhii{20. อชฺฌตฺตตฺติก}\csromanheadernomark{2}{1. เหตุทุก 20. อชฺฌตฺตตฺติก}
\proseitem{520}{\csromanfolio{317}นเหตุยา เทฺว, นอารมฺมเณ ตีณิ, นอธิปติยา นว, นอนนฺตเร ตีณิ, นสมนนฺตเร ตีณิ, นอญฺญมญฺเญ ตีณิ, นอุปนิสฺสเย ตีณิ, นปุเรชาเต นว, นปจฺฉาชาเต นว, นอาเสวเน นว, นกมฺเม ตีณิ, นวิปาเก นว, นอาหาเร เอกํ, นอินฺทฺริเย เอกํ, นฌาเน เอกํ, นมคฺเค เอกํ, นสมฺปยุตฺเต ตีณิ, นวิปฺปยุตฺเต นว, โนนตฺถิยา ตีณิ, โนวิคเต ตีณิ. (สํขิตฺตํ.)}

\prose{เหตุปจฺจยา นอารมฺมเณ ตีณิ. (สํขิตฺตํ.)}

\prose{นเหตุปจฺจยา อารมฺมเณ เทฺว. (สํขิตฺตํ.)}

\prose{(สหชาตวาโรปิ ปจฺจยวาโรปิ นิสฺสยวาโรปิ สํสฏฺฐ\-วาโรปิ สมฺปยุตฺตวาโรปิ ปฏิจฺจ\-วาร\-สทิสา วิตฺถาเรตพฺพา.)}

\csromanmarkh{2. พหิทฺธาปท}\csromanmarkhii{7. ปญฺหาวาร}\csromanheadernomark{2}{2. \textbf{พหิทฺธาปท 7. ปญฺหาวาร}}
\prose{ปจฺจยจตุกฺกเหตฺวาทิ}

\proseitem{521}{เหตุ พหิทฺธา ธมฺโม เหตุสฺส พหิทฺธา ธมฺมสฺส เหตุปจฺจเยน ปจฺจโย. ตีณิ.}

\prose{เหตุ พหิทฺธา ธมฺโม เหตุสฺส พหิทฺธา ธมฺมสฺส อารมฺมณ\-ปจฺจเยน ปจฺจโย. นว.}

\prose{เหตุ พหิทฺธา ธมฺโม เหตุสฺส พหิทฺธา ธมฺมสฺส อธิปติ\-ปจฺจเยน ปจฺจโย. อารมฺมณา\-ธิปติ, สหชาตา\-ธิปติ. ตีณิ.}

\prose{พหิทฺธา ธมฺโม นเหตุสฺส พหิทฺธา ธมฺมสฺส อธิปติ\-ปจฺจเยน ปจฺจโย. อารมฺมณา\-ธิปติ, สหชาตา\-ธิปติ. ตีณิ.}

\prose{เหตุ พหิทฺธา จ นเหตุ พหิทฺธา จ ธมฺมา เหตุสฺส พหิทฺธา ธมฺมสฺส อธิปติ\-ปจฺจเยน ปจฺจโย. อารมฺมณา\-ธิปติ. ตีณิ. (สํขิตฺตํ.)}

\proseitem{522}{เหตุยา ตีณิ, อารมฺมเณ นว, อธิปติยา นว ฯเปฯ ปุเรชาเต ตีณิ, อาเสวเน นว, กมฺเม ตีณิ, วิปาเก นว ฯเปฯ อวิคเต นว. (สํขิตฺตํ.)}

\csromanheader{2}{ธมฺมานุโลม ทุกติก\-ปฏฺฐาน\-ปาฬิ ปจฺจนียุ\-ทฺธาร}
\proseitem{523}{\csromanfolio{318}เหตุ พหิทฺธา ธมฺโม เหตุสฺส พหิทฺธา ธมฺมสฺส อารมฺมณ\-ปจฺจเยน ปจฺจโย. สหชาต\-ปจฺจเยน ปจฺจโย. อุปนิสฺสย\-ปจฺจเยน ปจฺจโย. (สํขิตฺตํ.)}

\proseitem{524}{นเหตุยา นว, นอารมฺมเณ นว. (สํขิตฺตํ.)}

\prose{เหตุปจฺจยา นอารมฺมเณ ตีณิ. (สํขิตฺตํ.)}

\prose{นเหตุปจฺจยา อารมฺมเณ นว. (สํขิตฺตํ.)}

\prosewithcloser{(ยถา กุสลตฺติเก ปญฺหาวารสฺส อนุโลมมฺปิ ปจฺจนียมฺปิ อนุโลมปจฺจนียมฺปิ ปจฺจนียานุโลมมฺปิ คณิตํ, เอวํ คเณตพฺพํ. อชฺฌตฺต\-พหิทฺธา น ลพฺภนฺติ.)}{เหตุ\-ทุก\-อชฺฌตฺต\-ตฺติกํ นิฏฺฐิตํ.}
\csromanmatikaanchor{mtk.52}
\csromantocmarkref{section}{21. อชฺฌตฺตารมฺมณตฺติก}{mtk.52}
\bookmark[dest={mtk.52},level=1]{21. อชฺฌตฺตารมฺมณตฺติก}
\csromanmarkh{1. เหตุทุก}\csromanmarkhii{21. อชฺฌตฺตา\-รมฺมณ\-ตฺติก}\csromanheadernomark{1}{1. เหตุทุก 21. อชฺฌตฺตา\-รมฺมณ\-ตฺติก}
\csromanmarkh{1. อชฺฌตฺตา\-รมฺมณปท}\csromanmarkhii{1. ปฏิจฺจวาร}\csromanheadernomark{2}{1. อชฺฌตฺตา\-รมฺมณปท 1. ปฏิจฺจวาร}
\prose{ปจฺจยจตุกฺกเหตุ}

\proseitem{525}{เหตุํ อชฺฌตฺตา\-รมฺมณํ ธมฺมํ ปฏิจฺจ เหตุ อชฺฌตฺตา\-รมฺมโณ ธมฺโม อุปฺปชฺชติ เหตุปจฺจยา. ตีณิ.}

\prose{นเหตุํ อชฺฌตฺตา\-รมฺมณํ ธมฺมํ ปฏิจฺจ นเหตุ อชฺฌตฺตา\-รมฺมโณ ธมฺโม อุปฺปชฺชติ เหตุปจฺจยา. ตีณิ.}

\prose{เหตุํ อชฺฌตฺตา\-รมฺมณญฺจ นเหตุํ อชฺฌตฺตา\-รมฺมณญฺจ ธมฺมํ ปฏิจฺจ เหตุ อชฺฌตฺตา\-รมฺมโณ ธมฺโม อุปฺปชฺชติ เหตุปจฺจยา. ตีณิ. (สํขิตฺตํ.)}

\proseitem{526}{เหตุยา นว, อารมฺมเณ นว, อธิปติยา นว ฯเปฯ อวิคเต นว. (สํขิตฺตํ.)}

\csromanmarkh{1. เหตุทุก}\csromanmarkhii{21. อชฺฌตฺตา\-รมฺมณ\-ตฺติก นเหตุ นอธิปติ}\csromanheadernomark{2}{1. เหตุทุก 21. อชฺฌตฺตา\-รมฺมณ\-ตฺติก นเหตุ นอธิปติ}
\proseitem{527}{\csromanfolio{319}นเหตุํ อชฺฌตฺตา\-รมฺมณํ ธมฺมํ ปฏิจฺจ นเหตุ อชฺฌตฺตา\-รมฺมโณ ธมฺโม อุปฺปชฺชติ นเหตุปจฺจยา.  นเหตุํ อชฺฌตฺตา\-รมฺมณํ ธมฺมํ ปฏิจฺจ เหตุ อชฺฌตฺตา\-รมฺมโณ ธมฺโม อุปฺปชฺชติ นเหตุปจฺจยา. (2)}

\prose{เหตุํ อชฺฌตฺตา\-รมฺมณํ ธมฺมํ ปฏิจฺจ เหตุ อชฺฌตฺตา\-รมฺมโณ ธมฺโม อุปฺปชฺชติ นอธิปติ\-ปจฺจยา. (สํขิตฺตํ.)}

\proseitem{528}{นเหตุยา เทฺว, นอธิปติยา นว, นปุเรชาเต นว, นปจฺฉาชาเต นว, นอาเสวเน นว, นกมฺเม ตีณิ, นวิปาเก นว, นฌาเน เอกํ, นมคฺเค เอกํ, นวิปฺปยุตฺเต นว. (สํขิตฺตํ.)}

\prose{เหตุปจฺจยา นอธิปติยา นว. (สํขิตฺตํ.)}

\prose{นเหตุปจฺจยา อารมฺมเณ เทฺว. (สํขิตฺตํ.)}

\prose{(สหชาตวาโรปิ ปจฺจยวาโรปิ นิสฺสยวาโรปิ สํสฏฺฐ\-วาโรปิ สมฺปยุตฺตวาโรปิ ปฏิจฺจ\-วาร\-สทิสา วิตฺถาเรตพฺพา.)}

\csromanmarkh{1. อชฺฌตฺตา\-รมฺมณปท}\csromanmarkhii{7. ปญฺหาวาร}\csromanheadernomark{2}{1. อชฺฌตฺตา\-รมฺมณปท 7. ปญฺหาวาร}
\prose{ปจฺจยจตุกฺกเหตุ}

\proseitem{529}{เหตุ อชฺฌตฺตา\-รมฺมโณ ธมฺโม เหตุสฺส อชฺฌตฺตา\-รมฺมณสฺส ธมฺมสฺส เหตุปจฺจเยน ปจฺจโย. (สํขิตฺตํ.)}

\proseitem{530}{เหตุยา ตีณิ, อารมฺมเณ นว, อธิปติยา นว ฯเปฯ อวิคเต นว. (สํขิตฺตํ.)}

\proseitem{531}{เหตุ อชฺฌตฺตา\-รมฺมโณ ธมฺโม เหตุสฺส อชฺฌตฺตา\-รมฺมณสฺส ธมฺมสฺส อารมฺมณ\-ปจฺจเยน ปจฺจโย. สหชาต\-ปจฺจเยน ปจฺจโย. อุปนิสฺสย\-ปจฺจเยน ปจฺจโย. (สํขิตฺตํ.)}

\vspace{\tipitakaparskip}
\csromancenter{\textbf{นเหตุ}ยา นว, นอารมฺมเณ นว. (สํขิตฺตํ.)}
\prose{\csromanfolio{320}เหตุปจฺจยา นอารมฺมเณ ตีณิ. (สํขิตฺตํ.)}

\prose{นเหตุปจฺจยา อารมฺมเณ นว. (สํขิตฺตํ.)}

\prose{(ยถา กุสลตฺติเก ปญฺหาวารสฺส อนุโลมมฺปิ ปจฺจนียมฺปิ อนุโลมปจฺจนียมฺปิ ปจฺจนียานุโลมมฺปิ คณิตํ, เอวํ คเณตพฺพํ.)}

\csromanmarkh{2. พหิทฺธา\-รมฺมณปท}\csromanmarkhii{1. ปฏิจฺจวาร}\csromanheadernomark{2}{2. พหิทฺธา\-รมฺมณปท 1. ปฏิจฺจวาร}
\prose{ปจฺจยจตุกฺกเหตุ}

\proseitem{533}{เหตุํ พหิทฺธา\-รมฺมณํ ธมฺมํ ปฏิจฺจ เหตุ พหิทฺธา\-รมฺมโณ ธมฺโม อุปฺปชฺชติ เหตุปจฺจยา. (สํขิตฺตํ.)}

\proseitem{534}{เหตุยา นว, อารมฺมเณ นว ฯเปฯ อวิคเต นว. (สํขิตฺตํ.)}

\vspace{\tipitakaparskip}
\csromancenter{\textbf{นเหตุ}}
\proseitem{535}{นเหตุํ พหิทฺธา\-รมฺมณํ ธมฺมํ ปฏิจฺจ นเหตุ พหิทฺธา\-รมฺมโณ ธมฺโม อุปฺปชฺชติ นเหตุปจฺจยา. (สํขิตฺตํ.)}

\vspace{\tipitakaparskip}
\csromancenter{\textbf{นเหตุ}ยา เทฺว, นอธิปติยา นว ฯเปฯ นปุเรชาเต นว, นปจฺฉาชาเต นว, นอาเสวเน นว, นกมฺเม ตีณิ, นวิปาเก นว, นฌาเน เอกํ, นมคฺเค เอกํ, นวิปฺปยุตฺเต นว. (สํขิตฺตํ.)}
\prose{เหตุปจฺจยา นอธิปติยา นว. (สํขิตฺตํ.)}

\prose{\textbf{นเหตุ}ปจฺจยา อารมฺมเณ เทฺว. (สํขิตฺตํ.)}

\prose{(สหชาตวาโรปิ ปจฺจยวาโรปิ นิสฺสยวาโรปิ สํสฏฺฐ\-วาโรปิ สมฺปยุตฺตวาโรปิ ปฏิจฺจ\-วาร\-สทิสา วิตฺถาเรตพฺพา.)}

\csromanmarkh{2. พหิทฺธา\-รมฺมณปท}\csromanmarkhii{7. ปญฺหาวาร}\csromanheadernomark{2}{2. พหิทฺธา\-รมฺมณปท 7. ปญฺหาวาร}
\prose{ปจฺจยจตุกฺกเหตฺวาทิ}

\proseitem{537}{เหตุ พหิทฺธา\-รมฺมโณ ธมฺโม เหตุสฺส พหิทฺธา\-รมฺมณสฺส ธมฺมสฺส เหตุปจฺจเยน ปจฺจโย. ตีณิ.}

\csromanmatikaanchor{mtk.53}
\csromantocmarkref{section}{22. สนิทสฺสนสปฺปฏิฆตฺติก}{mtk.53}
\bookmark[dest={mtk.53},level=1]{22. สนิทสฺสนสปฺปฏิฆตฺติก}
\csromanmarkh{1. เหตุทุก}\csromanmarkhii{22. สนิทสฺสนสปฺปฏิฆ\-ตฺติก}\csromanheadernomark{1}{1. เหตุทุก 22. สนิทสฺสนสปฺปฏิฆ\-ตฺติก}
\prose{\csromanfolio{321}เหตุ พหิทฺธา\-รมฺมโณ ธมฺโม เหตุสฺส พหิทฺธา\-รมฺมณสฺส ธมฺมสฺส อารมฺมณ\-ปจฺจเยน ปจฺจโย. นว.}

\prose{เหตุ พหิทฺธา\-รมฺมโณ ธมฺโม เหตุสฺส พหิทฺธา\-รมฺมณสฺส ธมฺมสฺส อธิปติ\-ปจฺจเยน ปจฺจโย. สหชาตา\-ธิปติ. ตีณิ.}

\prose{พหิทฺธา\-รมฺมโณ ธมฺโม นเหตุสฺส พหิทฺธา\-รมฺมณสฺส ธมฺมสฺส อธิปติ\-ปจฺจเยน ปจฺจโย. สหชาตา\-ธิปติ. ตีณิ. (สํขิตฺตํ.)}

\proseitem{538}{เหตุยา ตีณิ, อารมฺมเณ นว, อธิปติยา ฉ, กมฺเม ตีณิ, วิปาเก นว, อาหาเร ตีณิ ฯเปฯ อวิคเต นว. (สํขิตฺตํ.)}

\proseitem{539}{เหตุ พหิทฺธา\-รมฺมโณ ธมฺโม เหตุสฺส พหิทฺธา\-รมฺมณสฺส ธมฺมสฺส อารมฺมณ\-ปจฺจเยน ปจฺจโย. สหชาต\-ปจฺจเยน ปจฺจโย. อุปนิสฺสย\-ปจฺจเยน ปจฺจโย. (สํขิตฺตํ.)}

\proseitem{540}{นเหตุยา นว, นอารมฺมเณ นว. (สํขิตฺตํ.)}

\prose{เหตุปจฺจยา นอารมฺมเณ ตีณิ. (สํขิตฺตํ.)}

\prose{นเหตุปจฺจยา อารมฺมเณ นว. (สํขิตฺตํ.)}

\prosewithcloser{(ยถา กุสลตฺติเก ปญฺหาวารสฺส อนุโลมมฺปิ ปจฺจนียมฺปิ อนุโลมปจฺจนียมฺปิ ปจฺจนียานุโลมมฺปิ คณิตํ, เอวํ คเณตพฺพํ.)}{เหตุทุกอชฺฌตฺตารมฺมณตฺติกํ นิฏฺฐิตํ.}
\csromanmarkh{1. เหตุทุก}\csromanmarkhii{22. สนิทสฺสนสปฺปฏิฆ\-ตฺติก}\csromanheadernomark{2}{1. เหตุทุก 22. สนิทสฺสนสปฺปฏิฆ\-ตฺติก}
\csromanmarkh{1. อนิทสฺสนสปฺปฏิฆปท}\csromanmarkhii{1. ปฏิจฺจวาร}\csromanheadernomark{2}{1. \textbf{อนิทสฺสนสปฺปฏิฆปท 1. ปฏิจฺจวาร}}
\prose{ปจฺจยจตุกฺกเหตฺวาทิ}

\proseitem{541}{นเหตุํ อนิทสฺสน\-สปฺปฏิฆํ ธมฺมํ ปฏิจฺจ นเหตุ อนิทสฺสน\-สปฺปฏิโฆ ธมฺโม อุปฺปชฺชติ เหตุปจฺจยา. อธิปติ\-ปจฺจยา. สหชาต\-ปจฺจยา. \csromanfolio{322}อญฺญมญฺญ\-ปจฺจยา. นิสฺสยปจฺจยา. กมฺมปจฺจยา. วิปากปจฺจยา. อาหารปจฺจยา. อินฺทฺริยปจฺจยา. ฌานปจฺจยา. มคฺคปจฺจยา. วิปฺปยุตฺต\-ปจฺจยา. อตฺถิปจฺจยา. อวิคตปจฺจยา.}

\proseitem{542}{เหตุยา เอกํ, อธิปติยา เอกํ, สหชาเต เอกํ, อญฺญมญฺเญ เอกํ, นิสฺสเย เอกํ, กมฺเม เอกํ, วิปาเก เอกํ, มคฺเค เอกํ, วิปฺปยุตฺเต เอกํ, อตฺถิยา เอกํ, อวิคเต เอกํ. (สํขิตฺตํ.)}

\proseitem{543}{นเหตุยา เอกํ, นอารมฺมเณ เอกํ, นอธิปติยา เอกํ, (สพฺเพ ปจฺจยา กาตพฺพา.) ฯเปฯ โนวิคเต เอกํ. (สํขิตฺตํ.)}

\csromanmarkh{1. อนิทสฺสนสปฺปฏิฆปท}\csromanmarkhii{7. ปญฺหาวาร}\csromanheadernomark{2}{1. \textbf{อนิทสฺสนสปฺปฏิฆปท 7. ปญฺหาวาร}}
\prose{ปจฺจยจตุกฺกสหชาตาทิ}

\proseitem{544}{อนิทสฺสน\-สปฺปฏิโฆ ธมฺโม นเหตุสฺส อนิทสฺสน\-สปฺปฏิฆสฺส ธมฺมสฺส สหชาต\-ปจฺจเยน ปจฺจโย. อญฺญมญฺญ\-ปจฺจเยน ปจฺจโย. นิสฺสย\-ปจฺจเยน ปจฺจโย. อตฺถิปจฺจเยน ปจฺจโย. อวิคต\-ปจฺจเยน ปจฺจโย. (สพฺพตฺถ เอกํ.)}

\prose{2. อนิทสฺสนอปฺปฏิฆปท 1. ปฏิจฺจวาร}

\prose{ปจฺจยจตุกฺกเหตุ}

\proseitem{545}{เหตุํ อนิทสฺสน\-อปฺปฏิฆํ ธมฺมํ ปฏิจฺจ เหตุ อนิทสฺสน\-อปฺปฏิโฆ ธมฺโม อุปฺปชฺชติ เหตุปจฺจยา. ตีณิ.}

\prose{นเหตุํ อนิทสฺสน\-อปฺปฏิฆํ ธมฺมํ ปฏิจฺจ นเหตุ อนิทสฺสน\-อปฺปฏิโฆ ธมฺโม อุปฺปชฺชติ เหตุปจฺจยา. ตีณิ.}

\prose{เหตุํ อนิทสฺสนอปฺปฏิฆญฺจ นเหตุํ อนิทสฺสนอปฺปฏิฆญฺจ ธมฺมํ ปฏิจฺจ เหตุ อนิทสฺสน\-อปฺปฏิโฆ ธมฺโม อุปฺปชฺชติ เหตุปจฺจยา. ตีณิ. (สํขิตฺตํ.)}

\csromanmarkh{1. เหตุทุก}\csromanmarkhii{22. สนิทสฺสนสปฺปฏิฆ\-ตฺติก}\csromanheadernomark{2}{1. เหตุทุก 22. สนิทสฺสนสปฺปฏิฆ\-ตฺติก}
\proseitem{546}{\csromanfolio{323}เหตุยา นว, อารมฺมเณ นว ฯเปฯ กมฺเม นว, วิปาเก นว, อาหาเร นว ฯเปฯ อวิคเต นว. (สํขิตฺตํ.)}

\prose{นเหตุ นอธิปติ}

\proseitem{547}{นเหตุํ อนิทสฺสน\-อปฺปฏิฆํ ธมฺมํ ปฏิจฺจ นเหตุ อนิทสฺสน\-อปฺปฏิโฆ ธมฺโม อุปฺปชฺชติ นเหตุปจฺจยา.  นเหตุํ อนิทสฺสน\-อปฺปฏิฆํ ธมฺมํ ปฏิจฺจ เหตุ อนิทสฺสน\-อปฺปฏิโฆ ธมฺโม อุปฺปชฺชติ นเหตุปจฺจยา. (2)}

\prose{เหตุํ อนิทสฺสน\-อปฺปฏิฆํ ธมฺมํ ปฏิจฺจ นเหตุ อนิทสฺสน\-อปฺปฏิโฆ ธมฺโม อุปฺปชฺชติ นอารมฺมณ\-ปจฺจยา. ตีณิ.}

\prose{เหตุํ อนิทสฺสน\-อปฺปฏิฆํ ธมฺมํ ปฏิจฺจ เหตุ อนิทสฺสน\-อปฺปฏิโฆ ธมฺโม อุปฺปชฺชติ นอธิปติ\-ปจฺจยา. (สํขิตฺตํ.)}

\proseitem{548}{นเหตุยา เทฺว, นอารมฺมเณ ตีณิ, นอธิปติยา นว ฯเปฯ นปุเรชาเต นว, นปจฺฉาชาเต นว, นอาเสวเน นว, นกมฺเม ตีณิ, นวิปาเก นว, นอาหาเร เอกํ, นฌาเน เอกํ, นมคฺเค เอกํ, นสมฺปยุตฺเต ตีณิ, นวิปฺปยุตฺเต นว, โนนตฺถิยา ตีณิ, โนวิคเต ตีณิ. (สํขิตฺตํ.)}

\prose{เหตุปจฺจยา นอารมฺมเณ ตีณิ. (สํขิตฺตํ.)}

\prose{นเหตุปจฺจยา อารมฺมเณ เทฺว. (สํขิตฺตํ.)}

\prose{(สหชาตวาโรปิ ปจฺจยวาโรปิ นิสฺสยวาโรปิ สํสฏฺฐ\-วาโรปิ สมฺปยุตฺตวาโรปิ ปฏิจฺจ\-วาร\-สทิสา วิตฺถาเรตพฺพา.)}

\prose{2. อนิทสฺสนอปฺปฏิฆปท 7. ปญฺหาวาร}

\prose{ปจฺจยจตุกฺกเหตฺวาทิ}

\proseitem{549}{เหตุ อนิทสฺสน\-อปฺปฏิโฆ ธมฺโม เหตุสฺส อนิทสฺสน\-อปฺปฏิฆสฺส ธมฺมสฺส เหตุปจฺจเยน ปจฺจโย. ตีณิ.}

\prose{เหตุ อนิทสฺสน\-อปฺปฏิโฆ ธมฺโม เหตุสฺส อนิทสฺสน\-อปฺปฏิฆสฺส ธมฺมสฺส อารมฺมณ\-ปจฺจเยน ปจฺจโย. นว.}

\prose{\csromanfolio{324}เหตุ อนิทสฺสน\-อปฺปฏิโฆ ธมฺโม เหตุสฺส อนิทสฺสน\-อปฺปฏิฆสฺส ธมฺมสฺส อธิปติ\-ปจฺจเยน ปจฺจโย. อารมฺมณา\-ธิปติ, สหชาตา\-ธิปติ. ตีณิ.}

\prose{อนิทสฺสน\-อปฺปฏิโฆ ธมฺโม นเหตุสฺส อนิทสฺสน\-อปฺปฏิฆสฺส ธมฺมสฺส อธิปติ\-ปจฺจเยน ปจฺจโย. อารมฺมณา\-ธิปติ, สหชาตา\-ธิปติ. ตีณิ.}

\prose{เหตุ อนิทสฺสน\-อปฺปฏิโฆ จ นเหตุ อนิทสฺสน\-อปฺปฏิโฆ จ ธมฺมา เหตุสฺส อนิทสฺสน\-อปฺปฏิฆสฺส ธมฺมสฺส อธิปติ\-ปจฺจเยน ปจฺจโย. อารมฺมณา\-ธิปติ. ตีณิ. (สํขิตฺตํ.)}

\proseitem{550}{เหตุยา ตีณิ, อารมฺมเณ นว, อธิปติยา นว, อนนฺตเร นว, สมนนฺตเร นว, สหชาเต นว, อญฺญมญฺเญ นว, นิสฺสเย นว, อุปนิสฺสเย นว, ปุเรชาเต ตีณิ, ปจฺฉาชาเต ตีณิ, อาเสวเน นว, กมฺเม ตีณิ, วิปาเก นว, อาหาเร ตีณิ, อินฺทฺริเย นว, ฌาเน ตีณิ, มคฺเค นว, สมฺปยุตฺเต นว, วิปฺปยุตฺเต ปญฺจ, อตฺถิยา นว, นตฺถิยา นว, วิคเต นว, อวิคเต นว. (สํขิตฺตํ.)}

\proseitem{551}{เหตุ อนิทสฺสน\-อปฺปฏิโฆ ธมฺโม เหตุสฺส อนิทสฺสน\-อปฺปฏิฆสฺส ธมฺมสฺส อารมฺมณ\-ปจฺจเยน ปจฺจโย. สหชาตปจฺจเยนปจฺจโย. อุปนิสฺสย\-ปจฺจเยน ปจฺจโย. (สํขิตฺตํ.)}

\proseitem{552}{นเหตุยา นว, นอารมฺมเณ นว. (สํขิตฺตํ.)}

\prose{เหตุปจฺจยา นอารมฺมเณ ตีณิ. (สํขิตฺตํ.)}

\prose{นเหตุปจฺจยา อารมฺมเณ นว. (สํขิตฺตํ.)}

\prosewithcloser{(ยถา กุสลตฺติเก ปญฺหาวารสฺส อนุโลมมฺปิ ปจฺจนียมฺปิ อนุโลมปจฺจนียมฺปิ ปจฺจนียานุโลมมฺปิ คณิตํ, เอวํ คเณตพฺพํ.)}{เหตุ\-ทุก\-สนิทสฺสนสปฺปฏิฆ\-ตฺติกํ นิฏฺฐิตํ.}
\csromanmatikaanchor{mtk.54}
\csromanmatikaanchor{mtk.55}
\csromanmatikaanchor{mtk.58}
\csromantocmarknopage{chapter}{2. สเหตุกทุก}{mtk.54}
\bookmark[dest={mtk.54},level=0]{2. สเหตุกทุก}
\csromantocmarkref{section}{1. กุสลตฺติก ...}{mtk.55}
\bookmark[dest={mtk.55},level=1]{1. กุสลตฺติก ...}
\csromanmarkh{2. สเหตุกทุก}\csromanmarkhii{1. กุสลตฺติก 2. สเหตุกทุก 1. กุสลตฺติก}\csromanheadernomark{5}{2. สเหตุกทุก 1. กุสลตฺติก \textbf{2. สเหตุกทุก 1. กุสลตฺติก}}
\csromanmarkh{1. กุสลปท}\csromanmarkhii{1. ปฏิจฺจวาร}\csromanheadernomark{2}{1. \textbf{กุสลปท 1. ปฏิจฺจวาร}}
\prose{\csromanfolio{325}ปจฺจยจตุกฺกเหตุ อารมฺมณ}

\proseitem{1}{สเหตุกํ กุสลํ ธมฺมํ ปฏิจฺจ สเหตุโก กุสโล ธมฺโม อุปฺปชฺชติ เหตุปจฺจยา. (1)}

\prose{สเหตุกํ กุสลํ ธมฺมํ ปฏิจฺจ สเหตุโก กุสโล ธมฺโม อุปฺปชฺชติ อารมฺมณ\-ปจฺจยา. (1) (สํขิตฺตํ.)}

\proseitem{2}{เหตุยา เอกํ, อารมฺมเณ เอกํ, อธิปติยา เอกํ, อนนฺตเร เอกํ, สมนนฺตเร เอกํ, สหชาเต เอกํ, อญฺญมญฺเญ เอกํ, นิสฺสเย เอกํ, อุปนิสฺสเย เอกํ, ปุเรชาเต เอกํ, อาเสวเน เอกํ, กมฺเม เอกํ, อาหาเร เอกํ ฯเปฯ อวิคเต เอกํ. (สํขิตฺตํ.)}

\prose{ปจฺจนียนอธิปติ}

\proseitem{3}{สเหตุกํ กุสลํ ธมฺมํ ปฏิจฺจ สเหตุโก กุสโล ธมฺโม อุปฺปชฺชติ นอธิปติ\-ปจฺจยา. (สํขิตฺตํ.)}

\proseitem{4}{นอธิปติยา เอกํ, นปุเรชาเต เอกํ, นปจฺฉาชาเต เอกํ, นอาเสวเน เอกํ, นกมฺเม เอกํ, นวิปาเก เอกํ, นวิปฺปยุตฺเต เอกํ. (สํขิตฺตํ.)}

\prose{เหตุปจฺจยา นอธิปติยา เอกํ. (สํขิตฺตํ.)}

\prose{นอธิปติ\-ปจฺจยา เหตุยา เอกํ. (สํขิตฺตํ.)}

\prose{(สหชาตวาโรปิ ปจฺจยวาโรปิ นิสฺสยวาโรปิ สํสฏฺฐ\-วาโรปิ สมฺปยุตฺตวาโรปิ ปฏิจฺจ\-วาร\-สทิสา วิตฺถาเรตพฺพา.)}

\csromanmarkh{1. กุสลปท}\csromanmarkhii{7. ปญฺหาวาร}\csromanheadernomark{2}{1. \textbf{กุสลปท 7. ปญฺหาวาร}}
\csromanheader{2}{\textbf{ปจฺจยจตุกฺก\csromandash{}เหตุ อารมฺมณ}}
\proseitem{5}{สเหตุโก กุสโล ธมฺโม สเหตุกสฺส กุสลสฺส ธมฺมสฺส เหตุปจฺจเยน ปจฺจโย. (1)}

\prose{\csromanfolio{326}สเหตุโก กุสโล ธมฺโม สเหตุกสฺส กุสลสฺส ธมฺมสฺส อารมฺมณ\-ปจฺจเยน ปจฺจโย. (1) (สํขิตฺตํ.)}

\proseitem{6}{เหตุยา เอกํ, อารมฺมเณ เอกํ, อธิปติยา เอกํ, อนนฺตเร เอกํ, สมนนฺตเร เอกํ, สหชาเต เอกํ, อญฺญมญฺเญ เอกํ, นิสฺสเย เอกํ, อุปนิสฺสเย เอกํ, อาเสวเน เอกํ, กมฺเม เอกํ, อาหาเร เอกํ, อินฺทฺริเย เอกํ, ฌาเน เอกํ, มคฺเค เอกํ, สมฺปยุตฺเต เอกํ, อตฺถิยา เอกํ, นตฺถิยา เอกํ, วิคเต เอกํ, อวิคเต เอกํ. (สํขิตฺตํ.)}

\proseitem{7}{สเหตุโก กุสโล ธมฺโม สเหตุกสฺส กุสลสฺส ธมฺมสฺส อารมฺมณ\-ปจฺจเยน ปจฺจโย. สหชาต\-ปจฺจเยน ปจฺจโย. อุปนิสฺสย\-ปจฺจเยน ปจฺจโย.}

\proseitem{8}{นเหตุยา เอกํ, นอารมฺมเณ เอกํ. (สํขิตฺตํ.)}

\prose{เหตุปจฺจยา นอารมฺมเณ เอกํ. (สํขิตฺตํ.)}

\prose{นเหตุปจฺจยา อารมฺมเณ เอกํ. (สํขิตฺตํ.)}

\prose{(ยถา กุสลตฺติเก ปญฺหาวารสฺส อนุโลมมฺปิ ปจฺจนียมฺปิ อนุโลมปจฺจนียมฺปิ ปจฺจนียานุโลมมฺปิ คณิตํ, เอวํ คเณตพฺพํ.)}

\csromanmarkh{2. อกุสลปท}\csromanmarkhii{1. ปฏิจฺจวาร}\csromanheadernomark{2}{2. \textbf{อกุสลปท 1. ปฏิจฺจวาร}}
\prose{ปจฺจยจตุกฺกเหตุ อารมฺมณ}

\proseitem{9}{สเหตุกํ อกุสลํ ธมฺมํ ปฏิจฺจ สเหตุโก อกุสโล ธมฺโม อุปฺปชฺชติ เหตุปจฺจยา. (1)}

\prose{อเหตุกํ อกุสลํ ธมฺมํ ปฏิจฺจ สเหตุโก อกุสโล ธมฺโม อุปฺปชฺชติ เหตุปจฺจยา. (1)}

\prose{สเหตุกํ อกุสลญฺจ อเหตุกํ อกุสลญฺจ ธมฺมํ ปฏิจฺจ สเหตุโก อกุสโล ธมฺโม อุปฺปชฺชติ เหตุปจฺจยา. (1)}

\csromanmarkh{2. สเหตุกทุก}\csromanmarkhii{1. กุสลตฺติก}\csromanheadernomark{2}{2. สเหตุกทุก 1. กุสลตฺติก}
\prose{\csromanfolio{327}สเหตุกํ อกุสลํ ธมฺมํ ปฏิจฺจ สเหตุโก อกุสโล ธมฺโม อุปฺปชฺชติ อารมฺมณ\-ปจฺจยา.  สเหตุกํ อกุสลํ ธมฺมํ ปฏิจฺจ อเหตุโก อกุสโล ธมฺโม อุปฺปชฺชติ อารมฺมณ\-ปจฺจยา.  สเหตุกํ อกุสลํ ธมฺมํ ปฏิจฺจ สเหตุโก อกุสโล จ อเหตุโก อกุสโล จ ธมฺมา อุปฺปชฺชนฺติ อารมฺมณ\-ปจฺจยา. (3)}

\prose{อเหตุกํ อกุสลํ ธมฺมํ ปฏิจฺจ สเหตุโก อกุสโล ธมฺโม อุปฺปชฺชติ อารมฺมณ\-ปจฺจยา. (1)}

\prose{สเหตุกํ อกุสลญฺจ อเหตุกํ อกุสลญฺจ ธมฺมํ ปฏิจฺจ สเหตุโก อกุสโล ธมฺโม อุปฺปชฺชติ อารมฺมณ\-ปจฺจยา. (1)}

\prose{สเหตุกํ อกุสลํ ธมฺมํ ปฏิจฺจ สเหตุโก อกุสโล ธมฺโม อุปฺปชฺชติ อธิปติ\-ปจฺจยา. (สํขิตฺตํ.)}

\proseitem{10}{เหตุยา ตีณิ, อารมฺมเณ ปญฺจ, อธิปติยา เอกํ, อนนฺตเร ปญฺจ, สมนนฺตเร ปญฺจ, สหชาเต ปญฺจ, อญฺญมญฺเญ ปญฺจ, นิสฺสเย ปญฺจ, อุปนิสฺสเย ปญฺจ, ปุเรชาเต ปญฺจ, อาเสวเน ปญฺจ, กมฺเม ปญฺจ, อาหาเร ปญฺจ, อินฺทฺริเย ปญฺจ, ฌาเน ปญฺจ, มคฺเค ปญฺจ, สมฺปยุตฺเต ปญฺจ, วิปฺปยุตฺเต ปญฺจ, อตฺถิยา ปญฺจ, นตฺถิยา ปญฺจ, วิคเต ปญฺจ, อวิคเต ปญฺจ. (สํขิตฺตํ.)}

\prose{นเหตุ นอธิปติ}

\proseitem{11}{สเหตุกํ อกุสลํ ธมฺมํ ปฏิจฺจ อเหตุโก อกุสโล ธมฺโม อุปฺปชฺชติ นเหตุปจฺจยา. (1)}

\prose{สเหตุกํ อกุสลํ ธมฺมํ ปฏิจฺจ สเหตุโก อกุสโล ธมฺโม อุปฺปชฺชติ นอธิปติ\-ปจฺจยา. (สํขิตฺตํ.)}

\proseitem{12}{นเหตุยา เอกํ, นอธิปติยา ปญฺจ, นปุเรชาเต ปญฺจ, นปจฺฉาชาเต ปญฺจ, นอาเสวเน ปญฺจ, นกมฺเม ตีณิ, นวิปาเก ปญฺจ, นวิปฺปยุตฺเต ปญฺจ. (สํขิตฺตํ.)}

\prose{เหตุปจฺจยา นอธิปติยา ตีณิ. (สํขิตฺตํ.)}

\prose{นเหตุปจฺจยา อารมฺมเณ เอกํ. (สํขิตฺตํ.)}

\prose{(สหชาตวาโรปิ ปจฺจยวาโรปิ นิสฺสยวาโรปิ สํสฏฺฐ\-วาโรปิ สมฺปยุตฺตวาโรปิ ปฏิจฺจ\-วาร\-สทิสา วิตฺถาเรตพฺพา.)}

\csromanmarkh{ธมฺมานุโลม ทุกติก\-ปฏฺฐาน\-ปาฬิ}\csromanmarkhii{2. อกุสลปท 7. ปญฺหาวาร}\csromanheadernomark{2}{ธมฺมานุโลม ทุกติก\-ปฏฺฐาน\-ปาฬิ 2. อกุสลปท 7. ปญฺหาวาร}
\prose{\csromanfolio{328}ปจฺจยจตุกฺกเหตุ อารมฺมณาทิ}

\proseitem{13}{สเหตุโก อกุสโล ธมฺโม สเหตุกสฺส อกุสลสฺส ธมฺมสฺส เหตุปจฺจเยน ปจฺจโย. เทฺว.}

\prose{สเหตุโก อกุสโล ธมฺโม สเหตุกสฺส อกุสลสฺส ธมฺมสฺส อารมฺมณ\-ปจฺจเยน ปจฺจโย. ตีณิ.}

\prose{อเหตุโก อกุสโล ธมฺโม อเหตุกสฺส อกุสลสฺส ธมฺมสฺส อารมฺมณ\-ปจฺจเยน ปจฺจโย. ตีณิ.}

\prose{สเหตุโก อกุสโล จ อเหตุโก อกุสโล จ ธมฺมา สเหตุกสฺส อกุสลสฺส ธมฺมสฺส อารมฺมณ\-ปจฺจเยน ปจฺจโย. ตีณิ.}

\prose{สเหตุโก อกุสโล ธมฺโม สเหตุกสฺส อกุสลสฺส ธมฺมสฺส อธิปติ\-ปจฺจเยน ปจฺจโย. อารมฺมณา\-ธิปติ, สหชาตา\-ธิปติ. (สํขิตฺตํ.)}

\proseitem{14}{เหตุยา เทฺว, อารมฺมเณ นว, อธิปติยา เอกํ, อนนฺตเร นว, สมนนฺตเร นว, สหชาเต ปญฺจ, อญฺญมญฺเญ ปญฺจ, นิสฺสเย ปญฺจ, อุปนิสฺสเย นว, อาเสวเน นว, กมฺเม ตีณิ, อาหาเร ตีณิ, อินฺทฺริเย ตีณิ, ฌาเน ตีณิ, มคฺเค ตีณิ, สมฺปยุตฺเต ปญฺจ, อตฺถิยา ปญฺจ, นตฺถิยา นว, วิคเต นว, อวิคเต ปญฺจ. (สํขิตฺตํ.)}

\proseitem{15}{สเหตุโก อกุสโล ธมฺโม สเหตุกสฺส อกุสลสฺส ธมฺมสฺส อารมฺมณ\-ปจฺจเยน ปจฺจโย. สหชาต\-ปจฺจเยน ปจฺจโย. อุปนิสฺสย\-ปจฺจเยน ปจฺจโย. (สํขิตฺตํ.)}

\proseitem{16}{นเหตุยา นว, นอารมฺมเณ นว. (สํขิตฺตํ.)}

\prose{เหตุปจฺจยา นอารมฺมเณ เทฺว. (สํขิตฺตํ.)}

\prose{นเหตุปจฺจยา อารมฺมเณ นว. (สํขิตฺตํ.)}

\prose{(ยถา กุสลตฺติเก ปญฺหาวารสฺส อนุโลมมฺปิ ปจฺจนียมฺปิ อนุโลมปจฺจนียมฺปิ ปจฺจนียานุโลมมฺปิ คณิตํ, เอวํ คเณตพฺพํ.)}

\csromanmarkh{2. สเหตุกทุก}\csromanmarkhii{1. กุสลตฺติก 3. อพฺยากตปท 1. ปฏิจฺจวาร}\csromanheadernomark{2}{2. สเหตุกทุก 1. กุสลตฺติก \textbf{3. อพฺยากตปท 1. ปฏิจฺจวาร}}
\prose{\csromanfolio{329}ปจฺจยจตุกฺกเหตุ}

\proseitem{17}{สเหตุกํ อพฺยากตํ ธมฺมํ ปฏิจฺจ สเหตุโก อพฺยากโต ธมฺโม อุปฺปชฺชติ เหตุปจฺจยา. ตีณิ.}

\prose{อเหตุกํ อพฺยากตํ ธมฺมํ ปฏิจฺจ อเหตุโก อพฺยากโต ธมฺโม อุปฺปชฺชติ เหตุปจฺจยา. ตีณิ.}

\prose{สเหตุกํ อพฺยากตญฺจ อเหตุกํ อพฺยากตญฺจ ธมฺมํ ปฏิจฺจ สเหตุโก อพฺยากโต ธมฺโม อุปฺปชฺชติ เหตุปจฺจยา. ตีณิ. (สํขิตฺตํ.)}

\proseitem{18}{เหตุยา นว, อารมฺมเณ จตฺตาริ, อธิปติยา ปญฺจ, อนนฺตเร จตฺตาริ, สมนนฺตเร จตฺตาริ, สหชาเต นว, อญฺญมญฺเญ ฉ ฯเปฯ ปุเรชาเต เทฺว, อาเสวเน เทฺว, กมฺเม นว, วิปาเก นว ฯเปฯ วิปฺปยุตฺเต นว ฯเปฯ อวิคเต นว. (สํขิตฺตํ.)}

\prose{นเหตุ นอารมฺมณ}

\proseitem{19}{อเหตุกํ อพฺยากตํ ธมฺมํ ปฏิจฺจ อเหตุโก อพฺยากโต ธมฺโม อุปฺปชฺชติ นเหตุปจฺจยา. (1)}

\prose{สเหตุกํ อพฺยากตํ ธมฺมํ ปฏิจฺจ อเหตุโก อพฺยากโต ธมฺโม อุปฺปชฺชติ นอารมฺมณ\-ปจฺจยา. (1)}

\prose{อเหตุกํ อพฺยากตํ ธมฺมํ ปฏิจฺจ อเหตุโก อพฺยากโต ธมฺโม อุปฺปชฺชติ นอารมฺมณ\-ปจฺจยา. (1)}

\prose{สเหตุกํ อพฺยากตญฺจ อเหตุกํ อพฺยากตญฺจ ธมฺมํ ปฏิจฺจ อเหตุโก อพฺยากโต ธมฺโม อุปฺปชฺชติ นอารมฺมณ\-ปจฺจยา. (1)}

\prose{สเหตุกํ อพฺยากตํ ธมฺมํ ปฏิจฺจ สเหตุโก อพฺยากโต ธมฺโม อุปฺปชฺชติ นอธิปติ\-ปจฺจยา. (สํขิตฺตํ.)}

\proseitem{20}{นเหตุยา เอกํ, นอารมฺมเณ ตีณิ, นอธิปติยา นว, นอนนฺตเร ตีณิ, นสมนนฺตเร ตีณิ, นอญฺญมญฺเญ ตีณิ, นอุปนิสฺสเย ตีณิ, นปุเรชาเต นว, นปจฺฉาชาเต นว, นอาเสวเน นว, นกมฺเม เทฺว,}

\prose{\csromanfolio{330}นวิปาเก ปญฺจ, นอาหาเร เอกํ, ไนนฺทฺริเย เอกํ, นฌาเน เอกํ, นมคฺเค เอกํ, นสมฺปยุตฺเต ตีณิ, นวิปฺปยุตฺเต เทฺว, โนนตฺถิยา ตีณิ, โนวิคเต ตีณิ. (สํขิตฺตํ.)}

\prose{เหตุปจฺจยา นอารมฺมเณ ตีณิ. (สํขิตฺตํ.)}

\prose{นเหตุปจฺจยา อารมฺมเณ เอกํ. (สํขิตฺตํ.)}

\prose{(สหชาตวาโรปิ ปจฺจยวาโรปิ นิสฺสยวาโรปิ สํสฏฺฐ\-วาโรปิ สมฺปยุตฺตวาโรปิ ปฏิจฺจ\-วาร\-สทิสา วิตฺถาเรตพฺพา.)}

\csromanmarkh{3. อพฺยากตปท}\csromanmarkhii{7. ปญฺหาวาร}\csromanheadernomark{2}{3. \textbf{อพฺยากตปท 7. ปญฺหาวาร}}
\prose{ปจฺจยจตุกฺกเหตุ อารมฺมณ}

\proseitem{21}{สเหตุโก อพฺยากโต ธมฺโม สเหตุกสฺส อพฺยากตสฺส ธมฺมสฺส เหตุปจฺจเยน ปจฺจโย. ตีณิ.}

\prose{สเหตุโก อพฺยากโต ธมฺโม สเหตุกสฺส อพฺยากตสฺส ธมฺมสฺส อารมฺมณ\-ปจฺจเยน ปจฺจโย. (สํขิตฺตํ.)}

\proseitem{22}{เหตุยา ตีณิ, อารมฺมเณ จตฺตาริ, อธิปติยา จตฺตาริ, อนนฺตเร จตฺตาริ, สมนนฺตเร จตฺตาริ, สหชาเต สตฺต, อญฺญมญฺเญ ฉ, นิสฺสเย สตฺต, อุปนิสฺสเย จตฺตาริ, ปุเรชาเต เทฺว, ปจฺฉาชาเต เทฺว, อาเสวเน เทฺว, กมฺเม จตฺตาริ, วิปาเก จตฺตาริ, อาหาเร จตฺตาริ อินฺทฺริเย จตฺตาริ, ฌาเน จตฺตาริ, มคฺเค ตีณิ, สมฺปยุตฺเต เทฺว, วิปฺปยุตฺเต ตีณิ, อตฺถิยา สตฺต ฯเปฯ อวิคเต สตฺต. (สํขิตฺตํ.)}

\proseitem{23}{สเหตุโก อพฺยากโต ธมฺโม สเหตุกสฺส อพฺยากตสฺส ธมฺมสฺส อารมฺมณ\-ปจฺจเยน ปจฺจโย. สหชาต\-ปจฺจเยน ปจฺจโย. อุปนิสฺสย\-ปจฺจเยน ปจฺจโย. (สํขิตฺตํ.)}

\proseitem{24}{นเหตุยา สตฺต, นอารมฺมเณ สตฺต. (สํขิตฺตํ.)}

\prose{เหตุปจฺจยา นอารมฺมเณ ตีณิ. (สํขิตฺตํ.)}

\prose{นเหตุปจฺจยา อารมฺมเณ จตฺตาริ. (สํขิตฺตํ.)}

\csromanmatikaanchor{mtk.56}
\csromanmatikaanchor{mtk.57}
\csromantocmarknopage{subsection}{3. เหตุสมฺปยุตฺตทุก}{mtk.56}
\bookmark[dest={mtk.56},level=2]{3. เหตุสมฺปยุตฺตทุก}
\csromantocmarkref{paragraph}{1. กุสลตฺติก ...}{mtk.57}
\bookmark[dest={mtk.57},level=4]{1. กุสลตฺติก ...}
\csromantocmarknopage{subparagraph}{4. เหตุ สเหตุกทุก}{mtk.58}
\bookmark[dest={mtk.58},level=5]{4. เหตุ สเหตุกทุก}
\csromanmarkh{3. เหตุ\-สมฺปยุตฺตทุก}\csromanmarkhii{1. กุสลตฺติก}\csromanheadernomark{4}{3. เหตุ\-สมฺปยุตฺตทุก 1. กุสลตฺติก}
\prosewithcloser{\csromanfolio{331}(ยถา กุสลตฺติเก ปญฺหาวารสฺส อนุโลมมฺปิ ปจฺจนียมฺปิ อนุโลมปจฺจนียมฺปิ ปจฺจนียานุโลมมฺปิ คณิตํ, เอวํ คเณตพฺพํ.)}{สเหตุก\-ทุก\-กุสล\-ตฺติกํ นิฏฺฐิตํ.}
\csromanmarkh{3. เหตุ\-สมฺปยุตฺตทุก}\csromanmarkhii{1. กุสลตฺติก}\csromanheadernomark{2}{3. เหตุ\-สมฺปยุตฺตทุก 1. กุสลตฺติก}
\csromanmarkh{1. กุสลปท}\csromanmarkhii{1. ปฏิจฺจวาร}\csromanheadernomark{2}{1. \textbf{กุสลปท 1. ปฏิจฺจวาร}}
\prose{ปจฺจยจตุกฺกเหตุ}

\proseitem{25}{เหตุ\-สมฺปยุตฺตํ กุสลํ ธมฺมํ ปฏิจฺจ เหตุสมฺปยุตฺโต กุสโล ธมฺโม อุปฺปชฺชติ เหตุปจฺจยา. (1) (สํขิตฺตํ.)}

\proseitem{26}{เหตุยา เอกํ, อารมฺมเณ เอกํ, อธิปติยา เอกํ ฯเปฯ อวิคเต เอกํ. (สํขิตฺตํ.)}

\csromanheader{2}{\textbf{ปจฺจนีย}}
\proseitem{27}{นอธิปติยา เอกํ, นปุเรชาเต เอกํ, นปจฺฉาชาเต เอกํ, นอาเสวเน เอกํ, นกมฺเม เอกํ, นวิปาเก เอกํ, นวิปฺปยุตฺเต เอกํ. (สํขิตฺตํ.)}

\prose{เหตุปจฺจยา นอธิปติยา เอกํ. (สํขิตฺตํ.)}

\prose{นอธิปติ\-ปจฺจยา เหตุยา เอกํ. (สํขิตฺตํ.)}

\prose{(สหชาตวาโรปิ ปจฺจยวาโรปิ นิสฺสยวาโรปิ สํสฏฺฐ\-วาโรปิ สมฺปยุตฺตวาโรปิ ปฏิจฺจ\-วาร\-สทิสา วิตฺถาเรตพฺพา.)}

\csromanmarkh{1. กุสลปท}\csromanmarkhii{7. ปญฺหาวาร}\csromanheadernomark{2}{1. \textbf{กุสลปท 7. ปญฺหาวาร}}
\prose{ปจฺจยจตุกฺกเหตุ}

\proseitem{28}{เหตุสมฺปยุตฺโต กุสโล ธมฺโม เหตุ\-สมฺปยุตฺตสฺส กุสลสฺส ธมฺมสฺส เหตุปจฺจเยน ปจฺจโย. (สํขิตฺตํ.)}

\prose{เหตุยา เอกํ, อารมฺมเณ เอกํ, อธิปติยา เอกํ, อนนฺตเร เอกํ, สมนนฺตเร เอกํ ฯเปฯ กมฺเม เอกํ ฯเปฯ อวิคเต เอกํ. (สํขิตฺตํ.)}

\csromanheader{2}{ธมฺมานุโลม ทุกติก\-ปฏฺฐาน\-ปาฬิ ปจฺจนียุ\-ทฺธาร}
\proseitem{29}{\csromanfolio{332}เหตุสมฺปยุตฺโต กุสโล ธมฺโม เหตุ\-สมฺปยุตฺตสฺส กุสลสฺส ธมฺมสฺส อารมฺมณ\-ปจฺจเยน ปจฺจโย. สหชาต\-ปจฺจเยน ปจฺจโย. อุปนิสฺสย\-ปจฺจเยน ปจฺจโย.}

\proseitem{30}{นเหตุยา เอกํ, นอารมฺมเณ เอกํ. (สํขิตฺตํ.)}

\prose{เหตุปจฺจยา นอารมฺมเณ เอกํ. (สํขิตฺตํ.)}

\prose{นเหตุปจฺจยา อารมฺมเณ เอกํ. (สํขิตฺตํ.)}

\prose{(ยถา กุสลตฺติเก ปญฺหาวารสฺส อนุโลมมฺปิ ปจฺจนียมฺปิ อนุโลมปจฺจนียมฺปิ ปจฺจนียานุโลมมฺปิ คณิตํ, เอวํ คเณตพฺพํ.)}

\csromanmarkh{2. อกุสลปท}\csromanmarkhii{1. ปฏิจฺจวาร}\csromanheadernomark{2}{2. \textbf{อกุสลปท 1. ปฏิจฺจวาร}}
\prose{ปจฺจยจตุกฺกเหตุ อารมฺมณ}

\proseitem{31}{เหตุ\-สมฺปยุตฺตํ อกุสลํ ธมฺมํ ปฏิจฺจ เหตุสมฺปยุตฺโต อกุสโล ธมฺโม อุปฺปชฺชติ เหตุปจฺจยา. (1)}

\prose{เหตุ\-วิปฺปยุตฺตํ อกุสลํ ธมฺมํ ปฏิจฺจ เหตุสมฺปยุตฺโต อกุสโล ธมฺโม อุปฺปชฺชติ เหตุปจฺจยา. (1)}

\prose{เหตุ\-สมฺปยุตฺตํ อกุสลญฺจ เหตุ\-วิปฺปยุตฺตํ อกุสลญฺจ ธมฺมํ ปฏิจฺจ เหตุสมฺปยุตฺโต อกุสโล ธมฺโม อุปฺปชฺชติ เหตุปจฺจยา. (1)}

\prose{เหตุ\-สมฺปยุตฺตํ อกุสลํ ธมฺมํ ปฏิจฺจ เหตุสมฺปยุตฺโต อกุสโล ธมฺโม อุปฺปชฺชติ อารมฺมณ\-ปจฺจยา. ตีณิ.}

\prose{เหตุ\-วิปฺปยุตฺตํ อกุสลํ ธมฺมํ ปฏิจฺจ เหตุสมฺปยุตฺโต อกุสโล ธมฺโม อุปฺปชฺชติ อารมฺมณ\-ปจฺจยา. (1)}

\prose{เหตุ\-สมฺปยุตฺตํ อกุสลญฺจ เหตุ\-วิปฺปยุตฺตํ อกุสลญฺจ ธมฺมํ ปฏิจฺจ เหตุสมฺปยุตฺโต อกุสโล ธมฺโม อุปฺปชฺชติ อารมฺมณ\-ปจฺจยา. (1)}

\prose{เหตุ\-สมฺปยุตฺตํ อกุสลํ ธมฺมํ ปฏิจฺจ เหตุสมฺปยุตฺโต อกุสโล ธมฺโม อุปฺปชฺชติ อธิปติ\-ปจฺจยา. (สํขิตฺตํ.)}

\csromanmarkh{3. เหตุ\-สมฺปยุตฺตทุก}\csromanmarkhii{1. กุสลตฺติก}\csromanheadernomark{2}{3. เหตุ\-สมฺปยุตฺตทุก 1. กุสลตฺติก}
\proseitem{32}{\csromanfolio{333}เหตุยา ตีณิ, อารมฺมเณ ปญฺจ, อธิปติยา เอกํ, อนนฺตเร ปญฺจ, สมนนฺตเร ปญฺจ, สหชาเต ปญฺจ, อญฺญมญฺเญ ปญฺจ, นิสฺสเย ปญฺจ, อุปนิสฺสเย ปญฺจ, ปุเรชาเต ปญฺจ, อาเสวเน ปญฺจ, กมฺเม ปญฺจ, อาหาเร ปญฺจ, อินฺทฺริเย ปญฺจ, ฌาเน ปญฺจ, มคฺเค ปญฺจ, สมฺปยุตฺเต ปญฺจ, วิปฺปยุตฺเต ปญฺจ, อตฺถิยา ปญฺจ ฯเปฯ อวิคเต ปญฺจ. (สํขิตฺตํ.)}

\prose{นเหตุ นอธิปติ}

\proseitem{33}{เหตุ\-สมฺปยุตฺตํ อกุสลํ ธมฺมํ ปฏิจฺจ เหตุวิปฺปยุตฺโต อกุสโล ธมฺโม อุปฺปชฺชติ นเหตุปจฺจยา. (1)}

\prose{เหตุ\-สมฺปยุตฺตํ อกุสลํ ธมฺมํ ปฏิจฺจ เหตุสมฺปยุตฺโต อกุสโล ธมฺโม อุปฺปชฺชติ นอธิปติ\-ปจฺจยา. (สํขิตฺตํ.)}

\proseitem{34}{นเหตุยา เอกํ, นอธิปติยา ปญฺจ, นปุเรชาเต ปญฺจ, นปจฺฉาชาเต ปญฺจ, นอาเสวเน ปญฺจ, นกมฺเม ตีณิ, นวิปาเก ปญฺจ, นวิปฺปยุตฺเต ปญฺจ. (สํขิตฺตํ.)}

\prose{เหตุปจฺจยา นอธิปติยา ตีณิ. (สํขิตฺตํ.)}

\prose{นเหตุปจฺจยา อารมฺมเณ เอกํ. (สํขิตฺตํ.)}

\prose{(สหชาตวาโรปิ ปจฺจยวาโรปิ นิสฺสยวาโรปิ สํสฏฺฐ\-วาโรปิ สมฺปยุตฺตวาโรปิ ปฏิจฺจ\-วาร\-สทิสา วิตฺถาเรตพฺพา.)}

\csromanmarkh{2. อกุสลปท}\csromanmarkhii{7. ปญฺหาวาร}\csromanheadernomark{2}{2. \textbf{อกุสลปท 7. ปญฺหาวาร}}
\prose{ปจฺจยจตุกฺกเหตุ อารมฺมณ}

\proseitem{35}{เหตุสมฺปยุตฺโต อกุสโล ธมฺโม เหตุ\-สมฺปยุตฺตสฺส อกุสลสฺส ธมฺมสฺส เหตุปจฺจเยน ปจฺจโย. เทฺว.}

\prose{เหตุสมฺปยุตฺโต อกุสโล ธมฺโม เหตุ\-สมฺปยุตฺตสฺส อกุสลสฺส ธมฺมสฺส อารมฺมณ\-ปจฺจเยน ปจฺจโย. ตีณิ.}

\prose{เหตุวิปฺปยุตฺโต อกุสโล ธมฺโม เหตุ\-วิปฺปยุตฺตสฺส อกุสลสฺส ธมฺมสฺส อารมฺมณ\-ปจฺจเยน ปจฺจโย. ตีณิ.}

\prose{\csromanfolio{334}เหตุสมฺปยุตฺโต อกุสโล จ เหตุวิปฺปยุตฺโต อกุสโล จ ธมฺมา เหตุ\-สมฺปยุตฺตสฺส อกุสลสฺส ธมฺมสฺส อารมฺมณ\-ปจฺจเยน ปจฺจโย. ตีณิ. (สํขิตฺตํ.)}

\proseitem{36}{เหตุยา เทฺว, อารมฺมเณ นว, อธิปติยา เอกํ, อนนฺตเร นว, สมนนฺตเร นว, สหชาเต ปญฺจ, อญฺญมญฺเญ ปญฺจ, นิสฺสเย ปญฺจ, อุปนิสฺสเย นว, อาเสวเน นว, กมฺเม ตีณิ, อาหาเร ตีณิ, อินฺทฺริเย ตีณิ, ฌาเน ตีณิ, มคฺเค ตีณิ, สมฺปยุตฺเต ปญฺจ, อตฺถิยา ปญฺจ ฯเปฯ อวิคเต ปญฺจ. (สํขิตฺตํ.)}

\proseitem{37}{เหตุสมฺปยุตฺโต อกุสโล ธมฺโม เหตุ\-สมฺปยุตฺตสฺส อกุสลสฺส ธมฺมสฺส อารมฺมณ\-ปจฺจเยน ปจฺจโย. สหชาต\-ปจฺจเยน ปจฺจโย. อุปนิสฺสย\-ปจฺจเยน ปจฺจโย. (สํขิตฺตํ.)}

\proseitem{38}{นเหตุยา นว, นอารมฺมเณ นว. (สํขิตฺตํ.)}

\prose{เหตุปจฺจยา นอารมฺมเณ เทฺว. (สํขิตฺตํ.)}

\prose{นเหตุปจฺจยา อารมฺมเณ นว. (สํขิตฺตํ.)}

\prose{(ยถา กุสลตฺติเก ปญฺหาวารสฺส อนุโลมมฺปิ ปจฺจนียมฺปิ อนุโลมปจฺจนียมฺปิ ปจฺจนียานุโลมมฺปิ คณิตํ, เอวํ คเณตพฺพํ.)}

\csromanmarkh{3. อพฺยากตปท}\csromanmarkhii{1. ปฏิจฺจวาร}\csromanheadernomark{2}{3. \textbf{อพฺยากตปท 1. ปฏิจฺจวาร}}
\prose{ปจฺจยจตุกฺกเหตุ}

\proseitem{39}{เหตุ\-สมฺปยุตฺตํ อพฺยากตํ ธมฺมํ ปฏิจฺจ เหตุสมฺปยุตฺโต อพฺยากโต ธมฺโม อุปฺปชฺชติ เหตุปจฺจยา. ตีณิ.}

\prose{เหตุ\-วิปฺปยุตฺตํ อพฺยากตํ ธมฺมํ ปฏิจฺจ เหตุวิปฺปยุตฺโต อพฺยากโต ธมฺโม อุปฺปชฺชติ เหตุปจฺจยา. ตีณิ.}

\prose{เหตุ\-สมฺปยุตฺตํ อพฺยากตญฺจ เหตุ\-วิปฺปยุตฺตํ อพฺยากตญฺจ ธมฺมํ ปฏิจฺจ เหตุสมฺปยุตฺโต อพฺยากโต ธมฺโม อุปฺปชฺชติ เหตุปจฺจยา. ตีณิ. (สํขิตฺตํ.)}

\csromanmarkh{3. เหตุ\-สมฺปยุตฺตทุก}\csromanmarkhii{1. กุสลตฺติก}\csromanheadernomark{2}{3. เหตุ\-สมฺปยุตฺตทุก 1. กุสลตฺติก}
\proseitem{40}{\csromanfolio{335}เหตุยา นว, อารมฺมเณ จตฺตาริ, อธิปติยา ปญฺจ, อนนฺตเร จตฺตาริ, สมนนฺตเร จตฺตาริ, สหชาเต นว, อญฺญมญฺเญ ฉ ฯเปฯ ปุเรชาเต เทฺว, อาเสวเน เทฺว, กมฺเม นว, วิปาเก นว ฯเปฯ อวิคเต นว. (สํขิตฺตํ.)}

\prose{นเหตุ นอารมฺมณ}

\proseitem{41}{เหตุ\-วิปฺปยุตฺตํ อพฺยากตํ ธมฺมํ ปฏิจฺจ เหตุวิปฺปยุตฺโต อพฺยากโต ธมฺโม อุปฺปชฺชติ นเหตุปจฺจยา. (1)}

\prose{เหตุ\-สมฺปยุตฺตํ อพฺยากตํ ธมฺมํ ปฏิจฺจ เหตุวิปฺปยุตฺโต อพฺยากโต ธมฺโม อุปฺปชฺชติ นอารมฺมณ\-ปจฺจยา. (สํขิตฺตํ.)}

\proseitem{42}{นเหตุยา เอกํ, นอารมฺมเณ ตีณิ, นอธิปติยา นว, นอนนฺตเร ตีณิ, นสมนนฺตเร ตีณิ, นอญฺญมญฺเญ ตีณิ, นอุปนิสฺสเย ตีณิ, นปุเรชาเต นว, นปจฺฉาชาเต นว, นอาเสวเน นว, นกมฺเม เทฺว, นวิปาเก ปญฺจ, นอาหาเร เอกํ, นอินฺทฺริเย เอกํ, นฌาเน เอกํ, นมคฺเค เอกํ, นสมฺปยุตฺเต ตีณิ, นวิปฺปยุตฺเต เทฺว, โนนตฺถิยา ตีณิ, โนวิคเต ตีณิ. (สํขิตฺตํ.)}

\prose{เหตุปจฺจยา นอารมฺมเณ ตีณิ. (สํขิตฺตํ.)}

\prose{นเหตุปจฺจยา อารมฺมเณ เอกํ. (สํขิตฺตํ.)}

\prose{(สหชาตวาโรปิ ปจฺจยวาโรปิ นิสฺสยวาโรปิ สํสฏฺฐ\-วาโรปิ สมฺปยุตฺตวาโรปิ ปฏิจฺจ\-วาร\-สทิสา วิตฺถาเรตพฺพา.)}

\csromanmarkh{3. อพฺยากตปท}\csromanmarkhii{7. ปญฺหาวาร}\csromanheadernomark{2}{3. \textbf{อพฺยากตปท 7. ปญฺหาวาร}}
\prose{ปจฺจยจตุกฺกเหตุ อารมฺมณ}

\proseitem{43}{เหตุสมฺปยุตฺโต อพฺยากโต ธมฺโม เหตุ\-สมฺปยุตฺตสฺส อพฺยากตสฺส ธมฺมสฺส เหตุปจฺจเยน ปจฺจโย. ตีณิ.}

\prose{เหตุสมฺปยุตฺโต อพฺยากโต ธมฺโม เหตุ\-สมฺปยุตฺตสฺส อพฺยากตสฺส ธมฺมสฺส อารมฺมณ\-ปจฺจเยน ปจฺจโย. (สํขิตฺตํ.)}

\proseitem{44}{\csromanfolio{336}เหตุยา ตีณิ, อารมฺมเณ จตฺตาริ, อธิปติยา จตฺตาริ, อนนฺตเร จตฺตาริ, กมฺเม จตฺตาริ, วิปาเก จตฺตาริ ฯเปฯ อวิคเต สตฺต. (สํขิตฺตํ.)}

\proseitem{45}{เหตุสมฺปยุตฺโต อพฺยากโต ธมฺโม เหตุ\-สมฺปยุตฺตสฺส อพฺยากตสฺส ธมฺมสฺส อารมฺมณ\-ปจฺจเยน ปจฺจโย. สหชาต\-ปจฺจเยน ปจฺจโย. อุปนิสฺสย\-ปจฺจเยน ปจฺจโย.}

\proseitem{46}{นเหตุยา สตฺต, นอารมฺมเณ สตฺต. (สํขิตฺตํ.)}

\prose{เหตุปจฺจยา นอารมฺมเณ ตีณิ. (สํขิตฺตํ.)}

\prose{นเหตุปจฺจยา อารมฺมเณ จตฺตาริ. (สํขิตฺตํ.)}

\prosewithcloser{(ยถา กุสลตฺติเก ปญฺหาวารสฺส อนุโลมมฺปิ ปจฺจนียมฺปิ อนุโลมปจฺจนียมฺปิ ปจฺจนียานุโลมมฺปิ คณิตํ, เอวํ คเณตพฺพํ.)}{เหตุ\-สมฺปยุตฺต\-ทุก\-กุสล\-ตฺติกํ นิฏฺฐิตํ.}
\csromanmatikaanchor{mtk.59}
\csromantocmarkref{subparagraph}{1. กุสลตฺติก ...}{mtk.59}
\bookmark[dest={mtk.59},level=5]{1. กุสลตฺติก ...}
\csromanmarkh{4. เหตุ\-สเหตุกทุก}\csromanmarkhii{1. กุสลตฺติก}\csromanheadernomark{6}{4. เหตุ\-สเหตุกทุก 1. กุสลตฺติก}
\csromanmarkh{1. กุสลปท}\csromanmarkhii{1. ปฏิจฺจวาร}\csromanheadernomark{2}{1. \textbf{กุสลปท 1. ปฏิจฺจวาร}}
\prose{ปจฺจยจตุกฺกเหตุ}

\proseitem{47}{เหตุญฺเจว สเหตุกญฺจ กุสลํ ธมฺมํ ปฏิจฺจ เหตุ เจว สเหตุโก จ กุสโล ธมฺโม อุปฺปชฺชติ เหตุปจฺจยา. ตีณิ.}

\prose{สเหตุกญฺเจว น จ เหตุํ กุสลํ ธมฺมํ ปฏิจฺจ สเหตุโก เจว น จ เหตุ กุสโล ธมฺโม อุปฺปชฺชติ เหตุปจฺจยา. ตีณิ.}

\prose{เหตุญฺเจว สเหตุกญฺจ กุสลญฺจ สเหตุกญฺเจว น จ เหตุํ กุสลญฺจ ธมฺมํ ปฏิจฺจ เหตุ เจว สเหตุโก จ กุสโล ธมฺโม อุปฺปชฺชติ เหตุปจฺจยา. ตีณิ. (สํขิตฺตํ.)}

\csromanmarkh{4. เหตุ\-สเหตุกทุก}\csromanmarkhii{1. กุสลตฺติก}\csromanheadernomark{2}{4. เหตุ\-สเหตุกทุก 1. กุสลตฺติก}
\proseitem{48}{\csromanfolio{337}เหตุยา นว, อารมฺมเณ นว ฯเปฯ กมฺเม นว, อาหาเร นว ฯเปฯ อวิคเต นว. (สํขิตฺตํ.)}

\prose{นอธิปติ}

\proseitem{49}{เหตุญฺเจว สเหตุกญฺจ กุสลํ ธมฺมํ ปฏิจฺจ เหตุ เจว สเหตุโก จ กุสโล ธมฺโม อุปฺปชฺชติ นอธิปติ\-ปจฺจยา. (สํขิตฺตํ.)}

\proseitem{50}{นอธิปติยา นว, นปุเรชาเต นว, นปจฺฉาชาเต นว, นอาเสวเน นว, นกมฺเม ตีณิ, นวิปาเก นว, นวิปฺปยุตฺเต นว. (สํขิตฺตํ.)}

\prose{เหตุปจฺจยา นอธิปติยา นว. (สํขิตฺตํ.)}

\prose{นอธิปติ\-ปจฺจยา เหตุยา นว. (สํขิตฺตํ.)}

\prose{(สหชาตวาโรปิ ปจฺจยวาโรปิ นิสฺสยวาโรปิ สํสฏฺฐ\-วาโรปิ สมฺปยุตฺตวาโรปิ ปฏิจฺจ\-วาร\-สทิสา วิตฺถาเรตพฺพา.)}

\csromanmarkh{1. กุสลปท}\csromanmarkhii{7. ปญฺหาวาร}\csromanheadernomark{2}{1. \textbf{กุสลปท 7. ปญฺหาวาร}}
\prose{ปจฺจยจตุกฺกเหตุ}

\proseitem{51}{เหตุ เจว สเหตุโก จ กุสโล ธมฺโม เหตุสฺส เจว สเหตุกสฺส จ กุสลสฺส ธมฺมสฺส เหตุปจฺจเยน ปจฺจโย. (สํขิตฺตํ.)}

\proseitem{52}{เหตุยา ตีณิ, อารมฺมเณ นว, อธิปติยา นว, อนนฺตเร นว, สมนนฺตเร นว, สหชาเต นว, อญฺญมญฺเญ นว, นิสฺสเย นว, อุปนิสฺสเย นว, อาเสวเน นว, กมฺเม ตีณิ, อาหาเร ตีณิ ฯเปฯ อวิคเต นว. (สํขิตฺตํ.)}

\proseitem{53}{เหตุ เจว สเหตุโก จ กุสโล ธมฺโม เหตุสฺส เจว สเหตุกสฺส จ กุสลสฺส ธมฺมสฺส อารมฺมณ\-ปจฺจเยน ปจฺจโย. สหชาต\-ปจฺจเยน ปจฺจโย. อุปนิสฺสย\-ปจฺจเยน ปจฺจโย. (สํขิตฺตํ.)}

\proseitem{54}{\csromanfolio{338}นเหตุยา นว, นอารมฺมเณ นว. (สํขิตฺตํ.)}

\prose{เหตุปจฺจยา นอารมฺมเณ ตีณิ. (สํขิตฺตํ.)}

\prose{นเหตุปจฺจยา อารมฺมเณ นว. (สํขิตฺตํ.)}

\prose{(ยถา กุสลตฺติเก ปญฺหาวารสฺส อนุโลมมฺปิ ปจฺจนียมฺปิ อนุโลมปจฺจนียมฺปิ ปจฺจนียานุโลมมฺปิ คณิตํ, เอวํ คเณตพฺพํ.)}

\csromanmarkh{2. อกุสลปท}\csromanmarkhii{1. ปฏิจฺจวาร}\csromanheadernomark{2}{2. \textbf{อกุสลปท 1. ปฏิจฺจวาร}}
\prose{ปจฺจยจตุกฺกเหตุ}

\proseitem{55}{เหตุญฺเจว สเหตุกญฺจ อกุสลํ ธมฺมํ ปฏิจฺจ เหตุ เจว สเหตุโก จ อกุสโล ธมฺโม อุปฺปชฺชติ เหตุปจฺจยา. (สํขิตฺตํ.)}

\proseitem{56}{เหตุยา นว, อารมฺมเณ นว ฯเปฯ กมฺเม นว, อาหาเร นว ฯเปฯ อวิคเต นว. (สํขิตฺตํ.)}

\prose{นอธิปติ}

\proseitem{57}{เหตุญฺเจว สเหตุกญฺจ อกุสลํ ธมฺมํ ปฏิจฺจ เหตุ เจว สเหตุโก จ อกุสโล ธมฺโม อุปฺปชฺชติ นอธิปติ\-ปจฺจยา. นว. (สํขิตฺตํ.)}

\proseitem{58}{นอธิปติยา นว, นปุเรชาเต นว, นปจฺฉาชาเต นว, นอาเสวเน นว, นกมฺเม ตีณิ, นวิปาเก นว, นวิปฺปยุตฺเต นว. (สํขิตฺตํ.)}

\prose{เหตุปจฺจยา นอธิปติยา นว. (สํขิตฺตํ.)}

\prose{นอธิปติ\-ปจฺจยา เหตุยา นว. (สํขิตฺตํ.)}

\prose{(สหชาตวาโรปิ ปจฺจยวาโรปิ นิสฺสยวาโรปิ สํสฏฺฐ\-วาโรปิ สมฺปยุตฺตวาโรปิ ปฏิจฺจ\-วาร\-สทิสา วิตฺถาเรตพฺพา.)}

\csromanmarkh{2. อกุสลปท}\csromanmarkhii{7. ปญฺหาวาร}\csromanheadernomark{2}{2. \textbf{อกุสลปท 7. ปญฺหาวาร}}
\prose{ปจฺจยจตุกฺกเหตุ}

\proseitem{59}{เหตุ เจว สเหตุโก จ อกุสโล ธมฺโม เหตุสฺส เจว สเหตุกสฺส จ อกุสลสฺส ธมฺมสฺส เหตุปจฺจเยน ปจฺจโย. (สํขิตฺตํ.)}

\csromanmarkh{4. เหตุ\-สเหตุกทุก}\csromanmarkhii{1. กุสลตฺติก}\csromanheadernomark{2}{4. เหตุ\-สเหตุกทุก 1. กุสลตฺติก}
\proseitem{60}{\csromanfolio{339}เหตุยา ตีณิ, อารมฺมเณ นว ฯเปฯ อาเสวเน นว, กมฺเม ตีณิ, อาหาเร ตีณิ, อินฺทฺริเย ฌาเน มคฺเค ตีณิ, สมฺปยุตฺเต นว ฯเปฯ อวิคเต นว. (สํขิตฺตํ.)}

\proseitem{61}{เหตุ เจว สเหตุโก จ อกุสโล ธมฺโม เหตุสฺส เจว สเหตุกสฺส จ อกุสลสฺส ธมฺมสฺส อารมฺมณ\-ปจฺจเยน ปจฺจโย. สหชาต\-ปจฺจเยน ปจฺจโย. อุปนิสฺสย\-ปจฺจเยน ปจฺจโย.}

\proseitem{62}{นเหตุยา นว, นอารมฺมเณ นว ฯเปฯ โนอวิคเต นว. (สํขิตฺตํ.)}

\prose{เหตุปจฺจยา นอารมฺมเณ ตีณิ. (สํขิตฺตํ.)}

\prose{นเหตุปจฺจยา อารมฺมเณ นว. (สํขิตฺตํ.)}

\prose{(ยถา กุสลตฺติเก ปญฺหาวารสฺส อนุโลมมฺปิ ปจฺจนียมฺปิ อนุโลมปจฺจนียมฺปิ ปจฺจนียานุโลมมฺปิ คณิตํ, เอวํ คเณตพฺพํ.)}

\csromanmarkh{3. อพฺยากตปท}\csromanmarkhii{1. ปฏิจฺจวาร}\csromanheadernomark{2}{3. \textbf{อพฺยากตปท 1. ปฏิจฺจวาร}}
\prose{ปจฺจยจตุกฺกเหตุ}

\proseitem{63}{เหตุญฺเจว สเหตุกญฺจ อพฺยากตํ ธมฺมํ ปฏิจฺจ เหตุ เจว สเหตุโก จ อพฺยากโต ธมฺโม อุปฺปชฺชติ เหตุปจฺจยา. (สํขิตฺตํ.)}

\proseitem{64}{เหตุยา นว, อารมฺมเณ นว ฯเปฯ กมฺเม นว, วิปาเก นว ฯเปฯ อวิคเต นว. (สํขิตฺตํ.)}

\prose{นอธิปติ}

\proseitem{65}{เหตุญฺเจว สเหตุกญฺจ อพฺยากตํ ธมฺมํ ปฏิจฺจ เหตุ เจว สเหตุโก จ อพฺยากโต ธมฺโม อุปฺปชฺชติ นอธิปติ\-ปจฺจยา. (สํขิตฺตํ.)}

\proseitem{66}{นอธิปติยา นว, นปุเรชาเต นว, นปจฺฉาชาเต นว, นอาเสวเน นว, นกมฺเม ตีณิ, นวิปาเก นว, นวิปฺปยุตฺเต นว. (สํขิตฺตํ.)}

\prose{เหตุปจฺจยา นอธิปติยา นว. (สํขิตฺตํ.)}

\prose{นอธิปติ\-ปจฺจยา เหตุยา นว. (สํขิตฺตํ.)}

\prose{(สหชาตวาโรปิ ปจฺจยวาโรปิ นิสฺสยวาโรปิ สํสฏฺฐ\-วาโรปิ สมฺปยุตฺตวาโรปิ ปฏิจฺจ\-วาร\-สทิสา วิตฺถาเรตพฺพา.)}

\csromanmarkh{ธมฺมานุโลม ทุกติก\-ปฏฺฐาน\-ปาฬิ}\csromanmarkhii{3. อพฺยากตปท 7. ปญฺหาวาร}\csromanheadernomark{2}{ธมฺมานุโลม ทุกติก\-ปฏฺฐาน\-ปาฬิ 3. อพฺยากตปท 7. ปญฺหาวาร}
\prose{\csromanfolio{340}ปจฺจยจตุกฺกเหตุ}

\proseitem{67}{เหตุ เจว สเหตุโก จ อพฺยากโต ธมฺโม เหตุสฺส เจว สเหตุกสฺส จ อพฺยากตสฺส ธมฺมสฺส เหตุปจฺจเยน ปจฺจโย. (สํขิตฺตํ.)}

\proseitem{68}{เหตุยา ตีณิ, อารมฺมเณ นว, อธิปติยา นว, อนนฺตเร นว, สมนนฺตเร นว, สหชาเต นว, อญฺญมญฺเญ นว, นิสฺสเย นว, อุปนิสฺสเย นว, อาเสวเน นว, กมฺเม ตีณิ, วิปาเก นว, อาหาเร ตีณิ ฯเปฯ อวิคเต นว. (สํขิตฺตํ.)}

\proseitem{69}{เหตุ เจว สเหตุโก จ อพฺยากโต ธมฺโม เหตุสฺส เจว สเหตุกสฺส จ อพฺยากตสฺส ธมฺมสฺส อารมฺมณ\-ปจฺจเยน ปจฺจโย. สหชาต\-ปจฺจเยน ปจฺจโย. อุปนิสฺสย\-ปจฺจเยน ปจฺจโย. (สํขิตฺตํ.)}

\proseitem{70}{นเหตุยา นว, นอารมฺมเณ นว. (สํขิตฺตํ.)}

\prose{เหตุปจฺจยา นอารมฺมเณ ตีณิ. (สํขิตฺตํ.)}

\prose{นเหตุปจฺจยา อารมฺมเณ นว. (สํขิตฺตํ.)}

\prosewithcloser{(ยถา กุสลตฺติเก ปญฺหาวารสฺส อนุโลมมฺปิ ปจฺจนียมฺปิ อนุโลมปจฺจนียมฺปิ ปจฺจนียานุโลมมฺปิ คณิตํ, เอวํ คเณตพฺพํ.)}{เหตุ\-สเหตุก\-ทุก\-กุสล\-ตฺติกํ นิฏฺฐิตํ.}
\csromanmatikaanchor{mtk.60}
\csromantocmarkref{subparagraph}{5. เหตุ เหตุสมฺปยุตฺตทุก 1. กุสลตฺติก ...}{mtk.60}
\bookmark[dest={mtk.60},level=5]{5. เหตุ เหตุสมฺปยุตฺตทุก 1. กุสลตฺติก ...}
\csromanmarkh{5. เหตุ\-เหตุ\-สมฺปยุตฺตทุก}\csromanmarkhii{1. กุสลตฺติก}\csromanheadernomark{6}{5. เหตุ\-เหตุ\-สมฺปยุตฺตทุก 1. กุสลตฺติก}
\csromanmarkh{1. กุสลปท}\csromanmarkhii{1. ปฏิจฺจวาร}\csromanheadernomark{2}{1. \textbf{กุสลปท 1. ปฏิจฺจวาร}}
\prose{ปจฺจยจตุกฺกเหตุ}

\proseitem{71}{เหตุญฺเจว เหตุ\-สมฺปยุตฺตญฺจ กุสลํ ธมฺมํ ปฏิจฺจ เหตุ เจว เหตุสมฺปยุตฺโต จ กุสโล ธมฺโม อุปฺปชฺชติ เหตุปจฺจยา.  }

\vspace{\tipitakaparskip}
\csromancenter{เหตุสเหตุกสมฺปยุตฺตทุก 1. กุสลตฺติก เหตุญฺเจว เหตุ\-สมฺปยุตฺตญฺจ กุสลํ ธมฺมํ ปฏิจฺจ เหตุสมฺปยุตฺโต เจว น จ เหตุ กุสโล ธมฺโม อุปฺปชฺชติ เหตุปจฺจยา. (สํขิตฺตํ.)}
\proseitem{72}{\csromanfolio{341}เหตุยา นว, อารมฺมเณ นว ฯเปฯ อวิคเต นว. (สํขิตฺตํ.)}

\prose{นอธิปติ}

\proseitem{73}{เหตุญฺเจว เหตุ\-สมฺปยุตฺตญฺจ กุสลํ ธมฺมํ ปฏิจฺจ เหตุ เจว เหตุสมฺปยุตฺโต จ กุสโล ธมฺโม อุปฺปชฺชติ นอธิปติ\-ปจฺจยา. (สํขิตฺตํ.)}

\proseitem{74}{นอธิปติยา นว, นปุเรชาเต นว, นปจฺฉาชาเต นว, นอาเสวเน นว, นกมฺเม ตีณิ, นวิปาเก นว, นวิปฺปยุตฺเต นว. (สํขิตฺตํ.)}

\prose{เหตุปจฺจยา นอธิปติยา นว. (สํขิตฺตํ.)}

\prose{อธิปติ\-ปจฺจยา เหตุยา นว. (สํขิตฺตํ.)}

\prose{(สหชาตวาโรปิ ปจฺจยวาโรปิ นิสฺสยวาโรปิ สํสฏฺฐ\-วาโรปิ สมฺปยุตฺตวาโรปิ ปฏิจฺจ\-วาร\-สทิสา วิตฺถาเรตพฺพา.)}

\csromanmarkh{1. กุสลปท}\csromanmarkhii{7. ปญฺหาวาร}\csromanheadernomark{2}{1. \textbf{กุสลปท 7. ปญฺหาวาร}}
\prose{ปจฺจยจตุกฺกเหตุ}

\proseitem{75}{เหตุ เจว เหตุสมฺปยุตฺโต จ กุสโล ธมฺโม เหตุสฺส เจว เหตุ\-สมฺปยุตฺตสฺส จ กุสลสฺส ธมฺมสฺส เหตุปจฺจเยน ปจฺจโย. (สํขิตฺตํ.)}

\proseitem{76}{เหตุยา ตีณิ, อารมฺมเณ นว, อธิปติยา นว, อนนฺตเร นว ฯเปฯ อาเสวเน นว, กมฺเม ตีณิ, อาหาเร ตีณิ ฯเปฯ อวิคเต นว. (สํขิตฺตํ.)}

\proseitem{77}{เหตุ เจว เหตุสมฺปยุตฺโต จ กุสโล ธมฺโม เหตุสฺส เจว เหตุ\-สมฺปยุตฺตสฺส จ กุสลสฺส ธมฺมสฺส อารมฺมณ\-ปจฺจเยน ปจฺจโย. สหชาต\-ปจฺจเยน ปจฺจโย. อุปนิสฺสย\-ปจฺจเยน ปจฺจโย. (สํขิตฺตํ.)}

\proseitem{78}{\csromanfolio{342}นเหตุยา นว, นอารมฺมเณ นว. (สํขิตฺตํ.)}

\prose{เหตุปจฺจยา นอารมฺมเณ ตีณิ. (สํขิตฺตํ.)}

\prose{นเหตุปจฺจยา อารมฺมเณ นว. (สํขิตฺตํ.)}

\prose{(ยถา กุสลตฺติเก ปญฺหาวารสฺส อนุโลมมฺปิ ปจฺจนียมฺปิ อนุโลมปจฺจนียมฺปิ ปจฺจนียานุโลมมฺปิ คณิตํ, เอวํ คเณตพฺพํ.)}

\csromanmarkh{2. อกุสลปท}\csromanmarkhii{1. ปฏิจฺจวาร}\csromanheadernomark{2}{2. \textbf{อกุสลปท 1. ปฏิจฺจวาร}}
\prose{ปจฺจยจตุกฺกเหตุ}

\proseitem{79}{เหตุญฺเจว เหตุ\-สมฺปยุตฺตญฺจ อกุสลํ ธมฺมํ ปฏิจฺจ เหตุ เจว เหตุสมฺปยุตฺโต จ อกุสโล ธมฺโม อุปฺปชฺชติ เหตุปจฺจยา. (สํขิตฺตํ.)}

\proseitem{80}{เหตุยา นว, อารมฺมเณ นว ฯเปฯ อวิคเต นว. (สํขิตฺตํ.)}

\prose{นอธิปติ}

\proseitem{81}{เหตุญฺเจว เหตุ\-สมฺปยุตฺตญฺจ อกุสลํ ธมฺมํ ปฏิจฺจ เหตุ เจว เหตุสมฺปยุตฺโต จ อกุสโล ธมฺโม อุปฺปชฺชติ นอธิปติ\-ปจฺจยา. (สํขิตฺตํ.)}

\proseitem{82}{นอธิปติยา นว, นปุเรชาเต นว, นปจฺฉาชาเต นว, นอาเสวเน นว, นกมฺเม ตีณิ, นวิปาเก นว, นวิปฺปยุตฺเต นว. (สํขิตฺตํ.)}

\prose{เหตุปจฺจยา นอธิปติยา ตีณิ. (สํขิตฺตํ.)}

\prose{นอธิปติ\-ปจฺจยา เหตุยา นว. (สํขิตฺตํ.)}

\prose{(สหชาตวาโรปิ ปฺจฺจยวาโรปิ นิสฺสยวาโรปิ สํสฏฺฐ\-วาโรปิ สมฺปยุตฺตวาโรปิ ปฏิจฺจ\-วาร\-สทิสา วิตฺถาเรตพฺพา.)}

\csromanmarkh{2. อกุสลปท}\csromanmarkhii{7. ปญฺหาวาร}\csromanheadernomark{2}{2. \textbf{อกุสลปท 7. ปญฺหาวาร}}
\prose{ปจฺจยจตุกฺกเหตุ}

\proseitem{83}{เหตุ เจว เหตุสมฺปยุตฺโต จ อกุสโล ธมฺโม เหตุสฺส เจว เหตุ\-สมฺปยุตฺตสฺส จ อกุสลสฺส ธมฺมสฺส เหตุปจฺจเยน ปจฺจโย. (สํขิตฺตํ.)}

\csromanmarkh{5. เหตุสเหตุกสมฺปยุตฺตทุก}\csromanmarkhii{1. กุสลตฺติก}\csromanheadernomark{2}{5. เหตุสเหตุกสมฺปยุตฺตทุก 1. กุสลตฺติก}
\prose{\csromanfolio{343}เหตุยา ตีณิ, อารมฺมเณ นว, อธิปติยา นว ฯเปฯ กมฺเม ตีณิ, อาหาเร ตีณิ ฯเปฯ อวิคเต นว. (สํขิตฺตํ.)}

\proseitem{84}{เหตุ เจว เหตุสมฺปยุตฺโต จ อกุสโล ธมฺโม เหตุสฺส เจว เหตุ\-สมฺปยุตฺตสฺส จ อกุสลสฺส ธมฺมสฺส อารมฺมณ\-ปจฺจเยน ปจฺจโย. สหชาต\-ปจฺจเยน ปจฺจโย. อุปนิสฺสย\-ปจฺจเยน ปจฺจโย. (สํขิตฺตํ.)}

\proseitem{85}{นเหตุยา นว, นอารมฺมเณ นว. (สํขิตฺตํ.)}

\prose{เหตุปจฺจยา นอารมฺมเณ ตีณิ. (สํขิตฺตํ.)}

\prose{นเหตุปจฺจยา อารมฺมเณ นว. (สํขิตฺตํ.)}

\prose{(ยถา กุสลตฺติเก ปญฺหาวารสฺส อนุโลมมฺปิ ปจฺจนียมฺปิ อนุโลมปจฺจนียมฺปิ ปจฺจนียานุโลมมฺปิ คณิตํ, เอวํ คเณตพฺพํ.)}

\csromanmarkh{3. อพฺยากตปท}\csromanmarkhii{1. ปฏิจฺจวาร}\csromanheadernomark{2}{3. \textbf{อพฺยากตปท 1. ปฏิจฺจวาร}}
\prose{ปจฺจยจตุกฺกเหตุ}

\proseitem{86}{เหตุญฺเจว เหตุ\-สมฺปยุตฺตญฺจ อพฺยากตํ ธมฺมํ ปฏิจฺจ เหตุ เจว เหตุสมฺปยุตฺโต จ อพฺยากโต ธมฺโม อุปฺปชฺชติ เหตุปจฺจยา. (สํขิตฺตํ.)}

\prose{เหตุยา นว, อารมฺมเณ นว ฯเปฯ กมฺเม นว, วิปาเก นว ฯเปฯ อวิคเต นว. (สํขิตฺตํ.)}

\prose{นอธิปติ}

\proseitem{87}{เหตุญฺเจว เหตุ\-สมฺปยุตฺตญฺจ อพฺยากตํ ธมฺมํ ปฏิจฺจ เหตุ เจว เหตุสมฺปยุตฺโต จ อพฺยากโต ธมฺโม อุปฺปชฺชติ นอธิปติ\-ปจฺจยา. (สํขิตฺตํ.)}

\proseitem{88}{นอธิปติยา นว, นปุเรชาเต นว, นปจฺฉาชาเต นว, นอาเสวเน นว, นกมฺเม ตีณิ, นวิปาเก นว, นวิปฺปยุตฺเต นว. (สํขิตฺตํ.)}

\prose{\csromanfolio{344}เหตุปจฺจยา นอธิปติยา นว. (สํขิตฺตํ.)}

\prose{นอธิปติ\-ปจฺจยา เหตุยา นว. (สํขิตฺตํ.)}

\prose{(สหชาตวาโรปิ ปจฺจยวาโรปิ นิสฺสยวาโรปิ สํสฏฺฐ\-วาโรปิ สมฺปยุตฺตวาโรปิ ปฏิจฺจ\-วาร\-สทิสา วิตฺถาเรตพฺพา.)}

\csromanmarkh{3. อพฺยากตปท}\csromanmarkhii{7. ปญฺหาวาร}\csromanheadernomark{2}{3. \textbf{อพฺยากตปท 7. ปญฺหาวาร}}
\prose{ปจฺจยจตุกฺกเหตุ}

\proseitem{89}{เหตุ เจว เหตุสมฺปยุตฺโต จ อพฺยากโต ธมฺโม เหตุสฺส เจว เหตุ\-สมฺปยุตฺตสฺส จ อพฺยากตสฺส ธมฺมสฺส เหตุปจฺจเยน ปจฺจโย. (สํขิตฺตํ.)}

\proseitem{90}{เหตุยา ตีณิ, อารมฺมเณ นว, อธิปติยา นว ฯเปฯ อาเสวเน นว, กมฺเม ตีณิ, วิปาเก นว, อาหาเร ตีณิ, อินฺทฺริเย นว, ฌาเน ตีณิ, มคฺเค สมฺปยุตฺเต นว ฯเปฯ อวิคเต นว. (สํขิตฺตํ.)}

\proseitem{91}{เหตุ เจว เหตุสมฺปยุตฺโต จ อพฺยากโต ธมฺโม เหตุสฺส เจว เหตุ\-สมฺปยุตฺตสฺส จ อพฺยากตสฺส ธมฺมสฺส อารมฺมณ\-ปจฺจเยน ปจฺจโย. สหชาต\-ปจฺจเยน ปจฺจโย. อุปนิสฺสย\-ปจฺจเยน ปจฺจโย. (สํขิตฺตํ.)}

\proseitem{92}{นเหตุยา นว, นอารมฺมเณ นว. (สํขิตฺตํ.)}

\prose{เหตุปจฺจยา นอารมฺมเณ ตีณิ. (สํขิตฺตํ.)}

\prose{นเหตุปจฺจยา อารมฺมเณ นว. (สํขิตฺตํ.)}

\prosewithcloser{(ยถา กุสลตฺติเก ปญฺหาวารสฺส อนุโลมมฺปิ ปจฺจนียมฺปิ อนุโลมปจฺจนียมฺปิ ปจฺจนียานุโลมมฺปิ คณิตํ, เอวํ คเณตพฺพํ.)}{เหตุ\-เหตุ\-สมฺปยุตฺต\-ทุก\-กุสล\-ตฺติกํ นิฏฺฐิตํ.}
\csromancenter{นเหตุ\-สเหตุกทุก 1. กุสลตฺติก 6. นเหตุ\-สเหตุกทุก 1. กุสลตฺติก}
\csromanheader{2}{1-2. กุลลากุสลปท 1-7. ปฏิจฺจาทิวาร}
\prose{\csromanfolio{345}ปจฺจยจตุกฺกเหตุ}

\proseitem{93}{นเหตุํ สเหตุกํ กุสลํ ธมฺมํ ปฏิจฺจ นเหตุ สเหตุโก กุสโล ธมฺโม อุปฺปชฺชติ เหตุปจฺจยา. (สํขิตฺตํ.)}

\proseitem{94}{เหตุยา เอกํ, อารมฺมเณ เอกํ, (สพฺพตฺถ เอกํ.) อวิคเต เอกํ.}

\prose{นอธิปติยา เอกํ, นปุเรชาเต เอกํ ฯเปฯ นวิปาเก เอกํ, นวิปฺปยุตฺเต เอกํ. (ปญฺหาวาเรปิ สพฺพตฺถ เอกํ.)}

\proseitem{95}{นเหตุํ สเหตุกํ อกุสลํ ธมฺมํ ปฏิจฺจ นเหตุ สเหตุโก อกุสโล ธมฺโม อุปฺปชฺชติ เหตุปจฺจยา. (สํขิตฺตํ.)}

\proseitem{96}{เหตุยา เอกํ, อารมฺมเณ เอกํ, (สพฺพตฺถ เอกํ.) อวิคเต เอกํ. (เหตุปจฺจโย นตฺถิ. ปญฺหาวาเรปิ สพฺพตฺถ เอกํ, ปญฺหาวาเรปิ อิเมสํ ติณฺณนฺนํ เหตุปจฺจโย นตฺถิ.)}

\csromanmatikaanchor{mtk.61}
\csromanmatikaanchor{mtk.62}
\csromantocmarknopage{subparagraph}{6. นเหตุ สเหตุกทุก}{mtk.61}
\bookmark[dest={mtk.61},level=5]{6. นเหตุ สเหตุกทุก}
\csromantocmarkref{subparagraph}{1. กุสลตฺติก ...}{mtk.62}
\bookmark[dest={mtk.62},level=5]{1. กุสลตฺติก ...}
\csromanmarkh{3. อพฺยากตปท}\csromanmarkhii{1. ปฏิจฺจวาร}\csromanheadernomark{6}{3. \textbf{อพฺยากตปท 1. ปฏิจฺจวาร}}
\prose{ปจฺจยจตุกฺกเหตุ}

\proseitem{97}{นเหตุํ สเหตุกํ อพฺยากตํ ธมฺมํ ปฏิจฺจ นเหตุ สเหตุโก อพฺยากโต ธมฺโม อุปฺปชฺชติ เหตุปจฺจยา.  นเหตุํ สเหตุกํ อพฺยากตํ ธมฺมํ ปฏิจฺจ นเหตุ อเหตุโก อพฺยากโต ธมฺโม อุปฺปชฺชติ เหตุปจฺจยา.  นเหตุํ สเหตุกํ อพฺยากตํ ธมฺมํ ปฏิจฺจ นเหตุ สเหตุโก อพฺยากโต จ นเหตุ อเหตุโก อพฺยากโต จ ธมฺมา อุปฺปชฺชนฺติ เหตุปจฺจยา. (3)}

\prose{นเหตุํ อเหตุกํ อพฺยากตํ ธมฺมํ ปฏิจฺจ นเหตุ อเหตุโก อพฺยากโต ธมฺโม อุปฺปชฺชติ เหตุปจฺจยา. ตีณิ.}

\prose{\csromanfolio{346}นเหตุํ สเหตุกํ อพฺยากตญฺจ นเหตุํ อเหตุกํ อพฺยากตญฺจ ธมฺมํ ปฏิจฺจ นเหตุ สเหตุโก อพฺยากโต ธมฺโม อุปฺปชฺชติ เหตุปจฺจยา. ตีณิ. (สํขิตฺตํ.)}

\proseitem{98}{เหตุยา นว, อารมฺมเณ จตฺตาริ, อธิปติยา ปญฺจ ฯเปฯ ปุเรชาเต อาเสวเน เทฺว, กมฺเม นว, วิปาเก นว ฯเปฯ อวิคเต นว. (สํขิตฺตํ.)}

\vspace{\tipitakaparskip}
\csromancenter{\textbf{นเหตุ}}
\proseitem{99}{\textbf{นเหตุ}ํ อเหตุกํ อพฺยากตํ ธมฺมํ ปฏิจฺจ นเหตุ อเหตุโก อพฺยากโต ธมฺโม อุปฺปชฺชติ นเหตุปจฺจยา. (สํขิตฺตํ.)}

\vspace{\tipitakaparskip}
\csromancenter{\textbf{นเหตุ}ยา เอกํ, นอารมฺมเณ ตีนิ, นอธิปติยา นว, นอนนฺตเร ตีณิ, นสมนนฺตเร ตีณิ, นอญฺญมญฺเญ ตีณิ, นฺเอาปนิสฺสเย ตีณิ, นปุเรชาเต นว, นปจฺฉาชาเต นว, นอาเสวเน นว, นกมฺเม เทฺว, นวิปาเก ปญฺจ, นอาหาเร เอกํ, ไนนฺทฺริเย เอกํ, นฌาเน เอกํ, นมคฺเค เอกํ, นสมฺปยุตฺเต ตีณิ, นวิปฺปยุตฺเต เทฺว, โนนตฺถิยา ตีณิ, โนวิคเต ตีณิ. (สํขิตฺตํ.)}
\prose{เหตุปจฺจยา นอารมฺมเณ ตีณิ. (สํขิตฺตํ.)}

\prose{\textbf{นเหตุ}ปจฺจยา อารมฺมเณ เอกํ. (สํขิตฺตํ.)}

\prose{(สหชาตวาโรปิ ปจฺจยวาโรปิ นิสฺสยวาโรปิ สํสฏฺฐ\-วาโรปิ สมฺปยุตฺตวาโรปิ ปฏิจฺจ\-วาร\-สทิสา วิตฺถาเรตพฺพา.)}

\csromanmarkh{3. อพฺยากตปท}\csromanmarkhii{7. ปญฺหาวาร}\csromanheadernomark{2}{3. \textbf{อพฺยากตปท 7. ปญฺหาวาร}}
\prose{ปจฺจยจตุกฺกอารมฺมณ}

\proseitem{101}{สเหตุโก อพฺยากโต ธมฺโม นเหตุสฺส สเหตุกสฺส อพฺยากตสฺส ธมฺมสฺส อารมฺมณ\-ปจฺจเยน ปจฺจโย. (สํขิตฺตํ.)}

\proseitem{102}{อารมฺมเณ จตฺตาริ, อธิปติยา จตฺตาริ, อนนฺตเร จตฺตาริ, สมนนฺตเร จตฺตาริ, สหชาเต สตฺต, อญฺญมญฺเญ ฉ, นิสฺสเย สตฺต,}

\vspace{\tipitakaparskip}
\csromancenter{สปฺปจฺจยทุก 1. กุสลตฺติก อุปนิสฺสเย จตฺตาริ, ปุเรชาเต เทฺว, ปจฺฉาชาเต เทฺว, อาเสวเน เทฺว, กมฺเม จตฺตาริ, วิปาเก จตฺตาริ, อาหาเร จตฺตาริ, อินฺทฺริเย จตฺตาริ, ฌาเน จตฺตาริ, มคฺเค ตีณิ, สมฺปยุตฺเต เทฺว, วิปฺปยุตฺเต ตีณิ, อตฺถิยา สตฺต, นตฺถิยา จตฺตาริ, วิคเต จตฺตาริ, อวิคเต สตฺต. (สํขิตฺตํ.)}
\proseitem{103}{\csromanfolio{347}สเหตุโก อพฺยากโต ธมฺโม นเหตุสฺส สเหตุกสฺส อพฺยากตสฺส ธมฺมสฺส อารมฺมณ\-ปจฺจเยน ปจฺจโย. สหชาต\-ปจฺจเยน ปจฺจโย. อุปนิสฺสย\-ปจฺจเยน ปจฺจโย. (สํขิตฺตํ.)}

\proseitem{104}{นเหตุยา สตฺต, นอารมฺมเณ สตฺต. (สํขิตฺตํ.)}

\prose{อารมฺมณ\-ปจฺจยา นเหตุยา จตฺตาริ. (สํขิตฺตํ.)}

\prose{นเหตุปจฺจยา อารมฺมเณ จตฺตาริ. (สํขิตฺตํ.)}

\prosewithcloser{(ยถา กุสลตฺติเก ปญฺหาวารสฺส อนุโลมมฺปิ ปจฺจนียมฺปิ อนุโลมปจฺจนียมฺปิ ปจฺจนียานุโลมมฺปิ คณิตํ, เอวํ คเณตพฺพํ.)}{นเหตุ\-สเหตุก\-ทุก\-กุสล\-ตฺติกํ นิฏฺฐิตํ.}
\nitthitam{เหตุโคจฺฉกํ นิฏฺฐิตํ.}
\csromanmatikaanchor{mtk.63}
\csromanmatikaanchor{mtk.64}
\csromantocmarknopage{subparagraph}{7. สปฺปจฺจยทุก}{mtk.63}
\bookmark[dest={mtk.63},level=5]{7. สปฺปจฺจยทุก}
\csromantocmarkref{subparagraph}{1. กุสลตฺติก ...}{mtk.64}
\bookmark[dest={mtk.64},level=5]{1. กุสลตฺติก ...}
\csromanmarkh{7. สปฺปจฺจยทุก}\csromanmarkhii{1. กุสลตฺติก}\csromanheadernomark{6}{7. สปฺปจฺจยทุก 1. กุสลตฺติก}
\csromanheader{2}{1-2. กุสลากุสลปท 1-7. ปฏิจฺจาทิวาร}
\prose{ปจฺจยจตุกฺกเหตุ}

\proseitem{103}{สปฺปจฺจยํ กุสลํ ธมฺมํ ปฏิจฺจ สปฺปจฺจโย กุสโล ธมฺโม อุปฺปชฺชติ เหตุปจฺจยา. (1) (สํขิตฺตํ.)}

\prose{เหตุยา เอกํ, อารมฺมเณ เอกํ, อธิปติยา เอกํ ฯเปฯ อวิคเต เอกํ. (สํขิตฺตํ.)}

\prose{นอธิปติ}

\proseitem{103}{สปฺปจฺจยํ กุสลํ ธมฺมํ ปฏิจฺจ สปฺปจฺจโย กุสโล ธมฺโม อุปฺปชฺชติ นอธิปติ\-ปจฺจยา. (สํขิตฺตํ.)}

\proseitem{4}{\csromanfolio{348}นอธิปติยา เอกํ, นปุเรชาเต เอกํ, นปจฺฉาชาเต เอกํ, นอาเสวเน เอกํ, นกมฺเม เอกํ, นวิปาเก เอกํ, นวิปฺปยุตฺเต เอกํ. (สํขิตฺตํ.)}

\prose{(สหชาตวาโรปิ ปจฺจยวาโรปิ นิสฺสยวาโรปิ สํสฏฺฐ\-วาโรปิ สมฺปยุตฺตวาโรปิ ปฏิจฺจ\-วาร\-สทิสา วิตฺถาเรตพฺพา.)}

\prose{ปจฺจยจตุกฺกเหตุ}

\proseitem{5}{สปฺปจฺจโย กุสโล ธมฺโม สปฺปจฺจยสฺส กุสลสฺส ธมฺมสฺส เหตุปจฺจเยน ปจฺจโย. (สํขิตฺตํ.)}

\proseitem{6}{เหตุยา เอกํ, อารมฺมเณ เอกํ ฯเปฯ อวิคเต เอกํ. (สํขิตฺตํ.)}

\proseitem{7}{สปฺปจฺจยํ อกุสลํ ธมฺมํ ปฏิจฺจ สปฺปจฺจโย อกุสโล ธมฺโม อุปฺปชฺชติ เหตุปจฺจยา. (1) (สํขิตฺตํ.)}

\proseitem{8}{เหตุยา เอกํ, อารมฺมเณ เอกํ ฯเปฯ อวิคเต เอกํ.}

\vspace{\tipitakaparskip}
\csromancenter{(สหชาตวาเรปิ ฯเปฯ ปญฺหาวาเรปิ สพฺพตฺถ เอกํ.)}
\csromanheader{2}{3. \textbf{อพฺยากตปท 1-7. ปฏิจฺจาทิวาร}}
\prose{ปจฺจยจตุกฺกเหตุ}

\proseitem{9}{สปฺปจฺจยํ อพฺยากตํ ธมฺมํ ปฏิจฺจ สปฺปจฺจโย อพฺยากโต ธมฺโม อุปฺปชฺชติ เหตุปจฺจยา. (สํขิตฺตํ.)}

\proseitem{10}{เหตุยา เอกํ, อารมฺมเณ เอกํ ฯเปฯ อวิคเต เอกํ. (สํขิตฺตํ.)}

\prose{นเหตุยา เอกํ, นอารมฺมเณ เอกํ, นอธิปติยา เอกํ. (สพฺพตฺถ เอกํ, สํขิตฺตํ.)}

\prose{ปจฺจยจตุกฺกเหตฺวาทิ}

\proseitem{11}{สปฺปจฺจโย อพฺยากโต ธมฺโม สปฺปจฺจยสฺส อพฺยากตสฺส ธมฺมสฺส เหตุปจฺจเยน ปจฺจโย. (1)}

\prose{สปฺปจฺจโย อพฺยากโต ธมฺโม สปฺปจฺจยสฺส อพฺยากตสฺส ธมฺมสฺส อารมฺมณ\-ปจฺจเยน ปจฺจโย. (1)}

\csromanmarkh{8. สงฺขตทุก}\csromanmarkhii{1. กุสลตฺติก}\csromanheadernomark{2}{8. สงฺขตทุก 1. กุสลตฺติก}
\prose{\csromanfolio{349}อปฺปจฺจโย อพฺยากโต ธมฺโม สปฺปจฺจยสฺส อพฺยากตสฺส ธมฺมสฺส อารมฺมณ\-ปจฺจเยน ปจฺจโย. (1)}

\prose{สปฺปจฺจโย อพฺยากโต ธมฺโม สปฺปจฺจยสฺส อพฺยากตสฺส ธมฺมสฺส อุปนิสฺสย\-ปจฺจเยน ปจฺจโย. (1)}

\prose{อปฺปจฺจโย อพฺยากโต ธมฺโม สปฺปจฺจยสฺส อพฺยากตสฺส ธมฺมสฺส อุปนิสฺสย\-ปจฺจเยน ปจฺจโย. อารมฺมณู\-ปนิสฺสโย. (1) (สํขิตฺตํ.)}

\prose{เหตุยา เอกํ, อารมฺมเณ เทฺว, อธิปติยา เทฺว, อนนฺตเร เอกํ ฯเปฯ นิสฺสเย เอกํ, อุปนิสฺสเย เทฺว, ปุเรชาเต เอกํ, ปจฺฉาชาเต เอกํ, (สพฺพตฺถ เอกํ.) อวิคเต เอกํ. (สํขิตฺตํ.)}

\proseitem{12}{สปฺปจฺจโย อพฺยากโต ธมฺโม สปฺปจฺจยสฺส อพฺยากตสฺส ธมฺมสฺส อารมฺมณ\-ปจฺจเยน ปจฺจโย. สหชาต\-ปจฺจเยน ปจฺจโย. อุปนิสฺสย\-ปจฺจเยน ปจฺจโย.}

\prose{อปฺปจฺจโย อพฺยากโต ธมฺโม สปฺปจฺจยสฺส อพฺยากตสฺส ธมฺมสฺส อารมฺมณ\-ปจฺจเยน ปจฺจโย. อุปนิสฺสย\-ปจฺจเยน ปจฺจโย. (สํขิตฺตํ.)}

\proseitem{13}{นเหตุยา เทฺว, นอารมฺมเณ เอกํ. (สํขิตฺตํ.)}

\prose{เหตุปจฺจยา นอารมฺมเณ เอกํ. (สํขิตฺตํ.)}

\prosewithcloser{นเหตุปจฺจยา อารมฺมเณ เทฺว. (สํขิตฺตํ.)}{สปฺปจฺจย\-ทุก\-กุสล\-ตฺติกํ นิฏฺฐิตํ.}
\csromanmarkh{8. สงฺขตทุก}\csromanmarkhii{1. กุสลตฺติก}\csromanheadernomark{2}{8. สงฺขตทุก 1. กุสลตฺติก}
\proseitemwithcloser{14}{สงฺขตํ กุสลํ ธมฺมํ ปฏิจฺจ สงฺขโต กุสโล ธมฺโม อุปฺปชฺชติ เหตุปจฺจยา. (สํขิตฺตํ. สปฺปจฺจย\-สทิสํ วิตฺถาเรตพฺพํ.)}{สงฺขต\-ทุก\-กุสล\-ตฺติกํ นิฏฺฐิตํ.}
\csromanmarkh{ธมฺมานุโลม ทุกติก\-ปฏฺฐาน\-ปาฬิ}\csromanmarkhii{9. สนิทสฺสนทุก 1. กุสลตฺติก}\csromanheadernomark{2}{ธมฺมานุโลม ทุกติก\-ปฏฺฐาน\-ปาฬิ 9. สนิทสฺสนทุก 1. กุสลตฺติก}
\csromanmarkh{1. กุสลปท}\csromanmarkhii{1. ปฏิจฺจวาร}\csromanheadernomark{2}{1. \textbf{กุสลปท 1. ปฏิจฺจวาร}}
\prose{\csromanfolio{350}ปจฺจยจตุกฺกเหตุ}

\proseitem{15}{อนิทสฺสนํ กุสลํ ธมฺมํ ปฏิจฺจ อนิทสฺสโน กุสโล ธมฺโม อุปฺปชฺชติ เหตุปจฺจยา. (สํขิตฺตํ.)}

\proseitem{16}{เหตุยา เอกํ, อารมฺมเณ เอกํ ฯเปฯ อวิคเต เอกํ.}

\vspace{\tipitakaparskip}
\csromancenter{(สหชาตวาเรปิ ฯเปฯ ปญฺหาวาเรปิ สพฺพตฺถ เอกํ.)}
\prose{2. อกุสลปทเหตุ}

\proseitem{17}{อนิทสฺสนํ อกุสลํ ธมฺมํ ปฏิจฺจ อนิทสฺสโน อกุสโล ธมฺโม อุปฺปชฺชติ เหตุปจฺจยา. (สํขิตฺตํ.)}

\proseitem{18}{เหตุยา เอกํ, อารมฺมเณ เอกํ ฯเปฯ อวิคเต เอกํ.}

\vspace{\tipitakaparskip}
\csromancenter{(สหชาตวาเรปิ ฯเปฯ ปญฺหาวาเรปิ สพฺพตฺถ เอกํ.)}
\prose{3. อพฺยากตปทเหตุ}

\proseitem{19}{อนิทสฺสนํ อพฺยากตํ ธมฺมํ ปฏิจฺจ อนิทสฺสโน อพฺยากโต ธมฺโม อุปฺปชฺชติ เหตุปจฺจยา. (สํขิตฺตํ.)}

\proseitem{20}{เหตุยา ตีณิ, อารมฺมเณ เอกํ, อธิปติยา ตีณิ ฯเปฯ อวิคเต ตีณิ. (สํขิตฺตํ.)}

\prose{นเหตุยา ตีณิ, (สพฺพตฺถ ตีณิ.) นกมฺเม ตีณิ, นวิปาเก ตีณิ, นอาหาเร ตีณิ ฯเปฯ โนวิคเต ตีณิ. (สํขิตฺตํ.)}

\prose{(สหชาตวาโรปิ ปจฺจยวาโรปิ นิสฺสยวาโรปิ สํสฏฺฐ\-วาโรปิ สมฺปยุตฺตวาโรปิ ปฏิจฺจ\-วาร\-สทิสา วิตฺถาเรตพฺพา.)}

\csromanmarkh{3. อพฺยากตปท}\csromanmarkhii{7. ปญฺหาวาร}\csromanheadernomark{2}{3. \textbf{อพฺยากตปท 7. ปญฺหาวาร}}
\prose{ปจฺจยจตุกฺกเหตฺวาทิ}

\proseitem{21}{อนิทสฺสโน อพฺยากโต ธมฺโม อนิทสฺสนสฺส อพฺยากตสฺส ธมฺมสฺส เหตุปจฺจเยน ปจฺจโย. ตีณิ.}

\csromanmarkh{10. สปฺปฏิฆทุก}\csromanmarkhii{1. กุสลตฺติก}\csromanheadernomark{2}{10. สปฺปฏิฆทุก 1. กุสลตฺติก}
\prose{\csromanfolio{351}สนิทสฺสโน อพฺยากโต ธมฺโม อนิทสฺสนสฺส อพฺยากตสฺส ธมฺมสฺส อารมฺมณ\-ปจฺจเยน ปจฺจโย. (1)}

\prose{อนิทสฺสโน อพฺยากโต ธมฺโม อนิทสฺสนสฺส อพฺยากตสฺส ธมฺมสฺส อารมฺมณ\-ปจฺจเยน ปจฺจโย. (1)}

\prose{สนิทสฺสโน อพฺยากโต ธมฺโม อนิทสฺสนสฺส อพฺยากตสฺส ธมฺมสฺส อธิปติ\-ปจฺจเยน ปจฺจโย. อารมฺมณา\-ธิปติ. (1)}

\prose{อนิทสฺสโน อพฺยากโต ธมฺโม อนิทสฺสนสฺส อพฺยากตสฺส ธมฺมสฺส อธิปติ\-ปจฺจเยน ปจฺจโย. อารมฺมณา\-ธิปติ, สหชาตา\-ธิปติ.  อนิทสฺสโน อพฺยากโต ธมฺโม สนิทสฺสนสฺส อพฺยากตสฺส จ อนิทสฺสนสฺส อพฺยากตสฺส จ ธมฺมสฺส อธิปติ\-ปจฺจเยน ปจฺจโย. สหชาตา\-ธิปติ. (2)}

\prose{สนิทสฺสโน อพฺยากโต ธมฺโม อนิทสฺสนสฺส อพฺยากตสฺส ธมฺมสฺส อุปนิสฺสย\-ปจฺจเยน ปจฺจโย. ปกตู\-ปนิสฺสโย. (1)}

\prose{อนิทสฺสโน อพฺยากโต ธมฺโม อนิทสฺสนสฺส อพฺยากตสฺส ธมฺมสฺส อุปนิสฺสย\-ปจฺจเยน ปจฺจโย. อารมฺมณู\-ปนิสฺสโย, อนนฺตรู\-ปนิสฺสโย, ปกตู\-ปนิสฺสโย ฯเปฯ. (1) (สํขิตฺตํ.)}

\proseitemwithcloser{22}{ปุเรชาเต ตีณิ, ปจฺฉาชาเต ตีณิ, อาเสวเน เอกํ, กมฺเม ตีณิ, วิปาเก ตีณิ, (สพฺพตฺถ ตีณิ.) สมฺปยุตฺเต เอกํ, วิปฺปยุตฺเต ตีณิ, อตฺถิยา ปญฺจ, นตฺถิยา เอกํ ฯเปฯ อวิคเต ปญฺจ. (สํขิตฺตํ.)}{สนิทสฺสน\-ทุก\-กุสล\-ตฺติกํ นิฏฺฐิตํ.}
\csromancenter{(อปฺปจฺจยมฺปิ อสงฺขตมฺปิ สนิทสฺสนมฺปิ น ลพฺภติ.)}
\csromanmarkh{10. สปฺปฏิฆทุก}\csromanmarkhii{1. กุสลตฺติก}\csromanheadernomark{2}{10. สปฺปฏิฆทุก 1. กุสลตฺติก}
\csromanmarkh{1-2. กุสลากุสลปท}\csromanmarkhii{1. ปฏิจฺจวาร}\csromanheadernomark{2}{1-2. กุสลากุสลปท 1. ปฏิจฺจวาร}
\prose{ปจฺจยจตุกฺกเหตุ}

\proseitem{23}{อปฺปฏิฆํ กุสลํ ธมฺมํ ปฏิจฺจ อปฺปฏิโฆ กุสโล ธมฺโม อุปฺปชฺชติ เหตุปจฺจยา. (สํขิตฺตํ.)}

\proseitem{24}{\csromanfolio{352}เหตุยา เอกํ, อารมฺมเณ เอกํ ฯเปฯ วิปฺปยุตฺเต เอกํ ฯเปฯ อวิคเต เอกํ. (สํขิตฺตํ.) (สหชาตวาเรปิ ฯเปฯ ปญฺหาวาเรปิ สพฺพตฺถ เอกํ.)}

\proseitem{25}{อปฺปฏิฆํ อกุสลํ ธมฺมํ ปฏิจฺจ อปฺปฏิโฆ อกุสโล ธมฺโม อุปฺปชฺชติ เหตุปจฺจยา. (สํขิตฺตํ.)}

\prose{เหตุยา เอกํ ฯเปฯ อวิคเต เอกํ. (ปญฺหาวาเรปิ สพฺพตฺถ เอกํ.)}

\csromanmarkh{3. อพฺยากตปท}\csromanmarkhii{1. ปฏิจฺจวาร}\csromanheadernomark{2}{3. \textbf{อพฺยากตปท 1. ปฏิจฺจวาร}}
\prose{ปจฺจยจตุกฺกเหตุ}

\proseitem{26}{สปฺปฏิฆํ อพฺยากตํ ธมฺมํ ปฏิจฺจ สปฺปฏิโฆ อพฺยากโต ธมฺโม อุปฺปชฺชติ เหตุปจฺจยา.  สปฺปฏิฆํ อพฺยากตํ ธมฺมํ ปฏิจฺจ อปฺปฏิโฆ อพฺยากโต ธมฺโม อุปฺปชฺชติ เหตุปจฺจยา.  สปฺปฏิฆํ อพฺยากตํ ธมฺมํ ปฏิจฺจ สปฺปฏิโฆ อพฺยากโต จ อปฺปฏิโฆ อพฺยากโต จ ธมฺมา อุปฺปชฺชนฺติ เหตุปจฺจยา. นว.}

\prose{อปฺปฏิฆํ อพฺยากตํ ธมฺมํ ปฏิจฺจ อปฺปฏิโฆ อพฺยากโต ธมฺโม อุปฺปชฺชติ อารมฺมณ\-ปจฺจยา. (สํขิตฺตํ.)}

\proseitem{27}{เหตุยา นว, อารมฺมเณ เอกํ, อธิปติยา นว ฯเปฯ อญฺญมญฺเญ ฉ, นิสฺสเย นว, อุปนิสฺสเย เอกํ, ปุเรชาเต อาเสวเน เอกํ, กมฺเม นว, วิปาเก นว ฯเปฯ อวิคเต นว. (สํขิตฺตํ.)}

\proseitem{28}{สปฺปฏิฆํ อพฺยากตํ ธมฺมํ ปฏิจฺจ สปฺปฏิโฆ อพฺยากโต ธมฺโม อุปฺปชฺชติ นเหตุปจฺจยา. (สํขิตฺตํ.)}

\proseitem{29}{\textbf{นเหตุ}ยา นว ฯเปฯ โนวิคเต นว. (สพฺพตฺถ นว.) (สํขิตฺตํ.)}

\csromanmarkh{3. อพฺยากตปท}\csromanmarkhii{7. ปญฺหาวาร}\csromanheadernomark{2}{3. \textbf{อพฺยากตปท 7. ปญฺหาวาร}}
\prose{ปจฺจยจตุกฺกเหตุ}

\proseitem{30}{อปฺปฏิโฆ อพฺยากโต ธมฺโม อปฺปฏิฆสฺส อพฺยากตสฺส ธมฺมสฺส เหตุปจฺจเยน ปจฺจโย. (3) (สํขิตฺตํ.)}

\csromanmarkh{11. รูปีทุก}\csromanmarkhii{1. กุสลตฺติก}\csromanheadernomark{2}{11. รูปีทุก 1. กุสลตฺติก}
\proseitem{31}{\csromanfolio{353}เหตุยา ตีณิ, อารมฺมเณ เทฺว, อธิปติยา ตีณิ, อนนฺตเร เอกํ, สมนนฺตเร เอกํ, สหชาเต นว, อญฺญมญฺเญ ฉ, นิสฺสเย นว, อุปนิสฺสเย เทฺว, ปุเรชาเต ตีณิ, ปจฺฉาชาเต ตีณิ, อาเสวเน เอกํ, กมฺเม ตีณิ, วิปาเก ตีณิ, อาหาเร ตีณิ, อินฺทฺริเย ปญฺจ, ฌาเน ตีณิ, มคฺเค ตีณิ, สมฺปยุตฺเต เอกํ, วิปฺปยุตฺเต จตฺตาริ, อตฺถิยา นว, นตฺถิยา เอกํ, วิคเต เอกํ, อวิคเต นว.}

\proseitem{32}{สปฺปฏิโฆ อพฺยากโต ธมฺโม สปฺปฏิฆสฺส อพฺยากตสฺส ธมฺมสฺส สหชาต\-ปจฺจเยน ปจฺจโย. (สํขิตฺตํ.)}

\proseitem{33}{นเหตุยา นว, นอารมฺมเณ นว. (สํขิตฺตํ.)}

\prose{เหตุปจฺจยา นอารมฺมเณ ตีณิ. (สํขิตฺตํ.)}

\prose{นเหตุปจฺจยา อารมฺมเณ เทฺว. (สํขิตฺตํ.)}

\prosewithcloser{(ยถา กุสลตฺติเก ปญฺหาวารสฺส อนุโลมมฺปิ ปจฺจนียมฺปิ อนุโลมปจฺจนียมฺปิ ปจฺจนียานุโลมมฺปิ คณิตํ, เอวํ คเณตพฺพํ.)}{สปฺปฏิฆ\-ทุก\-กุสล\-ตฺติกํ นิฏฺฐิตํ.}
\csromanmarkh{11. รูปีทุก}\csromanmarkhii{1. กุสลตฺติก}\csromanheadernomark{2}{11. รูปีทุก 1. กุสลตฺติก}
\csromanheader{2}{1-2. กุสลากุสลปท 1-7. ปฏิจฺจาทิวาร}
\prose{ปจฺจยจตุกฺกเหตุ}

\proseitem{34}{อรูปิํ กุสลํ ธมฺมํ ปฏิจฺจ อรูปี กุสโล ธมฺโม อุปฺปชฺชติ เหตุปจฺจยา. (สํขิตฺตํ.)}

\proseitem{35}{เหตุยา เอกํ, อารมฺมเณ เอกํ. (สํขิตฺตํ. สหชาตวาเรปิ ฯเปฯ ปญฺหาวาเรปิ สพฺพตฺถ เอกํ.)}

\proseitem{36}{อรูปิํ อกุสลํ ธมฺมํ ปฏิจฺจ อรูปี อกุสโล ธมฺโม อุปฺปชฺชติ เหตุปจฺจยา. (สํขิตฺตํ.)}

\proseitem{37}{เหตุยา เอกํ, อารมฺมเณ เอกํ ฯเปฯ อวิคเต เอกํ. (สํขิตฺตํ. สหชาตวาเรปิ ฯเปฯ ปญฺหาวาเรปิ สพฺพตฺถ เอกํ.)}

\csromanmarkh{ธมฺมานุโลม ทุกติก\-ปฏฺฐาน\-ปาฬิ}\csromanmarkhii{3. อพฺยากตปท 1-7. ปฏิจฺจาทิวาร}\csromanheadernomark{2}{ธมฺมานุโลม ทุกติก\-ปฏฺฐาน\-ปาฬิ 3. อพฺยากตปท 1-7. ปฏิจฺจาทิวาร}
\prose{\csromanfolio{354}ปจฺจยจตุกฺกเหตุ อารมฺมณ}

\proseitem{38}{รูปิํ อพฺยากตํ ธมฺมํ ปฏิจฺจ รูปี อพฺยากโต ธมฺโม อุปฺปชฺชติ เหตุปจฺจยา. ตีณิ.}

\prose{อรูปิํ อพฺยากตํ ธมฺมํ ปฏิจฺจ อรูปี อพฺยากโต ธมฺโม อุปฺปชฺชติ เหตุปจฺจยา. ตีณิ.}

\prose{รูปิํ อพฺยากตญฺจ อรูปิํ อพฺยากตญฺจ ธมฺมํ ปฏิจฺจ รูปี อพฺยากโต ธมฺโม อุปฺปชฺชติ เหตุปจฺจยา. ตีณิ.}

\prose{รูปิํ อพฺยากตํ ธมฺมํ ปฏิจฺจ อรูปี อพฺยากโต ธมฺโม อุปฺปชฺชติ อารมฺมณ\-ปจฺจยา. (สํขิตฺตํ.)}

\proseitem{39}{เหตุยา นว, อารมฺมเณ ตีณิ, อธิปติยา ปญฺจ, อนนฺตเร ตีณิ, สมนนฺตเร ตีณิ ฯเปฯ อญฺญมญฺเญ ฉ ฯเปฯ ปุเรชาเต เอกํ, อาเสวเน เอกํ, กมฺเม นว, วิปาเก นว ฯเปฯ อวิคเต นว. (สํขิตฺตํ.)}

\proseitem{40}{นเหตุยา นว, นอารมฺมเณ ตีณิ, นอธิปติยา นว ฯเปฯ นกมฺเม เทฺว, นวิปาเก ปญฺจ, นอาหาเร เอกํ, นอินฺทฺริเย เอกํ, นฌาเน เทฺว, นมคฺเค นว, นสมฺปยุตฺเต ตีณิ, นวิปฺปยุตฺเต เทฺว, โนนตฺถิยา ตีณิ, โนวิคเต ตีณิ. (สํขิตฺตํ. ปจฺจนียํ.)}

\prose{ปจฺจยจตุกฺกเหตุ}

\proseitem{41}{อรูปี อพฺยากโต ธมฺโม อรูปิสฺส อพฺยากตสฺส ธมฺมสฺส เหตุปจฺจเยน ปจฺจโย. (สํขิตฺตํ.)}

\proseitem{42}{เหตุยา ตีณิ, อารมฺมเณ เทฺว, อธิปติยา ตีณิ, อนนฺตเร เอกํ, สมนนฺตเร เอกํ, สหชาเต สตฺต, อญฺญมญฺเญ ฉ, นิสฺสเย สตฺต, อุปนิสฺสเย เทฺว, ปุเรชาเต เอกํ, ปจฺฉาชาเต เอกํ, อาเสวเน เอกํ, กมฺเม ตีณิ, วิปาเก ตีณิ, อาหาเร จตฺตาริ, อินฺทฺริเย ฉ, ฌาเน ตีณิ, มคฺเค ตีณิ, สมฺปยุตฺเต เอกํ, วิปฺปยุตฺเต เทฺว ฯเปฯ อวิคเต สตฺต. (สํขิตฺตํ.)}

\prose{นเหตุยา สตฺต, นอารมฺมเณ สตฺต. (สํขิตฺตํ.)}

\csromanmarkh{12. โลกิยทุก}\csromanmarkhii{1. กุสลตฺติก}\csromanheadernomark{2}{12. โลกิยทุก 1. กุสลตฺติก}
\prose{\csromanfolio{355}เหตุปจฺจยา นอารมฺมเณ ตีณิ. (สํขิตฺตํ.)}

\prose{นเหตุปจฺจยา อารมฺมเณ เทฺว. (สํขิตฺตํ.)}

\prosewithcloser{(ยถา กุสลตฺติเก ปญฺหาวารสฺส อนุโลมมฺปิ ปจฺจนียมฺปิ อนุโลมปจฺจนียมฺปิ ปจฺจนียานุโลมมฺปิ คณิตํ, เอวํ คเณตพฺพํ.)}{รูปีทุก\-กุสล\-ตฺติกํ นิฏฺฐิตํ.}
\csromanmarkh{12. โลกิยทุก}\csromanmarkhii{1. กุสลตฺติก}\csromanheadernomark{2}{12. โลกิยทุก 1. กุสลตฺติก}
\csromanheader{2}{1. \textbf{กุสลปท 1-7. ปฏิจฺจาทิวาร}}
\prose{ปจฺจยจตุกฺกเหตุ}

\proseitem{43}{โลกิยํ กุสลํ ธมฺมํ ปฏิจฺจ โลกิโย กุสโล ธมฺโม อุปฺปชฺชติ เหตุปจฺจยา. (1)}

\prose{โลกุตฺตรํ กุสลํ ธมฺมํ ปฏิจฺจ โลกุตฺตโร กุสโล ธมฺโม อุปฺปชฺชติ เหตุปจฺจยา. (1)}

\prose{โลกิยํ กุสลํ ธมฺมํ ปฏิจฺจ โลกิโย กุสโล ธมฺโม อุปฺปชฺชติ อารมฺมณ\-ปจฺจยา. (1)}

\prose{โลกุตฺตรํ กุสลํ ธมฺมํ ปฏิจฺจ โลกุตฺตโร กุสโล ธมฺโม อุปฺปชฺชติ อารมฺมณ\-ปจฺจยา. (1) (สํขิตฺตํ.)}

\proseitem{44}{เหตุยา เทฺว, อารมฺมเณ เทฺว, อธิปติยา เทฺว ฯเปฯ กมฺเม เทฺว ฯเปฯ อวิคเต เทฺว. (สํขิตฺตํ. อนุโลมํ.)}

\proseitem{45}{นอธิปติยา เทฺว ฯเปฯ นอาเสวเน เอกํ ฯเปฯ นวิปฺปยุตฺเต เทฺว. (สํขิตฺตํ.)}

\prose{(สหชาตวาโรปิ ปจฺจยวาโรปิ นิสฺสยวาโรปิ สํสฏฺฐ\-วาโรปิ สมฺปยุตฺตวาโรปิ ปฏิจฺจ\-วาร\-สทิสา วิตฺถาเรตพฺพา.)}

\prose{ปจฺจยจตุกฺกเหตฺวาทิ}

\proseitem{46}{โลกิโย กุสโล ธมฺโม โลกิยสฺส กุสลสฺส ธมฺมสฺส เหตุปจฺจเยน ปจฺจโย. (1)}

\prose{\csromanfolio{356}โลกุตฺตโร กุสโล ธมฺโม โลกุตฺตรสฺส กุสลสฺส ธมฺมสฺส เหตุปจฺจเยน ปจฺจโย. (1)}

\prose{โลกิโย กุสโล ธมฺโม โลกิยสฺส กุสลสฺส ธมฺมสฺส อารมฺมณ\-ปจฺจเยน ปจฺจโย. (1)}

\prose{โลกุตฺตโร กุสโล ธมฺโม โลกิยสฺส กุสลสฺส ธมฺมสฺส อารมฺมณ\-ปจฺจเยน ปจฺจโย. (1)}

\prose{โลกิโย กุสโล ธมฺโม โลกิยสฺส กุสลสฺส ธมฺมสฺส อธิปติ\-ปจฺจเยน ปจฺจโย. อารมฺมณา\-ธิปติ, สหชาตา\-ธิปติ. (1)}

\prose{โลกุตฺตโร กุสโล ธมฺโม โลกุตฺตรสฺส กุสลสฺส ธมฺมสฺส อธิปติ\-ปจฺจเยน ปจฺจโย. สหชาตา\-ธิปติ.  โลกุตฺตโร กุสโล ธมฺโม โลกิยสฺส กุสลสฺส ธมฺมสฺส อธิปติ\-ปจฺจเยน ปจฺจโย. อารมฺมณา\-ธิปติ. (2) (สํขิตฺตํ.)}

\proseitem{47}{เหตุยา เทฺว, อารมฺมเณ เทฺว, อธิปติยา ตีณิ, อนนฺตเร เทฺว, สมนนฺตเร เทฺว, สหชาเต เทฺว ฯเปฯ อุปนิสฺสเย จตฺตาริ, อาเสวเน เทฺว, กมฺเม เทฺว, อาหาเร เทฺว ฯเปฯ อวิคเต เทฺว. (สํขิตฺตํ.)}

\csromanheader{2}{2. \textbf{อกุสลปท 1-7. ปฏิจฺจาทิวาร}}
\prose{ปจฺจยจตุกฺกเหตุ}

\proseitem{48}{โลกิยํ อกุสลํ ธมฺมํ ปฏิจฺจ โลกิโย อกุสโล ธมฺโม อุปฺปชฺชติ เหตุปจฺจยา. (สํขิตฺตํ.)}

\proseitem{49}{เหตุยา เอกํ, อารมฺมเณ เอกํ ฯเปฯ อวิคเต เอกํ. (สํขิตฺตํ.)}

\prose{(สหชาตวาโรปิ ฯเปฯ ปญฺหาวาโรปิ สพฺพตฺถ สทิสา.)}

\csromanheader{2}{3. \textbf{อพฺยากตปท 1-7. ปฏิจฺจาทิวาร}}
\prose{ปจฺจยจตุกฺกเหตุ}

\proseitem{50}{โลกิยํ อพฺยากตํ ธมฺมํ ปฏิจฺจ โลกิโย อพฺยากโต ธมฺโม อุปฺปชฺชติ เหตุปจฺจยา. (1)}

\csromanmarkh{12. โลกิยทุก}\csromanmarkhii{1. กุสลตฺติก}\csromanheadernomark{2}{12. โลกิยทุก 1. กุสลตฺติก}
\prose{\csromanfolio{357}โลกุตฺตรํ อพฺยากตํ ธมฺมํ ปฏิจฺจ โลกุตฺตโร อพฺยากโต ธมฺโม อุปฺปชฺชติ เหตุปจฺจยา. ตีณิ.}

\prose{โลกิยํ อพฺยากตญฺจ โลกุตฺตรํ อพฺยากตญฺจ ธมฺมํ ปฏิจฺจ โลกิโย อพฺยากโต ธมฺโม อุปฺปชฺชติ เหตุปจฺจยา. (1) (สํขิตฺตํ.)}

\proseitem{51}{เหตุยา ปญฺจ, อารมฺมเณ เทฺว ฯเปฯ อาเสวเน เอกํ, กมฺเม ปญฺจ, วิปาเก ปญฺจ ฯเปฯ อวิคเต ปญฺจ. (สํขิตฺตํ. อนุโลมํ.)}

\proseitem{52}{นเหตุยา เอกํ, นอารมฺมเณ ตีณิ, นอธิปติยา เทฺว ฯเปฯ นปุเรชาเต จตฺตาริ, นปจฺฉาชาเต ปญฺจ ฯเปฯ นกมฺเม เอกํ, นวิปาเก เอกํ, นอาหาเร เอกํ ฯเปฯ นมคฺเค เอกํ, นสมฺปยุตฺเต ตีณิ, นวิปฺปยุตฺเต เทฺว, โนนตฺถิยา ตีณิ, โนวิคเต ตีณิ. (สํขิตฺตํ. ปจฺจนียํ.)}

\prose{(สหชาตวาโรปิ ปจฺจยวาโรปิ นิสฺสยวาโรปิ สํสฏฺฐ\-วาโรปิ สมฺปยุตฺตวาโรปิ ปฏิจฺจ\-วาร\-สทิสา วิตฺถาเรตพฺพา.)}

\prose{ปจฺจยจตุกฺกเหตุ}

\proseitem{53}{โลกิโย อพฺยากโต ธมฺโม โลกิยสฺส อพฺยากตสฺส ธมฺมสฺส เหตุปจฺจเยน ปจฺจโย. (1)}

\prose{โลกุตฺตโร อพฺยากโต ธมฺโม โลกุตฺตรสฺส อพฺยากตสฺส ธมฺมสฺส เหตุปจฺจเยน ปจฺจโย. ตีณิ. (สํขิตฺตํ.)}

\proseitem{54}{เหตุยา จตฺตาริ, อารมฺมเณ ตีณิ, อธิปติยา จตฺตาริ, อนนฺตเร จตฺตาริ, สหชาเต ปญฺจ, อญฺญมญฺเญ เทฺว, นิสฺสเย สตฺต, อุปนิสฺสเย จตฺตาริ, ปุเรชาเต เทฺว, ปจฺฉาชาเต เทฺว, อาเสวเน เอกํ, กมฺเม จตฺตาริ, วิปาเก จตฺตาริ, อาหาเร จตฺตาริ ฯเปฯ มคฺเค จตฺตาริ, สมฺปยุตฺเต เทฺว, วิปฺปยุตฺเต ตีณิ, อตฺถิยา สตฺต ฯเปฯ อวิคเต สตฺต. (สํขิตฺตํ.)}

\prosewithcloser{(ยถา กุสลตฺติเก ปญฺหาวารสฺส อนุโลมมฺปิ ปจฺจนียมฺปิ อนุโลมปจฺจนียมฺปิ ปจฺจนียานุโลมมฺปิ คณิตํ, เอวํ คเณตพฺพํ.)}{โลกิย\-ทุก\-กุสล\-ตฺติกํ นิฏฺฐิตํ.}
\csromanmarkh{ธมฺมานุโลม ทุกติก\-ปฏฺฐาน\-ปาฬิ}\csromanmarkhii{13. เกนจิ\-วิญฺเญยฺยทุก 1. กุสลตฺติก}\csromanheadernomark{2}{ธมฺมานุโลม ทุกติก\-ปฏฺฐาน\-ปาฬิ 13. เกนจิ\-วิญฺเญยฺยทุก 1. กุสลตฺติก}
\csromanheader{2}{1. \textbf{กุสลปท 1-7. ปฏิจฺจาทิวาร}}
\prose{\csromanfolio{358}ปจฺจยจตุกฺกเหตุ}

\proseitem{55}{เกนจิ วิญฺเญยฺยํ กุสลํ ธมฺมํ ปฏิจฺจ เกนจิ วิญฺเญยฺโย กุสโล ธมฺโม อุปฺปชฺชติ เหตุปจฺจยา.  เกนจิ วิญฺเญยฺยํ กุสลํ ธมฺมํ ปฏิจฺจ เกนจิ นวิญฺเญยฺโย กุสโล ธมฺโม อุปฺปชฺชติ เหตุปจฺจยา.  เกนจิ วิญฺเญยฺยํ กุสลํ ธมฺมํ ปฏิจฺจ เกนจิ วิญฺเญยฺโย กุสโล จ เกนจิ นวิญฺเญยฺโย กุสโล จ ธมฺมา อุปฺปชฺชนฺติ เหตุปจฺจยา. (3) (สํขิตฺตํ.)}

\proseitem{56}{เหตุยา นว, อารมฺมเณ นว ฯเปฯ อวิคเต นว. (สํขิตฺตํ.)}

\prose{ปจฺจนียนอธิปติ}

\proseitem{57}{เกนจิ วิญฺเญยฺยํ กุสลํ ธมฺมํ ปฏิจฺจ เกนจิ วิญฺเญยฺโย กุสโล ธมฺโม อุปฺปชฺชติ นอธิปติ\-ปจฺจยา. (สํขิตฺตํ.)}

\proseitem{58}{นอธิปติยา นว, นปุเรชาเต นว, นปจฺฉาชาเต นว, นอาเสวเน นว, นกมฺเม นว, นวิปาเก นว ฯเปฯ นวิปฺปยุตฺเต นว. (สํขิตฺตํ.)}

\prose{(สหชาตวาโรปิ ปจฺจยวาโรปิ นิสฺสยวาโรปิ สํสฏฺฐ\-วาโรปิ สมฺปยุตฺตวาโรปิ ปฏิจฺจ\-วาร\-สทิสา วิตฺถาเรตพฺพา.)}

\prose{ปจฺจยจตุกฺกเหตุ}

\proseitem{59}{เกนจิ วิญฺเญยฺโย กุสโล ธมฺโม เกนจิ วิญฺเญยฺยสฺส กุสลสฺส ธมฺมสฺส เหตุปจฺจเยน ปจฺจโย. (สํขิตฺตํ.)}

\proseitem{60}{เหตุยา นว, อารมฺมเณ นว, อธิปติยา นว, อนนฺตเร นว, สมนนฺตเร นว ฯเปฯ กมฺเม นว ฯเปฯ อวิคเต นว. (สํขิตฺตํ.)}

\prose{นเหตุยา นว, นอารมฺมเณ นว. (สํขิตฺตํ.)}

\prose{เหตุปจฺจยา นอารมฺมเณ นว. (สํขิตฺตํ.)}

\prose{นเหตุปจฺจยา อารมฺมเณ นว. (สํขิตฺตํ.)}

\csromanmatikaanchor{mtk.65}
\csromanmatikaanchor{mtk.66}
\csromantocmarknopage{subparagraph}{14. อาสวทุก}{mtk.65}
\bookmark[dest={mtk.65},level=5]{14. อาสวทุก}
\csromantocmarkref{subparagraph}{1. กุสลตฺติก ...}{mtk.66}
\bookmark[dest={mtk.66},level=5]{1. กุสลตฺติก ...}
\csromanmarkh{14. อาสวทุก}\csromanmarkhii{1. กุสลตฺติก}\csromanheadernomark{6}{14. อาสวทุก 1. กุสลตฺติก}
\prosewithcloser{\csromanfolio{359}(ยถา กุสลตฺติเก ปญฺหาวารสฺส อนุโลมมฺปิ ปจฺจนียมฺปิ อนุโลมปจฺจนียมฺปิ ปจฺจนียานุโลมมฺปิ คณิตํ, เอวํ คเณตพฺพํ. เกนจิ วิญฺเญยฺยํ อกุสลมฺปิ เกนจิ วิญฺเญยฺยํ อพฺยากตมฺปิ เกนจิ วิญฺเญยฺย\-กุสล\-สทิสํ วิตฺถาเรตพฺพํ.)}{เกนจิ วิญฺเญยฺยทุกกุสลตฺติกํ นิฏฺฐิตํ.}
\nitthitam{จูฬนฺตรทุกํ นิฏฺฐิตํ.}
\csromanmarkh{14. อาสวทุก}\csromanmarkhii{1. กุสลตฺติก}\csromanheadernomark{2}{14. อาสวทุก 1. กุสลตฺติก}
\csromanheader{2}{1. \textbf{กุสลปท 1-7. ปฏิจฺจา}ทิวารเหตุ}
\proseitem{1}{โนอสวํ กุสลํ ธมฺมํ ปฏิจฺจ โนอสโว กุสโล ธมฺโม อุปฺปชฺชติ เหตุปจฺจยา. (สํขิตฺตํ.)}

\proseitem{2}{เหตุยา เอกํ, อารมฺมเณ เอกํ ฯเปฯ กมฺเม เอกํ ฯเปฯ อวิคเต เอกํ. (สํขิตฺตํ.)}

\prose{(สหชาตวาเรปิ ฯเปฯ ปญฺหาวาเรปิ สพฺพตฺถ เอกํ.)}

\prose{2. อกุสลปทเหตุ}

\proseitem{3}{อาสวํ อกุสลํ ธมฺมํ ปฏิจฺจ อาสโว อกุสโล ธมฺโม อุปฺปชฺชติ เหตุปจฺจยา. (สํขิตฺตํ.)}

\proseitem{4}{เหตุยา นว, อารมฺมเณ นว ฯเปฯ อวิคเต นว. (สํขิตฺตํ.)}

\prose{นเหตุยา เอกํ, นอธิปติยา นว ฯเปฯ นกมฺเม ตีณิ, นวิปาเก นว, นวิปฺปยุตฺเต นว. (สํขิตฺตํ.)}

\prose{(สหชาตวาโรปิ ฯเปฯ สมฺปยุตฺตวาโรปิ สพฺพตฺถ วิตฺถาเรตพฺโพ.)}

\proseitem{5}{อาสโว อกุสโล ธมฺโม อาสวสฺส อกุสลสฺส ธมฺมสฺส เหตุปจฺจเยน ปจฺจโย.}

\proseitem{6}{เหตุยา สตฺต, อารมฺมเณ นว, อธิปติยา นว, (มชฺเฌ ติณฺณํ สหชาตา\-ธิปติ ลพฺภติ.) อนนฺตเร นว ฯเปฯ อุปนิสฺสเย นว, \csromanfolio{360}อาเสวเน นว, กมฺเม ตีณิ ฯเปฯ ฌาเน ตีณิ, มคฺเค นว ฯเปฯ อวิคเต นว. (สํขิตฺตํ.)}

\prose{3. อพฺยากตปทเหตุ}

\proseitem{7}{โนอสวํ อพฺยากตํ ธมฺมํ ปฏิจฺจ โนอสโว อพฺยากโต ธมฺโม อุปฺปชฺชติ เหตุปจฺจยา. (สํขิตฺตํ.)}

\proseitem{8}{เหตุยา เอกํ, อารมฺมเณ เอกํ ฯเปฯ อวิคเต เอกํ. (สํขิตฺตํ.)}

\prosewithcloser{(สหชาตวาเรปิ ปจฺจยวาเรปิ ฯเปฯ ปญฺหาวาเรปิ สพฺพตฺถ เอกํ.)}{อาสว\-ทุก\-กุสล\-ตฺติกํ นิฏฺฐิตํ.}
\csromanmarkh{15. สาสวทุก}\csromanmarkhii{1. กุสลตฺติก}\csromanheadernomark{2}{15. \textbf{สาสวทุก 1. กุสลตฺติก}}
\prose{1. กุสลปทเหตุ}

\proseitem{9}{สาสวํ กุสลํ ธมฺมํ ปฏิจฺจ สาสโว กุสโล ธมฺโม อุปฺปชฺชติ เหตุปจฺจยา. (1)}

\prose{อนาสวํ กุสลํ ธมฺมํ ปฏิจฺจ อนาสโว กุสโล ธมฺโม อุปฺปชฺชติ เหตุปจฺจยา. (1) (สํขิตฺตํ.)}

\proseitem{6}{เหตุยา เทฺว, อารมฺมเณ เทฺว, อธิปติยา เทฺว ฯเปฯ อวิคเต เทฺว. (สํขิตฺตํ. อนุโลมํ.)}

\proseitem{11}{นอธิปติยา เทฺว ฯเปฯ นอาเสวเน เอกํ ฯเปฯ นวิปฺปยุตฺเต เทฺว. (สํขิตฺตํ. ปจฺจนียํ.) (สหชาตวาเรปิ ฯเปฯ สมฺปยุตฺตวาเรปิ สพฺพตฺถ เทฺว.)}

\prose{ปญฺหาวารเหตุ}

\proseitem{12}{สาสโว กุสโล ธมฺโม สาสวสฺส กุสลสฺส ธมฺมสฺส เหตุปจฺจเยน ปจฺจโย. (1)}

\prose{อนาสโว กุสโล ธมฺโม อนาสวสฺส กุสลสฺส ธมฺมสฺส เหตุปจฺจเยน ปจฺจโย. (1) (สํขิตฺตํ.)}

\csromanmarkh{15. สาสวทุก}\csromanmarkhii{1. กุสลตฺติก}\csromanheadernomark{2}{15. สาสวทุก 1. กุสลตฺติก}
\proseitem{13}{\csromanfolio{361}เหตุยา เทฺว, อารมฺมเณ เทฺว, อธิปติยา ตีณิ, อนนฺตเร เทฺว ฯเปฯ นิสฺสเย เทฺว, อุปนิสฺสเย จตฺตาริ, อาเสวเน เทฺว ฯเปฯ อวิคเต เทฺว. (สํขิตฺตํ.)}

\prose{(ยถา กุสลตฺติเก ปญฺหาวารสฺส อนุโลมมฺปิ ปจฺจนียมฺปิ อนุโลมปจฺจนียมฺปิ ปจฺจนียานุโลมมฺปิ คณิตํ, เอวํ คเณตพฺพํ.)}

\prose{2. อกุสลปทเหตุ}

\proseitem{14}{สาสวํ อกุสลํ ธมฺมํ ปฏิจฺจ สาสโว อกุสโล ธมฺโม อุปฺปชฺชติ เหตุปจฺจยา. (สํขิตฺตํ.)}

\proseitem{15}{เหตุยา เอกํ, อารมฺมเณ เอกํ ฯเปฯ อวิคเต เอกํ. (สํขิตฺตํ.)}

\prose{(สหชาตวาเรปิ ฯเปฯ ปญฺหาวาเรปิ สพฺพตฺถ เอกํ.)}

\prose{3. อพฺยากตปทเหตุ}

\proseitem{16}{สาสวํ อพฺยากตํ ธมฺมํ ปฏิจฺจ สาสโว อพฺยากโต ธมฺโม อุปฺปชฺชติ เหตุปจฺจยา. (1)}

\prose{อนาสวํ อพฺยากตํ ธมฺมํ ปฏิจฺจ อนาสโว อพฺยากโต ธมฺโม อุปฺปชฺชติ เหตุปจฺจยา.  อนาสวํ อพฺยากตํ ธมฺมํ ปฏิจฺจ สาสโว อพฺยากโต ธมฺโม อุปฺปชฺชติ เหตุปจฺจยา.  อนาสวํ อพฺยากตํ ธมฺมํ ปฏิจฺจ สาสโว อพฺยากโต จ อนาสโว อพฺยากโต จ ธมฺมา อุปฺปชฺชนฺติ เหตุปจฺจยา. (3)}

\prose{สาสวํ อพฺยากตญฺจ อนาสวํ อพฺยากตญฺจ ธมฺมํ ปฏิจฺจ สาสโว อพฺยากโต ธมฺโม อุปฺปชฺชติ เหตุปจฺจยา. (1) (สํขิตฺตํ.)}

\proseitem{17}{เหตุยา ปญฺจ, อารมฺมเณ เทฺว, อธิปติยา ปญฺจ ฯเปฯ อาเสวเน เอกํ ฯเปฯ วิปาเก ปญฺจ ฯเปฯ อวิคเต ปญฺจ. (สํขิตฺตํ.)}

\prose{นเหตุยา เอกํ, นอารมฺมเณ ตีณิ, นอธิปติยา เทฺว, นปุเรชาเต จตฺตาริ, นปจฺฉาชาเต นอาเสวเน ปญฺจ, นกมฺเม ฯเปฯ นมคฺเค เอกํ ฯเปฯ นวิปฺปยุตฺเต เทฺว ฯเปฯ โนวิคเต ตีณิ. (สํขิตฺตํ.)}

\prose{(สหชาตวาโรปิ ฯเปฯ สมฺปยุตฺตวาโรปิ สพฺพตฺถ วิตฺถาเรตพฺโพ.)}

\csromanheader{2}{ธมฺมานุโลม ทุกติก\-ปฏฺฐาน\-ปาฬิ ปญฺหาวารเหตุ}
\proseitem{18}{\csromanfolio{362}สาสโว อพฺยากโต ธมฺโม สาสวสฺส อพฺยากตสฺส ธมฺมสฺส เหตุปจฺจเยน ปจฺจโย. (1)}

\prose{อนาสโว อพฺยากโต ธมฺโม อนาสวสฺส อพฺยากตสฺส ธมฺมสฺส เหตุปจฺจเยน ปจฺจโย. (1) (สํขิตฺตํ.)}

\proseitem{19}{เหตุยา จตฺตาริ, อารมฺมเณ ตีณิ, อธิปติยา อนนฺตเร สมนนฺตเร จตฺตาริ, สหชาเต ปญฺจ, อญฺญมญฺเญ เทฺว, นิสฺสเย สตฺต, อุปนิสฺสเย สตฺตาริ, ปุเรชาเต ปจฺฉาชาเต เทฺว, อาเสวเน เอกํ, กมฺเม ฯเปฯ มคฺเค จตฺตาริ, สมฺปยุตฺเต เทฺว, วิปฺปยุตฺเต ตีณิ ฯเปฯ อวิคเต สตฺต. (สํขิตฺตํ.)}

\prosewithcloser{(ยถา กุสลตฺติเก ปญฺหาวารสฺส อนุโลมมฺปิ ปจฺจนียมฺปิ อนุโลมปจฺจนียมฺปิ ปจฺจนียานุโลมมฺปิ คณิตํ, เอวํ คเณตพฺพํ.)}{สาสว\-ทุก\-กุสล\-ตฺติกํ นิฏฺฐิตํ.}
\csromanmarkh{16. อาสว\-สมฺปยุตฺตทุก}\csromanmarkhii{1. กุสลตฺติก}\csromanheadernomark{2}{16. อาสว\-สมฺปยุตฺตทุก 1. กุสลตฺติก}
\csromanheader{2}{1. \textbf{กุสลปท 1-7. ปฏิจฺจา}ทิวารเหตุ}
\proseitem{20}{อาสว\-วิปฺปยุตฺตํ กุสลํ ธมฺมํ ปฏิจฺจ อาสว\-วิปฺปยุตฺโต กุสโล ธมฺโม อุปฺปชฺชติ เหตุปจฺจยา. (สํขิตฺตํ.)}

\proseitem{21}{เหตุยา เอกํ, อารมฺมเณ เอกํ ฯเปฯ อวิคเต เอกํ. (สํขิตฺตํ. สหชาตวาเรปิ ฯเปฯ ปญฺหาวาเรปิ สพฺพตฺถ เอกํ.)}

\prose{2. อกุสลปทเหตุ}

\proseitem{22}{อาสว\-สมฺปยุตฺตํ อกุสลํ ธมฺมํ ปฏิจฺจ อาสว\-สมฺปยุตฺโต อกุสโล ธมฺโม อุปฺปชฺชติ เหตุปจฺจยา. ตีณิ.}

\prose{อาสว\-วิปฺปยุตฺตํ อกุสลํ ธมฺมํ ปฏิจฺจ อาสว\-สมฺปยุตฺโต อกุสโล ธมฺโม อุปฺปชฺชติ เหตุปจฺจยา. (1)}

\csromanmarkh{16. อาสว\-สมฺปยุตฺตทุก}\csromanmarkhii{1. กุสลตฺติก}\csromanheadernomark{2}{16. อาสว\-สมฺปยุตฺตทุก 1. กุสลตฺติก}
\prose{\csromanfolio{363}อาสว\-สมฺปยุตฺตํ อกุสลญฺจ อาสว\-วิปฺปยุตฺตํ อกุสลญฺจ ธมฺมํ ปฏิจฺจ อาสว\-สมฺปยุตฺโต อกุสโล ธมฺโม อุปฺปชฺชติ เหตุปจฺจยา. (1) (สํขิตฺตํ.)}

\proseitem{23}{เหตุยา ปญฺจ, อารมฺมเณ ปญฺจ ฯเปฯ อวิคเต ปญฺจ. (สํขิตฺตํ.)}

\prose{(สหชาตวาเรปิ ฯเปฯ สมฺปยุตฺตวาเรปิ สพฺพตฺถ ปญฺจ.)}

\prose{นเหตุยา เอกํ, นอธิปติยา ปญฺจ, นปุเรชาเต ปญฺจ ฯเปฯ นกมฺเม ตีณิ, นวิปาเก ปญฺจ, นวิปฺปยุตฺเต ปญฺจ. (สํขิตฺตํ.)}

\proseitem{24}{อาสว\-สมฺปยุตฺโต อกุสโล ธมฺโม อาสว\-สมฺปยุตฺตสฺส อกุสลสฺส ธมฺมสฺส เหตุปจฺจเยน ปจฺจโย. ตีณิ.}

\prose{อาสว\-วิปฺปยุตฺโต อกุสโล ธมฺโม อาสว\-สมฺปยุตฺตสฺส อกุสลสฺส ธมฺมสฺส เหตุปจฺจเยน ปจฺจโย. (1)}

\prose{อาสว\-สมฺปยุตฺโต อกุสโล จ อาสว\-วิปฺปยุตฺโต อกุสโล จ ธมฺมา อาสว\-สมฺปยุตฺตสฺส อกุสลสฺส ธมฺมสฺส เหตุปจฺจเยน ปจฺจโย. (1) (สํขิตฺตํ.)}

\proseitem{25}{เหตุยา ปญฺจ, อารมฺมเณ นว, อธิปติยา เอกํ, อนนฺตเร สมนนฺตเร นว, สหชาเต อญฺญมญฺเญ นิสฺสเย ปญฺจ, อุปนิสฺสเย อาเสวเน นว, กมฺเม อาหาเร อินฺทฺริเย ฌาเน มคฺเค ตีณิ, สมฺปยุตฺเต ปญฺจ ฯเปฯ อวิคเต ปญฺจ. (สํขิตฺตํ.)}

\prose{(ยถา กุสลตฺติเก ปญฺหาวารสฺส อนุโลมมฺปิ ปจฺจนียมฺปิ อนุโลมปจฺจนียมฺปิ ปจฺจนียานุโลมมฺปิ คณิตํ, เอวํ คเณตพฺพํ.)}

\prose{3. อพฺยากตปทเหตุ}

\proseitem{26}{อาสว\-วิปฺปยุตฺตํ อพฺยากตํ ธมฺมํ ปฏิจฺจ อาสว\-วิปฺปยุตฺโต อพฺยากโต ธมฺโม อุปฺปชฺชติ เหตุปจฺจยา. (สํขิตฺตํ.)}

\proseitem{27}{เหตุยา เอกํ, อารมฺมเณ เอกํ ฯเปฯ อวิคเต เอกํ. (สํขิตฺตํ.)}

\prose{นเหตุยา เอกํ, นอารมฺมเณ เอกํ, นอธิปติยา เอกํ ฯเปฯ โนวิคเต เอกํ. (สํขิตฺตํ. สหชาตวาเรปิ ฯเปฯ สมฺปยุตฺตวาเรปิ สพฺพตฺถ เอกํ.)}

\proseitem{28}{\csromanfolio{364}อาสว\-วิปฺปยุตฺโต อพฺยากโต ธมฺโม อาสว\-วิปฺปยุตฺตสฺส อพฺยากตสฺส ธมฺมสฺส เหตุปจฺจเยน ปจฺจโย. (สํขิตฺตํ.)}

\proseitem{29}{เหตุยา เอกํ, อารมฺมเณ เอกํ ฯเปฯ อวิคเต เอกํ. (สํขิตฺตํ.)}

\prosewithcloser{(ยถา กุสลตฺติเก ปญฺหาวารสฺส อนุโลมมฺปิ ปจฺจนียมฺปิ อนุโลมปจฺจนียมฺปิ ปจฺจนียานุโลมมฺปิ คณิตํ, เอวํ คเณตพฺพํ.)}{อาสว\-สมฺปยุตฺต\-ทุก\-กุสล\-ตฺติกํ นิฏฺฐิตํ.}
\csromanmarkh{17. อาสว\-สาสวทุก}\csromanmarkhii{1. กุสลตฺติก}\csromanheadernomark{2}{17. อาสว\-สาสวทุก 1. กุสลตฺติก}
\csromanheader{2}{1. \textbf{กุสลปท 1-7. ปฏิจฺจา}ทิวารเหตุ}
\proseitem{30}{สาสวญฺเจว โน จ อาสวํ กุสลํ ธมฺมํ ปฏิจฺจ สาสโว เจว โน จ อาสโว กุสโล ธมฺโม อุปฺปชฺชติ เหตุปจฺจยา. (สํขิตฺตํ.)}

\proseitem{31}{เหตุยา เอกํ, อารมฺมเณ เอกํ ฯเปฯ อวิคเต เอกํ. (สํขิตฺตํ.)}

\prose{(สหชาตวาเรปิ ฯเปฯ ปญฺหาวาเรปิ สพฺพตฺถ เอกํ.)}

\prose{2. อกุสลปทเหตุ}

\proseitem{32}{อาสวญฺเจว สาสวญฺจ อกุสลํ ธมฺมํ ปฏิจฺจ อาสโว เจว สาสโว จ อกุสโล ธมฺโม อุปฺปชฺชติ เหตุปจฺจยา. ตีณิ.}

\prose{สาสวญฺเจว โน จ อาสวํ อกุสลํ ธมฺมํ ปฏิจฺจ สาสโว เจว โน จ อาสโว อกุสโล ธมฺโม อุปฺปชฺชติ เหตุปจฺจยา. ตีณิ.}

\prose{อาสวญฺเจว สาสวญฺจ อกุสลญฺจ สาสวญฺเจว โน จ อาสวํ อกุสลญฺจ ธมฺมํ ปฏิจฺจ อาสโว เจว สาสโว จ อกุสโล ธมฺโม อุปฺปชฺชติ เหตุปจฺจยา. ตีณิ. (สํขิตฺตํ.)}

\proseitem{33}{เหตุยา นว, อารมฺมเณ นว ฯเปฯ อวิคเต นว. (สํขิตฺตํ.)}

\prose{ปจฺจนีย\-น\-เหตุ}

\proseitem{34}{สาสวญฺเจว โน จ อาสวํ อกุสลํ ธมฺมํ ปฏิจฺจ อาสโว เจว สาสโว จ อกุสโล ธมฺโม อุปฺปชฺชติ นเหตุปจฺจยา. (สํขิตฺตํ.)}

\csromanmarkh{17. อาสว\-สาสวทุก}\csromanmarkhii{1. กุสลตฺติก}\csromanheadernomark{2}{17. อาสว\-สาสวทุก 1. กุสลตฺติก}
\proseitem{35}{\csromanfolio{365}นเหตุยา เอกํ, นอธิปติยา นว, นปุเรชาเต นว, นปจฺฉาชาเต นว, นอาเสวเน นว, นกมฺเม ตีณิ, นวิปาเก นว, นวิปฺปยุตฺเต นว. (สํขิตฺตํ.)}

\prose{(สหชาตวาโรปิ ปจฺจยวาโรปิ นิสฺสยวาโรปิ สํสฏฺฐ\-วาโรปิ สมฺปยุตฺตวาโรปิ วิตฺถาเรตพฺโพ.)}

\prose{7. ปญฺหาวารเหตุ}

\proseitem{36}{อาสโว เจว สาสโวจ อกุสโล ธมฺโม อาสวสฺส เจว สาสวสฺส จ อกุสลสฺส ธมฺมสฺส เหตุปจฺจเยน ปจฺจโย. (สํขิตฺตํ.)}

\proseitem{37}{เหตุยา สตฺต, อารมฺมเณ นว, อธิปติยา นว ฯเปฯ อุปนิสฺสเย อาเสวเน นว, กมฺเม อาหาเร อินฺทฺริเย ฌาเน ตีณิ, มคฺเค สมฺปยุตฺเต นว ฯเปฯ อวิคเต นว. (สํขิตฺตํ.)}

\proseitem{38}{อาสโว เจว สาสโว จ อกุสโล ธมฺโม อาสวสฺส เจว สาสวสฺส จ อกุสลสฺส ธมฺมสฺส อารมฺมณ\-ปจฺจเยน ปจฺจโย. สหชาต\-ปจฺจเยน ปจฺจโย. อุปนิสฺสย\-ปจฺจเยน ปจฺจโย. (สํขิตฺตํ.)}

\proseitem{39}{นเหตุยา นว, นอารมฺมเณ นว. (สํขิตฺตํ.)}

\prose{เหตุปจฺจยา นอารมฺมเณ สตฺต. (สํขิตฺตํ.)}

\prose{นเหตุปจฺจยา อารมฺมเณ นว. (สํขิตฺตํ.)}

\prose{(ยถา กุสลตฺติเก ปญฺหาวารสฺส อนุโลมมฺปิ ปจฺจนียมฺปิ อนุโลมปจฺจนียมฺปิ ปจฺจนียานุโลมมฺปิ คณิตํ, เอวํ คเณตพฺพํ.)}

\prose{3. อพฺยากตปทเหตุ}

\proseitem{40}{สาสวญฺเจว โน จ อาสวํ อพฺยากตํ ธมฺมํ ปฏิจฺจ สาสโว เจว โน จ อาสโว อพฺยากโต ธมฺโม อุปฺปชฺชติ เหตุปจฺจยา. (สํขิตฺตํ.)}

\proseitem{41}{เหตุยา เอกํ, อารมฺมเณ เอกํ ฯเปฯ อวิคเต เอกํ. (สํขิตฺตํ.)}

\prosewithcloser{(สหชาตวาเรปิ ฯเปฯ ปญฺหาวาเรปิ สพฺพตฺถ เอกํ.)}{อาสว\-สาสว\-ทุก\-กุสล\-ตฺติกํ นิฏฺฐิตํ.}
\csromanmarkh{ธมฺมานุโลม ทุกติก\-ปฏฺฐาน\-ปาฬิ}\csromanmarkhii{18. อาสว\-อาสว\-สมฺปยุตฺตทุก 1. กุสลตฺติก}\csromanheadernomark{2}{ธมฺมานุโลม ทุกติก\-ปฏฺฐาน\-ปาฬิ 18. อาสว\-อาสว\-สมฺปยุตฺตทุก 1. กุสลตฺติก}
\prose{\csromanfolio{366}1. ปฏิจฺจวารอกุสลปทเหตุ}

\proseitem{42}{อาสวญฺเจว อาสว\-สมฺปยุตฺตญฺจ อกุสลํ ธมฺมํ ปฏิจฺจ อาสโว เจว อาสว\-สมฺปยุตฺโต จ อกุสโล ธมฺโม อุปฺปชฺชติ เหตุปจฺจยา. ตีณิ.}

\prose{อาสว\-สมฺปยุตฺต\-ญฺเจว โน จ อาสวํ อกุสลํ ธมฺมํ ปฏิจฺจ อาสว\-สมฺปยุตฺโต เจว โน จ อาสโว อกุสโล ธมฺโม อุปฺปชฺชติ เหตุปจฺจยา. ตีณิ.}

\prose{อาสวญฺเจว อาสว\-สมฺปยุตฺตญฺจ อกุสลญฺจ อาสว\-สมฺปยุตฺต\-ญฺเจว โน จ อาสวํ อกุสลญฺจ ธมฺมํ ปฏิจฺจ อาสโว เจว อาสว\-สมฺปยุตฺโต จ อกุสโล ธมฺโม อุปฺปชฺชติ เหตุปจฺจยา. ตีณิ. (สํขิตฺตํ.)}

\proseitem{43}{เหตุยา นว, อารมฺมเณ นว ฯเปฯ อวิคเต นว. (สํขิตฺตํ.)}

\prose{ปจฺจนียนอธิปติ}

\proseitem{44}{อาสวญฺเจว อาสว\-สมฺปยุตฺตญฺจ อกุสลํ ธมฺมํ ปฏิจฺจ อาสโว เจว อาสว\-สมฺปยุตฺโต จ อกุสโล ธมฺโม อุปฺปชฺชติ นอธิปติ\-ปจฺจยา.}

\proseitem{45}{นอธิปติยา นว, นปุเรชาเต นว, นปจฺฉาชาเต นว, นอาเสวเน นว, นกมฺเม ตีณิ, นวิปาเก นว, นวิปฺปยุตฺเต นว. (สํขิตฺตํ.)}

\prose{(สหชาตวาเรปิ สมฺปยุตฺตวาเรปิ สพฺพตฺถ นว.)}

\prose{7. ปญฺหาวาร ปจฺจยจตุกฺกเหตุ}

\proseitem{46}{อาสโว เจว อาสว\-สมฺปยุตฺโต จ อกุสโล ธมฺโม อาสวสฺส เจว อาสว\-สมฺปยุตฺตสฺส จ อกุสลสฺส ธมฺมสฺส เหตุปจฺจเยน ปจฺจโย. (สํขิตฺตํ.)}

\proseitem{47}{เหตุยา จตฺตาริ, อารมฺมเณ นว, อธิปติยา นว ฯเปฯ อวิคเต นว. (สํขิตฺตํ.)}

\csromanmarkh{19. อาสว\-วิปฺปยุตฺต\-สาสวทุก}\csromanmarkhii{1. กุสลตฺติก}\csromanheadernomark{2}{19. อาสว\-วิปฺปยุตฺต\-สาสวทุก 1. กุสลตฺติก}
\prose{\csromanfolio{367}นเหตุยา นว, นอารมฺมเณ นว. (สํขิตฺตํ.)}

\prose{เหตุปจฺจยา นอารมฺมเณ จตฺตาริ. (สํขิตฺตํ.)}

\prose{นเหตุปจฺจยา อารมฺมเณ นว. (สํขิตฺตํ.)}

\prosewithcloser{(ยถา กุสลตฺติเก ปญฺหาวารสฺส อนุโลมมฺปิ ปจฺจนียมฺปิ อนุโลมปจฺจนียมฺปิ ปจฺจนียานุโลมมฺปิ คณิตํ, เอวํ คเณตพฺพํ.)}{อาสว\-อาสว\-สมฺปยุตฺตทุกํ นิฏฺฐิตํ.}
\csromanmarkh{19. อาสว\-วิปฺปยุตฺต\-สาสวทุก}\csromanmarkhii{1. กุสลตฺติก}\csromanheadernomark{2}{19. อาสว\-วิปฺปยุตฺต\-สาสวทุก 1. กุสลตฺติก}
\prose{กุสลปทเหตุ}

\proseitem{48}{อาสว\-วิปฺปยุตฺตํ สาสวํ กุสลํ ธมฺมํ ปฏิจฺจ อาสว\-วิปฺปยุตฺโต สาสโว กุสโล ธมฺโม อุปฺปชฺชติ เหตุปจฺจยา.  อาสว\-วิปฺปยุตฺตํ อนาสวํ กุสลํ ธมฺมํ ปฏิจฺจ อาสว\-วิปฺปยุตฺโต อนาสโว กุสโล ธมฺโม อุปฺปชฺชติ เหตุปจฺจยา. (2) (สํขิตฺตํ.)}

\proseitem{49}{เหตุยา เทฺว, อารมฺมเณ เทฺว ฯเปฯ อวิคเต เทฺว. (สํขิตฺตํ.)}

\prose{(สหชาตวาโรปิ ฯเปฯ ปญฺหาวาโรปิ วิตฺถาเรตพฺโพ.)}

\prose{อพฺยากตปทเหตุ}

\proseitem{50}{อาสว\-วิปฺปยุตฺตํ สาสวํ อพฺยากตํ ธมฺมํ ปฏิจฺจ อาสว\-วิปฺปยุตฺโต สาสโว อพฺยากโต ธมฺโม อุปฺปชฺชติ เหตุปจฺจยา. (1) อาสว\-วิปฺปยุตฺตํ อนาสวํ อพฺยากตํ ธมฺมํ ปฏิจฺจ อาสว\-วิปฺปยุตฺโต อนาสโว อพฺยากโต ธมฺโม อุปฺปชฺชติ เหตุปจฺจยา.  อาสว\-วิปฺปยุตฺตํ อนาสวํ อพฺยากตํ ธมฺมํ ปฏิจฺจ อาสว\-วิปฺปยุตฺโต สาสโว อพฺยากโต ธมฺโม อุปฺปชฺชติ เหตุปจฺจยา.  อาสว\-วิปฺปยุตฺตํ อนาสวํ อพฺยากตํ ธมฺมํ ปฏิจฺจ อาสว\-วิปฺปยุตฺโต สาสโว อพฺยากโต จ อาสว\-วิปฺปยุตฺโต อนาสโว อพฺยากโต จ ธมฺมา อุปฺปชฺชนฺติ เหตุปจฺจยา. (3)}

\prose{อาสว\-วิปฺปยุตฺตํ สาสวํ อพฺยากตญฺจ อาสว\-วิปฺปยุตฺตํ อนาสวํ อพฺยากตญฺจ ธมฺมํ ปฏิจฺจ อาสว\-วิปฺปยุตฺโต สาสโว อพฺยากโต ธมฺโม อุปฺปชฺชติ เหตุปจฺจยา. (1) (สํขิตฺตํ.)}

\proseitem{51}{\csromanfolio{368}เหตุยา ปญฺจ, อารมฺมเณ เทฺว, อธิปติยา ปญฺจ ฯเปฯ อาเสวเน เอกํ, กมฺเม ปญฺจ, วิปาเก ปญฺจ ฯเปฯ อวิคเต ปญฺจ. (สํขิตฺตํ.)}

\prose{ปจฺจนีย\-น\-เหตุ}

\proseitem{52}{อาสว\-วิปฺปยุตฺตํ สาสวํ อพฺยากตํ ธมฺมํ ปฏิจฺจ อาสว\-วิปฺปยุตฺโต สาสโว อพฺยากโต ธมฺโม อุปฺปชฺชติ นเหตุปจฺจยา. (สํขิตฺตํ.)}

\proseitem{53}{นเหตุยา เอกํ, นอารมฺมเณ ตีณิ, นอธิปติยา เทฺว ฯเปฯ นปุเรชาเต จตฺตาริ, นปจฺฉาชาเต นอาเสวเน ปญฺจ, นกมฺเม นวิปาเก ฯเปฯ นมคฺเค เอกํ ฯเปฯ นวิปฺปยุตฺเต เทฺว ฯเปฯ โนวิคเต ตีณิ. (สํขิตฺตํ.)}

\prose{(สหชาตวาโรปิ ปจฺจยวาโรปิ นิสฺสยวาโรปิ สํสฏฺฐ\-วาโรปิ สมฺปยุตฺตวาโรปิ ปฏิจฺจ\-วาร\-สทิสา วิตฺถาเรตพฺพา.)}

\csromanheader{2}{7. \textbf{ปญฺหาวาร ปจฺจยจตุกฺก}}
\csromanheader{2}{\textbf{เหตุ อารมฺมณ}}
\proseitem{54}{อาสว\-วิปฺปยุตฺโต สาสโว อพฺยากโต ธมฺโม อาสว\-วิปฺปยุตฺตสฺส สาสวสฺส อพฺยากตสฺส ธมฺมสฺส เหตุปจฺจเยน ปจฺจโย. (1) อาสว\-วิปฺปยุตฺโต อนาสโว อพฺยากโต ธมฺโม อาสว\-วิปฺปยุตฺตสฺส อนาสวสฺส อพฺยากตสฺส ธมฺมสฺส เหตุปจฺจเยน ปจฺจโย.  อาสว\-วิปฺปยุตฺโต อนาสโว อพฺยากโต ธมฺโม อาสว\-วิปฺปยุตฺตสฺส สาสวสฺส อพฺยากตสฺส ธมฺมสฺส เหตุปจฺจเยน ปจฺจโย.  อาสว\-วิปฺปยุตฺโต อนาสโว อพฺยากโต ธมฺโม อาสว\-วิปฺปยุตฺตสฺส สาสวสฺส อพฺยากตสฺส จ อาสว\-วิปฺปยุตฺตสฺส อนาสวสฺส อพฺยากตสฺส จ ธมฺมสฺส เหตุปจฺจเยน ปจฺจโย. (3)}

\prose{อาสว\-วิปฺปยุตฺโต สาสโว อพฺยากโต ธมฺโม อาสว\-วิปฺปยุตฺตสฺส สาสวสฺส อพฺยากตสฺส ธมฺมสฺส อารมฺมณ\-ปจฺจเยน ปจฺจโย. (1)}

\csromanmatikaanchor{mtk.67}
\csromanmatikaanchor{mtk.68}
\csromantocmarknopage{subparagraph}{20. สญฺโญชนทุก}{mtk.67}
\bookmark[dest={mtk.67},level=5]{20. สญฺโญชนทุก}
\csromantocmarkref{subparagraph}{1. กุสลตฺติก ...}{mtk.68}
\bookmark[dest={mtk.68},level=5]{1. กุสลตฺติก ...}
\csromanmarkh{20. สญฺโญชนทุก}\csromanmarkhii{1. กุสลตฺติก}\csromanheadernomark{6}{20. สญฺโญชนทุก 1. กุสลตฺติก}
\prose{\csromanfolio{369}อาสว\-วิปฺปยุตฺโต อนาสโว อพฺยากโต ธมฺโม อาสว\-วิปฺปยุตฺตสฺส อนาสวสฺส อพฺยากตสฺส ธมฺมสฺส อารมฺมณ\-ปจฺจเยน ปจฺจโย.  อาสว\-วิปฺปยุตฺโต อนาสโว อพฺยากโต ธมฺโม อาสว\-วิปฺปยุตฺตสฺส สาสวสฺส อพฺยากตสฺส ธมฺมสฺส อารมฺมณ\-ปจฺจเยน ปจฺจโย. (2) (สํขิตฺตํ.)}

\proseitem{55}{เหตุยา จตฺตาริ, อารมฺมเณ ตีณิ, อธิปติยา จตฺตาริ, อนนฺตเร จตฺตาริ, สมนนฺตเร จตฺตาริ, สหชาเต ปญฺจ, อญฺญมญฺเญ เทฺว, นิสฺสเย สตฺต, อุปนิสฺสเย จตฺตาริ, ปุเรชาเต เทฺว, อาเสวเน เอกํ, กมฺเม จตฺตาริ, วิปาเก จตฺตาริ, อาหาเร จตฺตาริ ฯเปฯ อวิคเต สตฺต. (สํขิตฺตํ.)}

\proseitem{56}{อาสว\-วิปฺปยุตฺโต สาสโว อพฺยากโต ธมฺโม อาสว\-วิปฺปยุตฺตสฺส สาสวสฺส อพฺยากตสฺส ธมฺมสฺส อารมฺมณ\-ปจฺจเยน ปจฺจโย. สหชาต\-ปจฺจเยน ปจฺจโย. อุปนิสฺสย\-ปจฺจเยน ปจฺจโย. ปุเรชาต\-ปจฺจเยน ปจฺจโย. ปจฺฉาชาต\-ปจฺจเยน ปจฺจโย. อาหารปจฺจเยน ปจฺจโย. อินฺทฺริย\-ปจฺจเยน ปจฺจโย. (สํขิตฺตํ.)}

\proseitem{57}{นเหตุยา สตฺต, นอารมฺมเณ สตฺต. (สํขิตฺตํ.)}

\prose{เหตุปจฺจยา นอารมฺมเณ จตฺตาริ. (สํขิตฺตํ.)}

\prose{นเหตุปจฺจยา อารมฺมเณ ตีณิ. (สํขิตฺตํ.)}

\prosewithcloser{(ยถา กุสลตฺติเก ปญฺหาวารสฺส อนุโลมมฺปิ ปจฺจนียมฺปิ อนุโลมปจฺจนียมฺปิ ปจฺจนียานุโลมมฺปิ คณิตํ, เอวํ คเณตพฺพํ.)}{อาสว\-วิปฺปยุตฺต\-สาสว\-ทุก\-กุสล\-ตฺติกํ นิฏฺฐิตํ.}
\nitthitam{อาสวโคจฺฉกํ นิฏฺฐิตํ.}
\csromanmarkh{20. สญฺโญชนทุก}\csromanmarkhii{1. กุสลตฺติก}\csromanheadernomark{2}{20. สญฺโญชนทุก 1. กุสลตฺติก}
\prose{1. ปฏิจฺจวาร ปจฺจยจตุกฺกเหตุ}

\proseitem{1}{โนสญฺโญชนํ กุสลํ ธมฺมํ ปฏิจฺจ โนสญฺโญชโน กุสโล ธมฺโม อุปฺปชฺชติ เหตุปจฺจยา. (สํขิตฺตํ.)}

\proseitem{2}{\csromanfolio{370}เหตุยา เอกํ, อารมฺมเณ เอกํ ฯเปฯ อวิคเต เอกํ. (สํขิตฺตํ.) (สหชาตวาเรปิ ฯเปฯ ปญฺหาวาเรปิ สพฺพตฺถ เอกํ.)}

\proseitem{3}{สญฺโญชนํ อกุสลํ ธมฺมํ ปฏิจฺจ สญฺโญชโน อกุสโล ธมฺโม อุปฺปชฺชติ เหตุปจฺจยา.  สญฺโญชนํ อกุสลํ ธมฺมํ ปฏิจฺจ โนสญฺโญชโน อกุสโล ธมฺโม อุปฺปชฺชติ เหตุปจฺจยา.  สญฺโญชนํ อกุสลํ ธมฺมํ ปฏิจฺจ สญฺโญชโน อกุสโล จ โนสญฺโญชโน อกุสโล จ ธมฺมา อุปฺปชฺชนฺติ เหตุปจฺจยา. (3) (สํขิตฺตํ.)}

\proseitem{4}{เหตุยา นว, อารมฺมเณ นว, อธิปติยา นว ฯเปฯ อวิคเต นว. (สํขิตฺตํ.)}

\prose{นเหตุยา ตีณิ, นอธิปติยา นว ฯเปฯ นกมฺเม ตีณิ ฯเปฯ นวิปฺปยุตฺเต นว. (สํขิตฺตํ.)}

\prose{(สหชาตวาโรปิ ฯเปฯ สมฺปยุตฺตวาโรปิ ปฏิจฺจ\-วาร\-สทิสา วิตฺถาเรตพฺพา.)}

\prose{ปญฺหาวาร ปจฺจยจตุกฺกเหตุ}

\proseitem{5}{สญฺโญชโน อกุสโล ธมฺโม สญฺโญชนสฺส อกุสลสฺส ธมฺมสฺส เหตุปจฺจเยน ปจฺจโย. (สํขิตฺตํ.)}

\proseitem{6}{เหตุยา ตีณิ, อารมฺมเณ นว, อธิปติยา นว ฯเปฯ อุปนิสฺสเย อาเสวเน นว, กมฺเม ตีณิ, อาหาเร อินฺทฺริเย ฌาเน ตีณิ, มคฺเค สมฺปยุตฺเต นว ฯเปฯ อวิคเต นว. (สํขิตฺตํ.)}

\proseitem{7}{นเหตุยา นว, นอารมฺมเณ นว. (สํขิตฺตํ.)}

\prose{เหตุปจฺจยา นอารมฺมเณ ตีณิ. (สํขิตฺตํ.)}

\prose{นเหตุปจฺจยา อารมฺมเณ นว. (สํขิตฺตํ.)}

\prose{(ยถา กุสลตฺติเก ปญฺหาวารสฺส อนุโลมมฺปิ ปจฺจนียมฺปิ อนุโลมปจฺจนียมฺปิ ปจฺจนียานุโลมมฺปิ คณิตํ, เอวํ คเณตพฺพํ.)}

\prose{อพฺยากตปทเหตุ}

\proseitem{8}{โนสญฺโญชนํ อพฺยากตํ ธมฺมํ ปฏิจฺจ โนสญฺโญชโน อพฺยากโต ธมฺโม อุปฺปชฺชติ เหตุปจฺจยา. (สํขิตฺตํ.)}

\csromanmarkh{21. สญฺโญชนิยทุก}\csromanmarkhii{1. กุสลตฺติก}\csromanheadernomark{2}{21. สญฺโญชนิยทุก 1. กุสลตฺติก}
\proseitem{9}{\csromanfolio{371}เหตุยา เอกํ, อารมฺมเณ เอกํ ฯเปฯ อวิคเต เอกํ. (สํขิตฺตํ.)}

\vspace{\tipitakaparskip}
\csromancenter{(สหชาตวาเรปิ ฯเปฯ ปญฺหาวาเรปิ สพฺพตฺถ เอกํ.)}
\csromanmarkh{21. สญฺโญชนิยทุก}\csromanmarkhii{1. กุสลตฺติก}\csromanheadernomark{2}{21. สญฺโญชนิยทุก 1. กุสลตฺติก}
\prose{1. ปฏิจฺจวาร ปจฺจยจตุกฺกเหตุ}

\proseitem{10}{สญฺโญชนิยํ กุสลํ ธมฺมํ ปฏิจฺจ สญฺโญชนิโย กุสโล ธมฺโม อุปฺปชฺชติ เหตุปจฺจยา. (1)}

\prose{อสญฺโญชนิยํ กุสลํ ธมฺมํ ปฏิจฺจ อสญฺโญชนิโย กุสโล ธมฺโม อุปฺปชฺชติ เหตุปจฺจยา. (1) (สํขิตฺตํ.)}

\proseitem{11}{เหตุยา เทฺว, อารมฺมเณ เทฺว. (สํขิตฺตํ.)}

\prose{(ยถา จูฬนฺตรทุเก โลกิย\-ทุก\-คมนํ, เอวํ อิมมฺปิ ญาตพฺพํ. สหชาตวาโรปิ ฯเปฯ ปญฺหาวาโรปิ วิตฺถาเรตพฺพา.)}

\proseitem{12}{สญฺโญชนิยํ อกุสลํ ธมฺมํ ปฏิจฺจ สญฺโญชนิโย อกุสโล ธมฺโม อุปฺปชฺชติ เหตุปจฺจยา. (สํขิตฺตํ.)}

\proseitem{13}{เหตุยา เอกํ, อารมฺมเณ เอกํ ฯเปฯ อวิคเต เอกํ. (สํขิตฺตํ.) (สหชาตวาเรปิ ฯเปฯ ปญฺหาวาเรปิ สพฺพตฺถ เอกํ.)}

\prose{อพฺยากตปทเหตุ}

\proseitem{14}{สญฺโญชนิยํ อพฺยากตํ ธมฺมํ ปฏิจฺจ สญฺโญชนิโย อพฺยากโต ธมฺโม อุปฺปชฺชติ เหตุปจฺจยา. (1)}

\prose{อสญฺโญชนิยํ อพฺยากตํ ธมฺมํ ปฏิจฺจ อสญฺโญชนิโย อพฺยากโต ธมฺโม อุปฺปชฺชติ เหตุปจฺจยา.  อสญฺโญชนิยํ อพฺยากตํ ธมฺมํ ปฏิจฺจ สญฺโญชนิโย อพฺยากโต ธมฺโม อุปฺปชฺชติ เหตุปจฺจยา.  อสญฺโญชนิยํ อพฺยากตํ ธมฺมํ ปฏิจฺจ สญฺโญชนิโย อพฺยากโต จ อสญฺโญชนิโย อพฺยากโต จ ธมฺโม อุปฺปชฺชนฺติ เหตุปจฺจยา. (3)}

\prose{สญฺโญชนิยํ อพฺยากตญฺจ อสญฺโญชนิยํ อพฺยากตญฺจ ธมฺมํ ปฏิจฺจ สญฺโญชนิโย อพฺยากโต ธมฺโม อุปฺปชฺชติ เหตุปจฺจยา. (1)}

\proseitem{15}{\csromanfolio{372}เหตุยา ปญฺจ, อารมฺมเณ เทฺว, อธิปติยา ปญฺจ ฯเปฯ อวิคเต ปญฺจ. (สํขิตฺตํ.)}

\prose{จูฬนฺตรทุเก โลกิย\-ทุก\-สทิสํ.}

\vspace{\tipitakaparskip}
\csromancenter{(สหชาตวาโรปิ ฯเปฯ สมฺปยุตฺตวาโรปิ วิตฺถาเรตพฺพา.)}
\prose{ปญฺหาวารเหตุ}

\proseitem{16}{สญฺโญชนิโย อพฺยากโต ธมฺโม สญฺโญชนิยสฺส อพฺยากตสฺส ธมฺมสฺส เหตุปจฺจเยน ปจฺจโย. (1)}

\prose{อสญฺโญชนิโย อพฺยากโต ธมฺโม อสญฺโญ\-ชนิยสฺส อพฺยากตสฺส ธมฺมสฺส เหตุปจฺจเยน ปจฺจโย. ตีณิ. (สํขิตฺตํ.)}

\proseitem{17}{เหตุยา จตฺตาริ, อารมฺมเณ ตีณิ, อธิปติยา จตฺตาริ ฯเปฯ อวิคเต สตฺต. (สํขิตฺตํ.)}

\proseitem{18}{สญฺโญชนิโย อพฺยากโต ธมฺโม สญฺโญชนิยสฺส อพฺยากตสฺส ธมฺมสฺส อารมฺมณ\-ปจฺจเยน ปจฺจโย. สหชาต\-ปจฺจเยน ปจฺจโย. อุปนิสฺสย\-ปจฺจเยน ปจฺจโย. ปุเรชาต\-ปจฺจเยน ปจฺจโย. ปจฺฉาชาต\-ปจฺจเยน ปจฺจโย. อาหารปจฺจเยน ปจฺจโย. อินฺทฺริย\-ปจฺจเยน ปจฺจโย. (สํขิตฺตํ.)}

\proseitem{19}{นเหตุยา สตฺต, นอารมฺมเณ สตฺต. (สํขิตฺตํ.)}

\prose{เหตุปจฺจยา นอารมฺมเณ จตฺตาริ. (สํขิตฺตํ.)}

\prose{นเหตุปจฺจยา อารมฺมเณ ตีณิ. (สํขิตฺตํ.)}

\prose{(ยถา กุสลตฺติเก ปญฺหาวารสฺส อนุโลมมฺปิ ปจฺจนียมฺปิ อนุโลมปจฺจนียมฺปิ ปจฺจนียานุโลมมฺปิ คณิตํ, เอวํ คเณตพฺพํ.)}

\csromanmarkh{22. สญฺโญชน\-สมฺปยุตฺตทุก}\csromanmarkhii{1. กุสลตฺติก}\csromanheadernomark{2}{22. สญฺโญชน\-สมฺปยุตฺตทุก 1. กุสลตฺติก}
\prose{1. ปฏิจฺจวาร ปจฺจยจตุกฺกเหตุ}

\proseitem{20}{สญฺโญชน\-วิปฺปยุตฺตํ กุสลํ ธมฺมํ ปฏิจฺจ สญฺโญชน\-วิปฺปยุตฺโต กุสโล ธมฺโม อุปฺปชฺชติ เหตุปจฺจยา. (สํขิตฺตํ.)}

\csromanmarkh{22. สญฺโญชน\-สมฺปยุตฺตทุก}\csromanmarkhii{1. กุสลตฺติก}\csromanheadernomark{2}{22. สญฺโญชน\-สมฺปยุตฺตทุก 1. กุสลตฺติก}
\prose{\csromanfolio{373}เหตุยา เอกํ, อารมฺมเณ เอกํ ฯเปฯ อวิคเต เอกํ. (สํขิตฺตํ.) (สหชาตวาเรปิ ฯเปฯ ปญฺหาวาเรปิ สพฺพตฺถ เอกํ.) สญฺโญชน\-สมฺปยุตฺตํ อกุสลํ ธมฺมํ ปฏิจฺจ สญฺโญชน\-สมฺปยุตฺโต อกุสโล ธมฺโม อุปฺปชฺชติ เหตุปจฺจยา. ตีณิ. สญฺโญชน\-สมฺปยุตฺตํ อกุสลํ ธมฺมํ ปฏิจฺจ สญฺโญชน\-สมฺปยุตฺโต อกุสโล ธมฺโม อุปฺปชฺชติ อารมฺมณ\-ปจฺจยา. (สํขิตฺตํ.)}

\proseitem{22}{เหตุยา ตีณิ, อารมฺมเณ ปญฺจ, อธิปติยา เอกํ ฯเปฯ อวิคเต ปญฺจ. (สํขิตฺตํ. อนุโลมํ.)}

\prose{ปจฺจนีย\-น\-เหตุ สญฺโญชน\-สมฺปยุตฺตํ อกุสลํ ธมฺมํ ปฏิจฺจ สญฺโญชน\-สมฺปยุตฺโต อกุสโล ธมฺโม อุปฺปชฺชติ นเหตุปจฺจยา. วิจิกิจฺฉา\-สหคเต ขนฺเธ ปฏิจฺจ วิจิกิจฺฉา\-สหคโต โมโห.  สญฺโญชน\-สมฺปยุตฺตํ อกุสลํ ธมฺมํ ปฏิจฺจ สญฺโญชน\-วิปฺปยุตฺโต อกุสโล ธมฺโม อุปฺปชฺชติ นเหตุปจฺจยา. อุทฺธจฺจ\-สหคเต ขนฺเธ ปฏิจฺจ อุทฺธจฺจ\-สหคโต โมโห. (2) (สํขิตฺตํ.)}

\proseitem{24}{นเหตุยา เทฺว, นอธิปติยา ปญฺจ, นปุเรชาเต ปญฺจ ฯเปฯ นกมฺเม ตีณิ ฯเปฯ นวิปฺปยุตฺเต ปญฺจ. (สํขิตฺตํ. ปจฺจนียํ.)}

\prose{ปญฺหาวารเหตุ}

\proseitem{25}{สญฺโญชน\-สมฺปยุตฺโต อกุสโล ธมฺโม สญฺโญชน\-สมฺปยุตฺตสฺส อกุสลสฺส ธมฺมสฺส เหตุปจฺจเยน ปจฺจโย. (1)}

\prose{สญฺโญชน\-วิปฺปยุตฺโต อกุสโล ธมฺโม สญฺโญชน\-สมฺปยุตฺตสฺส อกุสลสฺส ธมฺมสฺส อารมฺมณ\-ปจฺจเยน ปจฺจโย. (1)}

\prose{สญฺโญชน\-สมฺปยุตฺโต อกุสโล ธมฺโม สญฺโญชน\-สมฺปยุตฺตสฺส อกุสลสฺส ธมฺมสฺส อารมฺมณ\-ปจฺจเยน ปจฺจโย. นว.}

\prose{สญฺโญชน\-สมฺปยุตฺโต อกุสโล ธมฺโม สญฺโญชน\-สมฺปยุตฺตสฺส อกุสลสฺส ธมฺมสฺส อธิปติ\-ปจฺจเยน ปจฺจโย. (1) (สํขิตฺตํ.)}

\proseitem{26}{\csromanfolio{374}อธิปติยา เอกํ, อนนฺตเร สมนนฺตเร นว, สหชาเต ปญฺจ, อุปนิสฺสเย อาเสวเน นว, กมฺเม ตีณิ, มคฺเค ตีณิ, สมฺปยุตฺเต ปญฺจ, อตฺถิยา ปญฺจ. (สํขิตฺตํ.)}

\prose{อพฺยากตปทเหตุ}

\proseitem{27}{สญฺโญชน\-วิปฺปยุตฺตํ อพฺยากตํ ธมฺมํ ปฏิจฺจ สญฺโญชน\-วิปฺปยุตฺโต อพฺยากโต ธมฺโม อุปฺปชฺชติ เหตุปจฺจยา. (สํขิตฺตํ.)}

\prose{เหตุยา เอกํ, อารมฺมเณ เอกํ ฯเปฯ อวิคเต เอกํ. (สํขิตฺตํ.)}

\vspace{\tipitakaparskip}
\csromancenter{(สหชาตวาเรปิ ฯเปฯ ปญฺหาวาเรปิ สพฺพตฺถ เอกํ.)}
\csromanmarkh{23. สญฺโญชน\-สญฺโญชนิยทุก}\csromanmarkhii{1. กุสลตฺติก}\csromanheadernomark{2}{23. สญฺโญชน\-สญฺโญชนิยทุก 1. กุสลตฺติก}
\csromanheader{2}{1-7. ปฏิจฺจาทิวาร ปจฺจยจตุกฺกเหตุ}
\proseitem{28}{สญฺโญชนิย\-ญฺเจว โน จ สญฺโญชนํ กุสลํ ธมฺมํ ปฏิจฺจ สญฺโญชนิโย เจว โน จ สญฺโญชโน กุสโล ธมฺโม อุปฺปชฺชติ เหตุปจฺจยา. (สํขิตฺตํ.)}

\proseitem{29}{เหตุยา เอกํ, อารมฺมเณ เอกํ ฯเปฯ อวิคเต เอกํ. (สํขิตฺตํ.) (สหชาตวาเรปิ ฯเปฯ ปญฺหาวาเรปิ สพฺพตฺถ เอกํ.)}

\proseitem{30}{สญฺโญชนญฺเจว สญฺโญชนิยญฺจ อกุสลํ ธมฺมํ ปฏิจฺจ สญฺโญชโน เจว สญฺโญชนิโย จ อกุสโล ธมฺโม อุปฺปชฺชติ เหตุปจฺจยา. ตีณิ. (สํขิตฺตํ.)}

\proseitem{31}{เหตุยา นว, อารมฺมเณ นว, อธิปติยา นว ฯเปฯ อวิคเต นว. (สํขิตฺตํ.)}

\prose{นเหตุยา ตีณิ, นอธิปติยา นว, นปุเรชาเต นว, นปจฺฉาชาเต นว, นอาเสวเน นว, นกมฺเม ตีณิ, นวิปาเก นว, นวิปฺปยุตฺเต นว. (สํขิตฺตํ.)}

\prose{(สหชาตวาโรปิ ฯเปฯ สมฺปยุตฺตวาโรปิ ปฏิจฺจ\-วาร\-สทิสา.)}

\prose{\csromanfolio{375}24. สญฺโญชน\-สญฺโญชน\-สมฺปยุตฺตทุก 1. กุสลตฺติก}

\proseitem{32}{สญฺโญชโน เจว สญฺโญชนิโย จ อกุสโล ธมฺโม สญฺโญชนสฺส เจว สญฺโญชนิยสฺส จ อกุสลสฺส ธมฺมสฺส เหตุปจฺจเยน ปจฺจโย. (สํขิตฺตํ.)}

\proseitem{33}{เหตุยา ตีณิ, อารมฺมเณ นว, อธิปติยา นว ฯเปฯ อุปนิสฺสเย อาเสวเน นว, กมฺเม อาหาเร อินฺทฺริเย ฌาเน ตีณิ, มคฺเค สมฺปยุตฺเต นว ฯเปฯ อวิคเต นว. (สํขิตฺตํ.)}

\proseitem{34}{นเหตุยา นว, นอารมฺมเณ นว. (สํขิตฺตํ.)}

\prose{เหตุปจฺจยา นอารมฺมเณ ตีณิ. (สํขิตฺตํ.)}

\prose{นเหตุปจฺจยา อารมฺมเณ นว. (สํขิตฺตํ.)}

\prose{(ยถา กุสลตฺติเก ปญฺหาวารสฺส อนุโลมมฺปิ ปจฺจนียมฺปิ อนุโลมปจฺจนียมฺปิ ปจฺจนียานุโลมมฺปิ คณิตํ, เอวํ คเณตพฺพํ.)}

\proseitem{35}{สญฺโญชนิย\-ญฺเจว โน จ สญฺโญชนํ อพฺยากตํ ธมฺมํ ปฏิจฺจ สญฺโญชนิโย เจว โน จ สญฺโญชโน อพฺยากโต ธมฺโม อุปฺปชฺชติ เหตุปจฺจยา. (สํขิตฺตํ.)}

\prose{36. เหตุยา เอกํ อารมฺมเณ เอกํ ฯเปฯ อวิคเต เอกํ. (สํขิตฺตํ.)}

\vspace{\tipitakaparskip}
\csromancenter{(สหชาตวาเรปิ ฯเปฯ ปญฺหาวาเรปิ สพฺพตฺถ เอกํ.)}
\csromanmarkh{24. สญฺโญชน\-สญฺโญชน\-สมฺปยุตฺตทุก}\csromanmarkhii{1. กุสลตฺติก}\csromanheadernomark{2}{24. สญฺโญชน\-สญฺโญชน\-สมฺปยุตฺตทุก 1. กุสลตฺติก}
\csromanheader{2}{1-7. ปฏิจฺจาทิวาร ปจฺจยจตุกฺกเหตุ}
\proseitem{37}{สญฺโญชนญฺเจว สญฺโญชน\-สมฺปยุตฺตญฺจ อกุสลํ ธมฺมํ ปฏิจฺจ สญฺโญชโน เจว สญฺโญชน\-สมฺปยุตฺโต จ อกุสโล ธมฺโม อุปฺปชฺชติ เหตุปจฺจยา. ตีณิ.}

\prose{สญฺโญชน\-สมฺปยุตฺต\-ญฺเจว โน จ สญฺโญชนํ อกุสลํ ธมฺมํ ปฏิจฺจ สญฺโญชน\-สมฺปยุตฺโต เจว โน จ สญฺโญชโน อกุสโล ธมฺโม อุปฺปชฺชติ เหตุปจฺจยา. ตีณิ.}

\prose{\csromanfolio{376}สญฺโญชนญฺเจว สญฺโญชน\-สมฺปยุตฺตํ อกุสลญฺจ สญฺโญชน\-สมฺปยุตฺต\-ญฺเจว โน จ สญฺโญชนํ อกุสลญฺจ ธมฺมํ ปฏิจฺจ สญฺโญชโน เจว สญฺโญชน\-สมฺปยุตฺโต จ อกุสโล ธมฺโม อุปฺปชฺชติ เหตุปจฺจยา. ตีณิ. (สํขิตฺตํ.)}

\proseitem{38}{เหตุยา นว, อารมฺมเณ นว, อธิปติยา นว ฯเปฯ อวิคเต นว. (สํขิตฺตํ.)}

\prose{นเหตุยา ตีณิ, นอธิปติยา นว, นปุเรชาเต นว, นปจฺฉาชาเต นว, นอาเสวเน นว, นกมฺเม ตีณิ, นวิปาเก นว, นวิปฺปยุตฺเต นว. (สํขิตฺตํ.)}

\prose{(สหชาตวาโรปิ สมฺปยุตฺตวาโรปิ ปฏิจฺจ\-วาร\-สทิสา วิตฺถาเรตพฺพา.)}

\proseitem{39}{สญฺโญชโน เจว สญฺโญชน\-สมฺปยุตฺโต จ อกุสโล ธมฺโม สญฺโญชนสฺส เจว สญฺโญชน\-สมฺปยุตฺตสฺส จ อกุสลสฺส ธมฺมสฺส เหตุปจฺจเยน ปจฺจโย. (สํขิตฺตํ.)}

\proseitem{40}{เหตุยา ตีณิ, อารมฺมเณ นว, อธิปติยา นว ฯเปฯ อุปนิสฺสเย อาเสวเน นว, กมฺเม อาหาเร อินฺทฺริเย ฌาเน ตีณิ, มคฺเค สมฺปยุตฺเต นว ฯเปฯ อวิคเต นว. (สํขิตฺตํ.)}

\proseitem{41}{นเหตุยา นว, นอารมฺมเณ นว, นอธิปติยา นว. (สํขิตฺตํ.)}

\prose{เหตุปจฺจยา นอารมฺมเณ ตีณิ. (สํขิตฺตํ.)}

\prose{นเหตุปจฺจยา อารมฺมเณ นว. (สํขิตฺตํ.)}

\prose{(ยถา กุสลตฺติเก ปญฺหาวารสฺส อนุโลมมฺปิ ปจฺจนียมฺปิ อนุโลมปจฺจนียมฺปิ ปจฺจนียานุโลมมฺปิ คณิตํ, เอวํ คเณตพฺพํ.)}

\prose{25. สญฺโญชน\-วิปฺปยุตฺต\-สญฺโญชนิยทุก 1. กุสลตฺติก}

\csromanheader{2}{1-7. ปฏิจฺจาทิวาร ปจฺจยจตุกฺกเหตุ}
\proseitem{42}{สญฺโญชน\-วิปฺปยุตฺตํ สญฺโญชนิยํ กุสลํ ธมฺมํ ปฏิจฺจ สญฺโญชน\-วิปฺปยุตฺโต สญฺโญชนิโย กุสโล ธมฺโม อุปฺปชฺชติ เหตุปจฺจยา. (1)}

\prose{\csromanfolio{377}25. สญฺโญชน\-วิปฺปยุตฺต\-สญฺโญชนิยทุก 1. กุสลตฺติก}

\prose{สญฺโญชน\-วิปฺปยุตฺตํ อสญฺโญชนิยํ กุสลํ ธมฺมํ ปฏิจฺจ สญฺโญชน\-วิปฺปยุตฺโต อสญฺโญชนิโย กุสโล ธมฺโม อุปฺปชฺชติ เหตุปจฺจยา. (1)}

\proseitem{43}{เหตุยา เทฺว, อารมฺมเณ เทฺว ฯเปฯ อวิคเต เทฺว. (สํขิตฺตํ.)}

\prose{(ยถา จูฬนฺตรทุเก โลกิย\-ทุก\-สทิสํ. สหชาตวาโรปิ ฯเปฯ ปญฺหาวาโรปิ สพฺพตฺถ วิตฺถาเรตพฺพา.)}

\proseitem{44}{สญฺโญชน\-วิปฺปยุตฺโต สญฺโญชนิโย อกุสโล ธมฺโม สญฺโญชน\-วิปฺปยุตฺตสฺส สญฺโญชนิยสฺส อกุสลสฺส ธมฺมสฺส อารมฺมณ\-ปจฺจเยน ปจฺจโย. (สํขิตฺตํ.)}

\proseitem{45}{อารมฺมเณ เอกํ. (สพฺพตฺถ เอกํ, สํขิตฺตํ.)}

\proseitem{46}{สญฺโญชน\-วิปฺปยุตฺตํ สญฺโญชนิยํ อพฺยากตํ ธมฺมํ ปฏิจฺจ สญฺโญชน\-วิปฺปยุตฺโต สญฺโญชนิโย อพฺยากโต ธมฺโม อุปฺปชฺชติ เหตุปจฺจยา. (1)}

\prose{สญฺโญชน\-วิปฺปยุตฺตํ อสญฺโญชนิยํ อพฺยากตํ ธมฺมํ ปฏิจฺจ สญฺโญชน\-วิปฺปยุตฺโต อสญฺโญชนิโย อพฺยากโต ธมฺโม อุปฺปชฺชติ เหตุปจฺจยา.  สญฺโญชน\-วิปฺปยุตฺตํ อสญฺโญชนิยํ อพฺยากตํ ธมฺมํ ปฏิจฺจ สญฺโญชน\-วิปฺปยุตฺโต สญฺโญชนิโย อพฺยากโต ธมฺโม อุปฺปชฺชติ เหตุปจฺจยา.  สญฺโญชน\-วิปฺปยุตฺตํ อสญฺโญชนิยํ อพฺยากตํ ธมฺมํ ปฏิจฺจ สญฺโญชน\-วิปฺปยุตฺโต สญฺโญชนิโย อพฺยากโต จ สญฺโญชน\-วิปฺปยุตฺโต อสญฺโญชนิโย อพฺยากโต จ ธมฺมา อุปฺปชฺชนฺติ เหตุปจฺจยา. (3)}

\prose{สญฺโญชน\-วิปฺปยุตฺตํ สญฺโญชนิยํ อพฺยากตญฺจ สญฺโญชน\-วิปฺปยุตฺตํ อสญฺโญชนิยํ อพฺยากตญฺจ ธมฺมํ ปฏิจฺจ สญฺโญชน\-วิปฺปยุตฺโต สญฺโญชนิโย อพฺยากโต ธมฺโม อุปฺปชฺชติ เหตุปจฺจยา. (1) (สํขิตฺตํ.)}

\proseitem{47}{เหตุยา ปญฺจ, อารมฺมเณ เทฺว, อธิปติยา ปญฺจ ฯเปฯ อาเสวเน เอกํ, กมฺเม ปญฺจ, วิปาเก ปญฺจ ฯเปฯ อวิคเต ปญฺจ. (สํขิตฺตํ.)}

\prose{\csromanfolio{378}นเหตุยา เอกํ, นอารมฺมเณ ตีณิ, นอธิปติยา เทฺว, นปุเรชาเต จตฺตาริ, นปจฺฉาชาเต ปญฺจ, นอาเสวเน ปญฺจ, นกมฺเม เอกํ, นวิปาเก เอกํ, นวิปฺปยุตฺเต เทฺว ฯเปฯ โนวิคเต ตีณิ. (สํขิตฺตํ.) (สหชาตวาโรปิ ฯเปฯ สมฺปยุตฺตวาโรปิ ปฏิจฺจ\-วาร\-สทิสา.)}

\proseitem{48}{สญฺโญชน\-วิปฺปยุตฺโต สญฺโญชนิโย อพฺยากโต ธมฺโม สญฺโญชน\-วิปฺปยุตฺตสฺส สญฺโญชนิยสฺส อพฺยากตสฺส ธมฺมสฺส เหตุปจฺจเยน ปจฺจโย. (สํขิตฺตํ.)}

\proseitem{49}{เหตุยา จตฺตาริ, อารมฺมเณ ตีณิ, อธิปติยา จตฺตาริ, อนนฺตเร จตฺตาริ ฯเปฯ อวิคเต สตฺต. (สํขิตฺตํ.)}

\proseitem{50}{นเหตุยา สตฺต, นอารมฺมเณ สตฺต. (สํขิตฺตํ.)}

\prose{เหตุปจฺจยา นอารมฺมเณ จตฺตาริ. (สํขิตฺตํ.)}

\prose{นเหตุปจฺจยา อารมฺมเณ ตีณิ. (สํขิตฺตํ.)}

\prosewithcloser{(ยถา กุสลตฺติเก ปญฺหาวารสฺส อนุโลมมฺปิ ปจฺจนียมฺปิ อนุโลมปจฺจนียมฺปิ ปจฺจนียานุโลมมฺปิ คณิตํ, เอวํ คเณตพฺพํ.)}{สญฺโญชน\-โคจฺฉก\-กุสล\-ตฺติกํ นิฏฺฐิตํ.}
\csromanmatikaanchor{mtk.69}
\csromanmatikaanchor{mtk.70}
\csromantocmarknopage{subparagraph}{26. คนฺถทุก}{mtk.69}
\bookmark[dest={mtk.69},level=5]{26. คนฺถทุก}
\csromantocmarkref{subparagraph}{1. กุสลตฺติก ...}{mtk.70}
\bookmark[dest={mtk.70},level=5]{1. กุสลตฺติก ...}
\csromanmarkh{26. คนฺถทุก}\csromanmarkhii{1. กุสลตฺติก}\csromanheadernomark{6}{26. \textbf{คนฺถทุก 1. กุสลตฺติก}}
\csromanheader{2}{1-7. ปฏิจฺจาทิวาร ปจฺจยจตุกฺก}
\proseitem{1}{โนคนฺถํ กุสลํ ธมฺมํ ปฏิจฺจ โนคนฺโถ กุสโล ธมฺโม อุปฺปชฺชติ เหตุปจฺจยา. (สํขิตฺตํ.)}

\proseitem{2}{เหตุยา เอกํ, อารมฺมเณ เอกํ ฯเปฯ อวิคเต เอกํ. (สํขิตฺตํ.)}

\vspace{\tipitakaparskip}
\csromancenter{(สหชาตวาเรปิ ฯเปฯ ปญฺหาวาเรปิ สพฺพตฺถ เอกํ.)}
\csromanmarkh{26. คนฺถทุก}\csromanmarkhii{1. กุสลตฺติก}\csromanheadernomark{2}{26. คนฺถทุก 1. กุสลตฺติก}
\proseitem{3}{\csromanfolio{379}คนฺถํ อกุสลํ ธมฺมํ ปฏิจฺจ คนฺโถ อกุสโล ธมฺโม อุปฺปชฺชติ เหตุปจฺจยา.  คนฺถํ อกุสลํ ธมฺมํ ปฏิจฺจ โนคนฺโถ อกุสโล ธมฺโม อุปฺปชฺชติ เหตุปจฺจยา.  คนฺถํ อกุสลํ ธมฺมํ ปฏิจฺจ คนฺโถ อกุสโล จ โนคนฺโถ อกุสโล จ ธมฺมา อุปฺปชฺชนฺติ เหตุปจฺจยา. (3)}

\prose{โนคนฺถํ อกุสลํ ธมฺมํ ปฏิจฺจ โนคนฺโถ อกุสโล ธมฺโม อุปฺปชฺชติ เหตุปจฺจยา. ตีณิ.}

\prose{คนฺถํ อกุสลญฺจ โนคนฺถํ อกุสลญฺจ ธมฺมํ ปฏิจฺจ คนฺโถ อกุสโล ธมฺโม อุปฺปชฺชติ เหตุปจฺจยา. ตีณิ. (สํขิตฺตํ.)}

\proseitem{4}{เหตุยา นว, อารมฺมเณ นว ฯเปฯ อวิคเต นว. (สํขิตฺตํ.)}

\prose{นเหตุยา เอกํ, นอธิปติยา นว, นปุเรชาเต นว, นปจฺฉาชาเต นว, นอาเสวเน นว, นกมฺเม ตีณิ, นวิปาเก นว, นวิปฺปยุตฺเต นว. (สํขิตฺตํ.)}

\prose{(สหชาตวาโรปิ ฯเปฯ สมฺปยุตฺตวาโรปิ ปฏิจฺจ\-วาร\-สทิสา วิตฺถาเรตพฺพา.)}

\proseitem{5}{คนฺโถ อกุสโล ธมฺโม คนฺถสฺส อกุสลสฺส ธมฺมสฺส เหตุปจฺจเยน ปจฺจโย. (สํขิตฺตํ.)}

\proseitem{6}{เหตุยา นว, อารมฺมเณ นว, อธิปติยา นว ฯเปฯ กมฺเม อาหาเร อินฺทฺริเย ฌาเน ตีณิ ฯเปฯ อวิคเต นว. (สํขิตฺตํ.)}

\proseitem{7}{นเหตุยา นว, นอารมฺมเณ นว. (สํขิตฺตํ.)}

\prose{เหตุปจฺจยา นอารมฺมเณ นว. (สํขิตฺตํ.)}

\prose{นเหตุปจฺจยา อารมฺมเณ นว. (สํขิตฺตํ.)}

\prose{(ยถา กุสลตฺติเก ปญฺหาวารสฺส อนุโลมมฺปิ ปจฺจนียมฺปิ อนุโลมปจฺจนียมฺปิ ปจฺจนียานุโลมมฺปิ คณิตํ, เอวํ คเณตพฺพํ.)}

\proseitem{8}{โนคนฺถํ อพฺยากตํ ธมฺมํ ปฏิจฺจ โนคนฺโถ อพฺยากโต ธมฺโม อุปฺปชฺชติ เหตุปจฺจยา. (สํขิตฺตํ.)}

\prose{9. เหตุยา เอกํ ฯเปฯ อวิคเต เอกํ. (สํขิตฺตํ.)}

\vspace{\tipitakaparskip}
\csromancenter{(สหชาตวาเรปิ ฯเปฯ ปญฺหาวาเรปิ สพฺพตฺถ เอกํ.)}
\csromanmarkh{ธมฺมานุโลม ทุกติก\-ปฏฺฐาน\-ปาฬิ}\csromanmarkhii{27. คนฺถนิยทุก 1. กุสลตฺติก}\csromanheadernomark{2}{ธมฺมานุโลม ทุกติก\-ปฏฺฐาน\-ปาฬิ 27. คนฺถนิยทุก 1. กุสลตฺติก}
\csromanheader{2}{1-7. ปฏิจฺจาทิวาร ปจฺจยจตุกฺกเหตุ}
\proseitem{10}{\csromanfolio{380}คนฺถนิยํ กุสลํ ธมฺมํ ปฏิจฺจ คนฺถนิโย กุสโล ธมฺโม อุปฺปชฺชติ เหตุปจฺจยา. (1)}

\prose{อคนฺถนิยํ กุสลํ ธมฺมํ ปฏิจฺจ อคนฺถนิโย กุสโล ธมฺโม อุปฺปชฺชติ เหตุปจฺจยา. (1) (สํขิตฺตํ.)}

\proseitem{11}{เหตุยา เทฺว, อารมฺมเณ เทฺว ฯเปฯ อวิคเต เทฺว. (สํขิตฺตํ. โลกิยโลกุตฺตร\-คมน\-สทิสํ. สหชาตวาโรปิ ฯเปฯ ปญฺหาวาโรปิ วิตฺถาเรตพฺพา.)}

\proseitem{12}{คนฺถนิยํ อกุสลํ ธมฺมํ ปฏิจฺจ คนฺถนิโย อกุสโล ธมฺโม อุปฺปชฺชติ เหตุปจฺจยา. (สํขิตฺตํ.)}

\proseitem{13}{เหตุยา เอกํ, อารมฺมเณ เอกํ ฯเปฯ อวิคเต เอกํ. (สํขิตฺตํ.) (สหชาตวาเรปิ ฯเปฯ ปญฺหาวาเรปิ สพฺพตฺถ เอกํ.)}

\proseitem{14}{คนฺถนิยํ อพฺยากตํ ธมฺมํ ปฏิจฺจ คนฺถนิโย อพฺยากโต ธมฺโม อุปฺปชฺชติ เหตุปจฺจยา. (1)}

\prose{อคนฺถนิยํ อพฺยากตํ ธมฺมํ ปฏิจฺจ อคนฺถนิโย อพฺยากโต ธมฺโม อุปฺปชฺชติ เหตุปจฺจยา. ตีณิ. (3)}

\prose{คนฺถนิยํ อพฺยากตญฺจ อคนฺถนิยํ อพฺยากตญฺจ ธมฺมํ ปฏิจฺจ คนฺถนิโย อพฺยากโต ธมฺโม อุปฺปชฺชติ เหตุปจฺจยา. (1) (สํขิตฺตํ.)}

\proseitem{15}{เหตุยา ปญฺจ, อารมฺมเณ เทฺว ฯเปฯ อวิคเต ปญฺจ. (สํขิตฺตํ. โลกิยโลกุตฺตร\-คมน\-สทิสํ. สหชาตวาโรปิ ฯเปฯ ปญฺหาวาโรปิ วิตฺถาเรตพฺพา.)}

\csromanmarkh{28. คนฺถ\-สมฺปยุตฺตทุก}\csromanmarkhii{1. กุสลตฺติก}\csromanheadernomark{2}{28. คนฺถ\-สมฺปยุตฺตทุก 1. กุสลตฺติก}
\csromanheader{2}{1-7. ปฏิจฺจาทิวาร ปจฺจยจตุกฺกเหตุ}
\proseitem{16}{คนฺถ\-วิปฺปยุตฺตํ กุสลํ ธมฺมํ ปฏิจฺจ คนฺถ\-วิปฺปยุตฺโต กุสโล ธมฺโม อุปฺปชฺชติ เหตุปจฺจยา. (สํขิตฺตํ.)}

\csromanmarkh{29. คนฺถ\-คนฺถนิยทุก}\csromanmarkhii{1. กุสลตฺติก}\csromanheadernomark{2}{29. คนฺถ\-คนฺถนิยทุก 1. กุสลตฺติก}
\proseitem{17}{\csromanfolio{381}เหตุยา เอกํ, อารมฺมเณ เอกํ ฯเปฯ อวิคเต เอกํ. (สํขิตฺตํ.) (สหชาตวาเรปิ ฯเปฯ ปญฺหาวาเรปิ สพฺพตฺถ เอกํ.)}

\proseitem{18}{คนฺถ\-สมฺปยุตฺตํ อกุสลํ ธมฺมํ ปฏิจฺจ คนฺถ\-สมฺปยุตฺโต อกุสโล ธมฺโม อุปฺปชฺชติ เหตุปจฺจยา. ตีณิ.}

\prose{คนฺถ\-วิปฺปยุตฺตํ อกุสลํ ธมฺมํ ปฏิจฺจ คนฺถ\-วิปฺปยุตฺโต อกุสโล ธมฺโม อุปฺปชฺชติ เหตุปจฺจยา. เทฺว.}

\prose{คนฺถ\-สมฺปยุตฺตํ อกุสลญฺจ คนฺถ\-วิปฺปยุตฺตํ อกุสลญฺจ ธมฺมํ ปฏิจฺจ คนฺถ\-สมฺปยุตฺโต อกุสโล ธมฺโม อุปฺปชฺชติ เหตุปจฺจยา. (1) (สํขิตฺตํ.)}

\proseitem{19}{เหตุยา ฉ, อารมฺมเณ ฉ, อธิปติยา ปญฺจ ฯเปฯ อวิคเต ฉ. (สํขิตฺตํ.)}

\prose{นเหตุยา เอกํ, นอธิปติยา ฉ, นปุเรชาเต ฉ ฯเปฯ นกมฺเม จตฺตาริ, นวิปฺปยุตฺเต ฉ. (สํขิตฺตํ.) (สหชาตวาโรปิ ฯเปฯ สมฺปยุตฺตวาโรปิ วิตฺถาเรตพฺพา.)}

\proseitem{20}{คนฺถ\-สมฺปยุตฺโต อกุสโล ธมฺโม คนฺถ\-สมฺปยุตฺตสฺส อกุสลสฺส ธมฺมสฺส เหตุปจฺจเยน ปจฺจโย. (สํขิตฺตํ.)}

\proseitem{21}{เหตุยา ฉ, อารมฺมเณ นว, อธิปติยา นว, อนนฺตเร นว, สหชาเต ฉ ฯเปฯ นิสฺสเย ฉ, อุปนิสฺสเย นว, อาเสวเน นว, กมฺเม จตฺตาริ ฯเปฯ มคฺเค จตฺตาริ, สมฺปยุตฺเต ฉ ฯเปฯ อวิคเต ฉ. (สํขิตฺตํ.)}

\proseitem{22}{คนฺถ\-วิปฺปยุตฺตํ อพฺยากตํ ธมฺมํ ปฏิจฺจ คนฺถ\-วิปฺปยุตฺโต อพฺยากโต ธมฺโม อุปฺปชฺชติ เหตุปจฺจยา. (สํขิตฺตํ.)}

\proseitem{23}{เหตุยา เอกํ, อารมฺมเณ เอกํ ฯเปฯ อวิคเต เอกํ. (สํขิตฺตํ.)}

\vspace{\tipitakaparskip}
\csromancenter{(สหชาตวาโรปิ ฯเปฯ ปญฺหาวาเรปิ สพฺพตฺถ เอกํ.)}
\csromanmarkh{29. คนฺถ\-คนฺถนิยทุก}\csromanmarkhii{1. กุสลตฺติก}\csromanheadernomark{2}{29. คนฺถ\-คนฺถนิยทุก 1. กุสลตฺติก}
\csromanheader{2}{1-7. ปฏิจฺจาทิวาร ปจฺจยจตุกฺกเหตุ}
\proseitem{24}{คนฺถนิย\-ญฺเจว โน จ คนฺถํ กุสลํ ธมฺมํ ปฏิจฺจ คนฺถนิโย เจว โน จ คนฺโถ กุสโล ธมฺโม อุปฺปชฺชติ เหตุปจฺจยา. (1) (สํขิตฺตํ.)}

\proseitem{25}{\csromanfolio{382}เหตุยา เอกํ, อารมฺมเณ เอกํ, อวิคเต เอกํ.}

\vspace{\tipitakaparskip}
\csromancenter{(สหชาตวาเรปิ ฯเปฯ ปญฺหาวาเรปิ สพฺพตฺถ เอกํ.)}
\proseitem{26}{คนฺถญฺเจว คนฺถนิยญฺจ อกุสลํ ธมฺมํ ปฏิจฺจ คนฺโถ เจว คนฺถนิโย จ อกุสโล ธมฺโม อุปฺปชฺชติ เหตุปจฺจยา. ตีณิ.}

\prose{คนฺถนิย\-ญฺเจว โน จ คนฺถํ อกุสลํ ธมฺมํ ปฏิจฺจ คนฺถนิโย เจว โน จ คนฺโถ อกุสโล ธมฺโม อุปฺปชฺชติ เหตุปจฺจยา. ตีณิ.}

\prose{คนฺถญฺเจว คนฺถนิยํ อกุสลญฺจ คนฺถนิย\-ญฺเจว โน จ คนฺถํ อกุสลญฺจ ธมฺมํ ปฏิจฺจ คนฺโถ เจว คนฺถนิโย จ อกุสโล ธมฺโม อุปฺปชฺชติ เหตุปจฺจยา. ตีณิ. (สํขิตฺตํ.)}

\proseitem{27}{เหตุยา นว, อารมฺมเณ นว, อธิปติยา นว ฯเปฯ อวิคเต นว. (สํขิตฺตํ.)}

\prose{นเหตุยา เอกํ, นอธิปติยา นว ฯเปฯ นกมฺเม ตีณิ ฯเปฯ นวิปฺปยุตฺเต นว. (สํขิตฺตํ.)}

\prose{(สหชาตวารมฺปิ ฯเปฯ สมฺปยุตฺตวารมฺปิ ปฏิจฺจ\-วาร\-สทิสํ วิตฺถาเรตพฺพํ.)}

\proseitem{28}{คนฺโถ เจว คนฺถนิโย จ อกุสโล ธมฺโม คนฺถสฺส เจว คนฺถนิยสฺส จ อกุสลสฺส ธมฺมสฺส เหตุปจฺจเยน ปจฺจโย. (สํขิตฺตํ.)}

\proseitem{29}{เหตุยา นว, อารมฺมเณ นว, อธิปติยา นว ฯเปฯ กมฺเม อาหาเร อินฺทฺริเย ฌาเน ตีณิ ฯเปฯ อวิคเต นว. (สํขิตฺตํ.)}

\proseitem{30}{นเหตุยา นว, นอารมฺมเณ นว. (สํขิตฺตํ.)}

\prose{เหตุปจฺจยา นอารมฺมเณ นว. (สํขิตฺตํ.)}

\prose{นเหตุปจฺจยา อารมฺมเณ นว. (สํขิตฺตํ.)}

\prose{(ยถา กุสลตฺติเก ปญฺหาวารสฺส อนุโลมมฺปิ ปจฺจนียมฺปิ อนุโลมปจฺจนียมฺปิ ปจฺจนียานุโลมมฺปิ คณิตํ, เอวํ คเณตพฺพํ.)}

\proseitem{31}{คนฺถนิย\-ญฺเจว โน จ คนฺถํ อพฺยากตํ ธมฺมํ ปฏิจฺจ คนฺถนิโย เจว โน จ คนฺโถ อพฺยากโต ธมฺโม อุปฺปชฺชติ เหตุปจฺจยา. (สํขิตฺตํ.)}

\proseitem{32}{เหตุยา เอกํ, อารมฺมเณ เอกํ ฯเปฯ อวิคเต เอกํ. (สํขิตฺตํ.) (สหชาตวาเรปิ ฯเปฯ ปญฺหาวาเรปิ สพฺพตฺถ เอกํ.)}

\vspace{\tipitakaparskip}
\csromancenter{คนฺถ\-วิปฺปยุตฺต\-คนฺถนิยทุก 1. กุสลตฺติก 30. คนฺถ\-คนฺถ\-สมฺปยุตฺตทุก 1. กุสลตฺติก}
\csromanheader{2}{1-7. ปฏิจฺจาทิวาร ปจฺจยจตุกฺกเหตุ}
\proseitem{33}{\csromanfolio{383}คนฺถญฺเจว คนฺถ\-สมฺปยุตฺตญฺจ อกุสลํ ธมฺมํ ปฏิจฺจ คนฺโถ เจว คนฺถ\-สมฺปยุตฺโต จ อกุสโล ธมฺโม อุปฺปชฺชติ เหตุปจฺจยา. (สํขิตฺตํ.)}

\proseitem{34}{เหตุยา นว, อารมฺมเณ นว ฯเปฯ อวิคเต นว. (สํขิตฺตํ.)}

\prose{นอธิปติยา นว, นปุเรชาเต นว, นกมฺเม ตีณิ ฯเปฯ นวิปฺปยุตฺเต นว. (สํขิตฺตํ.) (สหชาตวาโรปิ ฯเปฯ สมฺปยุตฺตวาโรปิ ปฏิจฺจ\-วาร\-สทิสา.)}

\proseitem{35}{คนฺโถ เจว คนฺถ\-สมฺปยุตฺโต จ อกุสโล ธมฺโม คนฺถสฺส เจว คนฺถ\-สมฺปยุตฺตสฺส จ อกุสลสฺส ธมฺมสฺส เหตุปจฺจเยน ปจฺจโย. (สํขิตฺตํ.)}

\proseitem{36}{เหตุยา นว, อารมฺมเณ นว ฯเปฯ กมฺเม อาหาเร อินฺทฺริเย ฌาเน ตีณิ ฯเปฯ อวิคเต นว. (สํขิตฺตํ.)}

\proseitem{37}{นเหตุยา นว, นอารมฺมเณ นว. (สํขิตฺตํ.)}

\prose{เหตุปจฺจยา นอารมฺมเณ นว. (สํขิตฺตํ.)}

\prose{นเหตุปจฺจยา อารมฺมเณ นว. (สํขิตฺตํ.)}

\prose{(ยถา กุสลตฺติเก ปญฺหาวารสฺส อนุโลมมฺปิ ปจฺจนียมฺปิ อนุโลมปจฺจนียมฺปิ ปจฺจนียานุโลมมฺปิ คณิตํ, เอวํ คเณตพฺพํ.)}

\csromanmarkh{31. คนฺถ\-วิปฺปยุตฺต\-คนฺถนิยทุก}\csromanmarkhii{1. กุสลตฺติก}\csromanheadernomark{2}{31. คนฺถ\-วิปฺปยุตฺต\-คนฺถนิยทุก 1. กุสลตฺติก}
\csromanheader{2}{1-7. ปฏิจฺจาทิวาร ปจฺจยจตุกฺกเหตุ}
\proseitem{38}{คนฺถ\-วิปฺปยุตฺตํ คนฺถนิยํ กุสลํ ธมฺมํ ปฏิจฺจ อนฺถวิปฺปยุตฺโต คนฺถนิโย กุสโล ธมฺโม อุปฺปชฺชติ เหตุปจฺจยา. (1)}

\prose{คนฺถ\-วิปฺปยุตฺตํ อคนฺถนิยํ กุสลํ ธมฺมํ ปฏิจฺจ คนฺถ\-วิปฺปยุตฺโต อคนฺถนิโย กุสโล ธมฺโม อุปฺปชฺชติ เหตุปจฺจยา. (1) (สํขิตฺตํ.)}

\proseitem{39}{\csromanfolio{384}เหตุยา เทฺว, อารมฺมเณ เทฺว ฯเปฯ อวิคเต เทฺว. (สํขิตฺตํ.) (สหชาตวาโรปิ ฯเปฯ ปญฺหาวาโรปิ สพฺพตฺถ วิตฺถาเรตพฺพา.)}

\proseitem{40}{คนฺถ\-วิปฺปยุตฺตํ คนฺถนิยํ อกุสลํ ธมฺมํ ปฏิจฺจ คนฺถ\-วิปฺปยุตฺโต คนฺถนิโย อกุสโล ธมฺโม อุปฺปชฺชติ เหตุปจฺจยา. (สํขิตฺตํ.)}

\proseitem{41}{เหตุยา เอกํ ฯเปฯ อวิคเต เอกํ. (สํขิตฺตํ.) (สหชาตวาเรปิ ฯเปฯ ปญฺหาวาเรปิ สพฺพตฺถ เอกํ.)}

\proseitem{42}{คนฺถ\-วิปฺปยุตฺตํ คนฺถนิยํ อพฺยากตํ ธมฺมํ ปฏิจฺจ คนฺถ\-วิปฺปยุตฺโต อคนฺถนิโย อพฺยากโต ธมฺโม อุปฺปชฺชติ เหตุปจฺจยา. (1)}

\prose{คนฺถ\-วิปฺปยุตฺตํ อคนฺถนิยํ อพฺยากตํ ธมฺมํ ปฏิจฺจ คนฺถ\-วิปฺปยุตฺโต อคนฺถนิโย อพฺยากโต ธมฺโม อุปฺปชฺชติ เหตุปจฺจยา. ตีณิ.}

\prose{คนฺถ\-วิปฺปยุตฺตํ คนฺถนิยํ อพฺยากตญฺจ คนฺถ\-วิปฺปยุตฺตํ อคนฺถนิยํ อพฺยากตญฺจ ธมฺมํ ปฏิจฺจ คนฺถ\-วิปฺปยุตฺโต คนฺถนิโย อพฺยากโต ธมฺโม อุปฺปชฺชติ เหตุปจฺจยา. (1) (สํขิตฺตํ.)}

\proseitem{43}{เหตุยา ปญฺจ, อารมฺมเณ เทฺว ฯเปฯ วิปาเก ปญฺจ ฯเปฯ อวิคเต ปญฺจ. (สํขิตฺตํ. สหชาตวารมฺปิ ฯเปฯ สมฺปยุตฺตวารมฺปิ ปฏิจฺจ\-วาร\-สทิสํ.)}

\proseitem{44}{คนฺถ\-วิปฺปยุตฺโต คนฺถนิโย อพฺยากโต ธมฺโม คนฺถ\-วิปฺปยุตฺตสฺส คนฺถนิยสฺส อพฺยากตสฺส ธมฺมสฺส เหตุปจฺจเยน ปจฺจโย. (1)}

\prose{คนฺถ\-วิปฺปยุตฺโต อคนฺถนิโย อพฺยากโต ธมฺโม คนฺถ\-วิปฺปยุตฺตสฺส อคนฺถนิยสฺส อพฺยากตสฺส ธมฺมสฺส เหตุปจฺจเยน ปจฺจโย. ตีณิ. (สํขิตฺตํ.)}

\proseitem{45}{เหตุยา จตฺตาริ, อารมฺมเณ ตีณิ, อธิปติยา จตฺตาริ ฯเปฯ อวิคเต สตฺต. (สํขิตฺตํ.)}

\proseitem{46}{นเหตุยา สตฺต, นอารมฺมเณ สตฺต. (สํขิตฺตํ.)}

\csromanmatikaanchor{mtk.71}
\csromanmatikaanchor{mtk.72}
\csromantocmarknopage{subparagraph}{44. นีวรณทุก}{mtk.71}
\bookmark[dest={mtk.71},level=5]{44. นีวรณทุก}
\csromantocmarkref{subparagraph}{1. กุสลตฺติก ...}{mtk.72}
\bookmark[dest={mtk.72},level=5]{1. กุสลตฺติก ...}
\csromanmarkh{44. นีวรณทุก}\csromanmarkhii{1. กุสลตฺติก}\csromanheadernomark{6}{44. นีวรณทุก 1. กุสลตฺติก}
\prose{\csromanfolio{385}เหตุปจฺจยา นอารมฺมเณ จตฺตาริ. (สํขิตฺตํ.)}

\prose{นเหตุปจฺจยา อารมฺมเณ ตีณิ. (สํขิตฺตํ.)}

\prosewithcloser{(ยถา กุสลตฺติเก ปญฺหาวารสฺส อนุโลมมฺปิ ปจฺจนียมฺปิ อนุโลมปจฺจนียมฺปิ ปจฺจนียานุโลมมฺปิ คณิตํ, เอวํ คเณตพฺพํ.)}{คนฺถ\-โคจฺฉก\-กุสล\-ตฺติกํ นิฏฺฐิตํ.}
\csromancenter{(โอฆโคจฺฉกมฺปิ โยคโคจฺฉกมฺปิ อาสวโคจฺฉก- กุสล\-ตฺติก\-สทิสํ.)}
\csromanmarkh{44. นีวรณทุก}\csromanmarkhii{1. กุสลตฺติก}\csromanheadernomark{2}{44. นีวรณทุก 1. กุสลตฺติก}
\csromanheader{2}{1-7. ปฏิจฺจาทิวาร ปจฺจยจตุกฺกเหตุ}
\proseitem{1}{โนนีวรณํ กุสลํ ธมฺมํ ปฏิจฺจ โนนีวรโณ กุสโล ธมฺโม อุปฺปชฺชติ เหตุปจฺจยา. (สํขิตฺตํ.)}

\proseitem{2}{เหตุยา เอกํ, อารมฺมเณ เอกํ ฯเปฯ อวิคเต เอกํ. (สํขิตฺตํ.)}

\prose{(สหชาตวาเรปิ ฯเปฯ ปญฺหาวาเรปิ สพฺพตฺถ เอกํ.)}

\proseitem{3}{นีวรณํ อกุสลํ ธมฺมํ ปฏิจฺจ นีวรโณ อกุสโล ธมฺโม อุปฺปชฺชติ เหตุปจฺจยา.  นีวรณํ อกุสลํ ธมฺมํ ปฏิจฺจ โนนีวรโณ อกุสโล ธมฺโม อุปฺปชฺชติ เหตุปจฺจยา.  นีวรณํ อกุสลํ ธมฺมํ ปฏิจฺจ นีวรโณ อกุสโล จ โนนีวรโณ อกุสโล จ ธมฺมา อุปฺปชฺชนฺติ เหตุปจฺจยา. (3) (สํขิตฺตํ.)}

\proseitem{4}{เหตุยา นว, อารมฺมเณ นว ฯเปฯ อวิคเต นว. (สํขิตฺตํ.)}

\prose{ปจฺจนีย\-น\-เหตุ}

\proseitem{5}{นีวรณํ อกุสลํ ธมฺมํ ปฏิจฺจ นีวรโณ อกุสโล ธมฺโม อุปฺปชฺชติ นเหตุปจฺจยา. วิจิกิจฺฉา\-นีวรณํ อุทฺธจฺจ\-นีวรณํ ปฏิจฺจ อวิชฺชา\-นีวรณํ. (1)}

\prose{โนนีวรณํ อกุสลํ ธมฺมํ ปฏิจฺจ นีวรโณ อกุสโล ธมฺโม อุปฺปชฺชติ นเหตุปจฺจยา. วิจิกิจฺฉา\-สหคเต อุทฺธจฺจ\-สหคเต ขนฺเธ ปฏิจฺจ อวิชฺชา\-นีวรณํ. (1)}

\prose{\csromanfolio{386}นีวรณญฺจ โนนีวรณญฺจ อกุสลํ ธมฺมํ ปฏิจฺจ นีวรโณ อกุสโล ธมฺโม อุปฺปชฺชติ นเหตุปจฺจยา. วิจิกิจฺฉานีวรณญฺจ อุทฺธจฺจ\-นีวรณญฺจ สมฺปยุตฺตเก จ ขนฺเธ ปฏิจฺจ อวิชฺชา\-นีวรณํ. (1) (สํขิตฺตํ.)}

\proseitem{6}{นเหตุยา ตีณิ, นอธิปติยา นว, ปุเรชาเต นว ฯเปฯ นกมฺเม ตีณิ ฯเปฯ นวิปฺปยุตฺเต นว. (สํขิตฺตํ. สหชาตวาเรปิ สมฺปยุตฺตวาเรปิ สพฺพตฺถ นว.)}

\proseitem{7}{นีวรโณ อกุสโล ธมฺโม นีวรณสฺส อกุสลสฺส ธมฺมสฺส เหตุปจฺจเยน ปจฺจโย. (สํขิตฺตํ.)}

\proseitem{8}{เหตุยา ตีณิ, อารมฺมเณ นว, อธิปติยา นว ฯเปฯ อาเสวเน นว, กมฺเม อาหาเร อินฺทฺริเย ฌาเน มคฺเค ตีณิ, สมฺปยุตฺเต นว ฯเปฯ อวิคเต นว. (สํขิตฺตํ.)}

\prose{นเหตุยา นว, นอารมฺมเณ นว. (สํขิตฺตํ.)}

\prose{เหตุปจฺจยา นอารมฺมเณ ตีณิ. (สํขิตฺตํ.)}

\prose{นเหตุปจฺจยา อารมฺมเณ นว. (สํขิตฺตํ.)}

\prose{(ยถา กุสลตฺติเก ปญฺหาวารสฺส อนุโลมมฺปิ ปจฺจนียมฺปิ อนุโลมปจฺจนียมฺปิ ปจฺจนียานุโลมมฺปิ คณิตํ, เอวํ คเณตพฺพํ.)}

\proseitem{9}{โนนีวรณํ อพฺยากตํ ธมฺมํ ปฏิจฺจ โนนีวรโณ อพฺยากโต ธมฺโม อุปฺปชฺชติ เหตุปจฺจยา. (สํขิตฺตํ.)}

\proseitem{10}{เหตุยา เอกํ, อารมฺมเณ เอกํ ฯเปฯ อวิคเต เอกํ. (สํขิตฺตํ.)}

\vspace{\tipitakaparskip}
\csromancenter{(สหชาตวาเรปิ ฯเปฯ ปญฺหาวาเรปิ สพฺพตฺถ เอกํ.)}
\csromanmarkh{45. นีวรณิยทุก}\csromanmarkhii{1. กุสลตฺติก}\csromanheadernomark{2}{45. \textbf{นีวรณิยทุก 1. กุสลตฺติก}}
\prose{1. ปฏิจฺจวาร ปจฺจยจตุกฺกเหตุ}

\proseitem{11}{นีวรณิยํ กุสลํ ธมฺมํ ปฏิจฺจ นีวรณิโย กุสโล ธมฺโม อุปฺปชฺชติ เหตุปจฺจยา. (1)}

\csromanmarkh{46. นีวรณ\-สมฺปยุตฺตทุก}\csromanmarkhii{1. กุสลตฺติก}\csromanheadernomark{2}{46. นีวรณ\-สมฺปยุตฺตทุก 1. กุสลตฺติก}
\prose{\csromanfolio{387}อนีวรณิยํ กุสลํ ธมฺมํ ปฏิจฺจ อนีวรณิโย กุสโล ธมฺโม อุปฺปชฺชติ เหตุปจฺจยา. (1) (สํขิตฺตํ.)}

\proseitem{12}{เหตุยา เทฺว, อารมฺมเณ เทฺว ฯเปฯ อวิคเต เทฺว. (สํขิตฺตํ.)}

\prose{(สหชาตวาโรปิ ฯเปฯ ปญฺหาวาโรปิ วิตฺถาเรตพฺพา.)}

\proseitem{13}{นีวรณิยํ อกุสลํ ธมฺมํ ปฏิจฺจ นีวรณิโย อกุสโล ธมฺโม อุปฺปชฺชติ เหตุปจฺจยา. (สํขิตฺตํ.)}

\proseitem{14}{เหตุยา เอกํ, อารมฺมเณ เอกํ ฯเปฯ อวิคเต เอกํ. (สํขิตฺตํ.) (สหชาตวาเรปิ ฯเปฯ ปญฺหาวาเรปิ สพฺพตฺถ เอกํ.)}

\proseitem{15}{นีวรณิยํ อพฺยากตํ ธมฺมํ ปฏิจฺจ นีวรณิโย อพฺยากโต ธมฺโม อุปฺปชฺชติ เหตุปจฺจยา. (1)}

\prose{อนีวรณิยํ อพฺยากตํ ธมฺมํ ปฏิจฺจ อนีวรณิโย อพฺยากโต ธมฺโม อุปฺปชฺชติ เหตุปจฺจยา. ตีณิ.}

\prose{นีวรณิยํ อพฺยากตญฺจ อนีวรณิยํ อพฺยากตญฺจ ธมฺมํ ปฏิจฺจ นีวรณิโย อพฺยากโต ธมฺโม อุปฺปชฺชติ เหตุปจฺจยา. (1) (สํขิตฺตํ.)}

\proseitem{16}{เหตุยา ปญฺจ, อารมฺมเณ เทฺว ฯเปฯ อาเสวเน เอกํ ฯเปฯ วิปาเก ปญฺจ ฯเปฯ อวิคเต ปญฺจ. (สํขิตฺตํ. สหชาตวาโรปิ ฯเปฯ ปญฺหาวาโรปิ วิตฺถาเรตพฺพา.)}

\csromanmarkh{46. นีวรณ\-สมฺปยุตฺตทุก}\csromanmarkhii{1. กุสลตฺติก}\csromanheadernomark{2}{46. นีวรณ\-สมฺปยุตฺตทุก 1. กุสลตฺติก}
\prose{1. ปฏิจฺจวาร ปจฺจยจตุกฺกเหตุ}

\proseitem{17}{นีวรณ\-วิปฺปยุตฺตํ กุสลํ ธมฺมํ ปฏิจฺจ นีวรณ\-วิปฺปยุตฺโต กุสโล ธมฺโม อุปฺปชฺชติ เหตุปจฺจยา. (สํขิตฺตํ.)}

\proseitem{18}{เหตุยา เอกํ, อารมฺมเณ เอกํ ฯเปฯ อวิคเต เอกํ. (สํขิตฺตํ.)}

\vspace{\tipitakaparskip}
\csromancenter{(สหชาตวาเรปิ ฯเปฯ ปญฺหาวาเรปิ สพฺพตฺถ เอกํ.)}
\proseitem{19}{\csromanfolio{388}นีวรณ\-สมฺปยุตฺตํ อกุสลํ ธมฺมํ ปฏิจฺจ นีวรณ\-สมฺปยุตฺโต อกุสโล ธมฺโม อุปฺปชฺชติ เหตุปจฺจยา. (สํขิตฺตํ.)}

\proseitem{20}{เหตุยา เอกํ, อารมฺมเณ เอกํ ฯเปฯ อวิคเต เอกํ. (สํขิตฺตํ.) (สหชาตวาเรปิ ฯเปฯ ปญฺหาวาเรปิ สพฺพตฺถ เอกํ.)}

\proseitem{21}{นีวรณ\-วิปฺปยุตฺตํ อพฺยากตํ ธมฺมํ ปฏิจฺจ นีวรณ\-วิปฺปยุตฺโต อพฺยากโต ธมฺโม อุปฺปชฺชติ เหตุปจฺจยา. (สํขิตฺตํ.)}

\proseitem{22}{เหตุยา เอกํ, อารมฺมเณ เอกํ ฯเปฯ อวิคเต เอกํ. (สํขิตฺตํ.)}

\vspace{\tipitakaparskip}
\csromancenter{(สหชาตวาเรปิ ฯเปฯ ปญฺหาวาเรปิ สพฺพตฺถ เอกํ.)}
\csromanmarkh{47. นีวรณ\-นีวรณิยทุก}\csromanmarkhii{1. กุสลตฺติก}\csromanheadernomark{2}{47. นีวรณ\-นีวรณิยทุก 1. กุสลตฺติก}
\prose{1. ปฏิจฺจวาร ปจฺจยจตุกฺกเหตุ}

\proseitem{23}{นีวรณิย\-ญฺเจว โน จ นีวรณํ กุสลํ ธมฺมํ ปฏิจฺจ นีวรณิโย เจว โน จ นีวรโณ กุสโล ธมฺโม อุปฺปชฺชติ เหตุปจฺจยา. (สํขิตฺตํ.)}

\proseitem{24}{เหตุยา เอกํ, อารมฺมเณ เอกํ ฯเปฯ อวิคเต เอกํ. (สํขิตฺตํ.) (สหชาตวาเรปิ ฯเปฯ ปญฺหาวาเรปิ สพฺพตฺถ เอกํ.)}

\proseitem{25}{นีวรณญฺเจว นีวรณิยญฺจ อกุสลํ ธมฺมํ ปฏิจฺจ นีวรโณ เจว นีวรณิโย จ อกุสโล ธมฺโม อุปฺปชฺชติ เหตุปจฺจยา. ตีณิ.}

\prose{นีวรณิย\-ญฺเจว โน จ นีวรณํ อกุสลํ ธมฺมํ ปฏิจฺจ นีวรณิโย เจว โน จ นีวรโณ อกุสโล ธมฺโม อุปฺปชฺชติ เหตุปจฺจยา. ตีณิ.}

\prose{นีวรณญฺเจว นีวรณิยํ อกุสลญฺจ นีวรณิย\-ญฺเจว โน จ นีวรณํ อกุสลญฺจ ธมฺมํ ปฏิจฺจ นีวรโณ จ นีวรณิโย จ อกุสโล ธมฺโม อุปฺปชฺชติ เหตุปจฺจยา. ตีณิ. (สํขิตฺตํ.)}

\proseitem{26}{เหตุยา นว, อารมฺมเณ นว ฯเปฯ อวิคเต นว. (สํขิตฺตํ.)}

\vspace{\tipitakaparskip}
\csromancenter{(สหชาตวาเรปิ ฯเปฯ ปญฺหาวาเรปิ สพฺพตฺถ นว.)}
\csromanmarkh{49. นีวรณ\-วิปฺปยุตฺต\-นีวรณิยทุก}\csromanmarkhii{1. กุสลตฺติก}\csromanheadernomark{2}{49. นีวรณ\-วิปฺปยุตฺต\-นีวรณิยทุก 1. กุสลตฺติก}
\proseitem{27}{\csromanfolio{389}นีวรณิย\-ญฺเจว โน จ นีวรณํ อพฺยากตํ ธมฺมํ ปฏิจฺจ นีวรณิโย เจว โน จ นีวรโณ อพฺยากโต ธมฺโม อุปฺปชฺชติ เหตุปจฺจยา. (สํขิตฺตํ.)}

\proseitem{28}{เหตุยา เอกํ, อารมฺมเณ เอกํ ฯเปฯ อวิคเต เอกํ. (สํขิตฺตํ.)}

\vspace{\tipitakaparskip}
\csromancenter{(สหชาตวาเรปิ ฯเปฯ ปญฺหาวาเรปิ สพฺพตฺถ เอกํ.)}
\csromanmarkh{48. นีวรณ\-นีวรณ\-สมฺปยุตฺตทุก}\csromanmarkhii{1. กุสลตฺติก}\csromanheadernomark{2}{48. นีวรณ\-นีวรณ\-สมฺปยุตฺตทุก 1. กุสลตฺติก}
\prose{1. ปฏิจฺจวาร ปจฺจยจตุกฺกเหตุ}

\proseitem{29}{นีวรณญฺเจว นีวรณ\-สมฺปยุตฺตญฺจ อกุสลํ ธมฺมํ ปฏิจฺจ นีวรโณ เจว นีวรณ\-สมฺปยุตฺโต จ อกุสโล ธมฺโม อุปฺปชฺชติ เหตุปจฺจยา. (สํขิตฺตํ.)}

\proseitem{30}{เหตุยา นว, อารมฺมเณ นว ฯเปฯ อวิคเต นว. (สํขิตฺตํ.)}

\vspace{\tipitakaparskip}
\csromancenter{(สหชาตวาเรปิ ฯเปฯ ปญฺหาวาเรปิ สพฺพตฺถ นว.)}
\csromanmarkh{49. นีวรณ\-วิปฺปยุตฺต\-นีวรณิยทุก}\csromanmarkhii{1. กุสลตฺติก}\csromanheadernomark{2}{49. นีวรณ\-วิปฺปยุตฺต\-นีวรณิยทุก 1. กุสลตฺติก}
\prose{1. ปฏิจฺจวาร ปจฺจยจตุกฺกเหตุ}

\proseitem{31}{นีวรณ\-วิปฺปยุตฺตํ นีวรณิยํ กุสลํ ธมฺมํ ปฏิจฺจ นีวรณ\-วิปฺปยุตฺโต นีวรณิโย กุสโล ธมฺโม อุปฺปชฺชติ เหตุปจฺจยา. (1)}

\prose{นีวรณ\-วิปฺปยุตฺตํ อนีวรณิยํ กุสลํ ธมฺมํ ปฏิจฺจ นีวรณ\-วิปฺปยุตฺโต อนีวรณิโย กุสโล ธมฺโม อุปฺปชฺชติ เหตุปจฺจยา. (1) (สํขิตฺตํ.)}

\proseitem{32}{เหตุยา เทฺว, อารมฺมเณ เทฺว ฯเปฯ อวิคเต เทฺว. (สํขิตฺตํ.) (สหชาตวาเรปิ ฯเปฯ ปญฺหาวาเรปิ วิตฺถาเรตพฺพา.)}

\proseitem{33}{นีวรณ\-วิปฺปยุตฺตํ นีวรณิยํ อพฺยากตํ ธมฺมํ ปฏิจฺจ นีวรณ\-วิปฺปยุตฺโต นีวรณิโย อพฺยากโต ธมฺโม อุปฺปชฺชติ เหตุปจฺจยา. (1)}

\prose{\csromanfolio{390}นีวรณ\-วิปฺปยุตฺตํ อนีวรณิยํ อพฺยากตํ ธมฺมํ ปฏิจฺจ นีวรณ\-วิปฺปยุตฺโต อนีวรณิโย อพฺยากโต ธมฺโม อุปฺปชฺชติ เหตุปจฺจยา. ตีณิ.}

\prose{นีวรณ\-วิปฺปยุตฺตํ นีวรณิยํ อพฺยากตญฺจ นีวรณ\-วิปฺปยุตฺตํ อนีวรณิยํ อพฺยากตญฺจ ธมฺมํ ปฏิจฺจ นีวรณ\-วิปฺปยุตฺโต นีวรณิโย อพฺยากโต ธมฺโม อุปฺปชฺชติ เหตุปจฺจยา. (1) (สํขิตฺตํ.)}

\proseitemwithcloser{34}{เหตุยา ปญฺจ, อารมฺมเณ เทฺว ฯเปฯ วิปาเก ปญฺจ. (สํขิตฺตํ. สหชาตวาเรปิ ฯเปฯ ปญฺหาวาเรปิ สพฺพตฺถ วิตฺถาเรตพฺพํ.)}{นีวรณ\-โคจฺฉก\-กุสล\-ตฺติกํ นิฏฺฐิตํ.}
\csromanmatikaanchor{mtk.73}
\csromanmatikaanchor{mtk.74}
\csromantocmarknopage{subparagraph}{50. ปรามาสทุก}{mtk.73}
\bookmark[dest={mtk.73},level=5]{50. ปรามาสทุก}
\csromantocmarkref{subparagraph}{1. กุสลตฺติก ...}{mtk.74}
\bookmark[dest={mtk.74},level=5]{1. กุสลตฺติก ...}
\csromanmarkh{50. ปรามาสทุก}\csromanmarkhii{1. กุสลตฺติก}\csromanheadernomark{6}{50. \textbf{ปรามาสทุก 1. กุสลตฺติก}}
\csromanheader{2}{1-7. ปฏิจฺจาทิวาร ปจฺจยจตุกฺกเหตุ}
\proseitem{1}{โนปรามาสํ กุสลํ ธมฺมํ ปฏิจฺจ โนปรามาโส กุสโล ธมฺโม อุปฺปชฺชติ เหตุปจฺจยา. (สํขิตฺตํ.)}

\proseitem{2}{เหตุยา เอกํ, อารมฺมเณ เอกํ ฯเปฯ อวิคเต เอกํ. (สํขิตฺตํ.) (สหชาตวาเรปิ ฯเปฯ ปญฺหาวาเรปิ สพฺพตฺถ เอกํ.)}

\proseitem{3}{ปรามาสํ อกุสลํ ธมฺมํ ปฏิจฺจ โนปรามาโส อกุสโล ธมฺโม อุปฺปชฺชติ เหตุปจฺจยา. (1)}

\prose{โนปรามาสํ อกุสลํ ธมฺมํ ปฏิจฺจ โนปรามาโส อกุสโล ธมฺโม อุปฺปชฺชติ เหตุปจฺจยา. ตีณิ.}

\prose{ปรามาสํ อกุสลญฺจ โนปรามาสํ อกุสลญฺจ ธมฺมํ ปฏิจฺจ โนปรามาโส อกุสโล ธมฺโม อุปฺปชฺชติ เหตุปจฺจยา. (1) (สํขิตฺตํ.)}

\proseitem{4}{เหตุยา ปญฺจ, อารมฺมเณ ปญฺจ ฯเปฯ อวิคเต ปญฺจ. (สํขิตฺตํ.)}

\prose{นเหตุยา เอกํ, นอธิปติยา ปญฺจ, นปุเรชาเต ปญฺจ, นปจฺฉาชาเต ปญฺจ ฯเปฯ. (สํขิตฺตํ. สหชาตวารมฺปิ ฯเปฯ สมฺปยุตฺตวารมฺปิ วิตฺถาเรตพฺพํ.)}

\csromanmarkh{51. ปรามฏฺฐทุก}\csromanmarkhii{1. กุสลตฺติก}\csromanheadernomark{2}{51. ปรามฏฺฐทุก 1. กุสลตฺติก}
\proseitem{5}{\csromanfolio{391}โนปรามาโส อกุสโล ธมฺโม โนปรามาสสฺส อกุสลสฺส ธมฺมสฺส เหตุปจฺจเยน ปจฺจโย. (สํขิตฺตํ.)}

\proseitem{6}{เหตุยา ตีณิ, อารมฺมเณ นว, อธิปติยา นว ฯเปฯ อวิคเต ปญฺจ. (สํขิตฺตํ.)}

\prose{(ยถา กุสลตฺติเก ปญฺหาวารสฺส อนุโลมมฺปิ ปจฺจนียมฺปิ อนุโลมปจฺจนียมฺปิ ปจฺจนียานุโลมมฺปิ คณิตํ, เอวํ คเณตพฺพํ.)}

\proseitem{7}{โนปรามาสํ อพฺยากตํ ธมฺมํ ปฏิจฺจ โนปรามาโส อพฺยากโต ธมฺโม อุปฺปชฺชติ เหตุปจฺจยา. (สํขิตฺตํ.)}

\proseitem{8}{เหตุยา เอกํ, อารมฺมเณ เอกํ ฯเปฯ อวิคเต เอกํ. (สํขิตฺตํ.)}

\vspace{\tipitakaparskip}
\csromancenter{(สหชาตวาเรปิ ฯเปฯ ปญฺหาวาเรปิ สพฺพตฺถ เอกํ.)}
\csromanmarkh{51. ปรามฏฺฐทุก}\csromanmarkhii{1. กุสลตฺติก}\csromanheadernomark{2}{51. ปรามฏฺฐทุก 1. กุสลตฺติก}
\prose{1. ปฏิจฺจวาร ปจฺจยจตุกฺกเหตุ}

\proseitem{9}{ปรามฏฺฐํ กุสลํ ธมฺมํ ปฏิจฺจ ปรามฏฺโฐ กุสโล ธมฺโม อุปฺปชฺชติ เหตุปจฺจยา. (1)}

\prose{อปรามฏฺฐํ กุสลํ ธมฺมํ ปฏิจฺจ อปรามฏฺโฐ กุสโล ธมฺโม อุปฺปชฺชติ เหตุปจฺจยา. (1) (สํขิตฺตํ.)}

\proseitem{10}{เหตุยา เทฺว, อารมฺมเณ เทฺว ฯเปฯ อวิคเต เทฺว. (สํขิตฺตํ.)}

\prose{(สหชาตวาโรปิ ฯเปฯ ปญฺหาวาโรปิ วิตฺถาเรตพฺพา.)}

\proseitem{11}{ปรามฏฺฐํ อกุสลํ ธมฺมํ ปฏิจฺจ ปรามฏฺโฐ อกุสโล ธมฺโม อุปฺปชฺชติ เหตุปจฺจยา. (สํขิตฺตํ.)}

\proseitem{12}{เหตุยา เอกํ, อารมฺมเณ เอกํ ฯเปฯ อวิคเต เอกํ. (สํขิตฺตํ.)}

\vspace{\tipitakaparskip}
\csromancenter{(สหชาตวาเรปิ ฯเปฯ ปญฺหาวาเรปิ สพฺพตฺถ เอกํ.)}
\proseitem{13}{\csromanfolio{392}ปรามฏฺฐํ อพฺยากตํ ธมฺมํ ปฏิจฺจ ปรามฏฺโฐ อพฺยากโต ธมฺโม อุปฺปชฺชติ เหตุปจฺจยา. (1)}

\prose{อปรามฏฺฐํ อพฺยากตํ ธมฺมํ ปฏิจฺจ อปรามฏฺโฐ อพฺยากโต ธมฺโม อุปฺปชฺชติ เหตุปจฺจยา. ตีณิ.}

\prose{ปรามฏฺฐํ อพฺยากตญฺจ อปรามฏฺฐํ อพฺยากตญฺจ ธมฺมํ ปฏิจฺจ ปรามฏฺโฐ อพฺยากโต ธมฺโม อุปฺปชฺชติ เหตุปจฺจยา. (1) (สํขิตฺตํ.)}

\proseitem{14}{เหตุยา ปญฺจ, อารมฺมเณ เทฺว ฯเปฯ อาเสวเน เอกํ ฯเปฯ อวิคเต ปญฺจ. (สํขิตฺตํ. สหชาตวารมฺปิ ฯเปฯ ปญฺหาวารมฺปิ วิตฺถาเรตพฺพํ.)}

\csromanmarkh{52. ปรามาส\-สมฺปยุตฺตทุก}\csromanmarkhii{1. กุสลตฺติก}\csromanheadernomark{2}{52. ปรามาส\-สมฺปยุตฺตทุก 1. กุสลตฺติก}
\csromanheader{2}{1-7. ปฏิจฺจาทิวาร ปจฺจยจตุกฺกเหตุ}
\proseitem{15}{ปรามาส\-วิปฺปยุตฺตํ กุสลํ ธมฺมํ ปฏิจฺจ ปรามาส\-วิปฺปยุตฺโต กุสโล ธมฺโม อุปฺปชฺชติ เหตุปจฺจยา. (สํขิตฺตํ.)}

\proseitem{16}{เหตุยา เอกํ, อารมฺมเณ เอกํ ฯเปฯ อวิคเต เอกํ. (สํขิตฺตํ.)}

\prose{(สหชาตวาเรปิ ฯเปฯ ปญฺหาวาเรปิ สพฺพตฺถ เอกํ.)}

\proseitem{17}{ปรามาส\-สมฺปยุตฺตํ อกุสลํ ธมฺมํ ปฏิจฺจ ปรามาส\-สมฺปยุตฺโต อกุสโล ธมฺโม อุปฺปชฺชติ เหตุปจฺจยา. (1)}

\prose{ปรามาส\-วิปฺปยุตฺตํ อกุสลํ ธมฺมํ ปฏิจฺจ ปรามาส\-วิปฺปยุตฺโต อกุสโล ธมฺโม อุปฺปชฺชติ เหตุปจฺจยา. (1) (สํขิตฺตํ.)}

\proseitem{18}{เหตุยา เทฺว, อารมฺมเณ เทฺว, (สพฺพตฺถ เทฺว.) ฯเปฯ อวิคเต เทฺว. (สํขิตฺตํ.)}

\prose{นเหตุยา เอกํ, นอธิปติยา เทฺว, (สพฺพตฺถ เทฺว.) ฯเปฯ นวิปฺปยุตฺเต เทฺว. (สํขิตฺตํ. สหชาตวาโรปิ ฯเปฯ สมฺปยุตฺตวาโรปิ วิตฺถาเรตพฺพา.)}

\proseitem{19}{ปรามาส\-สมฺปยุตฺโต อกุสโล ธมฺโม ปรามาส\-สมฺปยุตฺตสฺส อกุสลสฺส ธมฺมสฺส เหตุปจฺจเยน ปจฺจโย. (1)}

\csromanmarkh{53. ปรามาส\-ปรามฏฺฐทุก}\csromanmarkhii{1. กุสลตฺติก}\csromanheadernomark{2}{53. ปรามาส\-ปรามฏฺฐทุก 1. กุสลตฺติก}
\prose{\csromanfolio{393}ปรามาส\-วิปฺปยุตฺโต อกุสโล ธมฺโม ปรามาส\-วิปฺปยุตฺตสฺส อกุสลสฺส ธมฺมสฺส เหตุปจฺจเยน ปจฺจโย. (1) (สํขิตฺตํ.)}

\proseitem{20}{เหตุยา เทฺว, อารมฺมเณ จตฺตาริ, อธิปติยา จตฺตาริ, (มชฺฌิเม จ อนฺเต จ อารมฺมณา\-ธิปติ.) อนนฺตเร เทฺว ฯเปฯ สหชาเต เทฺว ฯเปฯ อุปนิสฺสเย จตฺตาริ, อาเสวเน เทฺว, กมฺเม เทฺว, อาหาเร เทฺว ฯเปฯ มคฺเค เทฺว ฯเปฯ อวิคเต เทฺว. (สํขิตฺตํ.)}

\proseitem{21}{ปรามาส\-วิปฺปยุตฺตํ อพฺยากตํ ธมฺมํ ปฏิจฺจ ปรามาส\-วิปฺปยุตฺโต อพฺยากโต ธมฺโม อุปฺปชฺชติ เหตุปจฺจยา. (สํขิตฺตํ.)}

\proseitem{22}{เหตุยา เอกํ, อารมฺมเณ เอกํ ฯเปฯ อวิคเต เอกํ. (สํขิตฺตํ.)}

\vspace{\tipitakaparskip}
\csromancenter{(สหชาตวาเรปิ ฯเปฯ ปญฺหาวาเรปิ สพฺพตฺถ เอกํ.)}
\csromanmarkh{53. ปรามาส\-ปรามฏฺฐทุก}\csromanmarkhii{1. กุสลตฺติก}\csromanheadernomark{2}{53. ปรามาส\-ปรามฏฺฐทุก 1. กุสลตฺติก}
\csromanheader{2}{1-7. ปฏิจฺจาทิวาร ปจฺจยจตุกฺก}
\proseitem{23}{ปรามฏฺฐ\-ญฺเจว โน จ ปรามาสํ กุสลํ ธมฺมํ ปฏิจฺจ ปรามฏฺโฐ เจว โน จ ปรามาโส กุสโล ธมฺโม อุปฺปชฺชติ เหตุปจฺจยา. (สํขิตฺตํ.)}

\proseitem{24}{เหตุยา เอกํ, อารมฺมเณ เอกํ ฯเปฯ อวิคเต เอกํ. (สํขิตฺตํ.) (สหชาตวาเรปิ ฯเปฯ ปญฺหาวาเรปิ สพฺพตฺถ เอกํ.)}

\proseitem{25}{ปรามาสญฺเจว ปรามฏฺฐญฺจ อกุสลํ ธมฺมํ ปฏิจฺจ ปรามฏฺโฐ เจว โน จ ปรามาโส อกุสโล ธมฺโม อุปฺปชฺชติ เหตุปจฺจยา. เอกํ.}

\prose{ปรามฏฺฐ\-ญฺเจว โน จ ปรามาสํ อกุสลํ ธมฺมํ ปฏิจฺจ ปรามฏฺโฐ เจว โน จ ปรามาโส อกุสโล ธมฺโม อุปฺปชฺชติ เหตุปจฺจยา. ตีณิ.}

\prose{ปรามาสญฺเจว ปรามฏฺฐํ อกุสลญฺจ ปรามฏฺฐ\-ญฺเจว โน จ ปรามาสํ อกุสลญฺจ ธมฺมํ ปฏิจฺจ ปรามฏฺโฐ เจว โน จ ปรามาโส อกุสโล ธมฺโม อุปฺปชฺชติ เหตุปจฺจยา. เอกํ. (สํขิตฺตํ.)}

\csromanhangingbegin
\hangingheaditem{26}{เหตุยา ปญฺจ, อารมฺมเณ ปญฺจ ฯเปฯ อวิคเต ปญฺจ. (สํขิตฺตํ.)}
\hangingbody{(สหชาตวารมฺปิ ฯเปฯ สมฺปยุตฺตวารมฺปิ วิตฺถาเรตพฺพํ.)}
\csromanhangingend

\proseitem{27}{\csromanfolio{394}ปรามฏฺโฐ เจว โน จ ปรามาโส อกุสโล ธมฺโม ปรามฏฺฐสฺส เจว โน จ ปรามาสสฺส อกุสลสฺส ธมฺมสฺส เหตุปจฺจเยน ปจฺจโย. (สํขิตฺตํ.)}

\proseitem{28}{เหตุยา ตีณิ, อารมฺมเณ นว, อธิปติยา นว ฯเปฯ อวิคเต ปญฺจ. (สํขิตฺตํ.)}

\prose{(ยถา กุสลตฺติเก ปญฺหาวารสฺส อนุโลมมฺปิ ปจฺจนียมฺปิ อนุโลมปจฺจนียมฺปิ ปจฺจนียานุโลมมฺปิ คณิตํ, เอวํ คเณตพฺพํ.)}

\proseitem{29}{ปรามฏฺฐ\-ญฺเจว โน จ ปรามาสํ อพฺยากตํ ธมฺมํ ปฏิจฺจ ปรามฏฺโฐ เจว โน จ ปรามาโส อพฺยากโต ธมฺโม อุปฺปชฺชติ เหตุปจฺจยา. (สํขิตฺตํ.)}

\proseitem{30}{เหตุยา เอกํ, อารมฺมเณ เอกํ ฯเปฯ อวิคเต เอกํ. (สํขิตฺตํ.)}

\vspace{\tipitakaparskip}
\csromancenter{(สหชาตวาเรปิ ฯเปฯ ปญฺหาวาเรปิ สพฺพตฺถ เอกํ.)}
\prose{54. ปรามาส\-วิปฺปยุตฺต\-ปรามฏฺฐทุก 1. กุสลตฺติก}

\prose{1. ปฏิจฺจวาร ปจฺจยจตุกฺกเหตุ}

\proseitem{31}{ปรามาส\-วิปฺปยุตฺตํ ปรามฏฺฐํ กุสลํ ธมฺมํ ปฏิจฺจ ปรามาส\-วิปฺปยุตฺโต ปรามฏฺโฐ กุสโล ธมฺโม อุปฺปชฺชติ เหตุปจฺจยา. (1)}

\prose{ปรามาส\-วิปฺปยุตฺตํ อปรามฏฺฐํ กุสลํ ธมฺมํ ปฏิจฺจ ปรามาส\-วิปฺปยุตฺโต อปรามฏฺโฐ กุสโล ธมฺโม อุปฺปชฺชติ เหตุปจฺจยา. (1) (สํขิตฺตํ.)}

\proseitem{32}{เหตุยา เทฺว, อารมฺมเณ เทฺว ฯเปฯ อวิคเต เทฺว. (สํขิตฺตํ. โลกิยโลกุตฺตร\-ทุก\-สทิสํ. สหชาตวาโรปิ ฯเปฯ ปญฺหาวาโรปิ วิตฺถาเรตพฺพา.)}

\proseitem{33}{ปรามาส\-วิปฺปยุตฺตํ ปรามฏฺฐํ อกุสลํ ธมฺมํ ปฏิจฺจ ปรามาส\-วิปฺปยุตฺโต ปรามฏฺโฐ อกุสโล ธมฺโม อุปฺปชฺชติ เหตุปจฺจยา. (สพฺพตฺถ เอกํ.)}

\proseitem{34}{ปรามาส\-วิปฺปยุตฺตํ ปรามฏฺฐํ อพฺยากตํ ธมฺมํ ปฏิจฺจ ปรามาส\-วิปฺปยุตฺโต ปรามฏฺโฐ อพฺยากโต ธมฺโม อุปฺปชฺชติ เหตุปจฺจยา. (1)}

\csromanmatikaanchor{mtk.75}
\csromanmatikaanchor{mtk.76}
\csromantocmarknopage{subparagraph}{55. สารมฺมณทุก}{mtk.75}
\bookmark[dest={mtk.75},level=5]{55. สารมฺมณทุก}
\csromantocmarkref{subparagraph}{1. กุสลตฺติก ...}{mtk.76}
\bookmark[dest={mtk.76},level=5]{1. กุสลตฺติก ...}
\csromanmarkh{55. สารมฺมณทุก}\csromanmarkhii{1. กุสลตฺติก}\csromanheadernomark{6}{55. สารมฺมณทุก 1. กุสลตฺติก}
\prose{\csromanfolio{395}ปรามาส\-วิปฺปยุตฺตํ อปรามฏฺฐํ อพฺยากตํ ธมฺมํ ปฏิจฺจ ปรามาส\-วิปฺปยุตฺโต อปรามฏฺโฐ อพฺยากโต ธมฺโม อุปฺปชฺชติ เหตุปจฺจยา. ตีณิ.}

\prose{ปรามาส\-วิปฺปยุตฺตํ ปรามฏฺฐํ อพฺยากตญฺจ ปรามาส\-วิปฺปยุตฺตํ อปรามฏฺฐํ อพฺยากตญฺจ ธมฺมํ ปฏิจฺจ ปรามาส\-วิปฺปยุตฺโต ปรามฏฺโฐ อพฺยากโต ธมฺโม อุปฺปชฺชติ เหตุปจฺจยา. (1) (สํขิตฺตํ.)}

\proseitemwithcloser{35}{เหตุยา ปญฺจ, อารมฺมเณ เทฺว ฯเปฯ อวิคเต ปญฺจ. (สํขิตฺตํ. โลกิย\-ทุก\-สทิสํ. สหชาตวารมฺปิ ฯเปฯ ปญฺหาวารมฺปิ วิตฺถาเรตพฺพํ.)}{ปรามาส\-โคจฺฉก\-กุสล\-ตฺติกํ นิฏฺฐิตํ.}
\csromanmarkh{55. สารมฺมณทุก}\csromanmarkhii{1. กุสลตฺติก}\csromanheadernomark{2}{55. สารมฺมณทุก 1. กุสลตฺติก}
\csromanheader{2}{1-7. ปฏิจฺจาทิวาร ปจฺจยจตุกฺกเหตุ}
\proseitem{1}{สารมฺมณํ กุสลํ ธมฺมํ ปฏิจฺจ สารมฺมโณ กุสโล ธมฺโม อุปฺปชฺชติ เหตุปจฺจยา.}

\vspace{\tipitakaparskip}
\csromancenter{(ยถา สปฺปจฺจยํ กุสลํ, เอวํ วิตฺถาเรตพฺพํ.)}
\proseitem{2}{สารมฺมณํ อกุสลํ ธมฺมํ ปฏิจฺจ สารมฺมโณ อกุสโล ธมฺโม อุปฺปชฺชติ เหตุปจฺจยา. (ยถา สปฺปจฺจยํ อกุสลํ, เอวํ กาตพฺพํ.)}

\vspace{\tipitakaparskip}
\csromancenter{(ยถา สปฺปจฺจยํ อกุสลํ, เอวํ กาตพฺพํ.)}
\proseitem{3}{สารมฺมณํ อพฺยากตํ ธมฺมํ ปฏิจฺจ สารมฺมโณ อพฺยากโต ธมฺโม อุปฺปชฺชติ เหตุปจฺจยา. ตีณิ.}

\prose{อนารมฺมณํ อพฺยากตํ ธมฺมํ ปฏิจฺจ อนารมฺมโณ อพฺยากโต ธมฺโม อุปฺปชฺชติ เหตุปจฺจยา. ตีณิ.}

\prose{สารมฺมณํ อพฺยากตญฺจ อนารมฺมณํ อพฺยากตญฺจ ธมฺมํ ปฏิจฺจ สารมฺมโณ อพฺยากโต ธมฺโม อุปฺปชฺชติ เหตุปจฺจยา. ตีณิ. (สํขิตฺตํ.)}

\proseitem{4}{เหตุยา นว, อารมฺมเณ ตีณิ, อธิปติยา ปญฺจ ฯเปฯ อญฺญมญฺเญ ฉ ฯเปฯ ปุเรชาเต เอกํ, อาเสวเน เอกํ. (สํขิตฺตํ.)}

\prose{นเหตุยา นว, นอารมฺมเณ ตีณิ, นอธิปติยา นว ฯเปฯ นกมฺเม เทฺว, นวิปาเก ปญฺจ, นอาหาเร เอกํ, นอินฺทฺริเย เอกํ, นฌาเน เทฺว, \csromanfolio{396}นมคฺเค นว, นสมฺปยุตฺเต ตีณิ, นวิปฺปยุตฺเต เทฺว, โนนตฺถิยา ตีณิ, โนวิคเต ตีณิ. (สํขิตฺตํ.)}

\csromanheader{2}{\textbf{เหตุ อารมฺมณ อธิปติ}}
\proseitem{5}{สารมฺมโณ อพฺยากโต ธมฺโม สารมฺมณสฺส อพฺยากตสฺส ธมฺมสฺส เหตุปจฺจเยน ปจฺจโย. ตีณิ.}

\proseitem{4}{สารมฺมโณ อพฺยากโต ธมฺโม สารมฺมณสฺส อพฺยากตสฺส ธมฺมสฺส อารมฺมณ\-ปจฺจเยน ปจฺจโย. (1)}

\prose{อนารมฺมโณ อพฺยากโต ธมฺโม สารมฺมณสฺส อพฺยากตสฺส ธมฺมสฺส อารมฺมณ\-ปจฺจเยน ปจฺจโย. (1)}

\proseitem{7}{สารมฺมโณ อพฺยากโต ธมฺโม สารมฺมณสฺส อพฺยากตสฺส ธมฺมสฺส อธิปติ\-ปจฺจเยน ปจฺจโย. อารมฺมณา\-ธิปติ, สหชาตา\-ธิปติ. ตีณิ. อนารมฺมโณ อพฺยากโต ธมฺโม สารมฺมณสฺส อพฺยากตสฺส ธมฺมสฺส อธิปติ\-ปจฺจเยน ปจฺจโย. อารมฺมณา\-ธิปติ. (สํขิตฺตํ.)}

\proseitem{8}{เหตุยา ตีณิ, อารมฺมเณ เทฺว, อธิปติยา จตฺตาริ, อนนฺตเร เอกํ ฯเปฯ สหชาเต สตฺต, อญฺญมญฺเญ ฉ, นิสฺสเย สตฺต, อุปนิสฺสเย เทฺว, ปุเรชาเต เอกํ, ปจฺฉาชาเต เอกํ, อาเสวเน เอกํ, กมฺเม ตีณิ, วิปาเก ตีณิ, สมฺปยุตฺเต เอกํ, วิปฺปยุตฺเต เทฺว. (สํขิตฺตํ.)}

\csromanmarkh{56. จิตฺตทุก}\csromanmarkhii{1. กุสลตฺติก}\csromanheadernomark{2}{56. \textbf{จิตฺตทุก 1. กุสลตฺติก}}
\csromanheader{2}{1-7. ปฏิจฺจาทิวาร ปจฺจยจตุกฺกเหตุ}
\proseitem{9}{จิตฺตํ กุสลํ ธมฺมํ ปฏิจฺจ โนจิตฺโต กุสโล ธมฺโม อุปฺปชฺชติ เหตุปจฺจยา. (1)}

\prose{โนจิตฺตํ กุสลํ ธมฺมํ ปฏิจฺจ โนจิตฺโต กุสโล ธมฺโม อุปฺปชฺชติ เหตุปจฺจยา.  โนจิตฺตํ กุสลํ ธมฺมํ ปฏิจฺจ จิตฺโต กุสโล ธมฺโม อุปฺปชฺชติ เหตุปจฺจยา.  โนจิตฺตํ กุสลํ ธมฺมํ ปฏิจฺจ จิตฺโต กุสโล จ โนจิตฺโต กุสโล จ ธมฺมา อุปฺปชฺชนฺติ เหตุปจฺจยา. (3)}

\csromanmarkh{56. จิตฺตทุก}\csromanmarkhii{1. กุสลตฺติก}\csromanheadernomark{2}{56. จิตฺตทุก 1. กุสลตฺติก}
\prose{\csromanfolio{397}จิตฺตํ กุสลญฺจ โนจิตฺตํ กุสลญฺจ ธมฺมํ ปฏิจฺจ โนจิตฺโต กุสโล ธมฺโม อุปฺปชฺชติ เหตุปจฺจยา. (1) (สํขิตฺตํ.)}

\proseitem{10}{เหตุยา ปญฺจ, อารมฺมเณ ปญฺจ. (สพฺพตฺถ ปญฺจ.) อวิคเต ปญฺจ. (สํขิตฺตํ.)}

\prose{นอธิปติยา ปญฺจ, นปุเรชาเต ปญฺจ, นปจฺฉาชาเต ปญฺจ, นอาเสวเน ปญฺจ, นกมฺเม ตีณิ, นวิปาเก ปญฺจ, นวิปฺปยุตฺเต ปญฺจ. (สํขิตฺตํ.) (สหชาตวารมฺปิ ฯเปฯ สมฺปยุตฺตวารมฺปิ วิตฺถาเรตพฺพํ.)}

\proseitem{11}{โนจิตฺโต กุสโล ธมฺโม โนจิตฺตสฺส กุสลสฺส ธมฺมสฺส เหตุปจฺจเยน ปจฺจโย.  โนจิตฺโต กุสโล ธมฺโม จิตฺตสฺส กุสลสฺส ธมฺมสฺส เหตุปจฺจเยน ปจฺจโย. โนจิตฺโต กุสโล ธมฺโม จิตฺตสฺส กุสลสฺส จ โนจิตฺตสฺส กุสลสฺส จ ธมฺมสฺส เหตุปจฺจเยน ปจฺจโย. (3) (สํขิตฺตํ.)}

\proseitem{12}{เหตุยา ตีณิ, อารมฺมเณ นว, อธิปติยา นว, อนนฺตเร นว ฯเปฯ สหชาเต อญฺญมญฺเญ นิสฺสเย ปญฺจ, อุปนิสฺสเย อาเสวเน นว, กมฺเม ตีณิ, อาหาเร อินฺทฺริเย ปญฺจ, ฌาเน มคฺเค ตีณิ, สมฺปยุตฺเต ปญฺจ ฯเปฯ อวิคเต ปญฺจ. (สํขิตฺตํ.)}

\prose{นเหตุยา นว, นอารมฺมเณ นว. (สํขิตฺตํ.)}

\prose{เหตุปจฺจยา นอารมฺมเณ ตีนิ. (สํขิตฺตํ.)}

\prose{นเหตุปจฺจยา อารมฺมเณ นว. (สํขิตฺตํ.)}

\prose{(ยถา กุสลตฺติเก ปญฺหาวารสฺส อนุโลมมฺปิ ปจฺจนียมฺปิ อนุโลมปจฺจนียมฺปิ ปจฺจนียานุโลมมฺปิ คณิตํ, เอวํ คเณตพฺพํ.)}

\proseitem{13}{จิตฺตํ อกุสลํ ธมฺมํ ปฏิจฺจ โนจิตฺโต อกุสโล ธมฺโม อุปฺปชฺชติ เหตุปจฺจยา. (1)}

\prose{โนจิตฺตํ อกุสลํ ธมฺมํ ปฏิจฺจ โนจิตฺโต อกุสโล ธมฺโม อุปฺปชฺชติ เหตุปจฺจยา.  โนจิตฺตํ อกุสลํ ธมฺมํ ปฏิจฺจ จิตฺโต อกุสโล ธมฺโม อุปฺปชฺชติ เหตุปจฺจยา.  โนจิตฺตํ อกุสลํ ธมฺมํ ปฏิจฺจ จิตฺโต อกุสโล จ โนจิตฺโต อกุสโล จ ธมฺมา อุปฺปชฺชนฺติ เหตุปจฺจโย. (3)}

\prose{\csromanfolio{398}จิตฺตํ อกุสลญฺจ โนจิตฺตํ อกุสลญฺจ ธมฺมํ ปฏิจฺจ โนจิตฺโต อกุสโล ธมฺโม อุปฺปชฺชติ เหตุปจฺจยา. (1) (สํขิตฺตํ.)}

\proseitem{14}{เหตุยา ปญฺจ, อารมฺมเณ ปญฺจ ฯเปฯ อวิคเต ปญฺจ. (สํขิตฺตํ. จิตฺต\-ทุก\-กุสล\-สทิสํ, นเหตุยาปิ กตฺตพฺพํ. สหชาตวารมฺปิ ฯเปฯ ปญฺหาวารมฺปิ กาตพฺพํ.)}

\proseitem{15}{จิตฺตํ อพฺยากตํ ธมฺมํ ปฏิจฺจ โนจิตฺโต อพฺยากโต ธมฺโม อุปฺปชฺชติ เหตุปจฺจยา. (1)}

\prose{โนจิตฺตํ อพฺยากตํ ธมฺมํ ปฏิจฺจ โนจิตฺโต อพฺยากโต ธมฺโม อุปฺปชฺชติ เหตุปจฺจยา.  โนจิตฺตํ อพฺยากตํ ธมฺมํ ปฏิจฺจ จิตฺโต อพฺยากโต ธมฺโม อุปฺปชฺชติ เหตุปจฺจยา.  โนจิตฺตํ อพฺยากตํ ธมฺมํ ปฏิจฺจ จิตฺโต อพฺยากโต จ โนจิตฺโต อพฺยากโต จ ธมฺมา อุปฺปชฺชนฺติ เหตุปจฺจยา. (3)}

\prose{จิตฺตํ อพฺยากตญฺจ โนจิตฺตํ อพฺยากตญฺจ ธมฺมํ ปฏิจฺจ โนจิตฺโต อพฺยากโต ธมฺโม อุปฺปชฺชติ เหตุปจฺจยา. (1) (สํขิตฺตํ.)}

\proseitem{16}{เหตุยา ปญฺจ, อารมฺมเณ ปญฺจ ฯเปฯ วิปาเก ปญฺจ ฯเปฯ อวิคเต ปญฺจ. (สํขิตฺตํ.)}

\prose{นเหตุยา ปญฺจ, นอารมฺมเณ ตีณิ, นอธิปติยา ปญฺจ ฯเปฯ นปุเรชาเต นปจฺฉาชาเต นอาเสวเน ปญฺจ, นกมฺเม ตีณิ, นวิปาเก ปญฺจ, นอาหาเร เอกํ, นอินฺทฺริเย เอกํ, นฌาเน ปญฺจ, นมคฺเค ปญฺจ, นสมฺปยุตฺเต ตีณิ, นวิปฺปยุตฺเต ปญฺจ ฯเปฯ โนวิคเต ตีณิ. (สํขิตฺตํ.)}

\prose{(สหชาตวารมฺปิ ฯเปฯ สมฺปยุตฺตวารมฺปิ วิตฺถาเรตพฺพํ.)}

\proseitem{17}{โนจิตฺโต อพฺยากโต ธมฺโม โนจิตฺตสฺส อพฺยากตสฺส ธมฺมสฺส เหตุปจฺจเยน ปจฺจโย. ตีณิ. (จิตฺต\-ทุก\-กุสล\-สทิสํ วิตฺถาเรตพฺพํ.)}

\proseitem{18}{เหตุยา ตีณิ, อารมฺมเณ นว, อธิปติยา นว, อนนฺตเร นว ฯเปฯ สหชาเต ปญฺจ, ปุเรชาเต ตีณิ, อาเสวเน นว, กมฺเม}

\vspace{\tipitakaparskip}
\csromancenter{เจตสิกทุก 1. กุสลตฺติก ตีณิ, วิปาเก ปญฺจ, วิปฺปยุตฺเต ปญฺจ, อตฺถิยา ปญฺจ, นตฺถิยา นว ฯเปฯ อวิคเต ปญฺจ. (สํขิตฺตํ.)}
\csromanmarkh{57. เจตสิกทุก}\csromanmarkhii{1. กุสลตฺติก}\csromanheadernomark{2}{57. เจตสิกทุก 1. กุสลตฺติก}
\csromanheader{2}{1-7. ปฏิจฺจาทิวาร ปจฺจยจตุกฺกเหตุ}
\proseitem{19}{\csromanfolio{399}เจตสิกํ กุสลํ ธมฺมํ ปฏิจฺจ เจตสิโก กุสโล ธมฺโม อุปฺปชฺชติ เหตุปจฺจยา.  เจตสิกํ กุสลํ ธมฺมํ ปฏิจฺจ อเจตสิโก กุสโล ธมฺโม อุปฺปชฺชติ เหตุปจฺจยา.  เจตสิกํ กุสลํ ธมฺมํ ปฏิจฺจ เจตสิโก กุสโล จ อเจตสิโก กุสโล จ ธมฺมา อุปฺปชฺชนฺติ เหตุปจฺจยา. (3)}

\prose{อเจตสิกํ กุสลํ ธมฺมํ ปฏิจฺจ เจตสิโก กุสโล ธมฺโม อุปฺปชฺชติ เหตุปจฺจยา. (1)}

\prose{เจตสิกํ กุสลญฺจ อเจตสิกํ กุสลญฺจ ธมฺมํ ปฏิจฺจ เจตสิโก กุสโล ธมฺโม อุปฺปชฺชติ เหตุปจฺจยา. (1) (สํขิตฺตํ.)}

\proseitem{20}{เหตุยา ปญฺจ, อารมฺมเณ ปญฺจ, อธิปติยา ปญฺจ ฯเปฯ อวิคเต ปญฺจ.}

\prose{นอธิปติยา ปญฺจ, นปุเรชาเต ปญฺจ ฯเปฯ นกมฺเม ตีณิ, นวิปาเก ปญฺจ ฯเปฯ นวิปฺปยุตฺเต ปญฺจ. (สํขิตฺตํ. สหชาตวาราทิ วิตฺถาเรตพฺโพ.)}

\proseitem{21}{เจตสิโก กุสโล ธมฺโม เจตสิกสฺส กุสลสฺส ธมฺมสฺส เหตุปจฺจเยน ปจฺจโย. (สํขิตฺตํ.)}

\proseitem{22}{เหตุยา ตีณิ, อารมฺมเณ นว, อธิปติยา นว, (อาทิมฺหิ ตีณิ, สหชาตา\-ธิปติ, มชฺฌิเมสุ ตีสุ มชฺฌิมา\-นุโลมิกาเยว ปญฺหา, สหชาตา\-ธิปติ,) อนนฺตเร นว ฯเปฯ สหชาเต ปญฺจ ฯเปฯ อุปนิสฺสเย นว, อาเสวเน นว, กมฺเม ตีณิ, อาหาเร ปญฺจ, อินฺทฺริเย ปญฺจ, ฌาเน ตีณิ, มคฺเค ตีณิ, สมฺปยุตฺเต ปญฺจ, อตฺถิยา ปญฺจ, นตฺถิยา นว ฯเปฯ อวิคเต ปญฺจ. (สํขิตฺตํ.)}

\proseitem{23}{\csromanfolio{400}เจตสิกํ อกุสลํ ธมฺมํ ปฏิจฺจ เจตสิโก อกุสโล ธมฺโม อุปฺปชฺชติ เหตุปจฺจยา.  เจตสิกํ อกุสลํ ธมฺมํ ปฏิจฺจ อเจตสิโก อกุสโล ธมฺโม อุปฺปชฺชติ เหตุปจฺจยา.  เจตสิกํ อกุสลํ ธมฺมํ ปฏิจฺจ เจตสิโก อกุสโล จ อเจตสิโก อกุสโล จ ธมฺมา อุปฺปชฺชนฺติ เหตุปจฺจยา. (3)}

\prose{อเจตสิกํ อกุสลํ ธมฺมํ ปฏิจฺจ เจตสิโก อกุสโล ธมฺโม อุปฺปชฺชติ เหตุปจฺจยา. (1)}

\prose{เจตสิกํ อกุสลญฺจ อเจตสิกํ อกุสลญฺจ ธมฺมํ ปฏิจฺจ เจตสิโก อกุสโล ธมฺโม อุปฺปชฺชติ เหตุปจฺจยา. (1) (สํขิตฺตํ.)}

\proseitem{24}{เหตุยา ปญฺจ, อารมฺมเณ ปญฺจ ฯเปฯ อวิคเต ปญฺจ. (สํขิตฺตํ.) (ยถา กุสลนยํ เอวํ นเหตุปจฺจยมฺปิ กาตพฺพํ, สหชาตวารมฺปิ ฯเปฯ ปญฺหาวารมฺปิ วิตฺถาเรตพฺพํ.)}

\proseitem{25}{เจตสิกํ อพฺยากตํ ธมฺมํ ปฏิจฺจ เจตสิโก อพฺยากโต ธมฺโม อุปฺปชฺชติ เหตุปจฺจยา.  เจตสิกํ อพฺยากตํ ธมฺมํ ปฏิจฺจ อเจตสิโก อพฺยากโต ธมฺโม อุปฺปชฺชติ เหตุปจฺจยา.  เจตสิกํ อพฺยากตํ ธมฺมํ ปฏิจฺจ เจตสิโก อพฺยากโต จ อเจตสิโก อพฺยากโต จ ธมฺมา อุปฺปชฺชนฺติ เหตุปจฺจยา. (3)}

\prose{อเจตสิกํ อพฺยากตํ ธมฺมํ ปฏิจฺจ อเจตสิโก อพฺยากโต ธมฺโม อุปฺปชฺชติ เหตุปจฺจยา. ตีณิ.}

\prose{เจตสิกํ อพฺยากตญฺจ อเจตสิกํ อพฺยากตญฺจ ธมฺมํ ปฏิจฺจ เจตสิโก อพฺยากโต ธมฺโม อุปฺปชฺชติ เหตุปจฺจยา. ตีณิ. (สํขิตฺตํ.)}

\proseitem{26}{เหตุยา นว, อารมฺมเณ นว, อธิปติยา นว ฯเปฯ ปุเรชาเต อาเสวเน ปญฺจ, กมฺเม นว, วิปาเก นว ฯเปฯ อวิคเต นว. (สํขิตฺตํ.)}

\prose{นเหตุยา นว, นอารมฺมเณ ตีณิ, นอธิปติยา นว, (สพฺเพ กตพฺพา) ฯเปฯ นกมฺเม จตฺตาริ, นวิปาเก นว, นอาหาเร เอกํ, นอินฺทฺริเย เอกํ, นฌาเน ฉ, นมคฺเค นว, นสมฺปยุตฺเต ตีณิ, นวิปฺปยุตฺเต ฉ ฯเปฯ โนวิคเต ตีณิ. (สํขิตฺตํ. สหชาตวาราทิ วิตฺถาเรตพฺโพ.)}

\csromanmarkh{58. จิตฺต\-สมฺปยุตฺตทุก}\csromanmarkhii{1. กุสลตฺติก}\csromanheadernomark{2}{58. จิตฺต\-สมฺปยุตฺตทุก 1. กุสลตฺติก}
\proseitem{27}{\csromanfolio{401}เจตสิโก อพฺยากโต ธมฺโม เจตสิกสฺส อพฺยากตสฺส ธมฺมสฺส เหตุปจฺจเยน ปจฺจโย. (สํขิตฺตํ.)}

\proseitem{28}{เหตุยา ตีณิ, อารมฺมเณ นว, อธิปติยา นว, (เอตฺถ ฉสุ สหชาตา\-ธิปติ.) อนนฺตเร นว ฯเปฯ สหชาเต นว, อญฺญมญฺเญ นว, นิสฺสเย นว, อุปนิสฺสเย นว, ปุเรชาเต ตีณิ, อาเสวเน นว, กมฺเม ตีณิ, วิปาเก นว, อาหาเร นว, อินฺทฺริเย นว, ฌาเน ตีณิ, มคฺเค ตีณิ, สมฺปยุตฺเต ปญฺจ, วิปฺปยุตฺเต ปญฺจ, อตฺถิยา นว ฯเปฯ อวิคเต นว. (สํขิตฺตํ.)}

\csromanmarkh{58. จิตฺต\-สมฺปยุตฺตทุก}\csromanmarkhii{1. กุสลตฺติก}\csromanheadernomark{2}{58. จิตฺต\-สมฺปยุตฺตทุก 1. กุสลตฺติก}
\csromanheader{2}{1-7. ปฏิจฺจาทิวาร ปจฺจยจตุกฺกเหตุ}
\proseitem{29}{จิตฺต\-สมฺปยุตฺตํ กุสลํ ธมฺมํ ปฏิจฺจ จิตฺต\-สมฺปยุตฺโต กุสโล ธมฺโม อุปฺปชฺชติ เหตุปจฺจยา. (สํขิตฺตํ.)}

\proseitem{30}{เหตุยา เอกํ, อารมฺมเณ เอกํ ฯเปฯ อวิคเต เอกํ. (สํขิตฺตํ.) (สหชาตวาเรปิ ฯเปฯ ปญฺหาวาเรปิ สพฺพตฺถ เอกํ.)}

\proseitem{31}{จิตฺต\-สมฺปยุตฺตํ อกุสลํ ธมฺมํ ปฏิจฺจ จิตฺต\-สมฺปยุตฺโต อกุสโล ธมฺโม อุปฺปชฺชติ เหตุปจฺจยา. (สํขิตฺตํ.)}

\proseitem{32}{เหตุยา เอกํ, อารมฺมเณ เอกํ ฯเปฯ อวิคเต เอกํ. (สํขิตฺตํ.) (สหชาตวาเรปิ ฯเปฯ ปญฺหาวาเรปิ สพฺพตฺถ เอกํ.)}

\prose{อพฺยากตปทเหตุ}

\proseitem{33}{จิตฺต\-สมฺปยุตฺตํ อพฺยากตํ ธมฺมํ ปฏิจฺจ จิตฺต\-สมฺปยุตฺโต อพฺยากโต ธมฺโม อุปฺปชฺชติ เหตุปจฺจยา.  จิตฺต\-สมฺปยุตฺตํ อพฺยากตํ ธมฺมํ ปฏิจฺจ จิตฺต\-วิปฺปยุตฺโต อพฺยากโต ธมฺโม อุปฺปชฺชติ เหตุปจฺจยา.  จิตฺต\-สมฺปยุตฺตํ อพฺยากตํ ธมฺมํ ปฏิจฺจ จิตฺต\-สมฺปยุตฺโต อพฺยากโต จ จิตฺต\-วิปฺปยุตฺโต อพฺยากโต จ ธมฺมา อุปฺปชฺชนฺติ เหตุปจฺจยา. (3)}

\prose{จิตฺต\-วิปฺปยุตฺตํ อพฺยากตํ ธมฺมํ ปฏิจฺจ จิตฺต\-วิปฺปยุตฺโต อพฺยากโต ธมฺโม อุปฺปชฺชติ เหตุปจฺจยา.  จิตฺต\-วิปฺปยุตฺตํ อพฺยากตํ ธมฺมํ ปฏิจฺจ จิตฺตธมฺมานุโลม \csromanfolio{402}ทุกติก\-ปฏฺฐาน\-ปาฬิ สมฺปยุตฺโต อพฺยากโต ธมฺโม อุปฺปชฺชติ เหตุปจฺจยา.  จิตฺต\-วิปฺปยุตฺตํ อพฺยากตํ ธมฺมํ ปฏิจฺจ จิตฺต\-สมฺปยุตฺโต อพฺยากโต จ จิตฺต\-วิปฺปยุตฺโต อพฺยากโต จ ธมฺมา อุปฺปชฺชนฺติ เหตุปจฺจยา. (3)}

\prose{จิตฺต\-สมฺปยุตฺตํ อพฺยากตญฺจ จิตฺต\-วิปฺปยุตฺตํ อพฺยากตญฺจ ธมฺมํ ปฏิจฺจ จิตฺต\-สมฺปยุตฺโต อพฺยากโต ธมฺโม อุปฺปชฺชติ เหตุปจฺจยา.  จิตฺต\-สมฺปยุตฺตํ อพฺยากตญฺจ จิตฺต\-วิปฺปยุตฺตํ อพฺยากตญฺจ ธมฺมํ ปฏิจฺจ จิตฺต\-วิปฺปยุตฺโต อพฺยากโต ธมฺโม อุปฺปชฺชติ เหตุปจฺจยา.  จิตฺต\-สมฺปยุตฺตํ อพฺยากตญฺจ จิตฺต\-วิปฺปยุตฺตํ อพฺยากตญฺจ ธมฺมํ ปฏิจฺจ จิตฺต\-สมฺปยุตฺโต อพฺยากโต จ จิตฺต\-วิปฺปยุตฺโต อพฺยากโต จ ธมฺมา อุปฺปชฺชนฺติ เหตุปจฺจยา. (3) (สํขิตฺตํ.)}

\proseitem{34}{เหตุยา นว, อารมฺมเณ ตีณิ, อธิปติยา ปญฺจ ฯเปฯ อญฺญมญฺเญ ฉ ฯเปฯ ปุเรชาเต เอกํ, อาเสวเน เอกํ, กมฺเม นว, วิปาเก นว ฯเปฯ อวิคเต นว. (สํขิตฺตํ.)}

\prose{นเหตุยา นว, นอารมฺมเณ ตีณิ, นอธิปติยา นว ฯเปฯ นปุเรชาเต นว ฯเปฯ นกมฺเม เทฺว, นวิปาเก ปญฺจ, นอาหาเร เอกํ, นอินฺทฺริเย เอกํ, นฌาเน เทฺว, นมคฺเค นว, นสมฺปยุตฺเต ตีณิ, นวิปฺปยุตฺเต เทฺว. (สํขิตฺตํ. สหชาตวาราทิ วิตฺถาเรตพฺโพ.)}

\prose{เหตุอารมฺมณ}

\proseitem{35}{จิตฺต\-สมฺปยุตฺโต อพฺยากโต ธมฺโม จิตฺต\-สมฺปยุตฺตสฺส อพฺยากตสฺส ธมฺมสฺส เหตุปจฺจเยน ปจฺจโย. ตีณิ.}

\prose{จิตฺต\-สมฺปยุตฺโต อพฺยากโต ธมฺโม จิตฺต\-สมฺปยุตฺตสฺส อพฺยากตสฺส ธมฺมสฺส อารมฺมณ\-ปจฺจเยน ปจฺจโย. (1)}

\prose{จิตฺต\-วิปฺปยุตฺโต อพฺยากโต ธมฺโม จิตฺต\-สมฺปยุตฺตสฺส อพฺยากตสฺส ธมฺมสฺส อารมฺมณ\-ปจฺจเยน ปจฺจโย. (1) (สํขิตฺตํ.)}

\proseitem{36}{เหตุยา ตีณิ, อารมฺมเณ เทฺว, อธิปติยา จตฺตาริ, อนนฺตเร เอกํ ฯเปฯ สหชาเต สตฺต, อญฺญมญฺเญ ฉ, นิสฺสเย สตฺต, อุปนิสฺสเย เทฺว, ปุเรชาเต เอกํ, ปจฺฉาชาเต เอกํ, อาเสวเน เอกํ, กมฺเม ตีณิ, วิปาเก ตีณิ, อาหาเร จตฺตาริ, อินฺทฺริเย ฉ, ฌาเน ตีณิ, มคฺเค ตีณิ, สมฺปยุตฺเต เอกํ, วิปฺปยุตฺเต เทฺว. (สํขิตฺตํ.)}

\vspace{\tipitakaparskip}
\csromancenter{จิตฺต\-สมฺปยุตฺตทุก 1. กุสลตฺติก 59. จิตฺต\-สํสฏฺฐทุก 1. กุสลตฺติก}
\prose{\csromanfolio{403}1. ปฏิจฺจวาร ปจฺจยจตุกฺกเหตุ}

\proseitem{37}{จิตฺต\-สํสฏฺฐํ กุสลํ ธมฺมํ ปฏิจฺจ จิตฺตสํสฏฺโฐ กุสโล ธมฺโม อุปฺปชฺชติ เหตุปจฺจยา. (สํขิตฺตํ.)}

\proseitem{38}{เหตุยา เอกํ, อารมฺมเณ เอกํ ฯเปฯ อวิคเต เอกํ. (สํขิตฺตํ.) (สหชาตวาเรปิ ฯเปฯ ปญฺหาวาเรปิ สพฺพตฺถ เอกํ.)}

\proseitem{39}{จิตฺต\-สํสฏฺฐํ อกุสลํ ธมฺมํ ปฏิจฺจ จิตฺตสํสฏฺโฐ อกุสโล ธมฺโม อุปฺปชฺชติ เหตุปจฺจยา. (สํขิตฺตํ.)}

\proseitem{40}{เหตุยา เอกํ, อารมฺมเณ เอกํ ฯเปฯ อวิคเต เอกํ. (สํขิตฺตํ.) (สหชาตวาเรปิ ฯเปฯ ปญฺหาวาเรปิ สพฺพตฺถ เอกํ.)}

\proseitem{41}{จิตฺต\-สํสฏฺฐํ อพฺยากตํ ธมฺมํ ปฏิจฺจ จิตฺตสํสฏฺโฐ อพฺยากโต ธมฺโม อุปฺปชฺชติ เหตุปจฺจยา. ตีณิ.}

\prose{โนจิตฺต\-สํสฏฺฐํ อพฺยากตํ ธมฺมํ ปฏิจฺจ โนจิตฺต\-สํสฏฺโฐ อพฺยากโต ธมฺโม อุปฺปชฺชติ เหตุปจฺจยา. ตีณิ.}

\prose{จิตฺต\-สํสฏฺฐํ อพฺยากตญฺจ โนจิตฺต\-สํสฏฺฐํ อพฺยากตญฺจ ธมฺมํ ปฏิจฺจ จิตฺตสํสฏฺโฐ อพฺยากโต ธมฺโม อุปฺปชฺชติ เหตุปจฺจยา. ตีณิ. (สํขิตฺตํ.)}

\proseitem{42}{เหตุยา นว, อารมฺมเณ ตีณิ ฯเปฯ อวิคเต นว. (สํขิตฺตํ.)}

\prose{(ยถา จิตฺต\-สมฺปยุตฺตทุกํ อพฺยากต\-สทิสํ. สหชาตวารมฺปิ ฯเปฯ ปญฺหาวารมฺปิ วิตฺถาเรตพฺพํ.)}

\csromanmarkh{60. จิตฺต\-สมุฏฺฐานทุก}\csromanmarkhii{1. กุสลตฺติก}\csromanheadernomark{2}{60. จิตฺต\-สมุฏฺฐานทุก 1. กุสลตฺติก}
\csromanheader{2}{1-7. ปฏิจฺจาทิวาร ปจฺจยจตุกฺกเหตุ}
\proseitem{43}{จิตฺต\-สมุฏฺฐานํ กุสลํ ธมฺมํ ปฏิจฺจ จิตฺต\-สมุฏฺฐาโน กุสโล ธมฺโม อุปฺปชฺชติ เหตุปจฺจยา.  จิตฺต\-สมุฏฺฐานํ กุสลํ ธมฺมํ ปฏิจฺจ โนจิตฺต\-สมุฏฺฐาโน กุสโล ธมฺโม อุปฺปชฺชติ เหตุปจฺจยา.  จิตฺต\-สมุฏฺฐานํ กุสลํ ธมฺมํ ปฏิจฺจ \csromanfolio{404}จิตฺต\-สมุฏฺฐาโน กุสโล จ โนจิตฺต\-สมุฏฺฐาโน กุสโล จ ธมฺมา อุปฺปชฺชนฺติ เหตุปจฺจยา. (3)}

\prose{โนจิตฺต\-สมุฏฺฐานํ กุสลํ ธมฺมํ ปฏิจฺจ จิตฺต\-สมุฏฺฐาโน กุสโล ธมฺโม อุปฺปชฺชติ เหตุปจฺจยา. (1)}

\prose{จิตฺต\-สมุฏฺฐานํ กุสลญฺจ โนจิตฺต\-สมุฏฺฐานํ กุสลญฺจ ธมฺมํ ปฏิจฺจ จิตฺต\-สมุฏฺฐาโน กุสโล ธมฺโม อุปฺปชฺชติ เหตุปจฺจยา. (1) (สํขิตฺตํ.)}

\proseitem{44}{เหตุยา ปญฺจ, อารมฺมเณ ปญฺจ, (สพฺพตฺถ ปญฺจ.) อวิคเต ปญฺจ. (สํขิตฺตํ.)}

\prose{นอธิปติยา ปญฺจ, นปุเรชาเต ปญฺจ ฯเปฯ นกมฺเม ตีณิ, นวิปาเก ปญฺจ ฯเปฯ นวิปฺปยุตฺเต ปญฺจ. (สํขิตฺตํ. สหชาตวาราทิ วิตฺถาเรตพฺโพ.)}

\proseitem{45}{จิตฺต\-สมุฏฺฐาโน กุสโล ธมฺโม จิตฺต\-สมุฏฺฐานสฺส กุสลสฺส ธมฺมสฺส เหตุปจฺจเยน ปจฺจโย. (สํขิตฺตํ.)}

\proseitem{46}{เหตุยา ตีณิ, อารมฺมเณ นว, อธิปติยา นว, อนนฺตเร นว, สมนนฺตเร นว, สหชาเต ปญฺจ ฯเปฯ อุปนิสฺสเย นว, อาเสวเน นว, กมฺเม ตีณิ, อาหาเร ปญฺจ, อินฺทฺริเย ปญฺจ, ฌาเน ตีณิ, มคฺเค ตีณิ, สมฺปยุตฺเต ปญฺจ, อตฺถิยา ปญฺจ ฯเปฯ อวิคเต ปญฺจ. (สํขิตฺตํ.)}

\proseitem{47}{จิตฺต\-สมุฏฺฐานํ อกุสลํ ธมฺมํ ปฏิจฺจ จิตฺต\-สมุฏฺฐาโน อกุสโล ธมฺโม อุปฺปชฺชติ เหตุปจฺจยา. ตีณิ.}

\prose{โนจิตฺต\-สมุฏฺฐานํ อกุสลํ ธมฺมํ ปฏิจฺจ จิตฺต\-สมุฏฺฐาโน อกุสโล ธมฺโม อุปฺปชฺชติ เหตุปจฺจยา. (1)}

\prose{จิตฺต\-สมุฏฺฐานํ อกุสลญฺจ โนจิตฺต\-สมุฏฺฐานํ อกุสลญฺจ ธมฺมํ ปฏิจฺจ จิตฺต\-สมุฏฺฐาโน อกุสโล ธมฺโม อุปฺปชฺชติ เหตุปจฺจยา. (1) (สํขิตฺตํ.)}

\proseitem{48}{เหตุยา ปญฺจ, อารมฺมเณ ปญฺจ ฯเปฯ อวิคเต ปญฺจ. (สํขิตฺตํ.)}

\prose{(กุสล\-ตฺติก\-สทิสํ สหชาตวารมฺปิ ฯเปฯ ปญฺหาวารมฺปิ วิตฺถาเรตพฺพํ.)}

\proseitem{49}{จิตฺต\-สมุฏฺฐานํ อพฺยากตํ ธมฺมํ ปฏิจฺจ จิตฺต\-สมุฏฺฐาโน อพฺยากโต ธมฺโม อุปฺปชฺชติ เหตุปจฺจยา. ตีณิ.}

\prose{โนจิตฺต\-สมุฏฺฐานํ อพฺยากตํ ธมฺมํ ปฏิจฺจ โนจิตฺต\-สมุฏฺฐาโน อพฺยากโต ธมฺโม อุปฺปชฺชติ เหตุปจฺจยา. ตีณิ.}

\csromanmarkh{61. จิตฺต\-สหภูทุก}\csromanmarkhii{1. กุสลตฺติก}\csromanheadernomark{2}{61. จิตฺต\-สหภูทุก 1. กุสลตฺติก}
\prose{\csromanfolio{405}จิตฺต\-สมุฏฺฐานํ อพฺยากตญฺจ โนจิตฺต\-สมุฏฺฐานํ อพฺยากตญฺจ ธมฺมํ ปฏิจฺจ จิตฺต\-สมุฏฺฐาโน อพฺยากโต ธมฺโม อุปฺปชฺชติ เหตุปจฺจยา. ตีณิ. (สํขิตฺตํ.)}

\proseitem{50}{เหตุยา นว, อารมฺมเณ นว, อธิปติยา ปญฺจ, (สพฺพตฺถ นว.) ปุเรชาเต ปญฺจ, อาเสวเน ปญฺจ ฯเปฯ อวิคเต นว. (สํขิตฺตํ.)}

\prose{นเหตุยา นว, นอารมฺมเณ ฉ, นอธิปติยา นว ฯเปฯ นกมฺเม ตีณิ, นวิปาเก ฉ, นอาหาเร เอกํ, นอินฺทฺริเย เอกํ, นฌาเน ฉ, นมคฺเค นว, นสมฺปยุตฺเต ฉ, นวิปฺปยุตฺเต ฉ, โนนตฺถิยา ฉ, โนวิคเต ฉ. (สํขิตฺตํ. สหชาตวารมฺปิ ฯเปฯ สมฺปยุตฺตวารมฺปิ วิตฺถาเรตพฺพํ.)}

\proseitem{51}{จิตฺต\-สมุฏฺฐาโน อพฺยากโต ธมฺโม จิตฺต\-สมุฏฺฐานสฺส อพฺยากตสฺส ธมฺมสฺส เหตุปจฺจเยน ปจฺจโย. ตีณิ.}

\prose{จิตฺต\-สมุฏฺฐาโน อพฺยากโต ธมฺโม จิตฺต\-สมุฏฺฐานสฺส อพฺยากตสฺส ธมฺมสฺส อารมฺมณ\-ปจฺจเยน ปจฺจโย. (สํขิตฺตํ.)}

\proseitem{52}{เหตุยา ตีณิ, อารมฺมเณ นว, อธิปติยา นว, อนนฺตเร นว, (สพฺพตฺถ นว.) ปุเรชาเต นว, ปจฺฉาชาเต นว, อาเสวเน นว, กมฺเม ตีณิ, วิปาเก นว, อาหาเร นว, (จิตฺต\-สมุฏฺฐานมูลํ โนจิตฺต\-สมุฏฺฐานสฺส กพฬีกาโร อาหาโร กาตพฺโพ. โนจิตฺต\-สมุฏฺฐาโน จิตฺต\-สมุฏฺฐานสฺส กพฬีกาโร อาหาโร ฆฏเน มชฺเฌ กพฬีกาโร อาหาโร.) อินฺทฺริเย นว, (รูปชีวตินฺทฺริยํ เอกํ) ฌาเน ตีณิ, มคฺเค ตีณิ, สมฺปยุตฺเต ปญฺจ, วิปฺปยุตฺเต นว, อตฺถิยา นว, นตฺถิยา นว, วิคเต นว, อวิคเต นว. (สํขิตฺตํ.)}

\prose{(ยถา กุสลตฺติเก ปญฺหาวารสฺส อนุโลมมฺปิ ปจฺจนียมฺปิ อนุโลมปจฺจนียมฺปิ ปจฺจนียานุโลมมฺปิ คณิตํ, เอวํ คเณตพฺพํ.)}

\csromanmarkh{61. จิตฺต\-สหภูทุก}\csromanmarkhii{1. กุสลตฺติก}\csromanheadernomark{2}{61. จิตฺต\-สหภูทุก 1. กุสลตฺติก}
\prose{1. ปฏิจฺจวาร ปจฺจยจตุกฺกเหตุ}

\proseitem{53}{จิตฺตสหภุํ กุสลํ ธมฺมํ ปฏิจฺจ จิตฺตสหภู กุสโล ธมฺโม อุปฺปชฺชติ เหตุปจฺจยา.  จิตฺตสหภุํ กุสลํ ธมฺมํ ปฏิจฺจ โนจิตฺตสหภู กุสโล \csromanfolio{406}ธมฺโม อุปฺปชฺชติ เหตุปจฺจยา.  จิตฺตสหภุํ กุสลํ ธมฺมํ ปฏิจฺจ จิตฺตสหภู กุสโล จ โนจิตฺตสหภู กุสโล จ ธมฺมา อุปฺปชฺชนฺติ เหตุปจฺจยา. (3)}

\prose{โนจิตฺต\-สหภุํ กุสลํ ธมฺมํ ปฏิจฺจ จิตฺตสหภู กุสโล ธมฺโม อุปฺปชฺชติ เหตุปจฺจยา. (1)}

\prose{จิตฺตสหภุํ กุสลญฺจ โนจิตฺต\-สหภุํ กุสลญฺจ ธมฺมํ ปฏิจฺจ จิตฺตสหภู กุสโล ธมฺโม อุปฺปชฺชติ เหตุปจฺจยา. (1) (สํขิตฺตํ.)}

\proseitem{54}{เหตุยา ปญฺจ, อารมฺมเณ ปญฺจ ฯเปฯ อวิคเต ปญฺจ. (สํขิตฺตํ.) (สหชาตวารมฺปิ ฯเปฯ ปญฺหาวารมฺปิ วิตฺถาเรตพฺพํ.)}

\proseitem{55}{จิตฺตสหภุํ อกุสลํ ธมฺมํ ปฏิจฺจ จิตฺตสหภู อกุสโล ธมฺโม อุปฺปชฺชติ เหตุปจฺจยา. ตีณิ.}

\prose{โนจิตฺต\-สหภุํ อกุสลํ ธมฺมํ ปฏิจฺจ จิตฺตสหภู อกุสโล ธมฺโม อุปฺปชฺชติ เหตุปจฺจยา. (1)}

\prose{จิตฺตสหภุํ อกุสลญฺจ โนจิตฺต\-สหภุํ อกุสลญฺจ ธมฺมํ ปฏิจฺจ จิตฺตสหภู อกุสโล ธมฺโม อุปฺปชฺชติ เหตุปจฺจยา. (1) (สํขิตฺตํ.)}

\proseitem{56}{เหตุยา ปญฺจ, อารมฺมเณ ปญฺจ ฯเปฯ อวิคเต ปญฺจ. (สํขิตฺตํ.) (สหชาตวารมฺปิ ฯเปฯ ปญฺหาวารมฺปิ วิตฺถาเรตพฺพํ.)}

\proseitem{57}{จิตฺตสหภุํ อพฺยากตํ ธมฺมํ ปฏิจฺจ จิตฺตสหภู อพฺยากโต ธมฺโม อุปฺปชฺชติ เหตุปจฺจยา. ตีณิ.}

\prose{โนจิตฺต\-สหภุํ อพฺยากตํ ธมฺมํ ปฏิจฺจ โนจิตฺตสหภู อพฺยากโต ธมฺโม อุปฺปชฺชติ เหตุปจฺจยา. ตีณิ.}

\prose{จิตฺตสหภุํ อพฺยากตญฺจ โนจิตฺต\-สหภุํ อพฺยากตญฺจ ธมฺมํ ปฏิจฺจ จิตฺตสหภู อพฺยากโต ธมฺโม อุปฺปชฺชติ เหตุปจฺจยา. ตีณิ. (สํขิตฺตํ.)}

\proseitem{58}{เหตุยา นว, อารมฺมเณ นว ฯเปฯ อวิคเต นว. (สํขิตฺตํ. ยถา เจตสิก\-ทุก\-มูลา ตีณิ คมนา, เอวํ อิเมปิ ตีณิ คมนา กาตพฺพา. สหชาตวารมฺปิ ฯเปฯ ปญฺหาวารมฺปิ วิตฺถาเรตพฺพํ.)}

\vspace{\tipitakaparskip}
\csromancenter{จิตฺตา\-นุปริวตฺติทุก 1. กุสลตฺติก 62. จิตฺตา\-นุปริวตฺติทุก 1. กุสลตฺติก}
\prose{\csromanfolio{407}1. ปฏิจฺจวาร ปจฺจยจตุกฺกเหตุ}

\proseitem{59}{จิตฺตา\-นุปริวตฺติํ กุสลํ ธมฺมํ ปฏิจฺจ จิตฺตา\-นุปริวตฺตี กุสโล ธมฺโม อุปฺปชฺชติ เหตุปจฺจยา. ตีณิ.}

\prose{โนจิตฺตา\-นุปริวตฺติํ กุสลํ ธมฺมํ ปฏิจฺจ จิตฺตา\-นุปริวตฺตี กุสโล ธมฺโม อุปฺปชฺชติ เหตุปจฺจยา. (1)}

\prose{จิตฺตา\-นุปริวตฺติํ กุสลญฺจ โนจิตฺตา\-นุปริวตฺติํ กุสลญฺจ ธมฺมํ ปฏิจฺจ จิตฺตา\-นุปริวตฺตี กุสโล ธมฺโม อุปฺปชฺชติ เหตุปจฺจยา. (1) (สํขิตฺตํ.)}

\proseitem{60}{เหตุยา ปญฺจ, อารมฺมเณ ปญฺจ ฯเปฯ อวิคเต ปญฺจ. (สํขิตฺตํ.) (สหชาตวารมฺปิ ฯเปฯ ปญฺหาวารมฺปิ วิตฺถาเรตพฺพํ.)}

\proseitem{61}{จิตฺตา\-นุปริวตฺติํ อกุสลํ ธมฺมํ ปฏิจฺจ จิตฺตา\-นุปริวตฺตี อกุสโล ธมฺโม อุปฺปชฺชติ เหตุปจฺจยา. ตีณิ.}

\prose{โนจิตฺตา\-นุปริวตฺติํ อกุสลํ ธมฺมํ ปฏิจฺจ จิตฺตา\-นุปริวตฺตี อกุสโล ธมฺโม อุปฺปชฺชติ เหตุปจฺจยา. (1)}

\prose{จิตฺตา\-นุปริวตฺติํ อกุสลญฺจ โนจิตฺตา\-นุปริวตฺติํ อกุสลญฺจ ธมฺมํ ปฏิจฺจ จิตฺตา\-นุปริวตฺตี อกุสโล ธมฺโม อุปฺปชฺชติ เหตุปจฺจยา. (1) (สํขิตฺตํ.)}

\proseitem{62}{เหตุยา ปญฺจ, อารมฺมเณ ปญฺจ ฯเปฯ อวิคเต ปญฺจ. (สํขิตฺตํ.) (สหชาตวารมฺปิ ฯเปฯ ปญฺหาวารมฺปิ วิตฺถาเรตพฺพํ.)}

\proseitem{63}{จิตฺตา\-นุปริวตฺติํ อพฺยากตํ ธมฺมํ ปฏิจฺจ จิตฺตา\-นุปริวตฺตี อพฺยากโต ธมฺโม อุปฺปชฺชติ เหตุปจฺจยา. (สํขิตฺตํ.)}

\proseitem{64}{เหตุยา นว, อารมฺมเณ นว ฯเปฯ อวิคเต นว. (สํขิตฺตํ. เจตสิก\-ทุก\-สทิสํ. สหชาตวารมฺปิ ฯเปฯ ปญฺหาวารมฺปิ วิตฺถาเรตพฺพํ.)}

\vspace{\tipitakaparskip}
\csromancenter{63. จิตฺต\-สํสฏฺฐ\-สมุฏฺฐานทุก 1. กุสลตฺติก}
\csromanheader{2}{1-7. ปฏิจฺจาทิวาร ปจฺจยจตุกฺกเหตุ}
\proseitem{65}{\csromanfolio{408}จิตฺต\-สํสฏฺฐ\-สมุฏฺฐานํ กุสลํ ธมฺมํ ปฏิจฺจ จิตฺต\-สํสฏฺฐ\-สมุฏฺฐาโน กุสโล ธมฺโม อุปฺปชฺชติ เหตุปจฺจยา.  จิตฺต\-สํสฏฺฐ\-สมุฏฺฐานํ กุสลํ ธมฺมํ ปฏิจฺจ โนจิตฺต\-สํสฏฺฐ\-สมุฏฺฐาโน กุสโล ธมฺโม อุปฺปชฺชติ เหตุปจฺจยา.  จิตฺต\-สํสฏฺฐ\-สมุฏฺฐานํ กุสลํ ธมฺมํ ปฏิจฺจ จิตฺต\-สํสฏฺฐ\-สมุฏฺฐาโน กุสโล จ โนจิตฺต\-สํสฏฺฐ\-สมุฏฺฐาโน กุสโล จ ธมฺมา อุปฺปชฺชนฺติ เหตุปจฺจยา. (3)}

\prose{โนจิตฺต\-สํสฏฺฐ\-สมุฏฺฐานํ กุสลํ ธมฺมํ ปฏิจฺจ จิตฺต\-สํสฏฺฐ\-สมุฏฺฐาโน กุสโล ธมฺโม อุปฺปชฺชติ เหตุปจฺจยา. (1)}

\prose{จิตฺต\-สํสฏฺฐ\-สมุฏฺฐานํ กุสลญฺจ โนจิตฺต\-สํสฏฺฐ\-สมุฏฺฐานํ กุสลญฺจ ธมฺมํ ปฏิจฺจ จิตฺต\-สํสฏฺฐ\-สมุฏฺฐาโน กุสโล ธมฺโม อุปฺปชฺชติ เหตุปจฺจยา. (1) (สํขิตฺตํ.)}

\proseitem{66}{เหตุยา ปญฺจ, อารมฺมเณ ปญฺจ ฯเปฯ อวิคเต ปญฺจ. (สํขิตฺตํ. ยถา มหนฺตรทุเก เจตสิก\-ทุก\-กุสล\-สทิสํ, ตตฺตกา เอว ปญฺหา. สหชาตวารมฺปิ ฯเปฯ ปญฺหาวารมฺปิ วิตฺถาเรตพฺพํ.)}

\proseitem{67}{จิตฺต\-สํสฏฺฐ\-สมุฏฺฐานํ อกุสลํ ธมฺมํ ปฏิจฺจ จิตฺต\-สํสฏฺฐ\-สมุฏฺฐาโน อกุสโล ธมฺโม อุปฺปชฺชติ เหตุปจฺจยา. ตีณิ.}

\prose{โนจิตฺต\-สํสฏฺฐ\-สมุฏฺฐานํ อกุสลํ ธมฺมํ ปฏิจฺจ จิตฺต\-สํสฏฺฐ\-สมุฏฺฐาโน อกุสโล ธมฺโม อุปฺปชฺชติ เหตุปจฺจยา. (1)}

\prose{จิตฺต\-สํสฏฺฐ\-สมุฏฺฐานํ อกุสลญฺจ โนจิตฺต\-สํสฏฺฐ\-สมุฏฺฐานํ อกุสลญฺจ ธมฺมํ ปฏิจฺจ จิตฺต\-สํสฏฺฐ\-สมุฏฺฐาโน อกุสโล ธมฺโม อุปฺปชฺชติ เหตุปจฺจยา. (1) (สํขิตฺตํ.)}

\proseitem{68}{เหตุยา ปญฺจ, อารมฺมเณ ปญฺจ ฯเปฯ อวิคเต ปญฺจ. (สํขิตฺตํ. เจตสิก\-ทุก\-อกุสล\-สทิสํ. สหชาตวารมฺปิ ฯเปฯ ปญฺหาวารมฺปิ วิตฺถาเรตพฺพํ.)}

\proseitem{69}{จิตฺต\-สํสฏฺฐ\-สมุฏฺฐานํ อพฺยากตํ ธมฺมํ ปฏิจฺจ จิตฺต\-สํสฏฺฐ\-สมุฏฺฐาโน อพฺยากโต ธมฺโม อุปฺปชฺชติ เหตุปจฺจยา. ตีณิ.}

\prose{\csromanfolio{409}64. จิตฺต\-สํสฏฺฐ\-สมุฏฺฐาน\-สหภูทุก 1. กุสลตฺติก}

\prose{โนจิตฺต\-สํสฏฺฐ\-สมุฏฺฐานํ อพฺยากตํ ธมฺมํ ปฏิจฺจ โนจิตฺต\-สํสฏฺฐ\-สมุฏฺฐาโน อพฺยากโต ธมฺโม อุปฺปชฺชติ เหตุปจฺจยา. ตีณิ.}

\prose{จิตฺต\-สํสฏฺฐ\-สมุฏฺฐานํ อพฺยากตญฺจ โนจิตฺต\-สํสฏฺฐ\-สมุฏฺฐานํ อพฺยากตญฺจ ธมฺมํ ปฏิจฺจ จิตฺต\-สํสฏฺฐ\-สมุฏฺฐาโน อพฺยากโต ธมฺโม อุปฺปชฺชติ เหตุปจฺจยา. ตีณิ. (สํขิตฺตํ.)}

\proseitem{70}{เหตุยา นว, อารมฺมเณ นว, อธิปติยา นว ฯเปฯ ปุเรชาเต ปญฺจ, อาเสวเน ปญฺจ, กมฺเม นว ฯเปฯ อวิคเต นว. (สํขิตฺตํ.)}

\prose{นเหตุยา นว, นอารมฺมเณ ตีณิ ฯเปฯ นกมฺเม จตฺตาริ, นวิปาเก นว, นอาหาเร เอกํ, นอินฺทฺริเย เอกํ, นฌาเน ฉ, นมคฺเค นว, นสมฺปยุตฺเต ตีณิ, นวิปฺปยุตฺเต ฉ, โนนตฺถิยา ตีณิ, โนวิคเต ตีณิ. (สํขิตฺตํ. สหชาตวารมฺปิ ฯเปฯ สมฺปยุตฺตวารมฺปิ วิตฺถาเรตพฺพํ.)}

\proseitem{71}{จิตฺต\-สํสฏฺฐ\-สมุฏฺฐาโน อพฺยากโต ธมฺโม จิตฺต\-สํสฏฺฐ\-สมุฏฺฐานสฺส อพฺยากตสฺส ธมฺมสฺส เหตุปจฺจเยน ปจฺจโย. (สํขิตฺตํ.)}

\proseitem{72}{เหตุยา ตีณิ, อารมฺมเณ นว, อธิปติยา นว, (เอตฺถ ฉสุ สหชาตา\-ธิปติ.) อนนฺตเร นว, (สพฺพตฺถ นว.) ปุเรชาเต ตีณิ, ปจฺฉาชาเต ตีณิ, อาเสวเน นว, กมฺเม ตีณิ, วิปาเก นว, อาหาเร นว, อินฺทฺริเย นว, ฌาเน ตีณิ, มคฺเค ตีณิ, สมฺปยุตฺเต ปญฺจ, วิปฺปยุตฺเต ปญฺจ, อตฺถิยา นว ฯเปฯ อวิคเต นว. (สํขิตฺตํ.)}

\prose{นเหตุยา นว, นอารมฺมเณ นว. (สํขิตฺตํ.)}

\prose{เหตุปจฺจยา นอารมฺมเณ ตีณิ. (สํขิตฺตํ.)}

\prose{นเหตุปจฺจยา อารมฺมเณ นว. (สํขิตฺตํ.)}

\prose{(ยถา กุสลตฺติเก ปญฺหาวารสฺส อนุโลมมฺปิ ปจฺจนียมฺปิ อนุโลมปจฺจนียมฺปิ ปจฺจนียานุโลมมฺปิ คณิตํ, เอวํ คเณตพฺพํ.)}

\prose{64. จิตฺต\-สํสฏฺฐ\-สมุฏฺฐาน\-สหภูทุก 1. กุสลตฺติก}

\prose{1. ปฏิจฺจวาร ปจฺจยจตุกฺกเหตุ}

\proseitem{73}{จิตฺต\-สํสฏฺฐ\-สมุฏฺฐาน\-สหภุํ กุสลํ ธมฺมํ ปฏิจฺจ จิตฺต\-สํสฏฺฐ\-สมุฏฺฐาน\-สหภู กุสโล ธมฺโม อุปฺปชฺชติ เหตุปจฺจยา. ตีณิ.}

\prose{\csromanfolio{410}โนจิตฺต\-สํสฏฺฐ\-สมุฏฺฐาน\-สหภุํ กุสลํ ธมฺมํ ปฏิจฺจ จิตฺต\-สํสฏฺฐ\-สมุฏฺฐาน\-สหภู กุสโล ธมฺโม อุปฺปชฺชติ เหตุปจฺจยา. (1)}

\prose{จิตฺต\-สํสฏฺฐ\-สมุฏฺฐาน\-สหภุํ กุสลญฺจ โนจิตฺต\-สํสฏฺฐ\-สมุฏฺฐาน\-สหภุํ กิสลญฺจ ธมฺมํ ปฏิจฺจ จิตฺต\-สํสฏฺฐ\-สมุฏฺฐาน\-สหภู กุสโล ธมฺโม อุปฺปชฺชติ เหตุปจฺจยา. (1) (สํขิตฺตํ.)}

\proseitem{74}{เหตุยา ปญฺจ, อารมฺมเณ ปญฺจ ฯเปฯ อวิคเต ปญฺจ. (สํขิตฺตํ.) (สหชาตวารมฺปิ ฯเปฯ ปญฺหาวารมฺปิ วิตฺถาเรตพฺพํ.)}

\proseitem{75}{จิตฺต\-สํสฏฺฐ\-สมุฏฺฐาน\-สหภุํ อกุสลํ ธมฺมํ ปฏิจฺจ จิตฺต\-สํสฏฺฐ\-สมุฏฺฐาน\-สหภู อกุสโล ธมฺโม อุปฺปชฺชติ เหตุปจฺจยา. ตีณิ.}

\prose{โนจิตฺต\-สํสฏฺฐ\-สมุฏฺฐาน\-สหภุํ อกุสลํ ธมฺมํ ปฏิจฺจ จิตฺต\-สํสฏฺฐ\-สมุฏฺฐาน\-สหภู อกุสโล ธมฺโม อุปฺปชฺชติ เหตุปจฺจยา. (1)}

\prose{จิตฺต\-สํสฏฺฐ\-สมุฏฺฐาน\-สหภุํ อกุสลญฺจ โนจิตฺต\-สํสฏฺฐ\-สมุฏฺฐาน\-สหภุํ อกุสลญฺจ ธมฺมํ ปฏิจฺจ จิตฺต\-สํสฏฺฐ\-สมุฏฺฐาน\-สหภู อกุสโล ธมฺโม อุปฺปชฺชติ เหตุปจฺจยา. (1) (สํขิตฺตํ.)}

\proseitem{76}{เหตุยา ปญฺจ, อารมฺมเณ ปญฺจ ฯเปฯ อวิคเต ปญฺจ. (สํขิตฺตํ.) (สหชาตวารมฺปิ ฯเปฯ ปญฺหาวารมฺปิ วิตฺถาเรตพฺพํ.)}

\proseitem{77}{จิตฺต\-สํสฏฺฐ\-สมุฏฺฐาน\-สหภุํ อพฺยากตํ ธมฺมํ ปฏิจฺจ จิตฺต\-สํสฏฺฐ\-สมุฏฺฐาน\-สหภู อพฺยากโต ธมฺโม อุปฺปชฺชติ เหตุปจฺจยา. (สํขิตฺตํ.)}

\proseitem{78}{เหตุยา นว, อารมฺมเณ นว ฯเปฯ อวิคเต นว. (สํขิตฺตํ. จิตฺต\-สํสฏฺฐ\-สมุฏฺฐาน\-ทุก\-อพฺยากต\-สทิสํ. สหชาตวารมฺปิ ฯเปฯ ปญฺหาวารมฺปิ วิตฺถาเรตพฺพํ.)}

\prose{65. จิตฺต\-สํสฏฺฐ\-สมุฏฺฐานา\-นุปริวตฺติทุก 1. กุสลตฺติก}

\prose{1. ปฏิจฺจวาร ปจฺจยจตุกฺกเหตุ}

\proseitem{79}{จิตฺตสํสฏฺฐสมุฏฺฐฺอานานุปริวตฺติํ กุสลํ ธมฺมํ ปฏิจฺจ จิตฺต\-สํสฏฺฐ\-สมุฏฺฐานา\-นุปริวตฺตี กุสโล ธมฺโม อุปฺปชฺชติ เหตุปจฺจยา. ตีณิ.}

\csromanmarkh{66. อชฺฌตฺติกทุก}\csromanmarkhii{1. กุสลตฺติก}\csromanheadernomark{2}{66. อชฺฌตฺติกทุก 1. กุสลตฺติก}
\prose{\csromanfolio{411}โนจิตฺต\-สํสฏฺฐ\-สมุฏฺฐานา\-นุปริวตฺติํ กุสลํ ธมฺมํ ปฏิจฺจ จิตฺต\-สํสฏฺฐ\-สมุฏฺฐานา\-นุปริวตฺตี กุสโล ธมฺโม อุปฺปชฺชติ เหตุปจฺจยา. (1)}

\prose{จิตฺต\-สํสฏฺฐ\-สมุฏฺฐานา\-นุปริวตฺติํ กุสลญฺจ โนจิตฺต\-สํสฏฺฐ\-สมุฏฺฐานา\-นุปริวตฺติํ กุสลญฺจ ธมฺมํ ปฏิจฺจ จิตฺต\-สํสฏฺฐ\-สมุฏฺฐานา\-นุปริวตฺตี กุสโล ธมฺโม อุปฺปชฺชติ เหตุปจฺจยา. (1) (สํขิตฺตํ.)}

\proseitem{80}{เหตุยา ปญฺจ, อารมฺมเณ ปญฺจ ฯเปฯ อวิคเต ปญฺจ. (สํขิตฺตํ.) (สหชาตวารมฺปิ ฯเปฯ ปญฺหาวารมฺปิ วิตฺถาเรตพฺพํ.)}

\proseitem{81}{จิตฺต\-สํสฏฺฐ\-สมุฏฺฐานา\-นุปริวตฺติํ อกุสลํ ธมฺมํ ปฏิจฺจ จิตฺต\-สํสฏฺฐ\-สมุฏฺฐานา\-นุปริวตฺตี อกุสโล ธมฺโม อุปฺปชฺชติ เหตุปจฺจยา. (สํขิตฺตํ.)}

\proseitem{82}{เหตุยา ปญฺจ, อรมฺมเณ ปญฺจ ฯเปฯ อวิคเต ปญฺจ. (สํขิตฺตํ.) (สหชาตวารมฺปิ ฯเปฯ ปญฺหาวารมฺปิ วิตฺถาเรตพฺพํ.)}

\proseitem{83}{จิตฺต\-สํสฏฺฐ\-สมุฏฺฐานา\-นุปริวตฺติํ อพฺยากตํ ธมฺมํ ปฏิจฺจ จิตฺต\-สํสฏฺฐ\-สมุฏฺฐานา\-นุปริวตฺตี อพฺยากโต ธมฺโม อุปฺปชฺชติ เหตุปจฺจยา. (สํขิตฺตํ.)}

\prose{เหตุยา นว, อารมฺมเณ นว ฯเปฯ อวิคเต นว. (สํขิตฺตํ. จิตฺต\-สํสฏฺฐ\-สมุฏฺฐาน\-ทุก\-อพฺยากต\-สทิสํ. สหชาตวารมฺปิ ฯเปฯ ปญฺหาวารมฺปิ วิตฺถาเรตพฺพํ.)}

\csromanmarkh{66. อชฺฌตฺติกทุก}\csromanmarkhii{1. กุสลตฺติก}\csromanheadernomark{2}{66. อชฺฌตฺติกทุก 1. กุสลตฺติก}
\csromanheader{2}{1-7. ปฏิจฺจาทิวาร ปจฺจยจตุกฺกเหตุ}
\proseitem{84}{อชฺฌตฺติกํ กุสลํ ธมฺมํ ปฏิจฺจ พาหิโร กุสโล ธมฺโม อุปฺปชฺชติ เหตุปจฺจยา. (1)}

\prose{พาหิรํ กุสลํ ธมฺมํ ปฏิจฺจ พาหิโร กุสโล ธมฺโม อุปฺปชฺชติ เหตุปจฺจยา.  พาหิรํ กุสลํ ธมฺมํ ปฏิจฺจ อชฺฌตฺติโก กุสโล ธมฺโม อุปฺปชฺชติ เหตุปจฺจยา.  พาหิรํ กุสลํ ธมฺมํ ปฏิจฺจ อชฺฌตฺติโก กุสโล จ พาหิโร กุสโล จ ธมฺมา อุปฺปชฺชนฺติ เหตุปจฺจยา. (3)}

\prose{\csromanfolio{412}อชฺฌตฺติกํ กุสลญฺจ พาหิรํ กุสลญฺจ ธมฺมํ ปฏิจฺจ พาหิโร กุสโล ธมฺโม อุปฺปชฺชติ เหตุปจฺจยา. (1) (สํขิตฺตํ.)}

\proseitem{85}{เหตุยา ปญฺจ, อารมฺมเณ ปญฺจ ฯเปฯ อวิคเต ปญฺจ. (สํขิตฺตํ. จิตฺต\-ทุก\-สทิสํ. สหชาตวารมฺปิ ฯเปฯ ปญฺหาวารมฺปิ วิตฺถาเรตพฺพํ.)}

\proseitem{86}{อชฺฌตฺติกํ อกุสลํ ธมฺมํ ปฏิจฺจ พาหิโร อกุสโล ธมฺโม อุปฺปชฺชติ เหตุปจฺจยา. (1)}

\prose{พาหิรํ อกุสลํ ธมฺมํ ปฏิจฺจ พาหิโร อกุสโล ธมฺโม อุปฺปชฺชติ เหตุปจฺจยา.  พาหิรํ อกุสลํ ธมฺมํ ปฏิจฺจ อชฺฌตฺติโก อกุสโล ธมฺโม อุปฺปชฺชติ เหตุปจฺจยา.  พาหิรํ อกุสลํ ธมฺมํ ปฏิจฺจ อชฺฌตฺติโก อกุสโล จ พาหิโร อกุสโล จ ธมฺมา อุปฺปชฺชนฺติ เหตุปจฺจยา. (3)}

\prose{อชฺฌตฺติกํ อกุสลญฺจ พาหิรํ อกุสลญฺจ ธมฺมํ ปฏิจฺจ พาหิโร อกุสโล ธมฺโม อุปฺปชฺชติ เหตุปจฺจยา. (1) (สํขิตฺตํ.)}

\proseitem{87}{เหตุยา ปญฺจ, อารมฺมเณ ปญฺจ ฯเปฯ อวิคเต ปญฺจ. (สํขิตฺตํ. จิตฺต\-ทุก\-สทิสํ. สหชาตวารมฺปิ ฯเปฯ ปญฺหาวารมฺปิ วิตฺถาเรตพฺพํ.)}

\proseitem{88}{อชฺฌตฺติกํ อพฺยากตํ ธมฺมํ ปฏิจฺจ อชฺฌตฺติโก อพฺยากโต ธมฺโม อุปฺปชฺชติ เหตุปจฺจยา. ตีณิ.}

\prose{พาหิรํ อพฺยากตํ ธมฺมํ ปฏิจฺจ พาหิโร อพฺยากโต ธมฺโม อุปฺปชฺชติ เหตุปจฺจยา. ตีณิ.}

\prose{อชฺฌตฺติกํ อพฺยากตญฺจ พาหิรํ อพฺยากตญฺจ ธมฺมํ ปฏิจฺจ อชฺฌตฺติโก อพฺยากโต ธมฺโม อุปฺปชฺชติ เหตุปจฺจยา. ตีณิ. (สํขิตฺตํ.)}

\proseitem{89}{เหตุยา นว, อารมฺมเณ ปญฺจ, อธิปติยา ปญฺจ ฯเปฯ อญฺญมญฺเญ ปญฺจ, นิสฺสเย นว, อุปนิสฺสเย ปญฺจ, ปุเรชาเต ปญฺจ, อาเสวเน ปญฺจ, กมฺเม ฯเปฯ มคฺเค นว, สมฺปยุตฺเต ปญฺจ, วิปฺปยุตฺเต นว ฯเปฯ อวิคเต นว. (สํขิตฺตํ.)}

\prose{นเหตุยา นว, นอารมฺมเณ นว, นอธิปติยา นว ฯเปฯ นกมฺเม ตีณิ, นวิปาเก ปญฺจ, นอาหาเร เอกํ, นอินฺทฺริเย เอกํ, นฌาเน ปญฺจ, นมคฺเค นว, นสมฺปยุตฺเต นว, นวิปฺปยุตฺเต ปญฺจ, โนนตฺถิยา นว, โนวิคเต นว. (สํขิตฺตํ. สหชาตวารมฺปิ ฯเปฯ สมฺปยุตฺตวารมฺปิ วิตฺถาเรตพฺพํ.)}

\csromanmarkh{67. อุปาทาทุก}\csromanmarkhii{1. กุสลตฺติก}\csromanheadernomark{2}{67. อุปาทาทุก 1. กุสลตฺติก}
\proseitem{90}{\csromanfolio{413}พาหิโร อพฺยากโต ธมฺโม พาหิรสฺส อพฺยากตสฺส ธมฺมสฺส เหตุปจฺจเยน ปจฺจโย. (สํขิตฺตํ.)}

\prose{พาหิโร อพฺยากโต ธมฺโม พาหิรสฺส อพฺยากตสฺส ธมฺมสฺส กมฺมปจฺจเยน ปจฺจโย. ตีณิ.}

\proseitem{91}{เหตุยา ตีณิ, อารมฺมเณ นว, อธิปติยา นว, (อชฺฌตฺติโก พาหิรสฺส คมนกาเล สหชาตา\-ธิปติ, มชฺเฌ ตีสุปิ สหชาตา\-ธิปติ.) อนนฺตเร นว, สหชาเต นว, อญฺญมญฺเญ ปญฺจ, นิสฺสเย นว, อุปนิสฺสเย นว, ปุเรชาเต นว, (อารมฺมณปุเรชาตมฺปิ วตฺถุปุเรชาตมฺปิ.) ปจฺฉาชาเต นว, อาเสวเน นว, กมฺเม ตีณิ, วิปาเก นว, อาหาเร นว, (ติณฺณนฺนํ กพฬีกาโร อาหาโร.) อินฺทฺริเย นว, (ติณฺณนฺนํ รูปชีวิตินฺทฺริยํ.) ฌาเน มคฺเค ตีณิ, สมฺปยุตฺเต ปญฺจ, วิปฺปยุตฺเต นว, อตฺถิยา นว ฯเปฯ อวิคเต นว. (สํขิตฺตํ.)}

\csromanmarkh{67. อุปาทาทุก}\csromanmarkhii{1. กุสลตฺติก}\csromanheadernomark{2}{67. อุปาทาทุก 1. กุสลตฺติก}
\csromanheader{2}{1-7. ปฏิจฺจาทิวาร ปจฺจยจตุกฺกเหตุ}
\proseitem{92}{โนอุปาทา กุสลํ ธมฺมํ ปฏิจฺจ โนอุปาทา กุสโล ธมฺโม อุปฺปชฺชติ เหตุปจฺจยา. (สํขิตฺตํ.)}

\prose{เหตุยา เอกํ, อารมฺมเณ เอกํ ฯเปฯ อวิคเต เอกํ. (สํขิตฺตํ.) (สหชาตวาเรปิ ฯเปฯ ปญฺหาวาเรปิ สพฺพตฺถ เอกํ.)}

\proseitem{93}{โนอุปาทา อกุสลํ ธมฺมํ ปฏิจฺจ โนอุปาทา อกุสโล ธมฺโม อุปฺปชฺชติ เหตุปจฺจยา. (สํขิตฺตํ.)}

\prose{เหตุยา เอกํ, อารมฺมเณ เอกํ ฯเปฯ อวิคเต เอกํ. (สํขิตฺตํ.) (สหชาตวาเรปิ ฯเปฯ ปญฺหาวาเรปิ สพฺพตฺถ เอกํ.)}

\proseitem{94}{อุปาทา อพฺยากตํ ธมฺมํ ปฏิจฺจ โนอุปาทา อพฺยากโต ธมฺโม อุปฺปชฺชติ เหตุปจฺจยา. (1)}

\prose{โนอุปาทา อพฺยากตํ ธมฺมํ ปฏิจฺจ โนอุปาทา อพฺยากโต ธมฺโม อุปฺปชฺชติ เหตุปจฺจยา.  โนอุปาทา อพฺยากตํ ธมฺมํ ปฏิจฺจ อุปาทา \csromanfolio{414}อพฺยากโต ธมฺโม อุปฺปชฺชติ เหตุปจฺจยา.  โนอุปาทา อพฺยากตํ ธมฺมํ ปฏิจฺจ อุปาทา อพฺยากโต จ โนอุปาทา อพฺยากโต จ ธมฺมา อุปฺปชฺชนฺติ เหตุปจฺจยา. (3)}

\prose{อุปาทา อพฺยากตญฺจ โนอุปาทา อพฺยากตญฺจ ธมฺมํ ปฏิจฺจ โนอุปาทา อพฺยากโต ธมฺโม อุปฺปชฺชติ เหตุปจฺจยา. (1) (สํขิตฺตํ.)}

\prose{เหตุยา ปญฺจ, อารมฺมเณ ตีณิ ฯเปฯ อญฺญมญฺเญ ปญฺจ ฯเปฯ ปุเรชาเต เอกํ, อาเสวเน เอกํ, กมฺเม ปญฺจ, วิปาเก ปญฺจ ฯเปฯ อวิคเต ปญฺจ. (สํขิตฺตํ.)}

\prose{นเหตุยา ปญฺจ, นอารมฺมเณ ตีณิ, นอธิปติยา ปญฺจ ฯเปฯ นปุเรชาเต ปญฺจ ฯเปฯ นกมฺเม ตีณิ, นวิปาเก ตีณิ, นอาหาเร ตีณิ, นอินฺทฺริเย ตีณิ, นฌาเน ตีณิ, นมคฺเค ปญฺจ, นสมฺปยุตฺเต ตีณิ, นวิปฺปยุตฺเต ตีณิ, โนนตฺถิยา ตีณิ, โนวิคเต ตีณิ. (สํขิตฺตํ. สหชาตวารมฺปิ ฯเปฯ สมฺปยุตฺตวารมฺปิ วิตฺถาเรตพฺพํ.)}

\csromanheader{2}{\textbf{เหตุ อารมฺมณาทิ}}
\proseitem{96}{โนอุปาทา อพฺยากโต ธมฺโม โนอุปาทา อพฺยากตสฺส ธมฺมสฺส เหตุปจฺจเยน ปจฺจโย. ตีณิ.}

\prose{อุปาทา อพฺยากโต ธมฺโม โนอุปาทา อพฺยากตสฺส ธมฺมสฺส อารมฺมณ\-ปจฺจเยน ปจฺจโย. (1)}

\prose{โนอุปาทา อพฺยากโต ธมฺโม โนอุปาทา อพฺยากตสฺส ธมฺมสฺส อารมฺมณ\-ปจฺจเยน ปจฺจโย. (1)}

\prose{โนอุปาทา อพฺยากโต ธมฺโม โนอุปาทา อพฺยากตสฺส ธมฺมสฺส อธิปติ\-ปจฺจเยน ปจฺจโย. ตีณิ. (ปฐเม เทฺวปิ อธิปตี. ทฺวีสุ สหชาตา\-ธิปติ.)}

\prose{อุปาทา อพฺยากโต ธมฺโม โนอุปาทา อพฺยากตสฺส ธมฺมสฺส อุปนิสฺสย\-ปจฺจเยน ปจฺจโย. (1)}

\prose{โนอุปาทา อพฺยากโต ธมฺโม โนอุปาทา อพฺยากตสฺส ธมฺมสฺส อุปนิสฺสย\-ปจฺจเยน ปจฺจโย. (1) (สํขิตฺตํ.)}

\csromanmatikaanchor{mtk.77}
\csromanmatikaanchor{mtk.78}
\csromantocmarknopage{subparagraph}{68. อุปาทินฺนทุก}{mtk.77}
\bookmark[dest={mtk.77},level=5]{68. อุปาทินฺนทุก}
\csromantocmarkref{subparagraph}{1. กุสลตฺติก ...}{mtk.78}
\bookmark[dest={mtk.78},level=5]{1. กุสลตฺติก ...}
\csromanmarkh{68. อุปาทินฺนทุก}\csromanmarkhii{1. กุสลตฺติก}\csromanheadernomark{6}{68. อุปาทินฺนทุก 1. กุสลตฺติก}
\proseitem{97}{\csromanfolio{415}อุปนิสฺสเย เทฺว, ปุเรชาเต ตีณิ, ปจฺฉาชาเต ตีณิ, อาเสวเน เอกํ, กมฺเม ตีณิ, วิปาเก ตีณิ, อาหาเร ฉ, อินฺทฺริเย สตฺต, ฌาเน ตีณิ, มคฺเค ตีณิ, สมฺปยุตฺเต เอกํ, วิปฺปยุตฺเต จตฺตาริ, อตฺถิยา นว ฯเปฯ อวิคเต นว. (สํขิตฺตํ.)}

\csromanmarkh{68. อุปาทินฺนทุก}\csromanmarkhii{1. กุสลตฺติก}\csromanheadernomark{2}{68. อุปาทินฺนทุก 1. กุสลตฺติก}
\csromanheader{2}{1-7. ปฏิจฺจาทิวาร ปจฺจยจตุกฺกเหตุ}
\proseitem{98}{อนุปาทินฺนํ กุสลํ ธมฺมํ ปฏิจฺจ อนุปาทินฺโน กุสโล ธมฺโม อุปฺปชฺชติ เหตุปจฺจยา. (สํขิตฺตํ.)}

\prose{เหตุยา เอกํ, อารมฺมเณ เอกํ ฯเปฯ อวิคเต เอกํ. (สพฺพตฺถ เอกํ. สํขิตฺตํ. สหชาตวาเรปิ ฯเปฯ ปญฺหาวาเรปิ เอกํ, สปฺปจฺจย\-กุสล\-สทิสํ.)}

\proseitem{99}{อนุปาทินฺนํ อกุสลํ ธมฺมํ ปฏิจฺจ อนุปาทินฺโน อกุสโล ธมฺโม อุปฺปชฺชติ เหตุปจฺจยา. (สํขิตฺตํ.)}

\prose{เหตุยา เอกํ ฯเปฯ อวิคเต เอกํ. (สํขิตฺตํ. สหชาตวาเรปิ ฯเปฯ ปญฺหาวาเรปิ สพฺพตฺถ เอกํ.}

\proseitem{100}{อุปาทินฺนํ อพฺยากตํ ธมฺมํ ปฏิจฺจ อุปาทินฺโน อพฺยากโต ธมฺโม อุปฺปชฺชติ เหตุปจฺจยา. ตีณิ.}

\prose{อนุปาทินฺนํ อพฺยากตํ ธมฺมํ ปฏิจฺจ อนุปาทินฺโน อพฺยากโต ธมฺโม อุปฺปชฺชติ เหตุปจฺจยา. (1)}

\prose{อุปาทินฺนํ อพฺยากตญฺจ อนุปาทินฺนํ อพฺยากตญฺจ ธมฺมํ ปฏิจฺจ อนุปาทินฺโน อพฺยากโต ธมฺโม อุปฺปชฺชติ เหตุปจฺจยา. อุปาทินฺเน ขนฺเธ จ มหาภูเต จ ปฏิจฺจ จิตฺต\-สมุฏฺฐานํ รูปํ. (1) (สํขิตฺตํ.)}

\proseitem{101}{เหตุยา ปญฺจ, อารมฺมเณ เทฺว, อธิปติยา เอกํ ฯเปฯ ปุเรชาเต เทฺว, อาเสวเน เอกํ, กมฺเม ปญฺจ, วิปาเก ปญฺจ ฯเปฯ อวิคเต ปญฺจ. (สํขิตฺตํ.)}

\prose{\csromanfolio{416}นเหตุยา ปญฺจ, นอารมฺมเณ จตฺตาริ, นอธิปติยา ปญฺจ ฯเปฯ นปุเรชาเต จตฺตาริ, นปจฺฉาชาเต นอาเสวเน ปญฺจ, นกมฺเม เอกํ, นวิปาเก เทฺว, นอาหาเร เทฺว, นอินฺทฺริเย เทฺว, นฌาเน เทฺว, นมคฺเค ปญฺจ, นสมฺปยุตฺเต จตฺตาริ, นวิปฺปยุตฺเต เทฺว, โนนตฺถิยา จตฺตาริ, โนวิคเต จตฺตาริ. (สํขิตฺตํ. สหชาตวาราทิ วิตฺถาเรตพฺโพ.)}

\csromanheader{2}{\textbf{เหตุ ปุเรชาต}}
\proseitem{102}{อุปาทินฺโน อพฺยากโต ธมฺโม อุปาทินฺนสฺส อพฺยากตสฺส ธมฺมสฺส เหตุปจฺจเยน ปจฺจโย. ตีณิ.}

\prose{อนุปาทินฺโน อพฺยากโต ธมฺโม อนุปาทินฺนสฺส อพฺยากตสฺส ธมฺมสฺส เหตุปจฺจเยน ปจฺจโย. (1)}

\prose{อนุปาทินฺโน อพฺยากโต ธมฺโม อนุปาทินฺนสฺส อพฺยากตสฺส ธมฺมสฺส อธิปติ\-ปจฺจเยน ปจฺจโย. อารมฺมณา\-ธิปติ, สหชาตา\-ธิปติ. เอกํ.}

\prose{อุปาทินฺโน อพฺยากโต ธมฺโม อุปาทินฺนสฺส อพฺยากตสฺส ธมฺมสฺส ปุเรชาต\-ปจฺจเยน ปจฺจโย. อารมฺมณ\-ปุเรชาตํ, วตฺถุ\-ปุเรชาตํ. (เทฺว ปญฺหา.)}

\prose{อนุปาทินฺโน อพฺยากโต ธมฺโม อนุปาทินฺนสฺส อพฺยากตสฺส ธมฺมสฺส ปุเรชาต\-ปจฺจเยน ปจฺจโย. เทฺว. (อารมฺมณปุเรชาตํเยว. ฆฏนา เทฺว, อารมฺมณปุเรชาตมฺปิ วตฺถุปุเรชาตมฺปิ.)}

\proseitem{103}{เหตุยา จตฺตาริ, อารมฺมเณ จตฺตาริ, (อุปาทินฺนมูลเก เทฺว. อนุปาทินฺน\-มูลเก เทฺว.) อธิปติยา เอกํ, อนนฺตเร สมนนฺตเร จตฺตาริ, สหชาเต ปญฺจ, อญฺญมญฺเญ เทฺว, นิสฺสเย ปญฺจ, อุปนิสฺสเย จตฺตาริ, ปุเรชาเต ฉ, ปจฺฉาชาเต ฉ, อาเสวเน เอกํ, กมฺเม จตฺตาริ, วิปาเก จตฺตาริ, อาหาเร นว, อินฺทฺริเย จตฺตาริ, ฌาเน จตฺตาริ, มคฺเค จตฺตาริ, สมฺปยุตฺเต เทฺว, วิปฺปยุตฺเต ฉ, อตฺถิยา นว, นตฺถิยา จตฺตาริ, วิคเต จตฺตาริ, อวิคเต นว. (สํขิตฺตํ.)}

\prose{นเหตุยา นว, นอารมฺมเณ นว. (สํขิตฺตํ.)}

\prose{เหตุปจฺจยา นอารมฺมเณ จตฺตาริ. (สํขิตฺตํ.)}

\prosewithcloser{นเหตุปจฺจยา อารมฺมเณ จตฺตาริ. (สํขิตฺตํ.)}{มหนฺตร\-ทุก\-กุสล\-ตฺติกํ นิฏฺฐิตํ.}
\csromanmatikaanchor{mtk.79}
\csromanmatikaanchor{mtk.80}
\csromantocmarknopage{subparagraph}{69. อุปาทานทุก}{mtk.79}
\bookmark[dest={mtk.79},level=5]{69. อุปาทานทุก}
\csromantocmarkref{subparagraph}{1. กุสลตฺติก ...}{mtk.80}
\bookmark[dest={mtk.80},level=5]{1. กุสลตฺติก ...}
\csromanmarkh{69. อุปาทานทุก}\csromanmarkhii{1. กุสลตฺติก 69. อุปาทานทุก 1. กุสลตฺติก}\csromanheadernomark{6}{69. อุปาทานทุก 1. กุสลตฺติก \textbf{69. อุปาทานทุก 1. กุสลตฺติก}}
\csromanheader{2}{1-7. ปฏิจฺจาทิวาร ปจฺจยจตุกฺก}
\proseitem{1}{\csromanfolio{417}โนอุปาทานํ กุสลํ ธมฺมํ ปฏิจฺจ โนอุปาทาโน กุสโล ธมฺโม อุปฺปชฺชติ เหตุปจฺจยา. (สํขิตฺตํ.)}

\prose{เหตุยา เอกํ, อารมฺมเณ เอกํ ฯเปฯ อวิคเต เอกํ. (สํขิตฺตํ.) (สหชาตวาเรปิ ฯเปฯ ปญฺหาวาเรปิ สพฺพตฺถ เอกํ.)}

\proseitem{2}{อุปาทานํ อกุสลํ ธมฺมํ ปฏิจฺจ อุปาทาโน อกุสโล ธมฺโม อุปฺปชฺชติ เหตุปจฺจยา. อุปาทานํ อกุสลํ ธมฺมํ ปฏิจฺจ โนอุปาทาโน อกุสโล ธมฺโม อุปฺปชฺชติ เหตุปจฺจยา.  อุปาทานํ อกุสลํ ธมฺมํ ปฏิจฺจ อุปาทาโน อกุสโล จ โนอุปาทาโน อกุสโล จ ธมฺมา อุปฺปชฺชนฺติ เหตุปจฺจยา. (3)}

\prose{โนอุปาทานํ อกุสลํ ธมฺมํ ปฏิจฺจ โนอุปาทาโน อกุสโล ธมฺโม อุปฺปชฺชติ เหตุปจฺจยา. ตีณิ.}

\prose{อุปาทานํ อกุสลญฺจ โนอุปาทานํ อกุสลญฺจ ธมฺมํ ปฏิจฺจ อุปาทาโน อกุสโล ธมฺโม อุปฺปชฺชติ เหตุปจฺจยา. ตีณิ. (สํขิตฺตํ.)}

\proseitem{3}{เหตุยา นว, อารมฺมเณ นว, (สพฺพตฺถ นว,) อวิคเต นว. (สํขิตฺตํ.)}

\proseitem{4}{โนอุปาทานํ อกุสลํ ธมฺมํ ปฏิจฺจ โนอุปาทาโน อกุสโล ธมฺโม อุปฺปชฺชติ นเหตุปจฺจยา. (สํขิตฺตํ.)}

\proseitem{5}{นเหตุยา เอกํ, นอธิปติยา นว, นปุเรชาเต นว ฯเปฯ นกมฺเม ตีณิ, นวิปาเก นว, นวิปฺปยุตฺเต นว. (สํขิตฺตํ. สหชาตวาราทิ วิตฺถาเตตพฺโพ.)}

\proseitem{6}{อุปาทาโน อกุสโล ธมฺโม อุปาทานสฺส อกุสลสฺส ธมฺมสฺส เหตุปจฺจเยน ปจฺจโย. (สํขิตฺตํ.)}

\proseitem{7}{เหตุยา นว, อารมฺมเณ นว, อธิปติยา นว, (โนอุปาทาน\-มูลเก ตีณิ. สหชาตา\-ธิปติ.) อนนฺตเร นว ฯเปฯ นิสฺสเย นว, \csromanfolio{418}อุปนิสฺสเย นว, อาเสวเน นว, กมฺเม ตีณิ, (โนอุปาทาน\-มูลเก.) อาหาเร ตีณิ, อินฺทฺริเย ตีณิ, ฌาเน ตีณิ, มคฺเค นว, สมฺปยุตฺเต นว, อตฺถิยา นว ฯเปฯ อวิคเต นว. (สํขิตฺตํ.)}

\proseitem{8}{โนอุปาทานํ อพฺยากตํ ธมฺมํ ปฏิจฺจ โนอุปาทาโน อพฺยากโต ธมฺโม อุปฺปชฺชติ เหตุปจฺจยา. (สํขิตฺตํ.)}

\proseitem{9}{เหตุยา เอกํ, อารมฺมเณ เอกํ ฯเปฯ อวิคเต เอกํ. (สํขิตฺตํ.) (สหชาตวาเรปิ ฯเปฯ ปญฺหาวาเรปิ สพฺพตฺถ เอกํ.)}

\csromanmarkh{70. อุปาทานิยทุก}\csromanmarkhii{1. กุสลตฺติก}\csromanheadernomark{2}{70. \textbf{อุปาทานิยทุก 1. กุสลตฺติก}}
\prose{1. ปฏิจฺจวาร ปจฺจยจตุกฺกเหตุ}

\proseitem{7}{อุปาทานิยํ กุสลํ ธมฺมํ ปฏิจฺจ อุปาทานิโย กุสโล ธมฺโม อุปฺปชฺชติ เหตุปจฺจยา. (1)}

\prose{อนุปาทานิยํ กุสลํ ธมฺมํ ปฏิจฺจ อนุปาทานิโย กุสโล ธมฺโม อุปฺปชฺชติ เหตุปจฺจยา. (1) (สํขิตฺตํ.)}

\proseitem{11}{เหตุยา เทฺว, อารมฺมเณ เทฺว ฯเปฯ อวิคเต เทฺว. (สํขิตฺตํ. โลกิยโลกุตฺตร\-ทุก\-กุสล\-สทิสํ. สหชาตวารมฺปิ ฯเปฯ ปญฺหาวารมฺปิ วิตฺถาเรตพฺพํ.)}

\proseitem{12}{อุปาทานิยํ อกุสลํ ธมฺมํ ปฏิจฺจ อุปาทานิโย อกุสโล ธมฺโม อุปฺปชฺชติ เหตุปจฺจยา. (สํขิตฺตํ.)}

\prose{เหตุยา เอกํ, อารมฺมเณ เอกํ ฯเปฯ อวิคเต เอกํ. (สํขิตฺตํ.) (สหชาตวาเรปิ ฯเปฯ ปญฺหาวาเรปิ สพฺพตฺถ เอกํ.)}

\proseitem{13}{อุปาทานิยํ อพฺยากตํ ธมฺมํ ปฏิจฺจ อุปาทานิโย อพฺยากโต ธมฺโม อุปฺปชฺชติ เหตุปจฺจยา. (1)}

\prose{อนุปาทานิยํ อพฺยากตํ ธมฺมํ ปฏิจฺจ อนุปาทานิโย อพฺยากโต ธมฺโม อุปฺปชฺชติ เหตุปจฺจยา.  อนุปาทานิยํ อพฺยากตํ ธมฺมํ ปฏิจฺจ อุปาทานิโย}

\vspace{\tipitakaparskip}
\csromancenter{อุปาทาน\-สมฺปยุตฺตทุก 1. กุสลตฺติก อพฺยากโต ธมฺโม อุปฺปชฺชติ เหตุปจฺจยา.  อนุปาทานิยํ อพฺยากตํ ธมฺมํ ปฏิจฺจ อุปาทานิโย อพฺยากโต จ อนุปาทานิโย อพฺยากโต จ ธมฺมา อุปฺปชฺชนฺติ เหตุปจฺจยา. (3)}
\prose{\csromanfolio{419}อุปาทานิยํ อพฺยากตญฺจ อนุปาทานิยํ อพฺยากตญฺจ ธมฺมํ ปฏิจฺจ อุปาทานิโย อพฺยากโต ธมฺโม อุปฺปชฺชติ เหตุปจฺจยา. (1) (สํขิตฺตํ.)}

\proseitem{14}{เหตุยา ปญฺจ, อารมฺมเณ เทฺว ฯเปฯ อาเสวเน เอกํ ฯเปฯ อวิคเต ปญฺจ. (สํขิตฺตํ. โลกิย\-ทุก\-อพฺยากต\-สทิสํ. สหชาตวารมฺปิ ฯเปฯ ปญฺหาวารมฺปิ วิตฺถาเรตพฺพํ.)}

\csromanmarkh{71. อุปาทาน\-สมฺปยุตฺตทุก}\csromanmarkhii{1. กุสลตฺติก}\csromanheadernomark{2}{71. อุปาทาน\-สมฺปยุตฺตทุก 1. กุสลตฺติก}
\csromanheader{2}{1-7. ปฏิจฺจาทิวาร ปจฺจยจตุกฺกเหตุ}
\proseitem{15}{อุปาทาน\-วิปฺปยุตฺตํ กุสลํ ธมฺมํ ปฏิจฺจ อุปาทาน\-วิปฺปยุตฺโต กุสโล ธมฺโม อุปฺปชฺชติ เหตุปจฺจยา. (สํขิตฺตํ.)}

\prose{เหตุยา เอกํ, อารมฺมเณ เอกํ ฯเปฯ อวิคเต เอกํ. (สํขิตฺตํ.) (สหชาตวาเรปิ ฯเปฯ ปญฺหาวาเรปิ สพฺพตฺถ เอกํ.)}

\proseitem{16}{อุปาทาน\-สมฺปยุตฺตํ อกุสลํ ธมฺมํ ปฏิจฺจ อุปาทาน\-สมฺปยุตฺโต อกุสโล ธมฺโม อุปฺปชฺชติ เหตุปจฺจยา. ตีณิ.}

\prose{อุปาทาน\-วิปฺปยุตฺตํ อกุสลํ ธมฺมํ ปฏิจฺจ อุปาทาน\-วิปฺปยุตฺโต อกุสโล ธมฺโม อุปฺปชฺชติ เหตุปจฺจยา. เทฺว.}

\prose{อุปาทาน\-สมฺปยุตฺตํ อกุสลญฺจ อุปาทาน\-วิปฺปยุตฺตํ อกุสลญฺจ ธมฺมํ ปฏิจฺจ อุปาทาน\-สมฺปยุตฺโต อกุสโล ธมฺโม อุปฺปชฺชติ เหตุปจฺจยา. (สํขิตฺตํ.)}

\proseitem{17}{เหตุยา ฉ, อารมฺมเณ ฉ, (สพฺพตฺถ ฉ,) อวิคเต ฉ. (สํขิตฺตํ.)}

\prose{นเหตุยา เอกํ, นอธิปติยา ฉ ฯเปฯ นกมฺเม จตฺตาริ, นวิปาเก ฉ. นวิปฺปยุตฺเต ฉ. (สํขิตฺตํ. สหชาตวาราทิ วิตฺถาเรตพฺโพ.)}

\proseitem{18}{\csromanfolio{420}อุปาทาน\-สมฺปยุตฺโต อกุสโล ธมฺโม อุปาทาน\-สมฺปยุตฺตสฺส อกุสลสฺส ธมฺมสฺส เหตุปจฺจเยน ปจฺจโย. ตีณิ.}

\prose{อุปาทาน\-วิปฺปยุตฺโต อกุสโล ธมฺโม อุปาทาน\-วิปฺปยุตฺตสฺส อกุสลสฺส ธมฺมสฺส เหตุปจฺจเยน ปจฺจโย. เทฺว.}

\prose{อุปาทาน\-สมฺปยุตฺโต อกุสโล จ อุปาทาน\-วิปฺปยุตฺโต อกุสโล จ ธมฺมา อุปาทาน\-สมฺปยุตฺตสฺส อกุสลสฺส ธมฺมสฺส เหตุปจฺจเยน ปจฺจโย. (สํขิตฺตํ.)}

\proseitem{19}{เหตุยา ฉ, อารมฺมเณ นว, อธิปติยา นว. (อุปาทาน\-สมฺปยุตฺต\-มูลเก ตีณิ, สหชาตา\-ธิปติ, อุปาทาน\-วิปฺปยุตฺเต เอกํ.) อนนฺตเร นว, สมนนฺตเร นว, สหชาเต ฉ, อญฺญมญฺเญ ฉ, นิสฺสเย ฉ, อุปนิสฺสเย นว, อาเสวเน นว, กมฺเม จตฺตาริ, อาหาเร จตฺตาริ, อินฺทฺริเย จตฺตาริ, ฌาเน จตฺตาริ, มคฺเค จตฺตาริ, สมฺปยุตฺเต ฉ, อตฺถิยา ฉ, นตฺถิยา นว, วิคเต นว, อวิคเต ฉ. (สํขิตฺตํ.)}

\proseitem{20}{อุปาทาน\-วิปฺปยุตฺตํ อพฺยากตํ ธมฺมํ ปฏิจฺจ อุปาทาน\-วิปฺปยุตฺโต อพฺยากโต ธมฺโม อุปฺปชฺชติ เหตุปจฺจยา. (สํขิตฺตํ.)}

\prose{เหตุยา เอกํ, อารมฺมเณ เอกํ ฯเปฯ อวิคเต เอกํ. (สํขิตฺตํ.)}

\vspace{\tipitakaparskip}
\csromancenter{(สหชาตวาเรปิ ฯเปฯ ปญฺหาวาเรปิ สพฺพตฺถ เอกํ.)}
\prose{72. อุปาทานอุปาทานิยทุก 1. กุสลตฺติก}

\prose{1. ปฏิจฺจวาร ปจฺจยจตุกฺกเหตุ}

\proseitem{21}{อุปาทานิยํ เจว โน จ อุปาทานํ กุสลํ ธมฺมํ ปฏิจฺจ อุปาทานิโย เจว โน จ อุปาทาโน กุสโล ธมฺโม อุปฺปชฺชติ เหตุปจฺจยา. (สํขิตฺตํ.)}

\prose{เหตุยา เอกํ, อารมฺมเณ เอกํ ฯเปฯ อวิคเต เอกํ. (สํขิตฺตํ.) (สหชาตวาเรปิ ฯเปฯ ปญฺหาวาเรปิ สพฺพตฺถ เอกํ.)}

\proseitem{22}{อุปาทานญฺเจว อุปาทานิยญฺจ อกุสลํ ธมฺมํ ปฏิจฺจ อุปาทาโน เจว อุปาทานิโย จ อกุสโล ธมฺโม อุปฺปชฺชติ เหตุปจฺจยา. ตีณิ.}

\prose{\csromanfolio{421}73. อุปาทานอุปาทานสมฺปยุตฺตทุก 1. กุสลตฺติก}

\prose{อุปาทานิยญฺเจว โน จ อุปาทานํ อกุสลํ ธมฺมํ ปฏิจฺจ อุปาทานิโย เจว โน จ อุปาทาโน อกุสโล ธมฺโม อุปฺปชฺชติ เหตุปจฺจยา. ตีณิ.}

\prose{อุปาทานญฺเจว อุปาทานิยํ อกุสลญฺจ อุปาทานิยญฺเจว โน จ อุปาทานํ อกุสลญฺจ ธมฺมํ ปฏิจฺจ อุปาทาโน เจว อุปาทานิโย จ อกุสโล ธมฺโม อุปฺปชฺชติ เหตุปจฺจยา. ตีณิ. (สํขิตฺตํ.)}

\proseitem{23}{เหตุยา นว, อารมฺมเณ นว ฯเปฯ อวิคเต นว. (สํขิตฺตํ. อิมํ คมนํ อุปาทาน\-ทุก\-อกุสล\-สทิสํ. สหชาตวารมฺปิ ฯเปฯ ปญฺหาวารมฺปิ วิตฺถาเรตพฺพํ.)}

\proseitem{24}{อุปาทานิยญฺเจว โน จ อุปาทานํ อพฺยากตํ ธมฺมํ ปฏิจฺจ อุปาทานิโย เจว โน จ อุปาทาโน อพฺยากโต ธมฺโม อุปฺปชฺชติ เหตุปจฺจยา. (สํขิตฺตํ.)}

\proseitem{25}{เหตุยา เอกํ, อารมฺมเณ เอกํ ฯเปฯ อวิคเต เอกํ. (สํขิตฺตํ.)}

\vspace{\tipitakaparskip}
\csromancenter{(สหชาตวาเรปิ ฯเปฯ ปญฺหาวาเรปิ สพฺพตฺถ เอกํ.)}
\prose{73. อุปาทานอุปาทานสมฺปยุตฺตทุก 1. กุสลตฺติก}

\prose{1. ปฏิจฺจวาร ปจฺจยจตุกฺกเหตุ}

\proseitem{26}{อุปาทานญฺเจว อุปาทาน\-สมฺปยุตฺตญฺจ อกุสลํ ธมฺมํ ปฏิจฺจ อุปาทาโน เจว อุปาทาน\-สมฺปยุตฺโต จ อกุสโล ธมฺโม อุปฺปชฺชติ เหตุปจฺจยา. ตีณิ.}

\prose{อุปาทาน\-สมฺปยุตฺต\-ญฺเจว โน จ อุปาทานํ อกุสลํ ธมฺมํ ปฏิจฺจ อุปาทาน\-สมฺปยุตฺโต เจว โน จ อุปาทาโน อกุสโล ธมฺโม อุปฺปชฺชติ เหตุปจฺจยา. ตีณิ.}

\prose{อุปาทานญฺเจว อุปาทาน\-สมฺปยุตฺตํ อกุสลญฺจ อุปาทาน\-สมฺปยุตฺต\-ญฺเจว โน จ อุปาทานํ อกุสลญฺจ ธมฺมํ ปฏิจฺจ อุปาทาโน เจว อุปาทาน\-สมฺปยุตฺโต จ อกุสโล ธมฺโม อุปฺปชฺชติ เหตุปจฺจยา. ตีณิ. (สํขิตฺตํ.)}

\proseitem{27}{\csromanfolio{422}เหตุยา นว, อารมฺมเณ นว, อธิปติยา นว ฯเปฯ อวิคเต นว. (สํขิตฺตํ. อิมํ คมนํ อุปาทาน\-ทุก\-อกุสล\-สทิสํ. อิธ ปน เหตุปจฺจยา อิทํ นานํ. สหชาตวาเรปิ ฯเปฯ ปญฺหาวาเรปิ สพฺพตฺถ นว.)}

\prose{74. อุปาทานวิปฺปยุตฺตอุปาทานิยทุก 1. กุสลตฺติก}

\prose{1. ปฏิจฺจวาร ปจฺจยจตุกฺกเหตุ}

\proseitem{28}{อุปาทาน\-วิปฺปยุตฺตํ อุปาทานิยํ กุสลํ ธมฺมํ ปฏิจฺจ อุปาทาน\-วิปฺปยุตฺโต อุปาทานิโย กุสโล ธมฺโม อุปฺปชฺชติ เหตุปจฺจยา. (1)}

\prose{อุปาทาน\-วิปฺปยุตฺตํ อนุปาทานิยํ กุสลํ ธมฺมํ ปฏิจฺจ อุปาทาน\-วิปฺปยุตฺโต อนุปาทานิโย กุสโล ธมฺโม อุปฺปชฺชติ เหตุปจฺจยา. (1) (สํขิตฺตํ.)}

\proseitem{29}{เหตุยา เทฺว, อารมฺมเณ เทฺว ฯเปฯ อวิคเต เทฺว. (สํขิตฺตํ. โลกิย\-ทุก\-กุสล\-สทิสํ. สหชาตวารมฺปิ ฯเปฯ ปญฺหาวารมฺปิ วิตฺถาเรตพฺพํ.)}

\proseitem{30}{อุปาทาน\-วิปฺปยุตฺตํ อุปาทานิยํ อกุสลํ ธมฺมํ ปฏิจฺจ อุปาทาน\-วิปฺปยุตฺโต อุปาทานิโย อกุสโล ธมฺโม อุปฺปชฺชติ เหตุปจฺจยา. (สพฺพตฺถ เอกํ.)}

\prose{อุปาทาน\-วิปฺปยุตฺตํ อุปาทานิยํ อพฺยากตํ ธมฺมํ ปฏิจฺจ อุปาทาน\-วิปฺปยุตฺโต อุปาทานิโย อพฺยากโต ธมฺโม อุปฺปชฺชติ เหตุปจฺจยา. (1)}

\prose{อุปาทาน\-วิปฺปยุตฺตํ อนุปาทานิยํ อพฺยากตํ ธมฺมํ ปฏิจฺจ อุปาทาน\-วิปฺปยุตฺโต อนุปาทานิโย อพฺยากโต ธมฺโม อุปฺปชฺชติ เหตุปจฺจยา. ตีณิ.}

\prose{อุปาทาน\-วิปฺปยุตฺตํ อุปาทานิยํ อพฺยากตญฺจ อุปาทาน\-วิปฺปยุตฺตํ อนุปาทานิยํ อพฺยากตญฺจ ธมฺมํ ปฏิจฺจ อุปาทาน\-วิปฺปยุตฺโต อุปาทานิโย อพฺยากโต ธมฺโม อุปฺปชฺชติ เหตุปจฺจยา. (1) (สํขิตฺตํ.)}

\proseitemwithcloser{31}{เหตุยา ปญฺจ, อารมฺมเณ เทฺว ฯเปฯ อวิคเต ปญฺจ. (สํขิตฺตํ. โลกิย\-ทุก\-อพฺยากต\-สทิสํ. สหชาตวารมฺปิ ฯเปฯ ปญฺหาวารมฺปิ สพฺพตฺถ วิตฺถาเรตพฺพํ.)}{อุปาทาน\-โคจฺฉก\-กุสล\-ตฺติกํ นิฏฺฐิตํ.}
\csromanmatikaanchor{mtk.81}
\csromanmatikaanchor{mtk.82}
\csromantocmarknopage{subparagraph}{75. กิเลสทุก}{mtk.81}
\bookmark[dest={mtk.81},level=5]{75. กิเลสทุก}
\csromantocmarkref{subparagraph}{1. กุสลตฺติก ...}{mtk.82}
\bookmark[dest={mtk.82},level=5]{1. กุสลตฺติก ...}
\csromanmarkh{75. กิเลสทุก}\csromanmarkhii{1. กุสลตฺติก 75. กิเลสทุก 1. กุสลตฺติก}\csromanheadernomark{6}{75. กิเลสทุก 1. กุสลตฺติก \textbf{75. กิเลสทุก 1. กุสลตฺติก}}
\csromanheader{2}{1-7. ปฏิจฺจาทิวาร ปจฺจยจตุกฺกเหตุ}
\proseitem{1}{\csromanfolio{423}โนกิเลสํ กุสลํ ธมฺมํ ปฏิจฺจ โนกิเลโส กุสโล ธมฺโม อุปฺปชฺชติ เหตุปจฺจยา. (สํขิตฺตํ.)}

\prose{เหตุยา เอกํ, อารมฺมเณ เอกํ ฯเปฯ อวิคเต เอกํ. (สํขิตฺตํ.)}

\prose{(สหชาตวาเรปิ ฯเปฯ ปญฺหาวาเรปิ สพฺพตฺถ เอกํ.)}

\proseitem{2}{กิเลสํ อกุสลํ ธมฺมํ ปฏิจฺจ กิเลโส อกุสโล ธมฺโม อุปฺปชฺชติ เหตุปจฺจยา. ตีณิ.}

\prose{โนกิเลสํ อกุสลํ ธมฺมํ ปฏิจฺจ โนกิเลโส อกุสโล ธมฺโม อุปฺปชฺชติ เหตุปจฺจยา. ตีณิ.}

\prose{กิเลสํ อกุสลญฺจ โนกิเลสํ อกุสลญฺจ ธมฺมํ ปฏิจฺจ กิเลโส อกุสโล ธมฺโม อุปฺปชฺชติ เหตุปจฺจยา. ตีณิ. (สํขิตฺตํ.)}

\proseitem{3}{เหตุยา นว, อารมฺมเณ นว, (สพฺพตฺถ นว,) อวิคเต นว. (สํขิตฺตํ.)}

\proseitem{4}{กิเลสํ อกุสลํ ธมฺมํ ปฏิจฺจ กิเลโส อกุสโล ธมฺโม อุปฺปชฺชติ นเหตุปจฺจยา. วิจิกิจฺฉํ ปฏิจฺจ วิจิกิจฺฉา\-สหคโต โมโห. อุทฺธจฺจํ ปฏิจฺจ อุทฺธจฺจ\-สหคโต โมโห. (1)}

\prose{โนกิเลสํ อกุสลํ ธมฺมํ ปฏิจฺจ กิเลโส อกุสโล ธมฺโม อุปฺปชฺชติ นเหตุปจฺจยา. วิจิกิจฺฉา\-สหคเต อุทฺธจฺจ\-สหคเต ขนฺเธ ปฏิจฺจ วิจิกิจฺฉา\-สหคโต อุทฺธจฺจ\-สหคโต โมโห. (1)}

\prose{กิเลสํ อกุสลญฺจ โนกิเลสํ อกุสลญฺจ ธมฺมํ ปฏิจฺจ กิเลโส อกุสโล ธมฺโม อุปฺปชฺชติ นเหตุปจฺจยา. วิจิกิจฺฉา\-สหคเต ขนฺเธ จ วิจิกิจฺฉญฺจ ปฏิจฺจ วิจิกิจฺฉา\-สหคโต โมโห. อุทฺธจฺจ\-สหคเต ขนฺเธ จ อุทฺธจฺจญฺจ ปฏิจฺจ อุทฺธจฺจ\-สหคโต โมโห. (1) (สํขิตฺตํ.)}

\proseitem{5}{นเหตุยา ตีณิ, นอธิปติยา นว, นปุเรชาเต นว, นปจฺฉาชาเต นว, นอาเสวเน นว, นกมฺเม ตีณิ, นวิปาเก นว, นวิปฺปยุตฺเต นว. (สํขิตฺตํ. สหชาตวาราทิ วิตฺถาเรตพฺโพ.)}

\proseitem{6}{\csromanfolio{424}กิเลโส อกุสโล ธมฺโม กิเลสสฺส อกุสลสฺส ธมฺมสฺส เหตุปจฺจเยน ปจฺจโย. (สํขิตฺตํ.)}

\proseitem{7}{เหตุยา ตีณิ, อารมฺมเณ นว, อธิปติยา นว, (มชฺเฌ ตีณิ, สหชาตา\-ธิปติ.) อนนฺตเร นว, สมนนฺตเร นว ฯเปฯ อุปนิสฺสเย นว, อาเสวเน นว, กมฺเม ตีณิ, อาหาเร ตีณิ, อินฺทฺริเย ตีณิ, ฌาเน ตีณิ, มคฺเค นว, สมฺปยุตฺเต นว ฯเปฯ อวิคเต นว. (สํขิตฺตํ.)}

\prose{โนกิเลสํ อพฺยากตํ ธมฺมํ ปฏิจฺจ โนกิเลโส อพฺยากโต ธมฺโม อุปฺปชฺชติ เหตุปจฺจยา. (สํขิตฺตํ.)}

\proseitem{8}{เหตุยา เอกํ, อารมฺมเณ เอกํ ฯเปฯ อวิคเต เอกํ. (สํขิตฺตํ.)}

\vspace{\tipitakaparskip}
\csromancenter{(สหชาตวาเรปิ ฯเปฯ ปญฺหาวาเรปิ สพฺพตฺถ เอกํ.)}
\csromanmarkh{76. สํกิเลสิกทุก}\csromanmarkhii{1. กุสลตฺติก}\csromanheadernomark{2}{76. \textbf{สํกิเลสิกทุก 1. กุสลตฺติก}}
\prose{1. ปฏิจฺจวาร ปจฺจยจตุกฺกเหตุ}

\proseitem{9}{สํกิเลสิกํ กุสลํ ธมฺมํ ปฏิจฺจ สํกิเลสิโก กุสโล ธมฺโม อุปฺปชฺชติ เหตุปจฺจยา. (1)}

\prose{อสํกิเลสิกํ กุสลํ ธมฺมํ ปฏิจฺจ อสํกิเลสิโก กุสโล ธมฺโม อุปฺปชฺชติ เหตุปจฺจยา. (1) (สํขิตฺตํ.)}

\proseitem{10}{เหตุยา เทฺว, อารมฺมเณ เทฺว ฯเปฯ อวิคเต เทฺว. (สํขิตฺตํ. โลกิย\-ทุก\-กุสล\-สทิสํ. สหชาตวารมฺปิ ฯเปฯ ปญฺหาวารมฺปิ วิตฺถาเรตพฺพํ.)}

\proseitem{11}{สํกิเลสิกํ อกุสลํ ธมฺมํ ปฏิจฺจ สํกิเลสิโก อกุสโล ธมฺโม อุปฺปชฺชติ เหตุปจฺจยา. (สํขิตฺตํ.)}

\prose{เหตุยา เอกํ, อารมฺมเณ เอกํ ฯเปฯ อวิคเต เอกํ. (สํขิตฺตํ.)}

\vspace{\tipitakaparskip}
\csromancenter{(สหชาตวาเรปิ ฯเปฯ ปญฺหาวาเรปิ สพฺพตฺถ เอกํ.)}
\proseitem{12}{สํกิเลสิกํ อพฺยากตํ ธมฺมํ ปฏิจฺจ สํกิเลสิโก อพฺยากโต ธมฺโม อุปฺปชฺชติ เหตุปจฺจยา. (1)}

\csromanmarkh{78. กิเลส\-สมฺปยุตฺตทุก}\csromanmarkhii{1. กุสลตฺติก}\csromanheadernomark{2}{78. กิเลส\-สมฺปยุตฺตทุก 1. กุสลตฺติก}
\prose{\csromanfolio{425}อสํกิเลสิกํ อพฺยากตํ ธมฺมํ ปฏิจฺจ อสํกิเลสิโก อพฺยากโต ธมฺโม อุปฺปชฺชติ เหตุปจฺจยา. ตีณิ.}

\prose{สํกิเลสิกํ อพฺยากตญฺจ อสํกิเลสิกํ อพฺยากตญฺจ ธมฺมํ ปฏิจฺจ สํกิเลสิโก อพฺยากโต ธมฺโม อุปฺปชฺชติ เหตุปจฺจยา. (สํขิตฺตํ.)}

\proseitem{13}{เหตุยา ปญฺจ, อารมฺมเณ เทฺว ฯเปฯ อวิคเต ปญฺจ. (สํขิตฺตํ. โลกิย\-ทุก\-อพฺยากต\-สทิสํ. สหชาตวาเรปิ ฯเปฯ ปญฺหาวาเรปิ สพฺพตฺถ วิตฺถาเรตพฺพํ.)}

\csromanmarkh{77. สํกิลิฏฺฐทุก}\csromanmarkhii{1. กุสลตฺติก}\csromanheadernomark{2}{77. \textbf{สํกิลิฏฺฐทุก 1. กุสลตฺติก}}
\csromanheader{2}{1. \textbf{ปฏิจฺจวาร ปจฺจยจตุกฺก}}
\proseitem{14}{อสํกิลิฏฺฐํ กุสลํ ธมฺมํ ปฏิจฺจ อสํกิลิฏฺโฐ กุสโล ธมฺโม อุปฺปชฺชติ เหตุปจฺจยา. (สํขิตฺตํ.)}

\prose{เหตุยา เอกํ, อารมฺมเณ เอกํ ฯเปฯ อวิคเต เอกํ. (สํขิตฺตํ.)}

\prose{(สหชาตวาเรปิ ฯเปฯ ปญฺหาวาเรปิ สพฺพตฺถ เอกํ.)}

\proseitem{15}{สํกิลิฏฺฐํ อกุสลํ ธมฺมํ ปฏิจฺจ สํกิลิฏฺโฐ อกุสโล ธมฺโม อุปฺปชฺชติ เหตุปจฺจยา. (สํขิตฺตํ.)}

\prose{เหตุยา เอกํ, อารมฺมเณ เอกํ ฯเปฯ อวิคเต เอกํ. (สํขิตฺตํ.)}

\prose{(สหชาตวาเรปิ ฯเปฯ ปญฺหาวาเรปิ สพฺพตฺถ เอกํ.)}

\proseitem{16}{อสํกิลิฏฺฐํ อพฺยากตํ ธมฺมํ ปฏิจฺจ อสํกิลิฏฺโฐ อพฺยากโต ธมฺโม อุปฺปชฺชติ เหตุปจฺจยา. (สํขิตฺตํ.)}

\prose{เหตุยา เอกํ, อารมฺมเณ เอกํ ฯเปฯ อวิคเต เอกํ. (สํขิตฺตํ.)}

\prose{(สหชาตวาเรปิ ฯเปฯ ปญฺหาวาเรปิ สพฺพตฺถ เอกํ.)}

\csromanmarkh{78. กิเลส\-สมฺปยุตฺตทุก}\csromanmarkhii{1. กุสลตฺติก}\csromanheadernomark{2}{78. กิเลส\-สมฺปยุตฺตทุก 1. กุสลตฺติก}
\csromanheader{2}{1. \textbf{ปฏิจฺจวาร ปจฺจยจตุกฺก}}
\proseitem{17}{กิเลส\-วิปฺปยุตฺตํ กุสลํ ธมฺมํ ปฏิจฺจ กิเลส\-วิปฺปยุตฺโต กุสโล ธมฺโม อุปฺปชฺชติ เหตุปจฺจยา. (สํขิตฺตํ.)}

\prose{\csromanfolio{426}เหตุยา เอกํ, อารมฺมเณ เอกํ ฯเปฯ อวิคเต เอกํ. (สํขิตฺตํ.)}

\prose{(สหชาตวาเรปิ ฯเปฯ ปญฺหาวาเรปิ สพฺพตฺถ เอกํ.)}

\proseitem{18}{กิเลส\-สมฺปยุตฺตํ อกุสลํ ธมฺมํ ปฏิจฺจ กิเลส\-สมฺปยุตฺโต อกุสโล ธมฺโม อุปฺปชฺชติ เหตุปจฺจยา. (สํขิตฺตํ.)}

\prose{เหตุยา เอกํ, อารมฺมเณ เอกํ ฯเปฯ อวิคเต เอกํ. (สํขิตฺตํ.)}

\prose{(สหชาตวาเรปิ ฯเปฯ ปญฺหาวาเรปิ สพฺพตฺถ เอกํ.)}

\proseitem{19}{กิเลส\-วิปฺปยุตฺตํ อพฺยากตํ ธมฺมํ ปฏิจฺจ กิเลส\-วิปฺปยุตฺโต อพฺยากโต ธมฺโม อุปฺปชฺชติ เหตุปจฺจยา. (สํขิตฺตํ.)}

\prose{เหตุยา เอกํ, อารมฺมเณ เอกํ ฯเปฯ อวิคเต เอกํ. (สํขิตฺตํ.)}

\prose{(สหชาตวาเรปิ ฯเปฯ ปญฺหาวาเรปิ สพฺพตฺถ เอกํ.)}

\csromanmarkh{79. สํกิเลส\-สํกิเลสิกทุก}\csromanmarkhii{1. กุสลตฺติก}\csromanheadernomark{2}{79. สํกิเลส\-สํกิเลสิกทุก 1. กุสลตฺติก}
\csromanheader{2}{1. \textbf{ปฏิจฺจวาร ปจฺจยจตุกฺก}}
\proseitem{20}{สํกิเลสิก\-ญฺเจว โน จ กิเลสํ กุสลํ ธมฺมํ ปฏิจฺจ สํกิเลสิโก เจว โน จ กิเลโส กุสโล ธมฺโม อุปฺปชฺชติ เหตุปจฺจยา. (สํขิตฺตํ.)}

\prose{เหตุยา เอกํ, อารมฺมเณ เอกํ ฯเปฯ อวิคเต เอกํ. (สํขิตฺตํ.)}

\prose{(สหชาตวาเรปิ ฯเปฯ ปญฺหาวาเรปิ สพฺพตฺถ เอกํ.)}

\proseitem{21}{กิเลสญฺเจว สํกิเลสิกญฺจ อกุสลํ ธมฺมํ ปฏิจฺจ กิเลโส เจว สํกิเลสิโก จ อกุสโล ธมฺโม อุปฺปชฺชติ เหตุปจฺจยา. ตีณิ.}

\prose{สํกิเลสิก\-ญฺเจว โน จ กิเลสํ อกุสลํ ธมฺมํ ปฏิจฺจ สํกิเลสิโก เจว โน จ กิเลโส อกุสโล ธมฺโม อุปฺปชฺชติ เหตุปจฺจยา. ตีณิ.}

\prose{กิเลสญฺเจว สํกิเลสิกํ อกุสลญฺจ สํกิเลสิก\-ญฺเจว โน จ กิเลสํ อกุสลญฺจ ธมฺมํ ปฏิจฺจ กิเลโส เจว สํกิเลสิโก จ อกุสโล ธมฺโม อุปฺปชฺชติ เหตุปจฺจยา. ตีณิ. (สํขิตฺตํ.)}

\csromanmarkh{81. กิเลส\-กิเลส\-สมฺปยุตฺตทุก}\csromanmarkhii{1. กุสลตฺติก}\csromanheadernomark{2}{81. กิเลส\-กิเลส\-สมฺปยุตฺตทุก 1. กุสลตฺติก}
\proseitem{22}{\csromanfolio{427}เหตุยา นว, อารมฺมเณ นว ฯเปฯ อวิคเต นว. (สํขิตฺตํ. กิเลส\-ทุก\-อกุสล\-สทิสํ. สหชาตวาเรปิ ฯเปฯ ปญฺหาวาเรปิ สพฺพตฺถ วิตฺถาเรตพฺพํ.)}

\proseitem{23}{สํกิเลสิก\-ญฺเจว โน จ กิเลสํ อพฺยากตํ ธมฺมํ ปฏิจฺจ สํกิเลสิโก เจว โน จ กิเลโส อพฺยากโต ธมฺโม อุปฺปชฺชติ เหตุปจฺจยา. (สํขิตฺตํ.)}

\prose{เหตุยา เอกํ, อารมฺมเณ เอกํ ฯเปฯ อวิคเต เอกํ. (สํขิตฺตํ.)}

\vspace{\tipitakaparskip}
\csromancenter{(สหชาตวาเรปิ ฯเปฯ ปญฺหาวาเรปิ สพฺพตฺถ เอกํ.)}
\csromanmarkh{80. กิเลส\-สํกิลิฏฺฐทุก}\csromanmarkhii{1. กุสลตฺติก}\csromanheadernomark{2}{80. กิเลส\-สํกิลิฏฺฐทุก 1. กุสลตฺติก}
\csromanheader{2}{1. \textbf{ปฏิจฺจวาร ปจฺจยจตุกฺก}}
\proseitem{24}{กิเลสญฺเจว สํกิลิฏฺฐญฺจ อกุสลํ ธมฺมํ ปฏิจฺจ กิเลโส เจว สํกิลิฏฺโฐ จ อกุสโล ธมฺโม อุปฺปชฺชติ เหตุปจฺจยา. ตีณิ.}

\prose{สํกิลิฏฺฐ\-ญฺเจว โน จ กิเลสํ อกุสลํ ธมฺมํ ปฏิจฺจ สํกิลิฏฺโฐ เจว โน จ กิเลโส อกุสโล ธมฺโม อุปฺปชฺชติ เหตุปจฺจยา. ตีณิ.}

\prose{กิเลสญฺเจว สํกิลิฏฺฐํ อกุสลญฺจ สํกิลิฏฺฐ\-ญฺเจว โน จ กิเลสํ อกุสลญฺจ ธมฺมํ ปฏิจฺจ กิเลโส เจว สํกิลิฏฺโฐ จ อกุสโล ธมฺโม อุปฺปชฺชติ เหตุปจฺจยา. ตีณิ. (สํขิตฺตํ.)}

\proseitem{25}{เหตุยา นว, อารมฺมเณ นว ฯเปฯ อวิคเต นว. (สํขิตฺตํ. กิเลส\-ทุก\-อกุสล\-สทิสํ. สหชาตวารมฺปิ ฯเปฯ ปญฺหาวารมฺปิ สพฺพตฺถ วิตฺถาเรตพฺพํ.)}

\csromanmarkh{81. กิเลส\-กิเลส\-สมฺปยุตฺตทุก}\csromanmarkhii{1. กุสลตฺติก}\csromanheadernomark{2}{81. กิเลส\-กิเลส\-สมฺปยุตฺตทุก 1. กุสลตฺติก}
\csromanheader{2}{1. \textbf{ปฏิจฺจวาร ปจฺจยจตุกฺก}}
\proseitem{26}{กิเลสญฺเจว กิเลส\-สมฺปยุตฺตญฺจ อกุสลํ ธมฺมํ ปฏิจฺจ กิเลโส เจว กิเลส\-สมฺปยุตฺโต จ อกุสโล ธมฺโม อุปฺปชฺชติ เหตุปจฺจยา. ตีณิ.}

\prose{\csromanfolio{428}กิเลส\-สมฺปยุตฺต\-ญฺเจว โน จ กิเลสํ อกุสลํ ธมฺมํ ปฏิจฺจ กิเลส\-สมฺปยุตฺโต เจว โน จ กิเลโส อกุสโล ธมฺโม อุปฺปชฺชติ เหตุปจฺจยา. ตีณิ.}

\prose{กิเลสญฺเจว กิเลส\-สมฺปยุตฺตํ อกุสลญฺจ กิเลส\-สมฺปยุตฺต\-ญฺเจว โน จ กิเลสํ อกุสลญฺจ ธมฺมํ ปฏิจฺจ กิเลโส เจว กิเลส\-สมฺปยุตฺโต จ อกุสโล ธมฺโม อุปฺปชฺชติ เหตุปจฺจยา. ตีณิ. (สํขิตฺตํ.)}

\proseitem{27}{เหตุยา นว, อารมฺมเณ นว ฯเปฯ อวิคเต นว. (สํขิตฺตํ. กิเลส\-ทุก\-อกุสล\-สทิสํ. สหชาตวาเรปิ ฯเปฯ ปญฺหาวาเรปิ สพฺพตฺถ วิตฺถาเรตพฺพํ.)}

\csromanmarkh{82. กิเลส\-วิปฺปยุตฺต\-สํกิเลสิกทุก}\csromanmarkhii{1. กุสลตฺติก}\csromanheadernomark{2}{82. กิเลส\-วิปฺปยุตฺต\-สํกิเลสิกทุก 1. กุสลตฺติก}
\csromanheader{2}{1. \textbf{ปฏิจฺจวาร ปจฺจยจตุกฺก}}
\proseitem{28}{กิเลส\-วิปฺปยุตฺตํ สํกิเลสิกํ กุสลํ ธมฺมํ ปฏิจฺจ กิเลส\-วิปฺปยุตฺโต สํกิเลสิโก กุสโล ธมฺโม อุปฺปชฺชติ เหตุปจฺจยา. (1)}

\prose{กิเลส\-วิปฺปยุตฺตํ อสํกิเลสิกํ กุสลํ ธมฺมํ ปฏิจฺจ กิเลส\-วิปฺปยุตฺโต อสํกิเลสิโก กุสโล ธมฺโม อุปฺปชฺชติ เหตุปจฺจยา. (สํขิตฺตํ.)}

\proseitem{29}{เหตุยา เทฺว, อารมฺมเณ เทฺว ฯเปฯ อวิคเต เทฺว. (สํขิตฺตํ. โลกิย\-ทุก\-กุสล\-สทิสํ. สหชาตวาเรปิ ฯเปฯ ปญฺหาวาเรปิ วิตฺถาเรตพฺพํ.)}

\proseitem{30}{กิเลส\-วิปฺปยุตฺตํ สํกิเลสิกํ อพฺยากตํ ธมฺมํ ปฏิจฺจ กิเลสวิปยุตฺโต สํกิเลสิโก อพฺยากโต ธมฺโม อุปฺปชฺชติ เหตุปจฺจยา. (1)}

\prose{กิเลส\-วิปฺปยุตฺตํ อสํกิเลสิกํ อพฺยากตํ ธมฺมํ ปฏิจฺจ กิเลส\-วิปฺปยุตฺโต อสํกิเลสิโก อพฺยากโต ธมฺโม อุปฺปชฺชติ เหตุปจฺจยา. ตีณิ.}

\csromanmatikaanchor{mtk.83}
\csromantocmarkref{subparagraph}{83. ทสฺสเนนปหาตพฺพทุก 1. กุสลตฺติก ...}{mtk.83}
\bookmark[dest={mtk.83},level=5]{83. ทสฺสเนนปหาตพฺพทุก 1. กุสลตฺติก ...}
\csromanmarkh{83. ทสฺสเนน\-ปหาตพฺพทุก}\csromanmarkhii{1. กุสลตฺติก}\csromanheadernomark{6}{83. ทสฺสเนน\-ปหาตพฺพทุก 1. กุสลตฺติก}
\prose{\csromanfolio{429}กิเลส\-วิปฺปยุตฺตํ สํกิเลสิกํ อพฺยากตญฺจ กิเลส\-วิปฺปยุตฺตํ อสํกิเลสิกํ อพฺยากตญฺจ ธมฺมํ ปฏิจฺจ กิเลส\-วิปฺปยุตฺโต สํกิเลสิโก อพฺยากโต ธมฺโม อุปฺปชฺชติ เหตุปจฺจยา. (1) (สํขิตฺตํ.)}

\proseitemwithcloser{31}{เหตุยา ปญฺจ, อารมฺมเณ เทฺว ฯเปฯ อาเสวเน เอกํ ฯเปฯ อวิคเต ปญฺจ. (สํขิตฺตํ. โลกิย\-ทุก\-อพฺยากต\-สทิสํ. สหชาตวาเรปิ ฯเปฯ ปญฺหาวาเรปิ สพฺพตฺถ วิตฺถาเรตพฺพํ.)}{กิเลส\-โคจฺฉก\-กุสล\-ตฺติกํ นิฏฺฐิตํ.}
\csromanmarkh{83. ทสฺสเนน\-ปหาตพฺพทุก}\csromanmarkhii{1. กุสลตฺติก}\csromanheadernomark{2}{83. ทสฺสเนน\-ปหาตพฺพทุก 1. กุสลตฺติก}
\csromanheader{2}{1-7. ปฏิจฺจาทิวาร ปจฺจยจตุกฺก}
\proseitem{1}{นทสฺสเนน ปหาตพฺพํ กุสลํ ธมฺมํ ปฏิจฺจ นทสฺสเนน ปหาตพฺโพ กุสโล ธมฺโม อุปฺปชฺชติ เหตุปจฺจยา. (สํขิตฺตํ.)}

\prose{เหตุยา เอกํ, อารมฺมเณ เอกํ, อธิปติยา เอกํ ฯเปฯ อวิคเต เอกํ. (สํขิตฺตํ. สหชาตวาเรปิ ฯเปฯ ปญฺหาวาเรปิ สพฺพตฺถ เอกํ.)}

\proseitem{2}{ทสฺสเนน ปหาตพฺพํ อกุสลํ ธมฺมํ ปฏิจฺจ ทสฺสเนน ปหาตพฺโพ อกุสโล ธมฺโม อุปฺปชฺชติ เหตุปจฺจยา. (สํขิตฺตํ.)}

\prose{เหตุยา เทฺว, อารมฺมเณ เทฺว, อธิปติยา เทฺว ฯเปฯ อวิคเต เทฺว. (สํขิตฺตํ.)}

\proseitem{3}{ทสฺสเนน ปหาตพฺพํ อกุสลํ ธมฺมํ ปฏิจฺจ ทสฺสเนน ปหาตพฺโพ อกุสโล ธมฺโม อุปฺปชฺชติ นเหตุปจฺจยา. วิจิกิจฺฉา\-สหคเต ขนฺเธ ปฏิจฺจ วิจิกิจฺฉา\-สหคโต โมโห. (1)}

\prose{นทสฺสเนน ปหาตพฺพํ อกุสลํ ธมฺมํ ปฏิจฺจ นทสฺสเนน ปหาตพฺโพ อกุสโล ธมฺโม อุปฺปชฺชติ นเหตุปจฺจยา. (1) (สํขิตฺตํ.)}

\proseitem{4}{นเหตุยา เทฺว, นอธิปติยา เทฺว, นปุเรชาเต เทฺว ฯเปฯ นกมฺเม เทฺว, นวิปาเก เทฺว, นวิปฺปยุตฺเต เทฺว. (สํขิตฺตํ. สหชาตวาเรปิ ฯเปฯ สมฺปยุตฺตวาเรปิ สพฺพตฺถ วิตฺถาเรตพฺพํ.)}

\proseitem{5}{\csromanfolio{430}ทสฺสเนน ปหาตพฺโพ อกุสโล ธมฺโม ทสฺสเนน ปหาตพฺพสฺส อกุสลสฺส ธมฺมสฺส เหตุปจฺจเยน ปจฺจโย. (1)}

\prose{นทสฺสเนน ปหาตพฺโพ อกุสโล ธมฺโม นทสฺสเนน ปหาตพฺพสฺส อกุสลสฺส ธมฺมสฺส เหตุปจฺจเยน ปจฺจโย. (1) (สํขิตฺตํ.)}

\proseitem{6}{เหตุยา เทฺว, อารมฺมเณ ตีณิ. (ทสฺสเน เอกํ, นทสฺสเน เทฺว,) อธิปติยา ตีณิ. (ทสฺสเนน ปหาตพฺพ\-มูลกํ เอกํ, นทสฺสเน เทฺว, อารมฺมณา\-ธิปติ, สหชาตา\-ธิปติ, เอการมฺมณา\-ธิปติ.) อนนฺตเร เทฺว, (ทสฺสนมูลกํ เอกํ, นทสฺสเน เอกํ.) สมนนฺตเร เทฺว, สหชาเต เทฺว ฯเปฯ อุปนิสฺสเย ตีณิ, อาเสวเน เทฺว, กมฺเม เทฺว, อาหาเร เทฺว ฯเปฯ สมฺปยุตฺเต เทฺว, อตฺถิยา เทฺว, นตฺถิยา เทฺว ฯเปฯ อวิคเต เทฺว. (สํขิตฺตํ.)}

\proseitem{7}{นทสฺสเนน ปหาตพฺพํ อพฺยากตํ ธมฺมํ ปฏิจฺจ นทสฺสเนน ปหาตพฺโพ อพฺยากโต ธมฺโม อุปฺปชฺชติ เหตุปจฺจยา. (สํขิตฺตํ.)}

\prose{เหตุยา เอกํ, อารมฺมเณ เอกํ ฯเปฯ อวิคเต เอกํ. (สํขิตฺตํ.)}

\vspace{\tipitakaparskip}
\csromancenter{(สหชาตวาเรปิ ฯเปฯ ปญฺหาวาเรปิ สพฺพตฺถ เอกํ.)}
\csromanmarkh{84. ภาวนาย\-ปหาตพฺพทุก}\csromanmarkhii{1. กุสลตฺติก}\csromanheadernomark{2}{84. ภาวนาย\-ปหาตพฺพทุก 1. กุสลตฺติก}
\csromanheader{2}{1-7. ปฏิจฺจาทิวาร ปจฺจยจตุกฺก}
\proseitem{8}{นภาวนาย ปหาตพฺพํ กุสลํ ธมฺมํ ปฏิจฺจ นภาวนาย ปหาตพฺโพ กุสโล ธมฺโม อุปฺปชฺชติ เหตุปจฺจยา. (สํขิตฺตํ.)}

\prose{เหตุยา เอกํ, อารมฺมเณ เอกํ ฯเปฯ อวิคเต เอกํ. (สํขิตฺตํ.)}

\vspace{\tipitakaparskip}
\csromancenter{(สหชาตวาเรปิ ฯเปฯ ปญฺหาวาเรปิ สพฺพตฺถ เอกํ.)}
\proseitem{9}{ภาวนาย ปหาตพฺพํ อกุสลํ ธมฺมํ ปฏิจฺจ ภาวนาย ปหาตพฺโพ อกุสโล ธมฺโม อุปฺปชฺชติ เหตุปจฺจยา. (สํขิตฺตํ.)}

\prose{เหตุยา เทฺว, อารมฺมเณ เทฺว ฯเปฯ อวิคเต เทฺว. (สํขิตฺตํ. ทสฺสเนน ปหาตพฺพ\-ทุก\-อกุสล\-สทิสํ. สหชาตวาเรปิ ฯเปฯ ปญฺหาวาเรปิ สพฺพตฺถ วิตฺถาเรตพฺพํ.)}

\csromanmarkh{85. ทสฺสเนน\-ปหาตพฺพ\-เหตุกทุก}\csromanmarkhii{1. กุสลตฺติก}\csromanheadernomark{2}{85. ทสฺสเนน\-ปหาตพฺพ\-เหตุกทุก 1. กุสลตฺติก}
\proseitem{10}{\csromanfolio{431}นภาวนาย ปหาตพฺพํ อพฺยากตํ ธมฺมํ ปฏิจฺจ นภาวนาย ปหาตพฺโพ อพฺยากโต ธมฺโม อุปฺปชฺชติ เหตุปจฺจยา. (สํขิตฺตํ.)}

\prose{เหตุยา เอกํ, อารมฺมเณ เอกํ ฯเปฯ อวิคเต เอกํ. (สํขิตฺตํ.)}

\vspace{\tipitakaparskip}
\csromancenter{(สหชาตวาเรปิ ฯเปฯ ปญฺหาวาเรปิ สพฺพตฺถ เอกํ.)}
\csromanmarkh{85. ทสฺสเนน\-ปหาตพฺพ\-เหตุกทุก}\csromanmarkhii{1. กุสลตฺติก}\csromanheadernomark{2}{85. ทสฺสเนน\-ปหาตพฺพ\-เหตุกทุก 1. กุสลตฺติก}
\csromanheader{2}{1-7. ปฏิจฺจาทิวาร ปจฺจยจตุกฺก}
\proseitem{11}{นทสฺสเนน ปหาตพฺพ\-เหตุกํ กุสลํ ธมฺมํ ปฏิจฺจ นทสฺสเนน ปหาตพฺพ\-เหตุโก กุสโล ธมฺโม อุปฺปชฺชติ เหตุปจฺจยา. (สํขิตฺตํ.)}

\prose{เหตุยา เอกํ, อารมฺมเณ เอกํ ฯเปฯ อวิคเต เอกํ. (สํขิตฺตํ.)}

\vspace{\tipitakaparskip}
\csromancenter{(สหชาตวาเรปิ ฯเปฯ ปญฺหาวาเรปิ สพฺพตฺถ เอกํ.)}
\proseitem{12}{ทสฺสเนน ปหาตพฺพ\-เหตุกํ อกุสลํ ธมฺมํ ปฏิจฺจ ทสฺสเนน ปหาตพฺพ\-เหตุโก อกุสโล ธมฺโม อุปฺปชฺชติ เหตุปจฺจยา. (1) นทสฺสเนน ปหาตพฺพ\-เหตุกํ อกุสลํ ธมฺมํ ปฏิจฺจ นทสฺสเนน ปหาตพฺพ\-เหตุโก อกุสโล ธมฺโม อุปฺปชฺชติ เหตุปจฺจยา.  นทสฺสเนน ปหาตพฺพ\-เหตุกํ อกุสลํ ธมฺมํ ปฏิจฺจ ทสฺสเนน ปหาตพฺพ\-เหตุโก อกุสโล ธมฺโม อุปฺปชฺชติ เหตุปจฺจยา. (2)}

\prose{ทสฺสเนน ปหาตพฺพ\-เหตุกํ อกุสลญฺจ นทสฺสเนน ปหาตพฺพ\-เหตุกํ อกุสลญฺจ ธมฺมํ ปฏิจฺจ ทสฺสเนน ปหาตพฺพ\-เหตุโก อกุสโล ธมฺโม อุปฺปชฺชติ เหตุปจฺจยา. (1)}

\prose{ทสฺสเนน ปหาตพฺพ\-เหตุกํ อกุสลํ ธมฺมํ ปฏิจฺจ ทสฺสเนน ปหาตพฺพ\-เหตุโก อกุสโล ธมฺโม อุปฺปชฺชติ อารมฺมณ\-ปจฺจยา. ตีณิ. (นทสฺสเน เทฺว, ฆฏเน เอกํ.) (สํขิตฺตํ.)}

\proseitem{13}{เหตุยา จตฺตาริ, อารมฺมเณ ฉ, อธิปติยา เทฺว, (สพฺพตฺถ ฉ,) อวิคเต ฉ. (สํขิตฺตํ.)}

\proseitem{14}{ทสฺสเนน ปหาตพฺพ\-เหตุกํ อกุสลํ ธมฺมํ ปฏิจฺจ นทสฺสเนน ปหาตพฺพ\-เหตุโก อกุสโล ธมฺโม อุปฺปชฺชติ นเหตุปจฺจยา.}

\prose{\csromanfolio{432}นทสฺสเนน ปหาตพฺพ\-เหตุกํ อกุสลํ ธมฺมํ ปฏิจฺจ นทสฺสเนน ปหาตพฺพ\-เหตุโก อกุสโล ธมฺโม อุปฺปชฺชติ นเหตุปจฺจยา. (สํขิตฺตํ.)}

\proseitem{15}{นเหตุยา เทฺว, นอธิปติยา ฉ ฯเปฯ นกมฺเม จตฺตาริ, นวิปาเก ฉ, นวิปฺปยุตฺเต ฉ. (สํขิตฺตํ. สหชาตวาราทิ วิตฺถาเรตพฺโพ.)}

\proseitem{16}{ทสฺสเนน ปหาตพฺพ\-เหตุโก อกุสโล ธมฺโม ทสฺสเนน ปหาตพฺพ\-เหตุกสฺส อกุสลสฺส ธมฺมสฺส เหตุปจฺจเยน ปจฺจโย. (สํขิตฺตํ.)}

\proseitem{17}{เหตุยา ตีณิ, อารมฺมเณ นว, อธิปติยา ตีณิ, (ทสฺสเน เอกํ, นทสฺสเน เทฺว.) อนนฺตเร นว, สมนนฺตเร นว, สหชาเต ฉ ฯเปฯ อุปนิสฺสเย นว, อาเสวเน นว, กมฺเม จตฺตาริ, อาหาเร จตฺตาริ, อินฺทฺริเย จตฺตาริ ฯเปฯ สมฺปยุตฺเต ฉ, อตฺถิยา ฉ, นตฺถิยา นว ฯเปฯ อวิคเต ฉ. (สํขิตฺตํ.)}

\proseitem{18}{นทสฺสเนน ปหาตพฺพ\-เหตุกํ อพฺยากตํ ธมฺมํ ปฏิจฺจ นทสฺสเนน ปหาตพฺพ\-เหตุโก อพฺยากโต ธมฺโม อุปฺปชฺชติ เหตุปจฺจยา. (สํขิตฺตํ.)}

\prose{เหตุยา เอกํ, อารมฺมเณ เอกํ ฯเปฯ อวิคเต เอกํ. (สํขิตฺตํ.)}

\vspace{\tipitakaparskip}
\csromancenter{(สหชาตวาเรปิ ฯเปฯ ปญฺหาวาเรปิ สพฺพตฺถ เอกํ.)}
\csromanmarkh{86. ภาวนาย\-ปหาตพฺพ\-เหตุกทุก}\csromanmarkhii{1. กุสลตฺติก}\csromanheadernomark{2}{86. ภาวนาย\-ปหาตพฺพ\-เหตุกทุก 1. กุสลตฺติก}
\csromanheader{2}{1. \textbf{ปฏิจฺจวาร ปจฺจยจตุกฺก}}
\proseitem{19}{นภาวนาย ปหาตพฺพ\-เหตุกํ กุสลํ ธมฺมํ ปฏิจฺจ นภาวนาย ปหาตพฺพ\-เหตุโก กุสโล ธมฺโม อุปฺปชฺชติ เหตุปจฺจยา. (สํขิตฺตํ.)}

\prose{เหตุยา เอกํ, อารมฺมเณ เอกํ ฯเปฯ อวิคเต เอกํ. (สํขิตฺตํ.)}

\vspace{\tipitakaparskip}
\csromancenter{(สหชาตวาเรปิ ฯเปฯ ปญฺหาวาเรปิ สพฺพตฺถ เอกํ.)}
\proseitem{20}{ภาวนาย ปหาตพฺพ\-เหตุกํ อกุสลํ ธมฺมํ ปฏิจฺจ ภาวนาย ปหาตพฺพ\-เหตุโก อกุสโล ธมฺโม อุปฺปชฺชติ เหตุปจฺจยา. (1)}

\csromanmarkh{87. สวิตกฺกทุก}\csromanmarkhii{1. กุสลตฺติก}\csromanheadernomark{2}{87. สวิตกฺกทุก 1. กุสลตฺติก}
\prose{\csromanfolio{433}นภาวนาย ปหาตพฺพ\-เหตุกํ อกุสลํ ธมฺมํ ปฏิจฺจ นภาวนาย ปหาตพฺพ\-เหตุโก อกุสโล ธมฺโม อุปฺปชฺชติ เหตุปจฺจยา. ตีณิ. (สํขิตฺตํ.)}

\proseitem{21}{เหตุยา จตฺตาริ, อารมฺมเณ ฉ, อธิปติยา เทฺว ฯเปฯ อวิคเต ฉ. (สํขิตฺตํ. ทสฺสเนน ปหาตพฺพ\-เหตุก\-ทุก\-อกุสล\-สทิสํ. สหชาตวาเรปิ ฯเปฯ ปญฺหาวาเรปิ สพฺพตฺถ วิตฺถาเรตพฺพํ.)}

\proseitem{22}{นภาวนาย ปหาตพฺพ\-เหตุกํ อพฺยากตํ ธมฺมํ ปฏิจฺจ นภาวนาย ปหาตพฺพ\-เหตุโก อพฺยากโต ธมฺโม อุปฺปชฺชติ เหตุปจฺจยา. (สํขิตฺตํ.)}

\prose{เหตุยา เอกํ, อารมฺมเณ เอกํ ฯเปฯ อวิคเต เอกํ. (สํขิตฺตํ.)}

\vspace{\tipitakaparskip}
\csromancenter{(สหชาตวาเรปิ ฯเปฯ ปญฺหาวาเรปิ สพฺพตฺถ เอกํ.)}
\csromanmarkh{87. สวิตกฺกทุก}\csromanmarkhii{1. กุสลตฺติก}\csromanheadernomark{2}{87. สวิตกฺกทุก 1. กุสลตฺติก}
\csromanheader{2}{1. \textbf{ปฏิจฺจวาร ปจฺจยจตุกฺก}}
\proseitem{23}{สวิตกฺกํ กุสลํ ธมฺมํ ปฏิจฺจ สวิตกฺโก กุสโล ธมฺโม อุปฺปชฺชติ เหตุปจฺจยา.  สวิตกฺกํ กุสลํ ธมฺมํ ปฏิจฺจ อวิตกฺโก กุสโล ธมฺโม อุปฺปชฺชติ เหตุปจฺจยา.  สวิตกฺกํ กุสลํ ธมฺมํ ปฏิจฺจ สวิตกฺโก กุสโล จ อวิตกฺโก กุสโล จ ธมฺมา อุปฺปชฺชนฺติ เหตุปจฺจยา. (3)}

\prose{อวิตกฺกํ กุสลํ ธมฺมํ ปฏิจฺจ อวิตกฺโก กุสโล ธมฺโม อุปฺปชฺชติ เหตุปจฺจยา.  อวิตกฺกํ กุสลํ ธมฺมํ ปฏิจฺจ สวิตกฺโก กุสโล ธมฺโม อุปฺปชฺชติ เหตุปจฺจยา. (2)}

\prose{สวิตกฺกํ กุสลญฺจ อวิตกฺกํ กุสลญฺจ ธมฺมํ ปฏิจฺจ สวิตกฺโก กุสโล ธมฺโม อุปฺปชฺชติ เหตุปจฺจยา. (1) (สํขิตฺตํ.)}

\proseitem{24}{เหตุยา ฉ, อารมฺมเณ ฉ, (สพฺพตฺถ ฉ.) อวิคเต ฉ. (สํขิตฺตํ.)}

\prose{นอธิปติยา ฉ, นปุเรชาเต ฉ ฯเปฯ นกมฺเม จตฺตาริ ฯเปฯ นวิปฺปยุตฺเต ฉ. (สํขิตฺตํ. สหชาตวารมฺปิ ฯเปฯ สมฺปยุตฺตวารมฺปิ วิตฺถาเรตพฺพํ.)}

\csromanheader{2}{ธมฺมานุโลม ทุกติก\-ปฏฺฐาน\-ปาฬิ ปญฺหาวารเหตุ อารมฺมณ}
\proseitem{25}{\csromanfolio{434}สวิตกฺโก กุสโล ธมฺโม สวิตกฺกสฺส กุสลสฺส ธมฺมสฺส เหตุปจฺจเยน ปจฺจโย.  สวิตกฺโก กุสโล ธมฺโม อวิตกฺกสฺส กุสลสฺส ธมฺมสฺส เหตุปจฺจเยน ปจฺจโย.  สวิตกฺโก กุสโล ธมฺโม สวิตกฺกสฺส กุสลสฺส จ อวิตกฺกสฺส กุสลสฺส จ ธมฺมสฺส เหตุปจฺจเยน ปจฺจโย. (3)}

\prose{อวิตกฺโก กุสโล ธมฺโม อวิตกฺกสฺส กุสลสฺส ธมฺมสฺส เหตุปจฺจเยน ปจฺจโย. (1)}

\prose{สวิตกฺโก กุสโล ธมฺโม สวิตกฺกสฺส กุสลสฺส ธมฺมสฺส อารมฺมณ\-ปจฺจเยน ปจฺจโย. ตีณิ.}

\prose{อวิตกฺโก กุสโล ธมฺโม อวิตกฺกสฺส กุสลสฺส ธมฺมสฺส อารมฺมณ\-ปจฺจเยน ปจฺจโย. ตีณิ.}

\prose{สวิตกฺโก กุสโล จ อวิตกฺโก กุสโล จ ธมฺมา สวิตกฺกสฺส กุสลสฺส ธมฺมสฺส อารมฺมณ\-ปจฺจเยน ปจฺจโย. ตีณิ. (สํขิตฺตํ.)}

\proseitem{26}{เหตุยา จตฺตาริ, อารมฺมเณ นว, อธิปติยา นว. (เหฏฺฐา ตีสุ สหชาตา\-ธิปติ, อวิตกฺเก เอกํ สหชาตา\-ธิปติ.) อนนฺตเร นว, สมนนฺตเร นว, สหชาเต ฉ ฯเปฯ อุปนิสฺสเย นว, อาเสวเน นว, กมฺเม จตฺตาริ, อาหาเร จตฺตาริ, อินฺทฺริเย จตฺตาริ, ฌาเน ฉ, มคฺเค ฉ, สมฺปยุตฺเต ฉ, อตฺถิยา ฉ, นตฺถิยา นว, วิคเต นว, อวิคเต ฉ. (สํขิตฺตํ.)}

\prose{นเหตุยา นว, นอารมฺมเณ นว. (สํขิตฺตํ.)}

\prose{เหตุปจฺจยา นอารมฺมเณ จตฺตาริ. (สํขิตฺตํ.)}

\prose{นเหตุปจฺจยา อารมฺมเณ นว. (สํขิตฺตํ.) (ยถา กุสลตฺติเก ปญฺหาวารํ, เอวํ วิตฺถาเรตพฺพํ.)}

\csromanheader{2}{\textbf{อกุสลปท}}
\proseitem{27}{สวิตกฺกํ อกุสลํ ธมฺมํ ปฏิจฺจ สวิตกฺโก อกุสโล ธมฺโม อุปฺปชฺชติ เหตุปจฺจยา.  สวิตกฺกํ อกุสลํ ธมฺมํ ปฏิจฺจ อวิตกฺโก}

\vspace{\tipitakaparskip}
\csromancenter{1. กุสลตฺติก อกุสโล ธมฺโม อุปฺปชฺชติ เหตุปจฺจยา.  สวิตกฺกํ อกุสลํ ธมฺมํ ปฏิจฺจ สวิตกฺโก อกุสโล จ อวิตกฺโก อกุสโล จ ธมฺมา อุปฺปชฺชนฺติ เหตุปจฺจยา. (3)}
\prose{\csromanfolio{435}อวิตกฺกํ อกุสลํ ธมฺมํ ปฏิจฺจ สวิตกฺโก อกุสโล ธมฺโม อุปฺปชฺชติ เหตุปจฺจยา. (1)}

\prose{สวิตกฺกํ อกุสลญฺจ อวิตกฺกํ อกุสลญฺจ ธมฺมํ ปฏิจฺจ สวิตกฺโก อกุสโล ธมฺโม อุปฺปชฺชติ เหตุปจฺจยา. (1) (สํขิตฺตํ.)}

\proseitem{28}{เหตุยา ปญฺจ, อารมฺมเณ ปญฺจ, (สพฺพตฺถ ปญฺจ.) ฯเปฯ อวิคเต ปญฺจ. (สํขิตฺตํ.)}

\proseitem{29}{สวิตกฺกํ อกุสลํ ธมฺมํ ปฏิจฺจ สวิตกฺโก อกุสโล ธมฺโม อุปฺปชฺชติ นเหตุปจฺจยา. (1)}

\prose{อวิตกฺกํ อกุสลํ ธมฺมํ ปฏิจฺจ สวิตกฺโก อกุสโล ธมฺโม อุปฺปชฺชติ นเหตุปจฺจยา. (1)}

\prose{สวิตกฺกํ อกุสลญฺจ อวิตกฺกํ อกุสลญฺจ ธมฺมํ ปฏิจฺจ สวิตกฺโก อกุสโล ธมฺโม อุปฺปชฺชติ นเหตุปจฺจยา. (1) (สํขิตฺตํ.)}

\proseitem{30}{นเหตุยา ตีณิ, นอธิปติยา ปญฺจ, นกมฺเม ตีณิ ฯเปฯ นวิปฺปยุตฺเต ปญฺจ. (สํขิตฺตํ. สหชาตวาราทิ วิตฺถาเรตพฺโพ.)}

\proseitem{31}{สวิตกฺโก อกุสโล ธมฺโม สวิตกฺกสฺส อกุสลสฺส ธมฺมสฺส เหตุปจฺจเยน ปจฺจโย. (สํขิตฺตํ.)}

\proseitem{32}{เหตุยา ตีณิ, อารมฺมเณ นว, อธิปติยา นว, อนนฺตเร นว ฯเปฯ สหชาเต ปญฺจ ฯเปฯ อุปนิสฺสเย นว, อาเสวเน นว, กมฺเม ตีณิ, อาหาเร ตีณิ ฯเปฯ ฌาเน ปญฺจ, มคฺเค ปญฺจ, สมฺปยุตฺเต ปญฺจ, อตฺถิยา ปญฺจ, นตฺถิยา นว ฯเปฯ อวิคเต ปญฺจ. (สํขิตฺตํ.)}

\csromanheader{2}{\textbf{อพฺยากตปท}}
\proseitem{33}{สวิตกฺกํ อพฺยากตํ ธมฺมํ ปฏิจฺจ สวิตกฺโก อพฺยากโต ธมฺโม อุปฺปชฺชติ เหตุปจฺจยา. ตีณิ.}

\prose{\csromanfolio{436}อวิตกฺกํ อพฺยากตํ ธมฺมํ ปฏิจฺจ อวิตกฺโก อพฺยากโต ธมฺโม อุปฺปชฺชติ เหตุปจฺจยา. ตีณิ.}

\prose{สวิตกฺกํ อพฺยากตญฺจ อวิตกฺกํ อพฺยากตญฺจ ธมฺมํ ปฏิจฺจ สวิตกฺโก อพฺยากโต ธมฺโม อุปฺปชฺชติ เหตุปจฺจยา. ตีณิ. (สํขิตฺตํ.)}

\proseitem{34}{เหตุยา นว, อารมฺมเณ นว ฯเปฯ ปุเรชาเต อาเสวเน ฉ ฯเปฯ อวิคเต นว. (สํขิตฺตํ.)}

\prose{นเหตุยา นว, นอารมฺมเณ ตีณิ, นอธิปติยา นว ฯเปฯ นปุเรชาเต นว ฯเปฯ นกมฺเม จตฺตาริ, นวิปาเก นว, นอาหาเร เอกํ, นอินฺทฺริเย เอกํ, นฌาเน เอกํ, นมคฺเค นว, นสมฺปยุตฺเต ตีณิ, นวิปฺปยุตฺเต ฉ ฯเปฯ โนวิคเต ตีณิ. (สํขิตฺตํ. สหชาตวารมฺปิ ฯเปฯ สมฺปยุตฺตวารมฺปิ วิตฺถาเรตพฺพํ.)}

\proseitem{35}{สวิตกฺโก อพฺยากโต ธมฺโม สวิตกฺกสฺส อพฺยากตสฺส ธมฺมสฺส เหตุปจฺจเยน ปจฺจโย. ตีณิ.}

\prose{อวิตกฺโก อพฺยากโต ธมฺโม อวิตกฺกสฺส อพฺยากตสฺส ธมฺมสฺส เหตุปจฺจเยน ปจฺจโย. (1)}

\prose{สวิตกฺโก อพฺยากโต ธมฺโม สวิตกฺกสฺส อพฺยากตสฺส ธมฺมสฺส อารมฺมณ\-ปจฺจเยน ปจฺจโย. (สํขิตฺตํ.)}

\proseitem{36}{เหตุยา จตฺตาริ, อารมฺมเณ นว, อธิปติยา นว, (สพฺพตฺถ นว.) อุปนิสฺสเย นว, ปุเรชาเต ตีณิ, ปจฺฉาชาเต ตีณิ, อาเสวเน นว, กมฺเม จตฺตาริ, วิปาเก นว, อาหาเร จตฺตาริ ฯเปฯ ฌาเน นว, มคฺเค นว, สมฺปยุตฺเต ฉ, วิปฺปยุตฺเต ปญฺจ, อตฺถิยา นว, นตฺถิยา นว, วิคเต นว, อวิคเต นว. (สํขิตฺตํ.)}

\prose{นเหตุยา นว, นอารมฺมเณ นว. (สํขิตฺตํ.)}

\prose{เหตุปจฺจยา นอารมฺมเณ จตฺตาริ. (สํขิตฺตํ.)}

\prose{นเหตุปจฺจยา อารมฺมเณ นว. (สํขิตฺตํ.)}

\prose{(ยถา กุสลตฺติเก ปญฺหาวารํ, เอวํ วิตฺถาเรตพฺพํ.)}

\csromanmarkh{88. สวิจารทุก}\csromanmarkhii{1. กุสลตฺติก 88. สวิจารทุก 1. กุสลตฺติก}\csromanheadernomark{2}{88. สวิจารทุก 1. กุสลตฺติก \textbf{88. สวิจารทุก 1. กุสลตฺติก}}
\csromanheader{2}{1-7. ปฏิจฺจาทิวาร ปจฺจยจตุกฺก}
\proseitem{37}{\csromanfolio{437}สวิจารํ กุสลํ ธมฺมํ ปฏิจฺจ สวิจาโร กุสโล ธมฺโม อุปฺปชฺชติ เหตุปจฺจยา. ตีณิ.}

\prose{อวิจารํ กุสลํ ธมฺมํ ปฏิจฺจ อวิจาโร กุสโล ธมฺโม อุปฺปชฺชติ เหตุปจฺจยา. เทฺว.}

\prose{สวิจารํ กุสลญฺจ อวิจารํ กุสลญฺจ ธมฺมํ ปฏิจฺจ สวิจาโร กุสโล ธมฺโม อุปฺปชฺชติ เหตุปจฺจยา. (1) (สํขิตฺตํ.)}

\proseitem{38}{เหตุยา ฉ, อารมฺมเณ ฉ ฯเปฯ อวิคเต ฉ. (สํขิตฺตํ. สวิตกฺก\-ทุก\-กุสล\-สทิสํ. สหชาตวารมฺปิ ฯเปฯ ปญฺหาวารมฺปิ วิตฺถาเรตพฺพํ.)}

\proseitem{39}{สวิจารํ อกุสลํ ธมฺมํ ปฏิจฺจ สวิจาโร อกุสโล ธมฺโม อุปฺปชฺชติ เหตุปจฺจยา. ตีณิ.}

\prose{อวิจารํ อกุสลํ ธมฺมํ ปฏิจฺจ สวิจาโร อกุสโล ธมฺโม อุปฺปชฺชติ เหตุปจฺจยา. (1)}

\prose{สวิจารํ อกุสลญฺจ อวิจารํ อกุสลญฺจ ธมฺมํ ปฏิจฺจ สวิจาโร อกุสโล ธมฺโม อุปฺปชฺชติ เหตุปจฺจยา. (1) (สํขิตฺตํ.)}

\proseitem{40}{เหตุยา ปญฺจ, อารมฺมเณ ปญฺจ ฯเปฯ อวิคเต ปญฺจ. (สํขิตฺตํ. สวิตกฺก\-ทุก\-อกุสล\-สทิสํ. สหชาตวารมฺปิ ฯเปฯ ปญฺหาวารมฺปิ วิตฺถาเรตพฺพํ.)}

\proseitem{41}{สวิจารํ อพฺยากตํ ธมฺมํ ปฏิจฺจ สวิจาโร อพฺยากโต ธมฺโม อุปฺปชฺชติ เหตุปจฺจยา. ตีณิ.}

\prose{อวิจารํ อพฺยากตํ ธมฺมํ ปฏิจฺจ อวิจาโร อพฺยากโต ธมฺโม อุปฺปชฺชติ เหตุปจฺจยา. ตีณิ.}

\prose{สวิจารํ อพฺยากตญฺจ อวิจารํ อพฺยากตญฺจ ธมฺมํ ปฏิจฺจ สวิจาโร อพฺยากโต ธมฺโม อุปฺปชฺชติ เหตุปจฺจยา. ตีณิ. (สํขิตฺตํ.)}

\proseitem{42}{เหตุยา นว, อารมฺมเณ นว ฯเปฯ วิปาเก นว ฯเปฯ อวิคเต นว. (สํขิตฺตํ. สวิตกฺก\-ทุก\-อพฺยากต\-สทิสํ. สหชาตวาโรปิ ฯเปฯ สมฺปยุตฺตวาโรปิ วิตฺถาเรตพฺพา.)}

\proseitem{43}{\csromanfolio{438}สวิจาโร อพฺยากโต ธมฺโม สวิจารสฺส อพฺยากตสฺส ธมฺมสฺส เหตุปจฺจเยน ปจฺจโย. (สํขิตฺตํ.)}

\proseitem{44}{เหตุยา จตฺตาริ, อารมฺมเณ นว ฯเปฯ มคฺเค จตฺตาริ ฯเปฯ อวิคเต นว. (สํขิตฺตํ.)}

\prose{นเหตุยา นว, นอารมฺมเณ นว. (สํขิตฺตํ.)}

\prose{เหตุปจฺจยา นอารมฺมเณ จตฺตาริ. (สํขิตฺตํ.)}

\prose{นเหตุปจฺจยา อารมฺมเณ นว. (สํขิตฺตํ.)}

\prose{(ยถา กุสลตฺติเก ปญฺหาวารํ, เอวํ วิตฺถาเรตพฺพํ.)}

\csromanmarkh{89. สปฺปีติกทุก}\csromanmarkhii{1. กุสลตฺติก}\csromanheadernomark{2}{89. \textbf{สปฺปีติกทุก 1. กุสลตฺติก}}
\csromanheader{2}{1-7. ปฏิจฺจาทิวาร ปจฺจยจตุกฺก}
\proseitem{45}{สปฺปีติกํ กุสลํ ธมฺมํ ปฏิจฺจ สปฺปีติโก กุสโล ธมฺโม อุปฺปชฺชติ เหตุปจฺจยา.  สปฺปีติกํ กุสลํ ธมฺมํ ปฏิจฺจ อปฺปีติโก กุสโล ธมฺโม อุปฺปชฺชติ เหตุปจฺจยา.  สปฺปีติกํ กุสลํ ธมฺมํ ปฏิจฺจ สปฺปีติโก กุสโล จ อปฺปีติโก กุสโล จ ธมฺมา อุปฺปชฺชนฺติ เหตุปจฺจยา. (3)}

\prose{อปฺปีติกํ กุสลํ ธมฺมํ ปฏิจฺจ อปฺปีติโก กุสโล ธมฺโม อุปฺปชฺชติ เหตุปจฺจยา.  อปฺปีติกํ กุสลํ ธมฺมํ ปฏิจฺจ สปฺปีติโก กุสโล ธมฺโม อุปฺปชฺชติ เหตุปจฺจยา. (2)}

\prose{สปฺปีติกํ กุสลญฺจ อปฺปีติกํ กุสลญฺจ ธมฺมํ ปฏิจฺจ สปฺปีติโก กุสโล ธมฺโม อุปฺปชฺชติ เหตุปจฺจยา. (1) (สํขิตฺตํ.)}

\proseitem{46}{เหตุยา ฉ, (สพฺพตฺถ ฉ.) ฯเปฯ อวิคเต ฉ. (สํขิตฺตํ.)}

\prose{นอธิปติยา ฉ, นปุเรชาเต ฉ ฯเปฯ นกมฺเม จตฺตาริ, นวิปฺปยุตฺเต ฉ. (สํขิตฺตํ. สหชาตวารมฺปิ ฯเปฯ สมฺปยุตฺตวารมฺปิ วิตฺถาเรตพฺพํ.)}

\proseitem{47}{สปฺปีติโก กุสโล ธมฺโม สปฺปีติกสฺส กุสลสฺส ธมฺมสฺส เหตุปจฺจเยน ปจฺจโย. ตีณิ.}

\csromanmarkh{89. สปฺปีติกทุก}\csromanmarkhii{1. กุสลตฺติก}\csromanheadernomark{2}{89. สปฺปีติกทุก 1. กุสลตฺติก}
\prose{\csromanfolio{439}อปฺปีติโก กุสโล ธมฺโม อปฺปีติกสฺส กุสลสฺส ธมฺมสฺส เหตุปจฺจเยน ปจฺจโย. (1)}

\prose{สปฺปีติโก กุสโล ธมฺโม สปฺปีติกสฺส กุสลสฺส ธมฺมสฺส อารมฺมณ\-ปจฺจเยน ปจฺจโย. (สํขิตฺตํ.)}

\proseitem{48}{เหตุยา จตฺตาริ, อารมฺมเณ นว, อธิปติยา นว, (จตฺตาริ สหชาตา\-ธิปติ.) อนนฺตเร นว, สมนนฺตเร นว, สหชาเต ฉ ฯเปฯ อุปนิสฺสเย นว, อาเสวเน นว, กมฺเม จตฺตาริ, อาหาเร จตฺตาริ, อินฺทฺริเย จตฺตาริ, ฌาเน ฉ, มคฺเค จตฺตาริ, สมฺปยุตฺเต ฉ, อตฺถิยา ฉ, นตฺถิยา นว, วิคเต นว, อวิคเต ฉ. (สํขิตฺตํ.)}

\prose{นเหตุยา นว, นอารมฺมเณ นว. (สํขิตฺตํ.)}

\prose{เหตุปจฺจยา นอารมฺมเณ จตฺตาริ. (สํขิตฺตํ.)}

\prose{นเหตุปจฺจยา อารมฺมเณ นว. (สํขิตฺตํ.) (ยถา กุสลตฺติเก ปญฺหาวารํ, เอวํ วิตฺถาเรตพฺพํ.)}

\prose{อกุสลปทเหตุ}

\proseitem{49}{สปฺปีติกํ อกุสลํ ธมฺมํ ปฏิจฺจ สปฺปีติโก อกุสโล ธมฺโม อุปฺปชฺชติ เหตุปจฺจยา. ตีณิ.}

\prose{อปฺปีติกํ อกุสลํ ธมฺมํ ปฏิจฺจ อปฺปีติโก อกุสโล ธมฺโม อุปฺปชฺชติ เหตุปจฺจยา. เทฺว.}

\prose{สปฺปีติกํ อกุสลญฺจ อปฺปีติกํ อกุสลญฺจ ธมฺมํ ปฏิจฺจ สปฺปีติโก อกุสโล ธมฺโม อุปฺปชฺชติ เหตุปจฺจยา. (1) (สํขิตฺตํ.)}

\proseitem{50}{เหตุยา ฉ, อารมฺมเณ ฉ ฯเปฯ อวิคเต ฉ. (สํขิตฺตํ. กุสลสทิสํ. สหชาตวาโรปิ ฯเปฯ ปญฺหาวาโรปิ วิตฺถาเรตพฺพา.)}

\prose{อพฺยากตปทเหตุ}

\proseitem{51}{สปฺปีติกํ อพฺยากตํ ธมฺมํ ปฏิจฺจ สปฺปีติโก อพฺยากโต ธมฺโม อุปฺปชฺชติ เหตุปจฺจยา. ตีณิ.}

\prose{อปฺปีติกํ อพฺยากตํ ธมฺมํ ปฏิจฺจ อปฺปีติโก อพฺยากโต ธมฺโม อุปฺปชฺชติ เหตุปจฺจยา. ตีณิ.}

\prose{\csromanfolio{440}สปฺปีติกํ อพฺยากตญฺจ อปฺปีติกํ อพฺยากตญฺจ ธมฺมํ ปฏิจฺจ สปฺปีติโก อพฺยากโต ธมฺโม อุปฺปชฺชติ เหตุปจฺจยา. ตีณิ. (สํขิตฺตํ.)}

\proseitem{52}{เหตุยา นว, อารมฺมเณ นว ฯเปฯ ปุเรชาเต ฉ, อาเสวเน ฉ ฯเปฯ อวิคเต นว. (สํขิตฺตํ.)}

\prose{นเหตุยา นว, นอารมฺมเณ ตีณิ, นอธิปติยา นว ฯเปฯ นกมฺเม จตฺตาริ, นวิปาเก นว, นอาหาเร เอกํ, นอินฺทฺริเย เอกํ, นฌาเน เอกํ, นมคฺเค นว, นสมฺปยุตฺเต ตีณิ, นวิปฺปยุตฺเต ฉ. (สํขิตฺตํ. สหชาตวาราทิ วิตฺถาเรตพฺโพ.)}

\proseitem{53}{สปฺปีติโก อพฺยากโต ธมฺโม สปฺปีติกสฺส อพฺยากตสฺส ธมฺมสฺส เหตุปจฺจเยน ปจฺจโย. (สํขิตฺตํ.)}

\proseitem{54}{เหตุยา จตฺตาริ, อารมฺมเณ นว, อธิปติยา นว, อนนฺตเร นว ฯเปฯ (สพฺพตฺถ นว.) อุปนิสฺสเย นว, ปุเรชาเต ตีณิ, ปจฺฉาชาเต ตีณิ, อาเสวเน นว, กมฺเม จตฺตาริ, วิปาเก นว, อาหาเร จตฺตาริ, อินฺทฺริเย จตฺตาริ, ฌาเน นว, มคฺเค จตฺตาริ, สมฺปยุตฺเต ฉ, วิปฺปยุตฺเต ปญฺจ, อตฺถิยา นว. (สํขิตฺตํ.)}

\csromanmarkh{90-91. ปีติสหคตา\-ทิทุก}\csromanmarkhii{1. กุสลตฺติก}\csromanheadernomark{2}{90-91. ปีติสหคตา\-ทิทุก 1. กุสลตฺติก}
\csromanheader{2}{1. \textbf{ปฏิจฺจวาร ปจฺจยจตุกฺก}}
\proseitem{55}{ปีติสหคตํ กุสลํ ธมฺมํ ปฏิจฺจ ฯเปฯ.}

\vspace{\tipitakaparskip}
\csromancenter{(กุสลมฺปิ อกุสลมฺปิ อพฺยากตมฺปิ สปฺปีติก\-ทุก\-สทิสํ.)}
\proseitem{56}{สุขสหคตํ กุสลํ ธมฺมํ ปฏิจฺจ ฯเปฯ. (กุสลมฺปิ อกุสลมฺปิ อพฺยากตมฺปิ สปฺปีติก\-ทุก\-สทิสํ. อกุสลํ ธมฺมํ ปฏิจฺจ ฯเปฯ. ปจฺจนีเย, นเหตุยา เอกํ.  อพฺยากตํ ธมฺมํ ปฏิจฺจ ฯเปฯ. ปจฺจนีเย, นเหตุยา นว ฯเปฯ นฌาเน ฉ, กาตพฺพา.  ปจฺจนีเย ปญฺหาวาเร กุสลากุสเล อินฺทฺริเย ฌาเน ฉ, ปญฺหาวาเร อพฺยากเต นว.)}

\vspace{\tipitakaparskip}
\csromancenter{1. กุสลตฺติก 92. อุเปกฺขาสหคตทุก 1. กุสลตฺติก}
\csromanheader{2}{1. \textbf{ปฏิจฺจวาร ปจฺจยจตุกฺก}}
\proseitem{57}{\csromanfolio{441}อุเปกฺขา\-สหคตํ กุสลํ ธมฺมํ ปฏิจฺจ ฯเปฯ. ฉ. (สปฺปีติก\-ทุก\-สทิสํ. “อุเปกฺขา”ติ นานาอุเปกฺขา อพฺยากเต. ปจฺจนีเย นเหตุยา นว, นฌาเน ฉ.)}

\proseitem{58}{อุเปกฺขา\-สหคตํ อกุสลํ ธมฺมํ ปฏิจฺจ อุเปกฺขา\-สหคโต อกุสโล ธมฺโม อุปฺปชฺชติ เหตุปจฺจยา. (สํขิตฺตํ.)}

\prose{เหตุยา ฉ. (สพฺพตฺถ ฉ. สํขิตฺตํ.)}

\proseitem{59}{อุเปกฺขา\-สหคตํ อกุสลํ ธมฺมํ ปฏิจฺจ อุเปกฺขา\-สหคโต อกุสโล ธมฺโม อุปชฺชติ นเหตุปจฺจยา. (1)}

\prose{นอุเปกฺขาสหคตํ อกุสลํ ธมฺมํ ปฏิจฺจ อุเปกฺขา\-สหคโต อกุสโล ธมฺโม อุปฺปชฺชติ นเหตุปจฺจยา. (1)}

\prose{อุเปกฺขาสหคตญฺจ นอุเปกฺขาสหคตญฺจ ธมฺมํ ปฏิจฺจ อุเปกฺขา\-สหคโต อกุสโล ธมฺโม อุปฺปชฺชติ นเหตุปจฺจยา. (1) (สํขิตฺตํ.)}

\proseitem{60}{นเหตุยา ตีณิ, นอธิปติยา ฉ ฯเปฯ นวิปฺปยุตฺเต ฉ. (สํขิตฺตํ. เอวํ สพฺพตฺถ อกุสลํ วิตฺถาเรตพฺพํ. สปฺปีติก\-ทุก\-สทิสํ.)}

\proseitem{61}{อุเปกฺขา\-สหคตํ อพฺยากตํ ธมฺมํ ปฏิจฺจ อุเปกฺขา\-สหคโต อพฺยากโต ธมฺโม อุปฺปชฺชติ เหตุปจฺจยา นว ปญฺหา. (สปฺปีติก\-ทุก\-อพฺยากต\-สทิสํ. ปญฺหาวาเร กุสลากุสเล อินฺทฺริเย ฌาเน ฉ, อพฺยากเต นว.)}

\csromanmarkh{93. กามาวจรทุก}\csromanmarkhii{1. กุสลตฺติก}\csromanheadernomark{2}{93. กามาวจรทุก 1. กุสลตฺติก}
\csromanheader{2}{1-7. ปฏิจฺจาทิวาร ปจฺจยจตุกฺก}
\proseitem{62}{กามาวจรํ กุสลํ ธมฺมํ ปฏิจฺจ กามาวจโร กุสโล ธมฺโม อุปฺปชฺชติ เหตุปจฺจยา. (1)}

\prose{นกามาวจรํ กุสลํ ธมฺมํ ปฏิจฺจ นกามาวจโร กุสโล ธมฺโม อุปฺปชฺชติ เหตุปจฺจยา. (1) (สํขิตฺตํ.)}

\proseitem{63}{\csromanfolio{442}เหตุยา เทฺว, อารมฺมเณ เทฺว, (สพฺพตฺถ เทฺว.) อวิคเต เทฺว. (สํขิตฺตํ.)}

\prose{นอธิปติยา เทฺว ฯเปฯ นวิปฺปยุตฺเต เทฺว. (สํขิตฺตํ.) (สหชาตวาราทิ วิตฺถาเรตพฺโพ.)}

\proseitem{64}{กามาวจโร กุสโล ธมฺโม กามาวจรสฺส กุสลสฺส ธมฺมสฺส เหตุปจฺจเยน ปจฺจโย. (1)}

\prose{นกามาวจโร กุสโล ธมฺโม นกามาวจรสฺส กุสลสฺส ธมฺมสฺส เหตุปจฺจเยน ปจฺจโย. (1) (สํขิตฺตํ.)}

\proseitem{65}{เหตุยา เทฺว, อารมฺมเณ จตฺตาริ, อธิปติยา ตีณิ, (กามาวจเร เอกํ, นกามาวจเร เทฺว.) อนนฺตเร ตีณิ, (กามาวจเร เทฺว, นกามาวจเร เอกํ.) ฯเปฯ สหชาเต เทฺว ฯเปฯ อุปนิสฺสเย จตฺตาริ, อาเสวเน ตีณิ, กมฺเม เทฺว, อาหาเร เทฺว ฯเปฯ นตฺถิยา ตีณิ ฯเปฯ อวิคเต เทฺว. (สํขิตฺตํ.)}

\proseitem{66}{กามาวจรํ อกุสลํ ธมฺมํ ปฏิจฺจ กามาวจโร อกุสโล ธมฺโม อุปฺปชฺชติ เหตุปจฺจยา. (สํขิตฺตํ.)}

\prose{เหตุยา เอกํ, อารมฺมเณ เอกํ ฯเปฯ อวิคเต เอกํ. (สํขิตฺตํ.) (สหชาตวาเรปิ ฯเปฯ ปญฺหาวาเรปิ สพฺพตฺถ เอกํ.)}

\proseitem{67}{กามาวจรํ อพฺยากตํ ธมฺมํ ปฏิจฺจ กามาวจโร อพฺยากโต ธมฺโม อุปฺปชฺชติ เหตุปจฺจยา. ตีณิ.}

\prose{นกามาวจรํ อพฺยากตํ ธมฺมํ ปฏิจฺจ นกามาวจโร อพฺยากโต ธมฺโม อุปฺปชฺชติ เหตุปจฺจยา. ตีณิ.}

\prose{กามาวจรํ อพฺยากตญฺจ นกามาวจรํ อพฺยากตญฺจ ธมฺมํ ปฏิจฺจ กามาวจโร อพฺยากโต ธมฺโม อุปฺปชฺชติ เหตุปจฺจยา. ตีณิ. (สํขิตฺตํ.)}

\proseitem{68}{เหตุยา นว, อารมฺมเณ จตฺตาริ, อธิปติยา ปญฺจ ฯเปฯ อญฺญมญฺเญ ฉ ฯเปฯ ปุเรชาเต เทฺว, อาเสวเน เทฺว ฯเปฯ อวิคเต นว. (สํขิตฺตํ.)}

\csromanmarkh{94. รูปาวจรทุก}\csromanmarkhii{1. กุสลตฺติก}\csromanheadernomark{2}{94. รูปาวจรทุก 1. กุสลตฺติก}
\prose{\csromanfolio{443}นเหตุยา เอกํ, นอารมฺมเณ ตีณิ, นอธิปติยา นว ฯเปฯ นกมฺเม เทฺว, นวิปาเก ปญฺจ, นอาหาเร เอกํ, นอินฺทฺริเย เอกํ, นฌาเน เอกํ, นมคฺเค เอกํ, นสมฺปยุตฺเต ตีณิ, นวิปฺปยุตฺเต เทฺว, โนนตฺถิยา ตีณิ, โนวิคเต ตีณิ. (สํขิตฺตํ. สหชาตวาราทิ วิตฺถาเรตพฺโพ.)}

\proseitem{69}{กามาวจโร อพฺยากโต ธมฺโม กามาวจรสฺส อพฺยากตสฺส ธมฺมสฺส เหตุปจฺจเยน ปจฺจโย. (1)}

\prose{นกามาวจโร อพฺยากโต ธมฺโม นกามาวจรสฺส อพฺยากตสฺส ธมฺมสฺส เหตุปจฺจเยน ปจฺจโย. ตีณิ. (สํขิตฺตํ.)}

\proseitem{70}{เหตุยา จตฺตาริ, อารมฺมเณ จตฺตาริ, อธิปติยา จตฺตาริ, (กามาวจเร เอกํ, นกามาวจเร ตีณิ, กามาวจเร สหชาตา\-ธิปติเยว.) อนนฺตเร จตฺตาริ, สมนนฺตเร จตฺตาริ, สหชาเต สตฺต, อญฺญมญฺเญ ฉ, นิสฺสเย สตฺต, อุปนิสฺสเย จตฺตาริ, ปุเรชาเต เทฺว, ปจฺฉาชาเต เทฺว, อาเสวเน ตีณิ, กมฺเม จตฺตาริ, วิปาเก จตฺตาริ, อาหาเร จตฺตาริ ฯเปฯ สมฺปยุตฺเต เทฺว, วิปฺปยุตฺเต ตีณิ, อตฺถิยา สตฺต, นตฺถิยา จตฺตาริ ฯเปฯ อวิคเต สตฺต. (สํขิตฺตํ.)}

\csromanmarkh{94. รูปาวจรทุก}\csromanmarkhii{1. กุสลตฺติก}\csromanheadernomark{2}{94. รูปาวจรทุก 1. กุสลตฺติก}
\csromanheader{2}{1-7. ปฏิจฺจาทิวาร ปจฺจยจตุกฺก}
\proseitem{71}{รูปาวจรํ กุสลํ ธมฺมํ ปฏิจฺจ รูปาวจโร กุสโล ธมฺโม อุปฺปชฺชติ เหตุปจฺจยา. (1)}

\prose{นรูปาวจรํ กุสลํ ธมฺมํ ปฏิจฺจ นรูปาวจโร กุสโล ธมฺโม อุปฺปชฺชติ เหตุปจฺจยา. (1) (สํขิตฺตํ.)}

\proseitem{72}{เหตุยา เทฺว, อารมฺมเณ เทฺว. (สพฺพตฺถ เทฺว, สํขิตฺตํ.)}

\prose{นอธิปติยา เทฺว ฯเปฯ นปุเรชาเต เอกํ, นอาเสวเน เอกํ ฯเปฯ นวิปฺปยุตฺเต เอกํ. (สํขิตฺตํ. สหชาตวาราทิ วิตฺถาเรตพฺโพ.)}

\proseitem{73}{รูปาวจโร กุสโล ธมฺโม รูปาวจรสฺส กุสลสฺส ธมฺมสฺส เหตุปจฺจเยน ปจฺจโย. (1)}

\prose{\csromanfolio{444}นรูปาวจโร กุสโล ธมฺโม นรูปาวจรสฺส กุสลสฺส ธมฺมสฺส เหตุปจฺจเยน ปจฺจโย. (1) (สํขิตฺตํ.)}

\proseitem{74}{เหตุยา เทฺว, อารมฺมเณ จตฺตาริ, อธิปติยา ตีณิ, (รูปาวจเร เอกํ, สหชาตา\-ธิปติเยว, นรูปาวจเร เทฺว.) อนนฺตเร ตีณิ, (รูปาวจเร เอกํ, นรูปาวจเร เทฺว.) สมนนฺตเร ตีณิ, สหชาเต เทฺว ฯเปฯ อุปนิสฺสเย จตฺตาริ, อาเสวเน ตีณิ, กมฺเม เทฺว ฯเปฯ อตฺถิยา เทฺว, นตฺถิยา ตีณิ ฯเปฯ อวิคเต เทฺว. (สํขิตฺตํ.)}

\proseitem{75}{นรูปาวจรํ อกุสลํ ธมฺมํ ปฏิจฺจ นรูปาวจโร อกุสโล ธมฺโม อุปฺปชฺชติ เหตุปจฺจยา. (สพฺพตฺถ เอกํ.)}

\prose{รูปาวจรํ อพฺยากตํ ธมฺมํ ปฏิจฺจ รูปาวจโร อพฺยากโต ธมฺโม อุปฺปชฺชติ เหตุปจฺจยา. ตีณิ.}

\prose{นรูปาวจรํ อพฺยากตํ ธมฺมํ ปฏิจฺจ นรูปาวจโร อพฺยากโต ธมฺโม อุปฺปชฺชติ เหตุปจฺจยา. ตีณิ.}

\prose{รูปาวจรํ อพฺยากตญฺจ นรูปวจรํ อพฺยากตญฺจ ธมฺมํ ปฏิจฺจ รูปาวจโร อพฺยากโต ธมฺโม อุปฺปชฺชติ เหตุปจฺจยา. ตีณิ. (สํขิตฺตํ.)}

\proseitem{76}{เหตุยา นว, อารมฺมเณ จตฺตาริ, อธิปติยา ปญฺจ ฯเปฯ ปุเรชาเต อาเสวเน เทฺว ฯเปฯ อวิคเต นว. (สํขิตฺตํ. ยถา กามาวจร\-ทุก\-อพฺยากต\-สทิสํ, เอตฺตกาเยว ปญฺหา เหฏฺฐุปริกํ ปริวตฺตนฺติ. สหชาตวารมฺปิ ฯเปฯ ปญฺหาวารมฺปิ สพฺพตฺถ วิตฺถาเรตพฺพํ.)}

\csromanmarkh{95. อรูปาวจรทุก}\csromanmarkhii{1. กุสลตฺติก}\csromanheadernomark{2}{95. \textbf{อรูปาวจรทุก 1. กุสลตฺติก}}
\csromanheader{2}{1-7. ปฏิจฺจาทิวาร ปจฺจยจตุกฺก}
\proseitem{77}{อรูปาวจรํ กุสลํ ธมฺมํ ปฏิจฺจ อรูปาวจโร กุสโล ธมฺโม อุปฺปชฺชติ เหตุปจฺจยา. (1)}

\prose{นอรูปาวจรํ กุสลํ ธมฺมํ ปฏิจฺจ นอรูปาวจโร กุสโล ธมฺโม อุปฺปชฺชติ เหตุปจฺจยา. (1) (สํขิตฺตํ.)}

\csromanmarkh{95. อรูปาวจรทุก}\csromanmarkhii{1. กุสลตฺติก}\csromanheadernomark{2}{95. อรูปาวจรทุก 1. กุสลตฺติก}
\proseitem{78}{\csromanfolio{445}เหตุยา เทฺว, อารมฺมเณ เทฺว ฯเปฯ อวิคเต เทฺว. (สํขิตฺตํ.)}

\prose{นอธิปติยา เทฺว ฯเปฯ นอาเสวเน เอกํ ฯเปฯ นวิปฺปยุตฺเต เทฺว. (สํขิตฺตํ. สหชาตวาราทิ วิตฺถาเรตพฺโพ.)}

\proseitem{79}{อรูปาวจโร กุสโล ธมฺโม อรูปาวจรสฺส กุสลสฺส ธมฺมสฺส เหตุปจฺจเยน ปจฺจโย. (1)}

\prose{นอรูปาวจโร กุสโล ธมฺโม นอรูปาวจรสฺส กุสลสฺส ธมฺมสฺส เหตุปจฺจเยน ปจฺจโย. (1) (สํขิตฺตํ.)}

\proseitem{80}{เหตุยา เทฺว, อารมฺมเณ ตีณิ, อธิปติยา ตีณิ, อนนฺตเร ตีณิ ฯเปฯ สหชาเต เทฺว ฯเปฯ อุปนิสฺสเย จตฺตาริ, อาเสวเน ตีณิ, กมฺเม เทฺว ฯเปฯ อตฺถิยา เทฺว, นตฺถิยา ตีณิ ฯเปฯ อวิคเต เทฺว. (สํขิตฺตํ.)}

\proseitem{81}{นอรูปาวจรํ อกุสลํ ธมฺมํ ปฏิจฺจ นอรูปาวจโร อกุสโล ธมฺโม อุปฺปชฺชติ เหตุปจฺจยา. เอกํ. (สพฺพตฺถ เอกํ. สํขิตฺตํ.)}

\prose{อรูปาวจรํ อพฺยากตํ ธมฺมํ ปฏิจฺจ อรูปาวจโร อพฺยากโต ธมฺโม อุปฺปชฺชติ เหตุปจฺจยา. ตีณิ.}

\prose{นอรูปาวจรํ อพฺยากตํ ธมฺมํ ปฏิจฺจ นอรูปาวจโร อพฺยากโต ธมฺโม อุปฺปชฺชติ เหตุปจฺจยา. (1)}

\prose{อรูปาวจรํ อพฺยากตญฺจ นอรูปาวจรํ อพฺยากตญฺจ ธมฺมํ ปฏิจฺจ นอรูปาวจโร อพฺยากโต ธมฺโม อุปฺปชฺชติ เหตุปจฺจยา. (1) (สํขิตฺตํ.)}

\proseitem{82}{เหตุยา ปญฺจ, อารมฺมเณ เทฺว, อธิปติยา ปญฺจ ฯเปฯ อวิคเต ปญฺจ. (สํขิตฺตํ.)}

\prose{นเหตุยา เอกํ, นอารมฺมเณ ตีณิ, นอธิปติยา เทฺว ฯเปฯ นปุเรชาเต จตฺตาริ, นปจฺฉาชาเต นอาเสวเน ปญฺจ, นกมฺเม เทฺว, นวิปาเก ปญฺจ, นอาหาเร เอกํ ฯเปฯ นวิปฺปยุตฺเต เทฺว ฯเปฯ โนวิคเต ตีณิ. (สํขิตฺตํ. สหชาตวาราทิ วิตฺถาเรตพฺโพ.)}

\proseitem{83}{อรูปาวจโร อพฺยากโต ธมฺโม อรูปาวจรสฺส อพฺยากตสฺส ธมฺมสฺส เหตุปจฺจเยน ปจฺจโย. ตีณิ.}

\prose{นอรูปาวจโร อพฺยากโต ธมฺโม นอรูปาวจรสฺส อพฺยากตสฺส ธมฺมสฺส เหตุปจฺจเยน ปจฺจโย. (1) (สํขิตฺตํ.)}

\proseitem{84}{\csromanfolio{446}เหตุยา จตฺตาริ, อารมฺมเณ ตีณิ, (อรูปาวจร\-มูเล เทฺว, นอรูปาวจเร เอกํ.) อธิปติยา จตฺตาริ, (อรูปาวจร\-มูเล ตีณิ, นอรูเป เอกํ.) อนนฺตเร จตฺตาริ ฯเปฯ สหชาเต ปญฺจ, อญฺญมญฺเญ เทฺว, นิสฺสเย สตฺต, อุปนิสฺสเย จตฺตาริ, ปุเรชาเต เทฺว, ปจฺฉาชาเต เทฺว, อาเสวเน ตีณิ, กมฺเม จตฺตาริ, วิปาเก เทฺว ฯเปฯ สมฺปยุตฺเต เทฺว, วิปฺปยุตฺเต ตีณิ, อตฺถิยา สตฺต, นตฺถิยา จตฺตาริ ฯเปฯ อวิคเต สตฺต. (สํขิตฺตํ.)}

\csromanmarkh{96. ปริยาปนฺนทุก}\csromanmarkhii{1. กุสลตฺติก}\csromanheadernomark{2}{96. \textbf{ปริยาปนฺนทุก 1. กุสลตฺติก}}
\csromanheader{2}{1. \textbf{ปฏิจฺจวาร ปจฺจยจตุกฺก}}
\proseitem{85}{ปริยาปนฺนํ กุสลํ ธมฺมํ ปฏิจฺจ ปริยาปนฺโน กุสโล ธมฺโม อุปฺปชฺชติ เหตุปจฺจยา. (1)}

\prose{อปริยาปนฺนํ กุสลํ ธมฺมํ ปฏิจฺจ อปริยาปนฺโน กุสโล ธมฺโม อุปฺปชฺชติ เหตุปจฺจยา. (1) (สํขิตฺตํ.)}

\proseitem{86}{เหตุยา เทฺว, อารมฺมเณ เทฺว ฯเปฯ อวิคเต เทฺว. (สํขิตฺตํ. โลกิย\-ทุก\-กุสล\-สทิสํ. สหชาตวาโรปิ ฯเปฯ ปญฺหาวาโรปิ วิตฺถาเรตพฺพา.)}

\proseitem{87}{ปริยาปนฺนํ อกุสลํ ธมฺมํ ปฏิจฺจ ปริยาปนฺโน อกุสโล ธมฺโม อุปฺปชฺชติ เหตุปจฺจยา. (สํขิตฺตํ.)}

\prose{เหตุยา เอกํ, อารมฺมเณ เอกํ ฯเปฯ อวิคเต เอกํ. (สํขิตฺตํ.) (สหชาตวาเรปิ ฯเปฯ ปญฺหาวาเรปิ สพฺพตฺถ เอกํ.)}

\proseitem{88}{ปริยาปนฺนํ อพฺยากตํ ธมฺมํ ปฏิจฺจ ปริยาปนฺโน อพฺยากโต ธมฺโม อุปฺปชฺชติ เหตุปจฺจยา. (1)}

\prose{อปริยาปนฺนํ อพฺยากตํ ธมฺมํ ปฏิจฺจ อปริยาปนฺโน อพฺยากโต ธมฺโม อุปฺปชฺชติ เหตุปจฺจยา. ตีณิ.}

\prose{ปริยาปนฺนํ อพฺยากตญฺจ อปริยาปนฺนํ อพฺยากตญฺจ ธมฺมํ ปฏิจฺจ ปริยาปนฺโน อพฺยากโต ธมฺโม อุปฺปชฺชติ เหตุปจฺจยา. (1) (สํขิตฺตํ.)}

\csromanmarkh{98. นิยตทุก}\csromanmarkhii{1. กุสลตฺติก}\csromanheadernomark{2}{98. นิยตทุก 1. กุสลตฺติก}
\proseitem{89}{\csromanfolio{447}เหตุยา ปญฺจ, อารมฺมเณ เทฺว ฯเปฯ อวิคเต ปญฺจ. (สํขิตฺตํ. โลกิย\-ทุก\-อพฺยากต\-สทิสํ. สหชาตวารมฺปิ ฯเปฯ ปญฺหาวารมฺปิ สพฺพตฺถ วิตฺถาเรตพฺพํ.)}

\csromanmarkh{97. นิยฺยานิกทุก}\csromanmarkhii{1. กุสลตฺติก}\csromanheadernomark{2}{97. \textbf{นิยฺยานิกทุก 1. กุสลตฺติก}}
\csromanheader{2}{1. \textbf{ปฏิจฺจวาร ปจฺจยจตุกฺก}}
\proseitem{90}{นิยฺยานิกํ กุสลํ ธมฺมํ ปฏิจฺจ นิยฺยานิโก กุสโล ธมฺโม อุปฺปชฺชติ เหตุปจฺจยา. (1)}

\prose{อนิยฺยานิกํ กุสลํ ธมฺมํ ปฏิจฺจ อนิยฺยานิโก กุสโล ธมฺโม อุปฺปชฺชติ เหตุปจฺจยา. (1) (สํขิตฺตํ.)}

\proseitem{91}{เหตุยา เทฺว, อารมฺมเณ เทฺว ฯเปฯ อวิคเต เทฺว. (สํขิตฺตํ. โลกิย\-ทุก\-สทิสํ. สหชาตวารมฺปิ ฯเปฯ ปญฺหาวารมฺปิ วิตฺถาเรตพฺพํ.)}

\proseitem{92}{อนิยฺยานิกํ อกุสลํ ธมฺมํ ปฏิจฺจ อนิยฺยานิโก อกุสโล ธมฺโม อุปฺปชฺชติ เหตุปจฺจยา. (1) (สํขิตฺตํ.)}

\prose{เหตุยา เอกํ, อารมฺมเณ เอกํ ฯเปฯ อวิคเต เอกํ. (สํขิตฺตํ.) (สหชาตวาเรปิ ฯเปฯ ปญฺหาวาเรปิ สพฺพตฺถ เอกํ.)}

\proseitem{93}{อนิยฺยานิกํ อพฺยากตํ ธมฺมํ ปฏิจฺจ อนิยฺยานิโก อพฺยากโต ธมฺโม อุปฺปชฺชติ เหตุปจฺจยา. (1) (สํขิตฺตํ.)}

\prose{เหตุยา เอกํ, อารมฺมเณ เอกํ ฯเปฯ อวิคเต เอกํ. (สํขิตฺตํ.)}

\vspace{\tipitakaparskip}
\csromancenter{(สหชาตวาเรปิ ฯเปฯ ปญฺหาวาเรปิ สพฺพตฺถ เอกํ.)}
\csromanmarkh{98. นิยตทุก}\csromanmarkhii{1. กุสลตฺติก}\csromanheadernomark{2}{98. นิยตทุก 1. กุสลตฺติก}
\csromanheader{2}{1-7. ปฏิจฺจาทิวาร ปจฺจยจตุกฺก}
\proseitem{94}{นิยตํ กุสลํ ธมฺมํ ปฏิจฺจ นิยโต กุสโล ธมฺโม อุปฺปชฺชติ เหตุปจฺจยา. (1)}

\prose{\csromanfolio{448}อนิยตํ กุสลํ ธมฺมํ ปฏิจฺจ อนิยโต กุสโล ธมฺโม อุปฺปชฺชติ เหตุปจฺจยา. (1) (สํขิตฺตํ.)}

\proseitem{95}{เหตุยา เทฺว, อารมฺมเณ เทฺว ฯเปฯ อวิคเต เทฺว. (สํขิตฺตํ. โลกิย\-ทุก\-กุสล\-สทิสํ. สหชาตวารมฺปิ ฯเปฯ ปญฺหาวารมฺปิ วิตฺถาเรตพฺพํ.)}

\proseitem{96}{นิยตํ อกุสลํ ธมฺมํ ปฏิจฺจ นิยโต อกุสโล ธมฺโม อุปฺปชฺชติ เหตุปจฺจยา. (1)}

\prose{อนิยตํ อกุสลํ ธมฺมํ ปฏิจฺจ อนิยโต อกุสโล ธมฺโม อุปฺปชฺชติ เหตุปจฺจยา. (1) (สํขิตฺตํ.)}

\proseitem{97}{เหตุยา เทฺว, อารมฺมเณ เทฺว, (สพฺพตฺถ เทฺว.) อวิคเต เทฺว. (สํขิตฺตํ.)}

\prose{นเหตุยา เอกํ, นอาธิปติยา เทฺว, นปุเรชาเต เอกํ ฯเปฯ นกมฺเม เทฺว ฯเปฯ นวิปฺปยุตฺเต เอกํ. (สํขิตฺตํ. สหชาตวาราทิ วิตฺถาเรตพฺโพ.)}

\proseitem{98}{นิยโต อกุสโล ธมฺโม นิยตสฺส อกุสลสฺส ธมฺมสฺส เหตุปจฺจเยน ปจฺจโย. (1)}

\prose{อนิยโต อกุสโล ธมฺโม อนิยตสฺส อกุสลสฺส ธมฺมสฺส เหตุปจฺจเยน ปจฺจโย. (1) (สํขิตฺตํ.)}

\proseitem{99}{เหตุยา เทฺว, อารมฺมเณ ตีณิ, อธิปติยา เทฺว, (นิยเต สหชาตา\-ธิปติ. ทุติเย อารมฺมณา\-ธิปติ, สหชาตา\-ธิปติ.) อนนฺตเร เทฺว ฯเปฯ นิสฺสเย เทฺว, อุปนิสฺสเย ตีณิ, อาเสวเน เทฺว, กมฺเม เทฺว ฯเปฯ อวิคเต เทฺว. (สํขิตฺตํ.)}

\proseitem{100}{อนิยตํ อพฺยากตํ ธมฺมํ ปฏิจฺจ อนิยโต อพฺยากโต ธมฺโม อุปฺปชฺชติ เหตุปจฺจยา. (สํขิตฺตํ.)}

\prose{เหตุยา เอกํ, อารมฺมเณ เอกํ ฯเปฯ อวิคเต เอกํ. (สํขิตฺตํ. สหชาตวาเรปิ ฯเปฯ ปญฺหาวาเรปิ สพฺพตฺถ เอกํ.)}

\csromanmarkh{100. สรณทุก}\csromanmarkhii{1. กุสลตฺติก 99. สอุตฺตรทุก 1. กุสลตฺติก}\csromanheadernomark{2}{100. สรณทุก 1. กุสลตฺติก 99. สอุตฺตรทุก 1. กุสลตฺติก}
\csromanheader{2}{1. \textbf{ปฏิจฺจวาร ปจฺจยจตุกฺก}}
\proseitem{101}{\csromanfolio{449}สอุตฺตรํ กุสลํ ธมฺมํ ปฏิจฺจ สอุตฺตโร กุสโล ธมฺโม อุปฺปชฺชติ เหตุปจฺจยา. (1)}

\prose{อนุตฺตรํ กุสลํ ธมฺมํ ปฏิจฺจ อนุตฺตโร กุสโล ธมฺโม อุปฺปชฺชติ เหตุปจฺจยา. (1) (สํขิตฺตํ.)}

\proseitem{102}{เหตุยา เทฺว, อารมฺมเณ เทฺว ฯเปฯ อวิคเต เทฺว. (สํขิตฺตํ. สหชาตวาโรปิ ฯเปฯ ปญฺหาวาโรปิ สพฺพตฺถ วิตฺถาเรตพฺพา.)}

\proseitem{103}{สอุตฺตรํ อกุสลํ ธมฺมํ ปฏิจฺจ สอุตฺตโร อกุสโล ธมฺโม อุปฺปชฺชติ เหตุปจฺจยา. (สํขิตฺตํ.)}

\prose{เหตุยา เอกํ, อารมฺมเณ เอกํ ฯเปฯ อวิคเต เอกํ. (สํขิตฺตํ. สหชาตวาเรปิ ฯเปฯ ปญฺหาวาเรปิ สพฺพตฺถ เอกํ.)}

\prose{สอุตฺตรํ อพฺยากตํ ธมฺมํ ปฏิจฺจ สอุตฺตโร อพฺยากโต ธมฺโม อุปฺปชฺชติ เหตุปจฺจยา. (1)}

\prose{อนุตฺตรํ อพฺยากตํ ธมฺมํ ปฏิจฺจ อนุตฺตโร อพฺยากโต ธมฺโม อุปฺปชฺชติ เหตุปจฺจยา. ตีณิ.}

\prose{สอุตฺตรํ อพฺยากตญฺจ อนุตฺตรํ อพฺยากตญฺจ ธมฺมํ ปฏิจฺจ สอุตฺตโร อพฺยากโต ธมฺโม อุปฺปชฺชติ เหตุปจฺจยา. (1) (สํขิตฺตํ.)}

\proseitem{104}{เหตุยา ปญฺจ, อารมฺมเณ เทฺว ฯเปฯ อวิคเต ปญฺจ. (สํขิตฺตํ. สหชาตวารมฺปิ ฯเปฯ ปญฺหาวารมฺปิ สพฺพตฺถ วิตฺถาเรตพฺพํ. โลกิย\-ทุก\-อพฺยากต\-สทิสํ.)}

\csromanmatikaanchor{mtk.84}
\csromanmatikaanchor{mtk.85}
\csromantocmarknopage{subparagraph}{100. สรณทุก}{mtk.84}
\bookmark[dest={mtk.84},level=5]{100. สรณทุก}
\csromantocmarkref{subparagraph}{1. กุสลตฺติก}{mtk.85}
\bookmark[dest={mtk.85},level=5]{1. กุสลตฺติก}
\csromanmarkh{100. สรณทุก}\csromanmarkhii{1. กุสลตฺติก}\csromanheadernomark{6}{100. สรณทุก 1. กุสลตฺติก}
\csromanheader{2}{1-7. ปฏิจฺจาทิวาร ปจฺจยจตุกฺก}
\proseitem{105}{อรณํ กุสลํ ธมฺมํ ปฏิจฺจ อรโณ กุสโล ธมฺโม อุปฺปชฺชติ เหตุปจฺจยา. (สํขิตฺตํ.)}

\prose{\csromanfolio{450}เหตุยา เอกํ, อารมฺมเณ เอกํ ฯเปฯ อวิคเต เอกํ. (สํขิตฺตํ. สหชาตวาเรปิ ฯเปฯ ปญฺหาวาเรปิ สพฺพตฺถ เอกํ.)}

\proseitem{106}{สรณํ อกุสลํ ธมฺมํ ปฏิจฺจ สรโณ อกุสโล ธมฺโม อุปฺปชฺชติ เหตุปจฺจยา. (สํขิตฺตํ.)}

\prose{เหตุยา เอกํ, อารมฺมเณ เอกํ ฯเปฯ อวิคเต เอกํ. (สํขิตฺตํ. สหชาตวาเรปิ ฯเปฯ ปญฺหาวาเรปิ สพฺพตฺถ เอกํ.)}

\proseitem{107}{อรณํ อพฺยากตํ ธมฺมํ ปฏิจฺจ อรโณ อพฺยากโต ธมฺโม อุปฺปชฺชติ เหตุปจฺจยา. (สํขิตฺตํ.)}

\prose{เหตุยา เอกํ, อารมฺมเณ เอกํ ฯเปฯ อวิคเต เอกํ. (สํขิตฺตํ. สหชาตวาเรปิ ฯเปฯ สมฺปยุตฺตวาเรปิ สพฺพตฺถ เอกํ.)}

\proseitem{108}{อรโณ อพฺยากโต ธมฺโม อรณสฺส อพฺยากตสฺส ธมฺมสฺส เหตุปจฺจเยน ปจฺจโย. (สํขิตฺตํ.)}

\prosewithcloser{เหตุยา เอกํ, อารมฺมเณ เอกํ ฯเปฯ อวิคเต เอกํ. (สํขิตฺตํ.) (ยถา กุสลตฺติเก ปญฺหาวารํ, เอวํ วิตฺถาเรตพฺพํ.)}{ปิฏฺฐิ\-ทุก\-กุสล\-ตฺติกํ นิฏฺฐิตํ.}
\csromanmatikaanchor{mtk.86}
\csromantocmarkref{subparagraph}{2. เวทนาตฺติก}{mtk.86}
\bookmark[dest={mtk.86},level=5]{2. เวทนาตฺติก}
\csromanmarkh{100. สรณทุก}\csromanmarkhii{2. เวทนาตฺติก}\csromanheadernomark{6}{100. \textbf{สรณทุก 2. เวทนาตฺติก}}
\csromanheader{2}{1. \textbf{ปฏิจฺจวาร ปจฺจยจตุกฺก}}
\proseitem{1}{สรณํ สุขาย เวทนาย สมฺปยุตฺตํ ธมฺมํ ปฏิจฺจ สรโณ สุขาย เวทนาย สมฺปยุตฺโต ธมฺโม อุปฺปชฺชติ เหตุปจฺจยา. (1)}

\prose{อรณํ สุขาย เวทนาย สมฺปยุตฺตํ ธมฺมํ ปฏิจฺจ อรโณ สุขาย เวทนาย สมฺปยุตฺโต ธมฺโม อุปฺปชฺชติ เหตุปจฺจยา. (1) (สํขิตฺตํ.)}

\proseitem{2}{เหตุยา เทฺว, อารมฺมเณ เทฺว ฯเปฯ กมฺเม เทฺว, วิปาเก เอกํ ฯเปฯ อวิคเต เทฺว. (สํขิตฺตํ. สหชาตวาโรปิ ฯเปฯ ปญฺหาวาโรปิ สพฺพตฺถ วิตฺถาเรตพฺพา.)}

\csromanmatikaanchor{mtk.87}
\csromantocmarkref{subparagraph}{3. วิปากตฺติก}{mtk.87}
\bookmark[dest={mtk.87},level=5]{3. วิปากตฺติก}
\csromanmarkh{100. สรณทุก}\csromanmarkhii{3. วิปากตฺติก}\csromanheadernomark{6}{100. สรณทุก 3. วิปากตฺติก}
\proseitem{3}{\csromanfolio{451}สรณํ ทุกฺขาย เวทนาย สมฺปยุตฺตํ ธมฺมํ ปฏิจฺจ สรโณ ทุกฺขาย เวทนาย สมฺปยุตฺโต ธมฺโม อุปฺปชฺชติ เหตุปจฺจยา. (1)}

\prose{สรณํ ทุกฺขาย เวทนาย สมฺปยุตฺตํ ธมฺมํ ปฏิจฺจ สรโณ ทุกฺขาย เวทนาย สมฺปยุตฺโต ธมฺโม อุปฺปชฺชติ อารมฺมณ\-ปจฺจยา. (1) (สํขิตฺตํ.)}

\proseitem{4}{เหตุยา เอกํ, อารมฺมเณ เทฺว ฯเปฯ อวิคเต เทฺว. (สํขิตฺตํ. สหชาตวารมฺปิ ฯเปฯ ปญฺหาวารมฺปิ วิตฺถาเรตพฺพํ.)}

\proseitem{5}{สรณํ อทุกฺขม\-สุขาย เวทนาย สมฺปยุตฺตํ ธมฺมํ ปฏิจฺจ สรโณ อทุกฺขม\-สุขาย เวทนาย สมฺปยุตฺโต ธมฺโม อุปฺปชฺชติ เหตุปจฺจยา. (1)}

\prose{อรณํ อทุกฺขม\-สุขาย เวทนาย สมฺปยุตฺตํ ธมฺมํ ปฏิจฺจ อรโณ อทุกฺขม\-สุขาย เวทนาย สมฺปยุตฺโต ธมฺโม อุปฺปชฺชติ เหตุปจฺจยา. (สํขิตฺตํ.)}

\proseitemwithcloser{6}{เหตุยา เทฺว, อารมฺมเณ เทฺว ฯเปฯ อวิคเต เทฺว. (สํขิตฺตํ. สหชาตวารมฺปิ ฯเปฯ ปญฺหาวารมฺปิ สพฺพตฺถ วิตฺถาเรตพฺพํ.)}{สรณ\-ทุก\-เวทนา\-ตฺติกํ นิฏฺฐิตํ.}
\csromanmarkh{100. สรณทุก}\csromanmarkhii{3. วิปากตฺติก}\csromanheadernomark{2}{100. สรณทุก 3. วิปากตฺติก}
\csromanheader{2}{1. \textbf{ปฏิจฺจวาร ปจฺจยจตุกฺก}}
\proseitem{7}{อรณํ วิปากํ ธมฺมํ ปฏิจฺจ อรโณ วิปาโก ธมฺโม อุปฺปชฺชติ เหตุปจฺจยา. (สํขิตฺตํ.)}

\prose{เหตุยา เอกํ, อารมฺมเณ เอกํ ฯเปฯ อวิคเต เอกํ. (สํขิตฺตํ. สหชาตวาเรปิ ฯเปฯ ปญฺหาวาเรปิ สพฺพตฺถ เอกํ.)}

\proseitem{8}{สรณํ วิปาก\-ธมฺม\-ธมฺมํ ปฏิจฺจ สรโณ วิปาก\-ธมฺม\-ธมฺโม อุปฺปชฺชติ เหตุปจฺจยา. (1)}

\prose{อรณํ วิปาก\-ธมฺม\-ธมฺมํ ปฏิจฺจ อรโณ วิปาก\-ธมฺม\-ธมฺโม อุปฺปชฺชติ เหตุปจฺจยา. (1) (สํขิตฺตํ.)}

\proseitem{9}{\csromanfolio{452}เหตุยา เทฺว, อารมฺมเณ เทฺว ฯเปฯ อวิคเต เทฺว. (สํขิตฺตํ. สหชาตวารมฺปิ ฯเปฯ ปญฺหาวารมฺปิ วิตฺถาเรตพฺพํ.)}

\prose{อรณํ เนว\-วิปาก\-น\-วิปากธมฺม\-ธมฺมํ ปฏิจฺจ อรโณ เนว\-วิปาก\-น\-วิปาก\-ธมฺม\-ธมฺโม อุปฺปชฺชติ เหตุปจฺจยา. (สํขิตฺตํ.)}

\proseitem{10}{เหตุยา เอกํ, อารมฺมเณ เอกํ ฯเปฯ อวิคเต เอกํ. (สํขิตฺตํ. สหชาตวาเรปิ ฯเปฯ ปญฺหาวาเรปิ สพฺพตฺถ เอกํ.)}

\csromanmatikaanchor{mtk.88}
\csromantocmarkref{subparagraph}{4. อุปทินฺนตฺติก}{mtk.88}
\bookmark[dest={mtk.88},level=5]{4. อุปทินฺนตฺติก}
\csromanmarkh{100. สรณทุก}\csromanmarkhii{4. อุปาทินฺนตฺติก}\csromanheadernomark{6}{100. \textbf{สรณทุก 4. อุปาทินฺนตฺติก}}
\csromanheader{2}{1. \textbf{ปฏิจฺจวาร ปจฺจยจตุกฺก}}
\proseitem{11}{อรณํ อุปาทินฺนุ\-ปาทานิยํ ธมฺมํ ปฏิจฺจ อรโณ อุปาทินฺนุ\-ปาทานิโย ธมฺโม อุปฺปชฺชติ เหตุปจฺจยา. (สํขิตฺตํ.)}

\prose{เหตุยา เอกํ, อารมฺมเณ เอกํ ฯเปฯ อวิคเต เอกํ. (สํขิตฺตํ. สหชาตวาเรปิ ฯเปฯ ปญฺหาวาเรปิ สพฺพตฺถ เอกํ.)}

\proseitem{12}{สรณํ อนุปาทินฺนุ\-ปาทานิยํ ธมฺมํ ปฏิจฺจ สรโณ อนุปาทินฺนุ\-ปาทานิโย ธมฺโม อุปฺปชฺชติ เหตุปจฺจยา. ตีณิ.}

\prose{อรณํ อนุปาทินฺนุ\-ปาทานิยํ ธมฺมํ ปฏิจฺจ อรโณ อนุปาทินฺนุ\-ปาทานิโย ธมฺโม อุปฺปชฺชติ เหตุปจฺจยา. (1)}

\prose{สรณํ อนุปาทินฺนุ\-ปาทานิยญฺจ อรณํ อนุปาทินฺนุ\-ปาทานิยญฺจ ธมฺมํ ปฏิจฺจ อรโณ อนุปาทินฺนุ\-ปาทานิโย ธมฺโม อุปฺปชฺชติ เหตุปจฺจยา. (1) (สํขิตฺตํ.)}

\proseitem{13}{เหตุยา ปญฺจ, อารมฺมเณ เทฺว ฯเปฯ อวิคเต ปญฺจ. (สํขิตฺตํ. สหชาตวารมฺปิ ฯเปฯ ปญฺหาวารมฺปิ สพฺพตฺถ วิตฺถาเรตพฺพํ.)}

\proseitem{14}{อรณํ อนุปาทินฺน\-อนุปาทานิยํ ธมฺมํ ปฏิจฺจ อรโณ อนุปาทินฺน\-อนุปาทานิโย ธมฺโม อุปฺปชฺชติ เหตุปจฺจยา. (สํขิตฺตํ.)}

\prose{เหตุยา เอกํ, อารมฺมเณ เอกํ ฯเปฯ อวิคเต เอกํ. (สํขิตฺตํ. สหชาตวาเรปิ ฯเปฯ ปญฺหาวาเรปิ สพฺพตฺถ เอกํ.)}

\vspace{\tipitakaparskip}
\csromancenter{6. วิตกฺกตฺติก 100. สรณทุก 5. สํกิลิฏฺฐ\-ตฺติก}
\csromanheader{2}{1. \textbf{ปฏิจฺจวาร ปจฺจยจตุกฺก}}
\csromanmatikaanchor{mtk.89}
\csromantocmarkref{subparagraph}{5-21. สํกิลิฏฺฐตฺติกาทิ}{mtk.89}
\bookmark[dest={mtk.89},level=5]{5-21. สํกิลิฏฺฐตฺติกาทิ}
\proseitem{15}{\csromanfolio{453}สรณํ สํกิลิฏฺฐ\-สํกิเลสิกํ ธมฺมํ ปฏิจฺจ สรโณ สํกิลิฏฺฐ\-สํกิเลสิโก ธมฺโม อุปฺปชฺชติ เหตุปจฺจยา. (สํขิตฺตํ.)}

\prose{เหตุยา เอกํ, อารมฺมเณ เอกํ ฯเปฯ อวิคเต เอกํ. (สํขิตฺตํ. สหชาตวาเรปิ ฯเปฯ ปญฺหาวาเรปิ สพฺพตฺถ เอกํ.)}

\proseitem{16}{อรณํ อสํกิลิฏฺฐ\-สํกิเลสิกํ ธมฺมํ ปฏิจฺจ อรโณ อสํกิลิฏฺฐ\-สํกิเลสิโก ธมฺโม อุปฺปชฺชติ เหตุปจฺจยา. (สํขิตฺตํ.)}

\prose{เหตุยา เอกํ, อารมฺมเณ เอกํ ฯเปฯ อวิคเต เอกํ. (สํขิตฺตํ. สหชาตวาเรปิ ฯเปฯ ปญฺหาวาเรปิ สพฺพตฺถ เอกํ.)}

\proseitem{17}{อรณํ อสํกิลิฏฺฐ\-อสํกิเลสิกํ ธมฺมํ ปฏิจฺจ อรโณ อสํกิลิฏฺฐ\-อสํกิเลสิโก ธมฺโม อุปฺปชฺชติ เหตุปจฺจยา. (สํขิตฺตํ.)}

\prose{เหตุยา เอกํ, อารมฺมเณ เอกํ ฯเปฯ อวิคเต เอกํ. (สํขิตฺตํ. สหชาตวาเรปิ ฯเปฯ ปญฺหาวาเรปิ สพฺพตฺถ เอกํ.)}

\csromanmarkh{100. สรณทุก}\csromanmarkhii{6. วิตกฺกตฺติก}\csromanheadernomark{2}{100. สรณทุก 6. วิตกฺกตฺติก}
\csromanheader{2}{1. \textbf{ปฏิจฺจวาร ปจฺจยจตุกฺก}}
\proseitem{18}{สรณํ สวิตกฺก\-สวิจารํ ธมฺมํ ปฏิจฺจ สรโณ สวิตกฺก\-สวิจาโร ธมฺโม อุปฺปชฺชติ เหตุปจฺจยา. (1)}

\prose{อรณํ สวิตกฺก\-สวิจารํ ธมฺมํ ปฏิจฺจ อรโณ สวิตกฺก\-สวิจาโร ธมฺโม อุปฺปชฺชติ เหตุปจฺจยา. (1) (สํขิตฺตํ.)}

\proseitem{19}{เหตุยา เทฺว, อารมฺมเณ เทฺว ฯเปฯ อวิคเต เทฺว. (สํขิตฺตํ. สหชาตวารมฺปิ ฯเปฯ ปญฺหาวารมฺปิ วิตฺถาเรตพฺพํ.)}

\proseitem{20}{อรณํ อวิตกฺก\-วิจารมตฺตํ ธมฺมํ ปฏิจฺจ อรโณ อวิตกฺก\-วิจารมตฺโต ธมฺโม อุปฺปชฺชติ เหตุปจฺจยา. (สํขิตฺตํ.)}

\prose{\csromanfolio{454}เหตุยา เอกํ ฯเปฯ อวิคเต เอกํ. (สํขิตฺตํ. สหชาตวาเรปิ ฯเปฯ ปญฺหาวาเรปิ สพฺพตฺถ เอกํ.)}

\proseitem{21}{อรณํ อวิตกฺก\-อวิจารํ ธมฺมํ ปฏิจฺจ อรโณ อวิตกฺก\-อวิจาโร ธมฺโม อุปฺปชฺชติ เหตุปจฺจยา. (สํขิตฺตํ.)}

\prose{เหตุยา เอกํ, (สพฺพตฺถ เอกํ) ฯเปฯ อวิคเต เอกํ. (สํขิตฺตํ.)}

\csromanmarkh{100. สรณทุก}\csromanmarkhii{7. ปีติตฺติก}\csromanheadernomark{2}{100. \textbf{สรณทุก 7. ปีติตฺติก}}
\csromanheader{2}{1. \textbf{ปฏิจฺจวาร ปจฺจยจตุกฺก}}
\proseitem{22}{สรณํ ปีติสหคตํ ธมฺมํ ปฏิจฺจ สรโณ ปีติสหคโต ธมฺโม อุปฺปชฺชติ เหตุปจฺจยา. (1)}

\prose{อรณํ ปีติสหคตํ ธมฺมํ ปฏิจฺจ อรโณ ปีติสหคโต ธมฺโม อุปฺปชฺชติ เหตุปจฺจยา. (1) (สํขิตฺตํ.)}

\proseitem{23}{เหตุยา เทฺว, อารมฺมเณ เทฺว ฯเปฯ อวิคเต เทฺว. (สํขิตฺตํ. สหชาตวาโรปิ ฯเปฯ ปญฺหาวาโรปิ สพฺพตฺถ วิตฺถาเรตพฺพา.)}

\proseitem{24}{สรณํ สุขสหคตํ ธมฺมํ ปฏิจฺจ สรโณ สุขสหคโต ธมฺโม อุปฺปชฺชติ เหตุปจฺจยา. (1)}

\prose{อรณํ สุขสหคตํ ธมฺมํ ปฏิจฺจ อรโณ สุขสหคโต ธมฺโม อุปฺปชฺชติ เหตุปจฺจยา. (1) (สํขิตฺตํ.)}

\proseitem{25}{เหตุยา เทฺว, อารมฺมเณ เทฺว ฯเปฯ อวิคเต เทฺว. (สํขิตฺตํ. สหชาตวาโรปิ ฯเปฯ ปญฺหาวาโรปิ วิตฺถาเรตพฺพา.)}

\proseitem{26}{สรณํ อุเปกฺขา\-สหคตํ ธมฺมํ ปฏิจฺจ สรโณ อุเปกฺขา\-สหคโต ธมฺโม อุปฺปชฺชติ เหตุปจฺจยา. (1)}

\prose{อรณํ อุเปกฺขา\-สหคตํ ธมฺมํ ปฏิจฺจ อรโณ อุเปกฺขา\-สหคโต ธมฺโม อุปฺปชฺชติ เหตุปจฺจยา. (1) (สํขิตฺตํ.)}

\proseitem{27}{เหตุยา เทฺว, อารมฺมเณ เทฺว ฯเปฯ อวิคเต เทฺว. (สํขิตฺตํ. สหชาตวาเรปิ ฯเปฯ ปญฺหาวาเรปิ วิตฺถาเรตพฺโพ.)}

\vspace{\tipitakaparskip}
\csromancenter{9. ทสฺสเนน\-ปหาตพฺพ\-เหตุ\-กตฺติก 100. สรณทุก 8. ทสฺสเนน\-ปหาตพฺพ\-ตฺติก}
\csromanheader{2}{1. \textbf{ปฏิจฺจวาร ปจฺจยจตุกฺก}}
\proseitem{28}{\csromanfolio{455}สรณํ ทสฺสเนน ปหาตพฺพํ ธมฺมํ ปฏิจฺจ สรโณ ทสฺสเนน ปหาตพฺโพ ธมฺโม อุปฺปชฺชติ เหตุปจฺจยา. (สํขิตฺตํ.)}

\prose{เหตุยา เอกํ, อารมฺมเณ เอกํ ฯเปฯ อวิคเต เอกํ. (สํขิตฺตํ. สหชาตวาเรปิ ฯเปฯ ปญฺหาวาเรปิ สพฺพตฺถ เอกํ.)}

\proseitem{29}{สรณํ ภาวนาย ปหาตพฺพํ ธมฺมํ ปฏิจฺจ สรโณ ภาวนาย ปหาตพฺโพ ธมฺโม อุปฺปชฺชติ เหตุปจฺจยา. (สํขิตฺตํ.)}

\prose{เหตุยา เอกํ, อารมฺมเณ เอกํ ฯเปฯ อวิคเต เอกํ. (สํขิตฺตํ. สหชาตวาเรปิ ฯเปฯ ปญฺหาวาเรปิ สพฺพตฺถ เอกํ.)}

\proseitem{30}{อรณํ เนวทสฺสเนน นภาวนาย ปหาตพฺพํ ธมฺมํ ปฏิจฺจ อรโณ เนวทสฺสเนน นภาวนาย ปหาตพฺโพ ธมฺโม อุปฺปชฺชติ เหตุปจฺจยา.}

\prose{เหตุยา เอกํ, (สพฺพตฺถ เอกํ) ฯเปฯ อวิคเต เอกํ. (สํขิตฺตํ.)}

\csromanmarkh{100. สรณทุก}\csromanmarkhii{9. ทสฺสเนน\-ปหาตพฺพ\-เหตุ\-กตฺติก}\csromanheadernomark{2}{100. สรณทุก 9. ทสฺสเนน\-ปหาตพฺพ\-เหตุ\-กตฺติก}
\csromanheader{2}{1. \textbf{ปฏิจฺจวาร ปจฺจยจตุกฺก}}
\proseitem{31}{สรณํ ทสฺสเนน ปหาตพฺพ\-เหตุกํ ธมฺมํ ปฏิจฺจ สรโณ ทสฺสเนน ปหาตพฺพ\-เหตุโก ธมฺโม อุปฺปชฺชติ เหตุปจฺจยา. (สํขิตฺตํ.)}

\prose{เหตุยา เอกํ, อารมฺมเณ เอกํ ฯเปฯ อวิคเต เอกํ. (สํขิตฺตํ. สหชาตวาเรปิ ฯเปฯ ปญฺหาวาเรปิ สพฺพตฺถ เอกํ.)}

\proseitem{32}{สรณํ ภาวนาย ปหาตพฺพ\-เหตุกํ ธมฺมํ ปฏิจฺจ สรโณ ภาวนาย ปหาตพฺพ\-เหตุโก ธมฺโม อุปฺปชฺชติ เหตุปจฺจยา. (สํขิตฺตํ.)}

\prose{เหตุยา เอกํ, อารมฺมเณ เอกํ ฯเปฯ อวิคเต เอกํ. (สํขิตฺตํ. สหชาตวาเรปิ ฯเปฯ ปญฺหาวาเรปิ สพฺพตฺถ เอกํ.)}

\proseitem{33}{\csromanfolio{456}สรณํ เนวทสฺสเนน นภาวนาย ปหาตพฺพ\-เหตุกํ ธมฺมํ ปฏิจฺจ อรโณ เนวทสฺสเนน นภาวนาย ปหาตพฺพ\-เหตุโก ธมฺโม อุปฺปชฺชติ เหตุปจฺจยา. (สํขิตฺตํ.)}

\prose{เหตุยา ตีณิ, อารมฺมเณ เอกํ. (สพฺพตฺถ วิตฺถาเรตพฺพํ.)}

\csromanmarkh{100. สรณทุก}\csromanmarkhii{10. อาจยคามิตฺติก}\csromanheadernomark{2}{100. \textbf{สรณทุก 1}0. อาจยคามิตฺติก}
\csromanheader{2}{1. \textbf{ปฏิจฺจวาร ปจฺจยจตุกฺก}}
\proseitem{34}{สรณํ อาจยคามิํ ธมฺมํ ปฏิจฺจ สรโณ อาจยคามี ธมฺโม อุปฺปชฺชติ เหตุปจฺจยา. (1)}

\prose{อรณํ อาจยคามิํ ธมฺมํ ปฏิจฺจ อรโณ อาจยคามี ธมฺโม อุปฺปชฺชติ เหตุปจฺจยา. (1) (สํขิตฺตํ.)}

\proseitem{35}{เหตุยา เทฺว, อารมฺมเณ เทฺว ฯเปฯ อวิคเต เทฺว. (สํขิตฺตํ.)}

\prose{(สหชาตวาเรปิ ฯเปฯ ปญฺหาวาเรปิ สพฺพตฺถ วิตฺถาเรตพฺโพ.)}

\proseitem{36}{อรณํ อปจยคามิํ ธมฺมํ ปฏิจฺจ อรโณ อปจยคามี ธมฺโม อุปฺปชฺชติ เหตุปจฺจยา. (สํขิตฺตํ.)}

\prose{เหตุยา เอกํ, อารมฺมเณ เอกํ ฯเปฯ อวิคเต เอกํ. (สํขิตฺตํ.)}

\prose{(สหชาตวาเรปิ ฯเปฯ ปญฺหาวาเรปิ สพฺพตฺถ เอกํ.)}

\proseitem{37}{อรณํ เนวา\-จยคามิ\-นา\-ปจยคามิํ ธมฺมํ ปฏิจฺจ อรโณ เนวา\-จยคามิ\-นา\-ปจยคามี ธมฺโม อุปฺปชฺชติ เหตุปจฺจยา. (สํขิตฺตํ.)}

\prose{เหตุยา เอกํ, อารมฺมเณ เอกํ ฯเปฯ อวิคเต เอกํ. (สํขิตฺตํ.)}

\prose{(สหชาตวาเรปิ ฯเปฯ ปญฺหาวาเรปิ สพฺพตฺถ เอกํ.)}

\csromanmarkh{100. สรณทุก}\csromanmarkhii{12. ปริตฺตตฺติก 100. สรณทุก 11. เสกฺขตฺติก}\csromanheadernomark{2}{100. สรณทุก 12. ปริตฺตตฺติก \textbf{100. สรณทุก 11. เสกฺขตฺติก}}
\csromanheader{2}{1. \textbf{ปฏิจฺจวาร ปจฺจยจตุกฺก}}
\proseitem{38}{\csromanfolio{457}อรณํ เสกฺขํ ธมฺมํ ปฏิจฺจ อรโณ เสกฺโข ธมฺโม อุปฺปชฺชติ เหตุปจฺจยา. (สํขิตฺตํ.)}

\prose{เหตุยา เอกํ, อารมฺมเณ เอกํ ฯเปฯ อวิคเต เอกํ. (สํขิตฺตํ.)}

\prose{(สหชาตวาเรปิ ฯเปฯ ปญฺหาวาเรปิ สพฺพตฺถ เอกํ.)}

\proseitem{39}{อรณํ อเสกฺขํ ธมฺมํ ปฏิจฺจ อรโณ อเสกฺโข ธมฺโม อุปฺปชฺชติ เหตุปจฺจยา. (สํขิตฺตํ.)}

\prose{เหตุยา เอกํ, อารมฺมเณ เอกํ ฯเปฯ อวิคเต เอกํ. (สํขิตฺตํ.)}

\prose{(สหชาตวาเรปิ ฯเปฯ ปญฺหาวาเรปิ สพฺพตฺถ เอกํ.)}

\proseitem{40}{สรณํ เนว\-เสกฺข\-นาเสกฺขํ ธมฺมํ ปฏิจฺจ สรโณ เนว\-เสกฺข\-นา\-เสกฺโข ธมฺโม อุปฺปชฺชติ เหตุปจฺจยา. ตีณิ.}

\prose{อรณํ เนว\-เสกฺข\-นาเสกฺขํ ธมฺมํ ปฏิจฺจ อรโณ เนว\-เสกฺข\-นา\-เสกฺโข ธมฺโม อุปฺปชฺชติ เหตุปจฺจยา. (1)}

\prose{สรณํ เนว\-เสกฺข\-นาเสกฺขญฺจ อรณํ เนว\-เสกฺข\-นาเสกฺขญฺจ ธมฺมํ ปฏิจฺจ อรโณ เนว\-เสกฺข\-นา\-เสกฺโข ธมฺโม อุปฺปชฺชติ เหตุปจฺจยา. (1) (สํขิตฺตํ.)}

\proseitem{41}{เหตุยา ปญฺจ, อารมฺมเณ เทฺว, อธิปติยา ปญฺจ ฯเปฯ อวิคเต ปญฺจ. (สํขิตฺตํ. สหชาตวาเรปิ ฯเปฯ ปญฺหาวาเรปิ สพฺพตฺถ วิตฺถาเรตพฺโพ.)}

\csromanmarkh{100. สรณทุก}\csromanmarkhii{12. ปริตฺตตฺติก}\csromanheadernomark{2}{100. สรณทุก 12. ปริตฺตตฺติก}
\csromanheader{2}{1. \textbf{ปฏิจฺจวาร ปจฺจยจตุกฺก}}
\proseitem{42}{สรณํ ปริตฺตํ ธมฺมํ ปฏิจฺจ สรโณ ปริตฺโต ธมฺโม อุปฺปชฺชติ เหตุปจฺจยา. ตีณิ.}

\prose{อรณํ ปริตฺตํ ธมฺมํ ปฏิจฺจ อรโณ ปริตฺโต ธมฺโม อุปฺปชฺชติ เหตุปจฺจยา. (1)}

\prose{\csromanfolio{458}สรณํ ปริตฺตญฺจ อรณํ ปริตฺตญฺจ ธมฺมํ ปฏิจฺจ อรโณ ปริตฺโต ธมฺโม อุปฺปชฺชติ เหตุปจฺจยา. (1) (สํขิตฺตํ.)}

\prose{เหตุยา ปญฺจ, อารมฺมเณ เทฺว ฯเปฯ อวิคเต ปญฺจ. (สํขิตฺตํ.)}

\prose{(สหชาตวาเรปิ ฯเปฯ ปญฺหาวาเรปิ สพฺพตฺถ วิตฺถาเรตพฺโพ.)}

\proseitem{43}{อรณํ มหคฺคตํ ธมฺมํ ปฏิจฺจ อรโณ มหคฺคโต ธมฺโม อุปฺปชฺชติ เหตุปจฺจยา. (สํขิตฺตํ.)}

\prose{เหตุยา เอกํ, อารมฺมเณ เอกํ ฯเปฯ อวิคเต เอกํ. (สํขิตฺตํ.)}

\prose{(สหชาตวาเรปิ ฯเปฯ ปญฺหาวาเรปิ สพฺพตฺถ เอกํ.)}

\proseitem{44}{อรณํ อปฺปมาณํ ธมฺมํ ปฏิจฺจ อรโณ อปฺปมาโณ ธมฺโม อุปฺปชฺชติ เหตุปจฺจยา. (สํขิตฺตํ.)}

\prose{เหตุยา เอกํ, อารมฺมเณ เอกํ ฯเปฯ อวิคเต เอกํ. (สํขิตฺตํ.)}

\prose{(สหชาตวาเรปิ ฯเปฯ ปญฺหาวาเรปิ สพฺพตฺถ เอกํ.)}

\csromanmarkh{100. สรณทุก}\csromanmarkhii{13. ปริตฺตา\-รมฺมณ\-ตฺติก}\csromanheadernomark{2}{100. สรณทุก 13. ปริตฺตา\-รมฺมณ\-ตฺติก}
\csromanheader{2}{1. \textbf{ปฏิจฺจวาร ปจฺจยจตุกฺก}}
\proseitem{45}{สรณํ ปริตฺตา\-รมฺมณํ ธมฺมํ ปฏิจฺจ สรโณ ปริตฺตารมฺมโณ ธมฺโม อุปฺปชฺชติ เหตุปจฺจยา. (1)}

\prose{อรณํ ปริตฺตา\-รมฺมณํ ธมฺมํ ปฏิจฺจ อรโณ ปริตฺตารมฺมโณ ธมฺโม อุปฺปชฺชติ เหตุปจฺจยา. (1) (สํขิตฺตํ.)}

\prose{เหตุยา เทฺว, อารมฺมเณ เทฺว ฯเปฯ อวิคเต เทฺว. (สํขิตฺตํ.)}

\prose{(สหชาตวาโรปิ ฯเปฯ ปญฺหาวาโรปิ สพฺพตฺถ วิตฺถาเรตพฺพา.)}

\proseitem{46}{สรณํ มหคฺคตา\-รมฺมณํ ธมฺมํ ปฏิจฺจ สรโณ มหคฺคตา\-รมฺมโณ ธมฺโม อุปฺปชฺชติ เหตุปจฺจยา. (1)}

\prose{อรณํ มหคฺคตา\-รมฺมณํ ธมฺมํ ปฏิจฺจ อรโณ มหคฺคตา\-รมฺมโณ ธมฺโม อุปฺปชฺชติ เหตุปจฺจยา. (1) (สํขิตฺตํ.)}

\csromanmarkh{100. สรณทุก}\csromanmarkhii{15. มิจฺฉตฺตตฺติก}\csromanheadernomark{2}{100. สรณทุก 15. มิจฺฉตฺตตฺติก}
\prose{\csromanfolio{459}เหตุยา เทฺว, อารมฺมเณ เทฺว ฯเปฯ อวิคเต เทฺว. (สํขิตฺตํ.)}

\prose{(สหชาตวาโรปิ ฯเปฯ ปญฺหาวาโรปิ วิตฺถาเรตพฺพา.)}

\proseitem{47}{อรณํ อปฺปมาณา\-รมฺมณํ ธมฺมํ ปฏิจฺจ อรโณ อปฺปมาณารมฺมโณ ธมฺโม อุปฺปชฺชติ เหตุปจฺจยา. (สํขิตฺตํ.)}

\prose{เหตุยา เอกํ, อารมฺมเณ เอกํ ฯเปฯ อวิคเต เอกํ. (สํขิตฺตํ.)}

\prose{(สหชาตวาเรปิ ฯเปฯ ปญฺหาวาเรปิ สพฺพตฺถ เอกํ.)}

\csromanmarkh{100. สรณทุก}\csromanmarkhii{14. หีนตฺติก}\csromanheadernomark{2}{100. \textbf{สรณทุก 1}4. หีนตฺติก}
\csromanheader{2}{1. \textbf{ปฏิจฺจวาร ปจฺจยจตุกฺก}}
\proseitem{48}{สรณํ หีนํ ธมฺมํ ปฏิจฺจ สรโณ หีโน ธมฺโม อุปฺปชฺชติ เหตุปจฺจยา.}

\prose{เหตุยา เอกํ, อารมฺมเณ เอกํ ฯเปฯ อวิคเต เอกํ. (สํขิตฺตํ.)}

\prose{(สหชาตวาเรปิ ฯเปฯ ปญฺหาวาเรปิ สพฺพตฺถ เอกํ.)}

\proseitem{49}{อรณํ มชฺฌิมํ ธมฺมํ ปฏิจฺจ อรโณ มชฺฌิโม ธมฺโม อุปฺปชฺชติ เหตุปจฺจยา.}

\prose{เหตุยา เอกํ, อารมฺมเณ เอกํ ฯเปฯ อวิคเต เอกํ. (สํขิตฺตํ.)}

\prose{(สหชาตวาเรปิ ฯเปฯ ปญฺหาวาเรปิ สพฺพตฺถ เอกํ.)}

\proseitem{50}{อรณํ ปณีตํ ธมฺมํ ปฏิจฺจ อรโณ ปณีโต ธมฺโม อุปฺปชฺชติ เหตุปจฺจยา. (สํขิตฺตํ.)}

\prose{เหตุยา เอกํ, อารมฺมเณ เอกํ ฯเปฯ อวิคเต เอกํ. (สํขิตฺตํ.)}

\prose{(สหชาตวาเรปิ ฯเปฯ ปญฺหาวาเรปิ สพฺพตฺถ เอกํ.)}

\csromanmarkh{100. สรณทุก}\csromanmarkhii{15. มิจฺฉตฺตตฺติก}\csromanheadernomark{2}{100. สรณทุก 15. มิจฺฉตฺตตฺติก}
\csromanheader{2}{1. \textbf{ปฏิจฺจวาร ปจฺจยจตุกฺก}}
\proseitem{51}{สรณํ มิจฺฉตฺต\-นิยตํ ธมฺมํ ปฏิจฺจ สรโณ มิจฺฉตฺต\-นิยโต ธมฺโม อุปฺปชฺชติ เหตุปจฺจยา. (สํขิตฺตํ.)}

\prose{\csromanfolio{460}เหตุยา เอกํ, อารมฺมเณ เอกํ ฯเปฯ อวิคเต เอกํ. (สํขิตฺตํ.)}

\prose{(สหชาตวาเรปิ ฯเปฯ ปญฺหาวาเรปิ สพฺพตฺถ เอกํ.)}

\proseitem{52}{อรณํ สมฺมตฺต\-นิยตํ ธมฺมํ ปฏิจฺจ อรโณ สมฺมตฺตนิยโต ธมฺโม อุปฺปชฺชติ เหตุปจฺจยา. (สํขิตฺตํ.)}

\prose{เหตุยา เอกํ, อารมฺมเณ เอกํ ฯเปฯ อวิคเต เอกํ. (สํขิตฺตํ.)}

\prose{(สหชาตวาเรปิ ฯเปฯ ปญฺหาวาเรปิ สพฺพตฺถ เอกํ.)}

\proseitem{53}{สรณํ อนิยตํ ธมฺมํ ปฏิจฺจ สรโณ อนิยโต ธมฺโม อุปฺปชฺชติ เหตุปจฺจยา. ตีณิ.}

\prose{อรณํ อนิยตํ ธมฺมํ ปฏิจฺจ อรโณ อนิยโต ธมฺโม อุปฺปชฺชติ เหตุปจฺจยา. (1)}

\prose{สรณํ อนิยตญฺจ อรณํ อนิยตญฺจ ธมฺมํ ปฏิจฺจ อรโณ อนิยโต ธมฺโม อุปฺปชฺชติ เหตุปจฺจยา. (1) (สํขิตฺตํ.)}

\proseitem{54}{เหตุยา ปญฺจ, อารมฺมเณ เทฺว ฯเปฯ อวิคเต ปญฺจ. (สํขิตฺตํ.)}

\prose{(สหชาตวาโรปิ ฯเปฯ ปญฺหาวาโรปิ สพฺพตฺถ วิตฺถาเรตพฺพา.)}

\csromanmarkh{100. สรณทุก}\csromanmarkhii{16. มคฺคา\-รมฺมณ\-ตฺติก}\csromanheadernomark{2}{100. สรณทุก 16. มคฺคา\-รมฺมณ\-ตฺติก}
\csromanheader{2}{1. \textbf{ปฏิจฺจวาร ปจฺจยจตุกฺก}}
\proseitem{55}{อรณํ มคฺคารมฺมณํ ธมฺมํ ปฏิจฺจ อรโณ มคฺคารมฺมโณ ธมฺโม อุปฺปชฺชติ เหตุปจฺจยา. (สํขิตฺตํ.)}

\prose{เหตุยา เอกํ, อารมฺมเณ เอกํ ฯเปฯ อวิคเต เอกํ. (สํขิตฺตํ.)}

\prose{(สหชาตวาเรปิ ฯเปฯ ปญฺหาวาเรปิ สพฺพตฺถ เอกํ.)}

\proseitem{56}{อรณํ มคฺคเหตุกํ ธมฺมํ ปฏิจฺจ อรโณ มคฺคเหตุโก ธมฺโม อุปฺปชฺชติ เหตุปจฺจยา. (สํขิตฺตํ.)}

\prose{เหตุยา เอกํ, อารมฺมเณ เอกํ ฯเปฯ อวิคเต เอกํ. (สํขิตฺตํ.)}

\prose{(สหชาตวาเรปิ ฯเปฯ ปญฺหาวาเรปิ สพฺพตฺถ เอกํ.)}

\csromanmarkh{100. สรณทุก}\csromanmarkhii{18. อตีตตฺติก}\csromanheadernomark{2}{100. สรณทุก 18. อตีตตฺติก}
\proseitem{57}{\csromanfolio{461}อรณํ มคฺคาธิปติํ ธมฺมํ ปฏิจฺจ อรโณ มคฺคาธิปติ ธมฺโม อุปฺปชฺชติ เหตุปจฺจยา. (สํขิตฺตํ.)}

\prose{เหตุยา เอกํ, อารมฺมเณ เอกํ ฯเปฯ อวิคเต เอกํ. (สํขิตฺตํ.)}

\prose{(สหชาตวาเรปิ ฯเปฯ ปญฺหาวาเรปิ สพฺพตฺถ เอกํ.)}

\csromanmarkh{100. สรณทุก}\csromanmarkhii{17. อุปฺปนฺนตฺติก}\csromanheadernomark{2}{100. \textbf{สรณทุก 1}7. อุปฺปนฺนตฺติก}
\csromanheader{2}{7. \textbf{ปญฺหาวาร ปจฺจยจตุกฺก}}
\proseitem{58}{สรโณ อุปฺปนฺโน ธมฺโม สรณสฺส อุปฺปนฺนสฺส ธมฺมสฺส เหตุปจฺจเยน ปจฺจโย.  สรโณ อุปฺปนฺโน ธมฺโม อรณสฺส อุปฺปนฺนสฺส ธมฺมสฺส เหตุปจฺจเยน ปจฺจโย.  สรโณ อุปฺปนฺโน ธมฺโม สรณสฺส อุปฺปนฺนสฺส จ อรณสฺส อุปฺปนฺนสฺส จ ธมฺมสฺส เหตุปจฺจเยน ปจฺจโย. (3)}

\prose{อรโณ อุปฺปนฺโน ธมฺโม อรณสฺส อุปฺปนฺนสฺส ธมฺมสฺส เหตุปจฺจเยน ปจฺจโย. (1)}

\prose{อรโณ อุปฺปนฺโน ธมฺโม อรณสฺส อุปฺปนฺนสฺส ธมฺมสฺส อารมฺมณ\-ปจฺจเยน ปจฺจโย. (สํขิตฺตํ.)}

\csromanhangingbegin
\hangingheaditem{59}{เหตุยา จตฺตาริ, อารมฺมเณ เทฺว ฯเปฯ อวิคเต สตฺต. (สํขิตฺตํ.)}
\hangingbody{(ยถา กุสลตฺติเก ปญฺหาวารํ, เอวํ วิตฺถาเรตพฺพํ.)}
\csromanhangingend

\csromanmarkh{100. สรณทุก}\csromanmarkhii{18. อตีตตฺติก}\csromanheadernomark{2}{100. สรณทุก 18. อตีตตฺติก}
\csromanheader{2}{7. \textbf{ปญฺหาวาร ปจฺจยจตุกฺก}}
\proseitem{60}{สรโณ ปจฺจุปฺปนฺโน ธมฺโม สรณสฺส ปจฺจุปฺปนฺนสฺส ธมฺมสฺส เหตุปจฺจเยน ปจฺจโย. ตีณิ.}

\prose{\csromanfolio{462}อรโณ ปจฺจุปฺปนฺโน ธมฺโม อรณสฺส ปจฺจุปฺปนฺนสฺส ธมฺมสฺส เหตุปจฺจเยน ปจฺจโย. (1) (สํขิตฺตํ.)}

\prose{เหตุยา จตฺตาริ, อารมฺมเณ เทฺว ฯเปฯ อวิคเต สตฺต. (สํขิตฺตํ.)}

\prose{(ยถา กุสลตฺติเก ปญฺหาวารํ, เอวํ วิตฺถาเรตพฺพํ.)}

\csromanmarkh{100. สรณทุก}\csromanmarkhii{19. อตีตา\-รมฺมณ\-ตฺติก}\csromanheadernomark{2}{100. สรณทุก 19. อตีตา\-รมฺมณ\-ตฺติก}
\csromanheader{2}{1. \textbf{ปฏิจฺจวาร ปจฺจยจตุกฺก}}
\proseitem{61}{สรณํ อตีตารมฺมณํ ธมฺมํ ปฏิจฺจ สรโณ อตีตารมฺมโณ ธมฺโม อุปฺปชฺชติ เหตุปจฺจยา. (1)}

\prose{อรณํ อตีตารมฺมณํ ธมฺมํ ปฏิจฺจ อรโณ อตีตารมฺมโณ ธมฺโม อุปฺปชฺชติ เหตุปจฺจยา. (1) (สํขิตฺตํ.)}

\prose{เหตุยา เทฺว, อารมฺมเณ เทฺว ฯเปฯ วิปาเก เอกํ ฯเปฯ อวิคเต เทฺว. (สํขิตฺตํ. สหชาตวาเรปิ ฯเปฯ ปญฺหาวาเรปิ สพฺพตฺถ วิตฺถาเรตพฺโพ.)}

\proseitem{62}{สรณํ อนาคตา\-รมฺมณํ ธมฺมํ ปฏิจฺจ สรโณ อนาคตารมฺมโณ ธมฺโม อุปฺปชฺชติ เหตุปจฺจยา. (1)}

\prose{อรณํ อนาคตา\-รมฺมณํ ธมฺมํ ปฏิจฺจ อรโณ อนาคตารมฺมโณ ธมฺโม อุปฺปชฺชติ เหตุปจฺจยา. (1) (สํขิตฺตํ.)}

\prose{เหตุยา เทฺว, อารมฺมเณ เทฺว ฯเปฯ วิปาเก เอกํ ฯเปฯ อวิคเต เทฺว. (สํขิตฺตํ. สหชาตวาเรปิ ฯเปฯ ปญฺหาวาเรปิ สพฺพตฺถ วิตฺถาเรตพฺโพ.)}

\proseitem{63}{สรณํ ปจฺจุปฺปนฺนา\-รมฺมณํ ธมฺมํ ปฏิจฺจ สรโณ ปจฺจุปฺปนฺนา\-รมฺมโณ ธมฺโม อุปฺปชฺชติ เหตุปจฺจยา. (1)}

\prose{อรณํ ปจฺจุปฺปนฺนา\-รมฺมณํ ธมฺมํ ปฏิจฺจ อรโณ ปจฺจุปฺปนฺนา\-รมฺมโณ ธมฺโม อุปฺปชฺชติ เหตุปจฺจยา. (1) (สํขิตฺตํ.)}

\prose{เหตุยา เทฺว, อารมฺมเณ เทฺว ฯเปฯ วิปาเก เอกํ ฯเปฯ อวิคเต เทฺว. (สํขิตฺตํ. สหชาตวาเรปิ ฯเปฯ ปญฺหาวาเรปิ สพฺพตฺถ วิตฺถาเรตพฺโพ.)}

\vspace{\tipitakaparskip}
\csromancenter{21. อชฺฌตฺตา\-รมฺมณ\-ตฺติก 100. สรณทุก 20. อชฺฌตฺตตฺติก}
\csromanheader{2}{1. \textbf{ปฏิจฺจวาร ปจฺจยจตุกฺก}}
\proseitem{64}{\csromanfolio{463}สรณํ อชฺฌตฺตํ ธมฺมํ ปฏิจฺจ สรโณ อชฺฌตฺโต ธมฺโม อุปฺปชฺชติ เหตุปจฺจยา. ตีณิ.}

\prose{อรณํ อชฺฌตฺตํ ธมฺมํ ปฏิจฺจ อรโณ อชฺฌตฺโต ธมฺโม อุปฺปชฺชติ เหตุปจฺจยา. (1)}

\prose{สรณํ อชฺฌตฺตญฺจ อรณํ อชฺฌตฺตญฺจ ธมฺมํ ปฏิจฺจ อรโณ อชฺฌตฺโต ธมฺโม อุปฺปชฺชติ เหตุปจฺจยา. (1) (สํขิตฺตํ.)}

\proseitem{65}{เหตุยา ปญฺจ, อรมฺมเณ เทฺว ฯเปฯ อวิคเต ปญฺจ. (สํขิตฺตํ.)}

\prose{(สหชาตวาเรปิ ฯเปฯ ปญฺหาวาเรปิ สพฺพตฺถ วิตฺถาเรตพฺโพ.)}

\proseitem{66}{สรณํ พหิทฺธา ธมฺมํ ปฏิจฺจ สรโณ พหิทฺธา ธมฺโม อุปฺปชฺชติ เหตุปจฺจยา. ตีณิ.}

\prose{อรณํ พหิทฺธา ธมฺมํ ปฏิจฺจ อรโณ พหิทฺธา ธมฺโม อุปฺปชฺชติ เหตุปจฺจยา.}

\prose{สรณํ พหิทฺธา จ อรณํ หิทฺธา จ ธมฺมํ ปฏิจฺจ อรโณ พหิทฺธา ธมฺโม อุปฺปชฺชติ เหตุปจฺจยา. (1) (สํขิตฺตํ.)}

\prose{เหตุยา ปญฺจ, อารมฺมเณ เทฺว ฯเปฯ อวิคเต ปญฺจ. (สํขิตฺตํ.)}

\prose{(สหชาตวาเรปิ ฯเปฯ ปญฺหาวาเรปิ สพฺพตฺถ วิตฺถาเรตพฺโพ.)}

\csromanmarkh{100. สรณทุก}\csromanmarkhii{21. อชฺฌตฺตา\-รมฺมณ\-ตฺติก}\csromanheadernomark{2}{100. สรณทุก 21. อชฺฌตฺตา\-รมฺมณ\-ตฺติก}
\csromanheader{2}{1. \textbf{ปฏิจฺจวาร ปจฺจยจตุกฺก}}
\proseitem{67}{สรณํ อชฺฌตฺตา\-รมฺมณํ ธมฺมํ ปฏิจฺจ สรโณ อชฺฌตฺตา\-รมฺมโณ ธมฺโม อุปฺปชฺชติ เหตุปจฺจยา. (1)}

\prose{อรณํ อชฺฌตฺตา\-รมฺมณํ ธมฺมํ ปฏิจฺจ อรโณ อชฺฌตฺตา\-รมฺมโณ ธมฺโม อุปฺปชฺชติ เหตุปจฺจยา. (1) (สํขิตฺตํ.)}

\prose{\csromanfolio{464}เหตุยา เทฺว, อารมฺมเณ เทฺว ฯเปฯ อวิคเต เทฺว. (สํขิตฺตํ.)}

\prose{(สหชาตวาเรปิ ฯเปฯ ปญฺหาวาเรปิ สพฺพตฺถ วิตฺถาเรตพฺโพ.)}

\proseitem{68}{สรณํ พหิทฺธา\-รมฺมณํ ธมฺมํ ปฏิจฺจ สรโณ พหิทฺธา\-รมฺมโณ ธมฺโม อุปฺปชฺชติ เหตุปจฺจยา. (1)}

\prose{อรณํ พหิทฺธา\-รมฺมณํ ธมฺมํ ปฏิจฺจ อรโณ พหิทฺธา\-รมฺมโณ ธมฺโม อุปฺปชฺชติ เหตุปจฺจยา. (1) (สํขิตฺตํ.)}

\prose{เหตุยา เทฺว, อารมฺมเณ เทฺว ฯเปฯ อวิคเต เทฺว. (สํขิตฺตํ.)}

\prose{(สหชาตวาเรปิ ฯเปฯ ปญฺหาวาเรปิ สพฺพตฺถ วิตฺถาเรตพฺโพ.)}

\csromanmatikaanchor{mtk.90}
\csromantocmarkref{subparagraph}{22. สนิทสฺสนตฺติก}{mtk.90}
\bookmark[dest={mtk.90},level=5]{22. สนิทสฺสนตฺติก}
\csromanmarkh{100. สรณทุก}\csromanmarkhii{22. สนิทสฺสน\-ตฺติก}\csromanheadernomark{6}{100. สรณทุก 22. สนิทสฺสน\-ตฺติก}
\csromanheader{2}{1-7. ปฏิจฺจาทิวาร ปจฺจยจตุกฺก}
\proseitem{69}{อรณํ อนิทสฺสน\-สปฺปฏิฆํ ธมฺมํ ปฏิจฺจ อรโณ อนิทสฺสน\-สปฺปฏิโฆ ธมฺโม อุปฺปชฺชติ เหตุปจฺจยา. (สํขิตฺตํ.)}

\prose{เหตุยา เอกํ, อธิปติยา เอกํ ฯเปฯ อวิคเต เอกํ. (สํขิตฺตํ.)}

\prose{(สหชาตวาเรปิ ฯเปฯ ปญฺหาวาเรปิ สพฺพตฺถ เอกํ.)}

\proseitem{70}{สรณํ อนิทสฺสน\-อปฺปฏิฆํ ธมฺมํ ปฏิจฺจ สรโณ อนิทสฺสน\-อปฺปฏิโฆ ธมฺโม อุปฺปชฺชติ เหตุปจฺจยา. ตีณิ.}

\prose{อรณํ อนิทสฺสน\-อปฺปฏิฆํ ธมฺมํ ปฏิจฺจ อรโณ อนิทสฺสน\-อปฺปฏิโฆ ธมฺโม อุปฺปชฺชติ เหตุปจฺจยา. (1)}

\prose{สรณํ อนิทสฺสนอปฺปฏิฆญฺจ อรณํ อนิทสฺสนอปฺปฏิฆญฺจ ธมฺมํ ปฏิจฺจ อรโณ อนิทสฺสน\-อปฺปฏิโฆ ธมฺโม อุปฺปชฺชติ เหตุปจฺจยา. (1) (สํขิตฺตํ.)}

\prose{เหตุยา ปญฺจ, อารมฺมเณ เทฺว ฯเปฯ อวิคเต ปญฺจ. (สํขิตฺตํ.)}

\prose{(สหชาตวาโรปิ ฯเปฯ สมฺปยุตฺตวาโรปิ วิตฺถาเรตพฺพา.)}

\csromanmarkh{100. สรณทุก}\csromanmarkhii{22. สนิทสฺสน\-ตฺติก}\csromanheadernomark{2}{100. สรณทุก 22. สนิทสฺสน\-ตฺติก}
\proseitem{71}{\csromanfolio{465}อรโณ อนิทสฺสน\-สปฺปฏิโฆ ธมฺโม อรณสฺส อนิทสฺสน\-สปฺปฏิฆสฺส ธมฺมสฺส สหชาต\-ปจฺจเยน ปจฺจโย. เอกํ. (อญฺญมญฺเญ เอกํ, นิสฺสเย เอกํ, อตฺถิยา เอกํ, อวิคเต เอกํ. สํขิตฺตํ.)}

\proseitem{72}{สรโณ อนิทสฺสน\-อปฺปฏิโฆ ธมฺโม สรณสฺส อนิทสฺสน\-อปฺปฏิฆสฺส ธมฺมสฺส เหตุปจฺจเยน ปจฺจโย. ตีณิ.}

\prose{อรโณ อนิทสฺสน\-อปฺปฏิโฆ ธมฺโม อรณสฺส อนิทสฺสน\-อปฺปฏิฆสฺส ธมฺมสฺส เหตุปจฺจเยน ปจฺจโย. (1) (สํขิตฺตํ.)}

\proseitem{73}{เหตุยา จตฺตาริ, อารมฺมเณ จตฺตาริ ฯเปฯ อวิคเต สตฺต. (สํขิตฺตํ.)}

\prose{นเหตุยา สตฺต, นอารมฺมเณ สตฺต. (สํขิตฺตํ.)}

\prose{เหตุปจฺจยา นอารมฺมเณ จตฺตาริ. (สํขิตฺตํ.)}

\prose{นเหตุยา อารมฺมเณ จตฺตาริ. (สํขิตฺตํ.)}

\prosewithcloser{(ยถา กุสลตฺติเก ปญฺหาวารสฺส อนุโลมมฺปิ ปจฺจนียมฺปิ อนุโลมปจฺจนียมฺปิ ปจฺจนียานุโลมมฺปิ คณิตํ, เอวํ คเณตพฺพํ.)}{ธมฺมานุโลเม ทุกติก\-ปฏฺฐานํ นิฏฺฐิตํ.}
\csromanpalirecto
\pitaka{อภิธมฺม\-ปิฏก}
\csromanheader{2}{\textbf{ธมฺมานุโลม}}
\csromanmatikaanchor{mtk.91}
\csromantocmarknopage{chapter}{ติกทุกปฏฺฐานปาฬิ}{mtk.91}
\bookmark[dest={mtk.91},level=0]{ติกทุกปฏฺฐานปาฬิ}
\csromantocmarknopage{chapter}{1. กุสลตฺติก}{mtk.92}
\bookmark[dest={mtk.92},level=0]{1. กุสลตฺติก}
\gambhira{ติกทุก\-ปฏฺฐาน\-ปาฬิ}
\namakkaram{นโม ตสฺส ภควโต อรหโต สมฺมา\-สมฺพุทฺธสฺส.}
\csromanmatikaanchor{mtk.93}
\csromantocmarkref{section}{1. เหตุทุก}{mtk.93}
\bookmark[dest={mtk.93},level=1]{1. เหตุทุก}
\csromanmarkh{1. กุสลตฺติก}\csromanmarkhii{1. เหตุทุก}\csromanheadernomark{1}{1. \textbf{กุสลตฺติก 1. เหตุทุก}}
\csromanmarkh{1. เหตุปท}\csromanmarkhii{1. ปฏิจฺจวาร}\csromanheadernomark{2}{1. \textbf{เหตุปท 1. ปฏิจฺจวาร}}
\prose{\csromanfolio{467}ปจฺจยจตุกฺกเหตุ อารมฺมณ}

\proseitem{1}{กุสลํ เหตุํ ธมฺมํ ปฏิจฺจ กุสโล เหตุ ธมฺโม อุปฺปชฺชติ เหตุปจฺจยา. (1)}

\prose{อกุสลํ เหตุํ ธมฺมํ ปฏิจฺจ อกุสโล เหตุ ธมฺโม อุปฺปชฺชติ เหตุปจฺจยา. (1)}

\prose{อพฺยากตํ เหตุํ ธมฺมํ ปฏิจฺจ อพฺยากโต เหตุ ธมฺโม อุปฺปชฺชติ เหตุปจฺจยา. (1)}

\proseitem{2}{กุสลํ เหตุํ ธมฺมํ ปฏิจฺจ กุสโล เหตุ ธมฺโม อุปฺปชฺชติ อารมฺมณ\-ปจฺจยา. (1)}

\prose{อกุสลํ เหตุํ ธมฺมํ ปฏิจฺจ อกุสโล เหตุ ธมฺโม อุปฺปชฺชติ อารมฺมณ\-ปจฺจยา. (1)}

\prose{อพฺยากตํ เหตุํ ธมฺมํ ปฏิจฺจ อพฺยากโต เหตุ ธมฺโม อุปฺปชฺชติ อารมฺมณ\-ปจฺจยา. (1) (สํขิตฺตํ.)}

\proseitem{3}{เหตุยา ตีณิ, อารมฺมเณ ตีณิ, อธิปติยา ตีณิ, อนนฺตเร ตีณิ, สมนนฺตเร ตีณิ, สหชาเต ตีณิ, อญฺญมญฺเญ ตีณิ, นิสฺสเย \csromanfolio{468}ตีณิ, อุปนิสฺสเย ตีณิ, ปุเรชาเต ตีณิ, อาเสวเน ตีณิ, กมฺเม ตีณิ, วิปาเก เอกํ, อาหาเร ตีณิ ฯเปฯ อวิคเต ตีณิ. (สํขิตฺตํ.)}

\prose{ปจฺจนียนอธิปติ}

\proseitem{4}{กุสลํ เหตุํ ธมฺมํ ปฏิจฺจ กุสโล เหตุ ธมฺโม อุปฺปชฺชติ นอธิปติ\-ปจฺจยา. (1)}

\prose{อกุสลํ เหตุํ ธมฺมํ ปฏิจฺจ อกุสโล เหตุ ธมฺโม อุปฺปชฺชติ นอธิปติ\-ปจฺจยา. (1)}

\prose{อพฺยากตํ เหตุํ ธมฺมํ ปฏิจฺจ อพฺยากโต เหตุ ธมฺโม อุปฺปชฺชติ นอธิปติ\-ปจฺจยา. (1) (สํขิตฺตํ.)}

\proseitem{5}{นอธิปติยา ตีณิ, นปุเรชาเต ตีณิ, นปจฺฉาชาเต ตีณิ, นอาเสวเน ตีณิ, นวิปาเก ตีณิ, นวิปฺปยุตฺเต ตีณิ. (สํขิตฺตํ.)}

\prose{เหตุปจฺจยา นอธิปติยา ตีณิ. (สํขิตฺตํ.)}

\prose{นอธิปติ\-ปจฺจยา เหตุยา ตีณิ. (สํขิตฺตํ.)}

\prose{(สหชาตวาเรปิ ปจฺจยวาเรปิ นิสฺสยวาเรปิ สํสฏฺฐ\-วาเรปิ สมฺปยุตฺตวาเรปิ สพฺพตฺถ ตีณิ.)}

\csromanheader{2}{7. \textbf{ปญฺหาวาร ปจฺจยจตุกฺก}}
\prose{เหตุอารมฺมณาทิ}

\proseitem{3}{กุสโล เหตุ ธมฺโม กุสลสฺส เหตุสฺส ธมฺมสฺส เหตุปจฺจเยน ปจฺจโย. (1)}

\prose{อกุสโล เหตุ ธมฺโม อกุสลสฺส เหตุสฺส ธมฺมสฺส เหตุปจฺจเยน ปจฺจโย. (1)}

\prose{อพฺยากโต เหตุ ธมฺโม อพฺยากตสฺส เหตุสฺส ธมฺมสฺส เหตุปจฺจเยน ปจฺจโย. (1)}

\proseitem{7}{กุสโล เหตุ ธมฺโม กุสลสฺส เหตุสฺส ธมฺมสฺส อารมฺมณ\-ปจฺจเยน ปจฺจโย.  กุสโล เหตุ ธมฺโม อกุสลสฺส เหตุสฺส}

\vspace{\tipitakaparskip}
\csromancenter{กุสลตฺติก 1. เหตุทุก ธมฺมสฺส อารมฺมณ\-ปจฺจเยน ปจฺจโย.  กุสโล เหตุ ธมฺโม อพฺยากตสฺส เหตุสฺส ธมฺมสฺส อารมฺมณ\-ปจฺจเยน ปจฺจโย. (3)}
\prose{\csromanfolio{469}อกุสโล เหตุ ธมฺโม อกุสลสฺส เหตุสฺส ธมฺมสฺส อารมฺมณ\-ปจฺจเยน ปจฺจโย. ตีณิ.}

\prose{อพฺยากโต เหตุ ธมฺโม อพฺยากตสฺส เหตุสฺส ธมฺมสฺส อารมฺมณ\-ปจฺจเยน ปจฺจโย. ตีณิ.}

\proseitem{8}{กุสโล เหตุ ธมฺโม กุสลสฺส เหตุสฺส ธมฺมสฺส อธิปติ\-ปจฺจเยน ปจฺจโย.  กุสโล เหตุ ธมฺโม อกุสลสฺส เหตุสฺส ธมฺมสฺส อธิปติ\-ปจฺจเยน ปจฺจโย.  กุสโล เหตุ ธมฺโม อพฺยากตสฺส เหตุสฺส ธมฺมสฺส อธิปติ\-ปจฺจเยน ปจฺจโย. ตีณิ.}

\prose{อกุสโล เหตุ ธมฺโม อกุสลสฺส เหตุสฺส ธมฺมสฺส อธิปติ\-ปจฺจเยน ปจฺจโย. (1)}

\prose{อพฺยากโต เหตุ ธมฺโม อพฺยากตสฺส เหตุสฺส ธมฺมสฺส อธิปติ\-ปจฺจเยน ปจฺจโย.  อพฺยากโต เหตุ ธมฺโม กุสลสฺส เหตุสฺส ธมฺมสฺส อธิปติ\-ปจฺจเยน ปจฺจโย.  อพฺยากโต เหตุ ธมฺโม อกุสลสฺส เหตุสฺส ธมฺมสฺส อธิปติ\-ปจฺจเยน ปจฺจโย. ตีณิ.}

\proseitem{9}{กุสโล เหตุ ธมฺโม กุสลสฺส เหตุสฺส ธมฺมสฺส อนนฺตร\-ปจฺจเยน ปจฺจโย. เทฺว.}

\prose{อกุสโล เหตุ ธมฺโม อกุสลสฺส เหตุสฺส ธมฺมสฺส อนนฺตร\-ปจฺจเยน ปจฺจโย. เทฺว.}

\prose{อพฺยากโต เหตุ ธมฺโม อพฺยากตสฺส เหตุสฺส ธมฺมสฺส อนนฺตร\-ปจฺจเยน ปจฺจโย. (1)}

\proseitem{10}{กุสโล เหตุ ธมฺโม กุสลสฺส เหตุสฺส ธมฺมสฺส สหชาต\-ปจฺจเยน ปจฺจโย. (1)}

\prose{อกุสโล เหตุ ธมฺโม อกุสลสฺส เหตุสฺส ธมฺมสฺส สหชาต\-ปจฺจเยน ปจฺจโย. (1)}

\csromanpalirecto
\gambhira{ธมฺมานุโลม ติกทุก\-ปฏฺฐาน\-ปาฬิ}
\prose{\csromanfolio{470}(อญฺญมญฺญ\-นิสฺสยปจฺจยา สหชาตปจฺจย\-สทิสา.)}

\csromanheader{2}{\textbf{อุปนิสฺสยาทิ}}
\proseitem{11}{กุสโล เหตุ ธมฺโม กุสลสฺส เหตุสฺส ธมฺมสฺส อุปนิสฺสย\-ปจฺจเยน ปจฺจโย.  กุสโล เหตุ ธมฺโม อกุสลสฺส เหตุสฺส ธมฺมสฺส อุปนิสฺสย\-ปจฺจเยน ปจฺจโย.  กุสโล เหตุ ธมฺโม อพฺยากตสฺส เหตุสฺส ธมฺมสฺส อุปนิสฺสย\-ปจฺจเยน ปจฺจโย. (3)}

\prose{อกุสโล เหตุ ธมฺโม อกุสลสฺส เหตุสฺส ธมฺมสฺส อุปนิสฺสย\-ปจฺจเยน ปจฺจโย.  อกุสโล เหตุ ธมฺโม กุสลสฺส เหตุสฺส ธมฺมสฺส อุปนิสฺสย\-ปจฺจเยน ปจฺจโย.  อกุสโล เหตุ ธมฺโม อพฺยากตสฺส เหตุสฺส ธมฺมสฺส อุปนิสฺสย\-ปจฺจเยน ปจฺจโย. ตีณิ.}

\prose{อพฺยากโต เหตุ ธมฺโม อพฺยากตสฺส เหตุสฺส ธมฺมสฺส อุปนิสฺสย\-ปจฺจเยน ปจฺจโย.  อพฺยากโต เหตุ ธมฺโม กุสลสฺส เหตุสฺส ธมฺมสฺส อุปนิสฺสย\-ปจฺจเยน ปจฺจโย.  อพฺยากโต เหตุ ธมฺโม อกุสลสฺส เหตุสฺส ธมฺมสฺส อุปนิสฺสย\-ปจฺจเยน ปจฺจโย. (3)}

\prose{อพฺยากโต เหตุ ธมฺโม อพฺยากตสฺส เหตุสฺส ธมฺมสฺส วิปาก\-ปจฺจเยน ปจฺจโย. (1) (สํขิตฺตํ.)}

\proseitem{12}{เหตุยา ตีณิ, อารมฺมเณ นว, อธิปติยา สตฺต, อนนฺตเร ปญฺจ, สมนนฺตเร ปญฺจ, สหชาเต ตีณิ, อญฺญมญฺเญ ตีณิ, นิสฺสเย ตีณิ, อุปนิสฺสเย นว, อาเสวเน ตีณิ, วิปาเก เอกํ, อินฺทฺริเย เทฺว, มคฺเค เทฺว, สมฺปยุตฺเต ตีณิ, อตฺถิยา ตีณิ, นตฺถิยา ปญฺจ, วิคเต ปญฺจ, อวิคเต ตีณิ. (สํขิตฺตํ.)}

\proseitem{13}{กุสโล เหตุ ธมฺโม กุสลสฺส เหตุสฺส ธมฺมสฺส อารมฺมณ\-ปจฺจเยน ปจฺจโย. สหชาต\-ปจฺจเยน ปจฺจโย. อุปนิสฺสย\-ปจฺจเยน ปจฺจโย.  กุสโล เหตุ ธมฺโม อกุสลสฺส เหตุสฺส ธมฺมสฺส อารมฺมณ\-ปจฺจเยน ปจฺจโย. อุปนิสฺสย\-ปจฺจเยน ปจฺจโย.  }

\vspace{\tipitakaparskip}
\csromancenter{กุสลตฺติก 1. เหตุทุก กุสโล เหตุ ธมฺโม อพฺยากตสฺส เหตุสฺส ธมฺมสฺส อารมฺมณ\-ปจฺจเยน ปจฺจโย. อุปนิสฺสย\-ปจฺจเยน ปจฺจโย. (3)}
\prose{\csromanfolio{471}อกุสโล เหตุ ธมฺโม อกุสลสฺส เหตุสฺส ธมฺมสฺส อารมฺมณ\-ปจฺจเยน ปจฺจโย. สหชาต\-ปจฺจเยน ปจฺจโย. อุปนิสฺสย\-ปจฺจเยน ปจฺจโย.  อกุสโล เหตุ ธมฺโม กุสลสฺส เหตุสฺส ธมฺมสฺส อารมฺมณ\-ปจฺจเยน ปจฺจโย. อุปนิสฺสย\-ปจฺจเยน ปจฺจโย.  อกุสโล เหตุ ธมฺโม อพฺยากตสฺส เหตุสฺส ธมฺมสฺส อารมฺมณ\-ปจฺจเยน ปจฺจโย. อุปนิสฺสย\-ปจฺจเยน ปจฺจโย. (3)}

\prose{อพฺยากโต เหตุ ธมฺโม อพฺยากตสฺส เหตุสฺส ธมฺมสฺส อารมฺมณ\-ปจฺจเยน ปจฺจโย. สหชาต\-ปจฺจเยน ปจฺจโย. อุปนิสฺสย\-ปจฺจเยน ปจฺจโย. (สํขิตฺตํ.)}

\proseitem{14}{นเหตุยา นว, นอารมฺมเณ นว. (สํขิตฺตํ.)}

\prose{เหตุปจฺจยา นอารมฺมเณ ตีณิ. (สํขิตฺตํ.)}

\prose{นเหตุปจฺจยา อารมฺมเณ นว. (สํขิตฺตํ.)}

\prose{(ยถา กุสลตฺติเก ปญฺหาวารสฺส อนุโลมมฺปิ ปจฺจนียมฺปิ อนุโลมปจฺจนียมฺปิ ปจฺจนียานุโลมมฺปิ คณิตํ, เอวํ คเณตพฺพํ.)}

\csromanmarkh{2. นเหตุปท}\csromanmarkhii{1. ปฏิจฺจวาร}\csromanheadernomark{2}{2. \textbf{นเหตุปท 1. ปฏิจฺจวาร}}
\prose{ปจฺจยจตุกฺกเหตุ อารมฺมณ}

\proseitem{15}{กุสลํ นเหตุํ ธมฺมํ ปฏิจฺจ กุสโล นเหตุ ธมฺโม อุปฺปชฺชติ เหตุปจฺจยา.  กุสลํ นเหตุํ ธมฺมํ ปฏิจฺจ อพฺยากโต นเหตุ ธมฺโม อุปฺปชฺชติ เหตุปจฺจยา.  กุสลํ นเหตุํ ธมฺมํ ปฏิจฺจ กุสโล นเหตุ จ อพฺยากโต นเหตุ จ ธมฺมา อุปฺปชฺชนฺติ เหตุปจฺจยา. (3)}

\prose{อกุสลํ นเหตุํ ธมฺมํ ปฏิจฺจ อกุสโล นเหตุ ธมฺโม อุปฺปชฺชติ เหตุปจฺจยา.  อกุสลํ นเหตุํ ธมฺมํ ปฏิจฺจ อพฺยากโต นเหตุ ธมฺโม อุปฺปชฺชติ เหตุปจฺจยา.  อกุสลํ นเหตุํ ธมฺมํ ปฏิจฺจ อกุสโล นเหตุ จ อพฺยากโต นเหตุ จ ธมฺมา อุปฺปชฺชนฺติ เหตุปจฺจยา. (3)}

\prose{\csromanfolio{472}อพฺยากตํ นเหตุํ ธมฺมํ ปฏิจฺจ อพฺยากโต นเหตุ ธมฺโม อุปฺปชฺชติ เหตุปจฺจยา. (1)}

\prose{กุสลํ นเหตุญฺจ อพฺยากตํ นเหตุญฺจ ธมฺมํ ปฏิจฺจ อพฺยากโต นเหตุ ธมฺโม อุปฺปชฺชติ เหตุปจฺจยา. (1)}

\prose{อกุสลํ นเหตุญฺจ อพฺยากตํ นเหตุญฺจ ธมฺมํ ปฏิจฺจ อพฺยากโต นเหตุ ธมฺโม อุปฺปชฺชติ เหตุปจฺจยา. (1)}

\proseitem{16}{กุสลํ นเหตุํ ธมฺมํ ปฏิจฺจ กุสโล นเหตุ ธมฺโม อุปฺปชฺชติ อารมฺมณ\-ปจฺจยา. (1)}

\prose{อกุสลํ นเหตุํ ธมฺมํ ปฏิจฺจ อกุสโล นเหตุ ธมฺโม อุปฺปชฺชติ อารมฺมณ\-ปจฺจยา. (1)}

\prose{อพฺยากตํ นเหตุํ ธมฺมํ ปฏิจฺจ อพฺยากโต นเหตุ ธมฺโม อุปฺปชฺชติ อารมฺมณ\-ปจฺจยา. (1) (สํขิตฺตํ.)}

\proseitem{17}{เหตุยา นว, อารมฺมเณ ตีณิ, อธิปติยา นว, อนนฺตเร ตีณิ, สมนนฺตเร ตีณิ, สหชาเต นว, อญฺญมญฺเญ ตีณิ, นิสฺสเย นว, อุปนิสฺสเย ตีณิ, ปุเรชาเต ตีณิ, อาเสวเน ตีณิ, กมฺเม นว, วิปาเก เอกํ, อาหาเร นว, อินฺทฺริเย นว, ฌาเน นว, มคฺเค นว, สมฺปยุตฺเต ตีณิ, วิปฺปยุตฺเต นว, อตฺถิยา นว, นตฺถิยา ตีณิ, วิคเต ตีณิ, อวิคเต นว. (สํขิตฺตํ.)}

\prose{นเหตุ นอารมฺมณ}

\proseitem{18}{อพฺยากตํ นเหตุํ ธมฺมํ ปฏิจฺจ อพฺยากโต นเหตุ ธมฺโม อุปฺปชฺชติ นเหตุปจฺจยา. (1)}

\prose{กุสลํ นเหตุํ ธมฺมํ ปฏิจฺจ อพฺยากโต นเหตุ ธมฺโม อุปฺปชฺชติ นอารมฺมณ\-ปจฺจยา. (1)}

\prose{อกุสลํ นเหตุํ ธมฺมํ ปฏิจฺจ อพฺยากโต นเหตุ ธมฺโม อุปฺปชฺชติ นอารมฺมณ\-ปจฺจยา. (1)}

\prose{อพฺยากตํ นเหตุํ ธมฺมํ ปฏิจฺจ อพฺยากโต นเหตุ ธมฺโม อุปฺปชฺชติ นอารมฺมณ\-ปจฺจยา. (1)}

\csromanmarkh{1. กุสลตฺติก}\csromanmarkhii{1. เหตุทุก}\csromanheadernomark{2}{1. กุสลตฺติก 1. เหตุทุก}
\prose{\csromanfolio{473}กุสลํ นเหตุญฺจ อพฺยากตํ นเหตุญฺจ ธมฺมํ ปฏิจฺจ อพฺยากโต นเหตุ ธมฺโม อุปฺปชฺชติ นอารมฺมณ\-ปจฺจยา. (1)}

\prose{อกุสลํ นเหตุญฺจ อพฺยากตํ นเหตุญฺจ ธมฺมํ ปฏิจฺจ อพฺยากโต นเหตุ ธมฺโม อุปฺปชฺชติ นอารมฺมณ\-ปจฺจยา. (1) (สํขิตฺตํ.)}

\proseitem{19}{นเหตุยา เอกํ, นอารมฺมเณ ปญฺจ, นอธิปติยา นว, นอนนฺตเร ปญฺจ, นสมนนฺตเร ปญฺจ, นอญฺญมญฺเญ ปญฺจ, นอุปนิสฺสเย ปญฺจ, นปุเรชาเต สตฺต, นปจฺฉาชาเต นว, นอาเสวเน นว, นกมฺเม ตีณิ, นวิปาเก นว, นอาหาเร เอกํ, นอินฺทฺริเย เอกํ, นฌาเน เอกํ, นมคฺเค เอกํ, นสมฺปยุตฺเต ปญฺจ, นวิปฺปยุตฺเต ตีณิ, โนนตฺถิยา ปญฺจ, โนวิคเต ปญฺจ. (สํขิตฺตํ.)}

\prose{เหตุปจฺจยา นอารมฺมเณ ปญฺจ. (สํขิตฺตํ.)}

\prose{นเหตุปจฺจยา อารมฺมเณ เอกํ. (สํขิตฺตํ.)}

\prose{(สหชาตวารมฺปิ ปจฺจยวารมฺปิ นิสฺสยวารมฺปิ สํสฏฺฐ\-วารมฺปิ สมฺปยุตฺตวารมฺปิ ปฏิจฺจ\-วาร\-สทิสา วิตฺถาเรตพฺพา.)}

\prose{7. ปญฺหาวาร ปจฺจยจตุกฺกอารมฺมณาทิ}

\proseitem{20}{กุสโล นเหตุ ธมฺโม กุสลสฺส นเหตุสฺส ธมฺมสฺส อารมฺมณ\-ปจฺจเยน ปจฺจโย.  กุสโล นเหตุ ธมฺโม อกุสลสฺส นเหตุสฺส ธมฺมสฺส อารมฺมณ\-ปจฺจเยน ปจฺจโย.  กุสโล นเหตุ ธมฺโม อพฺยากตสฺส นเหตุสฺส ธมฺมสฺส อารมฺมณ\-ปจฺจเยน ปจฺจโย. (3)}

\prose{อกุสโล นเหตุ ธมฺโม อกุสลสฺส นเหตุสฺส ธมฺมสฺส อารมฺมณ\-ปจฺจเยน ปจฺจโย.  อกุสโล นเหตุ ธมฺโม กุสลสฺส นเหตุสฺส ธมฺมสฺส อารมฺมณ\-ปจฺจเยน ปจฺจโย.  อกุสโล นเหตุ ธมฺโม อพฺยากตสฺส นเหตุสฺส ธมฺมสฺส อารมฺมณ\-ปจฺจเยน ปจฺจโย. (3)}

\prose{อพฺยากโต นเหตุ ธมฺโม อพฺยากตสฺส นเหตุสฺส ธมฺมสฺส อารมฺมณ\-ปจฺจเยน ปจฺจโย.  อพฺยากโต นเหตุ ธมฺโม กุสลสฺส นเหตุสฺส ธมฺมสฺส อารมฺมณ\-ปจฺจเยน ปจฺจโย.  อพฺยากโต นเหตุ ธมฺโม อกุสลสฺส นเหตุสฺส ธมฺมสฺส อารมฺมณ\-ปจฺจเยน ปจฺจโย. (3)}

\proseitem{21}{\csromanfolio{474}กุสโล นเหตุ ธมฺโม กุสลสฺส นเหตุสฺส ธมฺมสฺส อธิปติ\-ปจฺจเยน ปจฺจโย.  กุสโล นเหตุ ธมฺโม อกุสลสฺส นเหตุสฺส ธมฺมสฺส อธิปติ\-ปจฺจเยน ปจฺจโย.  กุสโล นเหตุ ธมฺโม อพฺยากตสฺส นเหตุสฺส ธมฺมสฺส อธิปติ\-ปจฺจเยน ปจฺจโย.  กุสโล นเหตุ ธมฺโม กุสลสฺส นเหตุสฺส จ อพฺยากตสฺส นเหตุสฺส จ ธมฺมสฺส อธิปติ\-ปจฺจเยน ปจฺจโย. จตฺตาริ.}

\prose{อกุสโล นเหตุ ธมฺโม อกุสลสฺส นเหตุสฺส ธมฺมสฺส อธิปติ\-ปจฺจเยน ปจฺจโย. ตีณิ.}

\prose{อพฺยากโต นเหตุ ธมฺโม อพฺยากตสฺส นเหตุสฺส ธมฺมสฺส อธิปติ\-ปจฺจเยน ปจฺจโย.  อพฺยากโต นเหตุ ธมฺโม กุสลสฺส นเหตุสฺส ธมฺมสฺส อธิปติ\-ปจฺจเยน ปจฺจโย. ตีณิ.}

\proseitem{22}{กุสโล นเหตุ ธมฺโม กุสลสฺส นเหตุสฺส ธมฺมสฺส อนนฺตร\-ปจฺจเยน ปจฺจโย.  กุสโล นเหตุ ธมฺโม อพฺยากตสฺส นเหตุสฺส ธมฺมสฺส อนนฺตร\-ปจฺจเยน ปจฺจโย. (2)}

\prose{อกุสโล นเหตุ ธมฺโม อกุสลสฺส นเหตุสฺส ธมฺมสฺส อนนฺตร\-ปจฺจเยน ปจฺจโย.  อกุสโล นเหตุ ธมฺโม อพฺยากตสฺส นเหตุสฺส ธมฺมสฺส อนนฺตร\-ปจฺจเยน ปจฺจโย. (2)}

\prose{อพฺยากโต นเหตุ ธมฺโม อพฺยากตสฺส นเหตุสฺส ธมฺมสฺส อนนฺตร\-ปจฺจเยน ปจฺจโย. ตีณิ ฯเปฯ.}

\csromanheader{2}{\textbf{อุปนิสฺสย}}
\proseitem{23}{กุสโล นเหตุ ธมฺโม กุสลสฺส นเหตุสฺส ธมฺมสฺส อุปนิสฺสย\-ปจฺจเยน ปจฺจโย.  กุสโล นเหตุ ธมฺโม อกุสลสฺส นเหตุสฺส ธมฺมสฺส อุปนิสฺสย\-ปจฺจเยน ปจฺจโย.  กุสโล นเหตุ ธมฺโม อพฺยากตสฺส นเหตุสฺส ธมฺมสฺส อุปนิสฺสย\-ปจฺจเยน ปจฺจโย. (3)}

\prose{อกุสโล นเหตุ ธมฺโม อกุสลสฺส นเหตุสฺส ธมฺมสฺส อุปนิสฺสย\-ปจฺจเยน ปจฺจโย.  อกุสโล นเหตุ ธมฺโม กุสลสฺส นเหตุสฺส ธมฺมสฺส อุปนิสฺสย\-ปจฺจเยน ปจฺจโย.  อกุสโล นเหตุ ธมฺโม อพฺยากตสฺส นเหตุสฺส ธมฺมสฺส อุปนิสฺสย\-ปจฺจเยน ปจฺจโย. (3)}

\csromanmatikaanchor{mtk.94}
\csromanmatikaanchor{mtk.95}
\csromantocmarknopage{section}{2. เวทนาตฺติก}{mtk.94}
\bookmark[dest={mtk.94},level=1]{2. เวทนาตฺติก}
\csromantocmarkref{subsection}{1. เหตุทุก}{mtk.95}
\bookmark[dest={mtk.95},level=2]{1. เหตุทุก}
\csromanmarkh{2. เวทนาตฺติก}\csromanmarkhii{1. เหตุทุก}\csromanheadernomark{2}{2. เวทนาตฺติก 1. เหตุทุก}
\prose{\csromanfolio{475}อพฺยากโต นเหตุ ธมฺโม อพฺยากตสฺส นเหตุสฺส ธมฺมสฺส อุปนิสฺสย\-ปจฺจเยน ปจฺจโย.  อพฺยากโต นเหตุ ธมฺโม กุสลสฺส นเหตุสฺส ธมฺมสฺส อุปนิสฺสย\-ปจฺจเยน ปจฺจโย.  อพฺยากโต นเหตุ ธมฺโม อกุสลสฺส นเหตุสฺส ธมฺมสฺส อุปนิสฺสย\-ปจฺจเยน ปจฺจโย. (3) (สํขิตฺตํ.)}

\proseitem{24}{อารมฺมเณ นว, อธิปติยา ทส, อนนฺตเร สตฺต, สมนนฺตเร สตฺต, สหชาเต นว, อญฺญมญฺเญ ตีณิ, นิสฺสเย เตรส, อุปนิสฺสเย นว, ปุเรชาเต ตีณิ, ปจฺฉาชาเต ตีณิ, อาเสวเน ตีณิ, กมฺเม สตฺต, วิปาเก เอกํ, อาหาเร สตฺต, อินฺทฺริเย สตฺต, ฌาเน สตฺต, มคฺเค สตฺต, สมฺปยุตฺเต ตีณิ, วิปฺปยุตฺเต ปญฺจ, อตฺถิยา เตรส, นตฺถิยา สตฺต, วิคเต สตฺต, อวิคเต เตรส. (สํขิตฺตํ.)}

\prose{นเหตุยา ปนฺนรส, นอารมฺมเณ ปนฺนรส. (สํขิตฺตํ.)}

\prose{อารมฺมณ\-ปจฺจยา นเหตุยา นว. (สํขิตฺตํ.)}

\prose{นเหตุปจฺจยา อารมฺมเณ นว. (สํขิตฺตํ.)}

\prose{(ยถา กุสลตฺติเก ปญฺหาวารํ, เอวํ วิตฺถาเรตพฺพํ.)}

\csromanmarkh{2. เวทนาตฺติก}\csromanmarkhii{1. เหตุทุก}\csromanheadernomark{2}{2. เวทนาตฺติก 1. เหตุทุก}
\csromanheader{2}{1. \textbf{ปฏิจฺจาทิวาร ปจฺจยจตุกฺก}}
\proseitem{25}{สุขาย เวทนาย สมฺปยุตฺตํ เหตุํ ธมฺมํ ปฏิจฺจ สุขาย เวทนาย สมฺปยุตฺโต เหตุ ธมฺโม อุปฺปชฺชติ เหตุปจฺจยา. (1)}

\prose{ทุกฺขาย เวทนาย สมฺปยุตฺตํ เหตุํ ธมฺมํ ปฏิจฺจ ทุกฺขาย เวทนาย สมฺปยุตฺโต เหตุ ธมฺโม อุปฺปชฺชติ เหตุปจฺจยา. (1)}

\prose{อทุกฺขม\-สุขาย เวทนาย สมฺปยุตฺตํ เหตุํ ธมฺมํ ปฏิจฺจ อทุกฺขม\-สุขาย เวทนาย สมฺปยุตฺโต เหตุ ธมฺโม อุปฺปชฺชติ เหตุปจฺจยา. (1) (สํขิตฺตํ.)}

\proseitem{26}{เหตุยา ตีณิ, อารมฺมเณ ตีณิ, อธิปติยา ตีณิ ฯเปฯ วิปาเก เทฺว ฯเปฯ อวิคเต ตีณิ. (สํขิตฺตํ.)}

\prose{\csromanfolio{476}นอธิปติยา ตีณิ, นปุเรชาเต เทฺว, นปจฺฉาชาเต ตีณิ, นอาเสวเน ตีณิ, นวิปาเก ตีณิ, นวิปฺปยุตฺเต เทฺว. (สํขิตฺตํ.)}

\prose{เหตุปจฺจยา นอธิปติยา ตีณิ. (สํขิตฺตํ.)}

\prose{นอธิปติ\-ปจฺจยา เหตุยา ตีณิ. (สํขิตฺตํ.) (สหชาตวารมฺปิ ฯเปฯ สมฺปยุตฺตวารมฺปิ วิตฺถาเรตพฺพํ.)}

\csromanheader{2}{\textbf{เหตุ อารมฺมณ}}
\proseitem{27}{สุขาย เวทนาย สมฺปยุตฺโต เหตุ ธมฺโม สุขาย เวทนาย สมฺปยุตฺตสฺส เหตุสฺส ธมฺมสฺส เหตุปจฺจเยน ปจฺจโย. (1)}

\prose{ทุกฺขาย เวทนาย สมฺปยุตฺโต เหตุ ธมฺโม ทุกฺขาย เวทนาย สมฺปยุตฺตสฺส เหตุสฺส ธมฺมสฺส เหตุปจฺจเยน ปจฺจโย. (1)}

\prose{อทุกฺขม\-สุขาย เวทนาย สมฺปยุตฺโต เหตุ ธมฺโม อทุกฺขม\-สุขาย เวทนาย สมฺปยุตฺตสฺส เหตุสฺส ธมฺมสฺส เหตุปจฺจเยน ปจฺจโย. (1)}

\proseitem{28}{สุขาย เวทนาย สมฺปยุตฺโต เหตุ ธมฺโม สุขาย เวทนาย สมฺปยุตฺตสฺส เหตุสฺส ธมฺมสฺส อารมฺมณ\-ปจฺจเยน ปจฺจโย. ตีณิ.}

\prose{ทุกฺขาย เวทนาย สมฺปยุตฺโต เหตุ ธมฺโม ทุกฺขาย เวทนาย สมฺปยุตฺตสฺส เหตุสฺส ธมฺมสฺส อารมฺมณ\-ปจฺจเยน ปจฺจโย. ตีณิ.}

\prose{อทุกฺขม\-สุขาย เวทนาย สมปยุตฺโต เหตุ ธมฺโม อทุกฺขม\-สุขาย เวทนาย สมฺปยุตฺตสฺส เหตุสฺส ธมฺมสฺส อารมฺมณ\-ปจฺจเยน ปจฺจโย. ตีณิ. (สํขิตฺตํ.)}

\proseitem{29}{เหตุยา ตีณิ, อารมฺมเณ นว, อธิปติยา จตฺตาริ, อนนฺตเร ฉ, สมนนฺตเร ฉ, สหชาเต ตีณิ, อญฺญมญฺเญ ตีณิ, นิสฺสเย ตีณิ, อุปนิสฺสเย นว, อาเสวเน ตีณิ, วิปาเก เทฺว, อินฺทฺริเย เทฺว, มคฺเค เทฺว, สมฺปยุตฺเต ตีณิ, อตฺถิยา ตีณิ, นตฺถิยา ฉ, วิคเต ฉ, อวิคเต ตีณิ. (สํขิตฺตํ.)}

\prose{นเหตุยา นว, นอารมฺมเณ นว. (สํขิตฺตํ.)}

\prose{เหตุปจฺจยา นอารมฺมเณ ตีณิ. (สํขิตฺตํ.)}

\prose{นเหตุปจฺจยา อารมฺมเณ นว. (สํขิตฺตํ.)}

\prose{(ยถา กุสลตฺติเก ปญฺหาวารํ, เอวํ วิตฺถาเรตพฺพํ.)}

\csromanmarkh{2. เวทนาตฺติก}\csromanmarkhii{1. เหตุทุก 1. ปฏิจฺจวาร ปจฺจยจตุกฺก}\csromanheadernomark{2}{2. เวทนาตฺติก 1. เหตุทุก \textbf{1. ปฏิจฺจวาร ปจฺจยจตุกฺก}}
\proseitem{30}{\csromanfolio{477}สุขาย เวทนาย สมฺปยุตฺตํ นเหตุํ ธมฺมํ ปฏิจฺจ สุขาย เวทนาย สมฺปยุตฺโต นเหตุ ธมฺโม อุปฺปชฺชติ เหตุปจฺจยา. (1)}

\prose{ทุกฺขาย เวทนาย สมฺปยุตฺตํ นเหตุํ ธมฺมํ ปฏิจฺจ ทุกฺขาย เวทนาย สมฺปยุตฺโต นเหตุ ธมฺโม อุปฺปชฺชติ เหตุปจฺจยา. (1)}

\prose{อทุกฺขม\-สุขาย เวทนาย สมฺปยุตฺตํ นเหตุํ ธมฺมํ ปฏิจฺจ อทุกฺขม\-สุขาย เวทนาย สมฺปยุตฺโต นเหตุ ธมฺโม อุปฺปชฺชติ เหตุปจฺจยา. (1) (สํขิตฺตํ.)}

\proseitem{31}{เหตุยา ตีณิ, อารมฺมเณ ตีณิ ฯเปฯ อวิคเต ตีณิ. (สํขิตฺตํ.)}

\prose{นเหตุยา ตีณิ, นอธิปติยา ตีณิ, นปุเรชาเต เทฺว, นปจฺฉาชาเต ตีณิ, นอาเสวเน ตีณิ, นกมฺเม ตีณิ, นวิปาเก ตีณิ ฯเปฯ นวิปฺปยุตฺเต เทฺว. (สํขิตฺตํ.)}

\prose{(สหชาตวารมฺปิ ฯเปฯ สมฺปยุตฺตวารมฺปิ วิตฺถาเรตพฺพํ.)}

\prose{7. ปญฺหาวาร ปจฺจยจตุกฺกอารมฺมณาทิ}

\proseitem{32}{สุขาย เวทนาย สมฺปยุตฺโต นเหตุ ธมฺโม สุขาย เวทนาย สมฺปยุตฺตสฺส นเหตุสฺส ธมฺมสฺส อารมฺมณ\-ปจฺจเยน ปจฺจโย.  สุขาย เวทนาย สมฺปยุตฺโต นเหตุ ธมฺโม ทุกฺขาย เวทนาย สมฺปยุตฺตสฺส นเหตุสฺส ธมฺมสฺส อารมฺมณ\-ปจฺจเยน ปจฺจโย.  สุขาย เวทนาย สมฺปยุตฺโต นเหตุ ธมฺโม อทุกฺขม\-สุขาย เวทนาย สมฺปยุตฺตสฺส นเหตุสฺส ธมฺมสฺส อารมฺมณ\-ปจฺจเยน ปจฺจโย. (3)}

\prose{ทุกฺขาย เวทนาย สมฺปยุตฺโต นเหตุ ธมฺโม ทุกฺขาย เวทนาย สมฺปยุตฺตสฺส นเหตุสฺส ธมฺมสฺส อารมฺมณ\-ปจฺจเยน ปจฺจโย.  ทุกฺขาย เวทนาย สมฺปยุตฺโต เนเหตุ ธมฺโม สุขาย เวทนาย สมฺปยุตฺตสฺส นเหตุสฺส ธมฺมสฺส อารมฺมณ\-ปจฺจเยน ปจฺจโย.  ทุกฺขาย เวทนาย สมฺปยุตฺโต นเหตุ ธมฺโม อทุกฺขม\-สุขาย เวทนาย สมฺปยุตฺตสฺส นเหตุสฺส ธมฺมสฺส อารมฺมณ\-ปจฺจเยน ปจฺจโย. (3)}

\prose{\csromanfolio{478}อทุกฺขม\-สุขาย เวทนาย สมฺปยุตฺโต นเหตุ ธมฺโม อทุกฺขม\-สุขาย เวทนาย สมฺปยุตฺตสฺส นเหตุสฺส ธมฺมสฺส อารมฺมณ\-ปจฺจเยน ปจฺจโย.  อทุกฺขม\-สุขาย เวทนาย สมฺปยุตฺโต นเหตุ ธมฺโม สุขาย เวทนาย สมฺปยุตฺตสฺส นเหตุสฺส ธมฺมสฺส อารมฺมณ\-ปจฺจเยน ปจฺจโย.  อทุกฺขม\-สุขาย เวทนาย สมฺปยุตฺโต นเหตุ ธมฺโม ทุกฺขาย เวทนาย สมฺปยุตฺตสฺส นเหตุสฺส ธมฺมสฺส อารมฺมณ\-ปจฺจเยน ปจฺจโย. (3)}

\proseitem{33}{สุขาย เวทนาย สมฺปยุตฺโต นเหตุ ธมฺโม สุขาย เวทนาย สมฺปยุตฺตสฺส นเหตุสฺส ธมฺมสฺส อธิปติ\-ปจฺจเยน ปจฺจโย.  สุขาย เวทนาย สมฺปยุตฺโต นเหตุ ธมฺโม อทุกฺขม\-สุขาย เวทนาย สมฺปยุตฺตสฺส นเหตุสฺส ธมฺมสฺส อธิปติ\-ปจฺจเยน ปจฺจโย. (2)}

\prose{ทุกฺขาย เวทนาย สมฺปยุตฺโต นเหตุ ธมฺโม ทุกฺขาย เวทนาย สมฺปยุตฺตสฺส นเหตุสฺส ธมฺมสฺส อธิปติ\-ปจฺจเยน ปจฺจโย. (1)}

\prose{อทุกฺขม\-สุขาย เวทนาย สมฺปยุตฺโต นเหตุ ธมฺโม อทุกฺขม\-สุขาย เวทนาย สมฺปยุตฺตสฺส นเหตุสฺส ธมฺมสฺส อธิปติ\-ปจฺจเยน ปจฺจโย.  อทุกฺขม\-สุขาย เวทนาย สมฺปยุตฺโต นเหตุ ธมฺโม สุขาย เวทนาย สมฺปยุตฺตสฺส นเหตุสฺส ธมฺมสฺส อธิปติ\-ปจฺจเยน ปจฺจโย. (2)}

\csromanheader{2}{\textbf{อนนฺตร อุปนิสฺสย}}
\proseitem{34}{สุขาย เวทนาย สมฺปยุตฺโต นเหตุ ธมฺโม สุขาย เวทนาย สมฺปยุตฺตสฺส นเหตุสฺส ธมฺมสฺส อนนฺตร\-ปจฺจเยน ปจฺจโย. เทฺว.}

\prose{ทุกฺขาย เวทนาย สมฺปยุตฺโต นเหตุ ธมฺโม ทุกฺขาย เวทนาย สมฺปยุตฺตสฺส นเหตุสฺส ธมฺมสฺส อนนฺตร\-ปจฺจเยน ปจฺจโย. เทฺว.}

\prose{อทุกฺขม\-สุขาย เวทนาย สมฺปยุตฺโต นเหตุ ธมฺโม อทุกฺขม\-สุขาย เวทนาย สมฺปยุตฺตสฺส นเหตุสฺส ธมฺมสฺส อนนฺตร\-ปจฺจเยน ปจฺจโย. ตีณิ ฯเปฯ.}

\proseitem{35}{สุขาย เวทนาย สมฺปยุตฺโต นเหตุ ธมฺโม สุขาย เวทนาย สมฺปยุตฺตสฺส นเหตุสฺส ธมฺมสฺส อุปนิสฺสย\-ปจฺจเยน ปจฺจโย. ตีณิ.}

\prose{ทุกฺขาย เวทนาย สมฺปยุตฺโต นเหตุ ธมฺโม ทุกฺขาย เวทนาย สมฺปยุตฺตสฺส นเหตุ ธมฺมสฺส อุปนิสฺสย\-ปจฺจเยน ปจฺจโย. ตีณิ.}

\csromanmatikaanchor{mtk.96}
\csromanmatikaanchor{mtk.97}
\csromantocmarknopage{paragraph}{3. วิปากตฺติก}{mtk.96}
\bookmark[dest={mtk.96},level=4]{3. วิปากตฺติก}
\csromantocmarkref{subparagraph}{1. เหตุทุก}{mtk.97}
\bookmark[dest={mtk.97},level=5]{1. เหตุทุก}
\csromanmarkh{3. วิปากตฺติก}\csromanmarkhii{1. เหตุทุก}\csromanheadernomark{5}{3. วิปากตฺติก 1. เหตุทุก}
\prose{\csromanfolio{479}อทุกฺขม\-สุขาย เวทนาย สมฺปยุตฺโต นเหตุ ธมฺโม อทุกฺขม\-สุขาย เวทนาย สมฺปยุตฺตสฺส นเหตุสฺส ธมฺมสฺส อุปนิสฺสย\-ปจฺจเยน ปจฺจโย. ตีณิ. (สํขิตฺตํ.)}

\proseitem{36}{อารมฺมเณ นว, อธิปติยา ปญฺจ, อนนฺตเร สตฺต, สมนนฺตเร สตฺต, สหชาเต ตีณิ, อญฺญมญฺเญ ตีณิ, นิสฺสเย ตีณิ, อุปนิสฺสเย นว, อาเสวเน ตีณิ, กมฺเม อฏฺฐ, วิปาเก ตีณิ, อาหาเร ตีณิ, อินฺทฺริเย ตีณิ, ฌาเน ตีณิ, มคฺเค ตีณิ, สมฺปยุตฺเต ตีณิ, อตฺถิยา ตีณิ, นตฺถิยา สตฺต, วิคเต สตฺต, อวิคเต ตีณิ. (สํขิตฺตํ.)}

\prose{นเหตุยา นว, นอารมฺมเณ นว. (สํขิตฺตํ.)}

\prose{อารมฺมณ\-ปจฺจยา นเหตุยา นว. (สํขิตฺตํ.)}

\prose{นเหตุปจฺจยา อารมฺมเณ นว. (สํขิตฺตํ.)}

\prose{(ยถา กุสลตฺติเก ปญฺหาวารํ, เอวํ วิตฺถาเรตพฺพํ.)}

\csromanmarkh{3. วิปากตฺติก}\csromanmarkhii{1. เหตุทุก}\csromanheadernomark{2}{3. วิปากตฺติก 1. เหตุทุก}
\csromanheader{2}{1. \textbf{ปฏิจฺจวาร ปจฺจยจตุกฺก}}
\proseitem{37}{วิปากํ เหตุํ ธมฺมํ ปฏิจฺจ วิปาโก เหตุ ธมฺโม อุปฺปชฺชติ เหตุปจฺจยา. (1)}

\prose{วิปาก\-ธมฺม\-ธมฺมํ เหตุํ ธมฺมํ ปฏิจฺจ วิปาก\-ธมฺม\-ธมฺโม เหตุ ธมฺโม อุปฺปชฺชติ เหตุปจฺจยา. (1)}

\prose{เนว\-วิปาก\-น\-วิปากธมฺม\-ธมฺมํ เหตุํ ธมฺมํ ปฏิจฺจ เนว\-วิปาก\-น\-วิปาก\-ธมฺม\-ธมฺโม เหตุ ธมฺโม อุปฺปชฺชติ เหตุปจฺจยา. (1) (สํขิตฺตํ.)}

\proseitem{38}{เหตุยา ตีณิ, อารมฺมเณ ตีณิ ฯเปฯ อาเสวเน เทฺว, กมฺเม ตีณิ, วิปาเก เอกํ ฯเปฯ อวิคเต ตีณิ. (สํขิตฺตํ.)}

\prose{นอธิปติยา ตีณิ, นปุเรชาเต ตีณิ, นปจฺฉาชาเต ตีณิ, นอาเสวเน ตีณิ, นวิปาเก เทฺว, นวิปฺปยุตฺเต ตีณิ. (สํขิตฺตํ.)}

\prose{(สหชาตวารมฺปิ ฯเปฯ สมฺปยุตฺตวารมฺปิ ปฏิจฺจ\-วาร\-สทิสํ วิตฺถาเรตพฺพํ.)}

\csromanmarkh{ธมฺมานุโลม ติกทุก\-ปฏฺฐาน\-ปาฬิ}\csromanmarkhii{7. ปญฺหาวาร ปจฺจยจตุกฺก}\csromanheadernomark{2}{ธมฺมานุโลม ติกทุก\-ปฏฺฐาน\-ปาฬิ 7. ปญฺหาวาร ปจฺจยจตุกฺก}
\proseitem{39}{\csromanfolio{480}วิปาโก เหตุ ธมฺโม วิปากสฺส เหตุสฺส ธมฺมสฺส เหตุปจฺจเยน ปจฺจโย. (1)}

\prose{วิปาก\-ธมฺม\-ธมฺโม เหตุ ธมฺโม วิปาก\-ธมฺม\-ธมฺมสฺส เหตุสฺส ธมฺมสฺส เหตุปจฺจเยน ปจฺจโย. (1)}

\prose{เนว\-วิปาก\-น\-วิปาก\-ธมฺม\-ธมฺโม เหตุ ธมฺโม เนว\-วิปาก\-น\-วิปากธมฺม\-ธมฺมสฺส เหตุสฺส ธมฺมสฺส เหตุปจฺจเยน ปจฺจโย. (1)}

\proseitem{40}{วิปาโก เหตุ ธมฺโม วิปากสฺส เหตุสฺส ธมฺมสฺส อารมฺมณ\-ปจฺจเยน ปจฺจโย.  วิปาโก เหตุ ธมฺโม วิปาก\-ธมฺม\-ธมฺมสฺส เหตุสฺส ธมฺมสฺส อารมฺมณ\-ปจฺจเยน ปจฺจโย.  วิปาโก เหตุ ธมฺโม เนว\-วิปาก\-น\-วิปากธมฺม\-ธมฺมสฺส เหตุสฺส ธมฺมสฺส อารมฺมณ\-ปจฺจเยน ปจฺจโย. (3)}

\prose{วิปาก\-ธมฺม\-ธมฺโม เหตุ ธมฺโม วิปาก\-ธมฺม\-ธมฺมสฺส เหตุสฺส ธมฺมสฺส อารมฺมณ\-ปจฺจเยน ปจฺจโย.  วิปาก\-ธมฺม\-ธมฺโม เหตุ ธมฺโม วิปากสฺส เหตุสฺส ธมฺมสฺส อารมฺมณ\-ปจฺจเยน ปจฺจโย.  วิปาก\-ธมฺม\-ธมฺโม เหตุ ธมฺโม เนว\-วิปาก\-น\-วิปากธมฺม\-ธมฺมสฺส เหตุสฺส ธมฺมสฺส อารมฺมณ\-ปจฺจเยน ปจฺจโย. (3)}

\prose{เนว\-วิปาก\-น\-วิปาก\-ธมฺม\-ธมฺโม เหตุ ธมฺโม เนว\-วิปาก\-น\-วิปากธมฺม\-ธมฺมสฺส เหตุสฺส ธมฺมสฺส อารมฺมณ\-ปจฺจเยน ปจฺจโย. ตีณิ.}

\proseitem{41}{วิปาโก เหตุ ธมฺโม วิปากสฺส เหตุสฺส ธมฺมสฺส อธิปติ\-ปจฺจเยน ปจฺจโย.  วิปาโก เหตุ ธมฺโม วิปาก\-ธมฺม\-ธมฺมสฺส เหตุสฺส ธมฺมสฺส อธิปติ\-ปจฺจเยน ปจฺจโย.  วิปาโก เหตุ ธมฺโม เนว\-วิปาก\-น\-วิปากธมฺม\-ธมฺมสฺส เหตุสฺส ธมฺมสฺส อธิปติ\-ปจฺจเยน ปจฺจโย. ตีณิ.}

\prose{วิปาก\-ธมฺม\-ธมฺโม เหตุ ธมฺโม วิปาก\-ธมฺม\-ธมฺมสฺส เหตุสฺส ธมฺมสฺส อธิปติ\-ปจฺจเยน ปจฺจโย. เทฺว.}

\prose{เนว\-วิปาก\-น\-วิปาก\-ธมฺม\-ธมฺโม เหตุ ธมฺโม เนว\-วิปาก\-น\-วิปากธมฺม\-ธมฺมสฺส เหตุสฺส ธมฺมสฺส อธิปติ\-ปจฺจเยน ปจฺจโย. (1) (สํขิตฺตํ.)}

\csromanmarkh{3. วิปากตฺติก}\csromanmarkhii{1. เหตุทุก}\csromanheadernomark{2}{3. วิปากตฺติก 1. เหตุทุก}
\proseitem{42}{\csromanfolio{481}เหตุยา ตีณิ, อารมฺมเณ นว, อธิปติยา ฉ, อนนฺตเร ปญฺจ, สมนนฺตเร ปญฺจ, สหชาเต ตีณิ, อญฺญมญฺเญ ตีณิ, นิสฺสเย ตีณิ, อุปนิสฺสเย นว, อาเสวเน เทฺว, วิปาเก เอกํ ฯเปฯ อวิคเต ตีณิ. (สํขิตฺตํ.)}

\prose{นเหตุยา นว, นอารมฺมเณ นว. (สํขิตฺตํ.)}

\prose{เหตุปจฺจยา นอารมฺมเณ ตีณิ. (สํขิตฺตํ.)}

\prose{นเหตุปจฺจยา อารมฺมเณ นว. (สํขิตฺตํ.) (ยถา กุสลตฺติเก ปญฺหาวารํ, เอวํ วิตฺถาเรตพฺพํ.)}

\csromanheader{2}{\textbf{นเหตุปทปจฺจยจตุกฺก}}
\proseitem{43}{วิปากํ นเหตุํ ธมฺมํ ปฏิจฺจ วิปาโก นเหตุ ธมฺโม อุปฺปชฺชติ เหตุปจฺจยา.  วิปากํ นเหตุํ ธมฺมํ ปฏิจฺจ เนว\-วิปาก\-น\-วิปาก\-ธมฺม\-ธมฺโม นเหตุ ธมฺโม อุปฺปชฺชติ เหตุปจฺจยา.  วิปากํ นเหตุํ ธมฺมํ ปฏิจฺจ วิปาโก นเหตุ จ เนว\-วิปาก\-น\-วิปาก\-ธมฺม\-ธมฺโม นเหตุ จ ธมฺมา อุปฺปชฺชนฺติ เหตุปจฺจยา. (3)}

\prose{วิปาก\-ธมฺม\-ธมฺมํ นเหตุํ ธมฺมํ ปฏิจฺจ วิปาก\-ธมฺม\-ธมฺโม นเหตุ ธมฺโม อุปฺปชฺชติ เหตุปจฺจยา.  วิปาก\-ธมฺม\-ธมฺมํ นเหตุํ ธมฺมํ ปฏิจฺจ เนว\-วิปาก\-น\-วิปาก\-ธมฺม\-ธมฺโม นเหตุ ธมฺโม อุปฺปชฺชติ เหตุปจฺจยา.  วิปาก\-ธมฺม\-ธมฺมํ นเหตุํ ธมฺมํ ปฏิจฺจ วิปาก\-ธมฺม\-ธมฺโม นเหตุ จ เนว\-วิปาก\-น\-วิปาก\-ธมฺม\-ธมฺโม นเหตุ จ ธมฺมา อุปฺปชฺชนฺติ เหตุปจฺจยา. (3)}

\prose{นววิปากนวิปากธมฺมธมฺมํ นเหตุํ ธมฺมํ ปฏิจฺจ เนว\-วิปาก\-น\-วิปาก\-ธมฺม\-ธมฺโม นเหตุ ธมฺโม อุปฺปชฺชติ เหตุปจฺจยา. ตีณิ.}

\prose{วิปากํ นเหตุญฺจ เนว\-วิปาก\-น\-วิปากธมฺม\-ธมฺมํ นเหตุญฺจ ธมฺมํ ปฏิจฺจ เนว\-วิปาก\-น\-วิปาก\-ธมฺม\-ธมฺโม นเหตุ ธมฺโม อุปฺปชฺชติ เหตุปจฺจยา. ตีณิ.}

\prose{วิปาก\-ธมฺม\-ธมฺมํ นเหตุญฺจ เนว\-วิปาก\-น\-วิปากธมฺม\-ธมฺมํ นเหตุญฺจ ธมฺมํ ปฏิจฺจ เนว\-วิปาก\-น\-วิปาก\-ธมฺม\-ธมฺโม นเหตุ ธมฺโม อุปฺปชฺชติ เหตุปจฺจยา. (สํขิตฺตํ.)}

\proseitem{44}{เหตุยา เตรส, อารมฺมเณ ปญฺจ, อธิปติยา นว, อนนฺตเร ปญฺจ, สมนนฺตเร ปญฺจ, สหชาเต เตรส, อญฺญมญฺเญ สตฺต, นิสฺสเย \csromanfolio{482}เตรส, อุปนิสฺสเย ปญฺจ, ปุเรชาเต ตีณิ, อาเสวเน เทฺว, กมฺเม เตรส, วิปาเก นว, อาหาเร เตรส ฯเปฯ สมฺปยุตฺเต ปญฺจ ฯเปฯ อวิคเต เตรส. (สํขิตฺตํ.)}

\prose{นเหตุ นอารมฺมณ}

\proseitem{45}{วิปากํ นเหตุํ ธมฺมํ ปฏิจฺจ วิปาโก นเหตุ ธมฺโม อุปฺปชฺชติ นเหตุปจฺจยา. นว.}

\prose{วิปากํ นเหตุํ ธมฺมํ ปฏิจฺจ เนว\-วิปาก\-น\-วิปาก\-ธมฺม\-ธมฺโม นเหตุ ธมฺโม อุปฺปชฺชติ นอารมฺมณ\-ปจฺจยา. (สํขิตฺตํ.)}

\proseitem{46}{นเหตุยา นว, นอารมฺมเณ ปญฺจ, นอธิปติยา เตรส, นอนนฺตเร ปญฺจ, นสมนนฺตเร ปญฺจ, นอญฺญมญฺเญ ปญฺจ, นอุปนิสฺสเย ปญฺจ, นปุเรชาเต ทฺวาทส, นปจฺฉาชาเต เตรส, นอาเสวเน เตรส, นกมฺเม เทฺว, นวิปาเก ปญฺจ ฯเปฯ นวิปฺปยุตฺเต ตีณิ. (สํขิตฺตํ.) (สหชาตวารมฺปิ ฯเปฯ สมฺปยุตฺตวารมฺปิ ปฏิจฺจ\-วาร\-สทิสํ วิตฺถาเรตพฺพํ.)}

\csromanheader{2}{7. \textbf{ปญฺหาวาร ปจฺจยจตุกฺก}}
\proseitem{47}{วิปาโก นเหตุ ธมฺโม วิปากสฺส นเหตุสฺส ธมฺมสฺส อารมฺมณ\-ปจฺจเยน ปจฺจโย.  วิปาโก นเหตุ ธมฺโม วิปาก\-ธมฺม\-ธมฺมสฺส นเหตุสฺส ธมฺมสฺส อารมฺมณ\-ปจฺจเยน ปจฺจโย.  วิปาโก นเหตุ ธมฺโม เนว\-วิปาก\-น\-วิปากธมฺม\-ธมฺมสฺส นเหตุสฺส ธมฺมสฺส อารมฺมณ\-ปจฺจเยน ปจฺจโย. (3)}

\prose{วิปาก\-ธมฺม\-ธมฺโม นเหตุ ธมฺโม วิปาก\-ธมฺม\-ธมฺมสฺส นเหตุสฺส ธมฺมสฺส อารมฺมณ\-ปจฺจเยน ปจฺจโย. ตีณิ.}

\prose{เนว\-วิปาก\-น\-วิปาก\-ธมฺม\-ธมฺโม นเหตุ ธมฺโม เนว\-วิปาก\-น\-วิปากธมฺม\-ธมฺมสฺส นเหตุสฺส ธมฺมสฺส อารมฺมณ\-ปจฺจเยน ปจฺจโย. ตีณิ.}

\proseitem{48}{วิปาโก นเหตุ ธมฺโม วิปากสฺส นเหตุสฺส ธมฺมสฺส อธิปติ\-ปจฺจเยน ปจฺจโย. จตฺตาริ.}

\prose{วิปาก\-ธมฺม\-ธมฺโม นเหตุ ธมฺโม วิปาก\-ธมฺม\-ธมฺมสฺส นเหตุสฺส ธมฺมสฺส อธิปติ\-ปจฺจเยน ปจฺจโย. ตีณิ.}

\csromanmatikaanchor{mtk.98}
\csromanmatikaanchor{mtk.99}
\csromantocmarknopage{subparagraph}{4. อุปาทินฺนตฺติก}{mtk.98}
\bookmark[dest={mtk.98},level=5]{4. อุปาทินฺนตฺติก}
\csromantocmarkref{subparagraph}{1. เหตุทุก}{mtk.99}
\bookmark[dest={mtk.99},level=5]{1. เหตุทุก}
\csromanmarkh{4. อุปาทินฺนตฺติก}\csromanmarkhii{1. เหตุทุก}\csromanheadernomark{6}{4. อุปาทินฺนตฺติก 1. เหตุทุก}
\prose{\csromanfolio{483}เนว\-วิปาก\-น\-วิปาก\-ธมฺม\-ธมฺโม นเหตุ ธมฺโม เนว\-วิปาก\-น\-วิปากธมฺม\-ธมฺมสฺส นเหตุสฺส ธมฺมสฺส อธิปติ\-ปจฺจเยน ปจฺจโย.  เนว\-วิปาก\-น\-วิปาก\-ธมฺม\-ธมฺโม นเหตุ ธมฺโม วิปากสฺส นเหตุสฺส ธมฺมสฺส อธิปติ\-ปจฺจเยน ปจฺจโย.  เนว\-วิปาก\-น\-วิปาก\-ธมฺม\-ธมฺโม นเหตุ ธมฺโม วิปาก\-ธมฺม\-ธมฺมสฺส นเหตุสฺส ธมฺมสฺส อธิปติ\-ปจฺจเยน ปจฺจโย. ตีณิ.}

\proseitem{49}{วิปาโก นเหตุ ธมฺโม วิปากสฺส นเหตุสฺส ธมฺมสฺส อนนฺตร\-ปจฺจเยน ปจฺจโย. เทฺว.}

\prose{วิปาก\-ธมฺม\-ธมฺโม นเหตุ ธมฺโม วิปาก\-ธมฺม\-ธมฺมสฺส นเหตุสฺส ธมฺมสฺส อนนฺตร\-ปจฺจเยน ปจฺจโย. เทฺว.}

\prose{เนว\-วิปาก\-น\-วิปาก\-ธมฺม\-ธมฺโม นเหตุ ธมฺโม เนว\-วิปาก\-น\-วิปากธมฺม\-ธมฺมสฺส นเหตุสฺส ธมฺมสฺส อนนฺตร\-ปจฺจเยน ปจฺจโย. ตีณิ. (สํขิตฺตํ.)}

\proseitem{50}{อารมฺมเณ นว, อธิปติยา ทส, อนนฺตเร สตฺต, สมนนฺตเร สตฺต, สหชาเต เอกาทส, อญฺญมญฺเญ สตฺต, นิสฺสเย เตรส, อุปนิสฺสเย นว, อาเสวเน เทฺว, กมฺเม นว, วิปาเก ตีณิ, อาหาเร สตฺต, อินฺทฺริเย นว, ฌาเน สตฺต, มคฺเค สตฺต, สมฺปยุตฺเต ตีณิ, วิปฺปยุตฺเต ปญฺจ, อตฺถิยา เตรส ฯเปฯ อวิคเต เตรส. (สํขิตฺตํ.)}

\prose{นเหตุยา โสฬส, นอารมฺมเณ โสฬส. (สํขิตฺตํ.)}

\prose{เหตุปจฺจยา นอารมฺมเณ นว. (สํขิตฺตํ.)}

\prose{นเหตุปจฺจยา อารมฺมเณ นว. (สํขิตฺตํ.)}

\prose{(ยถา กุสลตฺติเก ปญฺหาวารํ, เอวํ วิตฺถาเรตพฺพํ.)}

\csromanmarkh{4. อุปาทินฺนตฺติก}\csromanmarkhii{1. เหตุทุก}\csromanheadernomark{2}{4. อุปาทินฺนตฺติก 1. เหตุทุก}
\csromanheader{2}{1. \textbf{ปฏิจฺจวาร ปจฺจยจตุกฺก}}
\proseitem{51}{อุปาทินฺนุ\-ปาทานิยํ เหตุํ ธมฺมํ ปฏิจฺจ อุปาทินฺนุ\-ปาทานิโย เหตุ ธมฺโม อุปฺปชฺชติ เหตุปจฺจยา. (1)}

\prose{อนุปาทินฺนุ\-ปาทานิยํ เหตุํ ธมฺมํ ปฏิจฺจ อนุปาทินฺนุ\-ปาทานิโย เหตุ ธมฺโม อุปฺปชฺชติ เหตุปจฺจยา. (1)}

\prose{\csromanfolio{484}อนุปาทินฺน\-อนุปาทานิยํ เหตุํ ธมฺมํ ปฏิจฺจ อนุปาทินฺน\-อนุปาทานิโย เหตุ ธมฺโม อุปฺปชฺชติ เหตุปจฺจยา. (1) (สํขิตฺตํ.)}

\proseitem{52}{เหตุยา ตีณิ, อารมฺมเณ ตีณิ, อธิปติยา เทฺว ฯเปฯ วิปาเก เทฺว. (สํขิตฺตํ.)}

\prose{นอธิปติยา ตีณิ, นปุเรชาเต ตีณิ ฯเปฯ นอาเสวเน ตีณิ, นวิปาเก เทฺว, นวิปฺปยุตฺเต ตีณิ. (สํขิตฺตํ. สหชาตวารมฺปิ ฯเปฯ สมฺปยุตฺตวารมฺปิ วิตฺถาเรตพฺพํ.)}

\csromanheader{2}{7. \textbf{ปญฺหาวาร ปจฺจยจตุกฺก}}
\proseitem{53}{อุปาทินฺนุ\-ปาทานิโย เหตุ ธมฺโม อุปาทินฺนุ\-ปาทานิยสฺส เหตุสฺส ธมฺมสฺส เหตุปจฺจเยน ปจฺจโย. (1)}

\prose{อนุปาทินฺนุ\-ปาทานิโย เหตุ ธมฺโม อนุปาทินฺนุ\-ปาทานิยสฺส เหตุสฺส ธมฺมสฺส เหตุปจฺจเยน ปจฺจโย. (1)}

\prose{อนุปาทินฺน\-อนุปาทานิโย เหตุ ธมฺโม อนุปาทินฺน\-อนุปาทานิยสฺส เหตุสฺส ธมฺมสฺส เหตุปจฺจเยน ปจฺจโย. (1)}

\proseitem{54}{อุปาทินฺนุ\-ปาทานิโย เหตุ ธมฺโม อุปาทินฺนุ\-ปาทานิยสฺส เหตุสฺส ธมฺมสฺส อารมฺมณ\-ปจฺจเยน ปจฺจโย.  อุปาทินฺนุ\-ปาทานิโย เหตุ ธมฺโม อนุปาทินฺนุ\-ปาทานิยสฺส เหตุสฺส ธมฺมสฺส อารมฺมณ\-ปจฺจเยน ปจฺจโย. (2)}

\prose{อนุปาทินฺนุ\-ปาทานิโย เหตุ ธมฺโม อนุปาทินฺนุ\-ปาทานิยสฺส เหตุสฺส ธมฺมสฺส อารมฺมณ\-ปจฺจเยน ปจฺจโย.  อนุปาทินฺนุ\-ปาทานิโย เหตุ ธมฺโม อุปาทินฺนุ\-ปาทานิยสฺส เหตุสฺส ธมฺมสฺส อารมฺมณ\-ปจฺจเยน ปจฺจโย. (2)}

\prose{อนุปาทินฺน\-อนุปาทานิโย เหตุ ธมฺโม อนุปาทินฺนุ\-ปาทานิยสฺส เหตุสฺส ธมฺมสฺส อารมฺมณ\-ปจฺจเยน ปจฺจโย. (1)}

\proseitem{55}{อุปาทินฺนุ\-ปาทานิโย เหตุ ธมฺโม อนุปาทินฺนุ\-ปาทานิยสฺส เหตุสฺส ธมฺมสฺส อธิปติ\-ปจฺจเยน ปจฺจโย. (1)}

\prose{อนุปาทินฺนุ\-ปาทานิโย เหตุ ธมฺโม อนุปาทินฺนุ\-ปาทานิยสฺส เหตุสฺส ธมฺมสฺส อธิปติ\-ปจฺจเยน ปจฺจโย. (1)}

\prose{อนุปาทินฺน\-อนุปาทานิโย เหตุ ธมฺโม อนุปาทินฺน\-อนุปาทานิยสฺส เหตุสฺส ธมฺมสฺส อธิปติ\-ปจฺจเยน ปจฺจโย.  อนุปาทินฺน\-อนุปาทานิโย}

\vspace{\tipitakaparskip}
\csromancenter{อุปาทินฺนตฺติก 1. เหตุทุก เหตุ ธมฺโม อนุปาทินฺนุ\-ปาทานิยสฺส เหตุสฺส ธมฺมสฺส อธิปติ\-ปจฺจเยน ปจฺจโย. (2)}
\csromanheader{2}{\textbf{อนนฺตร อุปนิสฺสย}}
\proseitem{56}{\csromanfolio{485}อุปาทินฺนุ\-ปาทานิโย เหตุ ธมฺโม อุปาทินฺนุ\-ปาทานิยสฺส เหตุสฺส ธมฺมสฺส อนนฺตร\-ปจฺจเยน ปจฺจโย. (1)}

\prose{อนุปาทินฺนุ\-ปาทานิโย เหตุ ธมฺโม อนุปาทินฺนุ\-ปาทานิยสฺส เหตุสฺส ธมฺมสฺส อนนฺตร\-ปจฺจเยน ปจฺจโย. ตีณิ.}

\prose{อนุปาทินฺน\-อนุปาทานิโย เหตุ ธมฺโม อนุปาทินฺน\-อนุปาทานิยสฺส เหตุสฺส ธมฺมสฺส อนนฺตร\-ปจฺจเยน ปจฺจโย.  อนุปาทินฺน\-อนุปาทานิโย เหตุ ธมฺโม อุปาทินฺนุ\-ปาทานิยสฺส เหตุสฺส ธมฺมสฺส อนนฺตร\-ปจฺจเยน ปจฺจโย ฯเปฯ. (2)}

\proseitem{57}{อุปาทินฺนุ\-ปาทานิโย เหตุ ธมฺโม อุปาทินฺนุ\-ปาทานิยสฺส เหตุสฺส ธมฺมสฺส อุปนิสฺสย\-ปจฺจเยน ปจฺจโย.  อุปาทินฺนุ\-ปาทานิโย เหตุ ธมฺโม อนุปาทินฺนุ\-ปาทานิยสฺส เหตุสฺส ธมฺมสฺส อุปนิสฺสย\-ปจฺจเยน ปจฺจโย.  อุปาทินฺนุ\-ปาทานิโย เหตุ ธมฺโม อนุปาทินฺน\-อนุปาทานิยสฺส เหตุสฺส ธมฺมสฺส อุปนิสฺสย\-ปจฺจเยน ปจฺจโย. (3)}

\prose{อนุปาทินฺนุ\-ปาทานิโย เหตุ ธมฺโม อนุปาทินฺนุ\-ปาทานิยสฺส เหตุสฺส ธมฺมสฺส อุปนิสฺสย\-ปจฺจเยน ปจฺจโย. ตีณิ.}

\prose{อนุปาทินฺน\-อนุปาทานิโย เหตุ ธมฺโม อนุปาทินฺน\-อนุปาทานิยสฺส เหตุสฺส ธมฺมสฺส อุปนิสฺสย\-ปจฺจเยน ปจฺจโย. ตีณิ. (สํขิตฺตํ.)}

\proseitem{58}{เหตุยา ตีณิ, อารมฺมเณ ปญฺจ, อธิปติยา จตฺตาริ, อนนฺตเร ฉ, สมนนฺตเร ฉ, สหชาเต ตีณิ, อญฺญมญฺเญ ตีณิ, นิสฺสเย ตีณิ, อุปนิสฺสเย นว, อาเสวเน เทฺว, วิปาเก เทฺว ฯเปฯ อวิคเต ตีณิ. (สํขิตฺตํ.)}

\prose{นเหตุยา นว, นอารมฺมเณ นว. (สํขิตฺตํ.)}

\prose{เหตุปจฺจยา นอารมฺมเณ ตีณิ. (สํขิตฺตํ.)}

\prose{นเหตุปจฺจยา อารมฺมเณ ปญฺจ. (สํขิตฺตํ.)}

\prose{(ยถา กุสลตฺติเก ปญฺหาวารํ, เอวํ วิตฺถาเรตพฺพํ.)}

\csromanheader{2}{ธมฺมานุโลม ติกทุก\-ปฏฺฐาน\-ปาฬิ นเหตุปทปจฺจยจตุกฺก}
\proseitem{59}{\csromanfolio{486}อุปาทินฺนุ\-ปาทานิยํ นเหตุํ ธมฺมํ ปฏิจฺจ อุปาทินฺนุ\-ปาทานิโย นเหตุ ธมฺโม อุปฺปชฺชติ เหตุปจฺจยา.  อุปาทินฺนุ\-ปาทานิยํ นเหตุํ ธมฺมํ ปฏิจฺจ อนุปาทินฺนุ\-ปาทานิโย นเหตุ ธมฺโม อุปฺปชฺชติ เหตุปจฺจยา.  อุปาทินฺนุ\-ปาทานิยํ นเหตุํ ธมฺมํ ปฏิจฺจ อุปาทินฺนุ\-ปาทานิโย นเหตุ จ อนุปาทินฺนุ\-ปาทานิโย นเหตุ จ ธมฺมา อุปฺปชฺชนฺติ เหตุปจฺจยา. (3)}

\prose{อนุปาทินฺนุ\-ปาทานิยํ นเหตุํ ธมฺมํ ปฏิจฺจ อนุปาทินฺนุ\-ปาทานิโย นเหตุ ธมฺโม อุปฺปชฺชติ เหตุปจฺจยา. (1)}

\prose{อนุปาทินฺน\-อนุปาทานิยํ นเหตุํ ธมฺมํ ปฏิจฺจ อนุปาทินฺน\-อนุปาทานิโย นเหตุ ธมฺโม อุปฺปชฺชติ เหตุปจฺจยา.  อนุปาทินฺน\-อนุปาทานิยํ นเหตุํ ธมฺมํ ปฏิจฺจ อนุปาทินฺนุ\-ปาทานิโย นเหตุ ธมฺโม อุปฺปชฺชติ เหตุปจฺจยา.  อนุปาทินฺน\-อนุปาทานิยํ นเหตุํ ธมฺมํ ปฏิจฺจ อนุปาทินฺนุ\-ปาทานิโย นเหตุ จ อนุปาทินฺน\-อนุปาทานิโย นเหตุ จ ธมฺมา อุปฺปชฺชนฺติ เหตุปจฺจยา. (3)}

\prose{อุปาทินฺนุ\-ปาทานิยํ นเหตุญฺจ อนุปาทินฺนุ\-ปาทานิยํ นเหตุญฺจ ธมฺมํ ปฏิจฺจ อนุปาทินฺนุ\-ปาทานิโย นเหตุ ธมฺโม อุปฺปชฺชติ เหตุปจฺจยา. (1)}

\prose{อนุปาทินฺนุ\-ปาทานิยํ นเหตุญฺจ อนุปาทินฺน\-อนุปาทานิยํ นเหตุญฺจ ธมฺมํ ปฏิจฺจ อนุปาทินฺนุ\-ปาทานิโย นเหตุ ธมฺโม อุปฺปชฺชติ เหตุปจฺจยา. (1)}

\proseitem{60}{อุปาทินฺนุ\-ปาทานิยํ นเหตุํ ธมฺมํ ปฏิจฺจ อุปาทินฺนุ\-ปาทานิโย นเหตุ ธมฺโม อุปฺปชฺชติ อารมฺมณ\-ปจฺจยา. (1)}

\prose{อนุปาทินฺนุ\-ปาทานิยํ นเหตุํ ธมฺมํ ปฏิจฺจ อนุปาทินฺนุ\-ปาทานิโย นเหตุ ธมฺโม อุปฺปชฺชติ อารมฺมณ\-ปจฺจยา. (1)}

\prose{อนุปาทินฺน\-อนุปาทานิยํ นเหตุํ ธมฺมํ ปฏิจฺจ อนุปาทินฺน\-อนุปาทานิโย นเหตุ ธมฺโม อุปฺปชฺชติ อารมฺมณ\-ปจฺจยา. (1) (สํขิตฺตํ.)}

\proseitem{61}{เหตุยา นว, อารมฺมเณ ตีณิ, อธิปติยา ปญฺจ, อนนฺตเร ตีณิ, สมนนฺตเร ตีณิ, สหชาเต นว, อญฺญมญฺเญ ตีณิ, นิสฺสเย นว, อุปนิสฺสเย ตีณิ, ปุเรชาเต ตีณิ, อาเสวเน เทฺว, กมฺเม นว, วิปาเก นว ฯเปฯ อวิคเต นว. (สํขิตฺตํ.)}

\csromanmarkh{5. สํกิลิฏฺฐ\-ตฺติก}\csromanmarkhii{1. เหตุทุก ปจฺจนีย\-น\-เหตุ}\csromanheadernomark{2}{5. สํกิลิฏฺฐ\-ตฺติก 1. เหตุทุก ปจฺจนีย\-น\-เหตุ}
\proseitem{62}{\csromanfolio{487}อุปาทินฺนุ\-ปาทานิยํ นเหตุํ ธมฺมํ ปฏิจฺจ อุปาทินฺนุ\-ปาทานิโย นเหตุ ธมฺโม อุปฺปชฺชติ นเหตุปจฺจยา. ตีณิ.}

\prose{อนุปาทินฺนุ\-ปาทานิยํ นเหตุํ ธมฺมํ ปฏิจฺจ อนุปาทินฺนุ\-ปาทานิโย นเหตุ ธมฺโม อุปฺปชฺชติ นเหตุปจฺจยา. (1)}

\prose{อุปาทินฺนุ\-ปาทานิยํ นเหตุญฺจ อนุปาทินฺนุ\-ปาทานิยํ นเหตุญฺจ ธมฺมํ ปฏิจฺจ อนุปาทินฺนุ\-ปาทานิโย นเหตุ ธมฺโม อุปฺปชฺชติ นเหตุปจฺจยา. (1) (สํขิตฺตํ.)}

\proseitem{63}{นเหตุยา ปญฺจ, นอารมฺมเณ ฉ. (สํขิตฺตํ.)}

\prose{เหตุปจฺจยา นอารมฺมเณ ฉ. (สํขิตฺตํ.)}

\prose{นเหตุปจฺจยา อารมฺมเณ เทฺว. (สํขิตฺตํ.)}

\prose{(ยถา กุสลตฺติเก ปญฺหาวารํ, เอวํ วิตฺถาเรตพฺพํ.)}

\csromanmatikaanchor{mtk.100}
\csromanmatikaanchor{mtk.101}
\csromantocmarknopage{subparagraph}{5. สํกิลิฏฺฐตฺติก}{mtk.100}
\bookmark[dest={mtk.100},level=5]{5. สํกิลิฏฺฐตฺติก}
\csromantocmarkref{subparagraph}{1. เหตุทุก}{mtk.101}
\bookmark[dest={mtk.101},level=5]{1. เหตุทุก}
\csromanmarkh{5. สํกิลิฏฺฐ\-ตฺติก}\csromanmarkhii{1. เหตุทุก}\csromanheadernomark{6}{5. สํกิลิฏฺฐ\-ตฺติก 1. เหตุทุก}
\csromanheader{2}{1. \textbf{ปฏิจฺจวาร ปจฺจยจตุกฺก}}
\proseitem{64}{สํกิลิฏฺฐ\-สํกิเลสิกํ เหตุํ ธมฺมํ ปฏิจฺจ สํกิลิฏฺฐ\-สํกิเลสิโก เหตุ ธมฺโม อุปฺปชฺชติ เหตุปจฺจยา. (1)}

\prose{อสํกิลิฏฺฐ\-สํกิเลสิกํ เหตุํ ธมฺมํ ปฏิจฺจ อสํกิลิฏฺฐ\-สํกิเลสิโก เหตุ ธมฺโม อุปฺปชฺชติ เหตุปจฺจยา. (1)}

\prose{อสํกิลิฏฺฐ\-อสํกิเลสิกํ เหตุํ ธมฺมํ ปฏิจฺจ อสํกิลิฏฺฐ\-อสํกิเลสิโก เหตุ ธมฺโม อุปฺปชฺชติ เหตุปจฺจยา. (1) (สํขิตฺตํ.)}

\proseitem{65}{เหตุยา ตีณิ, อารมฺมเณ ตีณิ ฯเปฯ อวิคเต ตีณิ. (สํขิตฺตํ.)}

\prose{\csromanfolio{488}นอธิปติยา ตีณิ, นปุเรชาเต ตีณิ, นปจฺฉาชาเต ตีณิ, นอาเสวเน ตีณิ, นวิปาเก ตีณิ, นวิปฺปยุตฺเต ตีณิ. (สํขิตฺตํ.)}

\prose{(สหชาตวารมฺปิ ฯเปฯ สมฺปยุตฺตวารมฺปิ วิตฺถาเรตพฺพํ.)}

\csromanheader{2}{7. \textbf{ปญฺหาวาร ปจฺจยจตุกฺก}}
\prose{เหตุอารมฺมณาทิ}

\proseitem{66}{สํกิลิฏฺฐ\-สํกิเลสิโก เหตุ ธมฺโม สํกิลิฏฺฐ\-สํกิเลสิกสฺส เหตุสฺส ธมฺมสฺส เหตุปจฺจเยน ปจฺจโย. (1)}

\prose{อสํกิลิฏฺฐ\-สํกิเลสิโก เหตุ ธมฺโม อสํกิลิฏฺฐ\-สํกิเลสิกสฺส เหตุสฺส ธมฺมสฺส เหตุปจฺจเยน ปจฺจโย. (1)}

\prose{อสํกิลิฏฺฐ\-อสํกิเลสิโก เหตุ ธมฺโม อสํกิลิฏฺฐ\-อสํกิเลสิกสฺส เหตุสฺส ธมฺมสฺส เหตุปจฺจเยน ปจฺจโย. (1)}

\proseitem{67}{สํกิลิฏฺฐ\-สํกิเลสิโก เหตุ ธมฺโม สํกิลิฏฺฐ\-สํกิเลสิกสฺส เหตุสฺส ธมฺมสฺส อารมฺมณ\-ปจฺจเยน ปจฺจโย. เทฺว.}

\prose{อสํกิลิฏฺฐ\-สํกิเลสิโก เหตุ ธมฺโม อสํกิลิฏฺฐ\-สํกิเลสิกสฺส เหตุสฺส ธมฺมสฺส อารมฺมณ\-ปจฺจเยน ปจฺจโย. เทฺว.}

\prose{อสํกิลิฏฺฐ\-อสํกิเลสิโก เหตุ ธมฺโม อสํกิลิฏฺฐ\-สํกิเลสิกสฺส เหตุสฺส ธมฺมสฺส อารมฺมณ\-ปจฺจเยน ปจฺจโย. (1)}

\proseitem{68}{สํกิลิฏฺฐ\-สํกิเลสิโก เหตุ ธมฺโม สํกิลิฏฺฐ\-สํกิเลสิกสฺส เหตุสฺส ธมฺมสฺส อธิปติ\-ปจฺจเยน ปจฺจโย. (1)}

\prose{อสํกิลิฏฺฐ\-สํกิเลสิโก เหตุ ธมฺโม อสํกิลิฏฺฐ\-สํกิเลสิกสฺส เหตุสฺส ธมฺมสฺส อธิปติ\-ปจฺจเยน ปจฺจโย.  อสํกิลิฏฺฐ\-สํกิเลสิโก เหตุ ธมฺโม สํกิลิฏฺฐ\-สํกิเลสิกสฺส เหตุสฺส ธมฺมสฺส อธิปติ\-ปจฺจเยน ปจฺจโย. (2)}

\prose{อสํกิลิฏฺฐ\-อสํกิเลสิโก เหตุ ธมฺโม อสํกิลิฏฺฐ\-อสํกิเลสิกสฺส เหตุสฺส ธมฺมสฺส อธิปติ\-ปจฺจเยน ปจฺจโย. เทฺว.}

\csromanmarkh{5. สํกิลิฏฺฐ\-ตฺติก}\csromanmarkhii{1. เหตุทุก}\csromanheadernomark{2}{5. สํกิลิฏฺฐ\-ตฺติก 1. เหตุทุก}
\proseitem{69}{\csromanfolio{489}สํกิลิฏฺฐ\-สํกิเลสิโก เหตุ ธมฺโม สํกิลิฏฺฐ\-สํกิเลสิกสฺส เหตุสฺส ธมฺมสฺส อนนฺตร\-ปจฺจเยน ปจฺจโย. ฉ ฯเปฯ.}

\prose{สํกิลิฏฺฐ\-สํกิเลสิโก เหตุ ธมฺโม สํกิลิฏฺฐ\-สํกิเลสิกสฺส เหตุสฺส ธมฺมสฺส อุปนิสฺสย\-ปจฺจเยน ปจฺจโย.  สํกิลิฏฺฐ\-สํกิเลสิโก เหตุ ธมฺโม อสํกิลิฏฺฐ\-สํกิเลสิกสฺส เหตุสฺส ธมฺมสฺส อุปนิสฺสย\-ปจฺจเยน ปจฺจโย.  สํกิลิฏฺฐ\-สํกิเลสิโก เหตุ ธมฺโม อสํกิลิฏฺฐ\-อสํกิเลสิกสฺส เหตุสฺส ธมฺมสฺส อุปนิสฺสย\-ปจฺจเยน ปจฺจโย. (3)}

\prose{อสํกิลิฏฺฐ\-สํกิเลสิโก เหตุ ธมฺโม อสํกิลิฏฺฐ\-สํกิเลสิกสฺส เหตุสฺส ธมฺมสฺส อุปนิสฺสย\-ปจฺจเยน ปจฺจโย. ตีณิ.}

\prose{อสํกิลิฏฺฐ\-อสํกิเลสิโก เหตุ ธมฺโม อสํกิลิฏฺฐ\-อสํกิเลสิกสฺส เหตุสฺส ธมฺมสฺส อุปนิสฺสย\-ปจฺจเยน ปจฺจโย. เทฺว. (สํขิตฺตํ.)}

\proseitem{70}{เหตุยา ตีณิ, อารมฺมเณ ปญฺจ, อธิปติยา ปญฺจ, อนนฺตเร ฉ, สมนนฺตเร ฉ, สหชาเต ตีณิ, อญฺญมญฺเญ ตีณิ, นิสฺสเย ตีณิ, อุปนิสฺสเย อฏฺฐ, อาเสวเน ตีณิ, วิปาเก เทฺว, อินฺทฺริเย เทฺว ฯเปฯ อวิคเต ตีณิ. (สํขิตฺตํ.)}

\prose{นเหตุยา อฏฺฐ, นอารมฺมเณ อฏฺฐ. (สํขิตฺตํ.)}

\prose{เหตุปจฺจยา นอารมฺมเณ ตีณิ. (สํขิตฺตํ.)}

\prose{นเหตุปจฺจยา อารมฺมเณ ปญฺจ. (สํขิตฺตํ.)}

\prose{(ยถา กุสลตฺติเก ปญฺหาวารํ, เอวํ วิตฺถาเรตพฺพํ.)}

\csromanheader{2}{\textbf{นเหตุปทปจฺจยจตุกฺก}}
\prosecont{สํกิลิฏฺฐ\-สํกิเลสิกํ นเหตุํ ธมฺมํ ปฏิจฺจ สํกิลิฏฺฐ\-สํกิเลสิโก นเหตุ ธมฺโม อุปฺปชฺชติ เหตุปจฺจยา.  สํกิลิฏฺฐ\-สํกิเลสิกํ นเหตุํ ธมฺมํ ปฏิจฺจ อสํกิลิฏฺฐ\-สํกิเลสิโก นเหตุ ธมฺโม อุปฺปชฺชติ เหตุปจฺจยา.  สํกิลิฏฺฐ\-สํกิเลสิกํ นเหตุํ ธมฺมํ ปฏิจฺจ สํกิลิฏฺฐ\-สํกิเลสิโก นเหตุ จ อสํกิลิฏฺฐ\-สํกิเลสิโก นเหตุ จ ธมฺมา อุปฺปชฺชนฺติ เหตุปจฺจยา. (3)}

\proseitem{71}{อสํกิลิฏฺฐ\-สํกิเลสิกํ นเหตุํ ธมฺมํ ปฏิจฺจ อสํกิลิฏฺฐ\-สํกิเลสิโก นเหตุ ธมฺโม อุปฺปชฺชติ เหตุปจฺจยา. (1)}

\prose{\csromanfolio{490}อสํกิลิฏฺฐ\-อสํกิเลสิกํ นเหตุํ ธมฺมํ ปฏิจฺจ อสํกิลิฏฺฐ\-อสํกิเลสิโก นเหตุ ธมฺโม อุปฺปชฺชติ เหตุปจฺจยา. ตีณิ.}

\prose{สํกิลิฏฺฐ\-สํกิเลสิกํ นเหตุญฺจ อสํกิลิฏฺฐ\-สํกิเลสิกํ นเหตุญฺจ ธมฺมํ ปฏิจฺจ อสํกิลิฏฺฐ\-สํกิเลสิโก นเหตุ ธมฺโม อุปฺปชฺชติ เหตุปจฺจยา. (1)}

\prose{อสํกิลิฏฺฐ\-สํกิเลสิกํ นเหตุญฺจ อสํกิลิฏฺฐ\-อสํกิเลสิกํ นเหตุญฺจ ธมฺมํ ปฏิจฺจ อสํกิลิฏฺฐ\-สํกิเลสิโก นเหตุ ธมฺโม อุปฺปชฺชติ เหตุปจฺจยา. (1)}

\proseitem{72}{สํกิลิฏฺฐ\-สํกิเลสิกํ นเหตุํ ธมฺมํ ปฏิจฺจ สํกิลิฏฺฐ\-สํกิเลสิโก นเหตุ ธมฺโม อุปฺปชฺชติ อารมฺมณ\-ปจฺจยา. (1)}

\prose{อสํกิลิฏฺฐ\-สํกิเลสิกํ นเหตุํ ธมฺมํ ปฏิจฺจ อสํกิลิฏฺฐ\-สํกิเลสิโก นเหตุ ธมฺโม อุปฺปชฺชติ อารมฺมณ\-ปจฺจยา. (1)}

\prose{อสํกิลิฏฺฐ\-อสํกิเลสิกํ นเหตุํ ธมฺมํ ปฏิจฺจ อสํกิลิฏฺฐ\-อสํกิเลสิโก นเหตุ ธมฺโม อุปฺปชฺชติ อารมฺมณ\-ปจฺจยา. (1) (สํขิตฺตํ.)}

\proseitem{73}{เหตุยา นว, อารมฺมเณ ตีณิ, อธิปติยา นว, อนนฺตเร ตีณิ, สมนนฺตเร ตีณิ, สหชาเต นว, อญฺญมญฺเญ ตีณิ, นิสฺสเย นว, อุปนิสฺสเย ตีณิ, ปุเรชาเต ตีณิ, อาเสวเน ตีณิ, กมฺเม นว, วิปาเก ปญฺจ, อาหาเร นว ฯเปฯ อวิคเต นว. (สํขิตฺตํ.)}

\prose{ปจฺจนีย\-น\-เหตุ}

\proseitem{74}{อสํกิลิฏฺฐ\-สํกิเลสิกํ นเหตุํ ธมฺมํ ปฏิจฺจ อสํกิลิฏฺฐ\-สํกิเลสิโก นเหตุ ธมฺโม อุปฺปชฺชติ นเหตุปจฺจยา. (สํขิตฺตํ.)}

\prose{นเหตุยา เอกํ, นอารมฺมเณ ปญฺจ, นอธิปติยา ฉ, นปุเรชาเต สตฺต ฯเปฯ นกมฺเม ตีณิ ฯเปฯ นวิปฺปยุตฺเต ตีณิ ฯเปฯ โนวิคเต ปญฺจ. (สํขิตฺตํ.)}

\prose{(สหชาตวารมฺปิ ฯเปฯ สมฺปยุตฺตวารมฺปิ วิตฺถาเรตพฺพํ.)}

\csromanmarkh{6. วิตกฺกตฺติก}\csromanmarkhii{1. เหตุทุก 7. ปญฺหาวาร ปจฺจยจตุกฺก}\csromanheadernomark{2}{6. วิตกฺกตฺติก 1. เหตุทุก \textbf{7. ปญฺหาวาร ปจฺจยจตุกฺก}}
\prose{\csromanfolio{491}อารมฺมณ\-อธิปติ}

\proseitem{75}{สํกิลิฏฺฐ\-สํกิเลสิโก นเหตุ ธมฺโม สํกิลิฏฺฐ\-สํกิเลสิกสฺส นเหตุสฺส ธมฺมสฺส อารมฺมณ\-ปจฺจเยน ปจฺจโย. เทฺว.}

\prose{อสํกิลิฏฺฐ\-สํกิเลสิโก นเหตุ ธมฺโม อสํกิลิฏฺฐ\-สํกิเลสิกสฺส นเหตุสฺส ธมฺมสฺส อารมฺมณ\-ปจฺจเยน ปจฺจโย. เทฺว.}

\prose{อสํกิลิฏฺฐ\-อสํกิเลสิโก นเหตุ ธมฺโม อสํกิลิฏฺฐ\-อสํกิเลสิกสฺส นเหตุสฺส ธมฺมสฺส อารมฺมณ\-ปจฺจเยน ปจฺจโย. เทฺว.}

\proseitem{76}{สํกิลิฏฺฐ\-สํกิเลสิโก นเหตุ ธมฺโม สํกิลิฏฺฐ\-สํกิเลสิกสฺส นเหตุสฺส ธมฺมสฺส อธิปติ\-ปจฺจเยน ปจฺจโย. ตีณิ.}

\prose{อสํกิลิฏฺฐ\-สํกิเลสิโก นเหตุ ธมฺโม อสํกิลิฏฺฐ\-สํกิเลสิกสฺส นเหตุสฺส ธมฺมสฺส อธิปติเยน ปจฺจโย. เทฺว.}

\prose{อสํกิลิฏฺฐ\-อสํกิเลสิโก นเหตุ ธมฺโม อสํกิลิฏฺฐ\-อสํกิเลสิกสฺส นเหตุสฺส ธมฺมสฺส อธิปติ\-ปจฺจเยน ปจฺจโย. ตีณิ. (สํขิตฺตํ.)}

\proseitem{77}{อารมฺมเณ ฉ, อธิปติยา อฏฺฐ, อนนฺตเร สตฺต, สมนนฺตเร สตฺต, สหชาเต นว, อญฺญมญฺเญ ตีณิ, นิสฺสเย เตรส, อุปนิสฺสเย อฏฺฐ, อาเสวเน ตีณิ, กมฺเม สตฺต, วิปาเก จตฺตาริ, อาหาเร สตฺต ฯเปฯ อวิคเต เตรส. (สํขิตฺตํ.)}

\prose{นเหตุยา จุทฺทส, นอารมฺมเณ จุทฺทส. (สํขิตฺตํ.)}

\prose{อารมฺมณ\-ปจฺจยา นเหตุยา ฉ. (สํขิตฺตํ.)}

\prose{นเหตุปจฺจยา อารมฺมเณ ฉ. (สํขิตฺตํ.)}

\prose{(ยถา กุสลตฺติเก ปญฺหาวารํ, เอวํ วิตฺถาเรตพฺพํ.)}

\csromanmatikaanchor{mtk.102}
\csromanmatikaanchor{mtk.103}
\csromantocmarknopage{subparagraph}{6. วิตกฺกตฺติก}{mtk.102}
\bookmark[dest={mtk.102},level=5]{6. วิตกฺกตฺติก}
\csromantocmarkref{subparagraph}{1. เหตุทุก}{mtk.103}
\bookmark[dest={mtk.103},level=5]{1. เหตุทุก}
\csromanmarkh{6. วิตกฺกตฺติก}\csromanmarkhii{1. เหตุทุก}\csromanheadernomark{6}{6. วิตกฺกตฺติก 1. เหตุทุก}
\csromanheader{2}{1. \textbf{ปฏิจฺจวาร ปจฺจยจตุกฺก}}
\proseitem{78}{สวิตกฺก\-สวิจารํ เหตุํ ธมฺมํ ปฏิจฺจ สวิตกฺก\-สวิจาโร เหตุ ธมฺโม อุปฺปชฺชติ เหตุปจฺจยา. (1)}

\prose{\csromanfolio{492}อวิตกฺก\-วิจารมตฺตํ เหตุํ ธมฺมํ ปฏิจฺจ อวิตกฺก\-วิจารมตฺโต เหตุ ธมฺโม อุปฺปชฺชติ เหตุปจฺจยา. (1)}

\prose{อวิตกฺก\-อวิจารํ เหตุํ ธมฺมํ ปฏิจฺจ อวิตกฺก\-อวิจาโร เหตุ ธมฺโม อุปฺปชฺชติ เหตุปจฺจยา. (1) (สํขิตฺตํ.)}

\proseitem{79}{เหตุยา ตีณิ, อารมฺมเณ ตีณิ ฯเปฯ อวิคเต ตีณิ. (สํขิตฺตํ.)}

\prose{นอธิปติยา ตีณิ, นปุเรชาเต ตีณิ, นปจฺฉาชาเต ตีณิ, นอาเสวเน ตีณิ, นวิปาเก ตีณิ, นวิปฺปยุตฺเต ตีณิ. (สํขิตฺตํ.)}

\prose{(ยถา สหชาตวารมฺปิ สมฺปยุตฺตวารมฺปิ, เอวํ วิตฺถาเรตพฺพํ.)}

\csromanheader{2}{7. \textbf{ปญฺหาวาร ปจฺจยจตุกฺก}}
\proseitem{80}{สวิตกฺก\-สวิจาโร เหตุ ธมฺโม สวิตกฺก\-สวิจารสฺส เหตุสฺส ธมฺมสฺส เหตุปจฺจเยน ปจฺจโย. (1)}

\prose{อวิตกฺก\-วิจารมตฺโต เหตุ ธมฺโม อวิตกฺก\-วิจารมตฺตสฺส เหตุสฺส ธมฺมสฺส เหตุปจฺจเยน ปจฺจโย. (1)}

\prose{อวิตกฺก\-อวิจาโร เหตุ ธมฺโม อวิตกฺก\-อวิจารสฺส เหตุสฺส ธมฺมสฺส เหตุปจฺจเยน ปจฺจโย. (1)}

\proseitem{81}{สวิตกฺก\-สวิจาโร เหตุ ธมฺโม สวิตกฺก\-สวิจารสฺส เหตุสฺส ธมฺมสฺส อารมฺมณ\-ปจฺจเยน ปจฺจโย.  สวิตกฺก\-สวิจาโร เหตุ ธมฺโม อวิตกฺก\-อวิจารสฺส เหตุสฺส ธมฺมสฺส อารมฺมณ\-ปจฺจเยน ปจฺจโย. (2)}

\prose{อวิตกฺก\-วิจารมตฺโต เหตุ ธมฺโม สวิตกฺก\-สวิจารสฺส เหตุสฺส ธมฺมสฺส อารมฺมณ\-ปจฺจเยน ปจฺจโย.  อวิตกฺก\-วิจารมตฺโต เหตุ ธมฺโม อวิตกฺก\-อวิจารสฺส เหตุสฺส ธมฺมสฺส อารมฺมณ\-ปจฺจเยน ปจฺจโย. (2)}

\prose{อวิตกฺก\-อวิจาโร เหตุ ธมฺโม อวิตกฺก\-อวิจารสฺส เหตุสฺส ธมฺมสฺส อารมฺมณ\-ปจฺจเยน ปจฺจโย.  อวิตกฺก\-อวิจาโร เหตุ ธมฺโม สวิตกฺก\-สวิจารสฺส เหตุสฺส ธมฺมสฺส อารมฺมณ\-ปจฺจเยน ปจฺจโย. (2)}

\proseitem{82}{สวิตกฺก\-สวิจาโร เหตุ ธมฺโม สวิตกฺก\-สวิจารสฺส เหตุสฺส ธมฺมสฺส อธิปติ\-ปจฺจเยน ปจฺจโย. (1)}

\csromanmarkh{6. วิตกฺกตฺติก}\csromanmarkhii{1. เหตุทุก}\csromanheadernomark{2}{6. วิตกฺกตฺติก 1. เหตุทุก}
\prose{\csromanfolio{493}อวิตกฺก\-วิจารมตฺโต เหตุ ธมฺโม อวิตกฺก\-วิจารมตฺตสฺส เหตุสฺส ธมฺมสฺส อธิปติ\-ปจฺจเยน ปจฺจโย.  อวิตกฺก\-วิจารมตฺโต เหตุ ธมฺโม สวิตกฺก\-สวิจารสฺส เหตุสฺส ธมฺมสฺส อธิปติ\-ปจฺจเยน ปจฺจโย. (2)}

\prose{อวิตกฺก\-อวิจาโร เหตุ ธมฺโม อวิตกฺก\-อวิจารสฺส เหตุสฺส ธมฺมสฺส อธิปติ\-ปจฺจเยน ปจฺจโย.  อวิตกฺก\-อวิจาโร เหตุ ธมฺโม สวิตกฺก\-สวิจารสฺส เหตุสฺส ธมฺมสฺส อธิปติ\-ปจฺจเยน ปจฺจโย. (2)}

\proseitem{83}{สวิตกฺก\-สวิจาโร เหตุ ธมฺโม สวิตกฺก\-สวิจารสฺส เหตุสฺส ธมฺมสฺส อนนฺตร\-ปจฺจเยน ปจฺจโย. ตีณิ.}

\prose{อวิตกฺก\-วิจารมตฺโต เหตุ ธมฺโม อวิตกฺก\-วิจารมตฺตสฺส เหตุสฺส ธมฺมสฺส อนนฺตร\-ปจฺจเยน ปจฺจโย. ตีณิ.}

\prose{อวิตกฺก\-อวิจาโร เหตุ ธมฺโม อวิตกฺก\-อวิจารสฺส เหตุสฺส ธมฺมสฺส อนนฺตร\-ปจฺจเยน ปจฺจโย. ตีณิ ฯเปฯ.}

\prose{สวิตกฺก\-สวิจาโร เหตุ ธมฺโม สวิตกฺก\-สวิจารสฺส เหตุสฺส ธมฺมสฺส อุปนิสฺสย\-ปจฺจเยน ปจฺจโย. นว. (สํขิตฺตํ.)}

\proseitem{84}{เหตุยา ตีณิ, อารมฺมเณ ฉ, อธิปติยา ปญฺจ, อนนฺตเร นว, สมนนฺตเร นว, สหชาเต ตีณิ, อญฺญมญฺเญ ตีณิ, นิสฺสเย ตีณิ, อุปนิสฺสเย นว, อาเสวเน ปญฺจ, วิปาเก ตีณิ ฯเปฯ อวิคเต ตีณิ.}

\prose{นเหตุยา นว, นอารมฺมเณ นว. (สํขิตฺตํ.)}

\prose{เหตุปจฺจยา นอารมฺมเณ ตีณิ. (สํขิตฺตํ.)}

\prose{นเหตุปจฺจยา อารมฺมเณ ฉ. (สํขิตฺตํ.) (ยถา กุสลตฺติเก ปญฺหาวารํ, เอวํ วิตฺถาเรตพฺพํ.)}

\csromanmarkh{นเหตุปท}\csromanmarkhii{1. ปฏิจฺจวาร}\csromanheadernomark{2}{\textbf{นเหตุปท 1. ปฏิจฺจวาร}}
\prose{ปจฺจยจตุกฺกเหตุ}

\proseitem{85}{สวิตกฺก\-สวิจารํ นเหตุํ ธมฺมํ ปฏิจฺจ สวิตกฺก\-สวิจาโร นเหตุ ธมฺโม อุปฺปชฺชติ เหตุปจฺจยา. (เอกํ.). สวิตกฺก\-สวิจารํ นเหตุํ ธมฺมํ ปฏิจฺจ อวิตกฺก\-วิจารมตฺโต นเหตุ ธมฺโม อุปฺปชฺชติ เหตุปจฺจยา. เทฺว. \csromanfolio{494}สวิตกฺก\-สวิจารํ นเหตุํ ธมฺมํ ปฏิจฺจ อวิตกฺก\-อวิจาโร นเหตุ ธมฺโม อุปฺปชฺชติ เหตุปจฺจยา. (ตีณิ.). สวิตกฺก\-สวิจารํ นเหตุํ ธมฺมํ ปฏิจฺจ สวิตกฺก\-สวิจาโร นเหตุ จ อวิตกฺก\-อวิจาโร นเหตุ จ ธมฺมา อุปฺปชฺชนฺติ เหตุปจฺจยา. (จตฺตาริ.). สวิตกฺก\-สวิจารํ นเหตุํ ธมฺมํ ปฏิจฺจ อวิตกฺก\-วิจารมตฺโต นเหตุ จ อวิตกฺก\-อวิจาโร นเหตุ จ ธมฺมา อุปฺปชฺชนฺติ เหตุปจฺจยา. (ปญฺจ.). สวิตกฺก\-สวิจารํ นเหตุํ ธมฺมํ ปฏิจฺจ สวิตกฺก\-สวิจาโร นเหตุ จ อวิตกฺก\-วิจารมตฺโต นเหตุ จ ธมฺมา อุปฺปชฺชนฺติ เหตุปจฺจยา. (ฉ.). สวิตกฺก\-สวิจารํ นเหตุํ ธมฺมํ ปฏิจฺจ สวิตกฺก\-สวิจาโร นเหตุ จ อวิตกฺก\-วิจารมตฺโต นเหตุ จ อวิตกฺก\-อวิจาโร นเหตุ จ ธมฺมา อุปฺปชฺชนฺติ เหตุปจฺจยา. (สตฺต.) อวิตกฺก\-วิจารมตฺตํ นเหตุํ ธมฺมํ ปฏิจฺจ อวิตกฺก\-วิจารมตฺโต นเหตุ ธมฺโม อุปฺปชฺชติ เหตุปจฺจยา. (เอกํ.). อวิตกฺก\-วิจารมตฺตํ นเหตุํ ธมฺมํ ปฏิจฺจ สวิตกฺก\-สวิจาโร นเหตุ ธมฺโม อุปฺปชฺชติ เหตุปจฺจยา. (เทฺว.). อวิตกฺก\-วิจารมตฺตํ นเหตุํ ธมฺมํ ปฏิจฺจ อวิตกฺก\-อวิจาโร นเหตุ ธมฺโม อุปฺปชฺชติ เหตุปจฺจยา. (ตีณิ.). อวิตกฺก\-วิจารมตฺตํ นเหตุํ ธมฺมํ ปฏิจฺจ สวิตกฺก\-สวิจาโร นเหตุ จ อวิตกฺก\-อวิจาโร นเหตุ จ ธมฺมา อุปฺปชฺชนฺติ เหตุปจฺจยา. (จตฺตาริ.). อวิตกฺก\-วิจารมตฺตํ นเหตุํ ธมฺมํ ปฏิจฺจ อวิตกฺก\-วิจารมตฺโต นเหตุ จ อวิตกฺก\-อวิจาโร นเหตุ จ ธมฺมา อุปฺปชฺชนฺติ เหตุปจฺจยา. (ปญฺจ.)}

\proseitem{87}{อวิตกฺก\-อวิจารํ นเหตุํ ธมฺมํ ปฏิจฺจ อวิตกฺก\-อวิจาโร นเหตุ ธมฺโม อุปฺปชฺชติ เหตุปจฺจยา. (เอกํ.). อวิตกฺก\-อวิจารํ นเหตุํ ธมฺมํ ปฏิจฺจ สวิตกฺก\-สวิจาโร นเหตุ ธมฺโม อุปฺปชฺชติ เหตุปจฺจยา. (เทฺว.). อวิตกฺก\-อวิจารํ นเหตุํ ธมฺมํ ปฏิจฺจ อวิตกฺก\-วิจารมตฺโต นเหตุ ธมฺโม อุปฺปชฺชติ เหตุปจฺจยา. (ตีณิ.). อวิตกฺก\-อวิจารํ นเหตุํ ธมฺมํ ปฏิจฺจ สวิตกฺก\-สวิจาโร นเหตุ จ อวิตกฺก\-อวิจาโร นเหตุ จ ธมฺมา อุปฺปชฺชนฺติ เหตุปจฺจยา. (จตฺตาริ.). อวิตกฺก\-อวิจารํ นเหตุํ ธมฺมํ ปฏิจฺจ อวิตกฺก\-วิจารมตฺโต นเหตุ จ อวิตกฺก\-อวิจาโร นเหตุ จ ธมฺมา อุปฺปชฺชนฺติ เหตุปจฺจยา. (ปญฺจ.). อวิตกฺก\-อวิจารํ นเหตุํ ธมฺมํ ปฏิจฺจ สวิตกฺก\-สวิจาโร นเหตุ จ อวิตกฺก\-วิจารมตฺโต นเหตุ จ ธมฺมา อุปฺปชฺชนฺติ เหตุปจฺจยา. (ฉ.). อวิตกฺก\-อวิจารํ นเหตุํ ธมฺมํ ปฏิจฺจ สวิตกฺก\-สวิจาโร นเหตุ จ}

\vspace{\tipitakaparskip}
\csromancenter{วิตกฺกตฺติก 1. เหตุทุก อวิตกฺก\-วิจารมตฺโต นเหตุ จ อวิตกฺก\-อวิจาโร นเหตุ จ ธมฺมา อุปฺปชฺชนฺติ เหตุปจฺจยา. (สตฺต.)}
\proseitem{88}{\csromanfolio{495}สวิตกฺก\-สวิจารํ นเหตุญฺจ อวิตกฺก\-อวิจารํ นเหตุญฺจ ธมฺมํ ปฏิจฺจ สวิตกฺก\-สวิจาโร นเหตุ ธมฺโม อุปฺปชฺชติ เหตุปจฺจยา. (เอกํ.) สวิตกฺก\-สวิจารํ นเหตุญฺจ อวิตกฺก\-อวิจารํ นเหตุญฺจ ธมฺมํ ปฏิจฺจ อวิตกฺก\-วิจารมตฺโต นเหตุ ธมฺโม อุปฺปชฺชติ เหตุปจฺจยา. (เทฺว.) สวิตกฺก\-สวิจารํ นเหตุญฺจ อวิตกฺก\-อวิจารํ นเหตุญฺจ ธมฺมํ ปฏิจฺจ อวิตกฺก\-อวิจาโร นเหตุ ธมฺโม อุปฺปชฺชติ เหตุปจฺจยา. (ตีณิ.) สวิตกฺก\-สวิจารํ นเหตุญฺจ อวิตกฺก\-อวิจารํ นเหตุญฺจ ธมฺมํ ปฏิจฺจ สวิตกฺก\-สวิจาโร นเหตุ จ อวิตกฺก\-อวิจาโร นเหตุ จ ธมฺมา อุปฺปชฺชนฺติ เหตุปจฺจยา. (จตฺตาริ.). สวิตกฺก\-สวิจารํ นเหตุญฺจ อวิตกฺก\-อวิจารํ นเหตุญฺจ ธมฺมํ ปฏิจฺจ อวิตกฺก\-วิจารมตฺโต นเหตุ จ อวิตกฺก\-อวิจาโร นเหตุ จ ธมฺมา อุปฺปชฺชนฺติ เหตุปจฺจยา. (ปญฺจ.). สวิตกฺก\-สวิจารํ นเหตุญฺจ อวิตกฺก\-อวิจารํ นเหตุญฺจ ธมฺมํ ปฏิจฺจ สวิตกฺก\-สวิจาโร นเหตุ จ อวิตกฺก\-วิจารมตฺโต นเหตุ จ ธมฺมา อุปฺปชฺชนฺติ เหตุปจฺจยา. (ฉ.). สวิตกฺก\-สวิจารํ นเหตุญฺจ อวิตกฺก\-อวิจารํ นเหตุญฺจ ธมฺมํ ปฏิจฺจ สวิตกฺก\-สวิจาโร นเหตุ จ อวิตกฺก\-วิจารมตฺโต นเหตุ จ อวิตกฺก\-อวิจาโร นเหตุ จ ธมฺมา อุปฺปชฺชนฺติ เหตุปจฺจยา. (สตฺต.)}

\proseitem{89}{อวิตกฺก\-วิจารมตฺตํ นเหตุญฺจ อวิตกฺก\-อวิจารํ นเหตุญฺจ ธมฺมํ ปฏิจฺจ สวิตกฺก\-สวิจาโร นเหตุ ธมฺโม อุปฺปชฺชติ เหตุปจฺจยา. (เอกํ.). อวิตกฺก\-วิจารมตฺตํ นเหตุญฺจ อวิตกฺก\-อวิจารํ นเหตุญฺจ ธมฺมํ ปฏิจฺจ อวิตกฺก\-วิจารมตฺโต นเหตุ ธมฺโม อุปฺปชฺชติ เหตุปจฺจยา. (เทฺว.). อวิตกฺก\-วิจารมตฺตํ นเหตุญฺจ อวิตกฺก\-อวิจารํ นเหตุญฺจ ธมฺมํ ปฏิจฺจ อวิตกฺก\-อวิจาโร นเหตุ ธมฺโม อุปฺปชฺชติ เหตุปจฺจยา. (ตีณิ.). อวิตกฺก\-วิจารมตฺตํ นเหตุญฺจ อวิตกฺก\-อวิจารํ นเหตุญฺจ ธมฺมํ ปฏิจฺจ สวิตกฺก\-สวิจาโร นเหตุ จ อวิตกฺก\-อวิจาโร นเหตุ จ ธมฺมา อุปฺปชฺชนฺติ เหตุปจฺจยา. (จตฺตาริ.). อวิตกฺก\-วิจารมตฺตํ นเหตุญฺจ อวิตกฺก\-อวิจารํ นเหตุญฺจ ธมฺมํ ปฏิจฺจ อวิตกฺก\-วิจารมตฺโต นเหตุ จ อวิตกฺก\-อวิจาโร นเหตุ จ ธมฺมา อุปฺปชฺชนฺติ เหตุปจฺจยา. ปญฺจ.}

\proseitem{90}{\csromanfolio{496}สวิตกฺก\-สวิจารํ นเหตุญฺจ อวิตกฺก\-วิจารมตฺตํ นเหตุญฺจ ธมฺมํ ปฏิจฺจ สวิตกฺก\-สวิจาโร นเหตุ ธมฺโม อุปฺปชฺชติ เหตุปจฺจยา. (เอกํ.). สวิตกฺก\-สวิจารํ นเหตุญฺจ อวิตกฺก\-วิจารมตฺตํ นเหตุญฺจ ธมฺมํ ปฏิจฺจ อวิตกฺก\-อวิจาโร นเหตุ ธมฺโม อุปฺปชฺชติ เหตุปจฺจยา. (เทฺว.). สวิตกฺก\-สวิจารํ นเหตุญฺจ อวิตกฺก\-วิจารมตฺตํ นเหตุญฺจ ธมฺมํ ปฏิจฺจ สวิตกฺก\-สวิจาโร นเหตุ จ อวิตกฺก\-อวิจาโร นเหตุ จ ธมฺมา อุปฺปชฺชนฺติ เหตุปจฺจยา. (ตีณิ.)}

\proseitem{91}{สวิตกฺก\-สวิจารํ นเหตุญฺจ อวิตกฺก\-วิจารมตฺตํ นเหตุญฺจ อวิตกฺก\-อวิจารํ นเหตุญฺจ ธมฺมํ ปฏิจฺจ สวิตกฺก\-สวิจาโร นเหตุ ธมฺโม อุปฺปชฺชติ เหตุปจฺจยา. (เอกํ.). สวิตกฺก\-สวิจารํ นเหตุญฺจ อวิตกฺก\-วิจารมตฺตํ นเหตุญฺจ อวิตกฺก\-อวิจารํ นเหตุญฺจ ธมฺมํ ปฏิจฺจ อวิตกฺก\-อวิจาโร นเหตุ ธมฺโม อุปฺปชฺชติ เหตุปจฺจยา. (เทฺว.). สวิตกฺก\-สวิจารํ นเหตุญฺจ อวิตกฺก\-วิจารมตฺตํ นเหตุญฺจ อวิตกฺก\-อวิจารํ นเหตุญฺจ ธมฺมํ ปฏิจฺจ สวิตกฺก\-สวิจาโร นเหตุ จ อวิตกฺก\-อวิจาโร นเหตุ จ ธมฺมา อุปฺปชฺชนฺติ เหตุปจฺจยา. (ตีณิ.)}

\proseitem{92}{สวิตกฺก\-สวิจารํ นเหตุํ ธมฺมํ ปฏิจฺจ สวิตกฺก\-สวิจาโร นเหตุ ธมฺโม อุปฺปชฺชติ อารมฺมณ\-ปจฺจยา.  สวิตกฺก\-สวิจารํ นเหตุํ ธมฺมํ ปฏิจฺจ อวิตกฺก\-วิจารมตฺโต นเหตุ ธมฺโม อุปฺปชฺชติ อารมฺมณ\-ปจฺจยา.  สวิตกฺก\-สวิจารํ นเหตุํ ธมฺมํ ปฏิจฺจ สวิตกฺก\-สวิจาโร นเหตุ จ อวิตกฺก\-วิจารมตฺโต นเหตุ จ ธมฺมา อุปฺปชฺชนฺติ อารมฺมณ\-ปจฺจยา. ตีณิ.}

\prose{อวิตกฺก\-วิจารมตฺตํ นเหตุํ ธมฺมํ ปฏิจฺจ อวิตกฺก\-วิจารมตฺโต นเหตุ ธมฺโม อุปฺปชฺชติ อารมฺมณ\-ปจฺจยา. จตฺตาริ.}

\prose{อวิตกฺก\-อวิจารํ นเหตุํ ธมฺมํ ปฏิจฺจ อวิตกฺก\-อวิจาโร นเหตุ ธมฺโม อุปฺปชฺชติ อารมฺมณ\-ปจฺจยา. ปญฺจ.}

\prose{สวิตกฺก\-สวิจารํ นเหตุญฺจ อวิตกฺก\-อวิจารํ นเหตุญฺจ ธมฺมํ ปฏิจฺจ อวิตกฺก\-อวิจาโร นเหตุ ธมฺโม อุปฺปชฺชติ อารมฺมณ\-ปจฺจยา. ตีณิ.}

\prose{อวิตกฺก\-สวิจาโร นเหตุญฺจ อวิตกฺก\-อวิจารํ นเหตุญฺจ ธมฺมํ ปฏิจฺจ สวิตกฺก\-สวิจาโร นเหตุ ธมฺโม อุปฺปชฺชติ อารมฺมณ\-ปจฺจยา. จตฺตาริ.}

\csromanmarkh{6. วิตกฺกตฺติก}\csromanmarkhii{1. เหตุทุก}\csromanheadernomark{2}{6. วิตกฺกตฺติก 1. เหตุทุก}
\prose{\csromanfolio{497}สวิตกฺก\-สวิจารํ นเหตุญฺจ อวิตกฺก\-วิจารมตฺตํ นเหตุญฺจ ธมฺมํ ปฏิจฺจ สวิตกฺก\-สวิจาโร นเหตุ ธมฺโม อุปฺปชฺชติ อารมฺมณ\-ปจฺจยา. เอกํ.}

\prose{สวิตกฺก\-สวิจารํ นเหตุญฺจ อวิตกฺก\-วิจารมตฺตํ นเหตุญฺจ อวิตกฺก\-อวิจารํ นเหตุญฺจ ธมฺมํ ปฏิจฺจ สวิตกฺก\-สวิจาโร นเหตุ ธมฺโม อุปฺปชฺชติ อารมฺมณ\-ปจฺจยา. เอกํ. (สํขิตฺตํ.)}

\proseitem{93}{เหตุยา สตฺตติํส, อารมฺมเณ เอกวีส, อธิปติยา เตวีส, อนนฺตเร เอกวีส, สหชาเต สตฺตติํส, อญฺญมญฺเญ อฏฺฐวีส, นิสฺสเย สตฺตติํส, อุปนิสฺสเย เอกวีส, ปุเรชาเต เอกาทส, อาเสวเน เอกาทส, กมฺเม สตฺตติํส, วิปาเก อาหาเร อินฺทฺริเย ฌาเน มคฺเค สตฺตติํส, สมฺปยุตฺเต เอกวีส, วิปฺปยุตฺเต สตฺตติํส ฯเปฯ อวิคเต สตฺตติํส. (สํขิตฺตํ.)}

\prose{นเหตุยา เตตฺติํส, นอารมฺมเณ สตฺต. (สํขิตฺตํ.)}

\prose{เหตุปจฺจยา นอารมฺมเณ สตฺต. (สํขิตฺตํ.)}

\prose{นเหตุปจฺจยา อารมฺมเณ จุทฺทส. (สํขิตฺตํ.)}

\prose{(สหชาตวารมฺปิ ฯเปฯ สมฺปยุตฺตวารมฺปิ วิตฺถาเรตพฺพํ.)}

\csromanmarkh{นเหตุปท}\csromanmarkhii{7. ปญฺหาวาร}\csromanheadernomark{2}{\textbf{นเหตุปท 7. ปญฺหาวาร}}
\prose{ปจฺจยจตุกฺกอารมฺมณ}

\proseitem{94}{สวิตกฺก\-สวิจาโร นเหตุ ธมฺโม สวิตกฺก\-สวิจารสฺส นเหตุสฺส ธมฺมสฺส อารมฺมณ\-ปจฺจเยน ปจฺจโย.  สวิตกฺก\-สวิจาโร นเหตุ ธมฺโม อวิตกฺก\-วิจารมตฺตสฺส นเหตุสฺส ธมฺมสฺส อารมฺมณ\-ปจฺจเยน ปจฺจโย.  สวิตกฺก\-สวิจาโร นเหตุ ธมฺโม อวิตกฺก\-อวิจารสฺส นเหตุสฺส ธมฺมสฺส อารมฺมณ\-ปจฺจเยน ปจฺจโย.  สวิตกฺก\-สวิจาโร นเหตุ ธมฺโม สวิตกฺก\-สวิจารสฺส นเหตุสฺส จ อวิตกฺก\-วิจารมตฺตสฺส นเหตุสฺส จ ธมฺมสฺส อารมฺมณ\-ปจฺจเยน ปจฺจโย. จตฺตาริ.}

\proseitem{95}{\csromanfolio{498}อวิตกฺก\-วิจารมตฺโต นเหตุ ธมฺโม อวิตกฺก\-วิจารมตฺตสฺส นเหตุสฺส ธมฺมสฺส อารมฺมณ\-ปจฺจเยน ปจฺจโย.  อวิตกฺก\-วิจารมตฺโต เนเหตุ ธมฺโม สวิตกฺก\-สวิจารสฺส นเหตุสฺส ธมฺมสฺส อารมฺมณ\-ปจฺจเยน ปจฺจโย.  อวิตกฺก\-วิจารมตฺโต นเหตุ ธมฺโม อวิตกฺก\-อวิจารสฺส นเหตุสฺส ธมฺมสฺส อารมฺมณ\-ปจฺจเยน ปจฺจโย.  อวิตกฺก\-วิจารมตฺโต นเหตุ ธมฺโม สวิตกฺก\-สวิจารสฺส นเหตุสฺส จ อวิตกฺก\-วิจารมตฺตสฺส นเหตุสฺส จ ธมฺมสฺส อารมฺมณ\-ปจฺจเยน ปจฺจโย. จตฺตาริ.}

\proseitem{96}{อวิตกฺก\-อวิจาโร นเหตุ ธมฺโม อวิตกฺก\-อวิจารสฺส นเหตุสฺส ธมฺมสฺส อารมฺมณ\-ปจฺจเยน ปจฺจโย.  อวิตกฺก\-อวิจาโร นเหตุ ธมฺโม สวิตกฺก\-สวิจารสฺส นเหตุสฺส ธมฺมสฺส อารมฺมณ\-ปจฺจเยน ปจฺจโย.  อวิตกฺก\-อวิจาโร นเหตุ ธมฺโม อวิตกฺก\-วิจารมตฺตสฺส นเหตุสฺส ธมฺมสฺส อารมฺมณ\-ปจฺจเยน ปจฺจโย.  อวิตกฺก\-อวิจาโร นเหตุ ธมฺโม อวิตกฺก\-วิจารมตฺตสฺส นเหตุสฺส จ อวิตกฺก\-อวิจารสฺส นเหตุสฺส จ ธมฺมสฺส อารมฺมณ\-ปจฺจเยน ปจฺจโย.  อวิตกฺก\-อวิจาโร นเหตุ ธมฺโม สวิตกฺก\-สวิจารสฺส นเหตุสฺส จ อวิตกฺก\-วิจารมตฺตสฺส นเหตุสฺส จ ธมฺมสฺส อารมฺมณ\-ปจฺจเยน ปจฺจโย. ปญฺจ.}

\proseitem{97}{อวิตกฺก\-วิจารมตฺโต นเหตุ จ อวิตกฺก\-อวิจาโร นเหตุ จ ธมฺมา สวิตกฺก\-สวิจารสฺส นเหตุสฺส ธมฺมสฺส อารมฺมณ\-ปจฺจเยน ปจฺจโย.  อวิตกฺก\-วิจารมตฺโต นเหตุ จ อวิตกฺก\-อวิจาโร นเหตุ จ ธมฺมา อวิตกฺก\-วิจารมตฺตสฺส นเหตุสฺส ธมฺมสฺส อารมฺมณ\-ปจฺจเยน ปจฺจโย.  อวิตกฺก\-วิจารมตฺโต นเหตุ จ อวิตกฺก\-อวิจาโร นเหตุ จ ธมฺมา อวิตกฺก\-อวิจารสฺส นเหตุสฺส ธมฺมสฺส อารมฺมณ\-ปจฺจเยน ปจฺจโย.  อวิตกฺก\-วิจารมตฺโต นเหตุ จ อวิตกฺก\-อวิจาโร นเหตุ จ ธมฺมา สวิตกฺก\-สวิจารสฺส นเหตุสฺส จ อวิตกฺก\-วิจารมตฺตสฺส นเหตุสฺส จ ธมฺมสฺส อารมฺมณ\-ปจฺจเยน ปจฺจโย. จตฺตาริ.}

\proseitem{98}{สวิตกฺก\-สวิจาโร นเหตุ จ อวิตกฺก\-วิจารมตฺโต นเหตุ จ ธมฺมา สวิตกฺก\-สวิจารสฺส นเหตุสฺส ธมฺมสฺส อารมฺมณ\-ปจฺจเยน ปจฺจโย.  สวิตกฺก\-สวิจาโร นเหตุ จ อวิตกฺก\-วิจารมตฺโต นเหตุ จ ธมฺมา อวิตกฺก\-วิจารมตฺตสฺส นเหตุสฺส ธมฺมสฺส อารมฺมณ\-ปจฺจเยน ปจฺจโย.  }

\vspace{\tipitakaparskip}
\csromancenter{วิตกฺกตฺติก 1. เหตุทุก สวิตกฺก\-สวิจาโร นเหตุ จ อวิตกฺก\-วิจารมตฺโต นเหตุ จ ธมฺมา อวิตกฺก\-อวิจารสฺส นเหตุสฺส ธมฺมสฺส อารมฺมณ\-ปจฺจเยน ปจฺจโย.  สวิตกฺก\-สวิจาโร นเหตุ จ อวิตกฺก\-วิจารมตฺโต นเหตุ จ ธมฺมา สวิตกฺก\-สวิจารสฺส นเหตุสฺส จ อวิตกฺก\-วิจารมตฺตสฺส นเหตุสฺส จ ธมฺมสฺส อารมฺมณ\-ปจฺจเยน ปจฺจโย. จตฺตาริ.}
\csromanheader{2}{\textbf{อธิปติ}}
\proseitem{99}{\csromanfolio{499}สวิตกฺก\-สวิจาโร นเหตุ ธมฺโม สวิตกฺก\-สวิจารสฺส นเหตุสฺส ธมฺมสฺส อธิปติ\-ปจฺจเยน ปจฺจโย. สตฺต.}

\prose{อวิตกฺก\-วิจารมตฺโต นเหตุ ธมฺโม อวิตกฺก\-วิจารมตฺตสฺส นเหตุสฺส ธมฺมสฺส อธิปติ\-ปจฺจเยน ปจฺจโย. ปญฺจ.}

\prose{อวิตกฺก\-อวิจาโร นเหตุ ธมฺโม อวิตกฺก\-อวิจารสฺส นเหตุสฺส ธมฺมสฺส อธิปติ\-ปจฺจเยน ปจฺจโย. ปญฺจ.}

\prose{อวิตกฺก\-วิจารมตฺโต นเหตุ จ อวิตกฺก\-อวิจาโร นเหตุ จ ธมฺมา สวิตกฺก\-สวิจารสฺส นเหตุสฺส ธมฺมสฺส อธิปติ\-ปจฺจเยน ปจฺจโย. ตีณิ. สวิตกฺก\-สวิจาโร นเหตุ จ อวิตกฺก\-วิจารมตฺโต นเหตุ จ ธมฺมา สวิตกฺก\-สวิจารสฺส นเหตุสฺส ธมฺมสฺส อธิปติ\-ปจฺจเยน ปจฺจโย. ตีณิ.}

\proseitem{100}{สวิตกฺก\-สวิจาโร นเหตุ ธมฺโม สวิตกฺก\-สวิจารสฺส นเหตุสฺส ธมฺมสฺส อนนฺตร\-ปจฺจเยน ปจฺจโย ฯเปฯ. สมนนฺตร\-ปจฺจเยน ปจฺจโย.}

\prose{สวิตกฺก\-สวิจาโร นเหตุ ธมฺโม สวิตกฺก\-สวิจารสฺส นเหตุสฺส ธมฺมสฺส อุปนิสฺสย\-ปจฺจเยน ปจฺจโย. (สํขิตฺตํ.)}

\proseitem{101}{อารมฺมเณ เอกวีส, อธิปติยา เตวีส, อนนฺตเร ปญฺจวีส, สมนนฺตเร ปญฺจวีส, สหชาเต ติํส, อญฺญมญฺเญ อฏฺฐวีส, นิสฺสเย ติํส, อุปนิสฺสเย ปญฺจวีส, ปุเรชาเต ปญฺจ, ปจฺฉาชาเต ปญฺจ, อาเสวเน เอกวีส, กมฺเม เอกาทส, วิปาเก เอกวีส, อาหาเร เอกาทส ฯเปฯ อวิคเต ติํส. (สํขิตฺตํ.)}

\prose{\csromanfolio{500}นเหตุยา ปญฺจติํส, นอารมฺมเณ ปญฺจติํส, นอธิปติยา ปญฺจติํส. (สํขิตฺตํ.)}

\prose{อารมฺมณ\-ปจฺจยา นเหตุยา เอกวีส. (สํขิตฺตํ.)}

\prose{นเหตุปจฺจยา อารมฺมเณ เอกวีส. (สํขิตฺตํ.)}

\prose{(ยถา กุสลตฺติเก ปญฺหาวารํ, เอวํ วิตฺถาเรตพฺพํ.)}

\csromanmatikaanchor{mtk.104}
\csromanmatikaanchor{mtk.105}
\csromantocmarknopage{subparagraph}{7. ปีติตฺติก}{mtk.104}
\bookmark[dest={mtk.104},level=5]{7. ปีติตฺติก}
\csromantocmarkref{subparagraph}{1. เหตุทุก}{mtk.105}
\bookmark[dest={mtk.105},level=5]{1. เหตุทุก}
\csromanmarkh{7. ปีติตฺติก}\csromanmarkhii{1. เหตุทุก}\csromanheadernomark{6}{7. \textbf{ปีติตฺติก 1. เหตุทุก}}
\csromanheader{2}{1. \textbf{ปฏิจฺจวาร ปจฺจยจตุกฺก}}
\proseitem{102}{ปีติสหคตํ เหตุํ ธมฺมํ ปฏิจฺจ ปีติสหคโต เหตุ ธมฺโม อุปฺปชฺชติ เหตุปจฺจยา. ตีณิ.}

\prose{สุขสหคตํ เหตุํ ธมฺมํ ปฏิจฺจ สุขสหคโต เหตุ ธมฺโม อุปฺปชฺชติ เหตุปจฺจยา. ตีณิ.}

\prose{อุเปกฺขา\-สหคตํ เหตุํ ธมฺมํ ปฏิจฺจ อุเปกฺขา\-สหคโต เหตุ ธมฺโม อุปฺปชฺชติ เหตุปจฺจยา. เอกํ.}

\prose{ปีติสหคตํ เหตุญฺจ สุขสหคตํ เหตุญฺจ ธมฺมํ ปฏิจฺจ ปีติสหคโต เหตุ ธมฺโม อุปฺปชฺชติ เหตุปจฺจยา. ตีณิ. (สํขิตฺตํ.)}

\proseitem{103}{เหตุยา ทส, อารมฺมเณ ทส ฯเปฯ อวิคเต ทส. (สํขิตฺตํ.)}

\prose{นอธิปติยา ทส, นปุเรชาเต ทส, นปจฺฉาชาเต ทส, นอาเสวเน ทส, นวิปาเก ทส, นวิปฺปยุตฺเต ทส. (สํขิตฺตํ.)}

\prose{(สหชาตวารมฺปิ ฯเปฯ สมฺปยุตฺตวารมฺปิ วิตฺถาเรตพฺพํ.)}

\csromanheader{2}{7. \textbf{ปญฺหาวาร ปจฺจยจตุกฺก}}
\csromanheader{2}{\textbf{เหตุ อารมฺมณาทิ}}
\proseitem{104}{ปีติสหคโต เหตุ ธมฺโม ปีติสหคตสฺส เหตุสฺส ธมฺมสฺส เหตุปจฺจเยน ปจฺจโย. ตีณิ.}

\csromanmarkh{7. ปีติตฺติก}\csromanmarkhii{1. เหตุทุก}\csromanheadernomark{2}{7. ปีติตฺติก 1. เหตุทุก}
\prose{\csromanfolio{501}สุขสหคโต เหตุ ธมฺโม สุขสหคตสฺส เหตุสฺส ธมฺมสฺส เหตุปจฺจเยน ปจฺจโย. ตีณิ.}

\prose{อุเปกฺขา\-สหคโต เหตุ ธมฺโม อุเปกฺขา\-สหคตสฺส เหตุสฺส ธมฺมสฺส เหตุปจฺจเยน ปจฺจโย. เอกํ.}

\prose{ปีติสหคโต เหตุ จ สุขสหคโต เหตุ จ ธมฺมา ปีติสหคตสฺส เหตุสฺส ธมฺมสฺส เหตุปจฺจเยน ปจฺจโย. ตีณิ.}

\proseitem{105}{ปีติสหคโต เหตุ ธมฺโม ปีติสหคตสฺส เหตุสฺส ธมฺมสฺส อารมฺมณ\-ปจฺจเยน ปจฺจโย.  ปีติสหคโต เหตุ ธมฺโม สุขสหคตสฺส เหตุสฺส ธมฺมสฺส อารมฺมณ\-ปจฺจเยน ปจฺจโย.  ปีติสหคโต เหตุ ธมฺโม อุเปกฺขา\-สหคตสฺส เหตุสฺส ธมฺมสฺส อารมฺมณ\-ปจฺจเยน ปจฺจโย.  ปีติสหคโต เหตุ ธมฺโม ปีติสหคตสฺส เหตุสฺส จ สุขสหคตสฺส เหตุสฺส จ ธมฺมสฺส อารมฺมณ\-ปจฺจเยน ปจฺจโย. จตฺตาริ.}

\prose{สุขสหคโต เหตุ ธมฺโม สุขสหคตสฺส เหตุสฺส ธมฺมสฺส อารมฺมณ\-ปจฺจเยน ปจฺจโย.  สุขสหคโต เหตุ ธมฺโม ปีติสหคตสฺส เหตุสฺส ธมฺมสฺส อารมฺมณ\-ปจฺจเยน ปจฺจโย.  สุขสหคโต เหตุ ธมฺโม อุเปกฺขสหคตสฺส เหตุสฺส ธมฺมสฺส อารมฺมณ\-ปจฺจเยน ปจฺจโย.  สุขสหคโต เหตุ ธมฺโม ปีติสหคตสฺส เหตุสฺส จ สุขสหคตสฺส เหตุสฺส จ ธมฺมสฺส อารมฺมณ\-ปจฺจเยน ปจฺจโย. จตฺตาริ.}

\prose{อุเปกฺขา\-สหคโต เหตุ ธมฺโม อุเปกฺขา\-สหคตสฺส เหตุสฺส ธมฺมสฺส อารมฺมณ\-ปจฺจเยน ปจฺจโย.  อุเปกฺขา\-สหคโต เหตุ ธมฺโม ปีติสหคตสฺส เหตุสฺส ธมฺมสฺส อารมฺมณ\-ปจฺจเยน ปจฺจโย.  อุเปกฺขา\-สหคโต เหตุ ธมฺโม สุขสหคตสฺส เหตุสฺส ธมฺมสฺส อารมฺมณ\-ปจฺจเยน ปจฺจโย.  อุเปกฺขา\-สหคโต เหตุ ธมฺโม ปีติสหคตสฺส เหตุสฺส จ สุขสหคตสฺส เหตุสฺส จ ธมฺมสฺส อารมฺมณ\-ปจฺจเยน ปจฺจโย. จตฺตาริ.}

\prose{ปีติสหคโต เหตุ จ สุขสหคโต เหตุ จ ธมฺมา ปีติสหคตสฺส เหตุสฺส ธมฺมสฺส อารมฺมณ\-ปจฺจเยน ปจฺจโย. จตฺตาริ.}

\prose{\csromanfolio{502}ปีติสหคโต เหตุ ธมฺโม ปีติสหคตสฺส เหตุสฺส ธมฺมสฺส อธิปติ\-ปจฺจเยน ปจฺจโย.}

\proseitem{106}{เหตุยา ทส, อารมฺมเณ โสฬส, อธิปติยา โสฬส, อนนฺตเร โสฬส, สมนนฺตเร โสฬส, สหชาเต ทส, อญฺญมญฺเญ ทส, นิสฺสเย ทส, อุปนิสฺสเย โสฬส ฯเปฯ วิปาเก ทส ฯเปฯ อวิคเต ทส. (สํขิตฺตํ.)}

\prose{นเหตุยา โสฬส, นอารมฺมเณ โสฬส. (สํขิตฺตํ.)}

\prose{เหตุปจฺจยา นอารมฺมเณ ทส. (สํขิตฺตํ.)}

\prose{นเหตุปจฺจยา อารมฺมเณ โสฬส. (สํขิตฺตํ.)}

\prose{(ยถา กุสลตฺติเก ปญฺหาวารํ, เอวํ วิตฺถาเรตพฺพํ.)}

\csromanmarkh{นเหตุปท}\csromanmarkhii{1. ปฏิจฺจวาร}\csromanheadernomark{2}{\textbf{นเหตุปท 1. ปฏิจฺจวาร}}
\prose{ปจฺจยจตุกฺกเหตุ อารมฺมณ}

\proseitem{107}{ปีติสหคตํ นเหตุํ ธมฺมํ ปฏิจฺจ ปีติสหคโต นเหตุ ธมฺโม อุปฺปชฺชติ เหตุปจฺจยา.  ปีติสหคตํ นเหตุํ ธมฺมํ ปฏิจฺจ สุขสหคโต นเหตุ ธมฺโม อุปฺปชฺชติ เหตุปจฺจยา.  ปีติสหคตํ นเหตุํ ธมฺมํ ปฏิจฺจ ปีติสหคโต นเหตุ จ สุขสหคโต นเหตุ จ ธมฺมา อุปฺปชฺชนฺติ เหตุปจฺจยา. (3)}

\prose{สุขสหคตํ นเหตุํ ธมฺมํ ปฏิจฺจ สุขสหคโต นเหตุ ธมฺโม อุปฺปชฺชติ เหตุปจฺจยา.  สุขสหคตํ นเหตุํ ธมฺมํ ปฏิจฺจ ปีติสหคโต นเหตุ ธมฺโม อุปฺปชฺชติ เหตุปจฺจยา.  สุขสหคตํ นเหตุํ ธมฺมํ ปฏิจฺจ ปีติสหคโต นเหตุ จ สุขสหคโต นเหตุ จ ธมฺมา อุปฺปชฺชนฺติ เหตุปจฺจยา. (3)}

\prose{อุเปกฺขา\-สหคตํ นเหตุํ ธมฺมํ ปฏิจฺจ อุเปกฺขา\-สหคโต นเหตุ ธมฺโม อุปฺปชฺชติ เหตุปจฺจยา. (1)}

\prose{ปีติสหคตํ นเหตุญฺจ สุขสหคตํ นเหตุญฺจ ธมฺมํ ปฏิจฺจ ปีติสหคโต นเหตุ ธมฺโม อุปฺปชฺชติ เหตุปจฺจยา.  ปีติสหคตํ นเหตุญฺจ สุขสหคตํ นเหตุญฺจ ธมฺมํ ปฏิจฺจ สุขสหคโต นเหตุ}

\vspace{\tipitakaparskip}
\csromancenter{ปีติตฺติก 1. เหตุทุก ธมฺโม อุปฺปชฺชติ เหตุปจฺจยา.  ปีติสหคตํ นเหตุญฺจ สุขสหคตํ นเหตุญฺจ ธมฺมํ ปฏิจฺจ ปีติสหคโต นเหตุ จ สุขสหคโต นเหตุ จ ธมฺมา อุปฺปชฺชนฺติ เหตุปจฺจยา. (3)}
\proseitem{108}{\csromanfolio{503}ปีติสหคตํ นเหตุํ ธมฺมํ ปฏิจฺจ ปีติสหคโต นเหตุ ธมฺโม อุปฺปชฺชติ อารมฺมณ\-ปจฺจยา. ตีณิ.}

\prose{สุขสหคตํ นเหตุํ ธมฺมํ ปฏิจฺจ สุขสหคโต นเหตุ ธมฺโม อุปฺปชฺชติ อารมฺมณ\-ปจฺจยา. ตีณิ.}

\prose{อุเปกฺขา\-สหคตํ นเหตุํ ธมฺมํ ปฏิจฺจ อุเปกฺขา\-สหคโต นเหตุ ธมฺโม อุปฺปชฺชติ อารมฺมณ\-ปจฺจยา. เอกํ.}

\prose{ปีติสหคตํ นเหตุญฺจ สุขสหคตํ นเหตุญฺจ ธมฺมํ ปฏิจฺจ ปีติสหคโต นเหตุ ธมฺโม อุปฺปชฺชติ อารมฺมณ\-ปจฺจยา. ตีณิ. (สํขิตฺตํ.)}

\proseitem{109}{เหตุยา ทส, อารมฺมเณ ทส, อธิปติยา ทส, อนนฺตเร ทส ฯเปฯ อวิคเต ทส. (สํขิตฺตํ.)}

\prose{นเหตุยา ทส, นอธิปติยา ทส. (สํขิตฺตํ. สหชาตวารมฺปิ ฯเปฯ สมฺปยุตฺตวารมฺปิ วิตฺถาเรตพฺพํ.)}

\csromanmarkh{นเหตุปท}\csromanmarkhii{7. ปญฺหาวาร}\csromanheadernomark{2}{\textbf{นเหตุปท 7. ปญฺหาวาร}}
\prose{ปจฺจยจตุกฺกอารมฺมณาทิ}

\proseitem{110}{ปีติสหคโต นเหตุ ธมฺโม ปีติสหคตสฺส นเหตุสฺส ธมฺมสฺส อารมฺมณ\-ปจฺจเยน ปจฺจโย.  ปีติสหคโต นเหตุ ธมฺโม สุขสหคตสฺส นเหตุสฺส ธมฺมสฺส อารมฺมณ\-ปจฺจเยน ปจฺจโย.  ปีติสหคโต นเหตุ ธมฺโม อุเปกฺขา\-สหคตสฺส นเหตุสฺส ธมฺมสฺส อารมฺมณ\-ปจฺจเยน ปจฺจโย.  ปีติสหคโต นเหตุ ธมฺโม ปีติสหคตสฺส นเหตุสฺส จ สุขสหคตสฺส นเหตุสฺส จ ธมฺมสฺส อารมฺมณ\-ปจฺจเยน ปจฺจโย. (4)}

\prose{สุขสหคโต นเหตุ ธมฺโม สุขสหคตสฺส นเหตุสฺส ธมฺมสฺส อารมฺมณ\-ปจฺจเยน ปจฺจโย.  สุขสหคโต นเหตุ ธมฺโม ปีติสหคตสฺส นเหตุสฺส ธมฺมสฺส อารมฺมณ\-ปจฺจเยน ปจฺจโย.   \csromanfolio{504}สุขสหคโต นเหตุ ธมฺโม อุเปกฺขา\-สหคตสฺส นเหตุสฺส ธมฺมสฺส อารมฺมณ\-ปจฺจเยน ปจฺจโย.  สุขสหคโต นเหตุ ธมฺโม ปีติสหคตสฺส นเหตุสฺส จ สุขสหคตสฺส นเหตุสฺส จ ธมฺมสฺส อารมฺมณ\-ปจฺจเยน ปจฺจโย. (4)}

\prose{อุเปกฺขา\-สหคโต นเหตุ ธมฺโม อุเปกฺขา\-สหคตสฺส นเหตุสฺส ธมฺมสฺส อารมฺมณ\-ปจฺจเยน ปจฺจโย.  อุเปกฺขา\-สหคโต นเหตุ ธมฺโม ปีติสหคตสฺส นเหตุสฺส ธมฺมสฺส อารมฺมณ\-ปจฺจเยน ปจฺจโย.  อุเปกฺขา\-สหคโต นเหตุ ธมฺโม สุขสหคตสฺส นเหตุสฺส ธมฺมสฺส อารมฺมณ\-ปจฺจเยน ปจฺจโย.  อุเปกฺขา\-สหคโต นเหตุ ธมฺโม ปีติสหคตสฺส นเหตุสฺส จ สุขสหคตสฺส นเหตุสฺส จ ธมฺมสฺส อารมฺมณ\-ปจฺจเยน ปจฺจโย. (4)}

\prose{ปีติสหคโต นเหตุ จ สุขสหคโต นเหตุ จ ธมฺมา ปีติสหคตสฺส นเหตุสฺส ธมฺมสฺส อารมฺมณ\-ปจฺจเยน ปจฺจโย. จตฺตาริ.}

\proseitem{111}{ปีติสหคโต นเหตุ ธมฺโม ปีติสหคตสฺส นเหตุสฺส ธมฺมสฺส อธิปติ\-ปจฺจเยน ปจฺจโย. จตฺตาริ.}

\prose{สุขสหคโต นเหตุ ธมฺโม สุขสหคตสฺส นเหตุสฺส ธมฺมสฺส อธิปติ\-ปจฺจเยน ปจฺจโย. จตฺตาริ.}

\prose{อุเปกฺขา\-สหคโต นเหตุ ธมฺโม อุเปกฺขา\-สหคตสฺส นเหตุสฺส ธมฺมสฺส อธิปติ\-ปจฺจเยน ปจฺจโย. จตฺตาริ.}

\prose{ปีติสหคโต นเหตุ จ สุขสหคโต นเหตุ จ ธมฺมา ปีติสหคตสฺส นเหตุสฺส ธมฺมสฺส อธิปติ\-ปจฺจเยน ปจฺจโย. จตฺตาริ.}

\proseitem{112}{ปีติสหคโต นเหตุ ธมฺโม ปีติสหคตสฺส นเหตุสฺส ธมฺมสฺส อนนฺตร\-ปจฺจเยน ปจฺจโย ฯเปฯ อุปนิสฺสย\-ปจฺจเยน ปจฺจโย. จตฺตาริ.}

\prose{สุขสหคโต นเหตุ ธมฺโม สุขสหคตสฺส นเหตุสฺส ธมฺมสฺส อุปนิสฺสย\-ปจฺจเยน ปจฺจโย. จตฺตาริ.}

\prose{อุเปกฺขา\-สหคโต นเหตุ ธมฺโม อุเปกฺขา\-สหคตสฺส นเหตุสฺส ธมฺมสฺส อุปนิสฺสย\-ปจฺจเยน ปจฺจโย. จตฺตาริ.}

\csromanmatikaanchor{mtk.106}
\csromanmatikaanchor{mtk.107}
\csromantocmarknopage{subparagraph}{8. ทสฺสเนนปหาตพฺพตฺติก}{mtk.106}
\bookmark[dest={mtk.106},level=5]{8. ทสฺสเนนปหาตพฺพตฺติก}
\csromantocmarkref{subparagraph}{1. เหตุทุก}{mtk.107}
\bookmark[dest={mtk.107},level=5]{1. เหตุทุก}
\csromanmarkh{8. ทสฺสเนน\-ปหาตพฺพ\-ตฺติก}\csromanmarkhii{1. เหตุทุก}\csromanheadernomark{6}{8. ทสฺสเนน\-ปหาตพฺพ\-ตฺติก 1. เหตุทุก}
\prose{\csromanfolio{505}ปีติสหคโต นเหตุ จ สุขสหคโต นเหตุ จ ธมฺมา ปีติสหคตสฺส นเหตุสฺส ธมฺมสฺส อุปนิสฺสย\-ปจฺจเยน ปจฺจโย. จตฺตาริ. (สํขิตฺตํ.)}

\proseitem{113}{อารมฺมเณ โสฬส, อธิปติยา โสฬส, สหชาเต ทส, อญฺญมญฺเญ ทส, นิสฺสเย ทส, อุปนิสฺสเย โสฬส, อาหาเร ทส ฯเปฯ อวิคเต ทส. (สํขิตฺตํ.)}

\prose{นเหตุยา โสฬส, นอารมฺมเณ โสฬส, นอธิปติยา โสฬส. (สํขิตฺตํ.)}

\prose{อารมฺมณ\-ปจฺจยา นเหตุยา โสฬส. (สํขิตฺตํ.)}

\prose{นเหตุปจฺจยา อารมฺมเณ โสฬส. (สํขิตฺตํ.)}

\prose{(ยถา กุสลตฺติเก ปญฺหาวารํ, เอวํ วิตฺถาเรตพฺพํ.)}

\csromanmarkh{8. ทสฺสเนน\-ปหาตพฺพ\-ตฺติก}\csromanmarkhii{1. เหตุทุก}\csromanheadernomark{2}{8. ทสฺสเนน\-ปหาตพฺพ\-ตฺติก 1. เหตุทุก}
\csromanheader{2}{1-7. ปฏิจฺจาทิวาร ปจฺจยจตุกฺก}
\proseitem{114}{ทสฺสเนน ปหาตพฺพํ เหตุํ ธมฺมํ ปฏิจฺจ ทสฺสเนน ปหาตพฺโพ เหตุ ธมฺโม อุปฺปชฺชติ เหตุปจฺจยา. (1)}

\prose{ภาวนาย ปหาตพฺพํ เหตุํ ธมฺมํ ปฏิจฺจ ภาวนาย ปหาตพฺโพ เหตุ ธมฺโม อุปฺปชฺชติ เหตุปจฺจยา. (1)}

\prose{เนวทสฺสเนน นภาวนาย ปหาตพฺพํ เหตุํ ธมฺมํ ปฏิจฺจ เนวทสฺสเนน นภาวนาย ปหาตพฺโพ เหตุ ธมฺโม อุปฺปชฺชติ เหตุปจฺจยา. (สํขิตฺตํ.)}

\proseitem{115}{เหตุยา ตีณิ, (สพฺพตฺถ ตีณิ.) วิปาเก เอกํ ฯเปฯ อวิคเต ตีณิ.}

\prose{นอธิปติยา ตีณิ, นปุเรชาเต ตีณิ ฯเปฯ นวิปฺปยุตฺเต ตีณิ. (สํขิตฺตํ. สหชาตวารมฺปิ ฯเปฯ สมฺปยุตฺตวารมฺปิ วิตฺถาเรตพฺพํ.)}

\proseitem{116}{ทสฺสเนน ปหาตพฺโพ เหตุ ธมฺโม ทสฺสเนน ปหาตพฺพสฺส เหตุสฺส ธมฺมสฺส เหตุปจฺจเยน ปจฺจโย. ตีณิ.}

\prose{\csromanfolio{506}ทสฺสเนน ปหาตพฺโพ เหตุ ธมฺโม ทสฺสเนน ปหาตพฺพสฺส เหตุสฺส ธมฺมสฺส อารมฺมณ\-ปจฺจเยน ปจฺจโย. (สํขิตฺตํ.)}

\proseitem{117}{เหตุยา ตีณิ, อารมฺมเณ อฏฺฐ, อธิปติยา ฉ, อนนฺตเร ปญฺจ, สมนนฺตเร ปญฺจ, สหชาเต ตีณิ, อญฺญมญฺเญ ตีณิ, นิสฺสเย ตีณิ, อุปนิสฺสเย อฏฺฐ, อาเสวเน ตีณิ, วิปาเก เอกํ, อินฺทฺริเย เอกํ, มคฺเค เอกํ ฯเปฯ อวิคเต ตีณิ. (สํขิตฺตํ.)}

\prose{นเหตุยา อฏฺฐ, นอารมฺมเณ อฏฺฐ. (สํขิตฺตํ.)}

\prose{เหตุปจฺจยา นอารมฺมเณ ตีณิ. (สํขิตฺตํ.)}

\prose{นเหตุปจฺจยา อารมฺมเณ อฏฺฐ. (สํขิตฺตํ.)}

\prose{(ยถา กุสลตฺติเก ปญฺหาวารํ, เอวํ วิตฺถาเรตพฺพํ.)}

\csromanmarkh{นเหตุปท}\csromanmarkhii{1. ปฏิจฺจวาร}\csromanheadernomark{2}{\textbf{นเหตุปท 1. ปฏิจฺจวาร}}
\prose{ปจฺจยจตุกฺกเหตุ อารมฺมณ}

\proseitem{118}{ทสฺสเนน ปหาตพฺพํ นเหตุํ ธมฺมํ ปฏิจฺจ ทสฺสเนน ปหาตพฺโพ นเหตุ ธมฺโม อุปฺปชฺชติ เหตุปจฺจยา.  ทสฺสเนน ปหาตพฺพํ นเหตุํ ธมฺมํ ปฏิจฺจ เนวทสฺสเนน นภาวนาย ปหาตพฺโพ นเหตุ ธมฺโม อุปฺปชฺชติ เหตุปจฺจยา.  ทสฺสเนน ปหาตพฺพํ นเหตุํ ธมฺมํ ปฏิจฺจ ทสฺสเนน ปหาตพฺโพ นเหตุ จ เนวทสฺสเนน นภาวนาย ปหาตพฺโพ นเหตุ จ ธมฺมา อุปฺปชฺชนฺติ เหตุปจฺจยา. ตีณิ.}

\prose{ภาวนาย ปหาตพฺพํ นเหตุํ ธมฺมํ ปฏิจฺจ ภาวนาย ปหาตพฺโพ นเหตุ ธมฺโม อุปฺปชฺชติ เหตุปจฺจยา.  ภาวนาย ปหาตพฺพํ นเหตุํ ธมฺมํ ปฏิจฺจ เนวทสฺสเนน นภาวนาย ปหาตพฺโพ นเหตุ ธมฺโม อุปฺปชฺชติ เหตุปจฺจยา.  ภาวนาย ปหาตพฺพํ นเหตุํ ธมฺมํ ปฏิจฺจ ภาวนาย ปหาตพฺโพ นเหตุ จ เนวทสฺสเนน นภาวนาย ปหาตพฺโพ นเหตุ จ ธมฺมา อุปฺปชฺชนฺติ เหตุปจฺจยา. ตีณิ.}

\prose{เนวทสฺสเนน นภาวนาย ปหาตพฺพํ นเหตุํ ธมฺมํ ปฏิจฺจ เนวทสฺสเนน นภาวนาย ปหาตพฺโพ นเหตุ ธมฺโม อุปฺปชฺชติ เหตุปจฺจยา. (1)}

\csromanmarkh{8. ทสฺสเนน\-ปหาตพฺพ\-ตฺติก}\csromanmarkhii{1. เหตุทุก}\csromanheadernomark{2}{8. ทสฺสเนน\-ปหาตพฺพ\-ตฺติก 1. เหตุทุก}
\prose{\csromanfolio{507}ทสฺสเนน ปหาตพฺพํ นเหตุญฺจ เนวทสฺสเนน นภาวนาย ปหาตพฺพํ นเหตุญฺจ ธมฺมํ ปฏิจฺจ เนวทสฺสเนน นภาวนาย ปหาตพฺโพ นเหตุ ธมฺโม อุปฺปชฺชติ เหตุปจฺจยา. (1)}

\prose{ภาวนาย ปหาตพฺพํ นเหตุญฺจ เนวทสฺสเนน นภาวนาย ปหาตพฺพํ นเหตุญฺจ ธมฺมํ ปฏิจฺจ เนวทสฺสเนน นภาวนาย ปหาตพฺโพ นเหตุ ธมฺโม อุปฺปชฺชติ เหตุปจฺจยา. (1)}

\prose{ทสฺสเนน ปหาตพฺพํ นเหตุํ ธมฺมํ ปฏิจฺจ ทสฺสเนน ปหาตพฺโพ นเหตุ ธมฺโม อุปฺปชฺชติ อารมฺมณ\-ปจฺจยา. (1)}

\prose{ภาวนาย ปหาตพฺพํ นเหตุํ ธมฺมํ ปฏิจฺจ ภาวนาย ปหาตพฺโพ นเหตุ ธมฺโม อุปฺปชฺชติ อารมฺมณ\-ปจฺจยา. (1)}

\prose{เนวทสฺสเนน นภาวนาย ปหาตพฺพํ นเหตุํ ธมฺมํ ปฏิจฺจ เนวทสฺสเนน นภาวนาย ปหาตพฺโพ นเหตุ ธมฺโม อุปฺปชฺชติ อารมฺมณ\-ปจฺจยา. (สํขิตฺตํ.)}

\proseitem{119}{เหตุยา นว, อารมฺมเณ ตีณิ, อธิปติยา นว ฯเปฯ อวิคเต นว. (สํขิตฺตํ.)}

\prose{นเหตุ นอารมฺมณ}

\proseitem{120}{เนวทสฺสเนน นภาวนาย ปหาตพฺพํ นเหตุํ ธมฺมํ ปฏิจฺจ เนวทสฺสเนน นภาวนาย ปหาตพฺโพ นเหตุ ธมฺโม อุปฺปชฺชติ นเหตุปจฺจยา.}

\prose{ทสฺสเนน ปหาตพฺพํ นเหตุํ ธมฺมํ ปฏิจฺจ เนวทสฺสเนน นภาวนาย ปหาตพฺโพ นเหตุ ธมฺโม อุปฺปชฺชติ นอารมฺมณ\-ปจฺจยา. (1)}

\prose{ภาวนาย ปหาตพฺพํ นเหตุํ ธมฺมํ ปฏิจฺจ เนวทสฺสเนน นภาวนาย ปหาตพฺโพ นเหตุ ธมฺโม อุปฺปชฺชติ นอารมฺมณ\-ปจฺจยา. (1)}

\prose{เนวทสฺสเนน นภาวนาย ปหาตพฺพํ นเหตุํ ธมฺมํ ปฏิจฺจ เนวทสฺสเนน นภาวนาย ปหาตพฺโพ นเหตุ ธมฺโม อุปฺปชฺชติ นอารมฺมณ\-ปจฺจยา. (1)}

\prose{ทสฺสเนน ปหาตพฺพํ นเหตุญฺจ เนวทสฺสเนน นภาวนาย ปหาตพฺพํ นเหตุญฺจ ธมฺมํ ปฏิจฺจ เนวทสฺสเนน นภาวนาย ปหาตพฺโพ นเหตุ ธมฺโม อุปฺปชฺชติ นอารมฺมณ\-ปจฺจยา. (1)}

\prose{\csromanfolio{508}ภาวนาย ปหาตพฺพํ นเหตุญฺจ เนวทสฺสเนน นภาวนาย ปหาตพฺพํ นเหตุญฺจ ธมฺมํ ปฏิจฺจ เนวทสฺสเนน นภาวนาย ปหาตพฺโพ นเหตุ ธมฺโม อุปฺปชฺชติ นอารมฺมณ\-ปจฺจยา. (1) (สํขิตฺตํ.)}

\proseitem{121}{นเหตุยา เอกํ, นอารมฺมเณ ปญฺจ, นอธิปติยา นว ฯเปฯ โนวิคเต ปญฺจ. (สํขิตฺตํ.)}

\prose{เหตุปจฺจยา นอารมฺมเณ ปญฺจ. (สํขิตฺตํ.)}

\prose{นเหตุปจฺจยา อารมฺมเณ เอกํ. (สํขิตฺตํ.)}

\prose{(สหชาตวารมฺปิ ปจฺจยวารมฺปิ นิสฺสยวารมฺปิ สํสฏฺฐ\-วารมฺปิ สมฺปยุตฺตวารมฺปิ ปฏิจฺจ\-วาร\-สทิสํ วิตฺถาเรตพฺพํ.)}

\csromanheader{2}{7. \textbf{ปญฺหาวาร ปจฺจยจตุกฺก}}
\proseitem{122}{ทสฺสเนน ปหาตพฺโพ นเหตุ ธมฺโม ทสฺสเนน ปหาตพฺพสฺส นเหตุสฺส ธมฺมสฺส อารมฺมณ\-ปจฺจเยน ปจฺจโย.  ทสฺสเนน ปหาตพฺโพ เหตุ ธมฺโม เนวทสฺสเนน นภาวนาย ปหาตพฺพสฺส นเหตุสฺส ธมฺมสฺส อารมฺมณ\-ปจฺจเยน ปจฺจโย. เทฺว.}

\prose{ภาวนาย ปหาตพฺโพ นเหตุ ธมฺโม ภาวนาย ปหาตพฺพสฺส นเหตุสฺส ธมฺมสฺส อารมฺมณ\-ปจฺจเยน ปจฺจโย. ตีณิ.}

\prose{เนวทสฺสเนน นภาวนาย ปหาตพฺโพ นเหตุ ธมฺโม เนวทสฺสเนน นภาวนาย ปหาตพฺพสฺส นเหตุสฺส ธมฺมสฺส อารมฺมณ\-ปจฺจเยน ปจฺจโย. ตีณิ. (สํขิตฺตํ.)}

\proseitem{123}{อารมฺมเณ อฏฺฐ, อธิปติยา ทส ฯเปฯ อวิคเต เตรส. (สํขิตฺตํ.)}

\csromanmatikaanchor{mtk.108}
\csromanmatikaanchor{mtk.109}
\csromantocmarknopage{subparagraph}{9. ทสฺสเนนปหาตพฺพเหตุกตฺติก}{mtk.108}
\bookmark[dest={mtk.108},level=5]{9. ทสฺสเนนปหาตพฺพเหตุกตฺติก}
\csromantocmarkref{subparagraph}{1. เหตุทุก}{mtk.109}
\bookmark[dest={mtk.109},level=5]{1. เหตุทุก}
\csromanmarkh{9. ทสฺสเนน\-ปหาตพฺพ\-เหตุ\-กตฺติก}\csromanmarkhii{1. เหตุทุก}\csromanheadernomark{6}{9. ทสฺสเนน\-ปหาตพฺพ\-เหตุ\-กตฺติก 1. เหตุทุก}
\prose{\csromanfolio{509}นเหตุยา จุทฺทส, นอารมฺมเณ จุทฺทส. (สํขิตฺตํ.)}

\prose{อารมฺมณ\-ปจฺจยา นเหตุยา อฏฺฐ. (สํขิตฺตํ.)}

\prose{นเหตุปจฺจยา อารมฺมเณ อฏฺฐ. (สํขิตฺตํ.)}

\prose{(ยถา กุสลตฺติเก ปญฺหาวารํ, เอวํ วิตฺถาเรตพฺพํ.)}

\csromanmarkh{9. ทสฺสเนน\-ปหาตพฺพ\-เหตุ\-กตฺติก}\csromanmarkhii{1. เหตุทุก}\csromanheadernomark{2}{9. ทสฺสเนน\-ปหาตพฺพ\-เหตุ\-กตฺติก 1. เหตุทุก}
\csromanheader{2}{1-7. ปฏิจฺจาทิวาร ปจฺจยจตุกฺกเหตุ}
\proseitem{124}{ทสฺสเนน ปหาตพฺพ\-เหตุกํ เหตุํ ธมฺมํ ปฏิจฺจ ทสฺสเนน ปหาตพฺพ\-เหตุโก เหตุ ธมฺโม อุปฺปชฺชติ เหตุปจฺจยา. (1)}

\prose{ภาวนาย ปหาตพฺพ\-เหตุกํ เหตุํ ธมฺมํ ปฏิจฺจ ภาวนาย ปหาตพฺพ\-เหตุโก เหตุ ธมฺโม อุปฺปชฺชติ เหตุปจฺจยา. (1)}

\prose{เนวทสฺสเนน นภาวนาย ปหาตพฺพ\-เหตุกํ เหตุํ ธมฺมํ ปฏิจฺจ เนวทสฺสเนน นภาวนาย ปหาตพฺพ\-เหตุโก เหตุ ธมฺโม อุปฺปชฺชติ เหตุปจฺจยา. (1) (สํขิตฺตํ.)}

\proseitem{125}{เหตุยา ตีณิ, อารมฺมเณ ตีณิ ฯเปฯ อวิคเต ตีณิ. (สํขิตฺตํ.)}

\prose{นอธิปติยา ตีณิ, นปุเรชาเต ตีณิ ฯเปฯ นวิปฺปยุตฺเต ตีณิ. (สํขิตฺตํ. สหชาตวารมฺปิ ฯเปฯ สมฺปยุตฺตวารมฺปิ ปฏิจฺจ\-วาร\-สทิสํ วิตฺถาเรตพฺพํ.)}

\proseitem{126}{ทสฺสเนน ปหาตพฺพ\-เหตุโก เหตุ ธมฺโม ทสฺสเนน ปหาตพฺพ\-เหตุกสฺส เหตุสฺส ธมฺมสฺส เหตุปจฺจเยน ปจฺจโย. (สํขิตฺตํ.)}

\prose{เหตุยา ตีณิ, อารมฺมเณ อฏฺฐ, อธิปติยา ฉ ฯเปฯ อวิคเต ตีณิ. (สํขิตฺตํ.)}

\prose{\csromanfolio{510}นเหตุยา อฏฺฐ, นอารมฺมเณ อฏฺฐ. (สํขิตฺตํ.)}

\prose{เหตุปจฺจยา นอารมฺมเณ ตีณิ. (สํขิตฺตํ.)}

\prose{นเหตุปจฺจยา อารมฺมเณ อฏฺฐ. (สํขิตฺตํ.)}

\prose{(ยถา กุสลตฺติเก ปญฺหาวารํ, เอวํ วิตฺถาเรตพฺพํ.)}

\csromanheader{2}{\textbf{นเหตุปท 1-7. ปฏิจฺจาทิวาร}}
\prose{ปจฺจยจตุกฺกเหตุ}

\proseitem{127}{ทสฺสเนน ปหาตพฺพ\-เหตุกํ นเหตุํ ธมฺมํ ปฏิจฺจ ทสฺสเนน ปหาตพฺพ\-เหตุโก นเหตุ ธมฺโม อุปฺปชฺชติ เหตุปจฺจยา. ตีณิ.}

\prose{ภาวนาย ปหาตพฺพ\-เหตุกํ นเหตุํ ธมฺมํ ปฏิจฺจ ภาวนาย ปหาตพฺพ\-เหตุโก นเหตุ ธมฺโม อุปฺปชฺชติ เหตุปจฺจยา. ตีณิ.}

\prose{เนวทสฺสเนน นภาวนาย ปหาตพฺพ\-เหตุกํ นเหตุํ ธมฺมํ ปฏิจฺจ เนวทสฺสเนน นภาวนาย ปหาตพฺพ\-เหตุโก นเหตุ ธมฺโม อุปฺปชฺชติ เหตุปจฺจยา. ตีณิ. (สํขิตฺตํ.)}

\prose{เหตุยา นว, อารมฺมเณ ตีณิ ฯเปฯ อวิคเต นว. (สํขิตฺตํ.)}

\prose{ปจฺจนีย\-น\-เหตุ}

\proseitem{128}{เนวทสฺสเนน นภาวนาย ปหาตพฺพ\-เหตุกํ นเหตุํ ธมฺมํ ปฏิจฺจ เนวทสฺสเนน นภาวนาย ปหาตพฺพ\-เหตุโก นเหตุ ธมฺโม อุปฺปชฺชติ นเหตุปจฺจยา. (1) (สํขิตฺตํ.)}

\proseitem{129}{นเหตุยา เอกํ, นอารมฺมเณ ปญฺจ ฯเปฯ โนวิคเต ปญฺจ. (สํขิตฺตํ.)}

\prose{เหตุปจฺจยา นอารมฺมเณ ปญฺจ. (สํขิตฺตํ.)}

\prose{นเหตุปจฺจยา อารมฺมเณ เอกํ. (สํขิตฺตํ.)}

\prose{(สหชาตวารมฺปิ ฯเปฯ สมฺปยุตฺตวารมฺปิ ปฏิจฺจ\-วาร\-สทิสํ วิตฺถาเรตพฺพํ.)}

\csromanmarkh{10. อาจยคามิตฺติก}\csromanmarkhii{1. เหตุทุก อารมฺมณ}\csromanheadernomark{2}{10. อาจยคามิตฺติก 1. เหตุทุก \textbf{อารมฺมณ}}
\proseitem{130}{\csromanfolio{511}ทสฺสเนน ปหาตพฺพ\-เหตุโก นเหตุ ธมฺโม ทสฺสเนน ปหาตพฺพ\-เหตุกสฺส นเหตุสฺส ธมฺมสฺส อารมฺมณ\-ปจฺจเยน ปจฺจโย. เทฺว.}

\prose{ภาวนาย ปหาตพฺพ\-เหตุโก นเหตุ ธมฺโม ภาวนาย ปหาตพฺพ\-เหตุกสฺส นเหตุสฺส ธมฺมสฺส อารมฺมณ\-ปจฺจเยน ปจฺจโย. ตีณิ.}

\prose{เนวทสฺสเนน นภาวนาย ปหาตพฺพ\-เหตุโก นเหตุ ธมฺโม เนวทสฺสเนน นภาวนาย ปหาตพฺพ\-เหตุกสฺส นเหตุสฺส ธมฺมสฺส อารมฺมณ\-ปจฺจเยน ปจฺจโย. ตีณิ. (สํขิตฺตํ.)}

\proseitem{131}{อารมฺมเณ อฏฺฐ, อธิปติยา ทส ฯเปฯ อวิคเต เตรส. (สํขิตฺตํ.)}

\prose{นเหตุยา จุทฺทส นอารมฺมเณ จุทฺทส. (สํขิตฺตํ.)}

\prose{อารมฺมณ\-ปจฺจยา นเหตุยา อฏฺฐ. (สํขิตฺตํ.)}

\prose{นเหตุปจฺจยา อารมฺมเณ อฏฺฐ. (สํขิตฺตํ.)}

\prose{(ยถา กุสลตฺติเก ปญฺหาวารํ, เอวํ วิตฺถาเรตพฺพํ.)}

\csromanmatikaanchor{mtk.110}
\csromanmatikaanchor{mtk.111}
\csromantocmarknopage{subparagraph}{10. อาจยคามิตฺติก}{mtk.110}
\bookmark[dest={mtk.110},level=5]{10. อาจยคามิตฺติก}
\csromantocmarkref{subparagraph}{1. เหตุทุก}{mtk.111}
\bookmark[dest={mtk.111},level=5]{1. เหตุทุก}
\csromanmarkh{10. อาจยคามิตฺติก}\csromanmarkhii{1. เหตุทุก}\csromanheadernomark{6}{10. อาจยคามิตฺติก 1. เหตุทุก}
\csromanheader{2}{1-7. ปฏิจฺจาทิวาร ปจฺจยจตุกฺกเหตุ}
\proseitem{132}{อาจยคามิํ เหตุํ ธมฺมํ ปฏิจฺจ อาจยคามี เหตุ ธมฺโม อุปฺปชฺชติ เหตุปจฺจยา. (1)}

\prose{อปจยคามิํ เหตุํ ธมฺมํ ปฏิจฺจ อปจยคามี เหตุ ธมฺโม อุปฺปชฺชติ เหตุปจฺจยา. (1)}

\prose{เนวา\-จยคามิ\-นา\-ปจยคามิํ เหตุํ ธมฺมํ ปฏิจฺจ เนวา\-จยคามิ\-นา\-ปจยคามี เหตุ ธมฺโม อุปฺปชฺชติ เหตุปจฺจยา. (1) (สํขิตฺตํ.)}

\proseitem{133}{เหตุยา ตีณิ, อารมฺมเณ ตีณิ, อธิปติยา ตีณิ, (สพฺพตฺถ ตีณิ.)}

\prose{\csromanfolio{512}นอธิปติยา ตีณิ, นอาเสวเน เทฺว, นวิปาเก ตีณิ, นวิปฺปยุตฺเต ตีณิ. (สํขิตฺตํ.)}

\prose{เหตุปจฺจยา นอธิปติยา ตีณิ. (สํขิตฺตํ.)}

\prose{นอธิปติ\-ปจฺจยา เหตุยา ตีณิ. (สํขิตฺตํ.)}

\prose{(สหชาตวารมฺปิ ฯเปฯ สมฺปยุตฺตวารมฺปิ ปฏิจฺจ\-วาร\-สทิสํ วิตฺถาเรตพฺพํ.)}

\csromanheader{2}{\textbf{เหตุ อารมฺมณาทิ}}
\proseitem{134}{อาจยคามี เหตุ ธมฺโม อาจยคามิสฺส เหตุสฺส ธมฺมสฺส เหตุปจฺจเยน ปจฺจโย. (1)}

\prose{อปจยคามี เหตุ ธมฺโม อปจยคามิสฺส เหตุสฺส ธมฺมสฺส เหตุปจฺจเยน ปจฺจโย. (1)}

\prose{เนวา\-จยคามิ\-นา\-ปจยคามี เหตุ ธมฺโม เนวา\-จยคามิ\-นา\-ปจยคามิสฺส เหตุสฺส ธมฺมสฺส เหตุปจฺจเยน ปจฺจโย. (1)}

\prose{อาจยคามี เหตุ ธมฺโม อาจยคามิสฺส เหตุสฺส ธมฺมสฺส อารมฺมณ\-ปจฺจเยน ปจฺจโย.  อาจยคามี เหตุ ธมฺโม เนวา\-จยคามิ\-นา\-ปจยคามิสฺส เหตุสฺส ธมฺมสฺส อารมฺมณ\-ปจฺจเยน ปจฺจโย. (2)}

\prose{อปจยคามี เหตุ ธมฺโม อาจยคามิสฺส เหตุสฺส ธมฺมสฺส อารมฺมณ\-ปจฺจเยน ปจฺจโย.  อปจยคามี เหตุ ธมฺโม เนวา\-จยคามิ\-นา\-ปจยคามิสฺส เหตุสฺส ธมฺมสฺส อารมฺมณ\-ปจฺจเยน ปจฺจโย. (2)}

\prose{เนวา\-จยคามิ\-นา\-ปจยคามี เหตุ ธมฺโม เนวา\-จยคามิ\-นา\-ปจยคามิสฺส เหตุสฺส ธมฺมสฺส อารมฺมณ\-ปจฺจเยน ปจฺจโย.  เนวา\-จยคามิ\-นา\-ปจยคามี เหตุ ธมฺโม อาจยคามิสฺส เหตุสฺส ธมฺมสฺส อารมฺมณ\-ปจฺจเยน ปจฺจโย. (2)}

\proseitem{135}{อาจยคามี เหตุ ธมฺโม อาจยคามิสฺส เหตุสฺส ธมฺมสฺส อธิปติ\-ปจฺจเยน ปจฺจโย. อารมฺมณา\-ธิปติ, สหชาตา\-ธิปติ. (1)}

\prose{อปจยคามี เหตุ ธมฺโม อปจยคามิสฺส เหตุสฺส ธมฺมสฺส อธิปติ\-ปจฺจเยน ปจฺจโย. สหชาตา\-ธิปติ.  อปจยคามี เหตุ ธมฺโม อาจยคามิสฺส เหตุสฺส ธมฺมสฺส อธิปติ\-ปจฺจเยน ปจฺจโย.  อปจยคามี เหตุ ธมฺโม เนวา\-จยคามิ\-นา\-ปจยคามิสฺส เหตุสฺส ธมฺมสฺส อธิปติ\-ปจฺจเยน ปจฺจโย. อารมฺมณา\-ธิปติ. (3)}

\csromanmarkh{10. อาจยคามิตฺติก}\csromanmarkhii{1. เหตุทุก}\csromanheadernomark{2}{10. อาจยคามิตฺติก 1. เหตุทุก}
\prose{\csromanfolio{513}เนวา\-จยคามิ\-นา\-ปจยคามี เหตุ ธมฺโม เนวา\-จยคามิ\-นา\-ปจยคามิสฺส เหตุสฺส ธมฺมสฺส อธิปติ\-ปจฺจเยน ปจฺจโย.  เนวา\-จยคามิ\-นา\-ปจยคามี เหตุ ธมฺโม อาจยคามิสฺส เหตุสฺส ธมฺมสฺส อธิปติ\-ปจฺจเยน ปจฺจโย.}

\proseitem{136}{อาจยคามี เหตุ ธมฺโม อาจยคามิสฺส เหตุสฺส ธมฺมสฺส อนนฺตร\-ปจฺจเยน ปจฺจโย.  อาจยคามี เหตุ ธมฺโม อปจยคามิสฺส เหตุสฺส ธมฺมสฺส อนนฺตร\-ปจฺจเยน ปจฺจโย.  อาจยคามี เหตุ ธมฺโม เนวา\-จยคามิ\-นา\-ปจยคามิสฺส เหตุสฺส ธมฺมสฺส อนนฺตร\-ปจฺจเยน ปจฺจโย. (3)}

\prose{อปจยคามี เหตุ ธมฺโม เนวา\-จยคามิ\-นา\-ปจยคามิสฺส เหตุสฺส ธมฺมสฺส อนนฺตร\-ปจฺจเยน ปจฺจโย. (1)}

\prose{เนวา\-จยคามิ\-นา\-ปจยคามี เหตุ ธมฺโม เนวา\-จยคามิ\-นา\-ปจยคามิสฺส เหตุสฺส ธมฺมสฺส อนนฺตร\-ปจฺจเยน ปจฺจโย. (1) (สํขิตฺตํ.)}

\proseitem{137}{เหตุยา ตีณิ, อารมฺมเณ ฉ, อธิปติยา ฉ, อนนฺตเร ปญฺจ, สมนนฺตเร ปญฺจ, สหชาเต ตีณิ, อญฺญมญฺเญ ตีณิ, นิสฺสเย ตีณิ, อุปนิสฺสเย อฏฺฐ, อาเสวเน ตีณิ, วิปาเก เอกํ ฯเปฯ อวิคเต ตีณิ. (สํขิตฺตํ.)}

\prose{นเหตุยา อฏฺฐ, นอารมฺมเณ อฏฺฐ. (สํขิตฺตํ.)}

\prose{เหตุปจฺจยา นอารมฺมเณ ตีณิ. (สํขิตฺตํ.)}

\prose{นเหตุปจฺจยา อารมฺมเณ ฉ. (สํขิตฺตํ.)}

\prose{(ยถา กุสลตฺติเก ปญฺหาวารํ, เอวํ วิตฺถาเรตพฺพํ.)}

\csromanmarkh{นเหตุปท}\csromanmarkhii{1. ปฏิจฺจวาร}\csromanheadernomark{2}{\textbf{นเหตุปท 1. ปฏิจฺจวาร}}
\prose{ปจฺจยจตุกฺกเหตุ อารมฺมณ}

\proseitem{138}{อาจยคามิํ นเหตุํ ธมฺมํ ปฏิจฺจ อาจยคามี นเหตุ ธมฺโม อุปฺปชฺชติ เหตุปจฺจยา.  อาจยคามิํ นเหตุํ ธมฺมํ ปฏิจฺจ เนวา\-จยคามิ\-นา\-ปจยคามี นเหตุ ธมฺโม อุปฺปชฺชติ เหตุปจฺจยา.  อาจยคามิํ \csromanfolio{514}นเหตุํ ธมฺมํ ปฏิจฺจ อาจยคามี นเหตุ จ เนวา\-จยคามิ\-นา\-ปจยคามี นเหตุ จ ธมฺมา อุปฺปชฺชนฺติ เหตุปจฺจยา. (3)}

\prose{อปจยคามิํ นเหตุํ ธมฺมํ ปฏิจฺจ อปจยคามี นเหตุ ธมฺโม อุปฺปชฺชติ เหตุปจฺจยา.  อปจยคามิํ นเหตุํ ธมฺมํ ปฏิจฺจ เนวา\-จยคามิ\-นา\-ปจยคามี นเหตุ ธมฺโม อุปฺปชฺชติ เหตุปจฺจยา.  อปจยคามิํ นเหตุํ ธมฺมํ ปฏิจฺจ อปจยคามี นเหตุ จ เนวา\-จยคามิ\-นา\-ปจยคามี นเหตุ จ ธมฺมา อุปฺปชฺชนฺติ เหตุปจฺจยา. (3)}

\prose{เนวา\-จยคามิ\-นา\-ปจยคามิํ นเหตุํ ธมฺมํ ปฏิจฺจ เนวา\-จยคามิ\-นา\-ปจยคามี นเหตุ ธมฺโม อุปฺปชฺชติ เหตุปจฺจยา. (1)}

\prose{อาจยคามิํ นเหตุญฺจ เนวา\-จยคามิ\-นา\-ปจยคามิํ นเหตุญฺจ ธมฺมํ ปฏิจฺจ เนวา\-จยคามิ\-นา\-ปจยคามี นเหตุ ธมฺโม อุปฺปชฺชติ เหตุปจฺจยา. (1)}

\prose{อปจยคามิํ นเหตุญฺจ เนวา\-จยคามิ\-นา\-ปจยคามิํ นเหตุญฺจ ธมฺมํ ปฏิจฺจ เนวา\-จยคามิ\-นา\-ปจยคามี นเหตุ ธมฺโม อุปฺปชฺชติ เหตุปจฺจยา. (1)}

\proseitem{139}{อาจยคามิํ นเหตุํ ธมฺมํ ปฏิจฺจ อาจยคามี นเหตุ ธมฺโม อุปฺปชฺชติ อารมฺมณ\-ปจฺจยา. (1)}

\prose{อปจยคามิํ นเหตุํ ธมฺมํ ปฏิจฺจ อปจยคามี นเหตุ ธมฺโม อุปฺปชฺชติ อารมฺมณ\-ปจฺจยา. (1)}

\prose{เนวา\-จยคามิ\-นา\-ปจยคามิํ นเหตุํ ธมฺมํ ปฏิจฺจ เนวา\-จยคามิ\-นา\-ปจยคามี นเหตุ ธมฺโม อุปฺปชฺชติ อารมฺมณ\-ปจฺจยา. (1) (สํขิตฺตํ.)}

\prose{เหตุยา นว, อารมฺมเณ ตีณิ, อธิปติยา นว ฯเปฯ อวิคเต นว. (สํขิตฺตํ.)}

\prose{นเหตุ นอารมฺมณ}

\proseitem{140}{เนวา\-จยคามิ\-นา\-ปจยคามิํ นเหตุํ ธมฺมํ ปฏิจฺจ เนวา\-จยคามิ\-นา\-ปจยคามี นเหตุ ธมฺโม อุปฺปชฺชติ นเหตุปจฺจยา. (1)}

\prose{อาจยคามิํ นเหตุํ ธมฺมํ ปฏิจฺจ เนวา\-จยคามิ\-นา\-ปจยคามี นเหตุ ธมฺโม อุปฺปชฺชติ นอารมฺมณ\-ปจฺจยา. (สํขิตฺตํ.)}

\proseitem{141}{นเหตุยา เอกํ, นอารมฺมเณ ปญฺจ, นอธิปติยา ฉ ฯเปฯ โนวิคเต ปญฺจ. (สํขิตฺตํ.)}

\csromanmarkh{10. อาจยคามิตฺติก}\csromanmarkhii{1. เหตุทุก}\csromanheadernomark{2}{10. อาจยคามิตฺติก 1. เหตุทุก}
\prose{\csromanfolio{515}เหตุปจฺจยา นอารมฺมเณ ปญฺจ. (สํขิตฺตํ.)}

\prose{นเหตุปจฺจยา อารมฺมเณ เอกํ. (สํขิตฺตํ.)}

\prose{(สหชาตวารมฺปิ ฯเปฯ สมฺปยุตฺตวารมฺปิ ปฏิจฺจ\-วาร\-สทิสํ วิตฺถาเรตพฺพํ.)}

\csromanheader{2}{7. \textbf{ปญฺหาวาร ปจฺจยจตุกฺก}}
\proseitem{142}{อาจยคามี นเหตุ ธมฺโม อาจยคามิสฺส นเหตุสฺส ธมฺมสฺส อารมฺมณ\-ปจฺจเยน ปจฺจโย.  อาจยคามี นเหตุ ธมฺโม เนวา\-จยคามิ\-นา\-ปจยคามิสฺส นเหตุสฺส ธมฺมสฺส อารมฺมณ\-ปจฺจเยน ปจฺจโย. (2)}

\prose{อปจยคามี นเหตุ ธมฺโม อาจยคามิสฺส นเหตุสฺส ธมฺมสฺส อารมฺมณ\-ปจฺจเยน ปจฺจโย.  อปจยคามี นเหตุ ธมฺโม เนวา\-จยคามิ\-นา\-ปจยคามิสฺส นเหตุสฺส ธมฺมสฺส อารมฺมณ\-ปจฺจเยน ปจฺจโย. (2) เนวา\-จยคามิ\-นา\-ปจยคามี นเหตุ ธมฺโม เนวา\-จยคามิ\-นา\-ปจยคามิสฺส นเหตุสฺส ธมฺมสฺส อารมฺมณ\-ปจฺจเยน ปจฺจโย.  เนวา\-จยคามิ\-นา\-ปจยคามี นเหตุ ธมฺโม อาจยคามิสฺส นเหตุสฺส ธมฺมสฺส อารมฺมณ\-ปจฺจเยน ปจฺจโย.  เนวา\-จยคามิ\-นา\-ปจยคามี นเหตุ ธมฺโม อปจยคามิสฺส นเหตุสฺส ธมฺมสฺส อารมฺมณ\-ปจฺจเยน ปจฺจโย. (3)}

\proseitem{143}{อาจยคามี นเหตุ ธมฺโม อาจยคามิสฺส นเหตุสฺส ธมฺมสฺส อธิปติ\-ปจฺจเยน ปจฺจโย. ตีณิ.}

\prose{อปจยคามี นเหตุ ธมฺโม อปจยคามิสฺส นเหตุสฺส ธมฺมสฺส อธิปติ\-ปจฺจเยน ปจฺจโย.  อปจยคามี นเหตุ ธมฺโม อาจยคามิสฺส นเหตุสฺส ธมฺมสฺส อธิปติ\-ปจฺจเยน ปจฺจโย.  อปจยคามี นเหตุ ธมฺโม เนวา\-จยคามิ\-นา\-ปจยคามิสฺส นเหตุสฺส ธมฺมสฺส อธิปติ\-ปจฺจเยน ปจฺจโย.  อปจยคามี นเหตุ ธมฺโม อปจยคามิสฺส นเหตุสฺส จ เนวา\-จยคามิ\-นา\-ปจยคามิสฺส นเหตุสฺส จ ธมฺมสฺส อธิปติ\-ปจฺจเยน ปจฺจโย. (4) เนวา\-จยคามิ\-นา\-ปจยคามี นเหตุ ธมฺโม เนวา\-จยคามิ\-นา\-ปจยคามิสฺส นเหตุสฺส ธมฺมสฺส อธิปติ\-ปจฺจเยน ปจฺจโย. ตีณิ.}

\prose{\csromanfolio{516}อาจยคามี นเหตุ ธมฺโม อาจยคามิสฺส นเหตุสฺส ธมฺมสฺส อนนฺตร\-ปจฺจเยน ปจฺจโย.  อาจยคามี นเหตุ ธมฺโม อปจยคามิสฺส นเหตุสฺส ธมฺมสฺส อนนฺตร\-ปจฺจเยน ปจฺจโย.  อาจยคามี นเหตุ ธมฺโม เนวา\-จยคามิ\-นา\-ปจยคามิสฺส นเหตุสฺส ธมฺมสฺส อนนฺตร\-ปจฺจเยน ปจฺจโย. (3)}

\prose{อปจยคามี นเหตุ ธมฺโม เนวา\-จยคามิ\-นา\-ปจยคามิสฺส นเหตุสฺส ธมฺมสฺส อนนฺตร\-ปจฺจเยน ปจฺจโย. (1)}

\prose{เนวา\-จยคามิ\-นา\-ปจยคามี นเหตุ ธมฺโม เนวา\-จยคามิ\-นา\-ปจยคามิสฺส นเหตุสฺส ธมฺมสฺส อนนฺตร\-ปจฺจเยน ปจฺจโย.  เนวา\-จยคามิ\-นา\-ปจยคามี นเหตุ ธมฺโม อาจยคามิสฺส นเหตุสฺส ธมฺมสฺส อนนฺตร\-ปจฺจเยน ปจฺจโย. (2) ฯเปฯ.}

\prose{อาจยคามี นเหตุ ธมฺโม อาจยคามิสฺส นเหตุสฺส ธมฺมสฺส อุปนิสฺสย\-ปจฺจเยน ปจฺจโย. (สํขิตฺตํ.)}

\proseitem{145}{อารมฺมเณ สตฺต, อธิปติยา ทส, อนนฺตเร ฉ, สมนนฺตเร ฉ, สหชาเต นว, อญฺญมญฺเญ ตีณิ, นิสฺสเย เตรส, อุปนิสฺสเย นว ฯเปฯ อวิคเต เตรส. (สํขิตฺตํ.)}

\prose{นเหตุยา ปนฺนรส, นอารมฺมเณ ปนฺนรส. (สํขิตฺตํ.)}

\prose{อารมฺมณ\-ปจฺจยา นเหตุยา สตฺต. (สํขิตฺตํ.)}

\prose{นเหตุปจฺจยา อารมฺมเณ สตฺต. (สํขิตฺตํ.)}

\prose{(ยถา กุสลตฺติเก ปญฺหาวารํ, เอวํ วิตฺถาเรตพฺพํ.)}

\csromanmatikaanchor{mtk.112}
\csromanmatikaanchor{mtk.113}
\csromantocmarknopage{subparagraph}{11. เสกฺขตฺติก}{mtk.112}
\bookmark[dest={mtk.112},level=5]{11. เสกฺขตฺติก}
\csromantocmarkref{subparagraph}{1. เหตุทุก}{mtk.113}
\bookmark[dest={mtk.113},level=5]{1. เหตุทุก}
\csromanmarkh{11. เสกฺขตฺติก}\csromanmarkhii{1. เหตุทุก}\csromanheadernomark{6}{11. \textbf{เสกฺขตฺติก 1. เหตุทุก}}
\csromanheader{2}{1-7. ปฏิจฺจาทิวาร ปจฺจยจตุกฺกเหตุ}
\vspace{\tipitakaparskip}
\csromancenter{เสกฺขํ เหตุํ ธมฺมํ ปฏิจฺจ เสกฺโข เหตุ ธมฺโม อุปฺปชฺชติ เหตุปจฺจยา. (1)}
\csromanmarkh{11. เสกฺขตฺติก}\csromanmarkhii{1. เหตุทุก}\csromanheadernomark{2}{11. เสกฺขตฺติก 1. เหตุทุก}
\prose{\csromanfolio{517}อเสกฺขํ เหตุํ ธมฺมํ ปฏิจฺจ อเสกฺโข เหตุ ธมฺโม อุปฺปชฺชติ เหตุปจฺจยา. (1)}

\prose{เนว\-เสกฺข\-นาเสกฺขํ เหตุํ ธมฺมํ ปฏิจฺจ เนว\-เสกฺข\-นา\-เสกฺโข เหตุ ธมฺโม อุปฺปชฺชติ เหตุปจฺจยา. (1) (สํขิตฺตํ.)}

\proseitem{147}{เหตุยา ตีณิ, อารมฺมเณ ตีณิ, อธิปติยา ตีณิ ฯเปฯ อวิคเต ตีณิ. (สํขิตฺตํ.)}

\prose{นอธิปติยา ตีณิ, นปุเรชาเต ตีณิ, นปจฺฉาชาเต ตีณิ, นอาเสวเน ตีณิ, นวิปาเก เทฺว, นวิปฺปยุตฺเต ตีณิ. (สํขิตฺตํ.)}

\prose{เหตุปจฺจยา นอธิปติยา ตีณิ. (สํขิตฺตํ.)}

\prose{นอธิปติ\-ปจฺจยา เหตุยา ตีณิ. (สํขิตฺตํ.) (สหชาตวารมฺปิ ฯเปฯ สมฺปยุตฺตวารมฺปิ วิตฺถาเรตพฺพํ.)}

\prose{เหตุอารมฺมณาทิ}

\proseitem{148}{เสกฺโข เหตุ ธมฺโม เสกฺขสฺส เหตุสฺส ธมฺมสฺส เหตุปจฺจเยน ปจฺจโย. (1)}

\prose{อเสกฺโข เหตุ ธมฺโม อเสกฺขสฺส เหตุสฺส ธมฺมสฺส เหตุปจฺจเยน ปจฺจโย. (1)}

\prose{เนว\-เสกฺข\-นา\-เสกฺโข เหตุ ธมฺโม เนว\-เสกฺข\-นาเสกฺขสฺส เหตุสฺส ธมฺมสฺส เหตุปจฺจเยน ปจฺจโย. (1)}

\proseitem{149}{เสกฺโข เหตุ ธมฺโม เนว\-เสกฺข\-นาเสกฺขสฺส เหตุสฺส ธมฺมสฺส อารมฺมณ\-ปจฺจเยน ปจฺจโย. (1)}

\prose{อเสกฺโข เหตุ ธมฺโม เนว\-เสกฺข\-นาเสกฺขสฺส เหตุสฺส ธมฺมสฺส อารมฺมณ\-ปจฺจเยน ปจฺจโย. (1)}

\prose{เนว\-เสกฺข\-นา\-เสกฺโข เหตุ ธมฺโม เนว\-เสกฺข\-นาเสกฺขสฺส เหตุสฺส ธมฺมสฺส อารมฺมณ\-ปจฺจเยน ปจฺจโย. (1)}

\proseitem{150}{\csromanfolio{518}เสกฺโข เหตุ ธมฺโม เสกฺขสฺส เหตุสฺส ธมฺมสฺส อธิปติ\-ปจฺจเยน ปจฺจโย.  เสกฺโข เหตุ ธมฺโม เนว\-เสกฺข\-นาเสกฺขสฺส เหตุสฺส ธมฺมสฺส อธิปติ\-ปจฺจเยน ปจฺจโย. (2)}

\prose{อเสกฺโข เหตุ ธมฺโม อเสกฺขสฺส เหตุสฺส ธมฺมสฺส อธิปติ\-ปจฺจเยน ปจฺจโย.  อเสกฺโข เหตุ ธมฺโม เนว\-เสกฺข\-นาเสกฺขสฺส เหตุสฺส ธมฺมสฺส อธิปติ\-ปจฺจเยน ปจฺจโย. (2)}

\prose{เนว\-เสกฺข\-นา\-เสกฺโข เหตุ ธมฺโม เนว\-เสกฺข\-นาเสกฺขสฺส เหตุสฺส ธมฺมสฺส อธิปติ\-ปจฺจเยน ปจฺจโย. (1)}

\proseitem{151}{เสกฺโข เหตุ ธมฺโม เสกฺขสฺส เหตุสฺส ธมฺมสฺส อนนฺตร\-ปจฺจเยน ปจฺจโย.  เสกฺโข เหตุ ธมฺโม อเสกฺขสฺส เหตุสฺส ธมฺมสฺส อนนฺตร\-ปจฺจเยน ปจฺจโย.  เสกฺโข เหตุ ธมฺโม เนว\-เสกฺข\-นาเสกฺขสฺส เหตุสฺส ธมฺมสฺส อนนฺตร\-ปจฺจเยน ปจฺจโย. (3)}

\prose{อเสกฺโข เหตุ ธมฺโม อเสกฺขสฺส เหตุสฺส ธมฺมสฺส อนนฺตร\-ปจฺจเยน ปจฺจโย.  อเสกฺโข เหตุ ธมฺโม เนว\-เสกฺข\-นาเสกฺขสฺส เหตุสฺส ธมฺมสฺส อนนฺตร\-ปจฺจเยน ปจฺจโย. (2)}

\prose{เนว\-เสกฺข\-นา\-เสกฺโข เหตุ ธมฺโม เนว\-เสกฺข\-นาเสกฺขสฺส เหตุสฺส ธมฺมสฺส อนนฺตร\-ปจฺจเยน ปจฺจโย. เนว\-เสกฺข\-นา\-เสกฺโข เหตุ ธมฺโม เสกฺขสฺส เหตุสฺส ธมฺมสฺส อนนฺตร\-ปจฺจเยน ปจฺจโย.  เนว\-เสกฺข\-นา\-เสกฺโข เหตุ ธมฺโม อเสกฺขสฺส เหตุสฺส ธมฺมสฺส อนนฺตร\-ปจฺจเยน ปจฺจโย. (3) (สํขิตฺตํ.)}

\proseitem{152}{เหตุยา ตีณิ, อารมฺมเณ ตีณิ, อธิปติยา ปญฺจ, อนนฺตเร อฏฺฐ, สมนนฺตเร อฏฺฐ, สหชาเต ตีณิ, อญฺญมญฺเญ ตีณิ, นิสฺสเย ตีณิ, อุปนิสฺสเย อฏฺฐ ฯเปฯ อวิคเต ตีณิ. (สํขิตฺตํ.)}

\prose{นเหตุยา อฏฺฐ, นอารมฺมเณ อฏฺฐ. (สํขิตฺตํ.)}

\prose{เหตุปจฺจยา นอารมฺมเณ ตีณิ. (สํขิตฺตํ.)}

\prose{นเหตุปจฺจยา อารมฺมเณ ตีณิ. (สํขิตฺตํ.)}

\prose{(ยถา กุสลตฺติเก ปญฺหาวารํ, เอวํ วิตฺถาเรตพฺพํ.)}

\csromanmarkh{11. เสกฺขตฺติก}\csromanmarkhii{1. เหตุทุก นเหตุปท 1. ปฏิจฺจวาร}\csromanheadernomark{2}{11. เสกฺขตฺติก 1. เหตุทุก \textbf{นเหตุปท 1. ปฏิจฺจวาร}}
\prose{\csromanfolio{519}ปจฺจยจตุกฺก เหตุอารมฺมณ}

\proseitem{153}{เสกฺขํ นเหตุํ ธมฺมํ ปฏิจฺจ เสกฺโข นเหตุ ธมฺโม อุปฺปชฺชติ เหตุปจฺจยา.  เสกฺขํ นเหตุํ ธมฺมํ ปฏิจฺจ เนว\-เสกฺข\-นา\-เสกฺโข นเหตุ ธมฺโม อุปฺปชฺชติ เหตุปจฺจยา.  เสกฺขํ นเหตุํ ธมฺมํ ปฏิจฺจ เสกฺโข นเหตุ จ เนว\-เสกฺข\-นา\-เสกฺโข นเหตุ จ ธมฺมา อุปฺปชฺชนฺติ เหตุปจฺจยา. (3)}

\prose{อเสกฺขํ นเหตุํ ธมฺมํ ปฏิจฺจ อเสกฺโข นเหตุ ธมฺโม อุปฺปชฺชติ เหตุปจฺจยา. ตีณิ.}

\prose{เนว\-เสกฺข\-นาเสกฺขํ นเหตุํ ธมฺมํ ปฏิจฺจ เนว\-เสกฺข\-นา\-เสกฺโข นเหตุ ธมฺโม อุปฺปชฺชติ เหตุปจฺจยา. (1)}

\prose{เสกฺขํ นเหตุญฺจ เนว\-เสกฺข\-นาเสกฺขํ นเหตุญฺจ ธมฺมํ ปฏิจฺจ เนว\-เสกฺข\-นา\-เสกฺโข นเหตุ ธมฺโม อุปฺปชฺชติ เหตุปจฺจยา. (1)}

\prose{อเสกฺขํ นเหตุญฺจ เนว\-เสกฺข\-นาเสกฺขํ นเหตุญฺจ ธมฺมํ ปฏิจฺจ เนว\-เสกฺข\-นา\-เสกฺโข นเหตุ ธมฺโม อุปฺปชฺชติ เหตุปจฺจยา. (1)}

\proseitem{154}{เสกฺขํ นเหตุํ ธมฺมํ ปฏิจฺจ เสกฺโข นเหตุ ธมฺโม อุปฺปชฺชติ อารมฺมณ\-ปจฺจยา. (1)}

\prose{อเสกฺขํ นเหตุํ ธมฺมํ ปฏิจฺจ อเสกฺโข นเหตุ ธมฺโม อุปฺปชฺชติ อารมฺมณ\-ปจฺจยา. (1)}

\prose{เนว\-เสกฺข\-นาเสกฺขํ นเหตุํ ธมฺมํ ปฏิจฺจ เนว\-เสกฺข\-นา\-เสกฺโข นเหตุ ธมฺโม อุปฺปชฺชติ อารมฺมณ\-ปจฺจยา. (1) (สํขิตฺตํ.)}

\prose{เหตุยา นว, อารมฺมเณ ตีณิ, อธิปติยา นว ฯเปฯ อาเสวเน เทฺว ฯเปฯ อวิคเต นว. (สํขิตฺตํ.)}

\prose{ปจฺจนีย\-น\-เหตุ}

\proseitem{155}{เนว\-เสกฺข\-นาเสกฺขํ นเหตุํ ธมฺมํ ปฏิจฺจ เนว\-เสกฺข\-นา\-เสกฺโข นเหตุ ธมฺโม อุปฺปชฺชติ เหตุปจฺจยา. (สํขิตฺตํ.)}

\prose{\csromanfolio{520}นเหตุยา เอกํ, นอารมฺมเณ ปญฺจ, นอธิปติยา ตีณิ ฯเปฯ โนวิคเต ปญฺจ. (สํขิตฺตํ.)}

\prose{เหตุปจฺจยา นอารมฺมเณ ปญฺจ. (สํขิตฺตํ.)}

\prose{นเหตุปจฺจยา อารมฺมเณ เอกํ. (สํขิตฺตํ.)}

\prose{(สหชาตวารมฺปิ ฯเปฯ สมฺปยุตฺตวารมฺปิ ปฏิจฺจ\-วาร\-สทิสํ วิตฺถาเรตพฺพํ.)}

\csromanheader{2}{7. \textbf{ปญฺหาวาร ปจฺจยจตุกฺก}}
\proseitem{156}{เสกฺโข นเหตุ ธมฺโม เนว\-เสกฺข\-นาเสกฺขสฺส นเหตุสฺส ธมฺมสฺส อารมฺมณ\-ปจฺจเยน ปจฺจโย. (1)}

\prose{อเสกฺโข นเหตุ ธมฺโม เนว\-เสกฺข\-นาเสกฺขสฺส นเหตุสฺส ธมฺมสฺส อารมฺมณ\-ปจฺจเยน ปจฺจโย. (1)}

\prose{เนว\-เสกฺข\-นา\-เสกฺโข นเหตุ ธมฺโม เนว\-เสกฺข\-นาเสกฺขสฺส นเหตุสฺส ธมฺมสฺส อารมฺมณ\-ปจฺจเยน ปจฺจโย.  เนว\-เสกฺข\-นา\-เสกฺโข นเหตุ ธมฺโม เสกฺขสฺส นเหตุสฺส ธมฺมสฺส อารมฺมณ\-ปจฺจเยน ปจฺจโย.  เนว\-เสกฺข\-นา\-เสกฺโข นเหตุ ธมฺโม อเสกฺขสฺส นเหตุสฺส ธมฺมสฺส อารมฺมณ\-ปจฺจเยน ปจฺจโย. (3)}

\prose{เสกฺโข นเหตุ ธมฺโม เสกฺขสฺส นเหตุสฺส ธมฺมสฺส อธิปติ\-ปจฺจเยน ปจฺจโย. ตีณิ.}

\prose{อเสกฺโข นเหตุ ธมฺโม อเสกฺขสฺส นเหตุสฺส ธมฺมสฺส อธิปติ\-ปจฺจเยน ปจฺจโย. ตีณิ.}

\prose{เนว\-เสกฺข\-นา\-เสกฺโข นเหตุ ธมฺโม เนว\-เสกฺข\-นาเสกฺขสฺส นเหตุสฺส ธมฺมสฺส อธิปติ\-ปจฺจเยน ปจฺจโย. ตีณิ. (สํขิตฺตํ.)}

\proseitem{157}{อารมฺมเณ ปญฺจ, อธิปติยา นว, อนนฺตเร อฏฺฐ, สมนนฺตเร อฏฺฐ. สหชาเต นว, อญฺญมญฺเญ ตีณิ, นิสฺสเย เตรส, อุปนิสฺสเย อฏฺฐ ฯเปฯ อวิคเต เตรส. (สํขิตฺตํ.)}

\csromanmatikaanchor{mtk.114}
\csromanmatikaanchor{mtk.115}
\csromantocmarknopage{subparagraph}{12. ปริตฺตตฺติก}{mtk.114}
\bookmark[dest={mtk.114},level=5]{12. ปริตฺตตฺติก}
\csromantocmarkref{subparagraph}{1. เหตุทุก}{mtk.115}
\bookmark[dest={mtk.115},level=5]{1. เหตุทุก}
\csromanmarkh{12. ปริตฺตตฺติก}\csromanmarkhii{1. เหตุทุก}\csromanheadernomark{6}{12. ปริตฺตตฺติก 1. เหตุทุก}
\prose{\csromanfolio{521}นเหตุยา จุทฺทส, นอารมฺมเณ จุทฺทส. (สํขิตฺตํ.)}

\prose{อารมฺมณ\-ปจฺจยา นเหตุยา ปญฺจ. (สํขิตฺตํ.)}

\prose{นเหตุปจฺจยา อารมฺมเณ ปญฺจ. (สํขิตฺตํ.)}

\prose{(ยถา กุสลตฺติเก ปญฺหาวารํ, เอวํ วิตฺถาเรตพฺพํ.)}

\csromanmarkh{12. ปริตฺตตฺติก}\csromanmarkhii{1. เหตุทุก}\csromanheadernomark{2}{12. ปริตฺตตฺติก 1. เหตุทุก}
\csromanheader{2}{1-7. ปฏิจฺจาทิวาร ปจฺจยจตุกฺก}
\proseitem{158}{ปริตฺตํ เหตุํ ธมฺมํ ปฏิจฺจ ปริตฺโต เหตุ ธมฺโม อุปฺปชฺชติ เหตุปจฺจยา. (1)}

\prose{มหคฺคตํ เหตุํ ธมฺมํ ปฏิจฺจ มหคฺคโต เหตุ ธมฺโม อุปฺปชฺชติ เหตุปจฺจยา. (1)}

\prose{อปฺปมาณํ เหตุํ ธมฺมํ ปฏิจฺจ อปฺปมาโณ เหตุ ธมฺโม อุปฺปชฺชติ เหตุปจฺจยา. (1) (สํขิตฺตํ.)}

\prose{เหตุยา ตีณิ, อารมฺมเณ ตีณิ ฯเปฯ อวิคเต ตีณิ. (สํขิตฺตํ.)}

\prose{นอธิปติยา ตีณิ ฯเปฯ นวิปฺปยุตฺเต ตีณิ. (สํขิตฺตํ.)}

\prose{(สหชาตวารมฺปิ ฯเปฯ สมฺปยุตฺตวารมฺปิ ปฏิจฺจ\-วาร\-สทิสํ วิตฺถาเรตพฺพํ.)}

\prose{เหตุอารมฺมณาทิ}

\proseitem{159}{ปริตฺโต เหตุ ธมฺโม ปริตฺตสฺส เหตุสฺส ธมฺมสฺส เหตุปจฺจเยน ปจฺจโย. (1)}

\prose{มหคฺคโต เหตุ ธมฺโม มหคฺคตสฺส เหตุสฺส ธมฺมสฺส เหตุปจฺจเยน ปจฺจโย. (1)}

\prose{อปฺปมาโณ เหตุ ธมฺโม อปฺปมาณสฺส เหตุสฺส ธมฺมสฺส เหตุปจฺจเยน ปจฺจโย. (1)}

\prose{ปริตฺโต เหตุ ธมฺโม ปริตฺตสฺส เหตุสฺส ธมฺมสฺส อารมฺมณ\-ปจฺจเยน ปจฺจโย.  ปริตฺโต เหตุ ธมฺโม มหคฺคตสฺส เหตุสฺส ธมฺมสฺส อารมฺมณ\-ปจฺจเยน ปจฺจโย. (2)}

\prose{\csromanfolio{522}มหคฺคโต เหตุ ธมฺโม มหคฺคตสฺส เหตุสฺส ธมฺมสฺส อารมฺมณ\-ปจฺจเยน ปจฺจโย.  มหคฺคโต เหตุ ธมฺโม ปริตฺตสฺส เหตุสฺส ธมฺมสฺส อารมฺมณ\-ปจฺจเยน ปจฺจโย. (2)}

\prose{อปฺปมาโณ เหตุ ธมฺโม ปริตฺตสฺส เหตุสฺส ธมฺมสฺส อารมฺมณ\-ปจฺจเยน ปจฺจโย.  อปฺปมาโณ เหตุ ธมฺโม มหคฺคตสฺส เหตุสฺส ธมฺมสฺส อารมฺมณ\-ปจฺจเยน ปจฺจโย. (2)}

\proseitem{160}{ปริตฺโต เหตุ ธมฺโม ปริตฺตสฺส เหตุสฺส ธมฺมสฺส อธิปติ\-ปจฺจเยน ปจฺจโย. (1)}

\prose{มหคฺคโต เหตุ ธมฺโม มหคฺคตสฺส เหตุสฺส ธมฺมสฺส อธิปติ\-ปจฺจเยน ปจฺจโย.  มหคฺคโต เหตุ ธมฺโม ปริตฺตสฺส เหตุสฺส ธมฺมสฺส อธิปติ\-ปจฺจเยน ปจฺจโย. (2)}

\prose{อปฺปมาโณ เหตุ ธมฺโม อปฺปมาณสฺส เหตุสฺส ธมฺมสฺส อธิปติ\-ปจฺจเยน ปจฺจโย. เทฺว.}

\prose{ปริตฺโต เหตุ ธมฺโม ปริตฺตสฺส เหตุสฺส ธมฺมสฺส อนนฺตร\-ปจฺจเยน ปจฺจโย.  ปริตฺโต เหตุ ธมฺโม มหคฺคตสฺส เหตุสฺส ธมฺมสฺส อนนฺตร\-ปจฺจเยน ปจฺจโย.  ปริตฺโต เหตุ ธมฺโม อปฺปมาณสฺส เหตุสฺส ธมฺมสฺส อนนฺตร\-ปจฺจเยน ปจฺจโย. (3)}

\prose{มหคฺคโต เหตุ ธมฺโม มหคฺคตสฺส เหตุสฺส ธมฺมสฺส อนนฺตร\-ปจฺจเยน ปจฺจโย.  มหคฺคโต เหตุ ธมฺโม ปริตฺตสฺส เหตุสฺส ธมฺมสฺส อนนฺตร\-ปจฺจเยน ปจฺจโย.  มหคฺคโต เหตุ ธมฺโม อปฺปมาณสฺส เหตุสฺส ธมฺมสฺส อนนฺตร\-ปจฺจเยน ปจฺจโย. (3)}

\prose{อปฺปมาโณ เหตุ ธมฺโม อปฺปมาณสฺส เหตุสฺส ธมฺมสฺส อนนฺตร\-ปจฺจเยน ปจฺจโย.  อปฺปมาโณ เหตุ ธมฺโม ปริตฺตสฺส เหตุสฺส ธมฺมสฺส อนนฺตร\-ปจฺจเยน ปจฺจโย.  อปฺปมาโณ เหตุ ธมฺโม มหคฺคตสฺส เหตุสฺส ธมฺมสฺส อนนฺตร\-ปจฺจเยน ปจฺจโย. (3) ฯเปฯ.}

\prose{ปริตฺโต เหตุ ธมฺโม ปริตฺตสฺส เหตุสฺส ธมฺมสฺส อุปนิสฺสย\-ปจฺจเยน ปจฺจโย. (สํขิตฺตํ.)}

\csromanmarkh{12. ปริตฺตตฺติก}\csromanmarkhii{1. เหตุทุก}\csromanheadernomark{2}{12. ปริตฺตตฺติก 1. เหตุทุก}
\proseitem{161}{\csromanfolio{523}เหตุยา ตีณิ, อารมฺมเณ ฉ, อธิปติยา ปญฺจ, อนนฺตเร นว, สมนนฺตเร นว, สหชาเต ตีณิ, อญฺญมญฺเญ ตีณิ, นิสฺสเย ตีณิ, อุปนิสฺสเย นว ฯเปฯ อวิคเต ตีณิ. (สํขิตฺตํ.)}

\prose{นเหตุยา นว, นอารมฺมเณ นว. (สํขิตฺตํ.)}

\prose{เหตุปจฺจยา นอารมฺมเณ ตีณิ. (สํขิตฺตํ.)}

\prose{นเหตุปจฺจยา อารมฺมเณ ฉ. (สํขิตฺตํ.)}

\prose{(ยถา กุสลตฺติเก ปญฺหาวารํ, เอวํ วิตฺถาเรตพฺพํ.)}

\csromanmarkh{นเหตุปท}\csromanmarkhii{1. ปฏิจฺจวาร}\csromanheadernomark{2}{\textbf{นเหตุปท 1. ปฏิจฺจวาร}}
\prose{ปจฺจยจตุกฺกเหตุ}

\proseitem{162}{ปริตฺตํ นเหตุํ ธมฺมํ ปฏิจฺจ ปริตฺโต นเหตุ ธมฺโม อุปฺปชฺชติ เหตุปจฺจยา. ตีณิ.}

\prose{มหคฺคตํ นเหตุํ ธมฺมํ ปฏิจฺจ มหคฺคโต นเหตุ ธมฺโม อุปฺปชฺชติ เหตุปจฺจยา.  มหคฺคตํ นเหตุํ ธมฺมํ ปฏิจฺจ ปริตฺโต นเหตุ ธมฺโม อุปฺปชฺชติ เหตุปจฺจยา.  มหคฺคตํ นเหตุํ ธมฺมํ ปฏิจฺจ ปริตฺโต นเหตุ จ มหคฺคโต นเหตุ จ ธมฺมา อุปฺปชฺชนฺติ เหตุปจฺจยา. (3)}

\prose{อปฺปมาณํ นเหตุํ ธมฺมํ ปฏิจฺจ อปฺปมาโณ นเหตุ ธมฺโม อุปฺปชฺชติ เหตุปจฺจยา.  อปฺปมาณํ นเหตุํ ธมฺมํ ปฏิจฺจ ปริตฺโต นเหตุ ธมฺโม อุปฺปชฺชติ เหตุปจฺจยา.  อปฺปมาณํ นเหตุํ ธมฺมํ ปฏิจฺจ ปริตฺโต นเหตุ จ อปฺปมาโณ นเหตุ จ ธมฺมา อุปฺปชฺชนฺติ เหตุปจฺจยา. (3)}

\prose{ปริตฺตํ นเหตุญฺจ มหคฺคตํ นเหตุญฺจ ธมฺมํ ปฏิจฺจ ปริตฺโต นเหตุ ธมฺโม อุปฺปชฺชติ เหตุปจฺจยา. ตีณิ.}

\prose{ปริตฺตํ นเหตุญฺจ อปฺปมาณํ นเหตุญฺจ ธมฺมํ ปฏิจฺจ ปริตฺโต นเหตุ ธมฺโม อุปฺปชฺชติ เหตุปจฺจยา. (1) (สํขิตฺตํ.)}

\proseitem{163}{เหตุยา เตรส, อารมฺมเณ ปญฺจ, อธิปติยา นว ฯเปฯ วิคเต ปญฺจ, อวิคเต เตรส. (สํขิตฺตํ.)}

\csromanheader{2}{ธมฺมานุโลม ติกทุก\-ปฏฺฐาน\-ปาฬิ ปจฺจนีย\-น\-เหตุ}
\proseitem{164}{\csromanfolio{524}ปริตฺตํ นเหตุํ ธมฺมํ ปฏิจฺจ ปริตฺโต นเหตุ ธมฺโม อุปฺปชฺชติ นเหตุปจฺจยา. (1) (สํขิตฺตํ.)}

\prose{นเหตุยา เอกํ, นอารมฺมเณ ปญฺจ, นอธิปติยา ทส ฯเปฯ โนวิคเต ปญฺจ. (สํขิตฺตํ.)}

\prose{เหตุปจฺจยา นอารมฺมเณ ปญฺจ. (สํขิตฺตํ.)}

\prose{นเหตุปจฺจยา อารมฺมเณ เอกํ. (สํขิตฺตํ.)}

\prose{(สหชาตวารมฺปิ ฯเปฯ สมฺปยุตฺตวารมฺปิ ปฏิจฺจ\-วาร\-สทิสํ วิตฺถาเรตพฺพํ.)}

\csromanheader{2}{7. \textbf{ปญฺหาวาร ปจฺจยจตุกฺก}}
\proseitem{165}{ปริตฺโต นเหตุ ธมฺโม ปริตฺตสฺส นเหตุสฺส ธมฺมสฺส อารมฺมณ\-ปจฺจเยน ปจฺจโย.  ปริตฺโต นเหตุ ธมฺโม มหคฺคตสฺส นเหตุสฺส ธมฺมสฺส อารมฺมณ\-ปจฺจเยน ปจฺจโย. (2)}

\prose{มหคฺคโต นเหตุ ธมฺโม มหคฺคตสฺส นเหตุสฺส ธมฺมสฺส อารมฺมณ\-ปจฺจเยน ปจฺจโย.  มหคฺคโต นเหตุ ธมฺโม ปริตฺตสฺส นเหตุสฺส ธมฺมสฺส อารมฺมณ\-ปจฺจเยน ปจฺจโย. (2)}

\prose{อปฺปมาโณ นเหตุ ธมฺโม อปฺปมาณสฺส นเหตุสฺส ธมฺมสฺส อารมฺมณ\-ปจฺจเยน ปจฺจโย.  อปฺปมาโณ นเหตุ ธมฺโม ปริตฺตสฺส นเหตุสฺส ธมฺมสฺส อารมฺมณ\-ปจฺจเยน ปจฺจโย.  อปฺปมาโณ นเหตุ ธมฺโม มหคฺคตสฺส นเหตุสฺส ธมฺมสฺส อารมฺมณ\-ปจฺจเยน ปจฺจโย. (3)}

\prose{ปริตฺโต นเหตุ ธมฺโม ปริตฺตสฺส นเหตุสฺส ธมฺมสฺส อธิปติ\-ปจฺจเยน ปจฺจโย. (สํขิตฺตํ.)}

\proseitem{166}{อารมฺมเณ สตฺต, อธิปติยา สตฺต, อนนฺตเร นว, สมนนฺตเร นว, สหชาเต เอกาทส, อญฺญมญฺเญ สตฺต, นิสฺสเย เตรส, อุปนิสฺสเย นว ฯเปฯ อวิคเต เตรส. (สํขิตฺตํ.)}

\csromanmatikaanchor{mtk.116}
\csromanmatikaanchor{mtk.117}
\csromantocmarknopage{subparagraph}{13. ปริตฺตารมฺมณตฺติก}{mtk.116}
\bookmark[dest={mtk.116},level=5]{13. ปริตฺตารมฺมณตฺติก}
\csromantocmarkref{subparagraph}{1. เหตุทุก}{mtk.117}
\bookmark[dest={mtk.117},level=5]{1. เหตุทุก}
\csromantocmarknopage{subparagraph}{14. หีนตฺติก}{mtk.118}
\bookmark[dest={mtk.118},level=5]{14. หีนตฺติก}
\csromanmarkh{13. ปริตฺตา\-รมฺมณ\-ตฺติก}\csromanmarkhii{1. เหตุทุก}\csromanheadernomark{6}{13. ปริตฺตา\-รมฺมณ\-ตฺติก 1. เหตุทุก}
\prose{\csromanfolio{525}นเหตุยา ปนฺนรส, นอารมฺมเณ ปนฺนรส. (สํขิตฺตํ.)}

\prose{อารมฺมณ\-ปจฺจยา นเหตุยา สตฺต. (สํขิตฺตํ.)}

\prose{นเหตุปจฺจยา อารมฺมเณ สตฺต. (สํขิตฺตํ.)}

\prose{(ยถา กุสลตฺติเก ปญฺหาวารํ, เอวํ วิตฺถาเรตพฺพํ.)}

\csromanmarkh{13. ปริตฺตา\-รมฺมณ\-ตฺติก}\csromanmarkhii{1. เหตุทุก}\csromanheadernomark{2}{13. ปริตฺตา\-รมฺมณ\-ตฺติก 1. เหตุทุก}
\csromanheader{2}{1-7. ปฏิจฺจาทิวาร ปจฺจยจตุกฺกเหตุ}
\proseitem{167}{ปริตฺตา\-รมฺมณํ เหตุํ ธมฺมํ ปฏิจฺจ ปริตฺตารมฺมโณ เหตุ ธมฺโม อุปฺปชฺชติ เหตุปจฺจยา. (1)}

\prose{มหคฺคตา\-รมฺมณํ เหตุํ ธมฺมํ ปฏิจฺจ มหคฺคตา\-รมฺมโณ เหตุ ธมฺโม อุปฺปชฺชติ เหตุปจฺจยา. (1)}

\prose{อปฺปมาณา\-รมฺมณํ เหตุํ ธมฺมํ ปฏิจฺจ อปฺปมาณารมฺมโณ เหตุ ธมฺโม อุปฺปชฺชติ เหตุปจฺจยา. (1) (สํขิตฺตํ.)}

\prose{เหตุยา ตีณิ, อารมฺมเณ ตีณิ ฯเปฯ วิคเต ตีณิ, อวิคเต ตีณิ. (สํขิตฺตํ.)}

\prose{นอธิปติยา ตีณิ, นปุเรชาเต ตีณิ, นปจฺฉาชาเต ตีณิ, นอาเสวเน ตีณิ, นวิปาเก ตีณิ, นวิปฺปยุตฺเต ตีณิ. (สํขิตฺตํ.)}

\prose{(สหชาตวารมฺปิ ฯเปฯ สมฺปยุตฺตวารมฺปิ ปฏิจฺจ\-วาร\-สทิสํ วิตฺถาเรตพฺพํ.)}

\prose{เหตุอารมฺมณ}

\proseitem{168}{ปริตฺตารมฺมโณ เหตุ ธมฺโม ปริตฺตา\-รมฺมณสฺส เหตุสฺส ธมฺมสฺส เหตุปจฺจเยน ปจฺจโย. (1)}

\prose{มหคฺคตา\-รมฺมโณ เหตุ ธมฺโม มหคฺคตา\-รมฺมณสฺส เหตุสฺส ธมฺมสฺส เหตุปจฺจเยน ปจฺจโย. (1)}

\prose{อปฺปมาณารมฺมโณ เหตุ ธมฺโม อปฺปมาณา\-รมฺมณสฺส เหตุสฺส ธมฺมสฺส เหตุปจฺจเยน ปจฺจโย. (1)}

\prose{\csromanfolio{526}ปริตฺตารมฺมโณ เหตุ ธมฺโม ปริตฺตา\-รมฺมณสฺส เหตุสฺส ธมฺมสฺส อารมฺมณ\-ปจฺจเยน ปจฺจโย.  ปริตฺตารมฺมโณ เหตุ ธมฺโม มหคฺคตา\-รมฺมณสฺส เหตุสฺส ธมฺมสฺส อารมฺมณ\-ปจฺจเยน ปจฺจโย. (2)}

\prose{มหคฺคตา\-รมฺมโณ เหตุ ธมฺโม มหคฺคตา\-รมฺมณสฺส เหตุสฺส ธมฺมสฺส อารมฺมณ\-ปจฺจเยน ปจฺจโย.  มหคฺคตา\-รมฺมโณ เหตุ ธมฺโม ปริตฺตา\-รมฺมณสฺส เหตุสฺส ธมฺมสฺส อารมฺมณ\-ปจฺจเยน ปจฺจโย. (2)}

\prose{อปฺปมาณารมฺมโณ เหตุ ธมฺโม อปฺปมาณา\-รมฺมณสฺส เหตุสฺส ธมฺมสฺส อารมฺมณ\-ปจฺจเยน ปจฺจโย.  อปฺปมาณารมฺมโณ เหตุ ธมฺโม ปริตฺตา\-รมฺมณสฺส เหตุสฺส ธมฺมสฺส อารมฺมณ\-ปจฺจเยน ปจฺจโย.  อปฺปมาณารมฺมโณ เหตุ ธมฺโม มหคฺคตา\-รมฺมณสฺส เหตุสฺส ธมฺมสฺส อารมฺมณ\-ปจฺจเยน ปจฺจโย. (3) (สํขิตฺตํ.)}

\proseitem{169}{เหตุยา ตีณิ, อารมฺมเณ สตฺต, อธิปติยา สตฺต, อนนฺตเร นว, สมนนฺตเร นว, อุปนิสฺสเย นว ฯเปฯ อวิคเต ตีณิ. (สํขิตฺตํ.)}

\prose{นเหตุยา นว, นอารมฺมเณ นว. (สํขิตฺตํ.)}

\prose{เหตุปจฺจยา นอารมฺมเณ ตีณิ. (สํขิตฺตํ.)}

\prose{นเหตุปจฺจยา อารมฺมเณ สตฺต. (สํขิตฺตํ.)}

\prose{(ยถา กุสลตฺติเก ปญฺหาวารํ, เอวํ วิตฺถาเรตพฺพํ.)}

\csromanheader{2}{\textbf{นเหตุปท 1-7. ปฏิจฺจาทิวาร}}
\prose{ปจฺจยจตุกฺกเหตุ}

\proseitem{170}{ปริตฺตา\-รมฺมณํ นเหตุํ ธมฺมํ ปฏิจฺจ ปริตฺตารมฺมโณ นเหตุ ธมฺโม อุปฺปชฺชติ เหตุปจฺจยา. (สํขิตฺตํ.)}

\prose{เหตุยา ตีณิ, อารมฺมเณ ตีณิ, อธิปติยา ตีณิ ฯเปฯ อวิคเต ตีณิ. (สํขิตฺตํ.)}

\csromanmarkh{13. ปริตฺตา\-รมฺมณ\-ตฺติก}\csromanmarkhii{1. เหตุทุก}\csromanheadernomark{2}{13. ปริตฺตา\-รมฺมณ\-ตฺติก 1. เหตุทุก}
\proseitem{171}{\csromanfolio{527}ปริตฺตา\-รมฺมณํ นเหตุํ ธมฺมํ ปฏิจฺจ ปริตฺตารมฺมโณ นเหตุ ธมฺโม อุปฺปชฺชติ นเหตุปจฺจยา. (สํขิตฺตํ.)}

\prose{นเหตุยา ตีณิ, นอธิปติยา ตีณิ ฯเปฯ นวิปฺปยุตฺเต ตีณิ. (สํขิตฺตํ.)}

\prose{เหตุปจฺจยา นอธิปติยา ตีณิ. (สํขิตฺตํ.)}

\prose{นเหตุปจฺจยา อารมฺมเณ ตีณิ. (สํขิตฺตํ.)}

\prose{(สหชาตวารมฺปิ ฯเปฯ สมฺปยุตฺตวารมฺปิ ปฏิจฺจ\-วาร\-สทิสํ วิตฺถาเรตพฺพํ.)}

\prose{อารมฺมณ\-อธิปติ}

\proseitem{172}{ปริตฺตารมฺมโณ นเหตุ ธมฺโม ปริตฺตา\-รมฺมณสฺส นเหตุสฺส ธมฺมสฺส อารมฺมณ\-ปจฺจเยน ปจฺจโย.  ปริตฺตารมฺมโณ นเหตุ ธมฺโม มหคฺคตา\-รมฺมณสฺส นเหตุสฺส ธมฺมสฺส อารมฺมณ\-ปจฺจเยน ปจฺจโย. (2)}

\prose{มหคฺคตา\-รมฺมโณ นเหตุ ธมฺโม มหคฺคตา\-รมฺมณสฺส นเหตุสฺส ธมฺมสฺส อารมฺมณ\-ปจฺจเยน ปจฺจโย.  มหคฺคตา\-รมฺมโณ นเหตุ ธมฺโม ปริตฺตา\-รมฺมณสฺส นเหตุสฺส ธมฺมสฺส อารมฺมณ\-ปจฺจเยน ปจฺจโย. (2)}

\prose{อปฺปมาณารมฺมโณ นเหตุ ธมฺโม อปฺปมาณา\-รมฺมณสฺส นเหตุสฺส ธมฺมสฺส อารมฺมณ\-ปจฺจเยน ปจฺจโย.  อปฺปมาณารมฺมโณ นเหตุ ธมฺโม ปริตฺตา\-รมฺมณสฺส นเหตุสฺส ธมฺมสฺส อารมฺมณ\-ปจฺจเยน ปจฺจโย.  อปฺปมาณารมฺมโณ นเหตุ ธมฺโม มหคฺคตา\-รมฺมณสฺส นเหตุสฺส ธมฺมสฺส อารมฺมณ\-ปจฺจเยน ปจฺจโย. (3)}

\prose{ปริตฺตารมฺมโณ นเหตุ ธมฺโม ปริตฺตา\-รมฺมณสฺส นเหตุสฺส ธมฺมสฺส อธิปติ\-ปจฺจเยน ปจฺจโย. (สํขิตฺตํ.)}

\proseitem{173}{อารมฺมเณ สตฺต, อธิปติยา สตฺต, อนนฺตเร นว, สมนนฺตเร นว, อุปนิสฺสเย นว ฯเปฯ อวิคเต ตีณิ. (สํขิตฺตํ.)}

\prose{นเหตุยา นว, นอารมฺมเณ นว. (สํขิตฺตํ.)}

\prose{อารมฺมณ\-ปจฺจยา นเหตุยา สตฺต. (สํขิตฺตํ.)}

\prose{นเหตุปจฺจยา อารมฺมเณ สตฺต. (สํขิตฺตํ.)}

\prose{(ยถา กุสลตฺติเก ปญฺหาวารํ, เอวํ วิตฺถาเรตพฺพํ.)}

\csromanmatikaanchor{mtk.119}
\csromantocmarkref{subparagraph}{1. เหตุทุก}{mtk.119}
\bookmark[dest={mtk.119},level=5]{1. เหตุทุก}
\csromanmarkh{ธมฺมานุโลม ติกทุก\-ปฏฺฐาน\-ปาฬิ}\csromanmarkhii{14. หีนตฺติก 1. เหตุทุก}\csromanheadernomark{6}{ธมฺมานุโลม ติกทุก\-ปฏฺฐาน\-ปาฬิ 14. หีนตฺติก 1. เหตุทุก}
\csromanheader{2}{1-7. ปฏิจฺจาทิวาร ปจฺจยจตุกฺกเหตุ}
\proseitem{174}{\csromanfolio{528}หีนํ เหตุํ ธมฺมํ ปฏิจฺจ หีโน เหตุ ธมฺโม อุปฺปชฺชติ เหตุปจฺจยา.}

\prose{มชฺฌิมํ เหตุํ ธมฺมํ ปฏิจฺจ มชฺฌิโม เหตุ ธมฺโม อุปฺปชฺชติ เหตุปจฺจยา.}

\prose{ปณีตํ เหตุํ ธมฺมํ ปฏิจฺจ ปณีโต เหตุ ธมฺโม อุปฺปชฺชติ เหตุปจฺจยา. (3) (สํขิตฺตํ.)}

\proseitem{175}{เหตุยา ตีณิ, อารมฺมเณ ตีณิ ฯเปฯ อวิคเต ตีณิ. (สํขิตฺตํ.)}

\prose{นอธิปติยา ตีณิ, นปุเรชาเต ตีณิ, นปจฺฉาชาเต ตีณิ, นอาเสวเน ตีณิ, นวิปาเก ตีณิ, นวิปฺปยุตฺเต ตีณิ. (สํขิตฺตํ.)}

\prose{(สหชาตวารมฺปิ ฯเปฯ สมฺปยุตฺตวารมฺปิ ปฏิจฺจ\-วาร\-สทิสํ วิตฺถาเรตพฺพํ.)}

\prose{เหตุอารมฺมณาทิ}

\proseitem{176}{หีโน เหตุ ธมฺโม หีนสฺส เหตุสฺส ธมฺมสฺส เหตุปจฺจเยน ปจฺจโย. (1)}

\prose{มชฺฌิโม เหตุ ธมฺโม มชฺฌิมสฺส เหตุสฺส ธมฺมสฺส เหตุปจฺจเยน ปจฺจโย. (1)}

\prose{ปณีโต เหตุ ธมฺโม ปณีตสฺส เหตุสฺส ธมฺมสฺส เหตุปจฺจเยน ปจฺจโย. (1)}

\prose{หีโน เหตุ ธมฺโม หีนสฺส เหตุสฺส ธมฺมสฺส อารมฺมณ\-ปจฺจเยน ปจฺจโย.  หีโน เหตุ ธมฺโม มชฺฌิมสฺส เหตุสฺส ธมฺมสฺส อารมฺมณ\-ปจฺจเยน ปจฺจโย. (2)}

\prose{มชฺฌิโม เหตุ ธมฺโม มชฺฌิมสฺส เหตุสฺส ธมฺมสฺส อารมฺมณ\-ปจฺจเยน ปจฺจโย.  มชฺฌิโม เหตุ ธมฺโม หีนสฺส เหตุสฺส ธมฺมสฺส อารมฺมณ\-ปจฺจเยน ปจฺจโย. (2)}

\prose{ปณีโต เหตุ ธมฺโม มชฺฌิมสฺส เหตุสฺส ธมฺมสฺส อารมฺมณ\-ปจฺจเยน ปจฺจโย. (1)}

\csromanmarkh{14. หีนตฺติก}\csromanmarkhii{1. เหตุทุก}\csromanheadernomark{2}{14. หีนตฺติก 1. เหตุทุก}
\proseitem{177}{\csromanfolio{529}หีโน เหตุ ธมฺโม หีนสฺส เหตุสฺส ธมฺมสฺส อธิปติ\-ปจฺจเยน ปจฺจโย. (1)}

\prose{มชฺฌิโม เหตุ ธมฺโม มชฺฌิมสฺส เหตุสฺส ธมฺมสฺส อธิปติ\-ปจฺจเยน ปจฺจโย.  มชฺฌิโม เหตุ ธมฺโม หีนสฺส เหตุสฺส ธมฺมสฺส อธิปติ\-ปจฺจเยน ปจฺจโย. (2)}

\prose{ปณีโต เหตุ ธมฺโม ปณีตสฺส เหตุสฺส ธมฺมสฺส อธิปติ\-ปจฺจเยน ปจฺจโย.  ปณีโต เหตุ ธมฺโม มชฺฌิมสฺส เหตุสฺส ธมฺมสฺส อธิปติ\-ปจฺจเยน ปจฺจโย. (2)}

\prose{หีโน เหตุ ธมฺโม หีนสฺส เหตุสฺส ธมฺมสฺส อนนฺตร\-ปจฺจเยน ปจฺจโย. เทฺว.}

\prose{มชฺฌิโม เหตุ ธมฺโม มชฺฌิมสฺส เหตุสฺส ธมฺมสฺส อนนฺตร\-ปจฺจเยน ปจฺจโย. เทฺว.}

\prose{ปณีโต เหตุ ธมฺโม ปณีตสฺส เหตุสฺส ธมฺมสฺส อนนฺตร\-ปจฺจเยน ปจฺจโย. เทฺว. (สํขิตฺตํ.)}

\proseitem{178}{เหตุยา ตีณิ, อารมฺมเณ ปญฺจ, อธิปติยา ปญฺจ ฯเปฯ อวิคเต ตีณิ. (สํขิตฺตํ.)}

\prose{นเหตุยา อฏฺฐ, นอารมฺมเณ อฏฺฐ. (สํขิตฺตํ.)}

\prose{เหตุปจฺจยา นอารมฺมเณ ตีณิ. (สํขิตฺตํ.)}

\prose{นเหตุปจฺจยา อารมฺมเณ ปญฺจ. (สํขิตฺตํ.)}

\prose{(ยถา กุสลตฺติเก ปญฺหาวารํ, เอวํ วิตฺถาเรตพฺพํ.)}

\prose{นเหตุปทเหตุ}

\proseitem{179}{หีนํ นเหตุํ ธมฺมํ ปฏิจฺจ หีโน นเหตุ ธมฺโม อุปฺปชฺชติ เหตุปจฺจยา.}

\prose{(ยถา สํกิลิฏฺฐ\-ตฺติก\-เหตุทุกํ, เอวํ วิตฺถาเรตพฺพํ.)}

\csromanmatikaanchor{mtk.120}
\csromanmatikaanchor{mtk.121}
\csromantocmarknopage{subparagraph}{15. มิจฺฉตฺตนิยตตฺติก}{mtk.120}
\bookmark[dest={mtk.120},level=5]{15. มิจฺฉตฺตนิยตตฺติก}
\csromantocmarkref{subparagraph}{1. เหตุทุก}{mtk.121}
\bookmark[dest={mtk.121},level=5]{1. เหตุทุก}
\csromanmarkh{ธมฺมานุโลม ติกทุก\-ปฏฺฐาน\-ปาฬิ}\csromanmarkhii{15. มิจฺฉตฺตนิยต\-ตฺติก 1. เหตุทุก}\csromanheadernomark{6}{ธมฺมานุโลม ติกทุก\-ปฏฺฐาน\-ปาฬิ 15. มิจฺฉตฺตนิยต\-ตฺติก 1. เหตุทุก}
\csromanheader{2}{1-7. ปฏิจฺจาทิวาร ปจฺจยจตุกฺกเหตุ}
\proseitem{180}{\csromanfolio{530}มิจฺฉตฺต\-นิยตํ เหตุํ ธมฺมํ ปฏิจฺจ มิจฺฉตฺต\-นิยโต เหตุ ธมฺโม อุปฺปชฺชติ เหตุปจฺจยา. (1)}

\prose{สมฺมตฺต\-นิยตํ เหตุํ ธมฺมํ ปฏิจฺจ สมฺมตฺตนิยโต เหตุ ธมฺโม อุปฺปชฺชติ เหตุปจฺจยา. (1)}

\prose{อนิยตํ เหตุํ ธมฺมํ ปฏิจฺจ อนิยโต เหตุ ธมฺโม อุปฺปชฺชติ เหตุปจฺจยา. (1) (สํขิตฺตํ.)}

\proseitem{181}{เหตุยา ตีณิ, อารมฺมเณ ตีณิ, อธิปติยา ตีณิ ฯเปฯ วิปาเก เอกํ ฯเปฯ อวิคเต ตีณิ. (สํขิตฺตํ.)}

\prose{นอธิปติยา เทฺว, นปุเรชาเต เทฺว, นปจฺฉาชาเต ตีณิ, นอาเสวเน เอกํ, นวิปาเก ตีณิ, นวิปฺปยุตฺเต เทฺว. (สํขิตฺตํ.)}

\prose{(สหชาตวารมฺปิ ฯเปฯ สมฺปยุตฺตวารมฺปิ ปฏิจฺจ\-วาร\-สทิสํ วิตฺถาเรตพฺพํ.)}

\csromanheader{2}{\textbf{เหตุ อารมฺมณาทิ}}
\proseitem{182}{มิจฺฉตฺต\-นิยโต เหตุ ธมฺโม มิจฺฉตฺต\-นิยตสฺส เหตุสฺส ธมฺมสฺส เหตุปจฺจเยน ปจฺจโย. (1)}

\prose{สมฺมตฺตนิยโต เหตุ ธมฺโม สมฺมตฺต\-นิยตสฺส เหตุสฺส ธมฺมสฺส เหตุปจฺจเยน ปจฺจโย. (1)}

\prose{อนิยโต เหตุ ธมฺโม อนิยตสฺส เหตุสฺส ธมฺมสฺส เหตุปจฺจเยน ปจฺจโย. (1)}

\prose{มิจฺฉตฺต\-นิยโต เหตุ ธมฺโม อนิยตสฺส เหตุสฺส ธมฺมสฺส อารมฺมณ\-ปจฺจเยน ปจฺจโย. (1)}

\prose{สมฺมตฺตนิยโต เหตุ ธมฺโม อนิยตสฺส เหตุสฺส ธมฺมสฺส อารมฺมณ\-ปจฺจเยน ปจฺจโย. (1)}

\prose{อนิยโต เหตุ ธมฺโม อนิยตสฺส เหตุสฺส ธมฺมสฺส อารมฺมณ\-ปจฺจเยน ปจฺจโย. (1)}

\csromanmarkh{15. มิจฺฉตฺตนิยต\-ตฺติก}\csromanmarkhii{1. เหตุทุก}\csromanheadernomark{2}{15. มิจฺฉตฺตนิยต\-ตฺติก 1. เหตุทุก}
\prose{\csromanfolio{531}อนิยโต เหตุ ธมฺโม มิจฺฉตฺต\-นิยตสฺส เหตุสฺส ธมฺมสฺส อารมฺมณ\-ปจฺจเยน ปจฺจโย. (1)}

\prose{สมฺมตฺตนิยโต เหตุ ธมฺโม สมฺมตฺต\-นิยตสฺส เหตุสฺส ธมฺมสฺส อธิปติ\-ปจฺจเยน ปจฺจโย.  สมฺมตฺตนิยโต เหตุ ธมฺโม อนิยตสฺส เหตุสฺส ธมฺมสฺส อธิปติ\-ปจฺจเยน ปจฺจโย. (2)}

\prose{อนิยโต เหตุ ธมฺโม อนิยตสฺส เหตุสฺส ธมฺมสฺส อธิปติ\-ปจฺจเยน ปจฺจโย. (1)}

\proseitem{183}{มิจฺฉตฺต\-นิยโต เหตุ ธมฺโม อนิยตสฺส เหตุสฺส ธมฺมสฺส อนนฺตร\-ปจฺจเยน ปจฺจโย. (1)}

\prose{สมฺมตฺตนิยโต เหตุ ธมฺโม อนิยตสฺส เหตุสฺส ธมฺมสฺส อนนฺตร\-ปจฺจเยน ปจฺจโย. (1)}

\prose{อนิยโต เหตุ ธมฺโม อนิยตสฺส เหตุสฺส ธมฺมสฺส อนนฺตร\-ปจฺจเยน ปจฺจโย.  อนิยโต เหตุ ธมฺโม มิจฺฉตฺต\-นิยตสฺส เหตุสฺส ธมฺมสฺส อนนฺตร\-ปจฺจเยน ปจฺจโย.  อนิยโต เหตุ ธมฺโม สมฺมตฺต\-นิยตสฺส เหตุสฺส ธมฺมสฺส อนนฺตร\-ปจฺจเยน ปจฺจโย. (3) ฯเปฯ.}

\proseitem{184}{มิจฺฉตฺต\-นิยโต เหตุ ธมฺโม มิจฺฉตฺต\-นิยตสฺส เหตุสฺส ธมฺมสฺส อุปนิสฺสย\-ปจฺจเยน ปจฺจโย.  มิจฺฉตฺต\-นิยโต เหตุ ธมฺโม อนิยตสฺส เหตุสฺส ธมฺมสฺส อุปนิสฺสย\-ปจฺจเยน ปจฺจโย. (2)}

\prose{สมฺมตฺตนิยโต เหตุ ธมฺโม สมฺมตฺต\-นิยตสฺส เหตุสฺส ธมฺมสฺส อุปนิสฺสย\-ปจฺจเยน ปจฺจโย.  สมฺมตฺตนิยโต เหตุ ธมฺโม อนิยตสฺส เหตุสฺส ธมฺมสฺส อุปนิสฺสย\-ปจฺจเยน ปจฺจโย. (2)}

\prose{อนิยโต เหตุ ธมฺโม อนิยตสฺส เหตุสฺส ธมฺมสฺส อุปนิสฺสย\-ปจฺจเยน ปจฺจโย.  อนิยโต เหตุ ธมฺโม มิจฺฉตฺต\-นิยตสฺส เหตุสฺส ธมฺมสฺส อุปนิสฺสย\-ปจฺจเยน ปจฺจโย.  อนิยโต เหตุ ธมฺโม สมฺมตฺต\-นิยตสฺส เหตุสฺส ธมฺมสฺส อุปนิสฺสย\-ปจฺจเยน ปจฺจโย. (3) (สํขิตฺตํ.)}

\proseitem{185}{\csromanfolio{532}เหตุยา ตีณิ, อารมฺมเณ จตฺตาริ, อธิปติยา ตีณิ, อนนฺตเร ปญฺจ, สมนนฺตเร ปญฺจ, สหชาเต ตีณิ ฯเปฯ อุปนิสฺสเย สตฺต ฯเปฯ อวิคเต ตีณิ. (สํขิตฺตํ.)}

\prose{นเหตุยา สตฺต, นอารมฺมเณ สตฺต. (สํขิตฺตํ.)}

\prose{เหตุปจฺจยา นอารมฺมเณ ตีณิ. (สํขิตฺตํ.)}

\prose{นเหตุปจฺจยา อารมฺมเณ จตฺตาริ. (สํขิตฺตํ.)}

\prose{(ยถา กุสลตฺติเก ปญฺหาวารํ, เอวํ วิตฺถาเรตพฺพํ.)}

\csromanheader{2}{\textbf{นเหตุปท 1-7. ปฏิจฺจาทิวาร}}
\prose{ปจฺจยจตุกฺกเหตุ}

\proseitem{186}{มิจฺฉตฺต\-นิยตํ นเหตุํ ธมฺมํ ปฏิจฺจ มิจฺฉตฺต\-นิยโต นเหตุ ธมฺโม อุปฺปชฺชติ เหตุปจฺจยา. ตีณิ.}

\prose{สมฺมตฺต\-นิยตํ นเหตุํ ธมฺมํ ปฏิจฺจ สมฺมตฺตนิยโต นเหตุ ธมฺโม อุปฺปชฺชติ เหตุปจฺจยา.  สมฺมตฺต\-นิยตํ นเหตุํ ธมฺมํ ปฏิจฺจ อนิยโต นเหตุ ธมฺโม อุปฺปชฺชติ เหตุปจฺจยา.  สมฺมตฺต\-นิยตํ นเหตุํ ธมฺมํ ปฏิจฺจ สมฺมตฺตนิยโต นเหตุ จ อนิยโต นเหตุ จ ธมฺมา อุปฺปชฺชนฺติ เหตุปจฺจยา. (3)}

\prose{อนิยตํ นเหตุํ ธมฺมํ ปฏิจฺจ อนิยโต นเหตุ ธมฺโม อุปฺปชฺชติ เหตุปจฺจยา. (1)}

\prose{มิจฺฉตฺต\-นิยตํ นเหตุญฺจ อนิยตํ นเหตุญฺจ ธมฺมํ ปฏิจฺจ อนิยโต นเหตุ ธมฺโม อุปฺปชฺชติ เหตุปจฺจยา. (1)}

\prose{สมฺมตฺต\-นิยตํ นเหตุญฺจ อนิยตํ นเหตุญฺจ ธมฺมํ ปฏิจฺจ อนิยโต นเหตุ ธมฺโม อุปฺปชฺชติ เหตุปจฺจยา. (1) (สํขิตฺตํ.) เหตุยา นว, อารมฺมเณ ตีณิ, อธิปติยา นว ฯเปฯ วิปาเก เอกํ ฯเปฯ อวิคเต นว. (สํขิตฺตํ.)}

\csromanmarkh{15. มิจฺฉตฺตนิยต\-ตฺติก}\csromanmarkhii{1. เหตุทุก}\csromanheadernomark{2}{15. มิจฺฉตฺตนิยต\-ตฺติก 1. เหตุทุก}
\prose{\csromanfolio{533}ปจฺจนีย\-น\-เหตุ}

\proseitem{188}{อนิยตํ นเหตุํ ธมฺมํ ปฏิจฺจ อนิยโต นเหตุ ธมฺโม อุปฺปชฺชติ นเหตุปจฺจยา. (สํขิตฺตํ.)}

\prose{นเหตุยา เอกํ, นอารมฺมเณ ปญฺจ, นอธิปติยา ตีณิ ฯเปฯ โนวิคเต ปญฺจ. (สํขิตฺตํ.)}

\prose{เหตุปจฺจยา นอารมฺมเณ ปญฺจ. (สํขิตฺตํ.)}

\prose{นเหตุปจฺจยา อารมฺมเณ เอกํ. (สํขิตฺตํ.)}

\prose{(สหชาตวารมฺปิ ฯเปฯ สมฺปยุตฺตวารมฺปิ ปฏิจฺจ\-วาร\-สทิสํ วิตฺถาเรตพฺพํ.)}

\csromanheader{2}{\textbf{อารมฺมณ อธิปติ อนนฺตร}}
\proseitem{189}{มิจฺฉตฺต\-นิยโต นเหตุ ธมฺโม อนิยตสฺส นเหตุสฺส ธมฺมสฺส อารมฺมณ\-ปจฺจเยน ปจฺจโย. (1)}

\prose{สมฺมตฺตนิยโต นเหตุ ธมฺโม อนิยตสฺส นเหตุสฺส ธมฺมสฺส อารมฺมณ\-ปจฺจเยน ปจฺจโย. (1)}

\prose{อนิยโต นเหตุ ธมฺโม อนิยตสฺส นเหตุสฺส ธมฺมสฺส อารมฺมณ\-ปจฺจเยน ปจฺจโย.  อนิยโต นเหตุ ธมฺโม มิจฺฉตฺต\-นิยตสฺส นเหตุสฺส ธมฺมสฺส อารมฺมณ\-ปจฺจเยน ปจฺจโย.  อนิยโต นเหตุ ธมฺโม สมฺมตฺต\-นิยตสฺส นเหตุสฺส ธมฺมสฺส อารมฺมณ\-ปจฺจเยน ปจฺจโย. (3)}

\proseitem{190}{มิจฺฉตฺต\-นิยโต นเหตุ ธมฺโม มิจฺฉตฺต\-นิยตสฺส นเหตุสฺส ธมฺมสฺส อธิปติ\-ปจฺจเยน ปจฺจโย. ตีณิ. สมฺมตฺตนิยโต นเหตุ ธมฺโม สมฺมตฺต\-นิยตสฺส นเหตุสฺส ธมฺมสฺส อธิปติ\-ปจฺจเยน ปจฺจโย.  สมฺมตฺตนิยโต นเหตุ ธมฺโม อนิยตสฺส นเหตุสฺส ธมฺมสฺส อธิปติ\-ปจฺจเยน ปจฺจโย.  สมฺมตฺตนิยโต นเหตุ ธมฺโม สมฺมตฺต\-นิยตสฺส นเหตุสฺส จ อนิยตสฺส นเหตุสฺส จ ธมฺมสฺส อธิปติ\-ปจฺจเยน ปจฺจโย. (3)}

\prose{\csromanfolio{534}อนิยโต นเหตุ ธมฺโม อนิยตสฺส นเหตุสฺส ธมฺมสฺส อธิปติ\-ปจฺจเยน ปจฺจโย.  อนิยโต นเหตุ ธมฺโม สมฺมตฺต\-นิยตสฺส นเหตุสฺส ธมฺมสฺส อธิปติ\-ปจฺจเยน ปจฺจโย. (2)}

\prose{มิจฺฉตฺต\-นิยโต นเหตุ ธมฺโม อนิยตสฺส นเหตุสฺส ธมฺมสฺส อนนฺตร\-ปจฺจเยน ปจฺจโย. (1)}

\prose{สมฺมตฺตนิยโต นเหตุ ธมฺโม อนิยตสฺส นเหตุสฺส ธมฺมสฺส อนนฺตร\-ปจฺจเยน ปจฺจโย. (1)}

\prose{อนิยโต นเหตุ ธมฺโม อนิยตสฺส นเหตุสฺส ธมฺมสฺส อนนฺตร\-ปจฺจเยน ปจฺจโย.  อนิยโต นเหตุ ธมฺโม มิจฺฉตฺต\-นิยตสฺส นเหตุสฺส ธมฺมสฺส อนนฺตร\-ปจฺจเยน ปจฺจโย.  อนิยโต นเหตุ ธมฺโม สมฺมตฺต\-นิยตสฺส นเหตุสฺส ธมฺมสฺส อนนฺตร\-ปจฺจเยน ปจฺจโย. (3) (สํขิตฺตํ.)}

\proseitem{191}{อารมฺมเณ ปญฺจ, อธิปติยา อฏฺฐ, อนนฺตเร ปญฺจ, สมนนฺตเร ปญฺจ, สหชาเต นว, อญฺญมญฺเญ ตีณิ, นิสฺสเย เตรส, อุปนิสฺสเย สตฺต ฯเปฯ อวิคเต เตรส. (สํขิตฺตํ.)}

\prose{นเหตุยา เตรส, นอารมฺมเณ เตรส, นอธิปติยา เตรส. (สํขิตฺตํ.)}

\prose{อารมฺมณ\-ปจฺจยา นเหตุยา ปญฺจ. (สํขิตฺตํ.)}

\prose{นเหตุปจฺจยา อารมฺมเณ ปญฺจ. (สํขิตฺตํ.)}

\prose{(ยถา กุสลตฺติเก ปญฺหาวารํ, เอวํ วิตฺถาเรตพฺพํ.)}

\csromanmatikaanchor{mtk.122}
\csromanmatikaanchor{mtk.123}
\csromantocmarknopage{subparagraph}{16. มคฺคารมฺมณตฺติก}{mtk.122}
\bookmark[dest={mtk.122},level=5]{16. มคฺคารมฺมณตฺติก}
\csromantocmarkref{subparagraph}{1. เหตุทุก}{mtk.123}
\bookmark[dest={mtk.123},level=5]{1. เหตุทุก}
\csromanmarkh{16. มคฺคา\-รมฺมณ\-ตฺติก}\csromanmarkhii{1. เหตุทุก}\csromanheadernomark{6}{16. มคฺคา\-รมฺมณ\-ตฺติก 1. เหตุทุก}
\csromanheader{2}{1-7. ปฏิจฺจาทิวาร ปจฺจยจตุกฺกเหตุ}
\proseitem{192}{มคฺคารมฺมณํ เหตุํ ธมฺมํ ปฏิจฺจ มคฺคารมฺมโณ เหตุ ธมฺโม อุปฺปชฺชติ เหตุปจฺจยา. ตีณิ.}

\csromanmarkh{16. มคฺคา\-รมฺมณ\-ตฺติก}\csromanmarkhii{1. เหตุทุก}\csromanheadernomark{2}{16. มคฺคา\-รมฺมณ\-ตฺติก 1. เหตุทุก}
\prose{\csromanfolio{535}มคฺคเหตุกํ เหตุํ ธมฺมํ ปฏิจฺจ มคฺคเหตุโก เหตุ ธมฺโม อุปฺปชฺชติ เหตุปจฺจยา. ตีณิ.}

\prose{มคฺคาธิปติํ เหตุํ ธมฺมํ ปฏิจฺจ มคฺคาธิปติ เหตุ ธมฺโม อุปฺปชฺชติ เหตุปจฺจยา. ปญฺจ.}

\prose{มคฺคารมฺมณํ เหตุญฺจ มคฺคาธิปติํ เหตุญฺจ ธมฺมํ ปฏิจฺจ มคฺคารมฺมโณ เหตุ ธมฺโม อุปฺปชฺชติ เหตุปจฺจยา. ตีณิ.}

\prose{มคฺคเหตุกํ เหตุญฺจ มคฺคาธิปติํ เหตุญฺจ ธมฺมํ ปฏิจฺจ มคฺคเหตุโก เหตุ ธมฺโม อุปฺปชฺชติ เหตุปจฺจยา. ตีณิ. (สํขิตฺตํ.)}

\prose{193. เหตุยา สตฺตรส ฯเปฯ อวิคเต สตฺตรส. (สํขิตฺตํ.)}

\prose{นอธิปติยา สตฺตรส, นปุเรชาเต สตฺตรส ฯเปฯ นอาเสวเน นว, นวิปาเก สตฺตรส, นวิปฺปยุตฺเต สตฺตรส. (สํขิตฺตํ.)}

\prose{(สหชาตวารมฺปิ ฯเปฯ สมฺปยุตฺตวารมฺปิ ปฏิจฺจ\-วาร\-สทิสํ วิตฺถาเรตพฺพํ.)}

\csromanheader{2}{\textbf{เหตุ อารมฺมณาทิ}}
\proseitem{194}{มคฺคารมฺมโณ เหตุ ธมฺโม มคฺคา\-รมฺมณสฺส เหตุสฺส ธมฺมสฺส เหตุปจฺจเยน ปจฺจโย. ตีณิ.}

\prose{มคฺคเหตุโก เหตุ ธมฺโม มคฺคเหตุกสฺส เหตุสฺส ธมฺมสฺส เหตุปจฺจเยน ปจฺจโย. ตีณิ.}

\prose{มคฺคาธิปติ เหตุ ธมฺโม มคฺคา\-ธิปติสฺส เหตุสฺส ธมฺมสฺส เหตุปจฺจเยน ปจฺจโย. ปญฺจ.}

\prose{มคฺคารมฺมโณ เหตุ จ มคฺคาธิปติ เหตุ จ ธมฺมา มคฺคา\-รมฺมณสฺส เหตุสฺส ธมฺมสฺส เหตุปจฺจเยน ปจฺจโย. ตีณิ.}

\prose{มคฺคเหตุโก เหตุ จ มคฺคาธิปติ เหตุ จ ธมฺมา มคฺคเหตุกสฺส เหตุสฺส ธมฺมสฺส เหตุปจฺจเยน ปจฺจโย. ตีณิ.}

\prose{มคฺคเหตุโก เหตุ ธมฺโม มคฺคา\-รมฺมณสฺส เหตุสฺส ธมฺมสฺส อารมฺมณ\-ปจฺจเยน ปจฺจโย.  มคฺคเหตุโก เหตุ ธมฺโม มคฺคา\-ธิปติสฺส เหตุสฺส ธมฺมสฺส อารมฺมณ\-ปจฺจเยน ปจฺจโย.  มคฺคเหตุโก เหตุ ธมฺโม มคฺคา\-รมฺมณสฺส เหตุสฺส จ มคฺคา\-ธิปติสฺส เหตุสฺส จ ธมฺมสฺส อารมฺมณ\-ปจฺจเยน ปจฺจโย. (3)}

\prose{\csromanfolio{536}มคฺคาธิปติ เหตุ ธมฺโม มคฺคา\-ธิปติสฺส เหตุสฺส ธมฺมสฺส อารมฺมณ\-ปจฺจเยน ปจฺจโย. ตีณิ.}

\prose{มคฺคเหตุโก เหตุ จ มคฺคาธิปติ เหตุ จ ธมฺมา มคฺคา\-รมฺมณสฺส เหตุสฺส ธมฺมสฺส อารมฺมณ\-ปจฺจเยน ปจฺจโย. ตีณิ.}

\proseitem{195}{มคฺคารมฺมโณ เหตุ ธมฺโม มคฺคา\-รมฺมณสฺส เหตุสฺส ธมฺมสฺส อธิปติ\-ปจฺจเยน ปจฺจโย. ตีณิ.}

\prose{มคฺคเหตุโก เหตุ ธมฺโม มคฺคเหตุกสฺส เหตุสฺส ธมฺมสฺส อธิปติ\-ปจฺจเยน ปจฺจโย. ปญฺจ.}

\prose{มคฺคาธิปติ เหตุ ธมฺโม มคฺคา\-ธิปติสฺส เหตุสฺส ธมฺมสฺส อธิปติ\-ปจฺจเยน ปจฺจโย. ปญฺจ.}

\prose{มคฺคารมฺมโณ เหตุ จ มคฺคาธิปติ เหตุ จ ธมฺมา มคฺคา\-รมฺมณสฺส เหตุสฺส ธมฺมสฺส อธิปติ\-ปจฺจเยน ปจฺจโย. ตีณิ.}

\prose{มคฺคเหตุโก เหตุ จ มคฺคาธิปติ เหตุ จ ธมฺมา มคฺคา\-รมฺมณสฺส เหตุสฺส ธมฺมสฺส อธิปติ\-ปจฺจเยน ปจฺจโย. ปญฺจ.}

\prose{มคฺคารมฺมโณ เหตุ ธมฺโม มคฺคา\-รมฺมณสฺส เหตุสฺส ธมฺมสฺส อนนฺตร\-ปจฺจเยน ปจฺจโย. ตีณิ.}

\prose{มคฺคาธิปติ เหตุ ธมฺโม มคฺคา\-ธิปติสฺส เหตุสฺส ธมฺมสฺส อนนฺตร\-ปจฺจเยน ปจฺจโย. ตีณิ.}

\prose{มคฺคารมฺมโณ เหตุ จ มคฺคาธิปติ เหตุ จ ธมฺมา มคฺคา\-รมฺมณสฺส เหตุสฺส ธมฺมสฺส อนนฺตร\-ปจฺจเยน ปจฺจโย. ตีณิ. (สํขิตฺตํ.)}

\proseitem{196}{เหตุยา สตฺตรส, อารมฺมเณ นว, อธิปติยา เอกวีส, อนนฺตเร นว, สมนนฺตเร นว, สหชาเต สตฺตรส, อญฺญมญฺเญ สตฺตรส, นิสฺสเย สตฺตรส, อุปนิสฺสเย เอกวีส ฯเปฯ อวิคเต สตฺตรส. (สํขิตฺตํ.)}

\prose{นเหตุยา เอกวีส, นอารมฺมเณ สตฺตรส. (สํขิตฺตํ.)}

\prose{เหตุปจฺจยา นอารมฺมเณ สตฺตรส. (สํขิตฺตํ.)}

\prose{นเหตุปจฺจยา อารมฺมเณ นว. (สํขิตฺตํ.)}

\prose{(ยถา กุสลตฺติเก ปญฺหาวารํ, เอวํ วิตฺถาเรตพฺพํ.)}

\csromanmarkh{16. มคฺคา\-รมฺมณ\-ตฺติก}\csromanmarkhii{1. เหตุทุก นเหตุปท 1. ปฏิจฺจวาร}\csromanheadernomark{2}{16. มคฺคา\-รมฺมณ\-ตฺติก 1. เหตุทุก นเหตุปท 1. ปฏิจฺจวาร}
\prose{\csromanfolio{537}ปจฺจยจตุกฺกเหตุ มคฺคารมฺมณํ นเหตุํ ธมฺมํ ปฏิจฺจ มคฺคารมฺมโณ นเหตุ ธมฺโม อุปฺปชฺชติ เหตุปจฺจยา.  มคฺคารมฺมณํ นเหตุํ ธมฺมํ ปฏิจฺจ มคฺคาธิปติ นเหตุ ธมฺโม อุปฺปชฺชติ เหตุปจฺจยา.  มคฺคารมฺมณํ นเหตุํ ธมฺมํ ปฏิจฺจ มคฺคารมฺมโณ นเหตุ จ มคฺคาธิปติ นเหตุ จ ธมฺมา อุปฺปชฺชนฺติ เหตุปจฺจยา. (3) มคฺคเหตุกํ นเหตุํ ธมฺมํ ปฏิจฺจ มคฺคเหตุโก นเหตุ ธมฺโม อุปฺปชฺชติ เหตุปจฺจยา.  มคฺคเหตุกํ นเหตุํ ธมฺมํ ปฏิจฺจ มคฺคาธิปติ นเหตุ ธมฺโม อุปฺปชฺชติ เหตุปจฺจยา.  มคฺคเหตุกํ นเหตุํ ธมฺมํ ปฏิจฺจ มคฺคเหตุโก นเหตุ จ มคฺคาธิปติ นเหตุ จ ธมฺมา อุปฺปชฺชนฺติ เหตุปจฺจยา. (3)}

\proseitem{197}{มคฺคาธิปติํ นเหตุํ ธมฺมํ ปฏิจฺจ มคฺคาธิปติ นเหตุ ธมฺโม อุปฺปชฺชติ เหตุปจฺจยา.  มคฺคาธิปติํ นเหตุํ ธมฺมํ ปฏิจฺจ มคฺคารมฺมโณ นเหตุ ธมฺโม อุปฺปชฺชติ เหตุปจฺจยา.  มคฺคาธิปติํ นเหตุํ ธมฺมํ ปฏิจฺจ มคฺคเหตุโก นเหตุ ธมฺโม อุปฺปชฺชติ เหตุปจฺจยา.  มคฺคาธิปติํ นเหตุํ ธมฺมํ ปฏิจฺจ มคฺคารมฺมโณ นเหตุ จ มคฺคาธิปติ นเหตุ จ ธมฺมา อุปฺปชฺชนฺติ เหตุปจฺจยา.  มคฺคาธิปติํ นเหตุํ ธมฺมํ ปฏิจฺจ มคฺคเหตุโก นเหตุ จ มคฺคาธิปติ นเหตุ จ ธมฺมา อุปฺปชฺชนฺติ เหตุปจฺจยา. (5) มคฺคารมฺมณํ นเหตุญฺจ มคฺคาธิปติํ นเหตุญฺจ ธมฺมํ ปฏิจฺจ มคฺคารมฺมโณ นเหตุ ธมฺโม อุปฺปชฺชติ เหตุปจฺจยา.  มคฺคารมฺมณํ นเหตุญฺจ มคฺคาธิปติํ นเหตุญฺจ ธมฺมํ ปฏิจฺจ มคฺคาธิปติ นเหตุ ธมฺโม อุปฺปชฺชติ เหตุปจฺจยา. มคฺคารมฺมณํ นเหตุญฺจ มคฺคาธิปติํ นเหตุญฺจ ธมฺมํ ปฏิจฺจ มคฺคารมฺมโณ นเหตุ จ มคฺคาธิปติ นเหตุ จ ธมฺมา อุปฺปชฺชนฺติ เหตุปจยา. (3)}

\prose{มคฺคเหตุกํ นเหตุญฺจ มคฺคาธิปติํ นเหตุญฺจ ธมฺมํ ปฏิจฺจ มคฺคเหตุโก นเหตุ ธมฺโม อุปฺปชฺชติ เหตุปจฺจยา.  มคฺคเหตุกํ นเหตุญฺจ มคฺคาธิปติํ นเหตุญฺจ ธมฺมํ ปฏิจฺจ มคฺคาธิปติ นเหตุ ธมฺโม อุปฺปชฺชติ เหตุปจฺจยา. มคฺคเหตุกํ นเหตุญฺจ มคฺคาธิปติํ นเหตุญฺจ ธมฺมํ ปฏิจฺจ มคฺคเหตุโก นเหตุ จ มคฺคาธิปติ นเหตุ จ ธมฺมา อุปฺปชฺชนฺติ เหตุปจฺจยา. (3) (สํขิตฺตํ.)}

\proseitem{198}{\csromanfolio{538}เหตุยา สตฺตรส, อารมฺมเณ สตฺตรส ฯเปฯ อวิคเต สตฺตรส. (สํขิตฺตํ.)}

\prose{ปจฺจนีย\-น\-เหตุ}

\proseitem{199}{มคฺคารมฺมณํ นเหตุํ ธมฺมํ ปฏิจฺจ มคฺคารมฺมโณ นเหตุ ธมฺโม อุปฺปชฺชติ นเหตุปจฺจยา. (1) (สํขิตฺตํ.)}

\prose{นเหตุยา เอกํ, นอธิปติยา สตฺตรส, นปุเรชาเต สตฺตรส, นปจฺฉาชาเต สตฺตรส, นอาเสวเน นว, นวิปาเก สตฺตรส, นวิปฺปยุตฺเต สตฺตรส. (สํขิตฺตํ.)}

\prose{เหตุปจฺจยา นอธิปติยา สตฺตรส. (สํขิตฺตํ.)}

\prose{นเหตุปจฺจยา อารมฺมเณ เอกํ. (สํขิตฺตํ.)}

\prose{(สหชาตวารมฺปิ ฯเปฯ สมฺปยุตฺตวารมฺปิ ปฏิจฺจ\-วาร\-สทิสํ วิตฺถาเรตพฺพํ.)}

\csromanmarkh{นเหตุปท}\csromanmarkhii{7. ปญฺหาวาร}\csromanheadernomark{2}{\textbf{นเหตุปท 7. ปญฺหาวาร}}
\prose{ปจฺจยจตุกฺกอารมฺมณาทิ มคฺคเหตุโก นเหตุ ธมฺโม มคฺคา\-รมฺมณสฺส นเหตุสฺส ธมฺมสฺส อารมฺมณ\-ปจฺจเยน ปจฺจโย.  มคฺคเหตุโก นเหตุ ธมฺโม มคฺคา\-ธิปติสฺส นเหตุสฺส ธมฺมสฺส อารมฺมณ\-ปจฺจเยน ปจฺจโย.  มคฺคเหตุโก นเหตุ ธมฺโม มคฺคา\-รมฺมณสฺส นเหตุสฺส จ มคฺคา\-ธิปติสฺส นเหตุสฺส จ ธมฺมสฺส อารมฺมณ\-ปจฺจเยน ปจฺจโย. (3)}

\proseitem{200}{มคฺคาธิปติ นเหตุ ธมฺโม มคฺคา\-ธิปติสฺส นเหตุสฺส ธมฺมสฺส อารมฺมณ\-ปจฺจเยน ปจฺจโย.  มคฺคาธิปติ นเหตุ ธมฺโม มคฺคา\-รมฺมณสฺส นเหตุสฺส ธมฺมสฺส อารมฺมณ\-ปจฺจเยน ปจฺจโย.  มคฺคาธิปติ นเหตุ ธมฺโม มคฺคา\-รมฺมณสฺส นเหตุสฺส จ มคฺคา\-ธิปติสฺส นเหตุสฺส จ ธมฺมสฺส อารมฺมณ\-ปจฺจเยน ปจฺจโย. (3)}

\prose{มคฺคเหตุโก นเหตุ จ มคฺคาธิปติ นเหตุ จ ธมฺมา มคฺคา\-รมฺมณสฺส นเหตุสฺส ธมฺมสฺส อารมฺมณ\-ปจฺจเยน ปจฺจโย. ตีณิ.}

\proseitem{201}{มคฺคารมฺมโณ นเหตุ ธมฺโม มคฺคา\-รมฺมณสฺส นเหตุสฺส ธมฺมสฺส อธิปติ\-ปจฺจเยน ปจฺจโย. ตีณิ.}

\csromanmarkh{16. มคฺคา\-รมฺมณ\-ตฺติก}\csromanmarkhii{1. เหตุทุก}\csromanheadernomark{2}{16. มคฺคา\-รมฺมณ\-ตฺติก 1. เหตุทุก}
\prose{\csromanfolio{539}มคฺคเหตุโก นเหตุ ธมฺโม มคฺคเหตุกสฺส นเหตุสฺส ธมฺมสฺส อธิปติ\-ปจฺจเยน ปจฺจโย.  มคฺคเหตุโก นเหตุ ธมฺโม มคฺคา\-รมฺมณสฺส นเหตุสฺส ธมฺมสฺส อธิปติ\-ปจฺจเยน ปจฺจโย.  มคฺคเหตุโก นเหตุ ธมฺโม มคฺคา\-ธิปติสฺส นเหตุสฺส ธมฺมสฺส อธิปติ\-ปจฺจเยน ปจฺจโย.  มคฺคเหตุโก นเหตุ ธมฺโม มคฺคา\-รมฺมณสฺส นเหตุสฺส จ มคฺคา\-ธิปติสฺส นเหตุสฺส จ ธมฺมสฺส อธิปติ\-ปจฺจเยน ปจฺจโย.  มคฺคเหตุโก นเหตุ ธมฺโม มคฺคเหตุกสฺส นเหตุสฺส จ มคฺคา\-ธิปติสฺส นเหตุสฺส จ ธมฺมสฺส อธิปติ\-ปจฺจเยน ปจฺจโย. (5)}

\prose{มคฺคาธิปติ นเหตุ ธมฺโม มคฺคา\-ธิปติสฺส นเหตุสฺส ธมฺมสฺส อธิปติ\-ปจฺจเยน ปจฺจโย. ปญฺจ.}

\prose{มคฺคารมฺมโณ นเหตุ จ มคฺคาธิปติ นเหตุ จ ธมฺมา มคฺคา\-รมฺมณสฺส นเหตุสฺส ธมฺมสฺส อธิปติ\-ปจฺจเยน ปจฺจโย. ตีณิ.}

\prose{มคฺคเหตุโก นเหตุ จ มคฺคาธิปติ นเหตุ จ ธมฺมา มคฺคเหตุกสฺส นเหตุสฺส ธมฺมสฺส อธิปติ\-ปจฺจเยน ปจฺจโย. ปญฺจ.}

\csromanheader{2}{\textbf{อนนฺตร อุปนิสฺสย}}
\prosecont{มคฺคารมฺมโณ นเหตุ ธมฺโม มคฺคา\-รมฺมณสฺส นเหตุสฺส ธมฺมสฺส อนนฺตร\-ปจฺจเยน ปจฺจโย.  มคฺคารมฺมโณ นเหตุ ธมฺโม มคฺคา\-ธิปติสฺส นเหตุสฺส ธมฺมสฺส อนนฺตร\-ปจฺจเยน ปจฺจโย.  มคฺคารมฺมโณ นเหตุ ธมฺโม มคฺคา\-รมฺมณสฺส นเหตุสฺส จ มคฺคา\-ธิปติสฺส นเหตุสฺส จ ธมฺมสฺส อนนฺตร\-ปจฺจเยน ปจฺจโย. (3)}

\proseitem{202}{มคฺคาธิปติ นเหตุ ธมฺโม มคฺคา\-ธิปติสฺส นเหตุสฺส ธมฺมสฺส อนนฺตร\-ปจฺจเยน ปจฺจโย.  มคฺคาธิปติ นเหตุ ธมฺโม มคฺคา\-รมฺมณสฺส นเหตุสฺส ธมฺมสฺส อนนฺตร\-ปจฺจเยน ปจฺจโย.  มคฺคาธิปติ นเหตุ ธมฺโม มคฺคา\-รมฺมณสฺส นเหตุสฺส จ มคฺคา\-ธิปติสฺส นเหตุสฺส จ ธมฺมสฺส อนนฺตร\-ปจฺจเยน ปจฺจโย. (3)}

\prose{มคฺคารมฺมโณ นเหตุ จ มคฺคาธิปติ นเหตุ จ ธมฺมา มคฺคา\-รมฺมณสฺส นเหตุสฺส ธมฺมสฺส อนนฺตร\-ปจฺจเยน ปจฺจโย. ตีณิ ฯเปฯ. มคฺคารมฺมโณ นเหตุ ธมฺโม มคฺคา\-รมฺมณสฺส นเหตุสฺส ธมฺมสฺส อุปนิสฺสย\-ปจฺจเยน ปจฺจโย. ตีณิ.}

\prose{\csromanfolio{540}มคฺคเหตุโก นเหตุ ธมฺโม มคฺคเนตุกสฺส นเหตุสฺส ธมฺมสฺส อุปนิสฺสย\-ปจฺจเยน ปจฺจโย. ปญฺจ.}

\prose{มคฺคาธิปติ นเหตุ ธมฺโม มคฺคา\-ธิปติสฺส นเหตุสฺส ธมฺมสฺส อุปนิสฺสย\-ปจฺจเยน ปจฺจโย. ปญฺจ.}

\prose{มคฺคารมฺมโณ นเหตุ จ มคฺคาธิปติ นเหตุ จ ธมฺมา มคฺคา\-รมฺมณสฺส นเหตุสฺส ธมฺมสฺส อุปนิสฺสย\-ปจฺจเยน ปจฺจโย. ตีณิ.}

\prose{มคฺคเหตุโก นเหตุ จ มคฺคาธิปติ นเหตุ จ ธมฺมา มคฺคเหตุกสฺส นเหตุสฺส ธมฺมสฺส อุปนิสฺสย\-ปจฺจเยน ปจฺจโย. ปญฺจ. (สํขิตฺตํ.)}

\proseitem{204}{อารมฺมเณ นว, อธิปติยา เอกวีส, อนนฺตเร นว, สมนนฺตเร นว, สหชาเต สตฺตรส, อญฺญมญฺเญ สตฺตรส, นิสฺสเย สตฺตรส, อุปนิสฺสเย เอกวีส ฯเปฯ อวิคเต สตฺตรส. (สํขิตฺตํ.)}

\prose{นเหตุยา เอกวีส, นอารมฺมเณ สตฺตรส. (สํขิตฺตํ.)}

\prose{อารมฺมณ\-ปจฺจยา นเหตุยา นว. (สํขิตฺตํ.)}

\prose{นเหตุปจฺจยา อารมฺมเณ นว. (สํขิตฺตํ.)}

\prose{(ยถา กุสลตฺติเก ปญฺหาวารํ, เอวํ วิตฺถาเรตพฺพํ.)}

\csromanmatikaanchor{mtk.124}
\csromanmatikaanchor{mtk.125}
\csromantocmarknopage{subparagraph}{17. อุปฺปนฺนตฺติก}{mtk.124}
\bookmark[dest={mtk.124},level=5]{17. อุปฺปนฺนตฺติก}
\csromantocmarkref{subparagraph}{1. เหตุทุก}{mtk.125}
\bookmark[dest={mtk.125},level=5]{1. เหตุทุก}
\csromanmarkh{17. อุปฺปนฺนตฺติก}\csromanmarkhii{1. เหตุทุก}\csromanheadernomark{6}{17. \textbf{อุปฺปนฺนตฺติก 1. เหตุทุก}}
\prose{7. ปญฺหาวาร ปจฺจยจตุกฺกเหตฺวาทิ}

\proseitem{205}{อุปฺปนฺโน เหตุ ธมฺโม อุปฺปนฺนสฺส เหตุสฺส ธมฺมสฺส เหตุปจฺจเยน ปจฺจโย. อุปฺปนฺนา เหตู สมฺปยุตฺตกานํ เหตูนํ เหตุปจฺจเยน ปจฺจโย. ปฏิสนฺธิ\-กฺขเณ ฯเปฯ. (1)}

\prose{อนุปฺปนฺโน เหตุ ธมฺโม อุปฺปนฺนสฺส เหตุสฺส ธมฺมสฺส อารมฺมณ\-ปจฺจเยน ปจฺจโย. อนุปฺปนฺนา เหตู เจโตปริย\-ญาณสฺส, อนาคตํสญาณสฺส อารมฺมณ\-ปจฺจเยน ปจฺจโย.  อุปฺปาที เหตุ ธมฺโม อุปฺปนฺนสฺส เหตุสฺส ธมฺมสฺส อารมฺมณ\-ปจฺจเยน ปจฺจโย. อุปฺปาที เหตู เจโตปริย\-ญาณสฺส อนาคตํสญาณสฺส อารมฺมณ\-ปจฺจเยน ปจฺจโย. (2)}

\csromanmarkh{17. อุปฺปนฺนตฺติก}\csromanmarkhii{1. เหตุทุก}\csromanheadernomark{2}{17. อุปฺปนฺนตฺติก 1. เหตุทุก}
\prose{\csromanfolio{541}อุปฺปนฺโน เหตุ ธมฺโม อุปฺปนฺนสฺส เหตุสฺส ธมฺมสฺส อธิปติ\-ปจฺจเยน ปจฺจโย. ตีณิ.}

\prose{สหชาต\-ปจฺจเยน ปจฺจโย. อญฺญมญฺญ\-ปจฺจเยน ปจฺจโย. นิสฺสย\-ปจฺจเยน ปจฺจโย. อุปนิสฺสย\-ปจฺจเยน ปจฺจโย. วิปาก\-ปจฺจเยน ปจฺจโย. อินฺทฺริย\-ปจฺจเยน ปจฺจโย. มคฺคปจฺจเยน ปจฺจโย. สมฺปยุตฺต\-ปจฺจเยน ปจฺจโย. อตฺถิปจฺจเยน ปจฺจโย ฯเปฯ อวิคต\-ปจฺจเยน ปจฺจโย.}

\proseitem{206}{เหตุยา เอกํ, อารมฺมเณ เทฺว, อธิปติยา ตีณิ, สหชาเต เอกํ, อญฺญมญฺเญ เอกํ, นิสฺสเย เอกํ, อุปนิสฺสเย เทฺว ฯเปฯ อวิคเต เอกํ. (สํขิตฺตํ.)}

\prose{นเหตุยา เทฺว, นอารมฺมเณ เทฺว. (สํขิตฺตํ.)}

\prose{เหตุปจฺจยา นอารมฺมเณ เอกํ. (สํขิตฺตํ.)}

\prose{นเหตุปจฺจยา อารมฺมเณ เทฺว. (สํขิตฺตํ.)}

\prose{(ยถา กุสลตฺติเก ปญฺหาวารํ, เอวํ วิตฺถาเรตพฺพํ.)}

\prose{นเหตุปทอารมฺมณาทิ}

\proseitem{207}{อุปฺปนฺโน นเหตุ ธมฺโม อุปฺปนฺนสฺส นเหตุสฺส ธมฺมสฺส อารมฺมณ\-ปจฺจเยน ปจฺจโย. (1)}

\prose{อนุปฺปนฺโน นเหตุ ธมฺโม อุปฺปนฺนสฺส นเหตุสฺส ธมฺมสฺส อารมฺมณ\-ปจฺจเยน ปจฺจโย. (1)}

\prose{อุปฺปาที นเหตุ ธมฺโม อุปฺปนฺนสฺส นเหตุสฺส ธมฺมสฺส อารมฺมณ\-ปจฺจเยน ปจฺจโย. (1)}

\proseitem{208}{อุปฺปนฺโน นเหตุ ธมฺโม อุปฺปนฺนสฺส นเหตุสฺส ธมฺมสฺส อธิปติ\-ปจฺจเยน ปจฺจโย. (1)}

\prose{อนุปฺปนฺโน นเหตุ ธมฺโม อุปฺปนฺนสฺส นเหตุสฺส ธมฺมสฺส อธิปติ\-ปจฺจเยน ปจฺจโย. (1)}

\prose{อุปฺปาที นเหตุ ธมฺโม อุปฺปนฺนสฺส นเหตุสฺส ธมฺมสฺส อธิปติ\-ปจฺจเยน ปจฺจโย. (1)}

\prose{\csromanfolio{542}อุปฺปนฺโน นเหตุ ธมฺโม อุปฺปนฺนสฺส นเหตุสฺส ธมฺมสฺส สหชาต\-ปจฺจเยน ปจฺจโย. อญฺญมญฺญ\-ปจฺจเยน ปจฺจโย. นิสฺสย\-ปจฺจเยน ปจฺจโย. อุปนิสฺสย\-ปจฺจเยน ปจฺจโย.}

\prose{อนุปฺปนฺโน นเหตุ ธมฺโม อุปฺปนฺนสฺส นเหตุสฺส ธมฺมสฺส อุปนิสฺสย\-ปจฺจเยน ปจฺจโย.}

\prose{อุปฺปาที นเหตุ ธมฺโม อุปฺปนฺนสฺส นเหตุสฺส ธมฺมสฺส อุปนิสฺสย\-ปจฺจเยน ปจฺจโย. (สํขิตฺตํ.)}

\proseitem{209}{อารมฺมเณ ตีณิ, อธิปติยา ตีณิ, สหชาเต เอกํ, อญฺญมญฺเญ เอกํ, นิสฺสเย เอกํ, อุปนิสฺสเย ตีณิ ฯเปฯ อวิคเต เอกํ. (สํขิตฺตํ.)}

\prose{นเหตุยา ตีณิ, นอารมฺมเณ ตีณิ. (สํขิตฺตํ.)}

\prose{อารมฺมณ\-ปจฺจยา นเหตุยา ตีณิ. (สํขิตฺตํ.)}

\prose{นเหตุปจฺจยา อารมฺมเณ ตีณิ. (สํขิตฺตํ.)}

\csromanmatikaanchor{mtk.126}
\csromanmatikaanchor{mtk.127}
\csromantocmarknopage{subparagraph}{18. อตีตตฺติก}{mtk.126}
\bookmark[dest={mtk.126},level=5]{18. อตีตตฺติก}
\csromantocmarkref{subparagraph}{1. เหตุทุก}{mtk.127}
\bookmark[dest={mtk.127},level=5]{1. เหตุทุก}
\csromanmarkh{18. อตีตตฺติก}\csromanmarkhii{1. เหตุทุก}\csromanheadernomark{6}{18. \textbf{อตีตตฺติก 1. เหตุทุก}}
\csromanheader{2}{7. \textbf{ปญฺหาวาร ปจฺจยจตุกฺก}}
\proseitem{210}{ปจฺจุปฺปนฺโน เหตุ ธมฺโม ปจฺจุปฺปนฺนสฺส เหตุสฺส ธมฺมสฺส เหตุปจฺจเยน ปจฺจโย. (1)}

\prose{อตีโต เหตุ ธมฺโม ปจฺจุปฺปนฺนสฺส เหตุสฺส ธมฺมสฺส อารมฺมณ\-ปจฺจเยน ปจฺจโย.  อนาคโต เหตุ ธมฺโม ปจฺจุปฺปนฺนสฺส เหตุสฺส ธมฺมสฺส อารมฺมณ\-ปจฺจเยน ปจฺจโย. (2) (สํขิตฺตํ.)}

\proseitem{211}{เหตุยา เอกํ, อารมฺมเณ เทฺว, อธิปติยา ตีณิ ฯเปฯ สหชาเต เอกํ, อญฺญมญฺเญ เอกํ, นิสฺสเย เอกํ, อุปนิสฺสเย เทฺว ฯเปฯ วิปาเก เอกํ, อินฺทฺริเย เอกํ, มคฺเค เอกํ, สมฺปยุตฺเต เอกํ, อตฺถิยา เอกํ ฯเปฯ อวิคเต เอกํ. (สํขิตฺตํ.)}

\prose{นเหตุยา เทฺว, นอารมฺมเณ เทฺว. (สํขิตฺตํ.)}

\prose{เหตุปจฺจยา นอารมฺมเณ เอกํ. (สํขิตฺตํ.)}

\prose{นเหตุปจฺจยา อารมฺมเณ เทฺว. (สํขิตฺตํ.)}

\csromanmarkh{19. อตีตา\-รมฺมณ\-ตฺติก}\csromanmarkhii{1. เหตุทุก นเหตุปทอารมฺมณ}\csromanheadernomark{2}{19. อตีตา\-รมฺมณ\-ตฺติก 1. เหตุทุก นเหตุปทอารมฺมณ}
\proseitem{212}{\csromanfolio{543}อตีโต นเหตุ ธมฺโม ปจฺจุปฺปนฺนสฺส นเหตุสฺส ธมฺมสฺส อารมฺมณ\-ปจฺจเยน ปจฺจโย. (1)}

\prose{อนาคโต นเหตุ ธมฺโม ปจฺจุปฺปนฺนสฺส นเหตุสฺส ธมฺมสฺส อารมฺมณ\-ปจฺจเยน ปจฺจโย. (1)}

\prose{ปจฺจุปฺปนฺโน นเหตุ ธมฺโม ปจฺจุปฺปนฺนสฺส นเหตุสฺส ธมฺมสฺส อารมฺมณ\-ปจฺจเยน ปจฺจโย. (1) (สํขิตฺตํ.)}

\proseitem{213}{อารมฺมเณ ตีณิ, อธิปติยา ตีณิ ฯเปฯ สหชาเต เอกํ, อญฺญมญฺเญ เอกํ, นิสฺสเย เอกํ, อุปนิสฺสเย ตีณิ ฯเปฯ อาเสวเน เอกํ, กมฺเม เทฺว, วิปาเก เอกํ ฯเปฯ อินฺทฺริเย เอกํ, มคฺเค เอกํ, สมฺปยุตฺเต เอกํ, อตฺถิยา เอกํ ฯเปฯ อวิคเต เอกํ. (สํขิตฺตํ.)}

\prose{นเหตุยา ตีณิ, นอารมฺมเณ ตีณิ. (สํขิตฺตํ.)}

\prose{อารมฺมณ\-ปจฺจยา นเหตุยา ตีณิ. (สํขิตฺตํ.)}

\prose{นเหตุปจฺจยา อารมฺมเณ ตีณิ. (สํขิตฺตํ.)}

\csromanmatikaanchor{mtk.128}
\csromanmatikaanchor{mtk.129}
\csromantocmarknopage{subparagraph}{19. อตีตารมฺมณตฺติก}{mtk.128}
\bookmark[dest={mtk.128},level=5]{19. อตีตารมฺมณตฺติก}
\csromantocmarkref{subparagraph}{1. เหตุทุก}{mtk.129}
\bookmark[dest={mtk.129},level=5]{1. เหตุทุก}
\csromanmarkh{19. อตีตา\-รมฺมณ\-ตฺติก}\csromanmarkhii{1. เหตุทุก}\csromanheadernomark{6}{19. อตีตา\-รมฺมณ\-ตฺติก 1. เหตุทุก}
\csromanheader{2}{1-7. ปฏิจฺจาทิวาร ปจฺจยจตุกฺกเหตุ}
\proseitem{214}{อตีตารมฺมณํ เหตุํ ธมฺมํ ปฏิจฺจ อตีตารมฺมโณ เหตุ ธมฺโม อุปฺปชฺชติ เหตุปจฺจยา. (1)}

\prose{อนาคตา\-รมฺมณํ เหตุํ ธมฺมํ ปฏิจฺจ อนาคตารมฺมโณ เหตุ ธมฺโม อุปฺปชฺชติ เหตุปจฺจยา. (1)}

\prose{ปจฺจุปฺปนฺนา\-รมฺมณํ เหตุํ ธมฺมํ ปฏิจฺจ ปจฺจุปฺปนฺนา\-รมฺมโณ เหตุ ธมฺโม อุปฺปชฺชติ เหตุปจฺจยา. (1) (สํขิตฺตํ.)}

\proseitem{215}{เหตุยา ตีณิ, อารมฺมเณ ตีณิ ฯเปฯ อวิคเต ตีณิ. (สํขิตฺตํ.)}

\prose{นอธิปติยา ตีณิ, นปุเรชาเต ตีณิ, นปจฺฉาชาเต ตีณิ, นอาเสวเน ตีณิ, นวิปาเก ตีณิ, นวิปฺปยุตฺเต เทฺว. (สํขิตฺตํ.)}

\prose{(สหชาตวารมฺปิ ฯเปฯ สมฺปยุตฺตวารมฺปิ ปฏิจฺจ\-วาร\-สทิสํ วิตฺถาเรตพฺพํ.)}

\csromanheader{2}{ธมฺมานุโลม ติกทุก\-ปฏฺฐาน\-ปาฬิ เหตุอารมฺมณ}
\proseitem{216}{\csromanfolio{544}อตีตารมฺมโณ เหตุ ธมฺโม อตีตา\-รมฺมณสฺส เหตุสฺส ธมฺมสฺส เหตุปจฺจเยน ปจฺจโย. (1)}

\prose{อนาคตารมฺมโณ เหตุ ธมฺโม อนาคตา\-รมฺมณสฺส เหตุสฺส ธมฺมสฺส เหตุปจฺจเยน ปจฺจโย. (1)}

\prose{ปจฺจุปฺปนฺนา\-รมฺมโณ เหตุ ธมฺโม ปจฺจุปฺปนฺนา\-รมฺมณสฺส เหตุสฺส ธมฺมสฺส เหตุปจฺจเยน ปจฺจโย. (1)}

\prose{อตีตารมฺมโณ เหตุ ธมฺโม อตีตา\-รมฺมณสฺส เหตุสฺส ธมฺมสฺส อารมฺมณ\-ปจฺจเยน ปจฺจโย.  อตีตารมฺมโณ เหตุ ธมฺโม อนาคตา\-รมฺมณสฺส เหตุสฺส ธมฺมสฺส อารมฺมณ\-ปจฺจเยน ปจฺจโย.  อตีตารมฺมโณ เหตุ ธมฺโม ปจฺจุปฺปนฺนา\-รมฺมณสฺส เหตุสฺส ธมฺมสฺส อารมฺมณ\-ปจฺจเยน ปจฺจโย. (3)}

\prose{อนาคตารมฺมโณ เหตุ ธมฺโม อนาคตา\-รมฺมณสฺส เหตุสฺส ธมฺมสฺส อารมฺมณ\-ปจฺจเยน ปจฺจโย. ตีณิ.}

\prose{ปจฺจุปฺปนฺนา\-รมฺมโณ เหตุ ธมฺโม ปจฺจุปฺปนฺนา\-รมฺมณสฺส เหตุสฺส ธมฺมสฺส อารมฺมณ\-ปจฺจเยน ปจฺจโย. ตีณิ. (สํขิตฺตํ.)}

\proseitem{217}{เหตุยา ตีณิ, อารมฺมเณ นว, อธิปติยา สตฺต, อนนฺตเร ฉ, สมนนฺตเร ฉ, สหชาเต ตีณิ, อญฺญมญฺเญ ตีณิ, นิสฺสเย ตีณิ, อุปนิสฺสเย นว ฯเปฯ อวิคเต ตีณิ. (สํขิตฺตํ.)}

\prose{นเหตุยา นว, นอารมฺมเณ นว. (สํขิตฺตํ.)}

\prose{เหตุปจฺจยา นอารมฺมเณ ตีณิ. (สํขิตฺตํ.)}

\prose{นเหตุปจฺจยา อารมฺมเณ นว. (สํขิตฺตํ.)}

\prose{(ยถา กุสลตฺติเก ปญฺหาวารํ, เอวํ วิตฺถาเรตพฺพํ.)}

\csromanheader{2}{\textbf{นเหตุปท 1-7. ปฏิจฺจาทิวาร}}
\prose{ปจฺจยจตุกฺกเหตุ}

\proseitem{218}{อตีตารมฺมณํ นเหตุํ ธมฺมํ ปฏิจฺจ อตีตารมฺมโณ นเหตุ ธมฺโม อุปฺปชฺชติ เหตุปจฺจยา. (1)}

\csromanmarkh{19. อตีตา\-รมฺมณ\-ตฺติก}\csromanmarkhii{1. เหตุทุก}\csromanheadernomark{2}{19. อตีตา\-รมฺมณ\-ตฺติก 1. เหตุทุก}
\prose{\csromanfolio{545}อนาคตา\-รมฺมณํ นเหตุํ ธมฺมํ ปฏิจฺจ อนาคตารมฺมโณ นเหตุ ธมฺโม อุปฺปชฺชติ เหตุปจฺจยา. (1)}

\prose{ปจฺจุปฺปนฺนา\-รมฺมณํ นเหตุํ ธมฺมํ ปฏิจฺจ ปจฺจุปฺปนฺนา\-รมฺมโณ นเหตุ ธมฺโม อุปฺปชฺชติ เหตุปจฺจยา. (1) (สํขิตฺตํ.)}

\proseitem{219}{เหตุยา ตีณิ, อารมฺมเณ ตีณิ, อธิปติยา ตีณิ ฯเปฯ อวิคเต ตีณิ. (สํขิตฺตํ.)}

\prose{นเหตุยา ตีณิ, นอธิปติยา ตีณิ, นปุเรชาเต ตีณิ, นปจฺฉาชาเต ตีณิ, นอาเสวเน ตีณิ, นกมฺเม ตีณิ, นวิปาเก ตีณิ, นฌาเน เอกํ, นมคฺเค ตีณิ, นวิปฺปยุตฺเต เทฺว. (สํขิตฺตํ.)}

\prose{(สหชาตวารมฺปิ ฯเปฯ สมฺปยุตฺตวารมฺปิ ปฏิจฺจ\-วาร\-สทิสํ วิตฺถาเรตพฺพํ.)}

\proseitem{220}{อตีตารมฺมโณ นเหตุ ธมฺโม อตีตา\-รมฺมณสฺส นเหตุสฺส ธมฺมสฺส อารมฺมณ\-ปจฺจเยน ปจฺจโย. ตีณิ.}

\prose{อนาคตารมฺมโณ นเหตุ ธมฺโม อนาคตา\-รมฺมณสฺส นเหตุสฺส ธมฺมสฺส อารมฺมณ\-ปจฺจเยน ปจฺจโย. ตีณิ.}

\prose{ปจฺจุปฺปนฺนา\-รมฺมโณ นเหตุ ธมฺโม ปจฺจุปฺปนฺนา\-รมฺมณสฺส นเหตุสฺส ธมฺมสฺส อารมฺมณ\-ปจฺจเยน ปจฺจโย. ตีณิ. (สํขิตฺตํ.)}

\proseitem{221}{อารมฺมเณ นว, อธิปติยา สตฺต ฯเปฯ อวิคเต ตีณิ. (สํขิตฺตํ.)}

\prose{นเหตุยา นว, นอารมฺมเณ นว. (สํขิตฺตํ.)}

\prose{อารมฺมณ\-ปจฺจยา นเหตุยา นว. (สํขิตฺตํ.)}

\prose{นเหตุปจฺจยา อารมฺมเณ นว. (สํขิตฺตํ.)}

\prose{(ยถา กุสลตฺติเก ปญฺหาวารํ, เอวํ วิตฺถาเรตพฺพํ.)}

\csromanmatikaanchor{mtk.130}
\csromanmatikaanchor{mtk.131}
\csromantocmarknopage{subparagraph}{20. อชฺฌตฺตตฺติก}{mtk.130}
\bookmark[dest={mtk.130},level=5]{20. อชฺฌตฺตตฺติก}
\csromantocmarkref{subparagraph}{1. เหตุทุก}{mtk.131}
\bookmark[dest={mtk.131},level=5]{1. เหตุทุก}
\csromanmarkh{ธมฺมานุโลม ติกทุก\-ปฏฺฐาน\-ปาฬิ}\csromanmarkhii{20. อชฺฌตฺตตฺติก 1. เหตุทุก}\csromanheadernomark{6}{ธมฺมานุโลม ติกทุก\-ปฏฺฐาน\-ปาฬิ 20. อชฺฌตฺตตฺติก 1. เหตุทุก}
\csromanheader{2}{1-7. ปฏิจฺจาทิวาร ปจฺจยจตุกฺกเหตุ}
\proseitem{222}{\csromanfolio{546}อชฺฌตฺตํ เหตุํ ธมฺมํ ปฏิจฺจ อชฺฌตฺโต เหตุ ธมฺโม อุปฺปชฺชติ เหตุปจฺจยา. (1)}

\prose{พหิทฺธา เหตุํ ธมฺมํ ปฏิจฺจ พหิทฺธา เหตุ ธมฺโม อุปฺปชฺชติ เหตุปจฺจยา. (1) (สํขิตฺตํ.)}

\proseitem{223}{เหตุยา เทฺว, อารมฺมเณ เทฺว ฯเปฯ อวิคเต เทฺว. (สํขิตฺตํ.)}

\prose{นอธิปติยา เทฺว, นปุเรชาเต เทฺว, นปจฺฉาชาเต เทฺว, นอาเสวเน เทฺว, นวิปาเก เทฺว, นวิปฺปยุตฺเต เทฺว. (สํขิตฺตํ.)}

\prose{(สหชาตวารมฺปิ ฯเปฯ สมฺปยุตฺตวารมฺปิ ปฏิจฺจ\-วาร\-สทิสํ วิตฺถาเรตพฺพํ.)}

\csromanheader{2}{\textbf{เหตุ อารมฺมณาทิ}}
\proseitem{224}{อชฺฌตฺโต เหตุ ธมฺโม อชฺฌตฺตสฺส เหตุสฺส ธมฺมสฺส เหตุปจฺจเยน ปจฺจโย. (1)}

\prose{พหิทฺธา เหตุ ธมฺโม พหิทฺธา เหตุสฺส ธมฺมสฺส เหตุปจฺจเยน ปจฺจโย.}

\prose{อชฺฌตฺโต เหตุ ธมฺโม อชฺฌตฺตสฺส เหตุสฺส ธมฺมสฺส อารมฺมณ\-ปจฺจเยน ปจฺจโย.  อชฺฌตฺโต เหตุ ธมฺโม พหิทฺธา เหตุสฺส ธมฺมสฺส อารมฺมณ\-ปจฺจเยน ปจฺจโย. (2)}

\prose{พหิทฺธา เหตุ ธมฺโม พหิทฺธา เหตุสฺส ธมฺมสฺส อารมฺมณ\-ปจฺจเยน ปจฺจโย.  พหิทฺธา เหตุ ธมฺโม อชฺฌตฺตสฺส เหตุสฺส ธมฺมสฺส อารมฺมณ\-ปจฺจเยน ปจฺจโย. (2)}

\proseitem{225}{อชฺฌตฺโต เหตุ ธมฺโม อชฺฌตฺตสฺส เหตุสฺส ธมฺมสฺส อธิปติ\-ปจฺจเยน ปจฺจโย. (1)}

\prose{พหิทฺธา เหตุ ธมฺโม พหิทฺธา เหตุสฺส ธมฺมสฺส อธิปติ\-ปจฺจเยน ปจฺจโย. (1)}

\prose{อชฺฌตฺโต เหตุ ธมฺโม อชฺฌตฺตสฺส เหตุสฺส ธมฺมสฺส อนนฺตร\-ปจฺจเยน ปจฺจโย. (1)}

\csromanmarkh{20. อชฺฌตฺตตฺติก}\csromanmarkhii{1. เหตุทุก}\csromanheadernomark{2}{20. อชฺฌตฺตตฺติก 1. เหตุทุก}
\prose{\csromanfolio{547}พหิทฺธา เหตุ ธมฺโม พหิทฺธา เหตุสฺส ธมฺมสฺส อนนฺตร\-ปจฺจเยน ปจฺจโย. (1) (สํขิตฺตํ.)}

\proseitem{226}{เหตุยา เทฺว, อารมฺมเณ จตฺตาริ, อธิปติยา เทฺว, อนนฺตเร เทฺว, สมนนฺตเร เทฺว, สหชาเต เทฺว, อญฺญมญฺเญ เทฺว, นิสฺสเย เทฺว, อุปนิสฺสเย จตฺตาริ ฯเปฯ อวิคเต เทฺว. (สํขิตฺตํ.)}

\prose{นเหตุยา จตฺตาริ, นอารมฺมเณ จตฺตาริ. (สํขิตฺตํ.)}

\prose{เหตุปจฺจยา นอารมฺมเณ เทฺว. (สํขิตฺตํ.)}

\prose{นเหตุปจฺจยา อารมฺมเณ จตฺตาริ. (สํขิตฺตํ.)}

\prose{(ยถา กุสลตฺติเก ปญฺหาวารํ, เอวํ วิตฺถาเรตพฺพํ.)}

\csromanheader{2}{\textbf{นเหตุปท 1-7. ปฏิจฺจาทิวาร}}
\prose{ปจฺจยจตุกฺกเหตุ}

\proseitem{227}{อชฺฌตฺตํ นเหตุํ ธมฺมํ ปฏิจฺจ อชฺฌตฺโต นเหตุ ธมฺโม อุปฺปชฺชติ เหตุปจฺจยา. (1)}

\prose{พหิทฺธา นเหตุํ ธมฺมํ ปฏิจฺจ พหิทฺธา นเหตุ ธมฺโม อุปฺปชฺชติ เหตุปจฺจยา. (1) (สํขิตฺตํ.)}

\proseitem{228}{เหตุยา เทฺว, อารมฺมเณ เทฺว ฯเปฯ อวิคเต เทฺว. (สํขิตฺตํ.)}

\prose{นเหตุยา เทฺว, นอธิปติยา เทฺว ฯเปฯ นวิปฺปยุตฺเต เทฺว. (สํขิตฺตํ.)}

\prose{(สหชาตวารมฺปิ ฯเปฯ สมฺปยุตฺตวารมฺปิ ปฏิจฺจ\-วาร\-สทิสํ วิตฺถาเรตพฺพํ.)}

\proseitem{229}{อชฺฌตฺโต นเหตุ ธมฺโม อชฺฌตฺตสฺส นเหตุสฺส ธมฺมสฺส อารมฺมณ\-ปจฺจเยน ปจฺจโย.  อชฺฌตฺโต นเหตุ ธมฺโม พหิทฺธา นเหตุสฺส ธมฺมสฺส อารมฺมณ\-ปจฺจเยน ปจฺจโย. (2)}

\prose{พหิทฺธา นเหตุ ธมฺโม พหิทฺธา นเหตุสฺส ธมฺมสฺส อารมฺมณ\-ปจฺจเยน ปจฺจโย.  พหิทฺธา นเหตุ ธมฺโม อชฺฌตฺตสฺส นเหตุสฺส ธมฺมสฺส อารมฺมณ\-ปจฺจเยน ปจฺจโย. (2) (สํขิตฺตํ.)}

\proseitem{230}{\csromanfolio{548}อารมฺมเณ จตฺตาริ, อธิปติยา จตฺตาริ, อนนฺตเร เทฺว ฯเปฯ อุปนิสฺสเย จตฺตาริ ฯเปฯ อวิคเต ฉ. (สํขิตฺตํ.)}

\prose{นเหตุยา ฉ, นอารมฺมเณ ฉ. (สํขิตฺตํ.)}

\prose{อารมฺมณ\-ปจฺจยา นเหตุยา จตฺตาริ. (สํขิตฺตํ.)}

\prose{นเหตุปจฺจยา อารมฺมเณ จตฺตาริ. (สํขิตฺตํ.)}

\prose{(ยถา กุสลตฺติเก ปญฺหาวารํ, เอวํ วิตฺถาเรตพฺพํ.)}

\csromanmatikaanchor{mtk.132}
\csromanmatikaanchor{mtk.133}
\csromantocmarknopage{subparagraph}{21. อชฺฌตฺตารมฺมณตฺติก}{mtk.132}
\bookmark[dest={mtk.132},level=5]{21. อชฺฌตฺตารมฺมณตฺติก}
\csromantocmarkref{subparagraph}{1. เหตุทุก}{mtk.133}
\bookmark[dest={mtk.133},level=5]{1. เหตุทุก}
\csromanmarkh{21. อชฺฌตฺตา\-รมฺมณ\-ตฺติก}\csromanmarkhii{1. เหตุทุก}\csromanheadernomark{6}{21. อชฺฌตฺตา\-รมฺมณ\-ตฺติก 1. เหตุทุก}
\csromanheader{2}{1-7. ปฏิจฺจาทิวาร ปจฺจยจตุกฺกเหตุ}
\proseitem{231}{อชฺฌตฺตา\-รมฺมณํ เหตุํ ธมฺมํ ปฏิจฺจ อชฺฌตฺตา\-รมฺมโณ เหตุ ธมฺโม อุปฺปชฺชติ เหตุปจฺจยา. (1)}

\prose{พหิทฺธา\-รมฺมณํ เหตุํ ธมฺมํ ปฏิจฺจ พหิทฺธา\-รมฺมโณ เหตุ ธมฺโม อุปฺปชฺชติ เหตุปจฺจยา. (1) (สํขิตฺตํ.)}

\proseitem{232}{เหตุยา เทฺว, อารมฺมเณ เทฺว ฯเปฯ อวิคเต เทฺว. (สํขิตฺตํ.)}

\prose{นอธิปติยา เทฺว, นปุเรชาเต เทฺว, นปจฺฉาชาเต เทฺว, นอาเสวเน เทฺว, นวิปาเก เทฺว, นวิปฺปยุตฺเต เทฺว. (สํขิตฺตํ.)}

\prose{(สหชาตวารมฺปิ ฯเปฯ สมฺปยุตฺตวารมฺปิ ปฏิจฺจ\-วาร\-สทิสํ วิตฺถาเรตพฺพํ.)}

\prose{เหตุอารมฺมณ}

\proseitem{233}{อชฺฌตฺตา\-รมฺมโณ เหตุ ธมฺโม อชฺฌตฺตา\-รมฺมณสฺส เหตุสฺส ธมฺมสฺส เหตุปจฺจเยน ปจฺจโย. (1)}

\prose{พหิทฺธา\-รมฺมโณ เหตุ ธมฺโม พหิทฺธา\-รมฺมณสฺส เหตุสฺส ธมฺมสฺส เหตุปจฺจเยน ปจฺจโย. (1)}

\csromanmarkh{21. อชฺฌตฺตา\-รมฺมณ\-ตฺติก}\csromanmarkhii{1. เหตุทุก}\csromanheadernomark{2}{21. อชฺฌตฺตา\-รมฺมณ\-ตฺติก 1. เหตุทุก}
\prose{\csromanfolio{549}อชฺฌตฺตา\-รมฺมโณ เหตุ ธมฺโม อชฺฌตฺตา\-รมฺมณสฺส เหตุสฺส ธมฺมสฺส อารมฺมณ\-ปจฺจเยน ปจฺจโย.  อชฺฌตฺตา\-รมฺมโณ เหตุ ธมฺโม พหิทฺธา\-รมฺมณสฺส เหตุสฺส ธมฺมสฺส อารมฺมณ\-ปจฺจเยน ปจฺจโย. (2)}

\prose{พหิทฺธา\-รมฺมโณ เหตุ ธมฺโม พหิทฺธา\-รมฺมณสฺส เหตุสฺส ธมฺมสฺส อารมฺมณ\-ปจฺจเยน ปจฺจโย.  พหิทฺธา\-รมฺมโณ เหตุ ธมฺโม อชฺฌตฺตา\-รมฺมณสฺส เหตุสฺส ธมฺมสฺส อารมฺมณ\-ปจฺจเยน ปจฺจโย. (2) (สํขิตฺตํ.)}

\proseitem{234}{เหตุยา เทฺว, อารมฺมเณ จตฺตาริ, อธิปติยา ตีณิ, อนนฺตเร จตฺตาริ, สมนนฺตเร จตฺตาริ, อุปนิสฺสเย จตฺตาริ ฯเปฯ อวิคเต เทฺว. (สํขิตฺตํ.)}

\prose{นเหตุยา จตฺตาริ, นอารมฺมเณ จตฺตาริ. (สํขิตฺตํ.)}

\prose{เหตุปจฺจยา นอารมฺมเณ เทฺว. (สํขิตฺตํ.)}

\prose{นเหตุปจฺจยา อารมฺมเณ จตฺตาริ. (สํขิตฺตํ.)}

\prose{(ยถา กุสลตฺติเก ปญฺหาวารํ, เอวํ วิตฺถาเรตพฺพํ.)}

\csromanheader{2}{\textbf{นเหตุปท 1-7. ปฏิจฺจาทิวาร}}
\prose{ปจฺจยจตุกฺกเหตุ}

\proseitem{235}{อชฺฌตฺตา\-รมฺมณํ นเหตุํ ธมฺมํ ปฏิจฺจ อชฺฌตฺตา\-รมฺมโณ นเหตุ ธมฺโม อุปฺปชฺชติ เหตุปจฺจยา. (1)}

\prose{พหิทฺธา\-รมฺมณํ นเหตุํ ธมฺมํ ปฏิจฺจ พหิทฺธา\-รมฺมโณ นเหตุ ธมฺโม อุปฺปชฺชติ เหตุปจฺจยา. (1) (สํขิตฺตํ.)}

\prose{เหตุยา เทฺว, อารมฺมเณ เทฺว ฯเปฯ อวิคเต เทฺว. (สํขิตฺตํ.)}

\prose{(สหชาตวารมฺปิ ฯเปฯ สมฺปยุตฺตวารมฺปิ ปฏิจฺจ\-วาร\-สทิสํ วิตฺถาเรตพฺพํ.)}

\proseitem{236}{อชฺฌตฺตา\-รมฺมโณ นเหตุ ธมฺโม อชฺฌตฺตา\-รมฺมณสฺส นเหตุสฺส ธมฺมสฺส อารมฺมณ\-ปจฺจเยน ปจฺจโย. (สํขิตฺตํ.)}

\prose{\csromanfolio{550}อารมฺมเณ จตฺตาริ, อธิปติยา ตีณิ, อนนฺตเร จตฺตาริ, อุปนิสฺสเย จตฺตาริ ฯเปฯ อวิคเต เทฺว. (สํขิตฺตํ.)}

\prose{นเหตุยา จตฺตาริ, นอารมฺมเณ จตฺตาริ. (สํขิตฺตํ.)}

\prose{อารมฺมณ\-ปจฺจยา นเหตุยา จตฺตาริ. (สํขิตฺตํ.)}

\prose{นเหตุปจฺจยา อารมฺมเณ จตฺตาริ. (สํขิตฺตํ.)}

\prose{(ยถา กุสลตฺติเก ปญฺหาวารํ, เอวํ วิตฺถาเรตพฺพํ.)}

\csromanmatikaanchor{mtk.134}
\csromanmatikaanchor{mtk.135}
\csromantocmarknopage{subparagraph}{22. สนิทสฺสนสปฺปฏิฆตฺติก}{mtk.134}
\bookmark[dest={mtk.134},level=5]{22. สนิทสฺสนสปฺปฏิฆตฺติก}
\csromantocmarkref{subparagraph}{1. เหตุทุก}{mtk.135}
\bookmark[dest={mtk.135},level=5]{1. เหตุทุก}
\csromanmarkh{22. สนิทสฺสนสปฺปฏิฆ\-ตฺติก}\csromanmarkhii{1. เหตุทุก}\csromanheadernomark{6}{22. สนิทสฺสนสปฺปฏิฆ\-ตฺติก 1. เหตุทุก}
\csromanheader{2}{1-7. ปฏิจฺจาทิวาร ปจฺจยจตุกฺก}
\proseitem{237}{อนิทสฺสน\-อปฺปฏิฆํ เหตุํ ธมฺมํ ปฏิจฺจ อนิทสฺสน\-อปฺปฏิโฆ เหตุ ธมฺโม อุปฺปชฺชติ เหตุปจฺจยา. (1) (สํขิตฺตํ.)}

\prose{เหตุยา เอกํ, อารมฺมเณ เอกํ ฯเปฯ อวิคเต เอกํ. (สํขิตฺตํ.)}

\prose{(สหชาตวารมฺปิ ฯเปฯ สมฺปยุตฺตวารมฺปิ ปฏิจฺจ\-วาร\-สทิสํ.)}

\proseitem{238}{อนิทสฺสน\-อปฺปฏิโฆ เหตุ ธมฺโม อนิทสฺสน\-อปฺปฏิฆสฺส เหตุสฺส ธมฺมสฺส เหตุปจฺจเยน ปจฺจโย. (1)}

\prose{อนิทสฺสน\-อปฺปฏิโฆ เหตุ ธมฺโม อนิทสฺสน\-อปฺปฏิฆสฺส เหตุสฺส ธมฺมสฺส อารมฺมณ\-ปจฺจเยน ปจฺจโย. (1) (สํขิตฺตํ.)}

\csromanhangingbegin
\hangingheaditem{238}{เหตุยา เอกํ, อารมฺมเณ เอกํ ฯเปฯ อวิคเต เอกํ. (สพฺพตฺถ เอกํ.)}
\hangingbody{(ยถา กุสลตฺติเก ปญฺหาวารํ, เอวํ วิตฺถาเรตพฺพํ.)}
\csromanhangingend

\csromanmarkh{นเหตุปท}\csromanmarkhii{1. ปฏิจฺจวาร}\csromanheadernomark{2}{\textbf{นเหตุปท 1. ปฏิจฺจวาร}}
\csromanheader{2}{ปจฺจยจตุกฺก เหตุอารมฺมณ}
\prosecont{อนิทสฺสน\-สปฺปฏิฆํ นเหตุํ ธมฺมํ ปฏิจฺจ อนิทสฺสน\-สปฺปฏิโฆ นเหตุ ธมฺโม อุปฺปชฺชติ เหตุปจฺจยา.  อนิทสฺสน\-สปฺปฏิฆํ นเหตุํ ธมฺมํ ปฏิจฺจ สนิทสฺสน\-สปฺปฏิโฆ นเหตุ ธมฺโม อุปฺปชฺชติ เหตุปจฺจยา.  อนิทสฺสน\-สปฺปฏิฆํ นเหตุํ ธมฺมํ ปฏิจฺจ อนิทสฺสน\-อปฺปฏิโฆ นเหตุ ธมฺโม}

\vspace{\tipitakaparskip}
\csromancenter{สนิทสฺสนสปฺปฏิฆ\-ตฺติก 1. เหตุทุก อุปฺปชฺชติ เหตุปจฺจยา.  อนิทสฺสน\-สปฺปฏิฆํ นเหตุํ ธมฺมํ ปฏิจฺจ สนิทสฺสน\-สปฺปฏิโฆ นเหตุ จ อนิทสฺสน\-อปฺปฏิโฆ นเหตุ จ ธมฺมา อุปฺปชฺชนฺติ เหตุปจฺจยา.  อนิทสฺสน\-สปฺปฏิฆํ นเหตุํ ธมฺมํ ปฏิจฺจ อนิทสฺสน\-สปฺปฏิโฆ นเหตุ จ อนิทสฺสน\-อปฺปฏิโฆ นเหตุ จ ธมฺมา อุปฺปชฺชนฺติ เหตุปจฺจยา.  อนิทสฺสน\-สปฺปฏิฆํ นเหตุํ ธมฺมํ ปฏิจฺจ สนิทสฺสน\-สปฺปฏิโฆ นเหตุ จ อนิทสฺสน\-สปฺปฏิโฆ นเหตุ จ ธมฺมา อุปฺปชฺชนฺติ เหตุปจฺจยา.  อนิทสฺสน\-สปฺปฏิฆํ นเหตุํ ธมฺมํ ปฏิจฺจ สนิทสฺสน\-สปฺปฏิโฆ นเหตุ จ อนิทสฺสน\-สปฺปฏิโฆ นเหตุ จ อนิทสฺสน\-อปฺปฏิโฆ นเหตุ จ ธมฺมา อุปฺปชฺชนฺติ เหตุปจฺจยา. สตฺต.}
\proseitem{240}{\csromanfolio{551}อนิทสฺสน\-อปฺปฏิฆํ นเหตุํ ธมฺมํ ปฏิจฺจ อนิทสฺสน\-อปฺปฏิโฆ นเหตุ ธมฺโม อุปฺปชฺชติ เหตุปจฺจยา.  อนิทสฺสน\-อปฺปฏิฆํ นเหตุํ ธมฺมํ ปฏิจฺจ สนิทสฺสน\-สปฺปฏิโฆ นเหตุ ธมฺโม อุปฺปชฺชติ เหตุปจฺจยา.  อนิทสฺสน\-อปฺปฏิฆํ นเหตุํ ธมฺมํ ปฏิจฺจ อนิทสฺสน\-สปฺปฏิโฆ นเหตุ ธมฺโม อุปฺปชฺชติ เหตุปจฺจยา.  อนิทสฺสน\-อปฺปฏิฆํ นเหตุํ ธมฺมํ ปฏิจฺจ สนิทสฺสน\-สปฺปฏิโฆ นเหตุ จ อนิทสฺสน\-อปฺปฏิโฆ นเหตุ จ ธมฺมา อุปฺปชฺชนฺติ เหตุปจฺจยา.  อนิทสฺสน\-อปฺปฏิฆํ นเหตุํ ธมฺมํ ปฏิจฺจ อนิทสฺสน\-สปฺปฏิโฆ นเหตุ จ อนิทสฺสน\-อปฺปฏิโฆ นเหตุ จ ธมฺมา อุปฺปชฺชนฺติ เหตุปจฺจยา.  อนิทสฺสน\-อปฺปฏิฆํ นเหตุํ ธมฺมํ ปฏิจฺจ สนิทสฺสน\-สปฺปฏิโฆ นเหตุ จ อนิทสฺสน\-สปฺปฏิโฆ นเหตุ จ ธมฺมา อุปฺปชฺชนฺติ เหตุปจฺจยา.  อนิทสฺสน\-อปฺปฏิฆํ นเหตุํ ธมฺมํ ปฏิจฺจ สนิทสฺสน\-สปฺปฏิโฆ นเหตุ จ อนิทสฺสน\-สปฺปฏิโฆ นเหตุ จ อนิทสฺสน\-อปฺปฏิโฆ นเหตุ จ ธมฺมา อุปฺปชฺชนฺติ เหตุปจฺจยา. สตฺต.}

\proseitem{241}{อนิทสฺสน\-สปฺปฏิฆํ นเหตุญฺจ อนิทสฺสน\-อปฺปฏิฆํ นเหตุญฺจ ธมฺมํ ปฏิจฺจ สนิทสฺสน\-สปฺปฏิโฆ นเหตุ ธมฺโม อุปฺปชฺชติ เหตุปจฺจยา.  อนิทสฺสน\-สปฺปฏิฆํ นเหตุญฺจ อนิทสฺสน\-อปฺปฏิฆํ นเหตุญฺจ ธมฺมํ ปฏิจฺจ อนิทสฺสน\-สปฺปฏิโฆ นเหตุ ธมฺโม อุปฺปชฺชติ เหตุปจฺจยา.  อนิทสฺสน\-สปฺปฏิฆํ นเหตุญฺจ อนิทสฺสน\-อปฺปฏิฆํ นเหตุญฺจ ธมฺมํ ปฏิจฺจ อนิทสฺสน\-อปฺปฏิโฆ นเหตุ ธมฺโม อุปฺปชฺชติ เหตุปจฺจยา.  อนิทสฺสน\-สปฺปฏิฆํ นเหตุญฺจ อนิทสฺสน\-อปฺปฏิฆํ นเหตุญฺจ ธมฺมํ ปฏิจฺจ สนิทสฺสน\-สปฺปฏิโฆ นเหตุ จ อนิทสฺสน\-อปฺปฏิโฆ นเหตุ จ ธมฺมา อุปฺปชฺชนฺติ เหตุปจฺจยา.  อนิทสฺสน\-สปฺปฏิฆํ นเหตุญฺจ อนิทสฺสน\-อปฺปฏิฆํ นเหตุญฺจ ธมฺมํ ปฏิจฺจ \csromanfolio{552}อนิทสฺสน\-สปฺปฏิโฆ นเหตุ จ อนิทสฺสน\-อปฺปฏิโฆ นเหตุ จ ธมฺมา อุปฺปชฺชนฺติ เหตุปจฺจยา.  อนิทสฺสน\-สปฺปฏิฆํ นเหตุญฺจ อนิทสฺสน\-อปฺปฏิฆํ นเหตุญฺจ ธมฺมํ ปฏิจฺจ สนิทสฺสน\-สปฺปฏิโฆ นเหตุ จ อนิทสฺสน\-สปฺปฏิโฆ นเหตุ จ ธมฺมา อุปฺปชฺชนฺติ เหตุปจฺจยา.  อนิทสฺสน\-สปฺปฏิฆํ นเหตุญฺจ อนิทสฺสน\-อปฺปฏิฆํ นเหตุญฺจ ธมฺมํ ปฏิจฺจ สนิทสฺสน\-สปฺปฏิโฆ นเหตุ จ อนิทสฺสน\-สปฺปฏิโฆ นเหตุ จ อนิทสฺสน\-อปฺปฏิโฆ นเหตุ จ ธมฺมา อุปฺปชฺชนฺติ เหตุปจฺจยา. สตฺต.}

\prose{อนิทสฺสน\-อปฺปฏิฆํ นเหตุํ ธมฺมํ ปฏิจฺจ อนิทสฺสน\-อปฺปฏิโฆ นเหตุ ธมฺโม อุปฺปชฺชติ อารมฺมณ\-ปจฺจยา. (1) (สํขิตฺตํ.)}

\proseitem{242}{เหตุยา เอกวีส, อารมฺมเณ เอกํ, อธิปติยา เอกวีส ฯเปฯ อวิคเต เอกวีส. (สํขิตฺตํ.)}

\prose{นเหตุยา เอกวีส, นอารมฺมเณ เอกวีส ฯเปฯ โนวิคเต เอกวีส.}

\prose{เหตุปจฺจยา นอารมฺมเณ เอกวีส. (สํขิตฺตํ.)}

\prose{นเหตุปจฺจยา อารมฺมเณ เอกํ. (สํขิตฺตํ.)}

\prose{(สหชาตวารมฺปิ ฯเปฯ สมฺปยุตฺตวารมฺปิ ปฏิจฺจ\-วาร\-สทิสํ.)}

\prose{7. ปญฺหาวาร ปจฺจยจตุกฺกเหตุ}

\proseitem{243}{สนิทสฺสน\-สปฺปฏิโฆ นเหตุ ธมฺโม อนิทสฺสน\-อปฺปฏิฆสฺส นเหตุสฺส ธมฺมสฺส อารมฺมณ\-ปจฺจเยน ปจฺจโย. (สํขิตฺตํ.)}

\prose{อารมฺมเณ ตีณิ, อธิปติยา นว, อนนฺตเร เอกํ, สมนนฺตเร เอกํ, สหชาเต เอกวีส, อุปนิสฺสเย ตีณิ ฯเปฯ อวิคเต ปญฺจวีส. (สํขิตฺตํ.)}

\prose{นเหตุยา ปญฺจวีส, นอารมฺมเณ เอกวีส. (สํขิตฺตํ.)}

\prose{อารมฺมณ\-ปจฺจยา นเหตุยา ตีณิ. (สํขิตฺตํ.)}

\prosewithcloser{นเหตุปจฺจยา อารมฺมเณ ตีณิ. (สํขิตฺตํ.) (ยถา กุสลตฺติเก ปญฺหาวารํ, เอวํ วิตฺถาเรตพฺพํ.)}{สนิทสฺสนสปฺปฏิฆ\-ตฺติก\-เหตุทุกํ นิฏฺฐิตํ.}
\csromanmatikaanchor{mtk.136}
\csromanmatikaanchor{mtk.137}
\csromantocmarknopage{subparagraph}{1. กุสลตฺติก}{mtk.136}
\bookmark[dest={mtk.136},level=5]{1. กุสลตฺติก}
\csromantocmarkref{subparagraph}{2. สเหตุกทุก}{mtk.137}
\bookmark[dest={mtk.137},level=5]{2. สเหตุกทุก}
\csromanmarkh{1. กุสลตฺติก}\csromanmarkhii{2. สเหตุกทุก 1. กุสลตฺติก 2. สเหตุกทุก}\csromanheadernomark{6}{1. กุสลตฺติก 2. สเหตุกทุก \textbf{1. กุสลตฺติก 2. สเหตุกทุก}}
\csromanheader{2}{1-7. ปฏิจฺจาทิวาร ปจฺจยจตุกฺกเหตุ}
\proseitem{244}{\csromanfolio{553}กุสลํ สเหตุกํ ธมฺมํ ปฏิจฺจ กุสโล สเหตุโก ธมฺโม อุปฺปชฺชติ เหตุปจฺจยา. (1)}

\prose{อกุสลํ สเหตุกํ ธมฺมํ ปฏิจฺจ อกุสโล สเหตุโก ธมฺโม อุปฺปชฺชติ เหตุปจฺจยา. (1)}

\prose{อพฺยากตํ สเหตุกํ ธมฺมํ ปฏิจฺจ อพฺยากโต สเหตุโก ธมฺโม อุปฺปชฺชติ เหตุปจฺจยา. (1) (สํขิตฺตํ.)}

\prose{245. เหตุยา ตีณิ ฯเปฯ อวิคเต ตีณิ. (สํขิตฺตํ.)}

\prose{นอธิปติยา ตีณิ, นปุเรชาเต ตีณิ, นปจฺฉาชาเต ตีณิ, นอาเสวเน ตีณิ, นกมฺเม ตีณิ, นวิปาเก ตีณิ, นวิปฺปยุตฺเต ตีณิ.}

\prose{เหตุปจฺจยา นอธิปติยา ตีณิ. (สํขิตฺตํ.)}

\prose{นอธิปติ\-ปจฺจยา เหตุยา ตีณิ. (สํขิตฺตํ.) (สหชาตวารมฺปิ ฯเปฯ สมฺปยุตฺตวารมฺปิ ปฏิจฺจ\-วาร\-สทิสํ.)}

\prose{เหตุอารมฺมณ}

\proseitem{246}{กุสโล สเหตุโก ธมฺโม กุสลสฺส สเหตุกสฺส ธมฺมสฺส เหตุปจฺจเยน ปจฺจโย. (1)}

\prose{อกุสโล สเหตุโก ธมฺโม อกุสลสฺส สเหตุกสฺส ธมฺมสฺส เหตุปจฺจเยน ปจฺจโย. (1)}

\prose{อพฺยากโต สเหตุโก ธมฺโม อพฺยากตสฺส สเหตุกสฺส ธมฺมสฺส เหตุปจฺจเยน ปจฺจโย. (1)}

\prose{กุสโล สเหตุโก ธมฺโม กุสลสฺส สเหตุกสฺส ธมฺมสฺส อารมฺมณ\-ปจฺจเยน ปจฺจโย.  กุสโล สเหตุโก ธมฺโม อกุสลสฺส สเหตุกสฺส ธมฺมสฺส อารมฺมณ\-ปจฺจเยน ปจฺจโย.  กุสโล สเหตุโก ธมฺโม อพฺยากตสฺส สเหตุกสฺส ธมฺมสฺส อารมฺมณ\-ปจฺจเยน ปจฺจโย. (3)}

\prose{\csromanfolio{554}อกุสโล สเหตุโก ธมฺโม อกุสลสฺส สเหตุกสฺส ธมฺมสฺส อารมฺมณ\-ปจฺจเยน ปจฺจโย. ตีณิ.}

\prose{อพฺยากโต สเหตุโก ธมฺโม อพฺยากตสฺส สเหตุกสฺส ธมฺมสฺส อารมฺมณ\-ปจฺจเยน ปจฺจโย. ตีณิ.}

\prose{กุสโล สเหตุโก ธมฺโม กุสลสฺส สเหตุกสฺส ธมฺมสฺส อธิปติ\-ปจฺจเยน ปจฺจโย. (สํขิตฺตํ.)}

\proseitem{247}{เหตุยา ตีณิ, อารมฺมเณ นว, อธิปติยา สตฺต, อนนฺตเร ปญฺจ, สมนนฺตเร ปญฺจ, สหชาเต ตีณิ, อญฺญมญฺเญ ตีณิ, นิสฺสเย ตีณิ, อุปนิสฺสเย นว ฯเปฯ อวิคเต ตีณิ. (สํขิตฺตํ.)}

\prose{นเหตุยา นว, นอารมฺมเณ นว. (สํขิตฺตํ.)}

\prose{เหตุปจฺจยา นอารมฺมเณ ตีณิ. (สํขิตฺตํ.)}

\prose{นเหตุปจฺจยา อารมฺมเณ นว. (สํขิตฺตํ.)}

\prose{(ยถา กุสลตฺติเก ปญฺหาวารํ, เอวํ วิตฺถาเรตพฺพํ.)}

\csromanheader{2}{\textbf{อเหตุกปท 1-7. ปฏิจฺจาทิวาร}}
\prose{ปจฺจยจตุกฺกเหตุ อารมฺมณ}

\proseitem{248}{อกุสลํ อเหตุกํ ธมฺมํ ปฏิจฺจ อพฺยากโต อเหตุโก ธมฺโม อุปฺปชฺชติ เหตุปจฺจยา. (1)}

\prose{อพฺยากตํ อเหตุกํ ธมฺมํ ปฏิจฺจ อพฺยากโต อเหตุโก ธมฺโม อุปฺปชฺชติ เหตุปจฺจยา. (1)}

\prose{อกุสลํ อเหตุกญฺจ อพฺยากตํ อเหตุกญฺจ ธมฺมํ ปฏิจฺจ อพฺยากโต อเหตุโก ธมฺโม อุปฺปชฺชติ เหตุปจฺจยา. (1)}

\prose{อพฺยากตํ อเหตุกํ ธมฺมํ ปฏิจฺจ อพฺยากโต อเหตุโก ธมฺโม อุปฺปชฺชติ อารมฺมณ\-ปจฺจยา. (1) (สํขิตฺตํ.)}

\proseitem{249}{เหตุยา ตีณิ, อารมฺมเณ เอกํ, อธิปติยา เอกํ ฯเปฯ อวิคเต ตีณิ. (สํขิตฺตํ.)}

\csromanmarkh{1. กุสลตฺติก}\csromanmarkhii{2. สเหตุกทุก ปจฺจนีย\-น\-เหตุ นอารมฺมณ}\csromanheadernomark{2}{1. กุสลตฺติก 2. สเหตุกทุก ปจฺจนีย\-น\-เหตุ นอารมฺมณ}
\prose{\csromanfolio{555}อพฺยากตํ อเหตุกํ ธมฺมํ ปฏิจฺจ อพฺยากโต อเหตุโก ธมฺโม อุปฺปชฺชติ นเหตุปจฺจยา. (1)}

\prose{อกุสลํ อเหตุกํ ธมฺมํ ปฏิจฺจ อพฺยากโต อเหตุโก ธมฺโม อุปฺปชฺชติ นอารมฺมณ\-ปจฺจยา. (1)}

\prose{อพฺยากตํ อเหตุกํ ธมฺมํ ปฏิจฺจ อพฺยากโต อเหตุโก ธมฺโม อุปฺปชฺชติ นอารมฺมณ\-ปจฺจยา. (1)}

\prose{อกุสลํ อเหตุกญฺจ อพฺยากตํ อเหตุกญฺจ ธมฺมํ ปฏิจฺจ อพฺยากโต อเหตุโก ธมฺโม อุปฺปชฺชติ นอารมฺมณ\-ปจฺจยา. (1) (สํขิตฺตํ.)}

\proseitem{250}{นเหตุยา เอกํ, นอารมฺมเณ ตีณิ, นอธิปติยา ตีณิ ฯเปฯ โนวิคเต ตีณิ. (สํขิตฺตํ.)}

\prose{เหตุปจฺจยา นอารมฺมเณ ตีณิ. (สํขิตฺตํ.)}

\prose{นเหตุปจฺจยา อารมฺมเณ เอกํ. (สํขิตฺตํ.) (สหชาตวารมฺปิ ฯเปฯ สมฺปยุตฺตวารมฺปิ ปฏิจฺจ\-วาร\-สทิสํ.)}

\csromanheader{2}{\textbf{เหตุ อารมฺมณาทิ}}
\proseitem{251}{อกุสโล อเหตุโก ธมฺโม อพฺยากตสฺส อเหตุกสฺส ธมฺมสฺส เหตุปจฺจเยน ปจฺจโย. (1)}

\prose{อกุสโล อเหตุโก ธมฺโม อกุสลสฺส อเหตุกสฺส ธมฺมสฺส อารมฺมณ\-ปจฺจเยน ปจฺจโย.  อกุสโล อเหตุโก ธมฺโม อพฺยากตสฺส อเหตุกสฺส ธมฺมสฺส อารมฺมณ\-ปจฺจเยน ปจฺจโย. (2)}

\prose{อพฺยากโต อเหตุโก ธมฺโม อพฺยากตสฺส อเหตุกสฺส ธมฺมสฺส อารมฺมณ\-ปจฺจเยน ปจฺจโย.  อพฺยากโต อเหตุโก ธมฺโม อกุสลสฺส อเหตุกสฺส ธมฺมสฺส อารมฺมณ\-ปจฺจเยน ปจฺจโย. (2)}

\prose{อกุสโล อเหตุโก ธมฺโม อกุสลสฺส อเหตุกสฺส ธมฺมสฺส อนนฺตร\-ปจฺจเยน ปจฺจโย.  อกุสโล อเหตุโก ธมฺโม อพฺยากตสฺส อเหตุกสฺส ธมฺมสฺส อนนฺตร\-ปจฺจเยน ปจฺจโย. (2)}

\prose{\csromanfolio{556}อพฺยากโต อเหตุโก ธมฺโม อพฺยากตสฺส อเหตุกสฺส ธมฺมสฺส อนนฺตร\-ปจฺจเยน ปจฺจโย.  อพฺยากโต อเหตุโก ธมฺโม อกุสลสฺส อเหตุกสฺส ธมฺมสฺส อนนฺตร\-ปจฺจเยน ปจฺจโย ฯเปฯ. (2)}

\prose{อกุสโล อเหตุโก ธมฺโม อกุสลสฺส อเหตุกสฺส ธมฺมสฺส อุปนิสฺสย\-ปจฺจเยน ปจฺจโย. (สํขิตฺตํ.)}

\proseitem{252}{เหตุยา เอกํ, อารมฺมเณ จตฺตาริ, อนนฺตเร จตฺตาริ, สมนนฺตเร จตฺตาริ, สหชาเต ตีณิ, อญฺญมญฺเญ เอกํ, นิสฺสเย จตฺตาริ, อุปนิสฺสเย จตฺตาริ. (สํขิตฺตํ.)}

\prose{นเหตุยา ปญฺจ, นอารมฺมเณ ปญฺจ, นอธิปติยา ปญฺจ. (สํขิตฺตํ.)}

\prose{เหตุปจฺจยา นอารมฺมเณ เอกํ. (สํขิตฺตํ.)}

\prose{นเหตุปจฺจยา อารมฺมเณ จตฺตาริ. (สํขิตฺตํ.)}

\prose{(ยถา กุสลตฺติเก ปญฺหาวารํ, เอวํ วิตฺถาเรตพฺพํ.)}

\csromanmatikaanchor{mtk.138}
\csromantocmarkref{subparagraph}{3. เหตุสมฺปยุตฺตทุก}{mtk.138}
\bookmark[dest={mtk.138},level=5]{3. เหตุสมฺปยุตฺตทุก}
\csromanmarkh{1. กุสลตฺติก}\csromanmarkhii{3. เหตุ\-สมฺปยุตฺตทุก}\csromanheadernomark{6}{1. กุสลตฺติก 3. เหตุ\-สมฺปยุตฺตทุก}
\csromanheader{2}{1-7. ปฏิจฺจาทิวาร ปจฺจยจตุกฺกเหตุ}
\proseitem{253}{กุสลํ เหตุ\-สมฺปยุตฺตํ ธมฺมํ ปฏิจฺจ กุสโล เหตุสมฺปยุตฺโต ธมฺโม อุปฺปชฺชติ เหตุปจฺจยา. (1)}

\prose{อกุสลํ เหตุ\-สมฺปยุตฺตํ ธมฺมํ ปฏิจฺจ อกุสโล เหตุสมฺปยุตฺโต ธมฺโม อุปฺปชฺชติ เหตุปจฺจยา. (1)}

\prose{อพฺยากตํ เหตุ\-สมฺปยุตฺตํ ธมฺมํ ปฏิจฺจ อพฺยากโต เหตุสมฺปยุตฺโต ธมฺโม อุปฺปชฺชติ เหตุปจฺจยา. (1) (สํขิตฺตํ.)}

\proseitem{254}{เหตุยา ตีณิ, อารมฺมเณ ตีณิ ฯเปฯ อวิคเต ตีณิ. (สํขิตฺตํ.)}

\prose{นอธิปติยา ตีณิ ฯเปฯ นวิปฺปยุตฺเต ตีณิ. (สํขิตฺตํ.)}

\prose{(สหชาตวารมฺปิ ฯเปฯ สมฺปยุตฺตวารมฺปิ ปฏิจฺจ\-วาร\-สทิสํ.)}

\csromanmarkh{1. กุสลตฺติก}\csromanmarkhii{3. เหตุ\-สมฺปยุตฺตทุก เหตุอารมฺมณ}\csromanheadernomark{2}{1. กุสลตฺติก 3. เหตุ\-สมฺปยุตฺตทุก เหตุอารมฺมณ}
\proseitem{255}{\csromanfolio{557}กุสโล เหตุสมฺปยุตฺโต ธมฺโม กุสลสฺส เหตุ\-สมฺปยุตฺตสฺส ธมฺมสฺส เหตุปจฺจเยน ปจฺจโย. (1)}

\prose{อกุสโล เหตุสมฺปยุตฺโต ธมฺโม อกุสลสฺส เหตุ\-สมฺปยุตฺตสฺส ธมฺมสฺส เหตุปจฺจเยน ปจฺจโย. (1)}

\prose{อพฺยากโต เหตุสมฺปยุตฺโต ธมฺโม อพฺยากตสฺส เหตุ\-สมฺปยุตฺตสฺส ธมฺมสฺส เหตุปจฺจเยน ปจฺจโย. (1)}

\prose{กุสโล เหตุสมฺปยุตฺโต ธมฺโม กุสลสฺส เหตุ\-สมฺปยุตฺตสฺส ธมฺมสฺส อารมฺมณ\-ปจฺจเยน ปจฺจโย. ตีณิ.}

\prose{อกุสโล เหตุสมฺปยุตฺโต ธมฺโม อกุสลสฺส เหตุ\-สมฺปยุตฺตสฺส ธมฺมสฺส อารมฺมณ\-ปจฺจเยน ปจฺจโย. ตีณิ.}

\prose{อพฺยากโต เหตุสมฺปยุตฺโต ธมฺโม อพฺยากตสฺส เหตุ\-สมฺปยุตฺตสฺส ธมฺมสฺส อารมฺมณ\-ปจฺจเยน ปจฺจโย. ตีณิ. (สํขิตฺตํ.)}

\proseitem{256}{เหตุยา ตีณิ, อารมฺมเณ นว, อธิปติยา สตฺต, อนนฺตเร ปญฺจ, สมนนฺตเร ปญฺจ, สหชาเต ตีณิ, อุปนิสฺสเย นว ฯเปฯ อวิคเต ตีณิ. (สํขิตฺตํ.)}

\prose{นเหตุยา นว, นอารมฺมเณ นว. (สํขิตฺตํ.)}

\prose{เหตุปจฺจยา นอารมฺมเณ ตีณิ. (สํขิตฺตํ.)}

\prose{นเหตุปจฺจยา อารมฺมเณ นว. (สํขิตฺตํ.)}

\prose{(ยถา กสลตฺติเก ปญฺหาวารํ, เอวํ วิตฺถาเรตพฺพํ.)}

\csromanheader{2}{เหตุ\-วิปฺปยุตฺตปท 1-7. ปฏิจฺจาทิวาร}
\prose{ปจฺจยจตุกฺกเหตุ อารมฺมณ}

\proseitem{257}{อกุสลํ เหตุ\-วิปฺปยุตฺตํ ธมฺมํ ปฏิจฺจ อพฺยากโต เหตุวิปฺปยุตฺโต ธมฺโม อุปฺปชฺชติ เหตุปจฺจยา. (1)}

\prose{\csromanfolio{558}อพฺยากตํ เหตุ\-วิปฺปยุตฺตํ ธมฺมํ ปฏิจฺจ อพฺยากโต เหตุวิปฺปยุตฺโต ธมฺโม อุปฺปชฺชติ เหตุปจฺจยา. (1)}

\prose{อกุสลํ เหตุ\-วิปฺปยุตฺตญฺจ อพฺยากตํ เหตุ\-วิปฺปยุตฺตญฺจ ธมฺมํ ปฏิจฺจ อพฺยากโต เหตุวิปฺปยุตฺโต ธมฺโม อุปฺปชฺชติ เหตุปจฺจยา. (1)}

\prose{อพฺยากตํ เหตุ\-วิปฺปยุตฺตํ ธมฺมํ ปฏิจฺจ อพฺยากโต เหตุวิปฺปยุตฺโต ธมฺโม อุปฺปชฺชติ อารมฺมณ\-ปจฺจยา. (1) (สํขิตฺตํ.)}

\proseitem{258}{เหตุยา ตีณิ, อารมฺมเณ เอกํ, อธิปติยา เอกํ ฯเปฯ อวิคเต ตีณิ. (สํขิตฺตํ.)}

\prose{นเหตุ นอารมฺมณ}

\proseitem{259}{อพฺยากตํ เหตุ\-วิปฺปยุตฺตํ ธมฺมํ ปฏิจฺจ อพฺยากโต เหตุวิปฺปยุตฺโต ธมฺโม อุปฺปชฺชติ นเหตุปจฺจยา. (1)}

\prose{อกุสลํ เหตุ\-วิปฺปยุตฺตํ ธมฺมํ ปฏิจฺจ อพฺยากโต เหตุวิปฺปยุตฺโต ธมฺโม อุปฺปชฺชติ นอารมฺมณ\-ปจฺจยา. (1)}

\prose{อพฺยากตํ เหตุ\-วิปฺปยุตฺตํ ธมฺมํ ปฏิจฺจ อพฺยากโต เหตุวิปฺปยุตฺโต ธมฺโม อุปฺปชฺชติ นอารมฺมณ\-ปจฺจยา. (1)}

\prose{อกุสลํ เหตุ\-วิปฺปยุตฺตญฺจ อพฺยากตํ เหตุ\-วิปฺปยุตฺตญฺจ ธมฺมํ ปฏิจฺจ อพฺยากโต เหตุวิปฺปยุตฺโต ธมฺโม อุปฺปชฺชติ นอารมฺมณ\-ปจฺจยา. (สํขิตฺตํ.)}

\proseitem{260}{นเหตุยา เอกํ, นอารมฺมเณ ตีณิ, นอธิปติยา ตีณิ ฯเปฯ โนวิคเต ตีณิ. (สํขิตฺตํ. สหชาตวารมฺปิ ฯเปฯ สมฺปยุตฺตวารมฺปิ ปฏิจฺจ\-วาร\-สทิสํ.)}

\csromanheader{2}{\textbf{เหตุ อารมฺมณ อนนฺตร}}
\proseitem{261}{อกุสโล เหตุวิปฺปยุตฺโต ธมฺโม อพฺยากตสฺส เหตุ\-วิปฺปยุตฺตสฺส ธมฺมสฺส เหตุปจฺจเยน ปจฺจโย. (1)}

\prose{อกุสโล เหตุวิปฺปยุตฺโต ธมฺโม อกุสลสฺส เหตุ\-วิปฺปยุตฺตสฺส ธมฺมสฺส อารมฺมณ\-ปจฺจเยน ปจฺจโย.  อกุสโล เหตุวิปฺปยุตฺโต ธมฺโม อพฺยากตสฺส เหตุ\-วิปฺปยุตฺตสฺส ธมฺมสฺส อารมฺมณ\-ปจฺจเยน ปจฺจโย. (2)}

\csromanmatikaanchor{mtk.139}
\csromantocmarkref{subparagraph}{4. เหตุสเหตุกทุก}{mtk.139}
\bookmark[dest={mtk.139},level=5]{4. เหตุสเหตุกทุก}
\csromanmarkh{1. กุสลตฺติก}\csromanmarkhii{4. เหตุ\-สเหตุกทุก}\csromanheadernomark{6}{1. กุสลตฺติก 4. เหตุ\-สเหตุกทุก}
\prose{\csromanfolio{559}อพฺยากโต เหตุวิปฺปยุตฺโต ธมฺโม อพฺยากตสฺส เหตุ\-วิปฺปยุตฺตสฺส ธมฺมสฺส อารมฺมณ\-ปจฺจเยน ปจฺจโย.  อพฺยากโต เหตุวิปฺปยุตฺโต ธมฺโม อกุสลสฺส เหตุ\-วิปฺปยุตฺตสฺส ธมฺมสฺส อารมฺมณ\-ปจฺจเยน ปจฺจโย. (2)}

\prose{อกุสโล เหตุวิปฺปยุตฺโต ธมฺโม อกุสลสฺส เหตุ\-วิปฺปยุตฺตสฺส ธมฺมสฺส อนนฺตร\-ปจฺจเยน ปจฺจโย.  อกุสโล เหตุวิปฺปยุตฺโต ธมฺโม อพฺยากตสฺส เหตุ\-วิปฺปยุตฺตสฺส ธมฺมสฺส อนนฺตร\-ปจฺจเยน ปจฺจโย. (2)}

\prose{อพฺยากโต เหตุวิปฺปยุตฺโต ธมฺโม อพฺยากตสฺส เหตุ\-วิปฺปยุตฺตสฺส ธมฺมสฺส อนนฺตร\-ปจฺจเยน ปจฺจโย.  อพฺยากโต เหตุวิปฺปยุตฺโต ธมฺโม อกุสลสฺส เหตุ\-วิปฺปยุตฺตสฺส ธมฺมสฺส อนนฺตร\-ปจฺจเยน ปจฺจโย. (2) (สํขิตฺตํ.)}

\proseitem{262}{เหตุยา เอกํ, อารมฺมเณ จตฺตาริ, อนนฺตเร จตฺตาริ, สมนนฺตเร จตฺตาริ, สหชาเต ตีณิ, อญฺญมญฺเญ เอกํ, นิสฺสเย จตฺตาริ, อุปนิสฺสเย จตฺตาริ ฯเปฯ อวิคเต จตฺตาริ. (สํขิตฺตํ.)}

\prose{นเหตุยา ปญฺจ, นอารมฺมเณ ปญฺจ. (สํขิตฺตํ.)}

\prose{เหตุปจฺจยา นอารมฺมเณ เอกํ. (สํขิตฺตํ.)}

\prose{นเหตุปจฺจยา อารมฺมเณ จตฺตาริ. (สํขิตฺตํ.)}

\prose{(ยถา กุสลตฺติเก ปญฺหาวารํ, เอวํ วิตฺถาเรตพฺพํ.)}

\csromanmarkh{1. กุสลตฺติก}\csromanmarkhii{4. เหตุ\-สเหตุกทุก}\csromanheadernomark{2}{1. กุสลตฺติก 4. เหตุ\-สเหตุกทุก}
\csromanheader{2}{1-7. ปฏิจฺจาทิวาร ปจฺจยจตุกฺกเหตุ}
\proseitem{263}{กุสลํ เหตุญฺเจว สเหตุกญฺจ ธมฺมํ ปฏิจฺจ กุสโล เหตุ เจว สเหตุโก จ ธมฺโม อุปฺปชฺชติ เหตุปจฺจยา. (1)}

\prose{อกุสลํ เหตุญฺเจว สเหตุญฺจ ธมฺมํ ปฏิจฺจ อกุสโล เหตุ เจว สเหตุโก จ ธมฺโม อุปฺปชฺชติ เหตุปจฺจยา. (1)}

\prose{\csromanfolio{560}อพฺยากตํ เหตุญฺเจว สเหตุกญฺจ ธมฺมํ ปฏิจฺจ อพฺยากโต เหตุ เจว สเหตุโก จ ธมฺโม อุปฺปชฺชติ เหตุปจฺจยา. (1) (สํขิตฺตํ.)}

\proseitem{264}{เหตุยา ตีณิ, อารมฺมเณ ตีณิ ฯเปฯ อวิคเต ตีณิ. (สํขิตฺตํ.)}

\prose{นอธิปติยา ตีณิ, นปุเรชาเต ตีณิ, นปจฺฉาชาเต ตีณิ, นอาเสวเน ตีณิ, นวิปาเก ตีณิ, นวิปฺปยุตฺเต ตีณิ. (สํขิตฺตํ.) (สหชาตวารมฺปิ ฯเปฯ สมฺปยุตฺตวารมฺปิ ปฏิจฺจ\-วาร\-สทิสํ.)}

\csromanheader{2}{\textbf{เหตุ อารมฺมณ อธิปติ}}
\proseitem{265}{กุสโล เหตุ เจว สเหตุโก จ ธมฺโม กุสลสฺส เหตุสฺส เจว สเหตุกสฺส จ ธมฺมสฺส เหตุปจฺจเยน ปจฺจโย. (1)}

\prose{อกุสโล เหตุ เจว สเหตุโก จ ธมฺโม อกุสลสฺส เหตุสฺส เจว สเหตุกสฺส จ ธมฺมสฺส เหตุปจฺจเยน ปจฺจโย. (1)}

\prose{อพฺยากโต เหตุ เจว สเหตุโก จ ธมฺโม อพฺยากตสฺส เหตุสฺส เจว สเหตุกสฺส จ ธมฺมสฺส เหตุปจฺจเยน ปจฺจโย. (1)}

\prose{กุสโล เหตุ เจว สเหตุโก จ ธมฺโม กุสลสฺส เหตุสฺส เจว สเหตุกสฺส จ ธมฺมสฺส อารมฺมณ\-ปจฺจเยน ปจฺจโย.  กุสโล เหตุ เจว สเหตุโก จ ธมฺโม อกุสลสฺส เหตุสฺส เจว สเหตุกสฺส จ ธมฺมสฺส อารมฺมณ\-ปจฺจเยน ปจฺจโย.  กุสโล เหตุ เจว สเหตุโก จ ธมฺโม อพฺยากตสฺส เหตุสฺส เจว สเหตุกสฺส จ ธมฺมสฺส อารมฺมณ\-ปจฺจเยน ปจฺจโย. (3)}

\prose{อกุสโล เหตุ เจว สเหตุโก จ ธมฺโม อกุสลสฺส เหตุสฺส เจว สเหตุกสฺส จ ธมฺมสฺส อารมฺมณ\-ปจฺจเยน ปจฺจโย. ตีณิ.}

\prose{อพฺยากโต เหตุ เจว สเหตุโก จ ธมฺโม อพฺยากตสฺส เหตุสฺส เจว สเหตุกสฺส จ ธมฺมสฺส อารมฺมณ\-ปจฺจเยน ปจฺจโย. ตีณิ.}

\prose{กุสโล เหตุ เจว สเหตุโก จ ธมฺโม กุสลสฺส เหตุสฺส เจว สเหตุกสฺส จ ธมฺมสฺส อธิปติ\-ปจฺจเยน ปจฺจโย. ตีณิ.}

\csromanmarkh{1. กุสลตฺติก}\csromanmarkhii{4. เหตุ\-สเหตุกทุก}\csromanheadernomark{2}{1. กุสลตฺติก 4. เหตุ\-สเหตุกทุก}
\prose{\csromanfolio{561}อกุสโล เหตุ เจว สเหตุโก จ ธมฺโม อกุสลสฺส เหตุสฺส เจว สเหตุกสฺส จ ธมฺมสฺส อธิปติ\-ปจฺจเยน ปจฺจโย. (1)}

\prose{อพฺยากโต เหตุ เจว สเหตุโก จ ธมฺโม อพฺยากตสฺส เหตุสฺส เจว สเหตุกสฺส จ ธมฺมสฺส อธิปติ\-ปจฺจเยน ปจฺจโย. ตีณิ. (สํขิตฺตํ.)}

\proseitem{266}{เหตุยา ตีณิ, อารมฺมเณ นว, อธิปติยา สตฺต, อนนฺตเร ปญฺจ, สมนนฺตเร ปญฺจ, สหชาเต ตีณิ ฯเปฯ อุปนิสฺสเย นว ฯเปฯ อวิคเต ตีณิ. (สํขิตฺตํ.)}

\prose{นเหตุยา นว, นอารมฺมเณ นว. (สํขิตฺตํ.)}

\prose{เหตุปจฺจยา นอารมฺมเณ ตีณิ. (สํขิตฺตํ.)}

\prose{นเหตุปจฺจยา อารมฺมเณ นว. (สํขิตฺตํ.) (ยถา กุสลตฺติเก ปญฺหาวารํ, เอวํ วิตฺถาเรตพฺพํ.)}

\csromanheader{2}{สเหตุก\-น\-เหตุปท 1-7. ปฏิจฺจาทิวาร}
\prose{ปจฺจยจตุกฺกเหตุ}

\proseitem{267}{กุสลํ สเหตุกญฺเจว น จ เหตุํ ธมฺมํ ปฏิจฺจ กุสโล สเหตุโก เจว น จ เหตุ ธมฺโม อุปฺปชฺชติ เหตุปจฺจยา. (1)}

\prose{อกุสลํ สเหตุกญฺเจว น จ เหตุํ ธมฺมํ ปฏิจฺจ อกุสโล สเหตุโก เจว น จ เหตุ ธมฺโม อุปฺปชฺชติ เหตุปจฺจยา. (1)}

\prose{อพฺยากตํ สเหตุกญฺเจว น จ เหตุํ ธมฺมํ ปฏิจฺจ อพฺยากโต สเหตุโก เจว น จ เหตุ ธมฺโม อุปฺปชฺชติ เหตุปจฺจยา. (สํขิตฺตํ.)}

\proseitem{268}{เหตุยา ตีณิ, อารมฺมเณ ตีณิ ฯเปฯ อวิคเต ตีณิ. (สํขิตฺตํ.)}

\prose{นอธิปติยา ตีณิ ฯเปฯ นวิปฺปยุตฺเต ตีณิ. (สํขิตฺตํ.)}

\prose{(สหชาตวารมฺปิ ฯเปฯ สมฺปยุตฺตวารมฺปิ ปฏิจฺจ\-วาร\-สทิสํ.)}

\csromanheader{2}{ธมฺมานุโลม ติกทุก\-ปฏฺฐาน\-ปาฬิ อารมฺมณ}
\proseitem{269}{\csromanfolio{562}กุสโล สเหตุโก เจว น จ เหตุ ธมฺโม กุสลสฺส สเหตุกสฺส เจว น จ เหตุสฺส ธมฺมสฺส อารมฺมณ\-ปจฺจเยน ปจฺจโย. ตีณิ.}

\prose{อกุสโล สเหตุโก เจว น จ เหตุ ธมฺโม อกุสลสฺส สเหตุกสฺส เจว น จ เหตุสฺส ธมฺมสฺส อารมฺมณ\-ปจฺจเยน ปจฺจโย. ตีณิ.}

\prose{อพฺยากโต สเหตุโก เจว น จ เหตุ ธมฺโม อพฺยากตสฺส สเหตุกสฺส เจว น จ เหตุสฺส ธมฺมสฺส อารมฺมณ\-ปจฺจเยน ปจฺจโย. ตีณิ. (สํขิตฺตํ.)}

\proseitem{270}{อารมฺมเณ นว, อธิปติยา สตฺต, อนนฺตเร ปญฺจ ฯเปฯ สหชาเต ตีณิ ฯเปฯ อุปนิสฺสเย นว ฯเปฯ อวิคเต ตีณิ. (สํขิตฺตํ.)}

\prose{นเหตุยา นว, นอารมฺมเณ นว. (สํขิตฺตํ.)}

\prose{อารมฺมณ\-ปจฺจยา นเหตุยา นว. (สํขิตฺตํ.)}

\prose{นเหตุปจฺจยา อารมฺมเณ นว. (สํขิตฺตํ.)}

\prose{(ยถา กุสลตฺติเก ปญฺหาวารํ, เอวํ วิตฺถาเรตพฺพํ.)}

\csromanmatikaanchor{mtk.140}
\csromantocmarkref{subparagraph}{5. เหตุเหตุสมฺปยุตฺตทุก}{mtk.140}
\bookmark[dest={mtk.140},level=5]{5. เหตุเหตุสมฺปยุตฺตทุก}
\csromanmarkh{1. กุสลตฺติก}\csromanmarkhii{5. เหตุ\-เหตุ\-สมฺปยุตฺตทุก}\csromanheadernomark{6}{1. กุสลตฺติก 5. เหตุ\-เหตุ\-สมฺปยุตฺตทุก}
\csromanheader{2}{1-7. ปฏิจฺจาทิวาร ปจฺจยจตุกฺกเหตุ}
\proseitem{271}{กุสลํ เหตุญฺเจว เหตุ\-สมฺปยุตฺตญฺจ ธมฺมํ ปฏิจฺจ กุสโล เหตุ เจว เหตุสมฺปยุตฺโต จ ธมฺโม อุปฺปชฺชติ เหตุปจฺจยา. (1)}

\prose{อกุสลํ เหตุญฺเจว เหตุ\-สมฺปยุตฺตญฺจ ธมฺมํ ปฏิจฺจ อกุสโล เหตุ เจว เหตุสมฺปยุตฺโต จ ธมฺโม อุปฺปชฺชติ เหตุปจฺจยา. (1)}

\prose{อพฺยากตํ เหตุญฺเจว เหตุ\-สมฺปยุตฺตญฺจ ธมฺมํ ปฏิจฺจ อพฺยากโต เหตุ เจว เหตุสมฺปยุตฺโต จ ธมฺโม อุปฺปชฺชติ เหตุปจฺจยา. (สํขิตฺตํ.)}

\csromanmarkh{1. กุสลตฺติก}\csromanmarkhii{5. เหตุ\-เหตุ\-สมฺปยุตฺตทุก}\csromanheadernomark{2}{1. กุสลตฺติก 5. เหตุ\-เหตุ\-สมฺปยุตฺตทุก}
\prose{\csromanfolio{563}272. เหตุยา ตีณิ ฯเปฯ อวิคเต ตีณิ. (สํขิตฺตํ.)}

\prose{นอธิปติยา ตีณิ. (สํขิตฺตํ.)}

\prose{เหตุปจฺจยา นอธิปติยา ตีณิ. (สํขิตฺตํ.)}

\prose{นอธิปติ\-ปจฺจยา เหตุยา ตีณิ. (สํขิตฺตํ.) (สหชาตวารมฺปิ ฯเปฯ สมฺปยุตฺตวารมฺปิ ปฏิจฺจ\-วาร\-สทิสํ.)}

\csromanheader{2}{\textbf{เหตุ อารมฺมณ}}
\proseitem{273}{กุสโล เหตุ เจว เหตุสมฺปยุตฺโต จ ธมฺโม กุสลสฺส เหตุสฺส เจว เหตุ\-สมฺปยุตฺตสฺส จ ธมฺมสฺส เหตุปจฺจเยน ปจฺจโย. (1)}

\prose{อกุสโล เหตุ เจว เหตุสมฺปยุตฺโต จ ธมฺโม อกุสลสฺส เหตุสฺส เจว เหตุ\-สมฺปยุตฺตสฺส จ ธมฺมสฺส เหตุปจฺจเยน ปจฺจโย. (1)}

\prose{อพฺยากโต เหตุ เจว เหตุสมฺปยุตฺโต จ ธมฺโม อพฺยากตสฺส เหตุสฺส เจว เหตุ\-สมฺปยุตฺตสฺส จ ธมฺมสฺส เหตุปจฺจเยน ปจฺจโย. (1)}

\prose{กุสโล เหตุ เจว เหตุสมฺปยุตฺโต จ ธมฺโม กุสลสฺส เหตุสฺส เจว เหตุ\-สมฺปยุตฺตสฺส จ ธมฺมสฺส อารมฺมณ\-ปจฺจเยน ปจฺจโย. (สํขิตฺตํ.)}

\proseitem{274}{เหตุยา ตีณิ, อารมฺมเณ นว, อธิปติยา สตฺต, อนนฺตเร ปญฺจ, อุปนิสฺสเย นว ฯเปฯ อวิคเต ตีณิ. (สํขิตฺตํ.)}

\prose{นเหตุยา นว, นอารมฺมเณ นว. (สํขิตฺตํ.)}

\prose{เหตุปจฺจยา นอารมฺมเณ ตีณิ. (สํขิตฺตํ.)}

\prose{นเหตุปจฺจยา อารมฺมเณ นว. (สํขิตฺตํ.)}

\prose{(ยถา กุสลตฺติเก ปญฺหาวารํ, เอวํ วิตฺถาเรตพฺพํ.)}

\csromanheader{2}{ธมฺมานุโลม ติกทุก\-ปฏฺฐาน\-ปาฬิ เหตุ\-สมฺปยุตฺต\-น\-เหตุปท 1-7. ปฏิจฺจาทิวาร}
\prose{\csromanfolio{564}ปจฺจยจตุกฺกเหตุ}

\proseitem{275}{กุสลํ เหตุสมฺปยุตญฺเจว น จ เหตุํ ธมฺมํ ปฏิจฺจ กุสโล เหตุสมฺปยุตฺโต เจว น จ เหตุ ธมฺโม อุปฺปชฺชติ เหตุปจฺจยา. (1)}

\prose{อกุสลํ เหตุ\-สมฺปยุตฺต\-ญฺเจว น จ เหตุํ ธมฺมํ ปฏิจฺจ อกุสโล เหตุสมฺปยุตฺโต เจว น จ เหตุ ธมฺโม อุปฺปชฺชติ เหตุปจฺจยา. (1)}

\prose{อพฺยากตํ เหตุ\-สมฺปยุตฺต\-ญฺเจว น จ เหตุํ ธมฺมํ ปฏิจฺจ อพฺยากโต เหตุสมฺปยุตฺโต เจว น จ เหตุ ธมฺโม อุปฺปชฺชติ เหตุปจฺจยา. (สํขิตฺตํ.)}

\proseitem{276}{เหตุยา ตีณิ, อารมฺมเณ ตีณิ ฯเปฯ อวิคเต ตีณิ.}

\prose{นอธิปติยา ตีณิ. (สํขิตฺตํ.) (สหชาตวารมฺปิ ฯเปฯ สมฺปยุตฺตวารมฺปิ ปฏิจฺจ\-วาร\-สทิสํ.)}

\proseitem{277}{กุสโล เหตุสมฺปยุตฺโต เจว น จ เหตุ ธมฺโม กุสลสฺส เหตุ\-สมฺปยุตฺตสฺส เจว น จ เหตุสฺส ธมฺมสฺส อารมฺมณ\-ปจฺจเยน ปจฺจโย. ตีณิ.}

\prose{อกุสโล เหตุสมฺปยุตฺโต เจว น จ เหตุ ธมฺโม อกุสลสฺส เหตุ\-สมฺปยุตฺตสฺส เจว น จ เหตุสฺส ธมฺมสฺส อารมฺมณ\-ปจฺจเยน ปจฺจโย. ตีณิ.}

\prose{อพฺยากโต เหตุสมฺปยุตฺโต เจว น จ เหตุ ธมฺโม อพฺยากตสฺส เหตุ\-สมฺปยุตฺตสฺส เจว น จ เหตุสฺส ธมฺมสฺส อารมฺมณ\-ปจฺจเยน ปจฺจโย. ตีณิ. (สํขิตฺตํ.)}

\proseitem{278}{อารมฺมเณ นว, อธิปติยา สตฺต, อนนฺตเร ปญฺจ ฯเปฯ สหชาเต ตีณิ ฯเปฯ อุปนิสฺสเย นว ฯเปฯ อวิคเต ตีณิ. (สํขิตฺตํ.)}

\csromanmatikaanchor{mtk.141}
\csromantocmarkref{subparagraph}{6. นเหตุสเหตุกทุก}{mtk.141}
\bookmark[dest={mtk.141},level=5]{6. นเหตุสเหตุกทุก}
\csromanmarkh{1. กุสลตฺติก}\csromanmarkhii{6. นเหตุ\-สเหตุกทุก}\csromanheadernomark{6}{1. กุสลตฺติก 6. นเหตุ\-สเหตุกทุก}
\prose{\csromanfolio{565}นเหตุยา นว, นอารมฺมเณ นว. (สํขิตฺตํ.)}

\prose{อารมฺมณ\-ปจฺจยา นเหตุยา นว. (สํขิตฺตํ.)}

\prose{นเหตุปจฺจยา อารมฺมเณ นว. (สํขิตฺตํ.)}

\prose{(ยถา กุสลตฺติเก ปญฺหาวารํ, เอวํ วิตฺถาเรตพฺพํ.)}

\csromanmarkh{1. กุสลตฺติก}\csromanmarkhii{6. นเหตุ\-สเหตุกทุก}\csromanheadernomark{2}{1. กุสลตฺติก 6. นเหตุ\-สเหตุกทุก}
\csromanheader{2}{1-7. ปฏิจฺจาทิวาร ปจฺจยจตุกฺกเหตุ}
\proseitem{279}{กุสลํ นเหตุํ สเหตุกํ ธมฺมํ ปฏิจฺจ กุสโล นเหตุ สเหตุโก ธมฺโม อุปฺปชฺชติ เหตุปจฺจยา. (1)}

\prose{อกุสลํ นเหตุํ สเหตุกํ ธมฺมํ ปฏิจฺจ อกุสโล นเหตุ สเหตุโก ธมฺโม อุปฺปชฺชติ เหตุปจฺจยา. (1)}

\prose{อพฺยากตํ นเหตุํ สเหตุกํ ธมฺมํ ปฏิจฺจ อพฺยากโต นเหตุ สเหตุโก ธมฺโม อุปฺปชฺชติ เหตุปจฺจยา. (1) (สํขิตฺตํ.)}

\proseitem{280}{เหตุยา ตีณิ, อารมฺมเณ ตีณิ ฯเปฯ อวิคเต ตีณิ. (สํขิตฺตํ.)}

\prose{นอธิปติยา ตีณิ. (สํขิตฺตํ.) (สหชาตวารมฺปิ ฯเปฯ สมฺปยุตฺตวารมฺปิ ปฏิจฺจ\-วาร\-สทิสํ.)}

\vspace{\tipitakaparskip}
\csromancenter{\textbf{อารมฺมณ}}
\proseitem{281}{กุสโล นเหตุ สเหตุโก ธมฺโม กุสลสฺส นเหตุสฺส สเหตุกสฺส ธมฺมสฺส อารมฺมณ\-ปจฺจเยน ปจฺจโย. ตีณิ.}

\prose{อกุสโล นเหตุ สเหตุโก ธมฺโม อกุสลสฺส นเหตุสฺส สเหตุกสฺส ธมฺมสฺส อารมฺมณ\-ปจฺจเยน ปจฺจโย. ตีณิ.}

\prose{อพฺยากโต นเหตุ สเหตุโก ธมฺโม อพฺยากตสฺส นเหตุสฺส สเหตุกสฺส ธมฺมสฺส อารมฺมณ\-ปจฺจเยน ปจฺจโย. ตีณิ. (สํขิตฺตํ.)}

\csromanmatikaanchor{mtk.142}
\csromantocmarkref{subparagraph}{7-13. สปฺปจฺจยทุกาทิ}{mtk.142}
\bookmark[dest={mtk.142},level=5]{7-13. สปฺปจฺจยทุกาทิ}
\proseitem{282}{\csromanfolio{566}อารมฺมเณ นว, อธิปติยา สตฺต, อนนฺตเร ปญฺจ ฯเปฯ สหชาเต อญฺญมญฺเญ นิสฺสเย ตีณิ, อุปนิสฺสเย นว, อาเสวเน ตีณิ, กมฺเม ปญฺจ, วิปาเก เอกํ, อาหาเร ฯเปฯ สมฺปยุตฺเต ตีณิ ฯเปฯ อวิคเต ตีณิ. (สํขิตฺตํ.)}

\prose{นเหตุยา นว, นอารมฺมเณ นว. (สํขิตฺตํ.)}

\prose{อารมฺมณ\-ปจฺจยา นเหตุยา นว. (สํขิตฺตํ.)}

\prose{นเหตุปจฺจยา อารมฺมเณ นว. (สํขิตฺตํ.) (ยถา กุสลตฺติเก ปญฺหาวารํ, เอวํ วิตฺถาเรตพฺพํ.)}

\prose{นเหตุอเหตุกปทเหตุ}

\proseitem{283}{อพฺยากตํ นเหตุํ อเหตุกํ ธมฺมํ ปฏิจฺจ อพฺยากโต นเหตุ อเหตุโก ธมฺโม อุปฺปชฺชติ เหตุปจฺจยา. (สํขิตฺตํ.)}

\prose{เหตุยา เอกํ, อารมฺมเณ เอกํ ฯเปฯ อวิคเต เอกํ. (สํขิตฺตํ.)}

\prose{นเหตุยา เอกํ ฯเปฯ โนวิคเต เอกํ. (สํขิตฺตํ.) (สหชาตวารมฺปิ ฯเปฯ สมฺปยุตฺตวารมฺปิ ปฏิจฺจ\-วาร\-สทิสํ.)}

\proseitem{284}{อพฺยากโต นเหตุ อเหตุโก ธมฺโม อพฺยากตสฺส นเหตุสฺส อเหตุกสฺส ธมฺมสฺส อารมฺมณ\-ปจฺจเยน ปจฺจโย. (สํขิตฺตํ.)}

\prosewithcloser{อารมฺมเณ เอกํ, อธิปติยา เอกํ ฯเปฯ อวิคเต เอกํ. (สํขิตฺตํ.) (ยถา กุสลตฺติเก ปญฺหาวารํ, เอวํ วิตฺถาเรตพฺพํ.)}{เหตุโคจฺฉกํ นิฏฺฐิตํ.}
\csromanmarkh{1. กุสลตฺติก}\csromanmarkhii{7. สปฺปจฺจยทุก}\csromanheadernomark{6}{1. \textbf{กุสลตฺติก 7. สปฺปจฺจยทุก}}
\csromanheader{2}{1-7. ปฏิจฺจาทิวาร ปจฺจยจตุกฺก}
\proseitem{1}{กุสลํ สปฺปจฺจยํ ธมฺมํ ปฏิจฺจ กุสโล สปฺปจฺจโย ธมฺโม อุปฺปชฺชติ เหตุปจฺจยา.  กุสลํ สปฺปจฺจยํ ธมฺมํ ปฏิจฺจ อพฺยากโต สปฺปจฺจโย ธมฺโม อุปฺปชฺชติ เหตุปจฺจยา.  กุสลํ สปฺปจฺจยํ ธมฺมํ ปฏิจฺจ กุสโล}

\vspace{\tipitakaparskip}
\csromancenter{กุสลตฺติก 7. สปฺปจฺจยทุก สปฺปจฺจโย จ อพฺยากโต สปฺปจฺจโย จ ธมฺมา อุปฺปชฺชนฺติ เหตุปจฺจยา. (3)}
\prose{\csromanfolio{567}อกุสลํ สปฺปจฺจยํ ธมฺมํ ปฏิจฺจ อกุสโล สปฺปจฺจโย ธมฺโม อุปฺปชฺชติ เหตุปจฺจยา. ตีณิ.}

\prose{อพฺยากตํ สปฺปจฺจยํ ธมฺมํ ปฏิจฺจ อพฺยากโต สปฺปจฺจโย ธมฺโม อุปฺปชฺชติ เหตุปจฺจยา. (1)}

\prose{กุสลํ สปฺปจฺจยญฺจ อพฺยากตํ สปฺปจฺจยญฺจ ธมฺมํ ปฏิจฺจ อพฺยากโต สปฺปจฺจโย ธมฺโม อุปฺปชฺชติ เหตุปจฺจยา. (1)}

\prose{อกุสลํ สปฺปจฺจยญฺจ อพฺยากตํ สปฺปจฺจยญฺจ ธมฺมํ ปฏิจฺจ อพฺยากโต สปฺปจฺจโย ธมฺโม อุปฺปชฺชติ เหตุปจฺจยา. (1)}

\proseitem{2}{กุสลํ สปฺปจฺจยํ ธมฺมํ ปฏิจฺจ กุสโล สปฺปจฺจโย ธมฺโม อุปฺปชฺชติ อารมฺมณ\-ปจฺจยา. (1)}

\prose{อกุสลํ สปฺปจฺจยํ ธมฺมํ ปฏิจฺจ อกุสโล สปฺปจฺจโย ธมฺโม อุปฺปชฺชติ อารมฺมณ\-ปจฺจยา. (1)}

\prose{อพฺยากตํ สปฺปจฺจยํ ธมฺมํ ปฏิจฺจ อพฺยากโต สปฺปจฺจโย ธมฺโม อุปฺปชฺชติ อารมฺมณ\-ปจฺจยา. (1) (สํขิตฺตํ.)}

\proseitem{3}{เหตุยา นว, อารมฺมเณ ตีณิ, อธิปติยา นว ฯเปฯ อวิคเต นว. (สํขิตฺตํ.)}

\prose{ปจฺจนีย\-น\-เหตุ}

\proseitem{4}{อกุสลํ สปฺปจฺจยํ ธมฺมํ ปฏิจฺจ อกุสโล สปฺปจฺจโย ธมฺโม อุปฺปชฺชติ นเหตุปจฺจยา. (1)}

\prose{อพฺยากตํ สปฺปจฺจยํ ธมฺมํ ปฏิจฺจ อพฺยากโต สปฺปจฺจโย ธมฺโม อุปฺปชฺชติ นเหตุปจฺจยา. (1) (สํขิตฺตํ.)}

\proseitem{5}{นเหตุยา เทฺว, นอารมฺมเณ ปญฺจ, นอธิปติยา นว ฯเปฯ โนวิคเต ปญฺจ. (สํขิตฺตํ.)}

\prose{\csromanfolio{568}เหตุปจฺจยา นอารมฺมเณ ปญฺจ. (สํขิตฺตํ.)}

\prose{นเหตุปจฺจยา อารมฺมเณ เทฺว. (สํขิตฺตํ.) (สหชาตวารมฺปิ ฯเปฯ สมฺปยุตฺตวารมฺปิ ปฏิจฺจ\-วาร\-สทิสํ.)}

\csromanheader{2}{\textbf{เหตุ อารมฺมณ}}
\proseitem{6}{กุสโล สปฺปจฺจโย ธมฺโม กุสลสฺส สปฺปจฺจยสฺส ธมฺมสฺส เหตุปจฺจเยน ปจฺจโย. ตีณิ.}

\prose{อกุสโล สปฺปจฺจโย ธมฺโม อกุสลสฺส สปฺปจฺจยสฺส ธมฺมสฺส เหตุปจฺจเยน ปจฺจโย. ตีณิ.}

\prose{อพฺยากโต สปฺปจฺจโย ธมฺโม อพฺยากตสฺส สปฺปจฺจยสฺส ธมฺมสฺส เหตุปจฺจเยน ปจฺจโย. (1)}

\proseitem{7}{กุสโล สปฺปจฺจโย ธมฺโม กุสลสฺส สปฺปจฺจยสฺส ธมฺมสฺส อารมฺมณ\-ปจฺจเยน ปจฺจโย. ตีณิ.}

\prose{อกุสโล สปฺปจฺจโย ธมฺโม อกุสลสฺส สปฺปจฺจยสฺส ธมฺมสฺส อารมฺมณ\-ปจฺจเยน ปจฺจโย. ตีณิ.}

\prose{อพฺยากโต สปฺปจฺจโย ธมฺโม อพฺยากตสฺส สปฺปจฺจยสฺส ธมฺมสฺส อารมฺมณ\-ปจฺจเยน ปจฺจโย. ตีณิ. (สํขิตฺตํ.)}

\proseitem{8}{เหตุยา สตฺต, อารมฺมเณ นว, อธิปติยา ทส, อนนฺตเร สตฺต ฯเปฯ สหชาเต นว ฯเปฯ อุปนิสฺสเย นว ฯเปฯ อวิคเต เตรส. (สํขิตฺตํ.)}

\prose{นเหตุยา ปนฺนรส, นอารมฺมเณ ปนฺนรส. (สํขิตฺตํ.)}

\prose{เหตุปจฺจยา นอารมฺมเณ สตฺต. (สํขิตฺตํ.)}

\prose{นเหตุปจฺจยา อารมฺมเณ นว. (สํขิตฺตํ.)}

\prose{(ยถา กุสลตฺติเก ปญฺหาวารํ, เอวํ วิตฺถาเรตพฺพํ.)}

\vspace{\tipitakaparskip}
\csromancenter{กุสลตฺติก 9. สนิทสฺสนทุก 1. \textbf{ กุสลตฺติก 8.} สงฺขตทุก}
\csromanheader{2}{1. \textbf{ปฏิจฺจวาร ปจฺจยจตุกฺก}}
\proseitem{9}{\csromanfolio{569}กุสลํ สงฺขตํ ธมฺมํ ปฏิจฺจ กุสโล สงฺขโต ธมฺโม อุปฺปชฺชติ เหตุปจฺจยา. (สปฺปจฺจย\-ทุก\-สทิสํ.)}
\csromansectionrule

\csromanmarkh{1. กุสลตฺติก}\csromanmarkhii{9. สนิทสฺสนทุก}\csromanheadernomark{2}{1. กุสลตฺติก 9. สนิทสฺสนทุก}
\csromanheader{2}{1-7. ปฏิจฺจาทิวาร ปจฺจยจตุกฺก}
\proseitem{10}{กุสลํ อนิทสฺสนํ ธมฺมํ ปฏิจฺจ กุสโล อนิทสฺสโน ธมฺโม อุปฺปชฺชติ เหตุปจฺจยา. ตีณิ.}

\prose{อกุสลํ อนิทสฺสนํ ธมฺมํ ปฏิจฺจ อกุสโล อนิทสฺสโน ธมฺโม อุปฺปชฺชติ เหตุปจฺจยา. ตีณิ.}

\prose{อพฺยากตํ อนิทสฺสนํ ธมฺมํ ปฏิจฺจ อพฺยากโต อนิทสฺสโน ธมฺโม อุปฺปชฺชติ เหตุปจฺจยา. (1)}

\prose{กุสลํ อนิทสฺสนญฺจ อพฺยากตํ อนิทสฺสนญฺจ ธมฺมํ ปฏิจฺจ อพฺยากโต อนิทสฺสโน ธมฺโม อุปฺปชฺชติ เหตุปจฺจยา. (1)}

\prose{อกุสลํ อนิทสฺสนญฺจ อพฺยากตํ อนิทสฺสนญฺจ ธมฺมํ ปฏิจฺจ อพฺยากโต อนิทสฺสโน ธมฺโม อุปฺปชฺชติ เหตุปจฺจยา. (1)}

\proseitem{11}{กุสลํ อนิทสฺสนํ ธมฺมํ ปฏิจฺจ กุสโล อนิทสฺสโน ธมฺโม อุปฺปชฺชติ อารมฺมณ\-ปจฺจยา. (1)}

\prose{อกุสลํ อนิทสฺสนํ ธมฺมํ ปฏิจฺจ อกุสโล อนิทสฺสโน ธมฺโม อุปฺปชฺชติ อารมฺมณ\-ปจฺจยา. (1)}

\prose{อพฺยากตํ อนิทสฺสนํ ธมฺมํ ปฏิจฺจ อพฺยากโต อนิทสฺสโน ธมฺโม อุปฺปชฺชติ อารมฺมณ\-ปจฺจยา. (1) (สํขิตฺตํ.)}

\proseitem{12}{เหตุยา นว, อารมฺมเณ ตีณิ ฯเปฯ อวิคเต นว. (สํขิตฺตํ.)}

\csromanheader{2}{ธมฺมานุโลม ติกทุก\-ปฏฺฐาน\-ปาฬิ ปจฺจนีย\-น\-เหตุ}
\proseitem{13}{\csromanfolio{570}อกุสลํ อนิทสฺสนํ ธมฺมํ ปฏิจฺจ อกุสโล อนิทสฺสโน ธมฺโม อุปฺปชฺชติ นเหตุปจฺจยา. (1)}

\prose{อพฺยากตํ อนิทสฺสนํ ธมฺมํ ปฏิจฺจ อพฺยากโต อนิทสฺสโน ธมฺโม อุปฺปชฺชติ นเหตุปจฺจยา. (1) (สํขิตฺตํ.)}

\prose{นเหตุยา เทฺว, นอารมฺมเณ ปญฺจ, นอธิปติยา นว ฯเปฯ โนวิคเต ปญฺจ. (สํขิตฺตํ. สหชาตวารมฺปิ ฯเปฯ สมฺปยุตฺตวารมฺปิ ปฏิจฺจ\-วาร\-สทิสํ.)}

\csromanheader{2}{\textbf{เหตุ อารมฺมณ}}
\proseitem{14}{กุสโล อนิทสฺสโน ธมฺโม กุสลสฺส อนิทสฺสนสฺส ธมฺมสฺส เหตุปจฺจเยน ปจฺจโย. ตีณิ.}

\prose{อกุสโล อนิทสฺสโน ธมฺโม อกุสลสฺส อนิทสฺสนสฺส ธมฺมสฺส เหตุปจฺจเยน ปจฺจโย. ตีณิ.}

\prose{อพฺยากโต อนิทสฺสโน ธมฺโม อพฺยากตสฺส อนิทสฺสนสฺส ธมฺมสฺส เหตุปจฺจเยน ปจฺจโย. (1)}

\proseitem{15}{กุสโล อนิทสฺสโน ธมฺโม กุสลสฺส อนิทสฺสนสฺส ธมฺมสฺส อารมฺมณ\-ปจฺจเยน ปจฺจโย. ตีณิ.}

\prose{อกุสโล อนิทสฺสโน ธมฺโม อกุสลสฺส อนิทสฺสนสฺส ธมฺมสฺส อารมฺมณ\-ปจฺจเยน ปจฺจโย. ตีณิ.}

\prose{อพฺยากโต อนิทสฺสโน ธมฺโม อพฺยากตสฺส อนิทสฺสนสฺส ธมฺมสฺส อารมฺมณ\-ปจฺจเยน ปจฺจโย. ตีณิ. (สํขิตฺตํ.)}

\proseitem{16}{เหตุยา สตฺต, อารมฺมเณ นว, อธิปติยา ทส, อนนฺตเร สตฺต ฯเปฯ อุปนิสฺสเย นว ฯเปฯ อวิคเต เตรส. (สํขิตฺตํ.)}

\prose{(ยถา กุสลตฺติเก ปญฺหาวารํ, เอวํ วิตฺถาเรตพฺพํ.)}

\csromanmarkh{1. กุสลตฺติก}\csromanmarkhii{10. สปฺปฏิฆทุก}\csromanheadernomark{2}{1. \textbf{กุสลตฺติก} 10. สปฺปฏิฆทุก}
\csromanheader{2}{1-7. ปฏิจฺจาทิวาร ปจฺจยจตุกฺกเหตุ}
\proseitem{17}{อพฺยากตํ สปฺปฏิฆํ ธมฺมํ ปฏิจฺจ อพฺยากโต สปฺปฏิโฆ ธมฺโม อุปฺปชฺชติ เหตุปจฺจยา. (สํขิตฺตํ.)}

\csromanmarkh{1. กุสลตฺติก}\csromanmarkhii{10. สปฺปฏิฆทุก}\csromanheadernomark{2}{1. กุสลตฺติก 10. สปฺปฏิฆทุก}
\prose{\csromanfolio{571}เหตุยา เอกํ ฯเปฯ อวิคเต เอกํ. (สํขิตฺตํ.) (สหชาตวาเรปิ ฯเปฯ ปญฺหาวาเรปิ สพฺพตฺถ เอกํ.)}

\proseitem{18}{กุสลํ อปฺปฏิฆํ ธมฺมํ ปฏิจฺจ กุสโล อปฺปฏิโฆ ธมฺโม อุปฺปชฺชติ เหตุปจฺจยา. ตีณิ.}

\prose{อกุสลํ อปฺปฏิฆํ ธมฺมํ ปฏิจฺจ อกุสโล อปฺปฏิโฆ ธมฺโม อุปฺปชฺชติ เหตุปจฺจยา. ตีณิ.}

\prose{อพฺยากตํ อปฺปฏิฆํ ธมฺมํ ปฏิจฺจ อพฺยากโต อปฺปฏิโฆ ธมฺโม อุปฺปชฺชติ เหตุปจฺจยา. (1)}

\prose{กุสลํ อปฺปฏิฆญฺจ อพฺยากตํ อปฺปฏิฆญฺจ ธมฺมํ ปฏิจฺจ อพฺยากโต อปฺปฏิโฆ ธมฺโม อุปฺปชฺชติ เหตุปจฺจยา. (1)}

\prose{อกุสลํ อปฺปฏิฆญฺจ อพฺยากตํ อปฺปฏิฆญฺจ ธมฺมํ ปฏิจฺจ อพฺยากโต อปฺปฏิโฆ ธมฺโม อุปฺปชฺชติ เหตุปจฺจยา. (1)}

\proseitem{19}{กุสลํ อปฺปฏิฆํ ธมฺมํ ปฏิจฺจ กุสโล อปฺปฏิโฆ ธมฺโม อุปฺปชฺชติ อารมฺมณ\-ปจฺจยา. (1)}

\prose{อกุสลํ อปฺปฏิฆํ ธมฺมํ ปฏิจฺจ อกุสโล อปฺปฏิโฆ ธมฺโม อุปฺปชฺชติ อารมฺมณ\-ปจฺจยา. (1)}

\prose{อพฺยากตํ อปฺปฏิฆํ ธมฺมํ ปฏิจฺจ อพฺยากโต อปฺปฏิโฆ ธมฺโม อุปฺปชฺชติ อารมฺมณ\-ปจฺจยา. (1) (สํขิตฺตํ.)}

\proseitem{20}{เหตุยา นว, อารมฺมเณ ตีณิ, อธิปติยา นว ฯเปฯ อวิคเต นว. (สํขิตฺตํ.)}

\prose{ปจฺจนีย\-น\-เหตุ}

\proseitem{21}{อกุสลํ อปฺปฏิฆํ ธมฺมํ ปฏิจฺจ อกุสโล อปฺปฏิโฆ ธมฺโม อุปฺปชฺชติ นเหตุปจฺจยา. (1)}

\prose{อพฺยากตํ อปฺปฏิฆํ ธมฺมํ ปฏิจฺจ อพฺยากโต อปฺปฏิโฆ ธมฺโม อุปฺปชฺชติ นเหตุปจฺจยา. (1) (สํขิตฺตํ.)}

\proseitem{22}{\csromanfolio{572}นเหตุยา เทฺว, นอารมฺมเณ ปญฺจ, นอธิปติยา นว ฯเปฯ โนวิคเต ปญฺจ. (สํขิตฺตํ. สหชาตวารมฺปิ ฯเปฯ สมฺปยุตฺตวารมฺปิ ปฏิจฺจ\-วาร\-สทิสํ.)}

\proseitem{23}{กุสโล อปฺปฏิโฆ ธมฺโม กุสลสฺส อปฺปฏิฆสฺส ธมฺมสฺส เหตุปจฺจเยน ปจฺจโย. ตีณิ.}

\prose{อกุสโล อปฺปฏิโฆ ธมฺโม อกุสลสฺส อปฺปฏิฆสฺส ธมฺมสฺส เหตุปจฺจเยน ปจฺจโย. ตีณิ.}

\prose{อพฺยากโต อปฺปฏิโฆ ธมฺโม อพฺยากตสฺส อปฺปฏิฆสฺส ธมฺมสฺส เหตุปจฺจเยน ปจฺจโย. (1) (สํขิตฺตํ.)}

\proseitem{24}{\textbf{เหตุ}ยา สตฺต, อารมฺมเณ นว, อธิปติยา ทส ฯเปฯ อวิคเต เตรส. (สํขิตฺตํ. ยถา กุสลตฺติเก ปญฺหาวารํ, เอวํ วิตฺถาเรตพฺพํ.)}

\csromanmarkh{1. กุสลตฺติก}\csromanmarkhii{11. รูปีทุก}\csromanheadernomark{2}{1. \textbf{กุสลตฺติก} 11. รูปีทุก}
\csromanheader{2}{1-7. ปฏิจฺจาทิวาร ปจฺจยจตุกฺกเหตุ}
\proseitem{25}{อพฺยากตํ รูปิํ ธมฺมํ ปฏิจฺจ อพฺยากโต รูปี ธมฺโม อุปฺปชฺชติ เหตุปจฺจยา. (สํขิตฺตํ.)}

\prose{\textbf{เหตุ}ยา เอกํ, อารมฺมเณ เอกํ ฯเปฯ อวิคเต เอกํ. (สํขิตฺตํ.) (สหชาตวาเรปิ ฯเปฯ ปญฺหาวาเรปิ สพฺพตฺถ เอกํ.)}

\proseitem{26}{กุสลํ อรูปิํ ธมฺมํ ปฏิจฺจ กุสโล อรูปี ธมฺโม อุปฺปชฺชติ เหตุปจฺจยา. (1)}

\prose{อกุสลํ อรูปิํ ธมฺมํ ปฏิจฺจ อกุสโล อรูปี ธมฺโม อุปฺปชฺชติ เหตุปจฺจยา. (1)}

\prose{อพฺยากตํ อรูปิํ ธมฺมํ ปฏิจฺจ อพฺยากโต อรูปี ธมฺโม อุปฺปชฺชติ เหตุปจฺจยา. (1) (สํขิตฺตํ.)}

\csromanmarkh{1. กุสลตฺติก}\csromanmarkhii{12. โลกิยทุก}\csromanheadernomark{2}{1. กุสลตฺติก 12. โลกิยทุก}
\proseitem{27}{\csromanfolio{573}เหตุยา ตีณิ, อารมฺมเณ ตีณิ ฯเปฯ อวิคเต ตีณิ. (สํขิตฺตํ.)}

\prose{นเหตุยา เทฺว, นอธิปติยา ตีณิ. (สํขิตฺตํ.)}

\prose{(สหชาตวารมฺปิ ฯเปฯ สมฺปยุตฺตวารมฺปิ ปฏิจฺจ\-วาร\-สทิสํ.)}

\csromanheader{2}{\textbf{เหตุ อารมฺมณ}}
\proseitem{28}{กุสโล อรูปี ธมฺโม กุสลสฺส อรูปิสฺส ธมฺมสฺส เหตุปจฺจเยน ปจฺจโย. (1)}

\prose{อกุสโล อรูปี ธมฺโม อกุสลสฺส อรูปิสฺส ธมฺมสฺส เหตุปจฺจเยน ปจฺจโย. (1)}

\prose{อพฺยากโต อรูปี ธมฺโม อพฺยากตสฺส อรูปิสฺส ธมฺมสฺส เหตุปจฺจเยน ปจฺจโย. (1)}

\prose{กุสโล อรูปี ธมฺโม กุสลสฺส อรูปิสฺส ธมฺมสฺส อารมฺมณ\-ปจฺจเยน ปจฺจโย. ตีณิ.}

\prose{อกุสโล อรูปี ธมฺโม อกุสลสฺส อรูปิสฺส ธมฺมสฺส อารมฺมณ\-ปจฺจเยน ปจฺจโย. ตีณิ.}

\prose{อพฺยากโต อรูปี ธมฺโม อพฺยากตสฺส อรูปิสฺส ธมฺมสฺส อารมฺมณ\-ปจฺจเยน ปจฺจโย. (สํขิตฺตํ.)}

\proseitem{29}{เหตุยา ตีณิ, อารมฺมเณ นว, อธิปติยา สตฺต ฯเปฯ อวิคเต ตีณิ. (สํขิตฺตํ.)}

\prose{(ยถา กุสลตฺติเก ปญฺหาวารํ, เอวํ วิตฺถาเรตพฺพํ.)}

\csromanmarkh{1. กุสลตฺติก}\csromanmarkhii{12. โลกิยทุก}\csromanheadernomark{2}{1. กุสลตฺติก 12. โลกิยทุก}
\csromanheader{2}{1-7. ปฏิจฺจาทิวาร ปจฺจยจตุกฺกเหตุ}
\proseitem{30}{กุสลํ โลกิยํ ธมฺมํ ปฏิจฺจ กุสโล โลกิโย ธมฺโม อุปฺปชฺชติ เหตุปจฺจยา.  กุสลํ โลกิยํ ธมฺมํ ปฏิจฺจ อพฺยากโต \csromanfolio{574}โลกิโย ธมฺโม อุปฺปชฺชติ เหตุปจฺจยา.  กุสลํ โลกิยํ ธมฺมํ ปฏิจฺจ กุสโล โลกิโย จ อพฺยากโต โลกิโย จ ธมฺมา อุปฺปชฺชนฺติ เหตุปจฺจยา. (3)}

\prose{อกุสลํ โลกิยํ ธมฺมํ ปฏิจฺจ อกุสโล โลกิโย ธมฺโม อุปฺปชฺชติ เหตุปจฺจยา. ตีณิ.}

\prose{อพฺยากตํ โลกิยํ ธมฺมํ ปฏิจฺจ อพฺยากโต โลกิโย ธมฺโม อุปฺปชฺชติ เหตุปจฺจยา. (1)}

\prose{กุสลํ โลกิยญฺจ อพฺยากตํ โลกิยญฺจ ธมฺมํ ปฏิจฺจ อพฺยากโต โลกิโย ธมฺโม อุปฺปชฺชติ เหตุปจฺจยา. (1)}

\prose{อกุสโล โลกิยญฺจ อพฺยากตํ โลกิยญฺจ ธมฺมํ ปฏิจฺจ อพฺยากโต โลกิโย ธมฺโม อุปฺปชฺชติ เหตุปจฺจยา. (1) (สํขิตฺตํ.)}

\proseitem{31}{เหตุยา นว, อารมฺมเณ ตีณิ ฯเปฯ อวิคเต นว. (สํขิตฺตํ.)}

\prose{ปจฺจนีย\-น\-เหตุ}

\proseitem{32}{อกุสลํ โลกิยํ ธมฺมํ ปฏิจฺจ อกุสโล โลกิโย ธมฺโม อุปฺปชฺชติ นเหตุปจฺจยา. (1)}

\prose{อพฺยากตํ โลกิยํ ธมฺมํ ปฏิจฺจ อพฺยากโต โลกิโย ธมฺโม อุปฺปชฺชติ นเหตุปจฺจยา. (1) (สํขิตฺตํ.)}

\proseitem{33}{นเหตุยา เทฺว, นอารมฺมเณ ปญฺจ, นอธิปติยา นว ฯเปฯ โนวิคเต ปญฺจ. (สํขิตฺตํ. สหชาตวารมฺปิ ฯเปฯ สมฺปยุตฺตวารมฺปิ ปฏิจฺจ\-วาร\-สทิสํ.)}

\csromanheader{2}{\textbf{เหตุ อารมฺมณ}}
\proseitem{34}{กุสโล โลกิโย ธมฺโม กุสลสฺส โลกิยสฺส ธมฺมสฺส เหตุปจฺจเยน ปจฺจโย. ตีณิ.}

\prose{อกุสโล โลกิโย ธมฺโม อกุสลสฺส โลกิยสฺส ธมฺมสฺส เหตุปจฺจเยน ปจฺจโย. ตีณิ.}

\prose{อพฺยากโต โลกิโย ธมฺโม อพฺยากตสฺส โลกิยสฺส ธมฺมสฺส เหตุปจฺจเยน ปจฺจโย. (1)}

\csromanmarkh{1. กุสลตฺติก}\csromanmarkhii{12. โลกิยทุก}\csromanheadernomark{2}{1. กุสลตฺติก 12. โลกิยทุก}
\proseitem{35}{\csromanfolio{575}กุสโล โลกิโย ธมฺโม กุสลสฺส โลกิยสฺส ธมฺมสฺส อารมฺมณ\-ปจฺจเยน ปจฺจโย. ตีณิ.}

\prose{อกุสโล โลกิโย ธมฺโม อกุสลสฺส โลกิยสฺส ธมฺมสฺส อารมฺมณ\-ปจฺจเยน ปจฺจโย. ตีณิ.}

\prose{อพฺยากโต โลกิโย ธมฺโม อพฺยากตสฺส โลกิยสฺส ธมฺมสฺส อารมฺมณ\-ปจฺจเยน ปจฺจโย. ตีณิ. (สํขิตฺตํ.)}

\proseitem{36}{เหตุยา สตฺต, อารมฺมเณ นว, อธิปติยา นว ฯเปฯ อวิคเต เตรส. (สํขิตฺตํ. ยถา กุสลตฺติเก ปญฺหาวารํ, เอวํ วิตฺถาเรตพฺพํ.)}

\csromanheader{2}{\textbf{โลกุตฺตรปท 1-7. ปฏิจฺจาทิวาร}}
\prose{ปจฺจยจตุกฺกเหตุ}

\proseitem{37}{กุสลํ โลกุตฺตรํ ธมฺมํ ปฏิจฺจ กุสโล โลกุตฺตโร ธมฺโม อุปฺปชฺชติ เหตุปจฺจยา. (1)}

\prose{อพฺยากตํ โลกุตฺตรํ ธมฺมํ ปฏิจฺจ อพฺยากโต โลกุตฺตโร ธมฺโม อุปฺปชฺชติ เหตุปจฺจยา. (1) (สํขิตฺตํ.)}

\proseitem{38}{เหตุยา เทฺว, อารมฺมเณ เทฺว, อธิปติยา เทฺว ฯเปฯ อวิคเต เทฺว. (สํขิตฺตํ.)}

\prose{นอธิปติยา เทฺว. (สํขิตฺตํ. สหชาตวารมฺปิ ฯเปฯ สมฺปยุตฺตวารมฺปิ ปฏิจฺจ\-วาร\-สทิสํ.)}

\csromanheader{2}{\textbf{เหตุ อารมฺมณาทิ}}
\proseitem{39}{กุสโล โลกุตฺตโร ธมฺโม กุสลสฺส โลกุตฺตรสฺส ธมฺมสฺส เหตุปจฺจเยน ปจฺจโย. (1)}

\prose{อพฺยากโต โลกุตฺตโร ธมฺโม อพฺยากตสฺส โลกุตฺตรสฺส ธมฺมสฺส เหตุปจฺจเยน ปจฺจโย. (1)}

\prose{อพฺยากโต โลกุตฺตโร ธมฺโม อพฺยากตสฺส โลกุตฺตรสฺส ธมฺมสฺส อารมฺมณ\-ปจฺจเยน ปจฺจโย.  อพฺยากโต โลกุตฺตโร ธมฺโม กุสลสฺส โลกุตฺตรสฺส ธมฺมสฺส อารมฺมณ\-ปจฺจเยน ปจฺจโย. (2)}

\prose{\csromanfolio{576}กุสโล โลกุตฺตโร ธมฺโม กุสลสฺส โลกุตฺตรสฺส ธมฺมสฺส อธิปติ\-ปจฺจเยน ปจฺจโย. (1)}

\prose{อพฺยากโต โลกุตฺตโร ธมฺโม อพฺยากตสฺส โลกุตฺตรสฺส ธมฺมสฺส อธิปติ\-ปจฺจเยน ปจฺจโย.  อพฺยากโต โลกุตฺตโร ธมฺโม กุสลสฺส โลกุตฺตรสฺส ธมฺมสฺส อธิปติ\-ปจฺจเยน ปจฺจโย. (2)}

\prose{กุสโล โลกุตฺตโร ธมฺโม อพฺยากตสฺส โลกุตฺตรสฺส ธมฺมสฺส อนนฺตร\-ปจฺจเยน ปจฺจโย. (1)}

\prose{อพฺยากโต โลกุตฺตโร ธมฺโม อพฺยากตสฺส โลกุตฺตรสฺส ธมฺมสฺส อนนฺตร\-ปจฺจเยน ปจฺจโย. (1) (สํขิตฺตํ.)}

\proseitem{40}{เหตุยา เทฺว, อารมฺมเณ เทฺว, อธิปติยา ตีณิ, อนนฺตเร เทฺว, สมนนฺตเร เทฺว, สหชาเต เทฺว, อญฺญมญฺเญ เทฺว, นิสฺสเย เทฺว, อุปนิสฺสเย จตฺตาริ ฯเปฯ อวิคเต เทฺว. (สํขิตฺตํ.)}

\prose{(ยถา กุสลตฺติเก ปญฺหาวารํ, เอวํ วิตฺถาเรตพฺพํ.)}

\csromanmarkh{1. กุสลตฺติก}\csromanmarkhii{13. เกนจิ\-วิญฺเญยฺยทุก}\csromanheadernomark{2}{1. กุสลตฺติก 13. เกนจิ\-วิญฺเญยฺยทุก}
\csromanheader{2}{1-7. ปฏิจฺจาทิวาร ปจฺจยจตุกฺก}
\proseitem{41}{กุสลํ เกนจิ วิญฺเญยฺยํ ธมฺมํ ปฏิจฺจ กุสโล เกนจิ วิญฺเญยฺโย ธมฺโม อุปฺปชฺชติ เหตุปจฺจยา.  กุสลํ เกนจิ วิญฺเญยฺยํ ธมฺมํ ปฏิจฺจ อพฺยากโต เกนจิ วิญฺเญยฺโย ธมฺโม อุปฺปชฺชติ เหตุปจฺจยา.  กุสลํ เกนจิ วิญฺเญยฺยํ ธมฺมํ ปฏิจฺจ กุสโล เกนจิ วิญฺเญยฺโย จ อพฺยากโต เกนจิ วิญฺเญยฺโย จ ธมฺมา อุปฺปชฺชนฺติ เหตุปจฺจยา. (3)}

\prose{อกุสลํ เกนจิ วิญฺเญยฺยํ ธมฺมํ ปฏิจฺจ อกุสโล เกนจิ วิญฺเญยฺโย ธมฺโม อุปฺปชฺชติ เหตุปจฺจยา. ตีณิ.}

\prose{อพฺยากตํ เกนจิ วิญฺเญยฺยํ ธมฺมํ ปฏิจฺจ อพฺยากโต เกนจิ วิญฺเญยฺโย ธมฺโม อุปฺปชฺชติ เหตุปจฺจยา. (1)}

\prose{กุสลํ เกนจิ วิญฺเญยฺยญฺจ อพฺยากตํ เกนจิ วิญฺเญยฺยญฺจ ธมฺมํ ปฏิจฺจ อพฺยากโต เกนจิ วิญฺเญยฺโย ธมฺโม อุปฺปชฺชติ เหตุปจฺจยา. (1)}

\csromanmarkh{1. กุสลตฺติก}\csromanmarkhii{13. เกนจิ\-วิญฺเญยฺยทุก}\csromanheadernomark{2}{1. กุสลตฺติก 13. เกนจิ\-วิญฺเญยฺยทุก}
\prose{\csromanfolio{577}อกุสลํ เกนจิ วิญฺเญยฺยญฺจ อพฺยากตํ เกนจิ วิญฺเญยฺยญฺจ ธมฺมํ ปฏิจฺจ อพฺยากโต เกนจิ วิญฺเญยฺโย ธมฺโม อุปฺปชฺชติ เหตุปจฺจยา. (1)}

\proseitem{42}{กุสลํ เกนจิ วิญฺเญยฺยํ ธมฺมํ ปฏิจฺจ กุสโล เกนจิ วิญฺเญยฺโย ธมฺโม อุปฺปชฺชติ อารมฺมณ\-ปจฺจยา. (1)}

\prose{อกุสลํ เกนจิ วิญฺเญยฺยํ ธมฺมํ ปฏิจฺจ อกุสโล เกนจิ วิญฺเญยฺโย ธมฺโม อุปฺปชฺชติ อารมฺมณ\-ปจฺจยา. (1)}

\prose{อพฺยากตํ เกนจิ วิญฺเญยฺยํ ธมฺมํ ปฏิจฺจ อพฺยากโต เกนจิ วิญฺเญยฺโย ธมฺโม อุปฺปชฺชติ อารมฺมณ\-ปจฺจยา. (1) (สํขิตฺตํ.)}

\proseitem{43}{เหตุยา นว, อารมฺมเณ ตีณิ ฯเปฯ อวิคเต นว. (สํขิตฺตํ.)}

\prose{ปจฺจนีย\-น\-เหตุ}

\proseitem{44}{อกุสลํ เกนจิ วิญฺเญยฺยํ ธมฺมํ ปฏิจฺจ อกุสโล เกนจิ วิญฺเญยฺโย ธมฺโม อุปฺปชฺชติ นเหตุปจฺจยา. (1)}

\prose{อพฺยากตํ เกนจิ วิญฺเญยฺยํ ธมฺมํ ปฏิจฺจ อพฺยากโต เกนจิ วิญฺเญยฺโย ธมฺโม อุปฺปชฺชติ นเหตุปจฺจยา. (1) (สํขิตฺตํ.)}

\proseitem{45}{นเหตุยา เทฺว, นอารมฺมเณ ปญฺจ, นอธิปติยา นว ฯเปฯ โนวิคเต ปญฺจ. (สํขิตฺตํ. สหชาตวารมฺปิ ฯเปฯ สมฺปยุตฺตวารมฺปิ ปฏิจฺจ\-วาร\-สทิสํ.)}

\proseitem{46}{กุสโล เกนจิ วิญฺเญยฺโย ธมฺโม กุสลสฺส เกนจิ วิญฺเญยฺยสฺส ธมฺมสฺส เหตุปจฺจเยน ปจฺจโย. ตีณิ.}

\prose{อกุสโล เกนจิ วิญฺเญยฺโย ธมฺโม อกุสลสฺส เกนจิ วิญฺเญยฺยสฺส ธมฺมสฺส เหตุปจฺจเยน ปจฺจโย. ตีณิ.}

\prose{อพฺยากโต เกนจิ วิญฺเญยฺโย ธมฺโม อพฺยากตสฺส เกนจิ วิญฺเญยฺยสฺส ธมฺมสฺส เหตุปจฺจเยน ปจฺจโย. (1) (สํขิตฺตํ.)}

\proseitem{47}{\textbf{เหตุ}ยา สตฺต, อารมฺมเณ นว ฯเปฯ อวิคเต เตรส.}

\prose{(สํขิตฺตํ. ยถา กุสลตฺติเก ปญฺหาวารํ, เอวํ วิตฺถาเรตพฺพํ.)}

\csromanmatikaanchor{mtk.143}
\csromantocmarkref{subparagraph}{14-19. อาสวทุกาทิ}{mtk.143}
\bookmark[dest={mtk.143},level=5]{14-19. อาสวทุกาทิ}
\proseitemwithcloser{48}{\csromanfolio{578}กุสลํ นเกนจิ วิญฺเญยฺยํ ธมฺมํ ปฏิจฺจ กุสโล นเกนจิ วิญฺเญยฺโย ธมฺโม อุปฺปชฺชติ เหตุปจฺจยา. (เกนจิ\-วิญฺเญยฺย\-สทิสํ.)}{กุสล\-ตฺติก\-จูฬ\-นฺตรทุกํ นิฏฺฐิตํ.}
\csromanmarkh{1. กุสลตฺติก}\csromanmarkhii{14. อาสวทุก}\csromanheadernomark{6}{1. \textbf{กุสลตฺติก} 14. อาสวทุก}
\csromanheader{2}{1-7. ปฏิจฺจาทิวาร ปจฺจยจตุกฺกเหตุ}
\proseitem{1}{อกุสลํ อาสวํ ธมฺมํ ปฏิจฺจ อกุสโล อาสโว ธมฺโม อุปฺปชฺชติ เหตุปจฺจยา. (สํขิตฺตํ.)}

\prose{เหตุยา เอกํ ฯเปฯ อวิคเต เอกํ. (สํขิตฺตํ.)}

\prose{(สหชาตวาเรปิ ฯเปฯ สมฺปยุตฺตวาเรปิ ปญฺหาวาเรปิ สพฺพตฺถ เอกํ.)}

\proseitem{2}{กุสลํ โนอสวํ ธมฺมํ ปฏิจฺจ กุสโล โนอสโว ธมฺโม อุปฺปชฺชติ เหตุปจฺจยา.  กุสลํ โนอสวํ ธมฺมํ ปฏิจฺจ อพฺยากโต โนอสโว ธมฺโม อุปฺปชฺชติ เหตุปจฺจยา.  กุสลํ โนอสวํ ธมฺมํ ปฏิจฺจ กุสโล โนอสโว จ อพฺยากโต โนอสโว จ ธมฺมา อุปฺปชฺชนฺติ เหตุปจฺจยา. (3)}

\prose{อกุสลํ โนอสวํ ธมฺมํ ปฏิจฺจ อกุสโล โนอสโว ธมฺโม อุปฺปชฺชติ เหตุปจฺจยา. ตีณิ.}

\prose{อพฺยากตํ โนอสวํ ธมฺมํ ปฏิจฺจ อพฺยากโต โนอสโว ธมฺโม อุปฺปชฺชติ เหตุปจฺจยา. (1)}

\prose{กุสลํ โนอสวญฺจ อพฺยากตํ โนอสวญฺจ ธมฺมํ ปฏิจฺจ อพฺยากโต โนอสโว ธมฺโม อุปฺปชฺชติ เหตุปจฺจยา. (1)}

\prose{อกุสลํ โนอสวญฺจ อพฺยากตํ โนอสวญฺจ ธมฺมํ ปฏิจฺจ อพฺยากโต โนอสโว ธมฺโม อุปฺปชฺชติ เหตุปจฺจยา. (1)}

\proseitem{3}{กุสลํ โนอสวํ ธมฺมํ ปฏิจฺจ กุสโล โนอสโว ธมฺโม อุปฺปชฺชติ อารมฺมณ\-ปจฺจยา. (1)}

\csromanmarkh{1. กุสลตฺติก}\csromanmarkhii{15. สาสวทุก}\csromanheadernomark{2}{1. กุสลตฺติก 15. สาสวทุก}
\prose{\csromanfolio{579}อกุสลํ โนอสวํ ธมฺมํ ปฏิจฺจ อกุสโล โนอสโว ธมฺโม อุปฺปชฺชติ อารมฺมณ\-ปจฺจยา. (1)}

\prose{อพฺยากตํ โนอสวํ ธมฺมํ ปฏิจฺจ อพฺยากโต โนอสโว ธมฺโม อุปฺปชฺชติ อารมฺมณ\-ปจฺจยา. (1) (สํขิตฺตํ.)}

\proseitem{4}{เหตุยา นว, อารมฺมเณ ตีณิ, อธิปติยา นว ฯเปฯ อวิคเต นว. (สํขิตฺตํ.)}

\prose{นเหตุยา เอกํ, นอารมฺมเณ ปญฺจ, นอธิปติยา นว ฯเปฯ โนวิคเต ปญฺจ. (สํขิตฺตํ. สหชาตวารมฺปิ ฯเปฯ สมฺปยุตฺตวารมฺปิ ปฏิจฺจ\-วาร\-สทิสํ.)}

\prose{เหตุอารมฺมณ}

\proseitem{5}{กุสโล โนอสโว ธมฺโม กุสลสฺส โนอสวสฺส ธมฺมสฺส เหตุปจฺจเยน ปจฺจโย. ตีณิ.}

\prose{อกุสโล โนอสโว ธมฺโม อกุสลสฺส โนอสวสฺส ธมฺมสฺส เหตุปจฺจเยน ปจฺจโย. ตีณิ.}

\prose{อพฺยากโต โนอสโว ธมฺโม อพฺยากตสฺส โนอสวสฺส ธมฺมสฺส เหตุปจฺจเยน ปจฺจโย. (1)}

\prose{กุสโล โนอสโว ธมฺโม กุสลสฺส โนอสวสฺส ธมฺมสฺส อารมฺมณ\-ปจฺจเยน ปจฺจโย. (สํขิตฺตํ.)}

\proseitem{6}{เหตุยา สตฺต, อารมฺมเณ นว, อธิปติยา ทส ฯเปฯ อวิคเต เตรส.}

\prose{(สํขิตฺตํ. ยถา กุสลตฺติเก ปญฺหาวารํ, เอวํ วิตฺถาเรตพฺพํ.)}

\csromanmarkh{1. กุสลตฺติก}\csromanmarkhii{15. สาสวทุก}\csromanheadernomark{2}{1. กุสลตฺติก 15. สาสวทุก}
\csromanheader{2}{1-7. ปฏิจฺจาทิวาร ปจฺจยจตุกฺกเหตุ}
\proseitem{7}{กุสลํ สาสวํ ธมฺมํ ปฏิจฺจ กุสโล สาสโว ธมฺโม อุปฺปชฺชติ เหตุปจฺจยา. ตีณิ.}

\prose{อกุสลํ สาสวํ ธมฺมํ ปฏิจฺจ อกุสโล สาสโว ธมฺโม อุปฺปชฺชติ เหตุปจฺจยา. ตีณิ.}

\prose{\csromanfolio{580}อพฺยากตํ สาสวํ ธมฺมํ ปฏิจฺจ อพฺยากโต สาสโว ธมฺโม อุปฺปชฺชติ เหตุปจฺจยา. (1)}

\prose{กุสลํ สาสวญฺจ อพฺยากตํ สาสวญฺจ ธมฺมํ ปฏิจฺจ อพฺยากโต สาสโว ธมฺโม อุปฺปชฺชติ เหตุปจฺจยา. (1)}

\prose{อกุสลํ สาสวญฺจ อพฺยากตํ สาสวญฺจ ธมฺมํ ปฏิจฺจ อพฺยากโต สาสโว ธมฺโม อุปฺปชฺชติ เหตุปจฺจยา. (1) (สํขิตฺตํ.)}

\proseitem{8}{เหตุยา นว, อารมฺมเณ ตีณิ ฯเปฯ อวิคเต นว. (สํขิตฺตํ.)}

\prose{นเหตุยา เทฺว, นอารมฺมเณ ปญฺจ, นอธิปติยา นว. (สํขิตฺตํ. สหชาตวารมฺปิ ฯเปฯ สมฺปยุตฺตวารมฺปิ ปฏิจฺจ\-วาร\-สทิสํ.)}

\proseitem{9}{กุสโล สาสโว ธมฺโม กุสลสฺส สาสวสฺส ธมฺมสฺส เหตุปจฺจเยน ปจฺจโย. ตีณิ.}

\prose{อกุสโล สาสโว ธมฺโม อกุสลสฺส สาสวสฺส ธมฺมสฺส เหตุปจฺจเยน ปจฺจโย. ตีณิ.}

\prose{อพฺยากโต สาสโว ธมฺโม อพฺยากตสฺส สาสวสฺส ธมฺมสฺส เหตุปจฺจเยน ปจฺจโย. (1) (สํขิตฺตํ.)}

\prose{เหตุยา สตฺต, อารมฺมเณ นว ฯเปฯ อวิคเต เตรส. (สํขิตฺตํ.)}

\prose{(ยถา กุสลตฺติเก ปญฺหาวารํ, เอวํ วิตฺถาเรตพฺพํ.)}

\csromanheader{2}{\textbf{อนาสวปท 1-7. ปฏิจฺจาทิวาร}}
\prose{ปจฺจยจตุกฺกเหตุ}

\proseitem{10}{กุสลํ อนาสวํ ธมฺมํ ปฏิจฺจ กุสโล อนาสโว ธมฺโม อุปฺปชฺชติ เหตุปจฺจยา. (1)}

\prose{อพฺยากตํ อนาสวํ ธมฺมํ ปฏิจฺจ อพฺยากโต อนาสโว ธมฺโม อุปฺปชฺชติ เหตุปจฺจยา. (1) (สํขิตฺตํ.)}

\proseitem{11}{เหตุยา เทฺว ฯเปฯ อวิคเต เทฺว. (สํขิตฺตํ. สหชาตวารมฺปิ ฯเปฯ สมฺปยุตฺตวารมฺปิ ปฏิจฺจ\-วาร\-สทิสํ.)}

\csromanmarkh{1. กุสลตฺติก}\csromanmarkhii{16. อาสว\-สมฺปยุตฺตทุก เหตุ อารมฺมณาทิ}\csromanheadernomark{2}{1. กุสลตฺติก 16. อาสว\-สมฺปยุตฺตทุก เหตุ อารมฺมณาทิ}
\proseitem{12}{\csromanfolio{581}กุสโล อนาสโว ธมฺโม กุสลสฺส อนาสวสฺส ธมฺมสฺส เหตุปจฺจเยน ปจฺจโย. (1)}

\prose{อพฺยากโต อนาสโว ธมฺโม อพฺยากตสฺส อนาสวสฺส ธมฺมสฺส เหตุปจฺจเยน ปจฺจโย. (1)}

\prose{อพฺยากโต อนาสโว ธมฺโม อพฺยากตสฺส อนาสวสฺส ธมฺมสฺส อารมฺมณ\-ปจฺจเยน ปจฺจโย.  อพฺยากโต อนาสโว ธมฺโม กุสลสฺส อนาสวสฺส ธมฺมสฺส อารมฺมณ\-ปจฺจเยน ปจฺจโย. (2)}

\prose{กุสโล อนาสโว ธมฺโม กุสลสฺส อนาสวสฺส ธมฺมสฺส อธิปติ\-ปจฺจเยน ปจฺจโย. (1)}

\prose{อพฺยากโต อนาสโว ธมฺโม อพฺยากตสฺส อนาสวสฺส ธมฺมสฺส อธิปติ\-ปจฺจเยน ปจฺจโย.  อพฺยากโต อนาสโว ธมฺโม กุสลสฺส อนาสวสฺส ธมฺมสฺส อธิปติ\-ปจฺจเยน ปจฺจโย. (2)}

\prose{กุสโล อนาสโว ธมฺโม อพฺยากตสฺส อนาสวสฺส ธมฺมสฺส อนนฺตร\-ปจฺจเยน ปจฺจโย. (1)}

\prose{อพฺยากโต อนาสโว ธมฺโม อพฺยากตสฺส อนาสวสฺส ธมฺมสฺส อนนฺตร\-ปจฺจเยน ปจฺจโย. (1) (สํขิตฺตํ.)}

\proseitem{13}{เหตุยา เทฺว, อารมฺมเณ เทฺว, อธิปติยา ตีณิ, อนนฺตเร เทฺว, สมนนฺตเร เทฺว, สหชาเต เทฺว, อญฺญมญฺเญ เทฺว, นิสฺสเย เทฺว, อุปนิสฺสเย จตฺตาริ ฯเปฯ อวิคเต เทฺว. (สํขิตฺตํ. ยถา กุสลตฺติเก ปญฺหาวารํ, เอวํ วิตฺถาเรตพฺพํ.)}

\csromanmarkh{1. กุสลตฺติก}\csromanmarkhii{16. อาสว\-สมฺปยุตฺตทุก}\csromanheadernomark{2}{1. กุสลตฺติก 16. อาสว\-สมฺปยุตฺตทุก}
\csromanheader{2}{1-7. ปฏิจฺจาทิวาร ปจฺจยจตุกฺกเหตุ}
\proseitem{14}{อกุสลํ อาสว\-สมฺปยุตฺตํ ธมฺมํ ปฏิจฺจ อกุสโล อาสว\-สมฺปยุตฺโต ธมฺโม อุปฺปชฺชติ เหตุปจฺจยา. (สํขิตฺตํ.)}

\prose{เหตุยา เอกํ, อารมฺมเณ เอกํ. (สํขิตฺตํ.)}

\prose{(สหชาตวาเรปิ ฯเปฯ สมฺปยุตฺตวาเรปิ ปญฺหาวาเรปิ สพฺพตฺถ เอกํ.)}

\proseitem{15}{\csromanfolio{582}กุสลํ อาสว\-วิปฺปยุตฺตํ ธมฺมํ ปฏิจฺจ กุสโล อาสว\-วิปฺปยุตฺโต ธมฺโม อุปฺปชฺชติ เหตุปจฺจยา.  กุสลํ อาสว\-วิปฺปยุตฺตํ ธมฺมํ ปฏิจฺจ อพฺยากโต อาสว\-วิปฺปยุตฺโต ธมฺโม อุปฺปชฺชติ เหตุปจฺจยา.  กุสลํ อาสว\-วิปฺปยุตฺตํ ธมฺมํ ปฏิจฺจ กุสโล อาสว\-วิปฺปยุตฺโต จ อพฺยากโต อาสว\-วิปฺปยุตฺโต จ ธมฺมา อุปฺปชฺชนฺติ เหตุปจฺจยา. (3)}

\prose{อกุสลํ อาสว\-วิปฺปยุตฺตํ ธมฺมํ ปฏิจฺจ อพฺยากโต อาสว\-วิปฺปยุตฺโต ธมฺโม อุปฺปชฺชติ เหตุปจฺจยา. (1)}

\prose{อพฺยากตํ อาสว\-วิปฺปยุตฺตํ ธมฺมํ ปฏิจฺจ อพฺยากโต อาสว\-วิปฺปยุตฺโต ธมฺโม อุปฺปชฺชติ เหตุปจฺจยา. (1)}

\prose{กุสลํ อาสว\-วิปฺปยุตฺตญฺจ อพฺยากตํ อาสว\-วิปฺปยุตฺตญฺจ ธมฺมํ ปฏิจฺจ อพฺยากโต อาสว\-วิปฺปยุตฺโต ธมฺโม อุปฺปชฺชติ เหตุปจฺจยา. (1)}

\prose{อกุสลํ อาสว\-วิปฺปยุตฺตญฺจ อพฺยากตํ อาสว\-วิปฺปยุตฺตญฺจ ธมฺมํ ปฏิจฺจ อพฺยากโต อาสว\-วิปฺปยุตฺโต ธมฺโม อุปฺปชฺชติ เหตุปจฺจยา. (สํขิตฺตํ.)}

\proseitem{16}{เหตุยา สตฺต, อารมฺมเณ เทฺว, อธิปติยา สตฺต ฯเปฯ อวิคเต สตฺต. (สํขิตฺตํ.)}

\prose{(สหชาตวารมฺปิ ฯเปฯ สมฺปยุตฺตวารมฺปิ ปฏิจฺจ\-วาร\-สทิสํ.)}

\proseitem{17}{กุสโล อาสว\-วิปฺปยุตฺโต ธมฺโม กุสลสฺส อาสว\-วิปฺปยุตฺตสฺส ธมฺมสฺส เหตุปจฺจเยน ปจฺจโย. ตีณิ.}

\prose{อกุสโล อาสว\-วิปฺปยุตฺโต ธมฺโม อพฺยากตสฺส อาสว\-วิปฺปยุตฺตสฺส ธมฺมสฺส เหตุปจฺจเยน ปจฺจโย. (1)}

\prose{อพฺยากโต อาสว\-วิปฺปยุตฺโต ธมฺโม อพฺยากตสฺส อาสว\-วิปฺปยุตฺตสฺส ธมฺมสฺส เหตุปจฺจเยน ปจฺจโย. (สํขิตฺตํ.)}

\proseitem{18}{เหตุยา ปญฺจ, อารมฺมเณ นว, อธิปติยา ปญฺจ ฯเปฯ อวิคเต ทส. (สํขิตฺตํ.)}

\prose{(ยถา กุสลตฺติเก ปญฺหาวารํ, เอวํ วิตฺถาเรตพฺพํ.)}

\vspace{\tipitakaparskip}
\csromancenter{กุสลตฺติก 17. อาสว\-สาสวทุก 1. กุสลตฺติก 17. อาสว\-สาสวทุก}
\csromanheader{2}{1-7. ปฏิจฺจาทิวาร ปจฺจยจตุกฺกเหตุ}
\proseitem{19}{\csromanfolio{583}อกุสลํ อาสวญฺเจว สาสวญฺจ ธมฺมํ ปฏิจฺจ อกุสโล อาสโว เจว สาสโว จ ธมฺโม อุปฺปชฺชติ เหตุปจฺจยา. (สํขิตฺตํ.)}

\prose{เหตุยา เอกํ, อารมฺมเณ เอกํ ฯเปฯ อวิคเต เอกํ. (สํขิตฺตํ.)}

\prose{(สหชาตวาเรปิ ฯเปฯ สมฺปยุตฺตวาเรปิ ปญฺหาวาเรปิ สพฺพตฺถ เอกํ.)}

\proseitem{20}{กุสลํ สาสวญฺเจว โน จ อาสวํ ธมฺมํ ปฏิจฺจ กุสโล สาสโว เจว โน จ อาสโว ธมฺโม อุปฺปชฺชติ เหตุปจฺจยา. ตีณิ.}

\prose{อกุสลํ สาสวญฺเจว โน จ อาสวํ ธมฺมํ ปฏิจฺจ อกุสโล สาสโว เจว โน จ อาสโว ธมฺโม อุปฺปชฺชติ เหตุปจฺจยา. ตีณิ.}

\prose{อพฺยากตํ สาสวญฺเจว โน จ อาสวํ ธมฺมํ ปฏิจฺจ อพฺยากโต สาสโว เจว โน จ อาสโว ธมฺโม อุปฺปชฺชติ เหตุปจฺจยา. (1)}

\prose{กุสลํ สาสวญฺเจว โน จ อาสวญฺจ อพฺยากตํ สาสวญฺเจว โน จ อาสวญฺจ ธมฺมํ ปฏิจฺจ อพฺยากโต สาสโว เจว โน จ อาสโว ธมฺโม อุปฺปชฺชติ เหตุปจฺจยา. (1)}

\prose{อกุสลํ สาสวญฺเจว โน จ อาสวญฺจ อพฺยากตํ สาสวญฺเจว โน จ อาสวญฺจ ธมฺมํ ปฏิจฺจ อพฺยากโต สาสโว เจว โน จ อาสโว ธมฺโม อุปฺปชฺชติ เหตุปจฺจยา. (1) (สํขิตฺตํ.)}

\proseitem{21}{เหตุยา นว, อารมฺมเณ ตีณิ, อธิปติยา นว ฯเปฯ อวิคเต นว. (สํขิตฺตํ.)}

\prose{(สหชาตวารมฺปิ ฯเปฯ สมฺปยุตฺตวารมฺปิ ปฏิจฺจ\-วาร\-สทิสํ.)}

\proseitem{22}{กุสโล สาสโว เจว โน จ อาสโว ธมฺโม กุสลสฺส สาสวสฺส เจว โน จ อาสวสฺส ธมฺมสฺส เหตุปจฺจเยน ปจฺจโย. ตีณิ.}

\prose{อกุสโล สาสโว เจว โน จ อาสโว ธมฺโม อกุสลสฺส สาสวสฺส เจว โน จ อาสวสฺส ธมฺมสฺส เหตุปจฺจเยน ปจฺจโย. ตีณิ.}

\prose{\csromanfolio{584}อพฺยากโต สาสโว เจว โน จ อาสโว ธมฺโม อพฺยากตสฺส สาสวสฺส เจว โน จ อาสวสฺส ธมฺมสฺส เหตุปจฺจเยน ปจฺจโย. (1) (สํขิตฺตํ.)}

\proseitem{23}{เหตุยา สตฺต, อารมฺมเณ นว, อธิปติยา นว ฯเปฯ อวิคเต เตรส. (สํขิตฺตํ. ยถา กุสลตฺติเก ปญฺหาวารํ, เอวํ วิตฺถาเรตพฺพํ.)}

\prose{1. กุสลตฺติก 18. อาสว\-อาสว\-สมฺปยุตฺตทุก}

\csromanheader{2}{1-7. ปฏิจฺจาทิวาร ปจฺจยจตุกฺก}
\proseitem{24}{อกุสลํ อาสวญฺเจว อาสว\-สมฺปยุตฺตญฺจ ธมฺมํ ปฏิจฺจ อกุสโล อาสโว เจว อาสว\-สมฺปยุตฺโต จ ธมฺโม อุปฺปชฺชติ เหตุปจฺจยา. (สํขิตฺตํ.)}

\prose{เหตุยา เอกํ, อารมฺมเณ เอกํ ฯเปฯ อวิคเต เอกํ. (สํขิตฺตํ.)}

\prose{(สหชาตวาเรปิ ฯเปฯ สมฺปยุตฺตวาเรปิ ปญฺหาวาเรปิ สพฺพตฺถ เอกํ.)}

\proseitem{25}{อกุสลํ อาสว\-สมฺปยุตฺต\-ญฺเจว โน จ อาสวํ ธมฺมํ ปฏิจฺจ อกุสโล อาสว\-สมฺปยุตฺโต เจว โน จ อาสโว ธมฺโม อุปฺปชฺชติ เหตุปจฺจยา. (สํขิตฺตํ.)}

\prose{เหตุยา เอกํ, อารมฺมเณ เอกํ ฯเปฯ อวิคเต เอกํ. (สํขิตฺตํ.)}

\prose{(สหชาตวาเรปิ ฯเปฯ สมฺปยุตฺตวาเรปิ ปญฺหาวาเรปิ สพฺพตฺถ เอกํ.)}

\proseitem{26}{กุสลํ อาสว\-วิปฺปยุตฺตํ สาสวํ ธมฺมํ ปฏิจฺจ กุสโล อาสว\-วิปฺปยุตฺโต สาสโว ธมฺโม อุปฺปชฺชติ เหตุปจฺจยา. ตีณิ.}

\prose{อกุสลํ อาสว\-วิปฺปยุตฺตํ สาสวํ ธมฺมํ ปฏิจฺจ อพฺยากโต อาสว\-วิปฺปยุตฺโต สาสโว ธมฺโม อุปฺปชฺชติ เหตุปจฺจยา. (1)}

\prose{อพฺยากตํ อาสว\-วิปฺปยุตฺตํ สาสวํ ธมฺมํ ปฏิจฺจ อพฺยากโต อาสว\-วิปฺปยุตฺโต สาสโว ธมฺโม อุปฺปชฺชติ เหตุปจฺจยา. (1)}

\prose{กุสลํ อาสว\-วิปฺปยุตฺตํ สาสวญฺจ อพฺยากตํ อาสว\-วิปฺปยุตฺตํ สาสวญฺจ ธมฺมํ ปฏิจฺจ อพฺยากโต อาสว\-วิปฺปยุตฺโต สาสโว ธมฺโม อุปฺปชฺชติ เหตุปจฺจยา. (1)}

\csromanmarkh{1. กุสลตฺติก}\csromanmarkhii{19. อาสว\-วิปฺปยุตฺต\-สาสวทุก}\csromanheadernomark{2}{1. กุสลตฺติก 19. อาสว\-วิปฺปยุตฺต\-สาสวทุก}
\prose{\csromanfolio{585}อกุสลํ อาสว\-วิปฺปยุตฺตํ สาสวญฺจ อพฺยากตํ อาสว\-วิปฺปยุตฺตํ สาสวญฺจ ธมฺมํ ปฏิจฺจ อพฺยากโต อาสว\-วิปฺปยุตฺโต สาสโว ธมฺโม อุปฺปชฺชติ เหตุปจฺจยา. (1) (สํขิตฺตํ.)}

\proseitem{27}{เหตุยา สตฺต, อารมฺมเณ เทฺว ฯเปฯ อวิคเต สตฺต. (สํขิตฺตํ. สหชาตวารมฺปิ ฯเปฯ สมฺปยุตฺตวารมฺปิ ปฏิจฺจ\-วาร\-สทิสํ.)}

\proseitem{28}{กุสโล อาสว\-วิปฺปยุตฺโต สาสโว ธมฺโม กุสลสฺส อาสว\-วิปฺปยุตฺตสฺส สาสวสฺส ธมฺมสฺส เหตุปจฺจเยน ปจฺจโย. ตีณิ.}

\prose{อกุสโล อาสว\-วิปฺปยุตฺโต สาสโว ธมฺโม อพฺยากตสฺส อาสว\-วิปฺปยุตฺตสฺส สาสวสฺส ธมฺมสฺส เหตุปจฺจเยน ปจฺจโย. (1)}

\prose{อพฺยากโต อาสว\-วิปฺปยุตฺโต สาสโว ธมฺโม อพฺยากตสฺส อาสว\-วิปฺปยุตฺตสฺส สาสวสฺส ธมฺมสฺส เหตุปจฺจเยน ปจฺจโย. (สํขิตฺตํ.)}

\proseitem{29}{เหตุยา ปญฺจ, อารมฺมเณ นว, อธิปติยา จตฺตาริ ฯเปฯ อวิคเต ทส. (สํขิตฺตํ. ยถา กุสลตฺติเก ปญฺหาวารํ, เอวํ วิตฺถาเรตพฺพํ.)}

\csromanmarkh{1. กุสลตฺติก}\csromanmarkhii{19. อาสว\-วิปฺปยุตฺต\-สาสวทุก}\csromanheadernomark{2}{1. กุสลตฺติก 19. อาสว\-วิปฺปยุตฺต\-สาสวทุก}
\csromanheader{2}{1-7. ปฏิจฺจาทิวาร ปจฺจยจตุกฺกเหตุ}
\proseitem{30}{กุสลํ อาสว\-วิปฺปยุตฺตํ อนาสวํ ธมฺมํ ปฏิจฺจ กุสโล อาสว\-วิปฺปยุตฺโต อนาสโว ธมฺโม อุปฺปชฺชติ เหตุปจฺจยา. (1)}

\prose{อพฺยากตํ อาสว\-วิปฺปยุตฺตํ อนาสวํ ธมฺมํ ปฏิจฺจ อพฺยากโต อาสว\-วิปฺปยุตฺโต อนาสโว ธมฺโม อุปฺปชฺชติ เหตุปจฺจยา. (สํขิตฺตํ.)}

\proseitem{31}{เหตุยา เทฺว, อารมฺมเณ เทฺว ฯเปฯ อวิคเต เทฺว. (สํขิตฺตํ.)}

\prose{(สหชาตวารมฺปิ ฯเปฯ สมฺปยุตฺตวารมฺปิ ปฏิจฺจ\-วาร\-สทิสํ.)}

\csromanheader{2}{ธมฺมานุโลม ติกทุก\-ปฏฺฐาน\-ปาฬิ เหตุ อารมฺมณาทิ}
\proseitem{32}{\csromanfolio{586}กุสโล อาสว\-วิปฺปยุตฺโต อนาสโว ธมฺโม กุสลสฺส อาสว\-วิปฺปยุตฺตสฺส อนาสวสฺส ธมฺมสฺส เหตุปจฺจเยน ปจฺจโย. (1)}

\prose{อพฺยากโต อาสว\-วิปฺปยุตฺโต อนาสโว ธมฺโม อพฺยากตสฺส อาสว\-วิปฺปยุตฺตสฺส อนาสวสฺส ธมฺมสฺส เหตุปจฺจเยน ปจฺจโย. (1)}

\prose{อพฺยากโต อาสว\-วิปฺปยุตฺโต อนาสโว ธมฺโม อพฺยากตสฺส อาสว\-วิปฺปยุตฺตสฺส อนาสวสฺส ธมฺมสฺส อารมฺมณ\-ปจฺจเยน ปจฺจโย. (1)}

\prose{อพฺยากโต อาสว\-วิปฺปยุตฺโต อนาสโว ธมฺโม กุสลสฺส อาสว\-วิปฺปยุตฺตสฺส อนาสวสฺส ธมฺมสฺส อารมฺมณ\-ปจฺจเยน ปจฺจโย. (1)}

\prose{กุสโล อาสว\-วิปฺปยุตฺโต อนาสโว ธมฺโม กุสลสฺส อาสว\-วิปฺปยุตฺตสฺส อนาสวสฺส ธมฺมสฺส อธิปติ\-ปจฺจเยน ปจฺจโย. (1)}

\prose{อพฺยากโต อาสว\-วิปฺปยุตฺโต อนาสโว ธมฺโม อพฺยากตสฺส อาสว\-วิปฺปยุตฺตสฺส อนาสวสฺส ธมฺมสฺส อธิปติ\-ปจฺจเยน ปจฺจโย. (1)}

\prose{อพฺยากโต อาสว\-วิปฺปยุตฺโต อนาสโว ธมฺโม กุสลสฺส อาสว\-วิปฺปยุตฺตสฺส อนาสวสฺส ธมฺมสฺส อธิปติ\-ปจฺจเยน ปจฺจโย. (1)}

\prose{กุสโล อาสว\-วิปฺปยุตฺโต อนาสโว ธมฺโม อพฺยากตสฺส อาสว\-วิปฺปยุตฺตสฺส อนาสวสฺส ธมฺมสฺส อนนฺตร\-ปจฺจเยน ปจฺจโย. (1)}

\prose{อพฺยากโต อาสว\-วิปฺปยุตฺโต อนาสโว ธมฺโม อพฺยากตสฺส อาสว\-วิปฺปยุตฺตสฺส อนาสวสฺส ธมฺมสฺส อนนฺตร\-ปจฺจเยน ปจฺจโย. (1) (สํขิตฺตํ.)}

\proseitem{33}{เหตุยา เทฺว, อารมฺมเณ เทฺว, อธิปติยา ตีณิ, อนนฺตเร เทฺว, สมนนฺตเร เทฺว, สหชาเต เทฺว, อญฺญมญฺเญ เทฺว, นิสฺสเย เทฺว, อุปนิสฺสเย จตฺตาริ ฯเปฯ อวิคเต เทฺว. (สํขิตฺตํ.)}

\prosewithcloser{(ยถา กุสลตฺติเก ปญฺหาวารํ, เอวํ วิตฺถาเรตพฺพํ.)}{กุสล\-ตฺติก\-อาสว\-โคจฺฉกํ นิฏฺฐิตํ.}
\csromancenter{กุสลตฺติก 55. สารมฺมณทุก 1. กุสลตฺติก 20-54. สญฺโญชนา\-ทิทุก}
\csromanmatikaanchor{mtk.144}
\csromantocmarkref{subparagraph}{20-54. สญฺโญชนาทิทุก}{mtk.144}
\bookmark[dest={mtk.144},level=5]{20-54. สญฺโญชนาทิทุก}
\proseitem{34}{\csromanfolio{587}อกุสลํ สญฺโญชนํ ธมฺมํ ปฏิจฺจ ฯเปฯ. อกุสลํ คนฺถํ. โอฆํ. โยคํ. นีวรณํ.}

\prosewithcloser{กุสลํ โนปรามาสํ ธมฺมํ ปฏิจฺจ กุสโล โนปรามาโส ธมฺโม อุปฺปชฺชติ เหตุปจฺจยา. (อาสว\-โคจฺฉก\-สทิสํ.) ทิฏฺฐิํ ปฏิจฺจ ทิฏฺฐิ น อุปฺปชฺชติ.}{กุสลตฺติเก ฉ โคจฺฉกํ นิฏฺฐิตํ.}
\csromanmatikaanchor{mtk.145}
\csromantocmarkref{subparagraph}{55. สารมฺมณทุก}{mtk.145}
\bookmark[dest={mtk.145},level=5]{55. สารมฺมณทุก}
\csromanmarkh{1. กุสลตฺติก}\csromanmarkhii{55. สารมฺมณทุก}\csromanheadernomark{6}{1. กุสลตฺติก 55. สารมฺมณทุก}
\csromanheader{2}{1-7. ปฏิจฺจาทิวาร ปจฺจยจตุกฺก}
\proseitem{1}{กุสลํ สารมฺมณํ ธมฺมํ ปฏิจฺจ กุสโล สารมฺมโณ ธมฺโม อุปฺปชฺชติ เหตุปจฺจยา. (1)}

\prose{อกุสลํ สารมฺมณํ ธมฺมํ ปฏิจฺจ อกุสโล สารมฺมโณ ธมฺโม อุปฺปชฺชติ เหตุปจฺจยา. (1)}

\prose{อพฺยากตํ สารมฺมณํ ธมฺมํ ปฏิจฺจ อพฺยากโต สารมฺมโณ ธมฺโม อุปฺปชฺชติ เหตุปจฺจยา. (1) (สํขิตฺตํ.)}

\proseitem{2}{เหตุยา ตีณิ, อารมฺมเณ ตีณิ ฯเปฯ อวิคเต ตีณิ. (สํขิตฺตํ. สหชาตวารมฺปิ ฯเปฯ สมฺปยุตฺตวารมฺปิ ปฏิจฺจ\-วาร\-สทิสํ.)}

\prose{เหตุอารมฺมณ}

\proseitem{3}{กุสโล สารมฺมโณ ธมฺโม กุสลสฺส สารมฺมณสฺส ธมฺมสฺส เหตุปจฺจเยน ปจฺจโย. (1)}

\prose{อกุสโล สารมฺมโณ ธมฺโม อกุสลสฺส สารมฺมณสฺส ธมฺมสฺส เหตุปจฺจเยน ปจฺจโย. (1)}

\prose{อพฺยากโต สารมฺมโณ ธมฺโม อพฺยากตสฺส สารมฺมณสฺส ธมฺมสฺส เหตุปจฺจเยน ปจฺจโย. (1)}

\prose{\csromanfolio{588}กุสโล สารมฺมโณ ธมฺโม กุสลสฺส สารมฺมณสฺส ธมฺมสฺส อารมฺมณ\-ปจฺจเยน ปจฺจโย. (สํขิตฺตํ.)}

\proseitem{4}{เหตุยา ตีณิ, อารมฺมเณ นว, อธิปติยา สตฺต ฯเปฯ อวิคเต ตีณิ. (สํขิตฺตํ.)}

\prose{(ยถา กุสลตฺติเก ปญฺหาวารํ, เอวํ วิตฺถาเรตพฺพํ.)}

\csromanheader{2}{\textbf{เหตุสหชาตาทิ}}
\proseitem{5}{อพฺยากตํ อนารมฺมณํ ธมฺมํ ปฏิจฺจ อพฺยากโต อนารมฺมโณ ธมฺโม อุปฺปชฺชติ เหตุปจฺจยา. (สํขิตฺตํ.)}

\prose{เหตุยา เอกํ ฯเปฯ อวิคเต เอกํ. (สํขิตฺตํ.)}

\prose{นเหตุยา เอกํ. (สํขิตฺตํ.)}

\prose{(สหชาตวารมฺปิ ฯเปฯ สมฺปยุตฺตวารมฺปิ ปฏิจฺจ\-วาร\-สทิสํ.)}

\proseitem{6}{อพฺยากโต อนารมฺมโณ ธมฺโม อพฺยากตสฺส อนารมฺมณสฺส ธมฺมสฺส สหชาต\-ปจฺจเยน ปจฺจโย. อญฺญมญฺญ\-ปจฺจเยน ปจฺจโย. นิสฺสย\-ปจฺจเยน ปจฺจโย. อาหารปจฺจเยน ปจฺจโย. อินฺทฺริย\-ปจฺจเยน ปจฺจโย. อตฺถิปจฺจเยน ปจฺจโย. อวิคต\-ปจฺจเยน ปจฺจโย. (สพฺพตฺถ เอกํ. สํขิตฺตํ.)}

\csromanmarkh{1. กุสลตฺติก}\csromanmarkhii{56. จิตฺตทุก}\csromanheadernomark{2}{1. \textbf{กุสลตฺติก} 56. จิตฺตทุก}
\csromanheader{2}{1-7. ปฏิจฺจาทิวาร ปจฺจยจตุกฺก}
\proseitem{7}{กุสโล จิตฺโต ธมฺโม กุสลสฺส จิตฺตสฺส ธมฺมสฺส อารมฺมณ\-ปจฺจเยน ปจฺจโย. (สํขิตฺตํ.)}

\prose{อารมฺมเณ นว, อธิปติยา สตฺต. (สํขิตฺตํ.)}

\prose{นเหตุยา นว, นอารมฺมเณ นว. (สํขิตฺตํ.)}

\prose{(ยถา กุสลตฺติเก ปญฺหาวารํ, เอวํ วิตฺถาเรตพฺพํ.)}

\csromanmarkh{1. กุสลตฺติก}\csromanmarkhii{57. เจตสิกทุก}\csromanheadernomark{2}{1. กุสลตฺติก 57. เจตสิกทุก}
\proseitem{8}{\csromanfolio{589}กุสลํ โนจิตฺตํ ธมฺมํ ปฏิจฺจ กุสโล โนจิตฺโต ธมฺโม อุปฺปชฺชติ เหตุปจฺจยา. ตีณิ.}

\prose{อกุสลํ โนจิตฺตํ ธมฺมํ ปฏิจฺจ อกุสโล โนจิตฺโต ธมฺโม อุปฺปชฺชติ เหตุปจฺจยา. ตีณิ.}

\prose{อพฺยากตํ โนจิตฺตํ ธมฺมํ ปฏิจฺจ อพฺยากโต โนจิตฺโต ธมฺโม อุปฺปชฺชติ เหตุปจฺจยา. (1)}

\prose{กุสลํ โนจิตฺตญฺจ อพฺยากตํ โนจิตฺตญฺจ ธมฺมํ ปฏิจฺจ อพฺยากโต โนจิตฺโต ธมฺโม อุปฺปชฺชติ เหตุปจฺจยา. (1)}

\prose{อกุสลํ โนจิตฺตญฺจ อพฺยากตํ โนจิตฺตญฺจ ธมฺมํ ปฏิจฺจ อพฺยากโต โนจิตฺโต ธมฺโม อุปฺปชฺชติ เหตุปจฺจยา. (1)}

\prose{กุสลํ โนจิตฺตํ ธมฺมํ ปฏิจฺจ กุสโล โนจิตฺโต ธมฺโม อุปฺปชฺชติ อารมฺมณ\-ปจฺจยา. (สํขิตฺตํ.)}

\proseitem{9}{เหตุยา นว, อารมฺมเณ ตีณิ, อธิปติยา นว ฯเปฯ อวิคเต นว. (สํขิตฺตํ. สหชตวารมฺปิ ฯเปฯ สมฺปยุตฺตวารมฺปิ ปฏิจฺจ\-วาร\-สทิสํ.)}

\proseitem{10}{กุสโล โนจิตฺโต ธมฺโม กุสลสฺส โนจิตฺตสฺส ธมฺมสฺส เหตุปจฺจเยน ปจฺจโย. (สํขิตฺตํ.)}

\prose{เหตุยา สตฺต, อารมฺมเณ นว, อธิปติยา ทส ฯเปฯ อวิคเต เตรส. (สํขิตฺตํ.)}

\csromanmarkh{1. กุสลตฺติก}\csromanmarkhii{57. เจตสิกทุก}\csromanheadernomark{2}{1. กุสลตฺติก 57. เจตสิกทุก}
\csromanheader{2}{1-7. ปฏิจฺจาทิวาร ปจฺจยจตุกฺก}
\proseitem{11}{กุสลํ เจตสิกํ ธมฺมํ ปฏิจฺจ กุสโล เจตสิโก ธมฺโม อุปฺปชฺชติ เหตุปจฺจยา. (1)}

\prose{อกุสลํ เจตสิกํ ธมฺมํ ปฏิจฺจ อกุสโล เจตสิโก ธมฺโม อุปฺปชฺชติ เหตุปจฺจยา. (1)}

\prose{\csromanfolio{590}อพฺยากตํ เจตสิกํ ธมฺมํ ปฏิจฺจ อพฺยากโต เจตสิโก ธมฺโม อุปฺปชฺชติ เหตุปจฺจยา. (1) (สํขิตฺตํ.)}

\proseitem{12}{เหตุยา ตีณิ, อารมฺมเณ ตีณิ ฯเปฯ. (สํขิตฺตํ. สหชตวารมฺปิ ฯเปฯ สมฺปยุตฺตวารมฺปิ ปฏิจฺจ\-วาร\-สทิสํ.)}

\csromanheader{2}{\textbf{เหตุ อารมฺมณ}}
\proseitem{13}{กุสโล เจตสิโก ธมฺโม กุสลสฺส เจตสิกสฺส ธมฺมสฺส เหตุปจฺจเยน ปจฺจโย. (1)}

\prose{อกุสโล เจตสิโก ธมฺโม อกุสลสฺส เจตสิกสฺส ธมฺมสฺส เหตุปจฺจเยน ปจฺจโย. (1)}

\prose{อพฺยากโต เจตสิโก ธมฺโม อพฺยากตสฺส เจตสิกสฺส ธมฺมสฺส เหตุปจฺจเยน ปจฺจโย. (1)}

\prose{กุสโล เจตสิโก ธมฺโม กุสลสฺส เจตสิกสฺส ธมฺมสฺส อารมฺมณ\-ปจฺจเยน ปจฺจโย. ตีณิ.}

\prose{อกุสโล เจสิตโก ธมฺโม อกุสลสฺส เจตสิกสฺส ธมฺมสฺส อารมฺมณ\-ปจฺจเยน ปจฺจโย. ตีณิ.}

\prose{อพฺยากโต เจตสิโก ธมฺโม อพฺยากตสฺส เจตสิกสฺส ธมฺมสฺส อารมฺมณ\-ปจฺจเยน ปจฺจโย. ตีณิ. (สํขิตฺตํ.)}

\proseitem{14}{เหตุยา ตีณิ, อารมฺมเณ นว, อธิปติยา สตฺต ฯเปฯ อวิคเต ตีณิ. (สํขิตฺตํ.)}

\prose{(ยถา กุสลตฺติเก ปญฺหาวารํ, เอวํ วิตฺถาเรตพฺพํ.)}

\proseitem{15}{กุสลํ อเจตสิกํ ธมฺมํ ปฏิจฺจ อพฺยากโต อเจตสิโก ธมฺโม อุปฺปชฺชติ เหตุปจฺจยา. (1)}

\prose{อกุสลํ อเจตสิกํ ธมฺมํ ปฏิจฺจ อพฺยากโต อเจตสิโก ธมฺโม อุปฺปชฺชติ เหตุปจฺจยา. (1)}

\prose{อพฺยากตํ อเจตสิกํ ธมฺมํ ปฏิจฺจ อพฺยากโต อเจตสิโก ธมฺโม อุปฺปชฺชติ เหตุปจฺจยา. (1)}

\csromanmarkh{1. กุสลตฺติก}\csromanmarkhii{66. อชฺฌตฺติกทุก}\csromanheadernomark{2}{1. กุสลตฺติก 66. อชฺฌตฺติกทุก}
\prose{\csromanfolio{591}กุสลํ อเจตสิกญฺจ อพฺยากตํ อเจตสิกญฺจ ธมฺมํ ปฏิจฺจ อพฺยากโต อเจตสิโก ธมฺโม อุปฺปชฺชติ เหตุปจฺจยา. (1)}

\prose{อกุสลํ อเจตสิกญฺจ อพฺยากตํ อเจตสิกญฺจ ธมฺมํ ปฏิจฺจ อพฺยากโต อเจตสิโก ธมฺโม อุปฺปชฺชติ เหตุปจฺจยา. (1) (สํขิตฺตํ.)}

\proseitem{16}{เหตุยา ปญฺจ, อารมฺมเณ เอกํ ฯเปฯ อวิคเต ปญฺจ.}

\prose{นเหตุยา เอกํ, นอารมฺมเณ ปญฺจ, โนวิคเต ปญฺจ.}

\prose{(สหชาตวารมฺปิ ฯเปฯ ปญฺหาวารมฺปิ เอวํ วิตฺถาเรตพฺพํ.)}

\csromanmatikaanchor{mtk.146}
\csromantocmarkref{subparagraph}{58-65. จิตฺตสมฺปยุตฺตาทิทุก}{mtk.146}
\bookmark[dest={mtk.146},level=5]{58-65. จิตฺตสมฺปยุตฺตาทิทุก}
\csromanheader{6}{1. \textbf{กุสลตฺติก 58-}65. จิตฺตสมฺปยุตฺตาทิทุก}
\csromanheader{2}{1-7. ปฏิจฺจาทิวาร}
\proseitem{17}{กุสลํ จิตฺต\-สมฺปยุตฺตํ ธมฺมํ ปฏิจฺจ ฯเปฯ. กุสลํ จิตฺต\-สํสฏฺฐํ ธมฺมํ ปฏิจฺจ ฯเปฯ. กุสลํ จิตฺต\-สมุฏฺฐานํ ธมฺมํ ปฏิจฺจ ฯเปฯ. กุสลํ จิตฺตสหภุํ ธมฺมํ ปฏิจฺจ ฯเปฯ. กุสลํ จิตฺตา\-นุปริวตฺติํ ธมฺมํ ปฏิจฺจ ฯเปฯ. กุสลํ จิตฺต\-สํสฏฺฐ\-สมุฏฺฐานํ ธมฺมํ ปฏิจฺจ ฯเปฯ. กุสลํ จิตฺต\-สํสฏฺฐ\-สมุฏฺฐาน\-สหภุํ ธมฺมํ ปฏิจฺจ ฯเปฯ. กุสลํ จิตฺต\-สํสฏฺฐ\-สมุฏฺฐานา\-นุปริวตฺติํ ธมฺมํ ปฏิจฺจ กุสโล จิตฺต\-สํสฏฺฐ\-สมุฏฺฐานา\-นุปริวตฺติ ธมฺโม อุปฺปชฺชติ เหตุปจฺจยา. (สํขิตฺตํ.)}

\prose{เหตุยา ตีณิ, อารมฺมเณ ตีณิ ฯเปฯ อวิคเต ตีณิ. (สํขิตฺตํ. สหชาตวารมฺปิ ฯเปฯ สมฺปยุตฺตวารมฺปิ ปญฺหาวารมฺปิ เอวํ วิตฺถาเรตพฺพํ.)}

\csromanmarkh{1. กุสลตฺติก}\csromanmarkhii{66. อชฺฌตฺติกทุก}\csromanheadernomark{2}{1. กุสลตฺติก 66. อชฺฌตฺติกทุก}
\csromanheader{2}{1-7. ปฏิจฺจาทิวาร ปจฺจยจตุกฺก}
\proseitem{18}{อพฺยากตํ อชฺฌตฺติกํ ธมฺมํ ปฏิจฺจ อพฺยากโต อชฺฌตฺติโก ธมฺโม อุปฺปชฺชติ เหตุปจฺจยา. (สํขิตฺตํ.)}

\prose{เหตุยา เอกํ, สหชาเต เอกํ ฯเปฯ อวิคเต เอกํ. (สํขิตฺตํ.)}

\prose{(สหชาตวารมฺปิ ฯเปฯ สมฺปยุตฺตวารมฺปิ ปฏิจฺจ\-วาร\-สทิสํ.)}

\proseitem{19}{\csromanfolio{592}กุสโล อชฺฌตฺติโก ธมฺโม กุสลสฺส อชฺฌตฺติกสฺส ธมฺมสฺส อารมฺมณ\-ปจฺจเยน ปจฺจโย.  กุสโล อชฺฌตฺติโก ธมฺโม อกุสลสฺส อชฺฌตฺติกสฺส ธมฺมสฺส อารมฺมณ\-ปจฺจเยน ปจฺจโย.  กุสโล อชฺฌตฺติโก ธมฺโม อพฺยากตสฺส อชฺฌตฺติกสฺส ธมฺมสฺส อารมฺมณ\-ปจฺจเยน ปจฺจโย. (3)}

\prose{อกุสโล อชฺฌตฺติโก ธมฺโม อกุสลสฺส อชฺฌตฺติกสฺส ธมฺมสฺส อารมฺมณ\-ปจฺจเยน ปจฺจโย. ตีณิ.}

\prose{อพฺยากโต อชฺฌตฺติโก ธมฺโม อพฺยากตสฺส อชฺฌตฺติกสฺส ธมฺมสฺส อารมฺมณ\-ปจฺจเยน ปจฺจโย. ตีณิ. (สํขิตฺตํ.)}

\proseitem{20}{อารมฺมเณ นว ฯเปฯ ปุเรชาเต ตีณิ ฯเปฯ อวิคเต ปญฺจ. (สํขิตฺตํ. ยถา กุสลตฺติเก ปญฺหาวารํ, เอวํ วิตฺถาเรตพฺพํ.)}

\proseitem{21}{กุสลํ พาหิรํ ธมฺมํ ปฏิจฺจ กุสโล พาหิโร ธมฺโม อุปฺปชฺชติ เหตุปจฺจยา. (สํขิตฺตํ.)}

\prose{เหตุยา นว, อารมฺมเณ ตีณิ ฯเปฯ วิปาเก เอกํ ฯเปฯ อวิคเต นว. (สํขิตฺตํ. สหชาตวารมฺปิ ฯเปฯ สมฺปยุตฺตวารมฺปิ ปฏิจฺจ\-วาร\-สทิสํ.)}

\proseitem{22}{กุสโล พาหิโร ธมฺโม กุสลสฺส พาหิรสฺส ธมฺมสฺส เหตุปจฺจเยน ปจฺจโย. (สํขิตฺตํ.)}

\csromanhangingbegin
\hangingheaditem{22}{เหตุยา สตฺต, อารมฺมเณ นว ฯเปฯ อวิคเต เตรส. (สํขิตฺตํ.)}
\hangingbody{(ยถา กุสลตฺติเก ปญฺหาวารํ, เอวํ วิตฺถาเรตพฺพํ.)}
\csromanhangingend

\csromanmarkh{1. กุสลตฺติก}\csromanmarkhii{67. อุปาทาทุก}\csromanheadernomark{2}{1. \textbf{กุสลตฺติก} 67. อุปาทาทุก}
\csromanheader{2}{1-7. ปฏิจฺจาทิวาร ปจฺจยจตุกฺก}
\proseitem{23}{อพฺยากโต อุปาทา ธมฺโม อพฺยากตสฺส อุปาทา ธมฺมสฺส อาหารปจฺจเยน ปจฺจโย. อินฺทฺริย\-ปจฺจเยน ปจฺจโย. อตฺถิปจฺจเยน ปจฺจโย. อวิคต\-ปจฺจเยน ปจฺจโย.}

\prose{อาหาเร เอกํ, อินฺทฺริเย อตฺถิยา อวิคเต เอกํ. (สํขิตฺตํ.)}

\csromanmarkh{1. กุสลตฺติก}\csromanmarkhii{67. อุปาทาทุก}\csromanheadernomark{2}{1. กุสลตฺติก 67. อุปาทาทุก}
\proseitem{24}{\csromanfolio{593}กุสลํ โนอุปาทา ธมฺมํ ปฏิจฺจ กุสโล โนอุปาทา ธมฺโม อุปฺปชฺชติ เหตุปจฺจยา. ตีณิ.}

\prose{อกุสลํ โนอุปาทา ธมฺมํ ปฏิจฺจ อกุสโล โนอุปาทา ธมฺโม อุปฺปชฺชติ เหตุปจฺจยา. ตีณิ.}

\prose{อพฺยากตํ โนอุปาทา ธมฺมํ ปฏิจฺจ อพฺยากโต โนอุปาทา ธมฺโม อุปฺปชฺชติ เหตุปจฺจยา. (1)}

\prose{กุสลํ โนอุปาทา จ อพฺยากตํ โนอุปาทา จ ธมฺมํ ปฏิจฺจ อพฺยากโต โนอุปาทา ธมฺโม อุปฺปชฺชติ เหตุปจฺจยา. (1)}

\prose{อกุสลํ โนอุปาทา จ อพฺยากตํ โนอุปาทา จ ธมฺมํ ปฏิจฺจ อพฺยากโต โนอุปาทา ธมฺโม อุปฺปชฺชติ เหตุปจฺจยา. (1) (สํขิตฺตํ.)}

\proseitem{25}{เหตุยา นว. (สํขิตฺตํ. สหชาตวารมฺปิ ฯเปฯ สมฺปยุตฺตวารมฺปิ ปฏิจฺจ\-วาร\-สทิสํ.)}

\proseitem{26}{กุสโล โนอุปาทา ธมฺโม กุสลสฺส โนอุปาทา ธมฺมสฺส เหตุปจฺจเยน ปจฺจโย. ตีณิ.}

\prose{อกุสโล โนอุปาทา ธมฺโม อกุสลสฺส โนอุปาทา ธมฺมสฺส เหตุปจฺจเยน ปจฺจโย. ตีณิ.}

\prose{อพฺยากโต โนอุปาทา ธมฺโม อพฺยากตสฺส โนอุปาทา ธมฺมสฺส เหตุปจฺจเยน ปจฺจโย. (1)}

\proseitem{27}{กุสโล โนอุปาทา ธมฺโม กุสลสฺส โนอุปาทา ธมฺมสฺส อารมฺมณ\-ปจฺจเยน ปจฺจโย. ตีณิ.}

\prose{อกุสโล โนอุปาทา ธมฺโม อกุสลสฺส โนอุปาทา ธมฺมสฺส อารมฺมณ\-ปจฺจเยน ปจฺจโย. ตีณิ.}

\prose{อพฺยากโต โนอุปาทา ธมฺโม อพฺยากตสฺส โนอุปาทา ธมฺมสฺส อารมฺมณ\-ปจฺจเยน ปจฺจโย. ตีณิ. (สํขิตฺตํ.)}

\proseitem{28}{เหตุยา สตฺต, อารมฺมเณ นว, อธิปติยา ทส ฯเปฯ อวิคเต เอกาทส. (สํขิตฺตํ.)}

\prose{(ยถา กุสลตฺติเก ปญฺหาวารํ, เอวํ วิตฺถาเรตพฺพํ.)}

\csromanmarkh{ธมฺมานุโลม ติกทุก\-ปฏฺฐาน\-ปาฬิ}\csromanmarkhii{1. กุสลตฺติก 68. อุปาทินฺนทุก}\csromanheadernomark{2}{ธมฺมานุโลม ติกทุก\-ปฏฺฐาน\-ปาฬิ 1. กุสลตฺติก 68. อุปาทินฺนทุก}
\csromanheader{2}{1-7. ปฏิจฺจาทิวาร ปจฺจยจตุกฺก}
\proseitem{29}{\csromanfolio{594}อพฺยากตํ อุปาทินฺนํ ธมฺมํ ปฏิจฺจ อพฺยากโต อุปาทินฺโน ธมฺโม อุปฺปชฺชติ เหตุปจฺจยา. (1) (สํขิตฺตํ.)}

\prose{เหตุยา เอกํ, อารมฺมเณ เอกํ ฯเปฯ อวิคเต เอกํ. (สํขิตฺตํ.) (สหชาตวาเรปิ ฯเปฯ ปญฺหาวาเรปิ สพฺพตฺถ เอกํ.)}

\proseitem{30}{กุสลํ อนุปาทินฺนํ ธมฺมํ ปฏิจฺจ กุสโล อนุปาทินฺโน ธมฺโม อุปฺปชฺชติ เหตุปจฺจยา. ตีณิ.}

\prose{อกุสลํ อนุปาทินฺนํ ธมฺมํ ปฏิจฺจ อกุสโล อนุปาทินฺโน ธมฺโม อุปฺปชฺชติ เหตุปจฺจยา. ตีณิ.}

\prose{อพฺยากตํ อนุปาทินฺนํ ธมฺมํ ปฏิจฺจ อพฺยากโต อนุปาทินฺโน ธมฺโม อุปฺปชฺชติ เหตุปจฺจยา. (1)}

\prose{กุสลํ อนุปาทินฺนญฺจ อพฺยากตํ อนุปาทินฺนญฺจ ธมฺมํ ปฏิจฺจ อพฺยากโต อนุปาทินฺโน ธมฺโม อุปฺปชฺชติ เหตุปจฺจยา. (1)}

\prose{อกุสลํ อนุปาทินฺนญฺจ อพฺยากตํ อนุปาทินฺนญฺจ ธมฺมํ ปฏิจฺจ อพฺยากโต อนุปาทินฺโน ธมฺโม อุปฺปชฺชติ เหตุปจฺจยา. (1) (สํขิตฺตํ.)}

\proseitem{31}{เหตุยา นว, อารมฺมเณ ตีณิ ฯเปฯ อวิคเต นว. (สํขิตฺตํ. สหชาตวารมฺปิ ฯเปฯ สมฺปยุตฺตวารมฺปิ ปฏิจฺจ\-วาร\-สทิสํ.)}

\csromanheader{2}{\textbf{เหตุ อารมฺมณ}}
\proseitem{32}{กุสโล อนุปาทินฺโน ธมฺโม กุสลสฺส อนุปาทินฺนสฺส ธมฺมสฺส เหตุปจฺจเยน ปจฺจโย. ตีณิ.}

\prose{อกุสโล อนุปาทินฺโน ธมฺโม อกุสลสฺส อนุปาทินฺนสฺส ธมฺมสฺส เหตุปจฺจเยน ปจฺจโย. ตีณิ.}

\prose{อพฺยากโต อนุปาทินฺโน ธมฺโม อพฺยากตสฺส อนุปาทินฺนสฺส ธมฺมสฺส เหตุปจฺจเยน ปจฺจโย. (1)}

\csromanmarkh{1. กุสลตฺติก}\csromanmarkhii{75. กิเลสทุก}\csromanheadernomark{2}{1. กุสลตฺติก 75. กิเลสทุก}
\csromanmatikaanchor{mtk.147}
\csromanmatikaanchor{mtk.148}
\csromantocmarkref{subparagraph}{69-74. อุปาทานทุก}{mtk.147}
\bookmark[dest={mtk.147},level=5]{69-74. อุปาทานทุก}
\csromantocmarkref{subparagraph}{75-82. กิเลสทุกาทิ}{mtk.148}
\bookmark[dest={mtk.148},level=5]{75-82. กิเลสทุกาทิ}
\proseitem{33}{\csromanfolio{595}กุสโล อนุปาทินฺโน ธมฺโม กุสลสฺส อนุปาทินฺนสฺส ธมฺมสฺส อารมฺมณ\-ปจฺจเยน ปจฺจโย. ตีณิ.}

\prose{อกุสโล อนุปาทินฺโน ธมฺโม อกุสลสฺส อนุปาทินฺนสฺส ธมฺมสฺส อารมฺมณ\-ปจฺจเยน ปจฺจโย. ตีณิ.}

\prose{อพฺยากโต อนุปาทินฺโน ธมฺโม อพฺยากตสฺส อนุปาทินฺนสฺส ธมฺมสฺส อารมฺมณ\-ปจฺจเยน ปจฺจโย.  อพฺยากโต อนุปาทินฺโน ธมฺโม กุสลสฺส อนุปาทินฺนสฺส ธมฺมสฺส อารมฺมณ\-ปจฺจเยน ปจฺจโย.  อพฺยากโต อนุปาทินฺโน ธมฺโม อกุสลสฺส อนุปาทินฺนสฺส ธมฺมสฺส อารมฺมณ\-ปจฺจเยน ปจฺจโย. ตีณิ. (สํขิตฺตํ.)}

\proseitemwithcloser{34}{เหตุยา สตฺต, อารมฺมเณ นว, อธิปติยา ทส, อนนฺตเร ฉ ฯเปฯ สหชาเต นว, นิสฺสเย นว, อุปนิสฺสเย นว ฯเปฯ อวิคเต เอกาทส. (สํขิตฺตํ.)}{กุสล\-ตฺติก\-มหนฺตรทุกํ นิฏฺฐิตํ.}
\csromanheader{6}{1. \textbf{กุสลตฺติก} 69-74. อุปาทานทุก}
\proseitem{35}{อกุสลํ อุปาทานํ ธมฺมํ ปฏิจฺจ อกุสโล อุปาทาโน ธมฺโม อุปฺปชฺชติ เหตุปจฺจยา. (สํขิตฺตํ.)}

\prosewithcloser{เหตุยา เอกํ, อารมฺมเณ เอกํ ฯเปฯ อวิคเต เอกํ. (อุปาทาน\-โคจฺฉกํ วิตฺถาเรตพฺพํ.)}{กุสลตฺติกอุปาทานโคจฺฉกํ นิฏฺฐิตํ.}
\csromanmarkh{1. กุสลตฺติก}\csromanmarkhii{75. กิเลสทุก}\csromanheadernomark{2}{1. กุสลตฺติก 75. กิเลสทุก}
\csromanheader{2}{1-7. ปฏิจฺจาทิวาร ปจฺจยจตุกฺก}
\proseitem{1}{อกุสลํ กิเลสํ ธมฺมํ ปฏิจฺจ อกุสโล กิเลโส ธมฺโม อุปฺปชฺชติ เหตุปจฺจยา. (สํขิตฺตํ.)}

\prose{เหตุยา เอกํ, อารมฺมเณ เอกํ ฯเปฯ อวิคเต เอกํ. (สํขิตฺตํ.)}

\vspace{\tipitakaparskip}
\csromancenter{(สหชาตวาเรปิ ฯเปฯ ปญฺหาวาเรปิ สพฺพตฺถ เอกํ.)}
\proseitem{2}{\csromanfolio{596}กุสลํ โนกิเลสํ ธมฺมํ ปฏิจฺจ กุสโล โนกิเลโส ธมฺโม อุปฺปชฺชติ เหตุปจฺจยา. ตีณิ.}

\prose{อกุสลํ โนกิเลสํ ธมฺมํ ปฏิจฺจ อกุสโล โนกิเลโส ธมฺโม อุปฺปชฺชติ เหตุปจฺจยา. ตีณิ.}

\prose{อพฺยากตํ โนกิเลสํ ธมฺมํ ปฏิจฺจ อพฺยากโต โนกิเลโส ธมฺโม อุปฺปชฺชติ เหตุปจฺจยา. (1)}

\prose{กุสลํ โนกิเลสญฺจ อพฺยากตํ โนกิเลสญฺจ ธมฺมํ ปฏิจฺจ อพฺยากโต โนกิเลโส ธมฺโม อุปฺปชฺชติ เหตุปจฺจยา. (1)}

\prose{อกุสลํ โนกิเลสญฺจ อพฺยากตํ โนกิเลสญฺจ ธมฺมํ ปฏิจฺจ อพฺยากโต โนกิเลโส ธมฺโม อุปฺปชฺชติ เหตุปจฺจยา. (1) (สํขิตฺตํ.)}

\proseitem{3}{เหตุยา นว, อารมฺมเณ ตีณิ, อธิปติยา นว ฯเปฯ อวิคเต นว. (สํขิตฺตํ. สหชาตวารมฺปิ ฯเปฯ สมฺปยุตฺตวารมฺปิ ปฏิจฺจ\-วาร\-สทิสํ.)}

\proseitem{4}{กุสโล โนกิเลโส ธมฺโม กุสลสฺส โนกิเลสสฺส ธมฺมสฺส เหตุปจฺจเยน ปจฺจโย. ตีณิ.}

\prose{อพฺยากโต โนกิเลโส ธมฺโม อพฺยากตสฺส โนกิเลสสฺส ธมฺมสฺส เหตุปจฺจเยน ปจฺจโย. (1) (สํขิตฺตํ.)}

\proseitem{5}{เหตุยา จตฺตาริ, อารมฺมเณ นว ฯเปฯ อวิคเต เตรส. (สํขิตฺตํ. ยถา กุสลตฺติเก ปญฺหาวารํ, เอวํ วิตฺถาเรตพฺพํ.)}

\csromanmarkh{1. กุสลตฺติก}\csromanmarkhii{76. สํกิเลสิกทุก}\csromanheadernomark{2}{1. \textbf{กุสลตฺติก} 76. สํกิเลสิกทุก}
\csromanheader{2}{1-7. ปฏิจฺจาทิวาร ปจฺจยจตุกฺก}
\proseitem{6}{กุสลํ สํกิเลสิกํ ธมฺมํ ปฏิจฺจ กุสโล สํกิเลสิโก ธมฺโม อุปฺปชฺชติ เหตุปจฺจยา. ตีณิ.}

\prose{อกุสลํ สํกิเลสิกํ ธมฺมํ ปฏิจฺจ อกุสโล สํกิเลสิโก ธมฺโม อุปฺปชฺชติ เหตุปจฺจยา. ตีณิ.}

\csromanmarkh{1. กุสลตฺติก}\csromanmarkhii{76. สํกิเลสิกทุก}\csromanheadernomark{2}{1. กุสลตฺติก 76. สํกิเลสิกทุก}
\prose{\csromanfolio{597}อพฺยากตํ สํกิเลสิกํ ธมฺมํ ปฏิจฺจ อพฺยากโต สํกิเลสิโก ธมฺโม อุปฺปชฺชติ เหตุปจฺจยา. (1)}

\prose{กุสลํ สํกิเลสิกญฺจ อพฺยากตํ สํกิเลสิกญฺจ ธมฺมํ ปฏิจฺจ อพฺยากโต สํกิเลสิโก ธมฺโม อุปฺปชฺชติ เหตุปจฺจยา. (1)}

\prose{อกุสลํ สํกิเลสิกญฺจ อพฺยากตํ สํกิเลสิกญฺจ ธมฺมํ ปฏิจฺจ อพฺยากโต สํกิเลสิโก ธมฺโม อุปฺปชฺชติ เหตุปจฺจยา. (1)}

\proseitem{7}{เหตุยา นว, อารมฺมเณ ตีณิ ฯเปฯ อวิคเต นว. (สํขิตฺตํ. สหชาตวารมฺปิ ฯเปฯ สมฺปยุตฺตวารมฺปิ วิตฺถาเรตพฺพํ.)}

\proseitem{8}{กุสโล สํกิเลสิโก ธมฺโม อกุสลสฺส สํกิเลสิกสฺส ธมฺมสฺส เหตุปจฺจเยน ปจฺจโย. ตีณิ.}

\prose{อกุสโล สํกิเลสิโก ธมฺโม อกุสลสฺส สํกิเลสิกสฺส ธมฺมสฺส เหตุปจฺจเยน ปจฺจโย. ตีณิ.}

\prose{อพฺยากโต สํกิเลสิโก ธมฺโม อพฺยากตสฺส สํกิเลสิกสฺส ธมฺมสฺส เหตุปจฺจเยน ปจฺจโย. (1) (สํขิตฺตํ.)}

\proseitem{9}{เหตุยา สตฺต, อารมฺมเณ นว, อธิปติยา นว ฯเปฯ อวิคเต เตรส. (สํขิตฺตํ. ยถา กุสลตฺติเก ปญฺหาวารํ, เอวํ วิตฺถาเรตพฺพํ.)}

\proseitem{10}{กุสลํ อสํกิเลสิกํ ธมฺมํ ปฏิจฺจ กุสโล อสํกิเลสิโก ธมฺโม อุปฺปชฺชติ เหตุปจฺจยา. (1)}

\prose{อพฺยากตํ อสํกิเลสิกํ ธมฺมํ ปฏิจฺจ อพฺยากโต อสํกิเลสิโก ธมฺโม อุปฺปชฺชติ เหตุปจฺจยา. (1) (สํขิตฺตํ.)}

\proseitem{11}{เหตุยา เทฺว, อารมฺมเณ เทฺว ฯเปฯ อวิคเต เทฺว. (สํขิตฺตํ. สหชาตวารมฺปิ ฯเปฯ สมฺปยุตฺตวารมฺปิ ปฏิจฺจ\-วาร\-สทิสํ.)}

\proseitem{12}{กุสโล อสํกิเลสิโก ธมฺโม กุสลสฺส อสํกิเลสิกสฺส ธมฺมสฺส เหตุปจฺจเยน ปจฺจโย. (1)}

\prose{อพฺยากโต อสํกิเลสิโก ธมฺโม อพฺยากตสฺส อสํกิเลสิกสฺส ธมฺมสฺส เหตุปจฺจเยน ปจฺจโย. (1)}

\prose{\csromanfolio{598}อพฺยากโต อสํกิเลสิโก ธมฺโม อพฺยากตสฺส อสํกิเลสิกสฺส ธมฺมสฺส อารมฺมณ\-ปจฺจเยน ปจฺจโย.  อพฺยากโต อสํกิเลสิโก ธมฺโม กุสลสฺส อสํกิเลสิกสฺส ธมฺมสฺส อารมฺมณ\-ปจฺจเยน ปจฺจโย. (2) (สํขิตฺตํ.)}

\proseitem{13}{เหตุยา เทฺว, อารมฺมเณ เทฺว, อธิปติยา ตีณิ, อนนฺตเร เทฺว ฯเปฯ อุปนิสฺสเย จตฺตาริ ฯเปฯ อวิคเต เทฺว. (สํขิตฺตํ. ยถา กุสลตฺติเก ปญฺหาวารํ, เอวํ วิตฺถาเรตพฺพํ.)}

\csromanmarkh{1. กุสลตฺติก}\csromanmarkhii{77. สํกิลิฏฺฐทุก}\csromanheadernomark{2}{1. \textbf{กุสลตฺติก} 77. สํกิลิฏฺฐทุก}
\prose{1-7. ปฏิจฺจาทิวาร เหตุอารมฺมณ}

\proseitem{14}{อกุสลํ สํกิลิฏฺฐํ ธมฺมํ ปฏิจฺจ อกุสโล สํกิลิฏฺโฐ ธมฺโม อุปฺปชฺชติ เหตุปจฺจยา. (สํขิตฺตํ.)}

\prose{เหตุยา เอกํ, อารมฺมเณ เอกํ ฯเปฯ อวิคเต เอกํ. (สํขิตฺตํ.) (สหชาตวาเรปิ ฯเปฯ ปญฺหาวาเรปิ สพฺพตฺถ เอกํ.)}

\proseitem{15}{กุสลํ อสํกิลิฏฺฐํ ธมฺมํ ปฏิจฺจ กุสโล อสํกิลิฏฺโฐ ธมฺโม อุปฺปชฺชติ เหตุปจฺจยา.  กุสลํ อสํกิลิฏฺฐํ ธมฺมํ ปฏิจฺจ อพฺยากโต อสํกิลิฏฺโฐ ธมฺโม อุปฺปชฺชติ เหตุปจฺจยา.  กุสลํ อสํกิลิฏฺฐํ ธมฺมํ ปฏิจฺจ กุสโล อสํกิลิฏฺโฐ จ อพฺยากโต อสํกิลิฏฺโฐ จ ธมฺมา อุปฺปชฺชนฺติ เหตุปจฺจยา. (3)}

\prose{อพฺยากตํ อสํกิลิฏฺฐํ ธมฺมํ ปฏิจฺจ อพฺยากโต อสํกิลิฏฺโฐ ธมฺโม อุปฺปชฺชติ เหตุปจฺจยา.  กุสลํ อสํกิลิฏฺฐญฺจ อพฺยากตํ อสํกิลิฏฺฐญฺจ ธมฺมํ ปฏิจฺจ อพฺยากโต อสํกิลิฏฺโฐ ธมฺโม อุปฺปชฺชติ เหตุปจฺจยา. (2)}

\proseitem{16}{กุสลํ อสํกิลิฏฺฐํ ธมฺมํ ปฏิจฺจ กุสโล อสํกิลิฏฺโฐ ธมฺโม อุปฺปชฺชติ อารมฺมณ\-ปจฺจยา. (1)}

\prose{อพฺยากตํ อสํกิลิฏฺฐํ ธมฺมํ ปฏิจฺจ อพฺยากโต อสํกิลิฏฺโฐ ธมฺโม อุปฺปชฺชติ อารมฺมณ\-ปจฺจยา. (1) (สํขิตฺตํ.)}

\csromanmarkh{1. กุสลตฺติก}\csromanmarkhii{79. กิเลสสํกิเลสิกทุก}\csromanheadernomark{2}{1. กุสลตฺติก 79. กิเลสสํกิเลสิกทุก}
\proseitem{17}{\csromanfolio{599}เหตุยา ปญฺจ, อารมฺมเณ เทฺว, อธิปติยา ปญฺจ ฯเปฯ อวิคเต ปญฺจ. (สํขิตฺตํ. สหชาตวารมฺปิ ฯเปฯ สมฺปยุตฺตวารมฺปิ ปฏิจฺจ\-วาร\-สทิสํ.)}

\proseitem{18}{กุสโล อสํกิลิฏฺโฐ ธมฺโม กุสลสฺส อสํกิลิฏฺฐสฺส ธมฺมสฺส เหตุปจฺจเยน ปจฺจโย. ตีณิ.}

\prose{อพฺยากโต อสํกิลิฏฺโฐ ธมฺโม อพฺยากตสฺส อสํกิลิฏฺฐสฺส ธมฺมสฺส เหตุปจฺจเยน ปจฺจโย. (1)}

\prose{กุสโล อสํกิลิฏฺโฐ ธมฺโม กุสลสฺส อสํกิลิฏฺฐสฺส ธมฺมสฺส อารมฺมณ\-ปจฺจเยน ปจฺจโย. เทฺว.}

\prose{อพฺยากโต อสํกิลิฏฺโฐ ธมฺโม อพฺยากตสฺส อสํกิลิฏฺฐสฺส ธมฺมสฺส อารมฺมณ\-ปจฺจเยน ปจฺจโย. เทฺว. (สํขิตฺตํ.)}

\proseitem{19}{เหตุยา จตฺตาริ, อารมฺมเณ จตฺตาริ, อธิปติยา ปญฺจ, อนนฺตเร จตฺตาริ ฯเปฯ อุปนิสฺสเย จตฺตาริ ฯเปฯ อวิคเต สตฺต. (สํขิตฺตํ. ยถา กุสลตฺติเก ปญฺหาวารํ, เอวํ วิตฺถาเรตพฺพํ.)}

\csromanmarkh{1. กุสลตฺติก}\csromanmarkhii{78. กิเลส\-สมฺปยุตฺตทุก}\csromanheadernomark{2}{1. กุสลตฺติก 78. กิเลส\-สมฺปยุตฺตทุก}
\prose{1-7. ปฏิจฺจาทิวารเหตุ}

\proseitem{20}{อกุสลํ กิเลส\-สมฺปยุตฺตํ ธมฺมํ ปฏิจฺจ อกุสโล กิเลส\-สมฺปยุตฺโต ธมฺโม อุปฺปชฺชติ เหตุปจฺจยา. (สํกิลิฏฺฐ\-ทุก\-สทิสํ.)}

\csromanmarkh{1. กุสลตฺติก}\csromanmarkhii{79. กิเลสสํกิเลสิกทุก}\csromanheadernomark{2}{1. กุสลตฺติก 79. กิเลสสํกิเลสิกทุก}
\prose{1-7. ปฏิจฺจาทิวารเหตุ}

\proseitem{21}{อกุสลํ กิเลสญฺเจว สํกิเลสิกญฺจ ธมฺมํ ปฏิจฺจ อกุสโล กิเลโส เจว สํกิเลสิโก จ ธมฺโม อุปฺปชฺชติ เหตุปจฺจยา. (สํขิตฺตํ.)}

\prose{เหตุยา เอกํ ฯเปฯ อวิคเต เอกํ. (สํขิตฺตํ.)}

\vspace{\tipitakaparskip}
\csromancenter{(สหชาตวาเรปิ ฯเปฯ ปญฺหาวาเรปิ สพฺพตฺถ เอกํ.)}
\proseitem{22}{\csromanfolio{600}กุสลํ สํกิเลสิก\-ญฺเจว โน จ กิเลสํ ธมฺมํ ปฏิจฺจ กุสโล สํกิเลสิโก เจว โน จ กิเลโส ธมฺโม อุปฺปชฺชติ เหตุปจฺจยา. ตีณิ.}

\prose{อกุสลํ สํกิเลสิก\-ญฺเจว โน จ กิเลสํ ธมฺมํ ปฏิจฺจ อกุสโล สํกิเลสิโก เจว โน จ กิเลโส ธมฺโม อุปฺปชฺชติ เหตุปจฺจยา. ตีณิ.}

\prose{อพฺยากตํ สํกิเลสิก\-ญฺเจว โน จ กิเลสํ ธมฺมํ ปฏิจฺจ อพฺยากโต สํกิเลสิโก เจว โน จ กิเลโส ธมฺโม อุปฺปชฺชติ เหตุปจฺจยา.}

\prose{กุสลํ สํกิเลสิก\-ญฺเจว โน จ กิเลสญฺจ อพฺยากตํ สํกิเลสิก\-ญฺเจว โน จ กิเลสญฺจ ธมฺมํ ปฏิจฺจ อพฺยากโต สํกิเลสิโก เจว โน จ กิเลโส ธมฺโม อุปฺปชฺชติ เหตุปจฺจยา. (1)}

\prose{อกุสลํ สํกิเลสิก\-ญฺเจว โน จ กิเลสญฺจ อพฺยกตํ สํกิเลสิกญฺจว โน จ กิเลสญฺจ ธมฺมํ ปฏิจฺจ อพฺยากโต สํกิเลสิโก เจว โน จ กิเลโส ธมฺโม อุปฺปชฺชติ เหตุปจฺจยา. (สํขิตฺตํ.)}

\proseitem{23}{เหตุยา นว, อารมฺมเณ ตีณิ ฯเปฯ อวิคเต นว. (สํขิตฺตํ. สหชาตวารมฺปิ ฯเปฯ สมฺปยุตฺตวารมฺปิ ปฏิจฺจ\-วาร\-สทิสํ.)}

\proseitem{24}{กุสโล สํกิเลสิโก เจว โน จ กิเลโส ธมฺโม กุสลสฺส สํกิเลสิกสฺส เจว โน จ กิเลสสฺส ธมฺมสฺส เหตุปจฺจเยน ปจฺจโย. ตีณิ.}

\prose{อพฺยากโต สํกิเลสิโก เจว โน จ กิเลโส ธมฺโม อพฺยากตสฺส สํกิเลสิกสฺส เจว โน จ กิเลสสฺส ธมฺมสฺส เหตุปจฺจเยน ปจฺจโย. (1) (สํขิตฺตํ.)}

\proseitem{25}{เหตุยา จตฺตาริ, อารมฺมเณ นว, อธิปติยา นว ฯเปฯ อวิคเต เตรส. (สํขิตฺตํ. ยถา กุสลตฺติเก ปญฺหาวารํ, เอวํ วิตฺถาเรตพฺพํ.)}

\vspace{\tipitakaparskip}
\csromancenter{กุสลตฺติก 81. กิเลส\-กิเลส\-สมฺปยุตฺตทุก 1. กุสลตฺติก 80. กิเลส\-สํกิลิฏฺฐทุก}
\prose{\csromanfolio{601}1. ปฏิจฺจวารเหตุ}

\proseitem{26}{อกุสลํ กิเลสญฺเจว สํกิลิฏฺฐญฺจ ธมฺมํ ปฏิจฺจ อกุสโล กิเลโส เจว สํกิลิฏฺโฐ จ ธมฺโม อุปฺปชฺชติ เหตุปจฺจยา. (สํขิตฺตํ.)}

\prose{เหตุยา เอกํ, อารมฺมเณ เอกํ ฯเปฯ อวิคเต เอกํ. (สํขิตฺตํ.) (สหชาตวาเรปิ ฯเปฯ ปญฺหาวาเรปิ สพฺพตฺถ เอกํ.)}

\proseitem{27}{อกุสลํ สํกิลิฏฺฐ\-ญฺเจว โน จ กิเลสํ ธมฺมํ ปฏิจฺจ อกุสโล สํกิลิฏฺโฐ เจว โน จ กิเลโส ธมฺโม อุปฺปชฺชติ เหตุปจฺจยา.}

\prose{เหตุยา เอกํ ฯเปฯ อวิคเต เอกํ. (สํขิตฺตํ.)}

\vspace{\tipitakaparskip}
\csromancenter{(สหชาตวาเรปิ ฯเปฯ ปญฺหาวาเรปิ สพฺพตฺถ เอกํ.)}
\csromanmarkh{1. กุสลตฺติก}\csromanmarkhii{81. กิเลส\-กิเลส\-สมฺปยุตฺตทุก}\csromanheadernomark{2}{1. กุสลตฺติก 81. กิเลส\-กิเลส\-สมฺปยุตฺตทุก}
\prose{1. ปฏิจฺจวารเหตุ}

\proseitem{28}{อกุสลํ กิเลสญฺเจว กิเลส\-สมฺปยุตฺตญฺจ ธมฺมํ ปฏิจฺจ อกุสโล กิเลโส เจว กิเลส\-สมฺปยุตฺโต จ ธมฺโม อุปฺปชฺชติ เหตุปจฺจยา.}

\prose{เหตุยา เอกํ ฯเปฯ อวิคเต เอกํ. (สํขิตฺตํ.) (สหชาตวาเรปิ ฯเปฯ ปญฺหาวาเรปิ สพฺพตฺถ เอกํ.)}

\proseitem{29}{อกุสลํ กิเลส\-สมฺปยุตฺต\-ญฺเจว โน จ กิเลสํ ธมฺมํ ปฏิจฺจ อกุสโล กิเลส\-สมฺปยุตฺโต เจว โน จ กิเลโส ธมฺโม อุปฺปชฺชติ เหตุปจฺจยา. (สํขิตฺตํ.)}

\prose{เหตุยา เอกํ ฯเปฯ อวิคเต เอกํ. (สํขิตฺตํ.)}

\vspace{\tipitakaparskip}
\csromancenter{(สหชาตวาเรปิ ฯเปฯ ปญฺหาวาเรปิ สพฺพตฺถ เอกํ.)}
\csromancenter{1. กุสลตฺติก 82. กิเลส\-วิปฺปยุตฺต\-สํกิเลสิกทุก}
\prose{\csromanfolio{602}1-7. ปฏิจฺจาทิวารเหตุ}

\proseitem{30}{กุสลํ กิเลส\-วิปฺปยุตฺตํ สํกิเลสิกํ ธมฺมํ ปฏิจฺจ กุสโล กิเลส\-วิปฺปยุตฺโต สํกิเลสิโก ธมฺโม อุปฺปชฺชติ เหตุปจฺจยา.  กุสลํ กิเลส\-วิปฺปยุตฺตํ สํกิเลสิกํ ธมฺมํ ปฏิจฺจ อพฺยากโต กิเลส\-วิปฺปยุตฺโต สํกิเลสิโก ธมฺโม อุปฺปชฺชติ เหตุปจฺจยา.  กุสลํ กิเลส\-วิปฺปยุตฺตํ สํกิเลสิกํ ธมฺมํ ปฏิจฺจ กุสโล กิเลส\-วิปฺปยุตฺโต สํกิเลสิโก จ อพฺยากโต กิเลส\-วิปฺปยุตฺโต สํกิเลสิโก จ ธมฺมา อุปฺปชฺชนฺติ เหตุปจฺจยา. (3)}

\prose{อพฺยากตํ กิเลส\-วิปฺปยุตฺตํ สํกิเลสิกํ ธมฺมํ ปฏิจฺจ อพฺยากโต กิเลส\-วิปฺปยุตฺโต สํกิเลสิโก ธมฺโม อุปฺปชฺชติ เหตุปจฺจยา. (1)}

\prose{กุสลํ กิเลส\-วิปฺปยุตฺตํ สํกิเลสิกญฺจ อพฺยากตํ กิเลส\-วิปฺปยุตฺตํ สํกิเลสิกญฺจ ธมฺมํ ปฏิจฺจ อพฺยากโต กิเลส\-วิปฺปยุตฺโต สํกิเลสิโก ธมฺโม อุปฺปชฺชติ เหตุปจฺจยา. (1) (สํขิตฺตํ.)}

\proseitem{31}{เหตุยา ปญฺจ, อารมฺมเณ เทฺว ฯเปฯ อวิคเต ปญฺจ. (สํขิตฺตํ. สหชาตวารมฺปิ ฯเปฯ สมฺปยุตฺตวารมฺปิ ปฏิจฺจ\-วาร\-สทิสํ.)}

\proseitem{32}{กุสโล กิเลส\-วิปฺปยุตฺโต สํกิเลสิโก ธมฺโม กุสลสฺส กิเลส\-วิปฺปยุตฺตสฺส สํกิเลสิกสฺส ธมฺมสฺส เหตุปจฺจเยน ปจฺจโย. ตีณิ.}

\prose{อพฺยากโต กิเลส\-วิปฺปยุตฺโต สํกิเลสิโก ธมฺโม อพฺยากตสฺส กิเลส\-วิปฺปยุตฺตสฺส สํกิเลสิกสฺส ธมฺมสฺส เหตุปจฺจเยน ปจฺจโย. (1) (สํขิตฺตํ.)}

\proseitem{33}{เหตุยา จตฺตาริ, อารมฺมเณ จตฺตาริ ฯเปฯ อวิคเต สตฺต. (สํขิตฺตํ. ยถา กุสลตฺติเก ปญฺหาวารํ, เอวํ วิตฺถาเรตพฺพํ.)}

\proseitem{34}{กุสลํ กิเลส\-วิปฺปยุตฺตํ อสํกิเลสิกํ ธมฺมํ ปฏิจฺจ กุสโล กิเลส\-วิปฺปยุตฺโต อสํกิเลสิโก ธมฺโม อุปฺปชฺชติ เหตุปจฺจยา.}

\csromanmarkh{1. กุสลตฺติก}\csromanmarkhii{83. ทสฺสเนน\-ปหาตพฺพทุก}\csromanheadernomark{2}{1. กุสลตฺติก 83. ทสฺสเนน\-ปหาตพฺพทุก}
\csromanmatikaanchor{mtk.149}
\csromantocmarkref{subparagraph}{83-99. ทสฺสเนนปหาตพฺพทุกาทิ}{mtk.149}
\bookmark[dest={mtk.149},level=5]{83-99. ทสฺสเนนปหาตพฺพทุกาทิ}
\prose{\csromanfolio{603}อพฺยากตํ กิเลส\-วิปฺปยุตฺตํ อสํกิเลสิกํ ธมฺมํ ปฏิจฺจ อพฺยากโต กิเลส\-วิปฺปยุตฺโต อสํกิเลสิโก ธมฺโม อุปฺปชฺชติ เหตุปจฺจยา.}

\prose{เหตุยา เทฺว ฯเปฯ อวิคเต เทฺว. (สํขิตฺตํ. สหชาตวารมฺปิ ฯเปฯ สมฺปยุตฺตวารมฺปิ ปฏิจฺจ\-วาร\-สทิสํ.)}

\proseitem{35}{กุสโล กิเลส\-วิปฺปยุตฺโต อสํกิเลสิโก ธมฺโม กุสลสฺส กิเลส\-วิปฺปยุตฺตสฺส อสํกิเลสิกสฺส ธมฺมสฺส เหตุปจฺจเยน ปจฺจโย. (1)}

\prose{อพฺยากโต กิเลส\-วิปฺปยุตฺโต อสํกิเลสิโก ธมฺโม อพฺยากตสฺส กิเลส\-วิปฺปยุตฺตสฺส อสํกิเลสิกสฺส ธมฺมสฺส เหตุปจฺจเยน ปจฺจโย. (1) (สํขิตฺตํ.)}

\proseitemwithcloser{36}{เหตุยา เทฺว, อารมฺมเณ เทฺว, อธิปติยา ตีณิ, อนนฺตเร เทฺว ฯเปฯ อุปนิสฺสเย จตฺตาริ ฯเปฯ อวิคเต เทฺว. (สํขิตฺตํ. ยถา กุสลตฺติเก ปญฺหาวารํ, เอวํ วิตฺถาเรตพฺพํ.)}{กุสล\-ตฺติก\-กิเลส\-โคจฺฉกํ นิฏฺฐิตํ.}
\csromanmarkh{1. กุสลตฺติก}\csromanmarkhii{83. ทสฺสเนน\-ปหาตพฺพทุก}\csromanheadernomark{2}{1. กุสลตฺติก 83. ทสฺสเนน\-ปหาตพฺพทุก}
\prose{1-7. ปฏิจฺจาทิวารเหตุ}

\proseitem{1}{อกุสลํ ทสฺสเนน ปหาตพฺพํ ธมฺมํ ปฏิจฺจ อกุสโล ทสฺสเนน ปหาตพฺโพ ธมฺโม อุปฺปชฺชติ เหตุปจฺจยา. (สํขิตฺตํ.)}

\prose{เหตุยา เอกํ, อารมฺมเณ เอกํ ฯเปฯ อวิคเต เอกํ. (สํขิตฺตํ.) (สหชาตวาเรปิ ฯเปฯ ปญฺหาวาเรปิ สพฺพตฺถ เอกํ.)}

\proseitem{2}{กุสลํ นทสฺสเนน ปหาตพฺพํ ธมฺมํ ปฏิจฺจ กุสโล นทสฺสเนน ปหาตพฺโพ ธมฺโม อุปฺปชฺชติ เหตุปจฺจยา. ตีณิ.}

\prose{อกุสลํ นทสฺสเนน ปหาตพฺพํ ธมฺมํ ปฏิจฺจ อกุสโล นทสฺสเนน ปหาตพฺโพ ธมฺโม อุปฺปชฺชติ เหตุปจฺจยา. ตีณิ.}

\prose{อพฺยากตํ นทสฺสเนน ปหาตพฺพํ ธมฺมํ ปฏิจฺจ อพฺยากโต นทสฺสเนน ปหาตพฺโพ ธมฺโม อุปฺปชฺชติ เหตุปจฺจยา. (1)}

\prose{\csromanfolio{604}กุสลํ นทสฺสเนน ปหาตพฺพญฺจ อพฺยากตํ นทสฺสเนน ปหาตพฺพญฺจ ธมฺมํ ปฏิจฺจ อพฺยากโต นทสฺสเนน ปหาตพฺโพ ธมฺโม อุปฺปชฺชติ เหตุปจฺจยา. (1)}

\prose{อกุสลํ นทสฺสเนน ปหาตพฺพญฺจ อพฺยากตํ นทสฺสเนน ปหาตพฺพญฺจ ธมฺมํ ปฏิจฺจ อพฺยากโต นทสฺสเนน ปหาตพฺโพ ธมฺโม อุปฺปชฺชติ เหตุปจฺจยา. (1)}

\proseitem{3}{เหตุยา นว, อารมฺมเณ ตีณิ ฯเปฯ อวิคเต นว. (สํขิตฺตํ. สหชาตวารมฺปิ ฯเปฯ สมฺปยุตฺตวารมฺปิ ปฏิจฺจ\-วาร\-สทิสํ.)}

\proseitem{4}{กุสโล นทสฺสเนน ปหาตพฺโพ ธมฺโม กุสลสฺส นทสฺสเนน ปหาตพฺพสฺส ธมฺมสฺส เหตุปจฺจเยน ปจฺจโย. ตีณิ.}

\prose{อกุสโล นทสฺสเนน ปหาตพฺโพ ธมฺโม อกุสลสฺส นทสฺสเนน ปหาตพฺพสฺส ธมฺมสฺส เหตุปจฺจเยน ปจฺจโย. ตีณิ.}

\prose{อพฺยากโต นทสฺสเนน ปหาตพฺโพ ธมฺโม อพฺยากตสฺส นทสฺสเนน ปหาตพฺพสฺส ธมฺมสฺส เหตุปจฺจเยน ปจฺจโย. (1) (สํขิตฺตํ.)}

\proseitem{5}{เหตุยา สตฺต, อารมฺมเณ นว, อธิปติยา ทส, อนนฺตเร สตฺต ฯเปฯ อวิคเต เตรส. (สํขิตฺตํ.)}

\prose{(ยถา กุสลตฺติเก ปญฺหาวารํ, เอวํ วิตฺถาเรตพฺพํ.)}

\csromanmarkh{1. กุสลตฺติก}\csromanmarkhii{84. ภาวนาย\-ปหาตพฺพทุก}\csromanheadernomark{2}{1. กุสลตฺติก 84. ภาวนาย\-ปหาตพฺพทุก}
\prose{1. ปฏิจฺจวารเหตุ}

\proseitem{6}{อกุสลํ ภาวนาย ปหาตพฺพํ ธมฺมํ ปฏิจฺจ อกุสโล ภาวนาย ปหาตพฺโพ ธมฺโม อุปฺปชฺชติ เหตุปจฺจยา. (สํขิตฺตํ.)}

\prose{เหตุยา เอกํ, อารมฺมเณ เอกํ ฯเปฯ อวิคเต เอกํ. (สํขิตฺตํ.) (สหชาตวาเรปิ ฯเปฯ ปญฺหาวาเรปิ สพฺพตฺถ เอกํ.)}

\proseitem{7}{กุสลํ นภาวนาย ปหาตพฺพํ ธมฺมํ ปฏิจฺจ กุสโล นภาวนาย ปหาตพฺโพ ธมฺโม อุปฺปชฺชติ เหตุปจฺจยา. (สํขิตฺตํ.)}

\prose{เหตุยา นว, อารมฺมเณ ตีณิ ฯเปฯ อวิคเต นว. (สํขิตฺตํ.)}

\vspace{\tipitakaparskip}
\csromancenter{(สหชาตวารมฺปิ ฯเปฯ ปญฺหาวารมฺปิ วิตฺถาเรตพฺพํ.)}
\csromancenter{กุสลตฺติก 87. สวิตกฺกทุก 1. กุสลตฺติก 85. ทสฺสเนน\-ปหาตพฺพ\-เหตุกทุก}
\prose{\csromanfolio{605}1. ปฏิจฺจวารเหตุ}

\proseitem{8}{อกุสลํ ทสฺสเนน ปหาตพฺพ\-เหตุกํ ธมฺมํ ปฏิจฺจ อกุสโล ทสฺสเนน ปหาตพฺพ\-เหตุโก ธมฺโม อุปฺปชฺชติ เหตุปจฺจยา. (สํขิตฺตํ.)}

\prose{เหตุยา เอกํ, อารมฺมเณ เอกํ ฯเปฯ อวิคเต เอกํ. (สํขิตฺตํ.) (สหชาตวาเรปิ ฯเปฯ ปญฺหาวาเรปิ สพฺพตฺถ เอกํ.)}

\proseitem{9}{กุสลํ นทสฺสเนน ปหาตพฺพ\-เหตุกํ ธมฺมํ ปฏิจฺจ กุสโล นทสฺสเนน ปหาตพฺพ\-เหตุโก ธมฺโม อุปฺปชฺชติ เหตุปจฺจยา. (สํขิตฺตํ.)}

\prose{เหตุยา นว ฯเปฯ อวิคเต นว. (สํขิตฺตํ. สหชาตวารมฺปิ ฯเปฯ ปญฺหาวารมฺปิ วิตฺถาเรตพฺพํ.)}

\csromanmarkh{1. กุสลตฺติก}\csromanmarkhii{86. ภาวนาย\-ปหาตพฺพ\-เหตุกทุก}\csromanheadernomark{2}{1. กุสลตฺติก 86. ภาวนาย\-ปหาตพฺพ\-เหตุกทุก}
\csromanheader{2}{1. \textbf{ปฏิจฺจวาร ปจฺจยจตุกฺก}}
\proseitem{10}{อกุสลํ ภาวนาย ปหาตพฺพ\-เหตุกํ ธมฺมํ ปฏิจฺจ อกุสโล ภาวนาย ปหาตพฺพ\-เหตุโก ธมฺโม อุปฺปชฺชติ เหตุปจฺจยา. (สํขิตฺตํ.)}

\prose{เหตุยา เอกํ, อารมฺมเณ เอกํ ฯเปฯ อวิคเต เอกํ. (สํขิตฺตํ.) (สหชาตวาเรปิ ฯเปฯ ปญฺหาวาเรปิ สพฺพตฺถ เอกํ.)}

\proseitem{11}{กุสลํ นภาวนาย ปหาตพฺพ\-เหตุกํ ธมฺมํ ปฏิจฺจ กุสโล นภาวนาย ปหาตพฺพ\-เหตุโก ธมฺโม อุปฺปชฺชติ เหตุปจฺจยา.}

\prose{เหตุยา นว, อารมฺมเณ ตีณิ ฯเปฯ อวิคเต นว. (สํขิตฺตํ.)}

\vspace{\tipitakaparskip}
\csromancenter{(สหชาตวาเรปิ ฯเปฯ ปญฺหาวาเรปิ วิตฺถาเรตพฺพํ.)}
\csromanmarkh{1. กุสลตฺติก}\csromanmarkhii{87. สวิตกฺกทุก}\csromanheadernomark{2}{1. กุสลตฺติก 87. สวิตกฺกทุก}
\prose{1-7. ปฏิจฺจาทิวาร เหตุอารมฺมณ}

\proseitem{12}{กุสลํ สวิตกฺกํ ธมฺมํ ปฏิจฺจ กุสโล สวิตกฺโก ธมฺโม อุปฺปชฺชติ เหตุปจฺจยา. (1)}

\prose{\csromanfolio{606}อกุสลํ สวิตกฺกํ ธมฺมํ ปฏิจฺจ อกุสโล สวิตกฺโก ธมฺโม อุปฺปชฺชติ เหตุปจฺจยา. (1)}

\prose{อพฺยากตํ สวิตกฺกํ ธมฺมํ ปฏิจฺจ อพฺยากโต สวิตกฺโก ธมฺโม อุปฺปชฺชติ เหตุปจฺจยา. (1) (สํขิตฺตํ.)}

\proseitem{13}{เหตุยา ตีณิ, อารมฺมเณ ตีณิ ฯเปฯ อวิคเต ตีณิ. (สํขิตฺตํ. สหชาตวารมฺปิ ฯเปฯ สมฺปยุตฺตวารมฺปิ ปฏิจฺจ\-วาร\-สทิสํ.)}

\proseitem{14}{กุสโล สวิตกฺโก ธมฺโม กุสลสฺส สวิตกฺกสฺส ธมฺมสฺส เหตุปจฺจเยน ปจฺจโย. (1)}

\prose{อกุสโล สวิตกฺโก ธมฺโม อกุสลสฺส สวิตกฺกสฺส ธมฺมสฺส เหตุปจฺจเยน ปจฺจโย. (1)}

\prose{อพฺยากโต สวิตกฺโก ธมฺโม อพฺยากตสฺส สวิตกฺกสฺส ธมฺมสฺส เหตุปจฺจเยน ปจฺจโย. (1)}

\prose{กุสโล สวิตกฺโก ธมฺโม กุสลสฺส สวิตกฺกสฺส ธมฺมสฺส อารมฺมณ\-ปจฺจเยน ปจฺจโย. ตีณิ.}

\prose{อกุสโล สวิตกฺโก ธมฺโม อกุสลสฺส สวิตกฺกสฺส ธมฺมสฺส อารมฺมณ\-ปจฺจเยน ปจฺจโย. ตีณิ.}

\prose{อพฺยากโต สวิตกฺโก ธมฺโม อพฺยากตสฺส สวิตกฺกสฺส ธมฺมสฺส อารมฺมณ\-ปจฺจเยน ปจฺจโย. ตีณิ. (สํขิตฺตํ.)}

\proseitem{15}{เหตุยา ตีณิ, อารมฺมเณ นว, อธิปติยา สตฺต ฯเปฯ อวิคเต ตีณิ. (สํขิตฺตํ.)}

\prose{(ยถา กุสลตฺติเก ปญฺหาวารํ, เอวํ วิตฺถาเรตพฺพํ.)}

\prose{อวิตกฺกปทเหตุ}

\proseitem{16}{กุสลํ อวิตกฺกํ ธมฺมํ ปฏิจฺจ กุสโล อวิตกฺโก ธมฺโม อุปฺปชฺชติ เหตุปจฺจยา. (สํขิตฺตํ.)}

\csromanmarkh{1. กุสลตฺติก}\csromanmarkhii{89. สปฺปีติกทุก}\csromanheadernomark{2}{1. กุสลตฺติก 89. สปฺปีติกทุก}
\prose{\csromanfolio{607}เหตุยา สตฺต, อารมฺมเณ เทฺว ฯเปฯ วิปาเก เอกํ, อวิคเต สตฺต. (สํขิตฺตํ. สหชาตวารมฺปิ ฯเปฯ สมฺปยุตฺตวารมฺปิ ปฏิจฺจ\-วาร\-สทิสํ.)}

\proseitem{17}{กุสโล อวิตกฺโก ธมฺโม กุสลสฺส อวิตกฺกสฺส ธมฺมสฺส เหตุปจฺจเยน ปจฺจโย. ตีณิ.}

\prose{อพฺยากโต อวิตกฺโก ธมฺโม อพฺยากตสฺส อวิตกฺกสฺส ธมฺมสฺส เหตุปจฺจเยน ปจฺจโย. (1) (สํขิตฺตํ.)}

\proseitem{18}{เหตุยา จตฺตาริ, อารมฺมเณ นว ฯเปฯ อวิคเต ทส. (สํขิตฺตํ. ยถา กุสลตฺติเก ปญฺหาวารํ, เอวํ วิตฺถาเรตพฺพํ.)}

\csromanmarkh{1. กุสลตฺติก}\csromanmarkhii{88. สวิจารทุก}\csromanheadernomark{2}{1. \textbf{กุสลตฺติก} 88. สวิจารทุก}
\csromanheader{2}{1. \textbf{ปฏิจฺจวาร ปจฺจยจตุกฺก}}
\proseitem{19}{กุสลํ สวิจารํ ธมฺมํ ปฏิจฺจ กุสโล สวิจาโร ธมฺโม อุปฺปชฺชติ เหตุปจฺจยา. (กุสล\-สวิตกฺก\-สทิสํ.)}

\csromanmarkh{1. กุสลตฺติก}\csromanmarkhii{89. สปฺปีติกทุก}\csromanheadernomark{2}{1. กุสลตฺติก 89. สปฺปีติกทุก}
\csromanheader{2}{1-7. ปฏิจฺจาทิวาร ปจฺจยจตุกฺก}
\proseitem{20}{กุสลํ สปฺปีติกํ ธมฺมํ ปฏิจฺจ กุสโล สปฺปีติโก ธมฺโม อุปฺปชฺชติ เหตุปจฺจยา. (1)}

\prose{อกุสลํ สปฺปีติกํ ธมฺมํ ปฏิจฺจ อกุสโล สปฺปีติโก ธมฺโม อุปฺปชฺชติ เหตุปจฺจยา. (1)}

\prose{อพฺยากตํ สปฺปีติกํ ธมฺมํ ปฏิจฺจ อพฺยากโต สปฺปีติโก ธมฺโม อุปฺปชฺชติ เหตุปจฺจยา. (1) (สํขิตฺตํ.)}

\proseitem{21}{\csromanfolio{608}เหตุยา ตีณิ, อารมฺมเณ ตีณิ ฯเปฯ อวิคเต ตีณิ. (สํขิตฺตํ. สหชาตวารมฺปิ ฯเปฯ สมฺปยุตฺตวารมฺปิ ปฏิจฺจ\-วาร\-สทิสํ.)}

\proseitem{22}{กุสโล สปฺปีติโก ธมฺโม กุสลสฺส สปฺปีติกสฺส ธมฺมสฺส เหตุปจฺจเยน ปจฺจโย. (1)}

\prose{อกุสโล สปฺปีติโก ธมฺโม อกุสลสฺส สปฺปีติกสฺส ธมฺมสฺส เหตุปจฺจเยน ปจฺจโย. (1)}

\prose{อพฺยากโต สปฺปีติโก ธมฺโม อพฺยากตสฺส สปฺปีติกสฺส ธมฺมสฺส เหตุปจฺจเยน ปจฺจโย. (1) (สํขิตฺตํ.)}

\proseitem{23}{เหตุยา ตีณิ, อารมฺมเณ นว ฯเปฯ อวิคเต ตีณิ. (สํขิตฺตํ.)}

\prose{อปฺปีติกปทเหตุ}

\proseitem{24}{กุสลํ อปฺปีติกํ ธมฺมํ ปฏิจฺจ กุสโล อปฺปีติโก ธมฺโม อุปฺปชฺชติ เหตุปจฺจยา. ตีณิ.}

\prose{อกุสลํ อปฺปีติกํ ธมฺมํ ปฏิจฺจ อกุสโล อปฺปีติโก ธมฺโม อุปฺปชฺชติ เหตุปจฺจยา. ตีณิ.}

\prose{อพฺยากตํ อปฺปีติกํ ธมฺมํ ปฏิจฺจ อพฺยากโต อปฺปีติโก ธมฺโม อุปฺปชฺชติ เหตุปจฺจยา. (1)}

\prose{กุสลํ อปฺปีติกญฺจ อพฺยากตํ อปฺปีติกญฺจ ธมฺมํ ปฏิจฺจ อพฺยากโต อปฺปีติโก ธมฺโม อุปฺปชฺชติ เหตุปจฺจยา. (1)}

\prose{อกุสลํ อปฺปีติกญฺจ อพฺยากตํ อปฺปีติกญฺจ ธมฺมํ ปฏิจฺจ อพฺยากโต อปฺปีติโก ธมฺโม อุปฺปชฺชติ เหตุปจฺจยา. (1) (สํขิตฺตํ.)}

\proseitem{25}{เหตุยา นว, อารมฺมเณ ตีณิ ฯเปฯ อวิคเต นว. (สํขิตฺตํ. สหชาตวารมฺปิ ฯเปฯ สมฺปยุตฺตวารมฺปิ ปฏิจฺจ\-วาร\-สทิสํ.)}

\csromanmarkh{1. กุสลตฺติก}\csromanmarkhii{93. กามาวจรทุก}\csromanheadernomark{2}{1. กุสลตฺติก 93. กามาวจรทุก}
\proseitem{26}{\csromanfolio{609}กุสโล อปฺปีติโก ธมฺโม กุสลสฺส อปฺปีติกสฺส ธมฺมสฺส เหตุปจฺจเยน ปจฺจโย. (สํขิตฺตํ.)}

\prose{เหตุยา สตฺต, อารมฺมเณ นว ฯเปฯ อวิคเต เตรส. (สํขิตฺตํ. ยถา กุสลตฺติเก ปญฺหาวารํ, เอวํ วิตฺถาเรตพฺพํ.)}

\csromanheader{2}{1. กุสลตฺติก 90-92. ปีติสหคตา\-ทิทุก}
\csromanheader{2}{1. \textbf{ปฏิจฺจวาร ปจฺจยจตุกฺก}}
\proseitem{27}{กุสลํ ปีติสหคตํ ธมฺมํ ปฏิจฺจ กุสโล ปีติสหคโต ธมฺโม อุปฺปชฺชติ เหตุปจฺจยา.  กุสลํ นปีติสหคตํ ธมฺมํ ปฏิจฺจ ฯเปฯ. กุสลํ สุขสหคตํ ธมฺมํ ปฏิจฺจ ฯเปฯ. กุสลํ นสุขสหคตํ ธมฺมํ ปฏิจฺจ ฯเปฯ. กุสลํ อุเปกฺขา\-สหคตํ ธมฺมํ ปฏิจฺจ ฯเปฯ. กุสลํ นอุเปกฺขา\-สหคตํ ธมฺมํ ปฏิจฺจ. (สํขิตฺตํ.)}

\csromanmarkh{1. กุสลตฺติก}\csromanmarkhii{93. กามาวจรทุก}\csromanheadernomark{2}{1. กุสลตฺติก 93. กามาวจรทุก}
\csromanheader{2}{1-7. ปฏิจฺจาทิวาร ปจฺจยจตุกฺกเหตุ}
\proseitem{28}{กุสลํ กามาวจรํ ธมฺมํ ปฏิจฺจ กุสโล กามาวจโร ธมฺโม อุปฺปชฺชติ เหตุปจฺจยา.  กุสลํ กามาวจรํ ธมฺมํ ปฏิจฺจ อพฺยากโต กามาวจโร ธมฺโม อุปฺปชฺชติ เหตุปจฺจยา.  กุสลํ กามาวจรํ ธมฺมํ ปฏิจฺจ กุสโล กามาวจโร จ อพฺยากโต กามาวจโร จ ธมฺมา อุปฺปชฺชนฺติ เหตุปจฺจยา. (3)}

\prose{อกุสลํ กามาวจรํ ธมฺมํ ปฏิจฺจ อกุสโล กามาวจโร ธมฺโม อุปฺปชฺชติ เหตุปจฺจยา. ตีณิ.}

\prose{อพฺยากตํ กามาวจรํ ธมฺมํ ปฏิจฺจ อพฺยากโต กามาวจโร ธมฺโม อุปฺปชฺชติ เหตุปจฺจยา. (1)}

\prose{กุสลํ กามาวจรญฺจ อพฺยากตํ กามาวจรญฺจ ธมฺมํ ปฏิจฺจ อพฺยากโต กามาวจโร ธมฺโม อุปฺปชฺชติ เหตุปจฺจยา. (1)}

\prose{\csromanfolio{610}อกุสลํ กามาวจรญฺจ อพฺยากตํ กามาวจรญฺจ ธมฺมํ ปฏิจฺจ อพฺยากโต กามาวจโร ธมฺโม อุปฺปชฺชติ เหตุปจฺจยา. (1) (สํขิตฺตํ.)}

\proseitem{29}{เหตุยา นว, อารมฺมเณ ตีณิ ฯเปฯ อวิคเต นว. (สํขิตฺตํ. สหชาตวารมฺปิ ฯเปฯ สมฺปยุตฺตวารมฺปิ ปฏิจฺจ\-วาร\-สทิสํ.)}

\proseitem{30}{กุสโล กามาวจโร ธมฺโม กุสลสฺส กามาวจรสฺส ธมฺมสฺส เหตุปจฺจเยน ปจฺจโย. ตีณิ.}

\prose{อกุสโล กามาวจโร ธมฺโม กุสลสฺส กามาวจรสฺส ธมฺมสฺส เหตุปจฺจเยน ปจฺจโย. ตีณิ.}

\prose{อพฺยากโต กามาวจโร ธมฺโม อพฺยากตสฺส กามาวจรสฺส ธมฺมสฺส เหตุปจฺจเยน ปจฺจโย. (1) (สํขิตฺตํ.)}

\proseitem{31}{เหตุยา สตฺต, อารมฺมเณ นว, อธิปติยา นว ฯเปฯ อวิคเต เตรส. (สํขิตฺตํ. ยถา กุสลตฺติเก ปญฺหาวารํ, เอวํ วิตฺถาเรตพฺพํ.)}

\prose{นกามาวจรปทเหตุ อารมฺมณ}

\proseitem{32}{กุสลํ นกามาวจรํ ธมฺมํ ปฏิจฺจ กุสโล นกามาวจโร ธมฺโม อุปฺปชฺชติ เหตุปจฺจยา. (1)}

\prose{อพฺยากตํ นกามาวจรํ ธมฺมํ ปฏิจฺจ อพฺยากโต นกามาวจโร ธมฺโม อุปฺปชฺชติ เหตุปจฺจยา. (1) (สํขิตฺตํ.)}

\prose{เหตุยา เทฺว, อารมฺมเณ เทฺว ฯเปฯ อวิคเต เทฺว. (สํขิตฺตํ. สหชาตวารมฺปิ ฯเปฯ สมฺปยุตฺตวารมฺปิ ปฏิจฺจ\-วาร\-สทิสํ.)}

\proseitem{33}{กุสโล นกามาวจโร ธมฺโม กุสลสฺส นกามาวจรสฺส ธมฺมสฺส เหตุปจฺจเยน ปจฺจโย. (1)}

\prose{อพฺยากโต นกามาวจโร ธมฺโม อพฺยากตสฺส นกามาวจรสฺส ธมฺมสฺส เหตุปจฺจเยน ปจฺจโย. (1)}

\prose{กุสโล นกามาวจโร ธมฺโม กุสลสฺส นกามาวจรสฺส ธมฺมสฺส อารมฺมณ\-ปจฺจเยน ปจฺจโย.  กุสโล นกามาวจโร ธมฺโม อพฺยากตสฺส นกามาวจรสฺส ธมฺมสฺส อารมฺมณ\-ปจฺจเยน ปจฺจโย. (2)}

\csromanmarkh{1. กุสลตฺติก}\csromanmarkhii{94. รูปาวจรทุก}\csromanheadernomark{2}{1. กุสลตฺติก 94. รูปาวจรทุก}
\prose{\csromanfolio{611}อพฺยากโต นกามาวจโร ธมฺโม อพฺยากตสฺส นกามาวจรสฺส ธมฺมสฺส อารมฺมณ\-ปจฺจเยน ปจฺจโย.  อพฺยากโต นกามาวจโร ธมฺโม กุสลสฺส นกามาวจรสฺส ธมฺมสฺส อารมฺมณ\-ปจฺจเยน ปจฺจโย. (2)}

\proseitem{34}{เหตุยา เทฺว, อารมฺมเณ จตฺตาริ, อธิปติยา ตีณิ ฯเปฯ อวิคเต เทฺว. (สํขิตฺตํ.)}

\prose{(ยถา กุสลตฺติเก ปญฺหาวารํ, เอวํ วิตฺถาเรตพฺพํ.)}

\csromanmarkh{1. กุสลตฺติก}\csromanmarkhii{94. รูปาวจรทุก}\csromanheadernomark{2}{1. กุสลตฺติก 94. รูปาวจรทุก}
\prose{1-7. ปฏิจฺจาทิวารปจฺจยจตุกฺกเหตุ}

\proseitem{35}{กุสลํ รูปาวจรํ ธมฺมํ ปฏิจฺจ กุสโล รูปาวจโร ธมฺโม อุปฺปชฺชติ เหตุปจฺจยา. (1)}

\prose{อพฺยากตํ รูปาวจรํ ธมฺมํ ปฏิจฺจ อพฺยากโต รูปาวจโร ธมฺโม อุปฺปชฺชติ เหตุปจฺจยา. (1) (สํขิตฺตํ.)}

\prose{เหตุยา เทฺว ฯเปฯ อวิคเต เทฺว. (สํขิตฺตํ.) (สหชาตวารมฺปิ ฯเปฯ ปญฺหาวารมฺปิ วิตฺถาเรตพฺพํ.)}

\proseitem{36}{กุสลํ นรูปาวจรํ ธมฺมํ ปฏิจฺจ กุสโล นรูปาวจโร ธมฺโม อุปฺปชฺชติ เหตุปจฺจยา. ตีณิ.}

\prose{อกุสลํ นรูปาวจรํ ธมฺมํ ปฏิจฺจ อกุสโล นรูปาวจโร ธมฺโม อุปฺปชฺชติ เหตุปจฺจยา. ตีณิ.}

\prose{อพฺยากตํ นรูปาวจรํ ธมฺมํ ปฏิจฺจ อพฺยากโต นรูปาวจโร ธมฺโม อุปฺปชฺชติ เหตุปจฺจยา. (1)}

\prose{กุสลํ นรูปาวจรญฺจ อพฺยากตํ นรูปาวจรญฺจ ธมฺมํ ปฏิจฺจ อพฺยากโต นรูปาวจโร ธมฺโม อุปฺปชฺชติ เหตุปจฺจยา. (1)}

\prose{อกุสลํ นรูปาวจรญฺจ อพฺยากตํ นรูปาวจรญฺจ ธมฺมํ ปฏิจฺจ อพฺยากโต นรูปาวจโร ธมฺโม อุปฺปชฺชติ เหตุปจฺจยา. (1) (สํขิตฺตํ.)}

\proseitem{37}{\csromanfolio{612}เหตุยา นว, อารมฺมเณ ตีณิ ฯเปฯ อวิคเต นว. (สํขิตฺตํ. สหชาตวารมฺปิ ฯเปฯ สมฺปยุตฺตวารมฺปิ ปฏิจฺจ\-วาร\-สทิสํ.)}

\proseitem{38}{กุสโล นรูปาวจโร ธมฺโม กุสลสฺส นรูปาวจรสฺส ธมฺมสฺส เหตุปจฺจเยน ปจฺจโย. ตีณิ.}

\prose{อกุสโล นรูปาวจโร ธมฺโม อกุสลสฺส นรูปาวจรสฺส ธมฺมสฺส เหตุปจฺจเยน ปจฺจโย. ตีณิ.}

\prose{อพฺยากโต นรูปาวจโร ธมฺโม อพฺยากตสฺส นรูปาวจรสฺส ธมฺมสฺส เหตุปจฺจเยน ปจฺจโย. (1) (สํขิตฺตํ.)}

\proseitem{39}{เหตุยา สตฺต, อารมฺมเณ นว ฯเปฯ อวิคเต เตรส. (สํขิตฺตํ. ยถา กุสลตฺติเก ปญฺหาวารํ, เอวํ วิตฺถาเรตพฺพํ.)}

\csromanmarkh{1. กุสลตฺติก}\csromanmarkhii{95. อรูปาวจรทุก}\csromanheadernomark{2}{1. \textbf{กุสลตฺติก} 95. อรูปาวจรทุก}
\csromanheader{2}{1. \textbf{ปฏิจฺจวาร ปจฺจยจตุกฺก}}
\proseitem{40}{กุสลํ อรูปาวจรํ ธมฺมํ ปฏิจฺจ กุสโล อรูปาวจโร ธมฺโม อุปฺปชฺชติ เหตุปจฺจยา. (1)}

\prose{อพฺยากตํ อรูปาวจรํ ธมฺมํ ปฏิจฺจ อพฺยากโต อรูปาวจโร ธมฺโม อุปฺปชฺชติ เหตุปจฺจยา. (1)}

\prose{กุสลํ นอรูปาวจรํ ธมฺมํ ปฏิจฺจ กุสโล นอรูปาวจโร ธมฺโม อุปฺปชฺชติ เหตุปจฺจยา. (รูปาวจร\-ทุก\-สทิสํ.)}

\csromanmarkh{1. กุสลตฺติก}\csromanmarkhii{96. ปริยาปนฺนทุก}\csromanheadernomark{2}{1. \textbf{กุสลตฺติก} 96. ปริยาปนฺนทุก}
\csromanheader{2}{1-7. ปฏิจฺจาทิวาร ปจฺจยจตุกฺกเหตุ}
\proseitem{41}{กุสลํ ปริยาปนฺนํ ธมฺมํ ปฏิจฺจ กุสโล ปริยาปนฺโน ธมฺโม อุปฺปชฺชติ เหตุปจฺจยา. ตีณิ.}

\prose{อกุสลํ ปริยาปนฺนํ ธมฺมํ ปฏิจฺจ อกุสโล ปริยาปนฺโน ธมฺโม อุปฺปชฺชติ เหตุปจฺจยา. ตีณิ.}

\csromanmarkh{1. กุสลตฺติก}\csromanmarkhii{96. ปริยาปนฺนทุก}\csromanheadernomark{2}{1. กุสลตฺติก 96. ปริยาปนฺนทุก}
\prose{\csromanfolio{613}อพฺยากตํ ปริยาปนฺนํ ธมฺมํ ปฏิจฺจ อพฺยากโต ปริยาปนฺโน ธมฺโม อุปฺปชฺชติ เหตุปจฺจยา. (1)}

\prose{กุสลํ ปริยาปนฺนญฺจ อพฺยากตํ ปริยาปนฺนญฺจ ธมฺมํ ปฏิจฺจ อพฺยากโต ปริยาปนฺโน ธมฺโม อุปฺปชฺชติ เหตุปจฺจยา. (1)}

\prose{อกุสลํ ปริยาปนฺนญฺจ อพฺยากตํ ปริยาปนฺนญฺจ ธมฺมํ ปฏิจฺจ อพฺยากโต ปริยาปนฺโน ธมฺโม อุปฺปชฺชติ เหตุปจฺจยา. (1) (สํขิตฺตํ.)}

\proseitem{42}{เหตุยา นว, อารมฺมเณ ตีณิ ฯเปฯ วิปาเก เอกํ ฯเปฯ อวิคเต นว. (สํขิตฺตํ. สหชาตวารมฺปิ ฯเปฯ สมฺปยุตฺตวารมฺปิ ปฏิจฺจ\-วาร\-สทิสํ.)}

\proseitem{43}{กุสโล ปริยาปนฺโน ธมฺโม กุสลสฺส ปริยาปนฺนสฺส ธมฺมสฺส เหตุปจฺจเยน ปจฺจโย. ตีณิ.}

\prose{อกุสโล ปริยาปนฺโน ธมฺโม อกุสลสฺส ปริยาปนฺนสฺส ธมฺมสฺส เหตุปจฺจเยน ปจฺจโย. ตีณิ.}

\prose{อพฺยากโต ปริยาปนฺโน ธมฺโม อพฺยากตสฺส ปริยาปนฺนสฺส ธมฺมสฺส เหตุปจฺจเยน ปจฺจโย. (1) (สํขิตฺตํ.)}

\proseitem{44}{เหตุยา สตฺต, อารมฺมเณ นว, อธิปติยา นว ฯเปฯ อวิคเต เตรส. (สํขิตฺตํ. ยถา กุสลตฺติเก ปญฺหาวารํ, เอวํ วิตฺถาเรตพฺพํ.)}

\prose{อปริยาปนฺนปทเหตุ อารมฺมณ}

\proseitem{45}{กุสลํ อปริยาปนฺนํ ธมฺมํ ปฏิจฺจ กุสโล อปริยาปนฺโน ธมฺโม อุปฺปชฺชติ เหตุปจฺจยา. (1)}

\prose{อพฺยากตํ อปริยาปนฺนํ ธมฺมํ ปฏิจฺจ อพฺยากโต อปริยาปนฺโน ธมฺโม อุปฺปชฺชติ เหตุปจฺจยา. (1) (สํขิตฺตํ.)}

\prose{เหตุยา เทฺว, อารมฺมเณ เทฺว ฯเปฯ อวิคเต เทฺว. (สํขิตฺตํ.) (สหชาตวารมฺปิ ฯเปฯ สมฺปยุตฺตวารมฺปิ ปฏิจฺจ\-วาร\-สทิสํ.)}

\proseitem{46}{กุสโล อปริยาปนฺโน ธมฺโม กุสลสฺส อปริยาปนฺนสฺส ธมฺมสฺส เหตุปจฺจเยน ปจฺจโย. (1)}

\prose{อพฺยากโต อปริยาปนฺโน ธมฺโม อพฺยากตสฺส อปริยาปนฺนสฺส ธมฺมสฺส เหตุปจฺจเยน ปจฺจโย. (1)}

\prose{\csromanfolio{614}อพฺยากโต อปริยาปนฺโน ธมฺโม อพฺยากตสฺส อปริยาปนฺนสฺส ธมฺมสฺส อารมฺมณ\-ปจฺจเยน ปจฺจโย.  อพฺยากโต อปริยาปนฺโน ธมฺโม กุสลสฺส อปริยาปนฺนสฺส ธมฺมสฺส อารมฺมณ\-ปจฺจเยน ปจฺจโย. (2)}

\proseitem{47}{เหตุยา เทฺว, อารมฺมเณ เทฺว, อธิปติยา ตีณิ, อนนฺตเร เทฺว ฯเปฯ อุปนิสฺสเย จตฺตาริ ฯเปฯ อวิคเต เทฺว. (สํขิตฺตํ.)}

\prose{(ยถา กุสลตฺติเก ปญฺหาวารํ, เอวํ วิตฺถาเรตพฺพํ.)}

\csromanmarkh{1. กุสลตฺติก}\csromanmarkhii{97. นิยฺยานิกทุก}\csromanheadernomark{2}{1. \textbf{กุสลตฺติก} 97. นิยฺยานิกทุก}
\csromanheader{2}{1. \textbf{ปฏิจฺจวาร ปจฺจยจตุกฺก}}
\proseitem{48}{กุสลํ นิยฺยานิกํ ธมฺมํ ปฏิจฺจ กุสโล นิยฺยานิโก ธมฺโม อุปฺปชฺชติ เหตุปจฺจยา. (1) (สํขิตฺตํ.)}

\prose{เหตุยา เอกํ, อารมฺมเณ เอกํ ฯเปฯ อวิคเต เอกํ. (สํขิตฺตํ.) (สหชาตวาเรปิ ฯเปฯ ปญฺหาวาเรปิ สพฺพตฺถ เอกํ.)}

\proseitem{49}{กุสลํ อนิยฺยานิกํ ธมฺมํ ปฏิจฺจ กุสโล อนิยฺยานิโก ธมฺโม อุปฺปชฺชติ เหตุปจฺจยา. ตีณิ.}

\prose{อกุสลํ อนิยฺยานิกํ ธมฺมํ ปฏิจฺจ อกุสโล อนิยฺยานิโก ธมฺโม อุปฺปชฺชติ เหตุปจฺจยา. ตีณิ.}

\prose{อพฺยากตํ อนิยฺยานิกํ ธมฺมํ ปฏิจฺจ อพฺยากโต อนิยฺยานิโก ธมฺโม อุปฺปชฺชติ เหตุปจฺจยา. (1)}

\prose{กุสลํ อนิยฺยานิกญฺจ อพฺยากตํ อนิยฺยานิกญฺจ ธมฺมํ ปฏิจฺจ อพฺยากโต อนิยฺยานิโก ธมฺโม อุปฺปชฺชติ เหตุปจฺจยา. (1)}

\prose{อกุสลํ อนิยฺยานิกญฺจ อพฺยากตํ อนิยฺยานิกญฺจ ธมฺมํ ปฏิจฺจ อพฺยากโต อนิยฺยานิโก ธมฺโม อุปฺปชฺชติ เหตุปจฺจยา. (1) (สํขิตฺตํ.)}

\proseitem{50}{เหตุยา นว, อารมฺมเณ ตีณิ ฯเปฯ อวิคเต นว. (สํขิตฺตํ.)}

\vspace{\tipitakaparskip}
\csromancenter{(สหชาตวารมฺปิ ฯเปฯ ปญฺหาวารมฺปิ วิตฺถาเรตพฺพํ.)}
\prose{\csromanfolio{615}1. กุสลตฺติก 99. สฺเอาตฺตรทุก \textbf{1. กุสลตฺติก 98. นิยตทุก}}

\csromanheader{2}{1. \textbf{ปฏิจฺจวาร ปจฺจยจตุกฺก}}
\proseitem{51}{กุสลํ นิยตํ ธมฺมํ ปฏิจฺจ กุสโล นิยโต ธมฺโม อุปฺปชฺชติ เหตุปจฺจยา. (1)}

\prose{อกุสลํ นิยตํ ธมฺมํ ปฏิจฺจ อกุสโล นิยโต ธมฺโม อุปฺปชฺชติ เหตุปจฺจยา. (1) (สํขิตฺตํ.)}

\prose{เหตุยา เทฺว, อารมฺมเณ เทฺว ฯเปฯ อวิคเต เทฺว. (สํขิตฺตํ.) (สหชาตวารมฺปิ ฯเปฯ ปญฺหาวารมฺปิ เอวํ วิตฺถาเรตพฺพํ.)}

\proseitem{52}{กุสลํ อนิยตํ ธมฺมํ ปฏิจฺจ กุสโล อนิยโต ธมฺโม อุปฺปชฺชติ เหตุปจฺจยา. (สํขิตฺตํ.)}

\csromanhangingbegin
\hangingheaditem{52}{เหตุยา นว, อารมฺมเณ ตีณิ ฯเปฯ อวิคเต นว. (สํขิตฺตํ.)}
\hangingbody{(สหชาตวารมฺปิ ฯเปฯ ปญฺหาวารมฺปิ เอวํ วิตฺถาเรตพฺพํ.)}
\csromanhangingend

\prose{1. กุสลตฺติก 99. สอุตฺตรทุก}

\csromanheader{2}{1-7. ปฏิจฺจาทิวาร ปจฺจยจตุกฺก}
\proseitem{53}{กุสลํ สอุตฺตรํ ธมฺมํ ปฏิจฺจ กุสโล สอุตฺตโร ธมฺโม อุปฺปชฺชติ เหตุปจฺจยา. ตีณิ.}

\prose{อกุสลํ สอุตฺตรํ ธมฺมํ ปฏิจฺจ อกุสโล สอุตฺตโร ธมฺโม อุปฺปชฺชติ เหตุปจฺจยา. ตีณิ.}

\prose{อพฺยากตํ สอุตฺตรํ ธมฺมํ ปฏิจฺจ อพฺยากโต สอุตฺตโร ธมฺโม อุปฺปชฺชติ เหตุปจฺจยา. (1)}

\prose{กุสลํ สอุตฺตรญฺจ อพฺยากตํ สอุตฺตรญฺจ ธมฺมํ ปฏิจฺจ อพฺยากโต สอุตฺตโร ธมฺโม อุปฺปชฺชติ เหตุปจฺจยา. (1)}

\prose{อกุสลํ สอุตฺตรญฺจ อพฺยากตํ สอุตฺตรญฺจ ธมฺมํ ปฏิจฺจ อพฺยากโต สอุตฺตโร ธมฺโม อุปฺปชฺชติ เหตุปจฺจยา. (1) (สํขิตฺตํ.)}

\proseitem{54}{\csromanfolio{616}เหตุยา นว, อารมฺมเณ ตีณิ ฯเปฯ อวิคเต นว. (สํขิตฺตํ. สหชาตวารมฺปิ ฯเปฯ สมฺปยุตฺตวารมฺปิ ปฏิจฺจ\-วาร\-สทิสํ.)}

\proseitem{55}{กุสโล สอุตฺตโร ธมฺโม กุสลสฺส สอุตฺตรสฺส ธมฺมสฺส เหตุปจฺจเยน ปจฺจโย. ตีณิ.}

\prose{อกุสโล สอุตฺตโร ธมฺโม อกุสลสฺส สอุตฺตรสฺส ธมฺมสฺส เหตุปจฺจเยน ปจฺจโย. ตีณิ.}

\prose{อพฺยากโต สอุตฺตโร ธมฺโม อพฺยากตสฺส สอุตฺตรสฺส ธมฺมสฺส เหตุปจฺจเยน ปจฺจโย. (1) (สํขิตฺตํ.)}

\proseitem{56}{เหตุยา สตฺต, อารมฺมเณ นว, อธิปติยา นว ฯเปฯ อวิคเต เตรส. (สํขิตฺตํ.)}

\prose{อนุตฺตรปทเหตุ อารมฺมณ}

\proseitem{57}{กุสลํ อนุตฺตรํ ธมฺมํ ปฏิจฺจ กุสโล อนุตฺตโร ธมฺโม อุปฺปชฺชติ เหตุปจฺจยา. (1)}

\prose{อพฺยากตํ อนุตฺตรํ ธมฺมํ ปฏิจฺจ อพฺยากโต อนุตฺตโร ธมฺโม อุปฺปชฺชติ เหตุปจฺจยา. (1) (สํขิตฺตํ.)}

\prose{เหตุยา เทฺว, อารมฺมเณ เทฺว ฯเปฯ อวิคเต เทฺว. (สํขิตฺตํ. สหชาตวารมฺปิ ฯเปฯ สมฺปยุตฺตวารมฺปิ ปฏิจฺจ\-วาร\-สทิสํ.)}

\proseitem{58}{กุสโล อนุตฺตโร ธมฺโม กุสลสฺส อนุตฺตรสฺส ธมฺมสฺส เหตุปจฺจเยน ปจฺจโย. (1)}

\prose{อพฺยากโต อนุตฺตโร ธมฺโม อพฺยากตสฺส อนุตฺตรสฺส ธมฺมสฺส เหตุปจฺจเยน ปจฺจโย. (1)}

\prose{อพฺยากโต อนุตฺตโร ธมฺโม อพฺยากตสฺส อนุตฺตรสฺส ธมฺมสฺส อารมฺมณ\-ปจฺจเยน ปจฺจโย. (1)}

\prose{อพฺยากโต อนุตฺตโร ธมฺโม กุสลสฺส อนุตฺตรสฺส ธมฺมสฺส อารมฺมณ\-ปจฺจเยน ปจฺจโย. (1) (สํขิตฺตํ.)}

\proseitem{59}{เหตุยา เทฺว, อารมฺมเณ เทฺว, อธิปติยา ตีณิ, อนนฺตเร เทฺว ฯเปฯ อวิคเต เทฺว. (สํขิตฺตํ.)}

\prose{(ยถา กุสลตฺติเก ปญฺหาวารํ, เอวํ วิตฺถาเรตพฺพํ.)}

\csromanmatikaanchor{mtk.150}
\csromantocmarkref{subparagraph}{100. สรณทุก}{mtk.150}
\bookmark[dest={mtk.150},level=5]{100. สรณทุก}
\csromanmarkh{1. กุสลตฺติก}\csromanmarkhii{100. สรณทุก 1. กุสลตฺติก 100. สรณทุก}\csromanheadernomark{6}{1. กุสลตฺติก 100. สรณทุก \textbf{1. กุสลตฺติก 100. สรณทุก}}
\csromanheader{2}{1-7. ปฏิจฺจาทิวาร ปจฺจยจตุกฺก}
\proseitem{60}{\csromanfolio{617}อกุสลํ สรณํ ธมฺมํ ปฏิจฺจ อกุสโล สรโณ ธมฺโม อุปฺปชฺชติ เหตุปจฺจยา. (สํขิตฺตํ.)}

\prose{เหตุยา เอกํ ฯเปฯ อวิคเต เอกํ. (สํขิตฺตํ.) (สหชาตวาเรปิ ฯเปฯ ปญฺหาวาเรปิ สพฺพตฺถ เอกํ.)}

\proseitem{61}{กุสลํ อรณํ ธมฺมํ ปฏิจฺจ กุสโล อรโณ ธมฺโม อุปฺปชฺชติ เหตุปจฺจยา. ตีณิ.}

\prose{อพฺยากตํ อรณํ ธมฺมํ ปฏิจฺจ อพฺยากโต อรโณ ธมฺโม อุปฺปชฺชติ เหตุปจฺจยา. (1)}

\prose{กุสลํ อรณญฺจ อพฺยากตํ อรณญฺจ ธมฺมํ ปฏิจฺจ อพฺยากโต อรโณ ธมฺโม อุปฺปชฺชติ เหตุปจฺจยา. (1)}

\prose{กุสลํ อรณํ ธมฺมํ ปฏิจฺจ กุสโล อรโณ ธมฺโม อุปฺปชฺชติ อารมฺมณ\-ปจฺจยา. (1)}

\prose{อพฺยากตํ อรณํ ธมฺมํ ปฏิจฺจ อพฺยากโต อรโณ ธมฺโม อุปฺปชฺชติ อารมฺมณ\-ปจฺจยา. (1) (สํขิตฺตํ.)}

\proseitem{62}{เหตุยา ปญฺจ, อารมฺมเณ เทฺว ฯเปฯ อวิคเต ปญฺจ. (สํขิตฺตํ. สหชาตวารมฺปิ ฯเปฯ สมฺปยุตฺตวารมฺปิ ปฏิจฺจ\-วาร\-สทิสํ.)}

\proseitem{63}{กุสโล อรโณ ธมฺโม กุสลสฺส อรณสฺส ธมฺมสฺส เหตุปจฺจเยน ปจฺจโย. ตีณิ.}

\prose{อพฺยากโต อรโณ ธมฺโม อพฺยากตสฺส อรณสฺส ธมฺมสฺส เหตุปจฺจเยน ปจฺจโย. (1) (สํขิตฺตํ.)}

\proseitemwithcloser{64}{\csromanfolio{618}เหตุยา จตฺตาริ, อารมฺมเณ จตฺตาริ ฯเปฯ อวิคเต สตฺต. (สํขิตฺตํ.) (ยถา กุสลตฺติเก ปญฺหาวารํ, เอวํ วิตฺถาเรตพฺพํ.)}{กุสล\-ตฺติก\-ปิฏฺฐิทุกํ นิฏฺฐิตํ.}
\csromanmatikaanchor{mtk.151}
\csromanmatikaanchor{mtk.152}
\csromantocmarknopage{subparagraph}{2. เวทนาตฺติก}{mtk.151}
\bookmark[dest={mtk.151},level=5]{2. เวทนาตฺติก}
\csromantocmarkref{subparagraph}{100. สรณทุก}{mtk.152}
\bookmark[dest={mtk.152},level=5]{100. สรณทุก}
\csromanmarkh{2. เวทนาตฺติก}\csromanmarkhii{100. สรณทุก}\csromanheadernomark{6}{2. \textbf{เวทนาตฺติก} 100. สรณทุก}
\csromanheader{2}{1-7. ปฏิจฺจาทิวาร ปจฺจยจตุกฺก}
\proseitem{65}{สุขาย เวทนาย สมฺปยุตฺตํ สรณํ ธมฺมํ ปฏิจฺจ สุขาย เวทนาย สมฺปยุตฺโต สรโณ ธมฺโม อุปฺปชฺชติ เหตุปจฺจยา. (1)}

\prose{ทุกฺขาย เวทนาย สมฺปยุตฺตํ สรณํ ธมฺมํ ปฏิจฺจ ทุกฺขาย เวทนาย สมฺปยุตฺโต สรโณ ธมฺโม อุปฺปชฺชติ เหตุปจฺจยา. (1)}

\prose{อทุกฺขม\-สุขาย เวทนาย สมฺปยุตฺตํ สรณํ ธมฺมํ ปฏิจฺจ อทุกฺขม\-สุขาย เวทนาย สมฺปยุตฺโต สรโณ ธมฺโม อุปฺปชฺชติ เหตุปจฺจยา. (1) (สํขิตฺตํ.)}

\proseitem{66}{เหตุยา ตีณิ ฯเปฯ อวิคเต ตีณิ. (สํขิตฺตํ.) (สหชาตวารมฺปิ ฯเปฯ สมฺปยุตฺตวารมฺปิ ปฏิจฺจ\-วาร\-สทิสํ.)}

\proseitem{67}{สุขาย เวทนาย สมฺปยุตฺโต สรโณ ธมฺโม สุขาย เวทนาย สมฺปยุตฺตสฺส สรณสฺส ธมฺมสฺส เหตุปจฺจเยน ปจฺจโย. (1)}

\prose{ทุกฺขาย เวทนาย สมฺปยุตฺโต สรโณ ธมฺโม ทุกฺขาย เวทนาย สมฺปยุตฺตสฺส สรณสฺส ธมฺมสฺส เหตุปจฺจเยน ปจฺจโย. (1)}

\prose{อทุกฺขม\-สุขาย เวทนาย สมฺปยุตฺโต สรโณ ธมฺโม อทุกฺขม\-สุขาย เวทนาย สมฺปยุตฺตสฺส สรณสฺส ธมฺมสฺส เหตุปจฺจเยน ปจฺจโย. (1) (สํขิตฺตํ.)}

\csromanmarkh{3. วิปากตฺติก}\csromanmarkhii{100. สรณทุก}\csromanheadernomark{2}{3. วิปากตฺติก 100. สรณทุก}
\proseitem{68}{\csromanfolio{619}เหตุยา ตีณิ, อารมฺมเณ นว ฯเปฯ อวิคเต ตีณิ. (สํขิตฺตํ. ยถา กุสลตฺติเก ปญฺหาวารํ, เอวํ วิตฺถาเรตพฺพํ.)}

\prose{อรณปทเหตุ}

\proseitem{69}{สุขาย เวทนาย สมฺปยุตฺตํ อรณํ ธมฺมํ ปฏิจฺจ สุขาย เวทนาย สมฺปยุตฺโต อรโณ ธมฺโม อุปฺปชฺชติ เหตุปจฺจยา.}

\prose{อทุกฺขม\-สุขาย เวทนาย สมฺปยุตฺตํ อรณํ ธมฺมํ ปฏิจฺจ อทุกฺขม\-สุขาย เวทนาย สมฺปยุตฺโต อรโณ ธมฺโม อุปฺปชฺชติ เหตุปจฺจยา. (สํขิตฺตํ.)}

\prose{เหตุยา เทฺว, อารมฺมเณ ตีณิ ฯเปฯ อวิคเต ตีณิ. (สํขิตฺตํ.) (สหชาตวารมฺปิ ฯเปฯ สมฺปยุตฺตวารมฺปิ ปฏิจฺจ\-วาร\-สทิสํ.)}

\proseitem{70}{สุขาย เวทนาย สมฺปยุตฺโต อรโณ ธมฺโม สุขาย เวทนาย สมฺปยุตฺตสฺส อรณสฺส ธมฺมสฺส เหตุปจฺจเยน ปจฺจโย. (1)}

\prose{อทุกฺขม\-สุขาย เวทนาย สมฺปยุตฺโต อรโณ ธมฺโม อทุกฺขม\-สุขาย เวทนาย สมฺปยุตฺตสฺส อรณสฺส ธมฺมสฺส เหตุปจฺจเยน ปจฺจโย. (1) (สํขิตฺตํ.)}

\csromanhangingbegin
\hangingheaditem{71}{เหตุยา เทฺว, อารมฺมเณ ฉ ฯเปฯ อวิคเต ตีณิ. (สํขิตฺตํ.)}
\hangingbody{(ยถา กุสลตฺติเก ปญฺหาวารํ, เอวํ วิตฺถาเรตพฺพํ.)}
\csromanhangingend

\csromanmarkh{3. วิปากตฺติก}\csromanmarkhii{100. สรณทุก}\csromanheadernomark{2}{3. วิปากตฺติก 100. สรณทุก}
\csromanheader{2}{1-7. ปฏิจฺจาทิวาร ปจฺจยจตุกฺก}
\proseitem{72}{วิปาก\-ธมฺม\-ธมฺมํ สรณํ ธมฺมํ ปฏิจฺจ วิปาก\-ธมฺม\-ธมฺโม สรโณ ธมฺโม อุปฺปชฺชติ เหตุปจฺจยา. (สํขิตฺตํ.)}

\prose{เหตุยา เอกํ ฯเปฯ อวิคเต เอกํ. (สํขิตฺตํ.)}

\vspace{\tipitakaparskip}
\csromancenter{(สหชาตวาเรปิ ฯเปฯ ปญฺหาวาเรปิ สพฺพตฺถ เอกํ.)}
\proseitem{73}{\csromanfolio{620}วิปากํ อรณํ ธมฺมํ ปฏิจฺจ วิปาโก อรโณ ธมฺโม อุปฺปชฺชติ เหตุปจฺจยา. ตีณิ.}

\prose{วิปาก\-ธมฺม\-ธมฺมํ อรณํ ธมฺมํ ปฏิจฺจ วิปาก\-ธมฺม\-ธมฺโม อรโณ ธมฺโม อุปฺปชฺชติ เหตุปจฺจยา. ตีณิ.}

\prose{เนว\-วิปาก\-น\-วิปากธมฺม\-ธมฺมํ อรณํ ธมฺมํ ปฏิจฺจ เนว\-วิปาก\-น\-วิปาก\-ธมฺม\-ธมฺโม อรโณ ธมฺโม อุปฺปชฺชติ เหตุปจฺจยา. ตีณิ.}

\prose{วิปากํ อรณญฺจ เนว\-วิปาก\-น\-วิปากธมฺม\-ธมฺมํ อรณญฺจ ธมฺมํ ปฏิจฺจ วิปาโก อรโณ ธมฺโม อุปฺปชฺชติ เหตุปจฺจยา. ตีณิ.}

\prose{วิปาก\-ธมฺม\-ธมฺมํ อรณญฺจ เนว\-วิปาก\-น\-วิปากธมฺม\-ธมฺมํ อรณญฺจ ธมฺมํ ปฏิจฺจ เนว\-วิปาก\-น\-วิปาก\-ธมฺม\-ธมฺโม อรโณ ธมฺโม อุปฺปชฺชติ เหตุปจฺจยา. (1) (สํขิตฺตํ.)}

\proseitem{74}{เหตุยา เตรส, อารมฺมเณ ปญฺจ ฯเปฯ อวิคเต เตรส. (สํขิตฺตํ. สหชาตวารมฺปิ ฯเปฯ สมฺปยุตฺตวารมฺปิ ปฏิจฺจ\-วาร\-สทิสํ.)}

\proseitem{75}{วิปาโก อรโณ ธมฺโม วิปากสฺส อรณสฺส ธมฺมสฺส เหตุปจฺจเยน ปจฺจโย. ตีณิ.}

\prose{วิปาก\-ธมฺม\-ธมฺโม อรโณ ธมฺโม วิปาก\-ธมฺม\-ธมฺมสฺส อรณสฺส ธมฺมสฺส เหตุปจฺจเยน ปจฺจโย. ตีณิ.}

\prose{เนว\-วิปาก\-น\-วิปาก\-ธมฺม\-ธมฺโม อรโณ ธมฺโม เนว\-วิปาก\-น\-วิปากธมฺม\-ธมฺมสฺส อรณสฺส ธมฺมสฺส เหตุปจฺจเยน ปจฺจโย. (1) (สํขิตฺตํ.)}

\csromanhangingbegin
\hangingheaditem{76}{เหตุยา สตฺต, อรมฺมเณ นว ฯเปฯ อวิคเต เตรส. (สํขิตฺตํ.)}
\hangingbody{(ยถา กุสลตฺติเก ปญฺหาวารํ, เอวํ วิตฺถาเรตพฺพํ.)}
\csromanhangingend

\vspace{\tipitakaparskip}
\csromancenter{อุปาทินฺนตฺติก 100. สรณทุก 4. \textbf{ อุปาทินฺนตฺติก 100.} สรณทุก}
\csromanheader{2}{1-7. ปฏิจฺจาทิวาร ปจฺจยจตุกฺกเหตุ}
\proseitem{77}{\csromanfolio{621}อนุปาทินฺนุ\-ปาทานิยํ สรณํ ธมฺมํ ปฏิจฺจ อนุปาทินฺนุ\-ปาทานิโย สรโณ ธมฺโม อุปฺปชฺชติ เหตุปจฺจยา. (สํขิตฺตํ.)}

\prose{เหตุยา เอกํ ฯเปฯ อวิคเต เอกํ. (สํขิตฺตํ.) (สหชาตวาเรปิ ฯเปฯ ปญฺหาวาเรปิ สพฺพตฺถ เอกํ.)}

\proseitem{78}{อุปาทินฺนุ\-ปาทานิยํ อรณํ ธมฺมํ ปฏิจฺจ อุปาทินฺนุ\-ปาทานิโย อรโณ ธมฺโม อุปฺปชฺชติ เหตุปจฺจยา. ตีณิ.}

\prose{อนุปาทินฺนุ\-ปาทานิยํ อรณํ ธมฺมํ ปฏิจฺจ อนุปาทินฺนุ\-ปาทานิโย อรโณ ธมฺโม อุปฺปชฺชติ เหตุปจฺจยา. (1)}

\prose{อนุปาทินฺน\-อนุปาทานิยํ อรณํ ธมฺมํ ปฏิจฺจ อนุปาทินฺน\-อนุปาทานิโย อรโณ ธมฺโม อุปฺปชฺชติ เหตุปจฺจยา. ตีณิ.}

\prose{อุปาทินฺนุ\-ปาทานิยํ อรณญฺจ อนุปาทินฺนุ\-ปาทานิยํ อรณญฺจ ธมฺมํ ปฏิจฺจ อนุปาทินฺนุ\-ปาทานิโย อรโณ ธมฺโม อุปฺปชฺชติ เหตุปจฺจยา. (1)}

\prose{อนุปาทินฺนุ\-ปาทานิยํ อรณญฺจ อนุปาทินฺน\-อนุปาทานิยํ อรณญฺจ ธมฺมํ ปฏิจฺจ อนุปาทินฺนุ\-ปาทานิโย อรโณ ธมฺโม อุปฺปชฺชติ เหตุปจฺจยา. (1) (สํขิตฺตํ.)}

\proseitem{79}{เหตุยา นว, อารมฺมเณ ตีณิ ฯเปฯ อวิคเต นว. (สํขิตฺตํ.) (สหชาตวารมฺปิ ฯเปฯ สมฺปยุตฺตวารมฺปิ ปฏิจฺจ\-วาร\-สทิสํ.)}

\proseitem{80}{อุปาทินฺนุ\-ปาทานิโย อรโณ ธมฺโม อุปาทินฺนุ\-ปาทานิยสฺส อรณสฺส ธมฺมสฺส เหตุปจฺจเยน ปจฺจโย. ตีณิ.}

\prose{อนุปาทินฺนุ\-ปาทานิโย อรโณ ธมฺโม อนุปาทินฺนุ\-ปาทานิยสฺส อรณสฺส ธมฺมสฺส เหตุปจฺจเยน ปจฺจโย. (1)}

\prose{อนุปาทินฺน\-อนุปาทานิโย อรโณ ธมฺโม อนุปาทินฺน\-อนุปาทานิยสฺส อรณสฺส ธมฺมสฺส เหตุปจฺจเยน ปจฺจโย. ตีณิ. (สํขิตฺตํ.)}

\proseitem{81}{\csromanfolio{622}เหตุยา สตฺต, อารมฺมเณ ฉ ฯเปฯ อวิคเต เตวีส. (สํขิตฺตํ. ยถา กุสลตฺติเก ปญฺหาวารํ, เอวํ วิตฺถาเรตพฺพํ.)}

\csromanmarkh{5. สํกิลิฏฺฐ\-ตฺติก}\csromanmarkhii{100. สรณทุก}\csromanheadernomark{2}{5. สํกิลิฏฺฐ\-ตฺติก 100. สรณทุก}
\csromanheader{2}{1-7. ปฏิจฺจาทิวาร ปจฺจยจตุกฺกเหตุ}
\proseitem{82}{สํกิลิฏฺฐ\-สํกิเลสิกํ สรณํ ธมฺมํ ปฏิจฺจ สํกิลิฏฺฐ\-สํกิเลสิโก สรโณ ธมฺโม อุปฺปชฺชติ เหตุปจฺจยา. (สํขิตฺตํ.)}

\prose{เหตุยา เอกํ ฯเปฯ อวิคเต เอกํ. (สํขิตฺตํ.) (สหชาตวาเรปิ ฯเปฯ ปญฺหาวาเรปิ สพฺพตฺถ เอกํ.)}

\proseitem{83}{อสํกิลิฏฺฐ\-สํกิเลสิกํ อรณํ ธมฺมํ ปฏิจฺจ อสํกิลิฏฺฐ\-สํกิเลสิโก อรโณ ธมฺโม อุปฺปชฺชติ เหตุปจฺจยา. (1)}

\prose{อสํกิลิฏฺฐ\-อสํกิเลสิกํ อรณํ ธมฺมํ ปฏิจฺจ อสํกิลิฏฺฐ\-อสํกิเลสิโก อรโณ ธมฺโม อุปฺปชฺชติ เหตุปจฺจยา. ตีณิ.}

\prose{อสํกิลิฏฺฐ\-สํกิเลสิกํ อรณญฺจ อสํกิลิฏฺฐ\-อสํกิเลสิกํ อรณญฺจ ธมฺมํ ปฏิจฺจ อสํกิลิฏฺฐ\-สํกิเลสิโก อรโณ ธมฺโม อุปฺปชฺชติ เหตุปจฺจยา. (1) (สํขิตฺตํ.)}

\proseitem{84}{เหตุยา ปญฺจ, อารมฺมเณ เทฺว ฯเปฯ อวิคเต ปญฺจ. (สํขิตฺตํ.) (สหชาตวารมฺปิ ฯเปฯ สมฺปยุตฺตวารมฺปิ ปฏิจฺจ\-วาร\-สทิสํ.)}

\proseitem{85}{อสํกิลิฏฺฐ\-สํกิเลสิโก อรโณ ธมฺโม อสํกิลิฏฺฐ\-สํกิเลสิกสฺส อรณสฺส ธมฺมสฺส เหตุปจฺจเยน ปจฺจโย. (1)}

\prose{อสํกิลิฏฺฐ\-อสํกิเลสิโก อรโณ ธมฺโม อสํกิลิฏฺฐ\-อสํกิเลสิกสฺส อรณสฺส ธมฺมสฺส เหตุปจฺจเยน ปจฺจโย. ตีณิ. (สํขิตฺตํ.)}

\csromanmarkh{6. วิตกฺกตฺติก}\csromanmarkhii{100. สรณทุก}\csromanheadernomark{2}{6. วิตกฺกตฺติก 100. สรณทุก}
\csromanhangingbegin
\hangingheaditem{86}{\csromanfolio{623}เหตุยา จตฺตาริ, อารมฺมเณ ตีณิ ฯเปฯ อวิคเต สตฺต. (สํขิตฺตํ.)}
\hangingbody{(ยถา กุสลตฺติเก ปญฺหาวารํ, เอวํ วิตฺถาเรตพฺพํ.)}
\csromanhangingend

\csromanmarkh{6. วิตกฺกตฺติก}\csromanmarkhii{100. สรณทุก}\csromanheadernomark{2}{6. วิตกฺกตฺติก 100. สรณทุก}
\prose{1. ปฏิจฺจวาร ปจฺจยจตุกฺกเหตุ}

\proseitem{87}{สวิตกฺก\-สวิจารํ สรณํ ธมฺมํ ปฏิจฺจ สวิตกฺก\-สวิจาโร สรโณ ธมฺโม อุปฺปชฺชติ เหตุปจฺจยา.  สวิตกฺก\-สวิจารํ สรณํ ธมฺมํ ปฏิจฺจ อวิตกฺก\-วิจารมตฺโต สรโณ ธมฺโม อุปฺปชฺชติ เหตุปจฺจยา.  สวิตกฺก\-สวิจารํ สรณํ ธมฺมํ ปฏิจฺจ สวิตกฺก\-สวิจาโร สรโณ จ อวิตกฺก\-วิจารมตฺโต สรโณ จ ธมฺมา อุปฺปชฺชนฺติ เหตุปจฺจยา. (3)}

\prose{อวิตกฺก\-วิจารมตฺตํ สรณํ ธมฺมํ ปฏิจฺจ สวิตกฺก\-สวิจาโร สรโณ ธมฺโม อุปฺปชฺชติ เหตุปจฺจยา. (1)}

\prose{สวิตกฺก\-สวิจารํ สรณญฺจ อวิตกฺก\-วิจารมตฺตํ สรณญฺจ ธมฺมํ ปฏิจฺจ สวิตกฺก\-สวิจาโร สรโณ ธมฺโม อุปฺปชฺชติ เหตุปจฺจยา. (1) (สํขิตฺตํ.)}

\proseitem{88}{เหตุยา ปญฺจ ฯเปฯ อวิคเต ปญฺจ. (สํขิตฺตํ.) (สหชาตวาเรปิ ฯเปฯ ปญฺหาวาเรปิ สพฺพตฺถ วิตฺถาโร.)}

\proseitem{89}{สวิตกฺก\-สวิจารํ อรณํ ธมฺมํ ปฏิจฺจ สวิตกฺก\-สวิจาโร อรโณ ธมฺโม อุปฺปชฺชติ เหตุปจฺจยา. สตฺต.}

\prose{อวิตกฺก\-วิจารมตฺตํ อรณํ ธมฺมํ ปฏิจฺจ อวิตกฺก\-วิจารมตฺโต อรโณ ธมฺโม อุปฺปชฺชติ เหตุปจฺจยา. ปญฺจ.}

\prose{อวิตกฺก\-อวิจารํ อรณํ ธมฺมํ ปฏิจฺจ อวิตกฺก\-อวิจาโร อรโณ ธมฺโม อุปฺปชฺชติ เหตุปจฺจยา. (สํขิตฺตํ.)}

\proseitem{90}{\csromanfolio{624}เหตุยา สตฺตติํส ฯเปฯ อวิคเต สตฺตติํส. (สํขิตฺตํ. สหชาตวารมฺปิ ฯเปฯ ปญฺหาวารมฺปิ วิตฺถาเรตพฺพํ.)}

\csromanmarkh{7. ปีติตฺติก}\csromanmarkhii{100. สรณทุก}\csromanheadernomark{2}{7. ปีติตฺติก 100. สรณทุก}
\csromanheader{2}{1-7. ปฏิจฺจาทิวาร ปจฺจยจตุกฺก}
\proseitem{91}{ปีติสหคตํ สรณํ ธมฺมํ ปฏิจฺจ ปีติสหคโต สรโณ ธมฺโม อุปฺปชฺชติ เหตุปจฺจยา. ตีณิ.}

\prose{สุขสหคตํ สรณํ ธมฺมํ ปฏิจฺจ สุขสหคโต สรโณ ธมฺโม อุปฺปชฺชติ เหตุปจฺจยา. ตีณิ.}

\prose{อุเปกฺขา\-สหคตํ สรณํ ธมฺมํ ปฏิจฺจ อุเปกฺขา\-สหคโต สรโณ ธมฺโม อุปฺปชฺชติ เหตุปจฺจยา. (1)}

\prose{ปีติสหคตํ สรณญฺจ สุขสหคตํ สรณญฺจ ธมฺมํ ปฏิจฺจ ปีติสหคโต สรโณ ธมฺโม อุปฺปชฺชติ เหตุปจฺจยา. ตีณิ. (สํขิตฺตํ.)}

\prose{เหตุยา ทส ฯเปฯ อวิคเต ทส. (สํขิตฺตํ.) (สหชาตวารมฺปิ ฯเปฯ สมฺปยุตฺตวารมฺปิ ปฏิจฺจ\-วาร\-สทิสํ.)}

\proseitem{92}{ปีติสหคโต สรโณ ธมฺโม ปีติสหคตสฺส สรณสฺส ธมฺมสฺส เหตุปจฺจเยน ปจฺจโย. ตีณิ.}

\prose{สุขสหคโต สรโณ ธมฺโม สุขสหคตสฺส สรณสฺส ธมฺมสฺส เหตุปจฺจเยน ปจฺจโย. ตีณิ.}

\prose{อุเปกฺขา\-สหคโต สรโณ ธมฺโม อุเปกฺขา\-สหคตสฺส สรณสฺส ธมฺมสฺส เหตุปจฺจเยน ปจฺจโย. (1)}

\prose{ปีติสหคโต สรโณ จ สุขสหคโต สรโณ จ ธมฺมา ปีติสหคตสฺส สรณสฺส ธมฺมสฺส เหตุปจฺจเยน ปจฺจโย. ตีณิ. (สํขิตฺตํ.)}

\csromanmarkh{9. ทสฺสน\-เหตุ\-ตฺติก}\csromanmarkhii{100. สรณทุก}\csromanheadernomark{2}{9. ทสฺสน\-เหตุ\-ตฺติก 100. สรณทุก}
\proseitem{93}{\csromanfolio{625}เหตุยา ทส, อารมฺมเณ โสฬส ฯเปฯ อวิคเต ทส. (สํขิตฺตํ.) (ยถา กุสลตฺติเก ปญฺหาวารํ, เอวํ วิตฺถาเรตพฺพํ.)}

\proseitem{94}{ปีติสหคตํ อรณํ ธมฺมํ ปฏิจฺจ ปีติสหคโต อรโณ ธมฺโม อุปฺปชฺชติ เหตุปจฺจยา. (สํขิตฺตํ.)}

\prose{เหตุยา ทส ฯเปฯ อวิคเต ทส. (สํขิตฺตํ.)}

\vspace{\tipitakaparskip}
\csromancenter{(สหชาตวาเรปิ ฯเปฯ ปญฺหาวาเรปิ ตตฺตกาว ปญฺหา.)}
\csromanmarkh{8. ทสฺสนตฺติก}\csromanmarkhii{100. สรณทุก}\csromanheadernomark{2}{8. \textbf{ทสฺสนตฺติก} 100. สรณทุก}
\csromanheader{2}{1. \textbf{ปฏิจฺจวาร ปจฺจยจตุกฺก}}
\proseitem{95}{ทสฺสเนน ปหาตพฺพํ สรณํ ธมฺมํ ปฏิจฺจ ทสฺสเนน ปหาตพฺโพ สรโณ ธมฺโม อุปฺปชฺชติ เหตุปจฺจยา. (1)}

\prose{ภาวนาย ปหาตพฺพํ สรณํ ธมฺมํ ปฏิจฺจ ภาวนาย ปหาตพฺโพ สรโณ ธมฺโม อุปฺปชฺชติ เหตุปจฺจยา. (1) (สํขิตฺตํ.)}

\prose{เหตุยา เทฺว ฯเปฯ อวิคเต เทฺว. (สํขิตฺตํ.) (สหชาตวาเรปิ ฯเปฯ ปญฺหาวาเรปิ สพฺพตฺถ วิตฺถาโร.)}

\proseitem{96}{เนวทสฺสเนน นภาวนาย ปหาตพฺพํ อรณํ ธมฺมํ ปฏิจฺจ เนวทสฺสเนน นภาวนาย ปหาตพฺโพ อรโณ ธมฺโม อุปฺปชฺชติ เหตุปจฺจยา. (สํขิตฺตํ.)}

\prose{เหตุยา เอกํ ฯเปฯ อวิคเต เอกํ. (สํขิตฺตํ.)}

\vspace{\tipitakaparskip}
\csromancenter{(สหชาตวาเรปิ ฯเปฯ ปญฺหาวาเรปิ สพฺพตฺถ เอกํ.)}
\csromanmarkh{9. ทสฺสน\-เหตุ\-ตฺติก}\csromanmarkhii{100. สรณทุก}\csromanheadernomark{2}{9. ทสฺสน\-เหตุ\-ตฺติก 100. สรณทุก}
\prose{1. ปฏิจฺจวารเหตุ}

\proseitem{97}{ทสฺสเนน ปหาตพฺพ\-เหตุกํ สรณํ ธมฺมํ ปฏิจฺจ ฯเปฯ. ภาวนาย ปหาตพฺพ\-เหตุกํ สรณํ ธมฺมํ ปฏิจฺจ ฯเปฯ. เนวทสฺสเนน นภาวนาย \csromanfolio{626}ปหาตพฺพ\-เหตุกํ สรณํ ธมฺมํ ปฏิจฺจ เนวทสฺสเนน นภาวนาย ปหาตพฺพ\-เหตุโก สรโณ ธมฺโม อุปฺปชฺชติ เหตุปจฺจยา. (สํขิตฺตํ.)}

\csromanmarkh{10. อาจยคามิตฺติก}\csromanmarkhii{100. สรณทุก}\csromanheadernomark{2}{10. \textbf{อาจยคามิตฺติ}ก 100. สรณทุก}
\prose{1-7. ปฏิจฺจาทิวารเหตุ อารมฺมณ}

\proseitem{97}{อาจยคามิํ สรณํ ธมฺมํ ปฏิจฺจ อาจยคามี สรโณ ธมฺโม อุปฺปชฺชติ เหตุปจฺจยา. (สํขิตฺตํ.)}

\prose{เหตุยา เอกํ ฯเปฯ อวิคเต เอกํ. (สํขิตฺตํ.) (สหชาตวาเรปิ ฯเปฯ ปญฺหาวาเรปิ สพฺพตฺถ เอกํ.)}

\proseitem{99}{อาจยคามิํ อรณํ ธมฺมํ ปฏิจฺจ อาจยคามี อรโณ ธมฺโม อุปฺปชฺชติ เหตุปจฺจยา. ตีณิ.}

\prose{อปจยคามิํ อรณํ ธมฺมํ ปฏิจฺจ อปจยคามี อรโณ ธมฺโม อุปฺปชฺชติ เหตุปจฺจยา. ตีณิ.}

\prose{เนวา\-จยคามิ\-นา\-ปจยคามิํ อรณํ ธมฺมํ ปฏิจฺจ เนวา\-จยคามิ\-นา\-ปจยคามี อรโณ ธมฺโม อุปฺปชฺชติ เหตุปจฺจยา. (1)}

\prose{อาจยคามิํ อรณญฺจ เนวา\-จยคามิ\-นา\-ปจยคามิํ อรณญฺจ ธมฺมํ ปฏิจฺจ เนวา\-จยคามิ\-นา\-ปจยคามี อรโณ ธมฺโม อุปฺปชฺชติ เหตุปจฺจยา. (1)}

\prose{อปจยคามิํ อรณญฺจ เนวา\-จยคามิ\-นา\-ปจยคามิํ อรณญฺจ ธมฺมํ ปฏิจฺจ เนวา\-จยคามิ\-นา\-ปจยคามี อรโณ ธมฺโม อุปฺปชฺชติ เหตุปจฺจยา. (1) (สํขิตฺตํ.)}

\proseitem{97}{เหตุยา นว, อารมฺมเณ ตีณิ, อธิปติยา นว ฯเปฯ อวิคเต นว. (สํขิตฺตํ. สหชาตวารมฺปิ ฯเปฯ สมฺปยุตฺตวารมฺปิ ปฏิจฺจ\-วาร\-สทิสํ.)}

\proseitem{101}{อาจยคามี อรโณ ธมฺโม อาจยคามิสฺส อรณสฺส ธมฺมสฺส เหตุปจฺจเยน ปจฺจโย. ตีณิ.}

\csromanmarkh{11. เสกฺขตฺติก}\csromanmarkhii{100. สรณทุก}\csromanheadernomark{2}{11. เสกฺขตฺติก 100. สรณทุก}
\prose{\csromanfolio{627}อปจยคามี อรโณ ธมฺโม อปจยคามิสฺส อรณสฺส ธมฺมสฺส เหตุปจฺจเยน ปจฺจโย. ตีณิ.}

\prose{เนวา\-จยคามิ\-นา\-ปจยคามี อรโณ ธมฺโม เนวา\-จยคามิ\-นา\-ปจยคามิสฺส อรณสฺส ธมฺมสฺส เหตุปจฺจเยน ปจฺจโย. (1)}

\prose{อาจยคามี อรโณ ธมฺโม อาจยคามิสฺส อรณสฺส ธมฺมสฺส อารมฺมณ\-ปจฺจเยน ปจฺจโย. (สํขิตฺตํ.)}

\proseitem{102}{เหตุยา สตฺต, อารมฺมเณ สตฺต, อธิปติยา ทส, อนนฺตเร ฉ ฯเปฯ อวิคเต เตรส. (สํขิตฺตํ.)}

\prose{(ยถา กุสลตฺติเก ปญฺหาวารํ, เอวํ วิตฺถาเรตพฺพํ.)}

\csromanmarkh{11. เสกฺขตฺติก}\csromanmarkhii{100. สรณทุก}\csromanheadernomark{2}{11. เสกฺขตฺติก 100. สรณทุก}
\prose{1-7. ปฏิจฺจาทิวารเหตุ}

\proseitem{103}{เนว\-เสกฺข\-นาเสกฺขํ สรณํ ธมฺมํ ปฏิจฺจ เนว\-เสกฺข\-นา\-เสกฺโข สรโณ ธมฺโม อุปฺปชฺชติ เหตุปจฺจยา. (สํขิตฺตํ.)}

\prose{เหตุยา เอกํ ฯเปฯ อวิคเต เอกํ. (สํขิตฺตํ.) (สหชาตวาเรปิ ฯเปฯ ปญฺหาวาเรปิ สพฺพตฺถ เอกํ.)}

\proseitem{104}{เสกฺขํ อรณํ ธมฺมํ ปฏิจฺจ เสกฺโข อรโณ ธมฺโม อุปฺปชฺชติ เหตุปจฺจยา.  เสกฺขํ อรณํ ธมฺมํ ปฏิจฺจ เนว\-เสกฺข\-นา\-เสกฺโข อรโณ ธมฺโม อุปฺปชฺชติ เหตุปจฺจยา.  เสกฺขํ อรณํ ธมฺมํ ปฏิจฺจ เสกฺโข อรโณ จ เนว\-เสกฺข\-นา\-เสกฺโข อรโณ จ ธมฺมา อุปฺปชฺชนฺติ เหตุปจฺจยา. (3)}

\prose{อเสกฺขํ อรณํ ธมฺมํ ปฏิจฺจ อเสกฺโข อรโณ ธมฺโม อุปฺปชฺชติ เหตุปจฺจยา. ตีณิ.}

\prose{เนว\-เสกฺข\-นาเสกฺขํ อรณํ ธมฺมํ ปฏิจฺจ เนว\-เสกฺข\-นา\-เสกฺโข อรโณ ธมฺโม อุปฺปชฺชติ เหตุปจฺจยา. (1)}

\prose{เสกฺขํ อรณญฺจ เนว\-เสกฺข\-นาเสกฺขํ อรณญฺจ ธมฺมํ ฯเปฯ. เหตุปจฺจยา.}

\prose{\csromanfolio{628}อเสกฺขํ อรณญฺจ เนว\-เสกฺข\-นาเสกฺขํ อรณญฺจ ธมฺมํ ปฏิจฺจ เนว\-เสกฺข\-นา\-เสกฺโข อรโณ ธมฺโม อุปฺปชฺชติ เหตุปจฺจยา. (1) (สํขิตฺตํ.)}

\proseitem{105}{เหตุยา นว, อารมฺมเณ ตีณิ ฯเปฯ อวิคเต นว. (สํขิตฺตํ.) (สหชาตวารมฺปิ ฯเปฯ สมฺปยุตฺตวารมฺปิ ปฏิจฺจ\-วาร\-สทิสํ.)}

\proseitem{106}{เสกฺโข อรโณ ธมฺโม เสกฺขสฺส อรณสฺส ธมฺมสฺส เหตุปจฺจเยน ปจฺจโย. ตีณิ.}

\prose{อเสกฺโข อรโณ ธมฺโม อเสกฺขสฺส อรณสฺส ธมฺมสฺส เหตุปจฺจเยน ปจฺจโย. ตีณิ.}

\prose{เนว\-เสกฺข\-นา\-เสกฺโข อรโณ ธมฺโม เนว\-เสกฺข\-นาเสกฺขสฺส อรณสฺส ธมฺมสฺส เหตุปจฺจเยน ปจฺจโย. (1) (สํขิตฺตํ.)}

\proseitem{107}{เหตุยา สตฺต, อารมฺมเณ ปญฺจ, อธิปติยา นว, อนนฺตเร อฏฺฐ ฯเปฯ อวิคเต เตรส. (สํขิตฺตํ.)}

\prose{(ยถา กุสลตฺติเก ปญฺหาวารํ, เอวํ วิตฺถาเรตพฺพํ.)}

\csromanmarkh{12. ปริตฺตตฺติก}\csromanmarkhii{100. สรณทุก}\csromanheadernomark{2}{12. \textbf{ปริตฺตตฺติก} 100. สรณทุก}
\prose{1. ปฏิจฺจวารเหตุ}

\proseitem{108}{ปริตฺตํ สรณํ ธมฺมํ ปฏิจฺจ ปริตฺโต สรโณ ธมฺโม อุปฺปชฺชติ เหตุปจฺจยา. (สํขิตฺตํ.)}

\prose{เหตุยา เอกํ ฯเปฯ อวิคเต เอกํ. (สํขิตฺตํ.) (สหชาตวาเรปิ ฯเปฯ ปญฺหาวาเรปิ สพฺพตฺถ เอกํ.)}

\proseitem{109}{ปริตฺตํ อรณํ ธมฺมํ ปฏิจฺจ ปริตฺโต อรโณ ธมฺโม อุปฺปชฺชติ เหตุปจฺจยา. ตีณิ.}

\prose{มหคฺคตํ อรณํ ธมฺมํ ปฏิจฺจ มหคฺคโต อรโณ ธมฺโม อุปฺปชฺชติ เหตุปจฺจยา. ตีณิ.}

\csromanmarkh{13. ปริตฺตา\-รมฺมณ\-ตฺติก}\csromanmarkhii{100. สรณทุก}\csromanheadernomark{2}{13. ปริตฺตา\-รมฺมณ\-ตฺติก 100. สรณทุก}
\prose{\csromanfolio{629}อปฺปมาณํ อรณํ ธมฺมํ ปฏิจฺจ อปฺปมาโณ อรโณ ธมฺโม อุปฺปชฺชติ เหตุปจฺจยา. ตีณิ.}

\prose{ปริตฺตํ อรณญฺจ มหคฺคตํ อรณญฺจ ธมฺมํ ปฏิจฺจ ปริตฺโต อรโณ ธมฺโม อุปฺปชฺชติ เหตุปจฺจยา. ตีณิ.}

\prose{ปริตฺตํ อรณญฺจ อปฺปมาณํ อรณญฺจ ธมฺมํ ปฏิจฺจ ปริตฺโต อรโณ ธมฺโม อุปฺปชฺชติ เหตุปจฺจยา. (1) (สํขิตฺตํ.)}

\proseitem{110}{เหตุยา เตรส, อารมฺมเณ ปญฺจ ฯเปฯ อวิคเต เตรส. (สํขิตฺตํ. สหชาตวาเรปิ ฯเปฯ ปญฺหาวาเรปิ สพฺพตฺถ วิตฺถาโร.)}

\csromanmarkh{13. ปริตฺตา\-รมฺมณ\-ตฺติก}\csromanmarkhii{100. สรณทุก}\csromanheadernomark{2}{13. ปริตฺตา\-รมฺมณ\-ตฺติก 100. สรณทุก}
\prose{1. ปฏิจฺจวารเหตุ}

\proseitem{111}{ปริตฺตา\-รมฺมณํ สรณํ ธมฺมํ ปฏิจฺจ ปริตฺตารมฺมโณ สรโณ ธมฺโม อุปฺปชฺชติ เหตุปจฺจยา. (1)}

\prose{มหคฺคตา\-รมฺมณํ สรณํ ธมฺมํ ปฏิจฺจ มหคฺคตา\-รมฺมโณ สรโณ ธมฺโม อุปฺปชฺชติ เหตุปจฺจยา. (1) (สํขิตฺตํ.)}

\prose{เหตุยา เทฺว ฯเปฯ อวิคเต เทฺว. (สํขิตฺตํ.) (สหชาตวาเรปิ ฯเปฯ ปญฺหาวาเรปิ สพฺพตฺถ วิตฺถาโร.)}

\proseitem{112}{ปริตฺตา\-รมฺมณํ อรณํ ธมฺมํ ปฏิจฺจ ปริตฺตารมฺมโณ อรโณ ธมฺโม อุปฺปชฺชติ เหตุปจฺจยา. (1)}

\prose{มหคฺคตา\-รมฺมณํ อรณํ ธมฺมํ ปฏิจฺจ มหคฺคตา\-รมฺมโณ อรโณ ธมฺโม อุปฺปชฺชติ เหตุปจฺจยา. (1)}

\prose{อปฺปมาณา\-รมฺมณํ อรณํ ธมฺมํ ปฏิจฺจ อปฺปมาณารมฺมโณ อรโณ ธมฺโม อุปฺปชฺชติ เหตุปจฺจยา. (1) (สํขิตฺตํ.)}

\prose{เหตุยา ตีณิ ฯเปฯ อวิคเต ตีณิ. (สํขิตฺตํ.)}

\vspace{\tipitakaparskip}
\csromancenter{(สหชาตวาเรปิ ฯเปฯ ปญฺหาวาเรปิ สพฺพตฺถ วิตฺถาโร.)}
\csromanmarkh{ธมฺมานุโลม ติกทุก\-ปฏฺฐาน\-ปาฬิ}\csromanmarkhii{14. หีนตฺติก 100. สรณทุก}\csromanheadernomark{2}{ธมฺมานุโลม ติกทุก\-ปฏฺฐาน\-ปาฬิ 14. หีนตฺติก 100. สรณทุก}
\prose{\csromanfolio{630}1. ปฏิจฺจวารเหตุ}

\proseitem{113}{หีนํ สรณํ ธมฺมํ ปฏิจฺจ หีโน สรโณ ธมฺโม อุปฺปชฺชติ เหตุปจฺจยา. (สํขิตฺตํ.)}

\prose{เหตุยา เอกํ ฯเปฯ อวิคเต เอกํ. (สํขิตฺตํ.)}

\proseitem{114}{มชฺฌิมํ อรณํ ธมฺมํ ปฏิจฺจ มชฺฌิโม อรโณ ธมฺโม อุปฺปชฺชติ เหตุปจฺจยา. (1)}

\prose{ปณีตํ อรณํ ธมฺมํ ปฏิจฺจ ปณีโต อรโณ ธมฺโม อุปฺปชฺชติ เหตุปจฺจยา. ตีณิ.}

\prose{มชฺฌิมํ อรณญฺจ ปณีตํ อรณญฺจ ธมฺมํ ปฏิจฺจ มชฺฌิโม อรโณ ธมฺโม อุปฺปชฺชติ เหตุปจฺจยา. (1) (สํขิตฺตํ.)}

\prose{เหตุยา ปญฺจ ฯเปฯ อวิคเต ปญฺจ. (สํขิตฺตํ.)}

\vspace{\tipitakaparskip}
\csromancenter{(สหชาตวาเรปิ ฯเปฯ ปญฺหาวาเรปิ สพฺพตฺถ วิตฺถาโร.)}
\csromanmarkh{15. มิจฺฉตฺตนิยต\-ตฺติก}\csromanmarkhii{100. สรณทุก}\csromanheadernomark{2}{15. มิจฺฉตฺตนิยต\-ตฺติก 100. สรณทุก}
\prose{1. ปฏิจฺจวารเหตุ}

\proseitem{115}{มิจฺฉตฺต\-นิยตํ สรณํ ธมฺมํ ปฏิจฺจ มิจฺฉตฺต\-นิยโต สรโณ ธมฺโม อุปฺปชฺชติ เหตุปจฺจยา. (1)}

\prose{อนิยตํ สรณํ ธมฺมํ ปฏิจฺจ อนิยโต สรโณ ธมฺโม อุปฺปชฺชติ เหตุปจฺจยา. (1) (สํขิตฺตํ.)}

\prose{เหตุยา เทฺว ฯเปฯ อวิคเต เทฺว. (สํขิตฺตํ.) (สหชาตวาเรปิ ฯเปฯ ปญฺหาวาเรปิ สพฺพตฺถ วิตฺถาโร.)}

\proseitem{116}{สมฺมตฺต\-นิยตํ อรณํ ธมฺมํ ปฏิจฺจ สมฺมตฺตนิยโต อรโณ ธมฺโม อุปฺปชฺชติ เหตุปจฺจยา.  สมฺมตฺต\-นิยตํ อรณํ ธมฺมํ ปฏิจฺจ อนิยโต}

\vspace{\tipitakaparskip}
\csromancenter{มคฺคา\-รมฺมณ\-ตฺติก 100. สรณทุก อรโณ ธมฺโม อุปฺปชฺชติ เหตุปจฺจยา.  สมฺมตฺต\-นิยตํ อรณํ ธมฺมํ ปฏิจฺจ สมฺมตฺตนิยโต อรโณ จ อนิยโต อรโณ จ ธมฺมา อุปฺปชฺชนฺติ เหตุปจฺจยา. (3)}
\prose{\csromanfolio{631}อนิยตํ อรณํ ธมฺมํ ปฏิจฺจ อนิยโต อรโณ ธมฺโม อุปฺปชฺชติ เหตุปจฺจยา. (1)}

\prose{สมฺมตฺต\-นิยตํ อรณญฺจ อนิยตํ อรณญฺจ ธมฺมํ ปฏิจฺจ อนิยโต อรโณ ธมฺโม อุปฺปชฺชติ เหตุปจฺจยา. (1) (สํขิตฺตํ.)}

\prose{เหตุยา ปญฺจ, อารมฺมเณ เทฺว ฯเปฯ อวิคเต ปญฺจ. (สํขิตฺตํ.)}

\vspace{\tipitakaparskip}
\csromancenter{(สหชาตวารมฺปิ ฯเปฯ ปญฺหาวารมฺปิ วิตฺถาเรตพฺพํ.)}
\csromanmarkh{16. มคฺคา\-รมฺมณ\-ตฺติก}\csromanmarkhii{100. สรณทุก}\csromanheadernomark{2}{16. มคฺคา\-รมฺมณ\-ตฺติก 100. สรณทุก}
\prose{1-7. ปฏิจฺจาทิวารเหตุ}

\proseitem{117}{มคฺคารมฺมณํ อรณํ ธมฺมํ ปฏิจฺจ มคฺคารมฺมโณ อรโณ ธมฺโม อุปฺปชฺชติ เหตุปจฺจยา. ตีณิ.}

\prose{มคฺคเหตุกํ อรณํ ธมฺมํ ปฏิจฺจ มคฺคเหตุโก อรโณ ธมฺโม อุปฺปชฺชติ เหตุปจฺจยา. ตีณิ.}

\prose{มคฺคาธิปติํ อรณํ ธมฺมํ ปฏิจฺจ มคฺคาธิปติ อรโณ ธมฺโม อุปฺปชฺชติ เหตุปจฺจยา. ปญฺจ.}

\prose{มคฺคารมฺมณํ อรณญฺจ มคฺคาธิปติํ อรณญฺจ ธมฺมํ ปฏิจฺจ มคฺคารมฺมโณ อรโณ ธมฺโม อุปฺปชฺชติ เหตุปจฺจยา. ตีณิ.}

\prose{มคฺคเหตุกํ อรณญฺจ มคฺคาธิปติํ อารณญฺจ ธมฺมํ ปฏิจฺจ มคฺคเหตุโก อรโณ ธมฺโม อุปฺปชฺชติ เหตุปจฺจยา. ตีณิ. (สํขิตฺตํ.)}

\prose{เหตุยา สตฺตรส ฯเปฯ อวิคเต สตฺตรส. (สํขิตฺตํ.)}

\prose{(สหชาตวารมฺปิ ฯเปฯ สมฺปยุตฺตวารมฺปิ ปฏิจฺจ\-วาร\-สทิสํ.)}

\proseitem{118}{\csromanfolio{632}มคฺคารมฺมโณ อรโณ ธมฺโม มคฺคา\-รมฺมณสฺส อรณสฺส ธมฺมสฺส เหตุปจฺจเยน ปจฺจโย. ตีณิ.}

\prose{มคฺคเหตุโก อรโณ ธมฺโม มคฺคเหตุกสฺส อรณสฺส ธมฺมสฺส เหตุปจฺจเยน ปจฺจโย. ตีณิ.}

\prose{มคฺคาธิปติ อรโณ ธมฺโม มคฺคา\-ธิปติสฺส อรณสฺส ธมฺมสฺส เหตุปจฺจเยน ปจฺจโย. ปญฺจ.}

\prose{มคฺคารมฺมโณ อรโณ จ มคฺคาธิปติ อรโณ จ ธมฺมา มคฺคา\-รมฺมณสฺส อรณสฺส ธมฺมสฺส เหตุปจฺจเยน ปจฺจโย. ตีณิ.}

\prose{มคฺคเหตุโก อรโณ จ มคฺคาธิปติ อรโณ จ ธมฺมา มคฺคเหตุกสฺส อรณสฺส ธมฺมสฺส เหตุปจฺจเยน ปจฺจโย. ตีณิ. (สํขิตฺตํ.)}

\proseitem{119}{เหตุยา สตฺตรส, อารมฺมเณ นว ฯเปฯ อวิคเต สตฺตรส. (สํขิตฺตํ.)}

\prose{(ยถา กุสลตฺติเก ปญฺหาวารํ, เอวํ วิตฺถาเรตพฺพํ.)}

\csromanmarkh{17. อุปฺปนฺนตฺติก}\csromanmarkhii{100. สรณทุก}\csromanheadernomark{2}{17. \textbf{อุปฺปนฺนตฺติก} 100. สรณทุก}
\prose{7. ปญฺหาวารเหตุ อารมฺมณ}

\proseitem{120}{อุปฺปนฺโน สรโณ ธมฺโม อุปฺปนฺนสฺส สรณสฺส ธมฺมสฺส เหตุปจฺจเยน ปจฺจโย. (สํขิตฺตํ.)}

\prose{เหตุยา เอกํ, อารมฺมเณ เทฺว, อธิปติยา เทฺว, สหชาเต อญฺญมญฺเญ นิสฺสเย เอกํ, อุปนิสฺสเย เทฺว, กมฺเม ฯเปฯ สมฺปยุตฺเต เอกํ ฯเปฯ อวิคเต เอกํ. (สํขิตฺตํ.)}

\proseitem{121}{อุปฺปนฺโน อรโณ ธมฺโม อุปฺปนฺนสฺส อรณสฺส ธมฺมสฺส เหตุปจฺจเยน ปจฺจโย. (สํขิตฺตํ.)}

\csromanmarkh{19. อตีตา\-รมฺมณ\-ตฺติก}\csromanmarkhii{100. สรณทุก}\csromanheadernomark{2}{19. อตีตา\-รมฺมณ\-ตฺติก 100. สรณทุก}
\prose{\csromanfolio{633}เหตุยา เอกํ, อารมฺมเณ ตีณิ, อธิปติยา ตีณิ ฯเปฯ อุปนิสฺสเย ตีณิ ฯเปฯ อวิคเต เอกํ. (สํขิตฺตํ.)}

\csromanmarkh{18. อตีตตฺติก}\csromanmarkhii{100. สรณทุก}\csromanheadernomark{2}{18. \textbf{อตีตตฺติก 1}00. สรณทุก}
\prose{7. ปญฺหาวารเหตุ}

\proseitem{122}{ปจฺจุปฺปนฺโน สรโณ ธมฺโม ปจฺจุปฺปนฺนสฺส สรณสฺส ธมฺมสฺส เหตุปจฺจเยน ปจฺจโย. (สํขิตฺตํ.)}

\prose{เหตุยา เอกํ, อารมฺมเณ ตีณิ, อธิปติยา ตีณิ ฯเปฯ อุปนิสฺสเย ตีณิ ฯเปฯ อวิคเต เอกํ. (สํขิตฺตํ.)}

\proseitem{123}{ปจฺจุปฺปนฺโน อรโณ ธมฺโม ปจฺจุปฺปนฺนสฺส อรณสฺส ธมฺมสฺส เหตุปจฺจเยน ปจฺจโย. (สํขิตฺตํ.)}

\csromanhangingbegin
\hangingheaditem{123}{เหตุยา เอกํ, อารมฺมเณ ตีณิ ฯเปฯ อวิคเต เอกํ. (สํขิตฺตํ.)}
\hangingbody{(ยถา กุสลตฺติเก ปญฺหาวารํ, เอวํ วิตฺถาเรตพฺพํ.)}
\csromanhangingend

\csromanmarkh{19. อตีตา\-รมฺมณ\-ตฺติก}\csromanmarkhii{100. สรณทุก}\csromanheadernomark{2}{19. อตีตา\-รมฺมณ\-ตฺติก 100. สรณทุก}
\prose{1-7. ปฏิจฺจาทิวารเหตุ อารมฺมณ}

\proseitem{124}{อตีตารมฺมณํ สรณํ ธมฺมํ ปฏิจฺจ อตีตารมฺมโณ สรโณ ธมฺโม อุปฺปชฺชติ เหตุปจฺจยา. (1)}

\prose{อนาคตา\-รมฺมณํ สรณํ ธมฺมํ ปฏิจฺจ อนาคตารมฺมโณ สรโณ ธมฺโม อุปฺปชฺชติ เหตุปจฺจยา. (1)}

\prose{ปจฺจุปฺปนฺนา\-รมฺมณํ สรณํ ธมฺมํ ปฏิจฺจ ปจฺจุปฺปนฺนา\-รมฺมโณ สรโณ ธมฺโม อุปฺปชฺชติ เหตุปจฺจยา. (1) (สํขิตฺตํ.)}

\prose{เหตุยา ตีณิ ฯเปฯ อวิคเต ตีณิ. (สํขิตฺตํ.)}

\prose{(สหชาตวารมฺปิ ฯเปฯ ปญฺหาวารมฺปิ สพฺพตฺถ วิตฺถาเรตพฺพํ.)}

\proseitem{125}{\csromanfolio{634}อตีตารมฺมณํ อรณํ ธมฺมํ ปฏิจฺจ อตีตารมฺมโณ อรโณ ธมฺโม อุปฺปชฺชติ เหตุปจฺจยา. (1)}

\prose{อนาคตา\-รมฺมณํ อรณํ ธมฺมํ ปฏิจฺจ อนาคตารมฺมโณ อรโณ ธมฺโม อุปฺปชฺชติ เหตุปจฺจยา. (1)}

\prose{ปจฺจุปฺปนฺนา\-รมฺมณํ อรณํ ธมฺมํ ปฏิจฺจ ปจฺจุปฺปนฺนา\-รมฺมโณ อรโณ ธมฺโม อุปฺปชฺชติ เหตุปจฺจยา. (1) (สํขิตฺตํ.)}

\prose{เหตุยา ตีณิ ฯเปฯ อวิคเต ตีณิ. (สํขิตฺตํ.) (สหชาตวารมฺปิ ฯเปฯ สมฺปยุตฺตวารมฺปิ ปฏิจฺจ\-วาร\-สทิสํ.)}

\proseitem{126}{อตีตารมฺมโณ อรโณ ธมฺโม อตีตา\-รมฺมณสฺส อรณสฺส ธมฺมสฺส เหตุปจฺจเยน ปจฺจโย. (1)}

\prose{อนาคตารมฺมโณ อรโณ ธมฺโม อนาคตา\-รมฺมณสฺส อรณสฺส ธมฺมสฺส เหตุปจฺจเยน ปจฺจโย. (1)}

\prose{ปจฺจุปฺปนฺนา\-รมฺมโณ อรโณ ธมฺโม ปจฺจุปฺปนฺนา\-รมฺมณสฺส อรณสฺส ธมฺมสฺส เหตุปจฺจเยน ปจฺจโย. (1)}

\prose{อตีตารมฺมโณ อรโณ ธมฺโม อตีตา\-รมฺมณสฺส อรณสฺส ธมฺมสฺส อารมฺมณ\-ปจฺจเยน ปจฺจโย. (สํขิตฺตํ.)}

\proseitem{127}{เหตุยา ตีณิ, อารมฺมเณ นว ฯเปฯ อวิคเต ตีณิ. (สํขิตฺตํ. ยถา กุสลตฺติเก ปญฺหาวารํ, เอวํ วิตฺถาเรตพฺพํ.)}

\csromanmarkh{20. อชฺฌตฺตตฺติกฺอ}\csromanmarkhii{100. สรณทุก}\csromanheadernomark{2}{20. \textbf{อชฺฌตฺตตฺติกฺ}อ 100. สรณทุก}
\prose{1-7. ปฏิจฺจาทิวารเหตุ}

\proseitem{128}{อชฺฌตฺตํ สรณํ ธมฺมํ ปฏิจฺจ อชฺฌตฺโต สรโณ ธมฺโม อุปฺปชฺชติ เหตุปจฺจยา.}

\vspace{\tipitakaparskip}
\csromancenter{อชฺฌตฺตา\-รมฺมณ\-ตฺติก 100. สรณทุก พหิทฺธา สรณํ ธมฺมํ ปฏิจฺจ พหิทฺธา สรโณ ธมฺโม อุปฺปชฺชติ เหตุปจฺจยา.}
\prose{\csromanfolio{635}เหตุยา เทฺว ฯเปฯ อวิคเต เทฺว. (สํขิตฺตํ.) (สหชาตวาเรปิ ฯเปฯ ปญฺหาวาเรปิ สพฺพตฺถ วิตฺถาโร.)}

\proseitem{129}{อชฺฌตฺตํ อรณํ ธมฺมํ ปฏิจฺจ อชฺฌตฺโต อรโณ ธมฺโม อุปฺปชฺชติ เหตุปจฺจยา. (1)}

\prose{พหิทฺธา อรณํ ธมฺมํ ปฏิจฺจ พหิทฺธา อรโณ ธมฺโม อุปฺปชฺชติ เหตุปจฺจยา. (1) (สํขิตฺตํ.)}

\prose{เหตุยา เทฺว ฯเปฯ อวิคเต เทฺว. (สํขิตฺตํ.)}

\vspace{\tipitakaparskip}
\csromancenter{(สหชาตวาเรปิ ฯเปฯ ปญฺหาวาเรปิ สพฺพตฺถ วิตฺถาโร.)}
\csromanmarkh{21. อชฺฌตฺตา\-รมฺมณ\-ตฺติก}\csromanmarkhii{100. สรณทุก}\csromanheadernomark{2}{21. อชฺฌตฺตา\-รมฺมณ\-ตฺติก 100. สรณทุก}
\prose{1-7. ปฏิจฺจาทิวารเหตุ}

\proseitem{130}{อชฺฌตฺตา\-รมฺมณํ สรณํ ธมฺมํ ปฏิจฺจ อชฺฌตฺตา\-รมฺมโณ สรโณ ธมฺโม อุปฺปชฺชติ เหตุปจฺจยา. (1)}

\prose{พหิทฺธา\-รมฺมณํ สรณํ ธมฺมํ ปฏิจฺจ พหิทฺธา\-รมฺมโณ สรโณ ธมฺโม อุปฺปชฺชติ เหตุปจฺจยา. (1) (สํขิตฺตํ.)}

\prose{เหตุยา เทฺว ฯเปฯ อวิคเต เทฺว. (สํขิตฺตํ.) (สหชาตวาเรปิ ฯเปฯ ปญฺหาวาเรปิ สพฺพตฺถ วิตฺถาโร.)}

\proseitem{131}{อชฺฌตฺตา\-รมฺมณํ อรณํ ธมฺมํ ปฏิจฺจ อชฺฌตฺตา\-รมฺมโณ อรโณ ธมฺโม อุปฺปชฺชติ เหตุปจฺจยา. (1)}

\prose{พหิทฺธา\-รมฺมณํ อรณํ ธมฺมํ ปฏิจฺจ พหิทฺธา\-รมฺมโณ อรโณ ธมฺโม อุปฺปชฺชติ เหตุปจฺจยา. (1) (สํขิตฺตํ.)}

\prose{เหตุยา เทฺว ฯเปฯ อวิคเต เทฺว. (สํขิตฺตํ.)}

\vspace{\tipitakaparskip}
\csromancenter{(สหชาตวาเรปิ ฯเปฯ ปญฺหาวาเรปิ สพฺพตฺถ วิตฺถาโร.)}
\csromanmatikaanchor{mtk.153}
\csromanmatikaanchor{mtk.154}
\csromantocmarknopage{subparagraph}{22. สนิทสฺสนตฺติก}{mtk.153}
\bookmark[dest={mtk.153},level=5]{22. สนิทสฺสนตฺติก}
\csromantocmarkref{subparagraph}{100. สรณทุก}{mtk.154}
\bookmark[dest={mtk.154},level=5]{100. สรณทุก}
\csromanmarkh{ธมฺมานุโลม ติกทุก\-ปฏฺฐาน\-ปาฬิ}\csromanmarkhii{22. สนิทสฺสน\-ตฺติก 100. สรณทุก}\csromanheadernomark{6}{ธมฺมานุโลม ติกทุก\-ปฏฺฐาน\-ปาฬิ 22. สนิทสฺสน\-ตฺติก 100. สรณทุก}
\prose{\csromanfolio{636}1-7. ปฏิจฺจาทิวารเหตุ}

\proseitem{132}{อนิทสฺสน\-อปฺปฏิฆํ สรณํ ธมฺมํ ปฏิจฺจ อนิทสฺสน\-อปฺปฏิโฆ สรโณ ธมฺโม อุปฺปชฺชติ เหตุปจฺจยา. (สพฺพตฺถ เอกํ.)}

\prose{อนิทสฺสน\-อปฺปฏิฆํ อรณํ ธมฺมํ ปฏิจฺจ อนิทสฺสน\-อปฺปฏิโฆ อรโณ ธมฺโม อุปฺปชฺชติ เหตุปจฺจยา. สตฺต. (สํขิตฺตํ.)}

\proseitem{133}{เหตุยา เอกวีส ฯเปฯ อวิคเต เอกวีส. (สํขิตฺตํ.) (สหชาตวารมฺปิ ฯเปฯ สมฺปยุตฺตวารมฺปิ ปฏิจฺจ\-วาร\-สทิสํ.)}

\proseitem{134}{อนิทสฺสน\-อปฺปฏิโฆ อรโณ ธมฺโม อนิทสฺสน\-อปฺปฏิฆสฺส อรณสฺส ธมฺมสฺส เหตุปจฺจเยน ปจฺจโย. (สํขิตฺตํ.)}

\prosewithcloser{เหตุยา สตฺต, อารมฺมเณ ตีณิ ฯเปฯ อวิคเต ปญฺจวีส. (สํขิตฺตํ.) (ยถา กุสลตฺติเก ปญฺหาวารํ, เอวํ วิตฺถาเรตพฺพํ.)}{สนิทสฺสน\-ตฺติก\-สรณทุกํ นิฏฺฐิตํ.}
\nitthitam{ธมฺมานุโลเม ติกทุก\-ปฏฺฐานํ นิฏฺฐิตํ.}
